% !TeX program = xelatex
\documentclass[11pt,a4paper]{article}
\usepackage[margin=23mm]{geometry}
\usepackage{amsmath,amssymb,amsthm,mathtools}
\usepackage{enumitem,microtype,xcolor}
\usepackage[hyphens]{url}
\usepackage[colorlinks=true,linkcolor=blue!55!black,citecolor=blue!55!black,urlcolor=blue!55!black]{hyperref}
\hypersetup{
  pdftitle={Reading Notes on Four-Dimensional Yang--Mills Theory},
  pdfauthor={Zexi He (maintainer)},
  pdfsubject={AI-generated reading notes based on OpenAI's Lean proof},
  pdfkeywords={Yang--Mills theory, reading notes, Lean, AI assistance}
}
\allowdisplaybreaks
\setlength{\parskip}{.35em}
\setlength{\parindent}{0pt}
\setlist{topsep=4pt,itemsep=3pt}
\numberwithin{equation}{section}
\newtheorem{theorem}{Theorem}[section]
\newtheorem{proposition}[theorem]{Proposition}
\newtheorem{lemma}[theorem]{Lemma}
\newtheorem{corollary}[theorem]{Corollary}
\theoremstyle{definition}
\newtheorem{definition}[theorem]{Definition}
\newcommand{\R}{\mathbb R}
\newcommand{\C}{\mathbb C}
\newcommand{\N}{\mathbb N}
\newcommand{\Z}{\mathbb Z}
\newcommand{\dd}{\mathrm d}
\newcommand{\eps}{\varepsilon}
\newcommand{\supp}{\operatorname{supp}}
\title{Reading Notes on Four-Dimensional Yang--Mills Theory\\
  \large An AI-generated Working Exposition}
\author{Maintained by Zexi He}
\date{13 November 2026}

\begin{document}
\maketitle
\begin{quote}
\small
\textbf{Status statement.}
These notes are based on OpenAI's Lean proof of the four-dimensional
Yang--Mills construction~\cite{OpenAIYM26} and the related expositions cited
below. This substantially AI-generated draft is maintained by Zexi He
for personal study and discussion. It has been assembled
through its final construction argument, but has not been read or checked
line by line by the maintainer or subjected to an exhaustive independent
mathematical review. It has not been mechanically checked against Lean;
adapted arguments require their own justification. This is not an official
exposition or an endorsement by OpenAI or the source contributors.
Further source and review information appears in the accompanying README.
Corrections and missing references are welcome.
\end{quote}

\begin{abstract}
These reading notes aim to explain the four-dimensional Yang--Mills
construction formalised in OpenAI's Lean proof. For every closed subgroup
of a finite-dimensional unitary group whose real Lie algebra is nonabelian
and simple, the construction starts from one family of Wilson lattice
measures and local renormalised insertions.
For every requested increasing volume sequence, a refinement gives
one ordered continuum hierarchy. Its mixed correlations have a
full mesh limit at every suitable fixed period, followed by the
selected volume limit. The hierarchy has faithful finite-dimensional
field spaces, a compact metric phase decomposition with a probability
measure of full support, and a selected Wightman theory satisfying
W0--W3 with a unique translation-invariant vacuum. One coefficient
family gives every finite-order separated operator product expansion
in the mixture and all phases. The centred stress trace in that same
selected vacuum has a positive short-distance amplitude with an
inverse logarithmic correction, excluding the specified finite
mass-shell Wick profiles.

The exposition develops the stationary and Gaussian estimates, the
small-field and large-field activities, and the nonlinear
renormalisation-group induction with local bounds uniform in cutoff
and physical volume. Exact reassembly retains all masks, incoming
activities and normalisations. The continuum construction includes
position derivatives, the interacting stress and local charges,
full orthogonal covariance, phase selection, analytic continuation
and reconstruction. The final section assembles the construction and
records the common choices and limit order. The notes are intended for
readers with graduate measure theory and functional analysis, and aim to
supply the renormalisation-group background as needed. A positive mass
gap is outside the scope of the construction discussed here.
\end{abstract}
\clearpage
\tableofcontents

\section{The starting point and the target}
\label{sec:target}

The immediate problem is to construct continuum correlations from
finite-dimensional probability measures. Before discussing
renormalisation, we need to fix the group, the lattice measure, and the
order in which its cutoffs are removed. Otherwise an estimate for one
family of lattice observables could be mistaken for a statement about a
different family.

For historical context, the analysis of
\cite[Abstract and Sect.~I, printed pp.~325--326]{MRS93}
treats pure SU(2) Yang--Mills theory in the trivial topological
sector, with a fixed infrared cutoff and a regularised axial gauge.
This is a background comparison concerning the cutoff problem.
The Wilson family and the order of limits used here are specified below.

\subsection{The gauge group}

\begin{definition}[Admissible gauge group]
\label{def:group}
Fix a positive integer \(N\) and a closed subgroup \(G\subset U(N)\),
with its given faithful matrix representation. Set
\begin{equation}
 \mathfrak g
 =\{X\in \operatorname{Mat}_N(\C):\exp(tX)\in G\text{ for every }t\in\R\}.
 \label{eq:lie-algebra}
\end{equation}
The data include the assertion that this set is a real vector subspace.
Its bracket is the matrix commutator
\([X,Y]=XY-YX\). The simplicity assumption is that
\(\mathfrak g\) is nonabelian and that its only real Lie ideals are
\(\{0\}\) and \(\mathfrak g\). A Lie ideal is a real vector
subspace \(\mathfrak i\subset\mathfrak g\) satisfying
\([\mathfrak g,\mathfrak i]\subset\mathfrak i\).
\end{definition}

Here \(N\) is the dimension of the fixed representation.
Equation \eqref{eq:lie-algebra} is part of the group data:
the Lie algebra is tied to the group, rather than being an
independently chosen vector space. No connectedness or simple
connectedness assumption is added here.

Since \(U(N)\) is compact and \(G\) is closed, \(G\) is compact.
Write \(h_G\) for its probability Haar measure. Then
\begin{equation}
 \int_G f(g_0g)\,\dd h_G(g)
 =\int_G f(g)\,\dd h_G(g)
 =\int_G f(gg_0)\,\dd h_G(g)
 \label{eq:haar}
\end{equation}
for every integrable \(f\) and \(g_0\in G\).
For a compact group the probability Haar measure is both left and
right invariant. Indeed, right translation of a left Haar probability
measure is again a left Haar probability measure; uniqueness gives
right invariance.

\begin{lemma}[Matrix-group calculus from the stated group data]
\label{lem:group-data-calculus}
The space \(\mathfrak g\) is a closed subspace of the
skew-Hermitian matrices, is closed under commutators and
is invariant under conjugation by every element of \(G\).
Exponential is a real analytic chart from a neighbourhood
of zero in \(\mathfrak g\) onto a relative neighbourhood of
the identity in \(G\). In that chart, probability Haar
measure has a strictly positive real analytic density,
including a fixed positive normalisation constant.
These statements do not require \(G\) to be connected.
\end{lemma}
\begin{proof}
Differentiating \(\exp(tX)^*\exp(tX)=I\) at zero gives
\(X^*=-X\). Finite dimensionality makes the given vector
space \(\mathfrak g\) closed. For \(g\in G\),
\(\exp(tgXg^{-1})=g\exp(tX)g^{-1}\) belongs to \(G\);
this proves conjugation invariance. The difference
quotients of
\(\exp(tX)Y\exp(-tX)\in\mathfrak g\) tend to \([X,Y]\),
proving commutator closure.

Here is a local-chart argument which also excludes an
implicit connectedness hypothesis. Let
\(\mathcal A=\{X:X^*=-X\}\), with real inner product
\(\operatorname{Re}\operatorname{tr}(X^*Y)\) and norm
\(\|X\|_{\rm F}^2=\operatorname{Re}\operatorname{tr}(X^*X)\),
and take the orthogonal complement of
\(\mathfrak g\) in \(\mathcal A\). The matrix logarithm
near \(I\), given by its convergent power series
\(\log(I+Z)=\sum_{n\ge1}(-1)^{n+1}Z^n/n\), and the
inverse function theorem show that
\[
       (X,Y)\longmapsto\exp X\exp Y,\qquad
       X\in\mathfrak g,\quad Y\in\mathfrak g^\perp,
\]
is a real analytic chart near \((0,0)\); its differential
is \((X,Y)\mapsto X+Y\). Suppose there were group
elements tending to \(I\) whose chart coordinate \(Y_j\)
was nonzero. Since \(\exp X_j\in G\), also
\(\exp Y_j=\exp(-X_j)g_j\in G\).
Pass to a subsequence with
\(Y_j/\|Y_j\|_{\rm F}\to Y\in\mathfrak g^\perp\),
\(\|Y\|_{\rm F}=1\). For every real \(t\), choose integers
\(n_j\) with \(n_j\|Y_j\|_{\rm F}\to t\). Then
\(\exp(n_jY_j)\in G\) tends to \(\exp(tY)\).
Closedness gives \(\exp(tY)\in G\) for every \(t\),
so \(Y\in\mathfrak g\), a contradiction.
After shrinking the chart, its group points therefore
have \(Y=0\). Conversely every sufficiently small
\(\exp X\), \(X\in\mathfrak g\), belongs to \(G\).

Left translates give charts on all components. The
left-invariant metric induced by the Frobenius inner
product defines a left-invariant volume density there.
Compactness makes its total mass finite and positive,
so its probability normalisation is \(h_G\), by Haar
uniqueness. Left trivialisation of the exponential
differential is
\[
 V\longmapsto
      \int_0^1 e^{-sX}Ve^{sX}\,\dd s .
\]
Its determinant on \(\mathfrak g\) is analytic and equals
one at zero, hence is positive in a smaller chart.
Multiplication by the fixed normalising constant gives
the stated Haar density. In particular that constant
is part of the Haar integral; chart calculations do not
replace it by one.
\end{proof}

\subsection{What is to be constructed}

Each Wilson family comes with a chosen unbounded set of physical
side lengths; these are its admissible side lengths.
The construction theorem has the following quantifier order:
\begin{equation}
\begin{split}
 &\text{for every admissible simple }G,\quad
   \exists\, C\text{ a Wilson family},\\
 &C\text{ has Wilson regularity}
 \quad\text{and}\quad
  \forall\,(M_j)\text{ strictly increasing and admissible},\\
 &\quad\exists\,q:\N\to\N\text{ strictly increasing},\
 \exists\,T\text{ a continuum hierarchy for }C,\\
 &\quad\text{all four conclusions hold along }(M_{q(j)})_j.
\end{split}
\label{eq:public-target}
\end{equation}
The objects and the four conclusions are described in
Sections~\ref{sec:family} and \ref{sec:continuum}.

The Wilson family \(C\) is chosen before the requested volume
sequence \((M_j)\). Only \(q\) and \(T\) are chosen afterwards.
All four conclusions concern this same \(T\) and this same \(C\).

\section{The finite Wilson model}
\label{sec:wilson}

At finite cutoff there is no need to begin with a measure on
distribution-valued gauge potentials. The basic variables are compact
group elements on links. This gives finite-dimensional probability
measures while retaining exact gauge invariance. This lattice starting
point goes back to Wilson \cite{Wilson74}.

\subsection{Links, plaquettes, and the action}

For an integer \(\ell\ge1\), define
\begin{equation}
\begin{gathered}
 \Lambda_\ell=(\Z/\ell\Z)^4,\qquad
 E_\ell^+=\Lambda_\ell\times\{1,2,3,4\},\\
 \mathcal P_\ell^+
   =\{(x,i,j):x\in\Lambda_\ell,\ 1\le i<j\le4\},\qquad
 \mathcal U_\ell=G^{E_\ell^+}.
\end{gathered}
 \label{eq:lattice}
\end{equation}
The sets \(E_\ell^+\) and \(\mathcal P_\ell^+\) label the positively
oriented edges and plaquettes, respectively; \(\Lambda_\ell\) is the site set.
A configuration assigns \(U_e=U(x,i)\in G\) to the positively oriented
edge \(e=(x,i)\) from \(x\) to \(x+e_i\), where \(e_i\) is the \(i\)-th
coordinate vector. The reverse link carries \(U(x,i)^{-1}\).
Addition of sites is modulo \(\ell\).
The probability Haar prior on configurations is
\begin{equation}
 H_\ell=\bigotimes_{e\in E_\ell^+}h_G.
 \label{eq:prior}
\end{equation}
Both products are finite.

\begin{definition}[Plaquette and Wilson energy]
For \(p=(x,i,j)\in\mathcal P_\ell^+\), write
\(P_p(U)=P_{ij}(U;x)\), where
\begin{align}
 P_{ij}(U;x)
 &=U(x,i)U(x+e_i,j)
       U(x+e_j,i)^{-1}U(x,j)^{-1}, \label{eq:plaquette}\\
 W_\ell(U)
 &=\sum_{p\in\mathcal P_\ell^+}
       \left(1-\frac1N\operatorname{Re}\operatorname{tr}
                                       P_p(U)\right).
 \label{eq:energy}
\end{align}
The plaquette is the ordered group product around an elementary square.
\end{definition}

% !TeX root = 4DYM_notes.tex
\emph{The energy as a squared plaquette deviation.}
For a complex matrix \(B\), define its unnormalised Hilbert--Schmidt
norm by
\[
                 \|B\|_{\rm HS}^2:=\operatorname{tr}(B^*B).
\]
If \(P\in U(N)\), unitarity gives the exact calculation
\begin{align*}
 \|P-I_N\|_{\rm HS}^2
   &=\operatorname{tr}\bigl((P^*-I_N)(P-I_N)\bigr)\\
   &=\operatorname{tr}(2I_N-P-P^*)
     =2N-2\operatorname{Re}\operatorname{tr}P.
\end{align*}
Consequently the Wilson energy in \eqref{eq:energy} is
\[
 W_\ell(U)=\frac1{2N}\sum_{p\in\mathcal P_\ell^+}
                                      \|P_p(U)-I_N\|_{\rm HS}^2.
\]
It measures the total squared deviation of the plaquette holonomies
from the identity. This formula holds for every finite configuration.

\emph{Covariant differences and their commutator.}
Give the integer lattice a physical mesh \(\varepsilon>0\), so that
one lattice step has physical length \(\varepsilon\).
For an auxiliary test field \(\psi:\Lambda_\ell\to\C^N\), define
\[
 (D_{U,i}^{\varepsilon}\psi)(x)
       :=\frac{U(x,i)\psi(x+e_i)-\psi(x)}{\varepsilon}.
\]
The link matrix transports the neighbouring vector back to \(x\)
before taking its difference. Multiplying out two such operators gives
\begin{align*}
 (D_{U,i}^{\varepsilon}D_{U,j}^{\varepsilon}\psi)(x)
 =\varepsilon^{-2}\bigl\{&U(x,i)U(x+e_i,j)\psi(x+e_i+e_j)\\
             &-U(x,i)\psi(x+e_i)-U(x,j)\psi(x+e_j)+\psi(x)\bigr\}.
\end{align*}
Subtracting the same expression with \(i,j\) interchanged cancels
the last three terms. For the commutator
\([D_{U,i}^{\varepsilon},D_{U,j}^{\varepsilon}]
 =D_{U,i}^{\varepsilon}D_{U,j}^{\varepsilon}
  -D_{U,j}^{\varepsilon}D_{U,i}^{\varepsilon}\), this yields
\begin{align*}
 &[D_{U,i}^{\varepsilon},D_{U,j}^{\varepsilon}]\psi(x)\\
 &\quad=\varepsilon^{-2}
    \bigl\{U(x,i)U(x+e_i,j)-U(x,j)U(x+e_j,i)\bigr\}
                                           \psi(x+e_i+e_j)\\
 &\quad=\varepsilon^{-2}\bigl(P_{ij}(U;x)-I_N\bigr)
                 U(x,j)U(x+e_j,i)\psi(x+e_i+e_j).
\end{align*}
The final equality follows by right multiplying the plaquette
formula \eqref{eq:plaquette} by \(U(x,j)U(x+e_j,i)\).
Thus the commutator contains a shift by \(e_i+e_j\); its matrix
coefficient is the plaquette deviation followed by a two-link
parallel transport. Right multiplication by that unitary transport
preserves the Hilbert--Schmidt norm. The coefficient therefore has
squared norm
\[
                    \varepsilon^{-4}\|P_{ij}(U;x)-I_N\|_{\rm HS}^2.
\]
The deviation \(P-I_N\) is used here as a matrix-valued object.
The Lie-algebra-valued small-field coordinates used later are
defined separately by local group charts.

\emph{Why four dimensions have no additional mesh factor.}
In physical units each four-dimensional lattice cell has volume
\(\varepsilon^4\). Hence the preceding curvature-coefficient
calculation gives
\begin{align*}
 \varepsilon^4\sum_{p\in\mathcal P_\ell^+}
       \bigl\|\varepsilon^{-2}(P_p(U)-I_N)\bigr\|_{\rm HS}^2
   &=\sum_{p\in\mathcal P_\ell^+}\|P_p(U)-I_N\|_{\rm HS}^2\\
   &=2N\,W_\ell(U).
\end{align*}
These are exact finite-lattice identities. The cell-volume factor
and the squared-curvature factor cancel in four dimensions.
The coefficient multiplying \(W_\ell\) will nevertheless depend
on the cutoffs through the running-coupling law proved below.


\begin{lemma}[Elementary action bounds]
\label{lem:energy}
The function \(W_\ell\) is continuous and
\begin{equation}
 0\le W_\ell(U)\le12\ell^4.
 \label{eq:energy-bound}
\end{equation}
\end{lemma}
\begin{proof}
Every plaquette is unitary. For \(W\in U(N)\), each diagonal
entry satisfies \(|W_{aa}|\le1\), since a row has squared
Euclidean norm one. Hence
\begin{equation}
 -N\le\operatorname{Re}\operatorname{tr}W
       =\sum_{a=1}^N\operatorname{Re}W_{aa}\le N.
\end{equation}
Each energy summand lies in \([0,2]\). There are
\(\ell^4\binom42=6\ell^4\) summands, proving the bound.
Continuity follows from continuity of multiplication, inversion on
\(G\), and the trace.
\end{proof}

\begin{definition}[Finite Wilson probability measure]
For \(\beta\in\R\), set
\begin{align}
 Z_\ell(\beta)
 &=\int_{\mathcal U_\ell}e^{-\beta W_\ell(U)}\,\dd H_\ell(U),
 \label{eq:partition}\\
 \dd\mu_{\ell,\beta}(U)
 &=Z_\ell(\beta)^{-1}e^{-\beta W_\ell(U)}\,\dd H_\ell(U).
 \label{eq:wilson-measure}
\end{align}
\end{definition}

\emph{Normalisation of the Wilson coefficient at fixed matrix dimension.}
For a fixed constant \(c>0\), rescaling the energy by \(c\) and the
coefficient by \(c^{-1}\) leaves the weight unchanged:
\[
                  e^{-(\beta/c)(cW_\ell(U))}=e^{-\beta W_\ell(U)}.
\]
The Haar prior and the normalising integral are then the same, so this
is a reparameterisation of the same probability measure. We keep
the convention \eqref{eq:energy} throughout. Once it is fixed,
changing \(\beta\) generally changes the measure. Its cutoff-dependent
choice is part of the construction and is constrained by the
running-coupling law proved below.
The group and matrix dimension are fixed in this construction,
so no large-\(N\) scaling is involved. All later coefficients use
the normalisation in \eqref{eq:energy}.

\begin{proposition}[Finite-cutoff well-definedness]
\label{prop:finite-measure}
For every real \(\beta\), \(0<Z_\ell(\beta)<\infty\);
\(\mu_{\ell,\beta}\) is a probability measure; and every continuous
complex-valued observable is integrable. If \(\beta\ge0\), then
\begin{equation}
 e^{-12\beta\ell^4}\le Z_\ell(\beta)\le1.
 \label{eq:crude-partition}
\end{equation}
\end{proposition}
\begin{proof}
For arbitrary real \(\beta\), Lemma~\ref{lem:energy} implies
\begin{equation}
 e^{-12|\beta|\ell^4}\le e^{-\beta W_\ell(U)}
                              \le e^{12|\beta|\ell^4}.
\end{equation}
Integrating against the probability prior proves strict positivity
and finiteness. For \(\beta\ge0\), use the sharper pointwise bounds
\(e^{-12\beta\ell^4}\le e^{-\beta W_\ell(U)}\le1\).
Normalisation follows by dividing the partition integral by itself.
Finally, \(\mathcal U_\ell\) is a finite product of compact spaces.
A continuous observable is bounded and thus integrable against its
probability measure.
\end{proof}

The explicit dependence on \(\ell^4\) in
\eqref{eq:crude-partition} makes it insufficient for the continuum
construction. Uniform estimates require further analysis.

\subsection{Gauge invariance without gauge fixing}

\begin{definition}[Lattice gauge transformation]
For \(t:\Lambda_\ell\to G\), define
\begin{equation}
 U^t(x,i)=t(x)U(x,i)t(x+e_i)^{-1}.
 \label{eq:gauge}
\end{equation}
An observable \(F\) is gauge invariant if \(F(U^t)=F(U)\)
for every \(t,U\).
\end{definition}

\begin{proposition}[Gauge invariance of the finite measure]
\label{prop:gauge}
For every \(t\),
\begin{equation}
 P_{ij}(U^t;x)=t(x)P_{ij}(U;x)t(x)^{-1},
 \qquad W_\ell(U^t)=W_\ell(U).
 \label{eq:gauge-plaquette}
\end{equation}
For every integrable \(F\),
\begin{equation}
 \int F(U^t)\,\dd\mu_{\ell,\beta}(U)
       =\int F(U)\,\dd\mu_{\ell,\beta}(U).
 \label{eq:gauge-measure}
\end{equation}
\end{proposition}
\begin{proof}
Substitute \eqref{eq:gauge} into the four plaquette factors:
\begin{align}
 P_{ij}(U^t;x)
 ={}&t(x)U(x,i)t(x+e_i)^{-1}\\
 &{}\cdot t(x+e_i)U(x+e_i,j)t(x+e_i+e_j)^{-1}\notag\\
 &{}\cdot t(x+e_i+e_j)U(x+e_j,i)^{-1}t(x+e_j)^{-1}\notag\\
 &{}\cdot t(x+e_j)U(x,j)^{-1}t(x)^{-1}.\notag
\end{align}
The three intermediate pairs cancel. Conjugation invariance of
the trace proves energy invariance. For each link the transformation
is a fixed left and right translation on \(G\), so it preserves
Haar measure by \eqref{eq:haar}. The product transformation
preserves \(H_\ell\); invariance of the density then proves
\eqref{eq:gauge-measure}.
\end{proof}

Gauge fixing is not needed to define this finite measure.
It will enter later as a local analytic device for estimating
fluctuations.

\section{One Wilson family and its two cutoffs}
\label{sec:family}

\subsection{Physical size versus lattice spacing}

The matrix dimension \(N\) is fixed throughout. We use \(n\) for
the terminal ultraviolet level, \(m\) for an intermediate RG level,
and \(M\) for the physical side length. Lattice coordinates remain in lattice
units; multiplication by the mesh gives physical coordinates.
When the number of insertions occurs together with these indices,
we denote it by \(r\).

A Wilson family \(C\) chooses a strictly decreasing sequence of
mesh sizes \(\varepsilon_n\to0\), whose reciprocals are positive
integers, and an unbounded set of positive integer physical side
lengths \(M\). These chosen side lengths are called admissible;
every volume parameter below ranges over this set. The construction
chooses it to satisfy the blocking and reflection compatibility
conditions used later. Define
\begin{equation}
 \ell_{n,M}:=\frac{M}{\varepsilon_n}\in\N.
 \label{eq:two-cutoffs}
\end{equation}
Consequently,
\begin{equation}
 \varepsilon_n\ell_{n,M}=M.
 \label{eq:physical-side}
\end{equation}
The index \(n\) controls the mesh, while \(M\) is the actual side
length, measured in the fixed physical units of the construction.
An increasing enumeration of the admissible side lengths identifies
them with integer indices. Strictly increasing sequences and their
further subsequences correspond under this identification, and
uniform bounds are unchanged by it. We use \(\#\) for site count and
\(\operatorname{Vol}\) for physical Lebesgue volume. Thus
\[
 \#\Lambda_{\ell_{n,M}}=\ell_{n,M}^4,\qquad
 \operatorname{Vol}\bigl(\R^4/(M\Z^4)\bigr)
       =\varepsilon_n^4\#\Lambda_{\ell_{n,M}}=M^4.
\]

At a fixed terminal level \(n\), write \(\varepsilon=\varepsilon_n\).
A mesh subscript denotes the datum at that cutoff, with the side length first;
for example, \(\ell_{M,\varepsilon}=\ell_{n,M}\). The same convention
applies to couplings, insertions and correlations. Every limit
\(\varepsilon\downarrow0\) below is taken through the chosen mesh
sequence. We retain scale subscripts when comparing different RG levels.

The family also chooses \(\beta_{M,\varepsilon}\), a set \(\mathcal V\)
of insertion labels, and functions
\begin{equation}
 \mathcal O_{M,\varepsilon}(v,x;U)\in\C,\quad
 v\in\mathcal V,\
 x\in\Lambda_{\ell_{M,\varepsilon}},\
 U\in\mathcal U_{\ell_{M,\varepsilon}}.
 \label{eq:insertion}
\end{equation}
These functions and their normalisations are fixed as part of
the family, before any limit is taken.

\begin{definition}[Wilson regularity]
\label{def:regularity}
The required properties are as follows.
\begin{enumerate}
 \item There is a constant \(\operatorname{slope}>0\) such that
 \begin{equation}
 \forall\delta>0\ \exists n_0\ \forall n\ge n_0\ \forall M,\quad
 \left|\frac{\beta_{M,\varepsilon_n}}{\log\varepsilon_n^{-1}}
                        -\operatorname{slope}\right|<\delta.
 \label{eq:running-public}
 \end{equation}
 \item Each \(\mathcal O_{M,\varepsilon}(v,x;\cdot)\) is continuous and
 gauge invariant.
 \item For each \(v\) there is one finite set
 \(D_v\subset\Z^4\), its lattice stencil, such that the insertion
 only reads links whose two endpoints lie in \(x+D_v\), modulo
 \(\ell_{M,\varepsilon}\). Precisely, agreement of two configurations on
 all these links implies equality of their insertions.
\end{enumerate}
\end{definition}

The stencil is independent of \(\varepsilon,M\).
Its diameter in the physical torus metric is at most \(\varepsilon\)
times its fixed lattice diameter. The two-endpoint condition
specifies which links determine the insertion.

\subsection{Moments and smearing}

For labels \(v_1,\ldots,v_k\) and sites \(x_1,\ldots,x_k\), set
\begin{equation}
 C_{M,\varepsilon}(v;x)=
 \int\prod_{j=1}^k\mathcal O_{M,\varepsilon}(v_j,x_j;U)
          \,\dd\mu_{\ell_{M,\varepsilon},\beta_{M,\varepsilon}}(U).
 \label{eq:finite-moment}
\end{equation}
This is well defined by Proposition~\ref{prop:finite-measure}.
For \(k=0\), the empty product and its moment equal one.

For a site \(x\), let \(\widetilde x\) be its centered integer
representative, with each coordinate in
\((-\ell_{M,\varepsilon}/2,\ell_{M,\varepsilon}/2]\).
Its physical position is \(\varepsilon\widetilde x\).
This also fixes the seam convention when the period is even.
The separated configuration space is
\begin{equation}
 \operatorname{Conf}_k(\R^4)
 =\{(z_1,\ldots,z_k):z_i\ne z_j\text{ whenever }i\ne j\}.
\end{equation}
For \(f\in C_c^\infty(\operatorname{Conf}_k(\R^4);\C)\),
extended by zero, define
\begin{equation}
\begin{split}
 S_{M,\varepsilon}(v;f)
 =\varepsilon^{4k}\sum_{x\in\Lambda_{\ell_{M,\varepsilon}}^k}
 &f\bigl(\varepsilon\widetilde{x_1},\ldots,
                    \varepsilon\widetilde{x_k}\bigr)\\
 &{}\cdot C_{M,\varepsilon}(v;x).
\end{split}
\label{eq:smear}
\end{equation}
The factor \(\varepsilon^{4k}\) is the product cell volume for \(k\)
four-dimensional insertion points.

An admissible side length \(M\) is \emph{suitable} for \(f\) when
\begin{equation}
 \supp f\subset\left\{z:
 -\frac{M}2<(z_j)_i
 <\frac{M}2\text{ for every }j,i\right\}.
 \label{eq:suitable}
\end{equation}
Compact support gives suitability for all sufficiently large \(M\):
the finitely many coordinate functions are bounded on the support,
whereas the admissible side lengths are unbounded.

\begin{definition}[Ordered Wilson limit]
\label{def:ordered}
For a strictly increasing sequence of admissible side lengths
\((M_j)\) tending to infinity, the asserted
Wilson limit requires, for each fixed \(v,f\), numbers
\(S_{M'}(v;f)\) such that
\begin{align}
 &\lim_{\varepsilon\downarrow0}S_{M',\varepsilon}(v;f)
        =S_{M'}(v;f)
        &&\text{for every suitable fixed }M', \label{eq:row-limit}\\
 &\lim_{j\to\infty}S_{M_j}(v;f)
        =\int f(z)S(v;z)\,\dd z. \label{eq:volume-limit}
\end{align}
Here \(S\) denotes the unselected continuum correlation,
before a particular continuum phase is chosen.
\end{definition}

The first limit is a full mesh limit at each suitable fixed period.
The second may refine a requested volume sequence, as in
\eqref{eq:public-target}. A simultaneous diagonal limit would not
by itself give \eqref{eq:row-limit}.

\section{The continuum conclusions to be explained}
\label{sec:continuum}

This section specifies the continuum data and all conclusions of
the construction theorem. The subsequent sections construct them;
Section~\ref{sec:construction-assembly} assembles the proof with
the same family and the stated order of limits.

\subsection{Fields and phases}

The continuum hierarchy \(T\) contains a finite-dimensional complex
normed space \(F_d\) for each degree \(d\in\N\).
A tagged field is a pair \((d,v)\) with \(v\in F_d\).
There is a surjective map
\(\pi:\mathcal V\to\bigsqcup_{d\ge0}F_d\) from lattice labels
to tagged fields. For a finite list of lattice labels, the notation
\(S(v;z)\) in \eqref{eq:volume-limit} means the correlation with
the corresponding tagged fields \(\pi(v_i)\).
Each \(F_d\) embeds by an injective complex linear map
\(\iota_d:F_d\to\mathcal L\) into the Wightman label space
\(\mathcal L\), and \(\bigcup_{d\ge0}\iota_d(F_d)=\mathcal L\).

The phase space is a compact metric space \(\Sigma\), equipped
with a Borel probability measure \(\rho\) of full support and a
selected point \(\sigma_*\in\Sigma\).
Full support means that every nonempty open set has positive measure.
For a finite list of tagged fields \(a\) and separated locations \(z\),
write \(S^\sigma(a;z)\) for the correlation in phase \(\sigma\),
and \(S(a;z)\) for the unselected correlation. These correlations
are smooth in the separated locations and invariant under simultaneous
permutation of the labels and locations. They satisfy
\begin{equation}
 \Sigma\ni\sigma\longmapsto S^\sigma(a;z)
       \text{ is continuous},\qquad
 S(a;z)=\int_\Sigma S^\sigma(a;z)\,\dd\rho(\sigma).
\end{equation}

The selected phase supplies a complex Hilbert space, a common
dense invariant domain, a normalised cyclic vacuum, and tempered
smeared fields. Temperedness means continuity in a Schwartz test,
with the state in the domain fixed. Cyclicity means that finite
words of smeared fields applied to the vacuum span the domain.
Faithfulness means that each nonzero label gives a nonzero smeared
operator for some test.

To specify the spacetime terminology, write a point as
\(x=(\mathbf x,t)\in\R^3\times\R\), with time the final coordinate.
The closed forward light cone is
\(\{(\mathbf x,t):t\ge|\mathbf x|\}\).
A separation is spacelike when \(t^2-|\mathbf x|^2<0\).
The proper orthochronous Poincare group consists of affine maps
\(x\mapsto A x+b\), where \(A\) preserves this quadratic form,
has determinant one, and preserves the forward cone.

The remaining Wightman conditions require the joint
energy--momentum spectral measure to be supported in that cone;
uniqueness of the translation-invariant vacuum up to scalar;
covariance under the full group just defined; the smeared adjoint
identity; and commutativity at spacelike separated supports.
The hierarchy also has analytic continuation in complex times with
strictly ordered imaginary parts. Its ordered imaginary-time values
recover the selected Euclidean correlations; its smeared boundary
values recover Wightman vacuum moments.

\subsection{Separated OPE}

An operator product expansion (OPE) replaces a cluster of nearby
insertions by single insertions at its center, with a controlled
remainder. Here it is formulated in correlations with all spectators
retained.

For a finite family of tags \(a_i=(d_i,v_i)\), put
\(\operatorname{weight}(a)=\sum_i d_i\).
For each degree \(j\), let \(I_j\) be a finite index set and let
\(A_{j,\alpha}\in F_j\), \(\alpha\in I_j\), be the fields used
in the expansion. Write their cluster coefficients as
\(c_a^{j,\alpha}(x;o)\).

The cluster geometry is specified by
\begin{equation}
 |x_i-o|\le Kr,\qquad |x_i-x_{i'}|\ge\eta r\quad(i\ne i').
 \label{eq:cluster-geometry}
\end{equation}
The spectator locations obey
\begin{equation}
 |o-y_b|\ge c,\qquad |y_b-y_{b'}|\ge c\quad(b\ne b').
 \label{eq:spectator-geometry}
\end{equation}
Let \(\mathcal S\) stand for either \(S\) or one of the phase
correlations \(S^\sigma\). The notation
\(\mathcal S(a,x;b,y)\) means this single correlation with all
cluster and spectator insertions enumerated together.

For fixed \(K\ge0\), \(\eta,\eps>0\), \(j,\alpha\), and all
sufficiently small \(r>0\), the target requires
\begin{equation}
 |c_a^{j,\alpha}(x;o)|
       \le r^{j-\operatorname{weight}(a)-\eps}.
 \label{eq:ope-coefficient}
\end{equation}
For fixed cluster and spectator labels, \(K\ge0\),
\(\eta,c,\eps>0\), and every finite degree \(D\), it requires
\begin{align}
 \Big|\mathcal S(a,x;b,y)
 &-\sum_{j=0}^{D}\sum_{\alpha\in I_j}
       c_a^{j,\alpha}(x;o)\notag\\
 &\quad\cdot\mathcal S\bigl((j,A_{j,\alpha}),o;b,y\bigr)
       \Big|
       \le r^{D+1-\operatorname{weight}(a)-\eps}.
 \label{eq:ope-remainder}
\end{align}
The single tag at \(o\) is understood as a one-element cluster.
Both estimates apply to the geometries
\eqref{eq:cluster-geometry} and, for the remainder,
\eqref{eq:spectator-geometry}.
The threshold for \(r\) may depend on the fixed labels,
separation parameters, \(D,\eps\), but then works for all indicated
configurations and all phases.
The same coefficients precede both \(D\) and the phase choice.
These are finite-truncation estimates; no limit as \(D\to\infty\)
is asserted.

\subsection{The trace test and the construction theorem}

The hierarchy has a centered degree-four field \(\Theta\in F_4\).
Centering means that each smeared vacuum one-point expectation
is zero. Write its selected-phase two-point correlation as
\begin{equation}
 \tau(r)=S^{\sigma_*}\bigl((4,\Theta),(4,\Theta);
                       (\mathbf0,-r/2),(\mathbf0,r/2)\bigr),
 \qquad r>0.
\end{equation}
The required asymptotic is
\begin{equation}
 \lim_{r\downarrow0}r^8(\log r^{-1})^2
          \tau(r)
       =\frac{3\dim_\R\mathfrak g}{\pi^4}>0.
 \label{eq:trace}
\end{equation}
The test also excludes agreement near zero with finite sums of
mass-shell Wick profiles, defined as follows.
For \(m\ge0\), integers \(s,t\ge0\), and a continuous real function
\(h\) on \(\R^3\) satisfying
\(h(\lambda p)=\lambda^s h(p)\) for \(\lambda>0\), set
\begin{equation}
 k_{m,s,t,h}(r)=\int_{\R^3}
       \frac{E_m(p)^t h(p)}{2E_m(p)}
                  e^{-rE_m(p)}\,\dd p,\qquad
 E_m(p)=\sqrt{m^2+|p|^2}.
\end{equation}
The value of the integrand at \(E_m(p)=0\) is immaterial to
the integral. A finite mass-shell Wick profile is a finite
complex linear combination of finite products of these kernels.
The exclusion states that no such profile equals \(\tau(r)\)
for every sufficiently small \(r>0\).
This test does not assert exclusion of every generalised free field.

\begin{theorem}[Joint construction from one Wilson family]
\label{thm:target}
For every closed matrix group \(G\) satisfying
Definition~\ref{def:group}, there is one Wilson family \(C\)
with Wilson regularity, chosen independently of any requested
volume sequence, with the following property.
For every strictly increasing sequence of admissible side lengths
\(M_j\to\infty\), there are a
strictly increasing \(q:\N\to\N\) and one continuum hierarchy,
with all the field, phase and continuation data specified
above, such that:
\begin{enumerate}
 \item The mixed correlations are the ordered Wilson limit
 \eqref{eq:row-limit}--\eqref{eq:volume-limit} along \((M_{q(j)})_j\).
 The inner mesh limit exists at every suitable fixed period.
 \item The same hierarchy has finite-dimensional field spaces,
 a surjective lattice-label projection, injective exhaustive
 maps to its faithful Wightman label space, and a compact
 metric full-support probability phase decomposition.
 Its correlations have the stated smoothness, symmetry and
 character continuity. One selected phase has the specified
 analytic continuation and Wightman reconstruction, with a
 unique translation-invariant vacuum.
 \item One coefficient family gives
 \eqref{eq:ope-coefficient}--\eqref{eq:ope-remainder} for every
 finite truncation, with one remainder threshold for the
 mixture and all phases.
 \item The centred stress trace in that same selected vacuum
 satisfies \eqref{eq:trace} and the finite mass-shell
 Wick-profile exclusion defined above.
\end{enumerate}
The only group assumptions are those in
Definition~\ref{def:group}; the running law, analytic estimates,
phase properties and reconstruction are conclusions of the
construction.
\end{theorem}

The proof is completed in
Section~\ref{sec:construction-assembly}. A positive mass gap
is outside the assertion.

\section{Why RG acts on unnormalised densities}
\label{sec:rg}

The finite probability measure contains a global denominator,
the partition function. An exact recursion transports numerator and
denominator together. The following elementary argument explains
this choice before introducing model-specific effective actions.

\subsection{A normalised blocking kernel}

Let \((\mathcal U_f,H_f)\) and \((\mathcal U_c,H_c)\) be fine
and coarse configuration spaces with probability Haar priors.
A blocking kernel is a nonnegative measurable function \(K(U,V)\)
with
\begin{equation}
 \int_{\mathcal U_c}K(U,V)\,\dd H_c(V)=1
       \quad\text{for every }U.
 \label{eq:kernel-normalisation}
\end{equation}
It describes a random coarse configuration conditioned on a fine
one. For an integrable complex-valued \(F\), set
\begin{equation}
 (\mathcal PF)(V)=\int_{\mathcal U_f}
                      K(U,V)F(U)\,\dd H_f(U).
 \label{eq:block-push}
\end{equation}
Under these assumptions the integral exists for almost every \(V\).
If \(K\) is bounded, it exists for every \(V\), because
\(|K(U,V)F(U)|\le\|K\|_\infty|F(U)|\).

\begin{proposition}[Mass preservation and \(L^1\) control]
\label{prop:push}
For \(F\in L^1(H_f)\), \(\mathcal PF\in L^1(H_c)\) and
\begin{equation}
 \|\mathcal PF\|_{L^1(H_c)}\le\|F\|_{L^1(H_f)},\qquad
 \int\mathcal PF\,\dd H_c=\int F\,\dd H_f.
 \label{eq:push-mass}
\end{equation}
\end{proposition}
\begin{proof}
Tonelli's theorem gives
\begin{align}
 \int_{\mathcal U_c}\int_{\mathcal U_f}
       K(U,V)|F(U)|\,\dd H_f(U)\dd H_c(V)
 &=\int_{\mathcal U_f}|F(U)|
       \left(\int_{\mathcal U_c}K(U,V)\,\dd H_c(V)\right)
                                      \dd H_f(U)\notag\\
 &=\int_{\mathcal U_f}|F(U)|\,\dd H_f(U).
\end{align}
Thus the joint integrand is absolutely integrable.
The triangle inequality inside the \(U\)-integral proves the norm
bound. Fubini's theorem, justified by this absolute integrability,
and \eqref{eq:kernel-normalisation} prove the integral identity.
\end{proof}

\begin{corollary}[One denominator for a transported observable]
\label{cor:observable}
Suppose \(F\in L^1(H_f)\), \(F\ge0\), \(Z=\int F\,\dd H_f>0\),
and \(OF\in L^1(H_f)\). Then
\begin{equation}
 \frac{\int\mathcal P(OF)\,\dd H_c}
      {\int\mathcal PF\,\dd H_c}
       =\frac{\int OF\,\dd H_f}{Z}.
 \label{eq:source-push}
\end{equation}
\end{corollary}
\begin{proof}
Apply \eqref{eq:push-mass} to \(OF\) and \(F\), respectively.
\end{proof}

For several insertions, \(O\) here is their entire product.
Its transport generally cannot be obtained by multiplying
separately transported single insertions.
We must therefore transport products jointly, together with
the vacuum density. Dividing individual insertions by
different partition functions would lose \eqref{eq:source-push}.

\subsection{Exact recursion at fixed cutoff and volume}

Fix \(n,M\); let \(m=n,n-1,\ldots,0\) label intermediate scales,
with larger \(m\) meaning a finer lattice.
We seek unnormalised densities \(F_m\) satisfying
\begin{equation}
 F_n(U)=e^{-\beta_{n,M}W_{\ell_{n,M}}(U)},\qquad
 F_m=\mathcal P_mF_{m+1}.
 \label{eq:exact-recursion}
\end{equation}
The kernel can depend on the running coupling; its normalisation
still gives Proposition~\ref{prop:push} at each fixed step.

\begin{proposition}[Partition function at every scale]
\label{prop:history-mass}
If the densities in \eqref{eq:exact-recursion} are integrable,
then
\begin{equation}
 \int F_m\,\dd H_m=Z_{\ell_{n,M}}(\beta_{n,M})
                   \quad(0\le m\le n).
 \label{eq:history-partition}
\end{equation}
\end{proposition}
\begin{proof}
At \(m=n\) the identity is \eqref{eq:partition}.
If it holds at \(m+1\), then
\begin{equation}
 \int F_m\,\dd H_m
 =\int\mathcal P_mF_{m+1}\,\dd H_m
 =\int F_{m+1}\,\dd H_{m+1}
 =Z_{\ell_{n,M}}(\beta_{n,M}).
\end{equation}
Descending induction completes the proof.
\end{proof}

The construction must represent these densities by local
coefficients controlled under this exact operation. Mass preservation
identifies their measure; it does not estimate those coefficients.

\section{Running coupling: a detailed bookkeeping calculation}
\label{sec:coupling}

We now expand a consequence of the one-step estimates.
A small relative error per logarithmic scale suffices;
absolute summability over RG steps is not required.

\subsection{Scale conventions}

Choose odd block lengths \(L_m\ge3\) and define
\begin{equation}
 {\varepsilon}_0=1,\quad {\varepsilon}_{m+1}=\frac{{\varepsilon}_m}{L_m},\qquad
 s_m=\log {\varepsilon}_m^{-1}
                   =\sum_{i=0}^{m-1}\log L_i.
 \label{eq:scale}
\end{equation}
Thus \(s_m\to\infty\).
Let \(b_0>0\) be the leading coefficient of the running coupling,
and fix a positive coarse reference coupling \(\overline\beta\).
Set
\begin{equation}
 \widehat\beta_m=\overline\beta+b_0s_m.
 \label{eq:hat-beta}
\end{equation}

Write \(\beta_{n,M,m}\) for the intermediate coupling.
Its fine value \(\beta_{n,M,n}\) is the bare coefficient
\(\beta_{n,M}\) of the Wilson measure.
The construction fixes \(\beta_{n,M,0}=\overline\beta\).

\begin{proposition}[Uniform accumulation of drift]
\label{prop:drift}
Suppose nonnegative \(\delta_m\) satisfy
\begin{align}
 &|\beta_{n,M,m+1}-\beta_{n,M,m}-b_0\log L_m|
        \le\delta_m &&(m<n,\text{ every }M), \label{eq:drift-step}\\
 &\beta_{n,M,0}=\overline\beta,\qquad
       \frac{\delta_m}{\log L_m}\longrightarrow0.
       \label{eq:drift-small}
\end{align}
For every \(\eps>0\), there is \(m_0\) such that
\begin{equation}
 |\beta_{n,M,m}-\widehat\beta_m|
       \le\eps s_m
 \quad(m\ge m_0,\ n\ge m,\text{ every }M).
 \label{eq:running-internal}
\end{equation}
Consequently,
\begin{equation}
 \frac{\beta_{n,M,n}}{\log {\varepsilon}_n^{-1}}\longrightarrow b_0
                \quad\text{uniformly in }M.
 \label{eq:bare-running}
\end{equation}
\end{proposition}
\begin{proof}
Subtract the increment of \eqref{eq:hat-beta} from
\eqref{eq:drift-step}. Equality at \(m=0\) and repeated use
of the triangle inequality give
\begin{equation}
 |\beta_{n,M,m}-\widehat\beta_m|
                    \le\sum_{i=0}^{m-1}\delta_i.
 \label{eq:drift-sum}
\end{equation}
The upper bound contains neither \(n\) nor \(M\).

Fix \(\eps>0\). Choose \(j\) such that
\(\delta_i\le(\eps/2)\log L_i\) for every \(i\ge j\).
For \(m\ge j\), split the sum:
\begin{align}
 \sum_{i=0}^{m-1}\delta_i
 &\le\sum_{i=0}^{j-1}\delta_i
            +\frac\eps2\sum_{i=j}^{m-1}\log L_i\notag\\
 &=\sum_{i=0}^{j-1}\delta_i+
    \frac\eps2\bigl(s_m
                         -s_j\bigr).
 \label{eq:prefix-tail}
\end{align}
The first sum is finite and fixed. Increase \(m_0\ge j\) until
\(\sum_{i<j}\delta_i\le(\eps/2)s_m\)
for every \(m\ge m_0\).
Since \(s_j\ge0\), this proves
\eqref{eq:running-internal}.

Set \(m=n\) and divide by \(s_n>0\).
The remaining deterministic discrepancy is
\(\overline\beta/s_n\to0\).
This proves \eqref{eq:bare-running}, with a threshold independent
of \(M\).
\end{proof}

The drift bounds will be arranged by choosing a stage index \(J\)
and requiring
\begin{equation}
 \frac{\delta_m}{\log L_m}\le J(m+1)^{-2},\qquad
                                  J(m)\to\infty.
 \label{eq:clock-drift}
\end{equation}
Here \(J(m)\) is a separate stage index controlling the amount
of analytic regularity requested. It may remain constant for
many consecutive RG steps. Equation \eqref{eq:clock-drift}
implies \eqref{eq:drift-small}; the argument does not require
\(\sum_m J(m)^{-2}<\infty\).

\clearpage
\section{Effective densities and the small-field decomposition}
\label{sec:activities}

Fix one lattice \(\Lambda=\Lambda_\ell\), and suppress the cutoff and
volume indices temporarily. The immediate task is to represent a blocked
density by local functions whose bounds can be propagated to the next
scale. The representation must cover every configuration, including
those on which a local expansion about the identity is not valid.

We shall keep the rapidly varying principal action in a reference
function \(V(U)\), place regular local corrections in an exponential,
and represent the remaining contributions by compatible products of
activities. An activity consists of a label, a finite footprint, and a
function of the configuration. Different labels may have the same
footprint. None of the summations below identifies such labels.

For comparison, Ba{\l}aban's inductive density representation separates
small-field characteristic functions from operations on large-field
components; see \cite[Sect.~2, eqs.~(2.17)--(2.19),
printed pp.~257--258; setup on pp.~254--256]{Bal88b}.
This is a comparison of the organisation. The masks, labelled activities,
and derivative norms used here are defined below.

\subsection{A gauge-invariant measure of plaquette size}

On the matrix Lie algebra use the real inner product
\begin{equation}
 \langle X,Y\rangle_{\mathfrak g}
       =-\frac1r\operatorname{Re}\operatorname{tr}(XY),\qquad
 |X|_{\mathfrak g}=\frac{\|X\|_{\mathrm F}}{\sqrt r},
 \label{eq:physical-lie-norm}
\end{equation}
where \(\|\cdot\|_{\mathrm F}\) is the Frobenius matrix norm.
The matrices in \(\mathfrak g\) are skew Hermitian: differentiation of
\(\exp(tX)^*\exp(tX)=1\) at \(t=0\) gives \(X^*+X=0\).
Thus the two expressions for the norm agree, and conjugation by
\(G\) preserves it.

For \(W\in G\), define its normalised chord from the identity by
\begin{equation}
 d_G(W)=\frac{\|W-1\|_{\mathrm F}}{\sqrt r}.
 \label{eq:group-chord}
\end{equation}
Unitarity gives the useful exact identity
\begin{equation}
 d_G(W)^2
 =\frac1r\operatorname{Re}\operatorname{tr}
                   \bigl((W-1)^*(W-1)\bigr)
 =2\left(1-\frac1r\operatorname{Re}\operatorname{tr}W\right).
 \label{eq:chord-energy}
\end{equation}
In particular \(0\le d_G(W)\le2\), and \(d_G\) is continuous
and conjugation invariant.

Fix \(0<h\le1\). A plaquette is bad when
\(d_G(P_{ij}(U;x))>h\). Let \(\mathcal B_h(U)\) be the set of
all vertices incident to at least one bad plaquette. For a finite
set \(X\subset\Lambda\), put
\begin{equation}
 \chi_X(U)=\mathbf1_{\{X\cap\mathcal B_h(U)=\varnothing\}}.
 \label{eq:good-mask}
\end{equation}
Thus \(X\) is small-field precisely when \(\chi_X(U)=1\).
This tests plaquettes incident to \(X\); it is not a bound on the
individual links \(U(x,i)\). A gauge transformation can change those
links substantially while leaving every plaquette size unchanged.

\begin{lemma}[Measurability and symmetry of the decomposition]
\label{lem:mask-measurable}
The set \(\{U:x\in\mathcal B_h(U)\}\) is open for each site \(x\).
The set \(\{U:\chi_X(U)=1\}\) is closed. Exact flag events
\(\{U:\mathcal B_h(U)=B\}\), \(B\subset\Lambda\), are Borel.
All these events are gauge invariant.
\end{lemma}
\begin{proof}
For fixed \(x\), the bad-site event is a finite union, over
plaquettes incident to \(x\), of sets
\(\{U:d_G(P_{ij}(U;y))>h\}\). Each is open by continuity.
The good event is the finite intersection of the complements of
the bad-site events for \(x\in X\). An exact flag event is the
intersection of the bad-site events for \(x\in B\) and their
complements for \(x\notin B\).
Gauge invariance follows from \eqref{eq:gauge-plaquette} and
conjugation invariance of \(d_G\).
\end{proof}

Equality at the threshold belongs to the small-field region.
The functions \(\chi_X\) are indicators and need not be differentiable.
The derivative norms below will therefore be imposed on the local
function \emph{before} multiplication by \(\chi_X\).

\begin{proposition}[Wilson action cost of bad sites]
\label{prop:bad-cost}
For every configuration,
\begin{equation}
 W_\ell(U)\ge\frac{h^2}{8}|\mathcal B_h(U)|.
 \label{eq:bad-cost}
\end{equation}
Consequently, for \(\beta\ge0\) and \(0\le\theta\le1\),
\begin{equation}
 e^{-\beta W_\ell(U)}
 \le e^{-\theta\beta h^2|\mathcal B_h(U)|/8}
                    e^{-(1-\theta)\beta W_\ell(U)}.
 \label{eq:split-bad-cost}
\end{equation}
\end{proposition}
\begin{proof}
Let \(n_{\rm bad}\) be the number of bad plaquettes. Each has
energy \(d_G(P)^2/2>h^2/2\), by \eqref{eq:chord-energy}.
Every other plaquette has nonnegative energy. Also a plaquette
has at most four distinct vertices, so
\begin{equation}
 |\mathcal B_h(U)|\le4n_{\rm bad},\qquad
 W_\ell(U)\ge\frac{h^2}{2}n_{\rm bad}
              \ge\frac{h^2}{8}|\mathcal B_h(U)|.
\end{equation}
This includes the case \(n_{\rm bad}=0\). Split \(\beta W_\ell\)
into the fractions \(\theta\) and \(1-\theta\), and apply the
first inequality to the former.
\end{proof}

The constant \(1/8\) is independent of \(\ell,n,M\).
Equation \eqref{eq:split-bad-cost} is an estimate for the
unnormalised density. Passing to a probability estimate also involves
its partition function. In the RG argument the retained fraction of
the action is needed to perform that comparison.

\subsection{Footprints and the links actually read}

Say that \(U,V\) agree on \(X\) when they agree on every link
whose two endpoints belong to \(X\). A function is local on \(X\)
if such agreement determines its value. To determine the flags
in \(X\), one must also read plaquettes incident to \(X\).
Define
\begin{equation}
 X^+=X\cup\bigcup_{\substack{\text{plaquettes }f\\
                            \operatorname{Vert}(f)\cap X\ne\varnothing}}
                                      \operatorname{Vert}(f),
 \label{eq:inspected-footprint}
\end{equation}
where \(\operatorname{Vert}(f)\) is the vertex set of \(f\).

\begin{lemma}[Flags enlarge the read footprint]
\label{lem:flag-footprint}
If \(U,V\) agree on \(X^+\), then their flag sets agree in \(X\).
Consequently, a function determined by the links in \(X\) and the
flag set \(X\cap\mathcal B_h(U)\) is local on \(X^+\).
\end{lemma}
\begin{proof}
Every plaquette incident to \(X\) has all its vertices in \(X^+\).
Agreement on \(X^+\) gives equality of all its link factors, hence
of its plaquette matrix and its badness test. This determines all
flags in \(X\). It also gives agreement on the smaller set \(X\),
which proves the last assertion.
\end{proof}

We shall use \(X\) as a combinatorial footprint when testing whether
activities overlap. Its physical read footprint may be \(X^+\).
Disjoint combinatorial footprints alone do not imply independence of
the corresponding functions under Haar measure. Whenever an integral
is factored, the actual links read must be considered.

\subsection{Regular and forced activities}

Regular activities are a countable family of real functions \(E_i\),
with a footprint \(X_i\subset\Lambda\) and an anchor \(o_i\in X_i\).
Each \(E_i\) is Borel, gauge invariant, and local on \(X_i\).
It is smooth near every configuration good on \(X_i\), in the
precise finite-order sense defined in Section~\ref{sec:activity-norms}.
Their contribution to the action is
\begin{equation}
 E^\#(U)=\sum_i\chi_{X_i}(U)E_i(U).
 \label{eq:regular-action}
\end{equation}

A forced activity has a label \(\alpha\), a nonempty footprint
\(Y_\alpha\), a nonempty inventory \(B_\alpha\subset Y_\alpha\),
and an additional Borel, gauge-invariant test \(Q_\alpha(U)\).
The inventory specifies which bad sites this activity accounts for.
The real Borel function \(K_\alpha\) is gauge invariant and satisfies
\begin{equation}
 K_\alpha(U)\ne0\ \Longrightarrow\
 Y_\alpha\cap\mathcal B_h(U)=B_\alpha
                     \quad\text{and}\quad Q_\alpha(U).
 \label{eq:forced-eligibility}
\end{equation}
Both \(K_\alpha\) and \(Q_\alpha\) must be determined by the
links in \(Y_\alpha\) and the flags in \(Y_\alpha\).
Lemma~\ref{lem:flag-footprint} then makes them local on \(Y_\alpha^+\).

A finite set \(A\) of forced labels is admissible at \(U\) if
\begin{equation}
 Y_\alpha\cap Y_{\alpha'}=\varnothing\quad(\alpha\ne\alpha'\in A),
 \qquad \bigcup_{\alpha\in A}B_\alpha=\mathcal B_h(U).
 \label{eq:admissible-forced}
\end{equation}
Define its gas factor by
\begin{equation}
 Z_K(U)=\sum_{\substack{A\text{ finite}\\A\text{ admissible at }U}}
                                      \prod_{\alpha\in A}K_\alpha(U).
 \label{eq:forced-gas}
\end{equation}
Absolute convergence will follow from the norms below. The sum is
over finite sets of labels, so a label is not repeated.
The exact coverage condition in \eqref{eq:admissible-forced}
ensures that every bad site is accounted for.

\begin{lemma}[The all-good sector]
\label{lem:all-good-gas}
If \(\mathcal B_h(U)=\varnothing\), then \(Z_K(U)=1\).
\end{lemma}
\begin{proof}
The empty label set is admissible and contributes one.
Every nonempty label set has a nonempty union of inventories,
so cannot satisfy the coverage condition with an empty flag set.
\end{proof}

\begin{definition}[Effective density representation]
\label{def:effective-density}
Given a bounded real Borel reference action \(V\), a scalar
\(v\in\R\), and the activities above, the represented function is
\begin{equation}
 F(U)=
 \exp\bigl(v\#\Lambda-V(U)-E^\#(U)\bigr)Z_K(U).
 \label{eq:effective-density}
\end{equation}
The reference action contains the principal terms treated
nonperturbatively. For the original Wilson weight one may take
\(V=\beta W_\ell\); the reference action after blocking must be
constructed together with its local corrections.
\end{definition}

An identity with the blocked function \(F_m=\mathcal P_mF_{m+1}\)
is part of the RG step, not a consequence of this definition.
Individual forced activities may have either sign. Positivity of an
actual density follows from its exact identification with the push of
the positive Wilson weight. The activity bounds by themselves will
give integrability, not positivity.

\section{Local derivative norms and absolute activity bounds}
\label{sec:activity-norms}

\subsection{A geometric weight for footprints}

Let \(d_\infty(x,y)\) be the periodic maximum distance on
\((\Z/\ell\Z)^4\), computed using the smallest absolute coordinate
representatives. Sites are adjacent when their distance is at most one.
Fix an integer \(B\ge1\). For \(D\subset\Lambda\), let
\begin{equation}
 H_B(D)=\bigcup_{x\in D}\{y:d_\infty(x,y)\le B\}.
\end{equation}
For nonempty \(X\), define \(n_B(X)\) to be the smallest number
of centers in a set \(D\) for which \(H_B(D)\) is connected and
contains \(X\). Connectivity refers to the adjacency just specified.
Set \(n_B(\varnothing)=0\). Such a cover exists: take \(D=\Lambda\).
The centers need not belong to \(X\), and the union of their balls,
rather than the set of centers alone, must be connected.

Put \(D_B=(2B+1)^4\). Two facts will be used repeatedly:
\begin{align}
 |X|&\le D_B n_B(X), \label{eq:cover-cardinality}\\
 n_B\left(\bigcup_{i\in A}X_i\right)
   &\le\sum_{i\in A}n_B(X_i)
       \quad\text{if the overlap graph of }(X_i)_{i\in A}
                       \text{ is connected}. \label{eq:cover-connected}
\end{align}
The overlap graph has one vertex for each label and an edge when
the two footprints intersect.
For \eqref{eq:cover-cardinality}, a periodic \(B\)-ball has at most
\(D_B\) sites. For \eqref{eq:cover-connected}, choose minimal
connected covers of the individual sets. Covers of intersecting
sets intersect, so their union is connected along every path in
the overlap graph. It is a cover with at most the sum of the
numbers of centers. In particular the constants do not grow with
the torus volume.

\subsection{Left derivatives on the configuration space}

Fix integers \(q\ge6\), \(s\ge2\), and a support exponent
\(\kappa>0\). A direction is an assignment
\(A=(A_e)_e\in\mathfrak g^{\Lambda\times\{1,2,3,4\}}\).
For an unmasked function \(f\), define
\begin{equation}
 D_Af(U)=\left.\frac{\dd}{\dd t}
       f\bigl((e^{tA_e}U_e)_e\bigr)\right|_{t=0}.
 \label{eq:left-derivative}
\end{equation}
Higher derivatives are ordered compositions
\(D_{A_1}\cdots D_{A_k}f\); their order is retained.
Let \(\mathcal D_o\) be the set of directions obeying
\begin{equation}
 |A_{(x,i)}|_{\mathfrak g}
       \le h(1+d_\infty(x,o))^s
                  \quad\text{for all links }(x,i).
 \label{eq:weighted-directions}
\end{equation}
The factor \(h\) scales derivatives with the small-field amplitude.
The polynomial weight permits a direction to grow away from the
anchor while retaining control on a local footprint.

For \(f\) local on \(X\), define
\begin{equation}
 \|f\|_{q,h,s;o,X}
 =\max_{0\le k\le q}
   \sup_{\substack{U:\,\chi_X(U)=1\\A_1,\ldots,A_k\in\mathcal D_o}}
       |D_{A_1}\cdots D_{A_k}f(U)|.
 \label{eq:local-derivative-norm}
\end{equation}
At \(k=0\) the expression is \(|f(U)|\).
The required regularity is the existence of a \(C^q\) germ of
\(A\mapsto f((e^{A_e}U_e)_e)\) near \(A=0\) at every good \(U\),
including configurations with plaquette size exactly \(h\).
These germs must yield the ordered derivatives in
\eqref{eq:left-derivative}.

Only the unmasked function is differentiated.
For example, \(\chi_X f\) can jump as a plaquette crosses \(h\),
even when \(f\) is smooth. The norm
\eqref{eq:local-derivative-norm} makes no assertion of a derivative
bound across that jump.

\subsection{Local normalisation conditions}

The identity configuration has every link equal to \(1\in G\).
Regular corrections are normalised by \(E_i(1)=0\).
There is also a condition removing the local quadratic contribution
that belongs in the principal action.

For an anchor \(o\), let \(\widetilde{x-o}\) be the centered integer
lift of a site displacement. A direction is affine about \(o\) if
\begin{equation}
 A_{(x,i)}=a_i+\sum_{j=1}^4\widetilde{x-o}_j b_{ij},
                    \qquad a_i,b_{ij}\in\mathfrak g.
 \label{eq:affine-direction}
\end{equation}
A footprint \(X\) is nonwrapping about \(o\) when
\(10d_\infty(x,o)<\ell\) for every \(x\in X\).
On such a footprint these coordinates describe one local chart on
the cover without crossing a torus seam.
For two directions define the exponential-chart Hessian by
\begin{equation}
 \mathcal H_f(A,C)=
 \left.\partial_t\partial_u
       f\bigl((\exp(tA_e+uC_e))_e\bigr)\right|_{t=u=0}.
 \label{eq:chart-hessian}
\end{equation}
The additional normalisation is
\begin{equation}
 \mathcal H_{E_i}(A,C)=0
 \quad\text{for affine }A,C
 \text{ whenever }X_i\text{ is nonwrapping about }o_i.
 \label{eq:affine-subtraction}
\end{equation}
This is a renormalisation condition to be achieved when extracting
the principal terms. It does not follow merely from gauge invariance.
Nor does it say that the Hessian vanishes in all directions.
The chart Hessian and the ordered left derivatives are distinct
notions and are not interchanged here.

\subsection{The two activity norms}

With \(q,h,s,B\) fixed, define
\begin{align}
 \|E\|_{\mathcal E,\kappa}
 &=\sup_{x\in\Lambda}
     \sum_{i:o_i=x}e^{\kappa n_B(X_i)}
                    \|E_i\|_{q,h,s;o_i,X_i}, \label{eq:E-norm}\\
 \|K\|_{\mathcal K,\kappa}
 &=\sup_{x\in\Lambda}
     \sum_{\alpha:x\in Y_\alpha}e^{\kappa n_B(Y_\alpha)}
                    \|K_\alpha\|_\infty .
 \label{eq:K-norm}
\end{align}
The sums are sums of nonnegative numbers and may initially be
infinite. A finite bound includes their summability.
The first norm sums at a specified anchor; the second charges every
site in the footprint. This difference is part of the definition.
Regular activities already carry anchors and derivative information;
forced activities must also pay for every possible overlap.

Constants in an estimate may depend on the fixed profile
\((q,s,\kappa)\), the group, and specified block parameters.
Uniformity in cutoff and volume means that these constants do not
depend on \(n,M\). Increasing the derivative order is a separate
change of profile.

\begin{proposition}[Absolute convergence on a finite torus]
\label{prop:activity-convergence}
Suppose \(\|E\|_{\mathcal E,\kappa}\le\delta\) and
\(\|K\|_{\mathcal K,\kappa}\le w\).
Then both \eqref{eq:regular-action} and \eqref{eq:forced-gas}
converge absolutely, uniformly in \(U\), and are Borel functions.
Moreover,
\begin{equation}
 |E^\#(U)|\le\#\Lambda\delta,\qquad
 \sum_{\substack{A\text{ finite}\\A\text{ admissible at }U}}
          \prod_{\alpha\in A}|K_\alpha(U)|
                         \le e^{{\#\Lambda}w}.
 \label{eq:activity-global-bounds}
\end{equation}
For bounded \(V\), the function in \eqref{eq:effective-density}
is bounded and Haar integrable.
\end{proposition}
\begin{proof}
Put \(e_i=\|E_i\|_{q,h,s;o_i,X_i}\).
On the good set the order-zero bound gives \(|E_i(U)|\le e_i\);
off it the masked value is zero.
Since the weights are at least one,
\begin{equation}
 \sum_i e_i
 =\sum_{x\in\Lambda}\sum_{i:o_i=x}e_i
 \le\#\Lambda\delta.
 \label{eq:E-total}
\end{equation}
For \(k_\alpha=\|K_\alpha\|_\infty\), nonemptiness of
\(Y_\alpha\) gives
\begin{equation}
 \sum_\alpha k_\alpha
 \le\sum_{x\in\Lambda}\sum_{\alpha:x\in Y_\alpha}k_\alpha
 \le{\#\Lambda}w.
 \label{eq:K-total}
\end{equation}
For a finite label set \(J\), expansion of a product yields
\begin{equation}
 \sum_{A\subset J}\prod_{\alpha\in A}k_\alpha
     =\prod_{\alpha\in J}(1+k_\alpha)
     \le\exp\left(\sum_{\alpha\in J}k_\alpha\right).
\end{equation}
Increase \(J\) to the countable label set.
Monotone convergence proves summability of all finite-subset
products with bound \(e^{{\#\Lambda}w}\).
This summable majorant is independent of \(U\), so the
Weierstrass test proves uniform absolute convergence for the gas,
even after the admissibility indicators are inserted.
The argument for \(E^\#\) uses \eqref{eq:E-total}.
Each summand is Borel by Lemma~\ref{lem:mask-measurable},
so their limits are Borel.
Finally,
\begin{equation}
 |F(U)|\le
 \exp\bigl(|v|\#\Lambda+\|V\|_\infty
                         +\#\Lambda(\delta+w)\bigr),
\end{equation}
which is integrable against probability Haar measure.
\end{proof}

The extensive factors \(\#\Lambda\) in
\eqref{eq:activity-global-bounds} are appropriate for a total
action. The constants \(\delta,w\) in the local norms are the
quantities to control uniformly as the volume grows.

A long product should be estimated without paying a new
factor \(2^q\) at every multiplication. The useful idea is to
keep the derivative orders separately until the product has
been formed. This is the elementary version of the normed-algebra
organisation of Brydges and Slade
\cite[Proposition~3.7, PDF p.~15; Sect.~5.1, PDF p.~24]{BS15a}.
Their test-function seminorm is defined for a different field
algebra. We use only its product principle and prove the
version for our ordered left derivatives.

The \emph{finite jet} means the function and its derivatives
through order \(q\). Introduce the auxiliary seminorm
\begin{equation}
 \|f\|^{\rm alg}_{q,h,s;o,X}
 =\sum_{j=0}^q\frac1{j!}
   \sup_{\substack{U:\chi_X(U)=1\\
                   A_1,\ldots,A_j\in\mathcal D_o}}
           |D_{A_1}\cdots D_{A_j}f(U)|.
 \label{eq:finite-jet-algebra-seminorm}
\end{equation}
The \(j=0\) summand is the supremum of \(|f|\).
This is a temporary tool for estimating products;
the activity norms \eqref{eq:E-norm}--\eqref{eq:K-norm}
are unchanged.

\begin{lemma}[Factorial weights and products]
\label{lem:finite-jet-algebra}
For fixed \(q,h,s,o,X\),
\begin{align}
 \|f\|_{q,h,s;o,X}
 &\le q!\,\|f\|^{\rm alg}_{q,h,s;o,X}
 \le e q!\,\|f\|_{q,h,s;o,X}, \label{eq:finite-jet-equivalence}\\
 \|fg\|^{\rm alg}_{q,h,s;o,X}
 &\le\|f\|^{\rm alg}_{q,h,s;o,X}
                    \|g\|^{\rm alg}_{q,h,s;o,X},\label{eq:finite-jet-product}\\
 \|e^f-1\|^{\rm alg}_{q,h,s;o,X}
 &\le \exp\{\|f\|^{\rm alg}_{q,h,s;o,X}\}-1 .
 \label{eq:finite-jet-exponential}
\end{align}
No derivatives are commuted. All constants are independent
of the cutoff, volume, and number of factors.
\end{lemma}
\begin{proof}
Write \({\mathsf M}_j(f)\) for the order-\(j\) supremum in
\eqref{eq:finite-jet-algebra-seminorm}.
The factorial weights give
\[
 \max_{j\le q}{\mathsf M}_j(f)\le q!\sum_{j=0}^q{\mathsf M}_j(f)/j!,
 \qquad
 \sum_{j=0}^q{\mathsf M}_j(f)/j!\le e\max_{j\le q}{\mathsf M}_j(f).
 \]
Here \(j!\le q!\), and the exponential series bounds
the sum of reciprocal factorials by \(e\).
For a derivative of \(fg\), a subset of its ordered
positions acts on \(f\), and its complement acts on \(g\).
Keeping the order within both subsequences gives
\[
 {\mathsf M}_j(fg)\le\sum_{a+b=j}\frac{j!}{a!b!}{\mathsf M}_a(f){\mathsf M}_b(g).
 \]
Divide by \(j!\), sum up to \(q\), and enlarge the resulting
sum over \(a+b\le q\) to the product of the two nonnegative
sums. This proves \eqref{eq:finite-jet-product}.
Iteration gives \(\|f^n\|^{\rm alg}\le\|f\|_{\rm alg}^n\).
The ordinary scalar exponential of each \(C^q\) germ is
defined on the same neighbourhood. Its Taylor series and
its derivatives through \(q\) converge there; their
seminorms are bounded by
\(\sum_{n\ge1}\|f\|_{\rm alg}^n/n!\).
This proves \eqref{eq:finite-jet-exponential}.
\end{proof}

\begin{lemma}[One anchor change for a whole product]
\label{lem:finite-jet-anchor-product}
Let \(f_i\) be a finite nonempty family of functions
local on nonempty \(X_i\), with
\(o_i\in X_i\), and suppose their overlap graph is connected.
Put \(X=\bigcup_iX_i\) and \(K=\sum_i n_B(X_i)\).
For any \(o\in X\),
\begin{equation}
 \left\|\prod_i f_i\right\|_{q,h,s;o,X}
 \le q!(1+3BK)^{sq}
               \prod_i\|f_i\|^{\rm alg}_{q,h,s;o_i,X_i}.
 \label{eq:finite-jet-anchor-product}
\end{equation}
Consequently, for every \(\eta>0\) and \(\kappa\ge0\),
\begin{equation}
 e^{\kappa n_B(X)}
     \left\|\prod_i f_i\right\|_{q,h,s;o,X}
 \le C_{q,s,\eta}(1+B)^{sq}
     \prod_i e^{(\kappa+\eta)n_B(X_i)}
                       \|f_i\|^{\rm alg}_{q,h,s;o_i,X_i}.
 \label{eq:finite-jet-cover-reserve}
\end{equation}
The power of \(K\) does not grow with the number of factors.
If the output is kept in the algebraic seminorm, both
inequalities hold with the factor \(q!\) omitted.
\end{lemma}
\begin{proof}
The connected radius-\(B\) covers can be united into a
connected cover using at most \(K\) balls. Its diameter
in the periodic metric is at most \(3BK\).
For \(\tau=(1+3BK)^{-s}\), multiplying a direction
allowed at \(o\) by \(\tau\) makes it allowed at every
\(o_i\), since
\[
 1+d(x,o)\le(1+d(x,o_i))(1+d(o_i,o)).
 \]
Fix a field \(U\) good on \(X\). It is good on every
\(X_i\). The same factorial convolution as in
\eqref{eq:finite-jet-product}, at this fixed field, gives
\[
 \sum_{j=0}^q\frac{\tau^j}{j!}
   \sup_{A_1,\ldots,A_j\in\mathcal D_o}
    |D_{A_1}\cdots D_{A_j}(\prod_i f_i)(U)|
       \le\prod_i\|f_i\|^{\rm alg}_{q,h,s;o_i,X_i}.
 \]
The good-field domain is held fixed; only the directions
are multiplied by \(\tau\). In particular the cutoff
threshold has not been changed to \(\tau h\).
Taking the original good-field supremum and using
\(j!\tau^{-j}\le q!\tau^{-q}\) proves
\eqref{eq:finite-jet-anchor-product}.
For an algebraic output, sum the weighted derivative
orders instead and use \(\tau^{-j}\le\tau^{-q}\);
no \(q!\) conversion is then required.
Finally \(n_B(X)\le K\), and
\[
 \sup_{K\ge1}q!(1+3BK)^{sq}e^{-\eta K}
                         \le C_{q,s,\eta}(1+B)^{sq}.
 \]
Distributing \(e^{\eta K}\) over the factors proves
\eqref{eq:finite-jet-cover-reserve}.
\end{proof}

\begin{proposition}[Calculus of the local norm]
\label{prop:local-calculus}
On a fixed footprint with a fixed anchor, the local norm obeys
\begin{align}
 \|f+g\|_{q,h,s;o,X}&\le\|f\|_{q,h,s;o,X}+
                                      \|g\|_{q,h,s;o,X},\\
 \|fg\|_{q,h,s;o,X}&\le2^q\|f\|_{q,h,s;o,X}
                                      \|g\|_{q,h,s;o,X}.
 \label{eq:local-product}
\end{align}
There is a constant \(C_q\), depending only on \(q\), such that
if \(\|f\|_{q,h,s;o,X}\le t\le1\), then
\begin{equation}
 \|e^f-1\|_{q,h,s;o,X}\le C_q t .
 \label{eq:local-mayer}
\end{equation}
\end{proposition}
\begin{proof}
Linearity proves the first inequality. An ordered derivative of
order \(k\) of \(fg\) has one term for each subset of the \(k\)
positions: the derivatives in that subset act on \(f\), in their
original order, and the others act on \(g\).
Each term is bounded by the product of the two norms.
There are \(2^k\le2^q\) terms. This proof uses the Leibniz rule
for derivations, without commuting any left derivatives.

For the exponential use Lemma~\ref{lem:finite-jet-algebra}:
\[
 \|e^f-1\|_{q,h,s;o,X}
       \le q!(e^{et}-1)\le q!(e^e-1)t
                          \qquad(0\le t\le1).
 \]
The last inequality is convexity on \([0,1]\).
This gives an explicit permissible \(C_q\), and avoids
enumerating derivative partitions separately.
\end{proof}

Concatenating two anchored families also gives
\(\|E^{(1)}\cup E^{(2)}\|_{\mathcal E,\kappa}
\le\|E^{(1)}\|_{\mathcal E,\kappa}
+\|E^{(2)}\|_{\mathcal E,\kappa}\).
The normalisations \(E_i(1)=0\) and
\eqref{eq:affine-subtraction} are preserved under addition.
The product estimate alone makes no such assertion about a product
of normalised functions.

\section{A volume-uniform bound for connected families}
\label{sec:connected-bound}

One RG step produces many overlapping local terms. To collect them
into activities we need to sum connected families, with a bound
through each site rather than a bound proportional to the volume.
The incompatibility-graph description of polymer activities is
standard; see Fern\'andez and Procacci \cite[Sect.~2,
eq.~(2.2), PDF p.~3]{FP07}.
We prove one positive majorant with a variable covering reserve.
Its finite-family part sums distinct labels. Its rooted-tree part
allows repeated labels with exponential generating-function weights.
These are different sums, but the same exponential bound controls
both. This is the positive-tree organisation of
\cite[Theorem~1 and Sect.~4]{FP07}; the proof below supplies the
particular reserve and counting conventions used here.

\begin{lemma}[Connected families and rooted trees]
\label{lem:connected-bound}
Let \(L\) be a countable label set with nonempty footprints
\(X_i\subset\Lambda\) and nonnegative weights \(f_i\).
Suppose, for \(\kappa\ge0\), \(a>0\) and \(\gamma\ge0\),
\begin{equation}
 \sup_{x\in\Lambda}
  \sum_{i:x\in X_i}f_i e^{(\kappa+a)n_B(X_i)}
                 \le\gamma,\qquad D_B\gamma\le a.
 \label{eq:connected-input}
\end{equation}
Then, for every \(x\),
\begin{equation}
 \sum_{\substack{\varnothing\ne A\subset L\text{ finite}\\
            A\text{ connected},\ x\in\bigcup_{i\in A}X_i}}
 e^{\kappa n_B(\bigcup_{i\in A}X_i)}
                     \prod_{i\in A}f_i
                       \le e^\gamma-1.
 \label{eq:connected-output}
\end{equation}
Connectivity is that of the overlap graph.
More generally the sum of families whose every component meets
a set \(Y\), with individual weights
\(u_i=f_i e^{\kappa n_B(X_i)}\), is at most
\(e^{|Y|\gamma}\le e^{a n_B(Y)}\).
For labelled rooted trees with a prescribed root \(i\), omit
the root price and include \(1/m!\) for its unordered children.
The sum of their weights is at most \(e^{a n_B(X_i)}\), even
when a label occurs more than once in the tree.
\end{lemma}
\begin{proof}
First restrict to a finite available label set \(D\).
Put \(u_i=f_i e^{\kappa n_B(X_i)}\).
For a site set \(Y\), let
\begin{equation}
 W_D(Y)=
 \sum_{\substack{A\subset D\\
      \text{every overlap component of }A\text{ meets }Y}}
                                      \prod_{i\in A}u_i.
 \label{eq:rooted-family-sum}
\end{equation}
A component meets \(Y\) if one of its footprints intersects \(Y\).
The empty set contributes one.

We first prove the finite combinatorial inequality
\begin{equation}
 W_D(Y)\le
 \prod_{\substack{i\in D\\X_i\cap Y\ne\varnothing}}
                  \bigl(1+u_i W_{D\setminus\{i\}}(X_i)\bigr).
 \label{eq:rooted-recursion}
\end{equation}
Choose a fixed ordering of \(D\). For each component of a family
contributing on the left, choose its first label whose footprint
meets \(Y\), and call it the root. On deleting the root, every
remaining component has an edge to the deleted root: follow a path
in the original connected graph to that root.
Thus the remaining labels in that component form one term of
\(W_{D\setminus\{i\}}(X_i)\).
The product on the right can choose these roots and these terms.
Their union recovers the original family, so two different families
cannot yield the same chosen tuple.
The weight of that tuple is exactly the original product.
The right side permits additional choices, including overlapping
descendant families. All weights are nonnegative, which proves
the inequality. No equality or uniqueness of an unrestricted tree
expansion is being asserted.

We prove by strong induction on \(|D|\), simultaneously for all
\(Y\), that \(W_D(Y)\le e^{a n_B(Y)}\).
For \(D=\varnothing\), the left side is one.
If the bound holds for proper subsets of \(D\), then
\begin{align}
 W_D(Y)
 &\le\prod_{i:X_i\cap Y\ne\varnothing}
          \bigl(1+u_i e^{a n_B(X_i)}\bigr)\notag\\
 &\le\exp\left(
       \sum_{i:X_i\cap Y\ne\varnothing}
                    f_i e^{(\kappa+a)n_B(X_i)}\right)\notag\\
 &\le e^{|Y|\gamma}
 \le e^{a n_B(Y)}.
 \label{eq:rooted-induction}
\end{align}
The third inequality charges a label hitting \(Y\) to every site
of \(Y\) that it hits, which only increases the sum.
The last uses \eqref{eq:cover-cardinality} and
\(D_B\gamma\le a\).
This also shows the sharper estimate \(W_D(\{x\})\le e^\gamma\).

If every component meets the singleton \(\{x\}\), all components
would meet each other at \(x\); hence a nonempty such family has
only one component.
For any connected \(A\), \eqref{eq:cover-connected} gives
\begin{equation}
 e^{\kappa n_B(\bigcup_{i\in A}X_i)}\prod_{i\in A}f_i
 \le\prod_{i\in A}u_i .
\end{equation}
The finite version of \eqref{eq:connected-output} is therefore
bounded by \(W_D(\{x\})-1\le e^\gamma-1\).
Finally increase finite sets \(D\) to \(L\).
Every finite family eventually occurs, and all terms are nonnegative.
Monotone convergence proves the stated countable sum and its
summability.

For the last assertion define the price-free root sums of height
at most \(h\) by
\[
 \Phi_i^{(0)}=1,\qquad
 \Phi_i^{(h+1)}=
   \exp\!\left\{\sum_{j:X_j\cap X_i\ne\varnothing}
                                      u_j\Phi_j^{(h)}\right\}.
\]
The exponential puts a factor \(1/m!\) on \(m\) children.
Expanding it recursively enumerates rooted labelled trees;
the division by child factorials agrees with the usual
\(1/n!\) for the \(n\) labelled nonroot vertices after their
allocation to branches. No distinctness of species is imposed.
If \(\Phi_j^{(h)}\le e^{a n_B(X_j)}\), then
\[
 \sum_{j:X_j\cap X_i\ne\varnothing}u_j\Phi_j^{(h)}
 \le |X_i|\gamma\le a n_B(X_i).
\]
Induction and monotone convergence therefore give
\begin{equation}
       \Phi_i:=\lim_h\Phi_i^{(h)}\le e^{a n_B(X_i)}.
 \label{eq:common-pinned-tree-bound}
\end{equation}
This definition also covers \(u_i=0\): no division by a root
activity is made. Finite alphabets followed by monotone
convergence justify countably many labels in this calculation.
\end{proof}

The reserve \(a\) pays for the attached families in
\eqref{eq:rooted-induction}; the frequent choice \(a=1\)
recovers the unit-reserve criterion.
The factor \(D_B\) enters through the number of sites in a covering
ball; it is independent of the total number of sites.

\section{One-step reassembly of large-field corrections}
\label{sec:one-step-reassembly}

After fluctuation integration and local subtraction, some corrections
vanish in the small-field region. They should be absorbed into the
forced activities while preserving the represented density exactly.
We now prove that reassembly step, including its norm estimate.

\subsection{The preliminary local corrections}

Let the old forced activities be \(K_\alpha\), with footprints
\(Y_\alpha\), inventories \(B_\alpha\), and tests as above.
Let \(u_j\) be countably many real Borel corrections with nonempty
footprints \(S_j\). They are gauge invariant, determined by the
links and flags in \(S_j\), and satisfy
\begin{equation}
 u_j(U)=0\quad\text{if }S_j\cap\mathcal B_h(U)=\varnothing.
 \label{eq:attachment-bad}
\end{equation}
The function to be reassembled has the preliminary form
\begin{equation}
 e^{v\#\Lambda-V-E^\#}\, Z_K\,e^{\sum_j u_j}.
 \label{eq:preliminary-density}
\end{equation}
Define the Mayer factors \(m_j=e^{u_j}-1\).
This elementary exponential expansion, used in polymer methods,
replaces \(e^{\sum_j u_j}\) by products in which a label occurs
at most once.
Compare the subsequent grouping into large-field components with
\cite[Sect.~2, eq.~(2.19), printed p.~258]{Bal88b}.
The reassembly identity and its bounds below are derived for the
present labelled activities.

Suppose nonnegative scalar bounds \(a_\alpha,b_j\) obey
\begin{equation}
 |K_\alpha(U)|\le a_\alpha,\qquad |u_j(U)|\le b_j
 \label{eq:primitive-point-bounds}
\end{equation}
and, at exponent \(\kappa+1\),
\begin{align}
 \sup_x\sum_{\alpha:x\in Y_\alpha}
                  a_\alpha e^{(\kappa+1)n_B(Y_\alpha)}
                    &\le a, \label{eq:primitive-hard-bound}\\
 \sup_x\sum_{j:x\in S_j}
                  b_j e^{(\kappa+1)n_B(S_j)}
                    &\le b\le1. \label{eq:primitive-soft-bound}
\end{align}
Each \(b_j\le b\), by choosing a site in \(S_j\).
Thus
\begin{equation}
 |m_j(U)|\le e^{|u_j(U)|}-1\le2b_j .
 \label{eq:mayer-point-bound}
\end{equation}
Also \(\sum_\alpha a_\alpha\le{\#\Lambda}a\) and
\(\sum_j b_j\le{\#\Lambda}b\), by the same counting as
\eqref{eq:K-total}.
In particular all sums defining the preliminary expression converge.

\subsection{The exact connected reorganisation}

Use the disjoint union of the old forced labels and the correction
labels as a combined alphabet. Each combined label retains its
whole footprint. A new label is a finite connected combined family
\(C\) containing at least one old forced label, with the old
forced footprints within \(C\) pairwise disjoint.
Its footprint and inventory are
\begin{equation}
 Y_C=\bigcup_{\lambda\in C}X_\lambda,\qquad
 B_C=\bigcup_{\alpha\in C\text{ old}}B_\alpha,
 \label{eq:merged-footprint}
\end{equation}
where \(X_\lambda\) is \(Y_\alpha\) or \(S_j\) according to the
label type. In particular \(B_C\ne\varnothing\).
Set
\begin{equation}
 K'_C(U)=
 \mathbf1_{\{Y_C\cap\mathcal B_h(U)=B_C\}}
 \prod_{\alpha\in C\text{ old}}K_\alpha(U)
 \prod_{j\in C\text{ correction}}m_j(U).
 \label{eq:merged-activity}
\end{equation}
Its additional test can be the conjunction of the old additional
tests. A nonzero product already enforces them by
\eqref{eq:forced-eligibility}.
This gives Borel, gauge-invariant forced activities with the
required footprint and inventory rules.

\begin{proposition}[Exact reassembly identity]
\label{prop:exact-reassembly}
Under \eqref{eq:primitive-point-bounds}--\eqref{eq:primitive-soft-bound},
the new gas satisfies
\begin{equation}
 Z_{K'}(U)=Z_K(U)e^{\sum_j u_j(U)}.
 \label{eq:exact-reassembly}
\end{equation}
The identity is an equality of absolutely convergent series at each
configuration.
\end{proposition}
\begin{proof}
For a finite correction set \(J\),
\(\prod_{j\in J}(1+m_j)=e^{\sum_{j\in J}u_j}\).
Since \(\sum_j|m_j|\le2\sum_jb_j<\infty\), passing to all labels
gives
\begin{equation}
 e^{\sum_j u_j}
      =\sum_{J\text{ finite}}\prod_{j\in J}m_j,\qquad
 \sum_{J\text{ finite}}\prod_{j\in J}|m_j|
                    \le e^{2\sum_jb_j}.
 \label{eq:mayer-absolute}
\end{equation}
The first limit is justified by uniform convergence of \(\sum u_j\)
and absolute convergence of the product expansion.
Multiplying by the old gas gives an absolutely convergent double
series, bounded in absolute sum by
\(e^{\sum_\alpha a_\alpha+2\sum_jb_j}\).

Fix a nonzero term: an admissible set \(A\) of old labels and a
finite set \(J\) of correction labels. Decompose the overlap graph
of \(A\sqcup J\) into its connected components.
Every component contains an old label.
Indeed, a correction with \(m_j(U)\ne0\) has
\(S_j\cap\mathcal B_h(U)\ne\varnothing\), by
\eqref{eq:attachment-bad}. A bad site in that intersection belongs
to an inventory of an old label in \(A\), because \(A\) covers
the whole flag set. Its footprint therefore meets \(S_j\).

Different components have disjoint full footprints.
For a component \(C\), every bad site in \(Y_C\) belongs to an old
inventory from that same component: the old footprint containing
the site intersects \(Y_C\) and hence joins the component.
It follows that \(Y_C\cap\mathcal B_h(U)=B_C\).
Thus the components are admissible new labels, with the same
product of function values.

Conversely, flatten an admissible family of nonzero new activities.
The old labels are compatible, inside each component by construction
and between components by disjointness of the merged footprints.
Their inventories cover all flags. Correction labels cannot repeat
between components, since their nonempty footprint would then lie
in both merged footprints. The connected components of the
flattened family are precisely the chosen new labels.
The two operations are inverse and preserve each nonzero term.
This bijection and the absolute double-series bound justify
regrouping and prove \eqref{eq:exact-reassembly}.
\end{proof}

\subsection{The local norm after reassembly}

\begin{theorem}[One-step forced-activity estimate]
\label{thm:forced-reassembly}
Assume \eqref{eq:primitive-hard-bound} and
\eqref{eq:primitive-soft-bound}.
Choose \(\eta\ge0\) and put
\begin{equation}
 \gamma=e^\eta a+2b.
\end{equation}
If \(D_B\gamma\le1\), then
\begin{equation}
 \|K'\|_{\mathcal K,\kappa}
             \le e^{-\eta}(e^\gamma-1).
 \label{eq:forced-reassembly-bound}
\end{equation}
The regular action, reference action, and scalar \(v\) in
\eqref{eq:preliminary-density} are unchanged, and the density is
represented exactly by the new activities.
\end{theorem}
\begin{proof}
Assign combined scalar weights
\begin{equation}
 f_\alpha=e^\eta a_\alpha
       \quad\text{to old labels},\qquad
 f_j=2b_j\quad\text{to corrections}.
\end{equation}
Their local incidence sum at exponent \(\kappa+1\) is bounded by
\(\gamma\).
Every new connected label contains at least one old label, so
\eqref{eq:merged-activity} and \eqref{eq:mayer-point-bound} give
\begin{equation}
 \|K'_C\|_\infty
 \le e^{-\eta\,|C\cap\{\text{old labels}\}|}
                         \prod_{\lambda\in C}f_\lambda
 \le e^{-\eta}\prod_{\lambda\in C}f_\lambda.
\end{equation}
For each site \(x\), sum this bound over new footprints containing
\(x\), with weight \(e^{\kappa n_B(Y_C)}\).
The new labels form a subcollection of the connected families in
Lemma~\ref{lem:connected-bound}. Dropping their additional restrictions
only increases the nonnegative sum.
That lemma yields \eqref{eq:forced-reassembly-bound}.
Proposition~\ref{prop:exact-reassembly} proves the density identity.
All convergence and measurability requirements follow from the
same absolute majorants.
\end{proof}

\begin{lemma}[Choosing the reserve for a prescribed output bound]
\label{lem:reserve-choice}
Let \(w,P>0\), \(a,b\ge0\), and assume
\begin{equation}
 a\le\frac w8,\quad wP\le1,\quad b\le1,\quad
 Pb\le\frac18,\quad
 D_B\left(\frac1{8P}+2b\right)\le1.
 \label{eq:reserve-conditions}
\end{equation}
Then the choice \(\eta=-\log(wP)\) in
Theorem~\ref{thm:forced-reassembly} gives
\(\|K'\|_{\mathcal K,\kappa}\le3w/4\).
\end{lemma}
\begin{proof}
We have \(\eta\ge0\), \(e^{-\eta}=wP\), and
\begin{equation}
 0\le\gamma=\frac{a}{wP}+2b
              \le\frac1{8P}+2b,\qquad D_B\gamma\le1.
\end{equation}
Since \(D_B\ge1\), also \(\gamma\le1\), and convexity gives
\(e^\gamma-1\le2\gamma\).
Consequently,
\begin{equation}
 \|K'\|_{\mathcal K,\kappa}
 \le2wP\left(\frac1{8P}+2b\right)
 =w\left(\frac14+4Pb\right)\le\frac34w.
\end{equation}
\end{proof}

This calculation permits a very small output target \(w\).
The correction bound \(b\) need not itself be smaller than \(w\).
The old forced terms must instead have the bound \(a\le w/8\);
their presence in every connected component supplies the reserve.

\begin{corollary}[A useful large-coupling regime]
\label{cor:large-coupling-reserve}
Let \(\beta\to\infty\), and suppose
\begin{equation}
 B(\beta)\le C(1+\log\beta)^2,\qquad
 b\le\beta^{-1/8},\qquad
 a\le w/8,\qquad 0<w\le\beta^{-1/16}.
 \label{eq:large-coupling-regime}
\end{equation}
For all sufficiently large \(\beta\), reassembly gives
\(\|K'\|_{\mathcal K,\kappa}\le3w/4\).
The threshold depends on \(C\), not on the torus volume.
\end{corollary}
\begin{proof}
Take \(P=\beta^{1/16}\).
Then \(wP\le1\), \(Pb\le\beta^{-1/16}\to0\), and \(b\le1\)
eventually. Furthermore,
\begin{equation}
 D_{B(\beta)}
 \left(\frac18\beta^{-1/16}+2\beta^{-1/8}\right)\longrightarrow0,
\end{equation}
because every fixed power of \(\log\beta\) is smaller than
every positive power of \(\beta\) eventually.
Thus \eqref{eq:reserve-conditions} holds for large \(\beta\),
and Lemma~\ref{lem:reserve-choice} applies.
\end{proof}

\subsection{The regular budget and the remaining analytic step}

Regular corrections are combined by the triangle inequality at their
anchors. For example, suppose the transported regular family has
norm at most \(3\delta/16\), while further regular corrections have
total norm at most \(\delta/2\), with the same norm parameters.
Their union has norm at most \(11\delta/16<\delta\).
One way to obtain the second bound is the following elementary
power margin.

\begin{lemma}[A power margin for regular corrections]
\label{lem:regular-margin}
Fix a profile and \(c>0\), and let
\(\delta(\beta)=c\beta^{-11/10}\).
Suppose a regular remainder has, for every \(\nu>0\), a bound
\begin{equation}
 \|R(\beta)\|_{\mathcal E,\kappa}
                  \le C_\nu\beta^{-23/20+\nu}
 \label{eq:regular-power-input}
\end{equation}
for all sufficiently large \(\beta\), with constants independent
of cutoff and volume. Then
\(\|R(\beta)\|_{\mathcal E,\kappa}\le\delta(\beta)/2\)
for all sufficiently large \(\beta\), uniformly in those parameters.
\end{lemma}
\begin{proof}
Choose \(\nu=1/40\).
The exponent in \eqref{eq:regular-power-input} is then \(-45/40\),
while \(-11/10=-44/40\). Dividing the bound by \(\delta(\beta)\)
gives at most \((C_{1/40}/c)\beta^{-1/40}\to0\).
\end{proof}

Theorems~\ref{prop:exact-reassembly} and
\ref{thm:forced-reassembly} prove the exact large-field reassembly
part of a step. Lemma~\ref{lem:regular-margin} states precisely how
a sufficiently strong regular remainder enters the regular budget.
Their inputs concern the \emph{actual} preliminary local functions
and the stronger footprint weight \(\kappa+1\).

To complete a full RG step, the blocked Wilson density
\(\mathcal P_mF_{m+1}\) must still be put into
\eqref{eq:preliminary-density}, with the stated input bounds.
This requires the local stationary configuration, its fluctuation
Hessian, the associated Gaussian integral, and a controlled remainder.
The principal quadratic term must be extracted so that the new
regular family satisfies \eqref{eq:affine-subtraction}, and the
scalar and coupling must agree with the exact density identity.
The preceding reassembly theorem does not assume that these tasks
follow from a bound on the partition function.

The next analytic calculation will therefore concern that fluctuation
integral itself. The norms, support conventions, and volume-independent
estimates it must deliver have now been specified.
\section{The physical blocking integral}
\label{sec:physical-fluctuation}

We now retain the integration variables and all normalising constants.
The fluctuation calculation has three stages: an exact change of
variables, an estimate on a central Gaussian integral, and a comparison
with the integral outside its central region. These stages have different
uniformity requirements. This organisation follows the small-field
analysis of Ba{\l}aban \cite[Sect.~2, printed pp.~265--268,
especially eqs.~(2.12)--(2.15) on p.~268]{Bal87};
the finite-dimensional estimates below are proved directly.

\subsection{Noise normalisation and the finite branch sum}

Fix a single blocking step, from side \(f=Lc\) to side \(c\).
Here \(L\ge3\) is odd, and both sides are measured in lattice units.
Write
\[
 q_G(W)=1-N^{-1}\Re\operatorname{tr}W,\qquad b_c=4c^4 .
\]
Thus \(q_G(W)=d_G(W)^2/2\), and \(b_c\) is the number of positive
coarse links. Fix a positive noise parameter \(a\). Its normaliser is
\begin{equation}
 z_{\mathrm n}(\beta)=\int_G e^{-a\beta q_G(W)}\,\dd h_G(W)>0 .
 \label{eq:noise-normalizer}
\end{equation}
The finite branch data consist of a finite set \(\mathcal T\), Borel
weights \(\lambda_t(U)\ge0\) satisfying
\(\sum_{t\in\mathcal T}\lambda_t(U)=1\), and predicted coarse
connections \(Q_t(U)\in\mathcal U_c\).
A branch records the local frame and averaging choices at every
physical coarse site and link. There is one common choice at a shared
site; choices are not made separately by neighbouring blocks.

In terms of these data the physical kernel is
\begin{equation}
 K_\beta(U,V)=z_{\mathrm n}(\beta)^{-b_c}
 \sum_{t\in\mathcal T}\lambda_t(U)
 \exp\!\left\{-a\beta\sum_{e\in\Lambda_c\times\{1,\ldots,4\}}
             q_G\bigl(V_eQ_t(U)_e^{-1}\bigr)\right\}.
 \label{eq:physical-kernel}
\end{equation}
This formula specifies the kernel in terms of its finite branch maps.
The construction and estimates for those local maps are a separate
part of the blocking construction. The calculations in this section
apply to the fixed maps in \eqref{eq:physical-kernel}; they do not
change them with the input density.

\begin{lemma}[The physical prefactor]
\label{lem:noise-normalization}
The kernel \eqref{eq:physical-kernel} is nonnegative.
For every \(U\), it satisfies
\[
 \int K_\beta(U,V)\,\dd H_c(V)=1 .
 \]
Moreover,
\[
 e^{-2a|\beta|}\le z_{\mathrm n}(\beta)\le e^{2a|\beta|}.
\]
In particular its prefactor is finite at every finite volume.
\end{lemma}
\begin{proof}
Since \(0\le q_G\le2\), the bounds on \(z_{\mathrm n}\) follow by
integration against probability Haar measure.
For each fixed \(U,t\), right invariance of Haar measure gives
\[
 \int_G e^{-a\beta q_G(V_eQ_t(U)_e^{-1})}\,\dd h_G(V_e)
 =z_{\mathrm n}(\beta).
\]
Integrating the product over all \(b_c\) coarse links cancels
\(z_{\mathrm n}^{-b_c}\). The remaining sum is
\(\sum_t\lambda_t(U)=1\).
\end{proof}

The elementary bound on the prefactor is extensive in \(c^4\).
We will keep this prefactor exactly; it is not a bound on a local
activity. For an effective input in the form of
Section~\ref{sec:activities}, write its reference action as
\(\beta S_f+I_f\), and put
\[
 W_f(U)=e^{-E^\#(U)}Z_K(U).
\]
This notation only splits the already specified reference action;
it imposes no extra positivity on \(W_f\).
With \(n_f=f^4\), the exact blocked density is
\begin{align}
 (\mathcal P_\beta F)(V)
 &=e^{vn_f}z_{\mathrm n}(\beta)^{-b_c}
   \sum_t\int \lambda_t(U)W_f(U)
           e^{-\beta A_t(V,U)-I_f(U)}\,\dd H_f(U), \label{eq:actual-block-integral}\\
 A_t(V,U)&=S_f(U)+a\sum_e q_G\bigl(V_eQ_t(U)_e^{-1}\bigr).
 \label{eq:actual-stationary-action}
\end{align}
At the original Wilson level \(S_f=E_f\), \(I_f=0\),
\(v=0\), and \(W_f=1\). At later levels all the input activities
in \eqref{eq:actual-block-integral} remain inside the same integral.

\subsection{Tree coordinates and the variables left to integrate}

A rooted forest here means a collection of trees on fine sites,
one tree in each coarse block, with the coarse embedded site as root.
Let \(T\) be its set of tree links and \(I\) the complementary
set of independent positive fine links. A connection is in tree
gauge when its links on \(T\) are the identity. It is then determined
by its links in \(I\).

\begin{lemma}[Exact integration in tree gauge]
\label{lem:tree-haar}
Let \(J(U)\) be an integrable function invariant under fine gauge
transformations equal to \(1\) at every root. If \(U_T(u)\) denotes
the tree-gauge connection with free links \(u\in G^I\), then
\begin{equation}
 \int J(U)\,\dd H_f(U)=\int_{G^I}J(U_T(u))\,\dd h_G^I(u).
 \label{eq:tree-haar}
\end{equation}
\end{lemma}
\begin{proof}
For every \(U\) there is a unique root-fixed gauge transformation
making the tree links equal to \(1\). Starting at each root, its
value at a child is determined by its value at the parent and by the
link joining them. Thus the map from \(U\) to its root-fixed gauge
variables and its transformed free links is invertible.

Order the tree vertices by distance from the root. At each step the
coordinate change on a link is left or right multiplication by group
elements already determined at earlier steps. Haar invariance shows
that the product measure becomes the product of probability Haar
measures on the gauge variables and on \(G^I\).
Fubini applies first to nonnegative functions and then to the positive
and negative parts of \(J\). Gauge invariance removes the gauge
variables, whose total measure is \(1\).
\end{proof}

For the joint integrand at fixed \(V\), the invariance in this lemma
uses the covariance of the branch maps under root-fixed gauge
transformations. Full fine gauge invariance of \(F\) alone would
not imply it for \(K_\beta(U,V)F(U)\).

Choose the inner product on \(\mathfrak g\) fixed earlier and its
Lebesgue measure. There are \(r_G>0\), \(\zeta_G>0\), and a positive
smooth function \(j_G\), with \(j_G(0)=1\), such that
\begin{equation}
 \dd h_G(e^X)=\zeta_G j_G(X)\,\dd X,\qquad |X|_{\mathfrak g}<r_G.
 \label{eq:haar-chart-density}
\end{equation}
Only the chart in the identity component is used in this formula;
\(\zeta_G\) includes the normalisation of probability Haar measure
on the whole group.
For completeness, its local Jacobian is
\begin{equation}
 j_G(X)=
 \left|\det_{\mathfrak g}
      \int_0^1 e^{-t\operatorname{ad}X}\,\dd t\right|,
 \qquad \operatorname{ad}X(Y)=[X,Y].
 \label{eq:haar-chart-jacobian}
\end{equation}
Indeed, differentiating the matrix exponential and then translating
its tangent vector back to the identity gives
\(\int_0^1e^{-t\operatorname{ad}X}\,\dd t\).
The determinant of this map gives the pullback of left invariant
volume. At \(X=0\) it is \(1\), so it stays positive on a smaller
ball. Probability Haar measure is a positive constant multiple of
this volume. The exponential chart and these assertions use the
standard closed subgroup theorem for matrix groups.
Every fixed number of derivatives of \(j_G\), \(1/j_G\), and
\(\log j_G\) is bounded on \(|X|\le r_G/2\), by compactness.

To describe a sector of \eqref{eq:actual-block-integral}, split
\(I=I_{\mathrm a}\sqcup I_{\mathrm f}\). The subscripts mean
active and frozen: active links are in exponential charts, while
frozen links are still integrated against their original Haar
measure. If \(Y\in G^{I_{\mathrm f}}\) and
\(y\in\mathfrak g^{I_{\mathrm a}}\), write \(U(V,Y,y)\) for the
tree-gauge connection with active links
\[
 U_e=e^{y_e}R_e(V),
\]
where \(R_e(V)=1\) on internal replica links and equals the
corresponding coarse link on crossing replica links.
Let \(\theta(V,Y,y)\) be the actual sector multiplier.
It includes the sector indicator, the branch likelihood, and any
input activity multipliers not included in the smooth action.
Every such factor is retained in the identity
\begin{align}
 &\zeta_G^{|I_{\mathrm a}|}
 \int_{G^{I_{\mathrm f}}}\int_{\mathfrak g^{I_{\mathrm a}}}
 \theta(V,Y,y)
 e^{-\beta\varphi(V,Y,y)-\psi(V,Y,y)}\,\dd y\,\dd h_G^{I_{\mathrm f}}(Y),
 \label{eq:sector-chart-integral}\\
 &\psi(V,Y,y)
     =I_f(U(V,Y,y))-\sum_{e\in I_{\mathrm a}}\log j_G(y_e).
 \label{eq:smooth-action-with-haar}
\end{align}
Here \(\varphi=A_t(V,U(V,Y,y))\) on a branch where that action
is smooth, and the active chart cutoff is part of \(\theta\).
One obtains the full integral by summing the sectors of a Borel
partition of unity and retaining the common factors in
\eqref{eq:actual-block-integral}. Formula
\eqref{eq:sector-chart-integral} concerns the actual action and
Haar Jacobian, rather than just its quadratic Taylor polynomial.

\begin{lemma}[Gaussian change of variables on a sector]
\label{lem:sector-gaussian-change}
Put \(d=|I_{\mathrm a}|\dim_{\R}\mathfrak g\), \(g=\beta^{-1/2}\),
and let \(C\) be a positive definite \(d\times d\) covariance matrix.
For fixed \(V,Y\), let \(m=m(V,Y)\in\R^d\) be any measurable shift.
For every integrable sector integrand \(B(y)\),
\begin{equation}
 \int_{\R^d}B(y)\,\dd y
 =(2\pi g^2)^{d/2}(\det C)^{1/2}
 \int_{\R^d}
 e^{z^TC^{-1}z/2}B(m+gz)\,\dd\mu_C(z),
 \label{eq:exact-sector-gaussian-change}
\end{equation}
where
\(\dd\mu_C(z)=(2\pi)^{-d/2}(\det C)^{-1/2}
e^{-z^TC^{-1}z/2}\,\dd z\).
Absolute integrability is equivalent on the two sides.
\end{lemma}
\begin{proof}
Substitute \(y=m+gz\), whose Jacobian is \(g^d\), and insert
the displayed density of \(\mu_C\). The same calculation applied
to \(|B|\) proves the integrability assertion.
\end{proof}

In particular, the determinant counts only active real coordinates.
Frozen Haar variables do not contribute a factor of \(g\).
The multiplier \(\theta(V,Y,m+gz)\) in
\eqref{eq:sector-chart-integral} remains present after
\eqref{eq:exact-sector-gaussian-change}. A symmetric central
cutoff can be used in the next calculation only on a region
where the actual sector cutoffs admit that central region.

\section{A Gaussian remainder with its moving determinant}
\label{sec:gaussian-remainder}

\subsection{The central integral and its exact normalisation}

Fix the frozen variables and a smooth branch. Let \(x\) denote the
coarse exponential coordinates in a finite-dimensional real normed
space \(E\). The active coordinates form \(\R^d\) with its Euclidean
norm. Let \(\rho(x)\) be a stationary configuration:
\begin{equation}
 D_y\varphi(x,\rho(x))=0,\qquad
 {\mathsf H}(x)=D_y^2\varphi(x,\rho(x))>0 .
 \label{eq:stationary-graph}
\end{equation}
The existence and estimates for this particular stationary graph
must be proved for the action in
\eqref{eq:actual-stationary-action}. The estimates in this section
state their finite-dimensional inputs explicitly.

For a central radius \(D>0\), define
\begin{equation}
 I_D(g,x)=
 \int_{|y-\rho(x)|\le |g|D}
     e^{-g^{-2}\varphi(x,y)-\psi(x,y)}\,\dd y
 \quad(g\ne0).
 \label{eq:central-physical-integral}
\end{equation}
The group factor \(\zeta_G^{|I_{\mathrm a}|}\) and the noise
normaliser are still the factors already displayed in
\eqref{eq:sector-chart-integral} and
\eqref{eq:actual-block-integral}.
Introduce the Taylor integral
\begin{equation}
 V_g(x,z)=\int_0^1(1-t)
 D_y^2\varphi(x,\rho(x)+tgz)[z,z]\,\dd t .
 \label{eq:stationary-taylor-integral}
\end{equation}
Taylor's formula with integral remainder and
\eqref{eq:stationary-graph} give
\begin{equation}
 \varphi(x,\rho(x)+gz)-\varphi(x,\rho(x))=g^2V_g(x,z),
 \qquad V_0(x,z)=\tfrac12z^T{\mathsf H}(x)z .
 \label{eq:stationary-taylor-exact}
\end{equation}
Both formulas hold for either sign of \(g\).

Write \(\ell(x)=\frac12\log\det {\mathsf H}(x)\), and let
\(\mu_0=\mu_{{\mathsf H}(0)^{-1}}\).
The normalised integral we estimate is
\begin{align}
 Z_D(g,x)&=\int_{|z|\le D}e^{P(g,x,z)}\,\dd\mu_0(z), \label{eq:central-normalized}\\
 P(g,x,z)&=\ell(x)-\ell(0)+V_0(0,z)-V_g(x,z)\notag\\
 &\hspace{18mm}+\psi(x,\rho(x))-\psi(x,\rho(x)+gz).
 \label{eq:central-exponent}
\end{align}
It is defined also at \(g=0\).
Its Gaussian measure is fixed as \(x\) varies.
This keeps derivatives of the moving covariance inside the explicit
exponent \(P\).

\begin{proposition}[Exact stationary formula]
\label{prop:exact-stationary-formula}
For \(g>0\),
\begin{equation}
 I_D(g,x)=(2\pi)^{d/2}g^d
 e^{-g^{-2}\varphi(x,\rho(x))-\psi(x,\rho(x))-\ell(x)}
 Z_D(g,x).
 \label{eq:stationary-density-identity}
\end{equation}
Moreover,
\begin{equation}
 Z_D(-g,x)=Z_D(g,x),\qquad
 Z_D(0,x)=\mu_{{\mathsf H}(x)^{-1}}\{|z|\le D\}.
 \label{eq:parity-and-zero-coupling}
\end{equation}
\end{proposition}
\begin{proof}
Substitute \(y=\rho(x)+gz\) in
\eqref{eq:central-physical-integral} and use
\eqref{eq:stationary-taylor-exact}. The Lebesgue density of
\(\mu_0\) is
\[
 (2\pi)^{-d/2}e^{\ell(0)-V_0(0,z)}.
 \]
Multiplication by \(e^{P(g,x,z)}\) gives
\[
 (2\pi)^{-d/2}
 e^{\ell(x)-V_g(x,z)+\psi(x,\rho(x))
                              -\psi(x,\rho(x)+gz)} .
 \]
This proves \eqref{eq:stationary-density-identity}.

For parity, \eqref{eq:stationary-taylor-integral} gives
\(V_{-g}(x,-z)=V_g(x,z)\), and all the other terms in \(P\)
have the same property. Both the ball and the centered Gaussian
measure are invariant under \(z\mapsto-z\).
For \(g=0\), the same density calculation leaves exactly the
Gaussian density with precision \({\mathsf H}(x)\).
\end{proof}

The factor \(e^{-\ell(x)}\) in
\eqref{eq:stationary-density-identity} is
\((\det {\mathsf H}(x))^{-1/2}\).
Thus the moving determinant is part of the density.
Also \(Z_D(0,x)\) is a Gaussian probability of a finite ball:
it need not equal \(1\).

\subsection{Derivative norms and the constants used below}

For \(k\in\N\), an open set \(O\subset E\), and \(s>0\), put
\begin{equation}
 \|f\|_{C_s^k(O)}
 =\max_{0\le j\le k}s^j
       \sup_{x\in O}\|D^jf(x)\|_{\mathrm{op}} .
 \label{eq:finite-chart-jet-norm}
\end{equation}
The zeroth derivative is \(f\) itself. We omit \(s\) when \(s=1\).
All operator norms use the stated norm on \(E\) and the Euclidean
norm on active coordinates. On \(\R\times E\) we use the maximum
norm. This finite-chart norm will subsequently be compared with
the ordered left-derivative activity norm.

Here are explicit finite-chart assumptions for a fixed integer
\(k\ge0\). On \(\|x\|<w\), \(|g|<v\), \(|z|\le D\), and
\(0\le t\le1\), require:
\begin{equation}
 0<w,v\le1,\qquad {\mathsf H}(x)\ge\eta\,1\quad(\eta>0),
 \label{eq:chart-gap}
\end{equation}
and a common bound \(C\ge1\) for all derivatives, including order
zero, of \(\rho\) through order \(k+4\), and of
\(\varphi,\psi\) through order \(k+4\) at the points
\((x,\rho(x)+tgz)\).
All functions are defined and \(C^{k+4}\) on neighbourhoods of
these points. Finally, \(\rho(0)=0\) and
\eqref{eq:stationary-graph} holds throughout the coarse domain.
Bounds for \(\varphi\) and \(\psi\) themselves can be replaced by
bounds after subtracting their value at \((0,0)\); those constants
cancel in \(P\).

\begin{lemma}[A polynomial derivative bound]
\label{lem:central-polynomial-bound}
Under these assumptions there is a polynomial \(A_k(C,D,d,\eta^{-1})\),
with nonnegative coefficients depending only on \(k\), which is at
least \(1\), such that
\begin{equation}
 \sup_{|z|\le D}
 \|P(\,\cdot\,,\,\cdot\,,z)\|_{C^{k+2}}
 \le A_k,\qquad
 |P(g,x,z)|\le A_k\max(|g|,\|x\|).
 \label{eq:central-polynomial-bound}
\end{equation}
The first norm is on \((-v,v)\times B_E(0,w)\).
There is no dependence on \(\dim E\) apart from its occurrence
in the input operator norms.
\end{lemma}
\begin{proof}
The map \((g,x)\mapsto(x,\rho(x)+tgz)\) has derivatives bounded
by polynomials in \(1+C+D\), uniformly in \(t\in[0,1]\).
For a derivative of order \(j\), repeated chain rules are indexed
by partitions of the \(j\) differentiated directions. Each term
is one derivative of the outer function applied to derivatives
of this map; the number of terms depends only on \(j\).
Applying this observation to \(D_y^2\varphi\), and evaluating its
two final arguments at \(z\), adds a factor at most \(D^2\).
Integration in \(t\) proves the asserted polynomial bound for
\(V_g\). The same argument applies to the two terms involving
\(\psi\).

For the determinant, its first derivative is
\[
 D\ell(x)[u]=\tfrac12\operatorname{tr}
                  \bigl({\mathsf H}(x)^{-1}D{\mathsf H}(x)[u]\bigr).
 \]
Differentiation of an inverse gives
\(D({\mathsf H}^{-1})[u]=-{\mathsf H}^{-1}D{\mathsf H}[u]{\mathsf H}^{-1}\).
Repeated differentiation therefore expresses every derivative
of \(\ell\), of positive order at most \(k+2\), as a finite sum
of traces of products of derivatives of \({\mathsf H}\) and copies of
\({\mathsf H}^{-1}\). Their norms are bounded by a polynomial in
\(C,d,\eta^{-1}\): use \(|\operatorname{tr}B|\le d\|B\|\).
Derivatives of \({\mathsf H}\) are bounded by the preceding chain-rule
argument for \(D_y^2\varphi(x,\rho(x))\).
The zeroth order difference \(\ell(x)-\ell(0)\) is bounded
by integrating its first derivative along the segment to \(x\).

These calculations construct \(A_k\) by finitely many sums and
products; no dimension-dependent counting of coordinate
derivatives is used. Finally \(P(0,0,z)=0\), by
\eqref{eq:stationary-taylor-exact}. Integrating the first
derivative along the segment from \((0,0)\) to \((g,x)\)
gives the last bound, after increasing \(A_k\).
\end{proof}

\subsection{Gaussian tails and differentiation at zero coupling}

\begin{lemma}[A tail and moment bound]
\label{lem:gaussian-tail-moment}
If \({\mathsf H}\ge\eta1\) on \(\R^d\), \(d\ge1\), and
\(\mu=\mu_{{\mathsf H}^{-1}}\), then
\begin{align}
 \mu\{|z|>D\}&\le2d\exp\{-\eta D^2/(2d)\},\label{eq:gaussian-ball-tail}\\
 \int|z|^{2j}\,\dd\mu(z)
 &\le(2j-1)!!(d/\eta)^j\qquad(j\ge1).
 \label{eq:gaussian-moment}
\end{align}
\end{lemma}
\begin{proof}
Every coordinate has variance at most \(\eta^{-1}\). Its moment
generating function is
\(\int e^{t z_i}\,\dd\mu=\exp(t^2({\mathsf H}^{-1})_{ii}/2)\).
Markov's inequality, with \(t=\eta D/\sqrt d\), bounds each of
the two tails \(\pm z_i>D/\sqrt d\) by
\(\exp(-\eta D^2/(2d))\). If \(|z|>D\), at least one
coordinate exceeds this threshold in absolute value.
The union bound proves \eqref{eq:gaussian-ball-tail}.

For the moment, the convexity inequality
\((\sum_i a_i)^j\le d^{j-1}\sum_i a_i^j\), \(a_i\ge0\),
and the one-dimensional Gaussian moment formula give
\[
 \int\Bigl(\sum_i z_i^2\Bigr)^j\,\dd\mu
 \le d^{j-1}\sum_i(2j-1)!!({\mathsf H}^{-1})_{ii}^j .
 \]
This is \eqref{eq:gaussian-moment}. The one-dimensional formula
follows by \(j\) integrations by parts, or by differentiating
the displayed moment generating function \(2j\) times at zero.
\end{proof}

\begin{lemma}[Derivatives of the finite-cutoff Gaussian mass]
\label{lem:zero-coupling-tail}
Suppose \eqref{eq:chart-gap} holds, derivatives of \({\mathsf H}\) through
order \(k\) are bounded by \(C_{\mathsf H}\ge1\), and
the right side of \eqref{eq:gaussian-ball-tail} is at most \(1/2\).
There is a polynomial \(T_k(d,C_{\mathsf H},\eta^{-1})\), depending only
on \(k\), such that
\begin{equation}
 \|\log Z_D(0,\cdot)\|_{C^k(B_E(0,w))}
 \le T_k(d,C_{\mathsf H},\eta^{-1})
                  e^{-\eta D^2/(4d)} .
 \label{eq:zero-coupling-log-tail}
\end{equation}
\end{lemma}
\begin{proof}
The density \(p_x(z)\) of \(\mu_{{\mathsf H}(x)^{-1}}\) satisfies
\[
 D_xp_x(z)[u]=p_x(z)
 \left(D\ell(x)[u]-\tfrac12 z^TD{\mathsf H}(x)[u]z\right).
 \]
Iterating this identity and the determinant calculation in
Lemma~\ref{lem:central-polynomial-bound} gives, for \(j\le k\),
\[
 \|D_x^jp_x(z)\|
 \le p_x(z)\,T'_k(d,C_{\mathsf H},\eta^{-1})(1+|z|^{2k})
 \]
for a polynomial \(T'_k\).
Differentiation under the integral is legitimate locally in \(x\):
the spectral lower bound and upper bounds on \({\mathsf H}\) supply a
Gaussian majorant times a polynomial for all these derivatives.
The same argument applies on the complement of the ball.

Since \(\int p_x(z)\,\dd z=1\), derivatives of the truncated
mass of positive order are negatives of derivatives of its
tail mass. Cauchy--Schwarz, followed by
\eqref{eq:gaussian-moment} and \eqref{eq:gaussian-ball-tail},
therefore bounds all derivatives through order \(k\) of
\(Z_D(0,x)-1\) by a polynomial times
\(e^{-\eta D^2/(4d)}\).
Also \(1/2\le Z_D(0,x)\le1\).
For order zero, \(|\log t|\le2(1-t)\) on \([1/2,1]\).
For positive orders, differentiate \(\log Z_D\):
each term is a product of derivatives of \(Z_D\), divided by a
power of \(Z_D\). Each product contains at least one derivative,
and \(Z_D^{-1}\le2\).
The preceding tail bounds, together with their polynomial
bounds without the exponential factor, prove
\eqref{eq:zero-coupling-log-tail}.
\end{proof}

\begin{theorem}[Central Gaussian remainder]
\label{thm:central-gaussian-remainder}
Assume the finite-chart hypotheses above, with one additional
derivative if necessary when applying the result at order \(k+1\).
Suppose
\begin{equation}
 A_k\max(w,v)\le1,\qquad
 2d e^{-\eta D^2/(2d)}\le\tfrac12 .
 \label{eq:central-smallness}
\end{equation}
There are polynomials \(B_k(C,D,d,\eta^{-1})\) and
\(T_k'(C,d,\eta^{-1})\), depending only on \(k\), such that
\begin{align}
 \tfrac12e^{-1}\le Z_D(g,x)&\le e,\label{eq:central-positive}\\
 \|\log Z_D(g,\cdot)\|_{C^k}
 &\le B_k g^2+T_k'e^{-\eta D^2/(4d)} .
 \label{eq:central-remainder-estimate}
\end{align}
In particular there is no term of order \(g\).
\end{theorem}
\begin{proof}
Lemma~\ref{lem:central-polynomial-bound} and
\eqref{eq:central-smallness} give \(|P|\le1\).
Integration of \(e^{-1}\le e^P\le e\) on the ball gives
\eqref{eq:central-positive}, since the ball has \(\mu_0\)-mass
at least \(1/2\).

For every mixed derivative through order \(k+2\), the chain
rule for \(e^P\) is a finite sum indexed by partitions:
each term is \(e^P\) times a product of derivatives of \(P\).
Its norm is bounded by a polynomial in \(A_k\), times \(e\).
The domain \(\{|z|\le D\}\) and \(\mu_0\) are fixed, so these
uniform integrable bounds justify differentiation under the
integral. They bound \(Z_D\) and its mixed derivatives.
The lower bound \eqref{eq:central-positive} then bounds all
mixed derivatives of \(\log Z_D\) through order \(k+2\) by a
polynomial \(B'_k\).

Put \(L(g,x)=\log Z_D(g,x)\). Proposition
~\ref{prop:exact-stationary-formula} makes \(L\) even in \(g\),
so \(\partial_gD_x^jL(0,x)=0\).
For \(0\le j\le k\), the fundamental theorem of calculus twice
gives
\[
 D_x^jL(g,x)-D_x^jL(0,x)
 =g^2\int_0^1(1-t)
       \partial_g^2D_x^jL(tg,x)\,\dd t .
 \]
Taking operator norms yields \(B'_kg^2/2\).
Apply Lemma~\ref{lem:zero-coupling-tail} to the remaining
term \(L(0,x)\), and increase the polynomials.
\end{proof}

When \(d=0\), the active space has one point, its Gaussian mass
is \(1\), and \(P=0\). Thus \(Z_D=1\) and the remainder is
zero. No tail expression with a denominator \(d\) is needed
in that case.

\subsection{Coupling dependence and subtraction of the scalar term}

The following growth scale makes both the central cutoff and the
dimension dependence explicit. Let \(\widehat\beta>e\) be the
nominal fine-scale coupling and put
\begin{equation}
 P(\widehat\beta)
 =\exp\{(\log\log\widehat\beta)^2\},\qquad
 D=d_0P(\widehat\beta),\qquad
 w=\widehat\beta^{-7/16},
 \label{eq:monitor-and-chart-width}
\end{equation}
where \(d_0>0\) is fixed.
In this subsection \(P(\widehat\beta)\) is a scalar growth
function; the exponent \(P(g,x,z)\) in
\eqref{eq:central-exponent} always has three arguments.
For every \(\nu>0\) and every fixed \(m\),
\begin{equation}
 (1+\log\widehat\beta)^m
       =o(P(\widehat\beta)),\qquad
 P(\widehat\beta)^m=o(\widehat\beta^\nu).
 \label{eq:monitor-growth}
\end{equation}
Indeed, set \(s=\log\log\widehat\beta\).
The logarithms to compare are \(ms\), \(s^2\), and
\(\nu e^s\), respectively.

\begin{corollary}[Uniform finite-chart rates]
\label{cor:uniform-central-rates}
Fix \(G,L,k,d_0\) and \(\eta_0>0\).
Suppose \(\eta\ge\eta_0\), the auxiliary active dimension obeys
\begin{equation}
 d\le C_1(1+\log\widehat\beta)^{m_1},
 \label{eq:auxiliary-dimension-bound}
\end{equation}
and the derivative constant \(C\) is subpower in
\(\widehat\beta\): for every \(\nu>0\),
\(C\le C_\nu\widehat\beta^\nu\) eventually.
Assume the finite-chart hypotheses on
\(\|x\|<w\), \(|g|<w\), with the widths in
\eqref{eq:monitor-and-chart-width}.
Uniformly for
\(\widehat\beta/2\le\beta\le2\widehat\beta\) and
\(g=\beta^{-1/2}\), one has, for every \(\nu>0\),
\begin{align}
 \|\log Z_D(g,\cdot)\|_{C^k}
 &\le C_{k,\nu}\widehat\beta^{-1+\nu},
 \label{eq:gaussian-unscaled-rate}\\
 \|\log Z_D(g,\cdot)-\log Z_D(g,0)\|_{C_w^k}
 &\le C_{k,\nu}\widehat\beta^{-23/16+\nu}.
 \label{eq:gaussian-centered-rate}
\end{align}
For the second bound require the finite-chart hypotheses at
order \(k+1\).
The constants and lower threshold may depend on the stated
fixed data, \(C_1,m_1\), and the subpower bounds \(C_\nu\).
They are independent of any other cutoff or physical volume
provided those data are uniform.
\end{corollary}
\begin{proof}
All the polynomial constants in
Theorem~\ref{thm:central-gaussian-remainder} are subpower.
Since \(w=\widehat\beta^{-7/16}\), the first condition in
\eqref{eq:central-smallness} holds eventually, and
\(\beta^{-1/2}<w\).
Furthermore, \eqref{eq:auxiliary-dimension-bound} and
\eqref{eq:monitor-growth} imply
\[
 \frac{D^2}{d\log\widehat\beta}\longrightarrow\infty .
 \]
Thus a polynomial in the derivative data times
\(e^{-\eta D^2/(4d)}\) is at most
\(\widehat\beta^{-m}\) for every prescribed \(m>0\),
eventually. Formula \eqref{eq:central-remainder-estimate},
and \(g^2\le2\widehat\beta^{-1}\), prove
\eqref{eq:gaussian-unscaled-rate}.

For order zero in the centered norm, integrate the first
derivative along the segment from \(0\) to \(x\).
For order \(j\ge1\), use \(w^j\le w\). Consequently
\[
 \|f-f(0)\|_{C_w^k(B(0,w))}
 \le 2w\|f\|_{C^{k+1}(B(0,w))}.
 \]
Apply this to \(f=\log Z_D(g,\cdot)\).
\end{proof}

The estimate \eqref{eq:gaussian-centered-rate} concerns the
Gaussian part of the effective logarithm. The stationary action
and \(\frac12\log\det {\mathsf H}\) in
\eqref{eq:stationary-density-identity} have their own classical
and quadratic matching terms. Their subtraction is a separate
calculation.

\subsection{Returning from the central region to the full integral}

Let \(J_\beta(x)\) denote the pure-input integral in
\eqref{eq:actual-block-integral} after tree gauge, including its
branch sum and its actual cutoffs. Suppose its central part has
been identified with \(A_\beta(x)Z_D(g,x)\), where
\[
 A_\beta(x)=z_{\mathrm n}(\beta)^{-b_c}
 \zeta_G^{|I|}(2\pi)^{d/2}g^d
 e^{-\beta\varphi(x,\rho(x))-\psi(x,\rho(x))-\ell(x)}
 \]
in a sector with all free links active.
This identification requires that the ball is in the Haar chart,
that its fine configurations satisfy the actual small-field
cutoffs, and that the smooth central action is independent of
the branch label there. Under the last condition the finite
likelihood sum cancels by \(\sum_t\lambda_t=1\).
These conditions concern actual configurations, not just the
Hessian at the identity.

\begin{lemma}[Value comparison for a positive tail]
\label{lem:positive-integral-tail}
If \(A_\beta>0\) and
\begin{equation}
 0\le J_\beta-A_\beta Z_D\le A_\beta\epsilon_\beta,
 \label{eq:physical-tail-hypothesis}
\end{equation}
then \(J_\beta>0\), and
\begin{equation}
 0\le\log J_\beta-\log A_\beta-\log Z_D
          \le2e\,\epsilon_\beta .
 \label{eq:physical-log-tail}
\end{equation}
\end{lemma}
\begin{proof}
By \eqref{eq:central-positive}, \(Z_D\ge e^{-1}/2\).
Write \(J_\beta=A_\beta Z_D(1+t)\), where
\(0\le t\le\epsilon_\beta/Z_D\).
Now use \(0\le\log(1+t)\le t\).
\end{proof}

A small value bound alone does not control derivatives.
The following finite-difference argument supplies a short way
to obtain them when a coarse derivative bound is available.

\begin{lemma}[A small value and a coarse derivative bound]
\label{lem:finite-difference-interpolation}
Let \(f\) be real valued and \(C^{k+1}\) on \(B_E(x,r)\).
Suppose \(|f|\le\epsilon\) and
\(\|D^{k+1}f\|_{\mathrm{op}}\le {\mathsf M}\) there.
For \(k\ge1\), \(0<t<r/k\), one has
\begin{equation}
 \|D^kf(x)\|_{\mathrm{op}}
       \le2^k\epsilon t^{-k}+k{\mathsf M}t .
 \label{eq:finite-difference-interpolation}
\end{equation}
In particular, for any fixed \(q\), if all derivatives through
order \(q+1\) are bounded by a fixed power of \(\widehat\beta\)
on balls of radius \(\widehat\beta^{-a}\), and \(|f|\) is smaller
than every inverse power of \(\widehat\beta\) on those balls,
then the derivatives through order \(q\) at their centres have
the same latter property.
The constants are independent of \(\dim E\).
\end{lemma}
\begin{proof}
Fix unit vectors \(u_1,\ldots,u_k\). Repeated application of
the fundamental theorem of calculus gives
\[
 t^{-k}\Delta_{tu_1}\cdots\Delta_{tu_k}f(x)
 =\int_{[0,1]^k}
 D^kf\left(x+t\sum_i s_iu_i\right)
                      [u_1,\ldots,u_k]\,\dd s .
 \]
Here \(\Delta_vf(x)=f(x+v)-f(x)\).
The integral differs from \(D^kf(x)[u_1,\ldots,u_k]\)
by at most \(k{\mathsf M}t\), by the derivative bound along line
segments in the ball. Expanding the finite difference gives
\(2^k\) values, each of absolute value at most \(\epsilon\).
Taking the supremum over the \(u_i\) proves the estimate.

For the last assertion, suppose the coarse derivative bound is
\(\widehat\beta^b\), after increasing \(b\) to absorb a constant.
For a desired bound \(\widehat\beta^{-m}\), choose
\(t=\widehat\beta^{-s}\) with \(s>a+b+m+2\).
The term \(k{\mathsf M}t\) then has the desired decay.
Choose the inverse power in the value bound larger than
\(sk+m+2\); this bounds the other term.
Repeat for the finitely many \(k\le q\).
\end{proof}

For the actual pure density, the necessary estimate is
\eqref{eq:physical-tail-hypothesis} on a neighbourhood,
together with the derivative bounds for its logarithm.
Proving it requires coercivity outside the central region.
A spectral gap of \({\mathsf H}(x)\) only at \(\rho(x)\) does not by
itself give that coercivity.
The lemma explains precisely how a proved tail estimate enters
the Gaussian remainder; it does not supply the model-specific
tail estimate.

\begin{lemma}[Coarse derivatives of the actual positive integral]
\label{lem:actual-block-log-derivatives}
Let \(F\ge0\) be bounded, Borel, independent of the coarse
coordinates \(x\), and satisfy \(\int F\,\dd H_f>0\).
For \(V(x)_e=e^{x_e}\), equip \(E=\mathfrak g^{b_c}\) with
the maximum link norm.
For \(\beta>0\), put
\[
 J_\beta(x)=(\mathcal P_\beta F)(V(x)).
 \]
Then \(J_\beta>0\) and, on \(\|x\|<1\),
\begin{equation}
 \|D^j\log J_\beta(x)\|_{\mathrm{op}}
     \le C_{G,a,q}(1+\beta b_c)^j,
                \qquad 1\le j\le q+1 .
 \label{eq:actual-log-gross-derivative}
\end{equation}
No derivative of the Borel branch maps in the fine field
is needed for this bound. The factor \(b_c=4c^4\) is explicit.
\end{lemma}
\begin{proof}
Use \eqref{eq:physical-kernel}, and define
\[
 B_t(x,U)=a\beta\sum_e
             q_G(e^{x_e}Q_t(U)_e^{-1}).
 \]
Derivatives of the matrix exponential through any fixed
order are bounded on the unit ball in \(\mathfrak g\).
The fixed matrix representation and unitarity of \(Q_t(U)_e\)
give
\(\|D_x^jB_t(x,U)\|\le C_{G,a,q}\beta b_c\)
for \(j\le q+1\), uniformly in \(t,U\).
Also \(0\le B_t\le2a\beta b_c\).
Consequently all derivatives of \(e^{-B_t}\) through that
order have a common integrable majorant, multiplied by \(F\).
This justifies differentiation of the actual finite branch sum
and integral.

At fixed \(x\), normalise the positive measure on
\(\mathcal T\times\mathcal U_f\) with density
\(\lambda_t(U)F(U)e^{-B_t(x,U)}\).
Its mass is nonzero, since the likelihoods sum to \(1\)
and \(e^{-B_t}>0\).
The ratio \(D^jJ_\beta/J_\beta\) is the expectation of
\[
 \sum_{\pi}
 (-1)^{|\pi|}
       \prod_{A\in\pi}D^{|A|}B_t ,
 \]
evaluated on the corresponding groups of differentiated
directions; \(\pi\) ranges over partitions of
\(\{1,\ldots,j\}\).
Its norm is at most \(C_{G,a,q}(1+\beta b_c)^j\).
Repeated differentiation of \(\log J_\beta\) expresses its
\(j\)-th derivative as a finite sum of products of such
ratios whose total derivative order is \(j\).
This proves \eqref{eq:actual-log-gross-derivative}.
The noise normaliser is independent of \(x\), so it
contributes zero to these derivatives.
\end{proof}

This lemma supplies the coarse derivative bound required
by Lemma~\ref{lem:finite-difference-interpolation} for the
positive pure input. The other terms in the logarithmic
comparison, \(\log A_\beta\) and \(\log Z_D\), have the
derivative bounds already derived from the finite-chart
data. On auxiliary tori of polynomial logarithmic size,
the explicit \(b_c\) factor is a polynomial logarithmic
cost. On an arbitrary physical torus it remains a volume
factor, to be removed by the localisation argument.

\subsection{From chart derivatives to local activity derivatives}

Suppose a local function is written \(R(U)=r(x(U))\).
On its good region, assume every ordered left derivative
of the chart map through order \(q\), in directions allowed in
the norm of Section~\ref{sec:activity-norms}, obeys
\begin{equation}
 |D_{A_1}\cdots D_{A_j}x(U)|_E\le C_{\mathrm{ch}}w,
       \qquad 1\le j\le q .
 \label{eq:chart-left-derivative-bound}
\end{equation}
An ordered chain rule, proved by successively differentiating
the factors, gives
\begin{equation}
 \|R\|_{q,h,s;o,X}
 \le C_q(1+C_{\mathrm{ch}})^q
                 \|r\|_{C_w^q}.
 \label{eq:chart-to-activity}
\end{equation}
Indeed, a term with \(j\) derivatives of \(r\) contains exactly
\(j\) derivatives of the chart map, whose product is bounded
by \((C_{\mathrm{ch}}w)^j\).
Summing the finitely many partitions proves the assertion.
The masks are not differentiated.
For direct logarithmic link charts on a fixed compact subchart,
\eqref{eq:chart-left-derivative-bound} follows from bounded
derivatives of the matrix logarithm and
\[
 h(1+\operatorname{diam}X)^s\le Cw .
 \]
For charts obtained by a local gauge transformation, estimates
on the derivatives of that transformation are also needed.

There are two distinct remaining uses of locality.
First, the finite-chart dimension bound
\eqref{eq:auxiliary-dimension-bound} must concern auxiliary
tori surrounding the actual local footprint. Taking
\(d=4c^4\dim\mathfrak g\) for an arbitrarily large physical
torus would not satisfy it uniformly in volume.
Second, the resulting local corrections must be summed with
the footprint weight \(e^{\kappa n_B(X)}\), as required by
\(\|\cdot\|_{\mathcal E,\kappa}\).
Neither conclusion follows merely from
\eqref{eq:gaussian-unscaled-rate}.
The Gaussian estimate above is uniform once its local
dimension, spectral gap, and chart derivative inputs are
uniform. Sections~\ref{sec:actual-stationary}--%
\ref{sec:physical-interacting-extensions} supply these
inputs for the principal action and its local physical
extensions. Section~\ref{sec:nonlinear-induction-step}
uses the resulting local estimates in the full iteration.
\section{Quadratic extraction and coupling reset}
\label{sec:coupling-reset}

The Gaussian calculation produces logarithmic terms whose constant
and quadratic parts have to be assigned to the scalar and the
coupling. We first perform this operation on the actual local
functions, including their masks. We then compare it with the
principal action used in the density. The distinction matters because
the principal action contains cubic and quartic terms as well.
The general purpose of this extraction agrees with
\cite[Sect.~1, eqs.~(1.20)--(1.22), printed p.~264;
Sect.~5, eqs.~(5.1)--(5.15), pp.~292--293]{Bal87}.
For a comparison with local-polynomial extraction, see
\cite[Sect.~1, Proposition~1.5, preprint p.~7,
and Propositions~1.11--1.12, pp.~10--11]{BS15b}.
These references explain the purpose and organisation of the extraction.
The gauge-invariant affine extraction and its bounds are proved below
for the present lattice functions.

\subsection{The Hessian seen by affine fields}

Let \(f\) be a real, gauge invariant, local function with anchor
\(o\), footprint \(X\ni o\), and a \(C^q\) unmasked germ near
every good configuration, as in Section~\ref{sec:activity-norms}.
It need not yet have the two normalisations imposed on a regular
activity. In exponential coordinates put
\[
 \widetilde f(A)=f((e^{A_e})_e),\qquad
 H_f(A,B)=D^2\widetilde f(0)[A,B].
 \]
Both \(A\) and \(B\) are real \(\mathfrak g\)-valued link fields.
The Hessian is symmetric.

\begin{lemma}[The first two gauge identities at the identity]
\label{lem:gauge-hessian-identities}
For the simple Lie algebra in Definition~\ref{def:group},
\(D\widetilde f(0)=0\).
For a site field \(\omega:\Lambda_\ell\to\mathfrak g\), define
\[
 (\dd\omega)_i(x)=\omega(x+e_i)-\omega(x).
 \]
Then \(H_f(\dd\omega,A)=0\) for every link field \(A\).
The Hessian depends only on the values of its arguments on
links read by \(f\).
\end{lemma}
\begin{proof}
A constant gauge transformation conjugates every link by the
same element of \(G\). Hence, in each link slot, the linear
functional \(D\widetilde f(0)\) is invariant under
\(\operatorname{Ad}(G)\). Differentiating this invariance along
\(e^{tX}\) shows that it vanishes on every commutator \([X,Y]\).
The commutator subspace \([\mathfrak g,\mathfrak g]\) is a
nonzero ideal, and therefore equals \(\mathfrak g\).
Thus each such linear functional is zero.

Apply a site gauge transformation \(e^{t\omega(x)}\) to the
connection \(e^{sA_e}\). In exponential coordinates its first
order variation at \(s=t=0\) is \(sA-t\dd\omega\), with a
possible reversal of the sign according to the gauge convention.
Differentiate the gauge invariance identity once in \(s\) and
once in \(t\). The term containing the mixed derivative of the
coordinate change is multiplied by \(D\widetilde f(0)=0\).
The remaining term is \(H_f(A,\dd\omega)=0\).
Finally, locality implies that \(\widetilde f\) is constant in
all unused link coordinates, proving the last assertion.
\end{proof}

For a nonwrapping footprint use the local integer coordinates
\(\xi(x)=\operatorname{lift}(x-o)\) already used in the affine
normalisation. An affine link field has the form
\begin{equation}
 A_i(x)=a_i+\sum_{j=1}^4\xi_j(x)b_{ij},\qquad
 F_{ij}(A)=b_{ji}-b_{ij}\quad(i<j).
 \label{eq:affine-curl-reset}
\end{equation}
The six entries \(F_{ij}(A)\) form its constant linearised curl.
Write \(\mathcal F=\{(i,j):1\le i<j\le4\}\).

\begin{lemma}[Dependence on the affine curl]
\label{lem:affine-curl-hessian}
If \(X\) is nonwrapping about \(o\), \(H_f(A,B)\), for affine
fields \(A,B\), depends only on their curls
\(F(A),F(B)\in\mathfrak g^{\mathcal F}\).
\end{lemma}
\begin{proof}
It suffices to treat an affine field of zero curl.
Then \(b_{ij}=b_{ji}\). On the sites read by \(f\) and their
neighbours define
\[
 \omega(x)=\sum_i\xi_i(x)(a_i-\tfrac12b_{ii})
                  +\tfrac12\sum_{i,j}\xi_i(x)\xi_j(x)b_{ij}.
 \]
An explicit finite difference gives
\(\omega(x+e_i)-\omega(x)=A_i(x)\).
The nonwrapping condition guarantees
\(\xi(x+e_i)=\xi(x)+e_i\) at these links.
Extend \(\omega\) arbitrarily to the remaining torus sites.
The gradient agrees with \(A\) on every link read by \(f\);
it need not agree on the whole torus.
Apply Lemma~\ref{lem:gauge-hessian-identities} and locality.
Subtracting two affine fields with the same curl proves the
assertion in the first argument, and symmetry proves it in
the second.
\end{proof}

\begin{lemma}[Invariant real bilinear forms]
\label{lem:invariant-bilinear}
Every real bilinear form on the compact simple Lie algebra
\(\mathfrak g\) that is invariant under
\(\operatorname{Ad}(G)\) is a scalar multiple of
\(\langle\cdot,\cdot\rangle_{\mathfrak g}\).
\end{lemma}
\begin{proof}
Write the form as \(\langle TX,Y\rangle_{\mathfrak g}\).
Invariance implies that \(T\), and also its adjoint, commute
with all \(\operatorname{ad}Z\).
The eigenspaces of the self-adjoint part of \(T\) are ideals.
Simplicity makes that part a scalar.

Let \(A\) be the skew-adjoint part. The self-adjoint map
\(A^2\) again commutes with all \(\operatorname{ad}Z\), so
\(A^2=-\lambda^2 1\) for some \(\lambda\ge0\).
Commutation also implies
\([AX,Y]=A[X,Y]\), and hence
\(\operatorname{ad}(AX)=A\operatorname{ad}X\).
The Killing form
\(B(X,X)=\operatorname{tr}((\operatorname{ad}X)^2)\)
is negative for \(X\ne0\): \(\operatorname{ad}X\) is
skew-adjoint and nonzero, since the centre is a proper ideal.
If \(\lambda>0\), then
\[
 B(AX,AX)
 =\operatorname{tr}(A^2(\operatorname{ad}X)^2)
 =-\lambda^2B(X,X)>0 ,
 \]
which is impossible. Thus \(A^2=0\), and skew-adjointness gives
\(|AX|^2=-\langle A^2X,X\rangle=0\).
Therefore \(A=0\) and \(T\) is scalar.
\end{proof}

For \(\alpha=(i,j)\in\mathcal F\) and \(Z\in\mathfrak g\),
let \(A^{\alpha,Z}\) be the affine field whose only nonzero
component is \(A_j(x)=\xi_i(x)Z\).
Its curl is \(Z\) in entry \(\alpha\), and zero elsewhere.
The preceding lemmas give real coefficients
\(q_f(\alpha,\gamma)\) such that
\begin{equation}
 H_f(A,B)=
 \sum_{\alpha,\gamma\in\mathcal F}q_f(\alpha,\gamma)
        \langle F_\alpha(A),F_\gamma(B)\rangle_{\mathfrak g}
 \label{eq:affine-hessian-coefficients}
\end{equation}
for affine \(A,B\).
One can define them without choosing coordinates in colour
space: fix any unit \(Z\in\mathfrak g\), and set
\[
 q_f(\alpha,\gamma)
       =H_f(A^{\alpha,Z},A^{\gamma,Z}).
 \]
Lemma~\ref{lem:invariant-bilinear} proves independence of \(Z\).
If the enlarged footprint \(X\cup B_\infty(o,2)\) wraps, define
these extraction coefficients to be zero. No affine
normalisation is then asserted on that enlarged footprint.

\begin{lemma}[The extraction price]
\label{lem:quadratic-extraction-price}
If \(\|f\|_{q,h,s;o,X}\le u\), with \(q\ge2\) and \(s\ge2\), then
\begin{equation}
 h^2|q_f(\alpha,\alpha)|\le u,\qquad
 h^2|q_f(\alpha,\gamma)|\le3u .
 \label{eq:quadratic-extraction-price}
\end{equation}
\end{lemma}
\begin{proof}
The link field \(hA^{\alpha,Z}\) is an allowed direction:
its norm at \(x\) is at most \(h\,d_\infty(x,o)\), which is
at most \(h(1+d_\infty(x,o))^s\).
Along the curve \(e^{thA^{\alpha,Z}}\), the second derivative
at zero is both \(h^2q_f(\alpha,\alpha)\) and the ordered
left derivative \(D_{hA^{\alpha,Z}}^2f(1)\).
The local norm bounds it by \(u\).

For the mixed entry use the allowed direction
\(\frac h2(A^{\alpha,Z}+A^{\gamma,Z})\).
Its second derivative is
\[
 \frac{h^2}{4}
 \bigl(q_f(\alpha,\alpha)+q_f(\gamma,\gamma)
                                      +2q_f(\alpha,\gamma)\bigr).
 \]
Its absolute value is at most \(u\).
The two diagonal bounds then give
\(2h^2|q_f(\alpha,\gamma)|\le6u\).
The wrapping case has zero coefficients.
\end{proof}

\subsection{A local counterterm with the correct affine Hessian}

The curvature used in the counterterm is a smooth, based
plaquette quantity. Let \(\Pi_{\mathfrak g}\) be orthogonal
projection from skew-Hermitian matrices onto \(\mathfrak g\),
for the inner product fixed earlier, and define
\[
 \mathcal C(W)=\Pi_{\mathfrak g}\bigl((W-W^{-1})/2\bigr).
 \]
It satisfies
\(\mathcal C(1)=0\), \(D(\mathcal C\circ\exp)(0)=1\),
\(\mathcal C(W^{-1})=-\mathcal C(W)\), and
\(|\mathcal C(W)|_{\mathfrak g}\le d_G(W)\).
Projection commutes with conjugation by \(G\), since its
subspace and inner product are invariant.

For \(\epsilon,\epsilon'\in\{1,-1\}\), let
\(\gamma_{\epsilon i,\epsilon'j}\) be the based four-link loop
with steps \(\epsilon e_i,\epsilon'e_j,-\epsilon e_i,-\epsilon'e_j\).
Use the inverse of a positive link when traversing it backwards.
Define the clover curvature
\begin{equation}
 \Phi_{ij}(U;o)=\frac14
 \sum_{\epsilon,\epsilon'=\pm1}\epsilon\epsilon'
       \mathcal C\bigl(U(\gamma_{\epsilon i,\epsilon'j};o)\bigr).
 \label{eq:clover-curvature-reset}
\end{equation}
It reads links in \(B_\infty(o,2)\) and transforms by
conjugation at \(o\).
Every constituent plaquette is incident to \(o\).
Consequently
\begin{equation}
 |\Phi_\alpha(U;o)|\le h
       \quad\hbox{if }\chi_{\{o\}}(U)=1 .
 \label{eq:clover-good-bound}
\end{equation}

For an affine field \eqref{eq:affine-curl-reset}, differentiating
the four links of each loop at \(U=1\) gives
\[
 \left.\frac{\dd}{\dd t}\right|_{t=0}
 \Phi_\alpha((e^{tA_e})_e;o)=F_\alpha(A).
 \]
The constant terms in \(A\) cancel around the square, and the
two increments of its linear part give \(b_{ji}-b_{ij}\);
the factors \(\epsilon\epsilon'\) in
\eqref{eq:clover-curvature-reset} remove the orientation signs.
Define
\begin{equation}
 Q_f(U)=\frac12\sum_{\alpha,\gamma}
 q_f(\alpha,\gamma)
       \langle\Phi_\alpha(U;o),\Phi_\gamma(U;o)\rangle .
 \label{eq:local-quadratic-counterterm}
\end{equation}
Then \(Q_f(1)=0\), and its Hessian on affine fields is
exactly \eqref{eq:affine-hessian-coefficients}.

\begin{lemma}[Counterterm and residual norms]
\label{lem:local-reset-norms}
Put \(X'=X\cup B_\infty(o,2)\) and
\(\bar f=f-f(1)-Q_f\).
There is \(C_{\mathrm r}=C_{\mathrm r}(G,q,s)\), independent
of \(h\le1\), \(\ell\), and \(X\), such that
\begin{equation}
 \|\bar f\|_{q,h,s;o,X'}\le C_{\mathrm r}u .
 \label{eq:local-reset-regular-bound}
\end{equation}
The residual satisfies \(\bar f(1)=0\), and its Hessian on
affine fields is zero whenever \(X'\) is nonwrapping.
\end{lemma}
\begin{proof}
An ordered derivative of a four-link word is a sum of products
in which a derivative inserts one of the direction matrices
at a link, or differentiates such an inserted factor at a
subsequent step. At these links the allowed directions have
norm at most \(h3^s\).
For every positive derivative order at most \(q\), the norm
of the result is at most \(C_{G,q,s}h\), since \(h^j\le h\)
for \(j\ge1\). At order zero use
\eqref{eq:clover-good-bound}.
Thus every ordered derivative through order \(q\) of each
\(\Phi_\alpha\), on a good anchor, is bounded by
\(C_{G,q,s}h\).
The product rule gives
\[
 |D_{A_1}\cdots D_{A_j}\langle\Phi_\alpha,\Phi_\gamma\rangle(U)|
       \le C_{G,q,s}h^2\quad(0\le j\le q,\ \chi_{\{o\}}(U)=1).
 \]
Combine this with \eqref{eq:quadratic-extraction-price} and
the \(6^2\) terms in \eqref{eq:local-quadratic-counterterm}.
This bounds \(Q_f\) by \(C_{G,q,s}u\) on the good anchor.
On the good region of \(X'\), \(f\) retains its old norm
bound and \(|f(1)|\le u\).
The triangle inequality proves
\eqref{eq:local-reset-regular-bound}.
The normalisations follow from the definitions and the
equality of affine Hessians.
\end{proof}

\subsection{The mask discrepancy is retained exactly}

Let \((f_i,X_i,o_i)\) be the preliminary logarithmic family.
Define \(X_i'=X_i\cup B_\infty(o_i,2)\) and
\begin{equation}
 u_i(U)=
 \begin{cases}
 0,&\chi_{X_i'}(U)=1,\\
 f_i(1)+\chi_{\{o_i\}}(U)Q_{f_i}(U)
                       -\chi_{X_i}(U)f_i(U),&\chi_{X_i'}(U)=0.
 \end{cases}
 \label{eq:reset-mask-discrepancy}
\end{equation}
This is a Borel, gauge invariant function determined by the
links and flags in \(X_i'\). Its link-only read footprint can
be enlarged as in Section~\ref{sec:activities}.
It vanishes when \(X_i'\) is good.

\begin{lemma}[The exact masked subtraction]
\label{lem:exact-masked-subtraction}
For every configuration \(U\),
\begin{equation}
 \chi_{X_i}f_i
 =f_i(1)+\chi_{\{o_i\}}Q_{f_i}
                  -u_i+\chi_{X_i'}\bar f_i .
 \label{eq:exact-masked-subtraction}
\end{equation}
Also \(\|u_i\|_\infty\le C_{\mathrm r}'u_i^0\) if
\(\|f_i\|_{q,h,s;o_i,X_i}\le u_i^0\), where
\(C_{\mathrm r}'\) depends only on \(G,q,s\).
\end{lemma}
\begin{proof}
On \(\chi_{X_i'}=1\), both smaller masks equal \(1\).
The identity is then \(\bar f_i=f_i-f_i(1)-Q_{f_i}\).
On its complement the last term is zero and
\eqref{eq:reset-mask-discrepancy} is exactly the stated
identity.
In either case the value \(f_i(1)\) is bounded by \(u_i^0\).
The masked old term is bounded by \(u_i^0\), and the masked
counterterm is bounded by \(C_{G,q,s}u_i^0\) by the proof
of Lemma~\ref{lem:local-reset-norms}.
\end{proof}

Assume the family is invariant under translations and all
signed permutations of the four axes, by reindexing the
same family of labels. Assume also the stronger anchored
bound
\begin{equation}
 \sup_o\sum_{i:o_i=o}
 e^{(\kappa+2)n_B(X_i)}u_i^0\le H .
 \label{eq:preliminary-reset-budget}
\end{equation}
The stronger exponent will pay for enlargement of footprints
and passage from anchored sums to incidence sums.
Absolute convergence justifies summing the coefficients.
Put
\[
 q_o(\alpha,\gamma)=\sum_{i:o_i=o}q_{f_i}(\alpha,\gamma).
 \]

\begin{proposition}[One coupling coefficient on the same torus]
\label{prop:isotropic-reset-coefficient}
There is a real number \(\Delta\), independent of the anchor
and plaquette direction, such that
\[
 q_o(\alpha,\gamma)=\Delta\,1_{\{\alpha=\gamma\}},
 \qquad h^2|\Delta|\le H .
 \]
Consequently, with
\begin{equation}
 S_h(U)=\frac12\sum_o\chi_{\{o\}}(U)
                         \sum_{\alpha\in\mathcal F}|\Phi_\alpha(U;o)|^2,
 \label{eq:clover-kinetic-reset}
\end{equation}
one has the exact identity
\begin{equation}
 \sum_i\chi_{X_i}f_i
 =C_0+\Delta S_h-\sum_i u_i+\sum_i\chi_{X_i'}\bar f_i,
 \qquad C_0=\sum_i f_i(1).
 \label{eq:family-reset-splitting}
\end{equation}
\end{proposition}
\begin{proof}
Translations make \(q_o\) independent of \(o\).
Under reflection of axis \(j\), a curl entry changes sign
exactly when its index pair contains \(j\).
For distinct pairs \(\alpha,\gamma\), choose an axis in their
symmetric difference. Invariance then makes
\(q_o(\alpha,\gamma)\) equal to its negative, so it is zero.
Permutations of the axes make all diagonal entries equal.
These transformation rules also hold when positive links
are reoriented by a reflection: the resulting affine field
has an additional constant term, and
Lemma~\ref{lem:affine-curl-hessian} removes that term.
Nonwrapping is preserved by these symmetries.
This proves the stated scalar form before applying any masks.

Fix a pair \(\alpha\) and an anchor \(o\).
Lemma~\ref{lem:quadratic-extraction-price} gives
\[
 h^2|\Delta|
 \le\sum_{i:o_i=o}h^2|q_{f_i}(\alpha,\alpha)|
 \le\sum_{i:o_i=o}u_i^0\le H .
 \]
The same bounds prove absolute convergence of all coefficient
series. Sum \eqref{eq:exact-masked-subtraction}; at each
anchor the counterterms have the common mask
\(\chi_{\{o\}}\), so their sum is \(\Delta S_h\).
This proves \eqref{eq:family-reset-splitting}.
\end{proof}

\subsection{From anchored bounds to incidence bounds}

\begin{lemma}[A uniform geometric summation]
\label{lem:anchor-incidence}
Let \(B\ge2\), and suppose \(a_i\ge0\), \(o_i\in X_i\), and
\(\sup_o\sum_{i:o_i=o}e^{(\kappa+2)n_B(X_i)}a_i\le H\).
For \(X_i'=X_i\cup B_\infty(o_i,2)\),
\begin{align}
 \sup_o\sum_{i:o_i=o}e^{\kappa n_B(X_i')}a_i
       &\le e^\kappa H, \label{eq:anchor-enlargement}\\
 \sup_x\sum_{i:x\in X_i'}
        e^{(\kappa+1)n_B(X_i')}a_i
       &\le C_\kappa(1+B)^4H .
 \label{eq:anchor-to-incidence}
\end{align}
The constants are independent of the torus side.
\end{lemma}
\begin{proof}
An additional \(B\)-ball centred at \(o_i\) covers the added
halo and intersects a covering of \(X_i\).
Thus \(n_B(X_i')\le n_B(X_i)+1\), proving
\eqref{eq:anchor-enlargement}.
The diameter of a connected union of \(n\) lattice
\(B\)-balls is at most \((2B+1)n\).
Indeed, choose an intersection-or-neighbour chain of balls
between the two points; each step crosses at most \(2B+1\)
in the maximum metric.
For \(x\in X_i'\), this gives
\[
 d_\infty(x,o_i)\le(2B+1)n_B(X_i)+2 .
 \]
Reserve one unit of the exponent
\((\kappa+2)-(\kappa+1)=1\).
The left side of \eqref{eq:anchor-to-incidence}, at fixed \(x\),
is bounded by
\[
 e^{\kappa+1+2/(2B+1)}H
   \sum_o e^{-d_\infty(x,o)/(2B+1)} .
 \]
The number of torus sites at distance \(j\) is at most
\(C(1+j)^3\), uniformly in the side, by using their
representatives of minimum coordinate lengths.
The last sum is at most \(C(1+B)^4\), by comparison with
the corresponding exponential integral on \([0,\infty)\).
\end{proof}

Apply this lemma with \(a_i=u_i^0\). The residual regular
family is bounded in \(\mathcal E,\kappa\) by
\(C_{G,q,s,\kappa}H\).
The functions \(u_i\) in
\eqref{eq:reset-mask-discrepancy} have total weighted incidence
bound at exponent \(\kappa+1\) at most
\(C_{G,q,s,\kappa}(1+B)^4H\).
Both assertions are volume uniform. In addition,
\(|C_0|\le\ell^4H\), so the change of scalar per site
will be bounded by \(H\).

\subsection{The principal action that the coupling multiplies}

The coupling multiplies a particular principal action,
not just \(S_h\).
We describe the path data needed for the reset calculation.
A transported curvature slot, anchored at \(o\), is
\begin{equation}
 Z_{\gamma,\alpha}(U;o)=
 \operatorname{Ad}_{U(\gamma;o)}
       \mathcal C(P_\alpha(U;o+\operatorname{end}\gamma)).
 \label{eq:transported-curvature-slot}
\end{equation}
Its length is \(|\gamma|+4\), and its footprint consists of the
visited sites and the final plaquette. It is a covariant
\(\mathfrak g\)-valued function based at \(o\).
Let
\[
 \ell_b(U;o)=\sum_j r_{bj}Z_{\gamma_{bj},\alpha_{bj}}(U;o)
       \quad(b\in\mathcal F).
 \]
Fix \(\sigma>0\) and a path budget \(C_{\mathrm p}\) with
\begin{equation}
 \sum_j|r_{bj}|e^{\sigma(|\gamma_{bj}|+4)}\le C_{\mathrm p}.
 \label{eq:row-path-budget}
\end{equation}
For cubic and quartic terms take invariant multilinear
polynomials of three or four transported slots, with, at each
anchor, the sum
of the absolute values of their coefficients, multiplied by
the exponential of their total length, at most
\(C_{\mathrm p}\).
Denote these local polynomials by \(p_\nu\), their footprints
by \(Y_\nu\), and their degree by \(k_\nu\in\{3,4\}\).
All series below converge absolutely by these budgets.

Fix positive parameters \(\lambda,\lambda_{\mathrm l},r_0\).
The clipping of a plaquette of chord \(s\) is
\[
 c_h(s)=
 \begin{cases}
 \lambda s^2/2,&s\le h,\\
 \lambda_{\mathrm l}hs,&s>h.
 \end{cases}
 \]
For the row term use the predicate
\begin{equation}
 \mathcal H_o(U):
 \quad d_G(P_\alpha(U;x))
          \le h e^{d_\infty(o,x)/r_0}
          \quad\hbox{for all }x,\alpha .
 \label{eq:row-relaxed-cutoff}
\end{equation}
The principal action is
\begin{equation}
 S(U)=\Bigl\langle
 \sum_{o,\alpha}c_h(d_G(P_\alpha(U;o)))
 +\frac12\sum_{o,b}1_{\mathcal H_o(U)}|\ell_b(U;o)|^2
 +\sum_\nu\chi_{Y_\nu}(U)p_\nu(U)
 \Bigr\rangle_{\mathrm{sym}}.
 \label{eq:actual-reset-principal}
\end{equation}
The brackets mean the arithmetic average over the same fixed
signed permutations of the whole expression, including its
cutoffs. In the local-family formulas below, orientation
labels, the common averaging factor, and translation labels
are included in the indices and coefficients.

Define the zero-momentum row matrix
\[
 R_{b\alpha}=\sum_{j:\alpha_{bj}=\alpha}r_{bj},
 \]
and impose its metric normalisation
\begin{equation}
 \lambda\,1+R^TR=1
          \quad\hbox{on }\R^{\mathcal F}.
 \label{eq:row-metric-normalization}
\end{equation}
This is a normalisation of the principal path data.
It is a property to be maintained by their construction,
not a consequence of summability.

Expand the quadratic row squares in their two path labels.
Index their summands, and the direct plaquette summands, by
\(j\). Write their smooth unmasked functions as \(t_j\),
their path footprints as \(L_j\), and their anchors as \(o_j\).
Explicitly they are
\[
 \lambda q_G(P_\alpha(U;o)),\qquad
 \tfrac12r_{bi}r_{bj}
          \langle Z_{\gamma_{bi},\alpha_{bi}},
                  Z_{\gamma_{bj},\alpha_{bj}}\rangle .
 \]
Let \(\widetilde t_j\) be the corresponding actual summand
in the first two terms of \eqref{eq:actual-reset-principal}:
use the clipping in the first case and
\(1_{\mathcal H_o}\) in the second.
Finally, let \(c_j\) be its clover analogue:
\[
 \tfrac{\lambda}{2}|\Phi_\alpha(U;o)|^2,\qquad
 \tfrac12r_{bi}r_{bj}
                  \langle\Phi_{\alpha_{bi}}(U;o),
                         \Phi_{\alpha_{bj}}(U;o)\rangle ,
 \]
respectively. These \(c_j\) need only the common anchor
halo, whatever the lengths of their original paths.
Metric normalisation gives
\begin{equation}
 \sum_{j:o_j=o}c_j(U)
             =\tfrac12\sum_\alpha|\Phi_\alpha(U;o)|^2.
 \label{eq:clover-analogue-sum}
\end{equation}
To see it, expand the square
\(\sum_b|\sum_\alpha R_{b\alpha}\Phi_\alpha|^2\)
and use \eqref{eq:row-metric-normalization}.
Absolute convergence follows from
\eqref{eq:row-path-budget}.
This identity is preserved by the signed average.
Put
\begin{equation}
 S_{0,h}=\sum_j\chi_{L_j}t_j .
 \label{eq:masked-second-principal}
\end{equation}
It is this masked quadratic principal term that carries the
one-loop coefficient before the coupling reset.

\subsection{Exact transfer to the prescribed principal action}

Fix the already chosen one-loop coefficient \(\alpha\).
The same value of \(\alpha\) is used on all periods at the
given scale; no infinite-volume replacement is made during
the following algebra.
Define the regular correction
\begin{equation}
 R^\#_{\alpha,\Delta}
 =\Delta\sum_j\chi_{L_j'}(c_j-t_j)
          -(\alpha+\Delta)\sum_\nu\chi_{Y_\nu}p_\nu,
 \qquad L_j'=L_j\cup B_\infty(o_j,2).
 \label{eq:principal-reset-regular}
\end{equation}
Enlarging the mask from \(L_j\) to \(L_j'\) requires its own
correction. It is convenient to put it in the exact local
functions
\begin{align}
 v_j={}&(\alpha+\Delta)\widetilde t_j
       +\Delta\chi_{L_j'}(c_j-t_j)
       -\alpha\chi_{L_j}t_j
       -\Delta\chi_{\{o_j\}}c_j .
 \label{eq:principal-reset-forced}
\end{align}
Each \(v_j\) vanishes when all its read and flag footprints
are good.

\begin{lemma}[Exact principal-action transfer]
\label{lem:exact-principal-transfer}
For every \(U\),
\begin{equation}
 (\alpha+\Delta)S+R^\#_{\alpha,\Delta}
            =\alpha S_{0,h}+\Delta S_h+\sum_j v_j .
 \label{eq:exact-principal-transfer}
\end{equation}
Every regular summand in
\eqref{eq:principal-reset-regular} vanishes at \(1\) and has
zero Hessian on affine fields whenever its footprint is
nonwrapping.
\end{lemma}
\begin{proof}
Inserting \eqref{eq:actual-reset-principal} cancels its cubic
and quartic terms against the second term of
\eqref{eq:principal-reset-regular}.
The remaining identity is the sum of the definitions
\eqref{eq:principal-reset-forced}, using
\eqref{eq:clover-analogue-sum} and
\eqref{eq:masked-second-principal}.

At the identity all transported curvatures vanish.
Their first variations on affine fields are their plaquette
curls, independent of the transporting paths.
Thus \(c_j\) and \(t_j\) have identical quadratic variations.
Their difference has zero affine Hessian and value zero.
A polynomial of three or four curvature slots also has
value zero and Hessian zero there.
Polarisation gives the mixed affine Hessian assertion.
\end{proof}

\subsection{Uniform prices for this transfer}

We give the bounds including the extra cutoff
\eqref{eq:row-relaxed-cutoff}. Put
\begin{equation}
 R_h=\left\lceil r_0\log(2/h)\right\rceil+2 .
 \label{eq:relaxed-cutoff-radius}
\end{equation}
Since every plaquette chord is at most \(2\),
\(\mathcal H_o\) is determined by links in the halo of
\(B_\infty(o,R_h)\). This is a local radius depending on
the coupling threshold; it is not the physical torus side.
If that ball wraps, it is still interpreted as a subset
of the same torus, with no new period.

\begin{lemma}[Path derivative and footprint bounds]
\label{lem:path-reset-bounds}
Fix \(G,q,s,\sigma,\kappa\). Suppose
\eqref{eq:row-path-budget} and the analogous cubic and
quartic budgets hold, and \(r_0^{-1}<\sigma/4\).
Choose \(B\ge2\) sufficiently large in terms of
\(\sigma,\kappa\), and assume \(R_h\le B\).
There is \(C_{\mathrm p}'\), depending only on these fixed
parameters and \(C_{\mathrm p},\lambda,\lambda_{\mathrm l},r_0\),
such that
\begin{align}
 \|R_{\alpha,\Delta}\|_{\mathcal E,\kappa}
       &\le C_{\mathrm p}'(H+|\alpha|h^3),
       \label{eq:principal-transfer-regular-price}\\
 \sup_x\sum_{j:x\in V_j}
   e^{(\kappa+1)n_B(V_j)}\|v_j\|_\infty
       &\le C_{\mathrm p}'(1+B+R_h)^m
                            (|\alpha|h+H/h).
       \label{eq:principal-transfer-forced-price}
\end{align}
Here \(h^2|\Delta|\le H\), \(m\) is a fixed integer, and
\(V_j\) is the union of \(L_j'\), the actual plaquette read
footprint, and \(B_\infty(o_j,R_h+2)\).
The regular family in the first line consists of the
unmasked functions in \eqref{eq:principal-reset-regular}.
The constants are independent of the cutoff and the torus
side.
\end{lemma}
\begin{proof}
Let a slot have length \(l\). All its visited sites are at
distance at most \(l\) from its anchor. Every direction
inserted by an ordered left derivative has norm at most
\(h(1+l)^s\) at these links.
Differentiating a word of length at most \(2l\) through
order \(q\) gives at most \((2l)^q\) terms times a constant
depending on \(q\). Unitary factors have norm \(1\).
Derivatives of inverses and of conjugation have the same
type of bound.
On a good footprint the undifferentiated curvature has
norm at most \(h\); each positive derivative order has a
bound \(C_{G,q,s}(1+l)^m h\).
Thus, for a polynomial of \(k\) slots of total length \(l\)
and coefficient magnitude \(a\), all its derivatives
through order \(q\) are bounded by
\begin{equation}
 C_{G,q,s}a(1+l)^m h^k .
 \label{eq:path-polynomial-derivative-price}
\end{equation}
Every factor, whether differentiated or not, supplies one
power of \(h\). The same argument gives \(Ch^2\) for the
direct plaquette price and for every clover pair.

The union of a fixed number of based walks of total
length \(l\) has a connected covering by at most
\(C+l/B\) balls of radius \(B\): place centres successively
along each walk at intervals of at most \(B\) steps.
Added balls at their common anchor keep the cover
connected. Hence its weight at exponent \(\kappa+1\)
is at most \(C e^{(\kappa+1)l/B}\).
Choose \(B\) so that this length exponent is smaller than
\(\sigma/4\). Polynomial factors in \(l\) can be
absorbed in another \(e^{\sigma l/4}\).
For a row pair its coefficient sum has the exponential
bound supplied by the product of two row budgets.
For higher terms use their stated exponential budget.
These observations show that the sums of the prices in
\eqref{eq:path-polynomial-derivative-price}, with their
footprint weights, are bounded uniformly per anchor.

The low regular terms in
\eqref{eq:principal-reset-regular} therefore cost at most
\(C|\Delta|h^2\le CH\).
The cubic and quartic terms cost at most
\(C|\alpha+\Delta|h^3
 \le C(|\alpha|h^3+Hh)\).
This proves \eqref{eq:principal-transfer-regular-price}.

For the value bound on \(v_j\), the direct clipped
plaquette satisfies \(|\widetilde t_j|\le Ch\).
Every masked direct \(t_j\) and its masked clover
analogue is at most \(Ch^2\).
These terms contribute at most
\[
 C\bigl((|\alpha|+|\Delta|)h+|\Delta|h^2\bigr)
          \le C(|\alpha|h+H/h).
 \]
For a row pair of total length \(l\), on \(\mathcal H_o\)
the two curvatures are bounded by
\(h e^{l_1/r_0}\) and \(h e^{l_2/r_0}\), respectively.
Transport does not change their norms. Thus its actual
value is at most \(Ca h^2e^{l/r_0}\).
Masked smooth terms have bound \(Ca h^2\), and the same
holds for the clover analogue on the good anchor.
If a smaller mask vanishes, its term is simply zero.
The resulting bound on \(v_j\) is
\(Ca e^{l/r_0}(|\alpha|h^2+H)\).
This is bounded by the preceding displayed price,
multiplied by \(Ca e^{l/r_0}\).
The assumption \(r_0^{-1}<\sigma/4\) leaves an exponential
length reserve for summation.

The added ball in \(V_j\) has a covering by a fixed number
of \(B\)-balls, connected to the original walk cover,
because \(R_h\le B\) and \(B\ge2\). Its cost is therefore
a constant in the exponent. In particular,
\[
 n_B(V_j)\le C+l/B .
 \]
For fixed \(x\), summing over translations of a
footprint containing \(x\) costs at most its cardinality,
bounded by \(C((1+R_h)^4+(1+l)^m)\).
The length polynomial is paid by the same exponential
reserve; the radius polynomial is retained explicitly.
Signed orientation averaging costs a fixed factor
depending on \(q,s\) in the derivative bounds and no
new volume factor.
This proves \eqref{eq:principal-transfer-forced-price}.
\end{proof}

\begin{corollary}[Large-coupling prices for reset]
\label{cor:reset-large-coupling}
Let \(\widehat\beta\to\infty\), with fixed path and norm
parameters as in Lemma~\ref{lem:path-reset-bounds}.
Fix a decay exponent \(1<\vartheta<3/2\), and suppose
\[
 h=\widehat\beta^{-1/2}P(\widehat\beta),\quad
 B\asymp(1+\log\widehat\beta)^2,\quad
 |\alpha|\le C_\nu\widehat\beta^\nu,\quad
 H\le C_\nu\widehat\beta^{-\vartheta+\nu}
 \quad(\nu>0).
 \]
Here \(B\asymp(1+\log\widehat\beta)^2\) means two fixed
positive multiples bound their ratio, and
\(P(\widehat\beta)\) is defined in
\eqref{eq:monitor-and-chart-width}.
Then the new regular family, consisting of the residuals
\(\bar f_i\) and the transfer terms, has for every
\(\nu>0\) the bound
\begin{equation}
 \|E_{\mathrm{new}}\|_{\mathcal E,\kappa}
            \le C_\nu\widehat\beta^{-\vartheta+\nu}.
 \label{eq:reset-regular-rate}
\end{equation}
The combined soft corrections \(u_i,v_j\) have weighted
incidence norm at exponent \(\kappa+1\) at most
\(\widehat\beta^{-1/8}\), eventually.
Also
\begin{equation}
 |\Delta|\le H/h^2
       \le C_\nu\widehat\beta^{1-\vartheta+\nu}P(\widehat\beta)^{-2}.
 \label{eq:reset-coupling-rate}
\end{equation}
All lower thresholds and constants have the same
cutoff and volume independence as the input budgets.
\end{corollary}
\begin{proof}
Equation \eqref{eq:relaxed-cutoff-radius} gives
\(R_h=O(1+\log\widehat\beta)\), hence \(R_h\le B\)
eventually. Every radius polynomial in
\eqref{eq:principal-transfer-forced-price} is subpower.
The regular price is \(C(H+|\alpha|h^3)\), and
\(h^3=\widehat\beta^{-3/2}P^3\).
Use \eqref{eq:monitor-growth} to prove
\eqref{eq:reset-regular-rate}.
The soft price is bounded by a subpower factor times
\[
 H+|\alpha|h+H/h
 \le C_\nu\bigl(
 \widehat\beta^{-\vartheta+\nu}
 +\widehat\beta^{-1/2+\nu}P
 +\widehat\beta^{1/2-\vartheta+\nu}P^{-1}\bigr).
 \]
Its decay is stronger than \(\widehat\beta^{-1/8}\).
Finally divide \(h^2|\Delta|\le H\) by \(h^2\).
\end{proof}

For a new remainder of order
\(\widehat\beta^{-23/20+\nu}\), take \(\vartheta=23/20\).
If the preliminary family also contains transported input
activities at the budget \(\delta\asymp\widehat\beta^{-11/10}\),
their contribution is included in \(H\); it cannot simply be
assigned the stronger decay exponent. Their constant reserve
has to be estimated separately. The bounded, linear extraction
and the bound \(C(H+|\alpha|h^3)\) above allow this separation.
A power margin for newly generated remainders alone does
not prove contraction of the transported input budget.

\subsection{The exact reset density, scalar, and two tests}

\begin{theorem}[Coupling reset as an identity of densities]
\label{thm:exact-coupling-reset}
Suppose the actual preliminary density at side \(\ell\) is
\begin{equation}
 J(U)=
 e^{\tau\ell^4-\beta S(U)-I(U)
                 -\sum_i\chi_{X_i}(U)f_i(U)-\alpha S_{0,h}(U)}
 Z_K(U),
 \label{eq:pre-reset-actual-density}
\end{equation}
with the convergent families and symmetries described above.
Let \(\Delta,C_0\) be those in
Proposition~\ref{prop:isotropic-reset-coefficient}, and set
\begin{equation}
 \beta'=\beta+\alpha+\Delta,\qquad
 \tau'=\tau-C_0/\ell^4 .
 \label{eq:actual-coupling-and-scalar-reset}
\end{equation}
Take the new regular family to be the residuals
\(\bar f_i\) together with
\eqref{eq:principal-reset-regular}.
Let the soft functions be all \(u_i\) and \(v_j\).
Whenever the hypotheses of the exact reassembly theorem
are satisfied for these functions, its resulting forced
family \(K'\) gives
\begin{equation}
 J(U)=e^{\tau'\ell^4-\beta'S(U)-I(U)
                                -E_{\mathrm{new}}^\#(U)}Z_{K'}(U)
                  \quad\hbox{for every }U .
 \label{eq:exact-reset-density}
\end{equation}
The regular affine and constant normalisations hold on the
new footprints. Furthermore,
\[
 h^2|\beta'-\beta-\alpha|\le H,\qquad
 |\tau'-\tau|\le H .
 \]
\end{theorem}
\begin{proof}
Insert \eqref{eq:family-reset-splitting} into the exponent
of \eqref{eq:pre-reset-actual-density}.
The factor \(e^{-C_0}\) changes the scalar precisely to
\(\tau'\). Apply
\eqref{eq:exact-principal-transfer} with its signs:
\[
 -\alpha S_{0,h}-\Delta S_h
 =-(\alpha+\Delta)S-R^\#_{\alpha,\Delta}
                                      +\sum_jv_j .
 \]
Thus the exponent becomes
\[
 \tau'\ell^4-\beta'S-I-E_{\mathrm{new}}^\#
                                  +\sum_i u_i+\sum_jv_j .
 \]
All sums converge absolutely, by the bounds already
proved. The exact reassembly identity absorbs
\(\exp(\sum_i u_i+\sum_jv_j)\) into \(Z_K\), yielding
\eqref{eq:exact-reset-density}.
The regular normalisations follow from
Lemmas~\ref{lem:local-reset-norms} and
\ref{lem:exact-principal-transfer}.
The two numerical bounds were proved in
Proposition~\ref{prop:isotropic-reset-coefficient} and
after Lemma~\ref{lem:anchor-incidence}.
\end{proof}

The scalar in this theorem includes whatever factors have
already entered \(\tau\) from the Gaussian, Haar, and noise
normalisations. Those factors cannot be discarded when
identifying \(\tau'\).

\begin{proposition}[Two evaluations identify this reset]
\label{prop:two-test-reset}
Assume \(S(1)=I(1)=0\), all new regular terms vanish at \(1\),
and \(V_*\) is good with \(S(V_*)>0\).
Freeze the new regular and forced families at the values
constructed above. Among trial densities
\[
 J_{d,v}(U)=
 e^{v\ell^4-(\beta+\alpha+d)S(U)-I(U)
                               -E_{\mathrm{new}}^\#(U)}Z_{K'}(U),
 \]
the equalities
\(J_{d,v}(1)=J(1)\), \(J_{d,v}(V_*)=J(V_*)\)
hold if and only if \(v=\tau'\) and \(d=\Delta\).
\end{proposition}
\begin{proof}
At a good configuration the forced gas is \(1\).
In particular \(J_{d,v}(1)=e^{v\ell^4}\), so the first
equality gives \(v=\tau'\).
By \eqref{eq:exact-reset-density}, \(J(V_*)>0\).
After taking logarithms, the second equality becomes
\((d-\Delta)S(V_*)=0\), and hence \(d=\Delta\).
The converse is \eqref{eq:exact-reset-density}.
\end{proof}

This last argument varies only the coupling after freezing
the actual output activities. Recomputing the activities
with a new trial value would be a different calculation.
The coefficient \(\Delta\) is a bounded linear functional of
the preliminary local germs: its continuity, and the bound
for its change, follow from
\eqref{eq:quadratic-extraction-price} applied to a difference
of two families with common labels.
No monotonicity of the Gaussian integral in a trial
parameter is needed for this identification.

The present section proves the subtraction, mask correction,
principal-action transfer, and scalar and coupling identities,
with their local norm prices.
To close a full step one must still construct the preliminary
family in \eqref{eq:pre-reset-actual-density} from
\eqref{eq:actual-block-integral}, with the stronger budget
\eqref{eq:preliminary-reset-budget}, the same path data, and
the proved physical tail and localisation bounds.
The finite-chart Gaussian theorem alone does not supply
that construction.
\section{The stationary graph for the lattice action}
\label{sec:actual-stationary}

We now verify the stationary-point and Hessian inputs of
Section~\ref{sec:gaussian-remainder} for the principal lattice
action. The estimates concern one step with its block size
fixed. Constants may depend on that block size; they will be
independent of the number of blocks and of the ultraviolet
cutoff.

For historical context, \cite[Sect.~1, printed p.~260]{Bal87}
attributes the constrained background-field construction to
Theorem~1 of \cite{Bal85v}. The original paper \cite{Bal85v}
has not been checked here; this attribution and theorem locator
are taken from \cite{Bal87}. This is a historical comparison.
The stationary graph for the blocking action used here is
constructed in this section.

\subsection{Which part of the Wilson integral is being estimated}

There is an important distinction between the full effective
density and its positive small-field reference contribution.
For the fine threshold \(h\), set
\[
 \chi_{\mathrm{sf}}(U)=1_{\{\mathcal B_h(U)=\varnothing\}},
 \qquad
 F_{\mathrm{sf},\beta}(U)
   =\chi_{\mathrm{sf}}(U)e^{-\beta S_f^{\mathrm{un}}(U)-I_f(U)} .
 \]
Here \(S_f^{\mathrm{un}}\) is the unmasked principal action:
it uses the direct Wilson plaquette price, the row squares,
and the unmasked cubic and quartic path polynomials of
Section~\ref{sec:coupling-reset}. It agrees with the prescribed
clipped and masked action \(S_f\) on the all-good set.
\(I_f\) is the quadratic path-loop action independent of
\(\beta\).
At the initial Wilson measure the direct coefficient is
\(\lambda=1\), the row and higher-degree path terms are absent,
and \(I_f=0\). Subsequent profiles are those of the same
exact blocking recursion.

For the same effective density \(F\) as before, the following
is an exact decomposition:
\begin{align}
 e^{-vf^4}F
 ={}&F_{\mathrm{sf},\beta}
   +F_{\mathrm{sf},\beta}(e^{-E^\#}-1)\notag\\
 &+(1-\chi_{\mathrm{sf}})
             e^{-\beta S_f-I_f-E^\#}Z_K .
 \label{eq:pure-reference-splitting}
\end{align}
Indeed \(Z_K=1\) on the all-good set.
The blocking operator is applied to all three terms.
The first supplies the stationary and Gaussian reference
calculation. The second retains the incoming regular
activities, and the third retains the large-field sectors.
Thus restricting the reference integral does not replace the
Wilson family by a different measure.

In particular, the pure integral used for the coercivity
estimates below contains \(\chi_{\mathrm{sf}}\).
An estimate for it is not an estimate for the full density
with this indicator removed. This specifies the reference
input in the model-specific use of
\eqref{eq:physical-tail-hypothesis}.

\subsection{The local averaging maps on their regular branches}

We describe the part of the actual branch maps needed here.
Fix an odd block side \(L\), an integer \(J\ge4\), and the
positive noise and branch parameters of the kernel.
The patch about a coarse root is a translate of
\[
 \mathcal Q_R=\{z\in\Z^4:|z|_\infty\le R\},
 \qquad R=(J+3)L .
 \]
All paths in this paragraph are on this integer cover.
The periods are large enough that the fixed enlarged patches
used below lift without identifications.
For two sites \(z,z'\) and a permutation \(\pi\) of the four
axes, \(\gamma^\pi_{z,z'}\) is the coordinate path that changes
the coordinates in the order \(\pi\).

A frame is a collection \(h_z\in G\), \(z\in\mathcal Q_R\),
with \(h_0=1\). Its score is
\begin{equation}
 \mathcal A_U(h)=
 \frac{1}{|\mathcal Q_R|^2\,4!}
 \sum_{z,z',\pi}
 q_G\bigl(h_zU(\gamma^\pi_{z,z'})h_{z'}^{-1}\bigr).
 \label{eq:frame-score-actual}
\end{equation}
The regular branch uses its minimiser.
The other, axial branches make the kernel defined at all
configurations. They retain their original finite
likelihoods. We will show that sufficiently small plaquettes
put all the patches considered here on the regular branch,
so the axial branches do not enter the stationary equation.

The connector of a coarse link is the collection
\[
 Y_z=h^{\,x}_z\,
 U(\gamma_{Lx+z,L(x+e_i)+z})\,
 (h^{\,x+e_i}_z)^{-1}.
 \]
The middle path has \(L\) positive steps in direction \(i\);
the frames are based at the two coarse roots.
The weights \(w(z)\) are the \(J\)-fold convolution of the
uniform measure on the centered \(L\)-block. Thus
\(w(z)\ge0\), \(\sum_z w(z)=1\), and their support is in
\(\mathcal Q_R\).
The regular predicted link minimises
\begin{equation}
 \mathcal A_Y(u)=\sum_z w(z)q_G(Y_zu^{-1}),\qquad u\in G .
 \label{eq:predictor-score-actual}
\end{equation}
These definitions use one frame per coarse site, shared by
all its incident links.

Here are the precise open tests for these regular branches.
Let \(c_{\rm fr},c_{\rm pr}>0\) be the fixed frame and
predictor thresholds. For the root-based axial transport
\(a_z^\pi=U(\gamma^\pi_{0,z})\), the frame test is
\begin{equation}
 \max_{\pi,\pi',z,z'}
 d_G\!\left(a_z^\pi U(\gamma^{\pi'}_{z,z'})
                         (a_{z'}^\pi)^{-1}\right)<c_{\rm fr}.
 \label{eq:actual-frame-branch-test}
\end{equation}
If it holds, the frame is the rooted minimiser of
\eqref{eq:frame-score-actual}, independently of its
otherwise retained axial tag. The predictor test is
\begin{equation}
 \min_{u\in G}\mathcal A_Y(u)<c_{\rm pr}.
 \label{eq:actual-predictor-branch-test}
\end{equation}
If it holds, the predictor is the minimiser of
\eqref{eq:predictor-score-actual}; otherwise the predictor
tag selects one of the \(Y_z\).
The finite tag likelihoods are uniform over the unused
tags on the regular branches.
Section~\ref{sec:axial-fallback} specifies the global
fallback probabilities, proves that their eligible sets
are nonempty, and verifies their exact normalisation.

\begin{lemma}[A uniform regular branch on a fixed patch]
\label{lem:fixed-patch-branch}
For fixed \(G,L,J\) and branch parameters, there is
\(r_{\mathrm{br}}>0\) with the following properties.
When all links in the relevant patch have logarithms of
norm less than \(r_{\mathrm{br}}\), the frame and predictor
minimisers are unique, belong to their regular branches,
and are real analytic functions of these link logarithms.
Every fixed number of their derivatives is bounded by a
constant depending only on these fixed data and the
derivative order. These constants are independent of the
torus side.
\end{lemma}
\begin{proof}
At \(U=1\), the frame score is nonnegative, and it is zero
only if \(h_z=h_{z'}\) for every pair of sites. Since
\(h_0=1\), its unique minimiser is \(h=1\).
Put \(m=|\mathcal Q_R|\).
For a frame tangent vector \(\omega\) with \(\omega_0=0\),
its Hessian at this minimiser is
\[
 \frac1{m^2}\sum_{z,z'}|\omega_z-\omega_{z'}|^2
 =\frac2m\sum_z|\omega_z|^2
                   -\frac2{m^2}\left|\sum_z\omega_z\right|^2
 \ge\frac2{m^2}\sum_z|\omega_z|^2 .
 \]
The last inequality uses Cauchy--Schwarz on the \(m-1\)
unpinned sites. It is a strictly positive finite-dimensional
Hessian. The analytic implicit function theorem gives a
unique critical point in a neighbourhood of \(h=1\).

This critical point is the global minimiser for small
enough input links. To see this, choose a smaller
neighbourhood inside that uniqueness neighbourhood.
The score at \(U=1\) has a strictly positive minimum on its
compact complement. Uniform continuity of
\(\mathcal A_U(h)\), uniformly over this compact frame
space, keeps every minimiser inside the neighbourhood
when \(U\) is sufficiently close to \(1\).
All such minimisers are critical there and hence coincide.
The finite branch tests are strict at \(U=1\); shrink the
input neighbourhood to keep them on the regular branch.

For the predictor reduce first to \(Y_z=1\).
The score has the unique minimiser \(u=1\), with Hessian
\(\sum_z w(z)|\omega|^2=|\omega|^2\).
The same implicit-function and compactness argument applies.
The functions are analytic in finite-dimensional
exponential coordinates. On a smaller closed input ball,
their derivatives of each fixed order are bounded.
The patches have a fixed finite number of coordinates;
neither compactness argument involves the number of patches.
\end{proof}

Only a local regular branch is asserted in this lemma.
The compact group may be disconnected: the small chart is
in its identity component, and the compact-complement
argument excludes other components from these minimisers.

% !TeX root = 4DYM_notes.tex
\subsection{Global branch choices and the axial cost condition}
\label{sec:axial-fallback}

The regular tests do not suffice to define the integral at a
large-field configuration. We now specify its other branches.
The axial condition below is part of the definition of the
kernel, rather than an estimate deduced from the sum of its
likelihoods. We give the layer argument and the first-power
comparison in full. This follows
\cite[Unit III.1, definitions on PDF pp.~83--84,
Propositions III.1.4--III.1.6 on p.~85, and
Lemmas 4.7--4.16 on pp.~95--97]{ArmstrongNotes}; the
corresponding concise branch argument is
\cite[Lemma 11.4, PDF pp.~48--49]{ArmstrongDraft}.

For an integer box or a coordinate slice \(X\), define
\begin{equation}
 \begin{split}
 \operatorname{edge}_X(U)&=\sum_{\substack{z,i:\ z,z+e_i\in X}}d_G(U_{z,i}),\\
 \operatorname{face}_X(U)&=\sum_{\substack{p:\ \text{all four vertices of }p\text{ in }X}}
                                                    d_G(P_p(U)).
 \end{split}
 \label{eq:axial-edge-face-costs}
\end{equation}
Both sums use positive orientations; a reversed edge or
plaquette has the same chord. In particular \(\operatorname{face}_X\) is a
sum of chords, not the Wilson action, which is a sum of
half their squares. For \(H:X\to G\) put
\(U^H_{z,i}=H_zU_{z,i}H_{z+e_i}^{-1}\).
Then \(\operatorname{face}_X(U^H)=\operatorname{face}_X(U)\).
Here all sets and paths are on the integer cover.

A frame tag is \(t=(p,\pi)\in\mathcal Q_R\times S_4\).
Its axial transport and its rooted frame are
\[
 a_z^t=U(\gamma^\pi_{p,z}),\qquad
 h_z^t=(a_0^t)^{-1}a_z^t.
\]
Fix once and for all \(a_{\rm ax}\ge16\), before \(L\),
\(J\), the fine cutoff and the period. A tag is
\emph{eligible} if
\begin{equation}
 \operatorname{edge}_{\mathcal Q_R}(U^{a^t})
        \le a_{\rm ax}(2R+1)\operatorname{face}_{\mathcal Q_R}(U).
 \label{eq:axial-eligibility}
\end{equation}
Write \(\mathcal E(U)\) for the finite set of eligible tags.
Root normalisation conjugates every gauged edge by
\((a_0^t)^{-1}\), so
\(\operatorname{edge}_{\mathcal Q_R}(U^{h^t})
       =\operatorname{edge}_{\mathcal Q_R}(U^{a^t})\).

\begin{lemma}[Parallel layers and an eligible axial frame]
\label{lem:axial-layer-existence}
Let \(n=2R+1\). On a coordinate slice with \(k\) varying
axes, in any prescribed order, there is a base point \(p\)
whose axial gauge satisfies
\[
                         \operatorname{edge}_X(U^{a^{p,\pi}})\le4kn \operatorname{face}_X(U).
\]
Consequently \(\mathcal E(U)\ne\varnothing\) for every
configuration, and for each of the \(4!\) axis orders at
least one tag is eligible. The numerical constants are
independent of the configuration and period.
\end{lemma}
\begin{proof}
We first prove the comparison used in the induction.
Suppose \(X\) has one additional varying axis \(j\), all
its internal \(j\)-links equal \(1\), and \(X_z\) is the
layer with \(j\)-coordinate \(z\). For an internal link in
direction \(i\ne j\), the \((i,j)\)-square gives, up to
orientation and conjugation,
\[
 P_{i,j}(y)=U_{y,i}U_{y+e_j,i}^{-1},\qquad
 |d_G(U_{y+e_j,i})-d_G(U_{y,i})|
                                  \le d_G(P_{i,j}(y)).
\]
The second inequality follows twice from the chord triangle
inequality. It also applies when the plaquette orientation
is reversed. Put
\(f_X(y)=\sum_{p\text{ based at }y,\ p\subset X}d_G(P_p(U))\).
Summing the link inequality over a layer, allowing the
harmless factor four for the four directions, yields
\[
 |\operatorname{edge}_{X_{z+1}}(U)-\operatorname{edge}_{X_z}(U)|
                         \le4\sum_{y\in X_z}f_X(y).
\]
Edges absent from both layers contribute zero. Telescoping
from \(z\) to \(u\), in either order, uses each intermediate
layer at most once. Thus
\begin{equation}
               |\operatorname{edge}_{X_z}(U)-\operatorname{edge}_{X_u}(U)|\le4\operatorname{face}_X(U).
 \label{eq:axial-layer-comparison}
\end{equation}

For \(k=0\) there are no internal links. Suppose the
assertion is known for \(k-1\) axes and append axis \(j\)
last. Write the resulting slice as the disjoint union of
its \(n\) layers \(X_z\). Plaquettes internal to different
layers are distinct and all are counted in \(\operatorname{face}_X(U)\).
Choose a layer \(X_{z_0}\) of least plaquette cost. Then
\[
                     n\operatorname{face}_{X_{z_0}}(U)\le \operatorname{face}_X(U).
\]
By induction choose \(p\in X_{z_0}\) with
\[
 \operatorname{edge}_{X_{z_0}}(U^{a^{p,\pi}})
       \le4(k-1)n\operatorname{face}_{X_{z_0}}(U)\le4(k-1)\operatorname{face}_X(U).
\]
Extend this axial gauge along \(j\) last. The transport to
\(y+e_j\) is the transport to \(y\) followed by that last
edge. Hence all internal \(j\)-links become \(1\), and the
gauge on \(X_{z_0}\) is unchanged. Apply
\eqref{eq:axial-layer-comparison} and gauge invariance of
\(\operatorname{face}_X\) to every layer:
\[
 \begin{split}
 \operatorname{edge}_X(U^{a^{p,\pi}})
 &=\sum_z\operatorname{edge}_{X_z}(U^{a^{p,\pi}})\\
 &\le n\{\operatorname{edge}_{X_{z_0}}(U^{a^{p,\pi}})+4\operatorname{face}_X(U)\}\\
 &\le4kn\operatorname{face}_X(U).
 \end{split}
\]
For \(k=4\) the coefficient is \(16n\), proving
\eqref{eq:axial-eligibility}. This is an existence argument
by choosing a least-cost layer at each stage; no
small-field assumption is used. The proof applies also to
a periodic connection lifted to the cover, with every
lifted edge counted as specified in \eqref{eq:axial-edge-face-costs}.
\end{proof}

Here are the full local branch probabilities. On the open
frame test \eqref{eq:actual-frame-branch-test}, use its rooted
minimiser and give every unused tag mass
\((4!|\mathcal Q_R|)^{-1}\). Off that test, use \(h^t\) and
give tag \(t\) mass
\begin{equation}
 p^{\rm fr}(U,t)=\frac{1_{\{t\in\mathcal E(U)\}}}
                             {|\mathcal E(U)|}.
 \label{eq:axial-fallback-mass}
\end{equation}
For a predictor connector \(Y\), on
\eqref{eq:actual-predictor-branch-test} use its minimiser
and give every unused tag \(z\in\mathcal Q_R\) mass
\(|\mathcal Q_R|^{-1}\). Otherwise use \(Q=Y_z\) with mass
\(w(z)\), including the tags whose weight is zero.
Thus the predictor fallback is weighted by the original
convolution law. The weights are not replaced by a uniform
law in this case.

For completeness, the regular minimisers used in this
definition can be made unique with thresholds depending
only on \(G\). We prove that assertion next, together with
the estimate needed for the large-field branches.

\subsection{A first-power estimate for the synchronising frame}
\label{sec:frame-first-power}

Use the normalised real Frobenius norm
\(\|A\|_{\rm F}^2=N^{-1}\operatorname{Re}\operatorname{tr}(AA^*)\).
Then \(d_G(uv^{-1})=\|u-v\|_{\rm F}\).
There is a ball about \(1\) on which the nearest-point
projection \(\pi_G\) onto \(G\) is analytic and
\[
 \pi_G(1)=1,\qquad D\pi_G(1)=\operatorname{pr}_{\mathfrak g}.
\]
Indeed the normal-coordinate map \((u,Z)\mapsto u+Z\),
with \(Z\) orthogonal to the tangent space at \(u\), has
invertible derivative at \((u,0)\). The inverse function
theorem and compactness give a uniform tubular
neighbourhood. If a common neighbourhood failed to give
unique nearest points, a sequence of such points at
distances tending to zero would converge to the same
group point and contradict that local inverse. The
normal-coordinate inverse is analytic, since the matrix
group charts and their tangent spaces are analytic.
Shrinking a ball about \(1\), we can therefore arrange
\begin{equation}
 \begin{split}
 \|\pi_G(A)-\pi_G(B)\|_{\rm F}&\le2\|A-B\|_{\rm F},\\
 \|\pi_G(A)-\pi_G(B)
       -\operatorname{pr}_{\mathfrak g}(A-B)\|_{\rm F}
                                      &\le\tfrac14\|A-B\|_{\rm F}.
 \end{split}
 \label{eq:frame-tubular-estimates}
\end{equation}
To obtain the second inequality, integrate
\(D\pi_G-\operatorname{pr}_{\mathfrak g}\) along the segment
in this ball. This argument includes disconnected \(G\):
the ball is chosen in the chart of its identity component,
and compactness separates the other components.

\begin{lemma}[Synchronisation with a linear mean-error bound]
\label{lem:frame-first-power}
Let a finite nonempty index set carry group-valued route
matrices \(W_{ij,\pi}\) with
\(W_{ii,\pi}=1\) and
\(W_{ji,\bar\pi}=W_{ij,\pi}^{-1}\), where reversing
\(\pi\) permutes the route set. If \(h\) minimises
\(\operatorname{av}_{i,j,\pi}q_G(h_iW_{ij,\pi}h_j^{-1})\),
then for some \(b\in G\),
\begin{equation}
 \operatorname{av}_i d_G(h_ib^{-1})
       \le C_G\operatorname{av}_{i,j,\pi}d_G(W_{ij,\pi}).
 \label{eq:frame-linear-mean-bound}
\end{equation}
The constant depends only on the fixed matrix group,
not on either index cardinality or the routes.
\end{lemma}
\begin{proof}
Put \(\epsilon=\operatorname{av}_{i,j,\pi}d_G(W_{ij,\pi})\),
\(\bar h=\operatorname{av}_ih_i\), and
\(\epsilon_i=\operatorname{av}_{j,\pi}d_G(W_{ij,\pi})\).
Comparison with the constant frame and
\(d_G(u)^2\le2d_G(u)\) give
\[
 \operatorname{av}_{i,j,\pi}
           d_G(h_iW_{ij,\pi}h_j^{-1})^2\le2\epsilon.
\]
The chord triangle inequality then gives
\[
 \operatorname{av}_{i,j}\|h_i-h_j\|_{\rm F}^2
 \le2\operatorname{av}_{i,j,\pi}
       \{d_G(h_iW_{ij,\pi}h_j^{-1})^2+d_G(W_{ij,\pi})^2\}
 \le8\epsilon.
\]
The variance identity implies
\(\operatorname{av}_i\|h_i-\bar h\|_{\rm F}^2\le4\epsilon\).
In particular \(\operatorname{dist}(\bar h,G)\le2\sqrt\epsilon\).
Choose a nearest point \(b\) and left-rotate the frame so
that \(b=1\). This changes neither its score nor the
preceding distances.

Let
\(A_i=\operatorname{av}_{j,\pi}h_jW_{ij,\pi}^{-1}\)
be the feedback at site \(i\). Reversal symmetry pairs the
two terms containing \(i\). Coordinate minimality implies
that \(h_i\) is nearest to the feedback with \(j\ne i\).
Including \(j=i\) adds a positive multiple of \(h_i\):
when a competitor \(v\) replaces \(h_i\), the added
difference is a nonnegative multiple of
\(\|v-h_i\|_{\rm F}^2\). Hence \(h_i\) is also nearest
to \(A_i\). Moreover
\[
                         \|A_i-\bar h\|_{\rm F}\le\epsilon_i.
\]
Choose \(\eta>0\), depending only on \(G\), so that the
ball of radius \(\eta\) lies in
\eqref{eq:frame-tubular-estimates}. If
\(\epsilon<\eta^2/64\), then
\(\|\bar h-1\|_{\rm F}\le\eta/4\).
When \(\epsilon_i<\eta/4\), both feedbacks are in that
ball, and their nearest points satisfy
\[
                    d_G(h_i)\le2\|A_i-\bar h\|_{\rm F}
                                      \le2\epsilon_i.
\]
For the other sites use
\(d_G(h_i)\le2\le8\epsilon_i/\eta\).
Averaging proves the claimed linear bound in this case.
If \(\epsilon\ge\eta^2/64\), use instead
\(\operatorname{av}_id_G(h_i)\le2\le128\epsilon/\eta^2\).
Thus one may take
\(C_G=\max(2,8/\eta,128/\eta^2)\).
The final right rotation by \(b^{-1}\) has the same
chord as the left rotation used in the proof.
\end{proof}

\begin{lemma}[Global definition and normalisation of the branches]
\label{lem:global-branch-definition}
The frame and predictor thresholds can be chosen using
\(G\) alone so that the preceding rules define Borel
branch maps and probabilities at every connection.
Each local probability sums to one. The joint law is
\begin{equation}
 \lambda_t(U)=\prod_xp^{\rm fr}(U^{(x)},t_x)
      \prod_{e=(x,i)}p^{\rm pr}
                         (Y_e(U,t_x,t_{x+e_i}),z_e),
 \qquad \sum_t\lambda_t(U)=1.
 \label{eq:complete-branch-law}
\end{equation}
One frame variable at a site is shared by all incident
predictors. The kernel \eqref{eq:physical-kernel} is
therefore the same normalised integral with these
explicit choices at every field, including cutoff boundaries.
\end{lemma}
\begin{proof}
We first verify uniform uniqueness. On the frame
test, make one origin-based axial change of gauge, so all
route matrices have norm distance at most \(d=c_{\rm fr}\)
from \(1\). A rooted minimiser exists by compactness.
It is also a minimiser among all frames, since any frame
can be root-normalised by a common left multiple without
changing its score. Comparison
with a constant frame now gives
\(\operatorname{av}_{i,j}\|h_i-h_j\|_{\rm F}^2\le4d^2\).
Rotate it by a nearest group point to its centroid.
Then \(\|\bar h-1\|_{\rm F}\le2d\) and
\(\|A_i-1\|_{\rm F}\le3d\), uniformly in the number of sites.

Centre two minimisers \(h,k\) this way, and put
\(u=\max_i\|h_i-k_i\|_{\rm F}\).
Both centroids have nearest point \(1\), so their
difference is orthogonal to \(\mathfrak g\). Their
feedbacks satisfy
\[
 \|A_i(h)-A_i(k)\|_{\rm F}\le u,\qquad
 \|A_i(h)-A_i(k)-(\bar h-\bar k)\|_{\rm F}\le du.
\]
Choose \(d\le1/4\) so that the ball of radius \(3d\)
is contained in \eqref{eq:frame-tubular-estimates}.
Applying its second inequality to a site attaining \(u\)
gives \(u\le(d+1/4)u\le u/2\). The centred minimisers
coincide. Thus all minimisers differ by a common left
multiple, and their root normalisation makes them unique.

For the predictor write \(A=\sum_z w(z)Y_z\). Expanding
the square shows that its score is
\[
 \mathcal A_Y(v)=\tfrac12\sum_z w(z)\|Y_z-A\|_{\rm F}^2
                                  +\tfrac12\|A-v\|_{\rm F}^2.
\]
Its minimisers are nearest group points to \(A\).
If its minimum is below \(c_{\rm pr}\), then
\(\operatorname{dist}(A,G)<\sqrt{2c_{\rm pr}}\).
Choose this smaller than the uniform tubular radius.
This proves uniqueness independently of the weights and
their number.

The regular tests are open and their minimisers are
continuous on those open sets: a convergent sequence of
data and compactness give convergent subsequences of
minimisers, and uniqueness identifies their limits.
On the complementary sets the finitely many axial maps
and connector maps are continuous. Each eligibility test
is closed, since the two finite costs are continuous.
Their finite cardinality is a Borel integer-valued
function bounded below by one, by
Lemma~\ref{lem:axial-layer-existence}. Thus every branch
map and every stated mass is Borel. Equality in either
open regular test belongs to its fallback branch.

The frame masses sum to one on either branch. The
predictor masses sum to one because \(\sum_z w(z)=1\).
Sum predictor tags first in \eqref{eq:complete-branch-law},
with the shared frame tags fixed, and then sum those frame
tags. This proves its normalisation without declaring
the weighted frame and predictor variables independent.
Gauge changes conjugate all costs and send each rooted
frame to the corresponding rooted frame; uniqueness on
the regular branch and the explicit axial formula give
the same covariance there and on fallback. Thus these
choices retain the original gauge-covariant kernel.
\end{proof}

\subsection{Comparing every realised frame with one gauge}
\label{sec:all-branch-frame-comparison}

\begin{lemma}[A common first-power comparison for both frame branches]
\label{lem:all-branch-frame-comparison}
Let \(H\) be any gauge on \(\mathcal Q_R\), and let
\(h^t\) be the realised frame of a tag with positive mass.
There is \(c\in G\) such that, with \(n=2R+1\),
\begin{equation}
 \sum_{z\in\mathcal Q_R}d_G(h_z^tH_z^{-1}c^{-1})
 \le C_G n\{\operatorname{edge}_{\mathcal Q_R}(U^H)
                         +a_{\rm ax}n\operatorname{face}_{\mathcal Q_R}(U)\}.
 \label{eq:all-branch-frame-comparison}
\end{equation}
The constant may be enlarged from
Lemma~\ref{lem:frame-first-power}; it remains independent
of \(R\), \(L\), cutoff and period.
\end{lemma}
\begin{proof}
First record the elementary route count. For fixed order
\(\pi\), a fixed undirected edge can occur in a coordinate
route from \(z\) to \(z'\) for at most
\(n^5=|\mathcal Q_R|n\) endpoint pairs. The three
transverse coordinates at the edge fix one of the two
endpoint coordinates in each direction. The three other
endpoint coordinates have at most \(n^3\) choices;
the two longitudinal endpoints have at most \(n^2\)
choices. The estimate includes travel in either direction.
The chord triangle inequality consequently gives the
convenient, weaker bound
\begin{equation}
 \sum_{z,z'}\operatorname{av}_\pi
        d_G(U^H(\gamma^\pi_{z,z'}))
              \le4|\mathcal Q_R|n\operatorname{edge}_{\mathcal Q_R}(U^H).
 \label{eq:axial-route-incidence}
\end{equation}

On the minimising branch the frame
\(k_z=h_z^tH_z^{-1}\) minimises the synchronisation score
for \(U^H\), modulo a common left multiple. Apply
Lemma~\ref{lem:frame-first-power} to its route matrices
and multiply their mean by \(|\mathcal Q_R|\).
Equation~\eqref{eq:axial-route-incidence} bounds the result
by \(4C_G n\operatorname{edge}_{\mathcal Q_R}(U^H)\), as required.

On fallback put again \(k_z=h_z^tH_z^{-1}\).
For an edge \((z,i)\), the identity
\[
 U^{h^t}_{z,i}=k_zU^H_{z,i}k_{z+e_i}^{-1}
\]
and conjugation invariance imply
\[
 d_G(k_zk_{z+e_i}^{-1})
                  \le d_G(U^{h^t}_{z,i})+d_G(U^H_{z,i}).
\]
Apply the route count to the pure-gauge connection
\(k_zk_{z+e_i}^{-1}\). Its transport from \(z\) to \(z'\)
is \(k_zk_{z'}^{-1}\). Averaging the possible choices of
\(z'\), choose one of them and set \(c=k_{z'}\), to get
\[
 \sum_zd_G(k_zc^{-1})
 \le4n\{\operatorname{edge}_{\mathcal Q_R}(U^{h^t})+\operatorname{edge}_{\mathcal Q_R}(U^H)\}.
\]
Positive fallback mass implies precisely
\eqref{eq:axial-eligibility}, which bounds its first cost
by \(a_{\rm ax}n\operatorname{face}_{\mathcal Q_R}(U)\).
This proves \eqref{eq:all-branch-frame-comparison} for
every retained tag, not just an average over the tags.
\end{proof}

\begin{proposition}[Macro-box circulation for every positive branch]
\label{prop:all-branch-macro-circulation}
Fix \(G,J,a_{\rm ax},c_{\rm fr},c_{\rm pr}\), with the
thresholds just chosen. There is \(C\), chosen before
\(L\), such that for each coarse plaquette \(x\),
\begin{equation}
 d_G(P(Q_t(U);x))\le CL^{-2}\operatorname{face}_{\mathcal M_x}(U)
 \label{eq:proved-macro-circulation}
\end{equation}
for every \(t\) of positive likelihood.
Here \(\mathcal M_x\) is the integer box about \(Lx\)
of radius \(3(J+3)L\). It contains the four frame
patches and all connector paths for that plaquette.
The estimate is uniform in cutoff, period, the number
of coarse sites and the branch choices. Each fine
plaquette belongs to at most \(C_J\) such macro boxes.
\end{proposition}
\begin{proof}
Choose one axial gauge \(H\) on \(\mathcal M_x\) by
Lemma~\ref{lem:axial-layer-existence}. Since its side
is at most \(C_J L\),
\[
                     \operatorname{edge}_{\mathcal M_x}(U^H)
                              \le C_J L \operatorname{face}_{\mathcal M_x}(U).
\]
For each of the four corner roots \(v\), use
Lemma~\ref{lem:all-branch-frame-comparison} on its patch
to choose \(c_v\). Both its edge and face costs are
bounded by the macro-box costs. Therefore
\begin{equation}
 \sum_zd_G(c_v^{-1}h_z^{(v)}H_{Lv+z}^{-1})
                              \le C_{G,J,a_{\rm ax}}L^2
                                                   \operatorname{face}_{\mathcal M_x}(U).
 \label{eq:macro-frame-price}
\end{equation}
The left and right placements of \(c_v^{-1}\) have the
same chord by conjugation invariance.

The uniform law on one centred block has supremum
\(L^{-4}\), and convolution with a probability law
cannot increase the supremum. Thus \(w(z)\le L^{-4}\).
For a connector edge \(e=(v,i)\), put
\[
 T_z=H_{Lv+z}U(\gamma_{Lv+z,L(v+e_i)+z})
                                          H_{L(v+e_i)+z}^{-1}.
\]
A fixed fine edge can occur for at most \(L\) pairs
\((z,k)\), where \(k\) is its position in that straight
connector. Hence
\[
 \sum_z w(z)d_G(T_z)
 \le L^{-3}\operatorname{edge}_{\mathcal M_x}(U^H)
 \le C_JL^{-2}\operatorname{face}_{\mathcal M_x}(U).
\]
Use \eqref{eq:macro-frame-price} for the two ends, with
the weight bound \(L^{-4}\), and the chord triangle
inequality to conclude that
\begin{equation}
 \epsilon_e:=\sum_z w(z)
                 d_G(c_v^{-1}Y_{e,z}c_{v+e_i})
                         \le CL^{-2}\operatorname{face}_{\mathcal M_x}(U).
 \label{eq:macro-connector-price}
\end{equation}

We check both predictor branches. On the minimising
branch, transform by the two constants \(c_v,c_{v+e_i}\)
and let \(A\) be the transformed matrix mean. Its
nearest group point \(Q\) obeys
\[
 d_G(Q)\le\|Q-A\|_{\rm F}+\|A-1\|_{\rm F}
                \le2\|A-1\|_{\rm F}\le2\epsilon_e.
\]
On fallback its minimum score is at least
\(c_{\rm pr}\). Testing that score at \(1\) gives
\[
 c_{\rm pr}\le\tfrac12\sum_z w(z)
                d_G(c_v^{-1}Y_{e,z}c_{v+e_i})^2\le\epsilon_e.
\]
Every selected connector has chord at most two, so its
predicted chord is at most \(2\epsilon_e/c_{\rm pr}\).
Thus in either case
\[
 d_G(c_v^{-1}Q_t(U)_e c_{v+e_i})
                         \le\max(2,2/c_{\rm pr})\epsilon_e.
\]
Multiply the four transformed bonds around the coarse
plaquette. Their corner constants cancel, leaving a
conjugate of \(P(Q_t(U);x)\). Summing the four chord
bounds proves \eqref{eq:proved-macro-circulation}.
The macro-box centres are spaced by \(L\) in each fine
direction, so a fixed fine plaquette is in at most
\((6(J+3)+5)^4\) boxes, up to the fixed orientation
multiplicity. This number is chosen before \(L\).

All arguments are local on the cover. On the admissible
periods the required enlarged boxes have no identifications.
Their translates use exactly the same constants.
The scale cutoff and the total number of blocks do not
occur in any of the bounds. In particular no factor
depending on the total volume is absorbed into \(C\).
\end{proof}

\begin{corollary}[Adding the actual coarse noise]
\label{cor:macro-circulation-with-noise}
For \(q_e=d_G(V_eQ_t(U)_e^{-1})\),
\[
 d_G(P(V;x))\le CL^{-2}\operatorname{face}_{\mathcal M_x}(U)
                                  +\sum_{e\in\partial x}q_e.
\]
\end{corollary}
\begin{proof}
Replace the four predicted bonds successively by the
four actual coarse bonds. Each replacement changes the
plaquette chord by at most its \(q_e\), independently
of orientation. Apply the proposition.
\end{proof}

\subsection{From small plaquettes to tree-gauge coordinates}

The following elementary path argument is the geometric
input. It plays the role of the local gauge estimates in
Ba{\l}aban's analysis \cite[Sect.~B, eq.~(1.19), printed p.~79;
Lemma~1 and eqs.~(1.24)--(1.26), pp.~79--80]{Bal85g}.
This is a finite-patch adaptation of the axial-gauge path argument;
we prove the version used here.

\begin{lemma}[Finite-box path comparison]
\label{lem:box-path-comparison}
Let two lattice paths in a fixed integer box have the same
endpoints. Suppose their concatenation is reduced by a
specified sequence of backtrack cancellations and elementary
square moves, using \(A\) squares. Then
\[
 d_G\bigl(U(\gamma_1)U(\gamma_2)^{-1}\bigr)
       \le\sum_{\text{square moves }p}d_G(P_p(U)).
 \]
In particular it is at most \(A\delta\) if every plaquette
used has chord at most \(\delta\).
All the finitely many path comparisons in a fixed blocking
patch admit a common such bound.
\end{lemma}
\begin{proof}
The chord is invariant under inverses and conjugations and
satisfies \(d_G(AB)\le d_G(A)+d_G(B)\).
Replacing the order of two adjacent perpendicular steps
changes the holonomy by insertion of a conjugate plaquette
holonomy. A backtrack changes nothing.
Iterating the chord inequality proves the assertion.

For existence of the required square moves, use the
coordinate-ordered path from a fixed box root to each site.
The loop formed by such a path, one extra link, and the
reverse coordinate path can be reduced by commuting that
link past the later coordinate segments, followed by
backtrack cancellations. All the squares remain in the box.
Apply this to every link of a path.
For a path of length \(l\) in a box of radius \(R_*\), this
procedure uses at most \(C(1+R_*)l\) squares.
The tree paths inside an \(L\)-block have length at most
\(L^4-1\). The set of their possible shapes and of the
coordinate routes in a fixed patch is finite.
Taking the largest of these filling bounds gives a
constant depending on \(L,J\), not on the physical volume.
\end{proof}

Choose a fixed enlarged patch \(\mathcal P_x\) around every
coarse root. It contains the frame and connector paths just
used, the adjacent blocks, and their tree paths.
Its radius and its number \(m_*\) of plaquettes depend
only on \(L,J\). For a configuration \(U\) set
\[
 \epsilon_x(U)=\sum_{p\in\mathcal P_x}d_G(P_p(U)).
 \]
The precise enlargement can be chosen once by taking the
union of those finitely many paths and square fillings.

\begin{proposition}[The nonlinear chart-entry bound]
\label{prop:nonlinear-chart-entry}
There are \(C_{\mathrm g}<\infty\) and \(\delta_{\mathrm g}>0\),
depending only on \(G,L,J\) and the fixed kernel parameters,
such that the following holds on every admissible period.
Suppose the fine plaquette chords are at most
\(\delta_{\mathrm g}\), and \(U\) is in the prescribed
block tree gauge. Let \(Q(U)\) be the actual predicted
coarse connection. It is independent of the finite branch
label in this region. For a free fine link \(e\), with
parent block \(x\),
\begin{equation}
 d_G\bigl(U_eR_e(V)^{-1}\bigr)
 \le C_{\mathrm g}\epsilon_x(U)
 +\begin{cases}
 0,&e\text{ is internal},\\
 d_G\bigl(Q(U)_{(x,i)}V_{(x,i)}^{-1}\bigr),
                         &e\text{ crosses in direction }i .
 \end{cases}
 \label{eq:nonlinear-chart-entry}
\end{equation}
Here \(R(V)\) is the replica connection of
Section~\ref{sec:physical-fluctuation}.
\end{proposition}
\begin{proof}
First consider one frame patch. Gauge transform by the
coordinate-ordered transports from its root. The
transformed links are holonomies of the loops in
Lemma~\ref{lem:box-path-comparison}; hence each is within
\(C_{L,J}\epsilon_x(U)\) of \(1\).
Choose \(\delta_{\mathrm g}\) so that this lies inside
the input neighbourhood of
Lemma~\ref{lem:fixed-patch-branch}.

We also need a quantitative comparison with these axial
transports. Denote the axial frame by \(h^0\).
Every term in its score is at most
\(C_{L,J}\epsilon_x(U)^2\), by the path lemma and
\(q_G=d_G^2/2\). The score of the minimising frame is no
larger. Each nonnegative term in its finite average is
therefore bounded by the number of terms times this
quantity. Taking square roots, and using terms with one
endpoint equal to the root, gives
\[
 d_G(h_zh_z^{0\, -1})
              \le C_{L,J}\epsilon_x(U).
 \]
The same argument applies to the adjacent root's frame.
The uniqueness statement is unaffected by the axial gauge
change: the frame score transforms by
\(h_z\mapsto g_0h_zg_z^{-1}\).

Consequently every connector \(Y_z\) differs by at most
\(C_{L,J}\epsilon_x(U)\) from the holonomy \(W\) of any
of the specified paths joining the two coarse roots.
For the weighted predictor score, comparison with \(u=W\)
gives \(\mathcal A_Y(u_{\min})\le
C_{L,J}\epsilon_x(U)^2\).
Choose a point of positive weight; its weight has a
positive lower bound in this fixed finite averaging rule.
Its term in the score bounds
\(d_G(Y_zu_{\min}^{-1})\), and then the chord triangle
inequality bounds
\[
 d_G(u_{\min}W^{-1})
              \le C_{L,J}\epsilon_x(U).
 \]
The connectors are close enough to one another that their
actual branch test uses this unique minimiser.
The frame branch test also holds: every axial-frame loop
in that finite test is small by the path lemma.
Thus all finite labels select these same regular maps.

For an internal free link, join its endpoints to the
block root along the block tree. The resulting loop has
holonomy \(U_e\), because the tree links are \(1\).
The path lemma bounds its chord by
\(C_{L,J}\epsilon_x(U)\).
For a crossing link use the path from its first block
root to the link, across the link, and back along the
second block tree to the second root.
Its holonomy in tree gauge is \(U_e\).
The preceding predictor comparison bounds
\(d_G(U_eQ(U)_{(x,i)}^{-1})\).
The triangle inequality with \(V_{(x,i)}\) proves
\eqref{eq:nonlinear-chart-entry}.
All constants came from finite paths, finite scores, and
the fixed group chart.
\end{proof}

\begin{corollary}[A common small chart for a low-noise configuration]
\label{cor:actual-chart-entry}
For every prescribed small radius \(r_*>0\), there exists
\(\delta_*>0\), with the same parameter dependence as above,
such that the following is true.
If all fine plaquette chords and all blocking mismatches
\[
 d_G\bigl(Q(U)_eV_e^{-1}\bigr)
 \]
are at most \(\delta_*\), then the tree-gauge connection
can be written uniquely as \(U=U(V,y)\), with
\(\|y\|_\infty\le r_*\).
The norms here are maximum norms of physical Lie algebra
coordinates, not Euclidean norms over the entire torus.
\end{corollary}
\begin{proof}
There are at most \(m_*\) plaquettes in each
\(\mathcal P_x\).
The right side of \eqref{eq:nonlinear-chart-entry} is at
most \((C_{\mathrm g}m_*+1)\delta_*\).
On a sufficiently small group neighbourhood the logarithm
is unique and has norm at most \(C_G\) times the chord.
Choose \(\delta_*\le\delta_{\mathrm g}\) so that these
chords lie in that neighbourhood and
\(C_G(C_{\mathrm g}m_*+1)\delta_*\le r_*\).
Use \(y_e=\log(U_eR_e(V)^{-1})\) on free links.
\end{proof}

\subsection{The actual quadratic form and a volume-independent gap}

Write \(V(x)_e=e^{x_e}\), and let \(U(x,y)\) be the
tree-gauge connection with free links
\(e^{y_e}R_e(V(x))\).
On the common regular branch define
\begin{equation}
 \varphi(x,y)=S_f^{\mathrm{un}}(U(x,y))
       +a\sum_e q_G\bigl(V(x)_eQ(U(x,y))_e^{-1}\bigr).
 \label{eq:actual-unmasked-stationary}
\end{equation}
This is the function used in the central fluctuation
formula. The \(\beta\)-independent loop action and the Haar
Jacobian belong to \(\psi\), as in
\eqref{eq:smooth-action-with-haar}.

For a fine link field \(A\) zero on the block trees, define
the linearised average \(Q_0A\) by
\[
 Q(e^{tA})_e=\exp\{t(Q_0A)_e+O(t^2)\}.
 \]
It is a finite-range linear operator, since the averaging
map reads a fixed patch. Let \(\dd A\) denote the
linearised plaquette curl:
\[
 (\dd A)_{ij}(z)
 =A_j(z+e_i)-A_j(z)-A_i(z+e_j)+A_i(z).
 \]

\begin{lemma}[Block Poincaré estimate]
\label{lem:actual-block-poincare}
There is \(C_{\mathrm P}<\infty\), depending only on
\(G,L,J\) and the kernel parameters, such that
\begin{equation}
 \sum_{e\in I}|A_e|^2
 \le C_{\mathrm P}
 \left(\sum_p|(\dd A)_p|^2+\sum_b|(Q_0A)_b|^2\right).
 \label{eq:actual-block-poincare}
\end{equation}
The sums use counting measure. The constant is independent
of the number of blocks.
\end{lemma}
\begin{proof}
% !TeX root = 4DYM_notes.tex
We give the linear reconstruction explicitly. It also specifies
which rows are needed when holes are introduced later. All paths in
this paragraph lie in one enlarged, nonwrapping integer box about
the link being estimated. Its radius is fixed by \(L,J\), the
block trees and the predictor read, before either cutoff is varied.
Oriented sums of a Lie-algebra-valued link array use a minus sign
when a link is traversed backwards.

Choose a box root \(r\) and the coordinate-ordered path
\(\gamma_z:r\to z\). Define a site array and a link array by
\[
 \phi(z)=\sum_{e\in\gamma_z}A_e,\qquad
 (d_0\phi)_e=\phi(e^+)-\phi(e^-),\qquad
 R_e=A_e-(d_0\phi)_e.
\]
The closed path \(\gamma_{e^-}\,e\,\gamma_{e^+}^{-1}\)
can be reduced by the square moves of
Lemma~\ref{lem:box-path-comparison}. In the additive linear
calculation a square contributes exactly its curl, with its
orientation sign. Hence there are integers \(c_{ep}\), on
the fixed patch, such that
\begin{equation}
 R_e=\sum_p c_{ep}(\dd A)_p,
       \qquad \sup_e\sum_p|c_{ep}|\le N_*<\infty .
 \label{eq:coercivity-path-homotopy}
\end{equation}
For instance, commuting a step in direction \(i\) past a
segment of length \(m\) in direction \(j\) gives the sum
of the \(m\) intervening \(ij\)-squares. Iterating over
four coordinates gives the displayed identity; backtracks
contribute zero. Thus this is an identity of link arrays,
not a compactness assertion about a finite matrix.

Let \(r_x\) be the root of block \(x\), and put
\(c_x=\phi(r_x)\),
\(\psi(z)=\phi(z)-c_x\) for \(z\) in that block.
On its tree links \(A_e=0\), so
\(d_0\phi=-R\). Summing along the tree path from
\(r_x\) to \(z\), of length at most \(L^4-1\), gives
\[
 \psi(z)=-\sum_{e\in\tau_{x,z}}R_e,
       \qquad |\psi(z)|\le (L^4-1)N_*
                         \sum_{p\text{ in the patch}}|(\dd A)_p|.
\]
Extend \(c\) to fine sites by the block-constant value
\(c_x\), and set \(\widetilde R=R+d_0\psi\). Then
\begin{equation}
 A=d_0c+\widetilde R,\qquad
 |\widetilde R_e|\le C_R
                      \sum_{p\text{ in the patch}}|(\dd A)_p|,
 \qquad C_R=\{1+2(L^4-1)\}N_* .
 \label{eq:coercivity-block-root-split}
\end{equation}
All expressions needed on the predictor's input links are
included by enlarging the same patch once.

We now identify the remaining block-root mode. Predictor
covariance is rooted at the two coarse endpoints. Apply it
to the fine pure gauge with site values \(\exp(-tc_x)\),
constant on each block, and differentiate at \(t=0\).
Both the frame and connector transformations are included.
The resulting exact linear identity is
\[
              Q_0d_0c=d_{0,\rm c}c,
       \qquad (d_{0,\rm c}c)_{(x,i)}=c_{x+e_i}-c_x.
\]
In particular the noise read measures precisely the
difference of adjacent block-root values. With
\(Q_*\) the maximum absolute row sum of this fixed local
linear predictor, \eqref{eq:coercivity-block-root-split}
implies, for a crossing link with coarse bond \(b\),
\begin{align}
 A_e&=(Q_0A)_b+\widetilde R_e-(Q_0\widetilde R)_b,\notag\\
 |A_e|&\le |(Q_0A)_b|+(1+Q_*)C_R
                         \sum_{p\in\mathcal P_x}|(\dd A)_p|.
 \label{eq:coercivity-noise-root-identity}
\end{align}
For an internal link \(d_0c=0\), so the same bound holds
without its noise term. Enlarge the earlier
\(C_{\mathrm g}\), if necessary, to dominate
\((1+Q_*)C_R\). This gives the link estimate also obtained
by linearising \eqref{eq:nonlinear-chart-entry}.

As a useful limiting case, zero curl gives \(R=0\);
the tree condition then makes \(\phi\) constant in each
block. Zero noise makes the adjacent constants equal.
Thus every tested free link is zero. The quantitative
estimate above retains the constants that this kernel
argument alone would leave unspecified. They may depend
on the fixed ratio and patch, but not on the number of
blocks, the terminal cutoff or the physical side length.

By Cauchy--Schwarz the square of its first sum is at most
\(m_*\sum_{p\in\mathcal P_x}|(\dd A)_p|^2\).
There are at most \(4L^4\) fine links per block.
Each fine plaquette is in at most \({\mathsf M}_*\) enlarged
patches, where \({\mathsf M}_*\) depends only on their fixed radius.
A coarse link has at most \(L^3\) crossing fine links.
The inequality \((u+v)^2\le2u^2+2v^2\), summed over free
links, therefore proves the result, for example with
\[
 C_{\mathrm P}
 =1+8L^4C_{\mathrm g}^2m_*{\mathsf M}_*+2L^3 .
 \]
Translation across a period does not change these overlap
counts.
\end{proof}

\begin{proposition}[The flat precision matrix]
\label{prop:actual-flat-gap}
Let \(\lambda>0\) be the direct Wilson coefficient in the
principal action and \(a>0\) the noise coefficient.
For
\[
 H_0=D_y^2\varphi(0,0),\qquad
 \eta_0=\min(\lambda,a)/C_{\mathrm P},
 \]
one has
\begin{equation}
 H_0\ge\eta_0\,1
       \quad\hbox{in the Euclidean counting norm on free coordinates}.
 \label{eq:actual-flat-gap}
\end{equation}
Also \(\varphi(0,0)=0\) and \(D\varphi(0,0)=0\).
\end{proposition}
\begin{proof}
The Wilson plaquette has expansion
\[
 q_G(P_p(e^{tA}))
             =\tfrac12t^2|(\dd A)_p|^2+O(t^3).
 \]
Every transported curvature slot has zero value at \(1\)
and first variation equal to its plaquette curl.
The row-square contribution to the second variation
is consequently nonnegative. Cubic and quartic curvature
polynomials have zero second variation.
The signed average of the direct plaquette action is
the direct plaquette action itself: orientation reversal
uses \(q_G(W^{-1})=q_G(W)\), and translations and axis
permutations reindex the plaquettes.
The noise second variation at fixed coarse field \(x=0\)
is \(a\sum_b|(Q_0A)_b|^2\).
It follows that
\[
 \langle A,H_0A\rangle
 \ge\lambda\sum_p|(\dd A)_p|^2
                         +a\sum_b|(Q_0A)_b|^2 .
 \]
Use Lemma~\ref{lem:actual-block-poincare}.
The value and first derivative vanish because every
curvature and noise mismatch vanishes at the identity.
\end{proof}

\subsection{Local derivative bounds for the actual action}

For discussing locality, attach a coarse site to each
coordinate. A coarse coordinate is attached to its link's
initial site; a free fine coordinate is attached to its
parent block. Choose an orthonormal basis of
\(\mathfrak g\), so coordinates are real numbers.
The number of coordinates attached to any coarse site
is at most
\[
 m_0=4(L^4+1)\dim_\R\mathfrak g.
 \]
Let \(d(i,j)\) be the periodic maximum distance between
their attached coarse sites.

For a matrix \(A=(A_{ij})\), define its weighted row and
column norm by
\begin{equation}
 \|A\|_\theta=
 \max\left\{\sup_i\sum_j e^{\theta d(i,j)}|A_{ij}|,
             \sup_j\sum_i e^{\theta d(i,j)}|A_{ij}|\right\}.
 \label{eq:weighted-schur-norm}
\end{equation}
For a rooted tensor \(T_{i;j_1,\ldots,j_k}\), define
\begin{equation}
 |T|_{\theta,k}
 =\sup_i\sum_{j_1,\ldots,j_k}
 e^{\theta\,\mathsf t(i,j_1,\ldots,j_k)}
                         |T_{i;j_1,\ldots,j_k}|.
 \label{eq:rooted-tensor-norm}
\end{equation}
Here \(\mathsf t\) is the minimum total edge length of a
tree joining the indicated attached sites, with edge
length given by the periodic maximum metric.
Additional lattice sites are allowed as vertices of the
tree. This convention ensures that joining several such
trees never increases the sum of their lengths.
Repeated sites add no length.
These are norms on coefficient arrays, not scalar bounds
on the total extensive action.

\begin{lemma}[Uniform local bounds from the path budget]
\label{lem:actual-action-local-bounds}
Fix the group, block and branch parameters, positive path
width \(\sigma\), path budget \(C_{\mathrm p}\), and
\(\lambda\).
There are \(r_0,\theta_0,K>0\), independent of the period
and ultraviolet cutoff, such that \(\varphi\) is real
analytic on \(\max(\|x\|_\infty,\|y\|_\infty)<2r_0\) and
\begin{align}
 \|D_y^2\varphi(x,y)\|_{\theta_0}&\le K,\label{eq:hessian-local-bound}\\
 \|D_y^2\varphi(x,y)-H_0\|_{\theta_0}
       &\le K\max(\|x\|_\infty,\|y\|_\infty),\label{eq:hessian-local-perturbation}\\
 \|D_y\varphi(x,0)\|_\infty&\le K\|x\|_\infty .
 \label{eq:actual-gradient-forcing}
\end{align}
The mixed \(x,y\) derivatives of the full gradient
\(D\varphi\) of any fixed
order have a uniform bound in
\eqref{eq:rooted-tensor-norm}, after decreasing
\(\theta_0\) if needed.
There is also a constant \(K_{\mathrm a}\) and a positive
radius, with the same uniformity, for which the absolute
Taylor coefficients of \(D\varphi\) have a convergent
rooted-tree weighted majorant.
\end{lemma}
\begin{proof}
First consider a single path polynomial of total length
\(l\) and coefficient magnitude \(b\).
It reads at most \(C_G(1+l)\) link-colour coordinates,
allowing for their coarse replica coordinates.
Every derivative is obtained by inserting direction
matrices into its words and differentiating their
exponentials. Through any fixed order its absolute
coefficient sums are bounded by
\[
 C_{G,k}b(1+l)^{m_k}e^{C_G r_0l}
 \]
on the indicated coordinate ball.
All the attached sites read by the term can be connected
by the projections of its based paths, with total length
at most \(C_L(1+l)\).
Thus its contribution to a rooted tensor weighted by
\(e^{\theta_0\mathsf t}\) costs at most another
\(e^{C_L\theta_0(1+l)}\).
For a fixed output coordinate, only translates of the
term whose read footprint contains that coordinate
contribute. Their number is bounded by a polynomial in
\(1+l\).

Choose \(r_0\) and \(\theta_0\) so that
\(C_Gr_0+C_L\theta_0<\sigma/2\).
The remaining polynomial factors are absorbed by
\(e^{\sigma l/2}\).
The exponential path budgets then sum these bounds,
uniformly per output coordinate.
A row-square term is a double path sum; apply the product
of its two row budgets. The direct Wilson term has a
fixed four-link support.
These observations give the local derivative bounds for
the principal action.

The predictor terms are analytic functions of a fixed
finite patch by Lemma~\ref{lem:fixed-patch-branch}.
Their derivatives have fixed finite support and fixed
cost. Each coordinate belongs to only a fixed number
of translated predictor patches. They therefore obey
the same rooted bounds after increasing the constants.

We justify also the analytic and Taylor assertions,
rather than inferring analyticity from a finite number
of derivative estimates.
The matrix exponentials and inverses in a word extend
holomorphically to complexified Lie coordinates.
For the price use the extension
\(1-\operatorname{tr}(W+W^{-1})/(2r)\).
Real curvature polynomials extend by complex linearity.
The local regular predictor extends holomorphically
on a smaller fixed patch polydisc.
Choose its radius and the word radius so small that the
complex bound of a term is \(C b e^{\sigma l/4}\).
For a term reading \(D_l\le C_G(1+l)\) coordinates,
Cauchy's estimate bounds its absolute power series on a
smaller polydisc of radius \(s\) by the same bound times
\[
 (1-s/R_0)^{-D_l}
                \le e^{2D_ls/R_0}\quad(s\le R_0/2),
 \]
where \(R_0\) is the chosen complex polydisc radius.
The extra exponential, the tree weight, and the
polynomial translation and differentiation costs are
all absorbed by the remaining path-width reserve.
This proves locally uniform convergence of the
holomorphic series and of its rooted, absolute Taylor
majorant. It also justifies all termwise derivatives.

Apply the third-derivative rooted bound along the segment
from \(0\) to \((x,y)\). Summing the remaining input
index against its maximum coordinate norm proves
\eqref{eq:hessian-local-perturbation}.
The same argument with the mixed second derivatives,
and \(D_y\varphi(0,0)=0\), proves
\eqref{eq:actual-gradient-forcing}.
Row and column estimates agree by symmetry of the full
Hessian. This proves the first two matrix bounds as well.
\end{proof}

The constants in this lemma do not include the total
value of \(\varphi\). Its scalar value can grow as the
number of sites. What is uniformly bounded is each
rooted derivative sum.

\subsection{From the counting gap to a maximum-norm inverse}

The next argument is the finite-lattice Combes--Thomas
method. Its useful feature here is that it gives a
maximum-norm inverse without a factor depending on the
dimension of the full torus. We use the organisation of
\cite[Sect.~2.4, Lemma~2.12, preprint p.~14;
Appendix~A, pp.~36--37]{DST23}, and give the finite-matrix proof needed
for the present Hessian.

\begin{lemma}[Weighted inverse estimate]
\label{lem:actual-combes-thomas}
Suppose \(H\) is real symmetric,
\(H\ge\eta1\) in the counting Euclidean norm,
\(\|H\|_\theta\le K\), and there are at most \(m_0\)
indices per attached site.
There are \(\mu>0\) and \(C_{\mathrm{inv}}<\infty\),
depending only on \(\eta,\theta,K,m_0\), such that
\begin{equation}
 |(H^{-1})_{ij}|\le\frac2\eta e^{-\mu d(i,j)},\qquad
 \|H^{-1}\|_{\mu/2}\le C_{\mathrm{inv}} .
 \label{eq:actual-inverse-locality}
\end{equation}
These bounds are independent of the torus side. One may take
\[
 \mu=\min\{\theta/2,\eta\theta/(8K)\}.
 \]
The same constants apply to every principal matrix
\(H_{\Omega,\Omega}\), on its original index set.
If \(H\) depends \(C^q\) on parameters and, for prescribed
unit directions, its derivatives through order \(q\)
have weighted Schur bounds \(K_1,\ldots,K_q\), the
inverse derivatives have bounds depending only on
\(C_{\mathrm{inv}},K_1,\ldots,K_q,q\), at the smaller
rate \(\mu/2\). These bounds are also common to the
principal matrices when their index sets are held fixed.
\end{lemma}
\begin{proof}
Fix an index \(j\), put
\(w_i=e^{\mu d(i,j)}\), and form
\(H^w=\operatorname{diag}(w)H\operatorname{diag}(w)^{-1}\).
Since the distance from a fixed site is \(1\)-Lipschitz,
\[
 |w_i/w_l-1|\le
          \mu d(i,l)e^{\mu d(i,l)}
       \le (2\mu/\theta)e^{\theta d(i,l)}
 \]
when \(\mu\le\theta/2\).
The row and column sums of \(|H^w-H|\) are at most
\(2\mu K/\theta\le\eta/4\).
The elementary Schur bound gives its Euclidean operator
norm at most \(\eta/4\).
For completeness this bound follows by applying
Cauchy--Schwarz to
\(\sum_l|A_{il}||u_l|\), summing over \(i\), and using
both row and column sums.

Since \(\|H^{-1}\|_{2\to2}\le\eta^{-1}\), the Neumann
series for
\((1+H^{-1}(H^w-H))^{-1}\) shows
\(\|(H^w)^{-1}\|_{2\to2}\le2/\eta\).
But
\((H^w)^{-1}=\operatorname{diag}(w)H^{-1}
                  \operatorname{diag}(w)^{-1}\).
Its \(ij\)-entry, with \(w_j=1\), proves the first bound.

The number of lattice sites at maximum distance \(n\)
is at most \(C(1+n)^3\), uniformly on every torus.
Thus
\[
 \sup_i\sum_l e^{-\mu d(i,l)/2}
       \le C m_0(1+\mu^{-1})^4 .
 \]
The entry bound and symmetry prove the second assertion,
with \(C_{\mathrm{inv}}\le
C m_0\eta^{-1}(1+\mu^{-1})^4\).

A principal matrix has the same gap: extend its test
vector by zero before applying \(H\ge\eta1\).
Restriction decreases every absolute row and column
sum and every distance-shell count. The preceding
proof therefore applies on any such index set, even
when it has no regular boundary shape.

For \(1\le s\le q\), the single inverse identity needed
for all subsequent derivatives is
\begin{equation}
 D^sA^{-1}[v_1,\ldots,v_s]
  =\sum_{(B_1,\ldots,B_b)}(-1)^b
 A^{-1}D^{|B_1|}A[v_{B_1}]A^{-1}\cdots
                  D^{|B_b|}A[v_{B_b}]A^{-1}.
 \label{eq:selected-inverse-partition-formula}
\end{equation}
The sum is over ordered partitions of \(\{1,\ldots,s\}\)
into nonempty blocks. It follows by differentiating
\(AA^{-1}=1\): a derivative on an inverse factor creates
a new singleton block, and one on an existing derivative
joins that block. For example at order two it gives
\[
 -A^{-1}D_{12}A A^{-1}
 +A^{-1}D_1A A^{-1}D_2A A^{-1}
 +A^{-1}D_2A A^{-1}D_1A A^{-1}.
\]
The weighted Schur norm is submultiplicative: apply
\(d(i,j)\le d(i,l)+d(l,j)\) in the matrix product and
sum absolute values. Formula
\eqref{eq:selected-inverse-partition-formula} consequently
gives the explicit derivative bound
\[
 \sum_{b=1}^s\ 
 \sum_{\substack{r_1+\cdots+r_b=s\\r_v\ge1}}
 \frac{s!}{r_1!\cdots r_b!}
 C_{\mathrm{inv}}^{b+1}\prod_{v=1}^b K_{r_v}.
\]
Every principal derivative is a restriction of the
original derivative; thus it satisfies the same bound.
For \(H_t=(1-t)1+tH\), \(0\le t\le1\), use
\(\min(1,\eta)\), \(K+2\) and the same positive-order
derivative bounds. All these statements are independent
of the matrix dimension and torus side. Local coercivity
supplies the gap; this inverse lemma does not assume a
global gauge in which the background is close to one.
\end{proof}

\begin{theorem}[Stationary graph]
\label{thm:actual-stationary-graph}
Fix the profile data of Lemma~\ref{lem:actual-action-local-bounds}.
There are positive \(r_y,r_x,\eta,C_\rho\), chosen before the
ultraviolet cutoff and volume, such that on
\(\|x\|_\infty<r_x\) the equation
\[
 D_y\varphi(x,y)=0,\qquad \|y\|_\infty\le r_y
 \]
has a unique solution \(y=\rho(x)\). It satisfies
\begin{align}
 \rho(0)&=0,\qquad
 \|\rho(x)\|_\infty\le C_\rho\|x\|_\infty,\label{eq:actual-rho-size}\\
 D_y^2\varphi(x,y)&\ge\eta1
            \quad(\|x\|_\infty\le r_x,\ \|y\|_\infty\le r_y).
 \label{eq:actual-uniform-chart-gap}
\end{align}
The graph is real analytic. Its derivatives of each fixed
order have rooted locality bounds, with constants
independent of the cutoff and volume.
It is the same stationary graph for every finite branch
label on this common regular chart.
\end{theorem}
\begin{proof}
Let \(\mathcal G(x,y)=D_y\varphi(x,y)\).
By \eqref{eq:hessian-local-perturbation}, after choosing
\(r_y,r_x\le r_0\) small enough, the perturbation of
\(H_0\) has row and column sums at most \(\eta_0/2\).
The Schur bound then gives
\(D_y\mathcal G(x,y)\ge\eta_0/2\); set \(\eta=\eta_0/2\).
By Lemma~\ref{lem:actual-combes-thomas},
\(\|H_0^{-1}\|_{\infty\to\infty}\le C_{\mathrm{inv}}\).
Decrease the radii further so that
\[
 C_{\mathrm{inv}}K\max(r_x,r_y)\le\tfrac14,\qquad
 C_{\mathrm{inv}}Kr_x\le r_y/4 .
 \]
All these choices involve only the fixed profile.

For fixed \(x\) consider
\[
 \mathcal T_x(y)=y-H_0^{-1}\mathcal G(x,y).
 \]
Its Lipschitz constant on the closed maximum-norm
\(r_y\)-ball is at most \(1/4\), by the Hessian perturbation
bound. At \(y=0\) its norm is at most
\(C_{\mathrm{inv}}K\|x\|_\infty\le r_y/4\).
It maps that ball into the \(r_y/2\)-ball.
Starting at \(0\), the successive differences of its
iterates are bounded by a geometric series with ratio
\(1/4\). Completeness gives a fixed point, and the same
Lipschitz inequality proves uniqueness.
Its norm is at most
\((1-1/4)^{-1}C_{\mathrm{inv}}K\|x\|_\infty\);
increase \(C_\rho\) to \(2C_{\mathrm{inv}}K\).
At \(x=0\) the fixed point is \(0\).

The analytic implicit function theorem at each point
applies to the invertible matrix \(D_y\mathcal G\). It makes these
locally defined analytic graphs agree with the unique
fixed point. Thus \(\rho\) is analytic on the whole
coarse ball.
The quantitative locality of its derivatives will be
proved in Section~\ref{sec:stationary-localisation}, using
the rooted derivative bounds already established.
Finally, the frame and predictor maps on this chart are
the unique regular minimisers, independent of the finite
tags. Hence the equations for all tags coincide.
\end{proof}

The proof uses a counting-\(\ell^2\) spectral gap and a
weighted-kernel estimate to obtain an \(\ell^\infty\)
inverse. The estimate
\(\|u\|_2\le\sqrt d\,\|u\|_\infty\), with \(d\) equal to
the number of free coordinates, is never used to choose
the graph radius.

\section{Coercivity and the actual small-field tail}
\label{sec:actual-coercivity}

\subsection{Quadratic growth in the common chart}

\begin{proposition}[Uniform chart coercivity]
\label{prop:actual-chart-coercivity}
Under Theorem~\ref{thm:actual-stationary-graph}, if
\(\|x\|_\infty<r_x\) and \(\|y\|_\infty\le r_y\), then
\begin{equation}
 \varphi(x,y)-\varphi(x,\rho(x))
                \ge\frac{\eta}{2}\|y-\rho(x)\|_2^2 .
 \label{eq:actual-chart-coercivity}
\end{equation}
The norm on the right is Euclidean counting norm.
\end{proposition}
\begin{proof}
Both endpoints lie in the maximum-norm \(r_y\)-ball.
Its convexity keeps their segment inside the chart.
Taylor's integral formula and \(D_y\varphi(x,\rho(x))=0\)
give
\[
 \varphi(x,y)-\varphi(x,\rho(x))
 =\int_0^1(1-t)
 D_y^2\varphi(x,\rho(x)+t(y-\rho(x)))
               [y-\rho(x),y-\rho(x)]\,\dd t .
 \]
Use \eqref{eq:actual-uniform-chart-gap} and
\(\int_0^1(1-t)\,\dd t=1/2\).
\end{proof}

This lower bound holds on a full small chart.
To estimate configurations outside that chart, we use
the geometry of the actual blocking mismatch and the
small-field indicator.

\subsection{Stability of the principal action on the good set}

\begin{lemma}[Good-set stability and the loop price]
\label{lem:actual-good-stability}
For all sufficiently small \(h\), uniformly in the
cutoff and period,
\begin{equation}
 S_f^{\mathrm{un}}(U)\ge\tfrac{\lambda}{2}E_f(U)
       \quad(\mathcal B_h(U)=\varnothing).
 \label{eq:actual-good-stability}
\end{equation}
If the quadratic loop coefficients have absolute
exponential path budget \(\Lambda\), then
\begin{equation}
 |I_f(U)|\le C f^4\Lambda h^2
       \quad(\mathcal B_h(U)=\varnothing).
 \label{eq:actual-good-loop-bound}
\end{equation}
The constants depend on the group and fixed path budgets,
not on volume.
\end{lemma}
\begin{proof}
Each transported curvature has norm equal to that of its
terminal curvature, which is at most the corresponding
plaquette chord and hence at most \(h\).
For \(k=3,4\) nonnegative numbers \(a_i\le h\),
\[
 \prod_{i=1}^k a_i
       \le\frac{h^{k-2}}{k}\sum_{i=1}^k a_i^2 .
 \]
For \(k=3\), use \(abc\le
(a^2b+b^2c+c^2a)/3\) followed by \(a,b,c\le h\);
equivalently apply the arithmetic--geometric mean
inequality to \(a^2,b^2,c^2\) and their bounded factors.
The same arithmetic--geometric mean argument proves
the displayed bound for all \(k\ge2\).
Indeed, after dividing by \(h^k\), the product is at most
the product to the power \(2/k\), which is at most the
arithmetic mean of the squares.

Apply this bound to each cubic or quartic path
polynomial, after taking absolute coefficient values.
Summation over translates moves each terminal
plaquette through the torus without a volume loss.
The coefficient budgets therefore give
\[
 \left|\sum_\nu p_\nu(U)\right|
           \le C C_{\mathrm p}(h+h^2)E_f(U).
 \]
The row-square terms are nonnegative, and the direct
term is \(\lambda E_f(U)\).
Choose \(h\) so small that
\(C C_{\mathrm p}(h+h^2)\le\lambda/2\).
The signed average preserves the same bound.
For \(I_f\), each of its degree-two path terms has
absolute value at most its coefficient magnitude
times \(h^2\). Sum the budget at each of the \(f^4\)
anchors to obtain \eqref{eq:actual-good-loop-bound}.
\end{proof}

\subsection{The two ways to be outside the central region}

Use the radii from Theorem~\ref{thm:actual-stationary-graph}
and choose \(\delta_*>0\) in
Corollary~\ref{cor:actual-chart-entry} with \(r_*=r_y\).
For \(g=\beta^{-1/2}\) and a central radius \(D\), define
the central set in tree-gauge variables by
\[
 \mathcal C_x=\{U(x,\rho(x)+gz):|z|_2\le D\}.
 \]
For now this is a coordinate set, not an assertion that
every one of its configurations is good.

\begin{proposition}[Coercivity outside the central set]
\label{prop:actual-outside-coercivity}
Suppose \(h\le\delta_*\), \(h\) is in the stability regime
of Lemma~\ref{lem:actual-good-stability}, and
\(\|x\|_\infty<r_x\).
For a good fine configuration \(U\) outside \(\mathcal C_x\),
one of the following two estimates applies:
\begin{align}
 \beta\bigl(A(V(x),U)-\varphi(x,\rho(x))\bigr)
       &\ge\tfrac{\eta}{2}D^2,
 &&\text{if every mismatch is at most }\delta_*,
 \label{eq:tail-chart-case}\\
 A(V(x),U)&\ge a\delta_*^2/2,
 &&\text{if some mismatch exceeds }\delta_* .
 \label{eq:tail-noise-case}
\end{align}
Here \(A=S_f^{\mathrm{un}}+
a\sum_e q_G(V_eQ(U)_e^{-1})\) is the actual action in the
positive reference integral.
\end{proposition}
\begin{proof}
In the first case chart entry gives \(U=U(x,y)\) with
\(\|y\|_\infty\le r_y\).
All the actual regular branches agree there, so
\(A(V(x),U)=\varphi(x,y)\).
Since \(U\notin\mathcal C_x\),
\(\|y-\rho(x)\|_2>gD\).
Use Proposition~\ref{prop:actual-chart-coercivity} and
\(\beta g^2=1\).

In the second case one noise summand is larger than
\(a\delta_*^2/2\), because \(q_G=d_G^2/2\).
All the noise summands are nonnegative, and
\eqref{eq:actual-good-stability} makes the principal
action nonnegative on this fine configuration.
\end{proof}

This argument is specific to the positive small-field
reference contribution in
\eqref{eq:pure-reference-splitting}.
It does not infer a global convexity statement for the
Wilson action.

\subsection{A relative tail bound on auxiliary volumes}

Let \(n_f=f^4\) and \(d=|I|\dim_\R\mathfrak g\).
The Hessian and forcing bounds give
\begin{equation}
 |\varphi(x,\rho(x))|
       \le C n_f\|x\|_\infty^2 .
 \label{eq:stationary-extensive-price}
\end{equation}
To verify it, apply Taylor's formula at \((0,0)\).
The full Hessian has uniformly bounded row sums, so its
quadratic form is bounded by a constant times the number
of joint coordinates times their maximum norm squared.
That number is at most \(C_Gn_f\), and
\eqref{eq:actual-rho-size} bounds the stationary coordinates.

Let \(T_\beta(x)\) be the contribution to
\((\mathcal P_\beta F_{\mathrm{sf},\beta})(V(x))\)
from tree-gauge configurations outside \(\mathcal C_x\).
It is nonnegative. In the central prefactor put
\begin{align*}
 \psi_*(x)&=I_f(U(x,\rho(x)))
                     -\sum_{e\in I}\log j_G(\rho_e(x)),\\
 A_\beta(x)&=
 z_{\mathrm n}(\beta)^{-b_c}\zeta_G^{|I|}
 (2\pi)^{d/2}g^d
 e^{-\beta\varphi(x,\rho(x))-\psi_*(x)}
                         (\det {\mathsf H}(x))^{-1/2},\\
 {\mathsf H}(x)&=D_y^2\varphi(x,\rho(x)).
\end{align*}
These are precisely the stationary prefactors of
Section~\ref{sec:gaussian-remainder}.

\begin{theorem}[The actual reference tail on auxiliary tori]
\label{thm:actual-auxiliary-tail}
Fix the block, group, branch, and principal path data.
Let \(\widehat\beta\to\infty\), and assume
\[
 \widehat\beta/2\le\beta\le2\widehat\beta,\quad
 h=\widehat\beta^{-1/2}P(\widehat\beta),\quad
 D=d_0P(\widehat\beta),\quad
 \|x\|_\infty\le\widehat\beta^{-7/16},
 \]
with fixed \(d_0>0\).
Suppose
\[
 n_f\le C_1(1+\log\widehat\beta)^{m_1},\qquad
 \Lambda\le C_\nu\widehat\beta^\nu
                     \quad\text{for every }\nu>0
 \]
eventually. Then there is \(c>0\) such that
\begin{equation}
 0\le T_\beta(x)/A_\beta(x)\le e^{-cD^2}
 \label{eq:actual-relative-tail}
\end{equation}
eventually, uniformly on all these auxiliary tori.
In particular this ratio is smaller than every inverse
power of \(\widehat\beta\).
Constants and lower thresholds may depend on
\(G,L,J,\sigma,C_{\mathrm p},\lambda,a,d_0,C_1,m_1\),
the fixed branch parameters, and the stated bounds for
\(\Lambda\), but not on the physical cutoff or volume.
\end{theorem}
\begin{proof}
The first case of Proposition
~\ref{prop:actual-outside-coercivity} gives the exponent
loss \(\eta D^2/2\).
For the second case,
\eqref{eq:stationary-extensive-price} and the auxiliary
size bound imply
\[
 |\varphi(x,\rho(x))|
 \le C n_f\widehat\beta^{-7/8}
                   \le a\delta_*^2/4
 \]
eventually. Hence
\[
 \beta\bigl(A(V(x),U)-\varphi(x,\rho(x))\bigr)
                 \ge\beta a\delta_*^2/4.
 \]
Since \(D^2\) is subpower in \(\widehat\beta\), this is
at least \(\eta D^2/2\) eventually.
Thus both cases have a common loss \(c_0D^2\), with
\(c_0>0\) depending only on the fixed profile.

On the good set, \eqref{eq:actual-good-loop-bound} gives
\(|I_f(U)|\le C n_f\Lambda h^2\).
After the loss is factored out, the remaining tree Haar
measure has mass at most \(1\). All actual finite branch
likelihoods sum to \(1\).
Consequently
\[
 T_\beta(x)\le z_{\mathrm n}(\beta)^{-b_c}
 e^{-\beta\varphi(x,\rho(x))
                         +C n_f\Lambda h^2-c_0D^2}.
 \]
The noise normaliser cancels exactly on division by
\(A_\beta\).

We bound the other logarithmic costs.
The entries of \(\rho\) are \(O(\widehat\beta^{-7/16})\).
Every terminal plaquette at this configuration has chord
at most \(C\widehat\beta^{-7/16}\); transported curvatures
have the same norm bound regardless of path length.
Therefore
\[
 |I_f(U(x,\rho(x)))|
                \le C n_f\Lambda\widehat\beta^{-7/8}.
 \]
Bounded first derivatives of \(\log j_G\) on a fixed
chart, and \(\log j_G(0)=0\), give
\[
 |\psi_*(x)|\le
 Cn_f\bigl(\Lambda\widehat\beta^{-7/8}
                           +\widehat\beta^{-7/16}\bigr).
 \]
Both terms tend to zero after using any sufficiently
small exponent in the subpower bound for \(\Lambda\).
The eigenvalues of \({\mathsf H}(x)\) lie between \(\eta\) and a
fixed upper row-sum bound. Hence
\(|\log\det {\mathsf H}(x)|\le Cd\).
The remaining factors
\(\zeta_G^{|I|}(2\pi)^{d/2}g^d\) cost at most
\(Cd(1+\log\widehat\beta)\) in their logarithm.
Also \(d\le C_Gn_f\).

Altogether the logarithm of the tail ratio is at most
\[
 -c_0D^2+
 Cn_f\bigl(1+\log\widehat\beta+\Lambda h^2
                  +\Lambda\widehat\beta^{-7/8}
                  +\widehat\beta^{-7/16}\bigr).
 \]
The positive terms grow at most as a fixed power of
\(1+\log\widehat\beta\), plus terms tending to zero.
The function \(D^2=d_0^2
\exp(2(\log\log\widehat\beta)^2)\) dominates every such
power. This proves \eqref{eq:actual-relative-tail},
with \(c=c_0/2\), after increasing the threshold.
Finally \(D^2/\log\widehat\beta\to\infty\).
\end{proof}

\subsection{What remains for the central cutoff}

The tail theorem bounds the actual integral outside the
central coordinate set. To replace its inside part by
the unmasked \(A_\beta Z_D\), the central set must also
lie in the actual good set.
The size estimate \(\|\rho(x)\|_\infty\le C_\rho\|x\|_\infty\)
does not establish this: the chart width
\(\widehat\beta^{-7/16}\) is larger than the plaquette
threshold \(\widehat\beta^{-1/2}P(\widehat\beta)\).
Thus this inclusion is a separate, quantitative issue.

The following useful criterion requires only one
configuration with a reserve in its plaquette bounds.

\begin{lemma}[Transport of a good central configuration]
\label{lem:central-good-reserve}
Suppose \(U(x_0,\rho(x_0)+gz_0)\), with \(|z_0|_2\le D\),
has every plaquette chord at most \(h/2\).
There is \(C_{\mathrm{cut}}\), depending only on the
fixed profile, such that every central configuration at
\(x\) is good whenever
\begin{equation}
 C_{\mathrm{cut}}
    \bigl((1+C_\rho)\|x-x_0\|_\infty+2gD\bigr)\le h/2 .
 \label{eq:central-cutoff-reserve}
\end{equation}
\end{lemma}
\begin{proof}
On a fixed link chart, the map
\((x,y)\mapsto U(x,y)\) is Lipschitz in maximum link norm,
with a fixed constant. The chord of a four-link
plaquette changes by at most a fixed constant times
that link distance, by telescoping its four factors.
The stationary graph has Lipschitz constant \(C_\rho\),
after increasing it using its derivative estimate.
For \(|z|_2,|z_0|_2\le D\),
\(\|z-z_0\|_\infty\le2D\).
These facts bound the change of every plaquette chord
by the left side of
\eqref{eq:central-cutoff-reserve}.
Add its original bound \(h/2\).
\end{proof}

When \(D=d_0P(\widehat\beta)\), choosing \(d_0\) small
pays the \(gD/h\) part of this criterion uniformly.
A coarse neighbourhood of radius
\(\widehat\beta^{-1}\) pays the first part.
At the flat coarse field the initial configuration can
be chosen explicitly: \(x_0=0,z_0=0\), so that \(U=1\).
This gives the following unconditional comparison near
the identity.

\begin{corollary}[The actual central comparison near the identity]
\label{cor:actual-central-comparison}
Under Theorem~\ref{thm:actual-auxiliary-tail}, there are
\(c_{\rm in},d_{\rm in}>0\), depending only on the fixed
profile, such that the following holds. Choose
\(0<d_0\le d_{\rm in}\) and
\(\|x\|_\infty\le c_{\rm in}h\).
For all sufficiently large \(\widehat\beta\),
the entire central set is good, and
\begin{equation}
 0\le
 (\mathcal P_\beta F_{\mathrm{sf},\beta})(V(x))
                  -A_\beta(x)Z_D(g,x)
       \le A_\beta(x)e^{-cD^2}.
 \label{eq:actual-central-comparison}
\end{equation}
All prefactors and the function \(Z_D\) are those of
the actual reference integral.
\end{corollary}
\begin{proof}
Since \(\rho(0)=0\), the configuration at \(x_0=z_0=0\)
has all plaquette chords equal to zero.
Also \(gD/h\le\sqrt2\,d_0\) because
\(\beta\ge\widehat\beta/2\).
Choose the two constants so that
\[
 C_{\mathrm{cut}}
       \bigl((1+C_\rho)c_{\rm in}+2\sqrt2\,d_{\rm in}\bigr)
                                                    \le\tfrac12 .
\]
Lemma~\ref{lem:central-good-reserve} proves the required
central inclusion. All its coordinates lie in the
common regular chart eventually.
There the branch predictions coincide, the branch
likelihoods sum to \(1\), and the fine indicator is \(1\).
The exact Haar change of variables and the Gaussian
formula of Section~\ref{sec:gaussian-remainder} therefore
identify the central contribution with \(A_\beta Z_D\).
The remaining contribution is precisely \(T_\beta\).
Apply \eqref{eq:actual-relative-tail}.
\end{proof}

For a general physical good coarse field, after the
local gauge and extension operations used in localisation,
the existence of the initial configuration with threshold
\(h/2\) requires a physical background or lower-bound
construction. Sections~\ref{sec:microscopic-centre}--%
\ref{sec:general-background-comparison} supply this
construction on auxiliary charts with the explicitly
specified calibrated linear RG data.
It is not a consequence of chart coercivity alone.
The preceding near-flat corollary verifies the two-sided comparison
\eqref{eq:physical-tail-hypothesis} on the stated near-flat
domain; it does not enlarge that domain to all physical
good fields.
\section{Localisation of the stationary and quadratic Gaussian terms}
\label{sec:stationary-localisation}

The stationary equation has been solved on the actual finite
torus. We next ask how its solution depends on coarse fields
far from a specified block. A bound on the largest derivative
alone would not answer this question. We need to sum derivatives
over their input sites, with an exponential weight measuring
the length needed to connect those sites.

We use the norms in
\eqref{eq:weighted-schur-norm}--\eqref{eq:rooted-tensor-norm}.
In this section a tree may contain additional lattice vertices,
as specified there. This minor convention is useful when an
inverse Hessian joins two derivative tensors.
All derivatives are ordinary derivatives in the exponential
coordinates already defined. Their conversion to physical
left derivatives will be distinguished explicitly.

The separation of the stationary action, Gaussian normalisation, and
remaining interaction may be compared with
\cite[Sect.~1, eqs.~(1.3)--(1.7), printed pp.~260--261]{Bal87}.
Localisation of the fluctuation interaction precedes its connected
expansion in \cite[Sect.~1, Lemmas~1--2, printed pp.~9 and~11;
Sect.~2, eqs.~(2.11)--(2.13), p.~14]{Bal88a}.
These are organisational comparisons; the rooted-kernel and determinant
estimates in this section are proved for the present Hessian.

\subsection{Why the weighted tensor norms are stable under substitution}

\begin{lemma}[Joining rooted kernels]
\label{lem:rooted-kernel-composition}
Suppose a rooted tensor is substituted into each input slot
of another rooted tensor. Sum over the indices that are joined.
The rooted norm of the resulting tensor is at most the product
of the rooted norms of the tensors being joined, with the same
exponential weight.
In particular, for matrices,
\[
 \|AB\|_\theta\le\|A\|_\theta\|B\|_\theta .
\]
The analogous assertion holds for rectangular matrices with
different input and output coordinate types.
\end{lemma}
\begin{proof}
For fixed indices in a summand, take connecting trees for the
outer tensor and for all the tensors inserted into its slots.
They meet at the joined indices. Their union is connected
and contains the output root and all the remaining input
sites. Removing cycle edges gives a connecting tree.
Consequently
\[
 \mathsf t(\hbox{output root, remaining inputs})
 \le
 \mathsf t(\hbox{outer indices})
       +\sum_{\text{insertions}}\mathsf t(\hbox{inserted indices}).
\]
The exponential weight of the summand is therefore bounded
by the product of these weights.
Take absolute values. Sum first over the inputs of the
inserted tensors, using their uniform rooted bounds at
each joined index. Then sum over the outer tensor's inputs.
All summands are nonnegative, so these changes of order
require no conditional rearrangement.
The matrix row estimate is this argument with one input
per tensor; the column estimate follows by reversing the
product.
\end{proof}

Notice that a connected graph produced by this operation
need not itself be a tree. The proof uses a tree to bound
its connecting length. It does not identify different
connected graphs with a single labelled tree.

\subsection{Locality of derivatives of the actual stationary graph}

Write
\[
 {\mathsf H}(x)=\varphi_{yy}(x,\rho(x)),\qquad
 {\mathsf H_{yx}}(x)=\varphi_{yx}(x,\rho(x)).
\]
The subscripts denote ordinary partial derivatives.
Their indices include the link and its Lie algebra component.
By the gap and local bounds proved above, the weighted inverse
estimate applies to \({\mathsf H}(x)\), with constants independent of
the torus side. Reduce the exponential weight to one common
number \(\theta>0\), smaller than both the action's derivative
weight and the inverse's weight.

\begin{proposition}[Local response of the stationary graph]
\label{prop:stationary-response-locality}
On a smaller fixed coarse ball, for each fixed \(k\ge1\)
there is \(C_k<\infty\) such that
\begin{equation}
 \sup_i\sum_{j_1,\ldots,j_k}
 e^{\theta\mathsf t(i,j_1,\ldots,j_k)}
 \left|
   \partial_{x_{j_1}}\cdots\partial_{x_{j_k}}\rho_i(x)
 \right|\le C_k .
 \label{eq:stationary-response-locality}
\end{equation}
The output index is fine and the input indices are coarse.
The constants depend on the fixed group, block and branch
data, path width and budgets, and \(k\), but not on the
cutoff or number of blocks.
\end{proposition}
\begin{proof}
Differentiate \(\mathcal G(x,\rho(x))=0\), where
\(\mathcal G=\varphi_y\). At first order,
\begin{equation}
 D\rho(x)=-{\mathsf H}(x)^{-1}{\mathsf H_{yx}}(x).
 \label{eq:actual-linear-response}
\end{equation}
The weighted inverse bound and
Lemma~\ref{lem:rooted-kernel-composition} prove the assertion
for \(k=1\).

For higher orders put \(\Gamma(x)=(x,\rho(x))\).
For directions \(u_1,\ldots,u_k\), repeated use of the
ordinary chain rule gives
\[
 D^k(\mathcal G\circ\Gamma)[u_1,\ldots,u_k]
 =
 \sum_{\pi}
 D^{|\pi|}\mathcal G(\Gamma)
   \bigl[D^{|B|}\Gamma[u_j:j\in B]:B\in\pi\bigr].
\]
Here \(\pi\) ranges over the set partitions of
\(\{1,\ldots,k\}\), and \(B\) is a block of that partition.
This identity follows inductively: an extra derivative
either makes a new block by differentiating \(\mathcal G\), or joins
an existing block by differentiating that argument.
Every partition is thereby produced once.

For \(k\ge2\) the partition with one block contributes
\({\mathsf H} D^k\rho\), since \(D^k x=0\).
In every other partition each block has size less than \(k\).
Move these terms to the other side and multiply by
\({\mathsf H}^{-1}\).
All derivatives of \(\mathcal G\) have the rooted bounds from
Lemma~\ref{lem:actual-action-local-bounds}; the derivatives
of \(\Gamma\) occurring on the right have bounds by
induction. Lemma~\ref{lem:rooted-kernel-composition} bounds
each contraction of indices.
There are finitely many partitions, depending only on \(k\).
For example the resulting constants can be chosen recursively
by
\[
 C_1=C_{\rm inv}K_1,\qquad
 C_k=C_{\rm inv}
   \sum_{\substack{\pi\\|\pi|\ge2}}
       K_{|\pi|}
       \prod_{B\in\pi}\bigl(1_{\{|B|=1\}}+C_{|B|}\bigr).
\]
Increase \(K_j\) to bound all required mixed derivatives
of \(\mathcal G\). Every \(C_{|B|}\) in the second expression has
\(|B|<k\), so this is a genuine finite recursion.
\end{proof}

In particular a change of coarse field supported at distance
at least \(R\) from an output link changes that stationary
link by at most \(C_1e^{-\theta R}\) times the maximum
change, provided the connecting segment of coarse fields
stays in the chart. This follows by integrating the first
derivative and summing only over those distant inputs.
It is a statement about this same finite-period solution.

\subsection{An absolute analytic norm}

To turn derivative locality into an absolutely convergent
family of local functions, we keep all Taylor orders.
This is legitimate here because the path budget gave an
absolute analytic majorant, not only finitely many
smoothness bounds.

For a vector-valued analytic function \(F\) with an output
root, write
\[
 F_i(z)=\sum_{k\ge0}\sum_{j_1,\ldots,j_k}
       F^{[k]}_{i;j_1,\ldots,j_k}
                         z_{j_1}\cdots z_{j_k},
 \qquad F^{[k]}=\frac1{k!}D^kF(0).
\]
The inputs \(z\) can be the joint coordinates \((x,y)\).
Define
\begin{equation}
 \mathcal A_{r,\theta}(F)
       =\sum_{k\ge0}r^k|F^{[k]}|_{\theta,k}.
 \label{eq:analytic-rooted-norm}
\end{equation}
For \(k=0\) the rooted norm means \(\sup_i|F_i^{[0]}|\).
The radius \(r\) is a coordinate radius and is unrelated
to the exponent \(s\ge2\) in the physical derivative norm.
Repeated input indices are allowed. The factor \(1/k!\)
makes the displayed ordered-index expansion the usual
Taylor series.

For a matrix-valued function \(A(x)\) we use the corresponding
kernel norm
\begin{equation}
 \mathcal A_{r,\theta}(A)
 =\sum_{k\ge0}r^k
   \sup_i\sum_{j,j_1,\ldots,j_k}
   e^{\theta\mathsf t(i,j,j_1,\ldots,j_k)}
                  |A^{[k]}_{ij;j_1,\ldots,j_k}|.
 \label{eq:analytic-kernel-norm}
\end{equation}
The index \(j\) is a matrix index, rather than a field input.
At \(k=0\) this is a weighted row norm. For a constant
matrix it is bounded by \(\|A\|_\theta\).
Symmetry is unnecessary for the product estimate
\begin{equation}
 \mathcal A_{r,\theta}(AB)
       \le\mathcal A_{r,\theta}(A)
                         \mathcal A_{r,\theta}(B).
 \label{eq:analytic-kernel-product}
\end{equation}
Indeed expand both series, apply
Lemma~\ref{lem:rooted-kernel-composition}, and sum the
convolution of their degrees. Averaging over permutations
of the field inputs makes each coefficient symmetric and
does not increase the absolute norm.

\begin{lemma}[A uniform analytic majorant for the graph]
\label{lem:analytic-stationary-majorant}
There are \(r_{\rm a},C_{\rm a}>0\), chosen independently of
cutoff and volume, such that
\begin{equation}
 \mathcal A_{r,\theta}(\rho)\le C_{\rm a}r,
       \qquad 0\le r\le r_{\rm a}.
 \label{eq:analytic-rho-bound}
\end{equation}
After decreasing \(r_{\rm a}\), the matrix
\(E(x)={\mathsf H}(x)-H_0\) satisfies
\begin{equation}
 \mathcal A_{r,\theta}(E)\le C_{\rm a}r .
 \label{eq:analytic-hessian-bound}
\end{equation}
\end{lemma}
\begin{proof}
Multiply the Taylor series of \(\mathcal G(x,y)=0\) by
\(H_0^{-1}\) and isolate its linear \(y\)-term:
\begin{equation}
 y=B x+\mathcal R(x,y),\qquad
 B=-H_0^{-1}\mathcal G_x(0,0).
 \label{eq:majorant-stationary-equation}
\end{equation}
The series \(\mathcal R\) has no terms of degree zero or one.
The weighted inverse bound and the gradient's absolute
analytic majorant give constants \(a\ge1,b\ge1\) such that
\[
 |B|_{\theta,1}\le a,\qquad
 \sum_{k\ge2}u^k|\mathcal R^{[k]}|_{\theta,k}
                    \le\frac{bu^2}{1-bu}
              \quad(0\le u<b^{-1}).
\]
For clarity, such \(b\) follows from any finite bound
\(\sum_k R^k|\mathcal R^{[k]}|_{\theta,k}\le K\):
each coefficient is at most \(KR^{-k}\), and increasing
\(b\) bounds these numbers by \(b^{k-1}\) for \(k\ge2\).
Both \(K\) and \(R\) were chosen from local path and
inverse budgets.

Consider the scalar iteration
\[
 t_0(r)=0,\qquad
 t_{n+1}(r)=ar+
       \frac{b(r+t_n(r))^2}{1-b(r+t_n(r))}.
 \]
All its power-series coefficients are nonnegative and
increase with \(n\).
Choose \(r_{\rm a}>0\) so small that for
\(0\le r\le r_{\rm a}\), \(0\le t\le2ar\),
\[
 b(r+t)\le\tfrac12,\qquad
 \frac{b(r+t)^2}{1-b(r+t)}\le ar,\qquad
 \left|\partial_t\frac{b(r+t)^2}{1-b(r+t)}\right|
                                                 \le\tfrac12 .
 \]
For example all three inequalities follow by taking a
sufficiently small positive multiple of
\(b^{-1}(1+2a)^{-2}\).
Thus \(t_n(r_{\rm a})\le2ar_{\rm a}\).
The increasing coefficients have limits and their
absolute series is bounded there by the same number.
The limit \(t(r)\) solves the scalar fixed-point equation
and satisfies \(t(r)\le2ar\).

Now iterate the vector equation
\eqref{eq:majorant-stationary-equation} starting at \(0\)
in its Taylor series. Apply
Lemma~\ref{lem:rooted-kernel-composition} at each substitution.
It follows inductively that its rooted coefficient sums are dominated by those
of \(t_n\). Each fixed Taylor degree stabilises:
the nonlinear remainder has degree at least two, so a
change first appearing at degree \(k\) appears in the
next iterate only at degree \(k+1\) or higher.
The limiting series is absolutely convergent on the
smaller coordinate ball, solves the vector equation,
and has rooted norm at most \(t(r)\).
It lies in the stationary uniqueness ball when
\(r_{\rm a}\) is decreased once more.
It therefore equals the actual graph \(\rho\).

The Hessian \(\varphi_{yy}\) has an absolute rooted
analytic kernel majorant on a smaller radius, by the
same local analytic estimates used for the gradient.
Differentiating the analytic series costs a smaller
radius, and no spatial weight is lost.
Substitute \((x,\rho(x))\), using
\eqref{eq:analytic-rho-bound} and the joining lemma.
The constant term is \(H_0\). The remaining majorant
vanishes at radius zero and is bounded by \(Cr\) on
a sufficiently small fixed interval. This proves
\eqref{eq:analytic-hessian-bound}.
\end{proof}

This proof supplies a common positive analytic radius.
It avoids applying a finite-dimensional Cauchy bound
whose constants could grow with the number of blocks.
It also completes the quantitative derivative assertion
in Theorem~\ref{thm:actual-stationary-graph}.
For example, at \(\|x\|_\infty\le r_{\rm a}/2\),
termwise differentiation gives the bound
\[
 |D^k\rho(x)|_{\theta,k}
 \le C_{\rm a}r_{\rm a}^{\,1-k}
       \sup_{l\ge k}(l)_k\,2^{-(l-k)} ,
\]
where \((l)_k=l(l-1)\cdots(l-k+1)\).
The supremum is finite for each \(k\).

\subsection{The on-shell classical action}

Define
\[
 \Psi(x)=\varphi(x,\rho(x)).
 \]
The word on-shell means that the fine variable has been
set equal to its stationary value. It does not mean
that the coarse field solves its own stationary equation.
The total value of \(\Psi\) is extensive.

\begin{proposition}[Local classical response]
\label{prop:classical-local-response}
The gradient of \(\Psi\) has a uniformly summable rooted
analytic expansion, vanishes at \(0\), and satisfies
\begin{equation}
 \Psi_{xx}
   =\varphi_{xx}
           -\varphi_{xy}{\mathsf H}^{-1}\varphi_{yx}
       \quad\hbox{at }(x,\rho(x)).
 \label{eq:actual-schur-complement}
\end{equation}
In particular \(\Psi\) is an absolutely convergent sum
of rooted coordinate monomials of degree at least two.
Their rooted analytic coefficient norm is bounded by
\(Cr^2\), uniformly per coarse site, at small fixed
radius \(r\).
\end{proposition}
\begin{proof}
Stationarity gives
\[
 \partial_{x_i}\Psi(x)=\varphi_{x_i}(x,\rho(x)).
 \]
Apply the gradient's analytic path bound and
Lemma~\ref{lem:analytic-stationary-majorant}.
The resulting rooted analytic norm is at most \(Cr\),
since the value at \(0\) is zero.
Differentiating the preceding identity and using
\eqref{eq:actual-linear-response} proves
\eqref{eq:actual-schur-complement}.

Write \(\varphi_{x_i}(x,\rho(x))
=\sum_{k\ge1}\sum_{\mathbf j}a_{i;\mathbf j}x_{\mathbf j}\),
where \(x_{\mathbf j}=\prod_{\nu=1}^kx_{j_\nu}\).
Since \(\Psi(0)=0\), integration along the radial segment
gives the exact series
\begin{equation}
 \Psi(x)=
 \sum_i\sum_{k\ge1}\sum_{\mathbf j}
             \frac{a_{i;\mathbf j}}{k+1}\,x_i x_{\mathbf j}.
 \label{eq:rooted-classical-series}
\end{equation}
Its absolute rooted coefficient sum is at most \(Cr^2\).
There are only \(4\dim\mathfrak g\) possible coarse
coordinate roots per coarse site.
On each finite torus the sum over those sites is finite.
This proves both the identity and the stated uniform
per-site bound.
\end{proof}

\subsection{The quadratic Gaussian determinant on the same torus}

Put \(C_0=H_0^{-1}\) and
\[
 A(x)=C_0({\mathsf H}(x)-H_0)=C_0E(x).
 \]
Decrease \(r_{\rm a}\) so that
\(\mathcal A_{r_{\rm a},\theta}(A)\le1/2\).
This is possible by \eqref{eq:analytic-hessian-bound}.
The precision matrices and all inverses in what follows
are those of the actual finite torus. No infinite-volume
inverse is substituted.

\begin{proposition}[An absolutely convergent determinant expansion]
\label{prop:actual-determinant-localisation}
For \(\|x\|_\infty<r_{\rm a}\),
\begin{equation}
 \frac12\log\frac{\det {\mathsf H}(x)}{\det H_0}
  =\sum_i p_i(x),\qquad
 p_i(x)=\frac12\sum_{n\ge1}
              \frac{(-1)^{n+1}}n(A(x)^n)_{ii}.
 \label{eq:actual-logdet-series}
\end{equation}
Each rooted coefficient series is absolutely summable.
Its tree-weighted analytic norm is at most \(Cr\)
at radius \(0\le r\le r_{\rm a}\), uniformly per fine
coordinate root.
\end{proposition}
\begin{proof}
First prove the finite-matrix identity.
For \(0\le t\le1\) the matrix \(H_0+tE\) is positive
definite, since it is a convex combination of \(H_0\)
and \({\mathsf H}\). Jacobi's formula gives
\[
 \frac{\dd}{\dd t}\log\det(H_0+tE)
 =\operatorname{Tr}\bigl((1+tA)^{-1}A\bigr).
 \]
The absolute row norm of \(A(x)\) is at most \(1/2\).
Its Neumann series converges uniformly in \(t\).
Integrate this series from \(0\) to \(1\) and take the
finite trace. This proves the identity in
\eqref{eq:actual-logdet-series}.
The matrix \(A\) itself need not be symmetric; the
argument uses positive definiteness of \(H_0+tE\).

To obtain locality, expand each matrix factor into
its Taylor series. Write \(a(r)=
\mathcal A_{r,\theta}(A)\le Cr\), with \(a(r)\le1/2\).
By \eqref{eq:analytic-kernel-product}, the rooted norm
of the expansion of \(A^n\) is at most \(a(r)^n\).
Restricting the final matrix index to the initial
index \(i\) only decreases this absolute sum.
Therefore the rooted norm of \(p_i\) is bounded by
\[
 \frac12\sum_{n\ge1}\frac{a(r)^n}{n}
       \le a(r)\le Cr .
 \]
Every coefficient has a connected description:
the matrix indices form a closed chain, and the
field inputs of each factor are joined to its two
matrix indices by a connecting tree.
The sum of these lengths pays the connecting length
of their union, by the joining lemma.
More precisely, one may retain the summed matrix
indices as separate coefficient labels and weight
each label by the sum of the connecting-tree lengths
of its factors. The product estimate still applies:
its proof sums exactly the product of the individual
weighted absolute coefficients, before any matrix
indices are combined. Thus this more detailed labelled
series obeys the same majorant \(a(r)^n\).
The number of fine roots at one coarse site is
bounded by \(m_0\). This converts the last bound to
a uniform per-coarse-site estimate as well.
\end{proof}

The constant \(\frac12\log\det H_0\) is not bounded
uniformly as a scalar when the volume grows.
It belongs to the vacuum normalisation.
The proposition concerns the field-dependent difference
and its summability per site.

There is a similar elementary localisation of the
remaining stationary prefactor
\[
 \psi_*(x)=I_f(U(x,\rho(x)))
                     -\sum_{e\in I}\log j_G(\rho_e(x)).
 \]
The quadratic loop has the same based-path structure
as above, with coefficient budget \(\Lambda\).
Its value and first derivative vanish at the identity.
The loop's rooted analytic expansion is therefore
bounded by \(C\Lambda r^2\), uniformly per site, after
substitution of the stationary graph.
For the Haar term, \(j_G(0)=1\) and \(Dj_G(0)=0\).
Indeed the linear term of
\(\int_0^1e^{-t\,\operatorname{ad}X}\dd t\) is
\(-\operatorname{ad}X/2\), whose trace vanishes because
\(\operatorname{ad}X\) is skew adjoint.
Thus \(\log j_G(X)=O(|X|^2)\) analytically near \(0\).
Each such term initially reads one fine link.
The graph majorant and the bounded number of fine
links per block give a rooted bound \(Cr^2\).
Combining the two gives
\begin{equation}
 \mathcal A_{r,\theta}(\psi_*)_{\rm per\ site}
                         \le C(1+\Lambda)r^2 .
 \label{eq:stationary-loop-haar-localisation}
\end{equation}
Here the notation means the absolute sum of the
rooted scalar coefficient families just described;
it is not an absolute bound on the total scalar
function over the whole volume.

\subsection{Connected footprints and the covering weight}

We now check explicitly that these rooted expansions
give connected supports with the covering weight used
for activities. This is stronger than merely controlling
the value of a kernel far from a fixed root.

For each coefficient label, retain a chosen connecting
tree from the preceding estimates, including any summed
matrix indices. If its total metric length is \(l\),
replace each tree edge by a coordinate path on the
torus. Each path has at most four times that edge's
periodic maximum distance in nearest-neighbour edges.
Let \(X\) be the distance-one neighbourhood of the union
of these paths. This includes the two endpoints of each
coarse link read by the monomial.
It is a connected footprint with a specified root.

\begin{lemma}[The price of a connected path footprint]
\label{lem:tree-covering-price}
For \(B\ge2\), a footprint constructed in this way obeys
\begin{equation}
 n_B(X)\le2+\frac{16l}{B}.
 \label{eq:tree-covering-price}
\end{equation}
Consequently, if \(B\ge64\kappa/\theta\),
\[
 e^{\kappa n_B(X)}
                 \le e^{2\kappa}e^{\theta l/4}.
 \]
\end{lemma}
\begin{proof}
The coordinate paths have at most \(4l\) edges in total.
Their connected union can be traversed by a walk of
at most \(8l\) edges, starting at the root: traverse
each tree edge out and back along its coordinate path.
Choose centres successively on this walk at intervals
of \(B-1\) steps, including its first centre.
At most \(1+\lceil8l/(B-1)\rceil\) centres are needed.
Every point on the walk is within \(B-1\) of a centre.
Its distance-one neighbourhood is thus covered by
the radius-\(B\) balls. Consecutive balls intersect,
so their union is connected.
Since \(B-1\ge B/2\), their number is at most
\(2+16l/B\). For \(l=0\) a single radius-\(B\) ball
suffices.
\end{proof}

To state the norm bound precisely within the current
coordinate chart, define
\[
 \|f\|^{\rm chart}_{q,r,s;o,X}
 =\max_{0\le j\le q}
   \sup_{\max_e|x_e|_{\mathfrak g}\le r}
   \sup_{\substack{|u_{\nu,e}|\le
               r(1+d(e,o))^s\\1\le\nu\le j}}
         |D^jf(x)[u_1,\ldots,u_j]|.
 \]
Only coordinates read by \(f\) matter in the inner
supremum; \(d(e,o)\) is the distance from the initial
site of its coarse link to the root.
The outer supremum uses the maximum Lie algebra norm
over links. Its comparison with the maximum of the
individual colour components costs only a fixed group
constant; shrink the analytic radii accordingly.
This is a coordinate-chart norm.
Unlike the physical activity norm, its domain requires
each link itself to be close to \(1\).

\begin{proposition}[Weighted local families in the common chart]
\label{prop:stationary-chart-local-families}
Fix \(q,s,\kappa\), and choose
\(B\ge\max(2,64\kappa/\theta)\).
On a smaller fixed chart the field-dependent stationary
prefactor has an exact expansion
\begin{equation}
 \beta\Psi(x)+\psi_*(x)
          +\frac12\log\frac{\det {\mathsf H}(x)}{\det H_0}
                 =\sum_\alpha f_\alpha(x).
 \label{eq:stationary-local-family}
\end{equation}
Each \(f_\alpha\) reads a connected footprint \(X_\alpha\)
with root \(o_\alpha\), vanishes at \(0\), and
\begin{equation}
 \sup_o\sum_{\alpha:o_\alpha=o}
       e^{\kappa n_B(X_\alpha)}
       \|f_\alpha\|^{\rm chart}_{q,r,s;o,X_\alpha}
 \le C_{q,s,\kappa}
               \bigl(\beta r^2+(1+\Lambda)r^2+r\bigr).
 \label{eq:stationary-chart-family-bound}
\end{equation}
The constant is independent of volume and cutoff.
Its dependence on the fixed group, block, branch,
path and gap data remains as stated above.
\end{proposition}
\begin{proof}
Use the monomials in
\eqref{eq:rooted-classical-series}, the path and Haar
expansions in \eqref{eq:stationary-loop-haar-localisation},
and the matrix-chain expansions of
\eqref{eq:actual-logdet-series}.
Keep their index labels, even if two labels happen
to give the same monomial.
Assign the connected footprints just constructed.
All these assignments preserve their exact finite-torus
series identities.

For a degree-\(k\) monomial of coefficient magnitude
\(b\) and connecting length \(l\), its chart derivative
norm is at most
\[
 C_q b(1+k)^q r^k(2+4l)^{sq}.
 \]
In fact every read coordinate is at distance at most
\(4l\) from its root. At order \(j\), differentiation
chooses at most \((k)_j\) ordered factors, removes
\(j\) powers of \(r\), and supplies at most
\(r^j(1+4l)^{sj}\) from the allowed directions.
This proves the bound also for \(j=0\).

The footprint price costs at most
\(e^{2\kappa}e^{\theta l/4}\).
The polynomial in \(l\) is bounded by
\(C_{q,s,\theta}e^{\theta l/4}\).
Finally \((1+k)^q\le C_q2^k\).
The rooted analytic coefficient estimates at radius
\(2r\), with their full \(e^{\theta l}\) weights,
therefore bound all these sums on a smaller chart.
The classical, loop--Haar, and determinant series
respectively cost \(C\beta r^2\),
\(C(1+\Lambda)r^2\), and \(Cr\).
There are at most \(m_0\) possible coordinate roots
at a coarse site.
These observations prove
\eqref{eq:stationary-chart-family-bound} and absolute
convergence of all differentiated series through
order \(q\).
\end{proof}

\subsection{What the Taylor remainder already gains}

The width \(w=\widehat\beta^{-7/16}\) has a useful
quantitative purpose. The classical phase has a factor
\(\beta\), so its Taylor terms must be retained through
degree four before the remaining local norm is small
enough for the desired RG budget.

For an analytic function \(F\), let
\(T_{\le m}F\) denote its Taylor polynomial at \(0\)
through total degree \(m\).
For an extensive scalar this notation means the
rooted Taylor series constructed above, truncated
degree by degree on its finite torus.
Set
\[
 \mathcal L(x)=\frac12\log\frac{\det {\mathsf H}(x)}{\det H_0},
 \qquad
 \mathcal R_{\rm st}(x)=
 \beta(1-T_{\le4})\Psi(x)
        +(1-T_{\le2})(\psi_*(x)+\mathcal L(x)).
 \]
Here \(1-T_{\le m}\) is the remainder operator,
not multiplication by a scalar.

\begin{corollary}[The local Taylor remainder of the stationary prefactor]
\label{cor:stationary-taylor-remainder}
Under Proposition~\ref{prop:stationary-chart-local-families},
the remainder has an exact connected-footprint
expansion \(\mathcal R_{\rm st}=\sum_\alpha r_\alpha\).
For \(r\) sufficiently small,
\begin{equation}
 \sup_o\sum_{\alpha:o_\alpha=o}
 e^{\kappa n_B(X_\alpha)}
 \|r_\alpha\|^{\rm chart}_{q,r,s;o,X_\alpha}
       \le C_{q,s,\kappa}
                     \bigl(\beta r^5+(1+\Lambda)r^3\bigr).
 \label{eq:stationary-taylor-remainder}
\end{equation}
In particular, if
\(\beta\le2\widehat\beta\),
\(\Lambda\le C_\nu\widehat\beta^\nu\) for every \(\nu>0\),
and \(r=w=\widehat\beta^{-7/16}\), the right side is
\begin{equation}
 C_\nu\widehat\beta^{-19/16+\nu}.
 \label{eq:stationary-taylor-exponent}
\end{equation}
These bounds have no factor depending on the
physical volume.
\end{corollary}
\begin{proof}
Keep exactly the coefficient labels of the preceding
proposition but remove the classical terms of degrees
at most four, and the other terms of degrees at most two.
No identity, root, or footprint is changed.
The resulting classical labels have degree at least
five; the other labels have degree at least three.

For explicit control of constants, choose a fixed
analytic radius \(r_*>0\) at which all the rooted
majorants are finite.
For every \(k\ge m\) and \(2r\le r_*\),
\[
 (2r)^k\le(2r/r_*)^m r_*^k.
 \]
In the proof of
\eqref{eq:stationary-chart-family-bound}, the
polynomial degree factors cost the radius \(2r\),
and the spatial polynomial factors cost part of
the exponential reserve.
Apply the displayed inequality with \(m=5\) to
the classical coefficients, and \(m=3\) to the
loop--Haar and determinant coefficients.
Their remaining majorants at \(r_*\) are fixed
local constants, with the explicit multiplier
\(\Lambda\) for the loop part.
This proves \eqref{eq:stationary-taylor-remainder}.

At the stated width,
\[
 \beta w^5\le2\widehat\beta^{-19/16},\qquad
 (1+\Lambda)w^3
       \le C_\nu\widehat\beta^{-21/16+\nu}.
 \]
The second exponent gives a stronger decay than
the first. This proves
\eqref{eq:stationary-taylor-exponent}.
\end{proof}

Since \(19/16-23/20=3/80\), choosing
\(\nu<3/80\) gives a reserve beyond the target
exponent \(23/20\) for this stationary contribution.
This is a remainder estimate after Taylor truncation.
Identifying the retained degree-at-most-four
classical polynomial with the prescribed coarse
principal action is an additional matching problem.
Likewise the quadratic Gaussian and loop terms must
be extracted with the gauge invariant procedure of
Section~\ref{sec:coupling-reset}.
The corollary does not assert either matching identity
merely from the existence of the Taylor polynomials.

\subsection{A finite-box gauge pullback and the physical derivative norm}

There is a useful local bridge between the coordinate
calculation and physical good fields.
The price of this bridge depends on the size of the
local box, which we retain explicitly.
The axial-gauge path mechanism is the one compared above with
\cite[Sect.~B, Lemma~1 and eqs.~(1.24)--(1.26),
printed pp.~79--80]{Bal85g}; the quantitative pullback below
is derived for our chosen paths and physical derivative norm.

Let \(\mathcal Q\) be a nonwrapping coarse box of radius
\(R\ge1\), with root \(o\). For each vertex \(z\) choose
the coordinate-ordered path \(T_z\) from \(o\) to \(z\)
inside \(\mathcal Q\). For every link
\(e=(z,z+e_i)\) inside this box set
\begin{equation}
 W_e(V)=V(T_z)V_eV(T_{z+e_i})^{-1},\qquad
 \kappa_e(V)=\log W_e(V)
 \label{eq:local-axial-pullback}
\end{equation}
whenever the logarithm lies in the fixed small group
chart. This operation is a local gauge change.
The read footprint contains the box and a fixed
plaquette halo.

\begin{lemma}[A quantitative local gauge chart]
\label{lem:local-gauge-chart}
If all plaquettes in the box have chord at most \(h\),
then
\begin{equation}
 \max_e|\kappa_e(V)|\le C_G(1+R)^2h
 \label{eq:local-gauge-chart-value}
\end{equation}
provided the right side is sufficiently small.
For ordered physical left derivatives of order
\(1\le j\le q\), with the allowed directions in the
activity norm, one has
\begin{equation}
 \max_e
 \bigl|D_{A_j}\cdots D_{A_1}\kappa_e(V)\bigr|
       \le C_{G,q,s}
                     \bigl[h(1+R)^{s+1}\bigr]^j .
 \label{eq:local-gauge-chart-derivatives}
\end{equation}
These bounds are independent of the surrounding
physical period.
Under a physical gauge transformation \(u\),
\begin{equation}
 \kappa_e(u\cdot V)
       =\operatorname{Ad}_{u_o}\kappa_e(V).
 \label{eq:local-chart-covariance}
\end{equation}
\end{lemma}
\begin{proof}
The loop in \eqref{eq:local-axial-pullback} consists
of two coordinate paths and one link.
Moving that link past the later coordinate segments
fills it by at most \(C(1+R)^2\) squares in the box.
Lemma~\ref{lem:box-path-comparison} bounds its chord
by \(C(1+R)^2h\).
The logarithm is Lipschitz on a fixed smaller group
chart. This proves
\eqref{eq:local-gauge-chart-value}.

The word defining \(W_e\) has length at most \(C(1+R)\).
An ordered left derivative inserts one of its direction
matrices at one of those occurrences; differentiating
an inverse uses
\(D(U^{-1})=-U^{-1}(DU)U^{-1}\).
At order \(j\) there are at most a constant depending
on \(j\) times \((1+R)^j\) such terms.
All undifferentiated word factors are unitary.
On the read paths the allowed directions have norm
at most \(C_s h(1+R)^s\).
The word's \(j\)-th derivative is consequently bounded
by the right side of
\eqref{eq:local-gauge-chart-derivatives}.
The derivatives of the fixed-chart logarithm are
bounded through order \(q\).
The iterated chain rule preserves the total derivative
order \(j\), hence preserves the same factor
\([h(1+R)^{s+1}]^j\), up to \(C_{G,q,s}\).

A path transport transforms as
\(V(T_z)\mapsto u_oV(T_z)u_z^{-1}\).
Its endpoint factors cancel in \(W_e\), leaving
\(W_e\mapsto u_oW_eu_o^{-1}\).
The small logarithm commutes with conjugation.
This proves \eqref{eq:local-chart-covariance}.
\end{proof}

To use this map for a scalar function, the scalar must
be invariant under simultaneous conjugation of all its
link coordinates. This condition can be retained in
the expansions above.
Restrict to a smaller ball in the maximum link Lie algebra
norm, so conjugated fields and their stationary points
remain in the uniqueness domains.
The action satisfies
\[
 \varphi(\operatorname{Ad}_u x,\operatorname{Ad}_u y)
                            =\varphi(x,y).
 \]
Uniqueness implies
\(\rho(\operatorname{Ad}_u x)=\operatorname{Ad}_u\rho(x)\).
The Hessian transforms by orthogonal conjugation on
the fine colour indices.
Thus \(\Psi,\psi_*\) and the determinant difference are
invariant under simultaneous conjugation.
In their Taylor expansions, group together all colour
indices for each fixed list of link indices.
The adjoint action mixes colour indices at a link but
does not mix links. Uniqueness of Taylor coefficients
therefore makes these grouped scalar polynomials
invariant. For a determinant chain, sum its internal
colour indices as well; the adjacent orthogonal
matrices cancel. Grouping does not increase an
absolute coefficient bound and does not enlarge a
read footprint.

\begin{proposition}[Physical pullback of a bounded stencil]
\label{prop:finite-box-physical-pullback}
Let \(f\) be a scalar chart function.
Assume it is invariant under simultaneous conjugation
and reads links inside the nonwrapping box \(\mathcal Q\).
Suppose it is defined on a neighbourhood of the
closed maximum-link \(2w\)-ball, with \(0<w\le1\) and
the \(4w\)-ball contained in the fixed exponential-coordinate
chart. Suppose also that
\begin{equation}
 C_{G,q,s}h(1+R)^{s+1}\le w .
 \label{eq:local-gauge-price-condition}
\end{equation}
Increase this fixed constant to pay also
\eqref{eq:local-gauge-chart-value}.
Then \(f\circ\kappa\) has a gauge invariant extension
\(\widetilde f\) as a raw local function with that
box and its fixed halo as footprint \(X\).
On the physical good set of \(X\),
\(\widetilde f=f\circ\kappa\), and
\begin{equation}
 \|\widetilde f\|_{q,h,s;o,X}
          \le C_{G,q,s}
                  \|f\|^{\rm chart}_{q,w,s;o,X}.
 \label{eq:finite-box-physical-pullback}
\end{equation}
If \(R\) is bounded by a fixed power of
\(1+\log\widehat\beta\),
\(h=\widehat\beta^{-1/2}P(\widehat\beta)\), and
\(w=\widehat\beta^{-7/16}\),
condition \eqref{eq:local-gauge-price-condition}
holds eventually, with no dependence on the physical
cutoff or volume.
\end{proposition}
\begin{proof}
Use maximum norms of link Lie algebra vectors; these
norms are invariant under the adjoint action.
Their comparison with the component norms used in the
coefficient estimates costs only a constant depending
on \(G\).
By the preceding lemma the good set maps into the
closed \(w\)-ball, with a reserve inside the domain
of \(f\).
Where all logarithms in \eqref{eq:local-axial-pullback}
exist and have maximum-link norm less than \(2w\),
define \(\widetilde f=f\circ\kappa\); elsewhere set it
equal to zero. One may decrease this defining region
within the fixed logarithm chart without changing a
neighbourhood of the physical good set.
The region and its complement are conjugation invariant.
Equation~\eqref{eq:local-chart-covariance} and the
invariance of \(f\) make this extension gauge invariant.
It is local because its definition reads only the
specified paths and links. It is \(C^q\) on a
neighbourhood of the good set, which is the raw local
regularity used in the activity norm.

For completeness, consider an order-\(j\) chain-rule
term involving \(D^pf\), with
\(j_1+\cdots+j_p=j\), \(j_\nu\ge1\).
The vectors supplied by the chart map have norm at most
\[
 C_{G,q,s}\bigl[h(1+R)^{s+1}\bigr]^{j_\nu}.
 \]
The chart norm bounds \(D^pf\) on these vectors by
its norm times the product of their maximum norms
divided by \(w^p\).
After increasing the constant in
\eqref{eq:local-gauge-price-condition}, the remaining
factor is at most \(C_{G,q,s}w^{j-p}\le C_{G,q,s}\).
There are only finitely many chain-rule partitions
through order \(q\).
The order-zero estimate follows directly by composition.
This proves \eqref{eq:finite-box-physical-pullback}.
No mask is differentiated.

Finally
\[
 \frac{h(1+R)^{s+1}}w
 =\widehat\beta^{-1/16}P(\widehat\beta)(1+R)^{s+1}
                                      \longrightarrow0 .
 \]
This proves the last assertion.
\end{proof}

This proposition proves the physical norm estimate
for each specified bounded stencil. It keeps the
size of that stencil explicit.
If the box including its halo lies in a radius-\(B\)
ball, then \(n_B(X)=1\) and its covering price is
at most \(e^\kappa\).
For the entire unbounded family of tree footprints,
compatible local extensions must retain the connected-path
price. Section~\ref{sec:physical-interacting-extensions}
constructs them without a large filled-box price;
the physical reassembly is proved in
Section~\ref{sec:physical-reassembly}.

\subsection{The scope of this localisation step}

The argument has so far localised the classical stationary
action, the stationary loop and Haar factors, and the
quadratic Gaussian determinant. The inverse Hessian,
stationary graph, and determinant all came from the
same finite-period integral.
The constants were obtained from rooted sums, never
from a uniform bound on an extensive scalar.
The connected footprints explicitly pay the prescribed
covering weight.

Two additional arguments turn these chart families
into output regular activities of the full RG step.
First, the local functions are organised as gauge invariant
families on the physical good sets.
Proposition~\ref{prop:finite-box-physical-pullback} supplies
the local gauge map and the physical derivative norm estimate
for bounded nonwrapping stencils.
Section~\ref{sec:physical-interacting-extensions} supplies
compatible extensions for the entire family, including
long footprints.
The smallness of plaquette holonomies does not place
all links on a periodic torus near \(1\), so this step
cannot be replaced by a single global logarithm.

Second, the interacting factor \(\log Z_D\) from
Section~\ref{sec:gaussian-remainder} is localised by
the connected expansion in
Section~\ref{sec:full-interacting-remainder}.
The determinant series is the quadratic
Gaussian calculation; it does not expand the connected
correlations of the nonquadratic remainder.
That expansion supplies the separate cutoff control
for these correlations. The incoming activities and
the large-field terms in
\eqref{eq:pure-reference-splitting} are retained in
Sections~\ref{sec:finite-range-gaussian}--%
\ref{sec:integrated-regular-derivatives}.
Ba{\l}aban's small-field and cluster constructions
\cite[Sect.~1, eqs.~(1.3)--(1.7), printed pp.~260--261]{Bal87}
and \cite[Sect.~2, eqs.~(2.11)--(2.13), printed p.~14]{Bal88a}
motivate this separation.
We have not used a scalar-model cluster theorem as
a substitute for these gauge-specific tasks.

Accordingly \eqref{eq:stationary-chart-family-bound}
is a proved chart localisation bound with its explicit
dependence on \(\beta,\Lambda,r\).
Corollary~\ref{cor:stationary-taylor-remainder}
gives the small stationary Taylor remainder.
To identify it with the required physical output
remainder, the principal quadratic through quartic
terms must be matched and extracted, the local
extensions must be reassembled, and the interacting
remainder must be included.
The budget for transported input activities likewise
remains separate, as explained in
Section~\ref{sec:coupling-reset}.

Sections~\ref{sec:full-interacting-remainder}--%
\ref{sec:physical-interacting-extensions} subsequently
construct the full connected interacting remainder,
covariant stationary kernels on physical good backgrounds,
and compatible local interacting families.
The Gaussian-coordinate failure sectors and their
reassembly with incoming activities belong to the later
full physical RG step.
\section{Microscopic rows and a good Gaussian centre}
\label{sec:microscopic-centre}

The lower bound needed in Section~\ref{sec:actual-coercivity}
is a value estimate on auxiliary volumes. For this purpose we
can use Gaussian completion and a common small fluctuation box.
We will estimate the primitive microscopic factors before any
connected-component factorisation. This avoids using the
full polymer expansion merely to prove that one positive
column is nonzero.

The Gaussian organisation follows the fluctuation calculation
in \cite[Sect.~2, pp.~266--268]{Bal87}. All actions, branch
likelihoods, Haar densities and noise normalisers below are
those already specified in these notes.

\subsection{The linear data required at one RG step}

We first specify the additional linear RG data used in this
part. Summability of the principal paths by itself is not
a substitute for these data.

Let \(\mathcal C\) be the transported curvature of
\eqref{eq:transported-curvature-slot}. On the regular branch,
collect the following Lie algebra valued rows into a vector
\(\mathcal F(x,y)\):
the direct rows \(\sqrt\lambda\,\mathcal C(P_p(U(x,y)))\),
the row sums of the principal fine action, and the noise rows
\(\sqrt a\,\mathcal C(V_bQ(U(x,y))_b^{-1})\).
The finite symmetry weights are included by multiplying
their rows by the square roots of their weights.
Equivalently one may combine the identical direct Wilson
contributions first. We allow either convention and let
\(\lambda_*>0\) be a fixed lower price for one direct
curvature row at each fine plaquette.
For the first convention one can take
\(\lambda_*=\lambda/|\mathsf S|\), where
\(\mathsf S\) is the fixed finite symmetry group.
There is likewise a fixed positive direct noise-row price
\(a_*\). Both depend only on the chosen row convention
and the fixed step parameters.

All row and coordinate spaces use their physical Lie
algebra norms. Euclidean sums use counting measure,
with the displayed symmetry weights already in the rows.
Then the exact principal phase has the form
\begin{equation}
 \varphi(x,y)=\tfrac12\|\mathcal F(x,y)\|_2^2
                                      +\mathcal R(x,y).
 \label{eq:microscopic-square-decomposition}
\end{equation}
The function \(\mathcal R\) consists of the cubic and
quartic curvature polynomials and the direct plaquette
and noise defects
\[
 \tfrac{\lambda}{2}
    \bigl(d_G(W)^2-|\mathcal C(W)|^2\bigr),
 \qquad
 \tfrac a2
    \bigl(d_G(W)^2-|\mathcal C(W)|^2\bigr),
\]
with their actual symmetry weights.

To check the local size of these defects, put \(W=e^X\).
The skew part of \(e^X\) is \(X+O(|X|^3)\); projecting
onto \(\mathfrak g\) keeps its linear term.
Also \(d_G(e^X)^2=|X|^2+O(|X|^4)\).
Consequently the displayed defects vanish through degree
three in \(X\). On a fixed smaller group chart their first
derivatives have size at most \(C|X|^3\).
In particular, for a sufficiently small fixed chord radius,
\begin{equation}
 d_G(W)\le2|\mathcal C(W)|.
 \label{eq:microscopic-curvature-controls-chord}
\end{equation}
These assertions follow from the convergent matrix
exponential series and its termwise derivatives; their
constants depend only on the fixed representation.

Let \(\mathcal C_{\rm c}(V)\) denote the corresponding
vector of direct and row curvatures of the prescribed
coarse principal action. It does not contain a noise row.
Its path budget implies
\begin{equation}
 \|\mathcal C_{\rm c}(V)\|_\infty\le C_{\rm c}h_{\rm c}
                  \quad\hbox{if \(V\) is coarse good}.
 \label{eq:coarse-row-good-bound}
\end{equation}
Here \(h_{\rm c}\) is the prescribed coarse threshold.
For a row sum this follows by the triangle inequality:
every transported slot has norm at most \(h_{\rm c}\),
and its absolute coefficient sum is at most its
exponentially weighted path budget.

Define the actual flat derivative matrices
\[
 \mathsf A=D_y\mathcal F(0,0),\qquad
 \mathsf B=D_x\mathcal F(0,0).
 \]
The flat precision from
Proposition~\ref{prop:actual-flat-gap} is
\(H_0=\mathsf A^T\mathsf A\), since
\(\mathcal R\) has no quadratic term at the identity.

\begin{definition}[Calibrated linear block data]
\label{def:calibrated-linear-data}
At this step the linear RG datum consists of a fine
free-coordinate lift \(\mathcal H\) and a row lift
\(\Phi_0\) from coarse curvature rows to fine and noise
rows, satisfying
\begin{align}
 \mathsf B+\mathsf A\mathcal H
                       &=\Phi_0D\mathcal C_{\rm c}(1),
                                  \label{eq:flat-row-matching}\\
 \mathsf A^T\Phi_0&=0 .             \label{eq:flat-row-orthogonality}
\end{align}
In the first identity \(D\mathcal C_{\rm c}(1)\) means
the derivative of \(x\mapsto\mathcal C_{\rm c}(V(x))\)
at \(x=0\).

The row lift has translation kernels on the integer
cover, with a fixed exponential absolute coefficient
budget. Its covariant version \(\Phi(V)\) is obtained by
parallel transport along specified coordinate paths
before periodisation.
More precisely, if \(\alpha\) includes the output
fine residue and row type, \(\gamma\) is an input
coarse row type, and \(u\in\Z^4\) is the coarse
displacement, the required cover prices are
\[
 \sup_{\alpha\ {\rm fine}}\sum_{\gamma,u}
     |k_{\alpha\gamma}(u)|e^{\sigma_\Phi|u|_1}
                           \le A_\Phi L^{-2},\qquad
 \sup_{\alpha\ {\rm noise}}\sum_{\gamma,u}
     |k_{\alpha\gamma}(u)|e^{\sigma_\Phi|u|_1}
                           \le A_\Phi ,
 \]
for a fixed \(\sigma_\Phi>0\).
There is a price \(A_\Phi\), chosen before the block
side, such that its fine and noise rows obey
\begin{equation}
 \|\Phi_{\rm f}(V)\|_{\infty\to\infty}
       \le A_\Phi L^{-2},\qquad
 \|\Phi_{\rm n}(V)\|_{\infty\to\infty}
       \le A_\Phi .
 \label{eq:row-lift-prices}
\end{equation}
The same bounds hold with absolute coefficients,
before any cancellation or periodisation.
When the block side is to be chosen later, these
cover prices must hold uniformly in the candidate
block sides.
\end{definition}

These are conditions on the intermediate linear RG
data, not additional assumptions on the gauge group or
on the Wilson measures in the construction theorem.
Their simultaneous construction at all scales is proved in
Sections~\ref{sec:all-scale-linear}--%
\ref{sec:all-scale-calibration}.
The one-step estimates in this section state the
input explicitly. They do not infer
\eqref{eq:flat-row-matching} or the pre-block price in
\eqref{eq:row-lift-prices} from the path budget alone.

The two exact identities have a direct Gaussian meaning.
At quadratic order, the fine-variable minimiser is
\begin{equation}
 \mathcal Hx=-H_0^{-1}\mathsf A^T\mathsf Bx.
 \label{eq:flat-lift-from-completion}
\end{equation}
Indeed multiply \eqref{eq:flat-row-matching} by
\(\mathsf A^T\) and use
\eqref{eq:flat-row-orthogonality}.
The weighted inverse bound for \(H_0\) and the
local derivative bounds for \(\mathsf A,\mathsf B\)
then give a weighted row bound for \(\mathcal H\),
uniform in period and allowed to depend on \(L\).
The residual row
\(\mathsf Bx+\mathsf A\mathcal Hx\) is orthogonal to
the range of \(\mathsf A\). Thus, for every fine \(y\),
\begin{equation}
 \tfrac12\|\mathsf Bx+\mathsf Ay\|_2^2
 =\tfrac12\|\Phi_0D\mathcal C_{\rm c}(1)x\|_2^2
       +\tfrac12\|\mathsf A(y-\mathcal Hx)\|_2^2.
 \label{eq:flat-gaussian-block-diagonalisation}
\end{equation}
This is an exact finite-matrix identity, not a
stationary-phase approximation.

We also explain how the coefficient prices give
physical, volume-independent operator bounds.
A typical term of \(\Phi(V)\) is
\[
 k_{\alpha\gamma}(u)
       \operatorname{Ad}_{V(\gamma_{o,o+u})}c_{o+u,\gamma}.
 \]
The adjoint transport has norm one.
For a fixed output row, sum absolute coefficients over
the input row types and over the cover displacement \(u\).
Each displacement has exactly one endpoint after
periodisation. The sum is therefore bounded by the
same cover sum on every period.
For columns sum also over the finite fine residues in
one block; this introduces a fixed \(L\)-dependent
multiplicity, never a physical-volume factor.
On \(\|x\|_\infty\le w\), telescoping a path of length
at most \(C_L(1+|u|)\) gives
\[
 \|\operatorname{Ad}_{V(\gamma)}-1\|
                              \le C_L(1+|u|)w .
 \]
Its polynomial factor is absorbed by the exponential
coefficient reserve. Hence, for some fixed
\(\theta>0\),
\begin{equation}
 \|\Phi(V(x))-\Phi_0\|_\theta\le C_Lw .
 \label{eq:covariant-row-lift-perturbation}
\end{equation}
The norm is the row and column norm of
\eqref{eq:weighted-schur-norm}; rectangular versions
use the attached input and output sites.
The same word estimates give the fixed-order local
derivative bounds used below.

\subsection{A Gaussian centre adapted to a good coarse field}

Work first at a fixed good coarse field \(V=V(x)\),
\(\|x\|_\infty\le w\).
Define its guide rows by
\[
\mathbf g(V)=\Phi(V)\mathcal C_{\rm c}(V).
 \]
We call \(\mathbf g(V)\) the guide row vector.
Equations
\eqref{eq:coarse-row-good-bound} and
\eqref{eq:row-lift-prices} give
\begin{equation}
 \|\mathbf g_{\rm f}(V)\|_\infty
       \le A L^{-2}h_{\rm c},\qquad
 \|\mathbf g_{\rm n}(V)\|_\infty\le A h_{\rm c},
 \qquad A=C_{\rm c}A_\Phi .
 \label{eq:good-guide-prices}
\end{equation}
In particular \(A\) is a pre-block price.

Put
\[
 \mathsf A(V)=D_y\mathcal F(x,0),\quad
 b(V)=\mathcal F(x,0)-\mathbf g(V),\quad
 H(V)=\mathsf A(V)^T\mathsf A(V).
 \]
Thus \(\mathsf A(1)=\mathsf A\); the argument \(V\)
distinguishes the replica derivative from its flat value.
Define the affine Gaussian centre
\begin{equation}
 m(V)=-H(V)^{-1}\mathsf A(V)^Tb(V).
 \label{eq:microscopic-affine-centre}
\end{equation}
The inverse exists on the small coarse chart by the
following estimates.

\begin{lemma}[Bounds for the affine rows]
\label{lem:affine-row-bounds}
For all sufficiently small \(w\), uniformly in the
admissible period,
\begin{align}
 H(V)&\ge\tfrac12\eta_0\,1,
       &\|H(V)^{-1}\|_{\infty\to\infty}&\le C_L,
                              \label{eq:affine-precision}\\
 \|\mathcal F(x,0)\|_\infty&\le C_Lh_{\rm c},
       &\|m(V)\|_\infty&\le C_Lh_{\rm c},
                              \label{eq:affine-centre-size}\\
 \|b(V)+\mathsf A(V)m(V)\|_\infty&\le C_Lw^2 .
                              \label{eq:affine-row-residual}
\end{align}
The constants may depend on the fixed group, block,
kernel, path and linear-data prices, but not on the
number of blocks or the ultraviolet cutoff.
\end{lemma}
\begin{proof}
Local derivative estimates give
\(\|\mathsf A(V)-\mathsf A\|_\theta\le C_Lw\).
They can also be checked directly on the individual
rows: two derivatives of a transported slot cost a
fixed polynomial in its word length, summable by the
path budget. The predictor rows have a fixed patch
and the regular-branch derivative bounds.
Integrating those second derivatives along the
coarse segment from \(0\) to \(x\) gives the stated
perturbation estimate, including the column sums.
Consequently \(H(V)-H_0\) has weighted row and column
norm at most \(C_Lw\).
The Schur bound and the flat gap give the first
inequality in \eqref{eq:affine-precision}.
The inverse estimate
of Lemma~\ref{lem:actual-combes-thomas} gives the second.

At \(y=0\) the fine connection is the replica of \(V\).
Each of its plaquettes is either the identity or a
coarse plaquette. Indeed, a fine square crossing
neither or only one coarse boundary has cancelling
replica factors; a square crossing both boundaries
has exactly the four factors of the corresponding
coarse plaquette.
Its transported curvature rows are therefore
\(O(h_{\rm c})\).
The predicted link differs from the root-to-root
replica transport \(V_b\) by at most \(C_Lh_{\rm c}\):
use the frame-score and path comparison in
Proposition~\ref{prop:nonlinear-chart-entry}, before
its final triangle inequality with \(V_b\).
This bounds the noise rows as well.
It proves the first bound in
\eqref{eq:affine-centre-size}.
The guide bounds, the row and column bounds on
\(\mathsf A(V)\), and \eqref{eq:affine-precision} then give
the bound on \(m\).

For the residual use the trial fine direction
\(v=\mathcal Hx\).
Its maximum norm is at most \(C_Lw\).
Taylor expansion of the actual rows and of the guide
at the flat configuration gives
\[
 \mathcal F(x,0)=\mathsf Bx+O_\infty(w^2),\qquad
 \mathsf A(V)=\mathsf A+O_\theta(w),\qquad
 \mathbf g(V)=
       \Phi_0D\mathcal C_{\rm c}(1)x+O_\infty(w^2).
 \]
All remainders have uniform rooted derivative prices,
so the notation \(O_\infty\) contains no volume factor.
Equation~\eqref{eq:flat-row-matching} now implies
\[
 \|b(V)+\mathsf A(V)v\|_\infty\le C_Lw^2 .
 \]
Let
\[
 P(V)=1-\mathsf A(V)H(V)^{-1}\mathsf A(V)^T .
 \]
Its maximum-norm operator bound is uniform, and
\(P(V)\mathsf A(V)=0\).
From \eqref{eq:microscopic-affine-centre},
\[
 b(V)+\mathsf A(V)m(V)
       =P(V)b(V)=P(V)(b(V)+\mathsf A(V)v).
 \]
This proves \eqref{eq:affine-row-residual}.
\end{proof}

\begin{proposition}[Fine curvature and stationary residual at the centre]
\label{prop:good-affine-centre}
Suppose \(h_{\rm c}\le2h\), \(h/w\to0\) and \(w^2/h\to0\).
Choose \(L\) sufficiently large in terms of the pre-block
price \(A\) and the direct row price \(\lambda_*\).
Choose a sufficiently large fixed mismatch width \(W\).
Then eventually the configuration \(U(x,m(V))\) satisfies
\begin{equation}
 d_G(P_p(U(x,m(V))))\le h/8
             \quad\hbox{for every fine plaquette }p,
 \qquad
 d_G(V_bQ(U(x,m(V)))_b^{-1})\le Wh/2 .
 \label{eq:good-centre-cutoff-reserve}
\end{equation}
Moreover
\begin{equation}
 \|\varphi_y(x,m(V))\|_\infty\le C_Lw^2,\qquad
 \|m(V)-\rho(x)\|_\infty\le C_Lw^2 .
 \label{eq:good-centre-stationary-residual}
\end{equation}
\end{proposition}
\begin{proof}
Taylor expansion in \(y\) and
\eqref{eq:affine-centre-size} give
\[
 \mathcal F(x,m)=\mathbf g(V)+O_\infty(w^2+h_{\rm c}^2)
               =\mathbf g(V)+O_\infty(w^2).
 \]
The last equality uses \(h_{\rm c}\le2h=o(w)\).
Before estimating the direct rows by
\eqref{eq:microscopic-curvature-controls-chord}, their
plaquette holonomies must be in its small chord domain.
This follows independently from the replica plaquette
bound and the Lipschitz change caused by
\(\|m\|_\infty\le C_Lh_{\rm c}\).
The same argument puts the noise mismatch in that
fixed small group domain.
The direct fine and noise rows therefore imply
\begin{align*}
 \sup_p d_G(P_p(U(x,m)))
       &\le \frac2{\sqrt{\lambda_*}}
                     (A L^{-2}h_{\rm c}+C_Lw^2),\\
 \sup_b d_G(V_bQ(U(x,m))_b^{-1})
       &\le \frac2{\sqrt{a_*}}(Ah_{\rm c}+C_Lw^2).
 \end{align*}
For example choose \(L\) so that
\(4A/(\sqrt{\lambda_*}L^2)\le1/16\).
The remaining fine term is at most \(h/16\)
eventually because \(w^2/h\to0\).
Choose \(W>16A/\sqrt{a_*}+1\).
The noise remainder is eventually small enough
to give the stated noise reserve.
This proves \eqref{eq:good-centre-cutoff-reserve}.

For the gradient use the exact square decomposition.
At the centre, all fine curvatures and noise
mismatches have size \(O(h_{\rm c}+w^2)=O(h)\).
Termwise differentiation of the cubic and quartic
path polynomials therefore gives
\(\|\mathcal R_y(x,m)\|_\infty\le C_Lh^2\).
The quartic plaquette and noise defects give a
stronger \(O(h^3)\) bound.
The sums of all these derivatives converge by the
path budgets and their polynomial length reserves.

The derivative of \(\mathcal F\) at \(m\) is
\(\mathsf A(V)+O_\theta(h_{\rm c})\).
Also \eqref{eq:covariant-row-lift-perturbation} and
flat orthogonality give
\[
 \|\mathsf A(V)^T\mathbf g(V)\|_\infty
 =\|\mathsf A(V)^T\Phi(V)\mathcal C_{\rm c}(V)\|_\infty
                                      \le C_Lw h_{\rm c}.
 \]
Indeed the constant product
\(\mathsf A^T\Phi_0\) is zero; every remaining term
has one \(O(w)\) matrix factor.
Combining these estimates in
\[
 \varphi_y(x,m)=
       \mathcal F_y(x,m)^T\mathcal F(x,m)
                                      +\mathcal R_y(x,m)
 \]
gives the first assertion of
\eqref{eq:good-centre-stationary-residual}.

Both \(m\) and \(\rho(x)\) lie in the common
stationary uniqueness chart eventually.
Their gradient difference equals
\(\overline H(m-\rho(x))\), where
\[
 \overline H=
 \int_0^1\varphi_{yy}(x,\rho(x)+t(m-\rho(x)))\,\dd t .
 \]
This averaged matrix is symmetric, has the same
uniform positive gap, and has a uniform weighted
row and column bound.
Lemma~\ref{lem:actual-combes-thomas} applies to it.
Multiplying by its maximum-norm inverse proves
the second assertion.
\end{proof}

The fine-curvature improvement used the factor \(L^{-2}\)
before taking the large-coupling limit.
The stationary residual used the exact flat row
matching and orthogonality identities.
Neither conclusion follows from the rough estimate
\(\|\rho(x)\|_\infty\le C_\rho w\) alone.
The constants denoted by \(C_L\) were only chosen after
the fixed block and linear data; they need not be
uniform as \(L\) varies.

\section{Microscopic activities and the intact column}
\label{sec:microscopic-intact-column}

\subsection{The actual partition of unity}

Fix a smooth function \(\chi:\R\to[0,1]\), equal to \(1\)
on \((-\infty,1]\) and to \(0\) on \([4,\infty)\).
For a positive denominator \(d\) put
\[
 \chi_d(r)=\chi((r/d)^2).
 \]
Thus \(\chi_d(r)=1\) for \(0\le r\le d\);
if it is nonzero then \(r<2d\).
For a good coarse field define the intact weight
\begin{equation}
 \alpha_{0,t}(U,V)=
 \prod_p\chi_{h/4}(d_G(P_p(U)))
 \prod_b\chi_{Wh}(d_G(V_bQ_t(U)_b^{-1})).
 \label{eq:actual-intact-weight}
\end{equation}
The first product is over fine plaquettes, the second
over coarse links. In particular
\begin{equation}
 0\le\alpha_{0,t}\le1,\qquad
 \alpha_{0,t}\ne0\ \Longrightarrow\
                  d_G(P_p(U))<h/2\quad\hbox{for all }p.
 \label{eq:intact-half-good-support}
\end{equation}

This is the empty choice of failures in a genuine
microscopic partition. To check this, attach a factor
\(\chi_d(r)\) to a test declared intact and its
complement \(1-\chi_d(r)\) to a test declared failed.
Summing the products over all subsets of failed tests
gives
\[
 \prod_{\text{tests}}\bigl(\chi_d(r)+(1-\chi_d(r))\bigr)=1.
 \]
There is no modification of the Wilson action,
blocking likelihoods, or physical small-field cutoff.

The restricted contribution to the positive reference
push is
\begin{equation}
 J_{\rm int}(V)=z_{\rm n}(\beta)^{-b_c}
 \sum_t\int
 \alpha_{0,t}(U,V)\lambda_t(U)\chi_{\rm sf}(U)
       e^{-\beta A_t(V,U)-I_f(U)}\,\dd H_{\rm tree}(U).
 \label{eq:restricted-microscopic-column}
\end{equation}
All terms are nonnegative. Hence an integral over a
subregion on which \(\alpha_{0,t}=\chi_{\rm sf}=1\)
gives a lower bound for this same column.

\subsection{One Gaussian box with a reserve}

Use \(m=m(V)\) from the preceding section and let
\[
 K(V)=\varphi_{yy}(x,m),\qquad
 \nu_V=\mu_{K(V)^{-1}},\qquad
 d=|I|\dim_\R\mathfrak g .
 \]
The precision \(K(V)\) is positive with the uniform
chart gap \(\eta\).
For \(z\in\mathfrak g^I\) define
\begin{equation}
 \mathcal Q_D=\{z:\max_{e\in I}|z_e|_{\mathfrak g}\le D\}.
 \label{eq:microscopic-gaussian-box}
\end{equation}
This is a product-shaped set in the physical link
coordinates, not the Euclidean central ball.
The Gaussian measure need not factor over its links.

\begin{lemma}[The intact Gaussian box]
\label{lem:intact-gaussian-box}
Let \(h=\widehat\beta^{-1/2}P(\widehat\beta)\),
\(\widehat\beta/2\le\beta\le2\widehat\beta\),
and \(D=d_0P(\widehat\beta)\).
After fixing \(L,W\) as above, choose \(d_0>0\)
sufficiently small in terms of the fixed step data.
Then eventually, for every \(z\in\mathcal Q_D\),
the connection \(U(x,m+gz)\) lies in the common
regular chart and satisfies
\[
 d_G(P_p(U(x,m+gz)))\le h/4,\qquad
 d_G(V_bQ(U(x,m+gz))_b^{-1})\le Wh .
 \]
Therefore all its intact and fine small-field
factors are \(1\).
\end{lemma}
\begin{proof}
On a fixed smaller chart, changes in plaquette
chords and in the regular prediction map are bounded
by a fixed constant times the maximum change in
the fine coordinates. The plaquette constant
depends only on the group; the prediction constant
may also depend on \(L,J\) and the branch parameters.
Their changes from \(m\) are at most \(C_LgD\).
Since \(gD/h\le\sqrt2\,d_0\), choose \(d_0\) so that
this change is at most \(h/8\) for plaquettes and
\(Wh/2\) for noise mismatches.
Use \eqref{eq:good-centre-cutoff-reserve}.
The whole box is in the fixed coordinate chart
eventually, since \(h\to0\) and
\(\|m+gz\|_\infty\le C_Lh_{\rm c}+gD\).
On this chart every tag selects the same regular
prediction. The tag likelihoods still sum to \(1\).
\end{proof}

Let \(d_G^{\rm Lie}=\dim_\R\mathfrak g\).
The Gaussian box has the explicit probability estimate
\begin{equation}
 \nu_V(\mathcal Q_D)\ge
          1-2d\exp\{-\eta D^2/(2d_G^{\rm Lie})\}.
 \label{eq:microscopic-box-probability}
\end{equation}
Indeed, every real component has variance at most
\(\eta^{-1}\). If a link Lie vector has norm greater
than \(D\), one of its \(d_G^{\rm Lie}\) components
has absolute value greater than
\(D/\sqrt{d_G^{\rm Lie}}\).
Apply the scalar Gaussian tail bound and the union
bound over all \(d\) components.
No independence assertion is used.

\subsection{The microscopic value norm}

To estimate the remaining exponential, hold this
background \(V\) fixed.
A microscopic letter is a real function \(v_\alpha(V,z)\)
with a nonempty coarse footprint \(X_\alpha\) containing
all parents of the fluctuation links it reads.
Define the conditional value norm
\begin{equation}
 \|v\|_{\rm mic,\kappa}=
 \sup_{o\in\Lambda_c}
 \sum_{\alpha:o\in X_\alpha}
 e^{\kappa n_B(X_\alpha)}
       \sup_{z\in\mathcal Q_D}|v_\alpha(V,z)|.
 \label{eq:microscopic-value-incidence-norm}
\end{equation}
We will bound it uniformly over the stated good
backgrounds. The coefficients can depend on that
fixed background through \(m(V)\).
This norm asserts locality in the fluctuation variables;
it does not assert that those coefficients are already
local functions of the physical coarse field.
It is consequently distinct from both the physical
regular norm and the incoming forced-activity norm.
The present lower bound only needs this derivative-order-zero
estimate.

Write the unmasked principal phase as the absolutely
convergent local path sum \(\varphi=\sum_a\varphi_a\).
Expand row squares into their double path sums.
Finite predictor terms are also included among these
local summands. Write
\(\psi=\sum_a\psi_a\), using the loop terms and the
one-link Haar logarithms.
The derivative path budgets established above
justify these sums and their required derivatives.
Set \(\ell=\varphi_y(x,m)\), and introduce the exact
logarithmic letters
\begin{align}
 v_i^{\rm lin}(V,z)&=-g^{-1}\ell_i z_i,
                                      \label{eq:micro-linear-letter}\\
 v_a^{(3)}(V,z)&=
 -\frac g2\int_0^1(1-t)^2
       D_y^3\varphi_a(x,m+tgz)[z,z,z]\,\dd t,
                                      \label{eq:micro-cubic-letter}\\
 v_a^\psi(V,z)&=
            -\bigl(\psi_a(x,m+gz)-\psi_a(x,m)\bigr).
                                      \label{eq:micro-measure-letter}
\end{align}
The index \(i\) in the first line is a real fine
coordinate. Terms that happen to be zero may be retained.
Each third-order or measure letter keeps its actual
path footprint, enlarged by a fixed halo for the
replica coarse links it reads.
The linear coefficients have already been summed at
their output coordinate; their fluctuation footprint
can be chosen as its parent block.

Taylor's integral formula gives the exact identity
\begin{equation}
 \sum_\alpha v_\alpha(V,z)
 =-\beta\bigl(\varphi(x,m+gz)-\varphi(x,m)\bigr)
       +\tfrac12 z^TK(V)z
       -\psi(x,m+gz)+\psi(x,m).
 \label{eq:exact-microscopic-exponent}
\end{equation}
No perturbation series on unbounded Gaussian
coordinates is being substituted for this identity.

\begin{proposition}[Primitive microscopic activity bound]
\label{prop:primitive-microscopic-bound}
Fix \(\kappa>0\) and the one-step data above.
For \(B\) larger than a fixed constant depending on
\(\kappa\) and the path widths, the letters can be
assigned connected coarse footprints so that
\begin{equation}
 \|v\|_{\rm mic,\kappa}\le
 C_L\left(g^{-1}w^2D+gD^3
                         +(1+\Lambda)ghD\right)
                       =:\epsilon_\beta .
 \label{eq:microscopic-letter-price}
\end{equation}
The bound is uniform over good backgrounds
\(\|x\|_\infty\le w\) and over admissible periods.
The loop coefficient budget is \(\Lambda\).
\end{proposition}
\begin{proof}
For the linear letters use
\(\|\ell\|_\infty\le C_Lw^2\).
There are at most a fixed number, depending on \(G,L\),
of fine real coordinates in one coarse block.
This proves their contribution to
\eqref{eq:microscopic-letter-price}.

For a path summand of total word length \(l\) and
coefficient magnitude \(b\), its third derivative has
absolute coordinate sum at most
\(C_Lb(1+l)^{\mathsf M}\), for a fixed integer \({\mathsf M}\).
This is the insertion estimate of
Lemma~\ref{lem:actual-action-local-bounds};
only three derivatives are needed here.
Evaluation on three vectors of maximum norm at
most \(D\) costs \(C_GD^3\).
The integral in \eqref{eq:micro-cubic-letter} costs
at most a fixed constant.
The projected word paths and their fixed halos give
connected coarse footprints with
\[
 n_B(X_\alpha)\le C+\frac{Cl}{B}.
 \]
For a given incidence site, the number of translates
of this summand that can contain it is bounded by a
polynomial in \(1+l\).
Choose \(B\) large enough that the covering weight
costs at most \(C_\kappa e^{\sigma l/4}\).
The polynomial derivative and incidence factors are
absorbed by another fixed fraction of the exponential
path width. Sum the remaining coefficient budget.
A row square uses the product of two row budgets.
A prediction term reads a fixed patch and has a
fixed derivative price. This proves the \(C_LgD^3\)
bound for the third-order letters.

On every segment \(m+tgz\) under consideration,
the fine curvature chords and the coordinate size
are \(O(h)\), by
Lemma~\ref{lem:intact-gaussian-box} and the centre bound.
Differentiating a quadratic loop path term therefore
costs at most its coefficient times
\(C_Lh(1+l)^{\mathsf M}\).
One differentiated curvature supplies a bounded
insertion and the other supplies \(h\);
an insertion in a transport supplies the stronger
factor \(h^2\).
For the Haar term use
\(D\log j_G(y)=O(|y|)=O(h)\) on this fixed group chart.
Integrate these first derivatives along the segment
of length \(g z\), and repeat the same covering
and incidence summation.
The total contribution is
\(C_L(1+\Lambda)ghD\).
All changes of order are dominated by these
nonnegative, absolutely summable path majorants.
\end{proof}

For the associated Mayer letters
\[
 a_\alpha(V,z)=e^{v_\alpha(V,z)}-1
 \]
the same norm is at most \(2\epsilon_\beta\)
when \(\epsilon_\beta\le\log2\).
Indeed \(|e^u-1|\le e^{|u|}|u|\), and the
incidence bound controls each individual \(|u|\).
On a finite coarse torus with \(n_c=c^4\),
\begin{equation}
 \sum_\alpha\sup_{\mathcal Q_D}|v_\alpha|
          \le n_c\epsilon_\beta,\qquad
 \prod_\alpha(1+a_\alpha)
                =\exp\left(\sum_\alpha v_\alpha\right).
 \label{eq:microscopic-absolute-product}
\end{equation}
Both the product and its expansion over finite
subsets of letters converge absolutely and uniformly
on the box. This follows from
\(\sum_\alpha\sup|a_\alpha|\le2n_c\epsilon_\beta\).
These are genuine microscopic factors of the
original fluctuation integral on the specified box.
Their later connected factorisation is unnecessary
for the positive lower bound below.

\subsection{Exact normalisation and a positive lower bound}

Define
\begin{equation}
 \mathcal N_\beta(V)=
 z_{\rm n}(\beta)^{-b_c}\zeta_G^{|I|}
       (2\pi)^{d/2}g^d
 e^{-\beta\varphi(x,m)-\psi(x,m)}
                                 (\det K(V))^{-1/2}.
 \label{eq:microscopic-normalisation}
\end{equation}
It uses the same noise normaliser and Haar chart
constant as \(A_\beta\).
On the intact box all finite tag predictions coincide
and their original likelihoods sum to \(1\).
The exact Gaussian change of variables therefore
gives
\begin{equation}
 J_{\rm int}(V)\ge
 \mathcal N_\beta(V)
 \int_{\mathcal Q_D}
       \exp\left(\sum_\alpha v_\alpha(V,z)\right)
                                            \dd\nu_V(z).
 \label{eq:intact-column-gaussian-lower}
\end{equation}
Using \eqref{eq:microscopic-absolute-product} and
\eqref{eq:microscopic-box-probability}, we obtain
\begin{equation}
 J_{\rm int}(V)\ge
 \mathcal N_\beta(V)e^{-n_c\epsilon_\beta}
 \left(1-2d e^{-\eta D^2/(2d_G^{\rm Lie})}\right).
 \label{eq:restricted-sector-lower-bound}
\end{equation}
In particular the Gaussian probability in parentheses
has not been set equal to \(1\).
The volume factor \(n_c\) has likewise been retained.

\section{The general-background central comparison}
\label{sec:general-background-comparison}

\subsection{The prescribed thresholds and auxiliary sizes}

We now check all the eventual inequalities for the
actual coupling band, without changing either
small-field threshold.
Let the coarse nominal coupling be \(t\), and put
\[
 \widehat\beta=t+b_0\log L,\qquad
 h_{\rm c}=t^{-1/2}P(t),\qquad
 h=\widehat\beta^{-1/2}P(\widehat\beta).
 \]
The block \(L\) is fixed in this one-step argument.
Then \(h_{\rm c}/h\to1\).
For a direct verification, the logarithm of the
threshold is
\[
 f(t)=-\tfrac12\log t+(\log\log t)^2,\qquad
 f'(t)=-\frac1{2t}
          +\frac{2\log\log t}{t\log t}\longrightarrow0 .
 \]
The mean value theorem on the interval of fixed
length \(b_0\log L\) proves the assertion.
In particular \(h_{\rm c}\le2h\) eventually.

With \(w=\widehat\beta^{-7/16}\), one also has
\[
 \frac h w=\widehat\beta^{-1/16}P(\widehat\beta)\to0,
 \qquad
 \frac{w^2}h
       =\widehat\beta^{-3/8}P(\widehat\beta)^{-1}\to0 .
 \]
Thus the hypotheses of
Proposition~\ref{prop:good-affine-centre} hold with
the prescribed scale thresholds.

In the following theorem the coarse auxiliary
volume and the fine coordinate count satisfy
\begin{equation}
 n_c\le C_1(1+\log\widehat\beta)^{m_1},\qquad
 d\le C_{G,L}n_c .
 \label{eq:microscopic-auxiliary-size}
\end{equation}
The second bound is just the finite number of fine
coordinates in one block.
These hypotheses concern the auxiliary tori used
for local estimates. They place no upper bound on
the physical volume of the final construction.

\begin{lemma}[Prices on the coupling band]
\label{lem:microscopic-prices-vanish}
Suppose \(\widehat\beta/2\le\beta\le2\widehat\beta\),
\(D=d_0P(\widehat\beta)\), and for every \(\nu>0\)
the loop budget satisfies
\(\Lambda\le C_\nu\widehat\beta^\nu\) eventually.
Then the microscopic price obeys, for every
\(\nu>0\),
\begin{equation}
 \epsilon_\beta\le C_{\nu,L}
                    \widehat\beta^{-3/8+\nu}.
 \label{eq:microscopic-price-exponent}
\end{equation}
Under \eqref{eq:microscopic-auxiliary-size},
\[
 n_c\epsilon_\beta\longrightarrow0,\qquad
 2d e^{-\eta D^2/(2d_G^{\rm Lie})}\longrightarrow0 .
 \]
All convergences are uniform over the stated
backgrounds and auxiliary periods.
\end{lemma}
\begin{proof}
The coupling band gives \(g\le\sqrt2\widehat\beta^{-1/2}\)
and \(g^{-1}\le\sqrt{2\widehat\beta}\).
The three terms in
\eqref{eq:microscopic-letter-price} are at most fixed
multiples of
\[
 \widehat\beta^{-3/8}P,\qquad
 \widehat\beta^{-1/2}P^3,\qquad
 (1+\Lambda)\widehat\beta^{-1}P^2 .
 \]
Every fixed power of \(P\), and every fixed power of
\(1+\log\widehat\beta\), is subpower in
\(\widehat\beta\).
Distribute the chosen exponent \(\nu\) among these
factors and the loop budget. This proves
\eqref{eq:microscopic-price-exponent} and its
first consequence.
For the Gaussian probability use
\(D^2/\log\widehat\beta\to\infty\).
\end{proof}

\subsection{Comparison with the stationary normalisation}

The normalisation \(A_\beta(x)\) is centred at
\(\rho(x)\), whereas \(\mathcal N_\beta(V)\) is centred
at \(m(V)\). Both are now explicit functions of
the same finite-period integral.

\begin{lemma}[The normalisation price]
\label{lem:microscopic-normalisation-comparison}
For good backgrounds \(\|x\|_\infty\le w\),
\begin{equation}
 \left|\log\frac{\mathcal N_\beta(V(x))}{A_\beta(x)}\right|
 \le C_Ld\bigl(\beta w^4+(1+\Lambda)w^2\bigr)
                         =:\mathcal E_\beta .
 \label{eq:microscopic-normalisation-price}
\end{equation}
Under the coupling band and
\eqref{eq:microscopic-auxiliary-size},
\(\mathcal E_\beta\to0\) uniformly.
\end{lemma}
\begin{proof}
The noise factor, Haar constant, Gaussian factor
\((2\pi)^{d/2}g^d\), and number of free links are
identical in the two normalisations and cancel exactly.
Their logarithmic ratio is
\[
 -\beta\{\varphi(x,m)-\varphi(x,\rho)\}
 -\{\psi(x,m)-\psi(x,\rho)\}
 -\tfrac12\{\log\det K-\log\det {\mathsf H}(x)\}.
 \]
The first difference has no linear term because
\(\varphi_y(x,\rho)=0\).
The uniform Hessian row bound and
\(\|m-\rho\|_\infty\le C_Lw^2\) give
\[
 |\varphi(x,m)-\varphi(x,\rho)|\le C_Ld w^4 .
 \]
The loop and Haar derivatives have per-coordinate
absolute sums at most \(C_L(1+\Lambda)\) on the
fixed common chart. Integration along the same
segment therefore gives
\[
 |\psi(x,m)-\psi(x,\rho)|
                         \le C_Ld(1+\Lambda)w^2 .
 \]
Finally the third-derivative row bounds imply
\(\|K-{\mathsf H}(x)\|_{2\to2}\le C_Lw^2\).
Every matrix on the segment between these two
precisions has eigenvalues at least \(\eta\).
Jacobi's formula, integrated on that segment, gives
\[
 |\log\det K-\log\det {\mathsf H}(x)|
                 \le d\eta^{-1}\|K-{\mathsf H}(x)\|_{2\to2}
                 \le C_Ld w^2 .
 \]
Combining these estimates proves
\eqref{eq:microscopic-normalisation-price}.
On the coupling band,
\(\beta w^4\le2\widehat\beta^{-3/4}\).
The second term is
\((1+\Lambda)\widehat\beta^{-7/8}\).
Their products with the auxiliary coordinate count
tend to zero using any sufficiently small exponent
in the loop-budget bound.
\end{proof}

\begin{theorem}[The restricted-sector lower bound]
\label{thm:actual-restricted-sector-lower}
Fix a principal path profile and calibrated linear
block data as in
Definition~\ref{def:calibrated-linear-data}.
Choose the block side after its pre-block guide price
and the positive direct fine-row price, then choose
the mismatch width \(W\) and the positive constant
\(d_0\) as above.
Assume the prescribed coupling band, thresholds,
loop budget and auxiliary-size bounds.
Then, for all sufficiently large \(\widehat\beta\),
uniformly for every good coarse background in the
specified chart,
\begin{equation}
 J_{\rm int}(V)\ge\tfrac12\mathcal N_\beta(V)
       \ge\tfrac12A_\beta(x)e^{-\mathcal E_\beta}
       \ge\tfrac14A_\beta(x).
 \label{eq:actual-restricted-lower-final}
\end{equation}
All constants and eventual thresholds are independent
of the ultraviolet cutoff and physical period.
\end{theorem}
\begin{proof}
By Lemma~\ref{lem:microscopic-prices-vanish}, increase
the threshold until
\[
 n_c\epsilon_\beta\le\log(4/3),\qquad
 2d e^{-\eta D^2/(2d_G^{\rm Lie})}\le1/4.
 \]
Equation~\eqref{eq:restricted-sector-lower-bound}
then gives
\(J_{\rm int}\ge(3/4)^2\mathcal N_\beta
\ge\mathcal N_\beta/2\).
Apply Lemma~\ref{lem:microscopic-normalisation-comparison}
and increase the threshold until
\(\mathcal E_\beta\le\log2\).
\end{proof}

This proof did not replace the intact column by the
full positive integral. Nor did it use the incoming
forced gas, which equals \(1\) on a good coarse field.
The microscopic logarithmic and Mayer factors were
estimated directly on a single common Gaussian box.
For arbitrary physical volumes the factor
\(n_c\epsilon_\beta\) need not be small; the
connected expansion is then essential.

\subsection{Inclusion of the central cutoff}

Here a general good background means a good coarse
field anywhere in the stated auxiliary chart.
It need not obey the stronger near-identity
restriction \(\|x_0\|_\infty\le c_{\rm in}h\).
The reduction of physical good fields to such
auxiliary charts still belongs to the local gauge
and extension part of the construction; a global
logarithm on an arbitrary physical torus is not
being asserted.

\begin{theorem}[Central inclusion at a general good auxiliary background]
\label{thm:general-good-central-inclusion}
Under Theorem~\ref{thm:actual-restricted-sector-lower},
choose \(d_0\) smaller once more if necessary.
For every good coarse background \(V(x_0)\) with
\(\|x_0\|_\infty\le w\), the actual central set at
\(x_0\) contains a configuration whose every fine
plaquette chord is at most \(h/2\).
Moreover, for all sufficiently large
\(\widehat\beta\), every configuration
\[
 U(x,\rho(x)+gz),\qquad |z|_2\le D,
       \qquad \|x-x_0\|_\infty\le\widehat\beta^{-1},
 \]
is fine good. The nearby coarse field \(V(x)\)
need not itself be good.
All these conclusions have uniform constants on
the prescribed auxiliary volumes.
\end{theorem}
\begin{proof}
The outside-tail theorem gives
\[
 T_\beta(x_0)\le A_\beta(x_0)e^{-cD^2}
                              <\tfrac14A_\beta(x_0)
 \]
eventually.
If the central set at \(x_0\) contained no
\(h/2\)-good configuration, the support condition
\eqref{eq:intact-half-good-support} would make the
intact integrand vanish everywhere on that set.
Outside it the intact weight is at most \(1\), so
positivity would give
\[
 J_{\rm int}(V(x_0))\le T_\beta(x_0).
 \]
This contradicts
\eqref{eq:actual-restricted-lower-final}.
It proves existence of the required configuration.

Apply Lemma~\ref{lem:central-good-reserve} to this
configuration.
Choose \(d_0\) so small that
\(2C_{\rm cut}\sqrt2\,d_0\le1/4\).
Then
\[
 2C_{\rm cut}gD\le h/4.
 \]
For all the stated nearby \(x\), the other cost is
\[
 C_{\rm cut}(1+C_\rho)\|x-x_0\|_\infty
       \le C_{\rm cut}(1+C_\rho)\widehat\beta^{-1}
       \le h/4
 \]
eventually. The last inequality uses
\(\widehat\beta^{-1}/h
=\widehat\beta^{-1/2}/P(\widehat\beta)\to0\).
The transport criterion proves the assertion.
These nearby fields lie in the common fixed chart
because \(w+\widehat\beta^{-1}\to0\).
\end{proof}

The preceding result also has a direct check.
Equation~\eqref{eq:good-centre-stationary-residual}
implies
\[
 \frac{\|m(V(x_0))-\rho(x_0)\|_2}{gD}
 \le \frac{C_L}{d_0}\sqrt d\,
          \frac{\widehat\beta^{-3/8}}{P(\widehat\beta)}
 \longrightarrow0 .
 \]
Thus the explicitly constructed good centre \(m(V(x_0))\)
itself belongs to the Euclidean central ball eventually
on these auxiliary volumes.
The restricted-integral proof above additionally
establishes the quantitative lower bound needed in
the original positive-column formulation.

\begin{corollary}[The actual reference comparison near a good background]
\label{cor:general-good-reference-comparison}
Under the same conditions, in the indicated
\(\widehat\beta^{-1}\)-neighbourhood of each good base
background,
\begin{equation}
 0\le
 (\mathcal P_\beta F_{\rm sf,\beta})(V(x))
                       -A_\beta(x)Z_D(g,x)
       \le A_\beta(x)e^{-cD^2}
 \label{eq:general-good-reference-comparison}
\end{equation}
eventually, with uniform constants.
\end{corollary}
\begin{proof}
Central inclusion makes the fine small-field
indicator equal to \(1\) throughout the central
set. The actual regular predictions agree and the
tag likelihoods sum to \(1\).
The exact Gaussian and Haar change of variables
therefore identifies its contribution with
\(A_\beta(x)Z_D(g,x)\).
The complement is the nonnegative tail.

The outside-tail proof in
Theorem~\ref{thm:actual-auxiliary-tail} applies on
\(\|x\|_\infty\le2w\) with the same conclusion
after a fixed change of its constants:
its stationary price is bounded by
\(C n_f(2w)^2\), its loop and Haar prices have
the same bounds with these fixed factors, and
the remaining positive costs are still
polylogarithmic or tending to zero.
Since \(w+\widehat\beta^{-1}\le2w\) eventually,
this covers the entire neighbourhood in question.
\end{proof}

\subsection{Constant choices and the remaining construction}

The order of choices is part of the estimate.
Fix the group and the two principal path profiles.
The linear block construction must provide
\eqref{eq:flat-row-matching}--\eqref{eq:row-lift-prices}
with their pre-block coefficient prices.
Choose a sufficiently large odd \(L\) to obtain
the fine-curvature gain, then fix \(W\), the
stationary chart constants, and the sufficiently
small positive \(d_0\).
All constants denoted by \(C_L\) are allowed after
these choices. Finally increase the nominal
coupling threshold to pay the finitely many
inequalities proved above.
The ultraviolet cutoff and physical volume occur
after these choices.
The actual path and kernel arrays may vary with a
cutoff or a step, provided the displayed identities
and coefficient budgets hold with the same profile
prices. Every bound above is uniform over that
priced class of arrays.

The incidence bound
\eqref{eq:microscopic-letter-price} is independent
of volume before the auxiliary restriction is used.
Its extensive sum is \(n_c\epsilon_\beta\).
Only the explicitly imposed auxiliary-size bound
makes this sum small and the Gaussian box have
probability close to \(1\).
This does not assert that a single such box captures
the Gaussian field on an unrestricted physical volume.

These sections close the restricted-sector
lower bound and general-good-background central
comparison with the listed linear RG data.
Sections~\ref{sec:all-scale-linear}--%
\ref{sec:all-scale-calibration} construct those data
simultaneously over all finite histories and periods.
Section~\ref{sec:first-interacting-localisation}
begins the connected interacting calculation.
Sections~\ref{sec:full-interacting-remainder}--%
\ref{sec:physical-interacting-extensions} give the
full connected interacting remainder and a separate
construction of physical gauge-invariant local families.
Their physical derivative norm is proved directly;
the fixed-background microscopic value norm is not
used as that derivative norm.
The Gaussian-coordinate failure sectors and the
full reassembly with incoming activities remain part
of the subsequent nonlinear RG step.
\section{Linear RG data from an arbitrary finite history}
\label{sec:all-scale-linear}

We now construct the data of
Definition~\ref{def:calibrated-linear-data}. The point of the
construction is its order of quantifiers. Constants are
chosen for a fixed smoothing index before the block sides,
the number of RG steps, and the periods. A bound obtained
by multiplying a constant once per step would not suffice.

We use cumulative averaging filters and a Gaussian Schur
complement. The use of Fourier analyticity to obtain
exponential kernel bounds is also explained in
\cite[Sect.~3.5, Lemma~3.17 and its contour-shift proof,
preprint p.~32]{DST23}. This is a comparison with the scalar
kernel argument. Here the covariance has a gauge
kernel. We give the cochain calculation explicitly, rather
than applying a scalar Green function estimate to that kernel.
All linear operators below act identically on each colour
component.

\subsection{Cochains, Fourier conventions, and coefficient norms}

An \(r\)-cochain assigns a number to each positively oriented
elementary \(r\)-cell. Thus a \(1\)-cochain is a link field
and a \(2\)-cochain is a plaquette field. Write \(d_r\) for
the oriented coboundary from \(r\)- to \((r+1)\)-cochains.
Opposite faces cancel, so \(d_{r+1}d_r=0\).
The inner product is the sum over sites and cell types,
without division by volume.

For \(p\in\C^4\), put
\[
 q_i(p)=e^{ip_i}-1,\qquad
 \Delta_{\mathsf M}(p)={\mathsf M}^2\sum_{i=1}^4
                  q_i(p/{\mathsf M})q_i(-p/{\mathsf M}).
\]
Here \({\mathsf M}\) is a positive odd integer. On real momenta,
\(\Delta_{\mathsf M}(p)=4{\mathsf M}^2\sum_i\sin^2(p_i/(2{\mathsf M}))\).
Let \(d_r(p)\) denote exterior multiplication by the
vector \(q(p)\). For a matrix symbol \(A\), the notation
\[
 A^\dagger(p)=A(-p)^T
\]
means its analytic transpose; it is the Hermitian adjoint
on real momenta for the real kernels used below.
We shall also write \(A^\dagger\) when its argument is
unambiguous. The symbol of \(d_r\) is \(d_r(p)\).
Write \(d_r^{(L)}(p)=L d_r(p)\).

For a translation kernel \(k_{\alpha\gamma}(u)\),
\(u\in\Z^4\), use the coefficient norm
\begin{equation}
 |k|_\sigma=
 \max_\alpha\sum_{\gamma,u}
       e^{\sigma|u|_1}|k_{\alpha\gamma}(u)|.
 \label{eq:linear-cover-norm}
\end{equation}
The cell-type sets in this formula have bounded cardinality.
For a kernel from a coarse lattice to a fine lattice,
the output index includes one fine residue in a centered
\(L\)-block; take the maximum over this residue as well.
The displacement \(u\) is still measured in coarse units.
In particular, the norm does not sum over all output
residues.

\begin{lemma}[Fourier coefficients and periodisation]
\label{lem:linear-fourier-coefficients}
Suppose a \(2\pi\)-periodic matrix symbol is holomorphic
for \(|\operatorname{Im}p_i|<\delta\) and has norm at most
\(C\) there. For \(0<\sigma<\delta\), its Fourier
coefficients satisfy
\[
 |k|_\sigma\le
 C'\,C\left(\frac{1+e^{-(\delta-\sigma)/2}}
                  {1-e^{-(\delta-\sigma)/2}}\right)^4 ,
\]
where \(C'\) depends only on the cell types. The same
coefficient budget bounds the \(\ell^\infty\) operator
norm after periodisation on every torus. Products of
kernels satisfy \(|k_1k_2|_\sigma\le
|k_1|_\sigma|k_2|_\sigma\).
\end{lemma}
\begin{proof}
Shift the \(i\)-th Fourier contour to imaginary part
\(-\delta'\operatorname{sgn}u_i\), with
\(\delta'=(\delta+\sigma)/2\). This bounds the coefficient
by \(C e^{-\delta'|u|_1}\). Sum the four geometric
series. Absolute convergence justifies the contour
shifts and convolution. In a periodised row, every
cover displacement has one image. Summing absolute
coefficients before identifying images proves the
operator bound, also when images coincide.
\end{proof}
The holomorphic bounds below are first established on an
enlarged cube
\[
 |\operatorname{Re}p_i|<\pi+\delta,\qquad
 |\operatorname{Im}p_i|<\delta.
\]
Identities on overlapping translates will supply
periodicity. We never periodise an expression with an
unremoved pole and then assign a separate value at \(p=0\).

\subsection{Cumulative filters and the geometric history bounds}

Define the scalar line and centered block averages
\[
 r_{\mathsf M}(z)=\frac1{\mathsf M}\sum_{\nu=0}^{{\mathsf M}-1}e^{i\nu z},
 \qquad
 a_{\mathsf M}(p)=\prod_{i=1}^4
       \left[e^{-i({\mathsf M}-1)p_i/2}r_{\mathsf M}(p_i)\right].
\]
The centering is integral because \({\mathsf M}\) is odd. Directly
regrouping the sums gives
\begin{equation}
 r_{L{\mathsf M}}(z)=r_L({\mathsf M}z)r_{\mathsf M}(z),\qquad
 a_{L{\mathsf M}}(p)=a_L({\mathsf M}p)a_{\mathsf M}(p).
 \label{eq:cumulative-average-product}
\end{equation}

A finite history ending at index \(j\ge4\) is defined
recursively. The empty history has mass \({\mathsf M}=1\) and filter
\(s(p)=1\). A step of odd side \(L\ge3\) takes a fine
history with index \(j\) to a coarse history with index
\(t\), where \(t=j\) or \(t=j-1\), by
\begin{equation}
 {\mathsf M}'=L{\mathsf M},\qquad s'(p)=a_L({\mathsf M}p)^j s(p).
 \label{eq:cumulative-filter-recursion}
\end{equation}
Consequently a history with final index \(j\) has
\[
 s(p)=a_{\mathsf M}(p)^j e(p).
\]
The extra factor \(e\) is a product of centered averages
\(a_{{\mathsf M}_\nu}\), one for each decrease of index. Their
masses satisfy \({\mathsf M}_{\nu+1}\le {\mathsf M}_\nu/3\) when listed
from the coarsest end. Repetitions at the same index
produce no extra factor. This is an exact consequence
of \eqref{eq:cumulative-average-product}, not an estimate.

For a centered alias \(\ell\in
\{-({\mathsf M}-1)/2,\ldots,({\mathsf M}-1)/2\}^4\), put
\[
 p_\ell=\frac{p+2\pi\ell}{{\mathsf M}},\qquad
 w_j(\ell)=\prod_{i=1}^4(1+|\ell_i|)^{-j}.
\]
Aliases are the distinct fine momenta corresponding
to a fixed coarse momentum.

\begin{lemma}[Bounds independent of history length]
\label{lem:cumulative-filter-bounds}
For every \(j\ge4\), there are \(\delta_j,c_j,C_j>0\),
independent of the finite history and its mass, such that
on the enlarged complex cube of width \(\delta_j\),
\begin{equation}
 |s(p_\ell)|\le C_jw_j(\ell),\qquad
 |s(p/{\mathsf M})|\ge c_j\quad\hbox{for real }p.
 \label{eq:cumulative-filter-bounds}
\end{equation}
Also,
\begin{equation}
 \sum_\ell w_2(\ell)\le
 \left(\sum_{\nu\in\Z}(1+|\nu|)^{-2}\right)^4<\infty.
 \label{eq:alias-summability}
\end{equation}
\end{lemma}
\begin{proof}
The identity
\[
 r_{\mathsf M}(z)=e^{i({\mathsf M}-1)z/2}
             \frac{\sin({\mathsf M}z/2)}{{\mathsf M}\sin(z/2)}
\]
is interpreted continuously at its removable zeros.
For \(|\operatorname{Re}p_i|\le\pi+\delta\) and small
fixed \(\delta>0\), it gives
\[
 |r_{\mathsf M}((p_i+2\pi\ell_i)/{\mathsf M})|
       \le\frac{C}{1+|\ell_i|}.
\]
For \(\ell_i\ne0\), the denominator is bounded below
by \(c(1+|\ell_i|)\), whereas the numerator is bounded.
For \(\ell_i=0\), use the finite sum, whose norm is at
most \(e^\delta\). On the real principal alias,
the same sine quotient is bounded below by a positive
constant.

An average of centered phases of length \(m\) has norm
at most \(e^{C m|\operatorname{Im}p|_\infty}\).
Thus at \(p_\ell\), the extra factors have total upper
bound
\[
 \exp\!\left(C\delta\sum_\nu {\mathsf M}_\nu/{\mathsf M}\right)\le e^{3C\delta/2}.
\]
On the real principal alias, the sine quotient also
gives
\[
 a_m(p/{\mathsf M})\ge e^{-C m/{\mathsf M}}\quad(1\le m\le {\mathsf M}).
\]
Indeed the quotient is positive on this compact
range, tends to one as \(m/{\mathsf M}\to0\), and its logarithm
is bounded below by \(-C m/{\mathsf M}\). The product of these
lower bounds is at least \(e^{-3C/2}\).
Combining this with the \(j\)-th power of \(a_{\mathsf M}\)
proves \eqref{eq:cumulative-filter-bounds}.
Finally extend the finite alias sum to \(\Z^4\) and
factor the four positive sums.
\end{proof}

\subsection{The cumulative covariance and its noise terms}

Fix the positive Gaussian noise price \(a\).
For a history \(v\), write \({\mathsf M}_v,s_v\) for its mass and
filter and put
\[
 b_{v,i}(p)=s_v(p)r_{{\mathsf M}_v}(p_i).
\]
For a scalar function \(f\), define
\[
 (\mathcal Q_{v,i}f)(p)
  ={\mathsf M}_v^{-2}\sum_\ell
             b_{v,i}(p_\ell)b_{v,i}(-p_\ell)f(p_\ell).
\]
The factor is \({\mathsf M}_v^{-2}\), with a sum over \({\mathsf M}_v^4\)
aliases. This is the four-dimensional counting-measure
normalisation for a link field, including its scaling
dimension one.

Let \(\mathcal P_v\) be the following list of shorter
histories. It is empty for the empty history; after
one step it contains the empty coarse history followed
by every member of the old list with that step appended.
These are precisely the histories of the independent
noise contributions inserted at previous steps. Set,
away from \(p=0\),
\begin{equation}
 D_{v,i}(p)=
  (\mathcal Q_{v,i}\Delta_1^{-1})(p)
   +a^{-1}\sum_{u\in\mathcal P_v}
                         (\mathcal Q_{u,i}1)(p).
 \label{eq:cumulative-covariance}
\end{equation}
These are auxiliary diagonal covariance coefficients.
They are not the inverse of the gauge-degenerate
potential kinetic operator.
Here singular expressions are used only away from
their real poles; their regular combinations are
defined below.

At a step from index \(j\) to \(t\), let
\[
 f(q)=a_L(q)^j,\quad
 h_r(q)=f(q)\,
      \operatorname{diag}_{|I|=r}\prod_{i\in I}r_L(q_i),
 \quad q_\ell=(p+2\pi\ell)/L .
\]
Then
\begin{equation}
 D_{v',i}(p)=
 L^{-2}\sum_\ell
       f(q_\ell)f(-q_\ell)
       r_L(q_{\ell,i})r_L(-q_{\ell,i})D_{v,i}(q_\ell)
       +a^{-1}.
 \label{eq:exact-covariance-recursion}
\end{equation}
To verify this, use
\(b_{v',i}(q)=f({\mathsf M}_vq)r_L({\mathsf M}_vq_i)b_{v,i}(q)\)
and split an \(L{\mathsf M}_v\)-alias into an \(L\)-alias and an
\({\mathsf M}_v\)-alias. The two factors \(L^{-2}\) and
\({\mathsf M}_v^{-2}\) multiply to \((L{\mathsf M}_v)^{-2}\).
The new empty history contributes exactly
\(a^{-1}\).

There is no accumulation proportional to the number of
steps in the noise part. By
Lemma~\ref{lem:cumulative-filter-bounds},
\[
 |\mathcal Q_{u,i}1|\le C_j{\mathsf M}_u^{-2}.
\]
The list of masses contributing to
\(\mathcal P_v\) starts with \(1\) and subsequently
increases by at least a factor \(3\). Hence
\begin{equation}
 \sum_{u\in\mathcal P_v}{\mathsf M}_u^{-2}
       \le\sum_{\nu\ge0}3^{-2\nu}=\frac98 .
 \label{eq:cumulative-noise-geometric}
\end{equation}
All these estimates hold holomorphically on one
common smaller cube, with constants depending on
\(j,a\) only.

\subsection{Removing the pole before taking Fourier coefficients}

Separate the principal alias in
\eqref{eq:cumulative-covariance}:
\[
 D_{v,i}(p)=\frac{L_{v,i}(p)}{\Delta_{{\mathsf M}_v}(p)}
                         +R_{v,i}(p),\qquad
 L_{v,i}(p)=b_{v,i}(p/{\mathsf M}_v)b_{v,i}(-p/{\mathsf M}_v).
\]
The remainder \(R_{v,i}\) is the sum of the nonprincipal
aliases and the noise terms in
\eqref{eq:cumulative-covariance}. Define
\begin{align}
 G_{v,i}&=L_{v,i}+\Delta_{{\mathsf M}_v}R_{v,i},\notag\\
 U_v(p)&=s_v(p/{\mathsf M}_v)s_v(-p/{\mathsf M}_v),\notag\\
 b_v(p)&=U_v(p)^{-1}
  -\sum_i\frac{q_i(p)q_i(-p)R_{v,i}(p)}
                         {L_{v,i}(p)G_{v,i}(p)} .
 \label{eq:filled-metric-functions}
\end{align}
The letter \(b_v\) without an axis index is distinct
from the filter \(b_{v,i}\).

\begin{lemma}[Uniform regularity of the filled symbols]
\label{lem:filled-symbol-bounds}
For fixed \(j\ge4,a>0\), the functions
\(G_{v,i},b_v\) and their inverses are holomorphic and
bounded on a common enlarged cube. Their bounds and
the cube width are independent of history and mass.
They satisfy \(G_{v,i}(0)=b_v(0)=1\).
For real \(p\) in that cube,
\[
 0<c_j\le G_{v,i}(p),b_v(p)\le C_j.
\]
\end{lemma}
\begin{proof}
For a nonzero centered alias \(\ell\), the denominator
\[
 {\mathsf M}_v^2\Delta_1((p+2\pi\ell)/{\mathsf M}_v)
\]
has magnitude at least
\(c(1+|\ell|_\infty)^2\) on a sufficiently thin
common complex cube. On the real cube this is the
sine bound. On a complex cube its real part remains
positive: the change caused by the imaginary parts
is at most \(C\delta(1+|\ell|_\infty)\), which can
be absorbed into the real lower bound.
The numerator is bounded by \(C_jw_j(\ell)\).
Sum with \eqref{eq:alias-summability}.
Together with \eqref{eq:cumulative-noise-geometric}
this bounds \(R_{v,i}\) holomorphically.
The other functions occurring in
\eqref{eq:filled-metric-functions} are likewise bounded.

On real momenta \(R_{v,i}\ge0\). The principal-alias
lower bound gives \(G_{v,i}\ge L_{v,i}\ge c_j>0\).
The elementary line identity
\[
 q_i(p)=r_{{\mathsf M}_v}(p_i/{\mathsf M}_v)\,
                  {\mathsf M}_vq_i(p/{\mathsf M}_v)
\]
gives
\[
 q_i(p)q_i(-p)U_v
   ={\mathsf M}_v^2q_i(p/{\mathsf M}_v)q_i(-p/{\mathsf M}_v)L_{v,i}.
\]
Substitution into \eqref{eq:filled-metric-functions}
and cancellation of \(L_{v,i}\) proves
\begin{equation}
 \Delta_{{\mathsf M}_v}(p)b_v(p)
              =\sum_i\frac{q_i(p)q_i(-p)}{G_{v,i}(p)}.
 \label{eq:filled-laplacian-identity}
\end{equation}
For nonzero real \(p\), both
\(\Delta_{{\mathsf M}_v}(p)\) and \(\sum_i|q_i(p)|^2\)
are comparable to \(|p|^2\), with numerical
constants independent of \({\mathsf M}_v\). The upper and
lower bounds on \(G_{v,i}\) therefore give upper
and lower bounds on \(b_v\). At zero, the defining
regular formula gives \(b_v(0)=1\), and also
\(G_{v,i}(0)=1\).

To pass from real lower bounds to holomorphic
inverses, first use the already obtained bounds on a
slightly larger cube. Cauchy's estimate bounds all
first derivatives on a smaller cube, uniformly in
history. If its width is reduced further by a
fixed factor depending on this derivative bound and
the real lower bound, each function remains within
half its real lower bound. Its reciprocal is then
holomorphic and uniformly bounded there.
\end{proof}

For a list \(z=(z_1,\ldots,z_4)\), write
\(\operatorname{diag}_r(z)\) for the diagonal matrix
on \(r\)-cell types with entry \(\prod_{i\in I}z_i\)
at the type \(I\). Define the unmodified
\(2\)-cochain metric and its inverse by
\begin{equation}
 H_v(p)=b_v(p)^{-1}
       \operatorname{diag}_2(G_{v,1}^{-1},\ldots,G_{v,4}^{-1}),
 \quad
 B_v(p)=b_v(p)
       \operatorname{diag}_2(G_{v,1},\ldots,G_{v,4}).
 \label{eq:all-scale-metric}
\end{equation}
They are inverse, positive Hermitian on real momenta,
and equal to the identity at \(p=0\).
For the empty history they are identically the identity.

These symbols glue periodically. To check this without
assuming periodicity of their separate pole splits,
put
\(\rho_v(p)=\sum_iq_i(p)q_i(-p)/D_{v,i}(p)\).
Away from zero,
\[
 D_{v,i}=G_{v,i}/\Delta_{{\mathsf M}_v},\qquad
 \rho_v=\Delta_{{\mathsf M}_v}^2b_v,\qquad
 H_v=\rho_v^{-1}\operatorname{diag}_2(D_{v,i}^{-1}).
\]
The last expression is periodic by relabelling aliases
in \eqref{eq:cumulative-covariance}. Hence the filled
expressions agree on real overlaps away from zero.
Analytic continuation gives agreement on the whole
overlap, including zero. The reality relation follows
from the real averaging coefficients.
Lemma~\ref{lem:linear-fourier-coefficients} now gives
exponential coefficient bounds for \(H_v,B_v\),
uniform over every finite history with final index \(j\).

\subsection{Positive row factors without changing the kinetic operator}

The prescribed principal action may use a metric
different from \(H_v\) away from exact \(2\)-cochains.
This freedom does not change its quadratic action on
potentials. To describe it, choose a nonnegative
number \(l_j\), one for each index \(j\), before
choosing block lengths. At the empty Wilson history
use the initial coefficient zero, as specified in
Lemma~\ref{lem:packet-complementary-row}. For nonempty
histories use \(l_j\), and set
\begin{equation}
 \widetilde H_v=
 \begin{cases}
 I,&v\text{ is empty},\\
 \left[B_v+l_j(d_2B_v)^\dagger(d_2B_v)\right]^{-1},
                                      &v\text{ is nonempty}.
 \end{cases}
 \label{eq:positive-curvature-extension}
\end{equation}
Since \(B_v\) has a uniform positive lower bound,
the matrix inside the inverse does as well.
Its upper bound depends on \(j,a,l_j\) only.
The empty case is the identity and obeys the same
uniform bounds after enlarging the fixed constants.
The Cauchy argument in the preceding proof gives
a uniform holomorphic inverse. Moreover,
\begin{equation}
 \widetilde H_vd_1=H_vd_1,\qquad
 d_1^\dagger\widetilde H_v=d_1^\dagger H_v .
 \label{eq:metric-extension-on-curvatures}
\end{equation}
Indeed multiply \(H_vd_1\) by the matrix being
inverted. Its first term gives \(d_1\), and its
second term is zero because \(B_vH_v=1\) and
\(d_2d_1=0\). Transposition proves the second identity.
Thus the potential kinetic operator is always
\[
 K_v=d_1^\dagger H_vd_1
                  =d_1^\dagger\widetilde H_vd_1.
 \]
At the Wilson cutoff we may instead set
\(\widetilde H_{\rm top}=1\). This has the same
potential kinetic operator, since \(H_{\rm top}=1\),
and the same uniform estimates after including
the numerical identity-matrix bounds. All formulas
below allow this top-only choice.

Choose \(0<\lambda_j<\frac12\inf_{v,p}
\lambda_{\min}(\widetilde H_v(p))\).
The infimum is over histories at index \(j\) and real
momenta; its positivity was just proved. The matrix
\(\widetilde H_v-\lambda_j1\) has a uniform positive
square root \(R_v\), with exponential coefficients.
Here is a direct verification of the last assertion.
For a matrix \(A\) with spectrum in \([c,C]\), the
contour formula
\[
 A^{1/2}=\frac{1}{2\pi i}
          \int_\Gamma z^{1/2}(z-A)^{-1}\,\dd z
\]
uses one fixed contour surrounding \([c,C]\) in
the right half-plane. On real momenta its resolvent
is uniformly bounded on that contour. On a common
smaller complex strip this remains true by the
resolvent identity and Cauchy's derivative estimate.
The contour formula is therefore holomorphic and
bounded on that strip. It gives a real symmetric
square root on every real fibre; Fourier inversion
proves the asserted coefficient bounds.

Set
\begin{equation}
 S_v=\begin{pmatrix}\sqrt{\lambda_j}\,1\\R_v\end{pmatrix},
 \qquad S_v^TS_v=\widetilde H_v .
 \label{eq:all-scale-row-factor}
\end{equation}
Here and below a superscript \(T\) on an installed
real operator is its counting-measure transpose.
All the displayed identities hold on every period,
by the same Fourier coefficients. The bounds for
\(S_v,\widetilde H_v^{\pm1},H_v^{\pm1}\) depend only on
\(j,a,l_j,\lambda_j\), not on depth or period.
Finite symmetry copies of the rows, with square-root
weights, give the same identity and change these
prices only by fixed type multiplicities.

\section{The two linear columns and exact Gaussian iteration}
\label{sec:linear-columns}

\subsection{Columns before their apparent poles are removed}

Fix a step \(v\longmapsto v'\) of side \(L\), from
index \(j\) to \(t\in\{j,j-1\}\), both at least \(4\).
All symbols in this subsection are first evaluated at
nonzero real \(p\). Put
\[
 q_\ell=(p+2\pi\ell)/L,\quad
 F_\ell=\operatorname{diag}_1(L^{-2}D_{v,i}(q_\ell)),
 \quad D=\operatorname{diag}_1(D_{v',i}(p)).
\]
These are diagonal \(1\)-cochain matrices. Denote the
fine and coarse unmodified \(2\)-cochain metrics by
\(H_f=H_v(q_\ell)\) and \(H_c=H_{v'}(p)\).
For brevity write
\[
 \mathbf q_\ell=d_0^{(L)}(q_\ell),\quad
 \mathbf q=d_0(p),\quad
 c_\ell=d_1^{(L)}(q_\ell),\quad c=d_1(p),
 \quad K=c^\dagger H_c c.
\]
Thus \(\mathbf q_\ell,\mathbf q\) are gradient
columns and \(c_\ell,c\) are curl matrices.
At an alias a dagger reverses the full momentum:
\(c_\ell^\dagger=d_1^{(L)}(-q_\ell)^T\);
it does not keep \(\ell\) fixed and replace only \(p\)
by \(-p\). Define the scalars and row
\[
 \rho_\ell=\mathbf q_\ell^\dagger F_\ell^{-1}\mathbf q_\ell,
 \quad \rho=\mathbf q^\dagger D^{-1}\mathbf q,\quad
 \pi=\rho^{-1}\mathbf q^\dagger D^{-1},
 \quad Z=\sum_\ell f(q_\ell)f(-q_\ell)
                                      \operatorname{tr}F_\ell .
\]
The denominator \(Z\) is strictly positive for real
nonzero \(p\), since its principal summand is positive.
For each alias set
\[
 u_\ell=\frac{f(-q_\ell)\operatorname{tr}F_\ell}{Z},
 \qquad
 y_\ell=f(q_\ell)f(-q_\ell)\rho/\rho_\ell .
\]
Here \(f,h_r\) are the one-step filters defined above.
The rescaled columns are
\begin{align}
 \mathcal E_\ell&=F_\ell h_1(q_\ell)^\dagger K
                                  +u_\ell\mathbf q_\ell\pi,
                         \label{eq:linear-E-column}\\
 \mathcal T_\ell&=
 H_f c_\ell F_\ell h_1(q_\ell)^\dagger c^\dagger
 +y_\ell h_2(q_\ell)^\dagger
                     (1-H_c cDc^\dagger).
                         \label{eq:linear-T-column}
\end{align}
The second terms specify the extensions on gradient
and nonexact curvature directions. Neither \(\pi\)
nor the individual terms in
\eqref{eq:linear-T-column} should be assigned Fourier
coefficients separately at the pole.

The elementary exterior algebra identity
\begin{equation}
 c_\ell^\dagger H_fc_\ell
  =F_\ell^{-1}
    -F_\ell^{-1}\mathbf q_\ell\rho_\ell^{-1}
                           \mathbf q_\ell^\dagger F_\ell^{-1}
 \label{eq:diagonal-Maxwell-identity}
\end{equation}
follows by expanding a diagonal \(2\)-form metric:
its diagonal and off-diagonal entries are, respectively,
\[
 \sum_{k\ne i}\frac{|\mathbf q_{\ell,k}|^2}
                    {\rho_\ell F_{\ell,i}F_{\ell,k}},
 \qquad
 -\frac{\mathbf q_{\ell,i}\overline{\mathbf q_{\ell,k}}}
                    {\rho_\ell F_{\ell,i}F_{\ell,k}}.
\]
Use
\(\rho_\ell=\sum_k|\mathbf q_{\ell,k}|^2/F_{\ell,k}\).
The same identity holds on the coarse side, with
\(F_\ell,\mathbf q_\ell,\rho_\ell\) replaced by
\(D,\mathbf q,\rho\).
These are identities on real momenta; analytic
transposes give the corresponding holomorphic formulas.

\begin{lemma}[Exact column identities]
\label{lem:exact-linear-columns}
On nonzero real momenta,
\begin{align}
 (c_\ell^\dagger H_fc_\ell)\mathcal E_\ell
                          &=h_1(q_\ell)^\dagger K,\notag\\
 \sum_\ell h_1(q_\ell)\mathcal E_\ell
                          &=1-a^{-1}K,\notag\\
 c_\ell^\dagger\mathcal T_\ell
                          &=h_1(q_\ell)^\dagger c^\dagger,\notag\\
 \mathcal T_\ell H_c c&=H_fc_\ell\mathcal E_\ell .
 \label{eq:exact-rescaled-column-identities}
\end{align}
\end{lemma}
\begin{proof}
The line identity \(r_L(z)Lq_i(z)=q_i(Lz)\) gives
\[
 h_{r+1}(q_\ell)d_r^{(L)}(q_\ell)
                                  =d_r(p)h_r(q_\ell).
\]
In particular
\(h_1(q_\ell)\mathbf q_\ell=f(q_\ell)\mathbf q\).
It also implies
\(\mathbf q_\ell^\dagger h_1(q_\ell)^\dagger K=0\).
Multiplication of \eqref{eq:linear-E-column} by
the right side of
\eqref{eq:diagonal-Maxwell-identity} proves the
first identity.

The covariance recursion and the definition of \(u_\ell\)
give, respectively,
\[
 \sum_\ell h_1(q_\ell)F_\ell h_1(q_\ell)^\dagger
                      =D-a^{-1}1,\qquad
 \sum_\ell u_\ell f(q_\ell)=1.
\]
The coarse version of
\eqref{eq:diagonal-Maxwell-identity} gives
\(DK=1-\mathbf q\pi\).
These three identities give the second assertion.

For the third assertion, multiplication of the first
term of \eqref{eq:linear-T-column} by
\(c_\ell^\dagger\) gives
\(h_1(q_\ell)^\dagger c^\dagger\); the gradient
part in \eqref{eq:diagonal-Maxwell-identity}
vanishes by \(c\mathbf q=0\).
The second term vanishes because
\[
 c^\dagger(1-H_c cDc^\dagger)=0.
\]
Indeed \(KD=1-D^{-1}\mathbf q\rho^{-1}\mathbf q^\dagger\) and
\(\mathbf q^\dagger c^\dagger=0\).
Finally
\[
 (1-H_c cDc^\dagger)H_c c
       =H_c c(1-DK)=H_c c\mathbf q\pi=0.
\]
The first term of \eqref{eq:linear-T-column}
then gives the fourth identity, since
\(c_\ell\mathbf q_\ell=0\).
\end{proof}

\subsection{Regular formulas for the principal alias}

We give the cancellations needed for coefficient
estimates. This also fixes the harmonic values of
the columns. Let
\[
 s=\Delta_{L{\mathsf M}_v}(p),\quad q_0=p/L,\quad
 g_i=G_{v,i}(q_0),\quad G_i=G_{v',i}(p).
\]
Define regular alias tails by
\[
 \tau_i=a^{-1}
  +\sum_{\ell\ne0}f(q_\ell)f(-q_\ell)
                r_L(q_{\ell,i})r_L(-q_{\ell,i})F_{\ell,i},
 \quad
 H=\sum_{\ell\ne0}f(q_\ell)f(-q_\ell)
                                       \operatorname{tr}F_\ell .
\]
Here \(F_{\ell,i}\) denotes its diagonal entry.
The covariance recursion says exactly that
\begin{equation}
 G_i=f(q_0)f(-q_0)r_L(q_{0,i})r_L(-q_{0,i})g_i+s\tau_i.
 \label{eq:principal-alias-cancellation}
\end{equation}
Set
\[
 A=f(q_0)f(-q_0)\sum_i g_i+sH,\qquad
 W=b_{v'}^{-1}\mathbf q^\dagger\operatorname{diag}_1(G_i^{-1}).
\]
The scalar \(A\) is bounded below by a positive
constant on real momenta, including zero. It has
a uniformly bounded holomorphic inverse on a
smaller cube, by the same Cauchy argument as before.

Write
\[
 \gamma_{0,i}=f(q_0)^{-1}(\tau_i/G_i-H/A),\qquad
 \gamma_{\ell,i}=f(-q_\ell)
       \left(\operatorname{tr}F_\ell/A
            -F_{\ell,i}r_L(q_{\ell,i})r_L(-q_{\ell,i})/G_i\right)
 \quad(\ell\ne0).
\]
Then the following formulas contain no pole:
\begin{align}
 \mathcal E_0&=
   \operatorname{diag}_1(g_i)h_1(q_0)^\dagger
                       \operatorname{diag}_1(G_i^{-1})
                    +\operatorname{diag}_1(\gamma_{0,i})\mathbf q_0W,
                                                        \notag\\
 \mathcal E_\ell&=
   sF_\ell h_1(q_\ell)^\dagger
                       \operatorname{diag}_1(G_i^{-1})
                    +\operatorname{diag}_1(\gamma_{\ell,i})\mathbf q_\ell W
                       \quad(\ell\ne0).
 \label{eq:filled-E-column}
\end{align}
In these occurrences \(\operatorname{diag}_1\)
is the usual \(4\)-by-\(4\) diagonal matrix.
To check \eqref{eq:filled-E-column}, use
\[
 D^{-1}=s\operatorname{diag}_1(G_i^{-1}),\quad
 \rho=s^2b_{v'},\quad Z=A/s,\quad \pi=W/s.
\]
The apparently singular coefficient in the principal
alias is
\[
 \frac{f(-q_0)}s\left(\frac{\sum_i g_i}{A}
                   -\frac{g_i r_L(q_{0,i})r_L(-q_{0,i})}{G_i}\right).
\]
Equation~\eqref{eq:principal-alias-cancellation}
turns it into \(\gamma_{0,i}\), proving the cancellation.

For completeness the dual column has an equally
explicit regular formula. Put
\[
 B=h_2(q_0)^{-1},\quad x_i=\tau_i/G_i,\quad
 Y=f(q_0)f(-q_0)b_{v'}/b_v(q_0),
\]
and let \(X_2\) be the diagonal \(2\)-form matrix
with entry \(x_i+x_k\) on a plaquette of type
\(\{i,k\}\). Write \(G_1=\operatorname{diag}_1(G_i)\)
and \(\tau_1=\operatorname{diag}_1(\tau_i)\).
Then
\begin{align}
 \mathcal T_0={}&Yh_2(q_0)^\dagger
 +\left[
 H_f\{X_2-s\operatorname{diag}_2(x_i)\}BcG_1
                         -H_fBc\tau_1\right]c^\dagger, \notag\\
 \mathcal T_\ell={}&H_fc_\ell F_\ell
                                h_1(q_\ell)^\dagger c^\dagger
 +\frac{f(q_\ell)f(-q_\ell)b_{v'}}{\rho_\ell}
       h_2(q_\ell)^\dagger
                 \{s^2 1-sH_c cG_1c^\dagger\},
                                      \quad\ell\ne0 .
 \label{eq:filled-T-column}
\end{align}
All factors in these formulas have already been
defined. For \(\ell=0\), \(H_f\) means \(H_v(q_0)\).
To verify its formula, substitute
\eqref{eq:principal-alias-cancellation} into
\eqref{eq:linear-T-column} and use the identity,
on each \(2\)-cell type \(\{i,k\}\),
\[
 1-(1-sx_i)(1-sx_k)=s(x_i+x_k)-s^2x_ix_k.
\]
The terms containing \(s^{-1}\) cancel by this
identity. For \(\ell\ne0\), simply use
\(\rho=s^2b_{v'}\) and \(D=s^{-1}G_1\).
This proves agreement with the previous formulas
away from the real pole. Formulas
\eqref{eq:filled-E-column} and
\eqref{eq:filled-T-column} define the columns at
the pole and on the complex cube.

\begin{lemma}[Pre-block bounds for the columns]
\label{lem:pre-block-linear-columns}
For fixed \(j,t\ge4\), \(t\in\{j,j-1\}\), and
\(a>0\), there exist \(\delta,C>0\) such that
for every finite fine history of index \(j\),
every odd \(L\ge3\), and every centered \(L\)-alias,
\begin{equation}
 \|\mathcal E_\ell(p)\|+
 \|\mathcal T_\ell(p)\|\le Cw_j(\ell)
 \label{eq:column-alias-bound}
\end{equation}
holomorphically on the enlarged cube of width \(\delta\).
\end{lemma}
\begin{proof}
For nonprincipal aliases the filled metric bounds
and the sine estimates give
\[
 |F_{\ell,i}|\le C(1+|\ell|_\infty)^{-2},\quad
 |\rho_\ell^{-1}|\le C(1+|\ell|_\infty)^{-4},\quad
 |\mathbf q_\ell|+\|c_\ell\|\le C(1+|\ell|_\infty).
\]
For the second bound use
\(\rho_\ell=L^4\Delta_{{\mathsf M}_v}(q_\ell)^2b_v(q_\ell)\).
These bounds hold on one common complex cube;
the constants in them do not depend on \({\mathsf M}_v,L\).
Moreover
\(\|h_r(q_\ell)\|\le C_jw_j(\ell)\) for
\(r=0,1,2\).

The tails \(\tau_i,H\) are bounded by the convergent
sum of \(C_jw_j(\ell)(1+|\ell|_\infty)^{-2}\)
and the fixed noise price. The quantities
\(s,G_i,g_i,b_v,b_{v'},A\) and their needed inverses
have uniform bounds. So do \(h_2(q_0)^{-1}\) and
\(f(q_0)^{-1}\), by their principal-alias lower
bounds and the Cauchy argument.
Every term of the principal formulas is now bounded.

For a nonprincipal \(\mathcal E_\ell\), the first
term has bound
\(Cw_j(\ell)(1+|\ell|_\infty)^{-2}\).
The second has bound
\(Cw_j(\ell)(1+|\ell|_\infty)^{-1}\).
For \(\mathcal T_\ell\) the first term has bound
\(Cw_j(\ell)(1+|\ell|_\infty)^{-1}\);
the second has bound
\(Cw_j(\ell)^2(1+|\ell|_\infty)^{-4}\).
Each is at most \(Cw_j(\ell)\).
This proves the claim with constants chosen before \(L\).
\end{proof}

\subsection{Restoring counting measure and taking fine differences}

For a centered fine residue \(r\), the coarse-momentum
symbols of the installed columns \(E,T\) are
\begin{equation}
 E_r(p)=L^{-1}\sum_\ell e^{iq_\ell\cdot r}\mathcal E_\ell(p),
 \qquad
 T_r(p)=L^{-2}\sum_\ell e^{iq_\ell\cdot r}\mathcal T_\ell(p).
 \label{eq:physical-linear-columns}
\end{equation}
These factors follow from counting measure.
The linear average \(Q\) has alias row
\(Lh_1(q_\ell)\); its transpose has alias column
\(L^{-3}h_1(q_\ell)^\dagger\), since a fine period
has \(L^4\) times as many sites as its coarse period.
In particular a factor \(L^{-4}\) cannot be inserted
in place of either displayed factor.

For a word \(\alpha\) of fine forward differences,
each alias acquires
\(\prod_{\nu\in\alpha}q_\nu(q_\ell)\).
Its norm is bounded by
\[
 C^{|\alpha|}L^{-|\alpha|}
                      (1+|\ell|_\infty)^{|\alpha|}.
\]
Changing the residue across its block boundary also
changes the coarse base point; the phase in
\eqref{eq:physical-linear-columns} includes that change.
Thus these are actual fine differences, including
boundary crossings.

Aliases are permuted when \(p\) is translated by
\(2\pi k\). The original expressions
\eqref{eq:linear-E-column}--\eqref{eq:linear-T-column}
therefore show periodic agreement on real overlaps
away from zero. The filled formulas extend this
agreement holomorphically. Their reality relation
is obtained by \(p,\ell\mapsto-p,-\ell\).
Taking Fourier coefficients of the periodic symbols
is consequently legitimate.
For \(0\le|\alpha|\le j-2\),
\eqref{eq:alias-summability} and
Lemma~\ref{lem:pre-block-linear-columns} now give
\begin{equation}
 |\nabla_f^\alpha E|_\sigma\le C L^{-1-|\alpha|},
 \qquad
 |\nabla_f^\alpha T|_\sigma\le C L^{-2-|\alpha|}.
 \label{eq:all-scale-linear-difference-bounds}
\end{equation}
Indeed the alias sum is bounded by a constant times
\(\sum_\ell w_{j-|\alpha|}(\ell)\).
Lemma~\ref{lem:linear-fourier-coefficients}
then gives one common \(\sigma>0\).

\begin{proposition}[The same Gaussian recursion on every period]
\label{prop:all-scale-gaussian-recursion}
Let \(K_f=d_f^TH_fd_f\) and \(K_c=d_c^TH_cd_c\),
where \(d_f,d_c\) are the fine and coarse curl
operators. The installed real columns satisfy
\begin{equation}
 K_fE=Q^TK_c,\quad
 QE=1-a^{-1}K_c,\quad
 T^Td_f=d_cQ,\quad
 TH_cd_c=H_fd_fE .
 \label{eq:physical-linear-column-identities}
\end{equation}
These identities hold on every compatible pair
of finite periods, including their harmonic fibres.
The minimum Gaussian energy on the gauge quotient
is \(\frac12x^TK_cx\).
\end{proposition}
\begin{proof}
Insert the factors in
\eqref{eq:physical-linear-columns} into
Lemma~\ref{lem:exact-linear-columns}.
For example \(K_f\) has rescaled alias symbol
\(L^{-2}c_\ell^\dagger H_fc_\ell\);
its product with the \(E\)-column has factor
\(L^{-3}\), as does \(Q^TK_c\).
The other identities follow in the same way.
The filled symbols extend these equalities to zero.
The finite Fourier transform is an invertible
change of basis with Gram matrix equal to the
number of sites times the identity, so the fibre
equalities give the displayed physical identities.

The first two identities imply
\((K_f+a Q^TQ)E=a Q^T\).
Hence \(Ex\) is stationary for
\[
 \tfrac12u^TK_fu+\tfrac a2|x-Qu|^2 .
\]
The quadratic form is nonnegative; directions in
its kernel have zero cost and can be fixed by
the prescribed gauge coordinates. Its value at
\(Ex\) is exactly \(\frac12x^TK_cx\), since
\[
 E^TK_fE+a(1-QE)^T(1-QE)
       =\{(QE)^T+(1-QE)^T\}K_c=K_c .
\]
\end{proof}

\section{Calibration on the actual tree slice at every scale}
\label{sec:all-scale-calibration}

\subsection{The frame correction in the actual predictor}

The derivative of the regular predictor of
Section~\ref{sec:actual-stationary} is not just \(Q\).
We compute the correction, because using \(Q\)
directly on the tree slice would change the noise
quadratic form.

Let \(u\) be a fine link tangent field. On one frame
patch, write \(\int_{\gamma_{z,z'}^\pi}u\) for the
oriented sum along its coordinate path. The tangent
of the rooted minimising frame is
\[
 t_z-t_0,\qquad
 t_z=-\frac{1}{|\mathcal Q_R|\,4!}
          \sum_{z',\pi}\int_{\gamma_{z,z'}^\pi}u .
\]
Indeed the linearised Euler equation for the
quadratic frame score is
\(\omega_z=\overline\omega+t_z\); the root condition
\(\omega_0=0\) fixes
\(\overline\omega=-t_0\).
Taking the weighted mean of the connector tangent
therefore gives
\begin{equation}
 \widehat Q=Q+d_{0,c}P,\qquad
 (Pu)_x=t_0^x-\sum_z w(z)t_z^x .
 \label{eq:actual-predictor-linear-frame}
\end{equation}
The signs follow from the connector
\(h_z^xU(\gamma)(h_z^{x+e_i})^{-1}\):
its frame contribution is the weighted frame
tangent at \(x\) minus that at \(x+e_i\).
This proves \eqref{eq:actual-predictor-linear-frame}
with the displayed definition of \(P\).

Let \((A_0\omega)_x=\sum_z w(z)\omega_{Lx+z}\)
for a fine site field \(\omega\). Telescoping the
paths gives
\[
 Qd_{0,f}\omega=d_{0,c}A_0\omega,\qquad
 Pd_{0,f}\omega=\omega_{Lx}-A_0\omega.
\]
Consequently
\begin{equation}
 \widehat Qd_{0,f}\omega
          =d_{0,c}\bigl(\omega|_{\text{coarse roots}}\bigr),
 \qquad
 d_c\widehat Q=d_cQ.
 \label{eq:rooted-predictor-gradient}
\end{equation}
This calculation includes overlapping patches and
repeated edge reads. It requires no volume estimate.

For any fine link field \(u\), let \(\omega(u)\)
at a fine vertex \(z\) be the sum of \(u\) along
the fixed tree path from \(z\) back to its coarse
root, minus \((Pu)\) at that root. Define
\[
 \widehat E=E+d_{0,f}\omega(E).
\]
The tree-path telescoping identity shows that
\(\widehat E\) vanishes on every tree edge. At a
coarse root the backward path is empty, so
\(\omega(E)|_{\rm roots}=-PE\).
Equations~\eqref{eq:physical-linear-column-identities}
and \eqref{eq:rooted-predictor-gradient} give
\begin{equation}
 d_f\widehat E=d_fE,\qquad
 \widehat Q\widehat E=QE=1-a^{-1}K_c,\qquad
 K_f\widehat E=\widehat Q^TK_c .
 \label{eq:actual-tree-linear-minimiser}
\end{equation}
For the last identity use
\(K_fd_{0,f}=0\) and
\(d_{0,c}^TK_c=0\).
The operator \(\omega\) reads paths in one fixed
block or frame patch. Its norm may depend on \(L,j\);
its coefficient budget is independent of the number
of blocks and of the period.

Let \(\mathcal R_{\rm rep}\) denote the linear replica
map from coarse links to fine links, and let
\(\iota\) include the free tree coordinates in
all fine link coordinates.
Since both \(\widehat E\) and the replica vanish
on tree edges, there is a unique free-coordinate
operator \(\mathcal H\) such that
\begin{equation}
 \mathcal R_{\rm rep}+\iota\mathcal H=\widehat E .
 \label{eq:actual-free-coordinate-lift}
\end{equation}
Equivalently, restrict \(\widehat E-\mathcal R_{\rm rep}\)
to the free edges. This constructs the lift in
Definition~\ref{def:calibrated-linear-data}
in the actual coordinates used in the integral.

\subsection{Constructing the row lift}

Let \(S_f,S_c\) be the physical row factors
of \eqref{eq:all-scale-row-factor}, including any
of the finite symmetry copies used by the
principal actions. Put
\[
 T\widetilde H_cd_c=\widetilde H_fd_fE .
\]
This follows from
\eqref{eq:physical-linear-column-identities}
and \eqref{eq:metric-extension-on-curvatures}
on the two sides of the step.
Define a map from coarse row space to the fine
and noise row spaces by
\begin{equation}
 \Phi_0=
 \begin{pmatrix}
   S_f\widetilde H_f^{-1}TS_c^T\\[2pt]
   a^{-1/2}d_c^TS_c^T
 \end{pmatrix}.
 \label{eq:constructed-calibration-row-lift}
\end{equation}
The inverse in this formula is that of the actual
principal row metric \(\widetilde H_f=S_f^TS_f\).
The unmodified metric was used to construct \(T\);
the identity just proved allows the row lift
to be written directly in the principal row spaces.

\begin{theorem}[Calibration with pre-block coefficient bounds]
\label{thm:constructed-calibration}
For every fixed pair of adjacent smoothing indices
and their fixed positive profile prices, the
operators \(\mathcal H,\Phi_0\) constructed above
satisfy
\[
 \mathsf B+\mathsf A\mathcal H
        =\Phi_0D\mathcal C_{\rm c}(1),\qquad
 \mathsf A^T\Phi_0=0 .
\]
There are \(\sigma_\Phi,A_\Phi>0\), chosen before
the odd side \(L\ge3\), such that
\[
 |(\Phi_0)_{\rm f}|_{\sigma_\Phi}
                     \le A_\Phi L^{-2},\qquad
 |(\Phi_0)_{\rm n}|_{\sigma_\Phi}\le A_\Phi .
\]
These statements hold for every finite fine
history and every compatible physical period.
After the tree coordinate conversion,
\(\mathcal H\) has an exponential coefficient
bound allowed to depend on \(L\). All the bounds
remain valid for the covariant row lift
\(\Phi(V)\) used in the preceding sections,
with the stated row prices.
\end{theorem}
\begin{proof}
The derivative of the coarse curvature row vector
is \(S_cd_c\). For the actual fine and noise rows,
\[
 \mathsf A=
 \begin{pmatrix}
   S_fd_f\iota\\-\sqrt a\,\widehat Q\iota
 \end{pmatrix},\qquad
 \mathsf B=
 \begin{pmatrix}
   S_fd_f\mathcal R_{\rm rep}\\
   \sqrt a(1-\widehat Q\mathcal R_{\rm rep})
 \end{pmatrix}.
\]
Use \eqref{eq:actual-free-coordinate-lift} and
\eqref{eq:actual-tree-linear-minimiser}.
The two rows of
\(\mathsf B+\mathsf A\mathcal H\) are
\[
 S_fd_fE=S_f\widetilde H_f^{-1}T\widetilde H_cd_c,\qquad
 a^{-1/2}K_c .
\]
They equal \(\Phi_0S_cd_c\), because
\(S_c^TS_c=\widetilde H_c\) and
\(d_c^T\widetilde H_cd_c=K_c\).
This proves matching.

For orthogonality, multiply out
\(\mathsf A^T\Phi_0\) and use
\(S_f^TS_f=\widetilde H_f\):
\begin{align*}
 \mathsf A^T\Phi_0
 &=\iota^T\{d_f^T\widetilde H_f\widetilde H_f^{-1}T
                              -\widehat Q^Td_c^T\}S_c^T\\
 &=\iota^T\{d_f^TT-Q^Td_c^T\}S_c^T=0 .
\end{align*}
Here \((\widehat Q-Q)^Td_c^T
=P^Td_{0,c}^Td_c^T=0\).
The last equality is the transpose of the third
identity in \eqref{eq:physical-linear-column-identities}.

All coarse factors in
\eqref{eq:constructed-calibration-row-lift}
have uniform exponential coefficient bounds.
For the fine factors, consider a fine displacement
\(z\) in \(S_f\widetilde H_f^{-1}\). It changes the output
residue of \(T\) and its coarse base by at most
\(1+|z|_1/L\). Consequently convolution costs at
most \(e^{\sigma_\Phi(1+|z|_1)}\).
Choose \(\sigma_\Phi\) below half the available
fine coefficient width. Summing over \(z\)
then preserves the \(L^{-2}\) factor from
\eqref{eq:all-scale-linear-difference-bounds};
there is no sum over all fine residues.
This proves the pre-block fine price.
The noise row is an ordinary coarse convolution
of fixed-size cell types, so its price is bounded
by a constant independent of \(L\).

The tree conversion consists of finite path sums
and the already bounded \(E\); this proves the
stated, possibly \(L\)-dependent, lift bound.
Finally replace each cover coefficient in
\(\Phi_0\) by that coefficient times the specified
adjoint transport along its coarse coordinate
path. Adjoint transport is isometric in the
fixed Lie algebra norm. The absolute cover sums
are unchanged, exactly as in the proof following
\eqref{eq:row-lift-prices}. Periodisation can
only reduce their upper bound. This proves
the assertion for \(\Phi(V)\).
\end{proof}

\subsection{Uniformity over the cutoff and closure of the history}

Choose a nondecreasing sequence of indices \(j_m\ge4\)
with \(j_{m+1}-j_m\in\{0,1\}\). Choose the profile
prices \(l_j,\lambda_j\) as above for each index.
For each cutoff \(n\), start with the empty
history at level \(n\), and append the fixed steps
\(n-1,n-2,\ldots,m\) to obtain \(v_{n,m}\).
This definition makes sense for all \(n\ge m\),
without any bound on \(n-m\) or on the product
of the block sides.

For a level of index \(j\), bounds for its metric
depend on \(j\) only. For a step whose fine index
is \(j\), take the smaller of the two coefficient
widths and the larger of the two prices for
\((j,j)\) and \((j,j-1)\). Denote the resulting
positive width and price by \(\sigma_j,C_j\).
This is a finite minimum and maximum, taken
before all block choices. Theorem
\ref{thm:constructed-calibration} therefore
provides one \(A_{\Phi,j}\) valid for every
cutoff, history and step of that index.
Only afterwards choose the odd block side
sufficiently large for the fine-curvature
reserve in Proposition~\ref{prop:good-affine-centre}.
The \(L\)-dependent chart constants and the
coupling threshold are chosen after that side,
in the order specified in
Section~\ref{sec:general-background-comparison}.

The closure assertion is stronger than a list
of bounded metrics. Equation
\eqref{eq:exact-covariance-recursion} says that
the coarse datum at one step is the fine datum
at the next step of this very history.
Equation~\eqref{eq:physical-linear-column-identities}
says it is its actual Gaussian Schur complement.
The empty history has \(H=1\), hence the Wilson
quadratic form. Positive extensions preserve
\(K=d_1^THd_1\) by
\eqref{eq:metric-extension-on-curvatures}.
Tree extraction preserves that Schur complement
by \eqref{eq:actual-tree-linear-minimiser}.
No Gaussian family is restarted between scales.

One can also check uniqueness directly. If
\((E,K_c)\) and \((E',K_c')\) obey the first
two identities of
\eqref{eq:physical-linear-column-identities}
for the same \(K_f,Q,a\), then
\[
 E'^TK_fE+a(1-QE')^T(1-QE)=K_c .
\]
Substitute \(K_fE=Q^TK_c\) and
\(a(1-QE)=K_c\) to prove this identity.
Transposing the same calculation with primes
interchanged gives the same left side equal
to \(K_c'\); symmetry of the kinetic matrices
therefore yields \(K_c=K_c'\).
Thus changing an admissible extension or a
gauge representative cannot silently choose
a different Gaussian trajectory.

The uniformity here is profilewise: the constants
may depend on \(j\), but never on the depth or
period of a history ending at that \(j\).
One need not assume that their supremum over
all \(j\) is finite. The slow profile clock
can pay any prescribed finite list of thresholds
at each index.

\begin{corollary}[Choosing the profile clock after the prices]
\label{cor:linear-profile-clock}
Choose one admissible odd side \(L_j\ge3\)
for each fine index \(j\), and then a finite
nominal-coupling threshold \(\mathcal B_j\)
covering its same-index and adjacent-index
one-step inequalities. There is a profile
clock \(j_m\to\infty\), with
\(j_{m+1}-j_m\in\{0,1\}\), such that with
\(L_m=L_{j_{m+1}}\) the nominal coupling at
every step of fine index \(j\) is at least
\(\mathcal B_j\). The same clock and blocks
are used for all ultraviolet cutoffs.
\end{corollary}
\begin{proof}
Choose \(\overline\beta\ge\mathcal B_4\).
Let \(m_4=0\), and choose strictly increasing
integers \(m_j\), for \(j\ge5\), satisfying
\[
 \overline\beta+b_0(m_j-1)\log3\ge\mathcal B_j .
\]
Set \(j_m=\max\{j\ge4:m_j\le m\}\).
The strict increase gives the stated step bound.
Every block side is at least \(3\), hence
\(\widehat\beta_m
=\overline\beta+b_0\sum_{i<m}\log L_i
\ge\overline\beta+b_0m\log3\).
A step of fine index \(j\ge5\) has
\(m+1\ge m_j\), so its coarse nominal
coupling already exceeds \(\mathcal B_j\).
For \(j=4\) use the choice of
\(\overline\beta\). No construction in this
argument depends on \(n\).
\end{proof}

This corollary pays the nominal thresholds
after the coefficient and block choices.
The comparison of actual and nominal running
couplings is the separate nonlinear induction
obligation stated in Section~\ref{sec:coupling}.

The Wilson quadratic boundary condition does
not identify a nonlinear row action with the
Wilson action. At the literal Wilson starting
density one may use the representation with
\(\lambda=1\) and no row term described in
Section~\ref{sec:actual-stationary}.
If a strict profile value \(\lambda_j<1\) is
used instead, the top row factor can be
\((\sqrt{\lambda_j}\,1,\sqrt{1-\lambda_j}\,1)^T\).
Its quadratic form is still exactly Wilson.
The remaining nonlinear difference must be
retained in the regular and forced activities.
The identity
\[
 e^{-\beta S_{\rm W}}
 =e^{-\beta S_{\rm top}}\,
                    e^{-\beta(S_{\rm W}-S_{\rm top})}
\]
specifies this retention exactly.
Matching the fourth-order curvature polynomial
makes \(S_{\rm W}-S_{\rm top}\) vanish through
degree four at the identity; it does not make
the difference identically zero. The finite
jet matching affects the choice of principal
coordinates, not the underlying Wilson measure.
The all-scale construction proved here concerns
the linear data and removes their standing-input
status in Sections~\ref{sec:microscopic-centre}--%
\ref{sec:general-background-comparison}.
The estimates for the full evolving nonlinear
activities are supplied by
Theorem~\ref{thm:nonlinear-one-step}; Wilson
initialisation and shooting give
Theorem~\ref{thm:finite-cutoff-nonlinear-induction}.
Both use the same calibrated linear trajectory.
\section{The first interacting Gaussian correction}
\label{sec:first-interacting-localisation}

The linear construction supplies the missing calibration
in the central comparison. We can now return to
localisation of the actual fluctuation integral.
We begin with its first nonquadratic connected
coefficient. This coefficient already involves
two distinct interaction vertices and their
Gaussian correlations; the quadratic determinant
calculation alone does not contain it.

The separation into connected contributions is the
principle behind the gauge-theory cluster expansion
in \cite[Sect.~2, eqs.~(2.11)--(2.13), printed p.~14]{Bal88a}.
This identifies the connected-expansion principle being adapted.
At the finite order considered
here, Wick's formula and the weighted inverse
estimate give a shorter argument. We prove the
coefficient bound directly, without assuming a
convergent expansion of an unbounded cubic
exponential.

\subsection{A local square root of the moving Gaussian covariance}

Retain the stationary graph \(\rho(x)\) and the
precision \({\mathsf H}(x)=\varphi_{yy}(x,\rho(x))\).
Restrict to a fixed smaller analytic chart on
which the rooted estimates in
Section~\ref{sec:stationary-localisation} hold.
Write \(C(x)={\mathsf H}(x)^{-1}\). All matrix coordinates
below are free fine link coordinates, including
their Lie algebra components.

\begin{lemma}[Local Gaussian whitening]
\label{lem:local-gaussian-whitening}
There are \(\theta>0,C<\infty\), independent
of period and ultraviolet cutoff within one
fixed profile, such that
\[
 W(x)={\mathsf H}(x)^{-1/2},\qquad
 \|W(x)\|_\theta+\|C(x)\|_\theta\le C .
\]
The norm is the weighted row and column norm
of \eqref{eq:weighted-schur-norm}.
The coefficients and all fixed-order coarse-field
derivatives of \(W,C\) have the rooted tree
bounds of Section~\ref{sec:stationary-localisation}.
Consequently \(z=W(x)\xi\), with independent
standard Gaussian coordinates \(\xi\), has
covariance \(C(x)\).
\end{lemma}
\begin{proof}
For \(t\ge0\), the counting spectral lower bound
for \({\mathsf H}+t1\) is \(\eta+t\).
The proof of Lemma~\ref{lem:actual-combes-thomas}
applies with the same conjugation parameter:
the change caused by conjugating \({\mathsf H}\) is
unchanged when \(t1\) is added, and is at most
\(\eta/2\). It gives
\[
 |({\mathsf H}+t1)^{-1}_{ab}|
       \le\frac{2e^{-\theta_0d(a,b)}}{\eta+t}.
\]
The number of sites at distance at most \(R\)
is bounded by \(C_L(1+R)^4\), also on a finite
torus. Choose \(\theta<\theta_0\), sum the
exponential reserve over sites, and obtain
\[
 \|({\mathsf H}+t1)^{-1}\|_\theta\le\frac{C_L}{\eta+t}.
\]
The scalar identity
\[
 \lambda^{-1/2}=\frac1\pi\int_0^\infty
                   \frac{t^{-1/2}}{\lambda+t}\,\dd t
 \quad(\lambda>0)
\]
follows by \(t=\lambda u^2\).
Diagonalising the real symmetric matrix proves
the norm-convergent formula
\[
 W=\frac1\pi\int_0^\infty
                         t^{-1/2}({\mathsf H}+t1)^{-1}\,\dd t .
\]
The preceding bound is integrable at both
endpoints, with integral \(C_L\eta^{-1/2}\).
It proves the weighted estimate.

For derivatives use
\[
 D({\mathsf H}+t1)^{-1}[v]
       =-({\mathsf H}+t1)^{-1}(D{\mathsf H}[v])({\mathsf H}+t1)^{-1}
\]
and its iterates, summed over partitions of
the derivative directions. Their products
have the rooted bounds by the tree-union
estimate already proved for matrix inversion.
The \(t\)-integrals are bounded by constants
times
\(\int_0^\infty t^{-1/2}(\eta+t)^{-k-1}\dd t\).
They are finite for every fixed \(k\).
The same argument on a sufficiently small
complex chart proves the analytic coefficient
bounds by the resolvent Neumann series.
Finally \(WW^T={\mathsf H}^{-1}\), so the covariance
assertion follows by a linear change of variables.
\end{proof}

Whitening is just a change of integration
variables; it does not change the Wilson
measure. In particular, the exact central
domain becomes
\[
 \{\xi:|W(x)\xi|_2\le D\}.
\]
It is generally an ellipsoid and does not
factor over the independent coordinates.
It must be retained, or its removal must
be estimated. We shall retain it in the
finite central integral and estimate its
effect on the coefficient below.

\subsection{The actual interaction after quadratic completion}

Use the moving Gaussian measure
\(\mu_{C(x)}\). The fixed-measure definition
\eqref{eq:central-normalized} is exactly
\begin{equation}
 Z_D(g,x)=
 \int_{|z|_2\le D}e^{U(g,x,z)}\,\dd\mu_{C(x)}(z),
 \label{eq:moving-central-interaction}
\end{equation}
where
\[
 U(g,x,z)=
 -g^{-2}\{\varphi(x,\rho+gz)-\varphi(x,\rho)
                                    -\tfrac{g^2}{2}z^T{\mathsf H}z\}
 -\{\psi(x,\rho+gz)-\psi(x,\rho)\}.
\]
The apparent singularity at \(g=0\) is removable
by Taylor's integral formula. No Haar, noise,
or determinant normaliser has been dropped:
they are exactly the normalisers in
\eqref{eq:stationary-density-identity}.

At the point \((x,\rho(x))\), set
\[
 T_k=D_y^k\varphi(x,\rho(x)),\qquad
 P_k=D_y^k\psi(x,\rho(x)).
\]
These are symmetric multilinear forms on the
free-coordinate space; \(T_k[z^k]\) means
\(T_k[z,\ldots,z]\), with \(k\) entries.
Expansion on the bounded central domain gives
\[
 U(g,x,z)=g u_1(x,z)+g^2u_2(x,z)+O(g^3),
\quad
 u_1=-\tfrac16T_3[z^3]-P_1[z],\quad
 u_2=-\tfrac1{24}T_4[z^4]-\tfrac12P_2[z^2].
\]
For this displayed expansion the remainder is
a finite-domain Taylor remainder. Its local
activity norm is not being inferred from a
global extensive bound.

Let
\[
 p_D(x)=\mu_{C(x)}\{|z|_2\le D\}.
\]
The domain and the Gaussian measure are
invariant under \(z\mapsto-z\), so the first
derivative of \(Z_D\) at zero vanishes. Thus
\begin{equation}
 \log Z_D(g,x)=\log p_D(x)+g^2L_{2,D}(x)+O(g^4),
 \quad
 L_{2,D}=\frac{1}{p_D}
       \int_{|z|\le D}(u_2+\tfrac12u_1^2)\,\dd\mu_C.
 \label{eq:actual-central-second-coefficient}
\end{equation}
This identifies the second Taylor coefficient
of the finite integral. Its evenness follows
from \(U(-g,x,-z)=U(g,x,z)\).

Define the full Gaussian polynomial coefficient
\begin{equation}
 L_2(x)=\int(u_2+\tfrac12u_1^2)\,\dd\mu_{C(x)} .
 \label{eq:full-gaussian-second-coefficient}
\end{equation}
This integral is well defined because its
integrand is a polynomial of degree at most
six. We do not define an unrestricted
Gaussian integral of \(e^{g u_1}\).

\begin{lemma}[The central cutoff changes the coefficient only by a tail]
\label{lem:second-coefficient-central-tail}
On the auxiliary volumes of
\eqref{eq:microscopic-auxiliary-size}, with
the prescribed \(D=d_0P(\widehat\beta)\),
the fixed-order chart derivatives satisfy
\[
 \|\log p_D\|_{C_w^q}
      +\|L_{2,D}-L_2\|_{C_w^q}
                         \le C_{q,L}e^{-c_LD^2}
\]
for sufficiently large nominal coupling.
This is an auxiliary-volume estimate, not
a statement about a single central ball on
arbitrarily large physical volumes.
\end{lemma}
\begin{proof}
The eigenvalues of \(C\) are at most \(\eta^{-1}\).
Gaussian integration in an orthonormal eigenbasis
gives
\[
 \int e^{\eta|z|^2/4}\,\dd\mu_C\le 2^{d/2},
 \qquad
 1-p_D\le 2^{d/2}e^{-\eta D^2/4}.
\]
Use half this exponential to absorb each of the
finitely many polynomial powers of \(z\).
Their coefficients are bounded by fixed powers
of \(d\), the profile prices and \(1+\Lambda\).
The same holds for fixed-order chart derivatives:
differentiate the moving Gaussian density on
the fixed ball. Its derivatives are the density
times polynomials in \(z\) and derivatives of
\({\mathsf H}\) and \(\log\det {\mathsf H}\).
For derivatives of the tail integral use that
the integral of the density over all space is
one. Hence no derivative of a moving boundary
is required.

On the auxiliary volumes,
\(d\le C_L(\log\widehat\beta)^m\), whereas
\(D^2=d_0^2\exp(2(\log\log\widehat\beta)^2)\).
Thus \(d=o(D^2)\), and all the indicated
polynomial factors are absorbed by a smaller
Gaussian exponential. The stated subpower
loop bounds are absorbed as well. In
particular \(p_D\ge1/2\).
Subtract the tail from the polynomial
expectation in
\eqref{eq:full-gaussian-second-coefficient},
and account for \(p_D^{-1}-1\).
The product and chain rules prove the claim.
\end{proof}

\subsection{Wick contractions and connected footprints}

For a centered Gaussian vector with covariance
\(C\), Wick's formula says that the expectation
of an even product is the sum over pairings,
with one covariance entry for each pair;
odd products have expectation zero.
For completeness, Gaussian integration by
parts gives
\[
 \int z_a f(z)\,\dd\mu_C
                   =\sum_b C_{ab}\int\partial_b f(z)\,\dd\mu_C .
\]
Apply this identity to a monomial and use induction
on its degree to obtain the pairing formula.

Write \(T_{abc},T_{abcd},P_a,P_{ab}\) for the
components of \(T_3,T_4,P_1,P_2\).
With all indices summed over the actual finite
free-coordinate space,
\begin{align}
 L_2={}&-\tfrac18T_{abcd}C_{ab}C_{cd}
        -\tfrac12P_{ab}C_{ab}\notag\\
 &+\tfrac18T_{abc}T_{def}C_{ab}C_{de}C_{cf}
        +\tfrac1{12}T_{abc}T_{def}C_{ad}C_{be}C_{cf}\notag\\
 &+\tfrac12T_{abc}P_d C_{ab}C_{cd}
        +\tfrac12P_aP_bC_{ab}.
 \label{eq:explicit-second-connected-coefficient}
\end{align}
For example, among the fifteen pairings of six
indices from two cubic vertices, nine have
one pair between vertices and six have three.
Their factors \(9/72\) and \(6/72\) give
\(1/8\) and \(1/12\).
There are three pairings between a cubic and
a linear vertex, giving \(3/6=1/2\).

Here a vertex means one factor \(T_k\) or \(P_k\);
an edge means one covariance factor linking
two of its indices. Each term has one vertex,
or has two vertices joined by at least one
covariance edge. This is the meaning of
connectedness at this order.
The logarithm has removed any product with
separate components; equivalently,
\(\log\mathbb E e^{g u_1+g^2u_2}\) has second
coefficient
\(\mathbb E u_2+\frac12
(\mathbb E u_1^2-(\mathbb E u_1)^2)\),
and \(\mathbb E u_1=0\).

\begin{proposition}[A local bound for the first interacting coefficient]
\label{prop:first-interacting-local-family}
There is an exact expansion
\(L_2(x)=\sum_\alpha l_\alpha(x)\) with connected
labelled footprints and anchors such that
\begin{equation}
 \sup_o\sum_{\alpha:o_\alpha=o}
       e^{\kappa n_B(X_\alpha)}
       \|l_\alpha\|^{\rm chart}_{q,r,s;o,X_\alpha}
                  \le C_{q,s,L,\kappa}(1+\Lambda)^2
 \label{eq:first-interacting-local-norm}
\end{equation}
on a fixed sufficiently small chart. The
constants and chart radius do not depend
on the physical volume or the ultraviolet
cutoff within a fixed profile.
The coefficient norm and rooted labels
are those of
Section~\ref{sec:stationary-localisation}.
\end{proposition}
\begin{proof}
First keep the rooted coefficient labels of
the principal path action, the loop--Haar
action and the stationary graph. Differentiating
four times in free coordinates preserves their
rooted tree bounds after a fixed radius loss:
the extra degree factors are absorbed by the
same absolute analytic majorant.
The vertex budgets are consequently \(C_L\)
for \(T_3,T_4\), and \(C_L(1+\Lambda)\)
for \(P_1,P_2\).
For \(C\), use the weighted inverse
and its rooted derivatives from
Lemma~\ref{lem:local-gaussian-whitening}.

In each of the six terms of
\eqref{eq:explicit-second-connected-coefficient},
retain every vertex label and every covariance
endpoint. Attach to a covariance entry a
specified shortest coordinate path between
its endpoints. The union of these paths
and the vertex trees is connected, because
the contraction graph is connected.
The tree length of the output label is at
most the sum of the vertex tree lengths
and covariance path lengths.
Thus its exponential tree weight is
paid by the product of the input weights.

Here is why the free sums create no volume
factor. With one vertex, fix its anchor and
sum its other indices using the vertex
majorant; a covariance pair within that
vertex is bounded by its weighted supremum.
With two vertices, choose one covariance
edge joining them. Fix the first anchor.
Sum the first vertex indices, then the
chosen covariance endpoint by its weighted
row bound, and then the second vertex
indices by the vertex majorant rooted at
that endpoint. Any remaining covariance
entries are bounded by their weighted
suprema. This procedure bounds the sum
by a product of fixed vertex and covariance
prices. It applies to both the one-edge
and three-edge contractions.

Coarse derivatives use the product rule.
At each fixed order there are finitely
many allocations among the vertices and
covariances, all with the rooted bounds
just used. The spatial polynomial weights
cost part of the exponential reserve.
The covering weight is paid using exactly
the path-cover estimate of
Section~\ref{sec:stationary-localisation};
choose \(B\) sufficiently large relative
to the fixed inverse-decay rate.
There are at most two loop--Haar vertices,
which gives the factor \((1+\Lambda)^2\).
Finally sum the finitely many Wick pairings.
Absolute convergence of the labelled
coefficient sums justifies the rearrangements.
\end{proof}

Removing the Taylor polynomial through degree
two gives the useful stronger estimate
\begin{equation}
 \bigl\|g^2(1-T_{\le2})L_2\bigr\|^{\rm chart}_{\rm local}
       \le C_{q,s,L,\kappa}\,
                       g^2(1+\Lambda)^2r^3 .
 \label{eq:first-interacting-reset-remainder}
\end{equation}
The notation on the left is the anchored family
norm in \eqref{eq:first-interacting-local-norm}.
To prove it, expand the same local analytic
families at their fixed radius, retain degrees
at least three, and use the radius ratio
argument of Corollary~\ref{cor:stationary-taylor-remainder}.
At \(r=w=\widehat\beta^{-7/16}\) and
comparable actual and nominal couplings,
the right side is
\(C_{\nu,L}\widehat\beta^{-37/16+\nu}\).
The constant, linear and quadratic jets must
still be included in the exact scalar and
coupling bookkeeping of
Section~\ref{sec:coupling-reset}; this estimate
does not discard them.

\subsection{What is now localised}

The formula and local family bound above are
for the first interacting Taylor coefficient
of the actual central Gaussian integral.
Its central-cutoff error is controlled on
the auxiliary volumes, with the original
normalisation. The all-scale linear input
used in that central comparison has been
constructed in
Sections~\ref{sec:all-scale-linear}--%
\ref{sec:all-scale-calibration}.

The following sections construct the
connected expansion with a bound for the
whole small-\(g\) remainder and compatible
gauge-covariant local families on physical
good backgrounds; see
Sections~\ref{sec:full-interacting-remainder}
and~\ref{sec:physical-interacting-extensions}.
These are additional estimates beyond
the coefficient calculation.
In particular, the ellipsoid in the
whitened variables has not been factorised,
and the fixed-background microscopic norm
has not been used as a physical activity
norm. For bounded stencils the extension
already follows from
Proposition~\ref{prop:finite-box-physical-pullback};
unbounded labelled families are treated
by the connected construction and the
short/long footprint decomposition below.
\section{A convergent interacting expansion and its full remainder}
\label{sec:full-interacting-remainder}

We retain the finite integral
\eqref{eq:moving-central-interaction}, including its
stationary normalisation. A polynomial expansion of its
integrand does not by itself justify taking a logarithm
uniformly in volume. We first construct a convergent
logarithm for bounded Gaussian coordinates, and then
compare with the actual central integral on the auxiliary
volumes. All errors in that comparison are retained.

The combinatorial argument uses the hard-core polymer
expansion and its tree bound; see
\cite[Sect.~2, eqs.~(2.2)--(2.7), preprint pp.~3--4;
Sect.~4.1, Proposition~5 and eqs.~(4.3)--(4.6),
pp.~10--11]{FP07}. The gauge-theory organisation
into connected activities is the one used in
\cite[Sect.~1, Lemmas~1--2, printed pp.~9 and~11;
Sect.~2, eqs.~(2.11)--(2.13), p.~14]{Bal88a}.
The hard-core graph identity is used in the proof below; the comparison
with the gauge-theory expansion concerns its organisation.
We give the convergence argument needed
here, and prove the interaction bounds for the present
phase. No Gaussian independence is assumed for \(z\).

\subsection{The radius used in the norm}

At this stage choose
\begin{equation}
 B=\left\lceil(\log\widehat\beta)^2\right\rceil .
 \label{eq:interacting-B}
\end{equation}
The definitions in Section~\ref{sec:activity-norms}
apply for each value of this integer. Its cardinality
price \(D_B=(2B+1)^4\) is retained in the estimates.
We use the notation
\[
 C_B=C(1+B)^p .
\]
Here \(C\) and the finite nonnegative integer \(p\)
can depend on the fixed group, block, path profile,
\(q,s,\kappa\), and the available exponential
coefficient reserves. They never depend on period,
ultraviolet cutoff, or the number of RG steps.
The power \(p\) may be increased a finite number of
times within an estimate. Thus \(C_B\) denotes a
specified polynomial allowance, not a volume
allowance. In particular, for every \(\nu>0\),
\(C_B\le C_\nu\widehat\beta^\nu\) eventually.

We shall need the incidence version of the chart
family norm. For scalar letters \(v_\alpha(x,\xi)\)
with anchors \(o_\alpha\) and finite coarse
footprints \(X_\alpha\), it is
\begin{equation}
 \sup_o\sum_{\alpha:o\in X_\alpha}
 e^{(\kappa+6)n_B(X_\alpha)}
 \sup_\xi\|v_\alpha(\,\cdot\,,\xi)\|
             ^{\rm chart}_{q,r,s;o_\alpha,X_\alpha}.
 \label{eq:interacting-letter-incidence}
\end{equation}
The Gaussian variables read by a letter are also
projected to its footprint. A footprint is nonempty
and connected, and contains its anchor.
The reserve \(6\) is convenient; any sufficiently
large fixed reserve would work.

Rebasing a derivative norm from \(o_\alpha\) to a
root \(o\) costs a polynomial in the diameter of
the connected union. Indeed
\[
 1+d(e,o)\le(1+d(o,o_\alpha))(1+d(e,o_\alpha)).
\]
Through total derivative order \(q\), the cost is
at most a fixed power of \(1+\operatorname{diam}X\).
If \(X\) is covered by a connected chain of \(k\)
\(B\)-balls, then
\(\operatorname{diam}X\le C Bk\).
The bound \(k^m\le C_{m,\epsilon}e^{\epsilon k}\)
therefore absorbs rebasing into a polynomial
allowance in \(B\) and part of the displayed
covering reserve. This explains where such
allowances enter below.

\subsection{Bounded independent coordinates}

Let \(\xi_e\in\mathfrak g\), one for each free
fine link, be independent standard Gaussian
vectors in the given Lie algebra inner product.
For \(R\ge1\) define the radial projection
\[
 \xi_e^{\,R}=
 \begin{cases}
  \xi_e,&|\xi_e|\le R,\\
  R\xi_e/|\xi_e|,&|\xi_e|>R .
 \end{cases}
\]
It is odd and commutes with every orthogonal
transformation of \(\mathfrak g\), hence with
\(\operatorname{Ad}(G)\). Smoothness in \(\xi\)
is not needed: no derivative below differentiates
this radial projection.

With \(W(x)={\mathsf H}(x)^{-1/2}\), put
\[
 z^R(x,\xi)=W(x)\xi^R,\qquad
 \mathcal Z_R(g,x)=
       \mathbb E_\gamma e^{U(g,x,z^R(x,\xi))}.
\]
The expectation is over the independent \(\xi_e\).
The original Gaussian field \(z\) remains correlated.
The group coordinates, path data, loop allowance
\(\Lambda\), \(R\), and \(\widehat\beta\) are held
fixed when \(g\) is treated as a complex analytic
variable.

\begin{lemma}[Primitive interacting letters]
\label{lem:full-interacting-letters}
On a fixed sufficiently small coarse chart,
there is an absolutely convergent expansion
\[
 U(g,x,z^R)=\sum_\alpha v_\alpha(g,x,\xi)
\]
with finite connected footprints. For
\(|g|R\) sufficiently small, its incidence norm
\eqref{eq:interacting-letter-incidence} is at most
\begin{equation}
 C_B|g|H_R,\qquad
 H_R=R^3+(1+\Lambda)R .
 \label{eq:full-primitive-price}
\end{equation}
Each letter is holomorphic in \(g\), vanishes at
\(g=0\), and satisfies
\(v_\alpha(-g,x,\xi)=v_\alpha(g,x,-\xi)\).
The constants and the chart radius are independent
of the physical volume and ultraviolet cutoff
within a fixed profile.
\end{lemma}
\begin{proof}
At the stationary point the exact interaction has
the convergent series
\begin{align}
 U(g,x,z^R)
 ={}&-\sum_{k\ge3}\frac{g^{k-2}}{k!}
              D_y^k\varphi(x,\rho)[(W\xi^R)^k]\notag\\
    &-\sum_{k\ge1}\frac{g^k}{k!}
              D_y^k\psi(x,\rho)[(W\xi^R)^k].
 \label{eq:full-interaction-series}
\end{align}
The linear phase term is zero and the quadratic
phase term has already been included in \({\mathsf H}\).
These cancellations are made before taking
absolute coefficient bounds.

Keep the path labels of each derivative tensor,
the rooted labels of \(\rho,W\), and the endpoint
of each \(\xi^R\). The rooted analytic bounds of
Section~\ref{sec:stationary-localisation} and
Lemma~\ref{lem:local-gaussian-whitening} give,
for a fixed radius \(a>0\),
vertex prices at most \(C k!a^{-k}\) for the
phase and \(C(1+\Lambda)k!a^{-k}\) for \(\psi\).
Products of their trees connect every endpoint
to the original vertex. Attach the corresponding
paths, exactly as in
Proposition~\ref{prop:first-interacting-local-family}.
The coefficient bound for \(W\) pays the sums
over the endpoints. Each such endpoint contributes
at most \(R\).
The two series of prices are bounded by
\[
 C\sum_{k\ge3}|g|^{k-2}(CR)^k
       \le C|g|R^3,\qquad
 C(1+\Lambda)\sum_{k\ge1}(|g|CR)^k
       \le C(1+\Lambda)|g|R ,
\]
after \(|g|CR\le1/2\).

To charge incidence at a site rather than only
at an anchor, retain that site in the connected
tree and sum the possible anchor positions.
The exponential tree reserve pays this sum:
for any fixed power \(m\),
\(\sum_{u\in\Z^4}(1+|u|)^m e^{-c|u|}<\infty\).
The degree factors introduced by choosing a
marked vertex are absorbed by a further fixed
analytic-radius loss.
The covering estimate and derivative rebasing
described above give at most \(C_B\).
This proves \eqref{eq:full-primitive-price}.
Finally each phase letter has Gaussian degree
\(k\) and \(g\)-degree \(k-2\), while each
\(\psi\) letter has both degrees \(k\).
Oddness of the radial projection proves the
stated parity for each retained label.
\end{proof}

\subsection{From independent coordinates to connected polymers}

Let \(a_\alpha=e^{v_\alpha}-1\).
Use the algebraic seminorm of
Lemma~\ref{lem:finite-jet-algebra}, with the original
chart domain. If \(t_\alpha\) is its price for
\(v_\alpha\), the Mayer factor has price
\(b_\alpha=e^{t_\alpha}-1\le2t_\alpha\) once
\(t_\alpha\le1\).
An original maximum-norm price is enlarged by at most
\(e\) to obtain \(t_\alpha\).
Thus the incidence bound is at most \(2e\) times
\eqref{eq:full-primitive-price} when that quantity
is small. Form all products in this seminorm.
Absolute summability on a finite torus justifies
\[
 e^{\sum_\alpha v_\alpha}
       =\sum_{A\ {\rm finite}}\prod_{\alpha\in A}a_\alpha .
\]
Each primitive address is selected at most once
in this formula.

A polymer \(P\) is a nonempty finite set of
primitive addresses connected by overlap of
their footprints. Define
\[
 X(P)=\bigcup_{\alpha\in P}X_\alpha,\qquad
 \zeta(P)=\mathbb E_\gamma\prod_{\alpha\in P}a_\alpha.
\]
Disjoint footprint components read disjoint
independent Gaussian vectors, so their
expectations factor. It follows exactly that
\begin{equation}
 \mathcal Z_R
   =\sum_{\{P_1,\ldots,P_m\}\ {\rm disjoint}}
                               \prod_{j=1}^m\zeta(P_j).
 \label{eq:clipped-polymer-gas}
\end{equation}
Disjointness here refers to the footprints.
Using overlap of the larger declared footprints
may exclude more pairs than overlap of the
Gaussian variables alone; it still gives the
displayed component factorisation.
The expectation here is under the fixed white
Gaussian law \(\gamma\), so the algebraic seminorm
of an expectation is at most the expectation of
that seminorm. Rebase once per whole polymer using
the algebraic version of
Lemma~\ref{lem:finite-jet-anchor-product}.
Keep these algebraic prices through the logarithm;
the maximum derivative norm is recovered at the end.
This statement is not used to omit derivatives of
the background-dependent physical Gaussian law.

Lemma~\ref{lem:connected-bound}, applied to the
primitive prices and the available covering
reserve, gives a polymer incidence price
\[
 \sup_o\sum_{P:o\in X(P)}
      e^{(\kappa+2)n_B(X(P))}\|\zeta(P)\|^{\rm alg}_q
                         \le C_B|g|H_R=:\epsilon .
\]
The subscript \(q\) here denotes the algebraic chart
seminorm after rebasing; all its costs
are included in \(C_B\).

We recall the logarithm estimate, including
why it has no volume factor.

\begin{lemma}[Rooted logarithm bound]
\label{lem:rooted-interacting-log}
For polymers with nonempty footprints, suppose
their absolute derivative prices have incidence
sum \(\epsilon\) at covering exponent
\(\kappa+3\). A derivative price bounds the original
maximum-over-orders norm at the polymer's own anchor;
the weighted chart version is also allowed.
If \(D_B\epsilon\le1/2\), the
hard-core gas has a convergent logarithm with
connected footprints and anchored family norm
at exponent \(\kappa\) at most \(C_B\epsilon\).
The same bound controls incidence of the absolute
cluster prices. All estimates are uniform in the
finite volume.
Alternatively, algebraic derivative prices at
covering exponent \(\kappa+2\) suffice, with the same
smallness condition and final maximum-norm bound.
\end{lemma}
\begin{proof}
For an ordered tuple \(P_1,\ldots,P_n\), let
\(\mathcal U(P_1,\ldots,P_n)\) be
\[
 \sum_{\substack{G\ {\rm connected}\\
                   \text{on }\{1,\ldots,n\}}}
       \prod_{\{i,j\}\in E(G)}
                   (-1_{\{X(P_i)\cap X(P_j)\ne\varnothing\}}).
\]
For \(n=1\) this is \(1\). The logarithm series is
\begin{equation}
 \sum_{n\ge1}\frac1{n!}
       \sum_{P_1,\ldots,P_n}
       \mathcal U(P_1,\ldots,P_n)\prod_i\zeta(P_i).
 \label{eq:interacting-Ursell}
\end{equation}
Repeated polymers are permitted in this ordered
series; their factorial is essential.

For a fixed incompatibility graph, the absolute
value of \(\mathcal U\) is at most its number of
spanning trees. One proof partitions its connected
spanning subgraphs into intervals
\([T,\mathcal P(T)]\). Root at vertex \(1\),
use graph distance for the levels, and choose
as parent the smallest-labelled neighbour in
the preceding level. The allowed extra edges
join vertices on the same level, or join
adjacent levels without replacing the chosen
smallest parent. These rules give precisely
the interval associated to the resulting tree.
The alternating sum over the optional edges
is \(0\) or has absolute value \(1\).
Summing the intervals proves the tree bound.
This is the hard-core case of the partition
identity discussed in \cite[Sect.~4.1, Proposition~5 and
eqs.~(4.3)--(4.4), preprint p.~10]{FP07}.

We now use one pinned estimate for both scalar
summation and derivatives. This is the rooted expansion
of \cite[Sect.~2, eqs.~(2.6)--(2.7),
PDF p.~3; convergence conditions, PDF p.~4]{FP07}.
Let \(p(P)\) be the assumed derivative price and
write \(k_P=n_B(X(P))\). Put
\[
 \rho_P=e^{(\kappa+1)k_P}
                       \|\zeta(P)\|^{\rm alg}_q .
 \]
The auxiliary seminorm uses the same fields, directions,
and anchor as the original norm. For chart functions the
proof is identical, using ordinary directional derivatives;
the chart domain is held fixed.
For physical germs the derivative norm is taken on the
healthy cluster union. Coarse masks are restored after
this unmasked estimate.
Lemma~\ref{lem:finite-jet-algebra} gives
\(\|\zeta(P)\|^{\rm alg}_q\le e\,p(P)\).
Since \(k_P\ge1\),
\begin{equation}
 \sup_x\sum_{P:x\in X(P)}\rho_Pe^{k_P}
 \le\sup_x\sum_{P:x\in X(P)}
                   p(P)e^{(\kappa+3)k_P}\le\epsilon .
 \label{eq:pinned-log-price}
\end{equation}
For the algebraic-price version, the same inequality
follows directly from its \(\kappa+2\) input and
does not need the additional conversion unit.
Apply the rooted-tree part of Lemma~\ref{lem:connected-bound}
with output exponent zero, reserve one and weights \(\rho_P\).
Its hypothesis is \eqref{eq:pinned-log-price}, since
\(D_B\epsilon\le1/2\le1\). A polymer is incompatible with
itself, so the neighbour sum includes \(Q=P\). The
spanning-tree bound already proved and
\eqref{eq:common-pinned-tree-bound} give
\begin{equation}
 \Pi_P(\rho):=
 \sum_{n\ge0}\frac1{n!}\sum_{P_1,\ldots,P_n}
       |\mathcal U(P,P_1,\ldots,P_n)|
                      \prod_{j=1}^n\rho_{P_j}
                      \le e^{k_P}.
 \label{eq:pinned-log-bound}
\end{equation}
The common root sum omits the root price, so this also
holds when that price is zero.

It remains to check that derivatives cost only one
constant after all vertices are assembled.
Equation~\eqref{eq:finite-jet-cover-reserve}, with
\(\eta=1\), bounds the original norm of a cluster
at any anchor in its union, including its output
covering weight, by \(C_B\prod_i\rho_{P_i}\).
The factor \(q!\) and the polynomial in the total
cover size are paid once here. They are not charged
at every vertex.
For a fixed site \(x\), bound the indicator that a
cluster union contains \(x\) by the number of its
vertices containing \(x\). Symmetry of the ordered
tuple changes \(1/n!\) to \(1/(n-1)!\) when such a
vertex is selected. Its total price is consequently
at most
\[
 C_B\sum_{P:x\in X(P)}\rho_P\Pi_P(\rho)
      \le C_B\sum_{P:x\in X(P)}\rho_Pe^{k_P}
      \le C_B\epsilon .
 \]
This proves the incidence and anchored bounds
without a sum over all lattice sites. The three
units of covering reserve have paid, respectively,
the common-anchor polynomial, the rooted descendants,
and the conversion from the original derivative norm
to the algebraic seminorm. If algebraic prices have
already been supplied, only the first two units are
used.

For finite alphabets the equality between
\eqref{eq:interacting-Ursell} and the logarithm
follows by expanding the exponential: a graph
is the unordered collection of its connected
components. With a fugacity multiplying all
activities it holds near zero as an identity
of convergent power series. The rooted bounds
extend it throughout the stated smallness
region. Passing to countable alphabets uses
the same absolute majorants. The exponential
of the convergent logarithm equals the gas;
in particular that gas has no zero there.
\end{proof}

\begin{lemma}[Marked logarithmic clusters]
\label{lem:marked-interacting-log}
Split the polymer species into ordinary species and
marked species. Let \(b\) and \(f\) be their respective
derivative-price incidences at exponent \(\kappa+3\),
in the original maximum-over-orders norm.
If \(D_Bb\le1/4\) and \(D_Bf\le1/4\), the logarithmic
clusters with at least one marked occurrence have
absolute derivative incidence at exponent \(\kappa\)
bounded by \(2C_Bf\), with the \(C_B\) of the preceding
lemma. The original cluster functions and anchors
are retained.
\end{lemma}
\begin{proof}
If \(f=0\), every marked cluster has zero price.
Otherwise put \(t=(4D_Bf)^{-1}\ge1\) and multiply each
marked occurrence price by \(t\). The new total income
is at most \(b+tf\le(2D_B)^{-1}\).
The preceding lemma bounds its nonnegative absolute
cluster majorant by \(C_B(b+tf)\), uniformly in volume.
Only majorant prices are scaled; the physical kernel,
Gaussian law, and cluster functions are unchanged.
A cluster with \(r\ge1\) marked occurrences has acquired
the factor \(t^r\ge t\). Dividing this majorant by \(t\)
therefore bounds the original marked part by
\[
 C_B(b/t+f)\le2C_Bf.
 \]
All series here are the absolute majorants of the
Ursell expansion, with repeated species and factorials.
No distinct-label subset sum is substituted for a
logarithmic cluster sum.
For finite alphabets the marking argument is an
identity of nonnegative series. The uniform majorant
and monotone convergence give the countable case.
\end{proof}

\subsection{The full fourth-order remainder}

Increase \(C_B\) to include the finite powers
of \(D_B\) and all derivative rebasing costs.
There is a fixed \(c>0\) such that
\[
 r_g=\frac{c}{C_BH_R}
\]
is a common complex \(g\)-radius for the
convergent logarithm. It also ensures
\(|g|R\) lies in the analytic domain of
\eqref{eq:full-interaction-series}.

\begin{theorem}[Connected interacting remainder]
\label{thm:full-interacting-remainder}
For \(|g|\le r_g/2\), there is an exact connected
family expansion
\begin{equation}
 \log\mathcal Z_R(g,x)
       =g^2L_2^R(x)+\sum_\alpha r_\alpha^R(g,x),
 \qquad
 \|r^R\|_{\mathcal E,\kappa}^{\rm chart}
                         \le C_B g^4H_R^4 .
 \label{eq:full-interacting-fourth-remainder}
\end{equation}
The norm is the anchored chart family norm.
Here \(L_2^R\) is the second coefficient
computed with \(z=W\xi^R\). It satisfies
\begin{equation}
 \|L_2^R-L_2\|_{\mathcal E,\kappa}^{\rm chart}
       \le C_B(1+\Lambda)^2(1+R)^6e^{-cR^2}.
 \label{eq:clipped-second-coefficient-error}
\end{equation}
All bounds are uniform over finite periods
and ultraviolet cutoffs within a fixed profile.
\end{theorem}
\begin{proof}
The preceding lemmas give an analytic connected
family for \(\log\mathcal Z_R\) on \(|g|<r_g\),
with uniformly bounded family norm there.
The primitive parity, followed by the change
\(\xi\mapsto-\xi\), makes every polymer weight
even in \(g\). The same is true of each
term in \eqref{eq:interacting-Ursell}.
Each vanishes at zero.
Its second coefficient, summed absolutely,
is exactly \(L_2^R\), by differentiating
the bounded-coordinate integral twice.

Apply Cauchy's coefficient estimate label by
label on a fixed smaller circle, retaining
the absolute rooted majorant. Summing the
even powers of degree at least four gives
\[
 \sum_{k\ge2}C(|g|/r_g)^{2k}
       \le C|g|^4r_g^{-4}.
\]
Absorb the resulting fixed power of the
polynomial allowance into \(C_B\). This
proves \eqref{eq:full-interacting-fourth-remainder}.
Thus the estimate is for the whole analytic
remainder, not just a list of finite-order
Feynman coefficients.

The polynomial defining \(L_2^R\) has Gaussian
degree at most six. Expand its \(W\)-entries
with the same rooted labels as in Section
\ref{sec:first-interacting-localisation}.
A monomial reads at most six Gaussian
vectors. It agrees with the unprojected
monomial whenever all these vectors have
norm at most \(R\).
Gaussian tails and Cauchy--Schwarz give,
uniformly for every such monomial of degree
at most six,
\[
 \left|\mathbb E p(\xi^R)-\mathbb E p(\xi)\right|
     \le C_G(1+R)^6e^{-c_GR^2}
\]
times its coefficient norm. This also holds
for repeated vectors. Summing the coefficient
majorants and differentiating the coarse
coefficients proves
\eqref{eq:clipped-second-coefficient-error}.
No union bound over all sites of the torus
was used in this local coefficient estimate.
\end{proof}

\subsection{Returning to the actual central integral}

For this comparison only, impose the auxiliary
size bound \eqref{eq:microscopic-auxiliary-size}.
Let \(d_*(\widehat\beta)\) be a fixed polynomial
in \(1+\log\widehat\beta\) bounding the number
of free real coordinates on every allowed
auxiliary volume. Choose
\begin{equation}
 R=\frac{D}{4C_W\sqrt{d_*}},\qquad
 D=d_0P(\widehat\beta),
 \label{eq:white-radius-auxiliary}
\end{equation}
where \(C_W\) bounds \(\|W\|_{2\to2}\) and is
increased to at least \(1\). Eventually \(R\ge1\).
This choice is made from the common size bound,
not from a subsequently chosen physical volume.
Both \(R\) and \(H_R\) are subpower functions
of the nominal coupling.

\begin{lemma}[Central integral versus bounded-coordinate reference]
\label{lem:central-versus-clipped}
On the auxiliary volumes, for every fixed
derivative order \(q\),
\begin{equation}
 \|\log Z_D(g,\cdot)-\log\mathcal Z_R(g,\cdot)\|_{C_w^q}
                             \le C_{q,L}e^{-c_LR^2}
 \label{eq:central-clipped-comparison}
\end{equation}
at the physical \(g=\beta^{-1/2}\), for
sufficiently large nominal coupling.
\end{lemma}
\begin{proof}
On the event \(\max_e|\xi_e|\le R\), projection
does nothing, and
\[
 |W\xi|_2\le C_W\sqrt{d_*}R\le D/4.
\]
The exact central integrand and the reference
integrand therefore agree on this event.
Its complement has probability at most
\(C_Gd_*e^{-c_GR^2}\).

On the central ellipsoid, the analytic local
interaction bounds, summed over the auxiliary
sites, give
\[
 |U|\le C_L n_c|g|
             \{D^3+(1+\Lambda)D\}=o(1).
\]
The same assertion holds for the projected
integrand with \(D\) replaced by \(C_LR\).
Here \(n_c\) is polynomial in \(\log\widehat\beta\);
every other displayed factor apart from \(g\)
is subpower. The two expectations differ in
value by at most \(Ce^{-cR^2}\).
Both are bounded below by a positive constant.

For derivatives use the fixed-\(z\) form
\eqref{eq:moving-central-interaction}, rather
than differentiating the ellipsoid boundary.
The higher derivatives of the central integral
have the polynomial coupling and auxiliary-size
prices established in
Section~\ref{sec:gaussian-remainder}.
Those of \(\mathcal Z_R\) follow directly
from its bounded-coordinate integral.
Apply the finite-difference interpolation
argument of that section to their value
difference on a slightly larger chart.
For example a difference step
\(t=e^{-aR^2}\), with fixed sufficiently small
\(a>0\), makes both the value-error term
\(t^{-q}e^{-cR^2}\) and the higher-derivative
error term exponentially small.
All polynomial prices are absorbed because
\(R^2\) dominates every fixed power of
\(\log\widehat\beta\), and in particular
dominates their logarithms.
The chart has a reserve much larger than
this difference step.
The chain rule for the logarithm proves
\eqref{eq:central-clipped-comparison}.
\end{proof}

This comparison retains the fact that
\(Z_D(0,x)=p_D(x)\), whereas
\(\mathcal Z_R(0,x)=1\). Their difference is
one of the Gaussian tails just estimated.
The estimate also retains the original
noise and Haar normalisers through the
exact stationary formula.

\begin{corollary}[The complete central interacting error]
\label{cor:complete-central-interacting-error}
On the auxiliary volumes, the actual central
logarithm has the exact form
\[
 \log Z_D(g,x)=g^2L_2(x)
                    +\sum_\alpha r_\alpha(g,x)+e_D(x),
\]
where the connected family has norm at most
\(C_B g^4H_R^4\), plus the local exponential
error of
\eqref{eq:clipped-second-coefficient-error},
and \(e_D\) satisfies
\eqref{eq:central-clipped-comparison}.
In particular, after retaining the Taylor
jets through degree two specified in Section
\ref{sec:first-interacting-localisation},
the local chart remainder is bounded by
\begin{equation}
 C_B\{g^4H_R^4+
                  g^2(1+\Lambda)^2w^3\}
       +Ce^{-cR^2}
             \le C_{\nu,L}\widehat\beta^{-2+\nu}.
 \label{eq:complete-central-interacting-rate}
\end{equation}
The exponentially small auxiliary error may
be assigned its full auxiliary footprint
without losing this bound.
\end{corollary}
\begin{proof}
Combine
Theorem~\ref{thm:full-interacting-remainder},
Lemma~\ref{lem:central-versus-clipped}, and
\eqref{eq:first-interacting-reset-remainder}.
The two polynomial terms have orders
\(\widehat\beta^{-2+\nu}\) and
\(\widehat\beta^{-37/16+\nu}\), respectively.
For the last assertion,
\(n_B(X)\le |X|\le C(\log\widehat\beta)^m\)
on an auxiliary volume.
This covering weight and its derivative
rebasing costs are absorbed by \(e^{-cR^2}\):
by \eqref{eq:white-radius-auxiliary},
\(R^2=P^2/\operatorname{polylog}(\widehat\beta)\)
dominates every fixed power of
\(\log\widehat\beta\).
This argument is deliberately confined to
the auxiliary volumes.
\end{proof}

The reference \(\mathcal Z_R\) has a
volume-independent connected expansion on
arbitrary finite periods. Its identification
with one central ball was proved only on
the auxiliary volumes. On unrestricted
physical volumes the Gaussian-coordinate
failures have to be retained in the
large-field reassembly; a union bound for
one central ball does not replace that step.
We next construct the local functions in
this reference expansion directly on
physical good backgrounds.
\section{Covariant kernels on physical good backgrounds}
\label{sec:covariant-kernel-construction}

The chart calculation must be extended without choosing
one logarithm of the whole periodic gauge field.
We construct the required kernels from parallel transport
on the integer cover. Their local functions will then
have a single physical definition, rather than separate
definitions in unrelated local gauges.

We keep the actual unmasked phase and loop--Haar
factor of the fluctuation integral. Finite-range
approximations and Gaussian covariance changes are
another way of organising this step. Here a
covariant parametrix is shorter for the all-good
column: a parametrix is an explicitly constructed
approximate inverse, whose defect is small in a
specified operator norm. Its Neumann series will
give the exact inverse. The exponential flat-kernel
estimates proved above, in the spirit of
\cite[Sect.~2.4, Lemma~2.12, preprint p.~14;
Sect.~3.5, Lemma~3.17, p.~32]{DST23}, provide that parametrix.
The cited scalar estimates are methodological comparisons; the
covariant parametrix used here is constructed below.
No estimate for a scalar Laplacian is substituted
for the gauge-fixed Hessian.

\subsection{Backgrounds, coordinates, and the norm to be proved}

Use the same rooted tree in each block; the exact
Haar tree-gauge identity allows this choice.
Let \(V\) be a coarse configuration whose
plaquette chords are at most \(h_{\rm c}\).
Write
\[
 U(V,y)_e=\exp(y_e)\operatorname{rep}(V)_e
\]
on free edges, and set the tree edges to one,
as in the original integral.
The phase \(\varphi(V,y)\) means that very
principal phase, with its regular frame and
predictor branches, written in these coordinates.
It includes the original noise coefficient.
For the physical good-background estimates
in this section, impose the explicit geometric
reserve
\begin{equation}
 \ell_{\rm c}\ge C_{\rm geom}B ,
 \label{eq:good-background-period-reserve}
\end{equation}
where \(\ell_{\rm c}\) is the coarse lattice
period. The fixed constant includes the
frame patches and the short filling boxes
used below. All volume estimates are uniform
over periods satisfying this reserve.
At the end of the construction we show how
one fixed Wilson family guarantees it at
every step, before any volume subsequence
is chosen.
For sufficiently small \(h_{\rm c}\) and
\(\|y\|_\infty\), these branches exist uniformly:
on each fixed frame patch a local gauge makes
the replica links \(O(C_Lh_{\rm c})\), and the
finite-patch branch lemma then applies.
The infinite path sums are analytic in a common
smaller \(y\)-ball. Indeed a complex link
perturbation of size \(r\) costs at most
\(e^{Cr\ell}\) in a word of length \(\ell\);
choose \(Cr\) smaller than the path coefficient
reserve.

A coarse gauge transformation \(u\) acts on a
free tangent vector at \(e\) by
\[
 (\mathcal O_u y)_e
             =\operatorname{Ad}_{u_{\operatorname{parent}(e)}}y_e.
\]
The replica, rooted tree coordinates and
regular predictor obey
\begin{equation}
 \varphi(u\cdot V,\mathcal O_u y)=\varphi(V,y).
 \label{eq:physical-phase-covariance}
\end{equation}
This is an identity of the physical functions:
path transports telescope their endpoint gauge
matrices, and the branch minimisers transform
by uniqueness. Negative link orientations
transform their tangent at the appropriate
endpoint. Thus no replacement of the predictor
by an ordinary noncovariant block average
is implicit in \eqref{eq:physical-phase-covariance}.

Let \(a,b\) denote free link indices without
their colour coordinates, and set
\(d(a)=\operatorname{parent}(a)\).
For a local matrix letter \(K_{\alpha,ab}(V)\),
use the operator norm on its Lie algebra entries.
Its local derivative norm is the physical norm
of Section~\ref{sec:activity-norms}, with scalar
absolute value replaced by the operator norm,
anchored at \(d(a)\). A matrix family is priced by the
maximum of its row and column sums, with
the factor \(e^{\kappa n_B(X_\alpha)}\) in
each sum. Column sums are rebased to their
column anchor. Scalar and vector letters are
treated in the same way.

We also retain the native length \(\ell_\alpha\):
the sum of the lengths of the declared cover
paths, plus one per primitive vertex.
It is measured before any torus identifications.
At a reduced exponent \(\theta>0\), we retain
an additional factor \(e^{\theta\ell_\alpha}\)
in these family bounds. The read footprint can
contain a filled box for a short letter; its
native length still records the paths that
produced the letter. This extra bookkeeping
will control the long-family tail.
The added test plaquettes and boxes are
read in \(X_\alpha\); they are not counted
as additional transport words in
\(\ell_\alpha\).
More precisely, include both attached sites
in each \(X_\alpha\), and for letters with
output \(a_\alpha\) and input \(b_\alpha\) put
\begin{align}
 \|\mathcal K\|_{\rm family}
 =\max\Bigl\{&
 \sup_a\sum_{\alpha:a_\alpha=a}
 e^{\kappa n_B(X_\alpha)+\theta\ell_\alpha}
       \|K_\alpha\|_{q,h_{\rm c},s;d(a),X_\alpha},\notag\\
 &\sup_b\sum_{\alpha:b_\alpha=b}
 e^{\kappa n_B(X_\alpha)+\theta\ell_\alpha}
       \|K_\alpha\|_{q,h_{\rm c},s;d(b),X_\alpha}
                                                    \Bigr\}.
 \label{eq:covariant-kernel-family-norm}
\end{align}
The vector version has one attached index.
A single positive final length reserve
is fixed below all the finitely many
losses in the constructions below.
Unbounded chain and tree lengths are
summed inside that fixed reserve.

The polynomial allowances \(C_B\) have the
meaning fixed in Section
\ref{sec:full-interacting-remainder}. Matrix
products, tensor contractions and derivative
rebasing cost such allowances and fixed
exponential reserves. The reasons are explicit:
endpoints are summed using the row or column
budget, the union of the declared footprints
is connected, and the total number of derivative
insertions is bounded by \(q\). The resulting
powers of length are paid by a reduced native
exponential reserve.

\subsection{Covariantising a flat kernel before periodisation}

The flat precision
\[
 A_0=\varphi_{yy}(1,0)
\]
has scalar colour entries: the first
derivatives of curvature and of the predictor
are scalar cochain maps tensored with
\(1_{\mathfrak g}\). On the integer cover
write its scalar coefficients as
\(A_{0,ab}(z)\), where \(a,b\) now include
the finitely many free residues in a block
and \(z\in\Z^4\) is the coarse displacement.
The flat gap and the exponential kernel
bounds give constants \(\eta,K,\sigma>0\)
such that
\[
 \eta1\le A_0\le K1 .
\]
The constants are uniform over the cutoff
within a fixed profile and block.
Fourier analyticity, or the weighted resolvent
proof already given, supplies exponential
coefficients for
\[
 P_t^0=(A_0+t1)^{-1},\quad
 B_0=A_0^{-1/2},\quad Q_0=A_0^{1/2},
 \qquad
 |P_t^0|_\sigma\le\frac{C}{\eta+t}.
 \]
The two square-root prices are bounded by \(C\).
The same coefficients give these operators
on every period. Positivity on the real
Fourier cube follows from the flat bound
on all periods and density of their
Fourier grids.

Choose paths \(\gamma_{ab,z}\) between the
attached coarse sites, of length at most
\(C_L(1+|z|_1)\), with the reverse path for
the reversed ordered pair. For a scalar
kernel \(K\) define
\begin{equation}
 K[V]_{ab}(z)
      =K_{ab}(z)\operatorname{Ad}_{V(\gamma_{ab,z})}.
 \label{eq:covariantised-flat-kernel}
\end{equation}
The kernel is periodised only after this
operation. Symmetric kernels remain symmetric,
and all kernels are covariant at their two
endpoints. Absolute coefficient and derivative
prices follow from unitarity of adjoint
transport and ordered word insertions.

\begin{lemma}[Covariant product defect]
\label{lem:covariant-product-defect}
Suppose \(K,L\) have exponential cover
coefficients. The difference
\(K[V]L[V]-(KL)[V]\) has an exact local
letter expansion with physical row and
column family price
\begin{equation}
 C_B\bigl(h_{\rm c}+e^{-cB}\bigr)|K|_\sigma|L|_\sigma
 \label{eq:covariant-product-defect}
\end{equation}
on the good sets of its declared footprints.
It retains a positive native exponential
reserve. The constants are independent of
the surrounding period.
\end{lemma}
\begin{proof}
For each intermediate index and pair of cover
displacements, retain the paired letter
\[
 K_{ai}(z)L_{ib}(w)
 \left\{\operatorname{Ad}_{V(\gamma_{ai,z}\gamma_{ib,w})}
            -\operatorname{Ad}_{V(\gamma_{ab,z+w})}\right\}.
\]
This pairing is made before taking absolute
values. The two paths have the same cover
endpoints. Their quotient is a closed cover
loop of length at most \(C_L\ell\), where
\(\ell=1+|z|_1+|w|_1\).
Coordinate square moves fill it in a box
of radius \(C_L\ell\), using at most
\(C_L\ell^2\) plaquettes.
Consequently its adjoint defect is at most
\(C_G\ell^2h_{\rm c}\) when those plaquettes
are good.

Separate the letters into \(\ell\le B\)
and \(\ell>B\). The fixed geometric
reserve contains the short filling boxes.
For a short letter declare the filling box
and a fixed halo as its footprint.
This box is covered by a fixed number of
\(B\)-balls, independently of \(\ell,B\).
The value estimate just obtained is therefore
compatible with the covering weight.
For a derivative of positive order \(j\),
word insertion gives a bound
\[
 C_{G,q,s}h_{\rm c}^{\,j}
                         (1+\ell)^{(s+1)j+C}.
\]
Since \(h_{\rm c}\le1\), this is at most
\(C h_{\rm c}(1+\ell)^{C_{q,s}}\).
The available exponential coefficient
reserve sums these polynomial factors.

For a long letter retain only the two paths
and their fixed halos. Do not fill its box.
The adjoint difference has norm at most two,
and its derivatives have the same polynomial
word prices without a small-curvature factor.
Its footprint is covered by a chain of
at most \(C_L(1+\ell/B)\) balls.
For sufficiently large \(B\) the covering
weight costs less than half the coefficient
length reserve. The remaining long sum is
bounded by \(Ce^{-cB}\).
Incidence and derivative rebasing cost only
polynomial length factors.

All these arguments are made on the cover.
In particular the closed comparison loop has
zero cover displacement, even if the paths
wrap after periodisation. Repeated plaquettes
and links are counted with multiplicity.
Summing absolute cover coefficients before
identifying images proves the stated period
uniformity.
\end{proof}

This proof explains the short/long distinction.
Filling an arbitrarily large box would cost
roughly \(e^{C\kappa(\ell/B)^4}\), which cannot
be paid by an exponential in \(\ell\).
Only the short letters use filled boxes.
Long letters retain their path footprints
and their actual exponential coefficient tail.

\subsection{The actual background gradient and Hessian}

Put
\[
 b(V)=\varphi_y(V,0),\qquad A(V)=\varphi_{yy}(V,0).
\]
Their covariance follows from
\eqref{eq:physical-phase-covariance}.
The following is the model-specific estimate
needed before using a parametrix.

\begin{lemma}[Small-curvature comparison of the actual Hessian]
\label{lem:actual-background-Hessian-comparison}
The physical functions \(b(V)\) and
\(A(V)-A_0[V]\) have exact local letter
expansions with family prices
\begin{equation}
 C_B(h_{\rm c}+e^{-cB}).
 \label{eq:actual-background-Hessian-comparison}
\end{equation}
The phase derivatives of order \(k\ge2\)
on a common small \(y\)-ball have rooted
analytic prices \(C_B k!r_0^{-k}\).
The loop--Haar derivatives have prices
\(C_B(1+\Lambda)k!r_0^{-k}\).
All statements are uniform in cutoff and
period within a fixed profile.
\end{lemma}
\begin{proof}
Retain the actual direct, row and noise
curvature vector \(\mathcal F(V,y)\) and
the higher-order remainder in
\eqref{eq:microscopic-square-decomposition}.
The replica of a good coarse field has
fine plaquette chords at most \(h_{\rm c}\):
its plaquettes are either trivial or coarse
plaquettes. On a fixed predictor patch,
a local rooted gauge makes its links
\(O(C_Lh_{\rm c})\).
The predictor agrees with the coarse link
for a flat replica and is analytic on this
patch. Hence its mismatch at a good
replica is \(O(C_Lh_{\rm c})\).
The absolute row budgets consequently give
\(\|\mathcal F(V,0)\|_\infty\le C_Lh_{\rm c}\).

Differentiate the exact sum of squares:
\[
 b=\mathcal F_y^T\mathcal F+\mathcal R_y,\qquad
 A=\mathcal F_y^T\mathcal F_y+
       \sum_r\mathcal F_r\mathcal F_{r,yy}
                                    +\mathcal R_{yy}.
\]
The derivative insertion budgets bound the
row and column sums of \(\mathcal F_y\)
and the rooted higher tensors.
The direct Wilson-versus-curvature defect
starts in degree four, the cubic and quartic
terms have their displayed curvature degrees,
and every term in \(\mathcal R_y,\mathcal R_{yy}\)
therefore contains a curvature factor when
needed in the Hessian comparison. These
terms vanish to the asserted orders on a
flat replica.

For clarity, the comparison of the leading
Gram term is also made record by record.
A record is a pair of row slots, or a
direct plaquette, or one fixed predictor
patch. At zero curvature its derivative
is its flat scalar colour matrix with
parallel transports at the endpoints.
Subtract precisely that transported flat
record before taking absolute values.
For a short record, put all its paths and
the coordinate filling paths in its
declared box. In the resulting local gauge
the coarse links are \(O(\ell^2h_{\rm c})\).
Analyticity of the curvature words and of
the fixed-patch predictor gives a difference
bounded by
\[
 C\,(\text{record coefficient})\,
                         h_{\rm c}(1+\ell)^{C_{q,s}} .
\]
The same bound holds for the physical
derivative norm: each left derivative
inserts a direction of size at most
\(h_{\rm c}(1+\ell)^s\).
For a row-pair coefficient the price is
the product of its two row coefficients;
their exponential budgets pay the
polynomial factors.
The predictor record has fixed patch size
and its full frame derivative is retained.
It is not replaced by \(Q\).

Composition of the transported flat
records differs from \(A_0[V]\) by exactly
the product defects of
Lemma~\ref{lem:covariant-product-defect}.
For long records use their bounded word
and derivative prices and retain their
path footprints, as in that lemma.
This proves
\eqref{eq:actual-background-Hessian-comparison}.
The same short/long proof applies to \(b\),
whose flat records are zero.

Finally take a common complex \(y\)-radius
so small that its exponential word cost
uses less than a fixed part of the path
width. The analytic geometric series and
Cauchy estimates give the asserted
\(k!r_0^{-k}\) bounds. Loop vertices carry
their original multiplier \(1+\Lambda\);
the fixed-chart Haar logarithm has the
same bound without \(\Lambda\).
No extensive scalar was bounded as a
single rooted object in this calculation.
\end{proof}

Every short estimate in this proof requires
good plaquettes only in its declared
filling box and halo. Every long estimate
requires only its own read patches.
Thus these are estimates on the healthy
sets of the individual footprints, not
only estimates on globally good fields.

\subsection{Positivity and the exact covariant inverse}

Define \(P_t(V)=P_t^0[V]\) and set
\[
 E_t(V)=1-P_t(V)(A(V)+t1).
\]
Separate the difference into
\[
 P_t(V)\{A_0[V]-A(V)\}
       +\{1-P_t(V)(A_0[V]+t1)\}.
\]
The second bracket is a covariant product
defect; the \(t1\) term has no such defect.
The preceding lemmas give
\begin{equation}
 \|E_t\|_{\rm family}\le
   \frac{C_B(h_{\rm c}+e^{-cB})}{\eta+t},
 \qquad
 \|P_t\|_{\rm family}\le\frac{C_B}{\eta+t}.
 \label{eq:covariant-parametrix-defect}
\end{equation}
The norm here is the row and column local
family norm, with a reduced native reserve.
Eventually its defect bound is less than
\(1/4\), uniformly in \(t\ge0\).

We also need positivity, not just an inverse.
The covariantised \(B_0[V]\) and \(Q_0[V]\)
satisfy
\[
 B_0[V]Q_0[V]=1+O(h_{\rm c}+e^{-cB}),
 \qquad
 B_0[V]^TA(V)B_0[V]=1+O(h_{\rm c}+e^{-cB})
\]
in the ordinary Schur operator norm.
The first identity makes \(B_0[V]\)
invertible. The second is at least
\(\frac34 1\), for sufficiently small
\(h_{\rm c}\) and large \(B\).
Since \(\|B_0[V]\|_{2\to2}\le C\) by its
absolute coefficients, there are constants
\(\eta',K'>0\), independent of cutoff,
period and \(B\), such that
\begin{equation}
 \eta'1\le A(V)\le K'1.
 \label{eq:physical-background-Hessian-gap}
\end{equation}
This argument also covers flat backgrounds
with nontrivial periodic holonomy.
It uses contractible comparison loops on
the cover, not a periodic gauge taking
every link to the identity.

The exact resolvent is
\begin{equation}
 (A(V)+t1)^{-1}
                    =\sum_{n\ge0}E_t(V)^nP_t(V).
 \label{eq:covariant-Neumann-inverse}
\end{equation}
Indeed \(P_t(A+t1)=1-E_t\); the Neumann
series is a left inverse, hence an inverse
on the finite coordinate space.
Every summand retains its ordered intermediate
indices. It is a finite product of local
covariant functions and is therefore a
local covariant function on the union of
its declared footprints. The family
series converges absolutely in the stated
norm by \eqref{eq:covariant-parametrix-defect}.
The same argument works for each fixed
number of physical derivatives.

\subsection{A stationary graph over the physical good set}

The nonlinear stationary equation can now
be solved without a global coarse logarithm.
Write
\[
 \varphi_y(V,y)=b(V)+A(V)y+{\mathcal R_{\rm nl}}(V,y),\qquad
 {\mathcal R_{\rm nl}}(V,0)=0,\quad {\mathcal R_{\rm nl}}_y(V,0)=0.
\]
With \(P_0=P_0(V)\), it becomes
\begin{equation}
 y=-P_0b+E_0y-P_0{\mathcal R_{\rm nl}}(V,y).
 \label{eq:physical-stationary-fixed-point}
\end{equation}
The rooted analytic bounds give
\(\|{\mathcal R_{\rm nl}}(V,y)\|_\infty\le C\|y\|_\infty^2\)
and the corresponding Lipschitz estimate.
Shrink the fixed \(y\)-radius \(r_0\), if
necessary, so that throughout that ball
\[
 \|\varphi_{yy}(V,y)-A(V)\|_{2\to2}
                          \le C\|y\|_\infty\le\eta'/2.
\]
The third-derivative row and column
budgets give the first inequality.
The phase consequently has at most one
critical point in this whole ball:
integrate its positive Hessian along
the segment between two candidate points.
The value bound for \(b\) is \(Ch_{\rm c}\)
on a globally good field; for this value
bound all cover comparison plaquettes
are available and the full exponential
sum is finite.
Thus \eqref{eq:physical-stationary-fixed-point}
is a contraction on a ball of radius
\(C h_{\rm c}\), with constants independent
of volume.

\begin{proposition}[Covariant stationary and Gaussian data]
\label{prop:physical-covariant-stationary}
For sufficiently large nominal coupling
there is a unique small stationary point
\(\rho(V)\) over every physical good
background, with
\[
 \|\rho(V)\|_\infty\le C h_{\rm c}.
\]
It is covariant under \(\mathcal O_u\).
The actual Hessian
\({\mathsf H}(V)=\varphi_{yy}(V,\rho(V))\) has a uniform
positive lower bound. The operators
\({\mathsf H}(V)^{-1}\) and \(W(V)={\mathsf H}(V)^{-1/2}\),
and the stationary graph itself, have
absolutely convergent local covariant
family expansions with physical derivative
prices \(C_B\). The graph price has the
additional factor \(h_{\rm c}\).
These constructions agree with the earlier
chart constructions on their common domain.
\end{proposition}
\begin{proof}
The contraction proves existence and uniqueness.
The equation is covariant, so uniqueness
proves covariance. Its family expansion
is obtained by iterating
\eqref{eq:physical-stationary-fixed-point}
and retaining its rooted trees.
An explicit scalar majorant is
\[
 t\longmapsto C_Bh_{\rm c}
    +C_B(h_{\rm c}+e^{-cB})t
                  +C_B\frac{t^2}{r_0-t}.
\]
After increasing a finite polynomial
allowance, its fixed point lies below
\(2C_Bh_{\rm c}\), and its derivative
there is less than \(1/2\).
This proves convergence of the local
families and all the fixed-order derivative
majorants. Connected unions and native
lengths are retained in every product.
The sharper value bound \(Ch_{\rm c}\)
is the ordinary sup-norm contraction
already proved.

The difference \({\mathsf H}(V)-A(V)\) has operator
norm at most \(C\|\rho(V)\|_\infty\).
It preserves half the gap in
\eqref{eq:physical-background-Hessian-gap}.
Use the same parametrices \(P_t(V)\),
with \({\mathsf H}+t1\) in place of \(A+t1\).
Their defects again satisfy
\eqref{eq:covariant-parametrix-defect},
with a possibly larger \(C_B\).
This constructs the exact resolvent with
the same local families.
Integrating
\[
 W(V)=\frac1\pi\int_0^\infty
                       t^{-1/2}({\mathsf H}(V)+t1)^{-1}\,\dd t
\]
gives its local expansion.
The bound \(C_B/(\eta'+t)\) is integrable
with this weight. Each ordered path
label has the same finite footprint
throughout the integral, so integration
does not introduce a new read region.
Absolute convergence justifies the
interchanges and physical derivatives.

On a small coarse exponential chart both
stationary points solve the same equation
and lie in the fixed \(r_0\)-ball.
The Hessian uniqueness argument above
applies even if the earlier chart estimate
only bounds its point by \(C\|x\|_\infty\).
They therefore
coincide. The Hessian and its positive
square root are then the same operators,
not just operators with comparable bounds.
\end{proof}

\subsection{Calibration transfer beyond a global link chart}

The stationary point must also lie in the
actual intact small-field sector. This
requires the linear calibration, in addition
to the new Hessian estimates.

\begin{lemma}[Physical transfer of the calibrated residual]
\label{lem:physical-calibrated-residual}
For a globally good physical background
define
\[
 \Gamma(V)=\mathcal F(V,\rho(V))
                              -\Phi(V)\mathcal C_{\rm c}(V).
\]
Then, in the maximum row norm,
\begin{equation}
 |\Gamma(V)|_\infty
                     \le C_B(h_{\rm c}^2+e^{-cB}).
 \label{eq:physical-calibrated-residual}
\end{equation}
The constants have the same cutoff and
volume uniformity as the stationary data.
\end{lemma}
\begin{proof}
The vector is covariant at its row anchors.
Its local series has a positive native
length reserve: use the stationary series,
the curvature-word budgets, and the cover
kernels of \(\Phi(V)\).
Every read site of a rooted label is
within \(C_L\ell_\alpha\) of its anchor.
Consequently changing the background
outside a ball of radius \(r\) changes
its rooted value by at most \(C_Be^{-cr}\),
provided both backgrounds lie in the
small-curvature domains of the series.

At the identity \(\Gamma(1)=0\).
Differentiating the actual stationary
equation gives \(D\rho(1)=\mathcal H\).
Theorem~\ref{thm:constructed-calibration}
then gives
\[
 D\Gamma(1)=\mathsf B+\mathsf A\mathcal H
                         -\Phi_0D\mathcal C_{\rm c}(1)=0.
\]
These cancellations concern the sum of
the rooted labels; they are not assumed
for each separate Taylor monomial.

Around a fixed row anchor take a
nonwrapping box of radius \(c_0B\),
with fixed sufficiently small \(c_0>0\).
In its axial gauge the good background
has link logarithms bounded by
\(C B^2h_{\rm c}\).
Extend these logarithms by zero outside
a smaller box, without changing any
links in its central half, to obtain
a globally small link array \(x^{\rm loc}\).
Its maximum norm is at most \(C B^2h_{\rm c}\);
its boundary plaquette size may be
\(C B^2h_{\rm c}\), which is retained here.
The parametrix and contraction construction
is valid also with this larger curvature
bound: \(C_BB^2h_{\rm c}\to0\).
Thus it evaluates the same unmasked
functions on this completion.

On that global small chart the two
vanishing jets and the rooted analytic
bound give
\[
 |\Gamma(e^{x^{\rm loc}})|_\infty
                         \le C_B h_{\rm c}^2.
\]
The zero completion and the locally
gauged original field agree through
a fixed multiple of \(B\).
Their rooted series differ only in
labels of native length at least \(cB\).
The preceding tail bound is \(C_Be^{-cB}\).
Gauge covariance and the isometry of
adjoint transport prove
\eqref{eq:physical-calibrated-residual}.
\end{proof}

The guide bounds
\eqref{eq:row-lift-prices} and
\eqref{eq:coarse-row-good-bound} now imply
the physical version of the earlier
intact-centre conclusion. Choose the
odd block side from its pre-block price
so that \(A_\Phi C_{\rm c}L^{-2}h_{\rm c}\)
uses less than a fixed fraction of
the fine threshold, and choose the
noise reserve \(W\) afterwards.
Since \(h_{\rm c}\le2h\) eventually and
\[
 C_Bh_{\rm c}^2/h\longrightarrow0,\qquad
 C_Be^{-cB}/h\longrightarrow0,
\]
the centre \(U(V,\rho(V))\) has fine
plaquette chords less than \(h/8\)
and noise mismatch less than \(Wh/2\),
after slightly strengthening the fixed
reserves.
Here we used the direct row price
\(\lambda_*>0\) and the small-chord
equivalence of curvature and chord.
That equivalence applies because the
replica already has small plaquettes
and \(\|\rho(V)\|_\infty\le Ch_{\rm c}\).

The weighted row bound for \(W(V)\)
gives
\(\|gW(V)\xi^R\|_\infty\le CgR\).
Its ordinary weighted matrix value bound
is uniform in \(B\): the actual Hessian
has the fixed spectral gap and weighted
word bound, so the resolvent argument of
Lemma~\ref{lem:local-gaussian-whitening}
applies with constants depending only
on the fixed profile and block.
Increase the common \(C_W\) in
\eqref{eq:white-radius-auxiliary} to
cover both its \(\ell^2\) and
\(\ell^\infty\) operator norms.
With \eqref{eq:white-radius-auxiliary}
and a sufficiently small fixed \(d_0\),
this uses less than the remaining
fine and noise reserves, on every
physical volume. The fixed-patch
predictor Lipschitz bound is used
for the noise term.
Thus every projected configuration in
the reference lies in the actual
intact sector, and all its branch
likelihoods are the original regular
ones. This is an inclusion statement,
not a factorisation of the central ball.

\section{Physical extensions of the interacting families}
\label{sec:physical-interacting-extensions}

\subsection{Local definitions before a background is chosen}

The preceding construction retains a useful
distinction. Its kernels are global sums on
a globally good field, but each summand is
a function of a specified finite footprint.
We construct the latter functions first.

A primitive record is a direct or path
curvature expression, a fixed predictor
patch, or an adjoint path transport.
It is gauge covariant at its declared
Lie algebra slots. Its regular branch
is smooth on an open neighbourhood of
the good set of its read footprint.
If necessary, enlarge that footprint
by the fixed branch-test halo.
Choose a fixed small branch radius
\(h_*>0\), independent of the nominal
coupling, with all regular tests safely
inside their uniqueness domains.
For definiteness define the record by
this smooth formula where its relevant
plaquette chords are less than \(h_*\)
and these branch reserves hold.
Set it to zero elsewhere.
This is a Borel gauge-covariant extension.
It agrees with the original formula on
a neighbourhood of the closed healthy
set at threshold \(h_{\rm c}\).
Eventually \(2h_{\rm c}<h_*\), which
gives that neighbourhood; no derivative of its guard is taken
in the physical activity norm.

For a paired short defect its footprint
includes the filling box used in its
estimate. For a long defect it contains
only its actual paths and fixed branch
halos. Products use the union of these
footprints. Thus every finite tree or
chain in the parametrix and stationary
expansions gives a raw local function
defined before the background is tested.
Its physical good-set bound follows
from exactly the same local calculation.

Colour indices at a vertex are retained
as a complete tensor. Under a physical
gauge transformation, each internal
index is acted on by one orthogonal
adjoint matrix and its transpose.
They cancel in a contraction. An external
fine Gaussian vector transforms by
\(\mathcal O_u\).
The radial projection and its Gaussian
law commute with these independent
orthogonal transformations at the free
links. Consequently each fully contracted
Gaussian letter is a gauge-invariant
raw local scalar function.

This supplies compatibility between local
definitions. The equality of their sums
on a globally good field comes from the
exact Neumann and stationary equations.
We do not evaluate different Taylor
polynomials in independently chosen
gauges and then infer equality of the
resulting sums from global conjugation
invariance.

\subsection{The same connected logarithm in the physical norm}

For a physical good \(V\), define
\(\mathcal Z_R(g,V)\) with the actual
\(\rho(V)\), \(W(V)\), phase and
loop--Haar factor. These data agree
with Section~\ref{sec:full-interacting-remainder}
on its chart domain.
Substitute their local family expansions
in \eqref{eq:full-interaction-series}.
The coefficients now are the raw physical
functions just constructed. Their Gaussian
variables are still the fixed independent
\(\xi_e\); differentiating \(V\) does not
differentiate a Gaussian law.

\begin{theorem}[Gauge-invariant physical remainder families]
\label{thm:physical-interacting-remainder}
Fix the group, block and principal path
profiles, derivative profile \(q,s\), and
covering exponent \(\kappa\).
For sufficiently large nominal coupling,
on periods satisfying
\eqref{eq:good-background-period-reserve},
the reference logarithm on physical good
fields has an exact representation
\begin{equation}
 \log\mathcal Z_R(g,V)
       =g^2\sum_\alpha l_\alpha(V)
                           +\sum_\alpha r_\alpha(g,V).
 \label{eq:physical-interacting-family-identity}
\end{equation}
The functions are gauge-invariant raw
local functions, smooth near the good
sets of their declared footprints.
They satisfy the physical bounds
\begin{align}
 \|l\|_{\mathcal E,\kappa}
                      &\le C_B(1+\Lambda)^2,\notag\\
 \|r\|_{\mathcal E,\kappa}
   &\le C_B\left\{
       g^4H_R^4+
        g^2(1+\Lambda)^2(1+R)^6e^{-cR^2}\right\}.
 \label{eq:physical-interacting-family-prices}
\end{align}
The first family sums to the full Gaussian
polynomial coefficient \(L_2(V)\).
The constants are independent of ultraviolet
cutoff and physical volume within the
stated profile and period reserve.
Both families retain a positive native
length reserve.
\end{theorem}
\begin{proof}
The phase tensors at \(\rho(V)\) and the
entries of \(W(V)\) have local covariant
families by
Proposition~\ref{prop:physical-covariant-stationary}.
Their connected products give primitive
interaction letters with physical incidence
price
\[
 C_B |g|\{R^3+(1+\Lambda)R\}.
\]
The proof is the two geometric sums in
Lemma~\ref{lem:full-interacting-letters};
the local physical derivatives have
replaced its chart derivatives.
All the guards, paths and filling boxes
are included in the read footprints.
Each estimate therefore holds on that
footprint's own healthy set.

Products of Mayer letters, their Gaussian
expectations, and the ordered logarithm
series are now formed with exactly the
same labels as in
\eqref{eq:clipped-polymer-gas} and
\eqref{eq:interacting-Ursell}.
Disjoint declared footprints again read
disjoint independent Gaussian vectors.
The physical product norm and derivative
rebasing are the ones used in
Lemma~\ref{lem:rooted-interacting-log}.
Its proof applies verbatim to these
norms: the local support cardinality is
at most \(D_B n_B\), and every sum is
rooted before the volume is chosen.
Thus the logarithm converges for the
same type of radius
\(|g|<c/(C_BH_R)\).

Every contracted polymer is a
gauge-invariant scalar by the preceding
construction. Cauchy subtraction of its
second \(g\)-coefficient preserves that
invariance and its footprint.
The parity argument gives the full
fourth-order family bound.
Replacing the projected second coefficient
by \(L_2(V)\) uses only Gaussian moments
of degree at most six, with the same
local tail estimate as
\eqref{eq:clipped-second-coefficient-error}.
This proves the exact identity and both
prices.

For the native reserve, attach lengths
additively in a matrix chain, a stationary
tree, a primitive letter, and an ordered
logarithm cluster. Choose a smaller
positive exponent at the beginning.
The same convergent majorants then
include the additional exponential
factor at every step. Enlarging a short
footprint to its declared box costs a
fixed number of \(B\)-balls; long
footprints retain their path covers.
No uncontrolled filled-box price is
introduced.
\end{proof}

The identity is an identity of unmasked
functions on the globally good set.
The raw functions in it are also defined
at fields with bad sites elsewhere.
In an effective density they are used
with precisely the good-footprint masks
of Section~\ref{sec:activities}.
Their masked reassembly and the corresponding
forced corrections are algebraic operations
on the same density.

\subsection{Removing the interacting vacuum without a global gauge}

The first bound in
\eqref{eq:physical-interacting-family-prices}
alone would leave an order-\(g^2\) value
budget. Its flat value belongs in the
vacuum constant. The difference has a
small-curvature gain which we now prove.

\begin{lemma}[Subtracting the flat scalar in the physical norm]
\label{lem:physical-interacting-vacuum-subtraction}
The leading physical family can be
relabelled, without changing its sum,
so that
\begin{equation}
 \bigl\|(l_\alpha-l_\alpha(1))_\alpha\bigr\|
                         _{\mathcal E,\kappa}
 \le C_B(1+\Lambda)^2
                         (h_{\rm c}+e^{-cB}).
 \label{eq:physical-leading-vacuum-subtraction}
\end{equation}
Every relabelled footprint is part of
its declared physical good test.
\end{lemma}
\begin{proof}
Split the already convergent native
length families at
\(\ell_\alpha\le B/C_*\), where \(C_*\)
is a fixed geometric constant.
For a short label enlarge its footprint
to the rooted box containing all its
declared paths, short filling boxes and
halos. Its radius is at most
\(C_L(1+\ell_\alpha)\le C_L B/C_*\).
It is covered by a fixed number of
\(B\)-balls and is nonwrapping by
\eqref{eq:good-background-period-reserve}.

The local axial gauge in
Lemma~\ref{lem:local-gauge-chart} maps a
healthy field on this box into link
coordinates of size
\(C(1+\ell_\alpha)^2h_{\rm c}\).
This is in a common analytic chart,
since \(B\) is logarithmic and
\(B^p h_{\rm c}\to0\) for every fixed
power \(p\).
The scalar is gauge invariant as a
physical function. Its value in this
local gauge therefore equals its
original value.
Subtract its value at the identity.
The mean-value theorem and the analytic
coefficient reserve give the bound
\[
 C(\text{label price})(1+\ell_\alpha)^p h_{\rm c}.
\]
For positive-order physical derivatives,
word insertion and the chain rule give
the same bound with a possibly larger
fixed power. Indeed each derivative
of the local gauge chart contributes
a factor \(h_{\rm c}\) times a power
of \(1+\ell_\alpha\).
Its constant subtraction has no such
derivatives.
The unmasked local formulas admit
analytic continuation to a common
small link chart. Their majorants
come from the convergent parametrix
and stationary series on that chart.
The guarded functions agree with
these formulas on the healthy
arguments under consideration;
no claim about analyticity across
a guard boundary is used.

The additional incidence of an enlarged
box costs at most its cardinality,
a polynomial in \(1+\ell_\alpha\).
All these polynomials are paid by
the native exponential reserve.
This proves the short-family bound.
For long labels make no enlargement.
Their total physical family price,
including their flat values, is at
most \(C_B(1+\Lambda)^2e^{-cB}\)
by the retained native reserve.
These labels include possible winding
contributions, which are not declared
zero at a flat background.
\end{proof}

Let \(n_c\) be the number of coarse
sites. The signed scalar in the
effective action is
\[
 c_{\rm int}=-\log\mathcal Z_R(g,1).
\]
From
\eqref{eq:physical-interacting-family-identity}
we obtain the exact good-field identity
\begin{equation}
 -\log\mathcal Z_R(g,V)
                     =c_{\rm int}+\sum_\alpha f_\alpha(V),
 \qquad f_\alpha(1)=0,
 \label{eq:physical-interacting-vacuum-identity}
\end{equation}
where the nonconstant family has price
\begin{align}
 H_{\rm int}\le C_B\{&
  g^2(1+\Lambda)^2(h_{\rm c}+e^{-cB})
                  +g^4H_R^4\notag\\
 &+g^2(1+\Lambda)^2(1+R)^6e^{-cR^2}\}.
 \label{eq:physical-interacting-nonconstant-price}
\end{align}
The scalar is extensive only in the
expected sense:
\[
 |c_{\rm int}|/n_c
       \le C_B\{g^2(1+\Lambda)^2+g^4H_R^4\}.
\]
This follows by summing the per-anchor
prices after, rather than before,
proving them.

At the physical coupling, with the
prescribed thresholds and loop allowances,
\begin{equation}
 H_{\rm int}\le
                  C_{\nu,L}\widehat\beta^{-3/2+\nu}.
 \label{eq:physical-interacting-final-rate}
\end{equation}
Indeed \(g^2h_{\rm c}\) is of order
\(\widehat\beta^{-3/2}\) times a subpower
factor, the fourth-order term is
\(\widehat\beta^{-2+\nu}\), and the
two exponential terms are smaller than
every power.
This gives a reserve beyond
\(\widehat\beta^{-23/20}\).

The affine quadratic projection and
the clover transfer of Section
\ref{sec:coupling-reset} can now be
applied to the actual functions
\(f_\alpha\).
In particular their coupling correction
obeys
\[
 h_{\rm c}^2|\Delta_{\rm int}|
                                      \le C_BH_{\rm int}.
\]
The exact vacuum reset uses
\(c_{\rm int}\) with the sign in
\eqref{eq:physical-interacting-vacuum-identity};
the new vacuum parameter is decreased
by \(c_{\rm int}/n_c\).
All scalar, affine and mask-transfer
terms remain those in the exact reset
identity. The estimate here concerns
the newly generated interacting
contribution; transported incoming
activities retain their own budgets.

\subsection{The geometric reserve in one fixed Wilson family}

It remains to check that
\eqref{eq:good-background-period-reserve}
is compatible with the family chosen
before any infrared subsequence.
At scale \(m\), write \(s_m=\log {\varepsilon}_m^{-1}\).
The lattice period and nominal coupling
have the forms
\[
 \ell_{m,M}=Me^{s_m},
 \qquad
 \widehat\beta_m=\overline\beta+b_0s_m .
\]
For a fixed geometric constant the function
\[
 s\longmapsto
 e^{-s}\left\lceil
          \{\log(\overline\beta+b_0s)\}^2\right\rceil
\]
is bounded on \(s\ge0\) and tends to zero.
Choose the smallest admissible side length \(M_{\min}\)
above the required fixed multiple of this bound. Choose the
remaining admissible integer side lengths above it, including
the reflection divisibility condition imposed below.
This guarantees the reserve at that
profile, uniformly in \(M\).

When the geometric constants vary with
the profile \(j\), enlarge its entry
threshold \(\mathcal B_j\) in
Corollary~\ref{cor:linear-profile-clock}.
Since \(e^s/B(\overline\beta+b_0s)\to\infty\),
there is a finite threshold after which
the reserve for that profile holds
at every later scale.
The slow clock can include this
threshold as well as the analytic
ones. All these choices precede \(n\)
and the requested volume subsequence.
Thus the uniformity statements concern
the same Wilson family.
If a step uses the fine nominal coupling
in its radius, include the fixed shift
\(b_0\log L_j\) in the logarithm's
argument at that profile. The same
boundedness and entry-threshold argument
then applies; the fine/coarse shift
has not been omitted.

\subsection{Relation to the exact physical integral}

The complete interacting remainder of
the actual central integral, with its
original normalisation, was identified
on auxiliary volumes in
Corollary~\ref{cor:complete-central-interacting-error}.
The constructions here give compatible
gauge-invariant physical local functions
for its bounded-coordinate reference
and its full fourth-order remainder.
On the common chart their stationary
point, Hessian, and Gaussian variable
change are exactly the earlier ones.

The full physical effective density
also includes the
Gaussian-coordinate failure sectors,
the original fine-field indicators,
and the incoming regular and forced
activities. Their component estimates
and exact masked reassembly are proved in
Sections~\ref{sec:finite-range-gaussian}--%
\ref{sec:physical-reassembly}; packet integrals and
regular derivatives are treated in
Sections~\ref{sec:hole-packet-integrals}--%
\ref{sec:integrated-regular-derivatives}.
The regular estimates just proved do
not remove any of those terms.
In particular a small error for one
auxiliary central ball has not been
declared an error independent of volume
for a single ball in an unrestricted
physical system.
\clearpage
\section{An exact finite-range Gaussian representation}
\label{sec:finite-range-gaussian}

The bounded white-coordinate expansion is useful for the
analytic interaction. The incoming forced activities, however,
are merely bounded Borel functions. We must evaluate these
functions on physical links without Taylor expanding them.
We therefore construct a different Gaussian representation
of the same finite-dimensional integral. Its coordinates
are the physical fluctuation coordinates themselves, and its
covariance has finite range. All changes in the Gaussian
density are retained in the integrand.

The physical fluctuation integral is
\eqref{eq:actual-block-integral}; compare the organisation of
the fluctuation integral in \cite[Sect.~2, eq.~(2.12),
printed p.~268]{Bal87}. This comparison concerns the physical
integral and its normalisation. The selected-window covariance
below is constructed directly from local inverses.
The matrix and measure calculations needed below are given
in full. A small Gaussian error will always mean a bound
per coordinate or per connected component, never a bound
on an extensive determinant without its volume factor.

\subsection{Holes, active links, and the exact physical integrand}

Keep the same block step, forest, branch choices, and
principal paths. Let \(I\) be the free fine links in tree
gauge, and put \(d_G^{\rm Lie}=\dim_{\R}\mathfrak g\).
An orthonormal basis of \(\mathfrak g\) identifies the
coordinate set with
\[
 \mathcal I=I\times\{1,\ldots,d_G^{\rm Lie}\}.
\]
All Euclidean norms and determinants in this section use
this counting-coordinate convention. In particular, no
factor equal to the reciprocal of the period is inserted.

A \emph{sector} is a fixed choice of values of the
discrete local tests used to partition the physical
integral. Its weight, denoted by \(\omega_\sigma(V,U)\),
includes these tests; it need not include the incoming
activities. Sector weights may be hard indicators.
An exact sector partition means
\begin{equation}
 0\le\omega_\sigma\le1,\qquad
 \sum_\sigma\omega_\sigma(V,U)=1 .
 \label{eq:sector-partition-finite-range}
\end{equation}
For example, a finite collection of tests taking values
in \([0,1]\) gives such a partition by selecting either
the test or its complement at every address.
Tests depending on the coarse field are evaluated at
the actual \(V\), and tests depending on the fine field
are evaluated at the actual \(U\).

Let \(\mathcal H_\sigma\) be the coarse hole set of the
sector. It contains all coarse bad flags and the
additional marked sites prescribed by that sector.
The notation \(H_r\) for a periodic \(r\)-neighbourhood
is the one used in Section~\ref{sec:activity-norms}.
For a fine link \(e\), let \(\pi(e)\) be the set of
the coarse parents of its two endpoints.
Define the active free links by
\[
 I_\sigma=\{e\in I:
       \pi(e)\cap H_{10B}(\mathcal H_\sigma)=\varnothing\},
 \qquad J_\sigma=I\setminus I_\sigma .
\]
The links in \(J_\sigma\) will be integrated with their
original probability Haar measures. Write their values
as \(Y=(Y_e)_{e\in J_\sigma}\).
Tree links remain equal to the identity.

For active coordinates \(a_e\in\mathfrak g\), define
\begin{equation}
 \mathcal U_\sigma(V,Y,a)_e=
 \begin{cases}
  \exp(a_e)\operatorname{rep}(V)_e,&e\in I_\sigma,\\
  Y_e,&e\in J_\sigma,\\
  1,&e\text{ is a tree link}.
 \end{cases}
 \label{eq:finite-range-physical-installation}
\end{equation}
Thus changing an active coordinate changes that physical
link, rather than a nonlocal linear combination of links.

Fix an injective exponential chart
\(\exp:\{|a|<r_*\}\to G\).
On this chart write probability Haar measure as
\[
 dh_G(\exp(a))=\zeta_G j_G(a)\,da,\qquad
                  j_G(0)=1,\quad \zeta_G>0 .
\]
Both \(j_G\) and \(\zeta_G\) are the same as in the
earlier physical integral. The sector formula below
applies when \(\omega_\sigma\) is supported in this
chart on every active link. A chart partition can
always enforce this property: expand the product
of a local chart cutoff and its complement, and put
all links carrying a complement into the frozen set.
If a frozen neighbourhood is enlarged, its original
chart-cutoff factors are kept in \(\omega_\sigma\).
Enlargement does not replace any test by one.

For a fixed branch \(t\), abbreviate
\begin{equation}
 \mathcal B_\sigma(V,Y,a)=
 \omega_\sigma(V,\mathcal U_\sigma)
 \lambda_t(\mathcal U_\sigma)
 e^{-I_f(\mathcal U_\sigma)-E^\#(\mathcal U_\sigma)}
 Z_K(\mathcal U_\sigma)
 \prod_{e\in I_\sigma}j_G(a_e).
 \label{eq:finite-range-full-multiplier}
\end{equation}
Every factor on the right is evaluated on the single
field \(\mathcal U_\sigma(V,Y,a)\). The indicators
inside \(E^\#\), the admissibility conditions in
\(Z_K\), and all branch weights are retained.
The product is zero off the active chart support.
For each finite torus it is absolutely integrable
against the physical measure, by
Proposition~\ref{prop:activity-convergence}.
This assertion permits signed incoming activities.

\subsection{The operated affine rows}

Put
\[
 R_B=\lfloor B/100000\rfloor,\qquad
 r_B=R_B+R_{\rm pred}+2 ,
\]
where \(R_{\rm pred}\) is the fixed predictor reach
at this profile. Increasing its coupling threshold,
we require \(4r_B+10<B/4\). The native word cut
at \(R_B\) counts the original fine walk length,
including its final plaquette. No cut is made using
the shorter periodic distance.

The finite row list contains direct curvature rows,
cut principal curvature rows, and noise rows in
each prescribed orientation. Symmetry weights are
absorbed through their positive square roots.
At the identity, its linear maps are the curl rows
and the flat derivative of the actual predictor.
The latter includes its frame derivative, as in
Section~\ref{sec:all-scale-linear}; it is not replaced
by the scalar line average alone.

Here is the transport convention defining the row
operator away from the identity. Pull a coordinate
array into the orientation of the row, with the
usual negative-link sign and adjoint action for
reversed links. Transport the colour at each
coarse parent to the row anchor along the fixed
coordinate path. Apply the flat cut row formula
to this transported array. Denote the resulting
map by \(d_r(V)\). This definition uses only finite
paths and adjoint actions; it is defined even when
the background is rough.

Only rows whose coarse anchors lie outside
\(H_{8B}(\mathcal H_\sigma)\) are operated.
The matrix \(D=D_\sigma(V)\) is
\begin{equation}
 D_{re}=\mathbf1_{\{r\ {\rm operated}\}}
        \mathbf1_{\{e\in I_\sigma\}}\,(d_r(V))_e .
 \label{eq:operated-affine-D}
\end{equation}
Colour indices are implicit in this formula.
Its potential stencil, rather than only the entries
which happen to be nonzero at \(V\), belongs to
the declared read set.

The affine source retains the actual finite-branch
row values. To specify it, choose a bounded local
extension \(s_e(V,Y)\) of the relative logarithm of
each frozen link, equal to that logarithm when its
norm is at most \(u\), zero on active and tree links,
and of norm at most \(2u\). We use the scale
\(u=B h\); its extension is obtained by a radial
smooth cutoff inside the fixed chart. This is a
definition of the affine source only.
The physical link in
\eqref{eq:finite-range-physical-installation}
is still \(Y_e\).

For an operated row, take its actual cut energetic
row at the replica field, subtract the prescribed
coarse guide row, and add the full transported flat
row applied to \(s\). Call this value \(\alpha_r\).
Put \(\alpha_r=0\) on the other rows, and set
\begin{equation}
 b=D^T\alpha,\qquad
 \mathsf M=D^TD+P_{\rm in}.
 \label{eq:operated-precision-source}
\end{equation}
Here \(P_{\rm in}\) is the diagonal projection onto
inactive real coordinates. The coarse guide is
the cut row lift of the prescribed coarse curvature
rows, with their original tests. The subtraction
is the same for every factor in a sector.
All differences between these affine rows and the
actual nonlinear phase will remain explicitly below.
The exact identity does not assume that the affine
rows are derivatives of that phase at \(V\).

The matrix \(D\), and consequently \(\mathsf M\),
depend on \(V\), the forest, and the hole geometry.
They do not depend on \(Y\) or on a resampled branch
choice. The source \(b\) can depend on the actual
frozen variables and branch choices.

\begin{lemma}[Uniform precision in the operated coordinates]
\label{lem:operated-precision-gap}
Fix the group, block and path profile.
There are \(\eta,C,\theta>0\), independent of the
hole geometry, cutoff, and period, such that,
for sufficiently large nominal coupling,
\begin{equation}
 \eta\,1\le\mathsf M\le C\,1,\qquad
                    |\mathsf M|_\theta\le C .
 \label{eq:operated-precision-gap}
\end{equation}
The norm \(|\cdot|_\theta\) is the maximum of the
weighted row and column sums from
\eqref{eq:weighted-schur-norm}, using the coarse
home of a fine coordinate.
The period reserve is \(\ell_{\rm c}\ge100(1+B)\).
Coarse plaquettes outside the hole set are tested
at the prescribed threshold \(h_{\rm c}\).
\end{lemma}
\begin{proof}
Adjoint transports have norm one. The exponential
word prices, including colour multiplicity, bound
both the row and column sums of the cut row matrix
uniformly in \(R_B\). Products give the stated
Schur upper bound for \(D^TD+P_{\rm in}\).
The Schur test gives its operator upper bound.

We spell out why deletion near a hole does not
remove the lower bound. The flat block Poincare
argument of Proposition~\ref{prop:actual-flat-gap}
has a local version: for a free link \(e\), there
is a fixed neighbourhood of radius \(R_*\), and
a constant \(C_*\), such that every array vanishing
on tree links obeys
\[
 |v_e|\le C_*\max_{r\ {\rm in\ that\ neighbourhood}}
                              |d_r(1)v|.
\]
Only direct curl and noise rows are needed here.
Apply \eqref{eq:coercivity-path-homotopy}--%
\eqref{eq:coercivity-noise-root-identity} on the enlarged
patch. The direct rows give its curls, multiplied by
the fixed positive square root of the direct coefficient.
The noise rows give \(Q_0v\), multiplied by \(\sqrt a\).
Thus the displayed \(C_*\) can be chosen as the product
of the finite path-reconstruction constant, the number
of rows in that patch and
\(\max\{1,\lambda^{-1/2},a^{-1/2}\}\).
All are fixed-profile constants. In particular the
block-root difference is supplied by an actual noise
row; it is not estimated by a missing boundary value.

For active \(e\), this whole patch lies outside
the \(8B\) hole neighbourhood once \(B\) is large.
Its required rows have therefore not been deleted.
In a local axial gauge, the coarse link logarithms
on the patch have size at most \(C_*h_{\rm c}\).
The local path comparison proved in
Lemma~\ref{lem:local-gauge-chart}, followed by
the finite row formula, gives
\[
 |v_e|\le C_*\max_{r\ {\rm near}\ e}|(Dv)_r|
              +C_*h_{\rm c}\max_{f\ {\rm near}\ e}|v_f|
\]
for arrays supported on active coordinates.
The cut principal rows can be ignored in this
lower bound; the direct and noise rows are
unchanged by their word cut.

Square this inequality, use \((a+b)^2\le2a^2+2b^2\),
and sum over active \(e\).
The number of links and rows in the fixed patch,
and the number of patches containing one fixed
link or row, are bounded by a constant depending
only on the block. We obtain
\[
\|v\|_2^2\le C_1\|Dv\|_2^2
                        +C_2h_{\rm c}^2\|v\|_2^2 .
 \]
For explicit dependence, let \(\nu_r\) be the maximum
number of these link patches containing one row, and
\(\nu_e\) the analogous number for a coordinate. If the
last local perturbation inequality uses constants
\(C_*,C_*'\) in front of its row and coordinate maxima,
one can take \(C_1=2C_*^2\nu_r\) and
\(C_2=2(C_*')^2\nu_e\). The maxima squared are bounded
by the corresponding sums of squares before summing
over the link patches. The numbers \(\nu_r,\nu_e\)
depend only on their fixed radius and coordinate types,
even next to a hole or a periodic seam.
Choose the coupling threshold so that
\(C_2h_{\rm c}^2\le1/2\). Thus
\(\|Dv\|_2^2\ge(2C_1)^{-1}\|v\|_2^2\)
on the active subspace.

For a general array write \(v=v_{\rm ac}+v_{\rm in}\).
Their supports are disjoint, \(Dv=Dv_{\rm ac}\), and
\[
 \langle v,\mathsf Mv\rangle
           =\|Dv_{\rm ac}\|_2^2+\|v_{\rm in}\|_2^2 .
\]
This proves the lower bound with
\(\eta=\min\{1,(2C_1)^{-1}\}\).
It also covers the case with no active coordinates.
Every estimate used a fixed patch or an absolutely
summed word budget. None used the total volume.
\end{proof}

\subsection{Selected inverse columns and their error}

For each real coordinate \(j\) choose the prescribed
window
\[
 \Omega_j=\{(e,a):\pi(e)\cap H_{\lfloor B/4\rfloor}
                         (\pi(e_j))\ne\varnothing\}.
\]
It contains \(j\). If \(d(i,j)\) denotes the maximum
distance between the coarse homes, then
\[
 d(i,j)\le\lfloor B/4\rfloor\ \Longrightarrow\
 i\in\Omega_j,\qquad
 i\in\Omega_j\ \Longrightarrow\ d(i,j)\le B/4+2 .
 \]
The windows include every colour at every selected
link. Their geometry is fixed before evaluating a
field. Define
\begin{equation}
 J_{ij}=\mathbf1_{\{i\in\Omega_j\}}
             (\mathsf M_{\Omega_j,\Omega_j}^{-1})_{ij},
 \quad E=\mathsf MJ-1,\quad
 C_{\rm sel}=\tfrac12(J+J^T),\quad S=\mathsf MC_{\rm sel}-1,
 \quad \mu=-J^Tb .
 \label{eq:selected-window-data}
\end{equation}
The subscripts of a principal inverse use its original
coordinate labels. It is the principal submatrix of
the full operated precision; boundary rows of \(D\)
are not removed when forming it.
In particular,
\[
 \mathsf M_{\Omega,\Omega}
       =D_{*,\Omega}^{\,T}D_{*,\Omega}
                        +(P_{\rm in})_{\Omega,\Omega}.
\]
The star here means all operated rows.
Replacing it by rows anchored in \(\Omega\)
would define a different principal matrix.

\begin{proposition}[Finite range and a uniform Gaussian law]
\label{prop:selected-finite-range-law}
For some reduced \(\theta_1>0\),
\begin{align}
 |J|_{\theta_1}+|C_{\rm sel}|_{\theta_1}
                                      &\le C, \notag\\
 |E|_{\theta_1}
 +|J-\mathsf M^{-1}|_{\theta_1}
 +|C_{\rm sel}-\mathsf M^{-1}|_{\theta_1}
 +|S|_{\theta_1}&\le Ce^{-cB}=:\epsilon_B .
 \label{eq:selected-window-error}
\end{align}
For sufficiently large \(B\),
\begin{equation}
 (2C)^{-1}1\le C_{\rm sel}\le C_1\,1,\qquad
                         C_{\rm sel}^{-1}\ge C_1^{-1}1.
 \label{eq:selected-covariance-gap}
\end{equation}
Moreover \((C_{\rm sel})_{ij}=0\) if \(d(i,j)>B/4+2\).
The corresponding Gaussian is one normalised joint
Gaussian law on all coordinates.
\end{proposition}
\begin{proof}
The full and principal inverse bounds are exactly the
two versions of Lemma~\ref{lem:actual-combes-thomas},
with the gap and Schur constants in
\eqref{eq:operated-precision-gap}. They give a common
envelope \(Ce^{-\nu d(i,j)}\) on every window, using
the original coarse-home metric. Summing this envelope
over the at most \(C_L(1+r)^4\) coordinates in a coarse
ball gives the weighted bounds for \(\mathsf M^{-1}\)
and \(J\). This sum also controls rows of \(J\): the
envelope is common to all columns, although their
windows differ. There is no extra number-of-windows factor.

For \(i\in\Omega_j\), the \(ij\) entry of
\(\mathsf MJ\) is the corresponding entry of
\(\mathsf M_{\Omega_j,\Omega_j}
 \mathsf M_{\Omega_j,\Omega_j}^{-1}\).
Indeed the \(j\)-th selected inverse column is
zero outside \(\Omega_j\). Thus \(E_{ij}=0\)
through the inner radius of that window.
The unreduced weighted norm of \(E\) is bounded
by the product bound for \(\mathsf M J-1\).
Reducing the exponent spends the inner radius
and proves \(|E|_{\theta_1}\le Ce^{-cB}\).

The exact identity
\[
 J-\mathsf M^{-1}=\mathsf M^{-1}E
\]
gives the next estimate. Symmetrisation gives
the covariance estimate, and
\(S=\mathsf M(C_{\rm sel}-\mathsf M^{-1})\)
gives the last one.
In operator norm \(\mathsf M^{-1}\ge C^{-1}1\).
For \(\epsilon_B\le(2C)^{-1}\), the covariance
differs from it by at most this amount.
This proves its positive lower bound; its upper
bound follows from the Schur bound. Inverting
the upper bound proves the precision reserve.
Finally both \(J_{ij}\) and \(J_{ji}\) vanish
beyond the stated range.
\end{proof}

The same argument applies uniformly to
\[
 \mathsf M_t=(1-t)1+t\mathsf M,\qquad 0\le t\le1.
\]
Use the lower bound \(\min\{1,\eta\}\) and the
upper Schur bound \(C+2\).
This observation will justify the determinant
flow below on its whole interval.

% !TeX root = 4DYM_notes.tex
\subsection{Derivatives of the same windowed covariance}
\label{sec:selected-window-jets}

The estimates for inverse columns must hold after
background differentiation. In particular a derivative
must not replace the selected covariance by the full
inverse. We prove all the needed estimates with the
windows and operated rows held fixed. These are the
unmasked germs behind the original physical tests.
The inverse-decay argument is the finite-matrix version
of the Combes--Thomas method discussed in
\cite[Sect.~2]{DST23}. The window construction is also
described in \cite[Unit III.5, Proposition III.5.13,
PDF p.~163]{ArmstrongNotes}.
Only the elementary matrix identities below are needed
for its derivative estimates.

Let \(x_\ell\) be the left exponential coordinate of a
coarse link, including its colour index. Extend the
rooted norm \eqref{eq:rooted-tensor-norm} to a matrix
derivative by rooting at either matrix index:
\begin{equation}
 \|D^q A\|_\theta=
 \max\left\{
 \sup_i\sum_{j,\boldsymbol\ell}
  e^{\theta\mathsf t(i,j,\boldsymbol\ell)}
       |h_{\rm c}^q\partial_{\boldsymbol\ell}A_{ij}|,
 \sup_j\sum_{i,\boldsymbol\ell}
  e^{\theta\mathsf t(i,j,\boldsymbol\ell)}
       |h_{\rm c}^q\partial_{\boldsymbol\ell}A_{ij}|
                         \right\}.
 \label{eq:selected-matrix-jet-norm}
\end{equation}
Here \(\boldsymbol\ell=(\ell_1,\ldots,\ell_q)\), repeated
indices are allowed, and \(\partial_{\boldsymbol\ell}\)
is the indicated ordered ordinary chart derivative.
The threshold is held fixed. At order zero this is the
original weighted row and column norm. Distances are in
coarse units, including the attached homes of fine
indices. The factor \(h_{\rm c}^q\) is the same scaling
used in the regular activity norms.

These bounds also control ordered left derivatives.
In a fixed left exponential chart such a derivative
is an ordinary first derivative with analytic bounded
coefficients. A product of \(q\) left derivatives is
a finite sum of chart derivatives of orders \(s\le q\),
with coefficients bounded at the fixed chart radius.
After multiplication by \(h_{\rm c}^q\), its order-\(s\)
term has the additional factor \(h_{\rm c}^{q-s}\le1\).
Derivatives of a chart coefficient use the same link,
so their repeated indices add no spatial length to
the rooted weight. Thus the conversion is uniform
in the number of links.

For fixed \(q\), the operated row formula gives
\begin{equation}
             \|D^s\mathsf M\|_{\theta_0}\le K_s,
                          \qquad 0\le s\le q,
 \label{eq:selected-precision-jet-input}
\end{equation}
on the healthy coarse germs, at one sufficiently small
\(\theta_0>0\). We spell out the counting in this input.
A transported row is a sum of its flat coefficients
times adjoint transports along its declared paths.
A derivative chooses a position on one of those paths
and inserts a bounded adjoint generator there. At order
\(s\), summing all input indices and choices costs at
most \(C_s(1+\ell)^s\) on a path of length \(\ell\).
The tree joining its row, column and differentiated
positions has length at most a fixed multiple of the
declared path length. Choose \(\theta_0\) so this costs
less than half its native exponential weight. Then
\[
 (1+\ell)^s e^{-\sigma\ell/2}
       \le \sup_{u\ge0}(1+u)^s e^{-\sigma u/2}<\infty.
\]
The word budget sums the remaining coefficients.
Deleting terms by the fixed cut and the fixed operated
and active indicators can only decrease this absolute
majorant. The same argument applies to the transpose
row sum. The finite product rule for \(D^TD\), with
Lemma~\ref{lem:rooted-kernel-composition}, proves
\eqref{eq:selected-precision-jet-input}; the fixed
projection \(P_{\rm in}\) has no derivatives.
The constants depend on \(G,L,J\), the path budgets,
their widths and \(q\), but not on the word cut, the
holes, the cutoff or the number of blocks.

\begin{lemma}[Differentiated inverse columns and exponential residual]
\label{lem:selected-window-jets}
For every fixed \(q\) there are \(\theta,c,C_q>0\),
chosen from the data in
\eqref{eq:selected-precision-jet-input} and the precision
gap, such that for \(0\le s\le q\),
\begin{align}
 \|D^s\mathsf M^{-1}\|_\theta+
 \|D^sJ\|_\theta+\|D^sC_{\rm sel}\|_\theta+
 \|D^sC_{\rm sel}^{-1}\|_\theta&\le C_q,\notag\\
 \|D^sE\|_\theta+
 \|D^s(J-\mathsf M^{-1})\|_\theta+
 \|D^s(C_{\rm sel}-\mathsf M^{-1})\|_\theta+
 \|D^sS\|_\theta&\le C_qe^{-cB}.
 \label{eq:selected-window-jet-bounds}
\end{align}
They hold for all hole geometries in the operated
construction. The constants are uniform in cutoff and
period. The corresponding estimates for selected
inverse columns of \((1-t)1+t\mathsf M\) hold uniformly
for \(0\le t\le1\).
\end{lemma}
\begin{proof}
Use the common inverse partition formula
\eqref{eq:selected-inverse-partition-formula}.
All threshold factors multiply to \(h_{\rm c}^s\).
The principal-matrix part of
Lemma~\ref{lem:actual-combes-thomas} gives a common envelope
\(C_0e^{-\nu d(i,j)}\) for the full inverse and every
principal inverse, since each principal matrix has
the same gap and weighted row bound. Choose a rate
less than \(\min(\nu,\theta_0)/4\).
In each term of \eqref{eq:selected-inverse-partition-formula},
the inverse edges and derivative tensors join into a
connected graph. Their total connecting length bounds
\(\mathsf t(i,j,\boldsymbol\ell)\), by
Lemma~\ref{lem:rooted-kernel-composition}.
Sum the nonnegative absolute envelopes, including all
joined indices. The four-dimensional counting bound
sums every inverse edge at the smaller rate. A
sufficient constant at order \(s\) is consequently
\begin{equation}
 \sum_{b=1}^s\ 
 \sum_{\substack{r_1+\cdots+r_b=s\\r_v\ge1}}
 \frac{s!}{r_1!\cdots r_b!}
                C_0^{b+1}K_{r_1}\cdots K_{r_b},
 \label{eq:selected-inverse-jet-price}
\end{equation}
after enlarging \(C_0\) for those inverse-edge sums.
For the column selected in \(\Omega_j\), every internal
inverse index is restricted to that window. Its entry
envelope is the same \(C_0e^{-\nu d}\). Removing these
restrictions in the nonnegative bound gives exactly
the same sums. This proves bounds rooted at either
\(i\) or \(j\) for \(D^sJ\), although the window
depends on the column. It introduces no count of the
number of windows.

The identity \((\mathsf MJ)_{ij}=\delta_{ij}\) for
\(i\in\Omega_j\) holds throughout the unmasked germ,
with the windows fixed. Its every derivative therefore
vanishes there. The product rule gives an unreduced
rooted bound \(C_q\) for \(D^sE\). On its support,
\(d(i,j)>B/4-2\). Reducing the rate by a factor two
then gives
\[
 e^{\theta\mathsf t(i,j,\boldsymbol\ell)/2}
 \le e^{-\theta(B/4-2)/2}
                      e^{\theta\mathsf t(i,j,\boldsymbol\ell)}.
\]
This proves its exponential smallness, including all
input-index sums. Use the exact identities
\[
 J-\mathsf M^{-1}=\mathsf M^{-1}E,\qquad
 C_{\rm sel}-\mathsf M^{-1}
       =\tfrac12\{(J-\mathsf M^{-1})+
                                  (J-\mathsf M^{-1})^T\}
\]
and \(S=\mathsf M(C_{\rm sel}-\mathsf M^{-1})\).
Leibniz's rule preserves one exponentially small
factor in every differentiated product. This proves
the second line of \eqref{eq:selected-window-jet-bounds}.

The selected covariance has the spectral bounds
\eqref{eq:selected-covariance-gap} and the weighted
Schur bound just proved. Apply Combes--Thomas once
more, followed by
\eqref{eq:selected-inverse-partition-formula}, to
\(C_{\rm sel}^{-1}\). This proves the remaining
bound. Finally replace the gap by \(\min(1,\eta)\)
and the derivative prices by \(K_0+2,K_1,\ldots,K_q\)
for the interpolation in \(t\). These choices do not
depend on \(t\), completing the proof.
\end{proof}

The same proof applies to a covariance marginal
\(C_I=(C_{\rm sel})_{I,I}\) on any fixed set of read
coordinates \(I\), and then to \(C_I^{-1}\).
Indeed a principal submatrix retains both spectral
bounds and all absolute rooted derivative bounds.
The inverse partition formula sums only indices in
\(I\), which can again be enlarged for the majorant.
Thus the inverse and score constants do not depend
on the shape or cardinality of \(I\).

To locate the sole size factor in a Gaussian score,
write \(K=C_I^{-1}\). Let \(T\subset I\) consist of all
colour coordinates on the retained physical links;
the other coordinates remain Gaussian. Let \(z\) include the
physical coordinates \((a_T-\mu_T)/g\) together with
the still-held Gaussian coordinates, as in
\eqref{eq:regular-mixed-physical-Gaussian-density}.
The Gaussian part of that exact density is
\[
 p_I=(2\pi)^{-|I|/2}(\det C_I)^{-1/2}
                                        e^{-z^TKz/2}.
\]
For one threshold-scaled direction, a prime denoting
that derivative, direct differentiation gives
\begin{equation}
 (\log p_I)'=-\tfrac12\operatorname{Tr}(KC_I')
       +\tfrac12 z^TKC_I'Kz-z^TKz'.
 \label{eq:selected-marginal-score}
\end{equation}
The trace term is at most
\(|I|\|KC_I'\|_{2\to2}/2\); it is not a trace over
the whole torus. Higher log derivatives are finite
sums of the same traces and quadratic expressions,
by \eqref{eq:selected-inverse-partition-formula}.
Their matrix bounds are independent of \(I\).
On the retained physical chart the scaled
\(z_T\)-derivatives have prices \(C_B(1+P)^{m_q}\)
from the logarithmic coordinate and affine-source
formulas. Each trace thus costs at most \(C_q|I|\),
and the other terms are bounded by a fixed polynomial
in \(P,|I|\) and \(1+\sum_{i\notin T}z_i^2\).

For clarity, derivatives of a density itself use
only one such polynomial at each fixed total order:
\begin{equation}
 \frac{D^qp_I[v_1,\ldots,v_q]}{p_I}
   =\sum_{\pi\in\operatorname{Part}(q)}
          \prod_{B\in\pi}D^{|B|}\log p_I[v_B].
 \label{eq:selected-density-partition-formula}
\end{equation}
There are finitely many partitions, depending on
\(q\) alone. The original Haar Jacobians and the
retained broad cutoffs multiply this density as
specified in
\eqref{eq:regular-mixed-physical-Gaussian-density};
their finite derivatives add only their already
declared local prices. No physical Borel multiplier
is differentiated in this calculation.
For a full footprint \(X\), \(|I|\le C_B n_B(X)\).
These formulas supply precisely the polynomial
dependence on \(n_B(X)\) in
Lemma~\ref{lem:central-common-score}. Its Gaussian
moments are those of
Corollary~\ref{cor:finite-range-tilted-polynomial},
with the same single precision tilt.

\subsection{The exact effective reach of the original row cutoff}

One original test deserves explicit
attention before any locality claim.
For a curvature-row anchor \(x\),
threshold \(t\), and fixed relaxation
parameter \(\varrho>0\), its retained-row
test is
\begin{equation}
 k_{t,\varrho}(x;U)=
 \mathbf1_{\{
   d_G(P_p(U))\le
          t e^{d_\infty(x,o(p))/\varrho}
                              \ {\rm for\ every}\ p\}} .
 \label{eq:original-exponential-row-test}
\end{equation}
Here \(o(p)\) is the plaquette's base
site. The distance and \(\varrho\)
are measured on the lattice where
that row is evaluated.
This test multiplies the principal
row square. It is distinct from
the mask on an incoming regular
activity. Its direct plaquette
term uses the prescribed piecewise
quadratic and linear clipping.
Both are retained in
\(\mathcal D_\sigma\).

Although \eqref{eq:original-exponential-row-test}
quantifies over all plaquettes, it
has the exact finite effective reach
\[
 R_{\rm keep}(t,\varrho)
          =\left\lceil\varrho\log(2/t)\right\rceil+2,
                         \qquad 0<t\le1.
\]
Indeed \(d_G(P)\le2\) everywhere.
At any distance at least
\(\varrho\log(2/t)\), the right side
of its inequality is at least two,
so that inequality is automatic.
Only the remaining plaquettes,
and the links on their boundaries,
are actually read.
This reasoning preserves equality
at every cutoff threshold.

For \(t=h\) or \(h_{\rm c}\), the
effective radius is at most
\(C_{\rm profile}(1+\log\widehat\beta)\),
with the actual fine nominal shift
included when appropriate.
Increase the profile's entry
threshold so that these radii,
in coarse units, and the predictor
reach are less than \(B/100\).
All read bounds below include
these tests.
In particular an operated coarse
guide row outside the hole
neighbourhood has its original
retained-row test equal to one:
nearby plaquettes are healthy,
and the more distant inequalities
are automatic. On its direct
coarse plaquette the clipping
factor also has its quadratic
value. These observations justify
using its unmasked local formula
on such a read region.

\subsection{Physical locality of the selected centre}

An entry of \(\mathsf M\) reads the finite paths
in its two row stencils, the active-column tests,
and the operated-row tests. Each test must read
all possible nearby holes. These read regions
have radius \(C B\), with \(C\) fixed.
An entry of a principal inverse reads the entries
of its entire principal matrix. Hence \(J_{ij}\)
reads a neighbourhood of radius \(CB\) of its
window, even when its numerical value vanishes.

In
\[
 \mu_j=-\sum_{i\in\Omega_j}J_{ij}b_i
\]
the source \(b_i\) reads only rows meeting \(i\)
and their protected frozen links. It follows that
\(\mu_j\) reads \(V,Y,t\), and the sector queries,
in a fixed \(CB\)-neighbourhood of \(j\).
There is no read of distant frozen variables
through a full inverse.
All active links in a physical observable have
coordinates \(\mu_e+gz_e\). Consequently a
bounded Borel function of finitely many physical
links is a bounded Borel function of finitely
many \(z\)'s, frozen links, and branch variables,
with the enlarged read region just described.
No continuity of that observable is needed.

These formulas also prove joint measurability
in \(V,Y,z\): path products and adjoint actions
are continuous, the discrete tests are Borel,
and a positive principal inverse is a continuous
function of its matrix entries. Integrating
against the selected Gaussian can equivalently
be done using its explicit Lebesgue density,
so it is a Borel operation in these parameters.

On a fixed healthy read region the matrix
entries have the physical left derivatives
of Section~\ref{sec:activity-norms}.
Principal inverse derivatives follow from
\[
 D_A(\mathsf M_\Omega^{-1})
       =-\mathsf M_\Omega^{-1}(D_A\mathsf M_\Omega)
                                      \mathsf M_\Omega^{-1}.
\]
Iterating this identity through the fixed order
\(q\), and using the common resolvent bounds,
gives prices \(C_B\).
Derivative directions on a window grow by at
most a fixed power of \(1+B\); this is part of
that price. The differentiated error \(E\)
still vanishes identically in the interior
window entries. The same exponent reduction
therefore gives a derivative price
\(C_B e^{-cB}\) for the small errors in
\eqref{eq:selected-window-error}.
These are statements about the unmasked
matrix formulas on their healthy regions,
not derivatives of the sector indicators.

There is also exact covariance under a
coarse gauge transformation. In the tree
slice it acts on a fine colour coordinate
by the adjoint action at its parent.
Denote this block diagonal orthogonal map
by \(O\), and the corresponding row map
by \(O_{\rm row}\).
The transport convention gives
\[
 D(V^u)=O_{\rm row}D(V)O^T,\quad
 \alpha(V^u,Y^u)=O_{\rm row}\alpha(V,Y).
\]
The inactive projection commutes with \(O\).
Every window includes the complete colour
space at its selected links, so restriction
and inversion give
\[
 \mathsf M(V^u)=O\mathsf M(V)O^T,\quad
 J(V^u)=OJ(V)O^T,\quad
 C_{\rm sel}(V^u)=OC_{\rm sel}(V)O^T,\quad
 \mu(V^u,Y^u)=O\mu(V,Y).
\]
Consequently the selected Gaussian is
carried into the selected Gaussian by
this orthogonal change of coordinates.
Its determinant is unchanged.
Frozen Haar measure is invariant under
the corresponding left and right group
translations. Thus fully contracted
component integrals of covariant letters
are gauge invariant. This argument uses
the same local paths and shared tags in
all factors.

\subsection{Exact change of variables, including all constants}

Let \(n_{\rm ac}=|I_\sigma|\),
\(n_{\rm fr}=|J_\sigma|\), and \(n_0=|I|\).
Set
\[
 s_g=\zeta_G(\sqrt{2\pi}\,g)^{d_G^{\rm Lie}},
 \qquad r=\mathsf M\mu+b .
\]
The letter \(r\) here is a vector residual, not
the chart radius \(r_*\).
Inactive rows and columns of \(\mathsf M\) form
an identity block, \(b_{\rm in}=0\), and
\(\mu_{\rm in}=0\).
The selected covariance also equals the identity
on that block and has zero cross entries with it.
We may therefore include inactive standard
Gaussian coordinates; their integral is exactly one.

Define the exact phase difference, on the
physical chart support, by
\begin{equation}
 \mathcal D_\sigma(V,Y,a)
       = A_t(V,\mathcal U_\sigma(V,Y,a))
                            -\tfrac12\|Da+\alpha\|_2^2 .
 \label{eq:finite-range-exact-phase-difference}
\end{equation}
There is no estimate in this definition.
In particular it retains the actual fine principal
action, all its masks and uncut tails, the actual
noise action, and every difference from the
chosen cut affine rows.
Put
\begin{equation}
 \mathcal Q_\sigma(z)=
       -g^{-1}\langle r,z\rangle
       +\tfrac12\langle z,
                    (C_{\rm sel}^{-1}-\mathsf M)z\rangle .
 \label{eq:finite-range-gaussian-density-correction}
\end{equation}
For every bounded Borel multiplier \(\mathcal B(a)\)
supported in the active chart,
\begin{align}
 &\zeta_G^{n_{\rm ac}}\int
   \mathcal B(a)e^{-A_t(V,\mathcal U_\sigma)/g^2}\,da
       \notag\\
 &\quad =
 s_g^{n_{\rm ac}}(\det C_{\rm sel})^{1/2}
 e^{-\|D\mu+\alpha\|_2^2/(2g^2)}
 \mathbb E_{\mu_{C_{\rm sel}}}
 \left[
  \mathcal B(\mu+gz)
  e^{\mathcal Q_\sigma(z)
             -\mathcal D_\sigma(V,Y,\mu+gz)/g^2}
 \right].
 \label{eq:exact-selected-Gaussian-identity}
\end{align}
The integral on the left is over active real
coordinates only; \(\mathcal B\) ignores the
inactive Gaussian coordinates on the right.

\begin{proof}
Make the translation and scaling \(a=\mu+gz\)
on the active subspace. Their Jacobian is
\(g^{d_G^{\rm Lie}n_{\rm ac}}\).
Since \(Dz=Dz_{\rm ac}\),
\[
 \frac{\|D(\mu+gz)+\alpha\|_2^2}{2g^2}
  =\frac{\|D\mu+\alpha\|_2^2}{2g^2}
        +g^{-1}\langle r,z\rangle
        +\tfrac12\langle z,\mathsf Mz\rangle
\]
on that subspace. Insert the Gaussian density
\[
 (2\pi)^{-d_G^{\rm Lie}n_{\rm ac}/2}
 (\det C_{\rm sel})^{-1/2}
                   e^{-\langle z,C_{\rm sel}^{-1}z\rangle/2}.
\]
The determinants of the completed full matrices
equal those of their active blocks.
The constants and exponents give exactly
\eqref{eq:exact-selected-Gaussian-identity}.
Adding the inactive probability Gaussian proves
the stated full-coordinate version.
The identity holds first for nonnegative Borel
functions by the change-of-variables theorem.
Applying it to positive and negative parts
proves it for absolutely integrable signed
functions, including the multiplier in
\eqref{eq:finite-range-full-multiplier}.
\end{proof}

For clarity, the full sector contribution to
\eqref{eq:actual-block-integral} is
\begin{align}
 &e^{vn_f}z_{\rm n}(\beta)^{-b_c}
 s_g^{n_{\rm ac}}
 \sum_t\int_{G^{J_\sigma}}
 (\det C_{\rm sel})^{1/2}
       e^{-\|D\mu+\alpha\|_2^2/(2g^2)}
       \notag\\
 &\hspace{27mm}\cdot
 \mathbb E_{\mu_{C_{\rm sel}}}
 \left[
  \mathcal B_\sigma(V,Y,\mu+gz)
  e^{\mathcal Q_\sigma(z)
       -\mathcal D_\sigma(V,Y,\mu+gz)/g^2}
 \right]\,dH_{J_\sigma}(Y).
 \label{eq:full-physical-selected-Gaussian-sector}
\end{align}
Summing the exact partition returns the original
blocked density. Noise normalisation is outside
the branch sum; the Haar density \(j_G\), fine
indicators, incoming activities, and branch
likelihoods are inside the same expectation.

Let \(\mathsf M_0\) be the cut flat operated
precision with no holes, on this same torus
and forest. Factoring the common reference
\(s_g^{n_0}(\det\mathsf M_0)^{-1/2}\)
out of every sector leaves exactly
\begin{equation}
 s_g^{-n_{\rm fr}}\,
 \exp\left\{\tfrac12\left[
  \log\det\mathsf M_0-\log\det\mathsf M
                         +\log\det(1+S)\right]\right\}.
 \label{eq:finite-range-vacuum-ledger}
\end{equation}
This follows from
\(\mathsf MC_{\rm sel}=1+S\).
The factor \(s_g^{-n_{\rm fr}}\) is the
normalisation ratio for the frozen links.
It is not discarded or bounded by a universal
constant. Its logarithm is of order
\(n_{\rm fr}(1+|\log g|)\).
The constant square in
\eqref{eq:full-physical-selected-Gaussian-sector}
is also retained, with its displayed sign.

The common scalar before this ledger is
\[
 e^{vn_f}z_{\rm n}(\beta)^{-b_c}
             s_g^{n_0}(\det\mathsf M_0)^{-1/2}.
\]
For one rooted tree in each block,
\(n_0=4n_f-(n_f-n_{\rm c})
      =(3L^4+1)n_{\rm c}\), where
\(n_{\rm c}=c^4\), and \(b_c=4n_{\rm c}\).
Its logarithm per coarse site is therefore
\[
 vL^4-4\log z_{\rm n}(\beta)
 +(3L^4+1)\log s_g
             -\frac{\log\det\mathsf M_0}{2n_{\rm c}}.
\]
The last term is uniformly bounded at
a fixed block and profile: every
eigenvalue of \(\mathsf M_0\) lies
between the same positive constants.
Its total logarithm is extensive.
This scalar is kept in the vacuum
bookkeeping before any further
determinant or log-cluster extraction.

\subsection{Relation to the earlier central reference}

On a globally good background, in the common
intact central domain with no frozen variables,
write the old
coordinate as \(y=\rho+gz\), and the selected
coordinate as \(y=\mu+g\widetilde z\).
Then
\[
 z=\widetilde z+\delta,\qquad
                         \delta=g^{-1}(\mu-\rho).
\]
For the old covariance \({\mathsf H}^{-1}\) and any positive
selected covariance \(C\), the exact density ratio is
\begin{equation}
 (\det {\mathsf H}\,\det C)^{1/2}
 \exp\left\{
 -\tfrac12\langle\widetilde z+\delta,
                         {\mathsf H}(\widetilde z+\delta)\rangle
 +\tfrac12\langle\widetilde z,C^{-1}\widetilde z\rangle
 \right\}.
 \label{eq:old-selected-Gaussian-bridge}
\end{equation}
The same physical test function must be evaluated
at \(y=\mu+g\widetilde z\) on the selected side.
This proves the bridge for bounded Borel tests,
not only for analytic observables. The constant,
linear, quadratic, and determinant factors in
the ratio are all present in the direct identity
\eqref{eq:exact-selected-Gaussian-identity}.
On this intact domain the masked and unmasked
principal actions agree. Outside it the direct
sector identity retains their actual difference
in \(\mathcal D_\sigma\); the old unmasked central
interaction is not substituted there.

There is no assertion that the two covariances
or centres are equal. On a sector with holes,
the direct Haar identity constructs the selected
representation without requiring a stationary
Hessian on a rough coarse field.
Likewise the radial projection used to define
\(\mathcal Z_R\) is not inserted into this identity:
it changes the reference integrand on its failure
set. Comparisons with it must retain that set.

\section{Local Gaussian corrections and component estimates}
\label{sec:finite-range-components}

\subsection{Localising the covariance and determinant corrections}

The error \(S=\mathsf MC_{\rm sel}-1\) has
finite range at most \(B/4+2+2r_B\), and a
weighted Schur norm at most \(\epsilon_B\).
For \(\epsilon_B<1\),
\begin{equation}
 C_{\rm sel}^{-1}-\mathsf M
       =\sum_{n\ge1}(-S)^n\mathsf M,\qquad
 \log\det(1+S)
       =\sum_{n\ge1}\frac{(-1)^{n+1}}n\operatorname{Tr}S^n .
 \label{eq:selected-ratio-path-series}
\end{equation}
The first order in the first series is important:
\(C_{\rm sel}^{-1}=(1+S)^{-1}\mathsf M\).
It need not equal \(\mathsf M(1+S)^{-1}\).
Both series converge absolutely at each finite
volume. Their rooted prices are \(C\epsilon_B\);
their unrooted trace price is at most
\(|\mathcal I|C\epsilon_B\).
To see the rooted statement, retain all intermediate
indices in a product, and use the row or column
sum at every successive index. For a trace fix
its initial index before performing the sum.
Then sum the initial indices only for an
extensive scalar.

For the second series, apply Jacobi's
formula to \(1+tS\), expand its inverse
as a Neumann series, and integrate \(t\)
from zero to one. The matrices \(1+tS\)
are invertible since \(\|tS\|<1\), and
their real determinants remain positive
along that interval. This proves the
logarithm formula even though \(S\) need
not be symmetric.

There is also an exact finite-window expansion
of the main determinant difference in
\eqref{eq:finite-range-vacuum-ledger}.
Let \(J_t\) be the selected columns of
\(\mathsf M_t^{-1}\), and \(E_t=\mathsf M_tJ_t-1\).
Uniformly on \(0\le t\le1\),
\[
 \mathsf M_t^{-1}
      =J_t(1+E_t)^{-1}
      =\sum_{n\ge0}J_t(-E_t)^n .
\]
Jacobi's determinant formula and positivity give
\begin{equation}
 \log\det\mathsf M
  =\int_0^1\sum_{n\ge0}
       \operatorname{Tr}\bigl[
         (\mathsf M-1)J_t(-E_t)^n\bigr]\,dt .
 \label{eq:selected-main-determinant-flow}
\end{equation}
The zeroth term is a local finite-window scalar.
The remaining rooted terms have price
\(C\epsilon_B^n\). The flat determinant has
the identical formula with \(\mathsf M_0\);
the difference is paired before taking
absolute bounds. Near holes its zeroth term
is retained as a local scalar cost. It is
not replaced by its flat value.

Every path term in these series reads its
endpoints, the windows of its factors,
their row stencils, and all sector queries
used by those matrices. Joining these read
regions gives its full physical footprint.
A chain of \(n\) finite-range factors is
covered by at most \(C(n+1)\) balls of
radius \(B\), with \(C\) fixed.
Consequently the covering price is paid
by a geometric series with ratio
\(e^{C(\kappa+1)}C_B e^{-cB}\).
This is less than \(1/2\) after a threshold
depending only on the profile.
The same reasoning applies to fixed-order
physical derivatives of the unmasked
matrix formulas.

\subsection{One joint quadratic-growth bound}

We need estimates which permit several
factors to grow in the same Gaussian
coordinate. Independence of single factors
will not be used.

\begin{lemma}[One Gaussian integral with linear and quadratic sources]
\label{lem:one-Gaussian-source-bound}
Let \(z\) be a centred finite-dimensional Gaussian vector
with positive covariance \(C\). Let \(\ell\) be a real
source vector and \(K\) a symmetric nonnegative matrix.
Suppose
\[
 0\le C^{1/2}KC^{1/2}\le\theta1,\qquad 0\le\theta<1.
 \]
Then
\begin{align}
 \mathbb E_C e^{\ell\cdot z+z^TKz/2}
 &=\det(1-C^{1/2}KC^{1/2})^{-1/2}
   \exp\{\tfrac12\ell^T(C^{-1}-K)^{-1}\ell\}\notag\\
 &\le\exp\left\{
       \frac{\operatorname{Tr}(CK)+\ell^TC\ell}
                         {2(1-\theta)}\right\}.
 \label{eq:one-Gaussian-source-bound}
\end{align}
For several factors use \(\ell=\sum_\alpha\ell_\alpha\)
and \(K=\sum_\alpha K_\alpha\) in this same formula.
In particular all cross terms of the linear sources are
included. Its constants do not depend on the dimension.
\end{lemma}
\begin{proof}
The precision \(C^{-1}-K\) is positive. Complete its square
with mean \((C^{-1}-K)^{-1}\ell\).
The Gaussian normalising constants give the displayed
determinant exactly. If \(T=C^{1/2}KC^{1/2}\), then
\[
 (C^{-1}-K)^{-1}
       =C^{1/2}(1-T)^{-1}C^{1/2}
       \le(1-\theta)^{-1}C .
 \]
For an eigenvalue \(0\le t\le\theta\),
\(-\log(1-t)\le t/(1-\theta)\).
Summing this inequality and using
\(\operatorname{Tr}T=\operatorname{Tr}(CK)\)
proves the bound.
\end{proof}

For local source factors the common linear square can
be charged without discarding its cross terms.
Choose numbers \(w_{\alpha i}\ge0\), positive wherever
\(\ell_\alpha(i)\ne0\), with
\(\sup_i\sum_\alpha w_{\alpha i}\le W\).
If \(\|C\|\le {\mathsf M}\), then
\begin{equation}
 \left(\sum_\alpha\ell_\alpha\right)^TC
          \left(\sum_\alpha\ell_\alpha\right)
 \le {\mathsf M}W\sum_{\alpha,i}
                   \frac{|\ell_\alpha(i)|^2}{w_{\alpha i}} .
 \label{eq:one-Gaussian-local-source-charge}
\end{equation}
A summand with both numerator and denominator zero is zero.
Indeed, weighted Cauchy--Schwarz at each coordinate bounds
the square of \(\sum_\alpha\ell_\alpha(i)\) by
\((\sum_\alpha w_{\alpha i})
 \sum_\alpha|\ell_\alpha(i)|^2/w_{\alpha i}\).
Summing over coordinates proves the assertion.
The exponentially summable endpoint envelopes of our
matrix words supply such \(w_{\alpha i}\), with bounded
\(W\) independently of the volume.
This is the common-source organisation used for local
integration in \cite[Lemma~11.13, PDF p.~73]{ArmstrongDraft};
the needed Gaussian formula has been proved here.

\begin{lemma}[A local exponential moment]
\label{lem:finite-range-joint-tilt}
Let \(P=C_{\rm sel}^{-1}\), and suppose
\(P\ge(r_0+u_0)1\), with \(r_0>0\).
If \(d_i\ge0\) and \(d_i\le u_0\), then
\begin{equation}
 \mathbb E_{\mu_{C_{\rm sel}}}
      e^{\frac12\sum_i d_i z_i^2}
 =\left(\frac{\det P}{\det(P-\operatorname{diag}d)}
                                      \right)^{1/2}
 \le \exp\left\{\frac1{2r_0}\sum_i d_i\right\}.
 \label{eq:finite-range-joint-tilt}
\end{equation}
In particular only the coordinates actually charged
by \(d\) contribute to this bound.
\end{lemma}
\begin{proof}
Apply Lemma~\ref{lem:one-Gaussian-source-bound} with
\(\ell=0\), \(C=P^{-1}\), and \(K=\operatorname{diag}d\)
to obtain the equality. Its positivity hypothesis holds
because \(P-K\ge r_0 1\).
The sharper bound in terms of \(r_0\) follows directly
from the same precision interpolation:
For \(P_t=P-t\operatorname{diag}d\),
\[
 \log\det P-\log\det P_1
     =\int_0^1\operatorname{Tr}
                    (P_t^{-1}\operatorname{diag}d)\,dt
     \le r_0^{-1}\sum_i d_i .
\]
The last inequality uses
\((P_t^{-1})_{ii}\le r_0^{-1}\).
Taking half the exponent proves the bound.
\end{proof}

\begin{corollary}[Polynomial insertions under the same tilt]
\label{cor:finite-range-tilted-polynomial}
Under the hypotheses of
Lemma~\ref{lem:finite-range-joint-tilt}, let \(c_i\ge0\)
and let \(m\ge1\) be an integer. Then
\begin{equation}
 \mathbb E_{\mu_{C_{\rm sel}}}
 \left[
  e^{\frac12\sum_i d_i z_i^2}
              \left(\sum_i c_i z_i^2\right)^m\right]
 \le
 e^{\sum_i d_i/(2r_0)}
 \frac{(2m-1)!!}{r_0^m}\left(\sum_i c_i\right)^m ,
 \label{eq:finite-range-tilted-polynomial}
\end{equation}
where \((2m-1)!!=1\cdot3\cdots(2m-1)\).
The case \(m=0\) is the preceding lemma.
No independence of coordinates is assumed.
\end{corollary}
\begin{proof}
Completing the same square changes the covariance to
\((P-\operatorname{diag}d)^{-1}\le r_0^{-1}1\),
with exactly the determinant ratio of
\eqref{eq:finite-range-joint-tilt}.
Under that one probability law, every coordinate has
variance at most \(r_0^{-1}\), so its \(2m\)-th moment
is at most \((2m-1)!!r_0^{-m}\).
Minkowski's inequality gives
\[
 \left\|\sum_i c_i z_i^2\right\|_{L^m}
 \le\sum_i c_i\|z_i^2\|_{L^m}
 \le ((2m-1)!!)^{1/m}r_0^{-1}\sum_i c_i.
 \]
Raise to the \(m\)-th power and multiply by the
already charged determinant ratio.
\end{proof}

More generally, suppose selected factors \(f_\alpha\)
have nonnegative prices \(a_\alpha\) and nonnegative
coordinate charges \(d_{\alpha i}\), with
\[
 |f_\alpha(z)|\le
 a_\alpha e^{\frac12\sum_i d_{\alpha i}z_i^2}.
\]
For any finite selected set \(A\) obeying
\(\sum_{\alpha\in A}d_{\alpha i}\le u_0\) at every
coordinate, the preceding lemma gives
\begin{equation}
 \int\prod_{\alpha\in A}|f_\alpha|\,d\mu_{C_{\rm sel}}
 \le\prod_{\alpha\in A}
       \left[a_\alpha
        e^{(2r_0)^{-1}\sum_i d_{\alpha i}}\right].
 \label{eq:finite-range-word-bound}
\end{equation}
This proof charges the whole selected word once.
It remains valid when factors read overlapping
sets of coordinates. An independent bound for
each factor would not justify the same conclusion.

\subsection{Paying the selected Gaussian correction letters}

The small matrix errors really can be
used in the preceding joint bound.
We verify this before adding any
physical failure factor.

First, on an operated row the replica
curvatures and the coarse guide are
bounded by \(C h_{\rm c}\).
Their paths lie outside the hole
neighbourhood and have the funded
exponential word price. The frozen
extension has norm at most \(2u\),
and the full flat row matrix has a
uniform Schur bound. Hence
\[
 |\alpha|_\infty+|b|_\infty
                         \le C(h_{\rm c}+u)\le C_B h .
\]
This argument uses no smallness of an
unprotected frozen link: only its
bounded extension enters \(\alpha\).
The true frozen link is still present
in the exact phase difference.
Since \(\mu=-J^Tb\) and
\(J^T-\mathsf M^{-1}\) has the same
error bound as \(J-\mathsf M^{-1}\),
\begin{equation}
 r=-\mathsf M(J^T-\mathsf M^{-1})b,\qquad
                  g^{-1}|r|_\infty
       \le C_B P(\widehat\beta)e^{-cB}.
 \label{eq:selected-centre-residual-price}
\end{equation}
The ratio \(h/g\) is bounded by a fixed
multiple of the fine monitor, using
the actual coupling band. The fixed
fine/coarse nominal shift is included
in this monitor.

\begin{lemma}[A joint price for Gaussian correction letters]
\label{lem:selected-correction-word-prices}
Expand the quadratic form in
\eqref{eq:finite-range-gaussian-density-correction}
using the retained paths in
\eqref{eq:selected-ratio-path-series},
and keep one linear letter
\(-g^{-1}r_i z_i\) per coordinate.
Expand their exponentials by Mayer
factors, selecting each address at
most once. For every fixed covering
exponent, these factors admit envelopes
\[
 |f_\alpha(z)|
 \le a_\alpha
             e^{\frac12\sum_i d_{\alpha i}z_i^2},
 \qquad a_\alpha,d_{\alpha i}\ge0,
\]
such that
\begin{align}
 \sup_x\sum_{\alpha:x\in X_\alpha}
           e^{(\kappa+3)n_B(X_\alpha)}a_\alpha
                         &\le C_Be^{-c'B},\notag\\
 \sup_i\sum_\alpha d_{\alpha i}
                         &\le C_Be^{-c'B}.
 \label{eq:selected-correction-occupation}
\end{align}
The constants depend on the profile but
not on the number of coordinates.
The small scalar determinant tails
have the same incidence estimate
with zero quadratic charge.
\end{lemma}
\begin{proof}
For \(0\le c\le1\) and \(Q\ge0\),
\[
 |e^v-1|\le cQe^{cQ}
              \le \sqrt c\,e^{2\sqrt c\,Q}
                 \quad\text{if }|v|\le cQ.
\]
The second inequality follows from
\(\sqrt c\,Q\le e^{\sqrt c\,Q}\).
For a quadratic path letter take
\(Q=z_i^2+z_j^2\), increasing \(c\)
to bound its coefficient if necessary.
It has cap \(\sqrt c\) and a charge
at most \(4\sqrt c\) at each endpoint,
counting both occurrences if the
endpoints coincide.

To sum these caps, keep every
intermediate coordinate of the path.
There are uniform pointwise
coefficient envelopes
\[
 |S_{ij}|\le C\epsilon_Be^{-\theta_1d(i,j)},
 \qquad
 |\mathsf M_{ij}|\le Ce^{-\theta_1d(i,j)} .
\]
Their square roots have row and
column sums \(C\sqrt{\epsilon_B}\)
and \(C\), respectively, by the
coordinate density bound in coarse
balls. Therefore the sum of the
square roots of coefficients of
paths with \(n\ge1\) small factors,
rooted at either endpoint, is at
most \(C(C\sqrt{\epsilon_B})^n\).
It also bounds the total charge
at that endpoint.

For the incidence norm, a point
of the full footprint belongs to
one of the \(CB\)-windows along
the path. Choose that position
as root. This costs \(C_B(n+1)\).
The covering weight costs at most
\(e^{C(\kappa+3)(n+1)}\).
Their sum over \(n\ge1\) is a
convergent geometric series with
a polynomial allowance in \(B\),
and is bounded by \(C_Be^{-c'B}\).

For the linear letters use
\(Q=1+z_i^2\) and
\(c=C_BP e^{-cB}\) in the same
elementary inequality. The constant
part of the exponential goes into
the cap. There is one such letter
per coordinate, with a \(CB\)-read
region. Thus their cap incidence
and total charge are bounded by
\(C_B\sqrt{P}e^{-cB/2}\).
Because \(\log P=(\log\log\widehat\beta)^2=o(B)\),
this is again \(C_Be^{-c'B}\).

The scalar determinant tails have
absolutely summable coefficients
with rooted prices \(C_B\epsilon_B\).
Their Mayer factors are bounded
by twice those coefficients once
the coupling threshold is large.
They read no Gaussian coordinate
and require no quadratic charge.
The zeroth, local term in
\eqref{eq:selected-main-determinant-flow}
is not one of these small tails.
\end{proof}

By increasing the fixed-profile
threshold, the total correction
charge is at most \(\lambda/8\),
where \(C_{\rm sel}^{-1}\ge\lambda1\).
The joint failure bound below
spends at most \(\lambda/4\).
Together they leave at least
\(5\lambda/8\) of precision before
any further physical interaction
is charged. Bounded incoming
activity factors have zero charge.
Thus these three parts can be
estimated simultaneously under
the same Gaussian law.

\subsection{The selected centre away from holes}

We check the physical meaning of a small
Gaussian coordinate box. This check is
local; no stationary point on the whole
rough background is introduced.

\begin{lemma}[An intact box for the selected centre]
\label{lem:selected-centre-intact-box}
At the fixed profile choose the block
and the noise reserve as in
Lemma~\ref{lem:physical-calibrated-residual}.
There is \(\varepsilon_L>0\), independent
of cutoff and period, with the following
property. A physical test at a coarse
anchor outside \(H_{20B}(\mathcal H_\sigma)\)
uses the selected centre of this same
sector. If all its local Gaussian
coordinates obey
\[
 |z_i|\le H,\qquad
                         H=\varepsilon_Lh/g,
\]
then its fine plaquettes have chord
less than \(h\), its tested noise
mismatch is less than \(Wh\), and its
frame and predictor branches are
the original regular branches.
The broad chart cutoff at width
\(hB^{17}\) is also identically one
on this box. The local coordinate
set and the number of such tests
reading one coordinate have bounds
of the form \(C_B\).
\end{lemma}
\begin{proof}
We first compare two centres on a
local completion. The row and column
bounds above give the sharper value
bound
\[
 |\mu|_\infty\le C_L(h_{\rm c}+u)
                                      \le C_LBh .
\]
For fine-field and noise tests we
need the more precise calibrated
bound near their anchor.

Take a box of radius \(B\) around
that anchor, reducing a fixed
geometric fraction if needed.
It lies in the healthy region
and is nonwrapping by the period
reserve. Put \(V\) in its local
axial gauge and extend its link
logarithms by zero outside the
box, preserving a central
neighbourhood containing all
selected windows used by the
test. Denote the completion by
\(V^{\rm loc}\).
Its global link coordinates and
global curvature have size at
most \(C B^2h_{\rm c}\).
The stationary and inverse formulas
of Section~\ref{sec:covariant-kernel-construction}
are valid on this larger auxiliary
domain once \(C_BB^2h_{\rm c}\) is
small, as in the proof of
Lemma~\ref{lem:physical-calibrated-residual}.

Use unmasked local formulas on
this auxiliary completion.
They agree with the actual
selected records in the central
windows, where every original
good-field test equals one.
They need not agree with masked
records at the artificial boundary.
No factor in the physical integral
is changed by this comparison.

Write \(D_{\rm full},\alpha_{\rm full}\)
for these auxiliary, no-hole,
cut affine rows and source, and
put
\[
 \mathsf M_{\rm full}=D_{\rm full}^TD_{\rm full},
 \qquad b_{\rm full}=D_{\rm full}^T\alpha_{\rm full}.
\]
All coordinates here are free
tree-slice coordinates.
If the auxiliary curvature bound
is \(\delta\), the path comparisons
used in
Lemma~\ref{lem:actual-background-Hessian-comparison}
give
\begin{align*}
 |\mathsf M_{\rm full}
       -\varphi_{yy}(V^{\rm loc},0)|_0
                       &\le C_B(\delta+e^{-cB}),\\
 |b_{\rm full}
       -\varphi_y(V^{\rm loc},0)|_\infty
                       &\le C_B(\delta^2+e^{-cB}).
\end{align*}
Here \(|\cdot|_0\) is the ordinary
row and column Schur norm.
We give the cancellation behind
the second line. With \(F_0\)
the actual energetic row at
the replica, \(F_y\) its actual
derivative, and \(G\) the coarse
guide, these two gradients have
the forms
\[
 \varphi_y=F_y^TF_0+\mathcal R_y,\qquad
 b_{\rm full}=D_{\rm full}^T(F_0-G)
                                      +\text{cut tails}.
\]
The comparison \(F_y-D_{\rm full}\)
has size \(C_B\delta\),
\(|F_0|_\infty\le C\delta\), and
\(\mathcal R_y=O(C_B\delta^2)\);
the latter consists of the cubic
and higher curvature defects of
Section~\ref{sec:microscopic-centre}.
Furthermore
\(G=\Phi(V^{\rm loc})\mathcal C_{\rm c}(V^{\rm loc})\)
up to the retained cut tails.
The flat identity
\(\mathsf A^T\Phi_0=0\), transported
with the paired product-defect
estimate, gives
\[
 |D_{\rm full}^T
        \Phi(V^{\rm loc})
        \mathcal C_{\rm c}(V^{\rm loc})|_\infty
                            \le C_B(\delta^2+e^{-cB}).
\]
The cut tail is \(e^{-cR_B}\);
at this fixed profile it is
written as \(e^{-cB}\) with a
smaller positive \(c\).
This proves the second line
without a residual term of
order \(\delta\).
The first line follows from
the same derivative comparison
and the Hessian formula
\(F_y^TF_y+\sum_rF_rF_{r,yy}
                       +\mathcal R_{yy}\).

Let \(\rho^{\rm loc}\) be the
actual unmasked stationary
point on the completion.
Taylor expansion of its equation
about zero, using
\(|\rho^{\rm loc}|_\infty\le C\delta\),
now gives
\[
 |\mathsf M_{\rm full}\rho^{\rm loc}
                         +b_{\rm full}|_\infty
                          \le C_B(\delta^2+e^{-cB}).
\]
The inverse has a uniform
sup-norm bound. Comparing its
exact affine centre with the
selected inverse columns gives
\[
 |\mu^{\rm loc}-\rho^{\rm loc}|_\infty
                          \le C_B(\delta^2+e^{-cB}).
\]
Set \(\delta=CB^2h_{\rm c}\);
the powers of \(B\) are absorbed
into \(C_B\).

The actual sector and this
completion have the same
principal matrices and source
entries in the central selected
windows. No frozen link can
enter these windows this far
from the holes. Thus their
selected centres there are
equal, after the local gauge
transformation.
By the local calibration
argument of
Lemma~\ref{lem:physical-calibrated-residual},
the stationary completion
has fine chord less than \(h/8\)
and noise mismatch less than
\(Wh/2\) in the central test.
That argument is local: its
native tails also pay the
artificial boundary of the
completion. Since
\[
 C_B(h_{\rm c}^2+e^{-cB})/h\longrightarrow0,
\]
the same selected centre
has, for example, fine chord
less than \(h/4\) and noise
mismatch less than \(3Wh/4\).

Finally the physical field
uses \(a=\mu+gz\), with no
inverse square root acting
on \(z\).
In a fixed predictor stencil
the perturbation is at most
\(g\sqrt{d_G^{\rm Lie}}H\)
in Lie norm at each link.
The fixed-stencil Lipschitz
bound for plaquettes and the
actual regular predictor
therefore preserves the
preceding margins when
\(\varepsilon_L\) is small
enough. The local frame
and predictor chart then
have their regular branch
values, including their
uniform unused-tag masses.
The sharper value bound
\(|\mu|_\infty\le C_LBh\)
also gives
\[
 |\mu+gz|_\infty
       \le C_LBh+C_L\varepsilon_Lh
                                  <hB^{17}
\]
eventually. Since \(hB^{17}\to0\),
these are genuine active
exponential-chart points.
The coordinate counts follow
from the fixed test list,
the finite windows, and
the ball bound \(C_L(1+B)^4\).
\end{proof}

\subsection{Failure witnesses and their multiplicity}

A failure factor \(f_\alpha\) reads a finite
coordinate set \(T_\alpha\) if its Gaussian
dependence is confined to that set. Suppose
\[
 |f_\alpha|\le A_\alpha,\qquad
 f_\alpha=0\quad\hbox{when }
                     |z_i|\le H\ (i\in T_\alpha).
\]
Here the vanishing condition must be checked
for the actual physical tests when this lemma
is applied; it is not inferred from a name
assigned to a sector.

\begin{lemma}[Joint Gaussian failure estimate]
\label{lem:finite-range-failure-witness}
Suppose each coordinate lies in at most \(N_B\)
of the sets \(T_\alpha\), and \(P\ge\lambda 1\).
For every finite set \(A\) of distinct failure
labels,
\begin{equation}
 \mathbb E\prod_{\alpha\in A}|f_\alpha|
 \le\prod_{\alpha\in A}
 \left[A_\alpha |T_\alpha|
                e^{1/(4N_B)}
                e^{-\lambda H^2/(8N_B)}\right].
 \label{eq:finite-range-failure-word}
\end{equation}
The constants do not depend on the ambient
number of coordinates.
\end{lemma}
\begin{proof}
Put \(v=\lambda/(4N_B)\), and use
\[
 \mathbf1_{\{\max_{i\in T_\alpha}|z_i|>H\}}
 \le\sum_{i\in T_\alpha}
                      e^{v(z_i^2-H^2)/2}.
\]
Expand this product of finite sums, keeping
one witness \(i_\alpha\) for each selected
label. A coordinate is charged at most
\(N_Bv=\lambda/4\).
Apply Lemma~\ref{lem:finite-range-joint-tilt}
with \(r_0=\lambda/2\) and \(u_0=\lambda/2\).
Each witness contributes at most
\(e^{v/\lambda}e^{-vH^2/2}\).
There are \(\prod_\alpha |T_\alpha|\)
choices of witnesses. Substituting \(v\)
gives the stated bound.
\end{proof}

For coordinate boxes and the fixed finite
list of local tests, both \(|T_\alpha|\)
and \(N_B\) are bounded by a polynomial
\(C_B\); colour multiplicity and the block
volume belong to its constant.
Taking \(H=c_LP(\widehat\beta)\), or
\(H=P(\widehat\beta)/(C_L(1+B)^p)\) with
a fixed \(p\), gives the incidence price
\begin{equation}
 C_B\exp\{-P(\widehat\beta)^2/C_B\}
          \le C_{\mathsf M}\widehat\beta^{-{\mathsf M}}
                  \quad\text{for every fixed }{\mathsf M}>0.
 \label{eq:finite-range-failure-superpower}
\end{equation}
Indeed, with \(t=\log\log\widehat\beta\),
\(P^2=e^{2t^2}\), \(B^p=O(e^{2pt})\),
and \(\log\widehat\beta=e^t\).
Thus \(P^2/(B^p\log\widehat\beta)\to\infty\).
The covering weight of a fixed \(CB\)-box
and its polynomial incidence count are
absorbed into the same bound.

The estimate applies immediately to Gaussian
coordinate failures \(\max_{T_\alpha}|z_i|>H\).
For a physical fine-field, noise, branch,
or broad-chart failure outside the protected
hole neighbourhood,
Lemma~\ref{lem:selected-centre-intact-box}
provides precisely this vanishing condition.
The scaled regular branch masses are one
there, so their deviations also vanish.
Extra protection factors bounded by one
do not change the estimate.
Tests in the hole neighbourhood remain
inside its actual packet weight.
There is no differentiation of a hard test
in the witness argument.

\subsection{Component factorisation with the actual joint law}

Let a letter read coordinates \(T_\alpha\),
frozen links \(F_\alpha\), branch variables
\(L_\alpha\), and the stated coarse and sector
queries. Its declared coarse footprint includes
all these read sets, the windows defining its
centre and covariance entries, and a Gaussian
range enlargement.
Take the enlargement large enough that
disjoint footprints imply
\[
 (C_{\rm sel})_{ij}=0
 \quad(i\in T_\alpha,\ j\in T_\beta)
\]
for letters in different footprint components.
A fixed extra \(B\)-neighbourhood suffices.

\begin{proposition}[Physical component factorisation]
\label{prop:finite-range-physical-factorisation}
For fixed \(V\) and sector geometry, integrate
all letters with respect to one selected
Gaussian law, the original product Haar law
of frozen links, and the product law of the
shared local branch variables.
If a finite word has disjoint full-footprint
components \(A_1,\ldots,A_k\), and its absolute
integral is finite, then
\begin{equation}
 \int\prod_{\alpha\in A}f_\alpha
          =\prod_{\nu=1}^k
                    \int\prod_{\alpha\in A_\nu}f_\alpha .
 \label{eq:physical-component-factorisation}
\end{equation}
Each factor on the right uses the marginal
of that same joint law.
\end{proposition}
\begin{proof}
The Gaussian covariance on the union of
the components' coordinate sets is block
diagonal. Its characteristic function
\(\exp(-\langle t,C_{\rm sel}t\rangle/2)\)
is therefore the product of those of
the component marginals. This proves
their independence.
The frozen Haar variables and local branch
variables read by different components
are disjoint. Their product laws also
factor. The matrix \(D\), and hence the
Gaussian covariance, do not depend on
these frozen values or branch choices;
the centre's dependence on them has
already been included in each letter's
read set. The joint marginal is therefore
the claimed product.
Fubini's theorem, using the absolute
integrability hypothesis, proves the
displayed identity.
\end{proof}

Local branch weights in this proposition
are kept as factors. More explicitly, the
actual branch variable consists of one
frame tag \(t_x\) at each coarse site and
one predictor tag \(t_e\) at each coarse
link. Its likelihood has the form
\[
 \lambda_t(U)=
 \prod_x p_x(U,t_x)
 \prod_{e=(x,i)}
       p_e(U,t_x,t_{x+\widehat i},t_e).
\]
Here \(p_x\) is the prescribed frame
probability on its finite patch, and
\(p_e\) is the predictor probability on
its connector patch. For every \(U\),
\(\sum_{t_x}p_x(U,t_x)=1\), and for
every fixed pair of frame tags,
\(\sum_{t_e}p_e(U,t_x,t_{x+\widehat i},t_e)=1\).
This is a product of local conditional
weights. It is represented by a product
uniform law on all tags, with the
conditional weights inside the letters;
no independence of the weighted tags is
asserted. If branch choices
are integrated with uniform probability
law instead of summed, multiply each
local likelihood by the cardinality of
its choice set. Their original normalisation
is then unchanged. A shared frame choice
is one variable, including in neighbouring
predictor factors; it is not copied for
each letter.

There is a useful consistency consequence.
If two sectors agree on a component's
full read region, the principal matrices
needed for its covariance marginal are
identical, and so are its selected centre,
frozen-coordinate reads, and test values.
Its integral is therefore the same in
both sectors. This follows from the
finite-window formulas, rather than from
an assumption about a restricted precision
chosen afresh for that component.

\subsection{Uniformity and the remaining physical payments}

The parameters \(\eta,C,\theta,\lambda\) above
depend on the group, profile, block, and
funded path bounds. The polynomial \(C_B\)
also depends on the fixed derivative and
covering exponents. All coupling thresholds
and the period reserve are chosen before
the ultraviolet cutoff and physical period.
The slow-clock argument at the end of
Section~\ref{sec:physical-interacting-extensions}
places these finitely many additional
thresholds on the same Wilson family.

The bridge is now exact for Borel incoming
activities, and the finite-range mechanism
for their component factorisation is available.
It has not removed the physical costs in
\eqref{eq:finite-range-vacuum-ledger}.
In particular, integration of the nonlinear
phase difference together with the frozen
Haar variables remains a separate physical
estimate. The actual fine-field, noise,
branch, and chart failures outside holes
have now been paid by their selected-centre
coordinate boxes. The remaining packet payments must
precede an unconditional output-activity
bound. The reassembly calculation in the
next section states precisely how those
payments enter the represented density.
\section{Incoming activities and the physical reassembly identity}
\label{sec:physical-reassembly}

We now give the local form of the incoming factors
in the selected coordinates, and the exact
reassembly after connected components have been
integrated. The reassembly theorem includes its
norm criterion. Its use for the complete nonlinear
step still requires paying the actual hole packets;
that remaining estimate is identified explicitly
at the end of the section.

\subsection{Incoming masks and the fine-flag prohibition}

Write \(e_i=\|E_i\|_{q,h,s;o_i,X_i}\) and
\(k_\alpha=\|K_\alpha\|_\infty\).
The incoming regular Mayer factor is
\begin{equation}
 m_i(U)=e^{-\chi_{X_i}(U)E_i(U)}-1 .
 \label{eq:incoming-regular-Mayer}
\end{equation}
If \(e_i\le1\), then
\begin{equation}
 |m_i(U)|\le e^{e_i}-1\le2e_i .
 \label{eq:incoming-regular-point-price}
\end{equation}
On a bad footprint the masked activity is zero,
so the same inequality holds there.
This bound does not differentiate the mask.
Since \(\sum_i e_i<\infty\) on a finite torus,
\[
 e^{-E^\#(U)}
       =\sum_{A\ {\rm finite}}\prod_{i\in A}m_i(U)
\]
with uniform absolute convergence.

The exact-coverage condition in the incoming
forced gas is global in its notation, but can
be expressed by local tests depending on the
selected labels. Let \(\Gamma\) be a finite
set of incoming forced labels with pairwise
disjoint original supports, and put
\(Y_\Gamma=\bigcup_{\alpha\in\Gamma}Y_\alpha\).
For each fine site \(x\), define
\[
 q_x(U;\Gamma)=
 \begin{cases}
  1,&x\in Y_\Gamma,\\
  \chi_{\{x\}}(U),&x\notin Y_\Gamma .
 \end{cases}
\]
Then the incoming gas has the exact expression
\begin{equation}
 Z_K(U)=
 \sum_{\Gamma\ {\rm original\ supports\ disjoint}}
       \left(\prod_{\alpha\in\Gamma}K_\alpha(U)\right)
                         \prod_x q_x(U;\Gamma).
 \label{eq:incoming-gas-local-prohibition}
\end{equation}
Indeed, a nonzero selected activity already
imposes its own inventory and additional test.
On disjoint selected supports, these inventories
are exactly the fine flags in \(Y_\Gamma\).
The product of \(q_x\)'s prohibits every fine
flag outside that union. Together they are
equivalent to the original exact coverage.
Terms with a zero selected activity vanish
on both sides.

The factors \(q_x=1+(q_x-1)\) can therefore
be included in a local word expansion.
Their selection query asks whether a selected
original support contains \(x\).
A letter and every selected label which can
answer that query must belong to the same
full-footprint component.
This requirement is retained even when a
numerical test happens to equal one.

Original supports still determine the
incoming \(K\)-versus-\(K\) exclusion.
Their enlarged read footprints determine
which letters must be integrated in the
same component. Two compatible original
labels may have intersecting enlarged
footprints; in that case they merge into
one component. They are not discarded.

Substitute
\(U=\mathcal U_\sigma(V,Y,\mu+gz)\)
in these exact factors. By the finite-window
construction their complete reads are finite.
The bounds \(|K_\alpha|\le k_\alpha\) and
\eqref{eq:incoming-regular-point-price}
hold for every \(z,Y,V\) on the support
of the corresponding factor.
No derivative, analytic continuation, or
regularity of \(K_\alpha\) is used.

\subsection{Covering and incidence costs of the enlarged reads}

Let \(B_f\) denote the incoming covering radius,
and retain \(B_f\le2B\) from the adjacent-scale
schedule. The read enlargement has a fixed
radius \(c_*B\) in coarse units; it includes
windows, row stencils, flag tests, and the
Gaussian range enlargement.
Its constant can be increased finitely many
times to contain all these queries.

For a nonempty original footprint \(X\),
choose a fixed minimal connected \(B_f\)-ball
cover, include the union of that cover in
the declared geometry, project it to coarse
parents, and take its \(c_*B\)-neighbourhood.
Call the resulting footprint \(\widehat X\).
This harmless geometric enlargement also
works for an originally disconnected \(X\).

\begin{lemma}[The enlarged input income]
\label{lem:incoming-enlarged-income}
There is a geometric integer \(c_1\), independent
of volume and cutoff, such that
\[
 n_B(\widehat X)\le c_1 n_{B_f}(X).
\]
Suppose incoming activity bounds are available
at exponent \(\kappa_{\rm in}\), and choose
\[
 \kappa_{\rm in}\ge c_1(\kappa_{\rm out}+a)+1 ,
 \qquad a>0 .
\]
Then the incidence sums of the enlarged
incoming value prices, at exponent
\(\kappa_{\rm out}+a\), are bounded by
\begin{equation}
 C_B\|E\|_{\mathcal E,\kappa_{\rm in}},
 \qquad
 C_B\|K\|_{\mathcal K,\kappa_{\rm in}} ,
 \label{eq:incoming-enlarged-income}
\end{equation}
respectively. For regular Mayer factors,
multiply the first bound by two.
\end{lemma}
\begin{proof}
A fine \(B_f\)-ball projects into a coarse
ball of radius at most \(2B+2\), since
\(L\ge3\). Its further \(c_*B\)-enlargement
is covered by a fixed number of \(B\)-balls.
The projected union of the original connected
cover, together with this enlargement, is
connected. This proves the covering estimate.

Put \(k=n_{B_f}(X)\). If a coarse site \(x\)
belongs to \(\widehat X\), its distance from
the parent of any chosen site of \(X\) is
at most \(C Bk\). There are at most
\(C_L B^4(1+k)^4\) possible fine sites in
that region. Use the original anchor for
an \(E\)-label; for a \(K\)-label select one
site of its nonempty original support.
The new covering weight is at most the
old one times \(e^{-k}\).
For fixed \(k\), sum the original activity
prices at each of the possible fine sites.
The original anchored bound for \(E\), or
the original incidence bound for \(K\),
gives its corresponding norm.
Finally
\[
 \sum_{k\ge1}C_LB^4(1+k)^4e^{-k}\le C_B .
\]
This also bounds the sum over every label
multiplicity; no bound on the number of
distinct supports was assumed.
\end{proof}

These are value estimates for the incoming
factors in a word. They do not yet prove
the regular derivative contraction after
integration and coupling reset. In particular
one must not differentiate a masked incoming
factor in order to infer that contraction.

\subsection{An integrated component and its price}

Fix a finite selected word \(A\).
Its letters include the incoming factors,
Gaussian density-correction letters, local
physical failure factors, and hole-packet
factors. Its internal tests retain original
compatibility and all selection queries.
Let \(X_\alpha\) be the complete read
footprint of a letter.
A component is a nonempty connected set
of letters in the overlap graph of these
footprints, with their full labels retained.

For a component \(P\), define its amplitude
by integrating its product of letters
against the common marginal of
Section~\ref{sec:finite-range-components},
including its internal tests.
Write the amplitude as \(a_P(V)\), and its
footprint as \(X(P)=\bigcup_{\alpha\in P}X_\alpha\).
The consistency at the end of
Proposition~\ref{prop:finite-range-physical-factorisation}
allows this definition to be made independently
of distant components.

Here is the precise estimate used for the
Gaussian part. If, conditionally on frozen
variables and the common branch choices,
every selected product obeys
\[
 \left|\prod_{\alpha\in P}f_\alpha\right|
 \le\left(\prod_{\alpha\in P}a_\alpha\right)
          e^{\frac12\sum_i d_i z_i^2},
 \qquad 0\le d_i\le u_0,
\]
then its Gaussian absolute integral is at
most
\[
 \left(\prod_{\alpha\in P}a_\alpha\right)
                  e^{(2r_0)^{-1}\sum_i d_i}.
\]
Bounded incoming factors contribute \(2e_i\)
and \(k_\alpha\), with zero quadratic charge.
Failure-witness factors have the charges and
prices in
Lemma~\ref{lem:finite-range-failure-witness}.
Other Gaussian corrections must use the
remaining precision reserve.
Integration of the frozen-variable majorant
is still required; it is not replaced by a
Gaussian tail probability.

If the resulting per-letter prices have
incidence sum \(\varepsilon\) at a sufficiently
larger covering exponent, the connected-family
bound of Section~\ref{sec:connected-bound}
gives a rooted component sum \(C_B\varepsilon\).
The Gaussian failure-only income in this
calculation is
\(C_B\exp\{-P(\widehat\beta)^2/C_B\}\).
The incoming value incomes are
\eqref{eq:incoming-enlarged-income}.
An unpaid hole-packet majorant is not included
in \(\varepsilon\) by this statement.

Absolute convergence supplied by these prices
justifies the component factorisation,
summation over words, and integration.
For finite truncations of the alphabet the
same reorganisation is elementary. An
infinite alphabet requires the absolute
majorant, rather than a purely formal
exchange of series and integrals.

\begin{corollary}[Failure components with bounded incoming factors]
\label{cor:Gaussian-failure-incoming-components}
Consider the selected Gaussian correction
letters, the Gaussian-coordinate failure
letters, and the bounded incoming factors
just described. Keep all their masks and
internal tests. Do not include an unpaid
physical phase or hole-packet factor.
The failure letters may also be the
actual fine-field, noise, branch, and
broad-chart failures outside the
protected hole neighbourhood, with
their original bounded protections,
by Lemma~\ref{lem:selected-centre-intact-box}.
After the common quadratic charges have
been integrated, let \(b_0\) be the
incidence income of the nonfailure
letters at exponent \(\kappa+1\), and
let \(f_0\) be that of the failure
letters. If
\(D_Bb_0\le1/4\) and \(D_Bf_0\le1/4\),
then the components containing a
failure label obey
\begin{equation}
 \sup_x\sum_{\substack{P:x\in X(P)\\
                          P\ {\rm contains\ a\ failure}}}
       e^{\kappa n_B(X(P))}\|a_P\|_\infty
                                      \le4f_0 .
 \label{eq:Gaussian-failure-component-income}
\end{equation}
Here one may take
\[
 b_0\le C_B\left(e^{-cB}
       +\|E\|_{\mathcal E,\kappa_{\rm in}}
       +\|K\|_{\mathcal K,\kappa_{\rm in}}\right),
 \qquad
 f_0\le C_Be^{-P(\widehat\beta)^2/C_B}.
\]
All constants are uniform in the
cutoff and volume at a fixed profile.
The estimate is uniform for each fixed
sector geometry. Summing full sector
geometries requires the physical local
component dictionary and its prices;
no multiplicity of such geometries is
silently omitted from an income.
\end{corollary}
\begin{proof}
The joint occupation bounds proved
above permit the correction letters
and failure witnesses to be integrated
in one use of the precision reserve.
Their sale prices multiply the
bounded incoming prices \(2e_i\)
and \(k_\alpha\). Haar probability
and uniform tag probability have
mass one; the original local
likelihood factors, when present,
must be included in their own
priced letters. Internal zero-one
tests can only decrease an absolute
majorant.

If \(f_0=0\), every component in the
sum has zero price. Otherwise put
\(t=(4D_Bf_0)^{-1}\ge1\), and multiply
every failure price by \(t\).
The total modified incidence is at
most \(1/(2D_B)\).
Every component in the displayed
sum gains at least one factor \(t\).
Lemma~\ref{lem:connected-bound} gives
its original rooted sum at most
\[
 t^{-1}(e^{b_0+tf_0}-1)
 \le t^{-1}\,2(b_0+tf_0)
 \le4f_0.
\]
The bounds for \(b_0,f_0\) follow
from Lemma~\ref{lem:incoming-enlarged-income},
Lemma~\ref{lem:selected-correction-word-prices},
and the failure-witness estimate,
using the available covering reserves.
The same majorants justify all
countable sums and suprema in the
statement.
\end{proof}

The scope of this corollary matters.
It proves the Gaussian failure
component estimate jointly with
bounded incoming multipliers.
It does not estimate the same
components under a further,
unbounded and unpaid physical
weight. Such a weight has to
receive its own joint price
before the corollary is applied
to the complete density.

\subsection{Unforced components and forced components}

Let
\(\mathcal D=\mathcal B_{h_{\rm c}}(V)\)
be the actual coarse flag set.
Each component has a prescribed coarse
inventory \(I(P)\subset X(P)\), carried by
its mandatory coarse-flag letters.
Fine incoming inventories and microscopic
Gaussian failures are not automatically
coarse inventories.
A component with \(I(P)=\varnothing\) is
unforced; one with nonempty inventory is
forced.

Let the integrated unforced species have
amplitudes \(u_i(V)\) and full footprints
\(X_i\). Their occurrence is masked by
\(\mathbf1_{\{X_i\cap\mathcal D=\varnothing\}}\).
Let the forced species have amplitudes
\(d_j(V)\), footprints \(Y_j\), and nonempty
inventories \(I_j\subset Y_j\). Their
eligibility includes
\(Y_j\cap\mathcal D=I_j\).
When the physical word expansion and its
component prices have been established,
its mixed gas has the form
\begin{equation}
 \Xi(V)=\sum_{\Gamma\ {\rm forced\ admissible}}
       \prod_{j\in\Gamma}d_j(V)\,
           Z_{\rm u}(\Lambda\setminus Y_\Gamma;V),
 \quad Y_\Gamma=\bigcup_{j\in\Gamma}Y_j .
 \label{eq:physical-mixed-gas}
\end{equation}
Admissibility uses disjoint forced
footprints and exact coverage of
\(\mathcal D\).
The unforced hard-core gas on the right
contains only species disjoint from
\(Y_\Gamma\), and from \(\mathcal D\).

Formula \eqref{eq:physical-mixed-gas} is the
reassembly input, not an assumption that
all Gaussian failures are forced.
It permits Gaussian failures and incoming
fine forced factors to contribute to the
unforced connected logarithm when their
coarse inventory is empty.

\subsection{Both exclusion compensations in the exact density}

Suppose the unforced gas satisfies the
rooted logarithm criterion of
Lemma~\ref{lem:rooted-interacting-log}.
Index its log clusters by \(i\), retaining
repeated species and the usual Ursell
coefficient, and denote the signed
cluster function by \(\ell_i(V)\).
Its footprint, again denoted by \(X_i\),
is the union of all the cluster reads.
Then, for a forbidden set \(Y\),
\[
 \log Z_{\rm u}(\Lambda\setminus Y;V)
   =\sum_{\substack{i:X_i\cap Y=\varnothing\\
                        X_i\cap\mathcal D=\varnothing}}\ell_i(V).
\]
This is the ordinary hard-core logarithm
with species meeting \(Y\) set equal to
zero. Its convergent expansion is described
in \cite[Sect.~2, eqs.~(2.3) and (2.6)--(2.7),
preprint pp.~3--4]{FP07}; the rooted estimate
needed here was proved in
Section~\ref{sec:full-interacting-remainder}.

Put
\[
 c_i=\ell_i(1),\qquad
 E_i^{\rm new}=-\ell_i+c_i,\qquad
 v_{\rm log}=\sum_i c_i .
\]
Here \(v_{\rm log}\) is an extensive scalar;
the increment per coarse site is
\(v_{\rm log}/\#\Lambda\).
This step subtracts only the flat scalar.
The affine quadratic extraction is the
separate coupling reset of
Section~\ref{sec:coupling-reset}.
At this stage these are local value
functions. The identities below do not
yet assert their full regular-family
derivative norm. For that assertion
one must construct the unmasked
\(C^q\) germs on healthy footprints
and prove the physical measure
derivative estimates. In particular,
one cannot obtain those germs by
differentiating a coarse retained-row
indicator at its threshold.
Define the restoration function
\begin{equation}
 t_i(V)=
 \begin{cases}
   -\ell_i(V),&X_i\cap\mathcal D=\varnothing,\\
   -c_i,&X_i\cap\mathcal D\ne\varnothing .
 \end{cases}
 \label{eq:physical-two-restorations}
\end{equation}

\begin{proposition}[Exact physical exclusion restoration]
\label{prop:physical-exclusion-restoration}
Under absolute convergence of the mixed gas
and of its log expansion,
\begin{align}
 \Xi(V)
 &=e^{\,v_{\rm log}-(E^{\rm new})^\#(V)}
   \sum_{\Gamma\ {\rm forced\ admissible}}
      \prod_{j\in\Gamma}d_j(V)
      \exp\left\{\sum_{i:X_i\cap Y_\Gamma\ne\varnothing}
                                                 t_i(V)\right\}.
 \label{eq:physical-exact-restoration}
\end{align}
The identity retains both the healthy
exclusion and the unmasked vacuum scalar.
\end{proposition}
\begin{proof}
For a fixed admissible \(\Gamma\), its
unforced log equals the full healthy log
minus all healthy log clusters meeting
\(Y_\Gamma\).
Also
\[
 \sum_{i:X_i\cap\mathcal D=\varnothing}\ell_i
   =v_{\rm log}-(E^{\rm new})^\#
                   -\sum_{i:X_i\cap\mathcal D\ne\varnothing}c_i .
\]
Every bad \(X_i\) meets \(Y_\Gamma\):
choose a coarse flag in \(X_i\), then use
the exact inventory coverage by \(\Gamma\).
Thus its scalar correction is included
among precisely the second-branch terms
of \eqref{eq:physical-two-restorations}
meeting \(Y_\Gamma\).
The healthy exclusions are the first-branch
terms there. Combining the three sums
proves the exponential identity for this
\(\Gamma\). Absolute convergence justifies
summing it over \(\Gamma\).
\end{proof}

Keeping only \(-\ell_i\) on healthy
footprints would miss the scalar
\(-c_i\) on bad footprints.
Keeping only the scalar correction would
allow unforced log clusters to pass
through forced footprints. Neither
omission preserves the density.

\subsection{Merging the restorations and its local norm}

Let \(b_i\ge0\) bound the actual healthy
value \(|\ell_i(V)|\) and \(|c_i|\).
Suppose their incidence sum at exponent
\(\kappa+1\) is at most \(b\le1\).
Let \(a_j\ge0\) bound \(|d_j(V)|\),
with corresponding incidence sum at
most \(a\).
These are prices of the integrated
physical components, including their
Gaussian/Haar normalisation.

Expand the exponential in
\eqref{eq:physical-exact-restoration}
with Mayer factors \(e^{t_i}-1\).
Only restoration addresses meeting the
selected original forced family may
be selected. A new connected label \(C\)
is a connected finite collection of
forced labels and restoration labels,
such that its forced labels are
compatible and every restoration
meets at least one of those forced
labels. It contains at least one
forced label.
Its footprint and inventory are
\[
 Y_C=\bigcup_{j\in C\ {\rm forced}}Y_j
                      \cup\bigcup_{i\in C\ {\rm rest.}}X_i,
 \qquad
 I_C=\bigcup_{j\in C\ {\rm forced}}I_j .
\]
The new activity is its product of
forced amplitudes and restoration
Mayer factors, multiplied by its
internal tests and by
\(\mathbf1_{\{Y_C\cap\mathcal D=I_C\}}\).

\begin{theorem}[Physical merger with a prescribed income]
\label{thm:physical-merger-income}
For the preceding component prices,
choose \(\eta\ge0\), put
\(\gamma=e^\eta a+2b\), and assume
\(D_B\gamma\le1\). Then the new forced
gas represents the sum in
\eqref{eq:physical-exact-restoration}
exactly, and
\begin{equation}
 \|K^{\rm new}\|_{\mathcal K,\kappa}
                   \le e^{-\eta}(e^\gamma-1).
 \label{eq:physical-merger-income}
\end{equation}
The regular action and scalar in
\eqref{eq:physical-exact-restoration}
are unchanged by this merger.
\end{theorem}
\begin{proof}
Each restoration satisfies \(|t_i|\le b_i\),
so \(|e^{t_i}-1|\le2b_i\).
For finite alphabets, multiply the
forced gas by the product expansion
of its allowed restoration factors.
Decompose each selected term into
overlap components. Every component
contains an original forced label,
because each restoration was required
to meet one. Different components
have disjoint full footprints.

Every coarse flag in a component's
footprint lies in an original inventory
from that same component: its original
forced footprint meets the component
at that flag and hence joins it.
Conversely, flattening an admissible
new family recovers a compatible
original forced selection and its
allowed restoration selection.
No restoration can appear in two
different components, since its
nonempty footprint would intersect
both. These operations are inverse
and preserve every signed product.

Give a forced label weight \(e^\eta a_j\)
and a restoration weight \(2b_i\).
Every new component has a forced label,
so its absolute weight is at most
\(e^{-\eta}\) times the product of
these combined weights.
Their incidence income is \(\gamma\).
The connected-family estimate used in
Theorem~\ref{thm:forced-reassembly}
therefore gives
\(e^{-\eta}(e^\gamma-1)\).
It also gives an absolutely summable
majorant. Passing from finite
alphabets proves both the identity
and the bound for countable alphabets.
\end{proof}

The functions \(t_i\) on a healthy
\(X_i\) need not vanish. Accordingly
they are attached only to the selected
forced supports; they are not treated
as unconditional corrections supported
on the bad set. This is the additional
selection rule needed beyond the
simpler reassembly in Section
\ref{sec:one-step-reassembly}.

\subsection{What has been bridged and what must still be paid}

The preceding construction resolves the
locality obstruction for bounded Borel
incoming functions. The selected variables
are physical fluctuation coordinates,
their centre has finite reads, and the
one joint Gaussian law has finite range.
The original masks and incoming activities
appear in the exact sector formula
\eqref{eq:full-physical-selected-Gaussian-sector}.
The Gaussian component bound and the
two signed reassembly corrections are
now explicit.

For an unconditional one-step RG bound,
the remaining physical estimate is the
combined price of each hole packet.
It must include the frozen-link ratio
\(s_g^{-n_{\rm fr}}\), its actual nonlinear
row differences, the constant affine
square, branch likelihoods, and any
incoming multipliers meeting that packet.
In particular an estimate on the number
of holes alone does not control these
factors. One must prove their joint
Gaussian/Haar absolute-integral bound,
with enough incidence reserve for the
numbers \(a,b\) in
Theorem~\ref{thm:physical-merger-income}.
The physical failures outside the
protected hole neighbourhood have
been checked on local Gaussian
boxes using the selected centre
of that same sector. Tests inside
the hole neighbourhood belong to
its packet payment.

These estimates are not premises silently
added to the construction theorem.
The packet payment is established in
Section~\ref{sec:hole-packet-integrals}.
Section~\ref{sec:integrated-regular-derivatives}
then treats the regular derivatives of
components containing a physical packet.
The theorem above remains an exact
reassembly and a quantitative closure
criterion. Full nonlinear RG closure
also needs the sharp central transport
of incoming regular activities and
its extraction estimates.
\clearpage
\section{Hole packets and the joint Gaussian--Haar integral}
\label{sec:hole-packet-integrals}

The remaining price in Section~\ref{sec:physical-reassembly}
is local, but it is not a Gaussian tail probability alone.
A hole removes some of the Gaussian coordinates and leaves their
physical links under probability Haar measure. Relative to the common
vacuum normalisation this costs \(s_g^{-1}\) per removed link.
In addition, the nonlinear action and the affine reference action
have to be compared on the boundary of that hole.
We first keep these terms in a single exponent. This avoids estimating
the negative affine square separately from the physical energy that
pays it.

For comparison with large-field integration and component suppression
in lattice gauge renormalisation, see
\cite[Sect.~1, eqs.~(1.71)--(1.75), printed pp.~379--380;
large-field factors on pp.~380--383; eqs.~(1.98)--(1.100),
p.~390]{Bal89b}.
This comparison concerns the organisation of large-field integration,
energy suppression, and exponentiation.
All estimates below concern the present clipped action, predictor,
and sector-dependent Gaussian law. They do not import an estimate
for another action or for independent copies of its Gaussian field.

\subsection{Physical seeds, ownership, and the exact packet exponent}

Use the fine quantities \(h,\widehat\beta_{\rm f},P\), the actual
coupling \(\beta=g^{-2}\), and the coarse threshold \(h_{\rm c}\)
of the preceding sections. Thus
\[
 \tfrac12\widehat\beta_{\rm f}\le\beta
       \le2\widehat\beta_{\rm f},\qquad
 \tfrac12 P^2\le\beta h^2\le2P^2,\qquad
 h_{\rm c}\le2h
\]
after the already prescribed fine-scale nominal shift.
Here and below \(C_B\) means \(C(1+B)^p\), with the exponent and
coefficient fixed at the current profile and derivative order.
It never denotes a constant depending on the number of lattice sites.

Choose a smooth function \(\chi:[0,\infty)\to[0,1]\) equal to one
on \([0,1]\) and zero on \([2,\infty)\).
For a fine plaquette put \(s_p=d_G(P_p(U))\), and for a coarse
bond put \(q_b=d_G(V_bQ_t(U)_b^{-1})\).
Resolve unity using
\begin{equation}
 \chi(4s_p/h)+(1-\chi(4s_p/h))=1,\qquad
 \chi(q_b/(Wh))+(1-\chi(q_b/(Wh)))=1.
 \label{eq:packet-smooth-sector-tests}
\end{equation}
These functions are smooth functions of the matrix variables:
their only possible nonsmooth radial point lies where they are
constant. An exact coarse bad-site flag is still a Borel indicator;
it is a different test and is not replaced by \(\chi\).

A \emph{seed record} consists of a selected fine plaquette, a selected
noise bond, or an actual coarse bad site. Selection refers respectively
to the second fine factor, the second noise factor, and the exact
coarse flag. Records are distinct even if their sites coincide.
For a collection \(p\) let \(m(p)\) be its number of records.
Project the vertices of a fine record to their coarse parents,
use the endpoints of a noise record, and retain a coarse-site record
itself. These are the locations of \(p\).
Write \(p^{[r]}\) for their coarse \(r\)-neighbourhood.

Join seed records whose locations are at distance at most \(500B\),
and take connected components of this graph. A nonempty such
component is called a \emph{geometric hole packet}.
On its \(450B\)-neighbourhood every selected seed belongs to this
same component. In particular the local tests imply
\begin{equation}
 \begin{array}{lll}
 p\text{ selected fine}:&s_p\ge h/4,
     &p\text{ unselected fine}:s_p\le h/2,\\
 b\text{ selected noise}:&q_b\ge Wh,
     &b\text{ unselected noise}:q_b\le2Wh .
 \end{array}
 \label{eq:packet-physical-plot}
\end{equation}
The weak inequalities include all boundary points of the tests.
Only tags with positive likelihood need to be considered.
The exact incidence \(Y_\alpha\cap\mathcal B_h(U)=I_\alpha\)
of each selected incoming activity is retained.
In this physical-sector partition take \(\mathcal H_\sigma\)
to be the union of these record locations. The default
fine and noise tests outside its neighbourhood imply the
active-link chart condition by
Lemma~\ref{lem:packet-annulus-payment}. Thus no independent,
unpaid chart-complement holes are introduced here.
Gaussian-coordinate failures are subsequently optional
letters; they do not enlarge this frozen-link set.

Each energetic row has an anchor, including its orientation and
row type. Assign a row to the packet whose \(200B\)-neighbourhood
contains its anchor. Distinct geometric packets have disjoint
such neighbourhoods. The actual energetic rows will be denoted
by \(w_r\), and their cut affine comparison rows by
\(D_ra+\alpha_r\). A nonoperated row has comparison row zero.
The packet energy that occurs in the exact integral is
\begin{equation}
 {\cal R}_p(a,Y,V,t)
 =\sum_{r\ {\rm assigned\ to}\ p}
 \left\{w_r(U,V,t)
       -\tfrac12|D_ra+\alpha_r|^2
       +\tfrac12|D_r\mu+\alpha_r|^2\right\}.
 \label{eq:packet-completed-energy}
\end{equation}
The same orientation average as in the action is understood in
this formula. The rows \(w_r\) include the coarse subtraction
and the row-calibration defects, not just positive fine rows.
The last term in \eqref{eq:packet-completed-energy} is the local
part of the completed-square constant in the exact Gaussian
identity. It cannot be dropped before a lower bound is obtained.

Let \(J_p\) be the frozen free fine edges assigned to \(p\).
Tree edges remain fixed at one. Simple counting gives
\begin{equation}
 |J_p|\le C_L(1+B)^4m(p).
 \label{eq:packet-frozen-count}
\end{equation}
There are at most \(C_L\) fine edges per coarse cell, and the
union of the \(10B\)-balls about the record locations has at most
\(C(1+B)^4m(p)\) cells. This proves the assertion also for
overlapping balls. Enlarging a footprint later because an incoming
activity reads far away does not freeze those extra edges.

The exact packet multiplier has the form
\begin{equation}
 \Theta_p
 = {\bf1}_p\,L_p\,J_p^{\rm Haar}\,
   s_g^{-|J_p|}
   \exp\{-\beta{\cal R}_p+\Delta_p\}\,{\mathsf M}_p .
 \label{eq:packet-exact-multiplier}
\end{equation}
Here \({\bf1}_p\) denotes all original sector tests and fine-flag
prohibitions assigned to this packet, \(L_p\) is the product of
the actual scaled frame and predictor likelihoods,
\(J_p^{\rm Haar}\) is the product of active-link Haar Jacobians
with their broad chart factors, and \(\Delta_p\) is its assigned
paired local determinant contribution. The symbol \({\mathsf M}_p\)
denotes the selected incoming multipliers and the microscopic
terms meeting this packet. Their precise sums are bounded below.
All these factors are evaluated at the same physical
\(U=\mathcal U_\sigma(V,Y,\mu+gz)\).
This is a multiplier relative to the coarse principal
factor and the common vacuum factor
\[
 e^{v n_f}z_{\mathrm n}(\beta)^{-b_c}
                 s_g^{n_0}(\det\mathsf M_0)^{-1/2}
 \]
in the exact sector identity. Both stay outside the
packet gas with their original values. In particular
no power of the noise normaliser \(z_{\mathrm n}(\beta)\) has
been estimated away.

The probability measure of integration is
\begin{equation}
 \prod_{e\in J_\sigma}\dd h_G(Y_e)\ \dd\nu(t)\
                   \dd\mu_{C_{\rm sel}}(z).
 \label{eq:packet-one-joint-law}
\end{equation}
The tag law \(\nu\) is the product uniform law used to express
the branch sum; the scaled likelihoods remain in \(L_p\).
In particular the tags, the frozen links in the source
\(\alpha\), and the Gaussian field are not resampled for
different factors.

\subsection{The bills that a packet must pay}

We separate elementary integration from the model-specific
energy inequality. This also specifies exactly the latter's
required strength.

\begin{lemma}[Local bills]
\label{lem:packet-local-bills}
Fix the profile and \(L,W\). There are an integer \(n_0\ge0\)
and \(H_0>0\), independent of cutoff and volume, such that
for \(n\ge n_0\) and \(H\ge H_0\) the bill
\begin{equation}
 H(1+B)^n(1+\log\widehat\beta_{\rm f})m(p)
 \label{eq:packet-general-bill}
\end{equation}
can include the frozen-link ratio, all scaled local likelihoods,
all broad-chart Haar Jacobians, and the assigned local
determinant terms in \eqref{eq:packet-exact-multiplier}.
The choices of \(H,n\) are independent of ultraviolet cutoff
and volume.
\end{lemma}
\begin{proof}
The exact constant is
\(s_g=\zeta_G(\sqrt{2\pi}g)^{d_G^{\rm Lie}}\).
Consequently, for \(g<1\),
\[
 \log(s_g^{-|J_p|})
 \le C_G|J_p|(1+\log(1/g))
 \le C_B(1+\log\widehat\beta_{\rm f})m(p).
\]
The first inequality is an upper bound even when \(s_g>1\).
The logarithm of a scaled local likelihood is bounded above
by the logarithm of the size of its finite tag alphabet.
At a fixed profile and \(L\) this is a fixed constant.
There are at most \(C_Bm(p)\) factors.
The broad chart lies in a fixed injectivity neighbourhood;
the smooth Haar density is bounded above there.
Its product therefore costs at most \(\exp(C_Bm(p))\).

Finally use the diagonal local determinant representation of
Section~\ref{sec:finite-range-gaussian}. The principal window
matrices and their interpolation have a common spectral gap.
For a diagonal local trace its absolute value is bounded by
the product of the two weighted Schur bounds. The integral in
the interpolation parameter has length one. Only local
diagonal terms whose full read region meets the packet
are assigned to it; all other terms retain their
good-background letters. The number of assigned
zero-length anchors is \(C_Bm(p)\).
The exponentially small terms of positive path length have
an absolutely summable geometric majorant. Thus
\(|\Delta_p|\le C_Bm(p)\).
This argument pairs the flat and physical contributions
before summing their absolute values.
It introduces no determinant factor proportional to the
volume outside the packet.
\end{proof}
\begin{lemma}[Incoming occurrences at a packet]
\label{lem:packet-incoming-occurrences}
Let \(w\) be the incoming forced incidence norm and \(\delta\)
the incoming regular norm, with the footprint reserve used in
Section~\ref{sec:physical-reassembly}. The sum of absolute
products of selected nonzero incoming forced occurrences
attached to packets with a total of \(m\) seed records is at
most
\begin{equation}
                         \exp(4wm).
 \label{eq:packet-incoming-K-income}
\end{equation}
Regular and microscopic multipliers meeting the hole
neighbourhoods cost at most
\(\exp\{C_B[\delta+(1+\Lambda)h^2+
              \beta C_{\rm p}h^3+\beta e^{-cR_B}]m\}\).
Here \(\Lambda\) is the loop coefficient budget and
\(C_{\rm p}\) the fixed principal path budget from the
microscopic estimates. The allowed subpolynomial growth
of \(\Lambda\) is retained; \((1+\Lambda)h^2\to0\)
along the chosen clock. The last term keeps the omitted
native row tails of the original action. Here \(c>0\)
may depend on the fixed \(L\): coarse path cuts still
count native fine length, including the factor \(L\).
Original support exclusions and flag masks remain in the
signed integral.
\end{lemma}
\begin{proof}
Every nonzero selected \(K_\alpha\) has a nonempty bad
inventory \(I_\alpha\). Choose a bad vertex in it. An incident
fine plaquette has chord \(>h\), so its smooth sector factor
forces it to be a selected fine seed. There are at most
\(4m\) possible fine vertices among these seed plaquettes.
Assign the occurrence to the first such vertex in a fixed
ordering. At each vertex the sum of its absolute occurrence
prices is at most \(w\). Dropping exclusions only for an
upper bound, and using \(1+x\le e^x\), gives
\[
 \sum_{\Gamma}\prod_{\alpha\in\Gamma}\|K_\alpha\|_\infty
 \le\prod_{\alpha:\,Y_\alpha\ {\rm meets\ a\ chosen\ vertex}}
                 (1+\|K_\alpha\|_\infty)
 \le e^{4wm}.
\]
An occurrence that meets several packets is counted once.
In the component expansion it joins those packets; it is not
replaced by several copies.

For the regular family use its value incidence bound and the
anchor-to-incidence counting with exponential covering reserve
proved in Section~\ref{sec:physical-reassembly}.
The number of possible hole-neighbourhood sites is \(C_Bm\).
Hence the sum of \(|E_i|\) that must be placed in the packet
exponent is at most \(C_B\delta m\).
The microscopic bounds give \(C_B(1+\Lambda)h^2m\)
for loops and \(C_B\beta C_{\rm p}h^3m\) for the
masked higher-degree principal terms.
The absolute native row tail is at most \(C e^{-cR_B}\)
per anchor; multiplying by the actual \(\beta\) retains
the last displayed bill. All three terms tend to zero
after any fixed polynomial loss in \(B\).
They keep their original masks. Exponentiating proves
the second assertion. Factors disjoint from these neighbourhoods
retain their original local letters.
\end{proof}

\begin{lemma}[Joint payment from a completed energy bound]
\label{lem:packet-joint-payment}
Suppose the actual completed energies satisfy, on the support
of their retained tests,
\begin{equation}
 \beta{\cal R}_p
 \ge c_*P^2m(p)
      +H(1+B)^n(1+\log\widehat\beta_{\rm f})m(p),
 \label{eq:packet-required-energy-payment}
\end{equation}
where \(c_*>0\) is fixed. Choose \(H,n\) large enough for the
bills in the preceding lemmas. The incoming norms
range over their prescribed bounded
profilewise ball; its bound is included in this choice of \(H\).
Consider any finite family of packets under the one joint
law \eqref{eq:packet-one-joint-law}, and additional Gaussian
letters with a combined bound
\[
 |F_\gamma|\le A_\gamma
      \exp\{\tfrac12\sum_i d_{\gamma i}z_i^2\},
 \qquad d_{\gamma i}\ge0,
\]
assume \(r_0>0\), \(u_0\ge0\),
\(C_{\rm sel}^{-1}\ge(r_0+u_0)1\), and
\(\sum_\gamma d_{\gamma i}\le u_0\).
Then their joint absolute integral, including the sum of
incoming occurrences, is at most
\begin{equation}
 \exp\{-c_*P^2\textstyle\sum_p m(p)\}
 \prod_\gamma A_\gamma\
 \exp\left\{\frac1{2r_0}
                    \sum_{\gamma,i}d_{\gamma i}\right\}.
 \label{eq:packet-joint-integral-bound}
\end{equation}
One can replace \(c_*\) on the right by any smaller fixed
positive constant to leave unused energy reserve.
\end{lemma}
\begin{proof}
Keep every multiplier in the original integral.
Apply \eqref{eq:packet-required-energy-payment} and the
preceding two lemmas pointwise. For the sums of incoming
occurrences use Tonelli on this nonnegative majorant.
The resulting packet bound is uniform in the actual frozen
links, tags, and Gaussian coordinates wherever the retained
tests are nonzero; elsewhere that multiplier vanishes.
Integrating probability Haar measure and the probability tag
law therefore costs one.

Write \(d_i=\sum_\gamma d_{\gamma i}\) and
\(P_0=C_{\rm sel}^{-1}\).
Lemma~\ref{lem:finite-range-joint-tilt} now supplies the
exact determinant ratio and the factor
\(\exp\{\sum_i d_i/(2r_0)\}\) for the remaining Gaussian
integral. Polynomial insertions, when needed for
derivatives, are covered by
Corollary~\ref{cor:finite-range-tilted-polynomial}.
Both statements use the sum of all quadratic charges once.
It does not apply a separate Gaussian tilt to each letter.
\end{proof}

The lemma is an integration estimate with an explicit energy
input. Establishing \eqref{eq:packet-required-energy-payment}
for \eqref{eq:packet-completed-energy} is essential; it is
addressed next. In particular probability Haar mass one does
not by itself cancel \(s_g^{-|J_p|}\).

\subsection{Completing the calibrated row system}

The stationary calibration determines the row lift on exact
curvatures. The energy comparison also uses its complement on
all coarse two-form rows. We now make the previously free
off-exact choices in Section~\ref{sec:all-scale-linear}
so that this complement is positive.

In this subsection \(d_1,d_2,d_3\) are the lattice exterior
derivatives from one-forms to two-forms, from two-forms to
three-forms, and from three-forms to four-forms, respectively.
Their symbols are exterior multiplication by
\((e^{ip_i}-1)_{i=1}^4\). Transposes in position space become
Hermitian adjoints on the physical Fourier torus.
Let \(B_v=H_v^{-1}\) be the unmodified two-form covariance.
For an adjacent pair write \(T\) for the two-form column
constructed earlier and put
\[
 Z=H_{\rm c}
  (B_{\rm c}-T^TB_{\rm f}T-a^{-1}d_{1,\rm c}d_{1,\rm c}^T)
       H_{\rm c}.
\]
The column identities give \(Zd_{1,\rm c}=0\);
also \(Z^T=Z\).

\begin{lemma}[A positive complementary row]
\label{lem:packet-complementary-row}
The off-exact parameters \(l_j\), and subsequently the
admissible lower bounds on \(L\), can be chosen so that the
flat lift \(\Phi_0=(\Phi_F,\Phi_I)^T\) has a complementary
exponentially local row operator \({\Phi_{\perp}}\) satisfying
\begin{equation}
 \Phi_F^T\Phi_F+\Phi_I^T\Phi_I+{\Phi_{\perp}}^T{\Phi_{\perp}}=1,\qquad
       {\Phi_{\perp}} S_{\rm c}d_{1,\rm c}=0.
 \label{eq:packet-full-row-isometry}
\end{equation}
Its coefficient bounds are chosen before \(L\), uniformly
over the legal finite histories at the fixed final profiles.
The fine and noise prices of \(\Phi_0\) remain \(C/L^2\)
and \(C\).
\end{lemma}
\begin{proof}
We give the division argument, including the zero Fourier
mode. For an exponentially summable Laurent series, division
of \(z^n-1\) by \(z-1\) is the finite signed interval sum,
for either sign of \(n\). Its coefficient sum with weight
\(e^{\sigma|x|/2}\) is bounded by
\(C_\sigma e^{\sigma|x|}\).
Thus taking
\((f-f|_{z_i=1})/(z_i-1)\) preserves exponential summability
on a smaller fixed strip.

Here is the resulting exterior division. If a row \(A\) on
positive-degree \(k\)-forms satisfies \(A d_{k-1}=0\), split
its dual exterior form along the first coordinate as
\(e^1\wedge a+b\), with neither \(a\) nor \(b\) containing
that coordinate. Put \(b_0=b|_{z_1=1}\).
The cycle equation gives
\(\iota_{q'}b=-q_1a\) and
\(\iota_{q'}b_0=0\), where
\(q_i=z_i-1\) and \(\iota_{q'}\) denotes contraction by
the remaining three components.
Consequently
\[
 \iota_q\left(e^1\wedge\frac{b-b_0}{q_1}\right)
                         =e^1\wedge a+b-b_0 .
 \]
Repeat with the residual \(b_0\) in coordinates two,
three, and four. A positive-degree residual with no coordinate
left is zero. This constructs a right lift
\({\cal L}(A)d_k=A\), with the stated coefficient bound.
Transposition gives a left lift. All divisions are interval
sums, so the identities hold at \(q=0\) as well.

For completeness, two-sided division requires a correction.
Take a left lift \(W\) of \(Z\), so \(d_2^TW=Z\).
The row \(Wd_1\) is left-annihilated by \(d_2^T\);
take a left lift \(A\) with \(d_3^TA=Wd_1\).
Then \(A d_0=0\): its left image under \(d_3^T\) is zero,
and the same exterior division at top degree has zero
target in degree five. Choose \(X\) with \(X d_1=A\).
It follows that
\[
 (W-d_3^TX)d_1=0,\qquad
 d_2^T(W-d_3^TX)=Z .
 \]
A final right lift and symmetrisation give an exponentially
local symmetric \(F\) such that \(d_2^TFd_2=Z\).
The all-history coefficient bounds for \(Z\) established by
the alias estimates give a bound \(|F|_\sigma\le C\)
before \(L\).

The same right division applied to
\(d_{2,\rm f}B_{\rm f}T H_{\rm c}\) gives an operator \(A\)
with
\[
 d_{2,\rm f}B_{\rm f}T=A\,d_{2,\rm c}B_{\rm c}.
 \]
Its coefficient price per fine residue is \(CL^{-3}\):
one fine difference is applied to the \(L^{-2}\) column,
and exterior division uses coarse differences only.
There are \(L^4\) fine residues. The geometric-mean
row/column Schur bound therefore gives
\(\|A\|_{\ell^2\to\ell^2}\le C/L\).
The same estimate holds on a smaller fixed coefficient strip.

Now use the allowed covariances
\[
 \widetilde B_{\rm f}
  =B_{\rm f}+l_{\rm f}(d_{2,\rm f}B_{\rm f})^T
                              (d_{2,\rm f}B_{\rm f}),\qquad
 \widetilde B_{\rm c}
  =B_{\rm c}+l_{\rm c}(d_{2,\rm c}B_{\rm c})^T
                              (d_{2,\rm c}B_{\rm c}).
 \]
Put \(D_0=(1+l_{\rm c}d_2B_{\rm c}d_2^T)^{-1}\).
The inverse formula, or multiplication by
\(\widetilde B_{\rm c}\), gives
\(d_2B_{\rm c}\widetilde H_{\rm c}=D_0d_2\).
Hence
\begin{align}
 &\widetilde H_{\rm c}
  (\widetilde B_{\rm c}-T^T\widetilde B_{\rm f}T
                   -a^{-1}d_1d_1^T)\widetilde H_{\rm c}
       =d_2^T R d_2, \label{eq:packet-rebased-defect}\\
 &R=D_0^T(F+l_{\rm c}1-l_{\rm f}A^TA)D_0 . \notag
\end{align}
At the initial Wilson input take \(l_{\rm f}=0\);
its unmodified quadratic boundary condition is unchanged.
For subsequent inputs use their already chosen off-exact
parameter. Choose \(l_{\rm c}\ge C+2\) and then
\(L^2\ge C^2l_{\rm f}\).
The middle operator is at least \(1\).
It and \(D_0\) have uniform exponential coefficient bounds
and spectral bounds at these fixed profiles. The same
resolvent argument used for \(S_v\) gives such bounds for
\(R^{1/2}\).

The matrix
\(P_S=S_{\rm c}\widetilde B_{\rm c}S_{\rm c}^T\)
is an orthogonal projection. Define
\[
 {\Phi_{\perp}}=\begin{pmatrix}
      1-P_S\\
      R^{1/2}d_{2,\rm c}\widetilde B_{\rm c}S_{\rm c}^T
    \end{pmatrix}.
 \]
Substitution of \eqref{eq:packet-rebased-defect} and
\(S_{\rm c}^TS_{\rm c}=\widetilde H_{\rm c}\)
gives \({\Phi_{\perp}}^T{\Phi_{\perp}}=1-\Phi_0^T\Phi_0\).
Also \({\Phi_{\perp}} S_{\rm c}d_1=0\), since \(P_SS_{\rm c}=S_{\rm c}\)
and \(d_2d_1=0\). This proves
\eqref{eq:packet-full-row-isometry}.
Each profile has only the prescribed adjacent profile types.
Choose its \(l_j\) above their finitely many bounds on \(C\),
then choose their \(L\)-thresholds. Repetitions within one
profile satisfy the same \(L^2\)-condition.
Thus this choice closes along the same finite-history
iteration; it introduces no history-dependent coefficient.
\end{proof}

\subsection{An exact row decomposition of the physical energy}

Write \(F_{\rm f}(U)\) for the two-layer fine curvature rows:
the direct layer is \(\sqrt{\lambda_{\rm f}}\mathcal C(P)\),
and the second layer is the cut principal row \(\ell_{R_B}\).
Use \(F_{\rm c}(V)\) at the coarse scale. The map
\(\mathcal C:G\to\mathfrak g\) is the skew-Hermitian
projection of \(W-1\), with the fixed Lie-algebra projection.
It satisfies
\[
 |\mathcal C(W)|\le d_G(W),\qquad
 0\le d_G(W)^2-|\mathcal C(W)|^2
                           \le C d_G(W)^4
\]
near the identity; these follow by Taylor expansion in the
fixed exponential chart and orthogonal projection.
Let \(c_{\rm f}(U;p)\) be the cut fine crown, namely its
clipped direct price plus its original retained-row square,
and define \(c_{\rm c}(V;p)\) similarly.
No retained-row predicate is removed from either definition.

Define a thinned coarse vector \(v\) by retaining
\(\ell_{R_B}(V;p)\) only when its original predicate holds,
and replacing its direct layer by
\[
 \sqrt{\lambda_{\rm c}}\,
       \min(1,h_{\rm c}/d_G(P_p(V)))
                         \mathcal C(P_p(V)).
 \]
At chord zero the multiplier is one.
The retained-row coefficient bound gives \(|v_r|\le C h_{\rm c}\).
Transport and cut \(\Phi_F,\Phi_I,{\Phi_{\perp}}\) by the path convention
of the earlier covariant rows, and set
\[
 G_F=\Phi_F(V)v,\quad G_I=\Phi_I(V)v,\quad
 Z_b=\sqrt a\,\mathcal C(V_bQ_t(U)_b^{-1}).
 \]
Here all operators are cut at the same native length \(R_B\).
Their full, uncut flat counterparts satisfy
\eqref{eq:packet-full-row-isometry}.
The pre-block prices give
\(|G_F|_\infty\le C L^{-2}h\) and
\(|G_I|_\infty\le C h\).

Use the following rows, with all displayed inner products
summed over the relevant two-layer indices:
\begin{align}
 w_F(p)&=c_{\rm f}(U;p)+\tfrac12|G_F(p)|^2
                           -\langle G_F(p),F_{\rm f}(U;p)\rangle,
                                                       \notag\\
 w_I(b)&=\tfrac a2q_b^2+\tfrac12|G_I(b)|^2
                           -\langle G_I(b),Z_b\rangle, \notag\\
 w_C(p)&=\langle v_p,F_{\rm c}(V;p)\rangle
                            -\tfrac12|v_p|^2-c_{\rm c}(V;p),\notag\\
 w_L(p)&=\langle v_p,
       [\Phi_F^TF_{\rm f}+\Phi_I^TZ-F_{\rm c}]_p\rangle,\notag\\
 w_{\perp}(x)&=\tfrac12|({\Phi_{\perp}}v)_x|^2, \notag\\
 w_D(p)&=\tfrac12\langle v_p,
       [1-\Phi_F^T\Phi_F-\Phi_I^T\Phi_I-{\Phi_{\perp}}^T{\Phi_{\perp}}]v|_p\rangle .
 \label{eq:packet-six-rows}
\end{align}
Group \(C,L,D\) at their common coarse anchor; only \(F,I\)
have affine squares subtracted in
\eqref{eq:packet-completed-energy}.
Their anchor grouping is important near a hole boundary.

\begin{lemma}[Exact row accounting]
\label{lem:packet-exact-row-accounting}
For every physical configuration and every branch,
\begin{equation}
 \sum(w_F+w_I+w_C+w_L+w_{\perp}+w_D)
  =\sum c_{\rm f}+\tfrac a2\sum q_b^2-\sum c_{\rm c}.
 \label{eq:packet-six-row-identity}
\end{equation}
The equality holds for the cut transported operators,
without asserting that they are isometries.
\end{lemma}
\begin{proof}
The two terms
\(-\langle\Phi_Fv,F_{\rm f}\rangle\) and
\(\langle v,\Phi_F^TF_{\rm f}\rangle\) cancel;
the noise pair cancels in the same way.
The terms containing \(\langle v,F_{\rm c}\rangle\)
also cancel. The remaining extra squares are
\[
 \tfrac12\|\Phi_Fv\|^2+\tfrac12\|\Phi_Iv\|^2
 -\tfrac12\|v\|^2+\tfrac12\|{\Phi_{\perp}}v\|^2
 +\tfrac12\langle v,
       (1-\Phi_F^T\Phi_F-\Phi_I^T\Phi_I-{\Phi_{\perp}}^T{\Phi_{\perp}})v\rangle=0 .
 \]
This proves the identity, including the original crown masks.
Taking the prescribed common orientation average preserves it.
\end{proof}

\subsection{Local physical energy per seed}

For clarity, denote the direct and linear clipping coefficients
by \(\lambda_{\rm f},\lambda_{{\rm l},\rm f}\), and
\(\lambda_{\rm c},\lambda_{{\rm l},\rm c}\).
They are the fixed positive profile coefficients of
Section~\ref{sec:coupling-reset}, with the prescribed
small linear coefficient relative to the direct one.
In particular no assertion below uses a lower bound on a
large plaquette with coefficient
\(\lambda_{\rm f}/128\) if its actual linear coefficient
is smaller.

Assign fine and noise rows within distance \(2B\) of a
packet, and coarse and gain rows by the same incident-anchor
rule. Call their total \({\cal A}_p\), with the coarse
\(C,L,D\) rows kept together. Its number of anchors is
at most \(C_L(1+B)^4m(p)\).

\begin{lemma}[The physical seed charge]
\label{lem:packet-physical-seed-charge}
There is \(L_0\), chosen from pre-block coefficient bounds,
such that for each allowed \(L\ge L_0\) there is \(W_0(L)\)
with the following property. For \(W\ge W_0(L)\) one can
choose \(\kappa>0\) and a nominal-coupling threshold, all
independent of period and cutoff, such that on
\eqref{eq:packet-physical-plot},
\begin{equation}
                       {\cal A}_p\ge\kappa h^2m(p).
 \label{eq:packet-seed-energy}
\end{equation}
Only the physical tests within the owned neighbourhood and
the actual positive branch likelihood are used.
\end{lemma}
\begin{proof}
We give the estimates and their finite budget. Constants
denoted by \(C\) in the coefficients of \(L^{-2}\) below
are chosen before \(L\); other polynomial losses may depend
on the fixed \(L,W\).

First consider a principal row whose retained-row predicate
fails. A witness plaquette has chord
\(s_q>h e^{d(o,q)/r_0}\). Summing the resulting exponential
cones, and then the row coefficients, gives
\begin{equation}
 h\,\#\{o:\mathcal H_o\text{ fails}\}
 +\sum_{o:\mathcal H_o\text{ fails}}|\ell_{R_B}(U;o)|
 \le C\sum_{q:s_q>h}s_q .
 \label{eq:packet-dropped-row-bound}
\end{equation}
To verify this, bound the indicator of a failed test by
\(\sum_{q:s_q>h}(s_q/h)e^{-d(o,q)/r_0}\).
For the row sum split its input plaquettes into \(s_q\le h\)
and \(s_q>h\). The first part is at most \(Ch\) per
failed row. For the second use its column-sum bound
\(\sum_o|r_{oq}|\le C\).
The lattice sum of the exponential cone is finite in
four dimensions. Witnesses have the effective finite reach
already proved for the original predicates and hence belong
to the owned halo whenever this estimate is used locally.

Complete squares in every retained fine row.
For the direct layer use \(|G_F|\le CL^{-2}h\).
A selected fine plaquette with \(h/4\le s\le h\)
then has a price at least \(d_{\rm f}h^2\), where
one may take \(d_{\rm f}=\lambda_{\rm f}/128\)
after increasing \(L_0\).
For \(s>h\), its clipped price is
\(\lambda_{{\rm l},\rm f}hs\).
Its direct cross term and the dropped-row loss in
\eqref{eq:packet-dropped-row-bound} are absorbed by choosing
\(CL^{-2}\le\lambda_{{\rm l},\rm f}/4\).
Unselected direct and retained-row squares are nonnegative,
apart from a quartic \(C_Bh^4\) error.
Thus, after harmless enlargement of that error,
\begin{equation}
 \sum w_F\ge
 \tfrac12\lambda_{{\rm l},\rm f}h
                 \sum_{\text{fine seeds}:s>h}s
 +d_{\rm f}h^2\#\{\text{fine seeds}:s\le h\}
                         -C_Bh^4m(p).
 \label{eq:packet-fine-row-payment}
\end{equation}
The quartic error comes from replacing chord squares by
curvature squares; no term of order \(h^2\) is discarded.

For noise use \(|G_I|\le Ch\) and
\(|Z_b|\le\sqrt a\,q_b\).
The difference
\(\tfrac a2q_b^2-\tfrac12|Z_b|^2\) is nonnegative,
and completing the remaining square gives nonnegativity
on unselected bonds. On a selected bond \(q_b\ge Wh\).
For sufficiently large \(W\),
\begin{equation}
             \sum w_I\ge \tfrac a8
                         \sum_{\text{noise seeds}}q_b^2 .
 \label{eq:packet-noise-row-payment}
\end{equation}

The principal coarse-row part of \(w_C\) is exactly zero:
when retained it is
\(|\ell|^2-\tfrac12|\ell|^2-\tfrac12|\ell|^2\),
and when not retained all three terms are zero.
Fix a sufficiently small chart chord \(\delta_0>0\).
The direct part therefore satisfies
\begin{equation}
 \begin{array}{ll}
 w_C\ge-Ch_{\rm c}^4,&s\le h_{\rm c},\\
 w_C\ge(\lambda_{\rm c}/4)h_{\rm c}s,
                                      &h_{\rm c}<s\le\delta_0,\\
 w_C\ge-\lambda_{{\rm l},\rm c}h_{\rm c}s,
                                      &s>\delta_0 .
 \end{array}
 \label{eq:packet-coarse-row-payment}
\end{equation}
For the middle interval substitute
\(v=\sqrt{\lambda_{\rm c}}(h_{\rm c}/s)
\mathcal C(P)\).
The first two terms are
\(\lambda_{\rm c}(h_{\rm c}/s-h_{\rm c}^2/(2s^2))
|\mathcal C(P)|^2\).
Use the quartic error, \(s\le\delta_0\), and
\(\lambda_{{\rm l},\rm c}<\lambda_{\rm c}/16\).
For the last interval the factor in parentheses is
nonnegative, giving the stated lower bound directly.

We next locate the possible losses in \(w_L+w_D\).
Put
\[
 T=\widehat\beta_{\rm f}^{-49/100}.
 \]
It tends to zero, while \(h/T\) tends to zero faster than
any inverse power of \(B\). Call a fine or noise seed
exceptional if its chord exceeds \(T\), and write
\(m_{\rm exc}\) for their number.
The hit coarse rows are those within \(4B/5\) of an
exceptional record. Their number is at most
\(C(1+B)^4m_{\rm exc}\).

We spell out the geometric estimate needed here. If
\(\operatorname{port}_x(U)\) is the sum of fine plaquette
chords in the predictor's macro box about \(x\), then
\begin{equation}
 d_G(P(Q_t(U);x))
       \le C L^{-2}\operatorname{port}_x(U),\qquad
 d_G(P(V);x)\le
       C L^{-2}\operatorname{port}_x(U)
                  +\sum_{b\in\partial x}q_b .
 \label{eq:packet-macro-circulation}
\end{equation}
Each fine plaquette belongs to at most \(C\) of these
macro boxes, with \(C\) chosen before \(L\).
Take the macro box to be \(\mathcal M_x\) of
Proposition~\ref{prop:all-branch-macro-circulation}, so that
\(\operatorname{port}_x=\operatorname{face}_{\mathcal M_x}(U)\).
That proposition proves the first bound for each individual
positive tag, using the eligibility condition
\eqref{eq:axial-eligibility}, the layer induction, and the
common frame comparison \eqref{eq:all-branch-frame-comparison}.
The three explicit prices are \(CL^2\operatorname{port}_x\)
for each frame before its weight, \(L^{-4}\) for that
weight, and \(CL^{-2}\operatorname{port}_x\) for each
weighted straight connector. Its predictor calculation
includes both branches, with constant at most
\(\max(2,2/c_{\rm pr})\). Thus every coefficient here is
chosen before \(L\). Corollary~\ref{cor:macro-circulation-with-noise}
gives the second bound, with the actual noise chords and
no replacement of the branch probabilities. The same
proposition gives the stated overlap count.

% !TeX root = 4DYM_notes.tex
\emph{The norms and the two different summations.}
We now expand the two defect rows in
\eqref{eq:packet-six-rows}. The organisation agrees with
\cite[Unit III.10, Lemmas 4.4--4.8, PDF pp.~259--260]{ArmstrongNotes}.
Only absolute coefficient estimates are used on the hit set.
In particular the argument applies at rough configurations
and on every positive-likelihood branch.

At a plaquette, combine its direct and principal coordinates
in the Hilbert space \(\mathfrak g\oplus\mathfrak g\), with
the sum of the two squared Lie algebra norms. Other fixed
row types have the analogous direct-sum norm. The original
symmetry weights are included in the rows. For a block
matrix \(K\) between these spaces define
\[
 \|K\|_{\rm row}=\sup_i\sum_j\|K_{ij}\|_{\rm op},\qquad
 \|K\|_{\rm col}=\sup_j\sum_i\|K_{ij}\|_{\rm op}.
\]
Here \(\|\cdot\|_{\rm op}\) is the operator norm for the
displayed Lie algebra Hilbert spaces; all spatial sums
use counting measure. These definitions imply
\[
 \|Ku\|_\infty\le\|K\|_{\rm row}\|u\|_\infty,\qquad
 \|K^T\|_{\rm row}=\|K\|_{\rm col},\qquad
 \|K_1K_2\|_{\rm row}
            \le\|K_1\|_{\rm row}\|K_2\|_{\rm row}.
\]
The last inequality follows by summing
\(\|(K_1K_2)_{ij}\|\le\sum_k\|(K_1)_{ik}\|\|(K_2)_{kj}\|\).

Choose one \(A\ge1\), depending on the group, the fixed
adjacent profiles, their path budgets and clipping and noise
coefficients, but chosen before \(L,W\). Enlarge it a finite
number of times below using only those data. The constructed
row lifts and the complementary rows obey
\[
\begin{array}{c|cc}
 K&\|K\|_{\rm row}&\|K\|_{\rm col}\\ \hline
 \Phi_F&A L^{-2}&A L^2\\
 \Phi_I&A&A\\
 {\Phi_{\perp}}&A&A .
\end{array}
\]
For the fine lift, the first entry is the pre-block
coefficient bound in Theorem~\ref{thm:constructed-calibration}.
To obtain the second, fix an input coarse plaquette and
write an output fine plaquette as its coarse parent and
one of the \(L^4\) residues. At each residue, sum the
coarse displacement using the same exponential cover
budget. It costs at most \(C L^{-2}\). Summing the residues
costs \(L^4\), giving \(C L^2\). The fixed direction and
layer counts are included in \(A\). For \(\Phi_I,{\Phi_{\perp}}\)
both indices have a fixed number of types per coarse site,
so there is no fine-residue factor.

Cutting terms decreases these absolute bounds. Every
adjoint transport is an isometry. Periodisation also
preserves the upper bound: for each cover displacement
there is just one compatible periodic endpoint in a row
or column, even when different displacements have the
same endpoint. Thus this table holds for the actual
cut transported operators, with the same \(A\), for
all cutoffs, periods and legal finite histories.

The same native path budgets give nonnegative arrays
\(k^{\rm f}_{rq},k^{\rm c}_{rq}\) with row and column sums
at most \(A\), such that
\[
 |F_{\rm f}(U;r)|\le\sum_q k^{\rm f}_{rq}s_q,\qquad
 |F_{\rm c}(V;r)|\le\sum_q k^{\rm c}_{rq}s_q^{\rm c},
 \qquad s_q^{\rm c}=d_G(P_q(V)).
\]
For the direct layer the array is its fixed coefficient
times the diagonal. For the principal layer, replace every
native path coefficient by its absolute value, use
\(|\mathcal C(P_q)|\le s_q\), and sum paths with the same
endpoint. Translation invariance bounds columns by the
same native budget that bounds rows. This is a same-level
sum and contributes no \(L^4\). All original retained-row
tests in the crowns and in \(v\) remain in place.
Their value bound and the direct thinning give, after
enlarging \(A\),
\[
                         |v_p|\le Ah .
\]
The comparison \(h_{\rm c}\le2h\) has been included here.

\emph{The common excess arrays.}
Let \(\mathcal Q_{\rm hit}\) be the owned coarse
plaquettes within \(4B/5\) of an exceptional record;
orientation types are counted in this set. Its cardinality
satisfies
\[
        |\mathcal Q_{\rm hit}|
                       \le A(1+B)^4m_{\rm exc}.
\]
Indeed a ball of this radius contains at most
\(C(1+B)^4\) coarse plaquettes, and the union of such
balls has at most the sum of their cardinalities.
Every full read of a row tested on this set remains
inside the owned halo: its extra reach is less than
\(B/50\), with the fixed predictor reach included at
entry. Enlarging these balls by that reach still has
cardinality \(C(1+B)^4m_{\rm exc}\), rather than a
product of two such ball counts.
At fixed \(L,W\), entry is also taken so that
\(2Wh\le T\). Hence an unselected noise bond cannot
be an input above the exceptional threshold.

Define nonnegative excess arrays by
\[
 e_q=s_q\mathbf1_{\{q\text{ a fine seed},\ s_q>h\}},\qquad
 n_b=q_b\mathbf1_{\{b\text{ a noise seed}\}},\qquad
 \mathcal B_p=L^{-2}\sum_qe_q+\sum_bn_b.
\]
The sums range over the records of this packet. On every
read just specified, the physical sector tests imply
\[
                        s_q\le h+e_q,\qquad
                        q_b\le2Wh+n_b.
\]
For the first inequality a chord greater than \(h\)
must be a selected fine seed. For the second, an
unselected noise bond has chord at most \(2Wh\).
Selected bonds with chord between \(Wh\) and \(2Wh\)
are still included in \(n_b\). Every selected record
on these reads belongs to the same geometric packet.

Use \eqref{eq:packet-macro-circulation} to define, on the
needed coarse plaquettes, a nonnegative macro excess
\[
 H_q=C_{\rm mac}L^{-2}
                    \sum_{u\in\mathcal M_q}e_u
                         +\sum_{b\in\partial q}n_b .
\]
The coefficient \(C_{\rm mac}\) is the pre-block constant
in that estimate. A macro box contains at most \(C_JL^4\)
fine plaquettes. Substituting the preceding two scalar
inequalities therefore gives
\[
 s_q^{\rm c}\le r_{L,W}h+H_q,
 \qquad r_{L,W}=A(1+L^2+W),\qquad
                         \sum_qH_q\le A\mathcal B_p.
\]
Here \(A\) has been chosen to dominate
\(C_{\rm mac}C_J\), the fixed noise-boundary count,
and the macro-box overlap count. To verify the last
inequality, sum the displayed definition of \(H_q\):
each fine plaquette occurs in at most a fixed number
of macro boxes, and each noise bond in a fixed number
of coarse plaquette boundaries. Both multiplicities
precede \(L\). Extending the sum beyond the tested
coarse set only increases this nonnegative majorant.
The absence of exceptional inputs on a macro box
also gives \(s_q^{\rm c}\le C_L T\). Thus every owned
plaquette with \(s_q^{\rm c}>\delta_0\) belongs to
\(\mathcal Q_{\rm hit}\), once \(C_LT<\delta_0\) and
the fixed macro reach is less than \(4B/5\).

\emph{The fine mixed term.}
The fine input estimate now reads
\[
       |F_{\rm f}(U;r)|
                     \le Ah+\sum_q k^{\rm f}_{rq}e_q.
\]
The triangle inequality, the bound on \(v\), and finite
rearrangement of nonnegative sums give
\begin{align*}
 &\left|\sum_{p\in\mathcal Q_{\rm hit}}
                  \langle v_p,(\Phi_F^TF_{\rm f})_p\rangle\right|\\
 &\quad\le Ah\sum_{p\in\mathcal Q_{\rm hit},r}
                \|\Phi_F(r,p)\|
                     \left(Ah+\sum_qk^{\rm f}_{rq}e_q\right)\\
 &\quad\le A^3L^2h^2|\mathcal Q_{\rm hit}|
       +Ah\sum_qe_q\sum_rk^{\rm f}_{rq}
                       \sum_{p\in\mathcal Q_{\rm hit}}\|\Phi_F(r,p)\|\\
 &\quad\le A^3L^2h^2|\mathcal Q_{\rm hit}|
                              +A^3L^{-2}h\sum_qe_q .
\end{align*}
The uniform \(Ah\) input uses the column sum \(AL^2\).
The excess uses the row sum \(AL^{-2}\), after \(r\)
has been fixed. Its remaining \(r\)-sum is a column
sum of the same-level array \(k^{\rm f}\), bounded by
\(A\). Thus the numerator of the large-input
coefficient contains no residue count.

\emph{The noise and coarse mixed terms.}
Since \(|Z_b|\le\sqrt a\,q_b\), the same calculation,
now using both row and column sums of \(\Phi_I\), yields
\begin{align*}
 &\left|\sum_{p\in\mathcal Q_{\rm hit}}
                       \langle v_p,(\Phi_I^TZ)_p\rangle\right|\\
 &\quad\le2A^2\sqrt a\,W h^2|\mathcal Q_{\rm hit}|
                                      +A^2\sqrt a\,h\sum_bn_b .
\end{align*}
For the coarse input, use its native majorant and
\(s_q^{\rm c}\le r_{L,W}h+H_q\). Then
\begin{align*}
 &\left|\sum_{p\in\mathcal Q_{\rm hit}}
                            \langle v_p,F_{\rm c}(V;p)\rangle\right|\\
 &\quad\le A^2r_{L,W}h^2|\mathcal Q_{\rm hit}|
       +Ah\sum_qH_q\sum_{p\in\mathcal Q_{\rm hit}}k^{\rm c}_{pq}\\
 &\quad\le A^2r_{L,W}h^2|\mathcal Q_{\rm hit}|
                                                +A^3h\mathcal B_p .
\end{align*}
This bounds the negative term in \(w_L\) as well as
the two terms with a positive sign. No sign of an
individual mixed inner product has been assumed.

\emph{The completeness-defect term.}
For arbitrary \(V\), the row norms satisfy
\[
 \|\Phi_F^T\Phi_F\|_{\rm row}
       \le(AL^2)(AL^{-2})=A^2,
 \qquad
 \|\Phi_I^T\Phi_I\|_{\rm row},\ \|{\Phi_{\perp}}^T{\Phi_{\perp}}\|_{\rm row}\le A^2.
\]
Consequently
\[
 \left|\sum_{p\in\mathcal Q_{\rm hit}}w_D(p)\right|
       \le\tfrac12 A^2(1+3A^2)h^2|\mathcal Q_{\rm hit}|.
\]
This estimate uses only the definition of \(w_D\),
not an isometry assertion for cut transported kernels.

Combining the four mixed-term estimates proves
\[
 \sum_{\mathcal Q_{\rm hit}}(w_L+w_D)
       \ge-C_{\rm mix}h\mathcal B_p
                    -Z_{\rm mix}(L,W)h^2|\mathcal Q_{\rm hit}|,
\]
where one permissible explicit choice is
\begin{align*}
 C_{\rm mix}&=2A^3+A^2\sqrt a,\\
 Z_{\rm mix}(L,W)&=A^3L^2+2A^2\sqrt a\,W
                  +A^2r_{L,W}+\tfrac12A^2(1+3A^2).
\end{align*}

\emph{Paying the negative coarse rows.}
Their remaining bill is bounded by
\begin{align*}
 \lambda_{{\rm l},\rm c}h_{\rm c}
                    \sum_{s_q^{\rm c}>\delta_0}s_q^{\rm c}
 &\le2\lambda_{{\rm l},\rm c}h
           \left(r_{L,W}h|\mathcal Q_{\rm hit}|+\sum_qH_q\right)\\
 &\le2\lambda_{{\rm l},\rm c}r_{L,W}h^2|\mathcal Q_{\rm hit}|
                       +2\lambda_{{\rm l},\rm c}A h\mathcal B_p .
\end{align*}
Thus take
\[
 C_7=2A^3+A^2\sqrt a+2\lambda_{{\rm l},\rm c}A,
 \qquad
 Z_{L,W}=Z_{\rm mix}(L,W)+2\lambda_{{\rm l},\rm c}r_{L,W}.
\]
The constant \(C_7\) is chosen before \(L,W\).
Only \(Z_{L,W}\) contains \(L^2\) and \(W\).
Using the cardinality bound for \(\mathcal Q_{\rm hit}\)
gives the following inequality, with the two occurrences
of \(C\) equal to \(C_7\) and with
\(C_{L,W}=AZ_{L,W}\):
\begin{align}
 &\sum_{\rm hit}(w_L+w_D)
   +\frac C{L^2}h\sum_{\text{fine seeds}:s>h}s
   +Ch\sum_{\rm noise\ seeds}q_b
   +C_{L,W}(1+B)^4h^2m_{\rm exc} \notag\\
 &\hspace{35mm}\ge
       \lambda_{{\rm l},\rm c}h_{\rm c}
                    \sum_{\text{coarse}:s>\delta_0}s .
 \label{eq:packet-hit-row-budget}
\end{align}

% !TeX root = 4DYM_notes.tex
\emph{One chart for all reads of an unhit row.}
Let \(p\) be an owned coarse plaquette outside
\(\mathcal Q_{\rm hit}\). Every fine and noise chord
needed below is at most \(T\), since its location is
within \(4B/5\) of \(p\). We justify the simultaneous
fine and coarse chart used in this assertion's
linearisation; compare
\cite[Unit III.10, Lemmas 4.9--4.13,
PDF pp.~260--262]{ArmstrongNotes}.

First the block-tree comparison gives, on the smaller
neighbourhood of radius \(2B/5\),
\[
       d_G(U_e\operatorname{rep}(V)_e^{-1})\le C_LT .
\]
Here is the path calculation. Fix coordinate routes
from a block root to every point of its predictor patch.
Closing two such routes with the same endpoints gives
a loop of length at most \(C_JL\) in a box of diameter
\(C_JL\). The coordinate filling of
Lemma~\ref{lem:box-path-comparison} uses at most
\(C_JL^2\) plaquettes. Its chord is at most \(C_JL^2T\).
In these route coordinates every pair matrix in the
frame score is therefore within \(C_LT\) of the identity.
The feedback argument of
Lemma~\ref{lem:global-branch-definition} bounds every
centred feedback by this same error; the nearest-point
map is Lipschitz. After root normalisation each frame
differs from the chosen route transport by at most
\(C_LT\), uniformly over its patch points. At entry the
regular test holds. Every connector consequently
differs from a fixed root-to-root transport by at
most \(C_LT\). Its weighted mean and its nearest
group point satisfy that same bound, with another
fixed factor. The predictor test is also regular,
since its minimum score is \(O_L(T^2)\).

For an internal fine edge, close it with the two
block-tree paths. In the actual tree gauge those
paths equal one, so the edge has chord at most
\(C_LT\). For a crossing fine edge the analogous
closed route compares that edge with the predicted
coarse bond, again with error \(C_LT\). Add its
actual mismatch \(q_b\le T\). This proves the
displayed bound for the physical replica. All loops
and predictor patches used here lie in the unhit
neighbourhood, after the fixed profile reach is
smaller than the available margin.

For coarse plaquettes in a box of radius
\(\lfloor B/2\rfloor\) about \(p\), the macro estimate
gives
\[
          s_q^{\rm c}\le C L^{-2}(C_JL^4T)+4T\le C_LT.
\]
Its macro reads still lie within \(4B/5\) for large
\(B\). The axial layer lemma on this coarse box gives
a gauge whose total edge cost is at most its side
times a fixed multiple of its total plaquette cost.
There are \(O(B^4)\) plaquettes. Hence every coarse
bond in the \(2B/5\)-neighbourhood has chord at most
\(C_L(1+B)^5T\). Lift this gauge to the fine sites
by making it constant on each block. The replica is
equivariant under this lift, and the preceding
fine-to-replica bound is unchanged by conjugation.
Thus both fine and coarse bonds on all needed reads
have chord, and then logarithm norm, at most
\[
             \varepsilon_B=C_{L,W}(1+B)^{10}T .
\]
The harmless exponent ten leaves room for all fixed
enlargements. The inverse exponential chart is
Lipschitz on a fixed smaller group neighbourhood.
Since \(\varepsilon_B\to0\), this chart is available
at one threshold independent of period and cutoff.
The original, injective cover boxes are used throughout.
This gauge only estimates a scalar row at its actual
configuration; it does not change the sector integral.

\emph{The complete flat cancellations.}
Write the two logarithmic link arrays as
\(u_{\rm f},u_{\rm c}\). For this calculation extend
them by zero outside the chart neighbourhood.
All cut reads at \(p\) are unchanged. Use the full
uncut flat row lifts \(\Phi_{F,0},\Phi_{I,0},\Phi_{\perp,0}\)
and row factors \(S_f,S_c\). Let \(d_f,d_c\) denote
the fine and coarse curls. The linearised row inputs are
\[
 F_{\rm f}^{(1)}=S_fd_fu_{\rm f},\qquad
 F_{\rm c}^{(1)}=S_cd_cu_{\rm c},\qquad
 Z^{(1)}=\sqrt a\,(u_{\rm c}-\widehat Q u_{\rm f}).
\]
From
\eqref{eq:constructed-calibration-row-lift},
\(S_f^TS_f=\widetilde H_f\), and the full column
identity \(T^Td_f=d_cQ\), one obtains
\begin{align*}
 \Phi_{F,0}^TS_fd_f
  &=S_cT^T\widetilde H_f^{-1}S_f^TS_fd_f
    =S_cT^Td_f=S_cd_cQ,\\
 \sqrt a\,\Phi_{I,0}^T&=S_cd_c .
\end{align*}
Since \(\widehat Q=Q+d_{0,c}P\) and
\(d_cd_{0,c}=0\), these imply, for the arbitrary
extended arrays just chosen,
\begin{align*}
 &\Phi_{F,0}^TF_{\rm f}^{(1)}
       +\Phi_{I,0}^TZ^{(1)}-F_{\rm c}^{(1)}\\
 &\qquad=S_cd_cQ u_{\rm f}
            +S_cd_c(u_{\rm c}-\widehat Q u_{\rm f})
            -S_cd_cu_{\rm c}=0 .
\end{align*}
No stationary equation for the actual fine configuration
was used. Separately
\eqref{eq:packet-full-row-isometry} gives, for every
coarse row array \(w\),
\[
 w-(\Phi_{F,0}^T\Phi_{F,0}
          +\Phi_{I,0}^T\Phi_{I,0}+\Phi_{\perp,0}^T\Phi_{\perp,0})w=0 .
\]
This is an identity on the full row space. In particular
it applies to the actual bounded thinned array after
the local gauge change; it does not require that
array to be a curvature or a differentiable function
of the background.

\emph{The nonlinear and cut errors.}
For the estimates in this paragraph, exponentiate the
zero-extended logarithms to obtain auxiliary fine and
coarse configurations. Every link of these configurations
lies in the small chart. Evaluate the row and column
norms below on these auxiliary configurations. All cut
entries in every product contributing to the row at
\(p\) agree with their original values, since their
reads lie inside the chart neighbourhood. This permits
global matrix-norm estimates for the auxiliary objects
without a smallness assumption on the original fields
outside that neighbourhood.

For clarity, if a native term of length \(\ell\) has
coefficient \(a_\ell\), two derivatives of its transports
and curvature products cost at most
\(C(1+\ell)^2|a_\ell|\). The exponential word budget
sums these factors. A discarded term of length
greater than \(R_B\) instead gains
\(e^{-c_LR_B}\), by spending part of that budget.
For a coarse word its length is measured in fine
units, so \(c_L>0\) may depend on the fixed \(L\).
It is independent of history length, cutoff and period.
The predictor and its frame have their fixed-patch
analytic derivatives on the regular branch just
verified. Their finite constants may also depend on
\(L\). Write \(\|K\|_{\rm row,col}\) for the maximum
of the row and column norms defined above.
There are fixed \(C_L,s,c_L>0\) such
that the following row-norm estimates hold on this
chart:
\begin{align*}
 |F_{\rm f}-F_{\rm f}^{(1)}|_\infty
       +|F_{\rm c}-F_{\rm c}^{(1)}|_\infty
       +|Z-Z^{(1)}|_\infty
 &\le C_L(1+R_B)^s\varepsilon_B^2
                         +C_Le^{-c_LR_B}\varepsilon_B,\\
 \|\Phi_F-\Phi_{F,0}\|_{\rm row,col}
       +\|\Phi_I-\Phi_{I,0}\|_{\rm row,col}
       +\|{\Phi_{\perp}}-\Phi_{\perp,0}\|_{\rm row,col}
 &\le C_L\{(1+R_B)^s\varepsilon_B+e^{-c_LR_B}\}.
\end{align*}
In this paragraph the column factor \(L^2\) is part of
\(C_L\); none of
these post-ratio constants is used to choose \(L\).

These estimates follow directly by the finite product
rule and the integral Taylor remainder in the fixed
charts. For example, a derivative of an adjoint path
transport chooses one of its \(\ell\) positions; two
derivatives choose an ordered pair of positions. A
curvature vanishes at zero and has its displayed
curl as first derivative. Multiplying its remainder
by a transport changes only the quadratic error.
The same argument for the finite predictor uses its
full derivative \(\widehat Q\). First shrink the chart
so the branch tests remain strict, and then use
\((1+R_B)^s\varepsilon_B\to0\). This gives the stated
Taylor paths with constants uniform in the original
number of lattice sites. The auxiliary extension is
used only for this pointwise row comparison; locality
returns the resulting bounds to the original reads.

Put
\[
 \mathcal L=\Phi_F^TF_{\rm f}+\Phi_I^TZ-F_{\rm c},
 \qquad
 \mathcal D=1-\Phi_F^T\Phi_F-\Phi_I^T\Phi_I-{\Phi_{\perp}}^T{\Phi_{\perp}}.
\]
To see every mixed remainder, subtract the first
flat identity as follows:
\begin{align*}
 \mathcal L={}&(\Phi_F^T-\Phi_{F,0}^T)F_{\rm f}^{(1)}
       +\Phi_F^T(F_{\rm f}-F_{\rm f}^{(1)})\\
 &+(\Phi_I^T-\Phi_{I,0}^T)Z^{(1)}
       +\Phi_I^T(Z-Z^{(1)})
       -(F_{\rm c}-F_{\rm c}^{(1)}).
\end{align*}
Each linear input has norm at most \(C_L\varepsilon_B\).
The row and column bounds therefore imply
\[
 |\mathcal L_p|
       \le C_L(1+R_B)^{s'}\varepsilon_B^2
                       +C_Le^{-c_LR_B}\varepsilon_B .
\]
For each of \(K=\Phi_F,\Phi_I,{\Phi_{\perp}}\), use the exact identity
\[
 K^TK-K_0^TK_0=(K^T-K_0^T)K+K_0^T(K-K_0).
\]
It and the second flat identity give
\[
 |(\mathcal Dw)_p|
       \le C_L\{(1+R_B)^{s'}\varepsilon_B+e^{-c_LR_B}\}
                                                   |w|_\infty .
\]
The latter estimate holds for arbitrary bounded \(w\).
We apply it to the actual thinned \(v\), keeping all
its original masks fixed pointwise.

Since \(R_B\) is a fixed positive fraction of \(B\),
while \(B\) grows quadratically in the coupling
logarithm, entry can impose \(e^{-c_LR_B}\le T\).
Indeed
\[
 \log(e^{-c_LR_B}/T)
       =-c_LR_B+\tfrac{49}{100}\log\widehat\beta_{\rm f}
                                                   \longrightarrow-\infty.
\]
Substitute \(\varepsilon_B=C_{L,W}(1+B)^{10}T\)
and \(|v|_\infty\le Ah\). All polynomial factors
in \(B\) can be included in \(C_B\), giving
\[
             |\mathcal L_p|\le C_BT^2,\qquad
             |(\mathcal Dv)_p|
                          \le C_Bh(T+e^{-c_LR_B}).
\]
Thus the two scalar rows, with their original factor
\(1/2\) in \(w_D\), satisfy
\begin{align*}
 w_L(p)&\ge-|v_p|\,|\mathcal L_p|\ge-C_BhT^2,\\
 w_D(p)&\ge-\tfrac12|v_p|\,|(\mathcal Dv)_p|
                           \ge-C_Bh^2(T+e^{-c_LR_B}).
\end{align*}
There are at most \(A(1+B)^4m(p)\) owned coarse
plaquettes. Finally \(h\le T\) eventually, so
\(h^2T\le hT^2\). Summing and enlarging the fixed
polynomial \(C_B\) proves the following bound:
\begin{equation}
 \sum_{\rm unhit}(w_L+w_D)
 \ge -C_B\{hT^2+h^2e^{-c_LR_B}\}m(p)
                         =-o(h^2)m(p).
 \label{eq:packet-unhit-row-budget}
\end{equation}

% !TeX root = 4DYM_notes.tex
\emph{The complete energy absorption.}
We now specify how the hit-row bill is paid and how
much positive energy remains. Increase the existing
fixed-ratio entry so that
\[
          h/2\le h_{\rm c}\le2h,\qquad h_{\rm c}\le\delta_0,
          \qquad h\le T.
\]
The first comparison follows directly from
\(h(x)=x^{-1/2}\exp((\log\log x)^2)\): for the
fixed adjacent nominal shift \(b_0\log L\),
\(h(x+b_0\log L)/h(x)\to1\). It is an entry
condition, not an additional hypothesis on a field.

Choose the ratio large enough for the fine-row bound
already established and for
\[
                         C_7/L^2\le\lambda_{{\rm l},\rm f}/4.
\]
Next choose
\[
 W\ge\max\{1,16C_7/a\},
\]
together with the earlier noise-row requirement.
These choices use the pre-block \(C_7\).
Subtracting the two large-input bills in
\eqref{eq:packet-hit-row-budget} from
\eqref{eq:packet-fine-row-payment} and
\eqref{eq:packet-noise-row-payment} leaves
\begin{align*}
 &\tfrac14\lambda_{{\rm l},\rm f}h
                         \sum_{\mathrm{fine\ seeds}:s_q>h}s_q
       +d_{\rm f}h^2\#\{\mathrm{fine\ seeds}:s_q\le h\}
                                      -C_Bh^4m(p),\\
 &\tfrac a{16}\sum_{\mathrm{noise\ seeds}}q_b^2 .
\end{align*}
For the second line, each selected bond obeys
\(q_b\ge Wh\), and hence
\(C_7hq_b\le(C_7/W)q_b^2\le(a/16)q_b^2\).
The negative coarse rows with chord above
\(\delta_0\) are paid exactly by the right side of
\eqref{eq:packet-hit-row-budget}. The gain rows
\(w_{\perp}\) remain nonnegative.

Fix
\[
 0<\mu_0\le\min\{\lambda_{{\rm l},\rm f}/16,
                         d_{\rm f}/4,a/64,\lambda_{\rm c}/64,1\}.
\]
Let \(m_F,m_I,m_C\) be the numbers of fine, noise and
coarse-site records of this packet. Let \(K\) count
its owned coarse plaquettes with
\(h_{\rm c}<s_q^{\rm c}\le\delta_0\).
A fine seed with \(s_q>h\) has remaining energy at
least \(4\mu_0hs_q\); one with \(s_q\le h\) has at
least \(4\mu_0h^2\). A noise seed has remaining
energy at least \(4\mu_0q_b^2\). In particular the
fine and noise payments together are at least
\[
             2\mu_0h^2(m_F+m_I)+2\mu_0hT m_{\rm exc}.
\]
Indeed an exceptional fine seed has \(s_q>T\ge h\),
so \(4hs_q\ge2h^2+2hT\). An exceptional noise
seed has \(q_b>T\) and \(q_b\ge h\), so
\(4q_b^2\ge2h^2+2hT\). Nonexceptional seeds
already pay the first term with a factor two to spare.

Each intermediate coarse plaquette contributes at least
\[
 \tfrac14\lambda_{\rm c}h_{\rm c}s_q^{\rm c}
       \ge\tfrac1{16}\lambda_{\rm c}h^2
       \ge4\mu_0h^2.
\]
Every coarse-site flag meets a plaquette above
\(h_{\rm c}\); a plaquette meets at most four sites.
For a plaquette above \(\delta_0\), the macro estimate
forces a fine or noise input above \(T\) once
\(C_LT<\delta_0\). Assign it to the first such input.
The macro-box overlap and the fixed boundary incidence
show that an input receives at most \(A\) assignments,
after the same pre-block enlargement of \(A\).
Consequently
\[
                              m_C\le4K+4A m_{\rm exc}.
\]
This also identifies explicitly the separate flag
cost that an exceptional input must pay.

The small-input bill of the hit rows is
\(C_{L,W}(1+B)^4h^2m_{\rm exc}\).
At fixed \(L,W\), impose the single entry inequality
\[
 \{4A\mu_0+C_{L,W}(1+B)^4\}h^2
                                           \le\mu_0hT.
\]
It is possible because
\[
 \frac{(1+B)^4h}{T}
   =(1+B)^4\widehat\beta_{\rm f}^{-1/100}
                    \exp((\log\log\widehat\beta_{\rm f})^2)
                                                   \longrightarrow0.
\]
Thus the exceptional energy pays both that bill and
the coarse flags charged to exceptional inputs.
Adding the displayed payments gives
\begin{align*}
 &2\mu_0h^2(m_F+m_I)+4\mu_0h^2K
        +2\mu_0hT m_{\rm exc}
        -C_{L,W}(1+B)^4h^2m_{\rm exc}\\
 &\quad\ge2\mu_0h^2(m_F+m_I)
               +4\mu_0h^2(K+A m_{\rm exc})
               +\mu_0hT m_{\rm exc}\\
 &\quad\ge\mu_0h^2(m_F+m_I+m_C)
                          =\mu_0h^2m(p).
\end{align*}

Finally the fine and small-coarse quartic errors total
at most \(C_B(h^4+h_{\rm c}^4)m(p)\).
Dividing by \(h^2m(p)\) gives \(C_Bh^2\to0\).
The unhit loss in \eqref{eq:packet-unhit-row-budget},
divided by that same quantity, is at most
\[
 C_B\{T^2/h+e^{-c_LR_B}\}\longrightarrow0,
 \qquad
 \frac{T^2}{h}
     =\widehat\beta_{\rm f}^{-12/25}P^{-1}.
\]
Here \(C_B\) is always a fixed polynomial in \(1+B\);
also \(c_LR_B\) grows as a fixed positive multiple
of \((\log\widehat\beta_{\rm c})^2\).
Choose entry so that each of these two total errors
uses at most \(\mu_0h^2m(p)/4\). We obtain
\eqref{eq:packet-seed-energy} with
\(\kappa=\mu_0/2\).

All estimates are pointwise on the original physical
tests. No incoming activity, branch likelihood, Haar
factor or Gaussian normaliser was changed. The kernels
were bounded by cover coefficient sums before any
volume sum. The norms and finite halo counts are
uniform over compatible periods and ultraviolet
cutoffs, including all legal histories at the fixed
profiles. The choices are: fixed profile constants,
then \(L\), then \(W\), then entry. Polynomial losses
in \(B\) or constants depending on fixed \(L,W\)
occur only in the last choice. The original finite
orientation average preserves all displayed bounds.
\end{proof}

\begin{lemma}[The annulus uses actual physical coordinates]
\label{lem:packet-annulus-payment}
On the retained sector tests, each operated row assigned to
a packet has all its physical coordinate reads outside the
\(2B\)-neighbourhood of every seed. On these reads the
protected frozen logarithm equals the true relative logarithm,
and the actual active and frozen coordinates have norm at
most \(C_{L,W}h\).
There are \(c,C>0\) and a polynomial \(C_B\) such that
\begin{equation}
 {\cal R}_p\ge {\cal A}_p-e_\beta m(p),\qquad
 e_\beta=C_B\{h^3+h^2e^{-cR_B}\},\qquad
       (1+B)^n\beta e_\beta\longrightarrow0
 \label{eq:packet-annulus-error}
\end{equation}
for each fixed integer \(n\). All statements use the
same actual \(V,Y,t\) as in the integral.
\end{lemma}
\begin{proof}
An operated anchor avoids every seed by \(8B\), whereas
its paths, predictor reads, and retained-row queries have
radius less than \(B/100\), apart from the already
declared harmless fixed enlargement. Such reads lie outside
the \(2B\)-neighbourhood. The default tests there give fine
plaquette chord at most \(h/2\), noise chord at most \(2Wh\),
and healthy coarse plaquettes. All branch masses used there
are positive.

We recall why this controls the relative link coordinates,
rather than an arbitrarily chosen logarithm.
In one fixed block tree, a link is the ordered product of
the plaquettes in its tree-strip and of the appropriate
block-boundary links. In a coordinate gauge on the fixed
predictor patch, the strip estimate bounds the first
factors by \(C_Lh\). Positivity of the frame branch puts
the chosen frame in its regular chart there; its local
displacement is bounded by the same strip error.
The predicted boundary links differ from the prescribed
coarse links by at most \(2Wh\).
Multiplying these factors gives
\[
 d_G(U_e\operatorname{rep}(V)_e^{-1})\le C_{L,W}h .
 \]
The tree links satisfy this with coordinate zero.
For sufficiently small \(h\) all other products lie in
the fixed injective exponential chart. Its inverse is
Lipschitz on a smaller compact chart, so their actual
logarithms have norm at most \(C_{L,W}h\).
The same strip argument is used on every orientation,
without changing the original physical link.
Since \(B>C_{L,W}\) eventually, the protected radius
\(u=Bh\) includes these logarithms. Hence its smooth
extension has not changed any frozen link on these reads.

On a healthy stencil the fine crown is its unmasked
quadratic row square plus the nonnegative chord/curvature
correction. Expand its row and the actual noise row at the
replica. For \(|a_e|\le C_{L,W}h\), Taylor's formula and
the weighted path sums give row remainder \(C_Bh^2\).
The linear derivative at the replica differs from the
transported flat derivative defining \(D\) by \(C_Bh\).
Acting on these same coordinates gives another \(C_Bh^2\).
The actual row and its affine comparison have size \(C_Bh\).
Their squared difference is consequently bounded by
\(C_Bh^3\). This proves, on each operated fine or noise row,
\[
 w_r-\tfrac12|D_ra+\alpha_r|^2\ge-C_Bh^3 .
 \]
The affine term here includes the true frozen logarithms,
as just checked. No comparison with an unrelated frozen
configuration is made.

For an unoperated fine or noise row outside the core,
the corresponding completed physical square is
nonnegative. For a coarse group \(C,L,D\), the first
line of \eqref{eq:packet-coarse-row-payment} and the
same zero-jet calculation as in
\eqref{eq:packet-unhit-row-budget} give a loss
\(C_B(h^3+h^2e^{-cR_B})\).
The gain rows are nonnegative.
Core rows have no affine subtraction, since their anchors
lie within \(2B\) of a seed.
There are at most \(C_Bm(p)\) assigned anchors.
Finally the completed-square terms
\(\tfrac12|D_r\mu+\alpha_r|^2\) in
\eqref{eq:packet-completed-energy} are nonnegative;
they are retained in the identity and may now be discarded
for this lower bound. Summing proves
\eqref{eq:packet-annulus-error}.
For its last assertion use
\(\beta h^3\le2P^2h\to0\) faster than any inverse power
of \(B\), and
\(\beta h^2e^{-cR_B}\le2P^2e^{-cR_B}\), which has the
same property.
\end{proof}

\begin{theorem}[Actual joint packet payment]
\label{thm:actual-joint-packet-payment}
Choose the profiles, off-exact row completions, \(L\),
and \(W\) in the order just specified.
There is \(c_0>0\), independent of cutoff, period, and RG
history length, such that the following holds after
increasing the profile's coupling entry threshold.
For each prescribed \(H,n\), the actual packet energies
satisfy
\begin{equation}
 \beta{\cal R}_p\ge
 c_0P^2m(p)+
 H(1+B)^n(1+\log\widehat\beta_{\rm f})m(p).
 \label{eq:packet-paid-completed-energy}
\end{equation}
Consequently, together with any additional letters
satisfying the combined Gaussian majorant in
Lemma~\ref{lem:packet-joint-payment}, the joint
Gaussian--Haar integral satisfies
\eqref{eq:packet-joint-integral-bound}, with all original
packet masks, incoming occurrences, and the exact
frozen-link normalisation included.
The packet energy bound itself has no additional
energy hypothesis.
The coefficient of \(P^2m(p)\) can be decreased by a
fixed factor to pay any fixed number of subsequent
derivatives.
\end{theorem}
\begin{proof}
Lemmas~\ref{lem:packet-physical-seed-charge}
and~\ref{lem:packet-annulus-payment} imply, eventually,
\[
 \beta{\cal R}_p
 \ge \tfrac12\kappa\beta h^2m(p)
 \ge \tfrac14\kappa P^2m(p).
 \]
For every fixed \(n\),
\[
 \frac{(1+B)^n(1+\log\widehat\beta_{\rm f})}{P^2}
 \longrightarrow0 .
 \]
Indeed \(B\) is quadratic in \(\log\widehat\beta\) whereas
\(\log P=(\log\log\widehat\beta_{\rm f})^2\).
Set, for example, \(c_0=\kappa/16\), and use part of
the remaining budget for the prescribed \(H,n\).
This proves \eqref{eq:packet-paid-completed-energy}.
Apply Lemma~\ref{lem:packet-joint-payment}.
All its pointwise bounds refer to the original integrand;
Haar and tag integrations therefore remain probability
integrations of those same factors.
\end{proof}

\subsection{Covering weights and the same Wilson family}

We record how the preceding estimate enters the local
norm rather than just a single-packet integral.
A \(200B\)-neighbourhood of one seed is covered by a fixed
number of radius-\(B\) balls. A connected seed graph can
be covered with at most \(C m\) such balls after adding
the full read halos. At a root site its number of labelled
seed graphs of size \(m\) is at most \((C_B)^m\):
encode a spanning exploration tree, its bounded local
record types, and the next location within distance \(CB\)
at each step. The number of trees is bounded exponentially
by the usual depth-first parenthesis encoding.
Thus, for every fixed desired covering exponent \(\kappa_1\),
\begin{equation}
 \sup_o\sum_{p:\,o\in X_p}
 e^{\kappa_1 n_B(X_p)}\,e^{-c_0P^2m(p)}
       \le C_Be^{-c_1P^2}
 \label{eq:packet-covering-income}
\end{equation}
for some \(c_1>0\).
The possible root-to-first-record count is polynomial in
\(B,m\) and is included in the same bound.
An incoming support of large diameter is handled by its
own covering reserve from Section~\ref{sec:physical-reassembly};
it is not filled with new frozen links.

Equation \eqref{eq:packet-joint-integral-bound} gives the
same product majorants when several packets and Gaussian
letters occur in one connected component. Their combined
quadratic charge is bounded once under \(C_{\rm sel}\).
Different components factor only when their full read sets,
including the covariance halo and the frozen and tag reads,
are disjoint. This is exactly the independence criterion
already established for the finite-range representation.
It does not assert independence of touching packets.

All coupling thresholds in this section are fixed before
the ultraviolet cutoff and the physical side length. The period
reserve of the single Wilson family already gives
\(\ell_{\rm c}\ge100(1+B)\) at each allowed station.
Local window and chart arguments use only radii below
that reserve. Larger hole neighbourhoods are unions on
the torus; union counts and covering counts remain upper
bounds even if they wrap. Native walk cuts still count
the original walk length. No estimate is obtained by
replacing that length with periodic distance.
Increasing a profile's entry threshold to satisfy the
finitely many bounds above is permitted by the existing
slow profile clock. It does not select a different Wilson
family at different volumes.

\clearpage
\section{Derivatives after physical integration}
\label{sec:integrated-regular-derivatives}

An unforced component can contain fine large fields and incoming
forced activities even when its coarse footprint is healthy.
Its output must therefore have a regular germ in \(V\), although
its integrand need not be differentiable in the physical fine links.
This section proves the derivative estimate for components containing
a physical hole packet and explains its passage to the regular
logarithmic family. The leading small-field transport of an incoming
regular family requires a sharper estimate and is kept separate.

The measure differentiation used here is elementary finite-dimensional
calculus. Its purpose in the fluctuation method is discussed in
\cite[Sect.~2, eqs.~(2.12)--(2.14), printed p.~268]{Bal87}.
The calculation below retains the actual
finite-range Gaussian marginal, the probability Haar measures,
and the original fine multipliers.

\subsection{What is differentiated}

At a coarse anchor \(o\), directions satisfy
\[
 |A_e|\le h_{\rm c}(1+d(e,o))^s .
 \]
As before, \(D_Af(V)=
\left.\frac{\dd}{\dd t}f((e^{tA_e}V_e)_e)\right|_{t=0}\),
and derivatives are ordered compositions, not a symmetric
multilinear replacement. The local norm is
\(\|f\|_{q,h_{\rm c},s;o,X}\) of
\eqref{eq:local-derivative-norm}. A regular germ is a
\(C^q\) function on a neighbourhood of the given field,
agreeing with the specified unmasked local formula there.
We require such a germ also at equality in every coarse
good-field test.

Fix the discrete word and all of its declared read sets
before differentiating. On a healthy coarse footprint,
replace its coarse good-field indicators by the value one,
use the quadratic coarse clipping formula, and use its
unthinned coarse curvature vector. These choices give the
unmasked formula on an open neighbourhood, including a
neighbourhood of fields with chord exactly \(h_{\rm c}\).
They are made only for the regular extension. Its coarse
mask is multiplied back after the derivatives are taken.
Fine plaquette tests, incoming eligibility masks,
and fine clipped actions are still evaluated on the
physical integration variable \(U\).

There is an important ordering in the word expansion.
Write \(b_e(V,a_e)\in[0,1]\) for the original broad chart
factor of an active link. For any fixed selection \(\Gamma\)
of physical Borel letters let \(T(\Gamma)\) contain all
active physical links they read. Before expanding the chart
factors, use the exact identity
\begin{equation}
 \left(\prod_{\alpha\in\Gamma}F_\alpha\right)\prod_e b_e
 =
 \left(\prod_{\alpha\in\Gamma}F_\alpha\right)
       \prod_{e\in T(\Gamma)}b_e\
 \sum_{A\subset T(\Gamma)^c}
                         \prod_{e\in A}(b_e-1).
 \label{eq:regular-retained-chart-identity}
\end{equation}
At finite volume this is simply expansion of
\(\prod_{e\notin T(\Gamma)}[1+(b_e-1)]\).
Products of \(b_e\) on overlapping read sets occur once,
on their union. A chart-failure letter \(b_e-1\) is
therefore prohibited if \(e\in T(\Gamma)\).
This is a local selection rule. It preserves the identity
and the original factors; dropping it only for an absolute
upper bound gives the same connected-family majorants,
because \(|b_e|\le1\).
The same retention is made for the physical reads of a
packet. Full read-set factorisation still holds:
disjoint components have disjoint retained chart products.
There is no new normalising factor in this rearrangement.

In particular one must not write
\[
 D_A K_\alpha(\mathcal U_\sigma(V,Y,\mu+gz)).
 \]
A bounded Borel \(K_\alpha\) need not have that derivative.
Instead change the finite set of active variables read by
this factor back to physical Haar variables before
differentiating. This moves the derivative onto the
integration density.

\subsection{The physical density of a Gaussian marginal}

Let \(T\) be the finite set of active links read by one
component, including every physical read of its incoming
activities. Its ghost coordinates can be integrated out.
Let \(D_T=d_G^{\rm Lie}|T|\),
\(\mu_T\) be the selected mean on \(T\), and
\(C_T=(C_{\rm sel})_{T,T}\) its Gaussian marginal covariance.
The original, unconditional marginal is used; it is not
the covariance obtained by conditioning on independent
copies of the exterior variables. Its spectral bounds are
\[
                 c\,1\le C_T\le C\,1
 \]
by the corresponding bounds for \(C_{\rm sel}\).
They are independent of the number or shape of links in \(T\).

For physical links in the broad chart set
\[
 a_e(V,U)=\log(U_e\operatorname{rep}(V)_e^{-1}),\qquad
 r_T=a_T-\mu_T,\qquad P_T=C_T^{-1}.
 \]
The push-forward of the selected Gaussian has the following
density with respect to probability Haar measure on \(G^T\):
\begin{equation}
 k_T(V,U;Y,t)=
 \frac{\exp\{-r_T^TP_Tr_T/(2g^2)\}}
      {(\sqrt{2\pi}g)^{D_T}(\det C_T)^{1/2}
               \prod_{e\in T}\zeta_Gj_G(a_e(V,U))}.
 \label{eq:regular-physical-Gaussian-density}
\end{equation}
If the original integrand includes \(j_G(a_e)\), it cancels
the corresponding denominator only where that original
factor is present. Broad chart cutoffs are kept in the
integrand. The density formula does not discard a chart
cutoff, or change the frozen-link normalisation.

Indeed the Gaussian marginal of \(a_T=\mu_T+gz_T\)
has Lebesgue density
\[
 (\sqrt{2\pi}g)^{-D_T}(\det C_T)^{-1/2}
                       e^{-r_T^TP_Tr_T/(2g^2)} .
 \]
Right multiplication by \(\operatorname{rep}(V)_e^{-1}\)
preserves Haar measure, and
\(\dd h_G(U_e)=\zeta_Gj_G(a_e)\dd a_e\)
on this chart. These two identities prove
\eqref{eq:regular-physical-Gaussian-density}, including
all constants.

If a component also reads Gaussian coordinates not included
in \(T\), they must not be integrated out of the kernel.
Let \(R\) be those coordinates, \(S=T\cup R\), and
\(C_S=(C_{\rm sel})_{S,S}\). In mixed physical and
Gaussian variables the exact density is
\begin{equation}
 \frac{\exp\{-\tfrac12w^TC_S^{-1}w\}}
 {(2\pi)^{D_S/2}g^{D_T}(\det C_S)^{1/2}
               \prod_{e\in T}\zeta_Gj_G(a_e)},\qquad
          w=((a_T-\mu_T)/g,z_R).
 \label{eq:regular-mixed-physical-Gaussian-density}
\end{equation}
Its measure is probability Haar measure on \(G^T\)
times Lebesgue measure on the remaining colour coordinates
\(z_R\); \(D_S\) counts all real colour coordinates in \(S\).
For the frozen links use their original Haar measures.
The same positive covariance \(C_S\) is used across both
parts. In particular \eqref{eq:regular-mixed-physical-Gaussian-density}
does not replace the outside coordinates by an independent
Gaussian. It is the required version when a chart-failure
letter outside \(T\) joins the component through a
covariance-range edge.

\begin{lemma}[Measure differentiation with a Borel multiplier]
\label{lem:regular-Borel-measure-derivative}
Let \(F(U,Y,t)\) be bounded Borel and independent of \(V\)
when its physical variables are held fixed. Let \(G(V,U,Y,t)\)
be a smooth local factor, including the retained broad chart
cutoffs and any differentiated noise partition factors.
At a fixed cutoff and finite volume, suppose their selected
matrices have the positive spectral bounds above on a
neighbourhood of the background under consideration.
Then
\[
 {\cal I}(V)=\int F\,G\, k_T\,
                \dd h_G^T(U)\dd h_G(Y)\dd\nu(t)
 \]
has the corresponding \(C^q\) germ, and its ordered
derivatives are obtained by differentiating \(Gk_T\)
under this same integral. No derivative of \(F\) is required.
\end{lemma}
\begin{proof}
At a fixed \(g>0\), all smooth derivatives of the chart
density on the support of the broad cutoffs are bounded.
The Haar density is bounded away from zero on the fixed
chart. The matrices and their inverses are smooth on a
smaller neighbourhood with a common positive lower bound.
Frozen variables range over a compact group, the tag sum
is finite, and \(F\) is bounded. Thus the differentiated
integrands have an integrable majorant at this fixed cutoff.
The fundamental theorem of calculus for difference quotients,
followed by dominated convergence, proves the first derivative.
Repeat for the ordered derivatives through \(q\).
For countably many incoming labels first truncate their
absolutely summable majorants, and then use the uniform
bounds below to remove the truncation.
For the mixed version, derivatives are polynomial in
the remaining Gaussian coordinates. The common positive
precision reserve gives an integrable Gaussian majorant
on a smaller neighbourhood; hence the same argument
applies. Its moments are bounded explicitly below.
\end{proof}

This lemma proves differentiability first. Uniformity in
the cutoffs requires the following quantitative calculation;
it does not follow merely from compactness.

\subsection{Scores of the same selected law}

The \emph{score} in direction \(A\) means
\(D_A\log k_T\). Write \(a'=D_Aa_T\),
\(\mu'=D_A\mu_T\), and \(C'=D_AC_T\).
Direct differentiation gives
\begin{align}
 D_A\log k_T
 &=-\tfrac12\operatorname{Tr}(P_TC')
   +g^{-2}\langle P_Tr_T,\mu'-a'\rangle \notag\\
 &\quad+\tfrac1{2g^2}
       \langle P_Tr_T,C'P_Tr_T\rangle
   -\sum_{e\in T}D_A\log j_G(a_e).
 \label{eq:regular-selected-score}
\end{align}
There is no derivative of \(D_T\log g\):
the coupling and this finite word are fixed while
the background varies.

\begin{lemma}[Local score bounds]
\label{lem:regular-selected-score-bounds}
Let \(X\) be the full component footprint, including its
window, frozen, and tag reads, and let \(k=n_B(X)\ge1\).
For each fixed \(q,s\) there is \(p_q<\infty\) such that,
on the broad chart support and a healthy coarse footprint,
\begin{align}
 |D_{A_1}\cdots D_{A_j}\log k_T|
       &\le C_B(1+P^2)k^{p_q},
                                      &&1\le j\le q, \label{eq:regular-log-score-bound}\\
 |D_{A_1}\cdots D_{A_j}k_T|
       &\le k_T\,[C_B(1+P^2)]^q k^{p_q},
                                      &&0\le j\le q. \label{eq:regular-density-jet-bound}
\end{align}
The constants are independent of \(|T|\), cutoff,
and period; the dependence on its size is the displayed
polynomial in \(k\).
\end{lemma}
\begin{proof}
The cardinality and diameter of a connected radius-\(B\)
cover give
\[
 |T|\le C_B k,\qquad
 \max_{x\in X}(1+d(x,o))\le C_B k
 \]
for an anchor in its declared cover.
Thus the allowed directions restricted to the footprint
have supremum at most \(C_Bh_{\rm c}k^s\).
Use Lemma~\ref{lem:selected-window-jets} for the
precision, its selected columns, \(J\) and \(C_{\rm sel}\).
Its marginal argument gives the same bounds for
\(C_T\) and \(P_T=C_T^{-1}\), independently of the
shape and cardinality of \(T\). Restricting the
allowed polynomially growing directions to this
read multiplies a fixed-order rooted jet by at most
\(C_B k^{p_q}\), by the preceding diameter bound.
Thus one fixed window formula is used in each
entry; no sum over windows or exterior coordinates
is introduced. The inverse identities and their
ordered partitions have already been proved in that
lemma and are not a different Hessian calculation.

The protected frozen logarithm has radius \(u=Bh\).
A derivative of its smooth cutoff costs \(u^{-1}\)
and a coarse variation supplies \(h_{\rm c}\le2h\).
In particular, for fixed \(j\), its scaled derivative
is bounded by \(C_Bh\,k^{p_q}\), also where it differs
from the true logarithm. The affine source has the
same bound. The formula \(\mu=-J^Tb\) therefore gives
\[
 \|\mu_T\|_\infty\le C_Bh,\qquad
 \|D_{A_1}\cdots D_{A_j}\mu_T\|_\infty
                                      \le C_Bh k^{p_q}.
 \]
The actual broad chart has width \(hB^{17}\), so
\(\|a_T\|_\infty\le C_Bh\). The smooth inverse exponential
chart and the replica path bounds give the same scaled
derivative estimate for \(a_T\).

Insert these estimates into
\eqref{eq:regular-selected-score}.
The trace costs at most \(C_Bk^{p_q}\).
Every quadratic term costs at most
\(C_B g^{-2}h^2k^{p_q}\le C_BP^2k^{p_q}\).
Derivatives of \(\log j_G\) are bounded on the fixed chart.
Higher derivatives follow by the product rule and
the inverse formula. At most \(q\) factors are inserted;
their extra powers of \(k\) can be included in \(p_q\).
This proves \eqref{eq:regular-log-score-bound}.
Apply the common partition identity
\eqref{eq:selected-density-partition-formula} to
\(k_T=\exp(\log k_T)\), including its Haar denominator.
Its finite partitions of at most \(q\) ordered
positions give
\eqref{eq:regular-density-jet-bound}.
\end{proof}

For the mixed density \eqref{eq:regular-mixed-physical-Gaussian-density},
the same calculation gives, in place of
\eqref{eq:regular-density-jet-bound}, that density times
\begin{equation}
             C_B k^{p_q}
              (1+P^2+\textstyle\sum_{i\in R}z_i^2)^q .
 \label{eq:regular-mixed-score-polynomial}
\end{equation}
Indeed only the \(T\)-part of \(w\) is bounded by \(C_BP\);
the remaining quadratic terms are bounded by
\(C_B\sum_R z_i^2\), using the row and column Schur bounds.
Repeated differentiation raises this quadratic polynomial
to at most the fixed power \(q\).
This polynomial requires no exponential price per coordinate.
Corollary~\ref{cor:finite-range-tilted-polynomial} applies
to this insertion under the same tilted law. In particular
\[
 \mathbb E_{\rm tilted}
        \left(\sum_{i\in R}z_i^2\right)^q
 \le (2q-1)!!\,r_0^{-q}|R|^q .
 \]
The determinant ratio is exactly the one already charged
in \eqref{eq:packet-joint-integral-bound}.
Thus mixed scores add only a polynomial in the covering
size. They do not spend a second precision reserve.

The germ needed here can be defined without differentiating
a coarse flag. Use the unmasked formulas on the open set
where their principal window matrices and \(C_T\) retain,
say, half their healthy spectral bounds. This open set
contains every healthy background. A smooth spectral
cutoff, constant one near the healthy spectral interval,
extends the resulting finite local formula by zero
outside that open set if an everywhere defined Borel
function is desired. Such a cutoff depends only on the
declared matrices and is invariant under their orthogonal
gauge conjugations. Its derivatives vanish on the healthy
set. It therefore changes none of the norms above.

\subsection{The nonlinear factors and the rare-component bound}

At fixed physical \(U,Y,t\), the incoming activities and
their fine inventory tests do not depend on \(V\).
The actual predictor \(Q_t(U)\) and the fine action are
also held fixed. The noise term is smooth in \(V\),
even if \(Q_t\) is only Borel in \(U\). For example,
unitarity and the chord-square formula give
\[
 |D_A\{\tfrac a2d_G(V_bQ_t(U)_b^{-1})^2\}|
                         \le C_a|A_b|.
 \]
Each further ordered derivative inserts another bounded
matrix. On a component footprint, derivatives of its
noise exponent are therefore bounded by
\(C_B\beta h_{\rm c} k^{p_q}\) at fixed order.
The same bound suffices for all the physical row terms
of the packet exponent; smooth affine terms in fact
cost only \(C_BP^2k^{p_q}\).

A differentiated smooth partition factor is treated by
its derivative bound and its support, not by dividing
by a possibly zero factor. For the noise partition the
price is \(C_q(h_{\rm c}/(Wh))^j\) times the direction
weights. Its derivative support still satisfies the
appropriate weak inequalities in
\eqref{eq:packet-physical-plot}.
The broad chart and protected logarithm similarly have
ratios \(h_{\rm c}/(hB^{17})\) and \(h_{\rm c}/(hB)\).
All are bounded by \(C_B\) at fixed order.
Fine hard masks and fine retained-row predicates are
not differentiated at all.
Coarse hard masks occur only outside the regular germ.

\begin{proposition}[Derivatives of a component with a packet]
\label{prop:regular-hole-component-derivatives}
Fix \(q,s\). Let an unforced component have \(m\ge1\) physical
seed records in its packets and full covering size
\(k=n_B(X)\). On its healthy coarse footprint its
integrated unmasked amplitude has a \(C^q\) germ.
In the joint majorant, replacing its value by its local
derivative norm costs at most
\begin{equation}
 [C_B(1+P^2+\beta h_{\rm c})]^q k^{p_q}.
 \label{eq:regular-hole-derivative-price}
\end{equation}
In particular the packet factor in that norm can be
bounded by
\begin{equation}
                      C_B k^{p_q}e^{-c_2P^2m}
 \label{eq:regular-hole-paid-derivatives}
\end{equation}
for some \(c_2>0\), after decreasing the fixed
constant in the packet payment.
The other incoming and Gaussian letter prices may be
inflated by a fixed factor depending on \(q\), but retain
the same combined precision reserve.
\end{proposition}
\begin{proof}
Retain the broad factors on all Borel reads by
\eqref{eq:regular-retained-chart-identity}.
Change those active physical reads to Haar variables using
\eqref{eq:regular-physical-Gaussian-density}, or the mixed
version \eqref{eq:regular-mixed-physical-Gaussian-density}
when other Gaussian reads remain.
Apply Lemma~\ref{lem:regular-Borel-measure-derivative}.
Ordered product differentiation acts on the smooth
kernel, affine corrections, determinant factors,
and smooth partition factors. It does not act on
bounded Borel incoming multipliers. The score bounds,
the preceding physical noise bound, and the bounded
cutoff ratios give
\eqref{eq:regular-hole-derivative-price}.
There are at most \(q\) differentiated positions.
Link placements cost a polynomial in the covering size.
If there are \(r\) selected letters, the additional
factor \(r^q\) is bounded by \(C_{q,\eta}e^{\eta r}\)
for any fixed \(\eta>0\). Include \(e^\eta\) in each
letter price before applying the connected criterion.
This handles arbitrarily many species on one footprint;
their number is not assumed bounded by its covering size.
Polynomial factors in a matrix word's native length are
paid by its previously reserved exponential length weight.

For each derivative term, use its differentiated
partition support as an upper bound. It still imposes
the weak physical tests needed for the core and annulus
payments. Thus Theorem~\ref{thm:actual-joint-packet-payment}
applies to its original physical exponent, leaving a
fixed fraction of \(P^2m\) unused.
No new Gaussian density is substituted.
If a differentiated Gaussian correction has an additional
fixed-order polynomial in its coordinates, use its moments
under the same tilted Gaussian as in
\eqref{eq:regular-mixed-score-polynomial}.
The covariance bound gives a polynomial in its number of
declared coordinates. At most \(q\) derivative positions
are involved. This is included in \(k^{p_q}\);
it introduces neither a new tilt nor a determinant
factor exponential in the number of reads.

Finally
\[
 \log(1+P^2+\beta h_{\rm c})=O(1+\log\widehat\beta_{\rm f}),
 \qquad
        1+\log\widehat\beta_{\rm f}=o(P^2).
 \]
For fixed \(q\) this absorbs the coefficient in
\eqref{eq:regular-hole-derivative-price} into any
prescribed fixed fraction of \(P^2m\), since \(m\ge1\).
The footprint polynomial is left explicit for the next
summation. This proves
\eqref{eq:regular-hole-paid-derivatives}.
\end{proof}

This proof deliberately spends part of a rare packet's
energy. Applying the same rough score estimate to a
component without a packet would lose the useful
small-field order. It would not prove the leading
linear RG estimate for the regular family.

\subsection{Passage to the regular logarithmic family}

For products of integrated regular amplitudes use the
auxiliary seminorm of
Lemma~\ref{lem:finite-jet-algebra}. It has product
constant one. Lemma~\ref{lem:finite-jet-anchor-product}
moves all factors to the final anchor in a single step,
while keeping the original healthy domain fixed.
Conversion to the prescribed norm costs \(q!\) once and
a polynomial in the total covering size. The positive
covering reserve pays that polynomial as in
\eqref{eq:finite-jet-cover-reserve}.
Thus a long cluster does not accumulate one \(2^q\)
factor per vertex.

Lemma~\ref{lem:rooted-interacting-log} then treats the
ordinary logarithm and
Lemma~\ref{lem:marked-interacting-log} treats its packet
part. These statements concern ordered Ursell tuples
with their factorials, so the distinct-label estimate
for physical word components is not needed a second
time for the logarithm.

\begin{corollary}[The packet part of the regular family]
\label{cor:regular-hole-family}
Use the rooted connected-logarithm criterion with enough
covering reserve for the product rule through order \(q\).
Suppose the other local-letter incomes, including their
required derivatives, obey that criterion.
Then the sum of unforced logarithmic clusters containing
at least one physical packet defines an anchored regular
family \(E^{\rm hole}\) satisfying
\begin{equation}
       \|E^{\rm hole}\|_{\mathcal E,\kappa_1}
                                  \le C_Be^{-c_3P^2}.
 \label{eq:regular-hole-family-bound}
\end{equation}
Its normalised residual has zero value at the identity
and zero affine Hessian on nonwrapping footprints.
The accompanying scalar, quadratic extraction, and
masked reassembly corrections are retained as in the
exact coupling reset.
\end{corollary}
\begin{proof}
Attach the norm price from
\eqref{eq:regular-hole-paid-derivatives} to a marked
packet and the original prices to the other letters.
Equation \eqref{eq:packet-covering-income}, with the
additional reserve \(\eta\), gives marked incidence
income \(f_q\le C_Be^{-c_3P^2}\).
Use the original maximum-over-orders prices at the
three extra covering units of
Lemma~\ref{lem:rooted-interacting-log}.
The factorial-weighted product estimate pays the
algebra conversion and the common-anchor change within
this reserve. Polynomial derivative-placement costs
inside a single integrated component have already
been included in its own price.
If \(b_q\) is the remaining incidence income, a sufficient
explicit smallness condition is
\[
                 D_B b_q\le\tfrac14,\qquad
                 D_B f_q\le\tfrac14 .
 \]
Mark the unforced species containing a physical packet.
Lemma~\ref{lem:marked-interacting-log} gives derivative
income at most \(2C_Bf_q\).
It uses the same nonnegative pinned majorant as the
ordinary logarithm, with unchanged cluster anchors.
Absolute convergence justifies the differentiated
family sums and the logarithm expansion.

The component species are collected with the shared tag
integration and the prescribed gauge covariance, as in
the physical reassembly. The broad factors retained on
their reads are radial, so this collection preserves
gauge invariance. The raw logarithmic functions have the sign convention
\(-\ell_i\). Subtract \(-\ell_i(1)\), then their quadratic
counterterms. Lemma~\ref{lem:local-reset-norms} bounds
this operation in the same scaled regular norm, with
a fixed enlargement of the footprint. The extraction
coefficient itself has the \(h_{\rm c}^{-2}\) price in
Lemma~\ref{lem:quadratic-extraction-price}; its counterterm
in the density carries the compensating \(h_{\rm c}^2\).
The exact masked subtraction, including its discrepancy
on bad regions, is the one already proved for the
coupling reset. Thus no discarded constant or masked
quadratic term is hidden in \eqref{eq:regular-hole-family-bound}.
\end{proof}

\subsection{Scope of this derivative estimate}

The combined packet integral and its regular germs now
have the local rare price required by physical reassembly.
For forced coarse inventories only the bounded Borel
incidence norm is needed. For unforced inventories the
preceding argument also controls the ordered derivatives
after integration, without imposing differentiability
on incoming forced activities.

The full regular RG induction also needs the central
component estimates under the selected finite-range law
and the sharp transport of a single incoming regular
activity with its affine subtraction.
These estimates must preserve the leading
small-field order; they cannot be obtained by spending a
packet energy when there is no packet. The rare-component
bound above is one part of that induction, not a claim
that the nonlinear RG iteration or continuum theorem
has already been closed.

Sections~\ref{sec:central-strength-expansion}
and~\ref{sec:finite-period-localisation} separate the
remaining tasks more economically. A scalar-strength
argument reduces higher central coefficients to one
joint analytic estimate, with its derivative and
precision inputs explicit. An exact reflected-box
identity transfers an auxiliary total-action bound to
local functions and retains the mask discrepancies.
Section~\ref{sec:central-derivative-estimates} proves the
primitive central estimates and the complex-strength
precision reserve.
Section~\ref{sec:incoming-sharp-transport} proves
the leading transport and its exact flat-value and
Hessian subtraction. Section~\ref{sec:ordinary-matching}
proves the classical and one-loop matching calculations.
Their ratio-uniform drift is proved in
Section~\ref{sec:one-loop-drift}, and their combined
nonlinear closure in
Section~\ref{sec:nonlinear-induction-step}.
\par
\clearpage
\section{A shorter route through the central expansion}
\label{sec:central-strength-expansion}

The central estimate should separate products, higher connected
coefficients, and the leading transport of an incoming regular
function. Repeating a full diagram expansion for each of them is
unnecessary. This section proves a scalar-parameter reduction
that uses the original local law and the algebraic seminorm of
Lemma~\ref{lem:finite-jet-algebra}.
Its analytic inputs are stated explicitly.
Section~\ref{sec:central-derivative-estimates} subsequently
proves them for the central letters of the physical
finite-range construction. The reduction itself is kept
separate from that model-specific verification.

The useful separation of incoming terms into a principal
average and a comparison remainder follows
\cite[Lemma~13.5, PDF p.~95]{ArmstrongDraft}.
The scalar coefficient argument below is written out for the
norms of these notes. The next section treats the finite-period
localisation in \cite[Lemma~13.6, PDF p.~98]{ArmstrongDraft}
as a separate exact operation.

\subsection{What is fixed when a strength varies}

A \emph{central component} has its entire read set disjoint
from the selected physical seeds. Its Gaussian law is still
the selected physical-coordinate law of
Section~\ref{sec:finite-range-gaussian}. Touching factors do
not receive independent Gaussian fields. All masks, shared
tags, frozen Haar integrations outside the read set, and
normalising constants retain their prescribed meanings.

Write ordinary nonincoming logarithmic corrections as
\(e^{v_\alpha}\). These include smooth row corrections
and the original bounded masked fine path terms.
The latter keep their physical arguments and are held
fixed when coarse derivatives are moved to the density,
as explained in Section~\ref{sec:central-derivative-estimates}.
Introduce a complex scalar \(u\) by replacing them with
\(e^{uv_\alpha}\). A strength is an auxiliary scalar used to
count coefficients; it does not change the coupling or the
reference covariance. Incoming occurrences have the form
\[
 e^{-e_i(U)},\qquad e_i(U)=\chi_{X_i}(U)E_i(U).
\]
Their strength, denoted by \(v\), replaces this expression
with \(e^{-v e_i(U)}\). Scalar dependence can be entire even
if the physical multiplier is only bounded Borel.

Warning factors and physical tests are not scaled. In
particular the broad chart factors retained on Borel reads
by \eqref{eq:regular-retained-chart-identity} stay in place.
Any smooth extension outside those charts must agree with
the original smooth formula on their support. No extension
is made in the arguments of an incoming function.
This requirement is essential when strengths are complex.

Use \(\|\cdot\|_{\mathrm{raw},\kappa}\) for the sum in
\eqref{eq:E-norm}, with the same maximum-over-orders
derivative norm, before imposing zero flat value or zero
affine Hessian. This is not an additional induction norm.
All auxiliary estimates may first be made with the
factorial-weighted seminorm; convert to the prescribed
norm once after forming a connected family.

\subsection{The precise analytic input}

Let \(F(u,v)\) denote the warning-free connected logarithmic
family with these strengths, and put \(F(u)=F(u,0)\).
Fix positive bounds \(b,d\) for the smooth and incoming
incidence prices, including the chosen covering reserves.
When there are no incoming letters, omit \(v,d\).
The following input will be called the
\emph{joint-strength bound}:
there are constants \(C_B,D_B\), polynomial in \(1+B\),
and a sufficiently small fixed \(c>0\), such that \(F\)
is analytic on a neighbourhood of the closed bidisc
\begin{equation}
 |u|\le R=\frac{c}{D_Bb},\qquad
 |v|\le S=\frac{c}{D_Bd},\qquad R,S\ge2,
 \label{eq:central-strength-radii}
\end{equation}
and has an absolutely summable connected-family majorant
there satisfying
\begin{align}
 \sup_{|u|\le R}\|F(u,0)\|_{\mathrm{raw},\kappa}
       &\le C_B Rb,\notag\\
 \sup_{\substack{|u|\le R\\|v|\le S}}
       \|F(u,v)\|_{\mathrm{raw},\kappa}
       &\le C_B(Rb+Sd).
 \label{eq:central-joint-strength-input}
\end{align}
The majorant must also hold after every prescribed ordered
derivative and before the family labels are summed.
Thus the two estimates are local family estimates, not
bounds for one total integral growing with the volume.
Their constants must be uniform in cutoff, volume, and
finite history length.

Here is how the preceding calculus reduces the verification
of this input. A bounded logarithmic letter in the smooth
or held-physical coordinate convention satisfies
\[
 \|e^{uv_\alpha}-1\|_{\rm alg}
 \le e^{|u|\|v_\alpha\|_{\rm alg}}-1
 \le 2|u|\|v_\alpha\|_{\rm alg}
\]
when \(|u|\|v_\alpha\|_{\rm alg}\) is sufficiently small.
For an unbounded Gaussian letter with
\(|v_\alpha|\le c_\alpha Q_\alpha\), \(Q_\alpha\ge0\),
\begin{equation}
 |e^{uv_\alpha}-1|
 \le |u|\sqrt{c_\alpha}\,
       e^{(\sqrt{c_\alpha}+|u|c_\alpha)Q_\alpha}.
 \label{eq:central-complex-Mayer}
\end{equation}
Indeed \(|e^w-1|\le |w|e^{|w|}\) and
\(\sqrt{c_\alpha}Q_\alpha
\le e^{\sqrt{c_\alpha}Q_\alpha}\).
Sum these quadratic charges over the whole selected word,
then apply Lemma~\ref{lem:one-Gaussian-source-bound}.
Fixed-order polynomial insertions use
Corollary~\ref{cor:finite-range-tilted-polynomial}
under that same tilted law.
Finally apply Lemma~\ref{lem:rooted-interacting-log},
using Lemma~\ref{lem:finite-jet-anchor-product} once
for the common anchor.

These operations prove
\eqref{eq:central-joint-strength-input} whenever the
integrated primitive derivative prices and the common
precision reserve hold throughout the indicated bidisc.
The primitive bounds are proved in
Section~\ref{sec:central-derivative-estimates}.
In particular the rough score estimate used for a physical
packet in Section~\ref{sec:integrated-regular-derivatives}
cannot simply be transferred to a central component with
no rare energy available to pay for it.

\begin{proposition}[Only two ordinary coefficients need matching]
\label{prop:central-two-coefficients}
Assume the one-variable joint-strength bound. Then
\begin{equation}
 F(1)=F'(0)+\tfrac12F''(0)+R_{\ge3},\qquad
 \|R_{\ge3}\|_{\mathrm{raw},\kappa}\le C_B b^3.
 \label{eq:central-two-coefficients}
\end{equation}
Constants in the conclusion absorb fixed powers of \(D_B\).
They inherit the uniformity of the input.
\end{proposition}
\begin{proof}
There is no selected letter at \(u=0\), so \(F(0)=0\).
Apply Cauchy's coefficient formula to each cluster
function and each ordered derivative. Integrating
the nonnegative family majorants on \(|u|=R\)
justifies all sums and suprema. The coefficient of
\(u^n\) has norm at most \(C_B RbR^{-n}\).
Consequently
\[
 \|R_{\ge3}\|_{\mathrm{raw},\kappa}
 \le C_B Rb\,\frac{R^{-3}}{1-R^{-1}}
 \le 2C_B bR^{-2}\le C_Bb^3 .
\]
This proves the assertion in the original derivative norm.
\end{proof}

When the ordinary series is the logarithm of the protected
ordinary Gaussian integral, let \(V=\sum_\alpha v_\alpha\),
including deterministic terms and the original fine masks.
At each fixed finite volume,
\[
 F'(0)=\mathbb E V,\qquad
 \tfrac12F''(0)
 =\tfrac12\{\mathbb EV^2-(\mathbb EV)^2\}.
\]
These equalities identify total actions. Their decomposition
into local families uses the same labelled connected
expansion. Classical and one-loop matching concerns these
two coefficients; the scalar argument does not establish
their cancellations.

\begin{proposition}[A single incoming insertion and its remainder]
\label{prop:central-inherited-comparison}
Assume the bidisc joint-strength bound. Then
\begin{equation}
 F(1,1)-F(1,0)=\partial_vF(0,0)+R_{\rm cmp},\qquad
 \|R_{\rm cmp}\|_{\mathrm{raw},\kappa}
       \le C_B(bd+d^2).
 \label{eq:central-inherited-remainder}
\end{equation}
The derivative \(\partial_vF(0,0)\) contains precisely one
incoming occurrence and no smooth interaction letter.
It retains the original function and its mask.
\end{proposition}
\begin{proof}
Subtract \(F(u,0)+v\partial_vF(0,0)\) from \(F(u,v)\).
The remaining coefficients have either positive degree
in both variables, or zero \(u\)-degree and \(v\)-degree
at least two. If \({\mathsf M}=C_B(Rb+Sd)\), Cauchy's formula gives
\begin{align*}
 \|\text{mixed terms at }(1,1)\|_{\mathrm{raw},\kappa}
       &\le \frac{{\mathsf M}}{(R-1)(S-1)},\\
 \|\text{pure higher incoming terms at }(1,1)\|_{\mathrm{raw},\kappa}
       &\le {\mathsf M}\frac{S^{-2}}{1-S^{-1}}.
\end{align*}
Since \(Rb=Sd=c/D_B\) and \(R,S\ge2\), these are at most
\(C_Bbd\) and \(C_Bd^2\), respectively.
The same absolute majorants justify the family sums.
No physical derivative of \(E_i\) has been created by
the scalar coefficient extraction.
\end{proof}

For example, when the reference coordinates on the full
read set \(T_i\) have no frozen entries, this single
insertion is exactly
\begin{equation}
 \partial_vF(0,0)=
 -\sum_i\mathbb E_{C_{\rm sel}}
    \left[\left(\prod_{e\in T_i}b_e\right)
             \chi_{X_i}(U)E_i(U)\right].
 \label{eq:central-principal-incoming-average}
\end{equation}
Here \(b_e\) is the original broad chart factor and the
mean is included in the installed physical links \(U\).
The prescribed shared tag and frame average is understood;
it is unchanged on both sides. An inner Gaussian protector
may replace the displayed mask and broad factors only
after the difference has received its own joint warning
estimate, including derivatives. Such a replacement is
not implicit in Proposition~\ref{prop:central-inherited-comparison}.

\subsection{The role of the strength reduction}

The intended scale bounds are
\[
 b\le\widehat\beta_{\rm f}^{-1/2+o(1)},\qquad
 d\le\widehat\beta_{\rm f}^{o(1)}\delta,\qquad
 \delta\le\widehat\beta_{\rm f}^{-11/10}.
\]
With the joint-strength input proved in
Theorem~\ref{thm:central-joint-strength-bound},
the ordinary higher remainder is
\(\widehat\beta_{\rm f}^{-3/2+o(1)}\), while the incoming
comparison remainder is at most
\(\widehat\beta_{\rm f}^{-1/2+o(1)}\delta\).
Warnings are handled separately by
Lemma~\ref{lem:marked-interacting-log}; their factors have
not been rescaled.

This replaces repeated higher-order expansions by one
coefficient estimate. Beyond the primitive central bounds
proved in Section~\ref{sec:central-derivative-estimates},
the sharp transport of the single incoming insertion,
with its constant chosen before \(L\), is proved in
Section~\ref{sec:incoming-sharp-transport}.
Section~\ref{sec:ordinary-matching} identifies the first
two ordinary coefficients on the auxiliary periods.
The exact localisation in
Section~\ref{sec:finite-period-localisation} then
transfers the resulting auxiliary-action bound.
Together with the incoming estimate these are the
central inputs of
Theorem~\ref{thm:nonlinear-one-step}.
\clearpage
\section{Finite-period localisation and exact mask corrections}
\label{sec:finite-period-localisation}

A small estimate for the sum of local terms on an auxiliary
torus need not be a small estimate for each summand.
Estimating the summands separately can destroy a cancellation.
We therefore localise the total auxiliary action, then put
back the original masks exactly.
The construction is a fourth finite difference of reflected
boxes, as in \cite[Lemma~13.6, PDF p.~98]{ArmstrongDraft}.
We prove the combinatorial identity, its physical meaning,
and its norm bound here. The small auxiliary-action estimate
is an input, not a consequence of this identity.

\subsection{Local formulas and reflected connections}

Fix an integer \(R\ge1\). In this section \(R\) is a read
radius, unrelated to the strength radius in
\eqref{eq:central-strength-radii}.
Select a finite collection of translation-covariant local
term types \(f_a\), with supports \(X_a\).
Choose nonnegative shares \(p_{xa}\) satisfying
\(\sum_xp_{xa}=1\), and set
\begin{equation}
 F_x(V)=\sum_a p_{xa}f_a(V),\qquad c=F_x(1).
 \label{eq:finite-period-site-formula}
\end{equation}
Every term contributing at \(x\), including its mask,
reads only links and plaquettes in \(N_R(x)\).
The formulas are literal finite-stencil formulas:
transferring them to another torus changes only the
coordinates, not their coefficients or the rule of evaluation.
In particular this step does not transfer an arbitrary
volume-dependent covariance or an incoming function
of unbounded support to a smaller torus.

Assume the site formula is covariant under translations
and coordinate reflections. If necessary, average the
entire indexed family, its masks, and its shares over
the finite reflection group. Whenever the total masked
action has these symmetries, this operation preserves
that action exactly. The radius includes the resulting
finite stencil enlargement. Gauge invariance of all
formulas is also retained.

Choose integers \(w\ge4R+10\) and
\[
 L_i=w+i,\quad
 \Lambda_i=(\mathbb Z/(2L_i)\mathbb Z)^4,\quad
 i=0,1,2,3,4,\qquad e=w+R+5.
\]
The physical torus has side greater than \(10e\).
The notation \(F_y^{(i)}\) means the same site formula
on \(\Lambda_i\).
For a coordinate \(u\in\{0,\ldots,2L_i-1\}\), put
\[
 \pi_i(u)=\min(u,2L_i-u),
\]
and apply \(\pi_i\) coordinatewise to a site.

For a physical site \(z\), define a connection
\(V^{i,z}\) on \(\Lambda_i\) by folding into the physical
box \(z+[0,L_i]^4\). An auxiliary oriented edge maps to
a unit physical edge: if its image has positive orientation,
use that link; if negative, use the inverse of the
positively oriented link. There is no collapsed edge.
Each auxiliary plaquette maps to a physical unit
plaquette, possibly reversed and with a different starting
vertex. Its holonomy is therefore a conjugate of the
physical holonomy or its inverse. The plaquette chord,
and hence the good-field condition, is preserved.
Gauge transformations pull back by the same map.

\begin{lemma}[Fourth-difference folding identity]
\label{lem:fourth-difference-fold}
Let
\[
 \lambda_i=\frac{(-1)^i\binom4i}{16\cdot4!},\qquad
 H_{iy}^{z}(V)=F_y^{(i)}(V^{i,z})-c.
\]
Then, pointwise in the physical connection,
\begin{equation}
 \sum_{i=0}^4\lambda_i\sum_{y\in\Lambda_i}
       H_{iy}^{\,x-\pi_i(y)}(V)=F_x(V)-c.
 \label{eq:fourth-difference-fold}
\end{equation}
\end{lemma}
\begin{proof}
After the shift \(z=x-\pi_i(y)\), every read is rooted
at \(x\). Classify each coordinate of \(y\) by its
distance to the two fold faces and its orientation.
An interior coordinate has
\(R\le \pi_i(y_j)\le L_i-R\). There are
\(2(L_i-2R+1)\) such coordinates.
On an entirely interior stencil the fold is a
translation followed by coordinate reflections.
Its value is \(F_x(V)-c\) by the stated covariance.
The multiplicity of these stencils is
\[
             [2(w+i-2R+1)]^4.
\]

A boundary coordinate records a distance less than \(R\)
from one of the faces, together with its fold orientation.
For each such record the relative folded read stencil
at \(x\) is independent of \(i\).
For example, close to the lower face it is determined by
the map \(t\mapsto |u+t|-u\) at the recorded distance \(u\);
the upper face gives its reflected version.
There are finitely many records, depending on \(R\).
If at least one coordinate is a boundary coordinate,
the other interior-coordinate multiplicities form a
polynomial in \(i\) of degree at most three.
The value associated with that record is independent
of \(i\), even if the fold changes the local connection.

The identities
\[
 \sum_{i=0}^4(-1)^i\binom4i i^k=0\quad(0\le k<4),
 \qquad
 \sum_{i=0}^4(-1)^i\binom4i i^4=4!
\]
follow by taking the fourth forward difference of
the polynomial \(t^k\).
They cancel every boundary record. The interior
polynomial has leading coefficient \(16\), so its
weighted multiplicity is one. This proves
\eqref{eq:fourth-difference-fold}.
\end{proof}

\subsection{Putting every original mask back}

Define the unmasked local function
\begin{equation}
 G_z(V)=\sum_{i=0}^4\lambda_i
                \sum_{y\in\Lambda_i}H_{iy}^{z}(V).
 \label{eq:finite-period-local-function}
\end{equation}
It reads only \(N_e(z)\).
Use the original good-field mask \(\chi_X\) from
Section~\ref{sec:activities}, at the coarse threshold.
Set
\begin{equation}
 \begin{split}
 L_x(V)={}&\sum_a p_{xa}\chi_{X_a}(V)f_a(V)-c\\
 &-\sum_{i=0}^4\lambda_i\sum_{y\in\Lambda_i}
    \chi_{N_e(x-\pi_i(y))}(V)
                 H_{iy}^{\,x-\pi_i(y)}(V).
 \end{split}
 \label{eq:finite-period-mask-discrepancy}
\end{equation}
These products use zero outside their masked domains;
no derivative of a hard mask is intended.

\begin{proposition}[Exact masked localisation]
\label{prop:finite-period-masked-localisation}
The following identity is exact:
\begin{equation}
 \sum_a\chi_{X_a}f_a
  =c\#\Lambda+\sum_z\chi_{N_e(z)}G_z+\sum_x L_x.
 \label{eq:finite-period-masked-localisation}
\end{equation}
Each \(L_x\) is local on \(N_{2e}(x)\) and vanishes
when this neighbourhood is healthy.
Gauge invariance is preserved. All terms outside
the selected finite collection remain unchanged.
\end{proposition}
\begin{proof}
Sum \eqref{eq:finite-period-mask-discrepancy} over \(x\).
For each \(i,y\), the change of variables \(z=x-\pi_i(y)\)
is a bijection on the physical torus. The last double
sum is consequently
\(\sum_z\chi_{N_e(z)}G_z\). The shares sum to one,
which gives \eqref{eq:finite-period-masked-localisation}.

All centres \(x-\pi_i(y)\) are within distance \(w+4\)
of \(x\); their \(e\)-neighbourhoods lie in \(N_{2e}(x)\).
If this larger neighbourhood is healthy, all masks in
\eqref{eq:finite-period-mask-discrepancy} equal one.
Then \eqref{eq:fourth-difference-fold} makes \(L_x=0\).
Pullback commutes with gauge transformations, so every
displayed scalar term is gauge invariant.
\end{proof}

\subsection{A bound from the total auxiliary action}

Put \(S_i(V)=\sum_{y\in\Lambda_i}F_y^{(i)}(V)\).
The important estimate is for \(S_i-S_i(1)\), not for
each \(F_y^{(i)}-c\) separately.

\begin{lemma}[Local derivative bound after folding]
\label{lem:finite-period-fold-norm}
Suppose the auxiliary actions have \(C^q\) germs on their
good-field domains and satisfy
\[
 \|S_i-S_i(1)\|_{q,h_{\rm c},s;0,\Lambda_i}\le A_i.
\]
If \(h_{\rm c}(1+w+R)^s\le1\), then
\begin{equation}
 \|G_z\|_{q,h_{\rm c},s;z,N_e(z)}
             \le C_{q,s,G}\sum_{i=0}^4|\lambda_i|A_i .
 \label{eq:finite-period-fold-norm}
\end{equation}
The constant is independent of \(w\), the physical
volume, and the cutoff. In particular it carries
no factor \(\#\Lambda_i\).
\end{lemma}
\begin{proof}
Write
\[
 G_z=\sum_i\lambda_i
          [S_i(V^{i,z})-S_i(1)].
\]
A healthy physical \(N_e(z)\) gives a healthy auxiliary
connection by the plaquette comparison above.
For a positively oriented image edge the differential
of pullback just copies the physical left direction.
For a reversed edge, inversion sends it to
\(-\operatorname{Ad}_{V_e^{-1}}A_e\), where
\(\operatorname{Ad}_g A=gAg^{-1}\).
This operation preserves its norm.

For \(y\in\Lambda_i\), the auxiliary distance from \(y\)
to \(0\) is \(|\pi_i(y)|_\infty\).
The source of the physical edge is \(z+\pi_i(y)\)
or differs from it by one unit. Thus a permitted
physical direction pulls back to at most \(2^s\)
times a permitted auxiliary direction.
Repeated ordered differentiation also differentiates
the adjoint coefficients at reversed edges.
An \(r\)-fold coefficient is bounded by
\[
 C_{r,G}h_{\rm c}^{\,r}
            (1+|\pi_i(y)|_\infty)^{rs}.
\]
Relative to one allowed auxiliary direction this
costs at most
\(C_{r,s,G}[h_{\rm c}(1+w+R)^s]^{r-1}\).
The hypothesis bounds it uniformly.
The ordered chain rule has finitely many terms
depending only on \(q\). Applying the assumed
auxiliary derivative estimate to each term gives
\eqref{eq:finite-period-fold-norm}.
No sum over edges is needed: each coefficient field
is itself an allowed full link direction.
This also proves existence of the composed germs,
including at the boundary of the healthy domain.
\end{proof}

There is a separate rough bound for the mask discrepancy.
Suppose
\(\sum_a p_{xa}\|f_a\|_{\infty,\mathrm{good}}\le\delta_0\)
uniformly in \(x\), and the same bound holds for the
transferred formulas on all five auxiliary tori.
Then \(|c|\le\delta_0\); whenever a term
\(\chi_{N_e(z)}H_{iy}^z\) is present, the folded
connection is healthy and \(|H_{iy}^z|\le2\delta_0\).
Consequently
\begin{equation}
          \|L_x\|_\infty\le C(1+w)^4\delta_0 .
 \label{eq:finite-period-mask-price}
\end{equation}
If \(R\le r_jB\) and \(w\le C_jB\) at a fixed profile,
then \(e\le C'_jB\) and
\(n_B(N_{2e}(x))\le C''_j\).
Thus this rough price, including its fixed covering
weight, is of the allowed \(C_B\delta_0\) type.
The small bound for \(G_z\) instead comes from the
total-action estimate in
Lemma~\ref{lem:finite-period-fold-norm}, where the
low-order cancellations must already have occurred.

\subsection{Period margins along one Wilson family}

The physical and auxiliary periods have different roles:
the latter are finite comparison devices. Their use
must not select a new physical Wilson family at each
volume. The required margin is compatible with the
single family already fixed.

\begin{lemma}[Eventual polynomial period reserve]
\label{lem:finite-period-reserve}
Along the scale parameter \(s_m\to\infty\), suppose
the lattice side is \(\ell_p e^{s_m}\), with
\(\ell_p\ge\ell_{\min}>0\), and
\(\widehat\beta_m=\bar\beta+b_0s_m\), \(b_0>0\).
For \(B_m=\lceil(\log\widehat\beta_m)^2\rceil\),
every fixed \(C,k\) admits an entry scale, independent
of \(p\), after which
\[
                 \ell_p e^{s_m}\ge C(1+B_m)^k .
\]
\end{lemma}
\begin{proof}
The logarithm of the right-hand side is
\(O(1+\log\log(\bar\beta+b_0s_m))=o(s_m)\).
Subtracting it from \(s_m+\log\ell_{\min}\) tends
to infinity. The entry scale therefore works for
all \(p\) simultaneously.
\end{proof}

At a fixed profile the finitely many constants in
the read radius and the derivative bounds are chosen
first. Their period margins and
\(h_{\rm c}(1+w+R)^s\le1\) can then be included among
the entry conditions for that profile.
Only afterwards is the slow profile clock chosen.
Early levels require their separately prescribed base
period reserve; the lemma is an eventual statement.
It does not deduce a larger margin from an insufficient
numerical lower bound at those levels.

The exact identity above may be used for selected
finite-stencil low-order terms once their auxiliary
action bound is proved. It does not localise arbitrary
incoming regular activities by changing their
arguments. Those use the principal transport
proved in Section~\ref{sec:incoming-sharp-transport},
with their original physical arguments.
\clearpage
\section{Central derivative estimates and a common precision reserve}
\label{sec:central-derivative-estimates}

We now prove the joint-strength input of
Section~\ref{sec:central-strength-expansion} for the central
letters of the physical finite-range construction. The
argument concerns the full read supports, with the original
incoming functions. It is not a change to a white-coordinate
reference or to independent Gaussian fields at neighbouring
addresses. Its two ingredients are a bounded smooth remainder
after calibration and one integration estimate for the
derivatives of the original local law.

The organisation of simultaneous local integration follows
\cite[Lemma~11.13, PDF p.~73]{ArmstrongDraft}.
We use the elementary product calculus already proved here,
rather than importing a field-space norm from another model;
compare \cite[Proposition~3.7, PDF p.~15]{BS15a}.
The role of affine Gaussian subtraction in the fluctuation
integral is described in \cite[Sect.~2, pp.~265--268]{Bal87}.
All quantitative estimates needed below are established in
the present norms.

\subsection{Parameters, domains, and the price to be proved}

Fix the adjacent profiles, their odd block ratio \(L\), the
noise reserve \(W\), the derivative order \(q\), the direction
power \(s\), and the covering exponents used in the local
expansion. The actual coupling is \(\beta\), with
\(g=\beta^{-1/2}\), in its prescribed nominal coupling band.
Recall
\[
 P=\exp\{(\log\log\widehat\beta_{\rm f})^2\},\qquad
 h=\widehat\beta_{\rm f}^{-1/2}P,\qquad
 B=\lceil(\log\widehat\beta_{\rm c})^2\rceil .
\]
The adjacent nominal shift is fixed once this ratio is
fixed. In particular \(h_{\rm c}\le2h\),
\(h/g\le CP\), and \(g\) is comparable with
\(\widehat\beta_{\rm f}^{-1/2}\) in this band.
The input path bounds are those of the fixed profiles;
write \(\Lambda\) for their allowed quadratic-loop price.
The existing slow-clock choice gives
\(\log(1+\Lambda)=o(\log\widehat\beta_{\rm f})\).

As before, \(C_B\) denotes \(C(1+B)^p\), with fixed \(C,p\).
In this section these constants may depend on the fixed
ratio, \(W,q,s\), group representation, profiles, and
the specified covering reserves. They do not depend on
period, ultraviolet cutoff, sector, or finite history
length. No pre-ratio uniformity is asserted here.
That stronger requirement belongs to the principal
incoming transport estimate, not to the present remainder.

At a coarse root \(o\) the directions satisfy
\[
 |A_e|\le h_{\rm c}(1+d(e,o))^s.
\]
We use the ordered left derivatives and the good-field
norm of \eqref{eq:local-derivative-norm}. Hard coarse
masks are restored after estimating the unmasked germ.
Hard fine masks are held fixed in physical variables.

Choose a fixed unused covering reserve. In applying
Lemma~\ref{lem:incoming-enlarged-income}, use its incoming
exponent \(\kappa_{\rm in}\) with
\[
 \kappa_{\rm in}\ge c_1(\kappa+a)+1.
\]
Here \(\kappa\) is the output exponent of Section 32;
the fixed \(a>0\) pays the component summation, the
common-anchor polynomial, and the connected logarithm.
One may take \(a\ge8\) after the geometric enlargement
constants have been included in \(c_1\).
These are the covering reserves in the component input,
not a new differentiability assumption.
Alternatively, Section~\ref{sec:incoming-sharp-transport}
proves the sharper compression estimate
\eqref{eq:incoming-compressed-incidence}.
With its ratio lower bound at exponent \(\kappa+a\),
one may use any fixed positive \(\kappa_{\rm in}\);
in particular repeated profiles may keep equal
incoming and outgoing covering exponents.
In either case set
\[
                  \delta=\|E\|_{\mathcal E,\kappa_{\rm in}}.
\]
As in Section 32, the incoming scale regime is
\(\delta\le\widehat\beta_{\rm f}^{-11/10}\).
Only its value bound is needed for the estimates in
this section.

We shall construct a positive envelope
\begin{align}
 A_\beta&=C(1+B)^{p_0}(1+P)^{p_1}
                              (1+\Lambda)^{p_2}\ge1,\notag\\
 \epsilon_B&=C_Be^{-cB},\notag\\
 b&=A_\beta\{gP^3+(1+\Lambda)h^2+\epsilon_B\},
       \qquad d=A_\beta\delta .
 \label{eq:central-explicit-prices}
\end{align}
The constants are enlarged finitely many times below;
choose their final values before defining the strength
radii. In particular \(b\) is this chosen envelope,
not the possibly much smaller actual incidence sum.
It satisfies
\begin{equation}
 b=\widehat\beta_{\rm f}^{-1/2+o(1)},\qquad
 A_\beta=\widehat\beta_{\rm f}^{o(1)},\qquad
 b\ge gP^3\ge c\,\widehat\beta_{\rm f}^{-1/2}.
 \label{eq:central-envelope-rates}
\end{equation}
The last inequality will control the largest strength
radius. If \(\delta=0\), omit the incoming variable.

\subsection{Protect the smooth remainder after its cancellations}

The exact flat Gaussian factor is already in the scalar
ledger \eqref{eq:finite-range-vacuum-ledger}, including
the noise and Haar normalisations. Smooth scalar
corrections below are the paired field-dependent
differences from that reference, together with the
retained completed-square constant. The extensive flat
determinant is not called a small activity.

Let \(H=hB^{17}\). Choose a smooth radial map
\(\mathfrak g\to\mathfrak g\), denoted \(a\mapsto a^{\rm pr}\),
equal to \(a\) on the support of the original broad
chart test, bounded by \(C H\), and satisfying
\begin{equation}
       \|D^r a^{\rm pr}\|\le C_r H^{1-r}\quad(r\ge1).
 \label{eq:central-smooth-protection}
\end{equation}
For example use a smooth function of \(|a|^2/H^2\)
multiplying \(a\), with its transition outside that
support. It commutes with orthogonal adjoint transport.
Protect only smooth auxiliary row formulas.
Incoming functions and all original tests keep their
physical arguments.

More precisely, let \(\mathcal R_{\rm sm}(V,a)\)
be the combined smooth energy remainder after subtracting
the installed affine squares. Its extension is
\[
 \mathcal R_{\rm sm}^{\rm pr}(V,a)
                       =\mathcal R_{\rm sm}(V,a^{\rm pr}).
\]
Protect this difference as a whole, after its
cancellations. The reference quadratic
\(\tfrac12\|Da+\alpha\|^2\) is still evaluated at
the original \(a\); neither it nor the selected
Gaussian law is replaced by a protected quadratic.
On the retained chart \(a^{\rm pr}=a\), so adding
the reference quadratic to this extended remainder
gives the original energy there.

This replacement is made in a fixed central column
before expanding its broad and default tests.
On the joint support of those tests the retained row
predicates equal one, the fine direct clipping is
quadratic, and the original formula equals its smooth
extension. Multiplying by these same tests therefore
gives an exact identity for the column. Expanding them
afterwards produces the original warning factors.
No formula in a packet column is replaced in this
operation.

On a healthy coarse stencil the local axial gauge and
its path comparisons give coordinates of size
\(C_Bh_{\rm c}\), with the corresponding fixed-order
scaled derivative bounds. Together with
\eqref{eq:central-smooth-protection}, all arguments of
a protected row lie in a fixed chart of size \(C_Bh\).
This tends to zero. A background-dependent gauge causes
no unbounded-coordinate loss: equivariance lets one
differentiate the protected coordinate itself, whose
size is at most \(CH\).

We record the elementary calculus used here. If a
smooth function \(f\) has vanishing jets of degrees
less than \(m\) at zero, Taylor's integral formula gives
\[
 \|D^r f(y)\|\le C|y|^{(m-r)_+}\quad(0\le r\le q),
\]
on a smaller fixed chart; higher derivatives through
\(q+m\) bound the displayed constants. If every
positive-order scaled derivative of \(y\) is bounded
by \(C_Bh\), the ordered chain rule gives
\begin{equation}
           |D_{A_1}\cdots D_{A_j}f(y)|
                      \le C_Bh^m,\qquad 0\le j\le q.
 \label{eq:central-scaled-vanishing-jet}
\end{equation}
Indeed a term with \(r\) derivatives of \(f\) has \(r\)
argument factors. For \(r\le m\) their product with
\(|y|^{m-r}\) is \(O(C_Bh^m)\); for \(r>m\) it is
\(O(C_Bh^r)\), which has the same bound eventually.
Repeated derivatives of a protector cost powers
\(H^{1-r}\), paired with \(r\) argument variations of
size \(C_Bh\); the same conclusion follows, with
the extra powers of \(B\) in \(C_B\).

\begin{lemma}[Bounded central logarithmic prices]
\label{lem:central-smooth-derivative-prices}
The protected smooth row corrections, the original
masked microscopic path terms, and the paired
local determinant differences admit elementary
addresses \(\alpha\) with local logarithms \(v_\alpha\)
and nonnegative prices \(t_\alpha\) such that,
in mixed coordinates with all Borel physical reads
held fixed,
\begin{align}
 \|v_\alpha\|^{\rm alg}_{q,h_{\rm c},s;o_\alpha,X_\alpha}
      &\le t_\alpha,\notag\\
 \sup_x\sum_{\alpha:x\in X_\alpha}
       e^{(\kappa+a)n_B(X_\alpha)}t_\alpha
      &\le C_B\{gP^3+(1+\Lambda)h^2+\epsilon_B\}.
 \label{eq:central-smooth-derivative-prices}
\end{align}
The Gaussian coordinates not changed to physical
variables are parameters in this norm. The bounds
are uniform in those parameters and the shared tags.
For an original fine masked path term its entire
physical expression is held fixed; no smoothness
across its fine mask is claimed.
The prices are uniform in the physical period and
the finite history. Small Gaussian correction letters,
whose coordinates remain unbounded, are estimated
separately below.
\end{lemma}
\begin{proof}
First keep the exact row grouping
\eqref{eq:packet-six-rows}. For fine and noise rows
write their true residual as \(r+e\), where
\(r=Da+\alpha\) is the installed affine row.
The local chart expansion of its multiplication,
curvature and predictor gives
\[
 |r|\le C_Bh,\qquad |e|\le C_Bh^2 .
\]
These are also scaled jet bounds through order \(q\):
the remainder in the affine expansion has at least
two argument factors, and
\eqref{eq:central-scaled-vanishing-jet} applies.
Consequently, before taking absolute values,
\[
 \tfrac12|r+e|^2-\tfrac12|r|^2
                  =\langle r,e\rangle+\tfrac12|e|^2
\]
has jet price \(C_Bh^3\).
The direct chord/curvature defect starts at degree
four and has price \(C_Bh^4\).

Here is the cancellation for the other rows.
Use the uncut flat operators temporarily, write
\(v_1=D\mathcal C_{\rm c}(1)x\), and let the first
joint jet of the fine and noise rows be
\(F_1=\mathsf A a+\mathsf Bx\).
Theorem~\ref{thm:constructed-calibration} and the
full row completion in Section 30 imply
\[
 \Phi_0^T\mathsf A=0,\qquad
 \Phi_0^TF_1=v_1,\qquad \Phi_{\perp,0}v_1=0 .
\]
For the middle identity, use
\(\mathsf B+\mathsf A\mathcal H=\Phi_0D\mathcal C_{\rm c}(1)\)
and
\(\Phi_0^T\Phi_0+\Phi_{\perp,0}^T\Phi_{\perp,0}=1\).
Thus the linear jet in the bracket defining \(w_L\)
vanishes; \(w_L\) has degree at least three.
The vector \({\Phi_{\perp}}v\) has zero linear jet, so \(w_{\perp}\)
has degree at least four.
The coefficient bracket defining \(w_D\) has zero
flat value; \(w_D\) has degree at least three.
The coarse direct defect in \(w_C\) has degree four.
Native cuts leave an error \(C_Be^{-cB}\), whose
fixed-order derivatives have the same bound after
decreasing \(c>0\).
Multiplying by their two curvature factors does not
create a second-order term of order one.

It follows that the six-row energy after its affine
subtraction has jet price
\[
 C_B(h^3+e^{-cB})
\]
at each local address. Its Boltzmann exponent has
price \(C_B\beta h^3\le C_BgP^3\), with the actual
coupling-band constants included. Factors
\(\beta^{\mathsf M} e^{-cB}\) for any fixed \({\mathsf M}\) are still
bounded by \(C_Be^{-c'B}\); the entry threshold is
chosen after \({\mathsf M}\).
This estimates the combined correction, not the
separate large affine squares.

For the constant square at the selected mean,
flat matching and the mean equation give
\[
  |D\mu+\alpha|\le C_B(h_{\rm c}^2+e^{-cB}).
\]
One can see this directly by comparing \(\mu\)
with the stationary completion in
Lemma~\ref{lem:selected-centre-intact-box} and
using \eqref{eq:physical-calibrated-residual}.
For its scaled derivatives the same conclusion
follows by differentiating the local formulas:
their uncut flat first jet is
\(\mathsf B+\mathsf A\mathcal H
 -\Phi_0D\mathcal C_{\rm c}(1)=0\);
the cut first-jet error is exponentially small.
Taylor's formula and the window inverse bounds
therefore give this estimate through order \(q\).
Its constant-square exponent costs
\(C_B\beta h_{\rm c}^4+C_Be^{-cB}\), which is
dominated by the preceding price.
The square and its sign remain in the exact density.

The original fine cubic, quartic, and quadratic-loop
terms retain their fine masks and physical arguments.
On the support of such a mask each curvature slot has
norm at most \(h\), and its parallel transport is
isometric. A degree-\(m\) polynomial is therefore
bounded by \(h^m\) times its absolute coefficient sum;
off that support the masked term is zero.
Their respective value prices are
\(C_B\beta h^3\), \(C_B\beta h^4\), and
\(C_B(1+\Lambda)h^2\).
Declare every such read to be a Borel physical read.
In the mixed coordinates its masked expression is
constant under coarse differentiation. Retained
smooth protectors on those reads are differentiated
separately, with bounded ratios \(h_{\rm c}/H\).
For coarse path terms use the unmasked good-field
germ and the same scaled jet calculation as for
the coarse rows.
For a path of native length \(n\), a fixed number of
derivatives costs a fixed power of \(1+n\).
Reserve a smaller positive part of its original
exponential length weight; then
\((1+n)^m e^{-\eta n}\) is bounded for each fixed \(m\).
The original absolutely summable path coefficients
therefore sum the prices with their declared covering
weights. A long path returning to the same footprint
is paid by its native length, not by its diameter.
Fine and coarse native tails retain their original
clipped expressions. Their coefficient bounds give
\(C_B\beta^{\mathsf M} e^{-cR_B}\) for a fixed \({\mathsf M}\), also for
the coarse unmasked jets. Fine Borel reads are again
held fixed. Since \(R_B\) is a fixed positive fraction
of \(B\) eventually, these prices are included in
\(C_Be^{-c'B}\).

The logarithmic Haar factor has price \(C_Bh^2\),
since \(j_G(0)=1\) and \(j_G(-a)=j_G(a)\).
For the main determinant use the paired formula
\eqref{eq:selected-main-determinant-flow}. In a
local gauge
\(\mathsf M-\mathsf M_0=O(C_Bh_{\rm c})\)
on its read stencil. The inverse difference formula
and the positive interpolation show that each paired
zeroth window trace is \(O(C_Bh_{\rm c})\).
Its fixed scaled derivatives obey the same bound:
the value difference vanishes at the flat field,
whereas each positive-order derivative supplies
at least one coarse factor \(h_{\rm c}\).
The window trace has \(C_B\) entries, not the
number of sites in the physical torus.
The remaining determinant words retain their
small residual factor and are summable as in
Section~\ref{sec:finite-range-components}.
Since \(h_{\rm c}\le CgP\), these prices are also
dominated by the right-hand side of
\eqref{eq:central-smooth-derivative-prices}.

All assertions concern the local scalar traces or
the covariantly transported rows. Orthogonal colour
conjugation preserves their values and norms.
Local gauges have consequently introduced no
global small-link hypothesis. The original
row-count and path incidence bounds now give the
displayed incidence sum. Conversion to the
algebraic seminorm costs only a fixed \(q\)-dependent
constant already included in \(C_B\).
\end{proof}

\subsection{Incoming insertions and derivatives of the common law}

Let \(\Gamma\) be one nonempty selected primitive word,
with distinct elementary addresses, and let \(X\) be
its full connected read footprint. Put \(k=n_B(X)\).
Its number of physical and Gaussian coordinates is at
most \(C_Bk\); the spatial weights at its root cost
a fixed power of \(1+k\).
Distinct species on the same footprint are not assumed
to have bounded multiplicity.

For each incoming occurrence put
\(e_i(U)=\chi_{X_i}(U)E_i(U)\) and
\(e_i^*=\|E_i\|_{q,h,s;o_i,X_i}\).
Then
\begin{equation}
 |e^{-v e_i(U)}-1|
       \le 2|v|e_i^*
       \quad\hbox{if }|v|e_i^*\le c_0.
 \label{eq:central-incoming-complex-price}
\end{equation}
No derivative of this masked expression is taken.
Retain the original broad factors on the union \(T\)
of all Borel physical reads, including any selected
fine masked path terms, by
\eqref{eq:regular-retained-chart-identity}.
Change just these coordinates to physical Haar
variables. The other reads remain Gaussian and use
the mixed density
\eqref{eq:regular-mixed-physical-Gaussian-density}.

\begin{lemma}[Central score and multiplier estimate]
\label{lem:central-common-score}
In these mixed coordinates, derivatives through
order \(q\) of the common density and retained
chart factors have a nonnegative majorant equal
to that density times
\begin{equation}
 C_B(1+P)^{m_q}(1+k)^{m_q}
       \left(1+\sum_{j\in R}z_j^2\right)^{m_q},
 \label{eq:central-common-score}
\end{equation}
for a fixed integer \(m_q\). Here \(R\) is the set
of remaining Gaussian reads. Incoming multipliers
contribute their value prices
\eqref{eq:central-incoming-complex-price} and
no additional Gaussian charge. The majorant is
uniform in period and cutoff.
\end{lemma}
\begin{proof}
Use the exact score
\eqref{eq:regular-selected-score}.
On the retained chart,
\[
 |a_T|+|\mu_T|\le C_Bh,\qquad
 |(a_T-\mu_T)/g|\le C_BP.
\]
Their scaled derivatives have bounds of this type,
with a fixed power of \(1+k\).
Apply Lemma~\ref{lem:regular-selected-score-bounds}
and its mixed version
\eqref{eq:regular-mixed-score-polynomial} to this
same marginal. They use the precision and window
jets of Lemma~\ref{lem:selected-window-jets}; in
particular the trace is over the read coordinates
only. The trace and quadratic score terms have the
common bound
\[
 C_B(1+k)^{m_q}
                  (1+P^2+\sum_R z_j^2).
\]
The Haar logarithmic derivative is bounded on the
fixed chart. Derivatives of a retained broad factor
cost powers of \(h_{\rm c}/H\), and the number of
their placements is polynomial in \(|T|\).
Use the product rule without dividing by a chart
factor that may vanish.
The partition formula
\eqref{eq:selected-density-partition-formula} gives
\eqref{eq:central-common-score}, after increasing
\(m_q\). There are finitely many partitions of the
at most \(q\) ordered positions.

At fixed physical \(U\), \(e_i(U)\) is constant
under these derivatives. Its value bound is global
because its original mask makes it zero off its
healthy domain. Formula
\eqref{eq:central-incoming-complex-price} therefore
supplies its entire contribution.
The factors \((\det C_S)^{-1/2}\),
\(g^{-D_T}\), and \(\zeta_G^{-1}j_G(a)^{-1}\)
in the mixed density have all been retained.
Integrating back in the original coordinates gives
exactly the same probability Gaussian, not a
dimension-dependent comparison constant.
\end{proof}

The important difference from the rough physical
packet estimate is that the small smooth energy
has already been bounded after cancellation.
No term \(\beta h_{\rm c}\) is charged for an
uncancelled physical noise energy.
The remaining score coefficients are only fixed
powers of \(P,1+B,1+k\).

\subsection{Unbounded Gaussian letters at complex strengths}

The unbounded letters come from the original residual
linear term and the covariance correction series
\eqref{eq:selected-ratio-path-series}.
Retain every native intermediate index.
For an address \(\alpha\), let \(c_\alpha\ge0\)
be a coefficient envelope, enlarged to cover its
fixed-order local jets and the reserved native-length
powers. For a quadratic address put
\(Q_\alpha=1+z_i^2+z_j^2\); for a linear address use
\(Q_\alpha=1+z_i^2\).
The estimates of
Lemma~\ref{lem:selected-correction-word-prices},
with a smaller decay rate, give
\begin{align}
 |v_\alpha|&\le c_\alpha Q_\alpha,\notag\\
 \sup_x\sum_{\alpha:x\in X_\alpha}
     e^{(\kappa+a)n_B(X_\alpha)}\sqrt{c_\alpha}
      &\le\epsilon_B,\notag\\
 \sup_i\sum_{\alpha:i\text{ an endpoint of }\alpha}
                 \sqrt{c_\alpha}&\le\epsilon_B.
 \label{eq:central-correction-jet-occupation}
\end{align}
Repeated endpoints are counted with multiplicity.
We justify the derivative enlargement explicitly.
For an invertible matrix \(\mathsf H\), differentiating
\(\mathsf H\mathsf H^{-1}=1\) gives
\(D(\mathsf H^{-1})=-\mathsf H^{-1}(D\mathsf H)\mathsf H^{-1}\).
Differentiating a product
places at most \(q\) insertions among its native
factors. Every resulting term still contains its
small window residual or its differentiated residual.
Exponential entry bounds absorb the fixed spatial
weights and the intermediate-index sums.
Native placement powers are paid by a smaller
exponential reserve. The resulting sums are
\(C_B\sum_{n\ge1}(1+n)^{m_q}
 (C_Be^{-cB})^{n/2}\), with their covering factor,
and hence at most \(C_Be^{-c'B}\).
The residual linear coefficient includes \(g^{-1}\);
its fixed power of the coupling and the monitor is
absorbed by decreasing \(c'\).
This proves \eqref{eq:central-correction-jet-occupation}
for the actual selected coefficients, not for
coefficients of a different reference law.

Under the mixed change of coordinates a derivative
of \(z_T=(a_T-\mu_T)/g\) costs \(C_BP\) times
a footprint polynomial. Thus differentiating a
Gaussian letter adds a fixed-degree polynomial
in its endpoint coordinates, with subpolynomial
coefficients. It does not remove its small
coefficient envelope.

\begin{lemma}[One precision reserve throughout the strength bidisc]
\label{lem:central-complex-precision}
Let \(C_{\rm sel}^{-1}\ge\lambda_{\rm G}1\) be the
fixed-profile precision gap. Choose \(b,d\) by
\eqref{eq:central-explicit-prices}, and set
\[
 R=\frac{c_0}{D_Bb},\qquad S=\frac{c_0}{D_Bd}.
\]
For sufficiently large entry coupling these radii
are at least two. Throughout \(|u|\le R\),
\(|v|\le S\), all selected Gaussian correction
letters and all their derivatives through order
\(q\) admit a common quadratic majorant whose
total diagonal charge is at most
\(\lambda_{\rm G}/16\) at each coordinate.
One may in addition allow a charge
\(\lambda_{\rm G}/16\) for local differentiation
neighbourhoods and \(\lambda_{\rm G}/8\) for
warning witnesses. Even after doubling these
charges, the remaining precision is at least
\(\lambda_{\rm G}/2\).
\end{lemma}
\begin{proof}
For a Gaussian Mayer factor use
\eqref{eq:central-complex-Mayer}. Its cap is
\(|u|\sqrt{c_\alpha}\); its quadratic charge
is bounded by
\[
             C_q(\sqrt{c_\alpha}+|u|c_\alpha)
\]
at its endpoints. The constant part of \(Q_\alpha\)
contributes the same small scalar to its cap.
The finitely many differentiated positions can be
bounded with this charge and fixed-degree polynomial
insertions. To check this, apply the product rule
to \(e^{uv_\alpha}\). A term has between one and
\(q\) factors \(uD^jv_\alpha\).
After extracting the cap
\(|u|\sqrt{c_\alpha}\), its remaining coefficient
is bounded since \(|u|c_\alpha\le1\).
Its endpoint polynomial has degree bounded by
a number depending only on \(q\).
For the undifferentiated factor the elementary
Mayer inequality gives the cap directly.

Because \(c_\alpha\le1\) eventually, the last line
of \eqref{eq:central-correction-jet-occupation}
implies the occupation estimate
\begin{equation}
 \sup_i\sum_{\alpha\ {\rm selected}}d_{\alpha i}(u)
             \le C_q(1+|u|)\epsilon_B.
 \label{eq:central-complex-total-charge}
\end{equation}
The sum is over distinct primitive addresses,
not over powers of one address.
It is bounded by the all-address endpoint sum,
independently of the number of selected letters.

By \eqref{eq:central-envelope-rates},
\(R\le C\widehat\beta_{\rm f}^{1/2}\).
Since \(B\) is quadratic in the logarithm of the
nominal coupling,
\[
 (1+R)C_Be^{-cB}\longrightarrow0 .
\]
Choose the entry coupling so the right-hand side
of \eqref{eq:central-complex-total-charge}
is at most \(\lambda_{\rm G}/16\).
This choice is uniform over all finite words.
The incoming strength does not enter this calculation:
its masked factor is bounded by
\eqref{eq:central-incoming-complex-price} and
has zero quadratic charge.

For clarity, differentiating a density does not
require an extra non-small tilt. At the base field
its derivatives are the polynomial scores of
Lemma~\ref{lem:central-common-score}.
For dominated difference quotients, restrict the
background to a sufficiently small open neighbourhood
with the same unmasked window formulas.
The covariance difference is arbitrarily small
there by the inverse identity. In the retained
physical chart the logarithmic-coordinate shift,
divided by \(g\), is also arbitrarily small.
The ratio to the base density is a scalar times
a linear and quadratic Gaussian exponential.
At the base it is identically one.
For a finite coordinate set \(I\) its coefficients
can consequently be made small enough that
\[
 \left|\log\frac{k(V)}{k(V_0)}\right|
       \le \frac{\epsilon_B}{1+|I|}
                         \left(1+\sum_{i\in I}z_i^2\right)
\]
on the union of the nearby retained chart supports.
This follows by continuity of the coefficients;
use Young's inequality for the linear term and
then reduce the neighbourhood once more.
The extra diagonal charge has sum at most
\(2\epsilon_B\) and maximum at most
\(2\epsilon_B/(1+|I|)\).
It is below \(\lambda_{\rm G}/16\) eventually,
and its determinant ratio is bounded by two
independently of \(|I|\).
This neighbourhood may depend on the finite read
set; the derivative bounds at its base do not.
It provides germs, not a new norm or a new
uniformity assumption on all backgrounds.
At the base the extra likelihood is one:
the quantitative derivative estimate uses only
its polynomial scores, not a nonzero
per-coordinate likelihood charge.

Warning witnesses, if retained, can be estimated
with the smaller exponential-moment parameter
\(\lambda_{\rm G}/8\) in the proof of
Lemma~\ref{lem:finite-range-failure-witness}.
Their weaker rare constant is still sufficient.
The total allowed charge is therefore at most
\(\lambda_{\rm G}/4\), or
\(\lambda_{\rm G}/2\) after doubling.
The base precision is at least
\(\lambda_{\rm G}1\), proving the reserve.
\end{proof}

No independent tilt is applied at each vertex.
If \(d_i=\sum_\alpha d_{\alpha i}\), the determinant
ratio for the common tilt is exactly
\[
 \left(\frac{\det C_{\rm sel}^{-1}}
 {\det(C_{\rm sel}^{-1}-\operatorname{diag}d)}\right)^{1/2}.
\]
Its upper bound from Lemma~\ref{lem:finite-range-joint-tilt}
is distributed as
\(\prod_\alpha\exp\{C\sum_i d_{\alpha i}\}\).
Each correction has at most two endpoints; this
factor changes its cap by a bounded constant.
It does not create an unassigned exponential
in the total number of reads.
The version with linear sources is precisely
Lemma~\ref{lem:one-Gaussian-source-bound}, which
keeps their cross terms.

\subsection{Integrated word bounds and the connected logarithm}

\begin{lemma}[Central integrated derivative majorants]
\label{lem:central-integrated-word}
For every nonempty admitted central primitive word
\(\Gamma\), its integrated local amplitude has a
\(C^q\) germ. At each healthy base field its ordered
derivatives through order \(q\) are dominated by
an integrable function under the same selected
Gaussian, frozen Haar, and shared tag measures.
The integral of this majorant is bounded by
\begin{equation}
 A_\beta(1+n_B(X_\Gamma))^{m_q}
          \prod_{\alpha\in\Gamma}p_\alpha(u,v).
 \label{eq:central-integrated-word}
\end{equation}
Smooth bounded letters have prices \(C|u|t_\alpha\);
Gaussian correction letters have prices
\(C|u|\sqrt{c_\alpha}\); incoming letters have
prices \(C|v|e_i^*\).
Fixed native-length reserves are included in these
prices. The constants and the polynomial degree
are independent of the number of coordinates.
\end{lemma}
\begin{proof}
Use the retained-chart identity and mixed physical
coordinates for the Borel reads.
Each smooth bounded Mayer factor has algebraic
jet price at most \(2|u|t_\alpha\), once
\(|u|t_\alpha\) is sufficiently small.
All its arguments may be protected as above.
Incoming factors use
\eqref{eq:central-incoming-complex-price}.
The differentiated Gaussian factors have the
caps and quadratic majorants just proved.
The product rule acts on their smooth factors
and the common density, with no derivative of
an incoming physical multiplier.

There are at most \(q\) differentiated positions.
Their Gaussian polynomials and the score
\eqref{eq:central-common-score} are bounded by
a fixed power of
\(1+P+\sum_{i\in I_\Gamma}z_i^2\), where
\(|I_\Gamma|\le C_Bn_B(X_\Gamma)\).
Apply Corollary~\ref{cor:finite-range-tilted-polynomial}
under the one common tilt. Its covariance is at
most \(2\lambda_{\rm G}^{-1}1\); all fixed-order
moments are polynomial in \(|I_\Gamma|\).
This proves the displayed polynomial cost and a
fixed power of \(P,1+B\).
The determinant ratio has already been assigned
to the individual correction prices.

There can be arbitrarily many distinct species on
one footprint. If the direct product rule gives
an extra factor \(r^q\) for \(r=|\Gamma|\), use
\(r^q\le C_{q,\eta}e^{\eta r}\), and absorb
\(e^\eta\) in each primitive price.
Do not replace \(r\) by the covering number.
Alternatively the factorial-weighted product
calculus performs this summation directly.
Different local anchors are moved to the final
anchor once; Lemma~\ref{lem:finite-jet-anchor-product}
costs only another footprint polynomial and
keeps the good-field domain fixed.

The mixed change of variables was exact, so all
Haar and Gaussian normalising constants remain.
Probability integration over disjoint frozen
and shared tag variables cannot increase these
uniform bounds.
Dominated differentiation follows from the
precision neighbourhood just constructed and
Lemma~\ref{lem:regular-Borel-measure-derivative}.
This proves the germs, including at a coarse
good-field threshold equality.
\end{proof}

\begin{theorem}[The central joint-strength bound]
\label{thm:central-joint-strength-bound}
For the central word expansion of the original
physical finite-range construction, choose
\(A_\beta,b,d\) as in
\eqref{eq:central-explicit-prices}, with the
fixed covering reserves stated above.
There are \(C_B,D_B\), polynomial in \(1+B\),
and a fixed \(c_0>0\) such that, above a
fixed-profile entry coupling, the warning-free
connected logarithmic family satisfies
\begin{align}
 \sup_{|u|\le R}\|F(u,0)\|_{\mathrm{raw},\kappa}
                  &\le C_B Rb,\notag\\
 \sup_{\substack{|u|\le R\\|v|\le S}}
     \|F(u,v)\|_{\mathrm{raw},\kappa}
                  &\le C_B(Rb+Sd),\notag\\
 R=\frac{c_0}{D_Bb}&\ge2,\qquad
 S=\frac{c_0}{D_Bd}\ge2.
 \label{eq:central-proved-joint-strength}
\end{align}
The family is analytic on a neighbourhood of the
closed bidisc and has a summable majorant for its
ordered physical derivatives through order \(q\).
The estimates are uniform in volume, cutoff, and
finite history length at the fixed profiles.
\end{theorem}
\begin{proof}
Sum the smooth prices using
Lemma~\ref{lem:central-smooth-derivative-prices}.
The Gaussian correction prices have income at
most \(C_B\epsilon_B\).
Use Lemma~\ref{lem:incoming-enlarged-income} for
the incoming value prices, with its original
fine masks and declared read enlargement.
It gives income at most \(C_B\delta\).
The fixed support factors and the common
subpolynomial cost in
\eqref{eq:central-integrated-word} enter
\(A_\beta\).
For a nonempty word a cost \(A_\beta\ge1\)
paid once is bounded by the product obtained
by multiplying each of its letter prices
by \(A_\beta\).
The remaining footprint polynomial is paid
by one positive covering reserve.

The incidence income is consequently at most
\[
                    C(|u|b+|v|d).
\]
Choose \(D_B\) to include the overlap and
primitive-to-component constants, and choose
\(c_0\) so that their smallness criterion
has, for example, a factor two to spare.
The component summation uses connected
distinct primitive subsets. Their disjoint
full reads factor exactly by
Proposition~\ref{prop:finite-range-physical-factorisation},
also after the retained-chart rearrangement.
Thus the resulting activities form the same
hard-core component gas.
Local selection restrictions remain in its
activities; dropping them is used only to
increase the nonnegative majorant.

Apply Lemma~\ref{lem:rooted-interacting-log}
at the reserved covering exponent.
Keep algebraic derivative prices through the
products and convert once to the original
maximum-over-orders norm.
Its pinned majorant is linear in the
component income, and gives
\eqref{eq:central-proved-joint-strength}.
Repeated component species in the ordered
logarithm retain their factorials.
They are products of already integrated
activities; they do not ask for repeated
charges in one primitive Gaussian integral.
The occupation estimate was needed only
inside each distinct-address primitive word.

Every finite word is entire in its bounded
incoming strengths and analytic in its smooth
strengths on the positive-precision domain.
The uniform word and logarithm majorants
justify their sums and all ordered
derivatives. The strict smallness margins
and the precision reserve permit a slightly
larger bidisc. Normal convergence there
proves the asserted analyticity.
All sums have been rooted before any
physical volume sum, so no factor
\(\#\Lambda\) occurs.
\end{proof}

\subsection{Uniformity along the fixed family and the next task}

The proof adds only finitely many fixed-profile
constants and entry inequalities: local chart
smallness, the complex charge bound, and the
component and logarithm smallness conditions.
All involve \(B,P,\Lambda\), the actual coupling
band, and the existing path and inverse prices.
Their thresholds can therefore be included
in the slow-clock choice of
Section~\ref{sec:all-scale-calibration}.
The comparison periods satisfy the reserve
of Lemma~\ref{lem:finite-period-reserve}.
No threshold is chosen after a volume
subsequence, and no new Wilson family is
selected at an intermediate step.

Propositions~\ref{prop:central-two-coefficients}
and~\ref{prop:central-inherited-comparison}
now apply to these central letters with the
explicit prices \eqref{eq:central-explicit-prices}.
They give the cubic ordinary remainder and
the \(O(bd+d^2)\) incoming comparison
remainder in the original derivative norm.
The first incoming insertion is still the
exact masked average
\eqref{eq:central-principal-incoming-average}.
This section does not replace it by an inner
protector without a warning comparison.

The sharp transport of that principal insertion,
with its constant chosen before \(L\), is proved
in Section~\ref{sec:incoming-sharp-transport}.
The classical and one-loop matching of the
first two ordinary coefficients is proved in
Section~\ref{sec:ordinary-matching}.
The exact localisation identity
of Section~\ref{sec:finite-period-localisation}
then transfers the resulting auxiliary
action estimate. This matching is not supplied
by a rough central majorant, even though the
central derivative and precision inputs of
Section 32 have now been proved.
\clearpage
\section{Sharp transport of a single incoming regular activity}
\label{sec:incoming-sharp-transport}

The comparison remainder in Section 32 is small in the
coupling. Its leading incoming insertion needs a different
estimate: its constant must be fixed before the block ratio.
We prove this estimate on the same selected Gaussian marginal,
retaining the original function and keeping a separate
remainder for the auxiliary inner cutoff.
The inherited-term separation is described in
\cite[Lemma~13.5, PDF p.~95]{ArmstrongDraft}.
Here the additional inverse power of the ratio is obtained
by an explicit local gauge calculation and the existing
affine-Hessian condition.

\subsection{The two norms and the order of choices}

Write \(q,s_{\rm c}\) for the outgoing derivative order
and direction power, and \(q_{\rm f},s_{\rm f}\) for
the incoming ones. We use the profile choices
\[
 q_{\rm f}\ge\max(q,3),\qquad
 s_{\rm f}\ge s_{\rm c}\ge2 .
\]
In particular the prescribed orders at least six suffice.
Equality \(q_{\rm f}=q\ge3\) is allowed: this estimate
does not consume an extra derivative at every RG step.
These are choices of the construction's regularity
parameters, not assumptions on the Wilson measure.

For an incoming label \(i\), put
\[
 e_i=\|E_i\|_{q_{\rm f},h,s_{\rm f};o_i,X_i},\qquad
 d_i=n_{B_{\rm f}}(X_i),\qquad
 \bar o_i=\operatorname{parent}(o_i),
 \qquad B_{\rm f}\le2B .
\]
Its flat value is zero and its exponential-chart Hessian
vanishes on two affine fields whenever \(X_i\) is
nonwrapping. All incoming functions remain the original
gauge invariant functions of the physical fine connection.

Let \(\kappa_{\rm f},\kappa_{\rm c}>0\) be the incoming
and outgoing covering exponents, and set
\(\delta=\|E\|_{\mathcal E,\kappa_{\rm f}}\).
The outgoing norm below is the raw norm of Section 32,
anchored at \(\bar o_i\); flat-value and Hessian
subtraction are not part of the definition of the
principal average.

Fix the profiles, these orders and exponents, and the
kinetic and penalty data chosen before the ratio.
Constants \(C_j,{\mathsf M}_j\) are selected from these data
before choosing \(L\).
After \(L\) is fixed, the entry coupling and the
inner Gaussian width may depend on it.
The bounds are uniform in ultraviolet cutoff,
physical period, sector, and finite history at the
fixed profiles. Later threshold choices do not
alter the kernels or the Wilson family.

\subsection{Compression of a connected cover}

Use the fixed minimizing incoming cover \(D_i\)
of Section~\ref{sec:physical-reassembly}:
its radius-\(B_{\rm f}\) balls have connected union
and cover \(X_i\), and \(|D_i|=d_i\).
Their centres need not belong to \(X_i\).
Let \(Y_i=\widehat X_i\) be its already declared
full coarse read footprint. It is contained in
\[
       N_{rB}(\operatorname{parent}(D_i)),
\]
for a fixed integer \(r\). The declared footprint
includes the original tests, mean and covariance
windows, paths, tags, and the anchor \(\bar o_i\).
The same footprint and coarse occurrence mask
are used throughout this section.
The neighbourhoods are sets on the actual coarse
torus; they need not embed when the incoming cover
is large.

\begin{lemma}[Ratio gain in the covering number]
\label{lem:incoming-cover-compression}
There are \(K_r,C_r<\infty\), independent of \(B,L\)
and the periods, such that for \(B,L\) above fixed
geometric lower bounds,
\begin{equation}
 n_B(Y_i)\le K_r\left(1+C_r\frac{d_i}{L}\right).
 \label{eq:incoming-cover-compression}
\end{equation}
Thus, after choosing
\(L\ge2C_rK_r\kappa_{\rm c}/\kappa_{\rm f}\),
\begin{equation}
 \kappa_{\rm c}n_B(Y_i)
       \le\kappa_{\rm c}K_r+\tfrac12\kappa_{\rm f}d_i .
 \label{eq:incoming-cover-exponent}
\end{equation}
The same conclusion holds at any fixed larger
outgoing exponent, with its corresponding ratio bound.
\end{lemma}
\begin{proof}
Join two centres when their radius-\(B_{\rm f}\)
balls overlap or are adjacent. This graph is
connected by the definition of \(n_{B_{\rm f}}\).
An edge has periodic maximum length at most
\(2B_{\rm f}+1\). Traverse a spanning tree in both
directions. Choosing a shortest lift for each
edge gives a continuous walk of length at most
\(C B_{\rm f}d_i\), measured in the sum of the
four coordinate lengths.

Divide the walk's coordinates by \(L\).
The parent of a fine point differs from this
scaled point by at most a fixed lattice distance.
Every projected \(B_{\rm f}\)-ball lies within
distance \(B_{\rm f}/L+C\) of the scaled walk.
Choose successive coarse centres at walk-length
increments \(B/4\), and round to coarse lattice
sites. There are at most \(C(1+d_i/L)\) centres.
For large enough fixed lower bounds on \(B,L\),
their radius-\(B\) balls cover the projected
incoming cover and have connected union.

An \(rB\)-thickening of each such ball is covered
by \(K_r\) radius-\(B\) balls on a fixed grid of
spacing at most \(B/2\). The grid and the walk
cover can be joined within the same fixed
multiple of this number. This proves
\eqref{eq:incoming-cover-compression}, after
enlarging \(K_r,C_r\).
The walk may wrap or visit a site repeatedly.
Projection to the torus identifies balls and
does not increase their number. Compatibility
of the periods, \(f=Lc\), makes the parent map
consistent with these lifts.
Multiplying the bound by \(\kappa_{\rm c}\)
gives \eqref{eq:incoming-cover-exponent}.
\end{proof}

In particular no large-cover conclusion is
drawn from a large diameter alone. We shall
split directly according to \(d_i^3<L\) or
\(d_i^3\ge L\).

\subsection{The three scales in a lifted direction}

Use \(E\) for the smooth free interpolation column
constructed in Section~\ref{sec:all-scale-linear}.
Equation \eqref{eq:all-scale-linear-difference-bounds}
gives, through two fine differences,
\[
 |\nabla_f^r E|_\sigma\le C_jL^{-1-r},
                      \qquad 0\le r\le2,
\]
with \(C_j,\sigma>0\) chosen before \(L\).
These bounds already include fine residue changes
across a block boundary.

For a fine field restricted to the physical reads
of \(E_i\), use the dimensionless size
\[
 |u|_i=\sup_e
       \frac{|u_e|}{h(1+d_f(e,o_i))^{s_{\rm f}}}.
 \]
This is only notation for the direction bound
already in the incoming norm. It does not add an
induction space. On a nonwrapping patch write
\(\xi(x)=x-o_i\) in its local integer lift.

\begin{lemma}[Constant, affine, and remaining parts]
\label{lem:incoming-three-parts}
Suppose
\[
 |v(y)|\le C h(1+d_c(y,\bar o_i))^{s_{\rm c}}.
\]
On a nonwrapping fine read patch the field \(Ev\)
has a decomposition \(c+b+\rho\), where \(c\)
is constant and \(b\) is affine, such that
\begin{equation}
 |c_\mu|\le C_jhL^{-1},\qquad
 |b_{\mu\nu}|\le C_jhL^{-2},\qquad
 |\rho|_i\le C_jL^{-3}.
 \label{eq:incoming-three-parts}
\end{equation}
Here \(b_\mu(x)=\sum_\nu\xi_\nu(x)b_{\mu\nu}\).
The constants are independent of the read radius,
\(B,L\), history and period.
\end{lemma}
\begin{proof}
Absolute kernel summation with a polynomial
weight gives
\[
 |\nabla_f^r Ev(x)|
 \le C_jhL^{-1-r}
                  (1+d_f(x,o_i)/L)^{s_{\rm c}} .
\]
Indeed the ratio of coarse polynomial weights
at two sites is bounded by a fixed power of
their separation, and every such moment is
summable against \(e^{-\sigma|y-y'|}\).
The fine anchor may be any residue in its
coarse block; this changes the weight by a
fixed factor.

Let \(c_\mu=(Ev)_\mu(o_i)\), and let
\(b_{\mu\nu}=\nabla_{f,\nu}(Ev)_\mu(o_i)\).
For \(r=d_f(x,o_i)\le L\), telescoping first
differences along coordinate paths gives
\[
 |(Ev)_\mu(x)-c_\mu-b_\mu(x)|
                         \le C_jhL^{-3}(1+r)^2 .
\]
This is the discrete Taylor formula; its error
is a sum of at most \(C(1+r)^2\) second
differences. For \(r\ge L\), use instead the
value bound and the two subtracted terms:
\[
 |\rho_\mu(x)|\le C_jh
 \{L^{-1}(1+r/L)^{s_{\rm c}}
                       +L^{-1}+L^{-2}r\}.
\]
Divide by \((1+r)^{s_{\rm f}}\).
Since \(s_{\rm f}\ge s_{\rm c}\ge2\), every
term is bounded by \(C_jhL^{-3}\).
The same division handles the near region.
This proves \eqref{eq:incoming-three-parts}
without a factor depending on the read radius.
\end{proof}

\subsection{A local gauge calculation giving the fifth power}

The constant part above is of order \(L^{-1}\);
estimating it as an arbitrary fine direction
would lose the required contraction.
It can instead be removed by a local gauge
transformation. The first correction it creates
is an affine field and is covered by the
incoming Hessian normalisation.

\begin{lemma}[Fifth-order gain from the affine Hessian]
\label{lem:incoming-fifth-order-gain}
Let \(f=E_i\), with \(X_i\) nonwrapping.
In local exponential coordinates suppose a
parameter-dependent fine connection has logarithm
\[
             a_\mu(x)=c_\mu+b_\mu(x)+\rho_\mu(x).
\]
For a number \(0<\tau\le1\), assume that the
coefficients of \(c\), the affine coefficients
of \(b\), and the weighted size of \(\rho\),
as well as their parameter derivatives through
order \(q\), satisfy respectively
\[
 C h\tau,\qquad C h\tau^2,\qquad C\tau^3.
\]
Assume the connection has plaquette chord at
most \(\eta h\) on the tested fine plaquettes,
where \(\eta>0\) is a sufficiently small fixed
group-dependent number. If \(R_i\) bounds the
fine read radius, require
\[
 h(1+R_i)^{2s_{\rm f}+2q+10}
                      \le\varepsilon(C,q,G).
\]
Then every parameter derivative through order
\(q\) of \(f(\exp a)\) is bounded by
\begin{equation}
                         C_{j,q}e_i\tau^5.
 \label{eq:incoming-fifth-order-gain}
\end{equation}
The constant is independent of the read radius.
Only derivatives of \(f\) through
\(\max(q,3)\) are used.
\end{lemma}
\begin{proof}
Put \(S_x=\sum_\nu\xi_\nu(x)c_\nu\).
Use the site gauge \(g_x=e^{S_x}\) on the
read patch and extend it to the torus with
a fixed spatial cutoff outside a larger
nonwrapping patch. On every read edge the
transformed link is
\[
 e^{S_x}e^{a_\mu(x)}e^{-S_{x+e_\mu}}.
\]
The smallness condition puts these factors
in a fixed logarithmic chart. The
Baker--Campbell--Hausdorff formula, namely the
Taylor expansion of a logarithm of a product
of exponentials, through
degree two in \(c,b,\rho\), gives its logarithm
in the form
\begin{equation}
 \widetilde a_\mu(x)
    =b_\mu(x)+\tfrac12[S_x,c_\mu]
                         +\widetilde\rho_\mu(x),
 \qquad
 |\widetilde\rho|_i\le C_{j,q}\tau^3 .
 \label{eq:incoming-constant-gauge}
\end{equation}
The same bounds hold for its parameter jets.
The first two displayed terms form an affine
field, of weighted size \(C_{j,q}\tau^2\).

Here are the remainder bounds in this chart
calculation. The first omitted constant-only
commutators have size
\(C h^3\tau^3(1+|\xi|)^2\).
The first mixed terms with \(b\) have size
\(C h^2\tau^3(1+|\xi|)^2\).
After division by
\(h(1+|\xi|)^{s_{\rm f}}\) these are bounded
by \(C\tau^3\).
Terms containing \(\rho\) have its weighted
bound, multiplied by bounded analytic
coefficients. Higher constant commutators
are summed using \(\sup_x|S_x|\ll1\).
The fixed-order product and chain rules give
the same bounds after parameter derivatives.
Thus this calculation uses a convergent
small-chart formula, with its full remainder.

All radial segments \(t\widetilde a\),
\(0\le t\le1\), are fine healthy.
To verify this rather than assume convexity
of the good-field set, expand the four
factors of each plaquette:
\[
 P(\exp u)=1+(d_1u)_{\square}
                              +O(\|u\|_\infty^2).
\]
Gauge invariance preserves the endpoint
plaquette bound \(\eta h\), so this identity
bounds \(d_1\widetilde a\) by
\(C_G(\eta h+\|\widetilde a\|_\infty^2)\).
Apply the same identity to \(t\widetilde a\).
The stated smallness condition makes
\(\|\widetilde a\|_\infty^2/h\) as small as
needed; choose \(\eta\) first and then
\(\varepsilon(C,q,G)\).
Every such plaquette consequently has chord
less than \(h/2\).
Only plaquettes queried by the incoming mask
are needed. Their four edges are included
in the declared physical read set.

Write \(F(u)=f((e^{u_e})_e)\).
On these segments its ordinary derivatives
through \(\max(q,3)\) satisfy
\[
 |D^mF(u)[v_1,\ldots,v_m]|
                  \le C_q e_i\prod_\nu|v_\nu|_i .
\]
Indeed the differential of the exponential
converts a chart direction to a left
direction by an integral of adjoint
isometries. Higher conversions multiply
pointwise directions. Their extra factors
are bounded by powers of
\(h(1+R_i)^{s_{\rm f}}\le1\).
The number of ordered chain-rule terms
depends only on the derivative order;
there is no edge-count factor.

By Lemma~\ref{lem:gauge-hessian-identities},
\(DF(0)=0\), and \(F(0)=0\) by normalisation.
Put \(H_0=D^2F(0)\).
Its value on two affine fields is zero by
the incoming affine-Hessian condition.
Write \(\widetilde a=a_2+a_3\), with \(a_2\)
affine of size \(O(\tau^2)\) and \(a_3\)
of size \(O(\tau^3)\); the same is true of
their parameter jets.
For value, first derivative, and second
derivative use the formulas
\begin{align*}
 F(u)&=\tfrac12H_0(u,u)
       +\tfrac12\int_0^1(1-t)^2
                      D^3F(tu)[u,u,u]\,dt,\\
 DF(u)v&=H_0(u,v)
       +\int_0^1(1-t)D^3F(tu)[u,u,v]\,dt,\\
 D^2F(u)[v,w]&=H_0(v,w)
       +\int_0^1D^3F(tu)[u,v,w]\,dt .
\end{align*}
In each Hessian term the pair of affine
parts cancels. Every remaining pair has
size at most \(C\tau^5\); every cubic
term has size at most \(C\tau^6\).
An ordered parameter derivative of
\(F(\widetilde a)\) is a finite sum of
\(D^mF(\widetilde a)\) applied to parameter
jets of \(\widetilde a\).
For \(m=1,2\) apply the last two formulas;
for \(m\ge3\) use directly the bound
\(C_q e_i\tau^{2m}\).
This proves \eqref{eq:incoming-fifth-order-gain}.
We have not differentiated a Taylor integral
through \(q\) more orders, which would
unnecessarily require higher regularity.
Finally use the exact gauge invariance
\(f(\exp a)=f(\exp\widetilde a)\).
\end{proof}

\subsection{The actual central marginal and an inner protector}

We need a rough estimate for arbitrary supports, as well
as the sharper small-cover estimate. It is useful to
record why its dependence on \(L\) can be kept polynomial.
An unspecified constant chosen after \(L\) would not
suffice in the large-cover argument.

\begin{lemma}[Polynomial ratio bounds for central records]
\label{lem:incoming-polynomial-ratio-records}
At the fixed profiles there are \(C_j,{\mathsf M}_j\), chosen
before \(L\), with the following properties above
a ratio-dependent entry coupling. Central operated
matrices, their selected inverse columns, means,
covariances, and marginal inverses have polynomial
spatial moments bounded by \(C_jL^{{\mathsf M}_j}\).
For fixed-order coarse derivatives each insertion
supplies its direction factor \(h_{\rm c}\), with
the prescribed spatial weights.
The central mean has value bound \(C_jL^{{\mathsf M}_j}h\).
These bounds have no factor growing with \(B\).
Taking a trace or counting reads is a separate
operation and may introduce powers of \(B\).
\end{lemma}
\begin{proof}
Here only central records are used: their full reads
meet no frozen link or physical seed. Their source
therefore consists of the regular replica rows and
coarse guides. In local gauges these rows are
bounded by \(C_jL^m h\).
The same finite path formulas give positive-order
derivatives bounded by \(C_jL^m h_{\rm c}^r\),
times their spatial direction weights.
There is no derivative of a threshold decision.
All coarse threshold decisions have their unmasked
regular extension in this germ.

Every block tree has at most \(L^4\) vertices.
Its route lengths, route comparisons, and numbers
of elementary plaquette interchanges are bounded
by fixed powers of \(L\).
The block axial estimate used for
Lemma~\ref{lem:operated-precision-gap} consequently
has constants at most \(C_jL^m\).
The fixed positive row metrics and noise penalty
give a lower precision bound \(c_jL^{-m}\).
The actual regular rows have the same bound after
the local curvature perturbation is made smaller
than half this gap by the entry choice.
Inactive coordinates retain their identity rows.
This reasoning is blockwise and counts no
physical volume.

The kernels and their fixed-order derivatives
have coefficient bounds of the form
\[
                 C_jL^m e^{-\sigma_j d_c(a,b)},
\]
with
\(\sigma_j>0\) fixed before the ratio.
The native cut only removes tails from their
absolute budgets.
Conjugation with
\(e^{\theta d_c(\cdot,a)}\), where
\(\theta=c_jL^{-m'}\) is sufficiently small,
perturbs the matrix by less than half its gap.
The resolvent identity therefore bounds its
inverse entries by
\[
        C_jL^{m'}e^{-c_jL^{-m'}d_c(a,b)}.
\]
There are at most \(C_GL^4\) coordinates above
one coarse site. Summing any fixed polynomial
moment of this exponential costs another
fixed power of \(L\). The same argument applies
to principal windows and positive marginal
covariances.
The selected covariance differs from the full
inverse by an exponentially small window error;
choose its entry threshold so it has half the
same positive spectral margins.

Repeated inverse differentiation inserts at
most the fixed number of matrix derivatives.
Their direction weights are moved along the
exponentially decaying kernels using the
triangle inequality. Fixed polynomial moments
again give powers of \(L\).
This also handles the fixed-order derivatives
of the regular frame and predictor: their
finite-dimensional implicit equations have
block Hessian gaps bounded below by an inverse
power of \(L\), and the route derivatives have
the path bounds just counted.
Finally \(\mu=-J^Tb\) gives its value bound
from the central source \(b=O(C_jL^m h)\).
This proves the assertion, with a common larger
integer \({\mathsf M}_j\), independently of window size,
history length and period.
Enlarge that integer to cover the finite
replica, tree-slice and exponential coordinate
changes through the same prescribed order.
Their route formulas and fixed-order
products have the polynomial ratio bounds
just established.
\end{proof}

Include the four edges of every tested
plaquette in the physical reads of
\(\chi_{X_i}E_i\). Let \(T_i\) be their
active free-link coordinates.
Tree links retain their installed values
and receive no new Gaussian variable.
Integrate out all other Gaussian coordinates first.
For a fixed discrete sector, forest and shared
tag, let
\[
 C_i=(C_{\rm sel})_{T_i,T_i}
\]
be this original marginal covariance.
Its mean is included in the installed physical
field \(\mathcal U_\sigma(V,\mu+gz)\).
The number of real coordinates is at most
\(C_j(1+B)^4(1+d_i)\).
The mean and covariance records read only \(Y_i\).
Using a marginal rather than the full Gaussian
is essential: differentiation must count the
reads of this insertion, not the whole torus.

Choose a fixed smooth function
\(\theta:[0,\infty)\to[0,1]\), equal to one on
\([0,1/4]\) and zero on \([1,\infty)\).
With \({\mathsf M}_j\) from the preceding lemma choose
\[
 \eta_L=\eta_jL^{-{\mathsf M}_j-6},\qquad
 H_L=\eta_Lh/g,\qquad
 \Gamma_i(z)=\prod_{e\in T_i}
                    \theta(|z_e|^2/H_L^2).
\]
Each \(z_e\) is a Lie-algebra valued coordinate.
The radial factors are invariant under orthogonal
colour transport. For fixed discrete reads
\(\Gamma_i\) is independent of \(V\).
This is an auxiliary protector; the physical
cutoffs have not been narrowed.

Choose \(L\) large and then its entry coupling
so that on \(\Gamma_i\ne0\) the fine plaquettes
queried by \(X_i\) have chord less than
\(\eta h\), and all original broad factors on
these reads equal one.
Here is the quantitative verification.
The calibrated fine row is bounded by
\(C_jL^{-2}h_{\rm c}\).
The selected centre differs from its stationary
completion by \(C_{j,L}(1+B)^m
(h_{\rm c}^2+e^{-c_{j,L}B})\), as proved in
Section~\ref{sec:central-derivative-estimates}.
The added fine fluctuation is at most
\(C_jL^{{\mathsf M}_j}\eta_Lh\).
First make \(C_jL^{-2}h_{\rm c}/h\) and
\(C_jL^{{\mathsf M}_j}\eta_L\) smaller than \(\eta/4\);
then make the centre error smaller than
\(\eta h/4\).
This is a local argument at each queried
plaquette and also applies to wrapping supports.
The centre's central value bound, together
with \(B^{17}\to\infty\), puts these coordinates
in the region where the original broad
factors equal one.

Define the principal average
\begin{equation}
 P_i(V)=\operatorname{Av}
       \int\Gamma_i(z)
           E_i(\mathcal U_\sigma(V,\mu+gz))\,d\gamma_{C_i}(z).
 \label{eq:sharp-principal-incoming-average}
\end{equation}
The symbol \(\operatorname{Av}\) denotes the
original normalized shared tag and forest
average. It is the same average as in the
single insertion in Section 32.
No independent copy is assigned at overlapping
reads. The marginal is normalized with its
actual determinant, and all Haar factors in
the physical coordinate identity keep their
original normalisation.

\subsection{Derivative bounds for the principal average}

For this principal term we can use the incoming
function's regular derivatives, since the
protector puts its arguments strictly inside
its healthy domain. Keep \(z\) fixed and leave
the mean inside the installed field.
Only covariance scores arise from differentiating
the marginal density:
\begin{equation}
 D_A\log\varphi_{C_i}(z)
 =-\tfrac12\operatorname{Tr}(C_i^{-1}D_AC_i)
   +\tfrac12 z^TC_i^{-1}(D_AC_i)C_i^{-1}z .
 \label{eq:incoming-centred-covariance-score}
\end{equation}
Here \(\varphi_{C_i}\) is its Lebesgue density.
There is no mean divided by \(g\) in this score.
The mean is differentiated as part of the
fine field, in its incoming scaled norm.

\begin{lemma}[Two bounds for a principal insertion]
\label{lem:incoming-principal-address}
There are \(C_j,m_j\), chosen before \(L\),
such that above its entry threshold
\begin{align}
 \|P_i\|_{q,h_{\rm c},s_{\rm c};\bar o_i,Y_i}
       &\le C_jL^{-5}e_i,
                                      &&d_i^3<L,\notag\\
 \|P_i\|_{q,h_{\rm c},s_{\rm c};\bar o_i,Y_i}
       &\le C_jL^{m_j}(1+d_i)^{m_j}e_i,
                                      &&d_i^3\ge L .
 \label{eq:incoming-principal-two-cases}
\end{align}
Each average has a \(C^q\) germ on its healthy
coarse domain.
\end{lemma}
\begin{proof}
First consider \(d_i^3<L\).
A connected cover bounds the fine read radius
by \(C(1+B)d_i\). Its coarse projection has
radius at most \(CBd_i/L+C\).
For sufficiently large \(L\) it lies in a
fixed fraction of a radius-\(B\) box about
\(\bar o_i\). The full healthy read set \(Y_i\)
therefore contains the coarse box needed for
the local gauges and selected-window comparisons.
The period reserve also makes the fine support
nonwrapping: its radius is at most
\(C(1+B)L^{1/3}\), whereas the fine side is
\(L\) times the coarse side.

Choose a local coarse gauge at the base field
and keep that gauge fixed when differentiating.
The axial path formula bounds its logarithms
by \(C h_{\rm c}(1+d_c(y,\bar o_i))\).
The first variation is bounded by the prescribed
coarse direction weight. Higher variations
of the logarithmic chart have at least two
direction factors; on this finite box their
extra normalized size is
\(C_{j,L}(1+B)^m h\), which tends to zero.
In particular differentiating the gauge itself
does not add an unwanted spatial weight.

We spell out the approximation to the free
column used here. In these local coordinates
the stationary graph has first derivative
\(\mathcal H\), and
\(\rho-\mathcal Hx=O(C_{j,L}(1+B)^m h^2)\)
in scaled jets.
The selected mean differs from \(\rho\) by
the bound already proved above.
Equation \eqref{eq:actual-free-coordinate-lift}
identifies the replica plus this first jet
with \(Ex\), after a fine gauge change.
The tree-slice and frame corrections are
included in that identity.
Use this fine gauge change on the local
patch inside the evaluation of the gauge
invariant function \(E_i\).
The selected Gaussian marginal and its
normalising determinant remain unchanged.
Expanding the coordinate multiplication and
logarithm adds only second-order chart errors.
The protected native tails remain exponentially
small. Consequently the fine logarithm, and
each of its ordered coarse jets through
order \(q\), equal the corresponding free
interpolation jet, a fluctuation of size at
most \(C_jL^{{\mathsf M}_j}\eta_L\) in \(|\cdot|_i\),
and an error bounded by
\begin{equation}
 C_{j,L}(1+B)^m(1+d_i)^m
                    \{h+e^{-c_{j,L}B}/h\}.
 \label{eq:incoming-linear-approximation-error}
\end{equation}
This follows from Taylor's integral formula
for the graph, inverse difference identities
for its selected windows, and the finite
product rule for the coordinate changes.
Every nonlinear term has at least two
small coordinate or direction factors;
division by the incoming scale \(h\)
leaves the factor \(h\) displayed here.
On the protector the fluctuation is \(gz\),
not a covariance square root applied to
an unrelated Gaussian field.

Lemma~\ref{lem:incoming-three-parts} decomposes
the free jets into constant, affine, and
remaining parts, with bounds of orders
\(L^{-1}\), \(L^{-2}\), and \(L^{-3}\).
The fluctuation fits in the last part by
the choice of \(\eta_L\).
Since \(d_i<L^{1/3}\), the error in
\eqref{eq:incoming-linear-approximation-error}
is smaller than \(L^{-3}\) after raising
the entry threshold.
The smallness condition in
Lemma~\ref{lem:incoming-fifth-order-gain}
also holds uniformly in this class, because
the read radius is bounded by
\(C(1+B)L^{1/3}\).
That lemma, with \(\tau=L^{-1}\), now bounds
every derivative of the incoming composition
on the support of \(\Gamma_i\) by
\(C_jL^{-5}e_i\).

It remains to bound derivatives of the
Gaussian density, without losing this
pre-ratio constant. By
\eqref{eq:incoming-centred-covariance-score},
Lemma~\ref{lem:incoming-polynomial-ratio-records},
and the finite ordered partition formula,
every positive-order density score on
\(\Gamma_i\ne0\) is at most
\begin{equation}
 C_jL^{m_j}h(1+B)^{m_j}(1+P)^{m_j}
                                      (1+d_i)^{m_j}.
 \label{eq:incoming-small-density-scores}
\end{equation}
Indeed \(|z_e|\le H_L\le CP\), the number
of coordinates is \(C_j(1+B)^4(1+d_i)\),
and every covariance derivative supplies
at least one coarse direction factor
\(h_{\rm c}\le2h\).
For \(d_i<L^{1/3}\) the displayed quantity
is at most one after the entry choice:
\(h\) dominates every fixed power of
\(B,P\).
The zeroth-order density has integral one,
and \(\Gamma_i\) is independent of \(V\)
at fixed \(z\).
The product rule therefore preserves the
bound \(C_jL^{-5}e_i\), with a fixed
order-dependent constant, proving the
first line.

For \(d_i^3\ge L\), no affine condition or
global fine chart is used.
On the protector all queried fine plaquettes
are still healthy. Differentiate
\(E_i(\mathcal U_\sigma(V,\mu+gz))\)
using its incoming regular derivatives.
The first pulled-back left direction has
incoming size at most \(C_jL^{m_j}\):
the parent distance is at most a fixed
multiple of \(1+d_f/L\), and
\(s_{\rm f}\ge s_{\rm c}\).
Higher direction coefficients cost a fixed
power of the read diameter and extra
factors of \(h\). Since that diameter
is at most \(C(1+B)(1+d_i)\), impose
\(h(1+B)^m(1+P)^m\le1\) at the entry
and retain the remaining polynomial
in \(1+d_i\).
For example, the direction-weight part
of an \(r\)-fold coefficient, divided
by one allowed fine direction, is bounded by
\[
 C_r\frac{h_{\rm c}^r(1+\rho/L)^{rs_{\rm c}}}
                  {h(1+\rho)^{s_{\rm f}}}
 \le C_r h^{r-1}(1+\rho)^{(r-1)s_{\rm c}},
 \qquad \rho=d_f(e,o_i)\le C(1+B)(1+d_i).
\]
For \(r\ge2\) its power of \(B\) is paired
with a positive power of \(h\); for \(r=1\)
there is no read-radius loss.
Lemma~\ref{lem:incoming-polynomial-ratio-records}
then bounds all these coefficients by
\(C_jL^{m_j}(1+d_i)^{m_j}\).
The same argument applies to the density
scores in \eqref{eq:incoming-small-density-scores}.
This proves the second line by the ordered
chain and product rules. Wrapping of the
support creates no extra assumption in
this argument.

Finally, the fine healthy margins on the
compact support of \(\Gamma_i\), together
with the positive marginal covariance,
persist on an open coarse neighbourhood.
Differentiation under this finite marginal
integral gives the claimed germs, also
at equality in a coarse good-field test.
No incoming derivative beyond
\(q_{\rm f}\ge\max(q,3)\) has been used.
\end{proof}

\subsection{Restore the original masked insertion exactly}

Let \(A_i\) denote the exact single incoming
average of Section 32:
\[
 A_i(V)=\operatorname{Av}\int
       \left(\prod_{e\in T_i}b_e\right)
          \chi_{X_i}(U)E_i(U)\,d\gamma_{C_i}(z).
\]
Work on the healthy coarse domain of \(Y_i\).
The following identities are multiplied by
their original coarse occurrence masks in
the physical reassembly.
On the support of \(\Gamma_i\), all these
broad factors and the original fine mask
equal one. Thus \(A_i=P_i+R_i^{\rm cut}\)
with the exact formula
\begin{equation}
 R_i^{\rm cut}(V)=\operatorname{Av}\int
       \left(\prod_{e\in T_i}b_e\right)
          (1-\Gamma_i)\chi_{X_i}(U)E_i(U)
                                           \,d\gamma_{C_i}(z).
 \label{eq:incoming-exact-cutoff-comparison}
\end{equation}
This identity keeps the original arguments
of \(E_i\) and the original fine mask.

\begin{lemma}[The cutoff comparison is superpolynomially small]
\label{lem:incoming-rare-cutoff-comparison}
For every fixed \(J>0\), above a ratio-dependent
entry threshold,
\begin{equation}
 \|R_i^{\rm cut}\|_{q,h_{\rm c},s_{\rm c};\bar o_i,Y_i}
 \le C_{j,L,J}\widehat\beta_{\rm f}^{-J}
                                  (1+d_i)^{m_j}e_i .
 \label{eq:incoming-rare-cutoff-comparison}
\end{equation}
The exponent \(m_j\) can be chosen independently
of \(J,L,B\). Period and cutoff do not enter
the constants.
\end{lemma}
\begin{proof}
Change all Borel physical reads to physical
Haar coordinates using the exact density of
Section~\ref{sec:integrated-regular-derivatives}.
The multiplier \(\chi_{X_i}(U)E_i(U)\) is
then fixed and bounded by \(e_i\).
Derivatives fall on the density, the retained
broad factors, and the smooth protector.
Their score bounds give a fixed power of
\(1+B,1+P,1+d_i\), with a ratio-dependent
constant. Derivatives of the protector
cost powers of \(h_{\rm c}/(gH_L)\);
this is at most \(C/\eta_L\) and is allowed
in that ratio-dependent constant.

Every derivative term retains the event
that some \(|z_e|\ge H_L/2\).
For a differentiated protector this follows
from its derivative support; for the other
terms it follows from \(1-\Gamma_i\ne0\).
The broad factors bound the remaining
coordinates by \(C_{j,L}(1+B)^mP\), so
the fixed-order score polynomials can be
bounded on this support.
Each Gaussian coordinate has variance at
most \(C_jL^{{\mathsf M}_j}\). A one-coordinate
Gaussian tail and a union bound therefore
give
\[
 \Pr\{\exists e\in T_i:|z_e|\ge H_L/2\}
       \le C_j(1+B)^4(1+d_i)
                              e^{-c_{j,L}P^2}.
\]
The Lie colour dimension changes only the
fixed constant. No independence of the
coordinates is used.
Combining these bounds gives a fixed
polynomial times \(e^{-c_{j,L}P^2}e_i\).
Since \(P^2/\log\widehat\beta_{\rm f}\)
dominates every fixed power of \(B\),
this implies
\eqref{eq:incoming-rare-cutoff-comparison}.
The physical change of variables kept
its exact determinant and Haar density;
no coordinate-count normalisation was
estimated by an extensive constant.
\end{proof}

\subsection{The anchored contraction and the inherited remainder}

\begin{theorem}[Sharp principal incoming transport]
\label{thm:sharp-incoming-transport}
There are \(V_j<\infty\) and a lower ratio bound,
both chosen before \(L\), such that for every
allowed larger odd ratio and then every
sufficiently large entry coupling,
\begin{equation}
       \|P\|_{\mathrm{raw},\mathcal E,\kappa_{\rm c}}
                       \le\frac{V_j}{L}\,
                                 \|E\|_{\mathcal E,\kappa_{\rm f}} .
 \label{eq:sharp-incoming-transport}
\end{equation}
At the same addresses
\begin{equation}
 \|R^{\rm cut}\|_{\mathrm{raw},\mathcal E,\kappa_{\rm c}}
       \le C_{j,L,J}\widehat\beta_{\rm f}^{-J}\delta
                    \quad\hbox{for every fixed }J>0 .
 \label{eq:sharp-incoming-cutoff-family}
\end{equation}
The averages are local, gauge invariant and
covariant as families under the lattice
symmetries. The estimates are uniform in
cutoff, period and finite history.
\end{theorem}
\begin{proof}
Let \(m_j\) be the larger exponent in
Lemma~\ref{lem:incoming-principal-address}.
If \(d_i^3\ge L\), then
\(L^{m_j}\le d_i^{3m_j}\) and
\(L^5\le d_i^{15}\).
Together with the small-cover bound this gives,
for a fixed integer \(p_j\),
\[
 \|P_i\|_{q,h_{\rm c},s_{\rm c};\bar o_i,Y_i}
       \le C_jL^{-5}(1+d_i)^{p_j}e_i
\]
for every incoming support.
Use \eqref{eq:incoming-cover-exponent}.
The number
\[
 C_j e^{\kappa_{\rm c}K_r}
      \sup_{d\ge1}(1+d)^{p_j}
                                      e^{-\kappa_{\rm f}d/2}
\]
is finite and fixed before \(L\).
It bounds the coefficient of
\(L^{-5}e^{\kappa_{\rm f}d_i}e_i\)
in the outgoing weighted price.

At a specified coarse anchor there are
exactly \(L^4\) fine site anchors above it.
Sum the incoming anchored norm at each of
these sites. The factor \(L^4\) multiplies
the per-address \(L^{-5}\), giving
\eqref{eq:sharp-incoming-transport}.
Do not convert this principal bound to an
incidence sum first: its norm is the
anchored norm, and no factor counting all
sites of \(Y_i\) is needed.
The same exponential reserve sums the
polynomial in
Lemma~\ref{lem:incoming-rare-cutoff-comparison},
proving \eqref{eq:sharp-incoming-cutoff-family};
its \(L^4\) is part of the allowed
ratio-dependent constant.

Every integral reads only \(Y_i\).
Gauge transformations act orthogonally
on its colour coordinates, transform the
installed physical links by the original
gauge action, and preserve the radial
protector. Change these Gaussian
coordinates and use the gauge invariance
of \(E_i\). The shared tag and forest
average preserves the same identity.
Transforming these indexed data also
gives the lattice covariance of the
family. There is no extra tag-count
constant because their weights form
the original probability average.
\end{proof}

\begin{corollary}[Normalised incoming transport and its exact reset]
\label{cor:normalised-incoming-transport}
For the prescribed outgoing orders \(q\ge2\),
put
\[
 Y_i'=Y_i\cup B_\infty(\bar o_i,2),\qquad
 \bar P_i=P_i-P_i(1)-Q_{P_i},
\]
using the clover counterterm
\eqref{eq:local-quadratic-counterterm}.
The normalised family satisfies
\begin{equation}
       \|\bar P\|_{\mathcal E,\kappa_{\rm c}}
                   \le \frac{\widetilde V_j}{L}\delta ,
 \label{eq:normalised-incoming-contraction}
\end{equation}
where \(\widetilde V_j\) is chosen before \(L\).
It has zero flat value and zero affine Hessian
whenever its enlarged footprint is nonwrapping.
The subtraction has the exact masked form
\begin{equation}
 \chi_{Y_i}P_i
   =P_i(1)+\chi_{\{\bar o_i\}}Q_{P_i}
                 -u_i^{\rm reset}+\chi_{Y_i'}\bar P_i .
 \label{eq:incoming-normalisation-mask-identity}
\end{equation}
The Borel local function \(u_i^{\rm reset}\)
vanishes on a healthy \(Y_i'\) and has value
bound \(C_jL^{-5}(1+d_i)^{p_j}e_i\).
\end{corollary}
\begin{proof}
Lemma~\ref{lem:local-reset-norms} has constant
\(C_{\rm r}(G,q,s_{\rm c})\), independent of
the footprint, threshold, period and ratio.
It bounds the norm of \(\bar P_i\) by
\(C_{\rm r}\) times that of \(P_i\).
Since \(B\ge2\) and \(\bar o_i\in Y_i\),
\[
                  n_B(Y_i')\le n_B(Y_i)+1 .
\]
Consequently one may take
\(\widetilde V_j=C_{\rm r}e^{\kappa_{\rm c}}V_j\).
There is no new factor polynomial in \(B\)
in this operation.

For completeness define \(u_i^{\rm reset}=0\)
on \(\chi_{Y_i'}=1\), and otherwise put
\[
 u_i^{\rm reset}
 =P_i(1)+\chi_{\{\bar o_i\}}Q_{P_i}
                                      -\chi_{Y_i}P_i .
\]
Lemma~\ref{lem:exact-masked-subtraction}
proves
\eqref{eq:incoming-normalisation-mask-identity}
and its stated value bound. Use the already
constructed covariant local extensions for
the unmasked germs; the masked expression
depends only on their healthy values.
The scalar and quadratic terms are retained
in the scalar and coupling ledger.
Their coefficients obey
\eqref{eq:quadratic-extraction-price}, including
its explicit factor \(h_{\rm c}^{-2}\).
The normalisation conditions follow from
the equality of the clover and affine
Hessians.
\end{proof}

This compression also supplies the incidence
estimate needed for the incoming prices of
Section~\ref{sec:central-derivative-estimates}
without requiring successively larger fine
covering exponents.
For a fixed coarse site \(x\), an anchor of
an incoming footprint containing \(x\) must
lie within \(C_rB(1+d_i/L)\) in coarse units.
For \(d_i=d\) the number of possible fine
anchors is at most
\(C_rL^4(1+B)^4(1+d)^4\).
Use \eqref{eq:incoming-cover-exponent} at
the desired output reserve exponent and
sum \(e^{-\kappa_{\rm f}d/2}(1+d)^m\).
For any fixed output reserve exponent
\(\kappa_{\rm out}\), with the ratio chosen
as in Lemma~\ref{lem:incoming-cover-compression}
at that exponent, the resulting value incidence is
\begin{equation}
 \sup_x\sum_{i:x\in Y_i}
       e^{\kappa_{\rm out}n_B(Y_i)}e_i
               \le C_{j,L}(1+B)^4\delta .
 \label{eq:incoming-compressed-incidence}
\end{equation}
It has exactly the allowed type in
\eqref{eq:central-explicit-prices}.
Thus any fixed positive incoming exponent
is compatible with a fixed output exponent
and reserve once the ratio is sufficiently
large. Equal exponents at repeated profiles
are allowed.

Finally recall that \(b,d\) in Section 34
denote the strength incomes, not \(d_i\).
The exact identities proved above and
Proposition~\ref{prop:central-inherited-comparison}
give the warning-free decomposition
\begin{equation}
 F(1,1)-F(1,0)=-\sum_iP_i+R_{\rm inh},\qquad
 \|R_{\rm inh}\|_{\mathrm{raw},\kappa}
 \le C_B(bd+d^2)
               +C_{j,L,J}\widehat\beta_{\rm f}^{-J}\delta .
 \label{eq:sharp-incoming-full-comparison}
\end{equation}
In the scale regime of Section 34 this remainder
is \(\widehat\beta_{\rm f}^{-1/2+o(1)}\delta\).
The signs agree with the incoming exponential
\(e^{-\chi_{X_i}E_i}\). The cutoff discrepancy
has been retained in \(R_{\rm inh}\), not
discarded.

Choose \(V_j,\widetilde V_j\) and the
geometric lower ratio bound first, then
choose the ratio so both contraction
coefficients fit the prescribed incoming
fraction, and only then raise the entry
coupling for the covariance scores,
local chart conditions, and rare remainder.
There are finitely many allowed adjacent
profile pairs at each stage; take their
maximum before choosing that stage's ratio.
All these entry conditions can be included
in the same slow clock and period reserves.
They are independent of a subsequently
requested volume subsequence.

The principal estimate is first proved
in the raw norm.
Corollary~\ref{cor:normalised-incoming-transport}
then performs the constant and
affine-Hessian subtraction, preserving
the contraction and all mask corrections.
Applying the same bounded extraction
to the comparison family preserves its
small remainder, with its own scalar,
quadratic and masked terms retained.
The next model-specific
calculation is classical and one-loop
matching of the ordinary coefficients,
followed by their auxiliary-action bound
and Section 33's localisation.
The present contraction does not assert
closure of the full nonlinear RG step.

For input curves with the prescribed
fixed-address continuity and common
summable envelopes, continuity can also
be checked in the held-physical variables.
There the derivatives fall on the smooth
kernel and protector. At a fixed cutoff
their positive covariance gaps and compact
chart supports give a common integrable
bound over a compact parameter interval.
Dominated convergence proves continuity
of each finite-address integral and its
coarse derivatives; the weighted bounds
above remove the address truncation.
\clearpage
\section{The two ordinary coefficients and matching of the principal action}
\label{sec:ordinary-matching}

The estimate of Section~\ref{sec:central-strength-expansion}
leaves two ordinary strength coefficients. Their cancellation is a
statement about the total action on a comparison torus. An individual
connected summand need not have the cancellation. We first identify
the coefficients with their original chart factors, and then match
the stationary classical action through degree four and the one-loop
action through degree two. The resulting total-action estimate is
the input needed in Section~\ref{sec:finite-period-localisation}.

The organisation of the matching follows
\cite[Lemmas~12.5--12.6, PDF pp.~79--81, and Lemma~13.6,
PDF pp.~98--100]{ArmstrongDraft}. The stationary derivative formulas,
the division into curvatures, and the norm estimates used here are
given below. The background-field interpretation of the quadratic
coefficient is also discussed in
\cite[Sect.~1, eqs.~(1.20)--(1.22), printed p.~264;
Sect.~5, eqs.~(5.1)--(5.15), pp.~292--293]{Bal87}.
Its numerical logarithmic coefficient is a further calculation;
we do not assume its value in the present matching argument.

Fix the group, an adjacent pair of profiles, their derivative and
covering reserves, and an odd blocking ratio \(L\).
Except where independence of \(L\) is stated explicitly, constants
may depend on these fixed data and on the incoming exponential
path width and budget. They are independent of the cutoff,
the physical period, and the length of a history satisfying these
budgets. Write \(\beta=g^{-2}\), \(h=h_{\rm f}\), and
\(h_{\rm c}\le2h\), as in the preceding sections.
An estimate \(C_B\) has the same meaning as before: its additional
dependence on \(B\) is a fixed polynomial. No incoming regular
function is changed in this section. Its transport remains the
one proved in Section~\ref{sec:incoming-sharp-transport}.

\subsection{The coefficients with the retained chart factors}

At first take a finite collection of ordinary letters
\(\{v_\alpha\}_{\alpha\in A}\). Let \(T_\alpha\) be the
physical Gaussian coordinates on which its broad chart factors
are retained, and set
\[
 b_T=\prod_{e\in T}b_e,\qquad b_\varnothing=1.
\]
A smooth letter without such a read has \(T_\alpha=\varnothing\).
The expectation below includes the original joint Gaussian law
and the common tag and forest average. It introduces no independent
copy for touching letters. Define
\begin{equation}
 Z_A(u)=1+\sum_{\varnothing\ne J\subset A}
  \mathbb E\left[b_{\bigcup_{\alpha\in J}T_\alpha}
            \prod_{\alpha\in J}(e^{uv_\alpha}-1)\right],
 \qquad F_A(u)=\log Z_A(u),
 \label{eq:ordinary-retained-partition}
\end{equation}
with the logarithm equal to zero at \(u=0\).
The letters are labels, even when their formulas or footprints
coincide. In particular the union in the subscript occurs once;
it cannot be replaced by a product containing repeated powers
of a smooth chart factor.

\begin{lemma}[The first two retained coefficients]
\label{lem:ordinary-retained-coefficients}
With the preceding conventions, put
\begin{align}
 a_1&=\sum_{\alpha\in A}\mathbb E[b_{T_\alpha}v_\alpha],\notag\\
 a_2&=\sum_{\alpha\in A}\mathbb E[b_{T_\alpha}v_\alpha^2]
       +2\sum_{\alpha<\gamma}
         \mathbb E[b_{T_\alpha\cup T_\gamma}
                             v_\alpha v_\gamma].
 \label{eq:ordinary-retained-two-moments}
\end{align}
Here any ordering of the finite label set can be used. Then
\begin{equation}
             F_A'(0)=a_1,\qquad F_A''(0)=a_2-a_1^2.
 \label{eq:ordinary-retained-cumulants}
\end{equation}
These identities extend to the absolutely summable ordinary
family and to each of its prescribed physical derivatives.
\end{lemma}
\begin{proof}
In differentiating \eqref{eq:ordinary-retained-partition} once
at zero, only a singleton \(J\) remains. Differentiating twice
leaves either a singleton differentiated twice or two distinct
letters each differentiated once. These give \(a_1,a_2\).
Since \(Z_A(0)=1\), differentiation of its logarithm gives
\eqref{eq:ordinary-retained-cumulants}.
The joint-strength majorant of
Theorem~\ref{thm:central-joint-strength-bound} is absolute
before the labels are summed and after every prescribed derivative.
Cauchy's coefficient formula therefore allows passage through
finite subfamilies. This proves the last assertion without
differentiating a hard fine mask.
\end{proof}

On a common intact chart all the factors \(b_{T_\alpha}\)
are one and the microscopic fine masks are one. There the two
formulas are the usual first moment and second cumulant of
\(\sum_\alpha v_\alpha\). Away from that chart one must use
\eqref{eq:ordinary-retained-two-moments}. We will estimate the
difference on the auxiliary tori, rather than silently omit it.

\subsection{Exact third and fourth stationary derivatives}

Use the actual unmasked phase
\(\varphi(x,y)\) of \eqref{eq:actual-unmasked-stationary}.
The variable \(x\) is the coarse connection in exponential
coordinates and \(y\) consists of the free fine coordinates in
the chosen tree gauge. Let
\[
 S_{\rm cl}(x)=\varphi(x,\rho(x)),\qquad
 Q=\varphi_{yy}(0,0),\qquad
 Wv=(v,D\rho(0)v).
\]
Thus \(Q\) is the positive fine Hessian without its factor
\(\beta\), not a new reference covariance. For coarse directions
\(v_i,v_j\), define a fine linear functional by
\[
 B_{ij}[u]=D^3\varphi(0,0)[Wv_i,Wv_j,(0,u)].
\]
All these objects are derivatives of the same physical phase.

\begin{lemma}[Stationary jets with their symmetry factors]
\label{lem:ordinary-stationary-jets}
At the flat field,
\begin{align}
 D^3S_{\rm cl}(0)[v_1,v_2,v_3]
   &=D^3\varphi(0,0)[Wv_1,Wv_2,Wv_3],\notag\\
 D^4S_{\rm cl}(0)[v_1,v_2,v_3,v_4]
   &=D^4\varphi(0,0)[Wv_1,Wv_2,Wv_3,Wv_4]
       -\sum_{(ij|kl)}B_{ij}Q^{-1}B_{kl}.
 \label{eq:ordinary-classical-jets}
\end{align}
The sum is over the three unordered pairings
\((12|34),(13|24),(14|23)\). The cubic and quartic Taylor
coefficients are these expressions divided by \(3!\) and \(4!\).
\end{lemma}
\begin{proof}
Differentiating \(\varphi_y(x,\rho(x))=0\) once gives
\[
        D^2\varphi(0,0)[Wv,(0,u)]=0.
\]
A second differentiation gives
\[
           QD^2\rho(0)[v_i,v_j]=-B_{ij}.
\]
These are identities of fine linear functionals. To see all
coefficients, first set every coarse direction equal to \(v\)
and put \(w=D^2\rho(0)[v,v]\).
In differentiating \(\varphi(tv,\rho(tv))\) three times,
the terms with \(D^2\rho\) pair with \(Wv\) through the
Hessian and vanish. The term with \(D^3\rho\) pairs with
\(\varphi_y\) and vanishes as well. At fourth order the
remaining terms apart from \(D^4\varphi[Wv]^4\) are
\[
       6B_{vv}[w]+3\langle w,Qw\rangle
                       =-3B_{vv}Q^{-1}B_{vv}.
\]
Terms with \(D^3\rho,D^4\rho\) again vanish by the two
stationary identities. Polarisation gives the displayed three
pairings. The factorials are the ordinary Taylor factorials;
no Wick factor has yet been used.
\end{proof}

The quadratic identity is already fixed by the calibrated free
construction:
\begin{equation}
 T_2S_{\rm cl}(x)=\tfrac12\langle d_1x,H_{\rm c}d_1x\rangle.
 \label{eq:ordinary-classical-quadratic}
\end{equation}
Here \(T_k\) denotes the homogeneous Taylor term of degree \(k\).
Equations~\eqref{eq:actual-schur-complement} and
\eqref{eq:physical-linear-column-identities} prove this identity;
it is not an extra matching hypothesis.

\subsection{How a transverse Taylor kernel becomes curvature atoms}

We give the algebra needed to turn the preceding jets into an
action on physical connections. A curvature slot consists of a
path \(\gamma:0\to z\) and an oriented plaquette \(p\) based at
\(z\). Its value at an anchor \(a\) is
\begin{equation}
 \mathcal F_{(\gamma,p)}(U,a)=
  \operatorname{Ad}(U_{a+\gamma})\operatorname{pr}_{\mathfrak g}
                  \frac{P_p(U,a+z)-P_p(U,a+z)^{-1}}2.
 \label{eq:ordinary-curvature-slot}
\end{equation}
The projection is the real orthogonal projection in the fixed
matrix inner product. Its complexification is used for analytic
estimates. A slot has length \(|\gamma|+4\). A degree-\(k\)
curvature atom is an invariant real \(k\)-linear contraction of
\(k\) slots. Its length is the sum of their lengths, and its
coefficient size is the sum of the absolute colour coefficients
in the fixed orthonormal basis. In particular the path budget is
\(\sum_a |c_a|e^{\sigma\operatorname{length}(a)}\).
The paths and their lengths are on the integer cover, including
when their evaluations wrap around a torus.

For a momentum \(p\) put \(q_a(p)=e^{ip_a}-1\).
The linearisation of a plaquette is
\(d_1(p)x=q(p)\wedge x\). A \(k\)-leg Taylor kernel has
independent momenta \(p_1,\ldots,p_{k-1}\) and
\(p_0=-\sum_{i=1}^{k-1}p_i\). It is called transverse in leg
\(i\) if contraction with \(q(p_i)\) is zero. This expresses
annihilation of a linear gauge increment in that leg.
Momentum order at least \(s\) means that all Taylor coefficients
in these momenta of total degree less than \(s\) vanish at zero.

\begin{lemma}[Analytic division into curls]
\label{lem:ordinary-curl-division}
Let \(f\) be a bounded periodic holomorphic \(k\)-leg kernel
on a strip of positive width, transverse in every leg.
For \(k=3,4\), it can be written as a contraction of one curl
in every leg against a bounded periodic holomorphic kernel
on a smaller strip. Its exponential coefficient sum on a still
smaller strip is at most a fixed constant times the original
bound. The constants depend on the strip and \(k\), not on a
period or on the bound itself.
For \(k=2\) the same assertion holds if \(f\) has momentum
order at least \(s+2\), \(s\ge0\); the resulting two-curvature
kernel has momentum order at least \(s\).
\end{lemma}
\begin{proof}
We first describe division in an independent momentum. On an
exterior leg let \(\delta=\iota_q\), contraction with \(q\).
In coordinate \(a\) write an exterior tensor as
\(f=e_a\wedge u+v\), where \(u,v\) have no \(e_a\) component,
and define
\[
 D_av=\frac{v-v|_{p_a=0}}{e^{ip_a}-1},\quad
 H_af=e_a\wedge D_av,\quad R_af=v|_{p_a=0}.
\]
The quotient is holomorphic: its numerator vanishes on
\(p_a=0\), and periodicity gives the same removable zero on
each translate by \(2\pi\). If \(\delta f=0\), contraction
in the other coordinates gives \(\delta'v=-q_au\), whence
\[
                         \delta H_af=f-R_af.
\]
Repeat in four coordinates. In positive exterior degree the
final remainder is zero, because it has no remaining exterior
coordinate. In exterior degree zero it is the value at zero;
it vanishes when the scalar tensor vanishes there.
For estimates, the Fourier multiplier of division on a mode
\(e^{inp_a}\) is the finite geometric sum
\((e^{inp_a}-1)/(e^{ip_a}-1)\). Its coefficient sum is
\(|n|\), with the analogous inverse sum for \(n<0\).
This loss is bounded by shrinking the exponential width once.
The construction therefore has a bounded norm between fixed
strips. It preserves transverse conditions in other independent
legs, since their contractions commute with these operations.

The dependent momentum \(p_0\) requires one additional
argument. Suppose it has exterior degree \(d\ge1\) and
there are \(r\) remaining independent one-form legs.
If \(d+2r>4\), it has a primitive under \(\delta_0\) that
remains transverse in these \(r\) legs. Prove this by induction
on \(r\). For \(r=0\) the asserted range has exterior degree
greater than four, so the tensor is zero. For the induction
step use \(p_0\), instead of \(p_1\), as independent variable
and perform the preceding positive-degree division to obtain
\(\delta_0X=f\), preserving legs \(2,\ldots,r\).
Write \(c_1\) for contraction in leg 1. Since
\(\delta_0c_1X=0\), a second division gives
\(\delta_0Y=c_1X\). At \(p_1=0\), \(Y_0\) has dependent
degree \(d+2\) and \(r-1\) independent legs. The induction
hypothesis applies, since
\((d+2)+2(r-1)=d+2r>4\), and gives
\(\delta_0Z=Y_0\). Extend \(Z\) independently of \(p_1\).
Then \(Y'=Y-\delta_0Z\) vanishes at \(p_1=0\).
Scalar division in leg 1 gives \(c_1D=Y'\).
Consequently
\[
                  V=X-\delta_0D,\qquad
                  \delta_0V=f,\quad c_1V=0.
\]
All the other transverse conditions are preserved. Each shear
of momentum variables and each division uses a smaller fixed
strip. There are only finitely many such operations for the
values of \(k\) in the statement.

For \(k=3,4\), start with \(d=1,r=k-1\), for which the
inequality just proved holds. After division in leg 0, divide
the independent legs successively. This gives a two-form in
each leg. For \(k=2\), start instead with a kernel of order
\(s+2\). The first primitive \(X\) has order at least
\(s+1\). Hence \(c_1X\) has order at least \(s+2\), and
its primitive \(Y\) has order at least \(s+1\).
There are no other momenta at \(p_1=0\), so \(Y_0=0\).
The same correction works with \(Z=0\), and the final
independent-leg division leaves order at least \(s\).

Finally a bounded holomorphic function on a strip of width
\(\eta\) has Fourier coefficients bounded by
\(Ce^{-\eta'\sum|z_i|_1}\) for every \(\eta'<\eta\), by
shifting each coordinate contour. Summation at a smaller
exponent is finite. This proves the exponential coefficient
bound. Taking real parts, averaging the compact adjoint action,
and averaging the finite lattice symmetry group preserve the
identities and change their coefficient norms by fixed factors.
\end{proof}

Here contraction of a two-form coefficient with \(q\) is
dual to evaluation on \(q\wedge x\). Thus the last lemma
has exactly the curl interpretation stated, including its
signs once a plaquette orientation is fixed.
Choose coordinate paths to the Fourier offsets of the
two-form coefficients and use \eqref{eq:ordinary-curvature-slot}.
Their linear terms are the required curls. Path length is at
most a fixed multiple of the sum of the offsets, plus \(4k\).
Shrinking the coefficient exponent once more therefore gives
the exponential path budget claimed for these physical atoms.
This defines arrays on the integer cover before any torus
is chosen.

% !TeX root = 4DYM_notes.tex
\paragraph{A common finite-difference estimate.}
We record once the position-space calculation behind the
moment gains used here and in the current transfer. It is
the lattice Taylor principle of \cite[Sect.~3]{BS15b}; we
give the estimate in the present coefficient conventions.

Let \(K_{\boldsymbol z}\) be a scalar or fixed finite
tensor kernel, with \(\boldsymbol z=(z_1,\ldots,z_r)\),
\(z_i\in\Z^4\), and
\[
 \sum_{\boldsymbol z}|K_{\boldsymbol z}|
                 e^{\sigma\sum_i|z_i|_1}\le C_K,
                                                   \qquad \sigma>0.
\]
Tensor contraction and colour norms are included in its
absolute value. Suppose functions \(A_i\) obey
\(\|\Delta^bA_i\|_\infty\le C_iL^{-d_i-b}\) for
\(b=0,1,2\), where \(\Delta^b\) ranges over ordered
unit lattice differences. The numbers \(d_i\ge0\)
specify only the scaling of these inputs. For
\(s=1\) assume \(\sum K_{\boldsymbol z}=0\). For
\(s=2\) assume also
\(\sum z_i^aK_{\boldsymbol z}=0\) for every \(i,a\).
Then, with the displayed product taken in its fixed order,
\begin{equation}
 \left|\sum_{\boldsymbol z}K_{\boldsymbol z}
                           \prod_{i=1}^r A_i(x+z_i)\right|
       \le C_{r,s,\sigma} C_K\prod_iC_i\,
                                     L^{-\sum_i d_i-s}.
 \label{eq:ordinary-common-moment-estimate}
\end{equation}
The same inequality with \(s=0\) needs no moment
cancellation. Its constants precede \(L\).

Here is the complete difference calculation. Put
\(F(\boldsymbol z)=\prod_iA_i(x+z_i)\), and
\(\ell=\sum_i|z_i|_1\). A unit difference in the
\(i\)-th variable differentiates only that input.
Two differences either both fall on it or fall on
two different inputs. Consequently
\[
 \|\Delta F\|_\infty\le C_r\prod_iC_i L^{-\sum_i d_i-1},
 \qquad
 \|\Delta^2F\|_\infty\le C_r\prod_iC_i L^{-\sum_i d_i-2}.
\]
Telescoping along a coordinate path from zero to
\(\boldsymbol z\) gives
\(F(\boldsymbol z)-F(0)\) as a sum of \(\ell\)
first differences. Subtracting
\(\sum_{i,a}z_i^a\Delta_{i,a}F(0)\) leaves differences
of those first differences, at most
\(C_r\ell(1+\ell)\) second differences. For a negative
step use
\(F(y-e)-F(y)=-\Delta_eF(y-e)\); comparing this with
\(-\Delta_eF(0)\) has the same second-difference bound.
Thus signed coordinates require no smooth extension
of a lattice function.

The zero moment removes \(F(0)\), and the first moments
remove the displayed linear polynomial. The remaining
absolute sum is bounded by
\(C_r\sum|K_{\boldsymbol z}|(1+\ell)^s\) times the
appropriate difference bound. Finally
\((1+\ell)^s\le C_{s,\sigma}e^{\sigma\ell}\)
proves \eqref{eq:ordinary-common-moment-estimate}.
This proof also works for a prescribed matrix product,
by the product rule and the fixed finite-dimensional
norm bound on multiplication.

In subsequent applications no new Taylor estimate is needed:
first remove the matched polynomial moments, identify the
input dimensions \(d_i\), and use \(s=0,1\) or \(2\)
in \eqref{eq:ordinary-common-moment-estimate}.
For two curvature inputs of dimension two and two cancelled
moments, the two second-difference placements have prices
\[
       (CL^{-3})(CL^{-3}),\qquad
       (CL^{-2})(CL^{-4}),
\]
both \(CL^{-6}\). A subsequent volume rescaling \(L^D\)
therefore gives \(CL^{D-6}\): \(D=4\) for an action and
\(D=5\) for a current defect. This single calculation
explains their different powers without introducing
different localization norms. Curl division and the
verification of the moments remain necessary at each use.

For rooted coefficient sums, take absolute values
before summing the output and input indices. Use the
weighted row bounds of the free columns on the
differenced inputs. Maximum-distance weights give
the required \(1\)-distance width after reducing
their exponent by the fixed factor four.
A translated coarse weight uses
at most its width divided by \(L\) on a fine
displacement. The admitted ratio floor leaves half
the native width for the kernel moments above.
Tonelli's theorem then gives the identical bound in
those rooted sums. All displacements are on the
integer cover; periodic evaluation sums their
original coefficients. No period-dependent norm
equivalence or residue count is used.

\subsection{Classical matching and the order of the path bounds}

Let \(S_{\rm row}\) be the unmasked coarse row action with
the already constructed free form \(H_{\rm c}\). It consists
of the direct Wilson term with coefficient \(\lambda_{\rm c}\)
and the squares of the transported-curvature rows from
\eqref{eq:all-scale-row-factor}. In particular
\[
      T_2S_{\rm row}=T_2S_{\rm cl},\qquad
      S_{\rm row}(1)=0.
\]
We now construct the cubic and quartic atoms, rather than
postulate agreement with this quadratic action.

\begin{proposition}[Classical matching]
\label{prop:ordinary-classical-matching}
There are real invariant curvature-atom arrays \(P_3,P_4\)
such that
\begin{equation}
 S_{\rm c}=S_{\rm row}+P_3+P_4,\qquad
             T_{\le4}(S_{\rm cl}-S_{\rm c})=0.
 \label{eq:ordinary-classical-matching}
\end{equation}
The arrays are independent of \(\beta\) and the physical
period. For fixed incoming path width and budget their
exponential budgets are finite, uniformly in the cutoff
and the finite history. Moreover their output width and
budget can be fixed before \(L\), provided \(L\) is then
taken above a floor depending on the incoming path data.
On the common analytic chart their remaining connected family
satisfies
\begin{equation}
 \|\beta(S_{\rm cl}-S_{\rm c})\|^{\rm chart}_{q,r,s;\mathrm{local}}
                      \le C_B\beta r^5.
 \label{eq:ordinary-classical-remainder}
\end{equation}
Here the local norm means the anchored, exponentially weighted
sum of the coefficient families used in
Proposition~\ref{prop:stationary-chart-local-families}.
\end{proposition}
\begin{proof}
First fix \(L\). The stationary rooted expansions and
Lemma~\ref{lem:ordinary-stationary-jets} give holomorphic
translation-invariant cubic and quartic kernels with finite
exponential bounds. Gauge invariance supplies their needed
transverse identities as follows. If an invariant analytic
functional \(D\) starts in field degree \(k\), differentiate
\[
             DD(x)[d_0\omega+O(x\omega)]=0.
\]
The coefficient of degree \(k-1\) in \(x\) is
\(D^kD(0)[d_0\omega,x,\ldots,x]=0\). Polarisation gives
transversality in every leg of its leading kernel.

Apply this to \(D=S_{\rm cl}-S_{\rm row}\), which starts
at degree three. Lemma~\ref{lem:ordinary-curl-division} and
its path construction give \(P_3\) with the same cubic jet.
Use the entire nonlinear physical function \(P_3\), not its
cubic Taylor polynomial, in the next difference
\(D_4=S_{\rm cl}-S_{\rm row}-P_3\).
Its jets through degree three vanish, so its quartic kernel
is transverse in every leg. A second application gives
\(P_4\). This proves \eqref{eq:ordinary-classical-matching}.
Subtracting only the polynomial \(T_3P_3\) would lose
invariance of the residual at degree four and would not
justify this second division.

We give the estimates that permit the output budget to be
chosen before the ratio. They also explain why the construction
can be repeated. The smooth free lift \(E\) satisfies
\eqref{eq:all-scale-linear-difference-bounds}:
\[
       \|\nabla_f^aEv\|_\infty
           \le C_jL^{-1-a}\|v\|_\infty,\qquad 0\le a\le2.
\]
The constant precedes \(L\). The axial tangent \(W\) in
\eqref{eq:ordinary-classical-jets} is a different representative
of this lift. The two differ by a fine gauge increment,
including the induced frame increment. Work in joint fine-link
and centred frame coordinates to use the smooth representative.
Here centring removes the common frame rotation: at the linear
level the local frame increments have weighted mean zero,
and their common mean is retained as the separate coarse
frame increment. Nonlinearly the centre is the unique nearby
minimiser of the normalised group-distance score. Its flat
Hessian on the centre variable is the identity.
Gauge invariance changes none of the stationary values or jets.

For clarity, the inverse needed for the quartic pairing has
a uniform bound in block energy units, even though an axial
counting-norm inverse need not. Put \(\mu=L^{-1}\), let
\(\theta_z=(p+2\pi z)/L\) be the fine aliases of a real
coarse momentum, and set
\[
 \omega_z^2=\mu^2+\sum_a|e^{i\theta_{z,a}}-1|^2,\qquad
            \|u\|_{\rm bl}^2=\sum_z\omega_z^2|u_z|^2.
\]
The alias transform is unitary on one block; coarse frame
variables use their Euclidean norm. On their direct sum
take the quotient by the joint infinitesimal gauge directions.
The quadratic energy on this quotient is
\begin{equation}
 \langle d_1u,H_{\rm f}d_1u\rangle
                +a\|Q_0u-d_0\vartheta\|^2.
 \label{eq:ordinary-block-quotient-energy}
\end{equation}
Here \(\vartheta\) is the centred coarse frame increment
and \(Q_0\) is the linear average in these centred coordinates.
This is the same quadratic form as the axial fine Hessian
after taking a slice.

% !TeX root = 4DYM_notes.tex
We give the centring and the inverse estimate explicitly.
Let \(\vartheta_{\rm old}\) be the coarse-frame tangent
before centring, and let \(P\) be the frame correction
in \eqref{eq:actual-predictor-linear-frame}. Set
\[
                 \vartheta=\vartheta_{\rm old}-Pu.
\]
This is the weighted mean frame tangent: the local
rooted frame has tangent \(t_z-t_0\), whose weighted
mean is \(-Pu\). The actual framed noise becomes
\[
 \widehat Qu-d_{0,c}\vartheta_{\rm old}
       =(Q+d_{0,c}P)u-d_{0,c}\vartheta_{\rm old}
                                      =Qu-d_{0,c}\vartheta.
\]
Thus \(Q_0\) in
\eqref{eq:ordinary-block-quotient-energy} is the
straight average \(Q\) in these centred coordinates.
This change retains the separate frame variable.
It does not change the physical predictor.
For each fixed fine connection the corresponding
change of the group-valued coarse frame is a
Haar-preserving multiplication. The joint change
is triangular in the fine connection and frame
variables; its Gaussian Hessian changes by
congruence, with the same physical normalisation.

Use a unitary alias transform on one block, with
\(L^4\) residues. If \(u_z\) is one unitary alias
coordinate, its fine plane wave has amplitude
\(L^{-2}u_z\). Consequently the straight average
has alias column
\[
 B_z=\mu h_1(\theta_z),\qquad
 Q_0u=\sum_zB_zu_z,\qquad \mu=L^{-1}.
\]
This factor is \(L\cdot L^{-2}\); it is important
for the lower bound. The site average \(A_0\)
from \eqref{eq:rooted-predictor-gradient} has
alias column \(\mu^2f(\theta_z)\), where
\(f=a_L^j\). A fine gauge parameter \(\phi\)
acts in the centred coordinates by
\[
                  (u,\vartheta)\longmapsto
                  (u+d_{0,f}\phi,\vartheta+A_0\phi).
\]
The noise is unchanged because
\(Q_0d_{0,f}=d_{0,c}A_0\).

At a real coarse momentum use a gauge change to
make every nonprincipal alias transverse to
\(q(\theta_z)\), and then make \(\vartheta=0\).
For \(z\ne0\) the first step chooses the gauge
coefficient by division by \(|q(\theta_z)|^2\).
The remaining principal gauge coefficient sets
\(\vartheta=0\), since
\(\mu^2f(\theta_0)\ne0\). This is a section of
the joint gauge quotient, including at \(p=0\):
there the principal gauge parameter acts on the
frame although its fine gradient is zero.

For transverse nonprincipal aliases the exterior
algebra identity gives
\[
 \|d_1(\theta_z)u_z\|^2
    =|q(\theta_z)|^2|u_z|^2
       -|q(\theta_z)^*u_z|^2
    =|q(\theta_z)|^2|u_z|^2.
\]
The sine bound gives
\(|q(\theta_z)|\ge c\mu(1+|z|_\infty)\) for
\(z\ne0\). Hence \(\omega_z^2\le C|q(\theta_z)|^2\),
and the uniformly positive curvature metric bounds
\(c_j\sum_{z\ne0}\omega_z^2|u_z|^2\).
The principal averaging column is diagonal up
to its scalar filter and satisfies
\[
               |B_0v|\ge c_j\mu|v|.
\]
Indeed the principal scalar filter and every
principal line average have the sine-quotient
lower bound of
Lemma~\ref{lem:cumulative-filter-bounds}.
For the other columns,
\begin{equation}
 \sum_{z\ne0}\frac{\|B_z\|^2}{\omega_z^2}
 \le C_j\sum_{z\ne0}
       \frac{w_j(z)^2}{(1+|z|_\infty)^2}\le C_j.
 \label{eq:block-average-weighted-columns}
\end{equation}
This is an absolutely summable alias sum, not
a bound by the \(L^4\) aliases separately.
Since \(B_0u_0=Q_0u-\sum_{z\ne0}B_zu_z\),
Cauchy--Schwarz with the weights \(\omega_z\)
now gives
\begin{align*}
 \mu^2|u_0|^2
 &\le C_j\left\{|Q_0u|^2+
                      \left|\sum_{z\ne0}B_zu_z\right|^2\right\}\\
 &\le C_j\left\{|Q_0u|^2+
                         \sum_{z\ne0}\omega_z^2|u_z|^2\right\}.
\end{align*}
Also \(\omega_0\le C\mu\). This proves a lower
bound for the block norm on this section.
The quotient norm is the infimum of the block
norm over all joint gauges, so it has the same
coercive lower bound. Conversely
\(\|B_z/\omega_z\|\) is square-summable also
including \(z=0\), and the coarse gradient is
bounded. The full energy is at most
\(C_j(\|u\|_{\rm bl}^2+|\vartheta|^2)\).
Taking the infimum over gauges proves the upper
quotient bound as well. In particular this
section has norm at most a fixed multiple of
the quotient norm.

We also justify the complex strip used for local
kernel estimates. Freeze the weights and the
preceding transverse section at a real momentum
\(p\). For a complex increment \(i\eta\),
\[
 |q(\theta_z(p+i\eta))-q(\theta_z(p))|
                                      \le C\mu|\eta|.
\]
Since \(\omega_z\ge\mu\), the change in the
curvature quadratic operator, between block
energy units, has norm at most \(C_j|\eta|\).
The derivative bounds for the positive metric
follow from its fixed holomorphic strip. For
the averaging row use
\eqref{eq:block-average-weighted-columns} on a
slightly larger strip and Cauchy's estimate:
\[
 \left\|\left(B(p+i\eta)-B(p)\right)
                    \operatorname{diag}(\omega_z^{-1})\right\|
                                      \le C_j|\eta|.
\]
Thus its squared noise operator has the same
perturbation bound. This proves
\(\|\delta H\|_{\rm bl\to bl,*}\le C_j|\eta|\)
for the restricted quadratic operator.

The fixed section remains a complement to the
complex gauge directions: for nonprincipal
aliases its gauge-fixing coefficients change
by at most \(C\mu|\eta|/|q(\theta_z)|\le C|\eta|\),
and the principal site-average coefficient
remains nonzero. Choose \(|\eta|\) so small
that this property holds and the preceding
operator perturbation is less than half the
real coercive constant. The inverse Neumann
series gives a common dual-to-primal bound
\(E_j\) on this strip. All constants here
precede \(L\). Local choices of section give
the same pairings for gauge-annihilating
functionals, so these analytic formulas agree
on overlaps.

In particular a functional \(b\) on the
quotient has dual norm
\[
 \|b\|_{\rm bl,*}
    =\sup_{\|[u,\vartheta]\|_{\rm bl}\le1}
                                      |b([u,\vartheta])|.
\]
Its quadratic pairing is at most
\(E_j\|b\|_{\rm bl,*}^2\). Returning to an
axial slice uses the finite block-tree gauge
map. It preserves this pairing; any polynomial
residue-counting price belongs to component
estimates, rather than to the quotient inverse
constant \(E_j\).

The mixed cubic functional in
\eqref{eq:ordinary-classical-jets} annihilates joint gauge
directions. Indeed differentiating joint gauge invariance twice
along the stationary curve leaves its contraction with a gauge
direction; all other terms contain a stationary fine gradient
or the Hessian paired with one stationary tangent and a fine
increment. These vanish. Pairing this functional through the
quotient inverse therefore gives exactly its pairing through
the axial inverse \(Q^{-1}\). There is no additional gauge
determinant in this classical identity.

We next bound the vertices. Give a connection component and
a fine coordinate difference weight one. The leading part of
every transported plaquette curvature has weight two and is
\(dA+[A,A]\). Changing the based transport changes this
expression only in higher weight. The coefficient of its
weight-four square is
\[
 \lambda_{\rm f}\delta_{pr}+\sum_b s_{bp}s_{br}=\delta_{pr},
\]
where \(s_{bp}\) is the sum of the row coefficients of
orientation \(p\). This is \(H_{\rm f}(0)=I\), already
proved for the constructed free form. The cubic and quartic
curvature atoms start at weights six and eight. Thus the fine
leading action and Wilson action have the same terms of
weight at most four. Their degree-\(k\) difference has at
least \(5-k\) fine differences for \(k=2,3,4\).
This statement concerns the full field jets, including the
nonlinear expansion of row transports.

The assertion can be estimated without a continuum limit.
Expand each finite word through degree four, and write
\(e^{ip_ax_a}-1\) as a coordinate difference times its finite
geometric sum. Each removed low-momentum Taylor coefficient
allows one more such factor. The resulting losses are fixed
polynomials in the actual path lengths; their sums are bounded
by the incoming exponential budget at a smaller exponent.
Fine momentum conservation moves differences off an arbitrary
fine leg when needed. For an arbitrary leg of block norm one,
\[
             \|u\|_{1,\mathrm{block}}\le CL^3,
       \qquad\|\nabla_fu\|_{1,\mathrm{block}}\le CL^2,
\]
by Cauchy--Schwarz, \(\|u\|_2\le L\|u\|_{\rm bl}\),
and \(\|\nabla_fu\|_2\le C\|u\|_{\rm bl}\).
With all other legs equal to smooth lifts, these are the same
scale counts as summing \(L^4\) sites and assigning
\(L^{-1}\) to each smooth field and \(L^{-1}\) to each
fine difference. A degree-\(k\), momentum-degree-\(s\)
vertex therefore has scale
\[
                              L^{4-k-s}.
\]
Wilson's vertices have \(s\ge4-k\); the differences just
described have \(s\ge5-k\). Hence, on a fixed smaller
coarse strip, their all-smooth and mixed dual block bounds
have the form
\begin{equation}
 C_3\le c_{3,j}+L^{-1}A_{3,j},\qquad
 C_4\le c_{4,j}+L^{-1}A_{4,j}.
 \label{eq:ordinary-pre-ratio-vertices}
\end{equation}
The \(c_{k,j}\) depend only on the free and geometric profile
data; \(A_{k,j}\) may also depend on the incoming path
width and budget. The geometric penalty has a fixed contribution
to \(c_{k,j}\): after centroid normalisation the linear frame
equation is identity on the mean-zero variables, while every
higher coefficient is a normalised average of words.
Differentiation through order four uses only finitely many
products of these averages. Route lengths cancel the
\(L^{-1}\) of a smooth field. Their bounds therefore contain
no number of patch sites. Centroid normalisation is legitimate
because the score and the predictor are unchanged by the common
frame gauge. This is also the geometry behind the linear
frame identity in \eqref{eq:actual-predictor-linear-frame}.

These estimates hold on a strip fixed before the incoming
amplitude is chosen. A path of length \(n\) moves at most
\(C_j+n/L\) coarse units. Its factor on that strip is bounded
by \(e^{C_j\eta}e^{\eta n/L}\); first taking
\(L\ge2\eta/\sigma\) pays it from the incoming width
\(\sigma\). All remaining fixed powers of \(n\) use the
other half of that width. This explains the ratio-dependent
floor and avoids a loss of a fixed fraction of the spatial
width at every successive step.

Choose \(L\) so that \(L^{-1}A_{3,j},L^{-1}A_{4,j}\le1\).
The factorial-normalised stationary kernels then have bounds
\[
       C_3/6,\qquad (C_4+3E_jC_3^2)/24,
\]
by \eqref{eq:ordinary-classical-jets} and the quotient inverse
estimate. The row-action kernels have fixed bounds as well.
Curl division first for the cubic difference and then for
the quartic difference therefore gives a fixed output width
and budget. The fourth jet of the newly constructed full
\(P_3\) has a fixed bound, using one further reserve in its
path exponent. This establishes the asserted order of choices.

Finally \(S_{\rm cl}-S_{\rm c}\) starts at degree five.
Apply the rooted analytic majorant and the monomial estimate
in Corollary~\ref{cor:stationary-taylor-remainder} to this
difference. The coarse path action has the same analytic
path bounds, so the proof applies to it as well and gives
\eqref{eq:ordinary-classical-remainder}.
\end{proof}

The common factor \(\beta\) in the leading fine action and
noise penalty cancels from the stationary equation. Thus none
of these cubic or quartic arrays has been tuned to an actual
coupling or to a volume subsequence. At the Wilson top the
same division, with \(S_{\rm cl}\) replaced by the Wilson
energy and \(H(0)=I\), matches the chosen top row action
through degree four. The remaining exact difference belongs
to the initial activity decomposition, including its masks;
the Wilson measure itself is unchanged.

\subsection{The one-loop action, including the Haar density}

Use the same stationary configuration \(U_{\rho(x)}\), and set
\begin{equation}
 {\mathsf H}(x)=\varphi_{yy}(x,\rho(x)),\qquad
 T(x)=I_f(U_{\rho(x)})-\sum_{e\in I}\log j_G(\rho_e(x))
                         +\tfrac12\log\det {\mathsf H}(x).
 \label{eq:ordinary-one-loop-action}
\end{equation}
The Haar factor is the one in
\eqref{eq:haar-chart-density}; \(I\) is the original set of
free tree-gauge links. Thus \(T=\psi_*+\tfrac12\log\det {\mathsf H}\).
The terms \((\dim y)\log\beta\), the noise normaliser,
and the constant Haar chart multiplier belong to the vacuum
scalar. They do not contribute a field derivative, but are
retained in the density identity below.

\begin{lemma}[The invariant one-loop quadratic jet]
\label{lem:ordinary-loop-hessian}
The combination in \eqref{eq:ordinary-one-loop-action} is
invariant under a change of fine gauge slice, and is gauge
invariant as a function of the coarse connection. Its Hessian
is
\begin{align}
 D_{ij}T={}&D_{ij}(I_f(U_\rho))
          -D_{ij}\sum_e\log j_G(\rho_e)\notag\\
 &+\tfrac12\operatorname{Tr}
       \{{\mathsf H}^{-1}D_{ij}{\mathsf H}-{\mathsf H}^{-1}(D_i{\mathsf H}){\mathsf H}^{-1}(D_j{\mathsf H})\}.
 \label{eq:ordinary-loop-hessian}
\end{align}
In particular it is determined by the leading phase jets
through order four, the inherited loop Hessian, and the
quadratic Haar jet. It is not determined by the determinant
alone.
\end{lemma}
\begin{proof}
For a change of free coordinates \(y=y(\widetilde y)\)
at the stationary point put \(J=D_{\widetilde y}y\).
Stationarity removes the terms involving second derivatives
of this change from the Hessian transformation, so
\(\widetilde {\mathsf H}=J^T{\mathsf H}J\). The Haar coordinate density transforms
as \(\widetilde j=j|\det J|\). Consequently
\[
       \tfrac12\log\det\widetilde {\mathsf H}-\log\widetilde j
                    =\tfrac12\log\det {\mathsf H}-\log j.
\]
Tree Haar integration has unit gauge volume by
Lemma~\ref{lem:tree-haar}. This proves slice independence.
A coarse gauge change induces just such a transformation
of the fine integration variables, and \(I_f\) is invariant.
It follows that \(T\) is coarse gauge invariant.

Differentiate \(\log\det {\mathsf H}\) once as
\(D\log\det {\mathsf H}=\operatorname{Tr}({\mathsf H}^{-1}D{\mathsf H})\), and use
\(D {\mathsf H}^{-1}=-{\mathsf H}^{-1}(D{\mathsf H}){\mathsf H}^{-1}\) for the second derivative.
This gives \eqref{eq:ordinary-loop-hessian} with its sign and
factor \(1/2\). All derivatives are along the stationary graph.
For example, at zero,
\begin{align*}
 D_i{\mathsf H}&=\varphi_{yyx}[v_i]+\varphi_{yyy}[\rho_i'],\\
 D_{ij}{\mathsf H}&=\varphi_{yyxx}[v_i,v_j]
       +\varphi_{yyyx}[\rho_i',v_j]
       +\varphi_{yyyx}[\rho_j',v_i]
       +\varphi_{yyyy}[\rho_i',\rho_j']
       +\varphi_{yyy}[\rho_{ij}''].
\end{align*}
Here subscripts \(i,j\) denote differentiation in the coarse
directions, and \(\rho_i',\rho_{ij}''\) are their stationary
derivatives. Lemma~\ref{lem:ordinary-stationary-jets} computes
\(\rho_{ij}''\) from the third phase jet.
Moreover \(DI_f(1)=0\), because its atoms have two curvature
slots, and \(D\log j_G(0)=0\). Thus their second derivatives
in \eqref{eq:ordinary-loop-hessian} require only the first
stationary lift and their own second jets. No fifth phase
derivative is hidden in the formula.
\end{proof}

The quadratic coefficient is chosen on the integer cover.
A closed trace on a torus can lift to a path whose total
displacement is zero, or a nonzero integer multiple of the
period. Call the first type \emph{nonwinding} and the second
type \emph{winding}. The cover kernel below uses traces of
the first type; the others are retained separately.
Let \(Y(p)\) be the translation-invariant kernel of \(T_2T\),
including its Taylor factor \(1/2!\), so that a quadratic
functional is written as \(\langle x,Yx\rangle\).
Write \(Y_{\rm kin}=\tfrac12d_1^*H_{\rm c}d_1\).
For a unit colour basis vector indexed by \(b\), put
\begin{equation}
 \alpha=2\operatorname{Re}[t^2]Y_{(2,b),(2,b)}(te_1),
       \qquad [t^2]f(t)=\tfrac12f''(0).
 \label{eq:ordinary-alpha-definition}
\end{equation}
The same operation gives one on \(Y_{\rm kin}\), since
\(H_{\rm c}(0)=I\). The coefficient \(\alpha\) is independent
of \(b\): an invariant symmetric bilinear form on the simple
compact Lie algebra is a scalar multiple of its fixed metric.
The finite signed-permutation average is understood here as
in the principal action.

\begin{proposition}[One-loop matching]
\label{prop:ordinary-one-loop-matching}
There is a quadratic curvature-atom action \(I_c\), defined
on the integer cover, for which the nonwinding one-loop action
satisfies
\begin{equation}
       T_{\le2}\{T-T(1)-\alpha S_{\rm c}-I_c\}=0.
 \label{eq:ordinary-one-loop-matching}
\end{equation}
Its curvature kernel has exponentially summable coefficients
\(a_{pr}^{bc}(z)\). For each pair of orientations and colours,
\begin{equation}
        \sum_z a_{pr}^{bc}(z)=0,\qquad
        \sum_z z_i a_{pr}^{bc}(z)=0\quad(1\le i\le4).
 \label{eq:ordinary-loop-moments}
\end{equation}
At fixed incoming width its budget can be bounded by
\begin{equation}
              C_jL^{-2}\Lambda_f+C_{j,L}.
 \label{eq:ordinary-loop-budget}
\end{equation}
The first constant is chosen before \(L\). The second may
depend on \(L\) and the fixed classical path budgets; it is
independent of \(\Lambda_f\), cutoff, period, and history.
The coefficient \(\alpha\) is independent of \(\beta\),
the period, and the incoming loop budget, and obeys
\(|\alpha|\le C_{j,L}\).
On a coarse torus of lattice side \(\ell_{\rm c}\), the same matching has an additional
winding family with chart norm per anchor at most
\(C_{j,L}e^{-c_{j,L}\ell_{\rm c}}\). The remaining local family obeys
\begin{equation}
 \|T-T(1)-\alpha S_{\rm c}-I_c\|^{\rm chart}_{q,r,s;\mathrm{local}}
       \le C_B(1+\Lambda_f+|\alpha|)r^3
                   +C_{j,L}e^{-c_{j,L}\ell_{\rm c}}.
 \label{eq:ordinary-loop-remainder}
\end{equation}
\end{proposition}
\begin{proof}
% !TeX root = 4DYM_notes.tex
\emph{The actual second derivative.}
Put \(R_0=\mathsf H(0)^{-1}\), and let \(u,v\) be two
coarse link increments. Write \(\rho'_u=D\rho(0)u\) and
\(\zeta_u=(u,\rho'_u)\) in the joint coarse/fine chart.
Differentiating \(\varphi_y(x,\rho(x))=0\) twice gives
\[
 \rho''_{uv}=-R_0\varphi^{(3)}(0)[(0,\cdot),\zeta_u,\zeta_v].
\]
Here the last expression denotes the fine covector obtained
by leaving its indicated slot open. In particular only the
third phase jet is used. The two derivatives of the actual
conditional precision are
\begin{align*}
 \mathsf H_u={}&\varphi_{yyx}[u]+\varphi_{yyy}[\rho'_u],\\
 \mathsf H_{uv}={}&\varphi_{yyxx}[u,v]
       +\varphi_{yyyx}[\rho'_u,v]+\varphi_{yyyx}[\rho'_v,u]\\
       &+\varphi_{yyyy}[\rho'_u,\rho'_v]
                                      +\varphi_{yyy}[\rho''_{uv}].
\end{align*}
Every phase derivative in these lines is evaluated at zero.
Let \(\mathcal L u\) be the full first fine-link variation
of \(U_{\rho(x)}\), including the replica and frame terms,
and put \(\mathsf J_G=D^2\log j_G(0)\). Since
\(DI_f(1)=D\log j_G(0)=0\), the complete quadratic read is
\begin{align}
 D^2T(0)[u,v]={}&D^2I_f(1)[\mathcal L u,\mathcal L v]
       -\sum_{e\in I}\mathsf J_G[\rho'_{u,e},\rho'_{v,e}]
       \notag\\
 &+\tfrac12\operatorname{Tr}(R_0\mathsf H_{uv})
       -\tfrac12\operatorname{Tr}(R_0\mathsf H_uR_0\mathsf H_v).
 \label{eq:ordinary-loop-complete-second-read}
\end{align}
Thus the last trace has two cubic insertions, and the
preceding one includes the quartic insertion and the
second stationary derivative. The Haar term has its own
minus sign. Retaining only the two-cubic trace would not
give this Hessian. This is the stationary Gaussian
coefficient organization of
\cite[Sect.~1, eqs.~(1.3)--(1.7), printed pp.~260--261]{Bal87},
with the present Haar and coordinate conventions written out.

For a real source \(x\), its quadratic Taylor functional
is \(D^2T(0)[x,x]/2\). Consequently the kernel \(Y\)
in \eqref{eq:ordinary-alpha-definition} includes this
factor \(1/2\); \(\alpha\) includes the compensating
factor two. The term \(T(1)\), the Gaussian dimension
factor, the constant Haar multiplier and the noise
normalizer stay in the exact vacuum scalar. Taking
field derivatives removes their contribution to this
read without removing them from the physical density.

\emph{A rooted bound before taking the period limit.}
All inverse and phase vertices in
\eqref{eq:ordinary-loop-complete-second-read} have the
rooted exponential bounds of
Section~\ref{sec:stationary-localisation}. To see where a
trace is paid, attach one external coarse leg to its
anchor. Every internal index is joined to that anchor
by inverse edges and vertex paths. Sum their absolute
coefficients using Lemma~\ref{lem:rooted-kernel-composition}.
The trace closes the last index; it does not free its
coarse anchor for an additional volume sum. There are
at most \(4\dim\mathfrak g\,L^4\) real fine free
coordinates per coarse block. Their residue count and
the finite sums of exponential distance shells are
included in a constant \(C_{j,L}\).

For example the second trace is summed in the order
\[
 \sum_{a,b,c,d}(R_0)_{ab}(\mathsf H_u)_{bc}
                         (R_0)_{cd}(\mathsf H_v)_{da}.
\]
Rooting at the external leg of \(\mathsf H_u\), the
four factors form a connected loop. Their summed
connecting length dominates the distance between the
two external coarse legs. Keeping a smaller exponential
weight therefore gives an exponentially summable kernel
in that distance. The other trace has the same argument
with fewer factors. Absolute convergence permits summing
the indices in this order. Periodization can identify
indices, but each lifted coefficient still occurs with
its original absolute price. This proves a common strip
bound at fixed \(L\), uniform in cutoff, period and
admissible history. It does not multiply a per-entry
bound by the number of blocks.

\emph{The kinetic coefficient and the two cancelled moments.}
Slice independence and coarse gauge invariance were proved
in Lemma~\ref{lem:ordinary-loop-hessian}. Their first
derivative vanishes by simplicity. The quadratic kernel
therefore annihilates a fine linear gauge increment in
either leg. Write its momentum expansion at zero as
\(Y^{[0]}+Y^{[1]}+Y^{[2]}+\cdots\). The first Ward
coefficient is \(p^TY^{[0]}=0\) for every \(p\), hence
\(Y^{[0]}=0\). Central reflection kills the possible
leading linear term: its two link signs cancel, while
the changes of link base have translation multipliers
equal to one at zero.

For clarity, signed coordinate invariance leaves the
following possibilities at degree two. A diagonal entry
is a multiple of \(\sum_{k\ne i}p_k^2\), plus a multiple
of \(p_i^2\); an off-diagonal \(ij\)-entry is a multiple
of \(p_ip_j\). Permutation symmetry makes the constants
independent of the distinct indices. Thus
\[
            Y^{[2]}_{ij}
               =a\delta_{ij}|p|^2+b p_ip_j+c\delta_{ij}p_i^2.
\]
Contracting the first index with \(p_i\) gives
\((a+b)|p|^2p_j+cp_j^3=0\). Taking first \(p_j=0\)
after factoring \(p_j\), and then arbitrary \(p\), gives
\(b=-a\) and \(c=0\). The invariant colour form is a
scalar multiple of the fixed metric on the simple Lie
algebra. Thus one transverse ray reads the entire
kinetic coefficient. Since \(H_{\rm c}(0)=I\),
\eqref{eq:ordinary-alpha-definition} gives one on
\(Y_{\rm kin}\), with exactly its Taylor normalization.

Subtract \(\alpha Y_{\rm kin}\). Its residual has no
degree-three term: a leading odd degree would change
sign under central reflection, whereas the two link
signs cancel; link-base phases multiply lower terms
which have already vanished. The residual has order
at least four. The two-leg curl division of
Lemma~\ref{lem:ordinary-curl-division} removes two powers
and gives an exponentially summable curvature kernel
\(a(p)\) with momentum order at least two. Fourier
evaluation and differentiation are absolutely convergent,
so
\[
 a(0)=\sum_z a(z)=0,\qquad
 \partial_i a(0)=i\sum_z z_i a(z)=0.
\]
They are precisely \eqref{eq:ordinary-loop-moments}.
Complete this kernel by actual transported curvature
atoms, taking the first at the anchor and the second
at offset \(z\). This defines the gauge-invariant
nonlinear function \(I_c\), rather than just its
quadratic polynomial. Its lower jets and those of
\(\alpha S_{\rm c}\) give
\eqref{eq:ordinary-one-loop-matching}.

\emph{The inherited budget contracts with a pre-ratio constant.}
The first term of \eqref{eq:ordinary-loop-complete-second-read}
is the inherited fine loop kernel pulled back by the full
linear lift. A gauge part of that lift is killed by its
transversality. Write fine aliases in centered residue
coordinates as \(\theta_z=(p+2\pi z)/L\). The two fine
curvature moments imply that its link kernel vanishes
through order three. On the chosen fixed coarse strip,
its absolute bound is therefore
\[
             C_j\Lambda_f L^{-4}(1+|z|_1)^4.
\]
Near the principal alias this is Taylor's remainder;
at aliases comparable to \(L\) it is its original
bounded strip norm. Each lift leg has bound
\(C_jL^{-1}w_j(z)\), where
\(w_j(z)=\prod_a(1+|z_a|)^{-j}\).
The fine cell sum contributes \(L^4\). Hence
\begin{align}
 \|Y^{\rm inherited}\|_{\rm strip}
 &\le C_j\Lambda_f\,
          \underbrace{L^4}_{\text{cell sum}}
          \underbrace{L^{-2}}_{\text{two lift legs}}
          \underbrace{L^{-4}}_{\text{fine kernel}}
       \sum_z(1+|z|_1)^4w_j(z)^2\notag\\
 &\le C_jL^{-2}\Lambda_f .
 \label{eq:ordinary-loop-inherited-factor-accounting}
\end{align}
The alias sum is bounded independently of \(L\), for
\(j\ge4\); extending its finite residue set to
\(\mathbb Z^4\) increases the majorant.

This inherited part has zero kinetic read. The principal
alias keeps the order-four zero of the fine kernel.
At a nonprincipal alias at least one averaging filter
vanishes to order \(j\). The pole-free filled-column
formulas \eqref{eq:filled-E-column}--%
\eqref{eq:filled-T-column} show that removing a fine
gradient costs at most two powers. The transverse part
of each such lift leg therefore vanishes to order at
least two. Two legs give order at least four. Thus
\(\alpha\) is independent of the inherited loop data.
Curl division at the fixed reduced widths preserves
the bound \(C_jL^{-2}\Lambda_f\). Its constant is
chosen before \(L\), as are the free-column constants.

The remaining Haar and determinant terms have the rooted
bound \(C_{j,L}\) proved above. Cauchy's formula on a
fixed external momentum circle bounds their second
coefficient by \(C_{j,L}\). Subtraction of its kinetic
kernel and the same curl division give
\eqref{eq:ordinary-loop-budget}. None of these terms
contains \(\Lambda_f\) or the actual coupling \(\beta\).
This proves the asserted independence of \(\alpha\).
Its error uniform in \(L\) is a further calculation
in Section~\ref{sec:one-loop-drift}; a fixed-ratio
bound does not replace that calculation.

\emph{Winding and higher field orders.}
The only failure of direct cover periodization in the
quadratic read is a closed internal trace with nonzero
period displacement. Such a lifted loop has connecting
length at least \(\ell_{\rm c}\). Keep half its native
exponential weight to sum its internal indices and
analytic degrees; the other half gives
\(e^{-c_{j,L}\ell_{\rm c}}\) per anchor. A fixed number
of external derivatives adds a length polynomial, paid
by one further fixed reduction of the width. Stationary
trees and inherited substitutions have no free closed
internal displacement and periodize with their original
coefficients. Thus the winding family is retained as an
activity and never changes \(\alpha\).

Finally the matched nonwinding analytic function has
no field terms below degree three. On a fixed smaller
chart, its rooted Taylor series and each fixed scaled
derivative are bounded by
\[
 C_{j,L,q}(1+\Lambda_f+|\alpha|)
          \sum_{a\ge3}(a+1)^q(r/r_0)^a
       \le C_{j,L,q,r_0}(1+\Lambda_f+|\alpha|)r^3
\]
for \(0<r\le r_0/2\). The derivative directions use
the same radius scaling as the chart norm; a derivative
on a degree-\(a\) term therefore preserves its factor
\(r^a\). Spatial direction weights are paid by the
rooted exponential reserve, giving the stated \(C_B\).
Together with the retained winding family this proves
\eqref{eq:ordinary-loop-remainder}. All sums were
performed with one external anchor fixed, and all
phase and path bounds were uniform over legal histories.
Consequently neither the terminal cutoff nor the
physical volume enters these constants.
\end{proof}

There is a useful closed choice for the inherited loop
allowance. One may take \(\Lambda_m=1+\log\widehat\beta_m\)
at a fixed pair of profiles. Choose the ratio so that the
constant in \eqref{eq:ordinary-loop-budget} satisfies
\(2C_jL^{-2}\le1/4\). At sufficiently late entry,
\(\Lambda_{m+1}\le2\Lambda_m\) and
\(C_{j,L}\le\Lambda_m/4\). Then the output budget is at
most \(\Lambda_m/2\). These choices precede the physical
volume and any subsequence. More generally the same proof
uses any prescribed subpower envelope with this reserve.
Together with \eqref{eq:ordinary-pre-ratio-vertices}, this
is the local path-budget closure used when the finitely many
requirements of a profile are installed in the slow clock.
It does not compute \(\alpha=-b_0\log L+O(1)\); that
uniform logarithmic evaluation remains a separate task.

\subsection{Comparison on the five auxiliary tori}

We use precisely the reflected tori of
Section~\ref{sec:finite-period-localisation}: the read radius
is \(R\le r_jB\), the integer \(w\) satisfies
\(4R+10\le w\le C_jB\), and the auxiliary sides are
\(n_i=2(w+i)\), \(0\le i\le4\).
Their numbers of fine and coarse integration coordinates are
bounded by a fixed polynomial in \(B\), with constants allowed
to depend on \(L\). For the following analytic estimate put
\begin{equation}
                 r_\beta=\widehat\beta_{\rm f}^{-7/16}.
 \label{eq:ordinary-auxiliary-width}
\end{equation}
This is an analytic chart width, distinct from the integer
\(w\) used in the reflection construction.

Only germs at folded good connections are needed for the
localisation. They have a small periodic gauge representative.
Indeed use the coordinate paths from the physical root to
\(\pi_i(y)\) to define the gauge at the auxiliary vertex \(y\).
The map \(\pi_i\) has the same value on the identified
endpoints, so this is a periodic gauge. Every transformed
link is a based loop filled by at most \(C(1+w+R)^2\)
good physical plaquettes. Lemma~\ref{lem:box-path-comparison}
and the fixed small logarithm chart give
\begin{equation}
       \max_e|\log(V^{i,z})_e^{\rm gauge}|
                       \le C_j(1+B)^2h_{\rm c}=o(r_\beta).
 \label{eq:ordinary-folded-small-chart}
\end{equation}
At such a connection fix this gauge while taking derivatives.
Gauge invariance allows this, so derivatives of a chosen gauge
map are unnecessary. A permitted direction has norm at most
\(C_jh_{\rm c}(1+B)^s=o(r_\beta)\) on the auxiliary torus.
The bounded derivatives of the exponential-coordinate chart
therefore convert its ordered physical derivatives into the
chart derivatives at radius \(r_\beta\), with at most a fixed
polynomial in \(B\). The same argument holds on a neighbourhood
of each folded good connection. It supplies the full derivative
germ, rather than only derivatives tangent to folded fields.

\begin{lemma}[An intact auxiliary chart with the original normalisation]
\label{lem:ordinary-auxiliary-comparison}
On these folded good-field germs, the warning-free ordinary
logarithm at strength one agrees, up to an error smaller than
any fixed inverse power of \(\widehat\beta_{\rm f}\), with
the normalised central fluctuation integral, including the
deterministic compensation for the matched coarse principal
action. The same assertion holds for its first two strength
coefficients and their prescribed physical derivatives.
The comparison keeps all original broad factors, masks,
Gaussian factors and shared tag averages before estimating
their difference. Its constants are independent of cutoff,
physical volume and history length.
\end{lemma}
\begin{proof}
% !TeX root = 4DYM_notes.tex
We separate an exact density identity from a comparison of
domains. The calculation is on one auxiliary torus and one
forest at a time. All dimensions below count real colour
coordinates, with counting measure. Constants can depend on
the fixed adjacent profiles, ratio and derivative order.
They are independent of the original physical volume,
ultraviolet cutoff and length of the legal history.

\emph{The exact identity on a common domain.}
There are no frozen links on the intact domain. Put
\(d=|I|\dim_{\R}\mathfrak g\), and retain the original
common constant
\[
 C_0=e^{vn_f}z_{\rm n}(\beta)^{-b_c}\zeta_G^{|I|}.
\]
An already extracted vacuum scalar is included in \(C_0\).
It is constant as the coarse connection varies. For the
moment keep one branch and its multiplier \(\theta(V,a)\),
including any retained physical chart tests. Write
\[
 \varphi(V,a)=A_t(V,\mathcal U(V,a)),\qquad
 \psi(V,a)=I_f(\mathcal U(V,a))
                              -\sum_{e\in I}\log j_G(a_e).
\]
Thus the physical expression to compare is exactly
\begin{equation}
 C_0\int\theta(V,a)
                   e^{-g^{-2}\varphi(V,a)-\psi(V,a)}\,\dd a.
 \label{eq:auxiliary-one-physical-integral}
\end{equation}
The factor \(\zeta_G^{|I|}\) occurs in \(C_0\), whereas
the nonconstant Haar density occurs in \(\psi\).
Neither is inserted a second time.

Let \(C=C_{\rm sel}\), let \(\mu\) be the selected affine
centre, and let \(\rho\) be the stationary centre. In this
proof \(\mathsf H=D_a^2\varphi(V,\rho)>0\) denotes the
stationary Hessian. The two covariance matrices
\(C\) and \(\mathsf H^{-1}\) need not agree. Define
\[
 \delta=g^{-1}(\mu-\rho),\qquad a=\mu+gz=\rho+g(z+\delta).
\]
Both translations have Lebesgue Jacobian \(g^d\).
The ratio of the stationary Gaussian density at
\(z+\delta\) to the selected density at \(z\) is
\begin{equation}
 \mathcal R(V,z)=
 (\det\mathsf H\det C)^{1/2}
 \exp\left\{-\tfrac12(z+\delta)^T\mathsf H(z+\delta)
                           +\tfrac12z^TC^{-1}z\right\}.
 \label{eq:auxiliary-exact-density-ratio}
\end{equation}
For every absolutely integrable physical test \(F(V,a)\),
including a Borel test, ordinary substitution gives
\[
 \mathbb E_{\mathsf H^{-1}}F(V,\rho+gw)
       =\mathbb E_C[\mathcal R(V,z)F(V,\mu+gz)].
\]
In particular, the argument of a physical test is changed
along with the density. This identity asserts no smallness
of \(\mathcal R-1\).

For completeness set
\[
 \mathcal V(V,w)=g^{-2}\{
        \varphi(V,\rho+gw)-\varphi(V,\rho)
                      -\tfrac12g^2 w^T\mathsf H w\},\qquad
 \Delta\psi(V,w)=\psi(V,\rho+gw)-\psi(V,\rho).
\]
The stationary equation removes the linear term. Expression
\eqref{eq:auxiliary-one-physical-integral} therefore equals
\begin{align}
 &C_0(2\pi)^{d/2}g^d(\det\mathsf H)^{-1/2}
              e^{-g^{-2}\varphi(V,\rho)-\psi(V,\rho)}
       \notag\\
 &\hspace{15mm}\cdot
 \mathbb E_{\mathsf H^{-1}}\left[
  \theta(V,\rho+gw)e^{-\mathcal V(V,w)-\Delta\psi(V,w)}
                         \right].
 \label{eq:auxiliary-stationary-density-expanded}
\end{align}
After \eqref{eq:auxiliary-exact-density-ratio}, its scalar
determinants cancel as
\[
 (\det\mathsf H)^{-1/2}
                    (\det\mathsf H\det C)^{1/2}=(\det C)^{1/2}.
\]
Its entire exponent, including the stationary constant, is
\[
 -g^{-2}\varphi(V,\mu+gz)-\psi(V,\mu+gz)
                                      +\tfrac12z^TC^{-1}z.
\]
This is exactly the selected-coordinate expression obtained
by applying Lemma~\ref{lem:sector-gaussian-change} directly.

One can also check every affine-row factor in
\eqref{eq:exact-selected-Gaussian-identity}. With
\(r=\mathsf M\mu+b\) and the exact phase difference
\(\mathcal D=\varphi-\|Da+\alpha\|^2/2\), its exponent is
\begin{align*}
 &-\frac{\|D\mu+\alpha\|^2}{2g^2}
       -g^{-1}\langle r,z\rangle
       +\tfrac12z^T(C^{-1}-\mathsf M)z
       -g^{-2}\mathcal D(V,\mu+gz)-\psi(V,\mu+gz)\\
 &\hspace{20mm}=
       -g^{-2}\varphi(V,\mu+gz)-\psi(V,\mu+gz)
                                          +\tfrac12z^TC^{-1}z.
\end{align*}
Here \(\psi\) replaces the separate loop and Haar factors
in \(\mathcal B_\sigma\); those factors are not also
included in \(\theta\). The residual \(r\), constant
square, covariance correction and determinant are all
needed in this cancellation.

The normalisations can now be read without an asymptotic
comparison:
\[
\begin{array}{c|c|c}
 &\text{selected coordinates}&\text{stationary coordinates}\\ \hline
 \text{physical point}&\mu+gz&\rho+gw\\
 \text{Lebesgue Jacobian}&g^d&g^d\\
 \text{Gaussian normaliser}&(2\pi)^{d/2}(\det C)^{1/2}
                  &(2\pi)^{d/2}(\det\mathsf H)^{-1/2}\\
 \text{constant phase}&\|D\mu+\alpha\|^2/(2g^2)
                  &g^{-2}\varphi(V,\rho)\\
 \text{Haar factor}&\zeta_G^{|I|}\prod_e j_G(a_e)
                  &\zeta_G^{|I|}\prod_e j_G(a_e).
\end{array}
\]
The two constant phases are accompanied by their respective
exact remainder exponents above. There is no additional
normalisation for a forest: the tree Haar identity is a
probability identity. On the intact domain every local
regular branch selects the same predictor; summing its
original conditional probabilities gives one. The site
frames remain shared in this sum.

\emph{The inner domain and its complement.}
Use the inner protector of
Lemma~\ref{lem:selected-centre-intact-box}, on the full
union of reads of the auxiliary torus. Its width is a
fixed multiple \(\eta_LP\) in each selected coordinate.
The torus has at most \(C_{j,L}(1+B)^{p_*}\) such coordinates,
where \(p_*\) is a fixed integer enlarged below to
dominate the finite polynomial losses.
The stationary curvature term has size \(C_jL^{-2}h_c\).
Choose \(L\) first to make this less than \(h/16\), then
choose \(\eta_L\) for the Gaussian displacement. The
nonlinear and cut errors use at most another \(h/16\)
at entry. All queried fine plaquettes then have chord
less than \(h/4\), and all original broad tests are one
on a smaller inner box. The exact density calculation
above applies there before any factor is removed.

The centre comparison used for this inclusion is
\[
 \|\mu-\rho\|_\infty
       \le C_{j,L}(1+B)^{p_*}(h^2+h e^{-cB}),\qquad
 \frac{\|\delta\|_\infty}{P}
       \le C_{j,L}(1+B)^{p_*}(h+e^{-cB})\longrightarrow0.
\]
It follows either by inserting the free lift in the
stationary equation and applying its inverse, or from
the explicit intact-box calculation in
Lemma~\ref{lem:selected-centre-intact-box}. At fixed
profiles and ratio the Hessian square root and its
inverse have bounded absolute row sums, by the
exponential inverse estimates and functional calculus
already used for the stationary Gaussian. Thus this
translation also compares the inner boxes in the two
coordinate systems. The central Euclidean ball requires
an additional dimension factor. Its radius is
\(D=d_0P\), and the bounded-white-coordinate comparison
uses \(R_{\rm G}=D/(4C_W\sqrt{d_*})\), where \(d_*\)
is the fixed polynomial bound on the auxiliary dimension.
Choose one fixed integer \(p_*\) large enough and a
small fixed \(c_{j,L}>0\), and put
\[
                    \tau_B=\frac{c_{j,L}P}{(1+B)^{p_*}}.
\]
The box of radius \(\tau_B\) then fits in the original
selected inner protector, in the central Euclidean
ball and in the admitted white box after the linear
change of covariance. The preceding centre bound gives
\(\|\delta\|_\infty=o(\tau_B)\). Thus the same statement
holds with either centre. A smooth cutoff equal to one
on half this box and supported in this box can be
inserted together with its complement. This partitions
each original integral and keeps every original test.
The derivatives of the additional cutoff cost powers
of \(P/\tau_B\), hence only powers of \(1+B\).

Consequently a discrepancy between the retained domains
requires \(|z_i|>c\tau_B\) in one of a polynomial number
of selected or stationary white coordinates. Use only
one such witness at a time. If the precision gap is
\(\lambda>0\), then
\[
 1_{\{|z_i|>c\tau_B\}}
       \le e^{-\lambda c^2\tau_B^2/16}
                                      e^{\lambda z_i^2/16}.
\]
This consumes charge \(\lambda/8\) at that coordinate.
The common reserve of
Lemma~\ref{lem:central-complex-precision} leaves room
for it. Corollary~\ref{cor:finite-range-tilted-polynomial}
bounds all fixed-order score polynomials under that
same tilt. Summing the possible witnesses contributes
only a fixed power of \(1+B\).

We keep track also of a possible finite-alphabet or
normalisation bill before absorbing it. Its logarithm
is bounded by
\(C_{j,L}(1+B)^{p_*}(1+\log\widehat\beta_{\rm f})\), since
only auxiliary reads are involved. For every fixed
\(p_*,t\),
\[
 C_{j,L}(1+B)^{p_*}(1+\log\widehat\beta_{\rm f})
       +t\log(1+P)+t\log(1+\Lambda_f)=o(\tau_B^2).
\]
Hence the bill and all polynomial scores can be paid
by half the witness exponent. The resulting error,
including derivatives through order \(q\), is at most
\begin{equation}
 C_{j,L,q}\exp\left(-\frac{c_{j,L,q}P^2}{(1+B)^{2p_*}}\right)
                         =O_{j,L,q,J}(\widehat\beta_{\rm f}^{-J})
                      \quad\text{for every fixed }J.
 \label{eq:auxiliary-common-tail-estimate}
\end{equation}
Indeed
\(\log\{\tau_B^2/\log\widehat\beta_{\rm f}\}
=2(\log\log\widehat\beta_{\rm f})^2
 -(4p_*+1)\log\log\widehat\beta_{\rm f}+O_{j,L}(1)\to\infty\).
This also shows explicitly why the Euclidean dimension
loss does not change the required inverse-power bound.
The first comparison keeps the masks and branch weights
in the physical variables. Lemma~\ref{lem:regular-Borel-measure-derivative}
moves differentiation to the smooth common density;
its score formula and
\eqref{eq:selected-density-partition-formula} account
for all derivative placements. A differentiated smooth
protector is supported in the same tail. No derivative
of a hard fine indicator is used.

\emph{The strength coefficients and the logarithm.}
On the common intact chart the finite ordinary source is
\(\exp(u\sum_\alpha v_\alpha)\). Apply the preceding
change of variables to this literal source. Away from
that chart retain the union factors
\(b_{\cup T_\alpha}\) in
\eqref{eq:ordinary-retained-partition}. For
\(|u|\le3/2\), the central majorants and precision
reserve used above are uniform. The discrepancy of
the two physical source integrals therefore has the
bound \eqref{eq:auxiliary-common-tail-estimate} on
this circle, after every prescribed physical derivative.
Cauchy's formula gives the same bound, up to a fixed
factor, for their first two coefficients. Equivalently,
one can use the two finite moment formulas in
\eqref{eq:ordinary-retained-two-moments}.
The parameter \(u\) multiplies the retained letters;
it is not substituted for \(g\), and no new stationary
graph at a rescaled coupling is assumed.

After the completed quadratic factor and deterministic
coarse compensation have been extracted, let the two
relative integrals be \(Z_1,Z_2\). Their total letter
income on the auxiliary torus is \(C_Bb=o(1)\).
The connected majorant therefore gives
\(|Z_i-1|\le1/4\) at entry, also on the fixed strength
circle. In particular every point of the segment
between them has modulus at least \(1/2\). Thus
\[
 \log Z_1-\log Z_2
       =\int_0^1\frac{Z_1-Z_2}{Z_2+t(Z_1-Z_2)}\,\dd t.
\]
The finite derivative rule for an inverse denominator
and the common score bounds preserve
\eqref{eq:auxiliary-common-tail-estimate}. Flat
subtraction merely doubles it. Finally, changing the
inner stationary domain to the original central ball
uses Lemma~\ref{lem:central-versus-clipped}; its
\(e^{-cR_{\rm G}^2}\) bound has the same inverse-power
property. This proves the lemma with its original
normalisation and derivative germs.
\end{proof}

For clarity about the compensation in this lemma, the actual
central integral, before flat division, has the form
\begin{equation}
 I_D(V)=C_0(2\pi)^{d/2}g^d
       e^{-\beta S_{\rm cl}(V)-T(V)}Z_D(g,V),
 \label{eq:ordinary-central-exact-prefactor}
\end{equation}
where \(d=\dim y\) and \(C_0\) contains the original
\(z_{\rm n}(\beta)^{-b_c}\), \(\zeta_G^{|I|}\), and any
incoming vacuum scalar. The factor \(\det {\mathsf H}\) is in \(T\).
The exact value \(I_D(1)\) supplies the new flat vacuum
scalar. In the matched relative density the compensating
factor is \(e^{(\beta+\alpha)S_{\rm c}+I_c}\).
The \(\beta T_2S_{\rm c}\) part has already been accounted
for by the calibrated reference Gaussian. Its remaining
deterministic logarithmic letters, of sizes
\(C_B\beta h_c^3\), \(C_B|\alpha|h_c^2\), and
\(C_B\Lambda_c h_c^2\), are included in the ordinary
strength convention. They all fit the envelope
\eqref{eq:central-explicit-prices}. Thus the flat-subtracted
ordinary logarithm is, with the sign appropriate to a log
density, the negative of
\begin{align}
 &\beta(S_{\rm cl}-S_{\rm c})
       +T-T(1)-\alpha S_{\rm c}-I_c\notag\\
 &\hspace{35mm}-\{\log Z_D(g,V)-\log Z_D(g,1)\},
 \label{eq:ordinary-matched-central-action}
\end{align}
up to the auxiliary comparison error just bounded.
In particular shifting the coupling uses the full nonlinear
\(S_{\rm c}\), not only its quadratic Taylor polynomial.
Here \(\beta+\alpha\) is the matched reference coupling;
the later affine extraction retains its additional correction
in the coupling ledger. For fixed \(L\), late entry makes
\(|\alpha|/\beta\) small and keeps this reference in the
same comparison band.
No normalising constant has been discarded: field-independent
terms in \eqref{eq:ordinary-central-exact-prefactor} are
recorded in the vacuum scalar before taking this difference.

\begin{theorem}[The ordinary auxiliary-action bound]
\label{thm:ordinary-auxiliary-action}
Choose the classical and loop arrays by
Propositions~\ref{prop:ordinary-classical-matching} and
\ref{prop:ordinary-one-loop-matching}.
Let \(S_i\) be the total first-two-coefficient action on
the \(i\)-th comparison torus, with these deterministic
coarse terms included, and with an optional finite truncation
error at most \(\epsilon_{\rm tr}\) per anchor in the raw
derivative norm. The formulas transferred to the five tori
are the same labelled local formulas.
For every \(0<\nu<3/80\), on folded good-field germs,
\begin{equation}
 \|S_i-S_i(1)\|_{q,h_c,s;0,\Lambda_i}
       \le C_{\nu,j,L}\widehat\beta_{\rm f}^{-19/16+\nu}
                    +2\#\Lambda_i\epsilon_{\rm tr}.
 \label{eq:ordinary-auxiliary-action}
\end{equation}
In particular, after entry and a truncation with
\(\#\Lambda_i\epsilon_{\rm tr}\le
\widehat\beta_{\rm f}^{-23/20}/4\), the right side is
at most \(\widehat\beta_{\rm f}^{-23/20}\).
All entry thresholds are independent of the physical
volume and cutoff.
\end{theorem}
\begin{proof}
Apply the stationary comparison to the total ordinary action
at strength one, after flat subtraction. Its classical part
is bounded by \(C_B\beta r_\beta^5\) per anchor using
\eqref{eq:ordinary-classical-remainder}. Its one-loop part
is bounded by
\(C_B(1+\Lambda_f+|\alpha|)r_\beta^3\), including the
nonlinear coarse compensation. The total winding error is
\(C_Be^{-c_{j,L}B}\), since every auxiliary side is at
least a fixed positive multiple of \(B\).

It remains to retain the interacting Gaussian integral rather
than replace it by its determinant. The results of
Sections~\ref{sec:first-interacting-localisation} and
\ref{sec:full-interacting-remainder} give
\[
  \log Z_D(g,x)=g^2L_2(x)+\mathcal R(g,x),\qquad
  \|\mathcal R\|_{\mathrm{local}}
       \le C_Bg^4H_{R_{\rm G}}^4+e^{-cR_{\rm G}^2},
\]
with the exponentially small clipped-coefficient error
included. The polynomial \(L_2\) is the full connected
Wick expression of \eqref{eq:explicit-second-connected-coefficient},
including the cubic--cubic, cubic--loop, and linear--linear
contractions. Its analytic rooted bound is
\(C(1+\Lambda_f)^2\). Subtract its flat value and apply
the mean-value formula on the chart to obtain
\[
 \|g^2(L_2(x)-L_2(0))\|^{\rm chart}_{q,r_\beta,s;\mathrm{local}}
                  \le C_Bg^2(1+\Lambda_f)^2r_\beta.
\]
The fixed-order derivatives obey the same bound: every
nonconstant monomial supplies at least one power of
\(r_\beta\) in the scaled chart norm. Retaining the full
flat value of \(\log Z_D\) is essential here. A bound for
\(g^2L_2\) without flat subtraction would only give order
\(g^2\), which is insufficient for the required allowance.
We can thus keep the entire nonconstant interacting term
as a remainder in this argument; a further extraction of
its quadratic jet is unnecessary for this auxiliary bound.

Finally Proposition~\ref{prop:central-two-coefficients}
changes strength one to the first two coefficients with
error \(C_Bb^3\) per anchor. Flat subtraction doubles
this bound. Summing the rooted estimates over an auxiliary
torus and using Lemma~\ref{lem:ordinary-auxiliary-comparison}
gives the following complete list of errors, up to fixed
polynomial factors in \(B\):
\begin{align}
 \#\Lambda_i\{&\beta r_\beta^5
       +(1+\Lambda_f+|\alpha|)r_\beta^3
       +g^2(1+\Lambda_f)^2r_\beta
       +g^4H_{R_{\rm G}}^4+b^3\}\notag\\
 &\quad+C_Be^{-c_{j,L}B}
             +C_Be^{-cR_{\rm G}^2}
             +2\#\Lambda_i\epsilon_{\rm tr}.
 \label{eq:ordinary-auxiliary-error-ledger}
\end{align}
This is a sum of total-action bounds, not a claim that its
first two coefficient families are individually small.
In particular there is no need to assign the classical or
one-loop contribution to one particular strength coefficient.
Their sum is identified through the exact strength-one
integral, with the higher-strength error estimated above.

The coordinate and physical derivative norms are related by
\eqref{eq:ordinary-folded-small-chart} and the fixed-gauge
argument following it. The extra powers of \(B\), the
auxiliary volumes, \(P\), the loop envelope, and
\(H_{R_{\rm G}}\) are all subpower in \(\widehat\beta_{\rm f}\).
The five algebraic terms in braces have respective rates
\[
 \widehat\beta_{\rm f}^{-19/16},\quad
 \widehat\beta_{\rm f}^{-21/16+o(1)},\quad
 \widehat\beta_{\rm f}^{-23/16+o(1)},\quad
 \widehat\beta_{\rm f}^{-2+o(1)},\quad
 \widehat\beta_{\rm f}^{-3/2+o(1)}.
\]
Both exponential terms are smaller than any fixed inverse
power. Absorb all subpower losses in the prescribed \(\nu\).
This proves \eqref{eq:ordinary-auxiliary-action}.
Since \(19/16-23/20=3/80\), a fixed
\(\nu<3/80\) leaves a positive margin to absorb the fixed
constant at entry. Every bound used the fixed profile and
the polynomial auxiliary size, so that entry is uniform
in the physical volume and in the cutoff.
\end{proof}

\subsection{The consequence for finite-period localisation}

To apply the last theorem to literal finite stencils, truncate
the native path and row labels of the first two coefficient
family at total length \(C_jB\). Their exponential budgets,
with a spare width for the physical derivatives, bound the
discarded rooted family by \(C_Be^{-c_jB}\). Coefficients
of the remaining countable family can then be approximated
by a finite subfamily with any prescribed
\(\epsilon_{\rm tr}\); absolute convergence gives this
simultaneously for the five transfers. These choices use
the cover labels and their budgets before periodic evaluation.
They may depend on the cutoff's input arrays but do not depend
on the physical volume. Identical footprints are still
different labels. The resulting formulas have a common
complete read radius \(R\le r_jB\), as required in
Section~\ref{sec:finite-period-localisation}.

\begin{corollary}[A small ordinary regular contribution]
\label{cor:ordinary-finite-period-regular}
The exact finite-period redistribution of this truncated
first-two-coefficient family gives functions \(G_z\) with
\begin{equation}
       \|G_z\|_{q,h_c,s;z,N_e(z)}
                         \le C_j\widehat\beta_{\rm f}^{-23/20}.
 \label{eq:ordinary-finite-period-regular}
\end{equation}
Their footprints have bounded covering size at the fixed
profile. After the exact local reset their normalised
regular family is smaller than any prescribed fixed
fraction of \(\widehat\beta_c^{-11/10}\), at late entry.
The flat scalar, the extracted kinetic term, and the mask
discrepancies remain in the exact density identity.
\end{corollary}
\begin{proof}
Use the same \(\lambda_i\), fold maps and \(G_z\) as in
\eqref{eq:finite-period-fold-norm}. That lemma is stated
with a bound on the whole good auxiliary domain. Its proof
only uses the auxiliary derivative germs at the folded
connections that occur in \(G_z\). The bound just proved
on precisely these germs therefore gives the same conclusion
by that proof, with no additional auxiliary volume factor.
The factors \(\#\Lambda_i\) have already been paid in
\eqref{eq:ordinary-auxiliary-error-ledger}.

As \(e\le C_jB\), the number \(n_B(N_e(z))\) is bounded
by a fixed profile constant. Each anchor has only the fixed
number of transferred functions, so its covering weight
changes \eqref{eq:ordinary-finite-period-regular} by a
fixed factor. Lemma~\ref{lem:local-reset-norms} bounds
flat-value and affine-Hessian subtraction by a constant
independent of this footprint. The enlargement by the
plaquette halo has a fixed covering cost.
Comparability of \(\widehat\beta_f,\widehat\beta_c\) then
gives the stated allowance, since
\(23/20-11/10=1/20>0\).

The redistribution identity still contains the rough
functions \(L_x\) of
\eqref{eq:finite-period-mask-price}. The two-coefficient
raw price is at most \(C_Bb\), so their price is at most
\(C_Bb\); their vanishing on healthy \(N_{2e}(x)\) is
unchanged. The additional reset discrepancies are exactly
those of Lemma~\ref{lem:exact-masked-subtraction}.
Neither kind of discrepancy is removed when the regular
functions are normalised.
\end{proof}

The ordinary terms of degree at least three already have
price \(C_Bb^3\), and the incoming part has the contraction
and comparison estimates of
Section~\ref{sec:incoming-sharp-transport}. The present
section supplies the missing small auxiliary total action
for the first two ordinary coefficients. Assembling these
families with the retained mask discrepancies and forced
packets, and evaluating the uniform logarithmic drift of
\(\alpha\), are subsequent parts of the full nonlinear
iteration. The result here is not a completed construction
of the continuum hierarchy.
\clearpage
\section{The one-loop coefficient with an error uniform in the block ratio}
\label{sec:one-loop-drift}

We now evaluate the coefficient \(\alpha\) defined in
\eqref{eq:ordinary-alpha-definition}. A bound depending on a
fixed ratio would not suffice: the ratio is subsequently chosen
to make the relative drift error small. We therefore prove a
bounded error before choosing that ratio. The calculation uses
the same fine leading action, stationary branch and free lift
as Section~\ref{sec:ordinary-matching}.

The organisation into a covariant determinant comparison and
two elementary one-loop contractions follows
\cite[Lemmas~12.7--12.9, PDF pp.~82--85]{ArmstrongDraft}.
The interpretation of the quadratic coefficient as the unique
gauge-invariant kinetic coefficient is the background-field
organisation of \cite[Sect.~5, printed pp.~292--297]{Bal87}.
We prove the determinant identities, estimates and contractions
needed here. In particular no value of a beta function is an
input to the calculation.

\subsection{The metric fixes the numerical coefficient}

Let \((e_a)\) be an orthonormal basis of \(\mathfrak g\) in
\eqref{eq:physical-lie-norm}, and write
\(\operatorname{ad}(X)Y=[X,Y]\).
Define \(C_G\) by
\begin{equation}
 -\operatorname{Tr}_{\mathfrak g}
       (\operatorname{ad}(X)\operatorname{ad}(Y))
                      =C_G\langle X,Y\rangle_{\mathfrak g}.
 \label{eq:drift-casimir}
\end{equation}
This constant exists because the bilinear form on the left
is invariant. Its self-adjoint representing operator has
invariant eigenspaces, hence ideals; simplicity makes it
scalar. The adjoint operators are skew adjoint in the trace
metric. Thus the form is nonnegative, and it is positive
on a nonzero \(X\), since a simple nonabelian Lie algebra
has zero centre. In particular \(C_G>0\).
Equivalently
\[
 C_G=-\frac1{\dim\mathfrak g}
               \sum_a\operatorname{Tr}_{\mathfrak g}
                                  (\operatorname{ad}(e_a)^2).
\]
The coefficient denoted by \(b_0\) in
\eqref{eq:hat-beta} will be
\begin{equation}
                         b_0=\frac{11C_G}{24\pi^2}.
 \label{eq:drift-b0}
\end{equation}
It is the slope of the inverse coupling multiplying
\(S_f\), with density \(e^{-\beta S_f}\).
All factors below refer to this convention and to the
Taylor factor \(1/2!\) in \(Y\).

\subsection{Replacing the axial determinant by local operators}

To estimate traces uniformly, keep the local gauge directions
instead of solving for each one in axial coordinates.
Temporarily adjoin one group frame per coarse root, with
probability Haar measure. Use the convention in which it
compensates the root gauge transformation of the predictor.
The joint action is invariant under all fine-site gauge
transformations. Fixing these frames and the block trees
recovers the original integral, by the tree Haar identity.
There is no additional gauge volume in its physical
normalisation.

At its stationary configuration let \(J\) be the Hessian
in joint left coordinates, without its factor \(\beta\).
Let \(R\phi\) insert a fine infinitesimal gauge transformation,
where \(\phi\) is a \(\mathfrak g\)-valued site field.
Let \(Q\) embed the axial free coordinates. Stationarity
and gauge invariance give \(JR=0\).
Choose the regular local frame already constructed in
Section~\ref{sec:actual-stationary}, and use the centred
frame coordinates of Section~\ref{sec:ordinary-matching}.
Hessians and predictor second differentials are computed
in physical joint left increments. Centring changes the
tangent basis used for their estimates. At a stationary
point the full Hessian transforms by congruence, since
the full gradient is zero.
Its fine gauge frame \(\Psi\) transforms by
\(\Psi_z\mapsto\gamma_z\Psi_z\) under a joint fine gauge
change. For example, before centring it is the inverse
root-normalised local frame followed by the compensating
coarse root frame. Define \(X\) to be its left logarithmic
differential. Then
\begin{equation}
                              XR=I.
 \label{eq:drift-frame-right-inverse}
\end{equation}
This follows by differentiating the displayed transformation
law. Centring is a change of coordinates and preserves it.

For a connection \(U\), put
\[
 (d_U\phi)_e=\phi_{e^-}-\operatorname{Ad}(U_e)\phi_{e^+},
       \qquad \Delta_U=d_U^*d_U.
\]
Adjoints use counting measure and the fixed colour metric.
Set \(\mu=L^{-1}\), and introduce the positive regulator
\begin{equation}
 P_\mu=\mu^2(\Delta_U+\mu^2I)^{-1},\qquad
 m_\mu=\mu^2P_\mu^9.
 \label{eq:drift-regulator}
\end{equation}
Here \(\mu\) is an auxiliary inverse block length, not a
coupling or an ultraviolet lattice spacing.
Let \(P\) project joint increments to fine links. Define
\begin{equation}
 F_\mu=d_U^*P+m_\mu X,\qquad
 C_\mu=J+F_\mu^*F_\mu,\qquad
 G_\mu=\Delta_U+m_\mu.
 \label{eq:drift-completed-operators}
\end{equation}
Thus \(F_\mu R=G_\mu\). Adding \(F_\mu^*F_\mu\) is
\emph{gauge completion}: it makes the Hessian positive in
the gauge directions as well as on its physical slice.
The compensating determinant of \(G_\mu\), an operator on
site gauge parameters, is usually called the ghost
determinant. We use ordinary finite matrices throughout.

\begin{lemma}[The completed determinant identity]
\label{lem:drift-completed-determinant}
On the regular stationary branch, the one-loop action of
\eqref{eq:ordinary-one-loop-action} satisfies
\begin{align}
 T(V)-T(1)={}&I_f(U_{\rho(V)})-I_f(1)
       +\tfrac12\log\frac{\det C_\mu(V)}{\det C_\mu(1)}
       -\log\frac{\det G_\mu(V)}{\det G_\mu(1)}.
 \label{eq:drift-completed-determinant}
\end{align}
In particular the left side is independent of the auxiliary
regulator. The original flat scalar remains the one fixed
by \eqref{eq:ordinary-central-exact-prefactor}.
\end{lemma}
\begin{proof}
Use the gauge and slice columns \((R,Q)\). Ordering tree
vertices from their roots makes the gauge-coordinate block
triangular with diagonal entries \(\pm I\); slice coordinates
are unchanged free left increments. Orthogonal adjoint
changes have determinant of absolute value one. Hence
\(\det(R,Q)^2=1\). In these columns the completed form is
\[
 \begin{pmatrix}G_\mu&F_\mu Q\\0&I\end{pmatrix}^{*}
 \begin{pmatrix}I&0\\0&Q^*JQ\end{pmatrix}
 \begin{pmatrix}G_\mu&F_\mu Q\\0&I\end{pmatrix}.
\]
Consequently
\(\det C_\mu=(\det G_\mu)^2\det(Q^*JQ)\).
If \(s\) changes the original free exponential coordinates
to these left slice coordinates, then
\({\mathsf H}=(Ds)^TQ^*JQ\,Ds\), while the Haar chart ratio is
\(|\det Ds|\). Taking the determinant cancels this
Jacobian against the Haar logarithm in \(T\).
All determinants are positive near the flat field, so
their real logarithms and local analytic continuations
give \eqref{eq:drift-completed-determinant}.
Field-independent coordinate constants cancel in the flat
difference; their vacuum contribution is unchanged.
\end{proof}

\subsection{Block energy norms and the trace reserve}

We record quantitatively the norms in this determinant
comparison. At a real coarse momentum \(\xi\), let
\(k_z=(\xi+2\pi z)/L\) be its centred fine aliases and put
\[
 \omega_z^2=\mu^2+\sum_a|e^{ik_{z,a}}-1|^2,\qquad
 \|u\|_{\rm bl}^2=\sum_z\omega_z^2|u_z|^2.
\]
This is the block norm used in
\eqref{eq:ordinary-block-quotient-energy}. The alias
transform is unitary, and the coarse frame has its
Euclidean norm. Let \(\Omega\) be multiplication by
\(\omega_z\) on fine components and by one on frames.
For a form \(A\) and its inverse the relevant operator
norms are respectively
\[
 \|\Omega^{-1}A\Omega^{-1}\|_{\rm op},\qquad
                         \|\Omega A^{-1}\Omega\|_{\rm op}.
\]
The trace norm of a matrix is
\(\|B\|_1=\operatorname{Tr}(B^*B)^{1/2}\).
For a form use \(\|\Omega^{-1}A\Omega^{-1}\|_1\).
Between different momentum fibres use the output and
input weights in the corresponding positions. These
weighted norms are used only to estimate the finite
matrices; they do not change any activity norm.

All field insertions in the rest of the comparison are
normalised coarse directions lifted by the smooth free
map \(E\). By \eqref{eq:all-scale-linear-difference-bounds},
\[
 \|\nabla_f^aEv\|_\infty\le
 C_j\mu^{1+a}\|v\|_\infty,\qquad 0\le a\le2.
\]
For a quadratic jet it suffices to use this tangent:
every invariant scalar has zero first field derivative
at the flat field, so the second-derivative chain rule
adds no second derivative of the stationary curve.

The auxiliary fine form is
\(A_\mu=A_f(U)+d_Ud_U^*+m_\mu^{(1)}\), where \(A_f\)
is the left Hessian of the unmasked fine leading action,
excluding the predictor and inherited loop action.
The superscript \((1)\) means that \(m_\mu\) acts on
each of the four link directions. All matrices and
field derivatives in the next lemma are evaluated at
the flat field; its strip refers to momentum variables.
Trace norms there are on one Bloch fibre.

\begin{lemma}[Uniform inverse and insertion bounds]
\label{lem:drift-block-trace-bounds}
Fix the input stage bounds, including its path width,
and take \(L\) above the classical floor in
Proposition~\ref{prop:ordinary-classical-matching}.
There are \(\eta_j>0\), \({\mathsf M}_j<\infty\), chosen before
\(L\), with the following properties on a coarse momentum
strip of width \(\eta_j\). The flat completed inverse
and the auxiliary fine inverse have block dual-to-primal
norm at most \({\mathsf M}_j\). The auxiliary fine form and the
contact-removed completed form through two normalised
field insertions have norm bounded by a fixed polynomial
in \({\mathsf M}_j\). Moreover
\begin{equation}
 \|P_\mu^3\|_1\le {\mathsf M}_j,\qquad
       \|D^rP_\mu\|_{\rm op}\le {\mathsf M}_j\quad(0\le r\le2).
 \label{eq:drift-regulator-trace-bounds}
\end{equation}
For the geometric maps, through these same insertions,
\begin{equation}
 \|D^r(H\Omega^{-1})\|+
 \|D^r(d_U^*P\Omega^{-1})\|+
 \|\mu^2D^r(X\Omega^{-1})\|\le {\mathsf M}_j,
                         \qquad 0\le r\le2.
 \label{eq:drift-geometric-insertions}
\end{equation}
Here \(H\) is the differential of the framed logarithmic
predictor discrepancy. The assertion also holds for the
contact-removed completed form used below.
The constants are uniform in cutoff, period and history.
\end{lemma}
\begin{proof}
At zero field write \(\lambda(k)=\sum_a|e^{ik_a}-1|^2\).
For \(x\ge0\),
\begin{equation}
 2^{-10}(x+\mu^2)\le
 x+\mu^2\left(\frac{\mu^2}{x+\mu^2}\right)^9
                                      \le x+\mu^2.
 \label{eq:drift-mass-floor}
\end{equation}
If \(x\ge\mu^2\), use the first term for the lower bound;
otherwise the factor in parentheses is at least \(1/2\).
The identity
\(\|d_1u\|^2+\|d_0^*u\|^2=\lambda(k)\|u\|^2\)
and the fixed positive bounds of \(H_f\) therefore make
the fine completed form coercive in the block norm.

For the joint form, Section~\ref{sec:ordinary-matching}
already proves coercivity on the quotient by joint gauge
directions. Choose a representative \(w_\perp\) bounded
by that quotient energy and write \(w=w_\perp+R\phi\).
The alias averaging bounds give
\[
 \|R\phi\|_{\rm bl}\le C_j\|G_\mu\phi\|_2,
       \qquad \|F_\mu w\|_2\le C_j\|w\|_{\rm bl}.
\]
For the first inequality, a fine gradient has weighted
size at most \(\omega_z^2|\phi_z|\); its averaged frame
component obeys the same estimate by the summable alias
filter. The mass floor controls these sizes by \(G_\mu\).
For the second inequality the divergence is bounded by
the energy norm and \(m_\mu X\) is controlled by the
frame estimate proved in the next paragraph. Since
\(G_\mu\phi=F_\mu w-F_\mu w_\perp\), these inequalities
give
\[
 \|w\|_{\rm bl}\le
 C_j\bigl(\langle w,Jw\rangle^{1/2}+\|F_\mu w\|_2\bigr).
\]
This is the required coercivity of \(C_\mu\).

Here is the quantitative frame estimate. At the flat
field the centred local frame equations act as identity
on the mean-zero increments. Their inhomogeneous term
is a normalised average of route sums. Every route has
length at most \(C_jL\). The absolute row sum of its
linear operator is at most \(C_jL\). Its column sum has
the same bound: for a fixed edge in a four-dimensional
patch, at most \(C_jL^5\) ordered endpoint pairs have
a coordinate-ordered route through that edge, whereas
the endpoint average contributes \(C_jL^{-4}\).
To check the count, for an edge in direction \(a\), the
endpoint coordinates before \(a\) leave \(a-1\) free
coordinates, those after \(a\) leave \(4-a\), and the
two coordinates in direction \(a\) leave at most two.
The total exponent is five. Finite route permutations
and the fixed normalised averaging rules change only
the constant. The Schur estimate therefore bounds the
fine-link part of \(X\) on counting \(\ell^2\) by
\(C_jL\). Its coarse-frame embedding has norm \(L^2\).
As \(\|u\|_2\le L\|u\|_{\rm bl}\),
\[
                         \|Xw\|_2\le C_jL^2\|w\|_{\rm bl}.
\]
One or two smooth field insertions in a route add at
most \(C_jL\) possible positions per insertion and a
factor \(C_jL^{-1}\) at each position. Differentiating
the frame equations adds normalised averages and the
bounded mean-zero inverse. Thus the same bound holds
through two insertions. This proves the last term of
\eqref{eq:drift-geometric-insertions}.
The predictor differential is another normalised route
average, and its block dual norm follows by the same
count and the averaging bounds. A derivative of
\(d_U^*\) inserts a smooth field of size \(C_j\mu\);
the inverse block weight is at most \(\mu^{-1}\).
These observations prove the other terms there.

% !TeX root = 4DYM_notes.tex
We spell out the insertion bounds in these units. Let
\(u,w\) be arbitrary fine increments with block norm
at most one, and let \(v,v'\) be normalised smooth
lifted coarse increments. Then
\[
 \|u\|_2,\|w\|_2\le\mu^{-1},\quad
 \|\nabla_fu\|_2,\|\nabla_fw\|_2\le C,\quad
 \|v\|_\infty\le C_j\mu,\quad
                         \|\nabla_fv\|_\infty\le C_j\mu^2.
\]
The cubic Wilson terms are sums of contractions of
one difference with a commutator. Their two relevant
placements, before summing direction and colour, obey
\begin{align*}
 \left|\sum_x\langle\nabla_fu,[v,w]\rangle\right|
 &\le C_G\|\nabla_fu\|_2\|v\|_\infty\|w\|_2\le C_j,\\
 \left|\sum_x\langle\nabla_fv,[u,w]\rangle\right|
 &\le C_G\|\nabla_fv\|_\infty\|u\|_2\|w\|_2\le C_j.
\end{align*}
Here \(|[x,y]|\le C_G|x||y|\) follows from continuity
of the bracket in the fixed finite-dimensional algebra.
The quartic terms have two commutators. For example
\[
 \left|\sum_x\langle[v,u],[v',w]\rangle\right|
       \le C_G^2\|v\|_\infty\|v'\|_\infty
                                      \|u\|_2\|w\|_2\le C_j.
\]
Permuting the two arbitrary legs gives the other
placements, with only a finite combinatorial factor.
Transporting a leg along a path preserves its colour
norm; differentiation of the transport inserts a
smooth field. The native exponential path budget
pays its path-position factors. Thus these are
operator bounds for the forms, rather than bounds
obtained by counting fine sites.

For the fine-minus-Wilson vertices use the stronger
weighted estimate
\eqref{eq:ordinary-pre-ratio-vertices}. Its additional
momentum gives the insertion price \(\mu A_j'\).
The classical ratio floor makes this at most a fixed
constant. This is where dependence on the incoming
path amplitude is used; it is not put into a numerator
that is later required to precede \(L\).

The regulator estimates use the same scale cancellation.
Put \(\mathcal A=\Delta_U+\mu^2I\), and evaluate field
derivatives at the flat connection. The left gauge
gradient satisfies
\[
 \|D_vd_U\|_{\rm op}\le C_j\mu,\qquad
 \|D_vD_{v'}d_U\|_{\rm op}\le C_j\mu^2.
\]
Indeed a first derivative of
\(\operatorname{Ad}(U_e)\phi_{e^+}\) is
\([v_e,\operatorname{Ad}(U_e)\phi_{e^+}]\), and the
second is a sum of two nested brackets. These act
pointwise on links; incidence multiplicities are
bounded independently of the number of links.
Also
\(\|d_U\mathcal A^{-1/2}\|\le1\) and
\(\|\mathcal A^{-1/2}\|\le\mu^{-1}\).
The exact derivatives
\begin{align*}
 D_v\Delta_U&=(D_vd_U)^*d_U+d_U^*D_vd_U,\\
 D_vD_{v'}\Delta_U
   &=(D_vD_{v'}d_U)^*d_U+d_U^*D_vD_{v'}d_U\\
   &\quad +(D_vd_U)^*D_{v'}d_U+(D_{v'}d_U)^*D_vd_U
\end{align*}
therefore give
\[
 \|\mathcal A^{-1/2}D_v\Delta_U\mathcal A^{-1/2}\|
 +\|\mathcal A^{-1/2}D_vD_{v'}\Delta_U\mathcal A^{-1/2}\|
                                                       \le C_j.
\]
For example a product of the two first derivatives
costs \((C_j\mu\cdot\mu^{-1})^2\), and a second
derivative next to \(d_U\) costs at most
\(C_j\mu^2\cdot\mu^{-1}\). No inverse power of
\(\mu\) remains.

Writing \(R_\Delta=\mathcal A^{-1}\), differentiation gives
\begin{align*}
 D_vP_\mu&=-\mu^2R_\Delta(D_v\Delta_U)R_\Delta,\\
 D_vD_{v'}P_\mu
 &=\mu^2R_\Delta(D_v\Delta_U)R_\Delta(D_{v'}\Delta_U)R_\Delta\\
 &\quad+\mu^2R_\Delta(D_{v'}\Delta_U)R_\Delta(D_v\Delta_U)R_\Delta
       -\mu^2R_\Delta(D_vD_{v'}\Delta_U)R_\Delta.
\end{align*}
Insert \(\mathcal A^{-1/2}\) between the form factors.
The two outside factors \(\mu\mathcal A^{-1/2}\)
have norm at most one. The preceding relative bounds
then prove \(\|D^rP_\mu\|_{\rm op}\le C_j\) for
\(r=0,1,2\). All these estimates concern one real
momentum fibre or compatible shifted fibres. Their
complex-strip versions follow by the fixed relative
perturbation argument below.

At zero field,
\[
 \frac{\mu^2}{\lambda(k_z)+\mu^2}
       \le C(1+|z|^2)^{-1},\qquad
 \|P_\mu^3\|_1\le
 C\dim\mathfrak g\sum_{z\in\mathbb Z^4}(1+|z|^2)^{-3}<\infty.
\]
This is the trace bound in
\eqref{eq:drift-regulator-trace-bounds}.
Finally extend these estimates to a smaller complex
coarse strip. Freeze weights at real momenta; shifts
of a fine momentum are \(O(\mu|\operatorname{Im}\xi|)\),
so the weighted symbol changes by at most
\(C_j|\operatorname{Im}\xi|\). The free and geometric
symbols have the bounds just proved there, by their
alias formulas and normalised finite words. Choose
the strip so that a Neumann series retains half the
real coercivity. Weights for fibres shifted by a
bounded external coarse momentum are comparable,
since every \(\omega_z\ge\mu\).
All estimates therefore also apply between the shifted
fibres of a field insertion. The contact-removed form
is the sum of the fine form, the predictor differential
square and the completed gauge square; the preceding
geometric bounds give its two insertion bounds directly.
At zero its inverse equals the full completed inverse.
The inverse differentiation identities supply its
inverse insertion bounds as needed. This proves the lemma.
\end{proof}

\subsection{The contact term and the fine determinant comparison}

The Hessian of the predictor penalty has two parts. The
first is the square of the predictor differential. The
second multiplies its second differential by the first
derivative of the penalty. Call the latter the contact
term. In joint left coordinates write it as
\(\sum_\lambda d^\lambda(v)C_\lambda(v)\), where \(v\)
is the coarse field, \(\lambda\) consists of a coarse
link and a colour, and \(C_\lambda\) is the second
differential of the framed logarithmic discrepancy.
The source \(d\) includes the fixed noise coefficient.
Let \(C_\mu^{\rm rem}\) be \(C_\mu\) with this term
removed and let \(R_\mu^{\rm rem}=(C_\mu^{\rm rem})^{-1}\).

\begin{lemma}[The contact term has zero kinetic coefficient]
\label{lem:drift-contact-cancellation}
Removing this contact term leaves the value of \(\alpha\)
unchanged.
\end{lemma}
\begin{proof}
The free normal equation gives
\begin{equation}
 d(0)=0,\qquad Dd(0)=-K_c,\qquad K_c=d_1^*H_cd_1.
 \label{eq:drift-contact-source}
\end{equation}
Thus its first field jet has momentum order two.
To retain all chain-rule terms, use an independent
source \(u\) and set
\[
 C(u,v)=\sum_\lambda u^\lambda C_\lambda(v),\quad
 f(v,u)=\operatorname{Tr}(R_\mu^{\rm rem}(v)C(u,v)),\quad
 f_2(v,u)=\operatorname{Tr}[(R_\mu^{\rm rem}(v)C(u,v))^2].
\]
The logarithmic difference, through field degree two,
is exactly
\begin{equation}
 \tfrac12 f(v,d(v))-\tfrac14 f_2(v,d(v))+O(|v|^3).
 \label{eq:drift-contact-log-expansion}
\end{equation}
At zero the covector
\(y_\lambda(v)=\operatorname{Tr}(R_\mu^{\rm rem}(v)C_\lambda(v))\)
vanishes, by adjoint invariance and simplicity.
It transforms covariantly in its coarse source index:
the discrepancy transforms by conjugation at its initial
root, and changes in the internal fine indices cancel
inside the trace. The regular stationary branch has the
same covariance. Applying a coarse gauge change to the
flat connection therefore gives
\(Dy(0)d_0\phi=0\).

Let \(b_{i\lambda}\) be the Fourier kernel of
\(\partial_{v_i}y_\lambda(0)\). The last identity says
\(\sum_a q_a(-p)b_{(a,c),\lambda}(p)=0\) for each colour
\(c\). On \(p=te_a\), only the \(a\)-th coefficient
remains. It is zero for \(t\ne0\), and continuity gives
\(b_{(a,c),\lambda}(0)=0\). Thus \(b(p)=O(|p|)\).
These are analytic kernels: the inverse, local
discrepancy derivatives and trace bounds of the previous
lemma justify their sums and the finite-period limits.

For completeness the mixed derivative defining \(b\)
keeps both terms
\[
 \partial_{v_i}y_\lambda
 =\operatorname{Tr}\{-R_\mu^{\rm rem}
       (\partial_{v_i}C_\mu^{\rm rem})R_\mu^{\rm rem}C_\lambda
                   +R_\mu^{\rm rem}\partial_{v_i}C_\lambda\}.
\]
In \eqref{eq:drift-contact-log-expansion}, the possible
term \(y(0)D^2d(0)\) is zero. The remaining quadratic
terms from \(f\) pair \(Dd(0)\) with \(b\), in the two
orders of the field directions. They have momentum
order at least three. The terms from \(f_2\) contain
two factors \(Dd(0)\) and an analytic coefficient, so
have order at least four. They have no second momentum
coefficient on the tested transverse ray. This proves
the assertion, including the derivatives of the
contact insertion itself.
\end{proof}

Using the auxiliary fine form defined above, put
\begin{equation}
 \Pi_\mu(U)=\tfrac12\log\det A_\mu-\log\det G_\mu.
 \label{eq:drift-fine-determinant}
\end{equation}
Let \(Y_\mu\) be the infinite-volume
quadratic Taylor kernel of this scalar. Define
\(\alpha_b(Y_\mu)\) by the same extraction as in
\eqref{eq:ordinary-alpha-definition}, now in fine units.

\begin{proposition}[A comparison uniform in the ratio]
\label{prop:drift-block-to-fine}
There is \(d_j^{\rm cmp}<\infty\), fixed before \(L\),
such that
\begin{equation}
                 |\alpha-\alpha_b(Y_{1/L})|
                                      \le d_j^{\rm cmp}.
 \label{eq:drift-block-to-fine}
\end{equation}
It is uniform in cutoff, period, actual coupling and
finite history satisfying the fixed stage bounds.
\end{proposition}
\begin{proof}
Work first with finite matrices and commensurate external
momenta. Put both vector operators on the same joint
space. Let \(Z\) project to coarse frames, and adjoin
their identity to the fine operator:
\(A_0=P^*A_\mu P+Z^*Z\). Its determinant is \(\det A_\mu\).
Let \(A_1=C_\mu^{\rm rem}\). If \(B\) is the coarse
penalty Hessian, \(H\) its logarithmic differential,
\(E=d_U^*P\), and \(K=m_\mu X\), their exact difference is
\begin{align}
 \Theta=A_1-A_0={}&H^*BH-Z^*Z+E^*K+K^*E+K^*K
                                  -P^*m_\mu^{(1)}P.
 \label{eq:drift-trace-class-difference}
\end{align}
% !TeX root = 4DYM_notes.tex
The fine-action Hessians cancel before any estimate.
We detail both derivative placements and the trace
bound for the remaining terms. This expands the
comparison in \cite[Lemma~12.8, PDF pp.~82--84]{ArmstrongDraft}.
The use of weighted operator bounds is the same
finite-matrix method as in \cite[Sect.~2]{DST23}.

Freeze the block weights at a real internal momentum.
Every product below uses the output and input weights
of its respective fibres; an external coarse momentum
shifts a fine momentum by \(O(\mu)\). Their weights
are consequently comparable by a constant fixed
before \(L\). Adjoint signs in the formulas below
refer to real momenta; analytic continuation uses
the corresponding opposite-momentum symbols.

The factors \(H^*BH\) and \(Z^*Z\) pass through
coarse spaces of dimensions \(4\dim\mathfrak g\)
and \(\dim\mathfrak g\). If \(h=H\Omega^{-1}\),
the full two-insertion derivative of the first is
\begin{align*}
 D_vD_w(h^*Bh)={}&h_{vw}^*Bh+h^*B_{vw}h+h^*Bh_{vw}\\
 &+h_v^*B_wh+h_w^*B_vh+h_v^*Bh_w+h_w^*Bh_v\\
 &+h^*B_vh_w+h^*B_wh_v.
\end{align*}
Subscripts mean derivatives of that factor only.
In each term the intermediate coarse space is the
same fixed-dimensional space. If the norms of its
factors and their two derivatives are at most \(A_j\),
\[
 \|D^r(h^*Bh)\|_1
       \le 9(4\dim\mathfrak g)\max(1,A_j)^3,
                                                  \qquad 0\le r\le2.
\]
The constant projection \(Z\) has zero field
derivatives; its zeroth-order trace norm is
\(\dim\mathfrak g\). Thus these terms introduce
no fine-residue multiplicity.

For a regulator factor put \(Q=P_\mu^3\), so
\(P_\mu^9=QQQ\). The bounds just proved give
\(\|Q\|_1\le C_j\) and \(\|D^rQ\|_{\rm op}\le C_j\)
through order two. The latter also follows directly
by differentiating its three factors. Explicitly,
\begin{align*}
 D_vD_w(QQQ)={}&Q_{vw}QQ+QQ_{vw}Q+QQQ_{vw}\\
 &+Q_vQ_wQ+Q_wQ_vQ+Q_vQQ_w+Q_wQQ_v\\
 &+QQ_vQ_w+QQ_wQ_v.
\end{align*}
Every term has at least one undifferentiated \(Q\).
Put that factor in trace norm and the other two
in operator norm. The ideal inequality
\(\|ABC\|_1\le\|A\|\|B\|_1\|C\|\) then proves
\(\|D^rP_\mu^9\|_1\le C_j\), with the same argument
on the smaller complex strip. For reference its
uninserted trace is bounded by the convergent sum
\[
 \sum_z\left|\frac{\mu^2}{\lambda(k_z)+\mu^2}\right|^3
                   \le C_j\sum_{z\in\Z^4}(1+|z|^2)^{-3}.
\]
This is a sum over aliases, not its supremum times
\(L^4\).

The placement of the weights can be checked without
commuting them through a differentiated matrix. Put
\[
 e=E\Omega^{-1},\qquad x=\mu^2X\Omega^{-1},\qquad
                         b_a=\mu J_aP\Omega^{-1},
\]
where \(J_a\) selects fine direction \(a\). These
are only names for the existing factors. Their
operator norms through two insertions are bounded;
the \(b_a\) are field independent. The regulator
parts of \(\Omega^{-1}\Theta\Omega^{-1}\) are
\[
 e^*P_\mu^9x+x^*(P_\mu^9)^*e
       +x^*(P_\mu^9)^*P_\mu^9x
                         -\sum_a b_a^*P_\mu^9b_a.
\]
For \(e^*P_\mu^9x\), distribute zero, one or two
derivatives among its three factors. Differentiated
outer factors are bounded operators, while every
middle derivative is trace class by the preceding
nine-term calculation. In the third term there are
six cubes; two derivatives again leave an untouched
cube in trace norm. In the last term
\(\|\mu P\Omega^{-1}\|\le1\); the two factors of
\(\mu\) pay the two inverse fine weights exactly.
These statements also apply to shifted fibres.

This proves the complete list of scale prices:
\[
\begin{array}{c|c|c}
 \text{term}&\text{trace-class factor}&
                                      \text{remaining scale price}\\ \hline
 H^*BH&\text{coarse-link space}&1\\
 Z^*Z&\text{coarse-frame space}&1\\
 E^*K,\ K^*E&\text{one untouched cube}&
                                      \|\mu^2X\Omega^{-1}\|\\
 K^*K&\text{one untouched cube}&
                                      \|\mu^2X\Omega^{-1}\|^2\\
 P^*m_\mu^{(1)}P&\text{one untouched cube}&
                                      \|\mu P\Omega^{-1}\|^2.
\end{array}
\]
All entries in the last column are uniformly bounded.
Adding the finitely many product-rule terms gives
\begin{equation}
 \|\Omega^{-1}D^r\Theta\Omega^{-1}\|_1\le C_j,
                                                   \qquad 0\le r\le2.
 \label{eq:drift-difference-trace-bounds}
\end{equation}
Every constant here is chosen from the fixed profile
data before \(L\).

There is a shorter way to organise the remaining
logarithmic differentiation. In the weighted coordinates
write \(\bar A_i=\Omega^{-1}A_i\Omega^{-1}\) and
\(\bar R_i=\Omega A_i^{-1}\Omega\). Let \(F_*,A_*,B_*\ge1\)
bound, respectively, these inverses, the first two
insertions of \(\bar A_i\), and the three trace norms
in \eqref{eq:drift-difference-trace-bounds}. These are
fixed-profile constants. The inverse identity gives
\[
 \|\bar R_1-\bar R_0\|_1\le F_*^2B_*=:D_*.
\]
For a normalised direction \(v\), define
\(S_i(v)=\bar R_iD_v\bar A_i\). Then
\begin{align*}
 S_1(v)-S_0(v)
    &=(\bar R_1-\bar R_0)D_v\bar A_1
                              +\bar R_0D_v(\bar A_1-\bar A_0),\\
 \|S_1(v)-S_0(v)\|_1&\le D_*A_*+F_*B_*=:E_*,
               \qquad \|S_i(v)\|\le F_*A_*.
\end{align*}
The same \(E_*\) bounds
\(\bar R_1D_{vw}\bar A_1-\bar R_0D_{vw}\bar A_0\).
For the quadratic product, use just
\[
 S_1(v)S_1(w)-S_0(v)S_0(w)
   =[S_1(v)-S_0(v)]S_1(w)+S_0(v)[S_1(w)-S_0(w)].
\]
It follows that
\[
 \left|D_vD_w(\log\det A_1-\log\det A_0)\right|
                                  \le E_*(1+2F_*A_*).
\]
Each trace has one trace-class factor and bounded
remaining factors. This identity avoids estimating
the two extensive logarithms separately. The common
ghost determinant cancels exactly.

Take the trace per coarse cell. Its internal coarse
momentum sum is a normalised average, and all its
fine aliases have already been summed in
\eqref{eq:drift-difference-trace-bounds}. The resulting
difference kernel has a bound \(C_j'\) on a fixed
external coarse strip \(|\operatorname{Im}p|<\eta_j'\).
Passing from commensurate momenta to this kernel uses
the absolutely summable native kernels and normalised
finite-period sums; their uniform bounds justify the
limit. Cauchy's coefficient formula on
\(|t|=\eta_j'/2\) bounds twice its real \(t^2\)
coefficient by \(8C_j'/(\eta_j')^2\).
The contact lemma has removed no such coefficient.
The inherited loop action has a link kernel of order
four, so its pullback has zero \(\alpha\), as proved
in Proposition~\ref{prop:ordinary-one-loop-matching}.

It remains to relate fine and coarse units in the
auxiliary scalar \(\Pi_\mu\). It is gauge invariant
and has zero first jet. Its quadratic kernel has
momentum order at least two. Nonprincipal transverse
aliases of the free lift vanish to order at least two
in each external leg, and contribute no quadratic
momentum coefficient. The principal alias is, modulo
fine gauge terms,
\[
                    L^4L^{-2}Y_\mu(p/L)+O(|p|^3).
\]
The factors are cell volume and \(L^{-1}\) per
interpolated field. Extraction of \(p^2\) contributes
one more \(L^{-2}\), so their product is one.
Thus folding changes no \(\alpha_b(Y_\mu)\).
This proves \eqref{eq:drift-block-to-fine}, with a
constant chosen before the ratio.
\end{proof}

\subsection{The leading vertices and their remainder bounds}

It remains to calculate \(\alpha_b(Y_\mu)\).
Use fine internal momentum \(k\in[-\pi,\pi]^4\), external
momentum \(p=te_1\), and external connection in direction
2 and unit colour \(e_b\). Put \(r=|k|\), the Euclidean
norm. The flat vector and ghost inverses will both be
denoted by \(R_\mu\) in estimates that apply to either;
their matrix types specify which is meant in a product.

We first justify the leading expressions. The weight-four
fine action, with connection and coordinate difference
each of weight one, is
\(\tfrac14\sum_{i,j}|F_{ij}|^2\), where
\[
 F_{ij}=\partial_iA_j-\partial_jA_i+[A_i,A_j],\qquad
                    D_i=\partial_i+\operatorname{ad}(A_i).
\]
These are polynomials used to read the lowest momentum
coefficients of the lattice words. They introduce no
continuum approximation to the integral. Indeed the
word expansion and \(H_f(0)=I\) in
Proposition~\ref{prop:ordinary-classical-matching} prove
that all other degree-\(a\) terms have momentum order
at least \(5-a\), \(a=2,3,4\).

Varying the displayed polynomial gives
\(\delta F_{ij}=D_i a_j-D_j a_i\) and
\(\delta^2F_{ij}[a,c]=[a_i,c_j]+[c_i,a_j]\).
Using the invariant metric and summation by parts,
and adding the divergence square, its vector Hessian
is
\[
                  -D^2\delta_{ij}-2\operatorname{ad}(F_{ij}).
\]
The ghost operator has leading part \(-D^2\).
On the chosen transverse external field, their first
vertices are consequently
\begin{align}
 V_g(t,k)&=-2ik_2\operatorname{ad}(e_b),\notag\\
 V_v(t,k)&=\{-2ik_2I_4+2itS\}\operatorname{ad}(e_b),
       \qquad S=E_{21}-E_{12}.
 \label{eq:drift-leading-vertices}
\end{align}
Here \(E_{ij}\) is the matrix unit in direction indices.
For example the orbital term is the first variation
of \(-D^2\); \(\partial_2A_2=0\) on this ray, leaving
\(-2ik_2\operatorname{ad}(e_b)\).
The other term follows from
\(F_{12}=it e_b\). At the second vertex replace \(t\)
by \(-t\) and \(k\) by \(k+te_1\); its orbital term
is unchanged.

\begin{lemma}[Integrable errors in the fine one-loop integral]
\label{lem:drift-integrand-bounds}
There are constants \(A_j,D_j<\infty\), fixed before
\(\mu\in(0,1]\), such that the actual quadratic
coefficient has the representation
\begin{equation}
 \alpha_b(Y_\mu)=(2\pi)^{-4}
              \int_{[-\pi,\pi]^4}\operatorname{Re}k_{\mu,b}(k)\,dk,
 \label{eq:drift-integral-representation}
\end{equation}
where
\begin{align}
 |k_{\mu,b}(k)|&\le A_j(r+\mu)^{-4},\notag\\
 |k_{\mu,b}(k)-h_b(k)|&\le
            D_j\{r^{-3}+\mu^2r^{-6}\},\qquad r\ge\mu,
 \label{eq:drift-integrand-bounds}
\end{align}
and \(h_b\) is the contraction of the leading vertices
and massless leading inverses computed below.
All constants are uniform over the fixed fine-stage
path budgets and histories.
\end{lemma}
\begin{proof}
At fixed \(\mu>0\), the positive flat operators and the
analytic fine vertices make the logarithmic Hessians
ordinary differentiable finite momentum integrals.
The bounds below supply integrable majorants for their
derivatives. Their finite-period Fourier sums converge
to these integrals; this is the definition of the cover
kernel, without a physical volume subsequence.

The mass floor and the fixed ellipticity of \(H_f\) give
\[
 \|\partial_k^aR_\mu(k)\|\le C_j(r+\mu)^{-2-a},
                                      \qquad 0\le a\le2.
\]
One can obtain the derivatives either by the inverse
identities or by a complex disc of radius \(c_j(r+\mu)\)
around a real momentum. On that disc the quadratic
operator remains invertible relative to its value at
the centre. For \(r\ge\mu\), replacing its unregulated
symbol by \(r^2I\) makes an error \(O_j(r^3)\), with
derivatives of order \(a\) bounded by \(C_jr^{3-a}\).
The regulator and its momentum derivatives cost at most
\(C_j\mu^2r^{-a}\). The resolvent identity and product
rule therefore give
\begin{equation}
 \|\partial_k^a(R_\mu-r^{-2}I)\|
             \le C_j\{r^{-1-a}+\mu^2r^{-4-a}\},
                                     \qquad 0\le a\le2.
 \label{eq:drift-propagator-replacement}
\end{equation}
For the ghost the lattice error is even smaller; this
common bound suffices.

Here are the corresponding vertex estimates. If \(V\)
is a first field vertex and \(V^{\rm lead}\) is
\eqref{eq:drift-leading-vertices}, then at \(p=0\),
\begin{align}
 \|\partial_p^aV\|&\le C_j(r+\mu)^{1-a},\notag\\
 \|\partial_p^a(V-V^{\rm lead})\|
       &\le C_j\{r^{2-a}+\mu^2r^{-1-a}\},
           \quad r\ge\mu,\quad 0\le a\le2.
 \label{eq:drift-vertex-replacement}
\end{align}
The unregulated error has momentum order two, as just
derived from the cubic fine-minus-Wilson vertex and
the expansions of the lattice gradient and gauge square.
Taylor's formula on the fixed fine strip gives the
three displayed derivative bounds for that error.
For the regulator, differentiate
\(P_\mu=\mu^2(\Delta_U+\mu^2I)^{-1}\).
A first variation of \(\Delta_U\) has size
\(C_j(r+|p|)\); each external momentum derivative of
an inverse adds one factor \((r+\mu)^{-1}\).
Thus a first field variation of \(m_\mu\) has bound
\(C_j\mu^2(r+\mu)^{-1}\), and its first two momentum
derivatives have the corresponding additional inverse
powers. The product rule for its ninth power changes
only constants. This proves
\eqref{eq:drift-vertex-replacement}, including the
region \(r\le\mu\) in the first line.

In the logarithmic determinant, a bubble means the
term with two first field vertices and two inverses.
The vector weight is \(-1/4\), the ghost weight is
\(+1/2\); these follow from
\(\tfrac12\log\det A_\mu-\log\det G_\mu\) and the
quadratic term of the logarithm. The extraction of
\(\alpha_b\) subsequently multiplies by two.
Two external momentum derivatives distributed over
the two vertices and inverses give at worst
\((r+\mu)^{2-a}(r+\mu)^{-4-(2-a)}
=(r+\mu)^{-4}\). This proves the first bound for
the bubbles. On \(r\ge\mu\), telescope the product,
replacing its factors and their derivatives using
\eqref{eq:drift-propagator-replacement} and
\eqref{eq:drift-vertex-replacement}.
An unregulated replacement gains one power of \(r\)
relative to degree \(-4\); a regulator replacement
gives \(\mu^2/r^2\). This proves the second bound
for the bubbles.

The other terms have one second field vertex and one
inverse; these are often called tadpoles. Their
quadratic external-momentum coefficients are bounded
by
\begin{equation}
          C_j\{(r+\mu)^{-2}+\mu^2(r+\mu)^{-6}\}.
 \label{eq:drift-tadpole-bound}
\end{equation}
Indeed the unregulated second vertex has bounded
second external derivative, by the quartic path
budget. The regulator's two field variations and
two external derivatives give at most
\(C_j\mu^2(r+\mu)^{-4}\) before multiplication by
the inverse. The leading unregulated second vertex
has no external momentum, so contributes nothing to
\(h_b\). As \(r\le2\pi\) and \(\mu\le1\),
\eqref{eq:drift-tadpole-bound} is at most a constant
times \((r+\mu)^{-4}\). On \(r\ge\mu\) it is at most
\(C_j(2\pi r^{-3}+\mu^2r^{-6})\).
This proves both bounds also for these terms. Summing
the finite direction and colour indices gives
\eqref{eq:drift-integrand-bounds}.
\end{proof}

\subsection{The logarithmic contraction}

The logarithmic part can now be calculated without
the regulator or the higher lattice vertices.
By \eqref{eq:drift-casimir},
\(\operatorname{Tr}_{\mathfrak g}\operatorname{ad}(e_b)^2=-C_G\).
Since \(\operatorname{Tr}S=0\) and
\(\operatorname{Tr}S^2=-2\), the direction traces
of the paired vertices are
\[
                 4(-2ik_2)^2-8t^2\quad\hbox{(vector)},
             \qquad (-2ik_2)^2\quad\hbox{(ghost)}.
\]
The mixed orbital--spin terms have zero trace.
Include the colour trace, the two massless inverses,
the bubble weights, and the factor two from
\(\alpha_b\). Before taking the \(t^2\) coefficient
the result is
\begin{equation}
            -\frac{4C_G(k_2^2+t^2)}{r^2|k+te_1|^2}.
 \label{eq:drift-leading-bubble}
\end{equation}
Expanding the denominator gives
\[
 \frac1{r^2|k+te_1|^2}
   =r^{-4}-2tk_1r^{-6}
         +t^2(4k_1^2r^{-8}-r^{-6})+O(t^3).
\]
Consequently
\begin{equation}
 h_b(k)=4C_G\left(\frac{k_2^2}{r^6}
                  -\frac{4k_1^2k_2^2}{r^8}
                  -\frac1{r^4}\right).
 \label{eq:drift-homogeneous-pole}
\end{equation}
This expression is independent of the chosen unit
colour and is homogeneous of degree \(-4\).

For completeness its angular averages follow directly
from Euclidean rotations. On the unit sphere in
\(\mathbb R^4\), \(\mathbb E\theta_2^2=1/4\).
Write \(a=\mathbb E\theta_1^4\),
\(c=\mathbb E\theta_1^2\theta_2^2\).
Rotation by \(\pi/4\) in the first two coordinates
gives \(a=3c\), and
\((\sum_i\theta_i^2)^2=1\) gives \(4a+12c=1\).
Thus \(c=1/24\). The sphere has area \(2\pi^2\),
as follows, for example, by evaluating
\(\int_{\mathbb R^4}e^{-|x|^2}\,dx=\pi^2\) in polar
coordinates. Hence
\begin{align}
 (2\pi)^{-4}\int_{\mu<|k|\le1}h_b(k)\,dk
  &=\frac{8\pi^2C_G}{(2\pi)^4}
                \left(\frac14-\frac16-1\right)\log\mu^{-1}\notag\\
  &=-\frac{11C_G}{24\pi^2}\log\mu^{-1}.
 \label{eq:drift-pole-integral}
\end{align}

\begin{theorem}[Uniform logarithmic one-loop drift]
\label{thm:uniform-one-loop-drift}
For every fixed adjacent profile pair there is
\(d_j<\infty\), chosen before its blocking ratio,
such that the coefficient constructed in
Section~\ref{sec:ordinary-matching} satisfies
\begin{equation}
       |\alpha_{n,m}+b_0\log L_m|\le d_j,
                  \qquad b_0=\frac{11C_G}{24\pi^2}.
 \label{eq:uniform-one-loop-drift}
\end{equation}
The same constant works for every admissible ratio
above the profile's floor, every cutoff and finite
history, every physical period, and every actual
coupling in the stated band. The arrays and the
coefficient are chosen before the physical side length.
\end{theorem}
\begin{proof}
Split the integral of
Lemma~\ref{lem:drift-integrand-bounds} into the ball
\(r\le\mu\), the shell \(\mu<r\le1\), and the
remaining part of the cube. The first part has bound
\(A_j\mu^{-4}\cdot\pi^2\mu^4/2=A_j\pi^2/2\).
On the shell the two errors have respective bounds
\[
 2\pi^2D_j\int_\mu^1dr\le2\pi^2D_j,
 \qquad
 2\pi^2D_j\mu^2\int_\mu^1r^{-3}\,dr\le\pi^2D_j.
\]
Outside the unit ball the bound is
\(A_j(2\pi)^4\). Therefore
\[
 |\alpha_b(Y_\mu)+b_0\log\mu^{-1}|
 \le(2\pi)^{-4}
       \{A_j\pi^2/2+3\pi^2D_j+A_j(2\pi)^4\}.
\]
Combine this with
Proposition~\ref{prop:drift-block-to-fine} and set
\(\mu=L^{-1}\). Their sum is a possible \(d_j\).
Every constant on the right was fixed from the fine
profile, free-symbol bounds, and exponential path
budgets before \(L\). This proves the required
uniformity in the ratio and the history.

The physical finite-period winding terms have already
been retained as activities in
Proposition~\ref{prop:ordinary-one-loop-matching}.
They do not change the definition of \(\alpha_{n,m}\).
Their bounds may depend on the subsequently fixed
\(L\), and are made small by its common period reserve
and entry threshold. Thus no volume-dependent coupling
coefficient is introduced in
\eqref{eq:uniform-one-loop-drift}.
\end{proof}

In particular the ratio can now be chosen so that
\(d_j/\log L\) is as small as prescribed. Postponing
the coupling entry only pays constants from the
fixed-ratio remainder estimates. It is not used to
replace the uniform coefficient estimate just proved.
\clearpage
\section{Closing the nonlinear step and restricting the coupling window}
\label{sec:nonlinear-induction-step}

The preceding calculations now supply the analytic inputs
for a repeatable nonlinear step. We collect their prices
before applying the exact reset. The regular price after
the ordinary cancellations and the Borel price needed
for forced reassembly are different quantities. Keeping
both is necessary for the pointwise density identity.

The organisation of the exact density bookkeeping follows
\cite[Lemma~13.12, PDF pp.~107--108]{ArmstrongDraft}.
The activity allowances and calendar choices below are derived here
from the preceding estimates.

For level \(m\), retain the calendar quantities
\[
 \widehat\beta_m=\bar\beta+b_0s_m,\quad
 P_m=\exp\{(\log\log\widehat\beta_m)^2\},\quad
 B_m=\lceil(\log\widehat\beta_m)^2\rceil,\quad
 h_m=\widehat\beta_m^{-1/2}P_m.
\]
The two activity allowances and the loop allowance will be
\begin{equation}
 \delta_m=\widehat\beta_m^{-11/10},\qquad
 w_m=e^{-P_m^{1/5}},\qquad
 \Lambda_m=2P_m^{\theta_{J(m)}},
 \label{eq:induction-allowances}
\end{equation}
where \(0<\theta_j\le1/j\) is nonincreasing.
These are deterministic functions of the common calendar,
not of an actual running coupling or a volume subsequence.
All preceding estimates apply to this subpower loop
envelope. The norms remain
\(\mathcal E\) and \(\mathcal K\) of
Section~\ref{sec:activity-norms}.
At level \(m\) their derivative order and spatial
power are those of profile \(J(m)\). Write
\(\kappa(m)=\kappa_{J(m)}\) for its covering exponent.

In this section \(S_m\) denotes the full clipped and
masked principal action. Its unmasked small-field branch
is \(S_m^{\rm un}\); this is the action called \(S_c\)
in the stationary matching formulas of
Section~\ref{sec:ordinary-matching}.

An admissible input here means an exact effective
representation of the form already defined, with its
prescribed symmetries, masks, path bounds and regular
normalisations, and with
\(\|E_f\|_{\mathcal E,\kappa_f}\le\delta_{m+1}\),
\(\|K_f\|_{\mathcal K,\kappa_f}\le w_{m+1}\).
Its derivative order is at least \(\max(q_c,3)\) and
its spatial power is at least \(s_c\ge2\).
The stage orders can all be chosen at least three, so
this condition entails no repeated loss of derivatives.
The finite histories considered below start from the
actual Wilson density; the inductive assertion is that
its effective representations remain admissible.
The present section proves the propagation step and the
coupling-window argument. Its explicit starting
representation is proved in
Section~\ref{sec:wilson-start-shooting}.

At a step put \(j=J(m+1)\). Its coarse index is either
\(j\) or \(j-1\). All constants carrying \(j\) may be
replaced by the maximum over these two profile pairs.
Choose a fixed covering reserve \(a\ge8\) for the
following argument. A larger fixed reserve, if required
by the finite derivative order, has the same proof.
The compression estimate
\eqref{eq:incoming-compressed-incidence}, applied with
target exponent \(\kappa_c+a\), gives the central
incoming income using the original positive fine
exponent. It does not require increasing that exponent
at every repetition of a profile.

\subsection{The remaining warning components}

The physical packets have already been integrated with
their Haar ratios and normalisations in
Theorem~\ref{thm:actual-joint-packet-payment}.
We also need the Gaussian warnings that have no physical
packet. The following argument supplies their regular
derivatives without using a packet energy.

\begin{lemma}[Rare components and their derivative price]
\label{lem:induction-rare-components}
At fixed profiles and ratio, all unforced logarithmic
clusters containing a physical packet or a Gaussian
warning have raw regular norm at most
\begin{equation}
 C_B\exp\{-P_f^2/C_B\}
             \le C_{j,L,{\mathsf M}}\widehat\beta_f^{-{\mathsf M}}
                         \quad\hbox{for every fixed }{\mathsf M}.
 \label{eq:induction-rare-regular}
\end{equation}
The raw forced component family, before the healthy
exclusions and reset restorations are merged, has
incidence price
\begin{equation}
                         a_f\le C_Be^{-c_{\rm pkt}P_f^2}.
 \label{eq:induction-raw-forced}
\end{equation}
Here \(c_{\rm pkt}>0\) depends only on the fixed profiles
and ratio. These estimates keep the original masks, frozen Haar
integrations, tag averages and shared Gaussian law.
\end{lemma}
\begin{proof}
For components containing a physical packet, use the
paid derivative estimate
\eqref{eq:regular-hole-paid-derivatives} and its marked
connected-log bound. The combined smooth and incoming
letter income is \(b+d\), with
\(b=\widehat\beta_f^{-1/2+o(1)}\) and
\(d=\widehat\beta_f^{o(1)}\delta_f\) from
\eqref{eq:central-explicit-prices}. The converted
incoming forced income is at most
\(C_{j,L}B^p w_f\), by the same footprint compression
and incidence count, with no derivatives required of
a forced multiplier. Thus the smallness hypotheses
of the connected majorant and of
Corollary~\ref{cor:regular-hole-family} hold after
entry. Packet counts are summed using
\eqref{eq:packet-covering-income}.

For a component with Gaussian warnings and no packet,
retain its physical Borel reads and use
Lemma~\ref{lem:regular-Borel-measure-derivative}.
All derivatives act on the common density, smooth
rows or smooth test factors. A derivative of a warning
factor is supported on the same warning annulus,
with its threshold decreased by a fixed factor if
necessary. Hard physical masks are held fixed.
The witness construction of
Lemma~\ref{lem:finite-range-failure-witness} therefore
still applies. Its multiplicity is polynomial in
\(B\), and its threshold is \(P_f\) divided by at
most a fixed polynomial in \(B\). After spending
the common precision reserve, its price is
\(C_B\exp(-P_f^2/C_B)\).

The differentiated density contributes polynomial
scores and inverse-Jacobian factors. The score bounds
in Lemma~\ref{lem:regular-selected-score-bounds}
give at most a fixed power of
\(\widehat\beta_f,P_f,B\) and of the read covering
number for a fixed derivative order. Tilted Gaussian
moments are bounded under the same retained positive
precision. Any such power is absorbed by decreasing
the warning exponent by a fixed factor, since
\(P_f^2/(B^p\log\widehat\beta_f)\to\infty\) for every
fixed \(p\). The product-rule polynomial in the
number of letters is paid by the unused covering
reserve, as in
Section~\ref{sec:integrated-regular-derivatives}.
There is no new inverse-coupling price per repeated
incoming derivative. Marking the species that contain
a warning and applying
Lemma~\ref{lem:marked-interacting-log} proves
\eqref{eq:induction-rare-regular}.

Finally a nonempty coarse forced inventory contains
an actual coarse bad site. Such a site is a physical
seed in the sector construction and belongs to a
physical packet. Apply the packet estimate with that
packet marked, and the ordinary small incomes to all
other species. This retains the stronger price
\(e^{-c_{\rm pkt}P_f^2}\) of a physical packet and gives
\eqref{eq:induction-raw-forced}.
The extra bounded coarse compensation after matching
has price \(C_{j,L}h_c\) per affected anchor, by the
clipped principal-action bounds. It is absorbed by
the packet reserve on bad regions; its good-region
part is already an ordinary letter. These uses of
the packet bound include its frozen-link factor
\(s_g^{-n_{\rm fr}}\). Dropping this ratio would not
give the asserted component price.
\end{proof}

\subsection{The small regular family before the final reset}

Write \(H_m\) for the anchored raw derivative price of
the preliminary regular family at exponent
\(\kappa_c+a\), after the finite-period ordinary
redistribution and before flat and affine subtraction.
This is the quantity in
\eqref{eq:preliminary-reset-budget}, with additional
fixed reserve. The input allowance \(\delta_f\) below
is \(\delta_{m+1}\); \(\delta_c=\delta_m\).

\begin{proposition}[A preparatory family with a pre-ratio principal price]
\label{prop:induction-preparatory-family}
The exact physical expansion gives a preliminary
regular family with
\begin{equation}
 H_m\le\frac{V_j}{L}\delta_f+r_{j,L}(\widehat\beta_c),
               \qquad r_{j,L}(x)=o(x^{-11/10}).
 \label{eq:induction-preparatory-price}
\end{equation}
The coefficient \(V_j\) is chosen before \(L\).
At fixed \(L\), the little-oh estimate and its entry
are uniform in cutoff, volume and history. The
remaining Borel corrections and healthy-log
restorations have incidence price at most
\begin{equation}
 b_m^{\rm rest}\le C_B\{b+d+H_m/h_c+|\alpha|h_c\}
                                  =\widehat\beta_c^{-1/2+o(1)}.
 \label{eq:induction-restoration-price}
\end{equation}
The exact scalar and all bad-region mask terms are
retained in this expansion.
\end{proposition}
\begin{proof}
Separate the connected logarithm by strengths as in
Section~\ref{sec:central-strength-expansion}.
For the ordinary coefficients of degree at least
three, Theorem~\ref{thm:central-joint-strength-bound}
and Proposition~\ref{prop:central-two-coefficients}
give raw price \(C_Bb^3\). The first two coefficients
give the local functions of
Corollary~\ref{cor:ordinary-finite-period-regular},
with price \(C_B\widehat\beta_f^{-23/20}\).
Use common finite truncations for their five
transfers and retain their exponentially small
tails. These are local regular functions on the
physical good footprints; the folded auxiliary
germ estimate is sufficient for this conclusion.

For the incoming part, the one-insertion term has
the raw anchored estimate of
Theorem~\ref{thm:sharp-incoming-transport}, applied
with the present fixed output reserve. Its price is
\(V_jL^{-1}\delta_f\). The original masked insertion
differs from it by the superpolynomial family in
\eqref{eq:incoming-rare-cutoff-comparison}.
The higher comparison terms have price
\(C_B(bd+d^2)\) by
Proposition~\ref{prop:central-inherited-comparison}.
There are no missing mixed incoming coefficients.
All remaining marked clusters are bounded in
Lemma~\ref{lem:induction-rare-components}.
Consequently a possible upper bound for \(r_{j,L}\)
is
\begin{align}
 C_B\{&\widehat\beta_f^{-23/20}+b^3+bd+d^2\}
       +C_{j,L,{\mathsf M}}\widehat\beta_f^{-{\mathsf M}}(1+\delta_f),
 \label{eq:induction-new-remainders}
\end{align}
where \({\mathsf M}\) can be any fixed sufficiently large
number. The comparison of neighbouring calendar
couplings and all their subpower factors is uniform
at fixed \(L\). Each term in
\eqref{eq:induction-new-remainders} is therefore
\(o(\delta_c)\), proving
\eqref{eq:induction-preparatory-price}.

We also check what was not put into this small
regular price. Before cancellation, the ordinary
healthy logarithmic species have incidence income
at most \(C_B(b+d)\). They are the species used
for exclusions from a selected forced footprint in
\eqref{eq:physical-exact-restoration}.
The finite-period discrepancy \(L_x\) has the
rough bound \eqref{eq:finite-period-mask-price},
also at most \(C_B(b+d)\), and is zero when its
declared neighbourhood is healthy. Flat and affine
subtraction contribute \(C_BH_m\) in the same way.
The principal-action transfer has precisely the
Borel price of \eqref{eq:principal-transfer-forced-price},
bounded here by \(C_B(H_m/h_c+|\alpha|h_c)\).
This latter term also bounds any such transfer
made at the preliminary matching stage; when the
whole \(\alpha S_c\) has already been matched it
is unnecessary but harmless in the upper bound.
The relaxed-cutoff radius \(R_{h_c}\) is at most
\(B_c\) at entry, so its read halo is included in
these covering prices.

This proves \eqref{eq:induction-restoration-price}.
In fact \(H_m=O_j(\delta_c)\), hence
\(H_m/h_c=\widehat\beta_c^{-3/5+o(1)}\), whereas
\(b=\widehat\beta_c^{-1/2+o(1)}\).
The coefficient \(\alpha\) is fixed once \(L\)
is fixed. All these incomes tend to zero in the
connected-family criterion. The restorations are
kept with their selection rule: every restoration
from a healthy-log exclusion meets an original
forced support. A correction that is zero on its
healthy footprint meets such a support whenever
it is nonzero, because the original inventories
cover every coarse bad site. Thus their eventual
merger contains a forced label in every component,
as required in the exact reassembly theorems.
\end{proof}

\subsection{The final reset has a constant chosen before the ratio}

\begin{lemma}[Uniform price of the principal transfer]
\label{lem:induction-reset-constant}
There is \(C_j^{\rm reset}<\infty\), fixed before \(L\),
such that a preliminary family of price \(H_m\) gives
\begin{equation}
 \|E_c\|_{\mathcal E,\kappa_c}\le C_j^{\rm reset}H_m,
                  \qquad h_c^2|\Delta_m|\le H_m.
 \label{eq:induction-reset-price}
\end{equation}
The reference coupling before this reset is
\(\beta_f+\alpha_{n,m}\). After the reset it is
\(\beta_c=\beta_f+\alpha_{n,m}+\Delta_m\).
The scalar and forced corrections obey the exact
identities of Section~\ref{sec:coupling-reset}.
\end{lemma}
\begin{proof}
The ordinary matching has already compensated the full
\(\alpha S_c\); apply the exact reset with this reference
coupling and with no additional \(\alpha S_{0,h}\) term.
For each preliminary function, Lemma~\ref{lem:local-reset-norms}
bounds its flat and affine subtraction by a footprint-independent
constant. Lemma~\ref{lem:anchor-incidence} pays the added
halo and the required incidence sums using the reserve
in \(H_m\). Isotropy and
Proposition~\ref{prop:isotropic-reset-coefficient} give
one scalar \(\Delta_m\), with
\(h_c^2|\Delta_m|\le H_m\).

The remaining transfer from its masked clover term to
\(\Delta_m S_c\) is exactly
Lemma~\ref{lem:exact-principal-transfer}, now with its
parameter \(\alpha\) set to zero. The first estimate
of Lemma~\ref{lem:path-reset-bounds} is \(C_jH_m\).
Its constant depends on the coarse leading path width,
budget, clip coefficients, relaxed-cutoff parameter,
derivative order and output exponent. All these coarse
bounds and scalar parameters were chosen before \(L\) in
Proposition~\ref{prop:ordinary-classical-matching}.
In particular the unscaled loop allowance does not
enter this principal transfer. Adding the two regular
prices defines \(C_j^{\rm reset}\) before \(L\).
The accompanying Borel terms retain the price in
\eqref{eq:induction-restoration-price}; the scalar is
the sum of the extracted flat values. Thus neither
has been assigned to a regular error or discarded.
\end{proof}

Set, for example,
\begin{equation}
                 c_j=\{8(1+C_j^{\rm reset})\}^{-1}.
 \label{eq:induction-reset-reserve}
\end{equation}
Choose the odd ratio so that \(V_j/L\le c_j/4\).
After this choice, increase entry until
\(r_{j,L}(\widehat\beta_c)\le c_j\delta_c/4\).
Since \(\delta_f\le\delta_c\), the preparatory bound
gives \(H_m\le c_j\delta_c/2\). Consequently
\begin{equation}
 \|E_c\|_{\mathcal E,\kappa_c}\le\delta_c/16,
 \qquad
 |\Delta_m|\le\tfrac12c_j
                    \widehat\beta_c^{-1/10}P_c^{-2}\longrightarrow0.
 \label{eq:induction-regular-contraction}
\end{equation}
The spare fraction allows a further finite allocation
of truncation tails within the same entry threshold.
The essential point is the order of choices: the
coefficient multiplying the inherited \(\delta_f\)
is fixed before \(L\); only the new remainder is made
small by later entry.

\subsection{The forced allowance and the pointwise density}

\begin{proposition}[Forced closure with all restorations]
\label{prop:induction-forced-closure}
At sufficiently late entry the physical forced family,
the healthy exclusions, the finite-period mask
discrepancies and the final reset corrections merge
to a family satisfying
\begin{equation}
                         \|K_c\|_{\mathcal K,\kappa_c}
                                           \le3w_c/4.
 \label{eq:induction-forced-closure}
\end{equation}
The merger preserves the density at every coarse
configuration, including its bad regions.
\end{proposition}
\begin{proof}
Use \(a_f\) from \eqref{eq:induction-raw-forced} and
the combined restoration income
\(b_m^{\rm rest}\) from
\eqref{eq:induction-restoration-price}.
At fixed \(L\), the neighbouring calendar quantities
satisfy \(P_f/P_c\to1\), \(B_f/B_c\to1\), and
\(h_f/h_c\to1\). It follows that
\[
        a_f/w_c\le C_B
               \exp\{-c_{\rm pkt}P_f^2+P_c^{1/5}\}\longrightarrow0.
\]
Also \(b_m^{\rm rest}\le\widehat\beta_c^{-1/8}\)
and \(w_c\le\widehat\beta_c^{-1/16}\) at entry.
Indeed \(P_c^{1/5}/\log\widehat\beta_c\to\infty\).
Thus \(a_f\le w_c/8\), and all the numerical
conditions of Lemma~\ref{lem:reserve-choice} hold
with its polynomial reserve
\(P=\widehat\beta_c^{1/16}\). This auxiliary reserve
in that lemma is distinct from the monitor \(P_c\).

Apply the resulting numerical estimate to the bound
\eqref{eq:physical-merger-income}. It is the same
bound used in Lemma~\ref{lem:reserve-choice}.
Each merged component has at least one original
forced label. Exclusion restorations are selected
only when meeting such a label. Every other nonzero
mask discrepancy meets a covered coarse flag and
therefore also meets a forced label. Thus the
mandatory forced factor in that estimate is present
even for the restorations that can be nonzero on
healthy footprints. This proves
\eqref{eq:induction-forced-closure}.

For finite families this is the exact flattening
and overlap-component bijection proved in the
physical reassembly. The nonnegative connected
majorant is absolutely summable with the displayed
prices, so the same identity passes to countable
families. The internal eligibility tests and exact
coarse flag inventories are retained in the merged
labels. The restored exponential and the forced
gas therefore give the same function pointwise.
\end{proof}

The field-independent scalar can be specified without
an asymptotic convention. If \(f=Lc\) are the fine and
coarse sides, set
\begin{align}
 e^{v_c c^4}={}&e^{v_f f^4}z_{\rm n}(\beta_f)^{-4c^4}
  \sum_t\int\lambda_t(U)
   e^{-\beta_f A_t(1,U)-I_f(U)-E_f^\#(U)}
                                  Z_{K_f}(U)\,dH_f(U).
 \label{eq:induction-exact-vacuum}
\end{align}
This is the actual blocked density at the identity.
For the Wilson history it is positive. The Gaussian
change of variables gives exactly the same scalar,
with its original Haar chart factors, determinants
and powers of \(g\). Equation~\eqref{eq:induction-exact-vacuum}
keeps all those factors through the physical Haar
integral and the exact noise normaliser.

\begin{theorem}[One nonlinear step]
\label{thm:nonlinear-one-step}
Choose the profile bounds, ratio and entry threshold
in the order above, and construct the common path arrays
for that calendar. For an
admissible fine input with
\(\widehat\beta_f/2\le\beta_f\le2\widehat\beta_f\),
the same physical blocking operator has
the exact representation
\begin{equation}
 \mathcal P_mF_f(V)=
 e^{v_c c^4-\beta_cS_c(V)-I_c(V)-E_c^\#(V)}Z_{K_c}(V)
                           \quad\hbox{for every }V,
 \label{eq:induction-pointwise-step}
\end{equation}
where
\begin{align}
 \beta_c&=\beta_f+\alpha_{n,m}+\Delta_m,\notag\\
 \|E_c\|_{\mathcal E,\kappa_c}&\le\delta_c/2,
       \qquad\|K_c\|_{\mathcal K,\kappa_c}\le3w_c/4,\notag\\
 |\beta_f-\beta_c-b_0\log L_m|&\le d_j+1.
 \label{eq:induction-one-step-bounds}
\end{align}
The regular flat-value and affine-Hessian normalisations,
forced eligibility, locality and prescribed symmetries
hold. Constants and entries are uniform in cutoff,
volume and finite history. The next subsection specifies
how the running window is retained for a further step.
\end{theorem}
\begin{proof}
The physical component expansion and
\eqref{eq:physical-exact-restoration} are identities
for the original blocked integral. Inserting the
ordinary finite-period redistribution, the incoming
comparison and their masks gives the preparatory
family of Proposition~\ref{prop:induction-preparatory-family}.
Apply the exact local and principal reset of
Lemma~\ref{lem:induction-reset-constant}, and then
Proposition~\ref{prop:induction-forced-closure}.
Their identities are pointwise; convergence has
been proved with the same local prices, so there
is no exchange of conditional sums. This gives
\eqref{eq:induction-pointwise-step}.

The scalar is \eqref{eq:induction-exact-vacuum}, since
\(S_c(1)=I_c(1)=0\), every new regular function has
zero flat value, and the forced gas at a good field
is one. The reset proves the regular normalisations.
The merger proves the forced inventories and locality.
The original branch, tag and path constructions and
their reindexing under translations and signed
permutations preserve the prescribed symmetries.
All Borel fine multipliers have stayed in their
physical arguments.

The regular bound follows from
\eqref{eq:induction-regular-contraction} with room
for the stated tails; the forced bound is
\eqref{eq:induction-forced-closure}.
Late entry gives \(|\Delta_m|\le1\).
Combining \(\beta_c=\beta_f+\alpha+\Delta_m\) with
Theorem~\ref{thm:uniform-one-loop-drift} proves the
last line of \eqref{eq:induction-one-step-bounds}.
Integrating \eqref{eq:induction-pointwise-step}
preserves the partition function by
Proposition~\ref{prop:push}. There has been no
change of the Wilson density whose exact history
is represented.
\end{proof}

% !TeX root = 4DYM_notes.tex
\subsection{Explicit entry bounds for every new remainder}
\label{sec:induction-numerical-ledger}

We now check the limiting prices quantitatively. This
also distinguishes the coefficient that must be chosen
before the ratio from the errors that entry can pay.
No change of activity norm or physical measure is made.

The following ledger locates the choices used in this section.
The bounds use the original path, regular and forced norms.
\begin{center}
\begin{tabular}{@{}p{.20\linewidth}p{.43\linewidth}p{.30\linewidth}@{}}
\hline
Quantity & Available bound & Choice that closes it\\ \hline
\raggedright Inherited regular part & \((V_j/L)\delta_f\) &
\raggedright Choose \(L\) after the reset constant.\tabularnewline[3pt]
\raggedright New regular part & \(r_{j,L}(x)=o(\delta_c)\) &
\raggedright Choose entry after \(L\).\tabularnewline[3pt]
\raggedright Inherited loop part & \(C_jL^{-2}\Lambda_f+C_{j,L}\) &
\raggedright Choose the ratio for the first term, then entry for the second.\tabularnewline[3pt]
\raggedright Forced packets & \(C_B e^{-cP_f^2}\) &
\raggedright Compare with \(w_c=e^{-P_c^{1/5}}\) at entry.\tabularnewline[3pt]
\raggedright Mask restorations & \(b_m^{\rm rest}\le x^{-1/8}\) &
\raggedright Keep an original forced member in every component.\tabularnewline[3pt]
\raggedright Coupling reset & \(|\Delta_m|\le Cx^{-1/10}P_c^{-2}\) &
\raggedright Use the regular price divided by \(h_c^2\).\tabularnewline \hline
\end{tabular}
\end{center}
Thus entry makes the new errors small; it does not supply
the inherited regular contraction. The detailed calculations
below verify these comparisons and their uniform tails.

Write \(x=\widehat\beta_{\rm c}\) and
\(x_{\rm f}=x+b_0\log L\). The ratio \(L\) is fixed
in the following limits. The monitor has logarithm
\(\log P(x)=(\log\log x)^2\), so for every fixed
\(u,v,p\ge0\),
\begin{equation}
 \frac{\log\{(1+B)^p(1+P)^u(1+\Lambda)^v\}}{\log x}
                                                   \longrightarrow0.
 \label{eq:induction-finite-subpower-products}
\end{equation}
Indeed \(\log(1+B)=O(\log\log x)\),
\(\log(1+P)=O((\log\log x)^2)\), and the loop
allowance at a fixed profile is a fixed power of
\(P\). Each numerator is therefore
\(O_{j,L,q}((\log\log x)^2)\).
In particular every finite product of these factors
is at most \(C_{j,L,q,\varepsilon}x^\varepsilon\)
for each prescribed \(\varepsilon>0\). One entry
threshold pays all of the finitely many products
required at that profile and derivative order.

The use of adjacent fine quantities does not change
this conclusion. Direct differentiation gives
\[
 \frac{d}{dx}\log P(x)=\frac{2\log\log x}{x\log x},\qquad
 \frac{d}{dx}\log h(x)=-\frac1{2x}
                                  +\frac{2\log\log x}{x\log x}.
\]
The mean-value formula on \([x,x_{\rm f}]\) proves
\(P_{\rm f}/P_{\rm c}\to1\) and
\(h_{\rm f}/h_{\rm c}\to1\). Also
\[
 \frac{B_{\rm f}}{B_{\rm c}}
 =1+O\left(\frac1{\log^2x}
                         +\frac{\log L}{x\log x}\right)\longrightarrow1.
\]
The first error retains the integer rounding of
\(B\). Since \(x_{\rm f}\ge x\),
\(\delta_{\rm f}=x_{\rm f}^{-11/10}\le
\delta_{\rm c}=x^{-11/10}\).

\begin{lemma}[A uniform numerical bound for the regular birth]
\label{lem:induction-explicit-entry}
At fixed profiles and ratio, the remainder in
\eqref{eq:induction-preparatory-price} may be chosen so
that, above a cutoff- and volume-independent threshold,
\begin{equation}
 \frac{r_{j,L}(x)}{\delta_{\rm c}}
                                     \le C_{j,L}x^{-1/25}.
 \label{eq:induction-explicit-remainder-rate}
\end{equation}
Thus a sufficient further entry condition is
\(x\ge(4C_{j,L}/c_j)^{25}\).
The constant here is permitted to depend on \(L\);
it is not used to choose \(L\).
\end{lemma}
\begin{proof}
Use the chosen envelopes in
\eqref{eq:central-explicit-prices}, not their possibly
smaller actual incidence sums. After collecting all
fixed subpower products, their contributions to
\eqref{eq:induction-new-remainders}, divided by
\(x^{-11/10}\), have the following bounds for any
fixed \(\varepsilon>0\):
\begin{equation}
\begin{array}{c|c}
 \text{term before division}&
               \text{bound after division, up to }C_{j,L,\varepsilon}\\ \hline
 C_Bx_{\rm f}^{-23/20}&x^{-1/20+\varepsilon}\\
 C_Bb^3&x^{-2/5+\varepsilon}\\
 C_Bbd&x^{-1/2+\varepsilon}\\
 C_Bd^2&x^{-11/10+\varepsilon}\\
 C_{j,L,{\mathsf M}}x_{\rm f}^{-{\mathsf M}}(1+\delta_{\rm f})&x^{11/10-{\mathsf M}}.
\end{array}
\label{eq:induction-regular-rate-table}
\end{equation}
For example
\(bd\le x_{\rm f}^{-1/2-11/10+o(1)}
=x_{\rm f}^{-8/5+o(1)}\), while
\(d^2\le x_{\rm f}^{-11/5+o(1)}\).
Taking \(\varepsilon=1/100\) and \({\mathsf M}=2\), every
power in the right column is at most \(x^{-1/25}\).
There are five terms, so their constants can be added
to give \eqref{eq:induction-explicit-remainder-rate}.
All analytic constants in this calculation were local
per-anchor constants in the preceding estimates.
No physical volume or number of RG steps enters this
sum. The displayed entry condition makes its right
side at most \(c_j/4\).
\end{proof}

The Gaussian warnings satisfy the required algebraic
rates with room to spare. For every fixed \(p\ge0\)
and \(a>0\),
\begin{equation}
 \frac{P(x)^a}{(1+B)^p\log x}\longrightarrow\infty,
 \qquad
 \log\frac{P(x)^a}{(1+B)^p\log x}
  =a(\log\log x)^2-(2p+1)\log\log x+O(1).
 \label{eq:induction-monitor-dominance}
\end{equation}
Consequently multiplying
\(\exp\{-P_{\rm f}^2/C_B\}\) by any fixed power
of \(x,P,B\) still leaves a bound smaller than
\(C_{j,L,{\mathsf M}}x^{-{\mathsf M}}\), for every fixed \({\mathsf M}\).
This verifies the derivative and counting prices in
\eqref{eq:induction-rare-regular} without inserting
a new Gaussian precision reserve.

\begin{lemma}[All restoration prices fit the forced reserve]
\label{lem:induction-explicit-forced-entry}
After the regular entry above, the quantities in
Proposition~\ref{prop:induction-forced-closure} satisfy
\[
 a_f/w_c\longrightarrow0,\qquad
 b_m^{\rm rest}\le x^{-1/8},\qquad w_c\le x^{-1/16}
\]
at a further threshold depending only on the fixed
profiles and ratio. In particular that proposition
uses one threshold for all cutoffs and volumes.
\end{lemma}
\begin{proof}
The physical packet price gives
\[
 \log(a_f/w_c)\le O(\log(1+B))
                         -c_{\rm pkt}P_f^2+P_c^{1/5}.
\]
Here \(P_f/P_c\to1\),
\(P_c^{1/5}/P_f^2\to0\), and
\(\log(1+B)/P_f^2\to0\). The right side tends to
\(-\infty\). The same calculation, with a fixed
polynomial denominator, handles the weaker Gaussian
warning price if it occurs in a marked regular sum.

For the restoration income use
\eqref{eq:induction-restoration-price}. The regular
bound gives
\[
 H_m/h_c\le C_jx^{-3/5}P_c^{-1}.
\]
At this fixed ratio the kinetic coefficient
\(\alpha=-b_0\log L+O_j(1)\) is bounded independently
of cutoff and volume. Hence
\[
 b_m^{\rm rest}
 \le C_B\{x^{-1/2+o(1)}+x^{-11/10+o(1)}
                  +C_jx^{-3/5}P_c^{-1}+C_{j,L}x^{-1/2}P_c\}
 =x^{-1/2+o(1)}.
\]
Use \eqref{eq:induction-finite-subpower-products} and
increase entry to make it at most \(x^{-1/8}\).
Finally \(w_c=e^{-P_c^{1/5}}\), and
\eqref{eq:induction-monitor-dominance} gives
\(P_c^{1/5}\ge(\log x)/16\) eventually, proving
the stated bound on \(w_c\).

These are numerical prices for the original merger,
not independent forced factors. Each restoration
component retains the mandatory forced label and
its exact flag inventory, as proved in
Proposition~\ref{prop:induction-forced-closure}.
Thus none of the polynomial restoration prices has
been added by itself to the superpolynomial allowance
\(w_c\). Applying Lemma~\ref{lem:reserve-choice} with
its auxiliary numerical reserve \(x^{1/16}\) gives
the claimed forced closure.
\end{proof}

We can now state the order of choices without an
implicit limiting argument. First fix the coarse
reset constant \(C_j^{\rm reset}\), set
\(c_j=[8(1+C_j^{\rm reset})]^{-1}\), and retain the
pre-ratio incoming constant \(V_j\). Choose
\(L\ge4V_j/c_j\), together with the other ratio
floors. Only after this choice use
Lemma~\ref{lem:induction-explicit-entry} to set entry.
Then
\[
 H_m\le(c_j/4)\delta_f+(c_j/4)\delta_c
                \le(c_j/2)\delta_c,
 \qquad C_j^{\rm reset}H_m
         \le\frac{C_j^{\rm reset}}{16(1+C_j^{\rm reset})}
                       \delta_c\le\delta_c/16.
\]
Also
\[
 |\Delta_m|\le H_m/h_c^2
             \le(c_j/2)x^{-1/10}P_c^{-2}\longrightarrow0.
\]
The loop bound is equally explicit: the calendar
inequality \(C_j^{\rm pb}L^{-2}\le1/8\),
\(\Lambda_f\le2\Lambda_c\), and the later entry
condition \(C_{j,L}\le\Lambda_c/4\) give
\[
 \Lambda_{\rm out}
       \le C_j^{\rm pb}L^{-2}\Lambda_f+C_{j,L}
                                             \le\Lambda_c/2.
\]
The forced entry is then fixed by the preceding lemma.
Taking the maximum of these finitely many thresholds
still introduces no cutoff, volume or history
dependence. The common delayed calendar in the next
subsection admits them for every fixed request.

\subsection{Choosing one calendar after the estimates}

We spell out the choices because they are part of the
uniform induction. First fix the free kernel catalogue,
the derivative orders \(q_j\ge3\), spatial powers
\(s_j\ge2\), covering exponents and their reserves,
and the positive numbers \(\theta_j\). Fix the
finite marked support exponents between \(96\)
and \(100\) as part of this same catalogue. Include
the cover-compression floors with input exponent
\(\kappa_f-4\) and the prescribed output exponent
\(\kappa_c\). Their constants precede the block
ratio; enlarging its lower bound pays the finite
four-unit reserve at a physical source joining.
Fix the
off-exact row completions, then choose the coarse
classical widths and budgets, including a
stronger exponential width for their derivatives, by
Proposition~\ref{prop:ordinary-classical-matching}.
These output bounds depend on free and geometric
profile data; dependence on the incoming nonlinear
budget is assigned to the subsequent ratio floor.
They can therefore be chosen for all indices before
any of their ratios. Include the finite Wilson top
budget used in the next section. The clip coefficients
and relaxed-cutoff parameter are also fixed at this stage.

This fixes \(C_j^{\rm reset},c_j,V_j,d_j\).
Now choose odd nondecreasing ratios \(L_j\) satisfying
all the earlier geometric and path floors and
\begin{equation}
 \frac{V_j}{L_j}\le\frac{c_j}{4},\qquad
 C_j^{\rm pb}L_j^{-2}\le\frac18,\qquad
 \frac{d_j+1}{\log L_j}
                         \le\min(b_0/32,j^{-2}).
 \label{eq:induction-ratio-choices}
\end{equation}
Here \(C_j^{\rm pb}\) is the pre-ratio inherited-loop
constant of \eqref{eq:ordinary-loop-budget}.
Each requirement is a lower bound on a freely chosen
odd integer. A maximum with the preceding ratio
also makes the sequence nondecreasing.
Include here the finite-order current ratio floors
of Lemma~\ref{lem:trace-current-local-bounds}: its
exponential widths are chosen before the ratios,
and its diagonal contraction constants depend only
on those widths and the free profile data. In
particular impose \(L a_{j_f}\ge32a_{j_c}\) and the
lower bound for \(L\) in
\eqref{eq:trace-explicit-defect-matrix}.
For the finite degree-four and degree-five marked requests,
include the ratio floor \(L^{1/8}\ge4K_j\) of
Lemma~\ref{lem:ward-matching-coefficient-control}. Its
leading-map constant and the chosen coefficient tolerance
\(\xi\le1/(3K_j)\) are fixed before the ratio; the
coupling-dependent coefficient price is imposed only at
entry. These requirements supplement the existing
remainder-contraction floors and do not alter the
unmarked physical step.
Choose the noise-sector width as required by the
packet payment after fixing these ratios.

Only now choose the profile entry thresholds. Include
the smallness of the new remainder, the forced-merger
conditions, the relaxed-cutoff radius, the neighbouring
threshold comparisons and all read-period reserves.
The loop bound closes too: if the fine and coarse
indices are \(j_f\ge j_c\), then
\[
 \frac{\Lambda_f}{\Lambda_c}
       \le(P_f/P_c)^{\theta_{j_c}}\longrightarrow1.
\]
Thus at entry \(\Lambda_f\le2\Lambda_c\), and one
can require \(C_{j,L}\le\Lambda_c/4\), since
\(\Lambda_c\to\infty\) at each fixed profile.
Equation~\eqref{eq:ordinary-loop-budget} then gives
an output loop budget at most \(\Lambda_c/2\).
The large constant \(C_{j,L}\) has been used only
for this entry condition, never for selecting \(L\).
Likewise include the current entry conditions
\(\log\widehat\beta_m\ge64A_{j,L}\), enlarged for
their finite number of components and their bounded
normalised forcing. The constant \(A_{j,L}\) is the
finite-jet bound constructed in the proof of
Lemma~\ref{lem:trace-current-local-bounds}. It is
finite after \(L\) is fixed. These requests depend
on finite-volume local formulas, not on a continuum
hierarchy or a selected vacuum.
Include also the separated coefficient price
\eqref{eq:ward-separated-atom-prices}, the finite-source
normalisation request
\(4C(1+C_j)\le P^{1/12}\), and
\(C_j\le\log\widehat\beta\). Here \(C\) includes the
finite current and ordinary-source birth constants in
\eqref{eq:trace-joint-source-price}. Include also
the actual response bound
\eqref{eq:trace-actual-response-budget},
\(4(1+C_j)\le\sqrt P\), and
\(4(1+C_j)Ch^2P^{1/6}\le\widehat\beta^3\).
For the complete physical source joining impose
\eqref{eq:response-joining-smallness} with
\[
 c=w+r_B\sqrt w+\widehat\beta^{-1}
                              +r_B\widehat\beta^{-1}.
\]
At fixed profile these are genuine entry conditions:
\(D_B,r_B\) are fixed logarithmic powers,
\(w=e^{-P^{1/5}}\), and
\(C_j(1+\log\widehat\beta)^{A_j}\le\widehat\beta\)
eventually. Response functions use no derivatives
and consume no additional Gaussian precision.
Include the degree-three potential and degree-four
forward-difference births of
Lemma~\ref{lem:degree-three-potential-birth}.
Their complete polynomial representatives use
the already admitted current jets through order
four. Their higher healthy derivatives are
bounded by the finite polynomial product rule.
The finite constants \(A_j^{\rm pot},p_j^{\rm pot}\)
are fixed after the ratio and enter only the
requests
\(A_j^{\rm pot}(1+\log\widehat\beta)^{p_j^{\rm pot}}
\le\widehat\beta\) and
\eqref{eq:degree-three-potential-birth-price}.
The degree-three low read requires at most
four derivatives; it adds no derivative loss
along the RG iteration. These are requests
in the same initial calendar.
Include the single additional vacuum
derivative used for source birth: for every admitted
marked order \(q\), carry the unmarked regular-action
bound through \(q+1\), then use its restriction in the
marked step. No derivative reserve is consumed repeatedly
along the RG history. The longer spatial
cutoff \(C_j^{\rm cut}(1+B)^{61}\) used for the Ward
comparison, and its fixed derivative powers, are included
in the same local read and period reserves. Its logarithm
is \(O_j(\log\log\widehat\beta)\); it is a fixed
logarithmic power, so all these conditions hold after
entry. All fixed constants and exponents appearing in
this finite request can additionally be bounded by
\(\log\widehat\beta\) by delaying entry. No accuracy
depending on a later selected vacuum enters these choices.

For clarity all neighbouring comparisons follow from
fixed shifts of the calendar coupling. At fixed
\(L\), \(\widehat\beta_f=\widehat\beta_c+b_0\log L\).
The derivative of \((\log\log x)^2\) is
\(2\log\log x/(x\log x)\), so \(P_f/P_c\to1\).
Also the derivative of
\(\exp\{(\log\log x)^2/5\}\) tends to zero.
Hence \(w_f/w_c\to1\); this last assertion uses the
difference of their exponents, not merely the ratio
of their monitors. In particular the conditions
\(h_c\le2h_f\), \(B_f\le2B_c\), and
\(\widehat\beta_f/\widehat\beta_c\le2\) hold at a
common entry. The fixed factor two in covering
compression is included in the pre-ratio constants.

Include also the entry condition
\begin{equation}
              L_j\le\log\widehat\beta_m
       \quad\hbox{whenever }J(m+1)=j.
 \label{eq:induction-trace-calendar}
\end{equation}
For a fixed profile the left side is already fixed and the
right side increases without bound. Delaying its first entry
therefore imposes this condition for both adjacent profiles,
without changing any ratio or previously fixed local constant.
This choice will be used for comparison at arbitrary small
physical radii in Section~\ref{sec:trace-prescribed-radii}.
It is included in the initial common calendar, before the
Wilson family and any of its limits are chosen. For the
fixed profile constants \(A_j\) in the scalar-free vacuum
ratio \eqref{eq:insertion-vacuum-ratio}, include likewise
\(A_{J(m)}\le\log\widehat\beta_m\). This is another finite
entry condition at each profile. In particular that ratio's
local exponent is \(O((\log\widehat\beta_m)^2)\), with a
constant independent of the level. This reserve will be
used when a matching length grows as a power of
\(\log\widehat\beta_m\).

Every entry condition is eventually satisfied for
a fixed profile. Corollary~\ref{cor:linear-profile-clock}
therefore chooses one slow clock with
\(J(m+1)-J(m)\in\{0,1\}\), \(J(m)\to\infty\),
and \(L_m=L_{J(m+1)}\). Include the further entry
condition
\[
       s_m\ge j\log L_j
       \quad\hbox{whenever }J(m+1)=j\ge5.
\]
This is possible by delaying the first occurrence of
each profile, since \(s_m\ge m\log3\).
After its first occurrence the same condition persists
while that profile is repeated. Consequently
\(\log L_m/s_m\le1/J(m+1)\to0\).
This is the mesh condition needed to make each single
blocking ratio negligible relative to the accumulated scale.
The physical period conditions
can be included by Lemma~\ref{lem:finite-period-reserve}:
the physical sides grow exponentially in the scale,
while each required radius is a fixed polynomial
in the logarithm of the nominal coupling. The
finitely many initial period margins are imposed
when choosing the base volume family. These choices
are made once, before any cutoff, volume subsequence
or shooting parameter is selected. All earlier free
and local estimates were uniform over such calendars,
so this is a single choice for the entire construction.

\subsection{Continuity of a step on a compact parameter interval}

The weighted activity-sum argument adapts
\cite[Lemma~13.13, PDF p.~108]{ArmstrongDraft}.
The required majorants and the finite-volume integral are treated below
for the present curve.

A continuous activity curve means a curve with fixed
labels and footprints, continuous in the weighted
norms for differences of its regular and forced
families, and with continuous coupling and scalar.
For example, for a regular family the sum in
\eqref{eq:E-norm}, applied to
\(E_i(u)-E_i(u_0)\), tends to zero as \(u\to u_0\).
The hard tests and calendar thresholds are fixed
along the curve; they are not varied with its actual
coupling parameter.

\begin{proposition}[Continuity without extending a curve past its window]
\label{prop:induction-curve-continuity}
At fixed cutoff and volume, the nonlinear step sends
a continuous admissible Wilson-history curve on a
compact interval to a continuous output curve on
that same interval, while its couplings remain in
the permitted bands. All output labels can be chosen
independently of the parameter.
\end{proposition}
\begin{proof}
Use the held physical variables of
\eqref{eq:regular-mixed-physical-Gaussian-density}
for bounded Borel reads. The parameter then varies
the smooth kernel and Gaussian density and the
already continuous input multipliers. It does not
differentiate or move a hard physical flag test.
Gaussian density differences, with each fixed
coarse derivative, tend to zero in weighted
\(L^1\) on compact coupling bands. The positive
precision reserve and its score moment estimates
give a common integrable majorant. Frozen Haar
and tag integrations are unchanged probability
integrations. Thus every finite integrated word
and its regular germs depend continuously on the
parameter.

The word and connected-log majorants are uniform
on the whole compact interval. Consequently their
tails are summable in the corresponding local
norms, and finite partial families prove continuity
of their sums. The ordinary finite-stencil
approximations use the cover-kernel majorants
on the entire coupling band and are fixed
before the physical period and the parameter.
For other continuous input label families a
common finite approximation follows from
compactness: a
continuous map from a compact interval into the
weighted summable family norm has uniformly small
tails: take finite approximations at finitely many
parameter centres and use norm continuity on a
finite cover of the interval. The union of these
finite label sets works for all parameters.
Complete finite symmetry orbits when needed.

Flat-value and affine extraction are bounded
linear maps on these germs. Thus \(\Delta_m\) is
continuous; the principal transfer and the new
regular family are continuous too. The forced
merger has the same uniformly convergent
nonnegative majorant, which proves continuity in
its incidence norm. Finally the physical integral
in \eqref{eq:induction-exact-vacuum} is continuous
by domination on the finite compact Haar space,
and is positive on the Wilson history. Its
logarithm gives a continuous scalar.
All restrictions below are to closed subintervals
of this very interval, so no continuity outside
the valid input window is used.
\end{proof}

\subsection{A window that survives a step}

Write
\begin{equation}
 \rho_m=\bar\beta/16+
          \sum_{r=0}^{m-1}(d_{J(r+1)}+1),\qquad
 I_m=[\widehat\beta_m-\rho_m,
                         \widehat\beta_m+\rho_m].
 \label{eq:induction-running-windows}
\end{equation}
Here \(\rho_m\) is a scalar window radius, unrelated
to the stationary graph \(\rho(x)\).
By \eqref{eq:induction-ratio-choices},
\[
 \rho_m\le\bar\beta/16+(b_0/32)s_m
                         \le\widehat\beta_m/16.
\]
Include \(\widehat\beta_m\ge32(d_j+1)\) among the
entry conditions for a step of fine index \(j\).
If \(\beta_f\in I_{m+1}\), then
\begin{equation}
 |\beta_c-\widehat\beta_m|
       \le\rho_{m+1}+d_j+1
       =\rho_m+2(d_j+1)\le\widehat\beta_m/8.
 \label{eq:induction-output-band}
\end{equation}
Thus all those output couplings are in the required
broad band, before any further interval is selected.

\begin{lemma}[Restriction between two crossings]
\label{lem:induction-window-restriction}
Let \(F:[a,b]\to\mathbb R\) be continuous, and
\(\ell<u\) satisfy \(F(a)\le\ell\), \(F(b)\ge u\).
There is a closed subinterval \([a',b']\subset[a,b]\)
such that \(F(a')=\ell\), \(F(b')=u\), and
\(F([a',b'])\subset[\ell,u]\).
\end{lemma}
\begin{proof}
Let \(b'\) be the first point with \(F(b')=u\);
existence and a minimum follow from the intermediate
value theorem and compactness of the level set.
Let \(a'\) be the last point of \([a,b']\) with
\(F(a')=\ell\), obtained in the same way.
A value above \(u\) before \(b'\) would force an
earlier crossing of \(u\). A value below \(\ell\)
after \(a'\) would force a later crossing of
\(\ell\) before \(b'\). Both contradict the choices.
This proves the entire interval assertion.
\end{proof}

Suppose an input curve at level \(m+1\) stays in
\(I_{m+1}\), with its left and right endpoints equal
to the endpoints of that window. Apply the nonlinear
step on its whole parameter interval. At the left
endpoint its coarse coupling is at most
\[
 \widehat\beta_m-\rho_{m+1}+d_j+1
                         =\widehat\beta_m-\rho_m;
\]
at the right endpoint it is at least
\(\widehat\beta_m+\rho_m\). Continuity and
Lemma~\ref{lem:induction-window-restriction} give
a closed subinterval on which the new curve stays
in \(I_m\) and has both prescribed endpoint values.
Every finer representation already constructed is
simply restricted to that interval and keeps all
its previous bounds. This is the inductive window
argument; it uses no monotonicity of the reset map.
The next section supplies its actual Wilson starting
curve and completes the finite-cutoff induction.
\clearpage
\section{The Wilson starting curve and the finite-cutoff induction}
\label{sec:wilson-start-shooting}

We now supply the initial representation required in
Section~\ref{sec:nonlinear-induction-step}. Its density
is exactly the Wilson density on every configuration.
The regular part corrects the chosen principal action
on healthy plaquettes. A nonnegative unassigned action
pays for the forced gas on the remaining configurations.
The construction follows the local Wilson initialisation
in \cite[Lemma~13.21, PDF pp.~119--120]{ArmstrongDraft};
we include its pointwise identity, read sets, norms and
continuity.

At the empty history use the initial off-exact choice
\(l_f=0\) already specified in the proof of
Lemma~\ref{lem:packet-complementary-row}. Thus both
top curvature metrics are the identity. Choose the
following top path arrays as the initial data of the
descending construction in Section~\ref{sec:ordinary-matching}.
All lower arrays are then obtained from those data
before choosing a volume or an actual coupling.

\subsection{The local fourth-order polynomial}

For \(U\in G\), write
\[
 \mathcal C(U)=\operatorname{pr}_{\mathfrak g}
                         \frac{U-U^{-1}}2.
\]
This is the curvature used in the packet rows.
It is conjugation covariant and odd under inversion.
Also
\begin{equation}
       |\mathcal C(U)|\le d_G(U),\qquad
                                  d_G(U)^2=2q_G(U).
 \label{eq:startup-curvature-contraction}
\end{equation}
Indeed \((U-U^{-1})/2=(U-1)(1+U^{-1})/2\), the
second factor has operator norm at most one, and
orthogonal projection contracts the matrix norm.

The map \(X\mapsto\mathcal C(e^X)\) is real analytic
with derivative the identity at zero. Let \(\iota\)
be its local analytic inverse. One may obtain this
inverse directly: on a fixed small complex ball the
derivative of \(\mathcal C(e^X)-X\) has norm less
than \(1/2\), and
\(X=Y-(\mathcal C(e^X)-X)\) is a contraction for
smaller \(Y\). Uniform convergence of its analytic
iterates gives the analytic inverse.
Define the homogeneous quartic polynomial
\begin{equation}
 T_4(Y)=\frac1{4!}
     D^4\bigl[q_G(e^{\iota(\cdot)})\bigr](0)[Y,Y,Y,Y].
 \label{eq:startup-quartic}
\end{equation}
The inverse is conjugation covariant and odd, by
uniqueness, so \(q_G(e^{\iota(Y)})\) is invariant
and even. Its quadratic term is \(|Y|^2/2\).
Consequently the globally defined smooth function
\begin{equation}
 R(U)=q_G(U)-\tfrac12|\mathcal C(U)|^2
                                     -T_4(\mathcal C(U))
 \label{eq:startup-remainder}
\end{equation}
has zero jets through degree four at the identity.
Only the polynomial in \eqref{eq:startup-quartic}
uses the local inverse; \eqref{eq:startup-remainder}
is defined on all of \(G\). It is invariant under
conjugation and inversion.

Fix a top level \(n\). Put \(h=h_n\), \(B=B_n\),
\(\lambda=\lambda_{J(n)}\), and let
\(\lambda_{\rm l}\) be its linear clip coefficient,
with \(0<16\lambda_{\rm l}<\lambda\le1\).
For a plaquette \(q\) with vertex footprint \(X_q\)
and anchor \(o_q\), set
\[
 r_q=d_G(P_q(U)),\quad
 k_q=(1-\lambda)|\mathcal C(P_q(U))|^2/2,\quad
 \chi_q=\chi_{X_q}(U).
\]
Let \(\mathcal H_{o_q}\) be the original relaxed
row-health indicator of \eqref{eq:row-relaxed-cutoff}.
It implies \(r_q\le h\) and, at entry, reads within
distance \(B-2\) of its anchor, including its
plaquette halo. Increase entry if necessary to
obtain this margin from
\eqref{eq:relaxed-cutoff-radius}.

The top row is \(\sqrt{1-\lambda}\,\mathcal C(P_q)\),
there are no cubic or loop atoms, and the quartic
atom is \((1-\lambda)T_4(\mathcal C(P_q))\).
Thus the physical principal action is the prescribed
normalised signed-permutation average of
\begin{equation}
 \sum_q\{\operatorname{clip}_h(r_q)
       +\mathcal H_{o_q}k_q
       +(1-\lambda)\chi_qT_4(\mathcal C(P_q))\}.
 \label{eq:startup-principal-action}
\end{equation}
Here the clip is \(\lambda r^2/2\) for \(r\le h\)
and \(\lambda_{\rm l}hr\) for \(r>h\), with equality
assigned to the former branch. These are finite
path arrays with identity quadratic form and the
required Wilson fourth jet. Their finite exponential
budgets are included in the profile bounds before
the ratios are chosen.

\subsection{A nonnegative action on the flagged regions}

For each plaquette define
\begin{equation}
 W_q(U)=
 \begin{cases}
 q_G(P_q)-\lambda_{\rm l}hr_q,&r_q>h,\\
 \chi_qk_q+(1-\chi_q)(1-\lambda)q_G(P_q)
                         -\mathcal H_{o_q}k_q,&r_q\le h.
 \end{cases}
 \label{eq:startup-unassigned-charge}
\end{equation}
These functions are Borel and nonnegative. For the
second line, when \(\chi_q=1\) it is
\((1-\mathcal H_{o_q})k_q\); when \(\chi_q=0\)
use \(k_q\le(1-\lambda)q_G(P_q)\) from
\eqref{eq:startup-curvature-contraction}.
For the first line,
\begin{equation}
          W_q\ge r_q^2/4\ge h^2/4\qquad(r_q>h).
 \label{eq:startup-bad-plaquette-charge}
\end{equation}
This follows from \(q_G=r_q^2/2\), \(r_q>h\),
and \(\lambda_{\rm l}<1/16\).

In both cases direct substitution gives the exact
identity
\begin{align}
 q_G(P_q)={}&\operatorname{clip}_h(r_q)
       +\mathcal H_{o_q}k_q
       +(1-\lambda)\chi_qT_4(\mathcal C(P_q))\notag\\
       &+(1-\lambda)\chi_qR(P_q)+W_q.
 \label{eq:startup-plaquette-identity}
\end{align}
For \(r_q>h\), both masks multiplying the row and
the regular terms are zero. For \(r_q\le h\),
insert \eqref{eq:startup-remainder} and
\(\operatorname{clip}_h=\lambda q_G\) to verify
the identity. This also covers a healthy plaquette
next to a bad one, and a relaxed row test failed
because of a more distant bad plaquette.

Let \({\mathsf M}_{\rm sym}\) be the number of signed
permutations in the principal-action average, and
write \(\operatorname{Av}_\sigma\) for that normalised
finite average. Take regular labels \((q,\sigma)\)
with unmasked functions
\[
       \frac{\beta(1-\lambda)}{{\mathsf M}_{\rm sym}}
                                  R(P_q(\sigma^*U)),
\]
and the corresponding image plaquette footprints.
Different labels are retained even if their values
agree. Their masked sum, denoted by \(E_{n,\beta}^\#\),
is precisely the average of the fourth term in
\eqref{eq:startup-plaquette-identity}, multiplied
by \(\beta\).

To exponentiate the last term, let
\(D=\mathcal B_h(U)\) be the actual bad-site set.
Join two sites when their radius-\(2B\) balls
intersect. For every nonempty connected inventory
\(I\) in this fixed graph put
\(Y_I=N_{2B}(I)\). Define \(d_I(U)\) to be the
normalised signed average of all \(W_q\) whose
physical anchors, in that signed image, are within
distance \(B\) of \(I\). Set
\begin{equation}
        K_{I,\beta}(U)=
          \mathbf1_{\{Y_I\cap\mathcal B_h(U)=I\}}
                                         e^{-\beta d_I(U)}.
 \label{eq:startup-forced-activity}
\end{equation}

We verify the locality and assignment in this
definition. Every nonzero \(W_q\) has a bad site
within distance \(B\) of its anchor. For a bad
plaquette this is one of its own vertices. For a
failed \(\chi_q\) the same is true. A failed
relaxed test has a witnessing plaquette of size
greater than \(h\), and its witness lies within
distance \(B\) by the chosen read margin. The
complete link and flag reads of each charge
anchored within \(B\) of \(I\) lie in \(Y_I\).
Thus \(K_I\) is local on the declared links and
flags and has exactly the forced eligibility
condition. It is gauge invariant.

Distinct actual components of \(D\) have disjoint
\(Y_I\)'s. A charge within \(B\) of two components
would join them, so each nonzero charge is assigned
once. Conversely an admissible exact-cover product
of the functions in \eqref{eq:startup-forced-activity}
selects precisely these actual components. A
selected connected inventory lies within an actual
component. If a neighbouring flag belonged to a
different selected inventory, their radius-\(2B\)
supports would intersect, contradicting
compatibility. Exact coverage therefore forces
selection of each whole actual component.
This proves that the forced gas has exactly one
nonzero term, with value
\begin{equation}
       Z_{K_{n,\beta}}(U)=
       \exp\{-\beta\operatorname{Av}_\sigma\sum_qW_q(\sigma^*U)\}.
 \label{eq:startup-forced-identity}
\end{equation}
If \(D\) is empty the selected family is empty and
the right side is one, as required.

\begin{proposition}[Exact Wilson initialisation]
\label{prop:exact-wilson-initialisation}
The preceding top arrays and activities have zero
loop action and vacuum scalar, and for every \(U\)
and every positive \(\beta\) satisfy
\begin{equation}
 e^{-\beta W_{\ell_{n,M}}(U)}
      =e^{-\beta S_n(U)-E_{n,\beta}^\#(U)}Z_{K_{n,\beta}}(U).
 \label{eq:exact-wilson-initialisation}
\end{equation}
Here \(\ell_{n,M}\) is the side of the top torus.
The regular
activities vanish at the identity with their
affine Hessians. The representation has the
prescribed symmetries and exact forced inventories.
\end{proposition}
\begin{proof}
Sum \eqref{eq:startup-plaquette-identity}, take its
normalised signed average, and multiply by
\(\beta\). The Wilson energy is unchanged by this
average. Equations~\eqref{eq:startup-principal-action}
and \eqref{eq:startup-forced-identity} then give
\eqref{eq:exact-wilson-initialisation}.
The jets of \(R\) through degree four vanish at
the identity, so the regular flat and Hessian
normalisations hold on their plaquette footprints.
Conjugation and inversion invariance, and
reindexing the signed images and inventories,
give the stated symmetries. There is no scalar
factor: all actions vanish at the identity and
the forced gas there is one.
\end{proof}

\subsection{The initial activity norms and continuity}

\begin{lemma}[Uniform prices of the Wilson starting curve]
\label{lem:wilson-start-prices}
After the profile entry is chosen, uniformly in
\(n,M\) and
\(\widehat\beta_n/2\le\beta\le2\widehat\beta_n\),
\begin{equation}
 \|E_{n,\beta}\|_{\mathcal E,\kappa(n)}
       \le\delta_n/1024,\qquad
 \|K_{n,\beta}\|_{\mathcal K,\kappa(n)}
       \le w_n/1024.
 \label{eq:wilson-start-prices}
\end{equation}
Both families form continuous curves in their weighted
norms, with fixed labels, footprints and hard tests.
\end{lemma}
\begin{proof}
We first estimate the local regular remainder without
bounding the four individual links by \(h\).
Use the plaquette holonomy as one group coordinate
and the other three links as compact parameters;
the fourth link is recovered by a product and an
inverse. This is a smooth change of coordinates
with uniformly bounded word derivatives. At a
healthy plaquette \(|\log P_q|\le C_Gh\).
The jets of \(R(e^X)\) through order four vanish,
so Taylor's formula bounds its order-\(k\)
derivative by \(C_kh^{5-k}\) for \(k\le5\).
For \(k>5\) its derivatives are bounded on the
fixed chart. Each physical direction on this
four-link footprint has size at most \(C_sh\).
In the ordered chain rule a term with \(k\)
derivatives on \(R\) and total direction order
\(r\ge k\) supplies \(h^r\). It is therefore
bounded by \(C_{q,s}h^5\), also when \(k>5\),
since then \(r>5\) and \(h\le1\).
This proves the local norm bound at every fixed
derivative order. Each anchor has finitely many
plaquettes and signed images, and their covering
number is bounded. Thus
\[
             \|E_{n,\beta}\|_{\mathcal E,\kappa(n)}
                                    \le C_j\beta h_n^5.
\]
The top band gives
\[
 \frac{\widehat\beta_nh_n^5}{\delta_n}
             =\widehat\beta_n^{-2/5}P_n^5\longrightarrow0,
\]
which proves the regular bound after entry.

For the forced family take an eligible inventory
\(I\) of \(r\) sites. Every site in \(I\) touches
a bad plaquette. All its vertices are flags in
the same inventory, by eligibility, and a
plaquette has at most four vertices. There are
therefore at least \(r/4\) such distinct bad
plaquettes assigned to \(I\).
Retain just the identity signed image in the
nonnegative average defining \(d_I\). By
\eqref{eq:startup-bad-plaquette-charge},
\begin{equation}
        d_I(U)\ge\frac{h_n^2r}{16{\mathsf M}_{\rm sym}}
               \quad\hbox{whenever }Y_I\cap\mathcal B_h(U)=I.
 \label{eq:wilson-start-component-charge}
\end{equation}
Consequently \(\|K_{I,\beta}\|_\infty\le
e^{-c\beta h_n^2r}\), with
\(c=(16{\mathsf M}_{\rm sym})^{-1}>0\).

The overlap graph has degree at most
\(C(1+B)^4\). A connected inventory of size
\(r\), rooted at one of its sites, has a
spanning-tree traversal of length \(2(r-1)\).
Counting these walks bounds its number by
\([C(1+B)^4]^{2(r-1)}\). To have a fixed site
in \(Y_I\) leaves at most another \(C(1+B)^4\)
choices for the first inventory site. A
radius-\(2B\) ball has a cover by at most
\(5^4\) radius-\(B\) balls; the union of these
covers is connected along the inventory graph.
Thus \(n_B(Y_I)\le5^4r\), including periodic
unions. The incidence norm is bounded by
\begin{equation}
 C_j(1+B)^4\sum_{r\ge1}
       [C_j(1+B)^8e^{5^4\kappa(n)-c\beta h_n^2}]^r.
 \label{eq:wilson-start-forced-sum}
\end{equation}
As \(\beta h_n^2\ge P_n^2/2\), its base tends
to zero and the sum is at most
\(C_Be^{-c'P_n^2}\). This is smaller than
\(w_n/1024=e^{-P_n^{1/5}}/1024\) at entry.
There is no factor equal to the physical volume.

The tests and charges are independent of the
actual \(\beta\). The regular difference estimate
is \(C_j|\beta-\beta_0|h_n^5\).
For the forced functions the mean-value inequality
gives, with \(\beta_{\min}>0\) a lower bound on
the parameter interval,
\[
 |e^{-\beta d}-e^{-\beta_0d}|
       \le|\beta-\beta_0|\,d\,e^{-\beta_{\min}d}
       \le\frac{2|\beta-\beta_0|}{e\beta_{\min}}
                                      e^{-\beta_{\min}d/2}.
\]
Apply \eqref{eq:wilson-start-forced-sum} with half
the charge. It is still convergent on the whole
top band, so weighted difference norms tend to
zero. This proves the claimed continuity and
uniform prices.
\end{proof}

The preceding construction installs the entire Wilson
curve on the top band. Its restriction to \(I_n\)
is therefore available for the running-window
argument, with vacuum scalar zero and loop action
zero. No identification from fourth jets alone
has been used on its flagged configurations.

\subsection{The finite-cutoff Wilson history}

\begin{theorem}[A common Wilson family with uniformly controlled histories]
\label{thm:finite-cutoff-nonlinear-induction}
For the common calendar chosen in
Section~\ref{sec:nonlinear-induction-step}, and
every cutoff \(n\) and physical side length \(M\), there
is a bare coupling \(\beta_{n,M,n}>0\) and an
exact effective history \(F_{n,M,m}\),
\(0\le m\le n\), such that
\begin{align}
 F_{n,M,n}(U)&=e^{-\beta_{n,M,n}W_{\ell_{n,M}}(U)},\notag\\
 F_{n,M,m}&=\mathcal P_mF_{n,M,m+1},\qquad
                           \beta_{n,M,0}=\bar\beta.
 \label{eq:finite-cutoff-wilson-history}
\end{align}
At every level its path, regular and forced
families have the prescribed profile bounds and
\begin{equation}
 \|E_{n,M,m}\|_{\mathcal E,\kappa(m)}\le\delta_m,
       \qquad \|K_{n,M,m}\|_{\mathcal K,\kappa(m)}\le w_m,
       \qquad \beta_{n,M,m}\in I_m.
 \label{eq:finite-cutoff-uniform-bounds}
\end{equation}
The path arrays and one-loop coefficients are
chosen independently of \(M\) and the actual
couplings. The same calendar, bounds and entries
work for all \(n,M\). Moreover
\begin{equation}
       \frac{\beta_{n,M,n}}{\log {\varepsilon}_n^{-1}}
                          \longrightarrow b_0
                      \quad\hbox{uniformly in }M.
 \label{eq:constructed-bare-running}
\end{equation}
All scalar normalisations and the partition
function are retained in the exact histories.
\end{theorem}
\begin{proof}
Choose the catalogue and common calendar first.
For each \(n\), specify the canonical top arrays
of \eqref{eq:startup-principal-action} and
descend through the classical and one-loop
constructions of Section~\ref{sec:ordinary-matching}.
The pre-ratio classical bounds and the closed
loop allowance prove their profile bounds by
descending induction, independently of \(M\)
or a shooting parameter.

For fixed \(n,M\), start with the continuous
Wilson curve \(\beta\in I_n\) just proved.
Its endpoints are exactly the endpoints of the
running window. Suppose a continuous curve of
histories has been obtained down to level
\(m+1\), with that level's coupling staying in
\(I_{m+1}\) and taking both endpoint values.
All its finer-level representations retain their
activity bounds. Apply
Theorem~\ref{thm:nonlinear-one-step} and
Proposition~\ref{prop:induction-curve-continuity}
to the whole curve. Equation~\eqref{eq:induction-output-band}
keeps the coarse coupling in the broad band
required for the construction. The two endpoint
inequalities and
Lemma~\ref{lem:induction-window-restriction}
give a closed parameter subinterval with coarse
coupling in \(I_m\) and both endpoint values.
Restrict all finer data to it. This completes
the induction down to zero; \(n=0\) uses the
starting curve directly.

The final continuous coupling curve covers
\(I_0=[15\bar\beta/16,17\bar\beta/16]\).
It therefore takes the value \(\bar\beta\).
For definiteness choose the least original bare
parameter with this property in the final
compact interval. The preimage is nonempty
and closed, so the choice exists. Evaluating
the whole history at that parameter gives
\eqref{eq:finite-cutoff-wilson-history} and
\eqref{eq:finite-cutoff-uniform-bounds}.
At each step the scalar is the exact flat
integral \eqref{eq:induction-exact-vacuum}, and
mass preservation gives
\[
 \int F_{n,M,m}\,dH_m
             =Z_{\ell_{n,M}}(\beta_{n,M,n})
                         \quad(0\le m\le n).
\]
The endpoint selection has not altered a
represented density or its normalisation.

All constants in this induction are profile
constants. Each regular output is bounded by
the next level's allowance, using its pre-ratio
incoming contraction and its own new-remainder
entry. Each forced output uses a fresh marked
component bound and the exact merger. The next
step uses these local norm bounds and its own
conditional Gaussian law. Thus no factor
depending on the number of previous RG steps
enters the prices. The fixed calendar and the
period reserves are common to all volumes.

Finally the exact coupling increments obey
\[
 |\beta_{n,M,m+1}-\beta_{n,M,m}-b_0\log L_m|
                         \le d_{J(m+1)}+1,
 \qquad
 \frac{d_{J(m+1)}+1}{\log L_m}\le J(m+1)^{-2}\longrightarrow0.
\]
Apply Proposition~\ref{prop:drift}, with the
coarse value \(\bar\beta\). It gives
\eqref{eq:constructed-bare-running}, uniformly
in \(M\). Absolute summability of the step
errors is unnecessary for this conclusion.
\end{proof}

The choices in the last theorem define one Wilson
family for all cutoffs and all volume indices
before any volume subsequence is selected.
This establishes the finite-cutoff nonlinear
vacuum history and its uniform inverse-coupling
asymptotic. Construction of the continuum
correlation hierarchy, with the required
joint insertion estimates and volume
subsequence, remains a subsequent part of the
proof.
% !TeX root = 4DYM_notes.tex
\section{Local observables on a fixed stencil}
\label{sec:local-observable-sections}

The histories of Theorem~\ref{thm:finite-cutoff-nonlinear-induction}
represent the vacuum density. An observable needs an additional
normalisation: a bounded function of the lattice links can acquire
an unbounded coefficient as the spacing tends to zero. We first
construct the local functions to which that normalisation will be
applied. Neither a continuum field nor a continuum correlation
function is assumed in this construction.

The finite-measure argument of Section~\ref{sec:wilson} establishes
integrability at a fixed cutoff. The estimates below instead keep
track of their dependence on the small-field threshold. The local
section construction follows
\cite[Lemmas 10.2--10.3, PDF pp.~40--43]{ArmstrongDraft};
we give the construction and the derivative estimates explicitly.

\subsection{Polynomial labels and their two gradings}

Fix a nonnegative integer \(d\). For \(j+2\le d\), let
\(\mathcal J_j\) be the real vector space of \(\mathfrak g\)-valued
one-forms \(A_j=(A_{j,1},\ldots,A_{j,4})\) on \(\R^4\) whose
coefficients are homogeneous polynomials of degree \(j+1\) and obey
\begin{equation}
                  \sum_{a=1}^4 z_a A_{j,a}(z)=0.
 \label{eq:insertion-radial-gauge}
\end{equation}
This is the radial-gauge condition. A truncated jet is
\(A=(A_j)_{j+2\le d}\) in the finite-dimensional real space
\(\mathcal J_{\le d}=\bigoplus_{j+2\le d}\mathcal J_j\).
These are test potentials; they are not random lattice fields.
Equip this space with any coefficient inner product invariant under
the adjoint action of \(G\) and signed coordinate permutations.
Such a choice exists by averaging an inner product over the compact
group of these actions.

A coordinate of \(A_j\) has engineering weight \(j+2\).
Let \(V_d\) be the complex vector space of adjoint-invariant
polynomials in these coordinates with total engineering weight at
most \(d\). Polynomial degree, here called factor count, counts
the number of jet coordinates in a monomial. The two gradings are
different: a derivative of curvature increases engineering weight
without necessarily increasing factor count. Both gradings are
preserved by the adjoint and orthogonal coordinate actions.
There are only finitely many monomials of weight at most \(d\), so
\(V_d\) is finite-dimensional.

Choose a real bihomogeneous basis \((b_{d,i})_{i\in I_d}\), with
\(0\in I_d\), \(b_{d,0}=1\), and write \((w_i,k_i)\) for its
engineering weight and factor count. A real basis means that its
polynomials are real-valued on real jets. Such a basis exists:
complex conjugation preserves every invariant bihomogeneous
component, and its real fixed space complexifies to that component.
For \(x=\sum_i x_i b_{d,i}\), put
\[
 |x|_1=\sum_i|x_i|,\qquad x^\circ=x-x_0\,1,
 \qquad (\mathsf D_hx)_i=h^{k_i}x_i\quad(h>0).
\]
The diagonal map \(\mathsf D_h\) scales factor count. It is
invertible, fixes the unit, and commutes with the signed-permutation
and conjugation actions. Every nonconstant basis element has
\(k_i\ge1\), and \(2k_i\le w_i\le d\).

The spatial dilation of a test jet is
\[
       \operatorname{sc}_\varepsilon A
                 =(\varepsilon^{j+2}A_j)_{j+2\le d}.
\]
It acts on a polynomial of weight \(w_i\) by
\(b_{d,i}(\operatorname{sc}_\varepsilon A)
=\varepsilon^{w_i}b_{d,i}(A)\). Factor-count scaling and
spatial dilation must not be identified.

\subsection{Test holonomies and finite measurements}

Put \(A^{\le d}=\sum_{j+2\le d}A_j\). For a positive cover
link \((z,a)\), let \(U(\varepsilon,A)_{(z,a)}=u(1)\), where
\begin{equation}
 \frac{\dd u(t)}{\dd t}
   =u(t)\,\varepsilon A^{\le d}_a(\varepsilon(z+te_a)),
                  \qquad u(0)=1,
 \label{eq:insertion-test-holonomy}
\end{equation}
and negative links are read by inversion. The connection on the
right is identified with its skew-Hermitian matrix in the fixed
representation. The integral equation for \(u\) has its convergent
successive-integral series. On a bounded parameter set its term
of order \(l\) is at most \(C^l/l!\); hence the series and each
fixed number of parameter derivatives converge uniformly.
It gives existence and uniqueness on \([0,1]\). Alternatively,
uniqueness follows by differentiating \(u(t)v(t)^*\) for two
solutions. Approximating the coefficient by step functions gives
products of exponentials in \(G\); their uniform limit lies in
\(G\), since \(G\) is closed. Thus these are actual group-valued
links, and they are smooth functions of the finite jet parameters.
Homogeneity in \eqref{eq:insertion-test-holonomy} gives
\begin{equation}
 U(\varepsilon,A)=U(1,\operatorname{sc}_\varepsilon A),
                     \qquad U(0,A)=1.
 \label{eq:insertion-test-dilation}
\end{equation}

For a plaquette \(q\) based at \(z\), choose a finite cover path
\(\gamma\) from the origin to \(z\). With \(C\) as in
\eqref{eq:startup-curvature-contraction}, define the transported
measurement
\[
                  \sigma_{\gamma,q}(U)
                          =\operatorname{Ad}(U_\gamma)C(P_q(U)).
\]
The ordered product \(U_\gamma\) uses oriented links along that
path. Each measurement transforms by the adjoint action at the
origin. It is a real polynomial in real and imaginary matrix
entries of finitely many links: inverses of unitary matrices are
their adjoints, and \(C\) is a fixed real linear projection of
\((P-P^{-1})/2\).

\begin{lemma}[Finite curvature measurements determine the test jet]
\label{lem:insertion-finite-measurements}
There are finitely many measurements \(\sigma_1,\ldots,\sigma_{\mathsf M}\)
such that the map
\[
 \mathsf S:\mathcal J_{\le d}\longrightarrow\mathfrak g^{\mathsf M},
       \qquad \mathsf S(A)=(\sigma_l(U(1,A)))_{l=1}^{\mathsf M},
\]
has injective derivative at zero. The measurements and their paths
can be chosen before any lattice spacing, period or RG ratio.
\end{lemma}
\begin{proof}
The first variation of a link holonomy at the zero connection is
the integral of its varied connection along the link. The first
variation of a plaquette is therefore the ordinary circulation of
that one-form. Differentiating its transporting path adds zero,
because the plaquette curvature at the identity is zero.

Suppose every translated unit-square circulation of
\(A=\sum_{j=0}^{j_*}A_j\) is zero. For \(a<b\), write
\[
 {\mathsf F}_{j;ab}(z,t)=\int_0^1\!\int_0^1
 (\partial_a A_{j,b}-\partial_b A_{j,a})
                     (z+t(ue_a+ve_b))\,\dd u\,\dd v.
\]
The fundamental theorem of calculus on the four edges gives
\(\sum_{j\le j_*}{\mathsf F}_{j;ab}(z,1)=0\) at every integer \(z\).
Homogeneity gives
\({\mathsf F}_{j;ab}(tz,t)=t^j{\mathsf F}_{j;ab}(z,1)\). For \(t>0\), multiply
the vanishing circulation by \(t^{j_*}\) to obtain
\[
          \sum_{j\le j_*}t^{j_*-j}{\mathsf F}_{j;ab}(tz,t)=0.
\]
Given \(y\in\R^4\), choose the integer vector
\(z_t=(\lfloor y_a/t\rfloor)_a\); then \(tz_t\to y\).
Letting \(t\downarrow0\) proves
\(\partial_a A_{j_*,b}-\partial_b A_{j_*,a}=0\) at \(y\).
Differentiating \eqref{eq:insertion-radial-gauge} and using Euler's
identity for a homogeneous polynomial gives
\[
 (j_*+2)A_{j_*,b}(y)
    =\sum_a y_a(\partial_a A_{j_*,b}-\partial_b A_{j_*,a})(y)=0.
\]
Remove this homogeneous component and repeat. All components
vanish, so the intersection of the kernels of all the linearised
measurements is zero. Starting with \(\mathcal J_{\le d}\),
add one measurement whenever the current common kernel is nonzero.
Its dimension decreases strictly. This process terminates after
finitely many choices and proves the assertion.
\end{proof}

\subsection{A polynomial local section}

\begin{proposition}[Fixed-stencil Taylor sections]
\label{prop:insertion-local-sections}
For every \(d\) there are a finite box
\(\Sigma_d=[-r_d,r_d]^4\cap\Z^4\) and a complex-linear map
\(x\mapsto P_d[x]\) to polynomials in matrix entries and their
conjugates on links with both endpoints in \(\Sigma_d\), such that:
\begin{enumerate}
 \item \(P_d[x]\) is gauge invariant and \(P_d[1]=1\) identically.
 \item It respects complex conjugation and signed coordinate
       permutations.
 \item For every test jet \(A\),
 \begin{equation}
 \left.\partial_\varepsilon^l
             P_d[x](U(\varepsilon,A))\right|_{\varepsilon=0}
  =\left.\partial_\varepsilon^l
             x(\operatorname{sc}_\varepsilon A)
                                      \right|_{\varepsilon=0},
                      \qquad 0\le l\le d.
 \label{eq:insertion-section-matching}
 \end{equation}
 \item Each \(P_d[b_{d,i}]\) is a polynomial in transported
       curvatures with at least \(k_i\) factors in each monomial.
\end{enumerate}
The section, its coefficients and its stencil are independent of
cutoff, volume and the subsequent choice of blocking ratios.
\end{proposition}
\begin{proof}
For \(d<2\), take the constant section. Otherwise use
Lemma~\ref{lem:insertion-finite-measurements}. Choose invariant
Euclidean structures, and let \(C_0\) be the inverse of
\(D\mathsf S(0)\) on its range, composed with orthogonal projection
onto that range. It is adjoint equivariant and
\(C_0D\mathsf S(0)=I\).
The map \(C_0\mathsf S\) has a smooth inverse \(\varphi\)
near zero. For completeness, on a small ball its derivative differs
from \(I\) by at most \(1/2\); for a sufficiently small target
\(y\), the equation \(X=y-(C_0\mathsf S(X)-X)\) is a contraction
of that ball into itself. This proves local existence and uniqueness.
Differentiating this equation gives the smooth inverse and its
higher derivatives. Invariant source and target balls and uniqueness
make \(f=\varphi C_0\) equivariant, with
\(f\mathsf S(A)=A\) and \(f(0)=0\).

Let \(T_df\) be the ordinary Taylor polynomial of \(f\) at zero
through order \(d\). It is an equivariant polynomial with no
constant term. Define
\[
 \Psi_x(U)=\sum_i x_i\,
 b_{d,i}\bigl(T_df(\sigma_1(U),\ldots,\sigma_{\mathsf M}(U))\bigr).
\]
This is gauge invariant. Each of its nonconstant monomials in
the measurements has at least the original \(k_i\) factors.
Taylor's formula gives
\(T_df(\mathsf S(A))=A+O(|A|^{d+1})\). Substitution of
\(\operatorname{sc}_\varepsilon A\) and
\eqref{eq:insertion-test-dilation} proves the matching through
weight \(d\). All maps are real before complexification, so
\(\Psi\) respects conjugation.

Let \(\mathcal W\) denote the finite group of signed coordinate
permutations, acting on connections by pullback and on polynomial
labels by \((\rho x)(A)=x(\rho^{-1}A)\). Set
\[
 P_d[x](U)=\frac1{|\mathcal W|}
           \sum_{\rho\in\mathcal W}\Psi_{\rho^{-1}x}(\rho^*U).
\]
Changing the summation variable proves covariance. Pulling back
\eqref{eq:insertion-test-holonomy}, including reversal and inversion
on negative links, proves the corresponding transformation of test
holonomies by uniqueness of that equation. Each averaged term has
the same Taylor polynomial \(x(\operatorname{sc}_\varepsilon A)\).
A signed image of a transported curvature is again a transported
curvature, possibly with a sign and a changed boundary transport.
Consequently the factor-count assertion persists. Enclose all the
finitely many paths and plaquettes in one invariant box \(\Sigma_d\).
Expansion in matrix entries proves every stated property.
\end{proof}

Placement at a site \(o\) means evaluation on the periodic cover
extension \(U^{[o]}_{(z,a)}=U_{(o+[z],a)}\). We write the result
as \(P_d[x](o;U)\). Its combinatorial footprint is the torus image
of \(o+\Sigma_d\). Its unmasked physical reads are the links with
both endpoints in that image. The health indicator requires the
additional plaquette halo of Lemma~\ref{lem:flag-footprint}.
Even if cover links coincide after placement, the polynomial is
defined by the same cover formula; coincident arguments are simply
assigned the same link value.

\begin{lemma}[Threshold-scaled derivative bounds]
\label{lem:insertion-section-bounds}
For fixed \(G,d,q,s\) and the section just chosen, there are finite
constants \(C_{d,q,s}\) and \(C_d\), independent of period and
\(0<h\le1\), such that
\begin{align}
 \|P_d[b_{d,i}](o;\cdot)\|_{q,h,s;o,o+\Sigma_d}
                                      &\le C_{d,q,s}h^{k_i},\notag\\
 \|P_d[\mathsf D_{h^{-1}}x](o;\cdot)\|_{q,h,s;o,o+\Sigma_d}
                                      &\le C_{d,q,s}|x|_1,\notag\\
 \sup_{o,U}|P_d[\mathsf D_{h^{-1}}x](o;U)|
                                      &\le C_dh^{-d}|x|_1.
 \label{eq:insertion-section-bounds}
\end{align}
The norm in this display is precisely
\eqref{eq:local-derivative-norm}, extended to complex functions by
using complex absolute values.
\end{lemma}
\begin{proof}
The projection and adjoint transport are contractions, so a healthy
transported curvature has norm at most \(h\). An ordered derivative
of a fixed path product is a finite sum of products of unitary
matrices and direction matrices. At order \(l\ge1\), every term
has \(l\) direction factors. Their norms on the fixed cover paths
are at most
\(h\max_{z\in\Sigma_d}(1+|z|)^s\). The differentiated measurement
is therefore bounded by \(C_{d,l,s}h^l\le C_{d,l,s}h\).
This also covers repeated physical links: differentiate the cover
formula first, then assign the coincident directions.

In the product rule for a monomial with at least \(k_i\)
measurements, each differentiated or undifferentiated measurement
thus still contributes one factor \(h\). There are finitely many
terms, and the derivative order is bounded by \(q\). This proves
the first bound. Multiplication of the \(i\)-th column by
\(h^{-k_i}\), followed by the triangle inequality, proves the
second. On arbitrary configurations each unitary matrix entry has
absolute value at most one. The finite matrix-entry polynomial
therefore has a global bound \(C_d\) per column. Since \(k_i\le d\)
and \(h\le1\), its scaled version is bounded by
\(C_dh^{-d}|x|_1\). No volume cardinality was used.
\end{proof}

In particular, these sections give continuous gauge-invariant local
observables with the fixed stencil required by
Definition~\ref{def:regularity}. An eventual cutoff-dependent linear
normalisation changes their coefficients, but does not enlarge
their stencil.

\section{One insertion: the complete density and its initial bounds}
\label{sec:marked-density-start}

We now keep an insertion together with the exact vacuum density.
Fix \(n,M,d\) and one of the histories constructed in
Section~\ref{sec:wilson-start-shooting}. Suppress \(n,M\) in the
notation. On level \(m\), write
\[
 F_m=e^{\Phi_m}Z_{K_m},\qquad
 \Phi_m=v_m\#\Lambda_m-\beta_mS_m-I_m-E_m^\#,
\]
where \(\Lambda_m\) is the level-\(m\) site torus and \(H_m\)
remains its Haar probability measure. All original masks and
eligibility tests in \(F_m\) are retained. The insertion anchor is
zero until translations are considered.

\subsection{The marked format and its norm}

\begin{definition}[A one-insertion representation]
\label{def:one-insertion-format}
A one-insertion representation consists of a coordinate \(x\in V_d\),
regular functions \(R_i\) on footprints \(Z_i\) containing zero,
and forced functions \(Q_\alpha\) on forced addresses of the same
vacuum gas. Unused addresses may be added with unmarked activity
zero. Its density is
\begin{align}
 \mathcal M_m(U)={}&F_m(U)\left[x_0+
   \chi_{\Sigma_d}(U)P_d[\mathsf D_{h_m^{-1}}x^\circ](0;U)
                      +\sum_i\chi_{Z_i}(U)R_i(U)\right]\notag\\
 &+e^{\Phi_m(U)}
     \sum_{A\ {\rm admissible}}
           \sum_{\alpha\in A}Q_\alpha(U)
                      \prod_{\alpha'\in A\setminus\{\alpha\}}
                                             K_{m,\alpha'}(U).
 \label{eq:one-insertion-density}
\end{align}
Here a stencil in a health indicator means its placed torus image.
Each \(R_i\) is gauge invariant, Borel and regular in the sense of
the local derivative norm. It satisfies \(R_i(1)=0\). On a
nonwrapping \(Z_i\) its test-holonomy Taylor coefficients through
weight \(d\) vanish:
\begin{equation}
       \left.\partial_\varepsilon^l
       R_i(U(\varepsilon,A))\right|_{\varepsilon=0}=0,
                           \qquad 0\le l\le d.
 \label{eq:marked-jet-null}
\end{equation}
Only its cover reads enter this formula. Each \(Q_\alpha\) is
gauge invariant and Borel, uses the original forced physical reads,
and is zero unless that address is eligible. If it is nonzero as
a function, its footprint contains the anchor.
\end{definition}

The unit coefficient is unmasked. Masking it would change the
constant observable on bad configurations. Likewise, the last
line of \eqref{eq:one-insertion-density} replaces one forced
factor; it does not multiply the gas by an independently normalised
expectation.

For majorants \(e_i\ge\|R_i\|_{q_m,h_m,s_m;0,Z_i}\) and
\(k_\alpha\ge\|Q_\alpha\|_\infty\), define the anchored masses
\begin{align}
 E_m^{\rm ins}&=\sum_i e^{\kappa(m)n_{B_m}(Z_i)}e_i,\notag\\
 K_m^{\rm ins}&=\sum_\alpha
                e^{\kappa(m)n_{B_m}(Y_\alpha)}k_\alpha,\notag\\
 \|g_m\|_{\rm ins}&=E_m^{\rm ins}+w_m^{-1/2}K_m^{\rm ins},
                    \qquad g_m=(R,Q).
 \label{eq:marked-remainder-norm}
\end{align}
The infimum over such majorants is understood when the norm is
used without a chosen majorant. Here \(q_m,s_m\) are the profile's
derivative order and spatial power, and \(\kappa(m)\) and \(w_m\)
are the same support exponent and forced allowance as in
Section~\ref{sec:nonlinear-induction-step}. Unlike the unmarked
norm, this is a sum at one insertion anchor, with no spatial
supremum and no independent translation sum. The square-root
forced weight leaves half of the existing forced smallness
available for marked estimates.

\begin{lemma}[Absolute convergence of the one-insertion density]
\label{lem:marked-density-envelope}
At a fixed finite torus, a finite marked remainder norm makes
\eqref{eq:one-insertion-density} a bounded Borel function.
Its regular and forced sums converge absolutely, uniformly in
\(U\). In particular it is integrable against \(H_m\).
\end{lemma}
\begin{proof}
The masked regular sum is bounded by \(\sum_i e_i\).
For the forced sum drop disjointness, eligibility and the coverage
test after taking absolute values. For a fixed distinguished
address \(\alpha\), the remaining products have sum at most
\(\prod_{\alpha'\ne\alpha}(1+\|K_{m,\alpha'}\|_\infty)\).
The unmarked incidence bound gives
\(\sum_{\alpha'}\|K_{m,\alpha'}\|_\infty
\le \#\Lambda_mw_m\): sum the incidence bounds over sites and
use that every forced footprint is nonempty. Consequently the
whole marked gas is bounded by
\[
                  e^{\#\Lambda_mw_m}\sum_\alpha k_\alpha.
\]
This majorant proves uniform absolute convergence. Every term is
Borel, and the finite-volume scalar-free action and \(\Phi_m\)
are bounded. The section is a fixed continuous polynomial on a
compact configuration space. These observations prove the claim.
The bound here is an integrability bound, not a volume-uniform
bound on a normalised correlation.
\end{proof}

\subsection{Restoring the unhealthy part of the initial section}

The following exact construction uses only the forced incidence
bound, not positivity of individual forced activities.
\begin{proposition}[Birth of a local insertion]
\label{prop:marked-insertion-birth}
At the Wilson top, the unmasked density
\[
             \mathcal M_n(x;U)
                 =F_n(U)P_d[\mathsf D_{h_n^{-1}}x](0;U)
\]
has a representation \eqref{eq:one-insertion-density} with
coordinate \(x\), zero regular remainder, fixed marked addresses,
and
\begin{equation}
 \|g_n(x)\|_{\rm ins}
 \le C_d h_n^{-d}
 e^{\kappa(n)n_{B_n}(\Sigma_d)}w_n^{-1/2}
       (e^{|\Sigma_d|w_n}-1)|x^\circ|_1.
 \label{eq:marked-birth-price}
\end{equation}
This choice is complex linear in \(x\), vanishes on the unit, and
retains every original forced factor, inventory and auxiliary test.
For each fixed profile and admitted degree, its right-hand
coefficient is at most one after a uniform coupling entry.
\end{proposition}
\begin{proof}
Let
\(f_x=(1-\chi_{\Sigma_d})P_d[\mathsf D_{h_n^{-1}}x^\circ]\).
For an admissible vacuum family \(A\), write
\[
       A_\Sigma=\{\alpha\in A:Y_\alpha\cap\Sigma_d
                                                 \ne\varnothing\}.
\]
If \(A_\Sigma\) is empty, exact coverage of the bad inventory
implies that \(\Sigma_d\) is healthy, and \(f_x=0\).
Otherwise give this finite subfamily one new marked address with
footprint and inventory
\[
   Y=\Sigma_d\cup\bigcup_{\alpha\in A_\Sigma}Y_\alpha,
                   \qquad B=\bigcup_{\alpha\in A_\Sigma}B_\alpha.
\]
Its unmarked activity is zero. Its marked function is
\(f_x\prod_{\alpha\in A_\Sigma}K_{n,\alpha}\), restricted to
the original eligibility tests and the exact inventory event
\(Y\cap\mathcal B_{h_n}(U)=B\). The address records the original
subfamily, even when two subfamilies have identical supports.
All reads of the multiplied functions, health indicator and tests
are included; the flag halo remains a physical read.

Every remaining old footprint avoids \(\Sigma_d\) and all the
merged old footprints. Replacing \(A_\Sigma\) by the marked
address therefore gives an admissible marked family. Conversely,
expanding that recorded address recovers the unique original
family. Additional new marked addresses have unmarked activity
zero and cannot contribute to the one-mark sum. This bijection
proves that the new marked sum equals \(F_n f_x\) exactly.
It proves the density identity on all configurations, including
threshold equalities and wrapping stencils.

The connected-cover inequality gives
\[
 e^{\kappa(n)n_{B_n}(Y)}
 \le e^{\kappa(n)n_{B_n}(\Sigma_d)}
       \prod_{\alpha\in A_\Sigma}
                   e^{\kappa(n)n_{B_n}(Y_\alpha)}.
\]
Indeed every merged footprint meets \(\Sigma_d\), so their
covers are connected through a cover of the stencil. The total
weighted unmarked income of addresses meeting the stencil is at
most \(|\Sigma_d|w_n\). Dropping the disjointness requirement,
the sum of their products over nonempty finite subfamilies is at
most \(e^{|\Sigma_d|w_n}-1\). The global section bound of
Lemma~\ref{lem:insertion-section-bounds} now proves
\eqref{eq:marked-birth-price}.

Once \(B_n\ge r_d\), the stencil has covering number at most
one. Since \(w_n\to0\),
\(e^{|\Sigma_d|w_n}-1\le2|\Sigma_d|w_n\) after entry.
Also \(h_n^{-1}\le\widehat\beta_n\) after entry. Thus the
coefficient in \eqref{eq:marked-birth-price} is bounded by
\(2C_d|\Sigma_d|e^{\kappa(n)}
\widehat\beta_n^d w_n^{1/2}\), which tends to zero at each
fixed profile. This follows directly from
\(w_n=e^{-P_n^{1/5}}\) and
\(P_n^{1/5}/\log\widehat\beta_n\to\infty\).
There are finitely many admitted degrees per profile, so their
entries can be chosen together before any cutoff or volume.
\end{proof}

\subsection{Exact transport and the spacing factor}

For the moment \(x\) is a bare coordinate, not a renormalised
reference coordinate. Define the inserted densities at all lower
levels by exact transport:
\begin{equation}
       \mathcal M_m(x)=L_m^d\mathcal P_m\mathcal M_{m+1}(x),
                                  \qquad 0\le m<n.
 \label{eq:marked-exact-transport}
\end{equation}
The scalar \(L_m^d\) anticipates the physical normalisation
\({\varepsilon}_n^{-d}\) of an insertion. It is not another vacuum
partition function.

\begin{proposition}[One common denominator for an inserted density]
\label{prop:marked-exact-mass}
Every density in \eqref{eq:marked-exact-transport} is integrable,
is complex linear in \(x\), and satisfies
\begin{equation}
 \int\mathcal M_m(x)\,\dd H_m
  ={\varepsilon}_m^d Z_{\ell_{n,M}}(\beta_{n,M,n})
   \left\langle {\varepsilon}_n^{-d}
              P_d[\mathsf D_{h_n^{-1}}x](0;U)\right\rangle_{n,M}.
 \label{eq:marked-exact-mass}
\end{equation}
The bracket denotes expectation in the original Wilson probability
measure. The identity includes every vacuum scalar and every
Gaussian, Haar and blocking-noise normalisation.
\end{proposition}
\begin{proof}
At the top the section is bounded and continuous, so its product
with the Wilson density is integrable. Proposition~\ref{prop:push}
proves integrability at every subsequent step and preserves its
integral, apart from the explicit \(L_m^d\) in
\eqref{eq:marked-exact-transport}. Their product is
\(\prod_{i=m}^{n-1}L_i^d=({\varepsilon}_m/{\varepsilon}_n)^d\). The top integral is
\(Z_{\ell_{n,M}}(\beta_{n,M,n})\langle
P_d[\mathsf D_{h_n^{-1}}x]\rangle_{n,M}\), which proves the
formula. Linearity follows at the top and is preserved by each
integral operator. No field-dependent density was divided out.
\end{proof}

This identity does not yet bound a renormalised moment uniformly
in \(n,M\). That task requires bounds on the coordinates and
remainders of the transported insertion. Products of insertions
also require their joint transport: the product of separately
transported one-insertion densities is not their joint density.
% !TeX root = 4DYM_notes.tex
\section{The one-insertion RG estimate and reference normalisation}
\label{sec:insertion-transport}

The useful bound is in the current polynomial coordinate. A bound
only in the initial bare coordinate could acquire a new constant
at every RG step. We prove an inverse bound for the coordinate map
and then express the remainder in its new coordinate. This is the
coefficient-graph construction of
\cite[Lemmas 14.5--14.8, PDF pp.~129--132]{ArmstrongDraft}.

All statements in this section concern the same physical vacuum
step as Theorem~\ref{thm:nonlinear-one-step}. The numerator uses its
selected Gaussian law, tags, masks, incoming activities and Haar
normalisations. The denominator is its unmarked density. It is not
replaced by a Gaussian partition function.

\subsection{Additional choices for finitely many insertion degrees}

For this construction choose the common support exponent of the
previous calendar to be \(\kappa(m)=100\). The earlier estimates
allow this fixed choice; their covering reserves are chosen before
the ratios. We retain the marked estimates also at exponents
between \(96\) and \(100\), including the four-unit
physical-source joining reserve proved in
Lemma~\ref{lem:complete-source-response-joining}.
Fix nondecreasing sequences of derivative
orders, spatial powers and free averaging orders so that, at
profile \(j\),
\[
          q_j\ge j+4,\qquad s_j\ge j,
\]
and the free-difference estimates are available through order
\(j\). Enlarge the orders if an earlier vacuum estimate requires
more. Only finitely many orders and degrees are used at a fixed
profile. Choose a nondecreasing reference level \(\operatorname{ref}(d)\) so that
\(J(\operatorname{ref}(d))\ge d+3\), and use the marked construction only for
\(\operatorname{ref}(d)\le m\le n\).

Let \(K_j\ge1\) be a nondecreasing bound for the finite section
constants, basis recombinations, leading coordinate maps and their
inverses, for all degrees admitted through profile \(j\). The
leading-map bound is proved below and is fixed before the ratio.
Choose nondecreasing column bounds \(C_j\) with
\begin{equation}
                         C_j\ge64K_{j+1}^2.
 \label{eq:insertion-column-choice}
\end{equation}
This anticipates the adjacent finer profile. It is not a recursive
definition of \(K_j\) in terms of \(C_j\): the former comes from
the section and free kernels, before a remainder graph is chosen.
Further local prices depending on these columns are included in
the subsequent ratio floor. All these sequences are fixed before
any cutoff, period or insertion coefficient.

\begin{definition}[A coefficient graph]
\label{def:insertion-coefficient-graph}
At level \(m\), a coefficient graph is a complex-linear assignment
\(x\mapsto g_m(x)=(R(x),Q(x))\) with fixed addresses and supports,
the format of Definition~\ref{def:one-insertion-format}, and
\begin{equation}
                  \|g_m(x)\|_{\rm ins}
                                 \le C_{J(m)}|x^\circ|_1.
 \label{eq:insertion-graph-bound}
\end{equation}
It vanishes on the unit. The density represented by the graph
is \eqref{eq:one-insertion-density} with coordinate \(x\).
\end{definition}

The graph describes a representation, not a new choice of an
observable. Its value is required to equal the prescribed inserted
density after integration.

\subsection{Taylor reading and the leading map}

For a regular gauge-invariant germ \(f\) at the identity, substitute
the test links \(U(\varepsilon,A)\) and read the coefficients of
\(\varepsilon^l\), \(0\le l\le d\). Their sum is an element of
\(V_d\), denoted by \(\Pi_d f\). To see that these coefficients
are polynomials, expand the link integral equation and the Taylor
formula for \(f\) to order \(d\). Every jet occurrence contributes
its positive integer weight; only finitely many terms can have
weight at most \(d\). Gauge invariance makes them adjoint invariant.
The matching identity \eqref{eq:insertion-section-matching} gives
\(\Pi_dP_d[x]=x\). At threshold \(h\), the corresponding
dimensionless reading is
\[
                         \Pi_{d,h}f=\mathsf D_h\Pi_df.
\]
Thus \(P_d[\mathsf D_{h^{-1}}\Pi_{d,h}f]\) reproduces the
specified Taylor coefficients. Subtracting it gives
\eqref{eq:marked-jet-null}; the subtraction is on the unmasked
healthy germ. Its mask difference is restored by a marked forced
family exactly as in Proposition~\ref{prop:marked-insertion-birth}.

The reading has a finite derivative bound at each profile. Indeed,
the radial test directions of weight at most \(d\) grow at most
like \((1+|z|)^d\), which is allowed by \(s_m\ge d\).
An ordinary Taylor derivative of order \(l\) has at least \(l\)
jet factors. Its factor-count normalisation therefore leaves
\(h^{k-l}\le1\) when \(k\ge l\), after the derivative bound
has supplied \(h^{-l}\). The finite product rule and the chosen
basis give
\[
                   |\Pi_{d,h_m}f|_1
                   \le C_{j,d}\|f\|_{q_m,h_m,s_m;0,Z}.
\]
This bound concerns the nonwrapping cover reading. On a wrapping
support no such cancellation is imposed. Its entire regular germ
is retained with its covering weight; an empty bad inventory is
never relabelled as forced merely because a support wraps.

Put \(L=L_m\), \(f=m+1\), \(c=m\), and define
\[
              (\mathsf S_Lx)_i=L^{d-w_i}x_i.
\]
This is the engineering rescaling after the explicit \(L^d\) in
\eqref{eq:marked-exact-transport}. Its inverse has \(\ell^1\)
operator norm at most one. To define the leading map \(T_m\) on the cover,
substitute the free interpolation of the coarse radial test links
into a fine section, read its jet, and perform the two factor-count
normalisations. The engineering diagonal \(\mathsf S_L\) is
kept outside \(T_m\). More explicitly, if \(\mathsf R_m^{\rm free}\)
is the polynomial reading before either threshold normalisation,
then
\[
 T_m=\mathsf D_{h_c}\mathsf R_m^{\rm free}
             \operatorname{diag}(L^{w_i})\mathsf D_{h_f^{-1}},
 \qquad
 L^d\mathsf D_{h_c}\mathsf R_m^{\rm free}\mathsf D_{h_f^{-1}}
                                                    =T_m\mathsf S_L.
\]
The actual periodic integral is still used for the exact
transport. Its local winding and localisation differences are
retained among the comparison terms below; a polynomial test
potential is not assumed to be periodic.

\begin{lemma}[Uniform leading coordinate map]
\label{lem:insertion-leading-map}
The map \(T_m\) is triangular in increasing engineering weight,
fixes the unit, preserves the nonconstant subspace, and has
diagonal entries \((h_c/h_f)^{k_i}\). For fixed adjacent profiles,
\begin{equation}
                    \|T_m\|_{1\to1}+\|T_m^{-1}\|_{1\to1}
                                                       \le K_j,
                           \qquad j=J(m+1),
 \label{eq:insertion-leading-inverse}
\end{equation}
after enlarging \(K_j\) by a fixed factor. This constant is chosen
before \(L\) and is uniform in cutoff, period and history.
\end{lemma}
\begin{proof}
The free interpolation column has, in coarse block coordinates,
the bounds \eqref{eq:all-scale-linear-difference-bounds}. Through
order \(d\) these have the form
\[
 |\nabla_f^aE(Ly+{\varrho},y')|
       \le C_{j,a}L^{-1-|a|}e^{-c_j|y-y'|},
                                    \qquad |a|\le d.
\]
Its polynomial moments converge absolutely, since
\(\sum_z e^{-c_j|z|}(1+|z|)^d<\infty\). The zeroth normalisation
and the exact gradient identity
\eqref{eq:physical-linear-column-identities} reproduce the highest
homogeneous part of each polynomial test potential. Applying a
fine difference lowers its degree and adds one inverse power of
\(L\). Induction on that degree shows that the jet substitution
on an input label of weight \(w_i\) has coefficient
\(\delta_{bi}L^{-w_i}\) at every output weight \(w_b\le w_i\).
The remaining coefficients have \(w_b>w_i\), factor count at
least \(k_i\), and absolute value at most \(C_{j,d}L^{-w_i}\).
For the last bound expand the finitely many polynomial factors
and sum their moments with the displayed exponential majorant.
This also describes why lower spatial-degree terms of an
interpolated potential give higher, rather than lower, label
weights: their jet coefficients come from the higher-weight
coarse input.

Multiply the input column by \(L^{w_i}\), as already accounted
for in \(\mathsf S_L\), and by its two threshold powers.
If a higher output factor count occurs, the extra power of
\(h_c\le1\) only decreases the bound. The remaining threshold
ratio is bounded because \(1/2\le h_c/h_f\le2\) at entry.
This proves the triangular assertion and bounds all entries
before \(L\). Back substitution in the finite weight order
bounds the inverse as well: each diagonal inverse is at most
\(2^{k_i}\), and there are finitely many entries. The unit is
exactly constant throughout this substitution. Periodisation
sums the same exponential majorant over cover translates and
does not change these bounds. Differences between a local cover
reading and its physical-period integral have the exponentially
small tails from the read-period reserve; they enter the exact
comparison in Lemma~\ref{lem:insertion-local-prices}.
\end{proof}

\subsection{The local marked estimates}

We give the estimates before rebasing by the new coordinate.
They concern a single primary section, a single regular remainder
or a single forced remainder. A regular remainder is an entire
masked physical cylinder; a forced remainder replaces its actual
gas factor and retains its original eligibility. Neither is an
abstract polynomial evaluated at a different random field.

\begin{lemma}[Local marked prices before coordinate inversion]
\label{lem:insertion-local-prices}
Let \({\varrho}=\|g_{m+1}(x)\|_{\rm ins}\). The exact transport admits
a coarse reading \(u\) and coarse remainder of norm \({\varrho}'\) such
that
\begin{align}
 |u-T_m\mathsf S_Lx|_1
                              &\le\eta_m(|x^\circ|_1+{\varrho}),\notag\\
 {\varrho}'&\le K_j|\mathsf S_Lx^\circ|_1+\eta_m {\varrho}.
 \label{eq:insertion-atom-prices}
\end{align}
Here, at fixed adjacent profiles,
\begin{equation}
          \eta_m\le\frac{A_j}{L}
                         +\omega_{j,L}(\widehat\beta_m),
     \qquad \omega_{j,L}(b)\longrightarrow0\quad(b\to\infty),
 \label{eq:insertion-pre-ratio-price}
\end{equation}
with \(A_j\) fixed before the ratio. The same construction applied
to the difference of two incoming remainder families has both
output prices at most \(\eta_m\|g-\widetilde g\|_{\rm ins}\).
It uses at most \(\max(q_m,d+1)\) regular input derivatives.
All identities and estimates are uniform in the physical period,
cutoff and finite history.
\end{lemma}
\begin{proof}
Introduce a commuting symbol \(z\) with \(z^2=0\). Multiply the
regular exponential by
\[
 1+z\left(x_0+\chi_{\Sigma_d}P_d[\mathsf D_{h_f^{-1}}x^\circ]
                                      +\sum_i\chi_{Z_i}R_i\right)
\]
and replace every forced factor by \(K_\alpha+zQ_\alpha\).
The coefficient of \(z\) is exactly
\eqref{eq:one-insertion-density}. The physical component expansion
and reassembly of Sections~\ref{sec:physical-reassembly}
and \ref{sec:nonlinear-induction-step} can be performed
coefficientwise. A word with one mark has
\[
       [z]\prod_{i=1}^l(a_i+zb_i)
                       =\sum_{i=1}^l b_i\prod_{j\ne i}a_j.
\]
Root its connected support at the distinguished factor.
Lemma~\ref{lem:rooted-interacting-log} bounds all other factors
by the same absolutely convergent rooted tree sums as in the
vacuum step. There is no sum over translates of the marked anchor.
The retained precision is the common reserve of
Section~\ref{sec:central-derivative-estimates}; the marked factor
does not select another Gaussian law.

We detail the size of the three possible distinguished factors.
For a primary section the zero-fluctuation, linear-interpolation
reading is \(T_m\mathsf S_Lx\). Its remaining local function has
norm at most \(K_j|\mathsf S_Lx^\circ|_1\), by the finite section
bounds, the interpolation differences and the finite product rule.
The coefficient error comes from nonlinear interpolation,
fluctuations, retained measure factors and cutoff comparisons.
On a fixed block the nonlinear interpolation differs from its
linear part by \(O_{j,L}(h_c^2)\), with the corresponding
threshold-scaled derivatives. Dividing a curvature factor by
\(h_f\) leaves \(O_{j,L}(h_c)\). A Gaussian fluctuation divided
by \(h_f\) has size \(O_{j,L}(P_f^{-1})\) in every fixed
Gaussian moment. A contraction of two such factors costs
\(O_{j,L}(P_f^{-2})\). Thus these coefficient errors tend to zero
at fixed \(L\). Polynomial interpolation errors left after the
degree-\(d\) reading have the additional inverse spatial scale
specified below. The unit has no such error: its push is the
unmarked push and its coordinate is exactly \(L^d\,1\).

For a nonwrapping regular remainder, first use a local site gauge
to put the slowly varying interpolated field in radial gauge at
the marked anchor. Its polynomial Taylor terms have weights
\(2,3,\ldots\); pure gauge constant potentials do not contribute
to a gauge-invariant function. The gauge is obtained by parallel
transport along radial segments; expanding its integral equation
gives the same finite-order radial jets as
\eqref{eq:insertion-test-holonomy}. The free-difference bounds
just used and the discrete Taylor formula give, for a term of
weight \(w\), a direction bounded by
\[
               C_{j,d}h_c L^{-w}(1+|z|)^{s_f}
\]
in fine anchored coordinates. The exponential kernel moments
make the constant independent of the support radius and \(B\).
The test-jet condition \eqref{eq:marked-jet-null} cancels every
contribution of total weight at most \(d\). The first retained
weight is at least \(d+1\). Hence its norm, after multiplication
by \(L^d\), is at most \(A_jL^{-1}\) times its input norm.

Here the derivative count matters. If an output derivative word
already contains total weight at least \(d+1\), estimate that
word directly. Otherwise apply Taylor's integral remainder only
to its missing weights. At most \(d+1\) input derivatives are
then used in total. A word of length greater than \(d+1\) is
estimated directly, with its corresponding interpolation factors.
Thus the required order is \(\max(q_m,d+1)\), not
\(q_m+d+1\). The nonlinear and Gaussian corrections to the
low-order reading have the vanishing factors described in the
preceding paragraph. High derivative words keep their inverse
interpolation powers and need no continuity modulus of the highest
input derivative. This is the higher-weight version of the ordered
jet argument in Lemma~\ref{lem:incoming-fifth-order-gain}.

All support estimates are made before summing the marked
addresses. The connected-cover compression of
Lemma~\ref{lem:incoming-cover-compression} converts a fine
anchored weight into the coarse one. Its finite enlargement
constant is included before \(L\). Extra powers of the covering
number from the Taylor product rule are paid from two fixed units
of the reserved exponent, using
\(\sup_{r\ge0}(1+r)^ae^{-2r}<\infty\) for each fixed \(a\).
For supports beyond the local reading radius, keep the whole
function and use that same unused exponential weight. Splitting
at a fixed positive power of \(L\) bounds this tail by
\(C_{j,d}L^{-1}\). This also handles wrapping supports without
inventing a bad inventory. These operations leave the exact
physical masks on the retained input.

For a forced remainder, use its actual forced multiplier in the
packet expansion, with the ordinary factor at that address
replaced rather than multiplied again. The joint packet estimate
and physical-measure differentiation of
Sections~\ref{sec:hole-packet-integrals} and
\ref{sec:integrated-regular-derivatives} apply to this bounded
multiplier with all its Borel reads held fixed. The physical seed
retains a Gaussian/Haar price \(C_{j,L}\exp(-c_{j,L}P_f^2)\).
An output assigned to a nonempty inventory keeps its eligibility;
an empty-inventory output enters the regular marked family.
Dividing the marked forced prices by \(\sqrt{w_c}\), and writing
the input price as \(k_\alpha/\sqrt{w_f}\), costs only
\[
 C_{j,L}\widehat\beta_f^aP_f^bB_f^c
 \exp\{-c_{j,L}P_f^2+(P_f^{1/5}+P_c^{1/5})/2\}
                                      \longrightarrow0
\]
for fixed derivative order. The packet payment retains the
original Gaussian/Haar ratios; no probability of a separately
selected event is inserted.

Cutoff and chart-comparison pieces have the same vanishing price.
For the regular mark, its smooth part is differentiated before
masking; a derivative of a held Borel test is instead transferred
to the physical measure, as in the earlier comparison argument.
The resulting polynomial scores are paid by the retained
exponential warning or packet charge. There is no derivative of
a hard indicator. Local Taylor subtractions and their mask
restorations are exact, with the same attachment bijection as at
birth. All marked functions vanish at the identity after reading,
and nonwrapping regular functions have their degree-\(d\) jet
removed.

After the rooted sums, the preceding bounds have the form
\eqref{eq:insertion-pre-ratio-price}. The finitely many losses
from basis recombination and symmetry averaging are included in
\(A_j,K_j\). Summing the individual address prices gives
\eqref{eq:insertion-atom-prices}. Applying the construction to a
difference cancels the primary term. Every resulting rooted word
then has a distinguished remainder difference, giving its stated
Lipschitz price. Absolute majorants justify all coefficient
extractions and sums. None of these arguments contains a sum over
all sites of the physical torus.
\end{proof}

\subsection{Closing the coefficient graph}

First enlarge the ratio floor so that \(A_j/L<\tfrac12L^{-2/5}\),
and then postpone entry until
\(\omega_{j,L}(\widehat\beta_m)\le\tfrac12L^{-2/5}\).
We can therefore use \(\eta_m\le L^{-2/5}\). Because all columns
and basis constants were fixed first, the same ratio may be made
large enough that, with \(s=C_{J(m+1)}\), \(S=C_{J(m)}\),
\begin{equation}
 K_j\eta_m(1+s)\le\tfrac14,
       \qquad 8K_j\eta_ms\le1,
       \qquad 4K_j^2\le S.
 \label{eq:insertion-graph-reserves}
\end{equation}
The last inequality is the column choice
\eqref{eq:insertion-column-choice}, including a profile transition.
The first two follow from the ratio floor and entry just specified.

\begin{theorem}[An exact, uniformly bounded one-insertion step]
\label{thm:one-insertion-step}
The transport of an admitted fine coefficient graph has a
complex-linear invertible coordinate map \(x\mapsto u=L_m^{\rm ins}x\)
and a coarse coefficient graph such that
\begin{align}
 |x|_1&\le C_{J(m+1)}|L_m^{\rm ins}x|_1,\notag\\
 |L_m^{\rm ins}x|_1&\le C_{J(m+1)}L_m^d|x|_1,\notag\\
 \|g_m(u)\|_{\rm ins}&\le C_{J(m)}|u^\circ|_1.
 \label{eq:one-insertion-step-bounds}
\end{align}
Its represented density is exactly
\(L_m^d\mathcal P_m\mathcal M_{m+1}(x)\). The maps commute
with conjugation and signed coordinate permutations; they send
\(1\) to \(L_m^d1\) with zero remainder. On a fixed unmarked
step, the induced map of coefficient graphs contracts their
normalised difference by a factor at most \(1/2\).
\end{theorem}
\begin{proof}
Write \(y=\mathsf S_Lx\), so \(|x|_1\le|y|_1\le L^d|x|_1\).
The atom prices and \({\varrho}\le s|x^\circ|_1\) give
\[
 |u-T_my|_1\le\eta_m(1+s)|x^\circ|_1,
       \qquad {\varrho}'\le K_j|y^\circ|_1+\eta_ms|x^\circ|_1.
\]
Apply \(T_m^{-1}\), and absorb the error using
\eqref{eq:insertion-graph-reserves}. Do this also after projection
onto the nonconstant subspace, which \(T_m\) preserves. We obtain
\[
 |y|_1\le2K_j|u|_1,\qquad
 |u|_1\le2K_j|y|_1,\qquad
 |y^\circ|_1\le2K_j|u^\circ|_1.
\]
It follows that
\({\varrho}'\le2K_j|y^\circ|_1\le4K_j^2|u^\circ|_1
\le S|u^\circ|_1\). The linear coordinate map is injective by
the first bound, hence invertible in the fixed finite-dimensional
space. This defines the new graph in coordinate \(u\) and proves
\eqref{eq:one-insertion-step-bounds}.

To make the choice linear with fixed addresses, perform the
construction on a finite basis, take the union of its address
sets and pad absent entries by zero. Recombine linearly. Average
over signed permutations and conjugation, whose finite
recombination constants have already been included in the local
prices. The exact density identity is linear, so it is preserved.
For the unit use the exact unmarked push, which has zero remainder.
This prevents an arbitrary decomposition of the constant density
from introducing a spurious marked term.

For the contraction let \(D\) be the distance between two fine
graphs, normalised by \(s|x^\circ|_1\). At fixed fine coordinate,
their output readings and remainder families differ by at most
\(\eta_msD|x^\circ|_1\). Rebase the second output to the first
coordinate \(u\). Its graph bound costs at most
\(S\eta_msD|x^\circ|_1\). Hence their coarse graph distance is
at most
\[
 \frac{\eta_ms(1+S)}{S}\,2K_jD\le4K_j\eta_msD\le D/2.
\]
This comparison uses the same unmarked step on both sides. It
does not assert convergence of different-cutoff vacuum histories.
\end{proof}

\subsection{The actual local insertions}

Iterate the theorem on the birth graph of
Proposition~\ref{prop:marked-insertion-birth}, down to \(\operatorname{ref}(d)\).
For \(n\ge \operatorname{ref}(d)\), define the terminal coordinate map
\[
       \mathsf Z_{n,M,d}
             =L_{\operatorname{ref}(d)}^{\rm ins}\cdots L_{n-1}^{\rm ins}.
\]
An empty product is the identity. This map is invertible and is
chosen at the origin before placing any translates. Define
\begin{equation}
 \mathcal O_{n,M}((d,v),o;U)
   ={\varepsilon}_n^{-d}P_d[\mathsf D_{h_n^{-1}}
                         \mathsf Z_{n,M,d}^{-1}v](o;U),
                  \qquad v\in V_d.
 \label{eq:constructed-local-insertion}
\end{equation}
For the finitely many \(n<\operatorname{ref}(d)\), choose any fixed invertible
terminal map, for example the identity. These finitely many
initial values do not affect a cutoff tail. The lattice label set
can now be taken to be the disjoint union of the \(V_d\); its
relation to continuum field quotients will be specified later.

\begin{corollary}[Reference reading and Wilson regularity]
\label{cor:local-insertion-regularity}
The insertions \eqref{eq:constructed-local-insertion} are complex
linear in \(v\), continuous, gauge invariant and local on the fixed
stencil \(\Sigma_d\). Their translated definitions use the same
terminal map. At the reference level their polynomial coordinate
is exactly \(v\), and their anchored remainder satisfies
\(\|g_{\operatorname{ref}(d)}(v)\|_{\rm ins}\le C_{J(\operatorname{ref}(d))}|v^\circ|_1\).
They respect conjugation and signed coordinate permutations.
For \(v={\varepsilon}_{\operatorname{ref}(d)}^d1\), the insertion equals one identically for
\(n\ge \operatorname{ref}(d)\). Together with the coupling asymptotic of
Theorem~\ref{thm:finite-cutoff-nonlinear-induction}, these
properties give Definition~\ref{def:regularity}.
\end{corollary}
\begin{proof}
The reading identity is the definition of \(\mathsf Z\); the
graph bound follows at each step from
Theorem~\ref{thm:one-insertion-step}. The local section and its
placement prove continuity, gauge invariance and locality. A
cutoff-dependent coefficient map does not alter the links read.
The symmetry assertions follow by composition and inversion of
the equivariant maps. Exact transport of the unit gives
\(\mathsf Z_{n,M,d}({\varepsilon}_n^d1)={\varepsilon}_{\operatorname{ref}(d)}^d1\).
Substitute its inverse in
\eqref{eq:constructed-local-insertion} and use
\(\mathsf D_h1=1\), \(P_d[1]=1\). Translation is literal
translation of the same polynomial; no coefficient is refitted at
a translated site.
\end{proof}

The new choices above are additional preassigned ratio floors and
coupling entries for the same calendar construction. At a fixed
profile there are finitely many of them, and each vanishing error
is postponed only after its ratio is fixed. They can therefore be
included in the common choice before any Wilson cutoff or volume
subsequence is selected. The resulting family has actual
finite-stencil renormalised insertions; volume-uniform bounds on
their products require the several-insertion argument next.
% !TeX root = 4DYM_notes.tex
\section{Joint insertions and volume-uniform separated moments}
\label{sec:joint-insertion-bounds}

One-insertion bounds do not bound a product by multiplying its
separately transported representatives. Integration creates terms
containing two or more insertion names. We retain those terms in a
finite coefficient algebra and estimate the complete source. The
rooted expansion is the same expansion as in the vacuum argument,
with one distinguished insertion and controlled attachments from
all the others. This organisation follows
\cite[Lemmas 14.12--14.13, PDF pp.~135--137]{ArmstrongDraft}.

\subsection{The coefficient algebra and the exact source}

Let \(I_0\) be a finite set of insertion names. Work in the
commutative algebra
\[
 \mathcal A_{I_0}=\C[z_i:i\in I_0]/(z_i^2:i\in I_0).
\]
For \(I\subset I_0\), write \(z_I=\prod_{i\in I}z_i\) and
\([z_I]a=a_I\) when \(a=\sum_Ia_Iz_I\). Distinct variables
commute. Their squares are zero to ensure that each named
insertion is counted at most once. For positive numbers \(\rho_i\)
put
\[
              \|a\|_\rho=\sum_{I\subset I_0}|a_I|
                                             \prod_{i\in I}\rho_i.
\]
Multiplication gives
\((ab)_I=\sum_{J\sqcup K=I}a_Jb_K\), whence
\begin{equation}
       \|ab\|_\rho\le\|a\|_\rho\|b\|_\rho,\qquad
       |[z_I]a|\prod_{i\in I}\rho_i\le\|a\|_\rho.
 \label{eq:joint-marker-norm}
\end{equation}
The nonconstant ideal is nilpotent. Its exponential and logarithm
are finite sums, and \(e^{z_if}=1+z_if\). Thus no analytic
continuation in a physical source parameter is needed.

For physical insertions \(O_i=\mathcal O_{n,M}((d_i,v_i),y_i)\),
define the exact source densities
\begin{equation}
 \mathcal F_n(z;U)=F_n(U)\prod_{i\in I_0}(1+z_iO_i(U)),
             \qquad \mathcal F_m=\mathcal P_m\mathcal F_{m+1}.
 \label{eq:joint-exact-source}
\end{equation}
Each coefficient is transported by the original integral operator.
For every \(I\subset I_0\), mass preservation gives
\begin{equation}
 \frac{\int[z_I]\mathcal F_m\,\dd H_m}
      {\int F_m\,\dd H_m}
                     =\left\langle\prod_{i\in I}O_i\right\rangle_{n,M}.
 \label{eq:joint-exact-mass}
\end{equation}
This is one denominator for the whole product.

% !TeX root = 4DYM_notes.tex
We will need signed replacements as well as the positive-sign
source in \eqref{eq:joint-exact-source}. Their common algebra
is the following finite identity. A density push \(\mathcal P\)
is a linear map from fine densities to coarse densities.
Write \(F_{\rm c}=\mathcal P F\). The fine functions are
\(O_i\), the coarse functions are \(V_i\), and all products
below are pointwise products on the indicated configuration
space.

\begin{lemma}[Complete subset coefficients and signed replacements]
\label{lem:complete-subset-identity}
For a finite name set \(I_0\), put
\[
 \mathcal D(z)=\mathcal P\!\left(F\prod_{i\in I_0}(1+z_iO_i)\right)
                                      \prod_{i\in I_0}(1-z_iV_i).
\]
For every \(I\subseteq I_0\),
\begin{equation}
 [z_I]\mathcal D
   =\sum_{J\subseteq I}(-1)^{|I\setminus J|}
       \mathcal P\!\left(F\prod_{j\in J}O_j\right)
                                      \prod_{i\in I\setminus J}V_i.
 \label{eq:common-signed-subset-coefficient}
\end{equation}
More generally, for a finite product of marker polynomials,
\begin{equation}
 [z_I]\prod_{\gamma\in\Gamma}a_\gamma(z)
  =\sum_{\substack{(I_\gamma)_{\gamma\in\Gamma}\text{ pairwise disjoint}\\
                           \bigcup_\gamma I_\gamma=I}}
                              \prod_{\gamma\in\Gamma}a_{\gamma,I_\gamma}.
 \label{eq:common-subset-product}
\end{equation}
Suppose also that \(\mathcal K\) is the probability pullback
adjoint to \(\mathcal P\), that \(Z=\int F\,\dd H_{\rm f}>0\),
and that for every \(A\subseteq I\)
\begin{equation}
       \mathcal K\!\left(\prod_{i\in A}V_i\right)
                              =\prod_{i\in A}\mathcal K V_i.
 \label{eq:common-pullback-product-hypothesis}
\end{equation}
Then the coefficient in
\eqref{eq:common-signed-subset-coefficient}, integrated and
divided by the original \(Z\), is exactly
\begin{equation}
 \frac1Z\int[z_I]\mathcal D\,\dd H_{\rm c}
       =\frac1Z\int F\prod_{i\in I}(O_i-\mathcal K V_i)\,\dd H_{\rm f}.
 \label{eq:common-signed-expectation}
\end{equation}
The density identity itself requires no product property of
\(\mathcal K\).
\end{lemma}
\begin{proof}
In the product defining \(\mathcal D\), a name in \(I\)
is chosen either from its fine factor or from its coarse
factor. Choosing it from both gives \(z_i^2=0\).
If \(J\) is the set chosen on the fine side, the remaining
names contribute the sign \((-1)^{|I\setminus J|}\).
Linearity of \(\mathcal P\) proves
\eqref{eq:common-signed-subset-coefficient}. The same
selection argument for an arbitrary collection of factors
proves \eqref{eq:common-subset-product}, including factors
whose zero coefficient is zero. No factor is divided by
its zero coefficient.

Integrating a term of
\eqref{eq:common-signed-subset-coefficient} and using
adjointness gives
\[
 \int F\Bigl(\prod_{j\in J}O_j\Bigr)
                 \mathcal K\!\Bigl(\prod_{i\in I\setminus J}V_i\Bigr)
                                                       \,\dd H_{\rm f}.
\]
Use \eqref{eq:common-pullback-product-hypothesis} in this
line, and sum over \(J\). The result is the ordinary
expansion of \(\prod_{i\in I}(O_i-\mathcal K V_i)\).
This proves \eqref{eq:common-signed-expectation} with the
single original denominator. The empty coefficient is
\(F_{\rm c}\).
\end{proof}

Coefficient extraction in this lemma has norm at most
\(\prod_{i\in I}\rho_i^{-1}\), by
\eqref{eq:joint-marker-norm}; multiplication is
submultiplicative in that same norm. The estimates have no
additional factor counting subsets or reflected cells.
Their analytic inputs remain the previously proved activity,
integration and connected-component bounds. For an absolutely
convergent labelled expansion, apply
\eqref{eq:common-subset-product} to each finite term and
then use its absolute majorant to regroup the series.
The signed sum is formed before taking its absolute value.

In particular, an incoming forced member is still one factor
in each original term, even when several marked reads meet
it. Its full marker polynomial is used in
\eqref{eq:common-subset-product}; its mask, inventory and
allowance occur once. This observation also applies to
restored higher modes. Disjoint full read cones, when available,
prove \eqref{eq:common-pullback-product-hypothesis} by an
additional kernel argument; they are not a consequence of the
finite marker algebra. These identities hold at every cutoff
and period and alter neither the Wilson measure nor any
Gaussian or Haar normalisation.

In the local representation, the source has the form
\begin{align}
 \mathcal F_m(z;U)={}&
 \exp\left\{\Phi_m(U)+\sum_i z_if_{i,m}(U)
                              +\sum_X H_{m,X}(z;U)\right\}\notag\\
 &\quad\cdot\sum_{A\ {\rm admissible}}
  \prod_{\alpha\in A}\left(K_{m,\alpha}(U)
      +\sum_i z_iq_{i,m,\alpha}(U)
                              +Q^{\ge2}_{m,\alpha}(z;U)\right).
 \label{eq:joint-local-source}
\end{align}
Here \(f_{i,m}\) is the complete singleton exponential function
and \(q_{i,m,\alpha}\) its complete forced replacement. In terms
of the one-insertion representation, both include the factor
\({\varepsilon}_m^{-d_i}\). This follows from
\({\varepsilon}_m^{-d_i}L_m^{d_i}={\varepsilon}_{m+1}^{-d_i}\).
The polynomials \(H_{m,X}\) and \(Q^{\ge2}_{m,\alpha}\) contain
only modes with at least two names. A nonzero mode has a support
containing all of its named anchors. Regular modes carry their
health masks; forced modes retain their exact inventory and every
auxiliary eligibility test. Addresses of unmarked activity zero
are allowed as before.

Formula \eqref{eq:joint-local-source} is produced by the physical
connected expansion, not by taking a logarithm of the full
possibly signed forced gas. At the unmasked initial source its
higher modes are zero. Restoring masks at birth can generate
higher forced modes; the attachment construction below includes
them. They are not set to zero after that restoration.

\subsection{Separation pays for changing source radii}

\begin{lemma}[Covering cost of several separated names]
\label{lem:joint-marked-cover}
Let \(B\ge8\). Suppose the distinct named anchors in a footprint
\(X\) have pairwise periodic maximum-distance at least \(B^2\).
If \(I\) is their name set and \(|I|\ge2\), then
\begin{equation}
                         B|I|\le64n_B(X).
 \label{eq:joint-marked-cover}
\end{equation}
The estimate is independent of the torus period.
\end{lemma}
\begin{proof}
Take a minimal connected cover with \({\mathsf M}=n_B(X)\) balls. Their
intersection-or-adjacency graph has a spanning tree whose edges
have length at most \(2B+1\). Its doubled traversal has length
at most \(6B({\mathsf M}-1)\). Assign each named anchor to a covering
centre. Different anchors have different centres, and two assigned
centres are separated by at least \(B^2-2B\). Visit the assigned
centres in their cyclic order in the doubled traversal. The triangle
inequality, including the return segment, gives
\[
              |I|(B^2-2B)\le6B({\mathsf M}-1).
\]
For \(B\ge8\) this implies the displayed, weaker bound. All
distances are periodic; neither a Euclidean lift of \(X\) nor a
nonwrapping assumption is used.
\end{proof}

In particular, removing \(\delta\) units of support exponent
allows every name radius in a higher mode to increase by
\(e^{\delta B/64}\). This is a statement about nonnegative
coefficient majorants, not about cancellations between modes.

Fix a maximal degree \(D\). For the finitely many sections through
that degree, choose an integer \({\varepsilon}_D\) large enough that their
physical evaluations, derivative products and marked restorations
cost at most \(A_{j,D}\widehat\beta_m^{{\varepsilon}_D}\) times their
dimensionless singleton coefficient size. Such an integer exists:
there are finitely many polynomial columns and derivatives at a
profile, \(h_m^{-1}\le\widehat\beta_m\), and the leftover
polynomials in \(B_m,P_m\) are subpowers of
\(\widehat\beta_m\). Threshold-scaled derivatives introduce
polynomials in \(P_m,B_m\), rather than an increasing power of
\(\widehat\beta_m\) for every derivative: a Gaussian translation
score has scale \(h_m\sqrt{\beta_m}\asymp P_m\). At each fixed profile,
its finitely many such polynomials can be paid by one extra
coupling power after entry. Thus \({\varepsilon}_D\) can be fixed with the
degree, while \(A_{j,D}\) and the entry retain the profile and
derivative-order dependence. Define
\[
                       \zeta_{D,m}
                                  =\widehat\beta_m^{-({\varepsilon}_D+4)}.
\]
At a starting level \(k\), choose source radii so that
\(\rho_{i,k}{\varepsilon}_k^{-d_i}|x_{i,k}|_1\le\zeta_{D,k}\), and at
finer levels put
\begin{equation}
                \rho_{i,m+1}=\widehat\beta_{m+1}^{-1}\rho_{i,m}.
 \label{eq:joint-radius-transport}
\end{equation}
These are estimating weights; the final coefficient is still
extracted at radius \(\rho_{i,k}\).

The inverse coordinate estimate
\eqref{eq:one-insertion-step-bounds} gives
\[
 {\varepsilon}_{m+1}^{-d_i}|x_{i,m+1}|_1
       \le L_m^{d_i}C_{J(m+1)}{\varepsilon}_m^{-d_i}|x_{i,m}|_1.
\]
At fixed profiles and ratio, entry can therefore ensure
\begin{equation}
 \widehat\beta_{m+1}^{-1}L_m^DC_{J(m+1)}\zeta_{D,m}
                                                   \le\zeta_{D,m+1}.
 \label{eq:joint-radius-entry}
\end{equation}
After division by \(\zeta_{D,m+1}\), its left side is a fixed
constant times \(\widehat\beta_{m+1}^{-1}
(\widehat\beta_{m+1}/\widehat\beta_m)^{{\varepsilon}_D+4}\), which tends
to zero. The radii-weighted singleton size is consequently bounded
by \(\zeta_{D,m}\) throughout the interval.

\begin{lemma}[Uniform attachment income for the full source]
\label{lem:joint-source-attachments}
Suppose the anchors have separation at least \(B_m^2\) at every
level \(k\le m\le n\), and their parents agree with the named
anchors at the next coarser level. Let \(v_{a,m}^{\rm src}\) be
the integrated, radii-weighted majorant of a marked source factor
on its full physical read support \(Z_{a,m}\). Derivative and
covering weights are included in this majorant. The parameters
can be chosen before cutoff and volume so that
\begin{equation}
 (2B_m+1)^4\sup_y\sum_{a:y\in Z_{a,m}}v_{a,m}^{\rm src}
                                                    \le\eta_m^{\rm att}/2
 \label{eq:joint-source-attachments}
\end{equation}
throughout the interval, independently of the number of names.
Here \(\eta_m^{\rm att}>0\) is a valid small attachment threshold
for the rooted component expansion. Together with the unmarked
income, the complete attachment sum is at most
\(\eta_m^{\rm att}\).
\end{lemma}
\begin{proof}
Choose the threshold so that it is at most \(1/4\), and make the
unmarked income at most half of it using the already established
factor rates. The marker norm \eqref{eq:joint-marker-norm} is
submultiplicative. Apply the rooted connected majorant with one
distinguished marked occurrence, retaining the usual separate
derivative and support estimates. This use does not assert that
the original unweighted derivative seminorm is an algebra norm;
the finite-jet product bound of
Lemma~\ref{lem:finite-jet-algebra} supplies its product rule, and
the profile columns pay its finite derivative factorial. A
physical marked multiplier, including its mask, is bounded by
its amplitude majorant. Bounding these multipliers does not
spend a new Gaussian precision reserve for each name; the common
positive reserve controls the whole physical word as before.

Consider first the sum over possible marked anchors for a support
which also contains \(y\). If that anchor is at distance \(R\)
from \(y\), connected covering implies
\(n_{B_m}(Z)\ge R/(2B_m+1)\), up to an inessential fixed
endpoint factor. One reserved unit of support exponent pays
\(Ce^{-cR/(2B_m+1)}\). Separation implies occupancy at most one,
and hence
\[
 \sum_{\text{anchors }o}Ce^{-c d_\infty(o,y)/(2B_m+1)}
                                              \le C'(1+B_m)^4.
\]
For a periodic torus, choose nearest representatives; their sum is
bounded by the full cover sum. The remaining rooted descendants
cost at most a factor two once their complete attachment income
is below the chosen threshold. A second reserved support unit
pays their root-covering factor. The local singleton bound and
\eqref{eq:joint-radius-entry} thus give the numerical income
\begin{equation}
 2A_{j,D}(2B_m+1)^4(1+B_m)^4
           \widehat\beta_m^{{\varepsilon}_D}\zeta_{D,m}
                                                    \le\eta_m^{\rm att}/2.
 \label{eq:joint-attachment-entry}
\end{equation}
At fixed ratio this is an entry requirement: its left side is a
fixed profile constant times a polynomial in \(\log\widehat\beta_m\)
times \(\widehat\beta_m^{-4}\).

There is a second requirement when transporting from the finer
source radii to the coarser ones. On a mode \(I\) with at least
two names the increase is \(\widehat\beta_{m+1}^{|I|}\).
Lemma~\ref{lem:joint-marked-cover} bounds this by
\[
 \exp\left\{\frac{64\log\widehat\beta_{m+1}}{B_m}
                                      n_{B_m}(X)\right\}.
\]
Its coefficient tends to zero, since
\(B_m\asymp(\log\widehat\beta_m)^2\). At entry make it at
most \(1/4\). Cover compression supplies the reserved exponent
at the next step, so it is not depleted once for every level of
the history. Singleton modes use their exact one-insertion maps
and \eqref{eq:joint-radius-entry}, rather than this several-name
covering estimate.

These observations close a descending induction. At the unmasked
top source there are only singleton exponential terms. Mask
restoration merges all forced factors meeting the marked stencil,
as in Proposition~\ref{prop:marked-insertion-birth}. If several
stencils belong to the same merged component, its mode contains
all their names and its support contains all their anchors. The
same covering estimate pays its coefficient weights. Original
forced factors occur once and all inventory tests are retained.
More explicitly, at a masked birth a component meeting \(b\)
stencils has weighted source multiplier at most
\((1+C\widehat\beta_n^{{\varepsilon}_D}\zeta_{D,n})^b\).
For \(b\ge2\), \eqref{eq:joint-marked-cover} bounds its logarithm
by \(64C\widehat\beta_n^{-4}n_{B_n}(Y)/B_n\), paid from the
same unused support reserve. The singleton case is
\eqref{eq:marked-birth-price}. This supplies the initial
majorants, including any higher modes
created by the restoration.

At a descending step the inductive complete income gives an
absolutely convergent rooted expansion. Its output coefficients
have the same named-support property. Enlarge their radii by the
just-paid factor and apply \eqref{eq:joint-attachment-entry} to
obtain the same bound for the output. The ordinary rooted-tree
majorant is at most twice its distinguished income under the
chosen smallness; this is exactly the factor two already allowed
there. Thus the induction includes previously formed connected
source terms, rather than treating all other marks as bare
singletons. All bounds apply to positive absolute majorants and
are uniform in finite truncations; monotone convergence passes
them to the full address families. This proves the claim.
\end{proof}

\subsection{A partition estimate with its volume cost explicit}

Let \(D_G=2N^2\) be the real dimension of the ambient matrix space
of the fixed \(N\)-dimensional unitary representation. It is only
an upper bound on the dimension of \(G\).

\begin{lemma}[A scalar-free vacuum ratio]
\label{lem:insertion-vacuum-ratio}
There are profile constants \(A_j\ge2\), fixed before the ratio,
such that at level \(m\), on \(V=\#\Lambda_m\) sites,
\begin{align}
 \frac{\displaystyle\int e^{\Phi_m(U)}
                \sum_{A\ {\rm admissible}}
                    \prod_{\alpha\in A}|K_{m,\alpha}(U)|\,\dd H_m(U)}
      {\displaystyle\int F_m\,\dd H_m}
                  &\le\exp\{VQ_m^{\rm vac}\},\notag\\
 Q_m^{\rm vac}&=A_{J(m)}+4D_G
                    \log(A_{J(m)}\sqrt{\widehat\beta_m}).
 \label{eq:insertion-vacuum-ratio}
\end{align}
The vacuum scalar cancels. The estimate itself still has a volume
cost; taking a reflected-cell root below removes that cost.
\end{lemma}
\begin{proof}
The leading clipped row action is nonnegative. The native cubic,
quartic and loop terms, and the masked regular sum, have total
absolute value at most \(A_jV\) after entry. This follows by
their path budgets: their local sizes are bounded by constants
times \(\beta_mh_m^3\), \(\beta_mh_m^4\),
\(\Lambda_mh_m^2\) and \(\delta_m\), respectively. Every one
tends to zero at a fixed profile; a fixed larger \(A_j\) also
covers the initial levels. Dropping forced constraints bounds
their absolute sum by \(e^{Vw_m}\). This bounds the numerator
after its scalar has been removed.

For the denominator restrict every link to a group ball of chord
radius \(g_m=(A_j\sqrt{\widehat\beta_m})^{-1}\). Increasing
\(A_j\) makes every configuration in this set healthy: each
plaquette has chord at most \(4g_m<h_m\). The forced gas is
then one. The same path bounds and the finite row Schur sums
give principal action at most \(A_jV\), since
\(\beta_mg_m^2\le2/A_j^2\). The other terms are bounded as
above, so the scalar-free density is at least \(e^{-A'_jV}\)
on this set.

Haar measure of a chord ball of radius \(0<g\le1\) is at least
\((g/C)^{D_G}\), with \(C\) depending only on the fixed
representation. To verify this elementary bound, take a maximal
\(g/2\)-separated subset of \(G\) in the ambient Euclidean
metric. Disjoint ambient balls and volume comparison bound its
cardinality by \((C/g)^{D_G}\). The corresponding group balls
cover \(G\); left invariance makes their Haar measures equal.
Their sum is at least one, proving the bound. There are \(4V\)
independent positive links. The restricted set therefore has
Haar measure at least \((g_m/C)^{4VD_G}\). Combine this lower
bound with the numerator upper bound and enlarge \(A_j\).
This proves \eqref{eq:insertion-vacuum-ratio}, without replacing
the denominator by one or assuming positivity of individual
forced activities.
\end{proof}

\subsection{Replicated sources and reflection positivity}

We impose the further calendar requirements
\begin{align}
 B_{m+1}^2&\le2B_m^2,\notag\\
 (L_mB_m^2)^4Q_{m+1}^{\rm vac}
                         &=o(s_{m+1}),\notag\\
 \frac{\log\mathcal C_j}{\log L_j}&\le j^{-2},
 \qquad 4\mid M.
 \label{eq:insertion-root-calendar}
\end{align}
Here \(\mathcal C_j\ge2\) bounds all column and source constants
through profile \(j\), and is fixed before \(L_j\). The last
divisibility condition is part of the initial choice of periods.
All ratios remain odd. These requirements are compatible with
the preceding calendar: choose the ratio after \(\mathcal C_j\),
then postpone profile entry. At a fixed profile the second
quotient is a fixed ratio power times a polynomial in
\(\log\widehat\beta_m\), divided by the accumulated scale, and
tends to zero. The first follows from the fixed-shift comparison
of neighbouring nominal couplings. The finite initial profile is
handled by its initial coupling and period margins.

For a bond tiling into congruent reflected cells, an interior
observable uses only links whose two endpoints are strictly
inside its cell. Its reflected copy in another cell includes
complex conjugation whenever an odd number of reflections is
used. Define its reflected-cell root by
\[
 \|O\|_{\mathcal T}=
       \left|\left\langle\prod_{C\in\mathcal T}(O)_C
                                       \right\rangle_{n,M}\right|^{1/|\mathcal T|}.
\]
For the centred tiling of physical side \({\varepsilon}_m\), denote this
root by \(\|O\|_m\). Its cells have odd lattice side
\({\varepsilon}_m/{\varepsilon}_n\), so the centred boundaries pass through link
midpoints. The number of cells per axis is an even integer by
\eqref{eq:insertion-root-calendar}.

\begin{lemma}[Wilson reflection and the chessboard inequality]
\label{lem:insertion-chessboard}
The Wilson probability measure is reflection positive across
bond planes. For bounded gauge-invariant interior cell functions,
\begin{equation}
 \left|\left\langle\prod_{C\in\mathcal T}(O_C)_C
                                      \right\rangle_{n,M}\right|
                                  \le\prod_{C\in\mathcal T}\|O_C\|_{\mathcal T}.
 \label{eq:insertion-chessboard}
\end{equation}
These roots satisfy the triangle inequality and \(\|1\|=1\).
\end{lemma}
\begin{proof}
We use the finite-volume reflection-positivity and chessboard
theorems in \cite[Theorems 2.1--2.2, printed pp.~14--16]{Sei82}.
Their hypotheses in the present model can be checked directly.
The coupling is positive, the group is compact, and the action
is a sum of Wilson plaquette terms in the fixed unitary
representation. After setting the links crossing a reflection
plane to the identity by gauge invariance, each crossing
plaquette couples a half-loop matrix entry to the conjugate
entry of its reflected half. Its Boltzmann factor is an
exponential of a sum with nonnegative coefficients of terms
\(f\,\theta f\), where \(\theta\) is reflected pullback followed
by conjugation. Expanding the exponential preserves that cone.
Integrating the two half-configurations gives sums of squared
absolute integrals, hence \(\langle F\theta F\rangle\ge0\).
Interior observables do not read the crossing links. This
justifies the gauge reduction and avoids a boundary-link
dependence in the positivity assertion.

Schwarz's inequality for this nonnegative sesquilinear form,
applied successively to the two halves in each axis, is the
chessboard theorem. Even cell counts supply the periodic paired
halves; translation invariance supplies its cyclicity. It gives
\eqref{eq:insertion-chessboard}. Expanding the reflected product
of \(O+O'\) and applying that inequality bounds its absolute
value by \((\|O\|+\|O'\|)^{|\mathcal T|}\), proving the
triangle inequality. The normalised expectation of one is one.
Bounded Borel interior functions follow by approximation; our
original insertions are already continuous.
\end{proof}

\begin{proposition}[A cutoff- and volume-uniform insertion root]
\label{prop:uniform-insertion-root}
Fix a degree \(d\) and a bounded coefficient ball in \(V_d\),
enlarged to contain its signed and conjugated images. There are
positive radii \(\rho_k\), independent of the period, such that
whenever \(\operatorname{ref}(d)\le m\le k\) and \(B_k^2\le {\varepsilon}_m/{\varepsilon}_k\),
\begin{equation}
 \|\mathcal O_{n,M}((d,v),0)\|_m
       \le\rho_k^{-1}
               \exp\{({\varepsilon}_m/{\varepsilon}_k)^4Q_k^{\rm vac}+6\}
 \label{eq:uniform-insertion-root}
\end{equation}
for every sufficiently large \(n\), uniformly in \(M,v\).
Moreover
\begin{equation}
                   -\log\rho_k
                         =ds_k
                                      +o(s_k).
 \label{eq:insertion-radius-asymptotic}
\end{equation}
The cutoff threshold can depend on \(d\), the ball, \(m,k\),
but not on \(M\).
\end{proposition}
\begin{proof}
The reference coordinate has a fixed bound. Iterating the inverse
estimate in \eqref{eq:one-insertion-step-bounds} bounds its
dimensionless level-\(k\) size by
\(R\prod_{i<k}C_{J(i+1)}\), uniformly in \(n,M\). Absorb the
finite initial segment into \(R\). Set
\[
 E_k=R C_{J(k)}\prod_{i<k}C_{J(i+1)},\qquad
       \rho_k=\frac{{\varepsilon}_k^d\widehat\beta_k^{-b_d}}{E_k},
\]
where the fixed integer \(b_d\) is chosen above \(a_d+4\) and
all polynomial losses in the singleton bounds. Then the full
weighted singleton size is at most \(\zeta_{d,k}\).

Insert one reflected copy in every cell of side \({\varepsilon}_m\).
At level \(k\) their spacing is \({\varepsilon}_m/{\varepsilon}_k\ge B_k^2\).
It remains at least \(B_i^2\) at every finer level: spacing
increases by at least three per step and \(B_i^2\) by at most
two. Their parents agree, and reflection changes only the finite
signed and conjugated coefficient ball. Thus
Lemma~\ref{lem:joint-source-attachments} applies on the whole
interval \([k,n]\), using \(\rho_k\) as its final radius.

At level \(k\), the weighted complete singleton exponential
function has bound at most three per name; its forced counterpart
has bound at most one. The connected regular and forced source
families have absolute sums at most one each per name, by the
rooted attachment majorant. In the marker norm,
exponentiation and the eligible forced sum are therefore bounded
by \(e^{6{\mathsf M}}\) times the absolute unmarked expression, where
\({\mathsf M}\) is the number of reflected cells. Higher modes are included
in these two connected sums. Lemma~\ref{lem:insertion-vacuum-ratio}
and coefficient extraction give
\[
 \rho_k^{\mathsf M}\left|\left\langle\prod_{C\in\mathcal T_m}(O)_C
                                \right\rangle_{n,M}\right|
                    \le e^{\#\Lambda_kQ_k^{\rm vac}+6{\mathsf M}}.
\]
Exact transport preserves this normalised coefficient. Since
\(\#\Lambda_k/{\mathsf M}=({\varepsilon}_m/{\varepsilon}_k)^4\), taking the \({\mathsf M}\)-th root proves
\eqref{eq:uniform-insertion-root}. In particular the number of
physical sites has cancelled.

Finally \eqref{eq:insertion-root-calendar}, summed over levels,
gives \(\sum_{i<k}\log C_{J(i+1)}=o(s_k)\).
To verify this, split the sum before any fixed late profile;
the remaining summands are at most that profile's inverse square
times \(\log L_i\), and the initial finite sum vanishes after
division by the accumulated scale. The final column is controlled
by the same bound and the gentle mesh condition.
Also \(\log\widehat\beta_k=o(s_k)\).
The formula for \(\rho_k\) proves
\eqref{eq:insertion-radius-asymptotic}.
\end{proof}

For each \(m\), let \(k(m)\) be the first \(k\ge m\) with
\(B_k^2\le {\varepsilon}_m/{\varepsilon}_k\). It exists because \({\varepsilon}_kB_k^2\) decreases
by at least a factor \(2/3\) per step. If \(k=m\) the following
bounds are immediate; otherwise its minimality gives
\[
       {\varepsilon}_m/{\varepsilon}_{k(m)}\le L_{k(m)-1}B_{k(m)-1}^2.
\]
Consequently \eqref{eq:insertion-root-calendar} implies
\[
 ({\varepsilon}_m/{\varepsilon}_{k(m)})^4Q_{k(m)}^{\rm vac}
                      =o(s_{k(m)}).
\]
The logarithm of this ratio is also negligible relative to that
scale, by gentleness and \(\log B_i=o(s_i)\).
Thus \(s_{k(m)}/s_m\to1\).
In particular, for every \(\epsilon>0\), sufficiently late
centred roots have the useful bound
\begin{equation}
                 \|\mathcal O_{n,M}((d,v),0)\|_m
                                            \le {\varepsilon}_m^{-d-\epsilon}
 \label{eq:insertion-root-power}
\end{equation}
on a cutoff tail uniform in volume and the fixed coefficient ball.
Finite earlier levels have finite uniform constants from
\eqref{eq:uniform-insertion-root}.

\subsection{Moving anchors into a common tiling}

\begin{lemma}[Interior roots at uncentred anchors]
\label{lem:insertion-interior-root}
Suppose \(m\le k\), every centred root of a signed or conjugated
insertion in cells of side \({\varepsilon}_k\) is at most \(B_*\), and its
physical stencil radius is at most \({\varepsilon}_k\). In any translated
bond tiling of side \({\varepsilon}_m\), an anchor at distance at least
\(3{\varepsilon}_k\) from its cell boundary is interior and has root at most
\(B_*\).
\end{lemma}
\begin{proof}
First take an anchor at a centre of the aligned fine tiling.
Its reflected coarse copies occupy distinct fine cells. Apply
\eqref{eq:insertion-chessboard} to those copies, putting one in
each unoccupied fine cell. Taking the coarse-cell root proves
the bound at all aligned fine centres.

To move between adjacent such centres in one coordinate, let
\(H(z)\) be the coarse-cell root after \(z\) fine-lattice
translations of the anchor. Translating the entire replicated
product by one lattice step places shifts \(z-1\) and \(z+1\)
in its alternating reflected cells. Translation invariance and
\eqref{eq:insertion-chessboard} give
\(H(z)^2\le H(z-1)H(z+1)\). A finite nonnegative sequence with
this property cannot exceed the larger endpoint value: a positive
interior maximum forces both neighbours to have the same maximum,
and iteration reaches an endpoint. This argument also handles
zero values. Bracket each anchor coordinate by the neighbouring
aligned fine centres. Throughout the resulting rectangle the
margin is at least \(2{\varepsilon}_k\), so the stencil remains interior.
Apply the one-dimensional bound successively in the four
coordinates. This proves the assertion.
\end{proof}

\begin{theorem}[Uniform separated insertion bounds]
\label{thm:uniform-separated-insertions}
Fix finitely many degrees, bounded coefficient balls, an arity
\(r\ge1\), and a positive torus separation \(\delta\). There are
\(C<\infty\) and \(n_0\), independent of volume, such that
\begin{equation}
 \left|\left\langle\prod_{i=1}^r
          \mathcal O_{n,M}((d_i,v_i),y_i)\right\rangle_{n,M}\right|
                                                      \le C
 \label{eq:uniform-separated-insertions}
\end{equation}
whenever \(n\ge n_0\) and the physical torus distances between
the anchors are at least \(\delta\). For \(r=1\) no separation
condition is needed. The constants may depend on \(G,r,\delta\),
the specified coefficient balls and the common calendar.
They do not depend on cutoff or physical volume.
\end{theorem}
\begin{proof}
Choose a fixed \(m\) with \({\varepsilon}_m<\delta/8\), and then a fixed
\(k\ge m\) with \({\varepsilon}_k<{\varepsilon}_m/(100r)\). For sufficiently large
\(n\), every insertion stencil has physical radius at most
\({\varepsilon}_k\). Proposition~\ref{prop:uniform-insertion-root} gives
one common bound \(B_*\) for their centred roots in cells of
side \({\varepsilon}_k\), including every signed and conjugated image.
Its cutoff threshold is independent of \(M\).

Choose a translate of the side-\({\varepsilon}_m\) bond tiling with every
anchor at least \(3{\varepsilon}_k\) from its cell boundary. Such a translate
exists on the fine lattice: in one coordinate, one anchor excludes
at most \(6{\varepsilon}_k/{\varepsilon}_n+2\) of the \({\varepsilon}_m/{\varepsilon}_n\) possible translations.
For large \(n\), the union over \(r\) anchors excludes less
than half the translations. Choose each of the four coordinates
separately. Distinct anchors lie in different cells because their
torus separation exceeds the cell diameter in the maximum metric.
Lemma~\ref{lem:insertion-interior-root} bounds every occupied
cell root by \(B_*\). Apply the chessboard inequality with the
constant one in the other cells. The resulting bound is
\(B_*^r\), as required. The argument for one anchor uses any
fixed \(m\) and otherwise is the same.
\end{proof}

For a compact separated physical configuration set, the same
bound is uniform over all sufficiently large suitable periods:
their torus separation is bounded below by a positive number.
At each fixed smaller suitable period it holds on every compact
set separated also in that torus. Multilinearity gives the
corresponding operator norm on each finite list of label spaces.
No assertion about coincident-point products is made here.
% !TeX root = 4DYM_notes.tex
\section{A common ultraviolet extraction and a weak correlation hierarchy}
\label{sec:weak-hierarchy-extraction}

The preceding estimates provide a compactness route which is
sufficient for the first construction of correlations. The family
in the target theorem is existential. We may therefore choose one
ultraviolet subsequence, common to all fixed periods, as part of
that family before any requested infinite-volume sequence is
given. A subsequent volume extraction then leaves a full mesh
limit at each fixed period of the chosen family.

This observation simplifies the present step: it does not require
uniqueness of all ultraviolet vacuum trajectories. The insertion
input has just been developed in
Section~\ref{sec:joint-insertion-bounds}, following
\cite[Lemmas~14.12--14.13, PDF pp.~135--137]{ArmstrongDraft};
the countable diagonal
construction is given directly below. We separate what follows
from these bounds from the further smoothness and reconstruction
estimates. The result here is a hierarchy of
distributions on separated configurations. It is not yet the full
continuum theorem of Section~\ref{sec:continuum}.

\subsection{Refining the mesh without changing the order of choices}

Keep the reference unit \(u_d={\varepsilon}_{\operatorname{ref}(d)}^d1\in V_d\) fixed from
the original internal RG calendar. A refinement changes the
terminal cutoffs, not this reference normalisation.

\begin{lemma}[A fixed ultraviolet refinement preserves Wilson regularity]
\label{lem:fixed-uv-refinement}
Let \(\nu:\N\to\N\) be strictly increasing. Replace every datum
at cutoff \(n\) by its value at cutoff \(\nu(n)\), leaving the
physical side length and the insertion labels unchanged. Continuity,
gauge invariance, the fixed-stencil property and the uniform
inverse-coupling slope in Definition~\ref{def:regularity} are
preserved. Uniform estimates on cutoff tails are preserved as well.
If \(\nu\) is selected before a volume sequence is requested,
this is one Wilson family in the quantifier order
\eqref{eq:public-target}.
\end{lemma}
\begin{proof}
The new mesh at cutoff \(n\) is \(\varepsilon_{\nu(n)}\).
It decreases strictly to zero, and its reciprocal remains a positive integer.
The local functions at each new cutoff are precisely old local
functions on the same stencil; their continuity and invariance
are unchanged. A uniform tail estimate remains valid once
\(\nu(n)\) exceeds its original threshold. In particular the
coupling ratio and its logarithmic denominator are both read at
\(\nu(n)\), which preserves the slope. The refinement is one
fixed sequence, independent of every later volume request.
The internal RG histories keep their original levels; no claim
that one new cutoff interval is a single original RG step is
needed.
\end{proof}

\subsection{Compactness at every fixed period}

For the physical period \(M\), define the open separated configuration domain
\[
 \Omega_{r,M}=
 \left\{z\in(-M/2,M/2)^{4r}:z_i\ne z_j\ (i\ne j)\right\}.
\]
For \(r=0\) the domain is one point. A compact subset of
\(\Omega_{r,M}\) has positive torus separation: a difference
of two distinct interior representatives cannot be a nonzero
period vector, and compactness bounds its distance from those
vectors away from zero.

For a degree list \(\mathbf d=(d_1,\ldots,d_r)\), label list
\(v_i\in V_{d_i}\), and \(f\in C_c^\infty(\Omega_{r,M})\),
use the scale-index notation
\(S_{n,M}(\mathbf v;f)=S_{M,\varepsilon_n}(\mathbf v;f)\)
for the lattice smearing
\eqref{eq:smear}, with the insertions
\eqref{eq:constructed-local-insertion}. In particular it uses
the centered placement convention and the factor \({\varepsilon}_n^{4r}\).

\begin{lemma}[Uniform bounds for fixed-period smearings]
\label{lem:uniform-row-smearings}
For every fixed \(M,r,\mathbf d\) and compact
\(K\Subset\Omega_{r,M}\), there is a cutoff tail and a constant
\(A_{K,M,\mathbf d}<\infty\) such that
\begin{equation}
 |S_{n,M}(\mathbf v;f)|
       \le A_{K,M,\mathbf d}\prod_i|v_i|_1\,\|f\|_\infty,
                  \qquad \supp f\subset K.
 \label{eq:uniform-row-smearings}
\end{equation}
For a fixed compact \(K\Subset\operatorname{Conf}_r(\R^4)\),
the constant and cutoff tail can be chosen uniformly over all
sufficiently large periods. For each fixed test \(f\), one also
has the sharper limiting upper bound
\begin{equation}
 \limsup_{n\to\infty}|S_{n,M}(\mathbf v;f)|
       \le C_{K,M,\mathbf d}\prod_i|v_i|_1
                                         \int|f(z)|\,\dd z.
 \label{eq:row-l1-upper-bound}
\end{equation}
\end{lemma}
\begin{proof}
Theorem~\ref{thm:uniform-separated-insertions} bounds every
moment evaluated at a point of \(K\) on a common cutoff tail.
Multilinearity extends its bounded-ball estimate to the displayed
product of coefficient norms. If \(K\subset[-R,R]^{4r}\), the
number of possible grid points is at most
\((2R/{\varepsilon}_n+2)^{4r}\). Multiplication by \({\varepsilon}_n^{4r}\), with
\({\varepsilon}_n\le1\), proves \eqref{eq:uniform-row-smearings}.

For large periods the support lies in the interior box and its
torus separation is bounded below by a fixed fraction of its
Euclidean separation. The uniform clause of the insertion theorem
then proves the stated uniformity in \(M\). At a fixed smaller
suitable period use its own positive compact torus separation;
this is a genuine fixed-period tail, not a statement only along
a diagonal of growing volumes and cutoffs.

Finally, for a fixed continuous compactly supported \(f\), the
Riemann sums \({\varepsilon}_n^{4r}\sum|f|\) converge to \(\int|f|\).
One elementary proof partitions its bounding box into lattice
cubes; uniform continuity bounds the error on each cube by the
modulus of continuity of \(|f|\) at diameter \(2{\varepsilon}_n\), times
the volume of the bounding box. The boundary cubes have total
volume tending to zero. Apply this to the moment bound before
taking the upper limit. This proves
\eqref{eq:row-l1-upper-bound}.
\end{proof}

We use the following elementary simultaneous-selection argument.
If \({\varepsilon}_{n,j}\) is a complex sequence eventually bounded for each
\(j\in\N\), successively select a subsequence on which the
first, second, and subsequent coordinates converge. The diagonal
choice, taking its \(n\)-th term from the \(n\)-th nested
subsequence beyond the preceding choice, is strictly increasing
and makes every fixed coordinate converge. Finite exceptional
terms can first be included in a larger closed disc. Thus the
argument requires a bounded tail for each fixed coordinate; a
doubly eventual bound only at large periods would be insufficient.

\begin{proposition}[One mesh sequence for every period and every test]
\label{prop:common-uv-distribution-rows}
There is a strictly increasing \(\nu:\N\to\N\), chosen before
any volume request, such that, for every \(M,r,\mathbf d\) and
\(\mathbf v\), there is a continuous complex-linear functional
\(S_M(\mathbf v;\cdot)\) on \(C_c^\infty(\Omega_{r,M})\) with
\begin{equation}
       S_{\nu(n),M}(\mathbf v;f)
                             \longrightarrow S_M(\mathbf v;f)
                       \quad\hbox{for every }f.
 \label{eq:common-uv-distribution-rows}
\end{equation}
It is multilinear in the label slots. On every compact set it
satisfies the bounds of Lemma~\ref{lem:uniform-row-smearings},
including the \(L^1\) test bound. Constants are uniform over all
sufficiently large suitable periods. This is a full mesh limit
after the single fixed refinement.
\end{proposition}
\begin{proof}
Choose a compact exhaustion of every open set \(\Omega_{r,M}\).
For example intersect the centered box of margin \(1/l\) with
the set of torus separations at least \(1/l\), and let \(l\)
increase. When approximating a test, use a slightly larger
exhaustion compact so that its support has an interior margin.
For each compact choose countably many smooth tests, supported
in that larger compact, dense for the uniform norm among tests
supported in the smaller compact. Such sets can be constructed
by piecewise polynomial approximations on rational cubes followed
by convolution with smooth mollifiers and multiplication by a
fixed compactly supported cutoff. Rational real and imaginary
coefficients and a countable choice of smoothing radii suffice.

Enumerate all periods, arities, degree lists, finite basis lists,
exhaustion indices and these tests. This is a countable list:
each \(V_d\) has a finite basis, and there are countably many
finite lists of integer degrees. Lemma~\ref{lem:uniform-row-smearings}
gives an eventually bounded sequence for each listed item. Apply
the simultaneous-selection argument to choose \(\nu\). It
depends on all fixed periods together, not on one later requested
volume sequence.

Convergence extends to every test as follows. Given \(f\),
choose approximating enumerated tests \(f_l\) with common compact
support and \(\|f-f_l\|_\infty\to0\). The smearing difference
is bounded by \(A_K\|f-f_l\|_\infty\), uniformly on the retained
cutoff tail. First choose \(l\) to make this error small, then
use convergence for \(f_l\). This is the Cauchy criterion for
\(S_{\nu(n),M}(f)\). It proves a limit independent of the
approximation and linear in \(f\). The compact uniform bound
makes it a distribution of order zero. The sharper \(L^1\)
bound follows by taking the limit in
\eqref{eq:row-l1-upper-bound}. Finite multilinearity extends
convergence from the basis lists to all complex label vectors.
The volume-uniform constants are those already proved in the
insertion estimate, so the extension does not introduce a
period-dependent choice of \(\nu\).
\end{proof}

We henceforth fix this ultraviolet refinement as part of the
Wilson family, using Lemma~\ref{lem:fixed-uv-refinement}. Thus
the first limit in Definition~\ref{def:ordered} is now available
as a distributional row limit, for every suitable fixed period.
Nothing has been selected in response to an infinite-volume
sequence. The stronger assertion of full convergence for every
cutoff of the unrefined family is not needed for this compactness
step, and is not asserted.

In the remaining statements \(S_{n,M}\) denotes the smearing at
the retained old cutoff \(\nu(n)\). The internal reference unit
\(u_d\) and the definition of its reference level are unchanged.

\subsection{The outer extraction and the separated hierarchy}

\begin{theorem}[An ordered weak correlation hierarchy]
\label{thm:ordered-weak-hierarchy}
For the fixed Wilson family with the preceding common mesh
refinement, every strictly increasing sequence of admissible side lengths
\(M_j\to\infty\)
has a strictly increasing refinement \(q(j)\) and a hierarchy
\(S(\mathbf v;\cdot)\) of distributions on
\(\operatorname{Conf}_r(\R^4)\) such that
\begin{align}
 \lim_{n\to\infty}S_{n,M'}(\mathbf v;f)
                                  &=S_{M'}(\mathbf v;f)
                &&\text{for every suitable fixed }M',\notag\\
 \lim_{j\to\infty}S_{M_{q(j)}}(\mathbf v;f)
                                  &=S(\mathbf v;f).
 \label{eq:ordered-weak-hierarchy}
\end{align}
The hierarchy is multilinear in its finite-dimensional polynomial
labels, symmetric under simultaneous permutation of points and
labels, respects conjugation, and has empty moment one. On each
separated compact set it has a uniform \(L^1\) test bound.
The same \(q\) works for all arities, degrees, labels and tests.
\end{theorem}
\begin{proof}
For each arity take a countable compact exhaustion of
\(\operatorname{Conf}_r(\R^4)\), countable uniformly dense
test sets as above, and all finite basis lists. Every listed test
is suitable for all sufficiently large \(M_j\). The row
functionals on this tail have a common bound by
Proposition~\ref{prop:common-uv-distribution-rows} and its
volume-uniform clause. Apply the same countable diagonal argument
to these row values, now choosing \(q\). Inactive early periods
can be filled arbitrarily; they play no role in a fixed test's
limit. Extend from dense tests and basis lists using the same
uniform error bound as in the preceding proof. This constructs
all the asserted distributions on one volume refinement.

Finite lattice moments are multilinear and symmetric, and their
complex conjugates are obtained by conjugating every label and
the test. These identities are identities of smearings and pass
through the two successive limits. The empty product is one at
every cutoff and period. Passing the compact \(L^1\) test bound
through the limits gives the asserted bound on \(S\).
The first line of \eqref{eq:ordered-weak-hierarchy} was fixed
before \((M_j)\) was requested; the second selects only side lengths.
Consequently this is an ordered construction, rather than a
single limit along pairs \((n_j,M_j)\).
\end{proof}

The compact \(L^1\) bound means, in particular, that no
off-diagonal singular measure can appear in these distributions.
Equivalently, on each relatively compact separated domain the
functional extends boundedly to \(L^1\), and its dual is an
\(L^\infty\) function, unique up to null sets. This is the usual
\(L^1\)--\(L^\infty\) duality for Lebesgue measure. No smooth
representative has yet been obtained, and values at individual
points are not defined by this argument.

\subsection{Symmetries and positivity already retained}

\begin{proposition}[Euclidean symmetries available at this stage]
\label{prop:weak-hierarchy-symmetries}
The distributional hierarchy is invariant under common continuous
translations of all points and is covariant under signed
coordinate permutations with their polynomial-label action.
Insertion of the label \(u_d\in V_d\) removes that slot,
with integration of the test over its position. These assertions
hold for the same hierarchy and the same volume refinement.
\end{proposition}
\begin{proof}
The unit assertion follows from
Corollary~\ref{cor:local-insertion-regularity} on a cutoff tail.
Integrating its position in the lattice smearing is the ordinary
Riemann sum, whose error tends to zero for a fixed smooth test;
the remaining compact moment bound controls that error. The
resulting test in the remaining variables is still smooth,
compactly supported and separated.

Signed covariance is the exact finite-lattice covariance of the
Wilson action and the constructed insertion, followed by change
of summation variables. For compact tests eventually away from
the torus seam the centered representative convention makes no
boundary contribution. Pass through the ordered limits.

For a fixed translation \(b\in\R^4\), choose lattice translations
\(b_n\in {\varepsilon}_n\Z^4\) with \(b_n\to b\). Literal translation of
all insertion polynomials and translation invariance of the
Wilson measure give the smearing identity for \(b_n\). Choose
the periods large enough to contain the supports of the test and
its translate, with a fixed margin; only finitely many unsuitable
periods are excluded. The difference between the tests translated
by \(b_n\) and by \(b\) tends uniformly to zero on one separated
compact. Lemma~\ref{lem:uniform-row-smearings} bounds the
corresponding smearing error. First take the mesh limit at every
such fixed period and then the volume limit. This proves common
continuous translation invariance.
\end{proof}

There is also a useful restricted positivity statement.
Keep time as the final coordinate, and let
\(\vartheta(\mathbf z,f)=(\mathbf z,-f)\). Write
\(\vartheta v\) be its signed-permutation action on a polynomial
label. A separated positive-time word consists of a degree and
label list \(\mathbf v_\alpha\) and a smooth test
\(f_\alpha\) compactly supported on mutually distinct points
with every time coordinate strictly positive. For two words,
write
\[
 (\vartheta\overline{f_\alpha}\otimes f_\beta)(z,w)
                   =\overline{f_\alpha(\vartheta z)}f_\beta(w),
 \qquad
 \vartheta\overline{\mathbf v_\alpha}
                  =(\vartheta\overline{v_{\alpha,i}})_i.
\]
The reflected and unreflected point sets are separated from each
other because their time coordinates have opposite signs and
compact positive distance from zero.

\begin{proposition}[Reflection positivity for separated words]
\label{prop:weak-separated-reflection-positivity}
For every finite list of separated positive-time words,
\begin{equation}
 \sum_{\alpha,\beta}
 S\bigl(\vartheta\overline{\mathbf v_\alpha},\mathbf v_\beta;
                  \vartheta\overline{f_\alpha}\otimes f_\beta\bigr)
                                                        \ge0.
 \label{eq:weak-separated-reflection-positivity}
\end{equation}
The sum is a real nonnegative number.
\end{proposition}
\begin{proof}
At finite cutoff smear each word over its positive-time lattice
anchors and add the resulting finitely many observables. For
sufficiently fine mesh the entire stencil of each insertion is
strictly on the positive side of a bond plane at time
\(-{\varepsilon}_n/2\), and for large suitable periods it also avoids the
opposite periodic plane. Denote this lattice observable by
\(F_{n,M}\). Lemma~\ref{lem:insertion-chessboard} gives
\(\langle\theta_nF_{n,M}\,F_{n,M}\rangle_{n,M}\ge0\), where
\(\theta_n\) includes conjugation and reflection about that
bond plane.

Expanding this nonnegative expression gives the finite sum of
double smearings in the proposition, with reflected time
\(-f-{\varepsilon}_n\) in place of \(-f\). Exact signed covariance of
the insertion gives the displayed reflected labels. The shifted
tests tend uniformly, with common compact separated supports, to
the tests in \eqref{eq:weak-separated-reflection-positivity}.
The uniform compact moment bound controls the difference. Take
the full mesh limit at each suitable fixed period and then the
selected volume limit. A finite sum of these convergent numbers
retains its reality and nonnegativity.
\end{proof}

This proves positivity on separated test words for the
ordered hierarchy just obtained. The regularity and
growth arguments needed for reconstruction are developed
in Sections~\ref{sec:smooth-weak-hierarchy}--%
\ref{sec:wightman-reconstruction}. Full orthogonal
covariance is supplied by
Theorem~\ref{thm:full-orthogonal-covariance}; it is not
inferred from signed permutations alone.
The complete construction retains the local insertion
family, its cutoff- and volume-uniform bounds and this
one ultraviolet choice before every volume request,
as made explicit in
Section~\ref{sec:construction-assembly}.
% !TeX root = 4DYM_notes.tex
\section{Position differences and local matching}
\label{sec:position-local-matching}

The moment bound of Theorem~\ref{thm:uniform-separated-insertions}
does not yet control position derivatives. A lattice translation
by one site has physical length \({\varepsilon}_n\), so its difference
quotient contains \({\varepsilon}_n^{-1}\). The purpose of this section is
to represent that quotient in a higher-degree label space. A
later finite-dimensional estimate will bound its physical
coefficient; that bound is not assumed in local matching.

All estimates are first proved on the original finite histories,
with terminal RG level \(n\). They remain valid on the mesh
refinement fixed in Section~\ref{sec:weak-hierarchy-extraction}.
Internal reference levels and the units \(u_d\) are unchanged.
For position differences at terminal cutoff \(n\), we write
\(\varepsilon=\varepsilon_n\).
We use the local comparison method of
\cite[Lemma~14.16, PDF pp.~139--141, and
Lemmas~14.20--14.21, PDF pp.~144--145]{ArmstrongDraft},
giving the norms and the backward matching argument explicitly.

\subsection{Differences, admissible coefficients and small residuals}

For a local observable family \(F_{n,M}(y)\) and an axis
\(a\in\{1,2,3,4\}\), write
\[
 \tau_aF(y)=F(y+e_a),\qquad
          \partial_a^{\varepsilon}F(y)=\frac{F(y+e_a)-F(y)}{\varepsilon}.
\]
Translations here are of the insertion anchor. For a finite axis
list \(\boldsymbol a=(a_1,\ldots,a_h)\), put
\(\partial_{\boldsymbol a}^{\varepsilon}=\partial_{a_1}^{\varepsilon}\cdots\partial_{a_h}^{\varepsilon}\).
For a moment, a slot is specified as well as an axis; differences
in different slots commute and act only on their respective
insertions.

A nonnegative family \(A_{n,M}\) has \emph{uniform subpower growth}
if, for every \(\epsilon>0\),
\[
      \sup_M A_{n,M}\le {\varepsilon}_n^{-\epsilon}
\]
on a full cutoff tail. Its threshold may depend on \(\epsilon\)
and the fixed request, but not on \(M\). At degree \(e\), call
a coefficient family \(c_{n,M}\in V_e\) \emph{admissible} if
\(\left|\mathsf Z_{n,M,e}^{-1}c_{n,M}\right|_1\) has this growth.
In particular every bounded reference-coordinate family is
admissible. Indeed Theorem~\ref{thm:one-insertion-step} gives
\begin{equation}
 \|\mathsf Z_{n,M,e}^{-1}\|_{1\to1}
       \le\prod_{i=\operatorname{ref}(e)}^{n-1}C_{J(i+1)}={\varepsilon}_n^{-o(1)},
 \label{eq:position-inverse-subpower}
\end{equation}
uniformly in \(M\), by the logarithmic calendar choice.

For a translation-covariant fixed-stencil family \(E_{n,M}\),
say it is \emph{root-small} if there is one
\(\epsilon_n\downarrow0\) such that, at every fixed centred
tiling level \(m\),
\begin{equation}
      \sup_M\max_{\rho,\,\text{conjugation}}
                       \|E_{n,M}^{\rho}(0)\|_m\le\epsilon_n
 \label{eq:position-root-small}
\end{equation}
on a full cutoff tail. The finite maximum includes signed images
and conjugation. The cutoff threshold may depend on \(m\).
The physical stencil is eventually interior in every such cell.
By Lemma~\ref{lem:insertion-interior-root}, this also bounds
uncentred interior roots after choosing a sufficiently fine fixed
cell inside the required margin.

\begin{lemma}[Testing a root-small residual]
\label{lem:position-small-spectators}
A root-small family has vanishing moments against every fixed
finite list of bounded-coefficient ordinary insertions at positive
physical separation, uniformly in cutoff tails and volume.
The same is true after smearing on a fixed separated compact.
\end{lemma}
\begin{proof}
Choose a fixed bond tiling with cell side smaller than the
separation and translate its boundaries away from all anchors,
as in Theorem~\ref{thm:uniform-separated-insertions}. The interior
root of the residual is at most \(\epsilon_n\), and the finitely
many spectator roots have uniform bounds. The chessboard
inequality bounds the moment by their product. The choice of
cells and the cutoff tail are independent of volume. For smearing,
multiply by the bounded test and sum its cell-volume weights;
the total volume of its support box is bounded independently of
the mesh. This proves both assertions.
\end{proof}

We shall not differentiate a root-small residual. An additional
difference quotient could multiply its bound by \({\varepsilon}_n^{-1}\).
Instead, every fixed iterated difference is matched anew in its
final degree.

\subsection{Starting an arbitrary local source}

Fix \(e\) and a family \(F_{n,M}(y)\), continuous and gauge
invariant on one fixed stencil. Suppose that its dimensionless
representative \(f_{n,M}={\varepsilon}_n^eF_{n,M}(0)\) has both a healthy
derivative norm and an all-configuration supremum bound of uniform
subpower growth. The derivative norm is the original norm
\eqref{eq:local-derivative-norm}, at the current profile order and
threshold. These are two different hypotheses: the healthy norm
pays Taylor reading; the global bound pays the unhealthy part.

Enlarge the stencil by \(\Sigma_e\). Read its degree-\(e\)
radial jet, subtract the corresponding local section on the
healthy part, and restore the unhealthy part by the exact forced
attachment of Proposition~\ref{prop:marked-insertion-birth}.
This produces an exact one-insertion representation with coordinate
\(v_n\), remainder norm \(r_n\), and
\begin{equation}
                         |v_n|_1+r_n\le A_{n,M},
 \label{eq:position-local-start}
\end{equation}
where \(A\) has uniform subpower growth. No graph inequality
\(r_n\le C|v_n^\circ|_1\) is required here.

To check \eqref{eq:position-local-start}, the Taylor reading is a
bounded linear map on the finite jet space. The section subtraction
has its finite derivative price. At sufficiently large \(B_n\)
the fixed enlarged stencil has covering number at most a fixed
integer. Its support factor and finite derivative constants can
be included in the preassigned profile constant \(\mathcal C_j\).
For the forced complement, its global amplitude is multiplied by
the original forced factors meeting that stencil. The weighted
sum of their nonempty products is at most
\(e^{|\Sigma|w_n}-1\). Dividing by \(\sqrt{w_n}\) leaves a
factor \(O(\sqrt{w_n})\), which is at most one. The global
amplitude itself need only have subpower growth; it is not divided
by \(w_n\) without those factors. This proves the assertion.

We include in each \(\mathcal C_j\) the finite section constants,
derivative factorials, and changes of anchor for the finitely many
degrees and stencil shifts requested at that profile. Thus these
constants remain subpower along the calendar. This is compatible
with the previous choices: they are fixed before its ratio, whose
logarithmic floor can then be enlarged.

\subsection{The backward cone calculation}

The local marked estimates are linear in a distinguished input
and apply to arbitrary finite remainder size. Only the unmarked
attachments have to be small. After finite symmetry recombination,
enlarge the ratio floors so that those estimates may be used with
\(\varepsilon_L=L^{-3/8}\). For one step write
\(u,v\) for its coarse and fine coordinates and \(R,T\) for
the corresponding remainder norms. Put
\(Y=|\mathsf S_Lv|_1\). Lemma~\ref{lem:insertion-local-prices}
and the leading inverse bound give
\begin{equation}
 Y\le K\{|u|_1+\varepsilon_L(Y+T)\},
                 \qquad R\le KY+\varepsilon_LT.
 \label{eq:position-match-cone}
\end{equation}
Here \(K\) is a fixed adjacent-profile constant. The same
inequalities hold for a difference of two representations with
the same unmarked history, by linearity and the marked difference
estimate.

Choose \(A_i=32(1+C_{J(i)})^2\). If
\(A_i|u|_1\le R\) and \(A_i\ge4K^2\), then
\[
                  4K|u|_1\le Y+\varepsilon_LT.
\]
Substitute this into the first inequality of
\eqref{eq:position-match-cone}. When \(4K\varepsilon_L\le1\),
\[
 Y\le\tfrac14Y+\tfrac14\varepsilon_LT
                     +K\varepsilon_LY+K\varepsilon_LT
       \le\tfrac12Y+(K+\tfrac14)\varepsilon_LT,
\]
and hence \(Y\le4K\varepsilon_LT\). Since
\(|v|_1\le Y\), the further ratio requirements
\begin{equation}
 4KA_{i+1}\varepsilon_L\le1,
       \qquad (4K^2+1)L^{1/8}\varepsilon_L\le1
 \label{eq:position-cone-reserves}
\end{equation}
give
\begin{equation}
             A_{i+1}|v|_1\le T,
                    \qquad L^{1/8}R\le T.
 \label{eq:position-cone-propagation}
\end{equation}
All constants in these requirements precede the ratio. The two
powers \(L^{-3/8}\) and \(L^{-1/4}\) tend to zero, so the
requirements are compatible at a repeated profile and at a
profile transition.

\begin{proposition}[Local matching without a physical coefficient bound]
\label{prop:position-local-matching}
For the local family above there are admissible coefficients
\(b_{n,M}\in V_e\) such that
\begin{equation}
          \mathcal O_{n,M}((e,b_{n,M}),y)-F_{n,M}(y)
                                      \quad\text{is root-small}.
 \label{eq:position-local-comparison}
\end{equation}
The choices are made on the same vacuum history and use literal
translates of the origin's coefficients. This assertion does not
bound \(|b_{n,M}|_1\) in the reference coordinate.
\end{proposition}
\begin{proof}
Transport the exact local source from
\eqref{eq:position-local-start} down to \(\operatorname{ref}(e)\), using the
linear local marked construction before graph inversion. Define
\(b_{n,M}\) to be its coordinate there. Subtract the selected
primary inserted density with reference coordinate \(b_{n,M}\).
The residual has coordinate zero at \(\operatorname{ref}(e)\), so its remainder
starts the cone \(A_{\operatorname{ref}(e)}|u|_1\le R\). Applying
\eqref{eq:position-cone-propagation} towards the fine scale gives
\[
 R_k\le R_n\prod_{i=k}^{n-1}L_i^{-1/8}
                   =R_n({\varepsilon}_n/{\varepsilon}_k)^{1/8},\qquad A_n|v_n^{\rm res}|_1\le R_n.
\]
At the top the selected graph has remainder at most
\(C_{J(n)}|\mathsf Z^{-1}b_{n,M}|_1\). Its difference from the
original top coordinate is \(v_n^{\rm res}\), so
\[
 R_n\le(1+C_{J(n)})A_{n,M}
                       +C_{J(n)}|v_n^{\rm res}|_1.
\]
The last cone inequality absorbs the second term. Thus, with a
uniformly subpower envelope \(E_{n,M}\),
\begin{align}
 R_n&\le E_{n,M},\qquad
 R_k\le E_{n,M}({\varepsilon}_n/{\varepsilon}_k)^{1/8},\notag\\
 |\mathsf Z^{-1}b_{n,M}|_1
                 &\le |v_n|_1+|v_n^{\rm res}|_1\le2E_{n,M}.
 \label{eq:position-backward-matching}
\end{align}
The coefficient is therefore admissible.

We check that the small coarse remainder also gives a small
physical root, retaining the full several-source density.
For this physical comparison retain the actual unhealthy
response family \(L_n\) from the birth construction above.
Its reserved price \(e(L_n)\) has uniform subpower growth:
the stencil is fixed, and its global function and section
amplitudes have that growth. Enlarge \(E_{n,M}\), if
necessary, to include \(\widehat\beta_n e(L_n)\).
The enlarged envelope still has uniform subpower growth.
This is a separate response allowance, not a consequence
of the divided singleton norm. The
cone bounds its total physical singleton size at every level
\(i\ge \operatorname{ref}(e)\) by
\[
                  B_i^{\rm env}
                   =2E_{n,M}{\varepsilon}_n^{1/8}{\varepsilon}_i^{-e-1/8}.
\]
This includes its coordinate, regular remainder and square-root
normalised forced remainder. Fix a centred tiling level \(m\)
and a finer stopping level \(k\ge \operatorname{ref}(e),m\) with
\(B_k^2\le {\varepsilon}_m/{\varepsilon}_k\). At \(k\), take the estimating radius
\[
       \rho_k=\widehat\beta_n^{-1}
                            \frac{\zeta_{e,k}}{4B_k^{\rm env}}.
\]
At finer levels use \(\rho_{i+1}=\widehat\beta_{i+1}^{-1}\rho_i\).
The entry requirement
\[
 \widehat\beta_{i+1}^{-1}L_i^{e+1}\zeta_{e,i}
                                                    \le\zeta_{e,i+1}
\]
keeps \(\rho_iB_i^{\rm env}\le\zeta_{e,i}/4\), since the
envelope grows by only \(L_i^{e+1/8}\). It is the same numerical
radius transport as \eqref{eq:joint-radius-entry}, with degree
increased by one. The attachment proof of
Lemma~\ref{lem:joint-source-attachments} uses the complete
singleton size. Its initial physical joining is justified
by the additional response allowance, using the finite
regrouping of Lemma~\ref{lem:complete-source-response-joining}
and Corollary~\ref{cor:physical-response-general-envelope},
with birth level \(n\). This argument is expanded there
also for a source born at an intermediate level. The
radius reduction costs only the subpower factor
\(\widehat\beta_n\). Higher modes and their original
inventories are retained. Their radius changes are paid
by the separated-name covering estimate as before.

The replicated coefficient bound and the scalar-free partition
ratio at this fixed stopping level therefore give
\begin{equation}
 \sup_M\|\mathcal O^{(e)}(b)-F\|_m
                  \le C_{m,k,e}\widehat\beta_n
                                      E_{n,M}{\varepsilon}_n^{1/8}.
 \label{eq:position-matching-root-price}
\end{equation}
Its constant can depend on \(m,k,e\), but not on cutoff or
period. The large estimating radius at \(k\) has not been used
at the fine scale without shrinking it: the preceding transport
is what verifies the fine-scale attachment hypotheses.
The product \(\widehat\beta_n E_{n,M}\) still has
uniform subpower growth. It absorbs, for example,
\({\varepsilon}_n^{-1/32}\); the
fixed constant is absorbed by another such factor on each
requested tail. We may consequently use the single tolerance
\(\epsilon_n={\varepsilon}_n^{1/32}\) in
\eqref{eq:position-root-small}. Include the finitely many signed
and conjugated inputs in the same construction. This proves the
proposition.
\end{proof}

\subsection{Corrected inclusion and corrected differences}

\begin{corollary}[Local comparison of degrees and differences]
\label{cor:position-corrected-differences}
For \(d\le e\) and \(v\in V_d\), there is an admissible
\(c_{n,M}\in V_e\) with
\(\mathcal O^{(e)}(c)-\mathcal O^{(d)}(v)\) root-small.
For each fixed axis list \(\boldsymbol a\) of length \(h\),
there is an admissible \(c^{\boldsymbol a}_{n,M}\in V_{d+h}\)
with
\begin{equation}
 \mathcal O^{(d+h)}(c^{\boldsymbol a}_{n,M},y)
                  -\partial_{\boldsymbol a}^{\varepsilon}\mathcal O^{(d)}(v,y)
                                      \quad\text{root-small}.
 \label{eq:position-corrected-difference}
\end{equation}
Choices on each finite-dimensional input space can be made linear
and real, with common fixed stencils. Inclusions can also be
chosen signed-equivariant. All signed images of a difference,
including the fine shift caused by reflection, obey the same
comparison estimates.
\end{corollary}
\begin{proof}
In the first case the dimensionless target is
\[
 {\varepsilon}_n^{e-d}P_d[\mathsf D_{h_n^{-1}}
                                \mathsf Z_{n,M,d}^{-1}v].
\]
In the second, multiplication by \({\varepsilon}_n^{d+h}\) cancels every
inverse spacing and leaves a finite alternating sum of translates
of that same section. Its stencil is enlarged by at most the
fixed list of unit shifts. Equation
\eqref{eq:position-inverse-subpower} and the section bounds give
uniform subpower derivative and supremum bounds. Moving an
anchor by a unit multiplies the spatial weight by at most
\(2^{s_{J(n)}}\) per derivative. These finite profile costs
were included in \(\mathcal C_j\) before its ratio and hence
are subpower too. Proposition~\ref{prop:position-local-matching}
applies in the stated final degree. Basis recombination and
conjugation averaging give linear real choices. Signed averaging
is available for inclusion. For differences use the translated
reflected stencils in the finite list: reflection changes a
forward difference to the negative of a backward difference,
not exactly to the negative forward difference. No previously
small error has been differentiated.
\end{proof}

\subsection{An auxiliary section with a bounded bare comparison}

We will need the same matching for another fixed Taylor-normalised
section \(P_e^{\rm aux}\). Apply it simultaneously to the columns
\({\varepsilon}_n^{-e}P_e^{\rm aux}[\mathsf D_{h_n^{-1}}x]\), and write the
matched coefficient as \(\mathsf Z_{n,M,e}\mathsf P_{n,M,e}x\).
The exact Taylor read at the top is \(x\). Its healthy derivative
price \(H_{j,e}\), including the finite symmetry recombination
and the factor \(4K\) in the last-step cone bound,
is fixed before the ratio; the bad complement is multiplied by
the original forced gas. The last step of the cone calculation
therefore gives
\[
 \|I-\mathsf P_{n,M,e}\|_{1\to1}
       \le2(1+C_j)H_{j,e}L_{n-1}^{-1/8},\qquad j=J(n).
\]
Include the ratio floor making this at most \(1/4\).
Then \(\mathsf P\) is invertible with
\(\|\mathsf P^{-1}\|_{1\to1}\le2\), by its convergent
Neumann series. For the unit use the exact representation, so
\(\mathsf P1=1\). This matrix is chosen before later cell and
spectator requests; its construction does not depend on them.
Its comparison residual obeys
\eqref{eq:position-matching-root-price} for every admissible
column family. This will transfer exact lattice equation
identities to the selected insertion coordinates.
% !TeX root = 4DYM_notes.tex
\section{Exact equation observables and finite Maxwell tests}
\label{sec:equation-sections-maxwell}

Local matching supplies a higher-degree coefficient, but does not
bound it. Some coefficient directions instead correspond to exact
equations of the lattice measure. We construct those directions
before using a free-field comparison. In particular an equation
term will be removed by Haar integration, not estimated as an
arbitrary large insertion.

The local equation construction follows
\cite[Lemma~10.4, PDF pp.~44--45]{ArmstrongDraft}.
The finite free tests below use its free positivity calculation
in Sect.~G. All statements are at finite cutoff until the later
passage to the hierarchy.

\subsection{The free equation subspace}

Inside the radial-potential space \(\mathcal J_j\) of
Section~\ref{sec:local-observable-sections}, put
\[
 \mathcal J_j^{\rm Max}=
 \left\{A_j:\ \sum_a\partial_a
       (\partial_a A_{j,b}-\partial_b A_{j,a})=0
                                        \text{ for every }b\right\}.
\]
These are the linear Maxwell equations. Define
\begin{equation}
 \mathcal I_d=
 \{p\in V_d:p|_{\bigoplus_{j+2\le d}\mathcal J_j^{\rm Max}}=0\},
                         \qquad V_d^{\rm Max}=V_d/\mathcal I_d.
 \label{eq:position-maxwell-quotient}
\end{equation}
An element of \(V_d^{\rm Max}\) is an invariant polynomial of
the on-shell jets, meaning jets satisfying those equations.
The radial-potential description already incorporates the linear
Bianchi identity for their curvatures. The equation subspace
\(\mathcal I_d\) is the weight-\(d\) part of the polynomial
ideal generated by the linear Maxwell equations.

To verify the last statement, choose projections
\(\pi_j:\mathcal J_j\to\mathcal J_j^{\rm Max}\), preserving
the real structure and the adjoint action, and put \(\pi=\oplus_j\pi_j\).
For a monomial in linear coordinates,
\begin{equation}
 \prod_{i=1}^k\ell_i-\prod_{i=1}^k(\ell_i\circ\pi)
  =\sum_{i=1}^k
      \left(\prod_{l<i}\ell_l\circ\pi\right)
        (\ell_i-\ell_i\circ\pi)\prod_{l>i}\ell_l.
 \label{eq:position-equation-ideal}
\end{equation}
Each distinguished linear factor annihilates one Maxwell
subspace. If \(p\in\mathcal I_d\), then \(p\circ\pi=0\),
so this identity proves the ideal assertion. Projections may be
chosen equivariantly by invariant inner products. Restriction
onto invariant polynomials of the Maxwell jets is surjective:
compose such a polynomial with \(\pi\). Consequently
\begin{equation}
                 \dim V_d-\dim\mathcal I_d=\dim V_d^{\rm Max}.
 \label{eq:position-maxwell-dimension}
\end{equation}
This dimension calculation does not use an interacting limit.

\subsection{One-link integration identities}

Write \(C_{ab}(U;x)\) for the projected plaquette curvature based
at \(x\), oriented first along \(a\) and then along \(b\).
Thus \(C_{ba}=-C_{ab}\) and \(C_{aa}=0\). For a positive link
\(e=(x,a)\), set
\begin{equation}
 \mathcal E_e(U)=\sum_{b\ne a}
   \left\{C_{ab}(U;x)
       -\operatorname{Ad}(U(x-e_b,b)^{-1})C_{ab}(U;x-e_b)\right\}.
 \label{eq:position-lattice-force}
\end{equation}
For the Wilson energy \(W_\ell\), direct differentiation
in the left link direction \(A\in\mathfrak g\) gives
\begin{equation}
                  D^L_{e,A}W_\ell(U)
                                  =\langle A,\mathcal E_e(U)\rangle.
 \label{eq:position-action-derivative}
\end{equation}
For the forward plaquette its derivative is
\(-N^{-1}\operatorname{Re}\operatorname{Tr}(AP_q)\), which
equals \(\langle A,C(P_q)\rangle\) in our metric. In the
backward plaquette differentiate the inverse occurrence of the
link and cycle its trace. This gives the negative of the
transported term in \eqref{eq:position-lattice-force}. Summing
the two incidences for each transverse axis proves the identity.

If a continuous multiplier \(g\) and a bounded Borel spectator
\(S\) are both independent of \(U_e\), left invariance of
probability Haar measure gives, for \(\beta\ne0\),
\begin{equation}
        \int \langle A,\mathcal E_e\rangle gS\,
                                     \dd\mu_{\ell,\beta}=0.
 \label{eq:position-exact-haar-equation}
\end{equation}
Indeed the integral of \(gS e^{-\beta W_\ell}\) is unchanged
when that link is left multiplied by \(e^{tA}\). Its derivative
is the preceding integral times \(-\beta Z\). Compactness bounds
the differentiated smooth density uniformly, so differentiation
is legitimate for the bounded measurable spectator.

Choose a path \(\gamma_e\) from the origin to \(e^-\) avoiding
the physical link \(e\). If \(g\) also avoids it, then
\begin{equation}
        \langle\xi,\operatorname{Ad}(U_{\gamma_e})
                                      \mathcal E_e(U)\rangle g(U)
 \label{eq:position-transported-equation}
\end{equation}
has the same zero exterior moment. Expand
\(\operatorname{Ad}(U_{\gamma_e}^{-1})\xi\) in a fixed real
colour basis and apply \eqref{eq:position-exact-haar-equation}
to each component. Its multiplier still avoids \(e\).

\subsection{Equation sections with fixed support}

Order the finitely many bidegrees of \(V_d\) lexicographically
by factor count first and engineering weight second. A linear
map which is the identity plus terms strictly later in this
order is invertible: its nonidentity part is nilpotent. We use
such a finite triangular correction below.

\begin{proposition}[Local sections with exact exterior equations]
\label{prop:position-equation-sections}
There are a Taylor-normalised polynomial section \(P_d^{\rm eq}\)
on a fixed box, an invertible real signed-equivariant map
\(\mathsf R_d\) fixing the unit, and fixed integers \(R_d,{\mathsf M}_d\),
such that
\begin{equation}
 \int P_d^{\rm eq}[\mathsf R_dp](y;U)S(U)\,
                     \dd\mu_{\ell,\beta}(U)=0,
       \qquad p\in\mathcal I_d,\quad \ell\ge {\mathsf M}_d,\quad\beta\ne0,
 \label{eq:position-exterior-equation}
\end{equation}
whenever \(S\) is bounded Borel and reads only links with both
endpoints at distance greater than \(R_d\) from \(y\).
The section has the same threshold-scaled derivative and global
bounds as Lemma~\ref{lem:insertion-section-bounds}, with its own
constants depending only on \(G,d,q,s\) and its fixed choices.
\end{proposition}
\begin{proof}
First, curvature measurements can be chosen to avoid any one
specified link. Replace each use of that link in a transporting
path by the other three sides of a transverse square. Only
finitely many translated plaquettes contain the forbidden link.
Removing them does not affect the highest-homogeneous-component
argument of Lemma~\ref{lem:insertion-finite-measurements}: the
shrinking-square lattice points avoid any fixed finite set for
nonzero limiting position, and continuity handles the origin.
Thus finitely many avoiding measurements still determine every
linear jet coordinate.

Second, linearised forces determine every linear equation
coordinate in its highest requested weight. On test potentials,
the linearisation of \eqref{eq:position-lattice-force} is
\[
 \sum_b\{{\mathsf F}_{j;ab}(x,1)-{\mathsf F}_{j;ab}(x-e_b,1)\},
\]
summed over the homogeneous components. Here \({\mathsf F}_{j;ab}\) is
the square average of the curl used in
Lemma~\ref{lem:insertion-finite-measurements}. The fundamental
theorem of calculus writes its difference as the average of
\(\partial_b(\partial_aA_{j,b}-\partial_bA_{j,a})\) over one
extra unit segment. If all such force measurements vanish on
jets through component \(j_*\), rescale a translated unit cube
by \(t\), multiply by the appropriate highest-degree power of
\(t\), and let its lattice origin approach an arbitrary point
as in that lemma. The resulting identity is
\[
              \sum_b\partial_b
                   (\partial_aA_{j_*,b}-\partial_bA_{j_*,a})=0.
\]
Hence every functional of the top component annihilating
\(\mathcal J_{j_*}^{\rm Max}\) vanishes on the common kernel
of these force measurements. Finite-dimensional duality places
it in their finite linear span. On the full larger jet space,
the same combination can have additional linear terms only in
higher engineering weight. This argument does not require
discrete and continuum harmonic polynomials to coincide in
lower weights.

Now decompose each equation monomial by
\eqref{eq:position-equation-ideal}. Lift its distinguished
linear equation factor by those finite force combinations.
For each varied link lift all the other linear factors by
curvature measurements avoiding that link, with their full
linear jet prescribed through weight \(d\). Choose its
transporting path to avoid the link as well. Their product
is a finite sum of expressions
\eqref{eq:position-transported-equation}, so has zero exterior
moment. Its Taylor polynomial is the original monomial plus
strictly later bidegrees: extra linear force terms have higher
weight, and every nonlinear correction adds a jet factor.
Every lifted monomial retains at least its original factor count.

Average over constant gauge transformations, signed permutations
and conjugation. The colour orbit lies in a finite-dimensional
space of polynomials in the same measurements, so its Haar
average is still one finite polynomial. These operations preserve
the equation ideal, the triangular Taylor identity and the
zero-moment identity. A reflected varied link may use a right
derivative; right invariance of Haar measure gives the identical
argument, without introducing the varied link into a multiplier.
Lift complementary basis vectors by
Proposition~\ref{prop:insertion-local-sections} and the unit by
one. This gives a real signed-equivariant linear lift
\(p\mapsto Y_p\).

Its Taylor reading is \(\mathsf R_dp\), with
\(\mathsf R_d-I\) strictly later in bidegree. If its nilpotence
order is \(m_d\), then
\[
      \mathsf R_d^{-1}
            =\sum_{l=0}^{m_d-1}(-1)^l(\mathsf R_d-I)^l,
             \qquad P_d^{\rm eq}[x]=Y_{\mathsf R_d^{-1}x}.
\]
This proves Taylor normalisation and
\eqref{eq:position-exterior-equation}. The inverse preserves
minimum factor count, so the section derivative proof applies
unchanged to the new finite curvature polynomials. Enclose all
the paths and varied links in one box. Choose \(R_d\) to contain
it and \({\mathsf M}_d\) so that distinct cover links remain distinct
after placement. Then all multipliers and all exterior spectators
are independent of their varied links, at every period and
placement in the statement.
\end{proof}

\subsection{The equation projection in the selected coordinates}

Use the bounded bare comparison matrix at the end of
Section~\ref{sec:position-local-matching} for \(P_d^{\rm eq}\),
and define
\begin{equation}
 B_{n,M,d}=\mathsf Z_{n,M,d}\mathsf P_{n,M,d}
                                \mathsf D_{h_n}\mathsf R_d.
 \label{eq:position-equation-coordinate-map}
\end{equation}
This is invertible. Its image of \(\mathcal I_d\) is the
finite-cutoff equation subspace. Choose a real signed-invariant
Hermitian inner product on the fixed complex coefficient space,
with real Gram matrix in its real basis, and let
\(\Pi_{n,M,d}\) be the orthogonal projection onto its orthogonal
complement. Norm equivalence in this fixed finite-dimensional
space gives
\begin{equation}
 \Pi^2=\Pi,\qquad \|\Pi\|_{1\to1}\le C_d,\qquad
 \ker\Pi=B_{n,M,d}\mathcal I_d,\qquad
                       \operatorname{rank}\Pi=\dim V_d^{\rm Max}.
 \label{eq:position-equation-projection}
\end{equation}
The norm bound does not assume that \(B\) is well conditioned.
The subspace is real and signed invariant, so the projection
commutes with those actions.

\begin{corollary}[Admissible equation directions]
\label{cor:position-equation-observables}
Let \(c_{n,M}\in\ker\Pi_{n,M,d}\) be an admissible family.
There is a fixed-stencil equation family
\(Q_{n,M}\), bounded at each finite cutoff and period, such that
\(\mathcal O^{(d)}(c)-Q\) is root-small. Its moment against
every fixed exterior spectator product is exactly zero on a
full cutoff tail. Its centred root is exactly zero at every
fixed cell size on such a tail.
\end{corollary}
\begin{proof}
Write \(c=B\pi\) with \(\pi\in\mathcal I_d\), and take
\(Q={\varepsilon}_n^{-d}P_d^{\rm eq}[\mathsf R_d\pi]\).
Admissibility and \(\|\mathsf P^{-1}\|\le2\) imply subpower
growth of \(\pi\): the additional inverse \(\mathsf D_{h_n}\)
costs only a fixed power of \(h_n^{-1}\), which is subpower.
The auxiliary-section comparison proves the root-small assertion.
The exact moment is \eqref{eq:position-exterior-equation}.
At positive physical separation every fixed spectator stencil is
exterior for all sufficiently large \(n\), uniformly in \(M\).
For the root take all other reflected copies of \(Q\) as the
spectator. A fixed physical cell eventually has lattice side
larger than twice the equation radius; those copies are exterior.
They are bounded at each finite cutoff, which suffices for the
exact Haar identity. No cutoff-uniform supremum bound for \(Q\)
is needed.
\end{proof}

\subsection{Free curvature correlations and their nondegeneracy}

The following free calculation is a test calculation, not an
identification of the interacting hierarchy with a Gaussian field.
Choose an orthonormal real colour basis. Different colours have
zero cross covariance. For the same colour let
\(\Gamma(z)=(4\pi^2|z|^2)^{-1}\) and define
\begin{equation}
 C_{ab,ij}(z)=
 -\delta_{bj}\partial_a\partial_i\Gamma
 +\delta_{bi}\partial_a\partial_j\Gamma
 +\delta_{aj}\partial_b\partial_i\Gamma
 -\delta_{ai}\partial_b\partial_j\Gamma,
                                  \qquad z\ne0.
 \label{eq:position-free-curvature-covariance}
\end{equation}
Curvature jets use its differentiated contractions. Define
\(S^{\rm Max}\) of polynomial labels by Wick pairings between
different insertion sites only; this is normal ordering at each
site. Every such pairing is a finite product of derivatives of
\(\Gamma\). These functions are smooth off coincidences.
Their equations follow by commuting derivatives and using
\(\Delta\Gamma=0\) there, so they descend to
\(V_d^{\rm Max}\).

For clarity, the radial potentials can equivalently be replaced
by homogeneous polynomial two-forms \(f_j=dA_j\) with
\(df_j=0\), \(\sum_a\partial_af_{j,ab}=0\). The inverse is
\(A_{j,b}(z)=(j+2)^{-1}\sum_a z_af_{j,ab}(z)\), by the radial
condition and Euler's identity. Give these forms the coefficient
inner product summing ordered form and derivative indices, with
Taylor coefficients \(D^j/j!\).

\begin{lemma}[Strict reflected free positivity]
\label{lem:position-free-strict-positivity}
Every nonzero on-shell invariant polynomial label has a strictly
positive reflected two-point Wick value at two distinct points
reflected across a coordinate plane.
\end{lemma}
\begin{proof}
We give the finite jet Gram calculation used in
\cite[Lemma~16.1, PDF pp.~169--171]{ArmstrongDraft}.
Let \(p_j\) be a homogeneous closed and coclosed polynomial
two-form of degree \(j\), and let \(R_{p_j}(y)\) be the
curvature covariance generated by its normalised order-\(j\)
jet probe at zero. Put
\[
 I(y)=y/|y|^2,\qquad
       R_y=I_4-2yy^{\mathsf T}/|y|^2.
\]
Differentiating \eqref{eq:position-free-curvature-covariance}
gives the Kelvin identity
\begin{equation}
 R_{p_j}(y)=s_j|y|^{-4}(R_y\otimes R_y)p_j(I(y)),
             \qquad s_j=\frac{(j+1)(j+2)(-2)^j}{\pi^2}.
 \label{eq:position-free-kelvin}
\end{equation}
Here harmonicity follows from \(df_j=\delta f_j=0\). In the
differentiation, contractions of derivative indices vanish by
harmonicity, and the terms containing divergence or a cyclic curl
vanish by the two equations. Euler's identity leaves
\(s_0=2/\pi^2\) and
\((j+1)s_{j+1}=-2(j+3)s_j\), proving the displayed coefficient
by induction.

For separation \(z\ne0\), replace an order-\(j\) probe by the
tower whose order-\(k\) entry is
\( (s_j/s_k)T_k[p_j(\cdot-I(z))]\), \(k\le j\), where
\(T_k\) is its homogeneous Taylor term at zero. This is an
invertible triangular change \(\sigma_z\) of each finite tower.
Translations preserve the closed and coclosed equations.
The generated covariance is the right side of
\eqref{eq:position-free-kelvin} with
\(p_j(I(y)-I(z))\). Its Taylor orders below \(j\) at \(z\)
vanish, and \(DI(z)=|z|^{-2}R_z\) determines its top order.
Symmetry of the reflected pairing removes the opposite
off-diagonal entries as well. Thus its quadratic form is
\begin{equation}
 B_z(\sigma_zp,\sigma_zq)
  =\sum_j\frac{(j+1)(j+2)2^j}{\pi^2|z|^{2j+4}}
                                               \langle p_j,q_j\rangle.
 \label{eq:position-free-positive-pivots}
\end{equation}
The first argument is conjugated on complexifying the real
coefficient pairing, so this is a Hermitian form.
The reflected probe acts on every form and derivative index.
Its positional sign \((-1)^j\) cancels that of \(s_j\).
Every pivot is positive. Taking \(z\) normal to a coordinate
plane makes this the required reflected jet pairing.

Diagonalise this finite positive form, with pivots \(b_l>0\).
For normally ordered polynomial coordinates the cross-site Wick
pairing is
\[
       \langle X^\mu,X^\nu\rangle
              =\mathbf1_{\mu=\nu}\mu!\prod_l b_l^{\mu_l}.
\]
It is strictly positive on every nonzero polynomial. Each label
uses only a finite jet tower; restriction to its invariant
subspace preserves this property. This proves the lemma.
\end{proof}

\begin{corollary}[A finite family of free separating tests]
\label{cor:position-finite-free-tests}
For every real basis \(v_1^{\rm Max}=1,\ldots,v_s^{\rm Max}\)
of \(V_e^{\rm Max}\), there are finitely many one- or two-point
rows, with fixed basis spectators, whose matrix of values on
these basis elements has rank \(s\). The rows may be chosen
with every coordinate separation nonzero.
\end{corollary}
\begin{proof}
For any nonzero label, strict reflected positivity supplies a
nonzero functional with the reflected conjugate label as
spectator. Expand that spectator in the finite basis; one basis
functional remains nonzero. The free correlations are smooth at
these distinct points, so perturb the points to make all their
coordinate separations nonzero without destroying that value.
The common kernel of all these functionals is zero. Add a row
whenever the current common kernel is nonzero. Its dimension
strictly decreases, so finitely many choices suffice.
\end{proof}

The same covariance formula is orthogonally covariant: \(\Gamma\)
is radial and its derivatives transform as tensors. Wick pairing
therefore gives a real continuous \(O(4)\) action on
\(V_e^{\rm Max}\) and orthogonal covariance of its free table.
This free action will be a comparison datum. It is not yet an
action making the interacting hierarchy orthogonally covariant.
% !TeX root = 4DYM_notes.tex
\section{Uniform Gaussian test rows on the Wilson family}
\label{sec:gaussian-position-tests}

The free test matrix of
Corollary~\ref{cor:position-finite-free-tests} must be compared
with actual finite-cutoff observables. A Gaussian polynomial
calculation by itself does not do this: its residual source and
its normalised integral must also be controlled on the same
Wilson history. We prove these two estimates separately.
Only one- and two-insertion rows are needed for detection.

The construction adapts
\cite[Lemmas~14.23--14.26, PDF pp.~146--152]{ArmstrongDraft}.

We keep the slow-background error explicitly and use an inner
interval with only two extra monitor powers. For the normalised integral
we retain the actual positive direct clipping in the principal
action; no identification of that whole action with a Wilson
plaquette sum is required.

\subsection{A sharper scalar-free normalisation estimate}

Put \(q_G=\dim_{\R}\mathfrak g\). At a fixed profile let
\(c_{\rm dir},c_{\rm lin}>0\) be the direct clipping
coefficients. The leading row action contains the nonnegative
term
\[
 \sum_q\operatorname{clip}_h(d_G(P_q)),\qquad
 \operatorname{clip}_h({\varrho})=
 \begin{cases}c_{\rm dir}{\varrho}^2/2,&{\varrho}\le h,\\
                       c_{\rm lin}h{\varrho},&{\varrho}>h.
 \end{cases}
\]
Its other leading row squares are nonnegative. Cubic, quartic,
loop and regular contributions have a bounded absolute price per
site after entry, as in
Lemma~\ref{lem:insertion-vacuum-ratio}.

We need the exact Lie-group dimension in this calculation.
There are constants \(s_H>0,b_H\ge1\), depending only on the
fixed group and metric, such that, for \(0<g\le g_0\),
\begin{equation}
 H_G\{d_G(U)\le g/4\}\ge s_Hg^{q_G},\qquad
 \int_Ge^{-d_G(U)/g}\,\dd H_G(U)\le b_Hg^{q_G}.
 \label{eq:position-haar-dimension}
\end{equation}
Here \(H_G\) is probability Haar measure. In an exponential
chart its density is a positive smooth function, bounded above
and below on a smaller chart, and chord and Lie-algebra distance
are comparable. Integration in polar coordinates gives the two
bounds there. Off that chart the chord has a positive lower
bound, and \(e^{-c/g}\le Cg^{q_G}\). This proves the second
bound on the whole compact group, including its other components.

\begin{lemma}[Triangular plaquette integration on a periodic lattice]
\label{lem:position-triangular-plaquettes}
On a four-dimensional torus of side \(r\ge2\), with \(V=r^4\)
sites, there is a spanning tree and \(3V-3\) distinct plaquettes
such that, for every nonnegative conjugation-invariant function
\(f\) on \(G\), their product integrates to
\[
                         \left(\int_G f\,\dd H_G\right)^{3V-3}.
\]
Tree links have first been fixed to one by gauge invariance.
The construction and the count are independent of coupling.
\end{lemma}
\begin{proof}
Use representatives \(x_a\in\{0,\ldots,r-1\}\). The tree
consists of positive links \((x,a)\) with every higher coordinate
zero and \(x_a<r-1\). Equivalently, a nonzero vertex returns
towards zero in its highest nonzero coordinate. Its number of
links is \(\sum_{a=1}^4(r-1)r^{a-1}=V-1\).
Gauge fixing therefore leaves \(4V-(V-1)=3V+1\) variables.
Retain the four wrapping links \((-e_a,a)\) as untreated variables.

Every other free link is of one of the following types.
If \((x,a)\) has all higher coordinates zero, it wraps in
coordinate \(a\). Put \(y=x+e_a\ne0\) and let \(b<a\) be
the highest nonzero coordinate of \(y\). Select the plaquette
\((b,a)\) based at \(x-e_b\). It contains the selected wrapping
link, another wrapping link with smaller lower-coordinate depth,
one transverse link, and one tree link.
Otherwise let \(k>a\) be the highest nonzero coordinate of
\(x\). Select the plaquette \((a,k)\) based at \(x-e_k\).
Its two \(k\)-links are tree links; its other \(a\)-link has
smaller \(k\)-coordinate or belongs to a lower-dimensional stage.

Process stages \(k=4,3,2,1\) in decreasing order. Within one
stage process the wrapping variables first, in decreasing
lower-coordinate depth, then transverse variables in decreasing
height. In every selected plaquette the distinguished variable
occurs once, and all its other non-tree variables are processed
later. The selected plaquettes are distinct: their higher axis
determines the stage; wrapping selections have base height
\(r-1\), whereas transverse selections have smaller height;
their base and axes then reconstruct the distinguished link.
There is one selection per treated free variable, hence
\((3V+1)-4=3V-3\) selections.

In this order each variable occurs in its own plaquette but in no
remaining selected plaquette. Haar invariance and inversion give
\(\int f(AU^{\pm1}B)\,\dd H_G(U)=\int f\,\dd H_G\).
Fubini integrates the variables successively and proves the
formula. Fixing a tree preserves the integral of any
gauge-invariant integrand: integrate its successive vertex gauge
variables with probability Haar measure. This contributes one,
not a gauge-volume factor.
\end{proof}

\begin{proposition}[Normalisation cost per coarse site]
\label{prop:position-sharp-positive-ratio}
For the actual unmarked representation at level \(k\), write
\(F_k=e^{\Phi_k}Z_{K_k}\), \(V=\#\Lambda_k\) and
\(g=\beta_k^{-1/2}\). There is a profile constant \(Q_j\),
fixed before entry, such that any integrable gauge-invariant
complex function \(A\) satisfying
\(z|A(U)|\le e^{\Phi_k(U)+VX}\), \(z\ge0\), obeys
\begin{equation}
 \frac{z|\int A\,\dd H_k|}{\int F_k\,\dd H_k}
                \le\exp\{V(Q_j+X)+4q_G\log g^{-1}\}.
 \label{eq:position-sharp-positive-ratio}
\end{equation}
The same estimate holds with a finite marker coefficient norm
in place of \(z|A|\). The constant is uniform in cutoff, period
and history. All factors of the original denominator are retained.
\end{proposition}
\begin{proof}
Choose \(\sigma\le\min(1,c_{\rm dir},c_{\rm lin})\).
Since \(g\le h\) after entry, in both branches of the clip,
\[
               \beta_k\operatorname{clip}_h({\varrho})
                                     \ge\sigma {\varrho}/g-1.
\]
For the quadratic branch minimise
\(c_{\rm dir}({\varrho}/g)^2/2-\sigma {\varrho}/g\); its minimum is at least
\(-1\). For the linear branch use \(h/g\ge1\).
The discontinuity at \({\varrho}=h\) causes no difficulty.
Drop the other positive row squares and bound all native
nonquadratic terms by \(C_jV\). Applying the lemma to the
selected plaquette factors and
\eqref{eq:position-haar-dimension} gives
\[
 z\left|\int A\right|
  \le e^{v_kV+V(C_j+6+X)}
                     \{b_H(g/\sigma)^{q_G}\}^{3V-3}.
\]
For the denominator in the same tree gauge, restrict its
\(3V+1\) free links to chord at most \(g/4\). Every plaquette
is healthy there and the exact forced gas is one. Path Schur
sums bound every leading row by \(C_jg\), so its action,
including the other native terms, is at most \(C'_jV\).
The denominator is consequently at least
\[
                  e^{v_kV-C'_jV}(s_Hg^{q_G})^{3V+1}.
\]
The scalar cancels. The powers of \(g\) differ by \(-4q_G\).
Absorb the remaining group, clip and path constants into \(Q_j\).
Absolute values and the marker coefficient norm obey the same
pointwise upper bound, proving the last clause. Positivity used
here is positivity of the full physical density, not of its
individual forced activities.
\end{proof}

\subsection{Comparison intervals and covariance estimates}

Fix a degree bound \(e\), a finite family of one- or two-point
rows, and a compact separated set of their rescaled positions.
Set \({\mathsf M}=P_t\), the monitor at a late internal level \(t\), and
choose a fixed integer \(b\ge4\), sufficiently large for the
finite degree and row request. At a constant-profile interval
with ratio \(L\), choose \(l<t\) and \(k<l\) with
\begin{equation}
 L^{t-l}\ge {\mathsf M}^{b+2},\qquad
 L^{l-k}\ge B_t^{10},\qquad
                      {\varrho}={\varepsilon}_t{\mathsf M}^b.
 \label{eq:position-comparison-window}
\end{equation}
Use the first integers giving these bounds. Then
\({\varrho}/{\varepsilon}_l\le {\mathsf M}^{-2}\), and all row points \({\varrho}y_i\), rounded
to the \({\varepsilon}_t\) grid, lie well inside their level-\(k\) cell.
The rounding error in their rescaled positions is \(O({\mathsf M}^{-b})\).
On these intervals nominal and actual couplings are comparable,
and the actual coupling changes by at most
\(C\log({\varepsilon}_k/{\varepsilon}_t)=O_j(\log {\mathsf M}+\log B_t)\). Its relative change
tends to zero. These facts follow from the exact drift bound,
not from an assumption that the running coupling is constant.

All such interval requirements are compatible with the calendar.
The profile constants, requested derivative orders and \(b\) are
fixed first. Include the Taylor order and the additional input
derivatives of Lemma~\ref{lem:actual-chart-Taylor} in that same
finite request. Choose its ratio with the logarithmic floor for their
finite linear prices, then postpone entry and lengthen the
constant-profile interval. Its length is \(O_j(\log P_t+\log B_t)\),
which is negligible compared with \(t\) on a fixed profile.
Choose one terminal interval for each successive finite request
list. They can have \(k\to\infty\), \({\varrho}\downarrow0\), and
\begin{equation}
           \mathcal C_j^{t-k+2}\le {\mathsf M}^{1/64},
         \qquad Q_j+C\le\frac{\log {\mathsf M}}{32}.
 \label{eq:position-window-prices}
\end{equation}
For the first bound use the two ceilings in
\eqref{eq:position-comparison-window},
\(B_t={\mathsf M}^{o(1)}\), and a ratio floor making
\(\log\mathcal C_j/\log L\) sufficiently small for the fixed
finite requests. One may impose this floor for the entire request
list before \(L\). The second is an entry condition after \(Q_j\)
is known. None of these choices depends on the physical period.

We record the actual covariance calculation. On the cover the
full quadratic curvature symbol has the Maxwell leading symbol,
with
\[
 |\partial^\alpha(K-H)(p)|\le C_{j,e,q}|p|^{1-|\alpha|},
                           \qquad 0<|p|\le1,
\]
through any fixed required order. This follows from the filled
inverse and gradient identities of
Section~\ref{sec:all-scale-linear}: curl factors remove the inverse
square singularity, and replacing each lattice difference by its
linear term gains one power of \(|p|\). The uniform symbol
bounds justify differentiation of the inverse. The part away
from zero is smooth with the same finite derivative bounds.

On the dyadic annulus \(|p|\asymp2^{-v}\), a position
derivative of order \(s\) and \(q>5+s\) integrations by parts
bound the error kernel at \(Rx\), uniformly on a separated
compact, by
\[
 C2^{-v(5+s)}\min\{1,(R2^{-v})^{-q}\}.
\]
Sum below and above \(2^v\asymp R\). The result is
\(CR^{-5-s}\), compared with the Maxwell homogeneity
\(R^{-4-s}\). Here \(R={\varrho}/{\varepsilon}_t={\mathsf M}^b\), so the normalised error
is \(O({\mathsf M}^{-b})\).

The covariance integrated from \(t\) down to \(l\) is the
conditional fast covariance, not the full massless covariance.
The exact Gaussian recursion expresses their difference as the
lifted level-\(l\) covariance. The averaging filter and its
derivatives bound its curvature Fourier density by
\(C|p|^s(1+{\varepsilon}_l|p|)^{-j}\), after the required curls, for any
fixed \(j>4+s\) available at a sufficiently late profile.
Its absolute Fourier integral is at most \(C{\varepsilon}_l^{-4-s}\).
Multiplication by \({\varrho}^{4+s}\) therefore gives the error
\begin{equation}
                                C({\varrho}/{\varepsilon}_l)^{4+s}.
 \label{eq:position-low-mode-error}
\end{equation}
The block mass gap also gives exponential decrease of the fast
kernel at distance measured in units of \({\varepsilon}_l\). Its physical
periodisation has exponentially small image tails since every
period is bounded below and \({\varepsilon}_l\to0\). Summing these tails,
or summing the positive Fourier majorant on the discrete momentum
grid, has the same constant at every period. In particular no
absolute sum of a full massless \(|x|^{-4}\) image kernel is
being used. Fast covariances between different level-\(k\)
cells have exponential row sums, because
\({\varepsilon}_k/{\varepsilon}_l\ge B_t^{10}\).

Localising the covariance adds \(C {\mathsf M}^A e^{-c_jB_t}\), with
fixed \(A\) for the polynomial degree and derivative request.
The factor follows by the inverse-Hessian localisation estimates,
with their full inspected reads. All these estimates concern
the quadratic family of the actual history; its leading Maxwell
coefficient is fixed by calibration and its bounds are uniform
over the terminal history.

\subsection{Actual Taylor data and finite Gaussian comparison}

For a finite polynomial in curvature jets and a covariance \(C\),
define the finite contraction operator
\[
 \mathsf H_C=\exp(\tfrac12\Delta_C),\qquad
       \Delta_C=\sum_{i,j}C_{ij}\partial_i\partial_j,
                      \qquad\mathsf H_C^{-1}=\mathsf H_{-C}.
\]
It is finite because each contraction removes two factors.
Use \(\mathsf H_{-C}\) for contractions internal to one
insertion; contractions between distinct sites are retained.
This gives Wick ordering at the current finite covariance,
before replacing a separated contraction by its Maxwell limit.

% !TeX root = 4DYM_notes.tex
\subsubsection*{Taylor polynomials of the actual insertion}

Fix the finite degree horizon \(e\) and a requested output
derivative order \(q\). All constructions in this subsection
are on the already chosen Wilson history with terminal level
\(n\), physical side \(M\), and \(n\ge t\). The current
coefficient at level \(t\) is denoted by \(x\in V_e\).
The regular and forced families are those of its actual
coefficient graph. In particular
\[
 \|g_t(x)\|_{\rm ins}\le C_{J(t)}|x^\circ|_1.
\]
No auxiliary evolution of coefficient graphs will be used.
The Taylor, curvature-realisation and finite-interval
organisation follows
\cite[Unit~III.11, Lemma~4.13]{ArmstrongNotes};
we establish the required bounds in the present norms below.

Set \(j=J(t)\), \(h=h_t\), \(g=\beta_{n,M,t}^{-1/2}\),
\(P=P_t\), and \(B=B_t\). Thus \(g/h\le C/P\), uniformly
in the prescribed coupling band. A constant denoted by
\(C_{j,L,e,D,q}\) may depend on the group, the fixed profiles
and ratio, the indicated finite orders, the chosen sections,
path budgets and unused support reserves. It is independent
of \(n,M\), and the preceding length of the history.
Write \(C_B=C_{j,L,e,D,q}(1+B)^a\), allowing finitely many
increases of the fixed exponent \(a\).

Let \(\mathcal A_t(x)\) be the whole unmasked regular
function: its primary part is
\[
 x_0+P_e[\mathsf D_{h^{-1}}x^\circ],
\]
and its other terms are the regular remainders. This notation
is used only on configurations on which all these germs are
healthy. It does not identify the physical inserted density
with this function outside that domain.

Cut its regular addresses at connected covering number \(B\).
The near reads lie in a centred cover box of radius \(CB^2\).
Their complements have summed healthy derivative price at most
\begin{equation}
 C_j e^{-cB}|x^\circ|_1.
 \label{eq:actual-Taylor-large-read-tail}
\end{equation}
Indeed each such term has \(n_B(Z)>B\), whereas its original
majorant is weighted by \(e^{100n_B(Z)}\). Keep a fixed
positive part of that weight for later attachment and use the
rest to obtain the displayed exponential. The same argument
applies to the forced majorants after division by \(\sqrt{w_t}\).
Far addresses retain their full original functions and records.

Choose a slightly larger symmetric box containing all near
reads. On its links put \(U_a=\exp(hz_a)\), and set all other
links to one. The components \(z_a\) are real Lie-algebra
coordinates. Their number, denoted by \(d_t\), satisfies
\[
                         d_t\le C_j(1+B)^8 .
\]
One can increase the fixed power if a declared read requires
another fixed halo. The box is nonwrapping on a common late
tail, because \(M/\varepsilon_t\ge M_{\min}/\varepsilon_t\).
All its small coordinate balls, including the outside links
equal to one, are globally healthy.

Evaluate the whole regular sum on this chart and call the
result \(F_t(x;z)\). Its countable sum and each of the finitely
many required derivatives converge uniformly on a smaller
chart. This follows from the original summable derivative
majorants. The far regular terms are included in this definition;
their price remains \eqref{eq:actual-Taylor-large-read-tail}.
In particular the Taylor polynomial is that of an actual
function, rather than a polynomial obtained by deleting
nonlinear terms in a chosen representation.

\begin{lemma}[Actual chart Taylor bounds]
\label{lem:actual-chart-Taylor}
Fix \(D\ge e\), and admit input derivatives through \(D+q+1\).
There is a fixed chart radius \(r_0>0\) such that the polynomial
\[
 T_{t,D}(x;z)=
       \sum_{r=0}^D\frac1{r!}
                  D^rF_t(x;0)[z,\ldots,z]
\]
is complex linear in \(x\), has ordinary degree at most \(D\),
and satisfies
\begin{align}
 |T_{t,D}(x)|_{\rm coeff}
 &\le C_{j,D,q}|x|_1
                     \sum_{r=0}^D\frac{d_t^r}{r!},\notag\\
 |F_t(x;z)-T_{t,D}(x;z)|
 &\le C_{j,D,q}|x|_1\|z\|_\infty^{D+1},
                    \qquad \|z\|_\infty\le r_0 .
 \label{eq:actual-chart-Taylor-bounds}
\end{align}
Here \(|p|_{\rm coeff}\) is the sum of the absolute values
of the coefficients in the fixed real coordinate basis.
After \(s\le q\) chart derivatives, the second bound has
the corresponding power
\(\|z\|_\infty^{\max(D+1-s,0)}\).
The unit polynomial is exactly one.
\end{lemma}
\begin{proof}
An ordered left derivative with direction of size at most \(h\)
is bounded by the original regular norm. In exponential
coordinates it differs from an ordinary scaled derivative by
analytic chart coefficients. Their derivatives through a fixed
order are bounded on a smaller chart. Products and commutators
are pointwise on links, so the conversion constants do not
count links. This gives
\[
 \|D^rF_t(x;z)\|_{\infty\to\mathbb C}
           \le C_{j,D,q}|x|_1,\qquad r\le D+q+1,
\]
where the norm is the multilinear operator norm for the
maximum coordinate norm. Summation over the original
majorants justifies differentiation of the regular sum.

There are \(d_t^r\) ordered choices of \(r\) coordinate
directions. Expansion of the symmetric multilinear form
therefore gives
\[
 \sum_{|\alpha|=r}
       \frac{|\partial^\alpha F_t(x;0)|}{\alpha!}
       \le C_{j,D,q}|x|_1\,\frac{d_t^r}{r!}.
\]
This proves the coefficient estimate. Taylor's integral
formula on the segment from zero to \(z\) gives
\[
 F_t(x;z)-T_{t,D}(x;z)
 =\frac1{D!}\int_0^1(1-u)^D
       D^{D+1}F_t(x;uz)[z,\ldots,z]\,\dd u.
\]
Its absolute value proves the second estimate. Apply the same
formula to each derivative for the last assertion; the admitted
order \(D+q+1\) is sufficient. Linearity follows before any
estimate is taken. The exact physical unit has zero remainder,
so its chart function and Taylor polynomial are both one.
\end{proof}

The order \(D\) is a Taylor order; it need not equal the
engineering horizon \(e\). This distinction is essential.
The section \(P_e\) can have curvature degree greater than
\(e\), whereas the polynomial just defined has degree at most
\(D\) by its definition. All orders used below form a finite
request at a fixed stage. When several window exponents
are needed for the same coefficient map, choose one \(D\)
for their entire finite list, or their specified compact
interval. The polynomial and calibration then remain
the same.

\subsubsection*{A local curvature realisation}

On the finite box use the same rooted axial tree as in the
local gauge estimates. Read plaquette curvatures after
transporting them to the origin. Write \(c(U)\) for their
finite vector of Lie-algebra components. The linear surface
sum along that tree is denoted by \(S_t\). The elementary
plaquette stripping and the product expansion used in
Lemma~\ref{lem:box-path-comparison} give
\begin{equation}
 \|a(U)-S_tc(U)\|_\infty\le C_B\tau^2,\qquad
 \|a(U)\|_\infty+\|S_tc(U)\|_\infty\le C_B\tau
 \label{eq:actual-Taylor-axial-replacement}
\end{equation}
when the plaquettes in the box have size at most \(\tau\).
Here \(a(U)\) are the logarithms of the tree-gauge links.
The constants are polynomial in the box radius.

To check the quadratic error explicitly, strip each link
into its ordered list of transported plaquettes. The linear
term in the logarithm of that product is their sum. Every
other term has at least two plaquette increments. For a list
of length \(m\le C_B\), its error is bounded by
\(C m^2\tau^2\) when \(m\tau\) is inside the fixed chart.
Replacing the smooth curvature contraction by the logarithm
adds another \(Cm\tau^2\). The finite path comparison controls
the changes of transport in the same way. These estimates
also follow by the second-order integral remainder for the
finite product. They prove \eqref{eq:actual-Taylor-axial-replacement}.
The same calculation with finitely many ordered derivatives
has a polynomial \(B\)-price.

For independent variables \(X\) in the curvature component
space define
\begin{equation}
 p_t(x;X)=T_{t,D}\bigl(x;(g/h)S_tX\bigr).
 \label{eq:actual-Taylor-curvature-polynomial}
\end{equation}
It is an actual Taylor polynomial followed by a linear
substitution. In particular it still has degree at most \(D\).
Its homogeneous part of degree \(r\) has coefficient bound
\begin{equation}
 |p_t(x)^{[r]}|_{\rm coeff}
       \le C_B|x|_1(g/h)^r,\qquad 0\le r\le D.
 \label{eq:actual-Taylor-homogeneous-price}
\end{equation}
No recovery from a completed remainder is involved.

The entire chart function is invariant under simultaneous
adjoint conjugation. Its Taylor polynomial has this same
invariance. Thus evaluation at \(X=c(U)/g\) is gauge
invariant: all its slots transform at the same origin.
Reality and signed equivariance can be retained by averaging
over the original finite real and signed actions. The
coefficient and error estimates only change by their fixed
finite constants. The whole healthy insertion has the
required real structure by its exact physical density
identity; the far regular tail in this averaging retains
the bound \eqref{eq:actual-Taylor-large-read-tail}.

For later estimates use the smaller comparison threshold
\[
                             \tau=gP^{1/4}.
\]
It is an auxiliary threshold for the comparison, not a
change of the original physical masks. After entry
\(C_B\tau/h\le P^{-2/3}\). On this chart, the preceding
Taylor and product estimates imply
\begin{align}
 |\mathcal A_t^{\rm near}(x;U)-p_t(x;c(U)/g)|
 \le C_B|x|_1
       \left\{P^{-2(D+1)/3}+\tau^2/h+e^{-cB}\right\}.
 \label{eq:actual-Taylor-physical-approximation}
\end{align}
Indeed use gauge invariance for the near part, set outside
links to one, and include the far chart sum with its original
exponential bound. The difference of the two polynomial
evaluations is bounded by the elementary telescoping product
rule, or by its first derivative on the segment joining the
two arguments. Since \(C_B\tau/h<1\), that derivative has
price \(C_B|x|_1\). Finally
\(\tau^2/h\le CgP^{-1/2}\); replacing this stronger small
quantity by a fixed power of \(P\) times \(g\) is harmless.
The estimates use only the original physical reads.

\subsubsection*{One finite Gaussian polynomial estimate}

We use the polynomial heat-operator calculus of
\cite[Proposition~2.6, PDF p.~5; Proposition~3.18, PDF p.~21]{BS15a}.
We give its elementary coefficient version to avoid changing
the present norms. The localisation estimates of
\cite[Propositions~1.11--1.12, PDF pp.~10--11]{BS15b}
motivate keeping the polynomial and remainder bounds together;
their test-function seminorm is not assumed here.

For a matrix \(Q\) put
\[
 \|Q\|_{\max}=\max_{a,b}|Q_{ab}|,\qquad
 \Delta_Q=\sum_{a,b}Q_{ab}\partial_a\partial_b,\qquad
 \mathsf H_Q=\sum_{s=0}^{\lfloor D/2\rfloor}
                            \frac{\Delta_Q^s}{2^ss!}.
\]
The last expression is an operator on polynomials of degree
at most \(D\). It is algebraically defined even when \(Q\)
is not a covariance. Whenever \(Q\) is a real positive
covariance it equals Gaussian expectation after translation.

\begin{lemma}[Finite Gaussian coefficient calculus]
\label{lem:finite-Gaussian-coefficient-calculus}
Let \(p\) have degree at most \(D\), and define
\[
 B_D(K)=\sum_{s=0}^{\lfloor D/2\rfloor}
                 \frac{D!}{(D-2s)!\,2^ss!}K^s .
\]
Then
\begin{align}
 |\mathsf H_Qp|_{\rm coeff}
 &\le B_D(\|Q\|_{\max})|p|_{\rm coeff},\notag\\
 |(\mathsf H_{Q_1}-\mathsf H_{Q_0})p|_{\rm coeff}
 &\le \tfrac12D(D-1)B_D(K)
       \|Q_1-Q_0\|_{\max}|p|_{\rm coeff},
 \label{eq:finite-Gaussian-coefficient-bounds}
\end{align}
where \(K=\max(\|Q_0\|_{\max},\|Q_1\|_{\max})\).
The second bound is zero for \(D<2\).
The constants do not depend on the number of variables.
\end{lemma}
\begin{proof}
On a monomial of degree \(r\), the sum of the absolute
coefficients of its second derivatives, summed over all
ordered pairs of indices, is \(r(r-1)\). Consequently
\[
 |\Delta_Qp|_{\rm coeff}
       \le D(D-1)\|Q\|_{\max}|p|_{\rm coeff}.
\]
Each further application lowers the degree by two.
Iteration and summation prove the first bound.
Constant-coefficient differential operators commute. With
\(Q_u=(1-u)Q_0+uQ_1\), finite differentiation therefore gives
\begin{equation}
 \mathsf H_{Q_1}-\mathsf H_{Q_0}
       =\tfrac12\int_0^1
                \mathsf H_{Q_u}\Delta_{Q_1-Q_0}\,\dd u.
 \label{eq:finite-Gaussian-difference-identity}
\end{equation}
Apply the first estimate and the second-derivative bound
inside this integral. This proves the second assertion.
All sums are finite.
\end{proof}

Linear substitutions and means are handled by the same
coefficient norm. For a list of linear forms their difference
is expanded using
\[
 \prod_{a=1}^r u_a-\prod_{a=1}^r v_a
   =\sum_{a=1}^r
         \left(\prod_{c<a}u_c\right)(u_a-v_a)
                              \left(\prod_{c>a}v_c\right).
\]
Thus a Gaussian step difference is the sum of its source,
mean, linear-substitution and covariance differences, each
times a finite polynomial in the corresponding bounds.
Applying a fixed derivative merely lowers polynomial degree.
This single calculation also proves the estimates for a
product of at most \(r\) named sources, in its square-zero
algebra. No replacement of joint coefficients by products
of singleton errors is made.

All contractions of size \(P^{-2}\) in threshold variables
are retained in \(\mathsf H_Q\). They are not assigned to an
inverse power of the coupling. The polynomial bound \(B_D\),
rather than an exponential bound in a large covariance norm,
will also be used after the fixed spatial rescalings.

\subsubsection*{Calibration at one actual level}

Let \(Q_t\) be the full cover quadratic covariance of the
curvature variables normalized by \(g\), on the linear data
of this same history. Its entries on the declared source
reads are bounded by \(C_j\). The full covariance is used
for internal contractions; the lower integration endpoint
does not enter this definition.

For a polynomial in these curvature variables, let
\(\mathcal R_e^{\rm lin}\) denote its reading on the
linear circulations of radial test potentials, through
engineering weight \(e\). This is a linear map to \(V_e\).
The potentials are the same finite radial jets used to
define \(\Pi_e\). In this reading each curvature variable
is \(d_1A\), with its declared finite differences; parallel
transports have their linear value. Define
\begin{equation}
 \mathsf B_t x
  =\mathsf D_{h/g}\mathcal R_e^{\rm lin}
                                \mathsf H_{Q_t}p_t(x).
 \label{eq:actual-Taylor-calibration-matrix}
\end{equation}
This is a finite matrix obtained from derivatives of
the actual insertion. Its construction does not depend
on the chosen row positions.

\begin{lemma}[A uniformly invertible local calibration]
\label{lem:actual-Taylor-calibration}
On a common late tail, uniformly in \(n,M\),
\[
 \|\mathsf B_t-1\|_{1\to1}
       \le C_B\{h+P^{-2}+e^{-cB}\}\le\tfrac14,
 \qquad
                  \|\mathsf B_t^{-1}\|_{1\to1}\le2 .
\]
The maps preserve the unit, real structure and signed
coordinate actions.
\end{lemma}
\begin{proof}
First omit \(\mathsf H_{Q_t}\).
For every near regular address its physical test jets
through \(e\) vanish. The section reading is exactly \(x\).
Truncation at ordinary degree \(D\ge e\) changes none of
these jets: every test coordinate has positive weight.
The far chart terms contribute at most \(C_Be^{-cB}|x|_1\).

Replacing the actual test holonomies and the axial product
by their linear expressions gives the remaining error
\(C_Bh|x|_1\). To see the threshold power, a term of
ordinary degree \(r\) has coefficient \(h^{-r}\) before
the chart normalization is undone. A nonlinear replacement
contains at least \(r+1\) test factors. Multiplication by
the output factor-count power therefore leaves at least
one factor \(h\). The finite product rule sums these
terms; its path and jet moments have a polynomial
\(B\)-price. This argument concerns only jets through
\(e\), and does not assert that an ordinary Taylor
polynomial is separately invariant under nonlinear
local gauge transformations.

Now keep all Gaussian contractions. By
\eqref{eq:actual-Taylor-homogeneous-price}, the input
degree-\(r\) coefficients have the factor \((g/h)^r\).
A term with \(s\ge1\) pairs has degree \(r-2s\).
The output normalization \(\mathsf D_{h/g}\) leaves
precisely \((g/h)^{2s}\). The remaining reading moments
and covariance entries are bounded by \(C_B\).
The finite sum in
Lemma~\ref{lem:finite-Gaussian-coefficient-calculus}
is consequently at most \(C_BP^{-2}|x|_1\).
This proves the matrix estimate.

The factor \(C_B\) is a fixed power of \(B\), whereas
\(P\) dominates every such power; \(h\) and \(e^{-cB}\)
are smaller still. Entry gives the strict \(1/4\) bound.
The Neumann series has norm at most \(4/3\), so the
stated bound two is valid. The unit has polynomial one
at every operation. All covariance and chart data have
the prescribed real and signed symmetries; the fixed
finite averaging preserves them and the estimates.
\end{proof}

If \(u\) is a chosen real free polynomial representative,
put \(R=\varrho/\varepsilon_t=P^b\), and let
\(\delta_R^{\rm wt}\) multiply a label of engineering
weight \(w\) by \(R^w\). The current coefficient
\begin{equation}
 x_t(u)=\mathsf B_t^{-1}
                 \mathsf D_{h/g}\delta_R^{\rm wt}u
 \label{eq:actual-Taylor-current-lift}
\end{equation}
then satisfies
\[
 \mathcal R_e^{\rm lin}\mathsf H_{Q_t}p_t(x_t(u))
                         =\delta_R^{\rm wt}u,\qquad
 |x_t(u)|_1\le C_jP^{e/2}R^e|u|_1.
\]
The common physical factor \(\varepsilon_t^e\) is kept
outside this finite matrix when these dimensionless
coefficients are converted to the original insertion.
Transport back through the fixed interval
\([\operatorname{ref}(e),t]\) gives the reference lifts.
Their bounds may depend on that fixed interval, but not
on \(n,M\). The unit lift is exactly the original \(u_e\).

The Wick symbol of the comparison polynomial is
\(\mathsf H_{Q_t}p_t(x_t(u))\); equivalently its evaluation
is the Wick evaluation of that symbol, using
\(\mathsf H_{-Q_t}\). Thus internal contractions have been
matched before any separated covariance is replaced by
its Maxwell limit. There is no fitted correlation matrix,
and no limiting Gaussian coefficient graph.

\subsubsection*{The comparison is made on a finite interval}

We now pass from these polynomial data to the complete source.
Fix the finite row family, its separated compact and the output
orders \(q,s\). Choose \(b\) large enough for the finite degrees,
as specified below. The order \(D\) is chosen afterwards.
At a sufficiently late profile admit the input orders
\(D+q+1\) and the corresponding spatial moments. These are
finite-order requests on the original smooth germs. They
are not analyticity assumptions, nor regularity assumptions
on any Borel forced multiplier.

For a complete source at level \(i\), its price below means
the sum of its scalar coefficient price per cell, the
anchored regular derivative majorants, and the forced sup
majorants divided by \(\sqrt{w_i}\). Covering weights and
the marker norm are those of
\eqref{eq:joint-marker-norm} and
\eqref{eq:joint-local-source}. At most two names occur
per comparison cell. The constants below may be enlarged
for any fixed finite arity.

\begin{lemma}[A finite-interval source comparison]
\label{lem:actual-Taylor-window-comparison}
Use the lifts \eqref{eq:actual-Taylor-current-lift}
and the interval \eqref{eq:position-comparison-window}.
For a sufficiently large fixed \(b\), then a sufficiently
large finite \(D\), the complete corrected source at \(l\)
has price at radius one at most
\begin{equation}
 \mathcal C_j^{\,t-l+2}
 \left[
 P^{-b/2}+P^{-3}
 +C_{j,L,e,D,q,s}P^{a_{e,D,q,s,b}}
       \{\widehat\beta_t^{-1/16}+e^{-c_jB}\}
 \right].
 \label{eq:actual-Taylor-window-price}
\end{equation}
The exponent \(a_{e,D,q,s,b}\) is finite. All bounds are
uniform in the terminal cutoff and physical side.
The constant \(\mathcal C_j\) contains only the finite
free transport, section and attachment prices chosen
before the ratio. Post-ratio perturbation constants
remain in the last term and in the entry conditions.
\end{lemma}
\begin{proof}
We give the construction and estimates before restoring
the physical masks.

\emph{1. Geometry and derivative prices.}
The source box has fine radius \(CB^2\). A selected
kernel read of radius \(CB\) at level \(i\), pulled to
level \(t\), has radius at most
\(CB^2L^{t-i}\). Summing over \(l\le i\le t\)
is a geometric sum. The union of the necessary finite
reads therefore has at most
\begin{equation}
 C_{j,L}(1+B)^aP^{4(b+2)}
 \label{eq:actual-Taylor-window-read-count}
\end{equation}
fine colour coordinates, after including the two
ceiling constants. This is a bound for the reads
attached to one specified cell, not for the whole torus.
Native long reads keep their exponential tails.

Here is also a direct bound for the covariance of the
finite Gaussian path, before completing any curvature
slots. Use the smooth free representative of
Section~\ref{sec:all-scale-linear} for this calculation.
Its exact gradient identities permit the subsequent
return to the declared local tree slices.
Normalize a level-\(i\) potential by its actual
coupling scale \(g_i\). Its conditional free interpolation
from the next coarser potential has row norm at most
\(2C_j/L\). Include \(L\ge4C_j\) in the ratio floor.
On the stacked path coordinates let \(\mathsf E\)
be this strictly triangular interpolation matrix,
and let \(\mathsf C\) be the block diagonal matrix
of the fresh conditional covariances. Their selected
gaps give
\[
 \|\mathsf E\|_{\infty\to\infty}\le\tfrac12,\qquad
                    c_*1\le\mathsf C\le C_*1 .
\]
The constants \(c_*,C_*>0\) may depend on the fixed
ratio and profiles. They do not depend on path length.
For \(d_{\rm win}\) retained coordinates, the joint
covariance is
\[
 \mathsf C_{\rm win}=(1-\mathsf E)^{-1}
                    \mathsf C(1-\mathsf E)^{-T}.
\]
The geometric series bounds the maximum row norm of
\((1-\mathsf E)^{-1}\) by two. A matrix with maximum
row norm \(K\) has Euclidean operator norm at most
\(K\sqrt{d_{\rm win}}\). Apply this to the matrix
and to \(1-\mathsf E\), whose row norm is at most
\(3/2\). It follows that
\[
 \|\mathsf C_{\rm win}\|_{2\to2}
 +\|\mathsf C_{\rm win}^{-1}\|_{2\to2}
                      \le C_{j,L}d_{\rm win}.
\]
To return to a tree slice, a potential is changed
by a gradient whose vertex values are sums on tree
paths. Each path has at most \(d_{\rm win}\) edges.
The change matrix and its inverse therefore have
a polynomial norm in \(d_{\rm win}\); the reverse
change uses the reverse path construction.
For a smooth representative the same assertion
follows by first taking the orthogonal representative
on the finite gauge quotient. The elementary
path Poincare inequality on a pinned connected
graph bounds its inverse by a fixed power of
the number of edges. Thus, in the declared slice,
the last displayed bound has a possibly larger
fixed power of \(d_{\rm win}\).
Every covariance marginal retains the resulting spectral
bounds. Native cuts, restricted coordinate charts
and fixed local tree slices use the same selected
principal estimates; their unretained tails have
already kept their exponential price. This calculation
does not require independent fields in different cells.

The fixed-order physical-measure differentiation of
Lemma~\ref{lem:regular-Borel-measure-derivative}
can be performed on these same reads. We spell out why
it has only a polynomial loss. A log-density derivative
is a sum of traces of products of inverse covariance
matrices and their derivatives, and quadratic or linear
expressions in normalized Gaussian coordinates.
At order \(q\) there are at most \(q\) differentiated
factors. The selected principal inverses have the
common gap of Lemma~\ref{lem:selected-window-jets}.
Undifferentiated propagation is controlled by the
all-scale free columns, including their explicit
spatial rescalings. Thus a derivative selects at most
\(q\) levels and \(q\) coordinate positions; it does
not multiply an inverse bound once for every level.
The level choices cost at most \((1+t-l)^q\).
Traces and Gaussian moments of the resulting
fixed-degree polynomial cost a fixed power of
the count \eqref{eq:actual-Taylor-window-read-count}.
The same statement holds for the required anchor
moments. Consequently the total score price is
\begin{equation}
                  C_B P^{a_{q,s}(b)}.
 \label{eq:actual-Taylor-window-score}
\end{equation}
The exponent can be fixed from the derivative order
and geometry before \(D\) is chosen. More explicitly,
the inverse derivative partition formula, the trace
bound and the fixed-degree Gaussian moment bound use
at most \(C(q+s+1)\) powers of \(d_{\rm win}\) and
at most \(C(q+1)\) powers of \(1+P\).
Insert \eqref{eq:actual-Taylor-window-read-count}
to choose a finite \(a_{q,s}(b)\) of this form.
No coefficient degree occurs in this score count.
All remaining
kernel and ratio constants occur in \(C_B\).
This is the finite-read version of the score calculation
\eqref{eq:selected-marginal-score}; no trace over the
exterior volume occurs. The physical variables in
each Borel multiplier are held fixed throughout.

\emph{2. Taylor errors and their order.}
On the smaller source chart the normalized potentials
have size at most \(P^{-2/3}\). The current lift has
price at most \(C_jP^{e/2+be}\). Thus
\eqref{eq:actual-Taylor-physical-approximation},
followed by the score bound, gives Taylor-error price
\begin{equation}
 C_B P^{e/2+be+a_{q,s}(b)-2(D+1)/3}.
 \label{eq:actual-Taylor-window-remainder}
\end{equation}
Choose \(D\ge e\) with
\[
 \tfrac23(D+1)>e/2+be+a_{q,s}(b)+b+4.
\]
All fixed powers of \(B\) are \(P^{o(1)}\).
The resulting term is at most \(P^{-b}\) after entry.
Its finitely many derivative statements are obtained
either by the differentiated Taylor remainder or by
the preceding physical-measure score. In the latter
case only the value bound of the remainder is used.
There is no differentiation of its auxiliary chart test.

The curvature-realisation error has instead the factor
\(\tau^2/h\). The current lift and all the finite
polynomial and score factors multiply it by a fixed
power of \(P\). It is therefore
\(\widehat\beta_t^{-1/2+o(1)}\) at a fixed request.
The large-read tail retains \(e^{-cB}\).
This explains why a high Taylor order is used to pay
the ordinary approximation error, whereas the
nonlinear curvature conversion is paid by the coupling.

\emph{3. Gaussian propagation and nonlinear differences.}
Compose the quadratic steps on this finite interval
before taking polynomial coefficient bounds.
Gaussian convolution gives the conditional fast
covariance and free interpolation of the same
linear history. It is a single finite polynomial
calculation, to which
Lemma~\ref{lem:finite-Gaussian-coefficient-calculus}
applies. In particular a loss for every individual
pairing at every step is not iterated.

For comparison with the physical integral, use the
exact Gaussian/Haar changes of variables and the
stationary-density identity
\eqref{eq:auxiliary-stationary-density-expanded}.
Pair the stationary prefactor, affine square,
determinant, noise constant and flat vacuum scalar
with the corresponding unmarked denominator before
estimating the marked coefficient. The normalized
mean difference is its physical difference divided
by \(g\). In a common small chart its price is
\(C_Bh^2/g\), in addition to its native cut tail;
the covariance difference has price \(C_Bh\).
These are Taylor expansions of the actual row and
predictor formulas. Their first jets agree with
the calibrated free data. The inverse difference
identity and Lemma~\ref{lem:selected-window-jets}
retain the native \(e^{-cB}\) errors.
Since \(h^2/g\le CgP^2\), these errors have a
positive inverse coupling power after fixed
polynomial losses.

For clarity, on a common positive Gaussian chart
the normalized log-density ratio is
\[
 \tfrac12\log\frac{\det P_1}{\det P_0}
 -\tfrac12 z^T(P_1-P_0)z
 +z^TP_1u-\tfrac12u^TP_1u,\qquad
             u=(\mu_1-\mu_0)/g.
 \]
Here \(P_\nu\) are the two precision matrices and
\(z\) uses the reference mean. The determinant
difference is the integral of
\(\operatorname{Tr}(P_v^{-1}(P_1-P_0))\).
The common precision reserve pays its quadratic
charge. Its traces count only the declared reads.
Every derivative of a ratio difference still
contains one data difference. This proves the
same inverse coupling price through the required
order, with Borel tests left in their physical
variables.

The remaining smooth difference letters have prices
\[
 gP^3,\qquad (1+\Lambda)h^2,\qquad
 \|E_i\|_{\mathcal E,100},\qquad
 \frac{1+|\alpha_{n,i}|}{\widehat\beta_{i+1}},
 \qquad C_Be^{-cB}.
\]
These are the combined cancelled letters of
\eqref{eq:central-explicit-prices} and the coupling
reset. The extensive constant square is not one
of the small letters. Interpolate a smooth
exponential with \(e^{-V}-1=-\int_0^1Ve^{-uV}\,\dd u\).
For an incoming masked activity keep its physical
mask fixed. For an integrated packet interpolate
the packet amplitude, retaining the seed energy
which paid it.

Take the marked connected coefficient before its
absolute majorant. A difference term has one
source occurrence and one distinguished difference
letter. Lemma~\ref{lem:marked-interacting-log}
then bounds its attachments by the original
rooted tree sum. Terms not attached to a source
cancel with the unmarked coefficient. Product
derivatives cost \((1+r)^q\) for \(r\) letters,
which the unused covering reserve pays.
There is one precision reserve per complete
physical word. Source-free components and
independent copies of a shared forced member
are not inserted.

Over the \(t-l\) steps the nonlinear errors can
have a fixed post-ratio factor
\(C_{j,L,e,D,q,s}^{\,t-l+2}P^a\).
It is important to keep that factor. At this
fixed ratio
\[
 t-l=O_{j,L,b}(\log P),\qquad
 \log P=o(\log\widehat\beta_t),
 \]
so it is \(\widehat\beta_t^{o(1)}\).
The errors just obtained are bounded by
\(C_{j,L,e,D,q,s}P^a
  (\widehat\beta_t^{-1/16}+e^{-c_jB})\)
on their admitted tail, after decreasing \(c_j\)
if necessary. This payment uses entry after the
ratio; it makes no assertion that a post-ratio
constant was chosen before that ratio.

For propagation of the remaining source majorants,
use the free-column moments with the engineering
factors left in the physical insertions. Their
finite section, coefficient and anchor constants
precede \(L\), as in
Lemma~\ref{lem:insertion-leading-map}.
An additional contraction in threshold variables
costs \(C_{j,L,D}P_i^{-2}\); every nonquadratic
extra letter has the vanishing prices above.
Entry makes their sum at most one. The
coefficientwise product rule and rooted sum
therefore give a linear bound \(\mathcal C_j\)
per step, with \(\mathcal C_j\) fixed before
the ratio. This is the finite-interval source
bound, not a comparison of two evolving graphs.

\emph{4. Separated polynomial moments.}
The matching identity
\eqref{eq:actual-Taylor-current-lift} specifies the
low engineering jets of the Wick symbol. Its
other coefficient kernels have the moment bounds
following from
\eqref{eq:actual-Taylor-homogeneous-price}.
We explain their spatial gain. A separated
curvature covariance and its derivatives satisfy
\[
 |D^\alpha K_t(Ry+a)|
             \le C_{\alpha}R^{-4-|\alpha|},
       \qquad |a|\le CB^2,\quad y\text{ in the fixed compact}.
\]
Taylor-expand in \(a/R\), through the required
finite spatial orders, in each cross-contraction.
A monomial with \(r\) curvature factors and
total difference order \(v\) has engineering
weight \(2r+v\). Its coefficient moments of
weight at most \(e\) are precisely the linear
radial reading used in \(\mathsf B_t\).
One can verify this directly from the radial
identity
\[
 A_b(y)=\sum_a y_a\int_0^1 u\,F_{ab}(uy)\,\dd u
\]
for a closed linear curvature jet. It recovers
the same finite radial test potentials. Thus
subtracting the prescribed low jets cancels
all these moments. A first uncancelled weight
is at least \(e+1\), and its contribution after
the factor \(R^e\) is \(O(R^{-1})\).

The Taylor integral remainder has exactly the
same bound: a derivative of order \(v+1\)
gives the additional inverse \(R\), and the
coefficient moment of the offsets is at most
a fixed power of \(B\). The input normalisation
adds at most \(P^{e/2}\) per source. In a product
of two sources, expand each Wick symbol into its
matched part and this higher-weight remainder.
A term containing a remainder consequently
costs at most \(C_BP^eR^{-1}\). The same
argument handles each fixed finite arity.
Choose \(b>2e+8\); then this is at most
\(P^{-b/2}\) on a late tail.

For the matched parts, replacing a separated
contraction by its Maxwell value costs
\(C R^{-1}\), by the preceding Fourier estimate.
The finite Gaussian coefficient lemma, or a
telescoping product of contractions, bounds all
one- and two-point errors by the same price.
For two points subtract the two singleton means
as well: if \(U\) bounds the limiting moments
and \(T,Z\) their singleton and other pair
errors, the resulting bound is
\[
                         T+Z+(U+T)T+TU .
\]
The mean subtraction is part of this calculation.

The full covariance used for internal Wick
ordering differs from the fast covariance by
the lifted slow covariance. Its normalized
jet contractions cost
\((\varrho/\varepsilon_l)^{w+w'}\), with
\(w,w'\ge2\). A remaining slow background jet
costs \(C_jP(\varrho/\varepsilon_l)^w\).
The matched polynomial is fixed, so these
errors are at most \(CP^{-8}\) and \(CP^{-3}\),
respectively. Higher-weight terms have already
received their extra inverse \(R\).
All these statements hold after the fixed
requested derivatives on a separated compact.
They give the first two terms of
\eqref{eq:actual-Taylor-window-price}.

\emph{5. Failed charts, forced norms and physical reassembly.}
The comparison threshold \(\tau\) is used only
to split the integrals. It replaces no physical
health test. Under the common Gaussian law,
failure of the source chart has a witness of
size \(P^{1/4}/C_B\); the slow background and
the nonlinear mean error are smaller than
half this margin after entry.
The finite-range joint tilt estimate gives
the value and fixed-order derivative price
\[
 C_B\widehat\beta_t^aP^{a'}
                          e^{-P^{1/2}/C_B}.
\]
The coordinate count in
\eqref{eq:actual-Taylor-window-read-count}
only adds another polynomial. Evaluation of
the degree-\(D\) curvature polynomial on
the compact group adds at most a fixed power
of \(\widehat\beta_t\), since \(g^{-1}\)
occurs only finitely many times.

Keep the exact eligibility record of each
such component. An empty bad inventory remains
regular, with its physical-measure derivative
bound. A nonempty inventory remains forced.
For an old forced mark of amplitude \(k_\alpha\),
write its input price as \(k_\alpha/\sqrt{w_t}\)
before applying the joint Gaussian/Haar packet
estimate. Division by the output allowance
costs at most the extra exponential
\[
               \exp\{(P_t^{1/5}+P_l^{1/5})/2\}.
\]
A new mask-restoration term contains its
actual coarse forced member and leaves
\(\sqrt{w_l}\) after division. Since adjacent
couplings stay in the stated band on this
finite interval,
\[
 P_t^{1/2}/C_B\gg P_t^{1/5}+P_l^{1/5}
                         +a\log\widehat\beta_t,
\]
the warning estimate and the original stronger
physical packet estimates pay these operations.
No pointwise division by a forced activity
or its gas is used.

Finally use the complete-source reassembly:
all proper subset coefficients are retained,
each old forced member occurs once, and the
old inventories, masks, predictor/frame
probabilities and noise factors stay in the
same normalized physical integral.
The support compression and its reserved
exponential pay the common-anchor moments.
Periodisation is dominated by the same
absolute cover sums. Wrapping reads belong
to the large-read tail, not to an invented
bad inventory. This proves the asserted
regular and divided forced prices together
with the scalar price per cell, with constants
independent of cutoff, volume and the number
of replicated cells.
\end{proof}

This construction uses the actual coefficient graph
only at its chosen starting level. The bounded inverse
in Lemma~\ref{lem:actual-Taylor-calibration} is a finite
reading of that graph, followed by Taylor approximation
and a controlled finite integration interval.
There is no inference from the norm of a completed
remainder to the coefficients of a previous polynomial.

\subsection{The full corrected source and its normalised integral}

For one cell containing at most two trial insertions \(F_1,F_2\),
let \(G_1,G_2,G_{12}\) be their prescribed free moments and put
\[
 C^{\rm Max}(z)=G_1z_1+G_2z_2+(G_{12}-G_1G_2)z_1z_2.
\]
For one slot omit the second slot. The corrected coefficients of
\(e^{-C^{\rm Max}(z)}\prod_i(1+z_iF_i)\) are
\begin{align}
 R_i&=F_i-G_i,\notag\\
 R_{12}&=F_1F_2-G_1F_2-G_2F_1+2G_1G_2-G_{12}.
 \label{eq:position-relative-pair-errors}
\end{align}
Thus the pair subtraction includes the singleton means. These
are identities in the square-zero algebra.

At level \(l\), the total corrected source price per cell,
including regular and normalised forced terms, is bounded by
\begin{equation}
 \mathcal C_j^{t-l+2}\mathcal E_t,\qquad
 \mathcal E_t={\mathsf M}^{-b/2}+{\mathsf M}^{-3}
       +C_{j,L,e,D,q,s}{\mathsf M}^{a_{e,D,q,s,b}}\{
           \widehat\beta_t^{-1/16}+e^{-c_jB_t}\}.
 \label{eq:position-gaussian-source-initial}
\end{equation}
Here and below the finite degree and derivative request is fixed.
Lemma~\ref{lem:actual-Taylor-window-comparison} proves this
estimate in the complete regular and divided forced norms.
The weaker first power includes the higher spatial moments
of the actual Taylor polynomial as well as rounding and the
lattice symbol error. For completeness,
a Wick polynomial moment is a finite polynomial in its
cross-contractions. If their values are bounded by \(U\) and
their errors by \(\epsilon\), telescoping a product of \(s\)
contractions bounds its error by
\(s\epsilon(U+\epsilon)^{s-1}\). For the pair cumulant also
subtract the two mean factors. If \(T\) bounds singleton errors
and \(Z\) the other pair error, its price is at most
\(T+Z+(U+T)T+TU\). These finite expressions prove the
polynomial part of \eqref{eq:position-gaussian-source-initial}.

The estimate is for the full source. Within each cell integrate
its one or two marks jointly as a finite polynomial packet; no
smallness of an uncontracted fine source is assumed there.
Different cells have fast-covariance row sums exponentially small
in \({\varepsilon}_k/{\varepsilon}_l\); long activity reads have their reserved covering
weight. These prices pay all cross-cell attachments uniformly in
the number of cells. Terms meeting physical failures keep the
original masks, incoming activities, inventories and Haar ratios;
their packet price absorbs the fixed source polynomial before
its forced allowance is divided out. Unforced failures stay
regular. Localisation and the full physical reassembly of the
earlier sections therefore add only the last terms already in
\eqref{eq:position-gaussian-source-initial}.

Use the constant source radius \(\alpha={\mathsf M}^{1/16}\) for each name.
At most two names occur per cell, so its initial price is
multiplied by at most \(\alpha^2={\mathsf M}^{1/8}\). Over the second
interval \([k,l]\), let \(a_i\) bound the complete corrected
regular and normalised forced income rooted at a cell, including
its scalar modes. The connected expansion gives
\begin{equation}
                   a_i\le K_j a_{i+1}
                                  +2{\mathsf M}^{1/2}a_{i+1}^2,
                  \qquad \mathcal C_j\ge K_j+2.
 \label{eq:position-source-iteration}
\end{equation}
The linear part uses the finite Gaussian maps with their
engineering factors kept in the physical sources; its remaining
operator bound is fixed before the ratio. Nonquadratic
post-ratio errors are made at most one by entry. The quadratic
part has two corrected source factors. Rooted summation and the
finite-jet product rule give its product of norms. Extra fixed
constants and powers of \(B_t={\mathsf M}^{o(1)}\) are at most \({\mathsf M}^{1/2}\)
after entry. Gaussian contraction powers are already in the
source norm; score losses on Borel failures retain their
exponential charge. They are not additional unpriced powers
of \({\mathsf M}\). At this stage occupancy is bounded by two names per
cell, rather than separated names. The radius is constant, so
the separated-name radius-increase lemma is not invoked.

Choose entry and then a full terminal-cutoff tail so that
\begin{equation}
 \alpha^2\mathcal C_j^{t-k+2}\mathcal E_t\le {\mathsf M}^{-3/4}.
 \label{eq:position-source-entrance}
\end{equation}
The first two terms have exponents at most
\(1/8+1/64-3<-3/4\).
The coupling and localisation terms are smaller than any fixed
negative power of \({\mathsf M}\), since
\(\log {\mathsf M}=o(\log\widehat\beta_t)\).
The construction starts from the actual graph at \(t\);
no terminal Gaussian graph or terminal-memory estimate is used.
The finite-cutoff admissibility threshold is uniform in \(M\).
Induction in \eqref{eq:position-source-iteration} gives
\(a_i\le\mathcal C_j^{l-i}a_l\le {\mathsf M}^{-3/4}\), because below
\({\mathsf M}^{-1/2}\) the quadratic term is at most \(2a_{i+1}\).
This also verifies the source attachment criterion: occupancy
and the spatial anchor sum add only a fixed polynomial in
\(B_t\), dominated by the small source income. The common
positive Gaussian precision is retained per physical word.

At level \(k\) every scalar, regular and forced source sum is
therefore at most a fixed constant per cell at radius \(\alpha\).
The interval margin in \eqref{eq:position-comparison-window}
keeps their inspected neighbourhoods inside the same cells;
the exact large-read and failed-sector terms remain in the
convergent source families. The complete coefficient norm is
bounded by \(e^{\Phi_k+CV}\), where \(V\) is both the number
of level-\(k\) sites and the number of replicated cells. The
absolute unmarked forced gas costs at most \(e^{Vw_k}\),
included in \(C\), not assumed to equal one.
Proposition~\ref{prop:position-sharp-positive-ratio} now gives
\[
 \left|\left\langle\prod_{C\in\mathcal T_k}(R_I)_C
                           \right\rangle_{n,M}\right|
       \le\exp\{Vq_t\}\alpha^{-|I|V},\qquad
                            q_t\le\log {\mathsf M}/32.
\]
Indeed \(4q_G\log g^{-1}/V\le1\) on late intervals, uniformly
in \(M\), since \(V\ge {\varepsilon}_k^{-4}\). The other terms are the
entry price \eqref{eq:position-window-prices}.
Taking the cell root gives the following result.

\begin{proposition}[Relative Gaussian test comparison]
\label{prop:position-relative-gaussian-tests}
For each fixed finite free test family there are intervals
\(k_\nu<l_\nu<t_\nu\), radii \(r_\nu\downarrow0\), rounded
positions \(y^{\nu,i}\to y^i\), and real common lifts
\(b_{j,n,M}^{\nu}\), bounded in \(n,M\) for each fixed \(\nu\),
such that on a full cutoff tail uniform in volume,
\begin{equation}
                         \|R_I\|_{k_\nu}\le\gamma_\nu,
                   \qquad \gamma_\nu=P_{t_\nu}^{-1/32}\longrightarrow0,
                          \quad I\ne\varnothing.
 \label{eq:position-relative-gaussian-tests}
\end{equation}
Here the row points are \(r_\nu y^{\nu,i}\), the free moments
are read at \(y^{\nu,i}\), and the errors are exactly
\eqref{eq:position-relative-pair-errors}. The unit lift is \(u_e\).
All masks, shared Gaussian laws, incoming activities and scalar
normalisations are included in this estimate.
\end{proposition}

The scalar and covariance calculations, followed by the source
induction, prove this proposition before any continuum rank or
position derivative assertion. In particular it can be used to
obtain those assertions without a circular appeal to them.
% !TeX root = 4DYM_notes.tex
\section{Uniform position derivative estimates}
\label{sec:position-derivative-bounds}

Local matching gives an admissible coefficient, not a bounded
physical coefficient. This distinction matters because the
renormalisation matrices need not have uniformly bounded inverses.
We now bound the component visible to exterior correlations by a
finite family of Gaussian tests. Exact equation directions supply
the complementary kernel. Averaging the variable test slot then
converts its differences into endpoint values, avoiding any
assumption about the derivative we are trying to bound.

This follows the finite-cutoff detection argument of
\cite[Lemma~14.27, PDF pp.~152--153]{ArmstrongDraft}.
For the undifferentiated displacement estimate we give a trace
H\"older argument which keeps the partition trace unchanged.
In position-difference operators at terminal cutoff \(n\),
we write \(\varepsilon=\varepsilon_n\).

\subsection{Displacing an undifferentiated correlation test}

A configuration with physical period \(M\) is
\emph{coordinate separated} with margin \(s>0\)
if
\[
 \min_{i\ne j}\min_{a=1,\ldots,4}\min_{q\in M\Z}
                         |x_i^a-x_j^a+q|\ge s.
\]
For one slot the minimum is \(+\infty\). This is stronger than
separation of points. It is used only for the finite test rows,
where it supplies empty slabs in every direction.

\begin{lemma}[Undifferentiated displacement]
\label{lem:position-test-displacement}
Suppose the insertions and their single-axis reflected conjugates
have centred roots at most \(A\) at one fixed cell level \(D\).
Choose \({\varrho}>{\varepsilon}_D\). If all coordinate hybrids of two grid
configurations have coordinate separation at least \(8{\varrho}\), and
each coordinate moves by at most \({\varrho}\), then
\begin{equation}
 |{\mathcal M}_{n,M}(\mathbf v;x/{\varepsilon}_n)-{\mathcal M}_{n,M}(\mathbf v;y/{\varepsilon}_n)|
       \le Cr A^r\frac{|x-y|_1}{{\varrho}}.
 \label{eq:position-test-displacement}
\end{equation}
Here a hybrid chooses each coordinate from either configuration.
The constant \(C\) is absolute. For bounded coefficient lists
the root hypothesis holds on a full cutoff tail uniform in volume.
\end{lemma}
\begin{proof}
We first explain the positive transfer operator for the original
Wilson measure, rather than for an intermediate clipped density.
Fix one direction as time. In temporal gauge, on the Hilbert
space of square-integrable spatial link configurations, one step
is
\[
       T={\mathsf M}_{e^{-\beta E_{\rm sp}/2}}\,P_{\rm gauge}K_\beta
                                      {\mathsf M}_{e^{-\beta E_{\rm sp}/2}}.
\]
Here \({\mathsf M}_f\) is multiplication by \(f\), \(E_{\rm sp}\)
is the spatial Wilson energy, \(P_{\rm gauge}\) averages the
spatial vertex gauges, and \(K_\beta\) is the product of
link convolution operators with kernel
\(f_\beta(UV^{-1})=\exp(c\beta\operatorname{Re}\operatorname{Tr}
(UV^{-1}))\), for the fixed positive Wilson coefficient \(c\).
This identity follows by separating spatial plaquettes into the
two adjacent half steps and integrating the temporal vertex
gauges with probability Haar measure.

The kernel is positive: expand its exponential into nonnegative
linear combinations of tensor products of the defining character
and its conjugate. Each such character is a sum of matrix
coefficient positive kernels. Convolution with \(f_\beta\)
is therefore positive. Gauge averaging is an orthogonal projection
commuting with this product convolution and with the spatial
energy. Consequently \(T\) is positive and self-adjoint.
Its continuous positive kernel on a compact configuration space
is trace class; its trace is the integral of the diagonal kernel,
as follows by positive finite-rank approximation. Rescale it to
\(0\le T\le1\). If the fine time period is \(n_t=M/{\varepsilon}_n\),
the partition function, up to the same cancelled scalar, is
\(Z=\operatorname{Tr}T^{n_t}>0\).

Put \(w={\varepsilon}_D/{\varepsilon}_n\) and \(m_t=n_t/w\), an even integer for
our tilings. An insertion contained in its centred width-\(w\)
slab defines a block operator \(A_i\) by integrating the
within-slab links. Reflection followed by conjugation replaces
it by \(A_i^*\). The alternating reflected time column has
normalised expectation
\(Z^{-1}\operatorname{Tr}|A_i|^{m_t}\). Applying the
chessboard inequality in the three remaining directions bounds
its \(m_t\)-th root by the full centred root. Thus
\begin{equation}
 \|A_i\|_{\mathcal S_{m_t}}\le Z^{1/m_t}A,
       \qquad \|B\|_{\mathcal S_q}
                       =(\operatorname{Tr}|B|^q)^{1/q}.
 \label{eq:position-block-schatten}
\end{equation}
The notation \(\mathcal S_q\) denotes this Schatten norm.
The unused part of a slab has empty evolution, so the block
identity applies also when the stencil occupies only its interior.

Order the marks around this time circle. Their slabs are disjoint.
Writing \(s_i\) for the empty gaps, their moment is the trace
of the cyclic product of \(A_iT^{s_i}\), divided by \(Z\),
and \(\sum_i s_i=n_t-rw\). A move of one block by a complex
amount \(z\), in fine sites, replaces its two neighbouring
gaps by \(s_-+z,s_+-z\). As long as their real parts remain
positive, this defines a holomorphic normalised trace. The
spectral definition of \(T^z\) is used on its positive spectrum;
its zero eigenspace contributes zero whenever the real exponent
is positive.

Trace H\"older gives the bound \(A^r\) throughout this strip.
Indeed use exponent \(m_t\) for each block and exponent
\(n_t/\operatorname{Re}s\) for each nonempty gap. Their
reciprocal exponents sum to
\[
               rw/n_t+(n_t-rw)/n_t=1.
\]
Moreover \(\|T^s\|_{\mathcal S_{n_t/s}}=Z^{s/n_t}\)
for real \(s>0\), and imaginary powers are contractions.
All the resulting powers of \(Z\) multiply to precisely \(Z\).
This proves the claimed bound without introducing a partition
function at a shorter period. Trace H\"older itself follows first
for finite-rank compressions from singular-value H\"older, and
then by convergence in the displayed Schatten norms.

The coordinate margin and \(w {\varepsilon}_n={\varepsilon}_D<{\varrho}\) leave, at every
point of the proposed real move, a complex disk of radius
\({\varrho}/{\varepsilon}_n\) inside this strip. Cauchy's formula bounds the
derivative by \(A^r{\varepsilon}_n/{\varrho}\). Integrating it over the real move
gives \(C A^r\) times its physical length divided by \({\varrho}\).
An arity-one moment is translation invariant and has zero change.
Negative moves use the reflected axis. Move the \(4r\)
coordinates in turn; the hybrid hypothesis supplies the same
margin at every step. This proves the assertion. All arguments
use the roots, not a cutoff-uniform supremum norm of the block.
The final clause is Theorem~\ref{thm:uniform-separated-insertions}
and its root estimates.
\end{proof}

\subsection{A detector before either limit}

% !TeX root = 4DYM_notes.tex
The elementary matrix argument used here will also be used for
continuum fields and scaling frames. We give its quantitative
form once. All the spaces in this paragraph have fixed bases;
their vector norms are coefficient \(\ell^1\) norms and their
matrix norms are the induced operator norms.

\begin{lemma}[Finite detection on a moving image]
\label{lem:common-finite-detection}
Let \(E\) be a finite-dimensional complex space, let
\(P:E\to E\) be a projection, and suppose
\(\dim\operatorname{ran}P\le s\). Let
\(X:\C^s\to E\), \(T:E\to\C^t\), and
\(H:\C^s\to\C^t\) be linear maps. Suppose that, for some
\(a>0\), \(0\le\delta<a\), and \(D>0\),
\begin{equation}
 \begin{split}
 |H u|_1&\ge a|u|_1,\qquad \|PX\|_{1\to1}\le D,\\
 \|TX-H\|_{1\to1}+\|T(I-P)X\|_{1\to1}&\le\delta.
 \end{split}
 \label{eq:common-detector-hypotheses}
\end{equation}
Then \(PX\) is injective, its image is
\(\operatorname{ran}P\), and
\begin{equation}
       |c|_1\le\frac{D}{a-\delta}|Tc|_1,
                     \qquad c\in\operatorname{ran}P.
 \label{eq:common-detector-bound}
\end{equation}
For a fixed domain, an injective matrix \(H\) therefore
continues to detect every vector under perturbations
\(T_\alpha\) with \(\|T_\alpha-H\|\le a/2\):
\begin{equation}
             |u|_1\le 2a^{-1}|T_\alpha u|_1.
 \label{eq:common-fixed-detector-bound}
\end{equation}
In particular this estimate applies to vectors depending on
\(\alpha\), without a prior bound on those vectors.
\end{lemma}
\begin{proof}
For every \(u\in\C^s\), the triangle inequality gives
\[
 \begin{split}
 |TPXu|_1
 &\ge |Hu|_1-| (TX-H)u|_1-|T(I-P)Xu|_1\\
 &\ge(a-\delta)|u|_1.
 \end{split}
\]
Thus \(PXu=0\) implies \(u=0\). Its image is contained
in \(\operatorname{ran}P\) and has dimension \(s\).
The assumed upper dimension bound makes the two images equal.
Write \(c=PXu\); then
\[
       |c|_1\le D|u|_1
                  \le\frac{D}{a-\delta}|TPXu|_1,
\]
which proves \eqref{eq:common-detector-bound}. If \(s=0\),
both images are zero and the conclusion is immediate.

For a fixed nonzero domain, injectivity of \(H\) makes
\(a=\min_{|u|_1=1}|Hu|_1\) positive, by compactness
of the unit sphere. Apply the same triangle inequality to
\(T_\alpha=H+(T_\alpha-H)\). It gives
\(|T_\alpha u|_1\ge(a/2)|u|_1\), which is
\eqref{eq:common-fixed-detector-bound}. This is also a
bounded-left-inverse assertion: invert on the image and extend
to a fixed finite-dimensional complement. Only the displayed
detection inequality will be needed below.
\end{proof}

For cutoff-dependent maps the lemma is applied pointwise on
the specified tails. A fixed weak trial first fixes \(H,a,D\)
and the finite rows; only then are its matrix errors made
smaller than \(a/2\). If those errors and the column bound
are uniform in the physical period, the resulting detector
constant is uniform in that period and in the terminal cutoff.
No uniformity in the unfixed weak-trial index is inferred.
The same argument applies on the ordered tails after taking
the fixed-period ultraviolet limits. All applications below
identify which of these two versions is being used.

Fix the degree \(e\), and put \(s=\dim V_e^{\rm Max}\).
Choose the finite free rows of
Corollary~\ref{cor:position-finite-free-tests}, retaining an
invertible \(s\)-by-\(s\) minor \(H_0\). In each row only
one label is variable; the other label, if present, is its fixed
basis spectator. For a sufficiently late but then fixed Gaussian
trial \(\nu\), use the common lifts
\(b_{j,n,M}^{\nu}\) from
Proposition~\ref{prop:position-relative-gaussian-tests}.
All their physical positions can be chosen inside a box smaller
than a quarter of the minimum allowed period. Their minimum
coordinate separation is therefore one positive number
\(\delta_0\), uniformly over periods.

The unaveraged matrix of actual moments on these lifts is
\(H_0+O(\gamma_\nu)+o_{n\to\infty}(1)\), uniformly in
volume. For pair rows this follows explicitly from
\eqref{eq:position-relative-pair-errors}:
\[
 \langle F_1F_2\rangle-G_{12}
  =\langle R_{12}\rangle
        +G_1\langle R_2\rangle+G_2\langle R_1\rangle.
\]
The chessboard inequality bounds each displayed residual moment
by its cell root. The free coefficients remain bounded for this
finite family. The rounded row points converge to the chosen
points and the free table is continuous there.

Project every variable lift to \(\Pi_{n,M,e}b_j^{\nu}\).
The removed coefficient is bounded, hence admissible, and belongs
to \(\ker\Pi\). Corollary~\ref{cor:position-equation-observables}
and Lemma~\ref{lem:position-small-spectators} show that this
changes the matrix by \(o(1)\), uniformly in volume. The
spectator lifts are left unchanged.

Fix an axis list \(\boldsymbol a\) of length \(h\). Choose
\(0<{\varrho}<\delta_0/32\), one fixed level \(D\) with \({\varepsilon}_D<{\varrho}\),
and then \(0<\rho<{\varrho}/[4(h+1)]\). For
\(Q_n=\lceil\rho/{\varepsilon}_n\rceil\), average the variable slot by
\begin{equation}
 \mathcal A_{\boldsymbol a,Q_n}
   =\prod_{j=1}^h\left(\frac1{Q_n}
                                \sum_{l=0}^{Q_n-1}\tau_{a_j}^{,l}\right).
 \label{eq:position-test-average}
\end{equation}
For \(h=0\) this is the identity. Denote by
\(\mathcal T_{n,M}:V_e\to\C^s\) the retained row tests with
this average applied to the variable slot. Every shift has total
physical size at most \(2h\rho<{\varrho}\), on a full cutoff tail.
All its coordinate hybrids retain the required slab margins.
The roots of the bounded lifts and their signed images are
bounded at the fixed level \(D\). Lemma~\ref{lem:position-test-displacement}
therefore changes the matrix by at most \(C_{\nu,D,h}\rho\).
The level and its root bounds were fixed before decreasing \(\rho\).

Choose first \(\nu\), then \(\rho\), then the full cutoff
tail so that the total matrix perturbation is less than
\((2\|H_0^{-1}\|)^{-1}\). Let \(X\) have the trial lifts
as columns and let \(D\) bound \(\|\Pi X\|_{1\to1}\)
uniformly in \(n,M\). Apply
Lemma~\ref{lem:common-finite-detection} with \(P=\Pi\),
\(T=\mathcal T_{n,M}\), and \(a=\|H_0^{-1}\|^{-1}\).
The projection error was included in the total perturbation
above. The lemma proves that the projected columns span the
actual image and gives \(C_{\rm det}=2D\|H_0^{-1}\|\).
Thus we have proved
\begin{lemma}[Uniform finite-cutoff detection]
\label{lem:position-finite-detector}
For the fixed degree and averaging list above there is a constant
\(C_{\rm det}<\infty\), independent of cutoff and volume,
such that on a full cutoff tail
\begin{equation}
        |c|_1\le C_{\rm det}|\mathcal T_{n,M}c|_1,
                           \qquad c\in\operatorname{ran}\Pi_{n,M,e}.
 \label{eq:position-finite-detector}
\end{equation}
\end{lemma}
In particular no rank assertion has been deduced by interchanging
a limit with a kernel. The inequality holds for the actual
finite-cutoff test matrix.

\subsection{Bounding the differentiated insertion}

\begin{proposition}[Roots of every fixed iterated difference]
\label{prop:position-difference-roots}
Fix \(d\), a bounded coefficient ball in \(V_d\), a finite
axis list \(\boldsymbol a\), and a sufficiently fine fixed
cell level \(m\). The centred roots of
\(\partial_{\boldsymbol a}^{\varepsilon}\mathcal O^{(d)}(v)\), and of all
their signed images and conjugates, are bounded on a full cutoff
tail, uniformly in volume and in that ball.
\end{proposition}
\begin{proof}
Put \(h=|\boldsymbol a|\), \(e=d+h\), and
\(F=\mathcal O^{(d)}(v)\). Local matching supplies an
admissible \(c=c_{n,M}\in V_e\) with
\(\mathcal O^{(e)}(c)-\partial_{\boldsymbol a}^{\varepsilon}F\) root-small.
Write \(c^{\rm pr}=\Pi c\), and first consider
\begin{equation}
                  z=\frac{c-c^{\rm pr}}{1+|c^{\rm pr}|_1}.
 \label{eq:position-normalised-equation}
\end{equation}
This family is admissible before \(c^{\rm pr}\) has been
bounded: the first numerator term divided by the denominator
is admissible by its original bound, and the second has norm
at most one, hence is admissible by
\eqref{eq:position-inverse-subpower}. It belongs to \(\ker\Pi\).
Its moments in the finite row tests tend to zero, uniformly in
volume, by the exact equation and root-small comparison. Averaging
does not change this conclusion, because all placements remain
in one fixed separated compact. Thus, on a full tail,
\begin{equation}
                  C_{\rm det}|\mathcal T_{n,M}z|_1\le\tfrac12.
 \label{eq:position-equation-absorption}
\end{equation}

The tests of \(c\) itself are bounded without a derivative
estimate. The exact averaging identity
\[
 \left(\frac1Q\sum_{l=0}^{Q-1}\tau_a^l\right)\partial_a^{\varepsilon} F
                         =\frac{\tau_a^QF-F}{Q{\varepsilon}_n}
\]
commutes with all other translations. Apply it once for every
entry of \(\boldsymbol a\). Since \(Q_n{\varepsilon}_n\ge\rho\),
the averaged difference is a sum of \(2^h\) ordinary insertion
tests, with total absolute coefficient at most \((2/\rho)^h\).
The undifferentiated moment bound applies to every endpoint
placement, with the same bounded spectator lifts and fixed margin.
The matching residual contributes \(o(1)\) by its root estimate.
Consequently
\[
                    |\mathcal T_{n,M}c|_1\le C_{\rm test}
\]
uniformly in \(n,M\). No differentiated root has been used.

Detection and \eqref{eq:position-equation-absorption} now give
\[
 |c^{\rm pr}|_1
 \le C_{\rm det}|\mathcal T c|_1
                  +(1+|c^{\rm pr}|_1)C_{\rm det}|\mathcal T z|_1
 \le C_{\rm det}C_{\rm test}+\tfrac12(1+|c^{\rm pr}|_1).
\]
Hence \(|c^{\rm pr}|_1\le2C_{\rm det}C_{\rm test}+1\).
Its ordinary insertion has bounded roots at the requested fixed
cell level. Now \(c-c^{\rm pr}\) is itself admissible, so its
equation observable has zero root and its comparison residual
has vanishing root. Add this residual and the local matching
residual to the bounded primary root. The triangle inequality for
the chessboard root proves the assertion. Finite basis
recombination makes all estimates uniform on the input ball.
The finite signed and conjugated lists are included in the same
choices and cutoff tail.
\end{proof}

\begin{theorem}[Cutoff- and volume-uniform position differences]
\label{thm:uniform-position-differences}
Fix an arity \(r\), degrees \(\mathbf d\), coefficient balls,
a finite list \(\boldsymbol l\) of slot-axis pairs, and
\(\delta>0\). There are \(B<\infty\) and \(n_0\), independent
of the period, such that
\begin{equation}
 \sup_M\ \sup_{\operatorname{sep}_{M}({\varepsilon}_n\mathbf y)\ge\delta}
            |\partial_{\boldsymbol l}^{\varepsilon}{\mathcal M}_{n,M}(\mathbf v;\mathbf y)|
                                                      \le B,
                       \qquad n\ge n_0.
 \label{eq:uniform-position-differences}
\end{equation}
Only suitable periods with nonempty indicated configuration sets
need be considered. The constants may depend on the fixed group,
degrees, orders, coefficient balls, separation and RG calendar.
They do not depend on the terminal history, cutoff or volume.
\end{theorem}
\begin{proof}
Distribute the differences among their indicated slots. Their
moment is exactly the expectation of the product of these
differentiated insertions; no product differentiation is involved
inside a slot. Their stencils are enlarged by only finitely many
fine shifts. Choose a fixed tiling side sufficiently smaller than
\(\delta\), and choose its translated boundaries away from
all anchors as in Theorem~\ref{thm:uniform-separated-insertions}.
In each direction the excluded strips have total width less than
one cell length if each margin is chosen below
\({\varepsilon}_m/[4(r+1)]\). A suitable fine-grid translate therefore
exists on a full tail, uniformly in period and placement.
Distinct anchors lie in distinct cells. Proposition~\ref{prop:position-difference-roots}
at a finer fixed centred level and
Lemma~\ref{lem:insertion-interior-root} bound each occupied
interior root. The chessboard inequality bounds their joint
moment by the product. Empty cells contain one. For \(r\le1\)
the same root bound suffices; nonempty position differences of
an arity-one moment are also zero by translation invariance.
\end{proof}

Coordinate separation was needed only to construct a finite
detector. The final bound uses ordinary separation of points,
including configurations lying on coordinate hyperplanes.
% !TeX root = 4DYM_notes.tex
\section{Smoothness of the ordered weak hierarchy}
\label{sec:smooth-weak-hierarchy}

We keep exactly the hierarchy and ordered limits constructed in
Theorem~\ref{thm:ordered-weak-hierarchy}. The difference estimates
upgrade their regularity; they do not define a second hierarchy.
In this section a retained cutoff means the original terminal
level \(\nu(n)\); its spacing is \({\varepsilon}_{\nu(n)}\). To keep the
summation formulas legible we denote that spacing by \({\varepsilon}\).
Every estimate used below was proved for a full tail of original
levels, and hence also holds for this retained sequence.

The elementary distributional argument below replaces, for this
regularity step, the need to prove uniqueness of all unrefined
ultraviolet trajectories. The difference bounds themselves were
proved before either limit and do not use this simplification.

\subsection{Summation by parts on separated compacts}

For a slot-axis list \(\boldsymbol l\) of length \(h\), write
\(\partial_{\boldsymbol l}\) for the corresponding ordinary
derivatives in \(\R^{4r}\). On tests define the backward
difference \(\partial_b^{\varepsilon,-}t(x)=[t(x)-t(x-{\varepsilon}e_b)]/{\varepsilon}\),
where \(b\) now denotes the specified coordinate of this
configuration space. The superscript \(-\) indicates a backward
difference; the adjoint of the forward difference is
\(-\partial_b^{\varepsilon,-}\).

Choose \(K\Subset U\Subset\Omega_{r,M}\), with positive
distance from the torus seam and the collision set. For a test
\(t\in C_c^\infty(K)\), its finitely many required fine shifts
are in \(U\) on a full tail. Periodic summation by parts gives
\begin{equation}
 {\varepsilon}^{4r}\sum_{\mathbf y}
       (\partial_{\boldsymbol l}^{\varepsilon}{\mathcal M}_{n,M})(\mathbf y)t({\varepsilon}\mathbf y)
  =(-1)^h {\varepsilon}^{4r}\sum_{\mathbf y}{\mathcal M}_{n,M}(\mathbf y)
                   (\partial_{\boldsymbol l}^{\varepsilon,-}t)({\varepsilon}\mathbf y).
 \label{eq:smooth-discrete-sbp}
\end{equation}
For one difference this is just the change of variable
\(\mathbf y\mapsto\mathbf y-e_b\) in its first summand;
iteration proves the formula. There is no boundary term.

The fundamental theorem of calculus writes each backward
difference as the average of the corresponding derivative over
an interval of length \({\varepsilon}\). Its iterated version proves
\[
       \|\partial_{\boldsymbol l}^{\varepsilon,-}t
                             -\partial_{\boldsymbol l}t\|_\infty\to0
\]
on one fixed compact in \(U\). The undifferentiated moment
bound controls the right-hand error in
\eqref{eq:smooth-discrete-sbp}: its absolute value is at most
the displayed supremum error times a fixed moment bound and
the total volume of the fine cubes meeting that compact. That
volume is bounded independently of the mesh. The weak row
limit consequently identifies the left-hand limit with
\((\partial_{\boldsymbol l}S_M)(\mathbf v;t)\).

On the other hand Theorem~\ref{thm:uniform-position-differences}
bounds the left-hand side in absolute value by
\(B_U {\varepsilon}^{4r}\sum|t({\varepsilon}\mathbf y)|\). The latter Riemann sum
converges to \(B_U\int|t|\). We obtain
\begin{equation}
 |(\partial_{\boldsymbol l}S_M)(\mathbf v;t)|
                   \le B_{U,\boldsymbol l}\int |t(x)|\,\dd x.
 \label{eq:smooth-distribution-derivative-bound}
\end{equation}
Its constant can be chosen uniformly over large periods for a
fixed Euclidean compact, by the volume-uniform difference bound.
Passing to the already selected volume subsequence gives the
same inequality for \(\partial_{\boldsymbol l}S\). This uses
only distributional convergence tested against
\(\partial_{\boldsymbol l}t\); it does not interchange a
pointwise derivative with a limit.

\subsection{Why all these weak derivatives give a smooth function}

\begin{lemma}[A smooth representative from bounded weak derivatives]
\label{lem:all-bounded-weak-derivatives}
Let \(T\) be a distribution on an open subset \(\Omega\) of
\(\R^m\). Suppose that for every multi-index \(\alpha\) and
every \(U\Subset\Omega\) there is \(B_{U,\alpha}<\infty\)
such that
\[
                     |\partial^\alpha T(t)|
                                 \le B_{U,\alpha}\|t\|_{L^1},
                    \qquad t\in C_c^\infty(U).
\]
Then \(T\) is represented by one unique smooth function \(f\).
On smaller compacts its derivatives obey the corresponding
local bounds.
\end{lemma}
\begin{proof}
By \(L^1\)--\(L^\infty\) duality, each weak derivative has
a locally bounded representative. Choose
\(U_0\Subset U_1\Subset\Omega\), a nonnegative smooth
mollifier \(\eta\) supported in the unit ball with integral
one, and \(\eta_\delta(x)=\delta^{-m}\eta(x/\delta)\).
For \(0<\delta<\operatorname{dist}(U_0,\partial U_1)\),
\(f_\delta=T*\eta_\delta\) is smooth on \(U_0\), and
\[
 \partial^\alpha f_\delta
          =(\partial^\alpha T)*\eta_\delta,
       \qquad\sup_{U_0}|\partial^\alpha f_\delta|
                                                   \le B_{U_1,\alpha}.
\]
One further derivative bounds its variation along sufficiently
short coordinate segments. Every finite derivative list is
therefore uniformly bounded and equicontinuous on a slightly
smaller compact. Arzel\`a--Ascoli and a diagonal extraction give
local uniform limits of all these derivatives. The fundamental
theorem of calculus on coordinate segments shows that the limit
of an order-\(q+1\) derivative is the derivative of its
order-\(q\) limit. Thus the limit \(f\) is smooth.
Convolution converges to \(T\) as a distribution, so \(f\)
represents \(T\). Two continuous representatives of the same
distribution agree everywhere: if their difference is nonzero
at a point, its real part after multiplication by a phase has
one strict sign on a small ball, contradicting testing against
a nonnegative function there. This also identifies all local
limits and glues them on overlaps. The supremum bounds pass to
the limit; shrinking the neighbourhoods proves their local form.
\end{proof}

\begin{theorem}[Smoothness of the existing ordered hierarchy]
\label{thm:smooth-ordered-hierarchy}
Every row \(S_M(\mathbf v;\cdot)\) has a unique smooth
representative on its periodic noncollision domain. The ordered
infinite-volume hierarchy \(S(\mathbf v;\cdot)\) has a unique
smooth representative on \(\operatorname{Conf}_r(\R^4)\).
For every separated compact \(K\) and finite derivative list,
\begin{equation}
 \sup_{x\in K}|\partial_{\boldsymbol l}S(\mathbf v;x)|
          \le B_{K,\mathbf d,\boldsymbol l}\prod_i|v_i|_1.
 \label{eq:smooth-hierarchy-derivative-bound}
\end{equation}
The analogous row bounds hold uniformly over sufficiently large
periods. The hierarchy retains the symmetries, units, reality and
separated-word reflection positivity already proved.
\end{theorem}
\begin{proof}
Apply Lemma~\ref{lem:all-bounded-weak-derivatives} to
\eqref{eq:smooth-distribution-derivative-bound}, at each fixed
period and then to the outer distribution. Every derivative
order is available from a separate finite difference request;
no one finite profile is being asked to provide infinitely many
derivatives. Multilinearity and finite-dimensional norm bounds
give the displayed coefficient factors. Existing distributional
symmetry identities are identities of smooth functions and hence
hold pointwise. Positivity on separated smooth tests is the same
inequality for the same distributions, so is unchanged.
\end{proof}

\subsection{Local uniform limits, including all fixed differences}

It is useful to record the stronger convergence now implied by
these bounds. It will let us pass finite-dimensional coefficient
maps through the already constructed hierarchy.

For each lattice function and each fixed difference list, use
its multilinear affine interpolation on fine cubes meeting a
separated compact. This means successively linear interpolation
in the \(4r\) coordinates; its weights are nonnegative and sum
to one. The interpolant is continuous, not asserted to be smooth
across cube faces. One more difference bound and telescoping
along coordinate lattice paths give
\begin{equation}
 |G_n(x)-G_n(y)|\le C|x-y|_1+C{\varepsilon}
 \label{eq:smooth-grid-equicontinuity}
\end{equation}
for grid points in a small convex box inside the domain. Its
interpolant is in fact Lipschitz there with constant at most a
fixed multiple of \(C\): on each cube a coordinate slope is
a convex combination of the adjacent forward differences.
Use a finite cover for a general compact.

These interpolants are locally precompact in the uniform norm.
Their distributional limits are already fixed. To verify this
last assertion, partition a compact test integral into fine cubes.
The difference between the interpolant and its lower vertex is
\(O({\varepsilon})\) uniformly; replacing the test by that vertex value
also costs \(O({\varepsilon})\) times the uniform bound. The resulting
integral is the lattice smearing, up to a vanishing boundary-strip
contribution. Equation~\eqref{eq:smooth-discrete-sbp} identifies
the limit for every difference list. Every uniform subsequential
limit is therefore the unique smooth representative of the
appropriate weak derivative. Precompactness and this uniqueness
show convergence of the full retained sequence: a contrary
subsequence would have a further uniformly convergent subsequence
with the prescribed limit, a contradiction.

The same reasoning applies in volume, now to the smooth row
functions. Their next derivative bounds supply equicontinuity
and their distributional limits were fixed by the outer
extraction. We have proved
\begin{corollary}[Ordered local uniform convergence]
\label{cor:ordered-smooth-local-convergence}
At every suitable fixed period, interpolated lattice moments
and each fixed interpolated difference converge locally uniformly
along the retained mesh family to \(S_M\) and its indicated
derivative. Along the selected volume subsequence,
\(S_M\to S\) in \(C^q\) on every separated compact for every
finite \(q\). The same conclusions hold with moving grid
points converging to separated points, and with convergent
coefficient lists in their fixed finite-dimensional spaces.
\end{corollary}

Coefficient variation in the last clause is controlled by
multilinearity and the uniform coefficient-ball bounds, not by
any bound on a bare inverse renormalisation matrix. All limits
keep their original order. Full orthogonal covariance has not
been used to establish regularity.
% !TeX root = 4DYM_notes.tex
\section{Continuum field quotients and their calculus}
\label{sec:continuum-field-calculus}

Smooth kernels allow us to identify labels which have the same
response against every exterior insertion. This gives actual
finite-dimensional continuum fields before choosing their
continuous rotation action. The construction follows
\cite[Lemma~15.3, PDF pp.~164--166]{ArmstrongDraft}.
We give the limit argument for the ordered hierarchy already
constructed here; it does not replace its two ordered limits
by a diagonal limit.

\subsection{Null labels and a faithful quotient}

For each \(d\), define the complex linear subspace
\begin{equation}
 \begin{aligned}
 \mathcal N_d=\{v\in V_d:\;&S((d,v),\mathbf w;x_0,\mathbf x)=0\\
       &\text{for every exterior list and separated configuration}\}.
 \end{aligned}
 \label{eq:continuum-null-labels}
\end{equation}
Here ``exterior'' only means that the other insertion points
differ from \(x_0\) and from one another. There is no fixed
bound on their arity or degrees. Put
\(\mathcal F_d=V_d/\mathcal N_d\), with quotient map \([v]_d\).
Give it the quotient of the fixed coefficient norm:
\(\|[v]_d\|=\inf_{z\in\mathcal N_d}|v+z|_1\).
Every subspace of a finite-dimensional space is closed, so this
is a norm. The hierarchy factors through these quotients in
every slot. It is faithful there: a field with zero response
against every exterior list is zero, by definition.
The physical unit is \(1_d=[u_d]_d\), where
\(u_d={\varepsilon}_{\operatorname{ref}(d)}^d1\) is the reference unit of
Corollary~\ref{cor:local-insertion-regularity}. In particular
the unscaled polynomial \(1\in V_d\) has class
\({\varepsilon}_{\operatorname{ref}(d)}^{-d}1_d\). This distinguishes polynomial
coefficients from the unit field represented by their
normalised physical insertion.
Conjugation and signed coordinate permutations preserve the
null spaces by the previously proved symmetry identities.

The dimension of this quotient is not being assumed. It follows
from the same finite-cutoff detector used for the differences.

\begin{proposition}[The equation kernel is exactly the continuum null space]
\label{prop:continuum-null-kernel}
For each degree there is a real signed-equivariant orthogonal
projection \(\Pi_d^\infty\) on \(V_d\) such that
\begin{equation}
 \mathcal N_d=\ker\Pi_d^\infty,
       \qquad \dim\mathcal F_d=\dim V_d^{\rm Max}.
 \label{eq:continuum-field-dimension}
\end{equation}
The projection is a convenient choice of representatives; the
quotient and its null space were defined independently of it.
\end{proposition}
\begin{proof}
Write \(M_j\) for the already selected increasing volume
sequence. We choose an auxiliary sequence of retained cutoffs
\(n_j\), solely to take limits of bounded finite matrices.
For the first \(j\) degrees, basis lists and compact
neighbourhoods of the finitely many detector rows, choose \(n_j\)
so large that the local uniform lattice-to-row errors at this
fixed \(M_j\) are at most \(1/j\). This is possible by
Corollary~\ref{cor:ordered-smooth-local-convergence}; the
coefficient-ball version handles all bounded columns involved.
Include the first \(j\) arbitrary exterior basis tests on a
countable compact exhaustion, and all previously requested
cutoff tails. The row-to-continuum errors on each such compact
already tend to zero as \(j\to\infty\).

The projections \(\Pi_{n_j,M_j,d}\) are uniformly bounded.
Their real Hermitian projection identities are closed under
matrix convergence. A countable diagonal extraction therefore
gives limits \(\Pi_d^\infty\) for every degree. Their ranks
are preserved: the trace is the fixed integer
\(\dim V_d^{\rm Max}\). Include in this extraction the bounded
trial-lift columns of one fixed, sufficiently accurate,
unaveraged detector at each degree. All these are bounded
finite-dimensional sequences. The extraction changes neither
the already defined \(S_M\) nor \(S\); every selected test
still converges to that same ordered hierarchy.

If \(v\in\ker\Pi_d^\infty\), put
\(v_j=(I-\Pi_{n_j,M_j,d})v\). This is a bounded admissible
equation coefficient, and \(v_j\to v\). Its observable is an
exact exterior-zero equation section plus a root-small residual.
Thus its response against every fixed bounded spectator list
tends to zero. Pass through the chosen local uniform limits,
and then use basis recombination and exhaustion. This gives
\(v\in\mathcal N_d\).

Conversely take \(v\in\operatorname{ran}\Pi_d^\infty
\cap\mathcal N_d\), and put \(v_j=\Pi_{n_j,M_j,d}v\).
The unaveraged inequality
\eqref{eq:position-finite-detector} holds with one fixed constant.
The retained spectator lifts converge along the same auxiliary
sequence. Their finite row points converge to distinct fixed
points. Consequently every row test of \(v_j\) converges to a
continuum response of \(v\), and is zero in the limit. Detection
gives \(|v_j|_1\to0\), so \(v=0\). Every null vector is a
sum of its projected part and a vector in the already null
kernel. Its projected part is therefore zero by this argument.
This proves equality of the two kernels and the dimension formula.
\end{proof}

\subsection{Inclusion and differentiation as field maps}

Corrected inclusions as well as corrected differences have
uniformly bounded projected coefficients. For clarity, the
inclusion proof uses the unaveraged detector: its matched source
is the original bounded lower-degree insertion, so its row tests
are bounded directly. For its equation part use
\((c-\Pi c)/(1+|\Pi c|_1)\) exactly as in
\eqref{eq:position-normalised-equation}; its tests tend to zero.
Detection and absorption bound \(|\Pi c|_1\). This is the same
calculation as the difference case, with no averaging or inverse
spacing factor. Finite basis recombination bounds the resulting
linear maps in operator norm.

For each \(d\le e\), let \(I_{n,M;d,e}:V_d\to V_e\) be
one such projected corrected inclusion. For each axis \(a\),
let \(D_{n,M;d,a}:V_d\to V_{d+1}\) be a projected corrected
first difference. These maps are real, and the inclusions can
be chosen signed-equivariant. A reflected forward difference
also shifts its anchor by one fine site; exact tensor covariance
is asserted below only for the limiting derivative, where it
follows from smooth covariance. At every fixed cell size the
comparison residuals are root-small. Include these countably
many bounded matrices in the auxiliary extraction above, and
write \(I_{d,e}\) and \(D_{d,a}\) for their limits.

\begin{theorem}[Faithful continuum field calculus]
\label{thm:continuum-field-calculus}
There are injective linear inclusions
\(\jmath_{d,e}:\mathcal F_d\to\mathcal F_e\), \(d\le e\),
and linear derivatives
\(\partial_a:\mathcal F_d\to\mathcal F_{d+1}\), with
\begin{align}
 S(\jmath_{d,e}v,\mathbf w;x_0,\mathbf x)
                         &=S(v,\mathbf w;x_0,\mathbf x),\notag\\
 S(\partial_av,\mathbf w;x_0,\mathbf x)
                         &=\partial_{x_0^a}S(v,\mathbf w;x_0,\mathbf x).
 \label{eq:continuum-field-calculus}
\end{align}
The inclusions compose, commute with differentiation in the
appropriate degrees, and preserve the unit. Derivatives commute
and annihilate the unit. All maps respect conjugation and the
signed coordinate action.
\end{theorem}
\begin{proof}
The corrected inclusion identity has a root-small error, hence
a vanishing response on each separated compact. For a corrected
derivative the same is true by the difference-root estimate and
local comparison. Along the auxiliary sequence the bounded
coefficient maps converge. The lattice functions and their fixed
differences converge to the smooth ordered hierarchy on all
required compacts. We obtain, first on basis lists and then on
every list,
\[
 S((e,I_{d,e}v),\mathbf w;x)=S((d,v),\mathbf w;x),
 \quad
 S((d+1,D_{d,a}v),\mathbf w;x)
                          =\partial_{x_0^a}S((d,v),\mathbf w;x).
\]
Thus the maps send null vectors to null vectors and descend to
the asserted maps on \(\mathcal F\). The first identity also
proves injectivity. Any two possible descended maps with these
responses coincide by faithfulness. The construction is
therefore independent of the auxiliary choices as a field map.

Testing a composition of inclusions gives the same response as
testing the direct inclusion, proving their composition law.
Testing either order of differentiation gives the same ordinary
mixed derivative of a smooth function, proving commutation.
The same reasoning proves compatibility with inclusions. The
unit has constant response in its position, so its derivative
is zero; inclusion preserves its response and hence the unit.
Conjugation and signed covariance follow by differentiating the
previously proved symmetry identities and invoking faithfulness.
Explicitly, if \(s\) is a signed coordinate permutation with
matrix \((s_{ba})\), then
\[
       s(\partial_av)=\sum_b s_{ba}\partial_b(sv).
\]
This also defines precisely the tensor convention for the axis
index. No continuous rotation covariance has entered this proof.
\end{proof}

The maps are bounded because their domain and range are finite
dimensional; the finite-cutoff proof gives their uniform
prelimit bounds. Their constants may depend on the degrees,
the fixed comparison rows and the RG calendar. A further
position cutoff or volume parameter is absent from the
continuum maps.
% !TeX root = 4DYM_notes.tex
\section{Preparing full orthogonal covariance}
\label{sec:rotation-ward-preparation}

The signed coordinate group does not contain arbitrary rotations.
Smoothness of its covariant kernels does not enlarge that group
by continuity, since the signed group is finite. We now identify
the finite-dimensional rotation comparison spaces and calculate
the free angular charge. This specifies the normalisation and
the interacting Ward identity needed next. The stress and charge
method is adapted from
\cite[Lemma~16.1, PDF pp.~169--171, and
Lemmas~16.9--16.11, PDF pp.~184--187]{ArmstrongDraft}.

\subsection{Comparison representations and the covariance to be proved}

Write \(H_4\) for the group of signed coordinate permutations.
An orthogonal matrix \(Q\) acts on a free curvature jet label by
\begin{equation}
 R_d^{\rm Max}(Q)F_{ab;c_1\cdots c_j}
  =\sum_{i,k,l_1,\ldots,l_j}
       Q_{ia}Q_{kb}\prod_{{\varrho}=1}^j Q_{l_{\varrho}c_{\varrho}}
                                      F_{ik;l_1\cdots l_j}.
 \label{eq:rotation-free-label-action}
\end{equation}
Here \(F_{ab;c_1\cdots c_j}\) denotes the order-\(j\)
curvature derivative coordinate; extend by substitution to
polynomials. The action preserves weight, the Maxwell equations
and colour invariance, so acts on \(V_d^{\rm Max}\). The tensor
formula proves the group law and continuity. The free covariance
formula gives
\begin{equation}
 S^{\rm Max}(R^{\rm Max}(Q)\mathbf v;Q\mathbf x)
                                      =S^{\rm Max}(\mathbf v;\mathbf x).
 \label{eq:rotation-free-covariance}
\end{equation}
For example the transformed two-point contraction cancels every
pair \(\sum_iQ_{ia}Q_{ib}=\delta_{ab}\); Wick products and
jet derivatives preserve this cancellation.

There are real \(H_4\)-equivariant vector-space isomorphisms
\(A_d:V_d^{\rm Max}\to\mathcal F_d\) preserving the unit.
Indeed choose the common trial lift columns in the proof of
Proposition~\ref{prop:continuum-null-kernel} by a real
\(H_4\)-equivariant section and its equivariant Gaussian
coefficient inverse. Their limits, projected by
\(\Pi_d^\infty\) and then quotiented, have an invertible
finite detector matrix and hence form a basis. Their unit column
is the exact \(u_d\). This proves the assertion.

Transporting \(R_d^{\rm Max}\) through \(A_d\) gives a
continuous real representation of \(O(4)\) on each
\(\mathcal F_d\), extending its known \(H_4\) action. This
is an algebraic comparison representation. It does not yet
assert covariance of interacting correlations, or compatibility
of these choices at different degrees. The charge construction
must select and identify the representations appropriate to
those correlations.

The desired covariance, in our label convention, is
\[
                   S(R(Q)\mathbf v;Q\mathbf x)
                                          =S(\mathbf v;\mathbf x).
\]
The two transformations occur together. In particular a tensor
label at a fixed point is not being declared a scalar.

\subsection{The free stress tensor and its local charge}

Use an orthonormal colour basis and let colons denote the
sitewise Wick ordering already used in the free table. Define
the degree-four free tensor
\begin{equation}
 T^{\rm Max}_{ab}
   =-\sum_{c,\alpha}:F_{ac}^{\alpha}F_{bc}^{\alpha}:
        +\frac14\delta_{ab}\sum_{i,j,\alpha}:(F_{ij}^{\alpha})^2:.
 \label{eq:rotation-free-stress}
\end{equation}
It is symmetric and traceless. Its divergence vanishes away
from other insertion points. To check conservation for a smooth
Maxwell two-form, expand its divergence: the terms containing
\(\sum_a\partial_aF_{ac}\) vanish, and Bianchi gives
\[
 \sum_{a,c}F_{ac}\partial_aF_{bc}
                    =\frac14\partial_b\sum_{i,j}F_{ij}^2.
\]
This cancels the derivative of the trace term. Applying the
same identity to each Wick contraction proves conservation of
the separated free correlation kernels. No product at coinciding
points of a random distribution is used in this calculation.

Fix one distinguished point \(x_0\) and a smooth function
\(\chi\) equal to one near it and zero outside a ball
containing no other insertion. For a constant vector field
\(v(y)=e_b\), or a rotational field
\(v(y)=K(y-x_0)\) with \(K^{\mathsf T}=-K\), define its free
shell charge by
\begin{equation}
 \mathcal Q_v^{\rm Max}(u,\mathbf w)
  =-\int\sum_{a,b}(\partial_a\chi)(y)v_b(y)
       S^{\rm Max}(T^{\rm Max}_{ab}(y),u(x_0),\mathbf w(\mathbf x))\,\dd y.
 \label{eq:rotation-free-shell}
\end{equation}
The integrand is smooth on the compact shell \(\supp\nabla\chi\).
Its integral is independent of the choice of \(\chi\) with
these properties: integrate conservation in the annular region
between two shells, using
\(\sum_{a,b}(\partial_av_b)T_{ab}=0\) for these vector fields.

\begin{lemma}[Normalisation of the free angular charge]
\label{lem:rotation-free-charge}
With this sign convention,
\begin{align}
 \mathcal Q_{e_b}^{\rm Max}(u,\mathbf w)
                     &=S^{\rm Max}(\partial_bu,\mathbf w),\notag\\
 \mathcal Q_{K(\cdot-x_0)}^{\rm Max}(u,\mathbf w)
                     &=S^{\rm Max}(L_d^{\rm Max}(K)u,\mathbf w),
 \label{eq:rotation-free-charge}
\end{align}
where \(L_d^{\rm Max}(K)=\left.\frac{\dd}{\dd t}\right|_{0}
R_d^{\rm Max}(e^{tK})\). The positions on the right are the
original distinguished and spectator positions.
\end{lemma}
\begin{proof}
First compute a single curvature probe. Polarise the tensor
between a smooth Maxwell wave \(f\) generated by exterior
probes and the wave \(g\) generated by \(F_{ij}(x_0)\).
Replace \(\Gamma\) in the latter by
\(\Gamma_\epsilon=\Gamma*\eta_\epsilon\), with
\(-\Delta\Gamma_\epsilon=\eta_\epsilon\). On the fixed
shell all its derivatives converge uniformly to the original
ones. The covariance formula gives
\[
 \sum_a\partial_ag_{ac}
   =\delta_{cj}\partial_i\eta_\epsilon(\cdot-x_0)
                       -\delta_{ci}\partial_j\eta_\epsilon(\cdot-x_0).
\]
Its Bianchi identity remains exact. The polarised tensor is
\[
 t_{ab}=-\sum_c(f_{ac}g_{bc}+g_{ac}f_{bc})
                         +\tfrac12\delta_{ab}\sum_{i,j}f_{ij}g_{ij}.
\]
Expanding derivatives and applying Bianchi as above leaves
\(\sum_a\partial_at_{ab}=-\sum_cf_{bc}\sum_a\partial_ag_{ac}\).
Since the contraction of \(\partial_av_b\) with this symmetric
tensor is zero, two integrations by parts give
\begin{align*}
 -\int(\partial_a\chi)v_bt_{ab}
   &= -\int\chi v_bf_{bc}
       (\delta_{cj}\partial_i\eta_\epsilon
                         -\delta_{ci}\partial_j\eta_\epsilon)\\
   &\longrightarrow
         \partial_i(v_bf_{bj})(x_0)-\partial_j(v_bf_{bi})(x_0).
\end{align*}
Repeated indices here are summed; the approximate identity is
centred at \(x_0\). No derivative of \(\chi\) survives at
that point. Bianchi rewrites the last expression as
\[
 v_b\partial_b f_{ij}
           +(\partial_iv_b)f_{bj}+(\partial_jv_b)f_{ib}.
\]
For a constant vector this is \(\partial_bf_{ij}\). For
\(v=K(y-x_0)\) it is
\(\sum_b(K_{bi}f_{bj}+K_{bj}f_{ib})(x_0)\), exactly the
generator of \eqref{eq:rotation-free-label-action}. For a
jet probe, differentiation gives in addition one term
\(\sum_bK_{bc_{\varrho}}F_{ij;c_1\cdots b\cdots c_j}\) for each
derivative index. This is also the derivative of the tensor
formula. Differentiation under the compact shell integral is
justified by uniform bounds on its separated kernels.

Finally expand a polynomial probe by Wick pairing. If neither
stress leg ends at the distinguished point, its shell is zero
by conservation through the enclosed ball. If both end there,
the shell is independent of its radius, whereas dilation of that
radius by two multiplies it by
\(2^{3+p-8-j-l}\), where \(j,l\ge0\) are the two jet
orders and \(p=0\) or \(1\) is the degree of \(v\).
This multiplier is never one, so that shell is zero. Exactly one
stress leg must therefore end there. The single-probe calculation
applies to each factor of the polynomial, with the ordinary
product rule. Summing the finitely many pairings proves the
two identities. Colour contractions commute with the calculation.
\end{proof}

These shell estimates use finite derivatives of explicit kernels
on a fixed compact separated from all other points. Their
constants depend on the degree list, shell radii, spectator
separation and group metric. There is no lattice cutoff or volume
in this free calculation. It fixes the sign and scale that an
interacting stress comparison must reproduce; conservation alone
would leave its overall scalar undetermined.

\subsection{The precise remaining angular identity}

For completeness we isolate the deterministic implication which
will close covariance once the interacting charges are controlled.
It must not be mistaken for that charge estimate.

\begin{proposition}[Angular Ward criterion for full covariance]
\label{prop:angular-ward-criterion}
Let \(R_d\) be real differentiable representations of \(O(4)\)
on \(\mathcal F_d\) extending its known \(H_4\) action, and
write \(L_d(K)=\left.\frac{\dd}{\dd t}\right|_0R_d(e^{tK})\)
for a skew-symmetric real matrix \(K\). Suppose the same smooth
hierarchy satisfies, on every separated configuration,
\begin{equation}
 \sum_i\left[
    S(v_1,\ldots,L_{d_i}(K)v_i,\ldots,v_r;\mathbf x)
           +(Kx_i)\cdot\nabla_{x_i}S(\mathbf v;\mathbf x)
          \right]=0.
 \label{eq:interacting-angular-ward-target}
\end{equation}
Then \(S\) is fully orthogonally covariant with these actions.
The inclusions and field derivatives have the corresponding
\(O(4)\) compatibility.
\end{proposition}
\begin{proof}
For fixed \(K\) differentiate
\(S(R(e^{tK})\mathbf v;e^{tK}\mathbf x)\). The chain rule
and \eqref{eq:interacting-angular-ward-target}, applied at the
transformed labels and points, give zero. Rotation preserves
all pair distances, so the entire path stays in the separated
domain. Its value is its value at zero. Every element of
\(SO(4)\) is a product of plane rotations: successive rotations
can eliminate the lower entries in each column of an orthogonal
matrix, leaving a diagonal sign matrix of determinant one, which
is itself a product of angle-\(\pi\) plane rotations. Thus
covariance holds on \(SO(4)\). A fixed coordinate reflection
belongs to \(H_4\); multiplying it by \(SO(4)\) covers the
other component of \(O(4)\). Its covariance is already known.

Test the inclusion identity
\eqref{eq:continuum-field-calculus} after rotation. For
derivatives, differentiate covariance in the insertion position
and use the tested derivative identity. Faithfulness gives
\[
 R_e(Q)\jmath_{d,e}=\jmath_{d,e}R_d(Q),\qquad
 R_{d+1}(Q)\partial_a=\sum_bQ_{ba}\partial_bR_d(Q).
\]
\end{proof}

The interacting identity
\eqref{eq:interacting-angular-ward-target} is the next estimate
to be proved, not an additional hypothesis of the construction
theorem. The free shell calculation does not establish it. Its
proof requires an interacting conserved stress tensor with the
local charge normalisation above and control of the exterior
boundary when local angular identities are combined. In
particular, off-diagonal smoothness supplies no decay at infinity
and cannot by itself discard that boundary. Neither full
orthogonal covariance nor this boundary estimate is claimed in
the present installment.
% !TeX root = 4DYM_notes.tex
\section{Stable Maxwell comparison at small radii}
\label{sec:stable-maxwell-comparison}

The finite rows used to prove position regularity do not yet
compare divergence maps or shell integrals. For those purposes
we need the same linear comparison maps for every word in a
fixed finite degree horizon, and convergence uniform over bounded
tests on separated compacts. We develop this strengthening first.
The argument follows the lattice and stable comparisons in
\cite[Lemmas~16.4--16.6, PDF pp.~176--181]{ArmstrongDraft}.
It uses no continuous rotation covariance of the interacting
hierarchy.

The input from Section~\ref{sec:gaussian-position-tests}
is Lemma~\ref{lem:actual-Taylor-window-comparison}, with
the actual Taylor bounds and local calibration of
Lemmas~\ref{lem:actual-chart-Taylor} and
\ref{lem:actual-Taylor-calibration}. It controls the
complete regular and divided forced source on each selected
interval. The finite coefficient inverse is defined from
the actual insertion at its starting level; it does not
require a Gaussian reference graph on the preceding history.

All norms on a fixed free label space are coefficient norms in
one fixed real basis. A \emph{comparison frame} at degree \(d\)
will mean a complex linear isomorphism
\(U_{k,d}:V_d^{\rm Max}\to\mathcal F_d\). The index \(k\)
labels a sequence of radii, not a new ultraviolet family or
volume extraction.

\subsection{All-arity corrected Gaussian sources}

Fix an arity bound \(r\ge2\), a degree bound \(D\), and a
compact set of configurations with positive point separation.
The source calculation of
Section~\ref{sec:gaussian-position-tests} extends to at most
\(r\) names per trial cell as follows. For \(m\le r\), let
\(F_i\) be its trial insertions, set \(G_\varnothing=1\) and
let \(G_I\) be the prescribed free moment on the subset \(I\).
In the square-zero algebra define
\begin{equation}
 G(z)=\sum_{I\subseteq\{1,\ldots,m\}}G_Iz_I,
 \qquad
 R(z)=G(z)^{-1}\prod_{i=1}^m(1+z_iF_i).
 \label{eq:stable-corrected-source}
\end{equation}
The inverse is finite: if \(G=1+H\), then
\(G^{-1}=\sum_{j=0}^m(-H)^j\). Its coefficient recursion is
\begin{equation}
 R_\varnothing=1,\qquad
 R_I=\prod_{i\in I}F_i-
                  \sum_{\varnothing\ne J\subseteq I}G_JR_{I\setminus J}.
 \label{eq:stable-residual-recursion}
\end{equation}
It subtracts the full free moment polynomial, including all
products of lower moments. In particular a noncentred pair has
the subtraction in \eqref{eq:position-relative-pair-errors}.

We spell out the changes in the norm estimates, so that this
extension does not amount to assuming independence of cells.
The coefficient norm sums over the \(2^m\) subsets and is
submultiplicative. All subset Gaussian moments and contractions
are finite polynomials of bounded degree on the chosen separated
compact. Telescoping their products gives their covariance-error
bound, with a constant depending on \(r,D\) and that compact.
Within a cell integrate all \(m\) marks as one polynomial
packet before estimating its corrected coefficients. Different
cells still have the exponential fast-covariance row bound.
Their long activity reads retain the original covering weight.
Neither summation counts the total number of replicated cells.

At a comparison level \(t\), write \({\mathsf M}=P_t\). Keep the
shorter interval \eqref{eq:position-comparison-window}, so
\({\varrho}/{\varepsilon}_l\le {\mathsf M}^{-2}\). The slow background costs \(C{\mathsf M}^{-3}\)
and the omitted low modes cost \(C{\mathsf M}^{-8}\). The total initial
error has the same form as
\eqref{eq:position-gaussian-source-initial}, with constants now
depending on \(r,D\). Give each name radius
\(\alpha_{\mathsf M}={\mathsf M}^{1/(8r)}\). Every mode is multiplied by at most
\(\alpha_{\mathsf M}^r={\mathsf M}^{1/8}\). The complete income \(a_i\), summed
over all subset modes, still satisfies
\[
 a_i\le K_{j,r,D}a_{i+1}+2{\mathsf M}^{1/2}a_{i+1}^2.
\]
Indeed the linear part contains exactly one corrected source;
all other connected terms contain at least two. Their rooted
finite-jet norm is bounded by the product of the two incomes.
The finitely many subset partitions add a fixed request constant;
the remaining powers of \(B_t\) are at most \({\mathsf M}^{1/2}\)
after entry. All Gaussian contraction powers have already been
included in the packet norm. On failures the polynomial is paid
before division by the forced square-root allowance, with the
same positive precision reserve for the complete word.

Choose the ratio floor giving
\(\mathcal C_j^{t-k+2}\le {\mathsf M}^{1/64}\) for this finite request.
Then
\({\mathsf M}^{1/8+1/64}({\mathsf M}^{-b/2}+{\mathsf M}^{-3})\le o({\mathsf M}^{-3/4})\) for
sufficiently large fixed \(b\). Coupling and localisation errors
are made smaller still by entry and a full cutoff tail. Thus
the same induction gives \(a_i\le {\mathsf M}^{-3/4}\). Choose also
\(Q_j+C\le\log {\mathsf M}/(16r)\) in the sharp normalisation bound.
Extracting a nonempty mode on every reflected cell and taking
its root gives
\begin{equation}
             \|R_I\|_{k}\le
        {\mathsf M}^{1/(16r)-|I|/(8r)}\le {\mathsf M}^{-1/(16r)}.
 \label{eq:stable-all-arity-root}
\end{equation}
Here the small cell subscript is the coarse end of that
comparison interval. The derivation retains all subset modes,
masks, incoming activities, inventories and Gaussian/Haar
normalisations. It holds on a full cutoff tail uniform in period,
with constants uniform over the prescribed compact position set.

\subsection{One coefficient map and three root normalisations}

The comparison coefficients must be common to all the rows just
described. They must also be common to three different root-cell
requests. This is why they are defined at the middle level,
rather than fitted to moments at any chosen positions.

Choose one Taylor order for the entire finite request,
including all arities and all window exponents used below.
The polynomial and calibration matrix at \(t\) are
then independent of the choice of that window exponent.
At level \(t\) first use the full cover quadratic covariance
for the internal Wick correction \(\mathsf H_{-C_t}\).
Curvature curls remove its zero-mode singularity, so each
same-site jet contraction is a finite integral with the symbol
bounds already proved. This correction depends on the linear
history and \(t\), but not on the lower integration endpoint.
Conditional integration down to \(l\) leaves the lifted slow
covariance; its contractions are included in the explicit
low-mode and slow-background errors above. This gives a common
coefficient inverse for all choices of the window exponent.

Choose real \(H_4\)-equivariant polynomial representatives of
\(V_d^{\rm Max}\). Let \(\delta_q^{\rm wt}\) multiply a
monomial of engineering weight \(w\) by \(q^w\). The
middle-level Gaussian inverse, its factor-count normalisation
and the finite marked transport determine one map
\(A_{t,n,M,d}\) such that the trial lift with exponent \(b\)
is
\begin{equation}
       \Lambda_{t,n,M,d}^{(b)}u
                       =A_{t,n,M,d}\delta_{{\mathsf M}^b}^{\rm wt}u.
 \label{eq:stable-common-lift}
\end{equation}
The inverse is supplied by
Lemma~\ref{lem:actual-Taylor-calibration} and the fixed
marked coordinate transport. Its norm may depend
on the fixed middle level, degree and history interval, but is
bounded uniformly in terminal cutoff and period. The factor
\({\varepsilon}_t^d\) is included on the physical-to-dimensionless side;
it is not an inverse-coupling loss. In particular
\(\Lambda^{(b)}1=u_d\). Linearity gives the exact relation
\[
       \Lambda^{(\alpha)}u
               =\Lambda^{(b)}\delta_{{\mathsf M}^{\alpha-b}}^{\rm wt}u.
\]
This relation concerns a common coefficient construction, not a
claim that two conditional covariances are identical.

Fix three sufficiently large exponents \(\alpha,\beta,\gamma\).
For the finite degree and arity horizon choose all of them
above \(2rD+16\), using the same Taylor order for all three.
Increasing them by a common constant preserves their
difference inequalities below and improves the inequality
containing the negative absolute exponent. In this
subsection \(\beta,\gamma\) denote comparison exponents;
running inverse couplings retain their full subscripts. For
each exponent \(b\), the comparison root cell has side at most
\begin{equation}
 {\varepsilon}_t L^2 {\mathsf M}^{b+2}B_t^{10}\le {\varepsilon}_t{\mathsf M}^{b+5}
 \label{eq:stable-root-cell-side}
\end{equation}
on a late constant-profile interval. The last inequality follows
by entry, since \(L\) is fixed and \(B_t={\mathsf M}^{o(1)}\).
Only this bound, not a window with exponent \(2b\), is needed.

Set \(r_t={\varepsilon}_t{\mathsf M}^\alpha\), \(s_t={\mathsf M}^{-\alpha}\), and
\(\sigma_t={\mathsf M}^\zeta\) with \(\zeta<0\). The normalised
grid has spacing \(s_t\), so \(r_ts_t={\varepsilon}_t\) exactly.
On the compact
\[
 K_{R,m}=\{\mathbf x:|\mathbf x|_\infty\le R,
                           \operatorname{sep}_\infty(\mathbf x)\ge R^{-1}\}
\]
Here \(|\mathbf x|_\infty=\max_{i,a}|x_i^a|\) and
\(\operatorname{sep}_\infty(\mathbf x)
=\min_{i\ne j}|x_i-x_j|_\infty\), with value \(+\infty\)
for at most one point. We obtain, for all basis words with \(m\le r\) and degrees
at most \(D\), the three bounds
\begin{align}
 |Q_t(\mathbf u;\mathbf x)|&\le\mathcal A_t^m,
       &\mathcal A_t&=C {\mathsf M}^{(\alpha-\beta)D},\notag\\
 |\nabla Q_t(\mathbf u;\mathbf x)|&\le
                           C m\mathcal B_t^m/\sigma_t,
       &\mathcal B_t&=C {\mathsf M}^{(\alpha-\gamma)D},\notag\\
 |Q_t-S^{\rm Max}(\mathbf u;\mathbf x)|&\le\epsilon_t\to0
       &&\text{on the grid with no zero coordinate differences}.
 \label{eq:stable-three-bounds}
\end{align}
Here \(Q_t\) is any continuum coefficient cluster limit of
the lifts \(\Lambda^{(\alpha)}\), read at \(r_t\mathbf x\).
The first bound holds on \(K_{R,m}\), the second on
\(K_{2R,m}\) where every coordinate separation exceeds
\(8\sigma_t\), and the third on the grid in \(K_{2R,m}\).

To prove the first two, apply
\eqref{eq:stable-all-arity-root} to singleton sources at the
\(\beta\) and \(\gamma\) normalisations. Their roots are
bounded by one plus their free singleton values. Relation
\eqref{eq:stable-common-lift} changes their norms by at most
the indicated weight powers. Their cell sides are smaller than
\(r_t\) and \(r_t\sigma_t\), respectively, provided
\begin{equation}
 \beta+6<\alpha,\qquad \gamma+5<\alpha+\zeta.
 \label{eq:stable-geometric-exponents}
\end{equation}
The interior-root/chessboard argument proves the first bound.
Lemma~\ref{lem:position-test-displacement}, followed by the
smooth row limits, proves the second, with the rescaled margin
\(\sigma_t\). The third follows by expanding
\eqref{eq:stable-residual-recursion} and applying its root
estimate. Its constants require only positive point separation;
nonzero coordinate differences are used to place the reflected
array on its aligned middle grid, not as a lower covariance
separation. All these bounds hold for every coefficient cluster
limit. For each fixed \(t\), take such limits using the
bounded-matrix auxiliary construction of
Proposition~\ref{prop:continuum-null-kernel}. It preserves the
same ordered hierarchy.

\subsection{Removing coordinate strips and making frames coherent}

Choose the exponents to satisfy, in addition,
\begin{equation}
  (\alpha-\beta)Dr+\zeta<0,\qquad
  (\alpha-\gamma)Dr-\alpha-\zeta<0.
 \label{eq:stable-strip-exponents}
\end{equation}
Here is an explicit simultaneous choice. Put \(H=\max(r,D)\),
\(v=7H(H+2)+1\), \(w=v+6\), and
\[
 \alpha=17+H+v+wH(H+2)+w,\quad
 \beta=\alpha-7,\quad\gamma=\alpha-w,\quad\zeta=-v.
\]
All three positive exponents exceed four.
The geometric inequalities have margins one and one. The last
two left sides are at most \(7H^2-v<0\) and
\(wH^2-\alpha+v<0\). Thus all four inequalities hold.
The common calendar can supply the three requested windows on
one late constant-profile interval: all constants and exponent
lists precede its ratio and entry, and their lengths are
\(O_j(\log {\mathsf M}+\log B_t)\). The gentle mesh permits
\(t\to\infty\), \(r_t\to0\).

The coordinate strips of width \(C\sigma_t\) in \(K_{R,m}\)
have volume at most \(C_{R,m}\sigma_t\), by integrating first
in the corresponding coordinate difference. Off these strips,
rounding to the \(s_t\) grid moves each coordinate by at most
\(s_t\) and keeps a fixed margin from the smaller strips.
Let \(\omega\) be the modulus of continuity of the free
table on \(K_{2R,m}\). For every smooth test supported in
\(K_{R,m}\), the three bounds give
\begin{align}
 \left|\int f(Q_t-S^{\rm Max})\right|
 \le C_{R,m}\|f\|_\infty\bigl[
       \epsilon_t+(\mathcal A_t^m+C)\sigma_t
       +\mathcal B_t^m s_t/\sigma_t+\omega(4ms_t)\bigr]
                                      \longrightarrow0.
 \label{eq:stable-tested-comparison}
\end{align}
The strip exponents make the second and third terms vanish.
No derivatives of the test enter this bound. Smooth approximation
to the phase of a locally integrable complex function then gives
local \(L^1\) convergence of the kernels. One can obtain that
approximation by first replacing the phase by
\(\overline g/(|g|+\epsilon)\), then approximating in the
finite measure \(|g|\,\dd x\) by bounded smooth functions.

Select a single sequence of these radii for increasing finite
arity, degree and compact requests and decreasing tolerances.
The maps initially obtained need not respect inclusions exactly.
Choose real \(H_4\)-equivariant left inverses \(L_d\) to the
free inclusions \(i_d:V_d^{\rm Max}\to V_{d+1}^{\rm Max}\),
using invariant inner products respecting the weight grading.
Replace them successively by
\begin{equation}
 \widetilde U_{k,0}=U_{k,0},\qquad
 \widetilde U_{k,d+1}
 =\jmath_{d,d+1}\widetilde U_{k,d}L_d
                         +U_{k,d+1}(I-i_dL_d).
 \label{eq:stable-coherent-frames}
\end{equation}
The two summands preserve the comparison, by expansion into
finitely many words and the tested inclusion identities. Reality
and signed covariance are preserved. All units follow from the
degree-zero unit. Finite free separating tested correlations make
these corrected maps eventually injective in each degree: their
matrices converge to an injective free matrix, so
\eqref{eq:common-fixed-detector-bound} applies. The equal dimensions
in \eqref{eq:continuum-field-dimension} make them isomorphisms.
We henceforth write \(U_{k,d}\) for these coherent frames.

\subsection{Derivatives of the comparison}

\begin{theorem}[Stable smooth Maxwell comparison]
\label{thm:stable-maxwell-comparison}
There is one sequence \(r_k\downarrow0\) of comparison radii
and real \(H_4\)-equivariant isomorphisms
\(U_{k,d}:V_d^{\rm Max}\to\mathcal F_d\), eventually coherent
with units and inclusions in every finite degree horizon, such
that
\begin{equation}
 S(U_{k,d_i}u_i;r_k\mathbf x)
              \longrightarrow S^{\rm Max}(u_i;\mathbf x)
       \quad\text{in }C^q(K)
 \label{eq:stable-smooth-comparison}
\end{equation}
for every fixed arity, degree list, separated compact \(K\)
and finite \(q\). The assertion remains true for convergent
free label lists. In operator norm,
\begin{equation}
 r_k U_{k,d+1}^{-1}\partial_aU_{k,d}
                                    \longrightarrow\partial_a^{\rm Max}.
 \label{eq:stable-derivative-comparison}
\end{equation}
The hierarchy and its two ordered physical limits are unchanged.
\end{theorem}
\begin{proof}
Only the derivative assertions remain. Choose finitely many free
tested responses injective on \(V_{d+1}^{\rm Max}\), with
smooth compact separated tests. Their matrices in the frames
converge by \eqref{eq:stable-tested-comparison}.
Equation~\eqref{eq:common-fixed-detector-bound} gives their
common left-inverse bound on a tail. Testing the left-hand
operator in \eqref{eq:stable-derivative-comparison} is, by
the field calculus and integration by parts, testing the
degree-\(d\) comparison kernel against the negative derivative
of the test. These values converge to those of
\(\partial_a^{\rm Max}\). The bounded left inverse proves
operator convergence, without presupposing boundedness of the
transported derivative.

Iteration and the chain rule give local \(L^1\) convergence of
all position derivatives, also for convergent label lists.
Choose a smooth cutoff equal to one near \(K\), supported in
a larger separated compact, and multiply a kernel difference
by it to obtain \(g\). In dimension \(4r\), integration by
parts bounds
\[
 |\widehat g(p)|\le C_m(1+|p|)^{-m}
                     \sum_{|\alpha|\le m}\|\partial^\alpha g\|_{L^1}.
\]
Take \(m>4r+q\). Fourier inversion and multiplication by
\(p\) through order \(q\) bound the \(C^q\) norm by a
finite sum of these \(L^1\) norms, which tend to zero.
This proves \eqref{eq:stable-smooth-comparison}. Every norm and
constant refers to a fixed finite request. All lattice inputs
were uniform in cutoff tails and period before passing through
the two physical limits. The final diagonal selects comparison
radii only.
\end{proof}
% !TeX root = 4DYM_notes.tex
\section{The interacting conserved line and local shell charges}
\label{sec:interacting-stress-charges}

We first construct a nonzero conserved tensor \(J\) on the
actual continuum field spaces. Its scale is initially unspecified.
We then prove that each of its local shell fluxes is represented
by a field of the required degree. Identifying the common scale
of these maps with positional differentiation comes after the
exterior boundary estimate. This order avoids using a global
Ward identity to define a local current.

The construction follows the conserved-line and annular arguments
in \cite[Lemmas~16.5, 16.7--16.8]{ArmstrongDraft}. We keep the complete
physical source in the annular estimate; a Ward anomaly at the
enclosed insertion is not assumed negligible.

\subsection{The finite free divergence calculation}

Use the free quadratic labels
\[
 Q_{ab,ij}=\sum_\alpha:F_{ab}^\alpha F_{ij}^\alpha:,
 \quad P_{ab}=\sum_cQ_{ac,bc},\quad Q_{\rm tot}=\sum_aP_{aa}.
\]
All have weight four. The real space of centred symmetric
\(H_4\)-tensorial degree-four arrays has the three-element basis
\begin{equation}
 \mathcal A_{ab}=\delta_{ab}P_{aa},\quad
 \mathcal B_{ab}=\delta_{ab}(Q_{\rm tot}-2P_{aa}),\quad
 \mathcal D_{ab}=(1-\delta_{ab})P_{ab}.
 \label{eq:stress-free-three-tensors}
\end{equation}
There is no sum in \(P_{aa}\) in the first two formulas.
Here is the counting argument. Simplicity makes an invariant
linear colour form zero and an invariant bilinear form a scalar
multiple of the fixed metric. Through weight four the remaining
nonconstant scalar labels are quadratic curvature labels.
For a diagonal tensor entry, coordinate signs leave two possible
index types: a curvature plane containing that coordinate and a
plane avoiding it. For an off-diagonal entry, its two free
indices must occur with odd parity and all other indices with
even parity, giving \(\sum_cF_{ac}F_{bc}\). Permutations
identify the coefficients within each type. The only constant
tensor is a multiple of \(\delta_{ab}1\), removed by centring.

Likewise the real \(H_4\)-vector space through weight five
is spanned by
\[
 V_b=\sum_c:F_{bc}\partial_bF_{bc}:,
 \qquad W_b=\sum_{c<d,\ c,d\ne b}:F_{cd}\partial_bF_{cd}:.
\]
Colour contraction is understood. To check the span, the parity
of a vector component permits a product \(F\partial F\)
with index multiplicities \(b^3c^2\) or \(b c^2d^2\).
If the derivative has index \(b\), these are exactly the two
displayed sums. Otherwise the possible contractions are reduced
by \(\sum_c\partial_cF_{cd}=0\) and by
\(2\sum_{c,d}F_{cd}\partial_cF_{bd}
=\sum_{c,d}F_{cd}\partial_bF_{cd}\).
There are no lower-weight vectors, by the same colour and parity
count. Thus this vector space has dimension at most two.

For explicit rank, test the first divergence component at
\(z=te_1\) against \(Q_{12,12}(0)\) and \(Q_{23,23}(0)\).
Divide by \(q_G/(4\pi^2)^2\) and put \(t=1\). Wick pairing
and \eqref{eq:position-free-curvature-covariance} give
\begin{equation}
 \begin{array}{c|rrr}
  &\operatorname{div}\mathcal A&\operatorname{div}\mathcal B
       &\operatorname{div}\mathcal D\\\hline
 Q_{12,12}&-256&0&128\\
 Q_{23,23}&0&-512&-128
 \end{array}
 \label{eq:stress-free-divergence-matrix}
\end{equation}
For example \(C_{1c,1c}(te_1)=-1/(\pi^2t^4)\) and
\(C_{cd,cd}(te_1)=1/(\pi^2t^4)\) for \(c,d\ne1\).
Thus the first two nonzero entries are the derivatives of
\(2q_G/(\pi^4t^8)\) and \(4q_G/(\pi^4t^8)\), respectively.
For the last column the terms \(a=3,4\) in
\(\partial_aP_{a1}\), tested against \(Q_{12,12}\), each
give \(4q_G/(\pi^4t^9)\). The terms \(a=2,3\), tested
against \(Q_{23,23}\), each give its negative. This proves
the displayed column. The two rows are independent, so divergence
is onto the vector space and the latter has dimension two.
Finally \(S^{\rm Max}(P_{12}(e_1),Q_{13,23}(0))
=-q_G/\pi^4\ne0\); this distinguishes the off-diagonal
tensor from the first two. All three tensors are independent.
The kernel is the line
\[
 \tfrac12\mathcal A-\tfrac14\mathcal B+\mathcal D=-T^{\rm Max}.
\]

\subsection{Transferring rank without assuming conservation}

For a tensor, centred means \(S_1(J_{ab})=0\) for all entries.
It is \(H_4\)-tensorial if
\(sJ_{ab}=\sum_{i,j}s_{ia}s_{jb}J_{ij}\).
Let \(\mathcal T\) be the real space of centred symmetric
such arrays in \(\mathcal F_4\), and \(\mathcal V\) the
real covariant vectors in \(\mathcal F_5\). Use fixed
finite-dimensional norms on whole arrays and vectors.

The comparison frames do not preserve one-point expectations
exactly. Correct this explicitly: transport a free centred
tensor by \(U_{k,4}\), then subtract each interacting
expectation times the unit. This is an isomorphism
\(C_k:\mathcal T^{\rm Max}\to\mathcal T\). Its inverse
transports back and subtracts the free expectation. The two
composites are identities because the frames preserve the unit;
their possible scalar error is removed by the final centring.
Transport by \(U_{k,5}\) gives an isomorphism
\(V_k:\mathcal V^{\rm Max}\to\mathcal V\).

The operator \(\operatorname{div}J\) is the vector with
entries \(\sum_a\partial_aJ_{ab}\), using the field calculus.
Equation~\eqref{eq:stable-derivative-comparison} implies
\begin{equation}
       V_k^{-1}(r_k\operatorname{div})C_k
                            \longrightarrow\operatorname{div}^{\rm Max}.
 \label{eq:stress-divergence-transfer}
\end{equation}
Derivatives kill the subtracted units, so centring introduces no
error. If \(R\) is a fixed right inverse of the free divergence,
the displayed operator composed with \(R\) is within \(1/2\)
of the identity for large \(k\), hence invertible by Neumann
series. Interacting divergence is therefore onto. Its domain and
target have dimensions three and two by the above isomorphisms.

\begin{proposition}[The interacting conserved line]
\label{prop:interacting-conserved-line}
The real space of centred symmetric \(H_4\)-tensorial arrays
\(J_{ab}\in\mathcal F_4\) satisfying
\(\sum_a\partial_aJ_{ab}=0\) is one-dimensional. One can
choose a nonzero member \(J\), one refinement of the comparison
radii and positive numbers \(\xi_k\) such that
\begin{equation}
 U_{k,4}^{-1}(\xi_k^{-1}r_k^4J_{ab})\longrightarrow T^{\rm Max}_{ab}.
 \label{eq:stress-normalised-preimage}
\end{equation}
\end{proposition}
\begin{proof}
The rank calculation proves the dimension assertion. Choose a
nonzero member \(B\) of the kernel, and put
\(\xi_k=\|U_{k,4}^{-1}(r_k^4B)\|/\|T^{\rm Max}\|\).
These numbers are positive. Their normalised preimages lie on
a fixed finite-dimensional sphere, so have a convergent
subsequence. The limit is real, symmetric and signed-tensorial.
Smooth comparison at arity one makes it free centred;
\eqref{eq:stable-derivative-comparison} makes its free
divergence zero. Its norm makes it \(T^{\rm Max}\) or
\(-T^{\rm Max}\). Choose \(J=B\) or \(-B\) once to fix
the sign. Only comparison radii have been refined.
\end{proof}

This constructs a tensor in the same hierarchy; it is not a
limit of unnormalised pointwise products of continuum fields.
Its correlations and their derivatives inherit the separated
compact bounds. Its local charge normalisation remains to be fixed.

\subsection{A quantitative annular matching estimate}

Fix \(A\in\mathcal F_d\), one component \(B\in\mathcal F_4\),
and \(\varkappa\in\{0,1\}\). Put
\(h=3+\varkappa\), \(D=d+1-\varkappa\), so \(D+h=d+4\).
Choose finitely many exterior response functionals
\[
 \mathcal L_i(v)=S(v(0),\mathbf B_i(\mathbf x_i)),
 \qquad H_i(y)=S(B(y),A(0),\mathbf B_i(\mathbf x_i)),
\]
at fixed coordinate-generic grid configurations. Their points
belong to some construction grids; their union is dense in the
separated domain. If \(g_i\in\C\) satisfy
\(\sum_i g_i\mathcal L_i=0\) on \(\mathcal F_D\), we need
\begin{equation}
 \sup_{1/4\le|z|\le2}
       \left|\sum_i g_i\rho_l^hH_i(\rho_lz)\right|\longrightarrow0
                         \quad\text{for some }\rho_l\downarrow0.
 \label{eq:stress-annular-annihilator}
\end{equation}
We now prove it, including the scale gain that permits degree
\(D\) rather than degree \(d+4\).

Write \(s_m=\log {\varepsilon}_m^{-1}\), and choose
\(\rho_m=128B_{m+1}^2{\varepsilon}_{m+1}\). Then
\(\log\rho_m=-s_m+o(s_m)\). The gentle mesh gives
\(s_{m+1}/s_m\to1\), so we can choose a finer \(j=j(m)\)
with
\[
 s_j/s_m\to t=1+\frac1{64(d+12)}.
\]
At a point \(y\in {\varepsilon}_j\Z^4\) in the enlarged annulus
\(\rho_m/16<|y|<4\rho_m\), the two insertions have separation
far larger than \({\varepsilon}_jB_j^2\). The complete two-source estimate
therefore starts at \(j\), uniformly in terminal cutoff and
period. Available marker amplitudes for the degree-\(d\)
source and the \(\rho_m^h\)-rescaled degree-four source obey
\[
 \log u_1=(-dt+o(1))s_m,\qquad
 \log u_2=(-4t+h+o(1))s_m.
\]
This follows from \eqref{eq:insertion-radius-asymptotic};
fixed threshold powers and coefficient columns contribute
\(o(s_m)\), uniformly in the annulus. The represented sources
include every incoming higher mode, not just the product of
their healthy primary functions.

Propagate to \(m+1\), without discarding either source. A
conservative estimate allows a factor \(L^4\) per mark and
step for the change of its anchor sum. In threshold-scaled
derivatives the remaining stationary pullbacks have fixed
profile Schur constants; they contribute \(e^{o(s_m)}\).
Thus the available joint amplitude can be reduced to
\(v_m=u_1u_2({\varepsilon}_j/{\varepsilon}_{m+1})^8e^{-o(s_m)}\), giving
\begin{align}
 \log v_m&=[-D-(d+12)(t-1)+o(1)]s_m,\notag\\
 E_m&\le\exp\{[(d+12)(t-1)+o(1)]s_m\}.
 \label{eq:stress-annular-source-price}
\end{align}
Here \(E_m\) is the dimensionless degree-\(D\) Taylor
remainder price after the two sources have entered one local
packet: its bound is \({\varepsilon}_{m+1}^D/v_m\) times the fixed jet
and support constants. To check the exponent, add
\(-(d+4)t+h\) and \(8(1-t)\), and use \(h=d+4-D\).
The packet diameter in level-\(m+1\) units is \(O(B_{m+1}^2)\).
Its fixed-order Taylor prices, derivative anchor changes and
covering weight have logarithm \(o(s_m)\), by the root-calendar
requirements. They are included in \(E_m\).

Choose a fixed centred cell disjoint from all spectator stencils,
and a fixed finer stopping level \(k_0\ge r(D)\). Match the
merged degree-\(D\) coordinate at its reference level, subtract
the selected primary, and apply the backward cone from
\(m+1\) to \(k_0\). The residual gains
\(({\varepsilon}_{m+1}/{\varepsilon}_{k_0})^{1/8}\). To obtain a quantitative root
bound, amplify it by \(c_m=e^{s_m/32}\) first. Its full
source price is at most
\begin{equation}
 c_mE_m e^{-(s_m-s_{k_0})/8}
       \le\exp\{[-5/64+o(1)]s_m\}\longrightarrow0,
 \label{eq:stress-annular-gain}
\end{equation}
since \(1/32+1/64-1/8=-5/64\). It fits the fixed source
allowance at \(k_0\). The replicated residual estimate then
bounds the unamplified root by \(C/c_m\), independently of
terminal cutoff and volume.

In the following physical-point notation,
\(O^{(e)}(v,y)\) abbreviates the original insertion anchored
at \(y/{\varepsilon}_n\). The estimate is for the literal difference
\(\rho_m^h O^{(4)}(b,y)O^{(d)}(a,0)-O^{(D)}(c(y),0)\)
on the same copies in every cell. To see that the subtraction
is retained, give a cell two product names and one counterterm
name. Permit the product pair together or the counterterm alone,
and forbid a mode containing a counterterm and a product name
in that cell. Its replicated coefficient is exactly the expansion
of the product of these differences, with the counterterm sign.
Merge and project only after physical propagation. The cone
therefore estimates the signed matched residual, rather than
adding absolute root bounds of its two large branches. Every
branch retains its original masks, incoming activities and
normalisation. Higher source modes use the same common Gaussian
precision reserve.

For fixed \(m,y\) all intermediate levels are fixed. Their
source prices and the inverse marked maps consequently bound
the reference coefficient \(c_{n,M}(y)\) uniformly in cutoff
tails and period, though not uniformly in \(m,y\). Take an
auxiliary coefficient cluster limit and project it to
\(\mathcal F_D\). Chessboard testing gives, simultaneously
for our finite spectator list,
\[
  |\rho_m^hH_i(y)-\mathcal L_i(c(y))|\le C_i/c_m.
\]
The same coefficient is used for all spectators. Multiplication
by \(g_i\) and summation cancels it exactly. This proves
\eqref{eq:stress-annular-annihilator} first at generic grid
points, with an enlarged annulus.

We justify passage through coordinate planes quantitatively.
For any fixed label of degree \(e\),
\eqref{eq:insertion-root-power} and the interior-root comparison
bound its roots in cells with logarithmic scale \(s_m+o(s_m)\)
by \(\exp[(e+2+o(1))s_m]\). These cells fit between annular
points separated by a fixed fraction of \(\rho_m\); their
cutoff tails are uniform in volume. A \(q\)-th position
derivative replaces the moving label by one of degree \(4+q\).
Consequently, for \(f_m(y)=\sum_i g_i\rho_m^hH_i(y)\),
\begin{equation}
 |\partial_a^qf_m(y)|\le\exp[(q+C_0+o(1))s_m],
           \quad \rho_m/16<|y|<4\rho_m,
 \label{eq:stress-annular-derivative-price}
\end{equation}
where \(C_0\) is independent of \(q\). The coefficient norm
of the differentiated field may depend on the fixed \(q\);
its logarithm divided by \(s_m\) tends to zero. Spectator
degrees contribute only to \(C_0\).

Choose \(q(t-1)>C_0\). In each coordinate choose \(q\)
distinct nonzero grid nodes within \(O_q({\varepsilon}_j)\) of a given
point. Polynomial interpolation of degree \(q-1\), with the
integral remainder, bounds its value by a fixed linear
combination of the node values plus
\(C_q {\varepsilon}_j^q\sup|\partial_a^qf_m|\).
The latter tends to zero by the choice of \(q\). Apply this
in each of the four axes. The enlarged annulus contains all
these nodes because \({\varepsilon}_j=o(\rho_m)\). The resulting node
values have all coordinates nonzero and already tend to zero.
This proves \eqref{eq:stress-annular-annihilator}, after
selecting decreasing radii. Constants depend on the fixed
fields, finite spectator tests, their margins and the calendar;
all physical estimates precede the ordered limits and are
uniform in cutoff tails and volume.

\subsection{Local charges as genuine field maps}

For a smooth window \(\chi\) equal to one near \(x_0\),
compactly supported away from the other insertions, define
\begin{equation}
 \mathcal Q_v^J(A,\mathbf B)
  =-\int\sum_{a,b}(\partial_a\chi)(y)v_b(y)
                      S(J_{ab}(y),A(x_0),\mathbf B(\mathbf x))\,\dd y,
 \label{eq:stress-interacting-shell}
\end{equation}
where \(v=e_b\) or \(v=K(y-x_0)\), \(K^{\mathsf T}=-K\).
This is an ordinary integral of a smooth kernel on a compact
separated shell. Conservation and symmetry make it independent
of the window, by integration by parts between two shells.

\begin{theorem}[Unnormalised local charges]
\label{thm:interacting-local-charges}
There are complex linear maps
\(p_b:\mathcal F_d\to\mathcal F_{d+1}\) and
\(Y_K:\mathcal F_d\to\mathcal F_d\), with real linear
dependence on \(K\), such that
\begin{equation}
 \mathcal Q_{e_b}^J(A,\mathbf B)=S(p_bA,\mathbf B),\qquad
 \mathcal Q_{K(\cdot-x_0)}^J(A,\mathbf B)=S(Y_KA,\mathbf B)
 \label{eq:stress-raw-charge-laws}
\end{equation}
for all separated configurations and windows. They are unique,
real, signed-covariant and compatible with inclusions. Moreover
\begin{equation}
 \partial_bp_a=p_a\partial_b,\qquad
 \partial_bY_K+\sum_iK_{ib}p_i=Y_K\partial_b.
 \label{eq:stress-shell-commutators}
\end{equation}
\end{theorem}
\begin{proof}
Translate the distinguished point to zero. For any finite grid
response family \(\mathcal L_i\), a linear relation annihilating
\(\mathcal F_D\) also annihilates the flux values. Indeed shrink
their common window to \(\chi(y/\rho_l)\) and change variables.
The sum of fluxes is a fixed compact integral of
\(\sum_i g_i\rho_l^{3+\varkappa}H_i(\rho_lz)\), times
\(-(\partial_a\chi)(z)v_b(z)\). It tends to zero by
\eqref{eq:stress-annular-annihilator} and is independent of
\(l\). Finite-dimensional duality therefore puts the flux
vector in the image of the response map on \(\mathcal F_D\).

Choose a finite separating set of grid responses. Such a set
exists by faithfulness, density of generic grids and continuity.
Its representing field is unique. Adjoin any further grid test;
the same duality and uniqueness give the same representing
field. Continuity of the smooth responses and compact shell
integrals extends the identity to every separated test, and
translation moves its centre to arbitrary \(x_0\).
Linearity follows by faithfulness. Conjugation, signed symmetry
and the tested inclusions give the corresponding field laws.

Keep one window fixed in a neighbourhood of a configuration and
differentiate its distinguished point. The compact shell remains
separated, so differentiation under its integral is justified.
For a constant vector field this gives the first commutator.
For \(K(y-x_0)\), differentiating the centre adds
\(-K_{ib}\) times the translation shell. Moving this term to
the other side gives the second commutator. Faithfulness turns
the tested identities into the displayed field identities.
\end{proof}

These are local field maps, with finite operator norms in every
fixed degree. At this point we have not identified \(p_b\)
with positional differentiation or \(Y_K\) with a rotation
generator. Those identifications require the global angular
identity. In particular conservation has removed inner annular
boundaries but has not removed the exterior boundary at infinity.
% !TeX root = 4DYM_notes.tex
\section{Coordinate reflection spaces and spectral bounds}
\label{sec:coordinate-spectral-bounds}

An exterior angular flux need not vanish from conservation alone.
We shall use positivity in each coordinate direction to control
its low-energy part. Continuous rotations are not assumed here.
The reflection and moment-subtraction method follows
\cite[Lemmas~17.1--17.3, PDF pp.~188--194]{ArmstrongDraft}.
We state the spectral variables and the subtraction formula
explicitly, including the possible zero-energy mass.

\subsection{Reflection Hilbert spaces and buffered operators}

Fix an axis \(d\) as time. A positive word is a finite product
of fields at distinct points with positive \(d\)-coordinate;
the empty word is allowed. For finite linear combinations put
\(\langle u,v\rangle=S(\theta_du\,v)\), where
\(\theta_d\) is coordinate reflection, conjugation of labels
and coefficients, and reversal of the word. Point reflection
positivity follows from the separated-test version by smooth
approximate identities and the smoothness theorem. Quotient by
zero norm and complete to obtain \(\mathcal H_d\).

Moving a positive word by \(te_d\), \(t\ge0\), defines
\(E_tv\). Reflection and translation give symmetry and
\(\|E_tv\|^2=\langle v,E_{2t}v\rangle\).
The globally uniform separated moment bound keeps
\(\|E_{2^kt}v\|\) bounded as \(k\to\infty\) for each fixed
word combination. Repeated Schwarz therefore gives
\(\|E_tv\|\le\|v\|\). This proves preservation of null
vectors and bounded extension. Gram continuity proves strong
continuity. Moreover \(E_t=E_{t/2}^2\) is positive, so the
spectral theorem gives \(E_t=e^{-tH_d}\), \(H_d\ge0\).
Transverse translations are commuting strongly continuous
unitaries, also commuting with \(E_t\). Write \(P_d\) for
their joint projection-valued measure on
\([0,\infty)\times\R^3\), with variables \((E,p)\).
This is the spectral theorem for commuting self-adjoint and
unitary operators; no assertion \(|p|\le E\) is used.

We also need insertions with positive buffers. A word of width
\(h\) and buffer \({\varrho}\) has mutual separation at least \({\varrho}\)
and times in \([{\varrho},h-{\varrho}]\). It defines on positive words the
insertion rule \(A_wv=w\,\tau_{he_d}v\).
For each fixed degree and coefficient bound there are \(C,Q\)
such that this rule extends to a bounded operator with
\begin{equation}
                    \|A_w\|\le C^r(r/{\varrho})^{Qr},
                         \quad 0<{\varrho}\le1.
 \label{eq:coordinate-buffered-operator}
\end{equation}
Here is the required estimate rather than an assumption of
bounded point fields. On cofinal construction grids choose a
tiling side below \({\varrho}/(100r)\) dividing the slab width.
Repeated reflected copies of the word have the same positions
modulo this tiling. Translate its boundaries away from the
finitely many positions in one copy. The interior roots of
these insertions are bounded by
\eqref{eq:insertion-root-power}. The gentle mesh allows a side
with inverse at most \(C(r/{\varrho})^2\); fixed earlier cells are
absorbed into \(C\). Chessboard bounds the repeated-word
contribution by \(D^{2j}\), \(D=C^r(r/{\varrho})^{Qr}\).
An exterior word contributes a constant \(K_v\) independent
of \(j\): group its factors in the occupied coarse cells and
bound their product roots with a finer aligned tiling. Odd
integer scale ratios preserve the reflection margins in that
tiling. The extra finite root bounds depend on the exterior
word only, so do not change \(D\). Its position modulo the
tiling is unchanged by the repeated slab width. Take ordered physical
limits first for every fixed \(j\), and then use smooth
approximation of positions and widths. Thus, with \(B\) the
reflected insertion and \(T=BA\),
\(\|T^jv\|\le K_vD^{2j}\).

The adjoint identity holds on the word span by the reflection
form. Put \(b_j=D^{-2j}\|T^jv\|\). Schwarz gives
\(b_j^2\le\|v\|b_{2j}\). Iteration and boundedness of this
sequence imply \(b_j\le\|v\|\), and hence
\(\|Av\|^2=\langle v,Tv\rangle\le D^2\|v\|^2\).
This proves \eqref{eq:coordinate-buffered-operator}. Extra
empty width is supplied by the contraction \(E_t\). Products
and adjoints agree with concatenation and reflected words.
Differentiated labels satisfy the same bound with a smaller
buffer. Twice integrating their matrix elements gives the
operator norm Taylor remainder, so buffered operators are smooth
in their insertion positions. In particular fixed positive-time
matrix elements are holomorphic in the additional time variable
\(z\), \(\operatorname{Re}z>0\), by the spectral calculus
for \(e^{-zH_d}\).

\subsection{Spectral measures and their finite moments}

For a word combination \(U\), translate it by \(ue_d\) into
positive times and call its vector \(U_u\). Define the possibly
infinite undamped positive measure by
\begin{equation}
        \dd\nu_d^U(E,p)=e^{2uE}\,\dd\langle U_u,P_d(E,p)U_u\rangle.
 \label{eq:coordinate-undamped-measure}
\end{equation}
The semigroup identity makes it independent of admissible \(u\).
It is finite on \(E\le {\mathsf M}\) for every finite \({\mathsf M}\).
Polarisation defines mixed measures \(\nu_d^{UV}\). Spectral
Schwarz gives
\begin{equation}
 \left|\int f\,\dd\nu_d^{UV}\right|^2
       \le\int g_1\,\dd\nu_d^U\int g_2\,\dd\nu_d^V,
                         \qquad |f|^2\le g_1g_2,
 \label{eq:coordinate-spectral-schwarz}
\end{equation}
for nonnegative integrable majorants \(g_1,g_2\).

For every admissible \(u>0\) and integer \(m\),
\begin{equation}
       \int e^{-2uE}(1+E+|p|)^m\,\dd\nu_d^U<\infty.
 \label{eq:coordinate-damped-moments}
\end{equation}
For a transverse component, the centred difference of order
\(2m\) in the reflected-pair kernel has the nonnegative
spectral integrand
\([2\sin(tp_j/2)/t]^{2m}e^{-2uE}\).
The integral remainder formula bounds it by the corresponding
smooth derivative on a separated compact. Fatou gives the
\(p_j^{2m}\) moment. Extra heat gives the energy moments.
If all word offsets lie perpendicular to \(d\), the same
argument at small \(u\), using the polynomial collision bound
from \eqref{eq:insertion-root-power}, gives some \(D_m\) with
\[
       \int(1+E)^{-D_m}|p_j|^{2m}\,\dd\nu_d^U<\infty.
\]
Indeed a bound \(Cu^{-D}\) can be integrated against
\(u^D\,\dd u\) on \((0,1)\) after increasing \(D\).
All constants depend on the fixed words, buffers and moment
orders. Their lattice ingredients were uniform in volume and
cutoff tails before taking the hierarchy.

\subsection{Opposite reflection signs}

Let \(Q\) be a nonempty set of coordinate axes. Suppose a
word combination has offsets zero in those axes and
\(\vartheta_jU=\varepsilon_jU\), \(\varepsilon_j=\pm1\),
with both signs occurring in \(Q\). Here \(\vartheta_j\)
is linear coordinate reflection on labels and positions.
For even nonnegative \(f\in C_c^\infty(\R)\), write
\(f_v(p)=v^{-1}f(p/v)\).

\begin{lemma}[Opposite-sign spectral bound]
\label{lem:coordinate-opposite-signs}
For \(i\in Q\), \(q<0\), fixed \(\beta_j\le0\), and
every \(\delta>0\),
\begin{equation}
 \int_{E\le1}\frac{\prod_{j\in Q\setminus\{i\}}
                         f_{j,\lambda^{\beta_j}}(p_j)}{E+\lambda^q}
                      \,\dd\nu_i^U=O_\delta(\lambda^\delta),
                         \qquad\lambda=e^s\to\infty.
 \label{eq:coordinate-opposite-signs}
\end{equation}
The estimate includes \(E=0\). The constant may depend on
the fixed profiles, exponents, word, axes and \(\delta\).
\end{lemma}
\begin{proof}
% !TeX root = 4DYM_notes.tex
\emph{Moment subtraction and its sign.}
Choose an even nonnegative \(\sigma\) of integral one,
supported in \(1\le|v|\le2\). Put \(a_j=j+2\) for
\(0\le j<m\), and let \(L_j\) be the Lagrange polynomials
with \(L_j(a_l^2)=\delta_{jl}\). For \(\rho>0\), define
the annular correction by its action on a test \(g\):
\[
 \int(\mathcal C_{m,\rho}f)(q)g(q)\,\dd q
 =\int\sigma(v)\int f(p)\sum_{j<m}
       L_j((p/(\rho v))^2)g(\rho va_j)\,\dd p\,\dd v.
\]
This is a smooth function supported in
\(2\rho\le|q|\le2(m+1)\rho\). It need not be positive.
For \(0\le l<m\), polynomial interpolation gives
\[
 \sum_{j<m}L_j(x)a_j^{2l}=x^l,
 \qquad
 \int q^{2l}\mathcal C_{m,\rho}f(q)\,\dd q
                              =\int p^{2l}f(p)\,\dd p.
\]
If \(f\) is even, the odd moments also vanish. Thus
\(f-\mathcal C_{m,\rho}f\) has all moments of orders
\(0,\ldots,2m-1\) equal to zero.

Here is also a useful bound for the signed correction. Write
\(L_j(x)=\sum_{l<m}b_{jl}x^l\) and
\(\mu_{2l}(f)=\int p^{2l}f(p)\,\dd p\). Changing
variables \(q=\rho va_j\) gives fixed smooth annular
functions \(A_{m,l}\) such that
\begin{equation}
 \mathcal C_{m,\rho}f_w(q)
  =\rho^{-1}\sum_{l<m}(w/\rho)^{2l}
                         \mu_{2l}(f)A_{m,l}(q/\rho),
 \quad
 A_{m,l}(z)=\sum_{j<m}b_{jl}a_j^{2l-1}z^{-2l}\sigma(z/a_j).
 \label{eq:coordinate-annular-majorant}
\end{equation}
The functions vanish near zero, so the displayed negative
powers are smooth on their supports. If \(w\le\rho\),
their sum is bounded in absolute value by a fixed even
nonnegative profile of width \(\rho\). If also \(\rho\ge1\),
then, for any integer \(r\ge0\),
\[
             |\mathcal C_{m,\rho}f_w(q)|
                              \le C_{m,r,f}\rho^{-2r}|q|^{2r}.
\]
Indeed on its support \(|q|\ge2\rho\), and its
supremum is at most \(C\rho^{-1}\). These bounds explain
both types of annular costs used below.

Use Fourier conventions for which
\[
 \operatorname{Lap}f(E)=\int e^{-E|t|}\widehat f_+(t)\,\dd t
                  =2E\int\frac{f(p)}{p^2+E^2}\,\dd p,
                         \qquad E>0.
\]
Interpolating the rational function \((x+t)^{-1}\) gives
\[
 \frac1{x+t}-\sum_{j<m}\frac{L_j(x)}{a_j^2+t}
       =\frac1{x+t}\prod_{j<m}\frac{a_j^2-x}{a_j^2+t}.
\]
Substitute \(x=(p/(\rho v))^2\),
\(t=(E/(\rho v))^2\), and integrate. The result is
\begin{equation}
 \mathcal R_{m,\rho}f(E)
 =2E\int\frac{f(p)}{p^2+E^2}\int\sigma(v)
       \prod_{j<m}\frac{(\rho va_j)^2-p^2}{(\rho va_j)^2+E^2}
                                      \,\dd v\,\dd p
 =\operatorname{Lap}(f-\mathcal C_{m,\rho}f)(E).
 \label{eq:coordinate-positive-residual}
\end{equation}
For \(f\ge0\) supported in \([-\rho,\rho]\), every
factor in the product is between zero and one. The residual
response is therefore nonnegative at every energy. Its value
at zero is \(2\pi f(0)\), by the approximate identity
\(E/(p^2+E^2)\); the annular correction has value zero
at the origin. This includes a possible zero-energy atom.

Choose an even bump equal to one near zero and supported in
\([-1,1]\). For \(0<w\le1\le\rho\), its rescaling obeys
\begin{equation}
 \mathcal R_{m,\rho}f_w(E)\ge\frac{c_m}{E+w}quad(0\le E\le1),
 \qquad
 \int\mathcal R_{m,\rho}f_w(E)|p_j|^{2r}\,\dd\nu_i^U
                                  \le C(w^{-1}+\rho^D).
 \label{eq:coordinate-residual-moment-cost}
\end{equation}
For the lower bound integrate over
\(0<p<wE/[2(E+w)]\). On that interval \(f_w=w^{-1}\),
the product is bounded below by a positive constant depending
on \(m\), and
\(\int_0^{wE/[2(E+w)]}E/(p^2+E^2)\,\dd p\ge cw/(E+w)\).
Continuity gives the same bound at zero.

For the upper bound, on \(E\le1\) use
\(\mathcal R_{m,\rho}f_w\le2\pi\|f_w\|_\infty\le Cw^{-1}\).
For \(E\ge1\), the product in
\eqref{eq:coordinate-positive-residual} gives
\[
 \mathcal R_{m,\rho}f_w(E)
          \le CE^{-1}\min\{1,(C_m\rho/E)^{2m}\}.
\]
The polynomial moment estimate proved above supplies an
integer \(D\) with
\(\int(1+E)^{-D}|p_j|^{2r}\,\dd\nu_i^U<\infty\).
Increase \(m\) beyond \(D/2\). Splitting at
\(E=C_m\rho\) bounds the last display by
\(C\rho^D(1+E)^{-D}\), proving the second assertion.
The same upper bound holds for a general fixed nonnegative
head profile of width \(w\le\rho/R_0\), after its fixed
support constant is included in \(R_0\). It also holds when
\(w\ge1\); its low-energy contribution is then bounded
by \(Cw^{-1}\).

\emph{Changing the temporal axis.}
For profiles in the axes of \(Q\), put
\[
 \operatorname{Resp}_i(g)
   =\int\operatorname{Lap}g_i(E)
                     \prod_{j\in Q\setminus\{i\}}g_j(p_j)\,\dd\nu_i^U.
\]
If the profiles in the two prospective temporal axes have
sufficiently many zero moments, their inverse Fourier
transforms vanish to a sufficiently high order at temporal
zero. Taylor's formula and the polynomial collision bound
then make the common position integral absolutely convergent.
Its spectral evaluation in axes \(i,j\) gives
\[
       \varepsilon_i\operatorname{Resp}_i(g)
                          =\varepsilon_j\operatorname{Resp}_j(g).
\]
For clarity, a collision in this calculation has temporal
coordinate zero in each prospective axis: all word offsets
in \(Q\) are zero. A residual with \(2m\) zero moments
contributes a factor \(O(|t|^{2m})\) there. Taking \(2m\)
above the collision exponent plus the integration dimension
justifies Fubini and the two spectral evaluations. For omitted
axes insert \(h(p/T)\), with \(h=1\) near zero, and
let \(T\to\infty\); its inverse Fourier transform is an
approximate identity. The same Taylor remainder controls
collisions, and \eqref{eq:coordinate-residual-moment-cost}
controls its spectral tails. This proves the identity for
the stated smaller set \(Q\) as well.

Let \(a_i=\mathcal C g_i\), \(a_j=\mathcal C g_j\)
be the respective annular corrections, allowing different
radii and subtraction orders, and put \(h_i=g_i-a_i\),
\(h_j=g_j-a_j\). For opposite signs the identity is
\(\operatorname{Resp}_i(h_i;h_j)=-\operatorname{Resp}_j(h_j;h_i)\).
Expanding just these two slots gives
\[
 \operatorname{Resp}_i(h_i;g_j)
  =\operatorname{Resp}_i(h_i;a_j)
    +\operatorname{Resp}_j(h_j;a_i)
    -\operatorname{Resp}_j(h_j;g_i).
\]
The final term is nonnegative, since its temporal residual
response and all transverse profiles are nonnegative. Hence
\begin{equation}
 \operatorname{Resp}_i(g_i-a_i;g_{\ne i})
 \le\operatorname{Resp}_i(g_i-a_i;g_j\mapsto a_j)
       +\operatorname{Resp}_j(g_j-a_j;g_i\mapsto a_i).
 \label{eq:coordinate-sign-transfer}
\end{equation}
When estimating the right side, bound each signed annular
correction by the nonnegative majorant in
\eqref{eq:coordinate-annular-majorant}. This uses positivity
of the residual response, without asserting positivity of
the correction itself.

\emph{The finite induction with all exponent choices.}
Choose a walk \(i_0,\ldots,i_L\) in \(Q\), starting
at the requested axis, with opposite signs at successive
visits. Require that every axis has occurred before the last
visit and that the last axis has already occurred. Such a
walk exists by inserting one axis of the opposite sign
between successive axes of the same sign. For example,
with signs \((-,-,+)\) one can use \(1,3,2,3,1\).
Put \(V_s=\{i_0,\ldots,i_{s-1}\}\).

Fix \(K\ge0\) and \(\delta>0\). The backward induction
asserts the existence of \(a_s>0\) such that the positive
residual response at head \(i_s\) is \(O(\lambda^\delta)\)
whenever
\[
 \begin{cases}
 0\le\omega_j\le a_s,&j\in V_s,\\
 -K\le\omega_j\le a_s,&j\in Q\setminus V_s,
 \end{cases}
 \quad
 0\le\eta\le a_s,\quad \omega_{i_s}\le\eta.
\]
Here \(\lambda^{\omega_j}\) is the width in axis \(j\),
and \(R_0\lambda^\eta\) is the correction radius of the
head, with fixed \(R_0\ge1\) large enough for its support.
The claim holds for all sufficiently large fixed subtraction
orders. Constants can depend on the fixed profiles,
exponents, \(R_0\), order, word and axes, but not on
\(\lambda\).

At \(s=L\), all axes belong to \(V_L\). The transverse
suprema are bounded, and
\eqref{eq:coordinate-residual-moment-cost} with \(r=0\)
bounds the response by
\(C(\lambda^{-\omega_{i_L}}+\lambda^{\eta D})\).
Choose \(a_LD<\delta\). Since \(\omega_{i_L}\ge0\),
this proves the terminal case.

Suppose the next case is known, put \(i=i_s\),
\(j=i_{s+1}\), and make the choices in the following order:
\[
 \begin{gathered}
 a'=a_{s+1},\qquad P=K|Q\setminus\{i,j\}|,\\
 2r a'>P+\max(K,\delta)+1,\\
 D\text{ is the moment exponent for }|p_j|^{2r},\qquad
 0<a_s<\min\{a',\delta/(D+1)\}.
 \end{gathered}
\]
Correct axis \(j\) at radius \(R_1\lambda^{a'}\),
with a fixed support constant \(R_1\). Its original width
exponent is at most \(a_s<a'\). Thus its correction
coefficients are bounded, and
\[
              |a_j(p_j)|\le C\lambda^{-2ra'}|p_j|^{2r}.
\]
The other transverse suprema cost at most \(C\lambda^P\).
The first term in \eqref{eq:coordinate-sign-transfer} is
therefore bounded by
\[
 C\lambda^{P-2ra'}
       (\lambda^{-\omega_i}+\lambda^{\eta D})
 \le C\lambda^{P-2ra'+\max(K,a_sD)}
                                      =O(\lambda^\delta).
\]
For the second term, replace \(a_i\) by its fixed
nonnegative annular majorant of width \(\lambda^\eta\).
Axis \(i\) has now been visited and its new exponent
\(\eta\) is nonnegative. All other exponents are unchanged,
and the next head correction radius has exponent \(a'\).
These are exactly the next induction hypotheses. They
bound this term by \(C\lambda^\delta\).

There are finitely many visits. Choose each required moment
and collision order before the corresponding subtraction
order, and then take maxima over the finite list. The
annular majorants are fixed profiles at each visit; their
profile-dependent constants are allowed by the induction.
Only finitely many responses are added, so their costs
remain \(O(\lambda^\delta)\). This completes the induction.

At its first visit choose \(K\ge\max(-q,\max_j(-\beta_j))\),
head width \(\lambda^q\), transverse widths
\(\lambda^{\beta_j}\), and correction radius \(R_0\).
The lower bound in
\eqref{eq:coordinate-residual-moment-cost} bounds the desired
low-energy resolvent integral by a constant times this
positive all-energy response. It proves
\eqref{eq:coordinate-opposite-signs}, including \(E=0\).
All spectral estimates belong to the same continuum
hierarchy obtained from the original uniform cutoff and
volume bounds. No rotation covariance or spectral gap has
entered this argument.
\end{proof}
% !TeX root = 4DYM_notes.tex
\section{Finite free Ward tests with bounded spectator degree}
\label{sec:free-ward-detection}

To identify an interacting charge by a finite matrix, its free
counterpart must have a kernel described using spectators of
bounded degree. The output of a vector map can have one larger
degree, so positivity tested only within the input degree does
not suffice by itself. We prove the needed uniform detection and
then the two Ward kernels, following
\cite[Lemmas~16.2--16.3, PDF pp.~172--175]{ArmstrongDraft}.
In the collision argument we use the exact two-leg Wick
decomposition. No assertion about arbitrary Taylor factors
preserving an improvement of singular order is needed.

\subsection{One spectator degree detecting every polynomial}

Let \(\mathcal P\) be the complex polynomial algebra on all
on-shell colour curvature jets, and \(\mathcal P^G\) its
colour-invariant part. Each polynomial uses finitely many jets.
For an integer \(b\), let \(\mathcal N_b^{\rm free}\) be
the polynomials whose free correlations vanish against every
finite exterior word with labels in \(V_b^{\rm Max}\).
This is a testing kernel, not a polynomial ideal.

\begin{lemma}[Uniform bounded-degree spectators]
\label{lem:free-uniform-spectators}
There is a finite \(b_0\), depending on the simple colour
algebra, such that
\(\mathcal N_{b_0}^{\rm free}\cap\mathcal P^G=\{0\}\).
There is no bound on the degree of the polynomial being tested;
the spectator arity is allowed to increase.
\end{lemma}
\begin{proof}
First restrict to jets of orders at most \(j\). Their
polynomial algebra has finitely many variables. It has the
Hilbert basis property: every ideal is finitely generated.
For completeness, choose a graded lexicographic monomial order.
The coordinatewise order on \(\N^r\) has no infinite
antichain: successively in each of its finitely many coordinates
an infinite sequence has a constant or strictly increasing
subsequence. Thus an upward closed set of monomials has finitely
many minimal elements. Apply this to leading monomials of an
ideal, choose polynomials attaining its minima, and cancel leading
terms successively. The well-order makes division terminate;
a nonzero remainder in the ideal would contradict minimality.
This proves the property.

Apply it to the ideal generated by all invariant polynomials of
positive ordinary degree. A finite subset of these invariants
generates the ideal; replace it by its ordinary-homogeneous
components. Average coefficients in any ideal representation
over probability Haar measure. For an invariant ordinary-
homogeneous polynomial this expresses it in terms of the chosen
invariants and invariant coefficients of strictly smaller
ordinary degree. Induction shows that finitely many invariant
generators generate the invariant algebra on these jets. Let
\(b_j\) exceed all their engineering weights.

These generators detect every nonzero invariant polynomial on
the jets through \(j\). Take its lowest nonzero ordinary part
\(P_r\), and express its reflected conjugate as a sum of
products of ordinary-homogeneous generators, each product of
ordinary degree \(r\). Put its factors at distinct nearby
spectator sites near one reflected point. Every leg at the
original site must contract to a spectator leg. Higher ordinary
parts have too many legs, and any internal spectator contraction
leaves too few legs; both contributions are zero. Let the
spectator sites coalesce. The remaining finite cross-contractions
give the strictly positive two-site pairing of \(P_r\) with
its reflected conjugate. Hence some bounded-degree spectator
word detects the original polynomial.

We remove the jet-order restriction by a collision calculation.
Put \(g=q_G\), choose \(b_0\ge\max(b_g,2g+5)\), and use
strong induction on engineering weight. Suppose an invariant
polynomial \(P\) in the testing kernel uses jets through order
\(j>g\). Insert a quadratic colour contraction with probe
orders \(k=g+1\) and \(l<g\), so its weight is at most
\(2g+5\le b_0\). Bring this insertion to zero along \(z\).
Its single contraction with the top order of \(P\) produces
the polynomial, with coefficients depending on \(z\),
\begin{equation}
 \sum_\alpha X_{l,\alpha}(q)
       \sum_\mu C_{k,j}(p,e_\mu;z)
                         \frac{\partial P}{\partial X_{j,\alpha}(e_\mu)}.
 \label{eq:free-ward-extracted-polynomial}
\end{equation}
Here \(X_{j,\alpha}(e_\mu)\) are jet coordinates in a real
colour basis and a basis of scalar on-shell probes; \(C_{k,j}\)
is their differentiated curvature covariance.

We justify extracting this term. After any fixed exterior test,
the exact two-leg Wick expansion has the finite form
\begin{equation}
 \sum_{s,t\le j}C_{k,s}(z)C_{l,t}(z)A_{st}
       +\sum_{s\le j}C_{k,s}(z)B_s(z)
       +\sum_{t\le j}C_{l,t}(z)D_t(z)+E(z).
 \label{eq:free-ward-two-leg-expansion}
\end{equation}
Tensor and colour indices are summed. The coefficients
\(A_{st}\) are independent of \(z\): both legs of the moving
quadratic have already been used at the original point. The
functions \(B_s,D_t,E\) are analytic near zero, since their
remaining moving legs connect only to fixed exterior points.
This distinction is essential.

Put \(r=4+k+j\). A contraction of order \(s\) is harmonic
homogeneous of degree \(-4-k-s\). In the single-contraction
terms only \(C_{k,j}(z)B_j(0)\) has degree \(-r\): lower
jet orders and positive Taylor terms have smaller singular
orders; the \(l\)-leg alone has order strictly below \(r\).
More singular double contractions have their own distinct radial
homogeneities and can first be removed from the identically zero
expansion. At homogeneity \(-r\), a double contraction is a
product of harmonic poles of orders \(r_1,r_2\) with
\(r_1+r_2=r\). Multiplying it by \(|z|^{2(r-2)}\) gives a
polynomial, because a harmonic pole of degree \(-r_i\) has
the form \(H_{r_i-2}(z)/|z|^{2(r_i-1)}\). No Taylor
multiplier occurs in this double-contraction term.

In contrast, a harmonic homogeneous \(K\) of degree \(-r\)
cannot have \(|z|^{2(r-2)}K\) analytic at zero unless \(K=0\).
Indeed
\(\Delta(|z|^{2a}K)=4a(a+1-r)|z|^{2a-2}K\);
successive Laplacians of an analytic extension give an analytic
extension of \(K\), whose negative homogeneity then makes it
zero. The calculation can be made explicit:
\[
 \begin{split}
 |z|^{2(r-2)}K_{r_1}(z)K_{r_2}(z)
                         &=H_{r_1-2}(z)H_{r_2-2}(z),\qquad r_1+r_2=r,\\
 \Delta^{r-2}(|z|^{2(r-2)}K)
                         &=\left(\prod_{a=1}^{r-2}4a(a+1-r)\right)K.
 \end{split}
\]
All factors in the second line are nonzero. Thus the first
line is polynomial, whereas a nonzero single harmonic pole
in the second line cannot acquire an analytic extension.
At the target radial order the zero correlation therefore
forces \(C_{k,j}(z)B_j(0)=0\). Wick removal rewrites this as
the vanishing exterior response of
\eqref{eq:free-ward-extracted-polynomial}. It belongs to
\(\mathcal N_{b_0}^{\rm free}\). Its weight is strictly
smaller than that of \(P\), so the induction hypothesis makes
it the zero polynomial.

For \(l=0,\ldots,g-1\), choose scalar probes \(q_l\).
The matrix \((X_{l,\alpha}(q_l))_{l,\alpha}\) has nonzero
polynomial determinant: assign a nonzero value in colour \(l\)
at order \(l\), normalized by \(q_l\), and zero in other
colours at that order, to make it the identity. Its adjugate
therefore shows that every colour coefficient in the extracted
identity is zero. The forms \(C_{k,j}(p,\cdot;z)\), as
\(p,z\) vary, span the dual of the order-\(j\) jets. Their
common annihilator would generate a Maxwell wave whose
order-\(k\) jets vanish everywhere; it would be a polynomial
of degree below \(k\), and its negative homogeneity makes it
zero. The Kelvin formula then makes its generating probe zero.
Thus every order-\(j\) derivative of \(P\) is zero. Descend
to order \(g\), where the bounded-jet detection already proved
applies. This closes the weight induction.
\end{proof}

\subsection{The vector Ward kernel}

Let \(\mathcal C_d\) be the real space of complex linear
tuples \(P_b:V_d^{\rm Max}\to V_{d+1}^{\rm Max}\) which
are real, kill the unit, preserve each lower filtration with
degree increase at most one, and are \(H_4\)-vector covariant.
Their Ward sums are
\[
 \mathcal E_b(P;\mathbf u;\mathbf x)
             =\sum_i S^{\rm Max}(u_1,\ldots,P_bu_i,\ldots;\mathbf x).
\]
\begin{lemma}[Finite vector tests]
\label{lem:finite-free-vector-tests}
For all sufficiently large \(d\), finitely many positional
derivatives of these sums on lines perpendicular to \(e_b\),
using labels of degree at most \(d\), define a real linear
map on \(\mathcal C_d\) with kernel
\(\R\partial^{\rm Max}\).
\end{lemma}
\begin{proof}
If every indicated jet vanishes, analyticity of free contractions
makes the Ward sums vanish near the testing lines. Every separated
configuration can be joined to such a line by a separated straight
segment: choose distinct line coordinates respecting the strict
order of the original coordinates, resolving ties in any fixed
order. Analytic continuation then gives all separated Ward sums.

First determine \(P_b\) on weight four. For a scalar curvature
probe \(p\), write \(U(p)=\sum_\alpha:X_{0,\alpha}(p)^2:\).
Colour simplicity, filtration and inversion parity make
\(P_bU(p)\) a contraction of one curvature and one first jet,
represented by a map \(A_p:\mathcal U_0\to\mathcal U_1^*\).
Here \(\mathcal U_j\) is the order-\(j\) on-shell curvature
jet space with its colour factor removed; these are not the
polynomial label spaces \(V_d\).
In the three-point Ward identity for \(U(p),U(q),U(w)\),
coalesce the first two points along \(tz\) and multiply by
\(t^5\). Only the two triangular Wick contractions survive;
the distant differentiated term has no such pole. Varying the
distant probe through the invertible order-zero covariance gives
\[
       C_{1,0}(A_pX,q;z)q(X)=C_{1,0}(A_qX,p;z)p(X)
\]
for every scalar jet value \(X\). If \(p(X)=0\), varying
\(q,z\) gives \(A_pX=0\). Hence
\(A_pX=p(X)f_bp\). Polarising the quadratic dependence on
\(p\) proves linearity of \(f_b\): use independent \(p,q\)
and values with \((p(X),q(X))=(1,0),(0,1)\) to compare
\(p+q\), then homogeneity. Nondegeneracy also gives
\(C_{1,0}(f_bp,q;z)=C_{1,0}(f_bq,p;z)\).

The signed-covariant index calculation leaves two constants.
For \(f_1F_{12}\), parity permits
\[
                  \partial_1F_{12},\quad\partial_3F_{23},
                                      \quad\partial_4F_{24};
\]
interchanging axes three and four and using Maxwell divergence
reduces this to \(\R\partial_1F_{12}\).
For \(f_1F_{23}\), the candidates
\[
                  \partial_1F_{23},\quad\partial_2F_{13},
                                      \quad\partial_3F_{12}
\]
reduce by antisymmetry and Bianchi to
\(\R\partial_1F_{23}\). They have the same constant because
\[
 C_{1,0}(\partial_1F_{12},F_{23};e_1+e_3)
                                      =-1/(2\pi^2)\ne0.
\]
Thus \(P_b=c\partial_b\) on weight
four; the factor two from differentiating the quadratic is
included in \(c\).

Put \(Q_b=P_b-c\partial_b\). Its Ward sums vanish and it
kills every free stress entry. Insert a free dilation stress
shell around just the distinguished point in this identity.
The same polarised calculation as
Lemma~\ref{lem:rotation-free-charge}, now with
\(v(y)=y-x_0\), gives the weight operator
\(\mathcal H X_{j,\alpha}=(j+2)X_{j,\alpha}\).
The extra derivative of \(v\) contracts with the trace of the
free stress and is zero. Subtract the Ward sum with label
\(\mathcal H u\). It follows that
\([\mathcal H,Q_b]u\in\mathcal N_d^{\rm free}\).
Its output may have degree \(d+1\); that is precisely why
Lemma~\ref{lem:free-uniform-spectators} is used. For \(d\ge b_0\)
it makes this commutator zero. Thus \(Q_b\) preserves each
weight. Inversion acts there by \((-1)^{\rm weight}\), while
vector covariance makes \(Q_b\) odd under inversion.
Consequently \(Q_b=-Q_b\) and \(Q_b=0\).

Conversely translations annihilate every test on
\(\partial^{\rm Max}\). The real space \(\mathcal C_d\)
is finite-dimensional. Select a test nonzero on the current
kernel modulo this line until its dimension no longer decreases.
This terminates and gives the claimed finite family.
\end{proof}

\subsection{The internal commutator kernel}

Let \(\mathcal B_d\), \(d\ge4\), consist of alternating
real-bilinear maps
\(\Delta:\operatorname{so}(4)^2\to\operatorname{End}_{\C}V_d^{\rm Max}\)
which are real, kill the unit, preserve each lower filtration,
and obey signed covariance for determinant-one signed matrices.
\begin{lemma}[Finite internal Ward tests]
\label{lem:finite-free-internal-tests}
Finitely many separated sums
\[
                 \sum_iS^{\rm Max}(\ldots,\Delta(K,L)u_i,\ldots;\mathbf x),
\]
with labels of degree at most \(d\), separate \(\mathcal B_d\).
\end{lemma}
\begin{proof}
Suppose all these sums vanish and fix \(Z=\Delta(K,L)\).
The one-point sum removes its possible constant part on centred
weight-four labels. Put \(b_x(u,v)=S_2^{\rm Max}(u,v;x,0)\).
The order-zero Kelvin formula gives, for equal nonzero radii,
\(b_x(u,v)=b_y(R^{\rm Max}(R_yR_x)u,v)\), with the
reflection matrices of \eqref{eq:position-free-kelvin}.
Apply the two-site Ward sum at \(x\) and \(y\) and use
nondegeneracy. It shows that \(Z\) commutes on weight four
with these rotations, and hence with every coordinate quarter
turn. Signed covariance then makes \(\Delta(K,L)\), on that
space, invariant when both arguments are conjugated by such
turns. An invariant alternating bilinear form on the six plane
generators is zero: overlapping distinct planes have opposite
signs under a determinant-one diagonal sign change, and disjoint
planes are exchanged by an even permutation. Thus \(Z\) kills
weight four, including the stress.

Insert a free angular stress shell in its Ward sum and subtract
the sum with the angularly differentiated distinguished label.
It gives \([L^{\rm Max}({\mathsf M}),Z]u\in\mathcal N_d^{\rm free}\).
This polynomial has degree at most \(d\), so strict two-site
positivity already makes it zero. Hence \(Z\) commutes with
all infinitesimal rotations and their quarter turns in every
degree. Repeat the alternating-plane argument to obtain
\(\Delta=0\). Finite-dimensional kernel reduction selects
finitely many sums with the same kernel.
\end{proof}
% !TeX root = 4DYM_notes.tex
\section{The exterior angular flux}
\label{sec:exterior-angular-flux}

We remove the exterior boundary first for insertions on a
coordinate line. Reflection signs supply stronger information
there than a general separated moment bound. The diagonal
stress component needs an additional use of conservation;
its crude spectral mass bound alone is insufficient.
This is the anisotropic-cutoff method of
\cite[Lemmas~16.9 and 17.3]{ArmstrongDraft}, with its scale
powers and the zero-energy contribution retained.
See also \cite[Unit VI.1, Lemma VI.1.e,
PDF pp.~449--450]{ArmstrongNotes}. The buffered transfer,
contact subtraction and diagonal positivity calculation are
proved explicitly below.

\subsection{The coordinate estimates used by the flux}

Fix different axes \(a,b\), a line axis \(l\), and
\(\lambda=e^s\), \(s\in\N\). Set
\[
 c_a=c_b=-100,\qquad c_l=-150\text{ if }l\notin\{a,b\},
 \qquad c_j=-300\text{ otherwise}.
\]
Choose an even nonnegative compactly supported profile \(h\),
equal to one on \([-1,1]\). For time axis \(d\), put
\[
 w_j=\lambda^{1+c_j},\quad
 H_\lambda^{(d)}(p)=\prod_{j\ne d}w_j^{-1}h(p_j/w_j),\quad
 I_d(U;q)=\int_{E\le1}\frac{H_\lambda^{(d)}(p)}{E+\lambda^q}
                                                \,\dd\nu_d^U.
\]
The letter \(c_j\) here denotes fixed anisotropy exponents,
not an insertion coefficient. For a word on \(\R e_l\),
let \(B\) be its projection to odd reflection parity in axes
\(a,b\) and even parity in the other two axes, using the
finite product of projections \((1\pm\vartheta_j)/2\).
It is a finite combination of collinear words.

\begin{lemma}[Coordinate low-energy bounds]
\label{lem:angular-coordinate-low-energy}
For the fixed word, tensor and profiles, and every \(\delta>0\),
\begin{align}
 I_d(B;-3000)&=O(\lambda^{\kappa_d}),
 &\kappa_d&=\begin{cases}1,&d=l,\\-c_l,&d\ne l,\end{cases}\notag\\
 I_d(J_{db};c_d)&=O_\delta(\lambda^\delta),&&d\ne b,\notag\\
 \int_{E\le1}\frac{EW^4}{(E+W)^4}H_\lambda^{(d)}
                       \,\dd\nu_d^{J_{dd}}
                  &=O_\delta(\lambda^\delta W^4),&&W=\lambda^{-99}.
 \label{eq:angular-coordinate-low-energy}
\end{align}
Also \(I_d(J_{dd};0)=O(\lambda^{900})\).
Constants may depend on these fixed data, the axes and
\(\delta\), but not on \(\lambda\).
\end{lemma}
\begin{proof}
For \(d\ne l\), omit the transverse \(l\)-profile and apply
Lemma~\ref{lem:coordinate-opposite-signs} on the other three
axes; both signs occur there. Restore its supremum
\(O(\lambda^{-1-c_l})\), and take \(\delta=1\). This gives
the first bound in this case. For a point field \(J_{db}\),
\(d\ne b\), the two indicated reflection signs are odd and
the others even, so the same lemma on all four axes gives the
second bound directly.

% !TeX root = 4DYM_notes.tex
For \(d=l\) the word has a nonzero temporal extent. We
give a buffered version of the sign transfer that avoids
testing an overlapping pair of words. Choose \(u>0\)
larger than their extent, and an even smooth compact cutoff
\(0\le\chi\le1\), equal to one on \([-2u,2u]\).
Take an even nonnegative compact convolution square \(f\)
whose Fourier transform \(\phi\) is nonnegative and has
\(\phi(0)=1\). One construction is to convolve an even
nonnegative bump with itself; its Fourier transform is the
square of a real Fourier transform. Put
\[
 R_v(E)=\int e^{-E|t|}(1-\chi(t))\phi(vt)\,\dd t,
                        \qquad v=\lambda^{-3000}.
\]
Both factors in its time profile are nonnegative. There are
constants depending on the fixed buffer, but not on \(v,E\),
such that
\begin{equation}
 \frac{c_u}{E+v}\le R_v(E)\quad(0\le E\le1),\qquad
 R_v(E)\le C_uv^{-1}e^{-2uE}\quad(E\ge0).
 \label{eq:angular-buffered-resolvent}
\end{equation}
For the lower bound choose \(R\) beyond the support of
\(\chi\), and integrate on
\([R,R+c/(E+v)]\). For fixed sufficiently small \(c>0\)
and then small \(v\), \(\phi(vt)\ge1/2\) there, and
\(e^{-Et}\ge e^{-R-c}\). Its length is \(c/(E+v)\).
For the upper bound use \(|t|\ge2u\) on the support and
\(\int\phi(vt)\,\dd t=v^{-1}\int\phi\). This also
proves both statements at zero energy.

In momentum variables the buffered head is \(f_v-a_v\),
where \(a_v\) is the Fourier transform of
\(\chi(t)\phi(vt)\), in the same convention. In particular
\[
                   \|a_v\|_\infty\le C_u
\]
uniformly for \(0<v\le1\), by its time \(L^1\) bound.
The buffered inverse Fourier transform vanishes on
\([-2u,2u]\), so the common position integral in the
sign-transfer identity involves only separated pairs in
axis \(l\).

The set \(Q'=\{j:j\ne l\}\) has both reflection signs
for the specified parity of \(B\). Choose \(j\in Q'\)
with sign opposite to that of \(l\). Apply the stronger
all-energy response bound proved during
Lemma~\ref{lem:coordinate-opposite-signs} on \(Q'\).
It permits a sufficiently small positive correction-radius
exponent \(a'\) for head \(j\). Correct its original
shrinking profile at radius \(R_0\lambda^{a'}\).
The exact sign-transfer expansion, with buffered head in
axis \(l\), now has two upper terms.

In the first term the transverse \(j\)-correction is at
most \(C\lambda^{-2ra'}|p_j|^{2r}\). If \(P\) is the
sum of the absolute values of the other negative width
exponents, its absolute value is at most
\[
 C\lambda^{P-2ra'}\int R_v(E)|p_j|^{2r}\,\dd\nu_l^B
 \le C_{u,r}\lambda^{P+3000-2ra'}.
\]
The last inequality uses the upper bound in
\eqref{eq:angular-buffered-resolvent} and the fixed damped
spectral moment. Choose \(r\) after \(a'\), with
\(2ra'>P+3001\), and take the subtraction order above
its moment and collision requirements. This term is
\(O(\lambda^{-1})\).

In the second term the head is the positive corrected
\(j\)-response, and the transverse \(l\)-profile is
\(a_v\). Replace it by its uniform supremum and omit
that axis. The all-energy bound on \(Q'\) gives
\(C_{u,\delta}\lambda^\delta\). The remaining dropped
term has positive buffered or corrected response and is
nonnegative, exactly as in
\eqref{eq:coordinate-sign-transfer}. Thus the original
positive buffered response is
\(O_\delta(\lambda^\delta)\). Its lower bound in
\eqref{eq:angular-buffered-resolvent} proves
\(I_l(B;-3000)=O(\lambda)\), taking \(\delta=1\).
Only the first transfer uses axis \(l\) as time; all
later transfers use axes perpendicular to the offsets of
the word. This proves the estimate with the same continuum
spectral measures and fixed buffer constants.

% !TeX root = 4DYM_notes.tex
For the diagonal component we spell out the contact
subtraction, the two conservation contractions and the
positivity calculation. Write
\(C_{ab,cd}(x)=S(J_{ab}(0)J_{cd}(x))\), \(x\ne0\).
Choose \(D_{\rm col}\) bounding the collision powers of
these finitely many pairs and their first derivatives.
Fix an even compact cutoff \(\chi=1\) near zero and
an integer \(M_0>D_{\rm col}+5\). For a product test
\(f(x)=\prod_j f_j(x_j)\), put
\[
 \Pi f_j(t)=\chi(t)\sum_{q<M_0}\frac{f_j^{(q)}(0)}{q!}t^q,
 \qquad
 R^0_{ab,cd}(f)=\int C_{ab,cd}(x)
                     \left(\prod_j f_j(x_j)-\prod_j\Pi f_j(x_j)\right)\dd x.
\]
The bracket is \(O(|x|^{M_0})\) at zero. Its product
with \(C\), and the differentiated tests needed for
divergence, is locally integrable. On general Schwartz
tests use the finite Taylor operator
\[
 \Pi_{\Box}f(x)=\prod_j\chi(x_j)
        \sum_{0\le\alpha_j<M_0}\frac{\partial^\alpha f(0)}{\alpha!}x^\alpha.
\]
It agrees with the preceding product operation. Taylor's
formula bounds \(R^0\) by finitely many Schwartz
seminorms, so this defines a tempered distribution.
Average under signed permutations and exchange of the two
tensor pairs; denote the result by \(R\). These averages
preserve its values away from zero.

\emph{Removing the first divergence jet.}
Conservation off zero implies that \(\sum_i\partial_iR_{ib,cd}\)
is supported at zero. A distribution of finite order
supported at one point is a finite sum of delta derivatives:
subtract its finite Taylor polynomial from a test to verify
this directly. Inversion makes this divergence odd, so only
odd orders occur. Put \(\chi_0(x)=\prod_j\chi(x_j)\) and
\[
           q_{ab,cd}=\sum_iR_{ib,cd}(\partial_i(x_a\chi_0)).
\]
The linear divergence jet evaluates on \(x_a\chi_0\)
as \(-q_{ab,cd}\). The divergence of
\(q_{ab,cd}\delta_0\) has the same value. Replacing
\(R_{ab,cd}\) by \(R_{ab,cd}-q_{ab,cd}\delta_0\)
therefore cancels this entire jet.

This subtraction retains the tensor symmetries. Indeed a
reflection makes \(q_{ab,cd}=0\) unless every axis index
occurs an even number of times. Its only possible types
are \(aaaa\), \(aabb\), and \(abab,abba\) with \(a\ne b\).
Coordinate permutations identify the coefficients within
each type, and symmetry in \(c,d\) identifies the last
two. Every resulting tensor is symmetric in \(a,b\), in
\(c,d\), and under exchange of those pairs.

Continue to denote the Fourier transform by \(R\), and set
\(F_{b,cd}=\sum_i p_iR_{ib,cd}\). These are polynomials
with no terms of degrees less than three. All their degrees
and coefficients are fixed by the continuum tensor and
the chosen extension, independently of \(\lambda\).

\emph{Two explicit conservation contractions.}
With all sums below over axes different from \(d\), put
\begin{align}
 \mathcal Q_0&=R_{dd,dd},\notag\\
 \mathcal Q_1&=\sum_{j,k\ne d}p_jp_kR_{kd,jd}
                               =p_d^2\mathcal Q_0+D_1,\notag\\
 \mathcal Q_2&=\sum_{j,k,r,s\ne d}p_jp_kp_rp_sR_{sk,rj}
                               =p_d^2\mathcal Q_1+D_2.
 \label{eq:angular-tensor-contractions}
\end{align}
Here multiplication by momentum means multiplication on
the distributional test. The polynomial errors are
\[
 \begin{split}
 D_1&=\sum_{j\ne d}p_jF_{d,jd}-p_dF_{d,dd},\\
 D_2&=\sum_{j,k,r\ne d}p_jp_kp_rF_{k,rj}
                         -p_d\sum_{j,k\ne d}p_jp_kF_{j,dk}.
 \end{split}
\]
For the first identity use
\(\sum_{k\ne d}p_kR_{kd,jd}=F_{d,jd}-p_dR_{dd,jd}\),
then contract in \(j\). For the second, first contract
in \(s\), exchange tensor pairs in the remaining term,
and contract in \(r\). The last term is exactly
\(p_d^2\mathcal Q_1\). Consequently every monomial in
\(D_1,D_2\) has degree at least four and six, respectively.

\emph{Fixed-annulus bounds by derivative labels.}
Fix a temporal momentum profile supported in an annulus
\(c\le|p_d|\le C\), with \(c>0\). Testing any entry
\(R_{ab,ce}\) against this profile and normalised shrinking
transverse profiles costs \(O_\delta(\lambda^\delta)\).
Here the component index \(e\) is unrelated to the time
axis \(d\). We prove the assertion using the same field
derivatives, rather than a quotient by spectral energy.

For each transverse profile \(g_j\), choose a fixed-radius
annular correction \(a_j\). For even profiles this is
\(\mathcal C_{m,\rho}g_j\). For odd profiles it is
\(p^{-1}\mathcal C_{m,\rho}(pg_j)\); its support avoids
zero, so it is smooth, and its odd moments agree with those
of \(g_j\) through order \(2m-1\). The other moments
are zero by parity. Take \(m\) above all contact and
collision orders used in this paragraph. Formula
\eqref{eq:coordinate-annular-majorant} shows that these
corrections are finite sums of fixed annular profiles with
bounded coefficients: a moment of a shrinking profile of
width \(w_j=\lambda^{\beta_j}\le1\) has a nonnegative
power of \(w_j\).

Order the three transverse axes, writing their profiles
as \(g_1,g_2,g_3\) just in the following identity. Then
\begin{equation}
 \begin{split}
 g_1g_2g_3={}&(g_1-a_1)g_2g_3
             +a_1(g_2-a_2)g_3\\
            &+a_1a_2(g_3-a_3)+a_1a_2a_3.
 \end{split}
 \label{eq:angular-three-axis-subtraction}
\end{equation}
The term with three corrections is bounded by continuity
of the fixed distribution \(R\): its tests lie in a
bounded subset of the finite span of fixed annular profiles.
It therefore costs \(O(1)\).

In each other term let \(i\ne d\) be its residual axis.
The inverse Fourier transform of \(g_i-a_i\) vanishes
to high order at \(x_i=0\); this removes the collision
and permits the spectral evaluation with axis \(i\) as
time. Let \(A,B\) be the two tensor labels. Put
\(A'=\partial_dA\), \(B'=\partial_dB\). The derivative
law and translation invariance give, off the collision,
\[
 S(A'(0)B'(x))=-\partial_d^2S(A(0)B(x)).
\]
Thus their distributional extensions obey
\(R^{A',B'}=p_d^2R^{A,B}\) up to a polynomial contact
term. Divide the fixed annular profile by \(p_d^2\).
This is a fixed smooth test; the residual in axis \(i\)
annihilates the polynomial difference. Consequently the
term is exactly the tested spectral pairing for \(A',B'\).
Since \(d\ne i\), differentiation in \(d\) is transverse
in this pairing, and equivalently
\(\dd\nu_i^{A',B'}=p_d^2\dd\nu_i^{A,B}\).

Every tensor component has either zero or two odd
reflection signs. Differentiation in the single axis
\(d\) changes that number to one or three. Both signs
therefore occur for each of \(A',B'\). The all-energy
sign-transfer bound applies to their positive self
measures for any choice of head axis \(i\).

Here is the bound on its head test. For an even residual,
the positive rational kernel in
\eqref{eq:coordinate-positive-residual} bounds its absolute
response by that of a fixed nonnegative even majorant at
the same width. For an odd profile \(f_v\), the odd
temporal kernel is controlled at low energy by
\[
 \int\frac{|p f_v(p)|}{p^2+E^2}\,\dd p
                               \le\frac{C_f}{E+v}.
\]
For \(E\le v\), change variables \(p=vq\) and use
\(|f(q)|\le C|q|\), giving \(C/v\). For \(E\ge v\),
use \(Cv/E^2\le C/E\). The odd correction was chosen
so that
\(p(f_v-a_v)=pf_v-\mathcal C_{m,\rho}(pf_v)\).
The same rational residual product therefore supplies
arbitrarily high fixed inverse powers above energy one.
Choose its subtraction order above the order of the
positive even majorant. The latter has the lower response
\(c/(E+v)\) below one and a fixed inverse-power lower
bound above one, directly from its positive product
formula. This proves absolute domination at every energy,
including zero.

Replace the other transverse factors by their even
nonnegative majorants, using their original shrinking
widths or the fixed correction width. Spectral Schwarz
bounds the mixed pairing by the geometric mean of two
positive residual responses for \(A'\) and \(B'\).
Each is \(O_\delta(\lambda^\delta)\) by the all-energy
bound in Lemma~\ref{lem:coordinate-opposite-signs}.
This bounds each of the three residual terms in
\eqref{eq:angular-three-axis-subtraction}. Together with
the all-annular term it proves the fixed-annulus estimate.
All profiles, division constants and derivative orders
were fixed first; no inverse power of a shrinking momentum
or spectral energy has been introduced.

On the transverse box \(|p_j|\le Cw_j\le CW\), each
momentum multiplier costs at most \(CW\). More precisely,
\(p_j^k w_j^{-1}h(p_j/w_j)\) is \(w_j^k\) times
a fixed normalised profile. Applying the preceding estimate
to these finitely many profiles yields fixed-annulus prices
\(O_\delta(\lambda^\delta W^{2r})\) for
\(\mathcal Q_r\), \(r=0,1,2\).

\emph{A positive difference and the diagonal gain.}
Use the annular bump \(\sigma\) above, and put
\[
 g_{v,r}(p)=v^{-1}(p/v)^{-2r}\sigma(p/v),\qquad
 c_r(t)=W^{-2r}\operatorname{Re}\mathcal Q_r(g_{tW,r}),
                          \quad t=1,2.
\]
The transverse test is always \(H_\lambda^{(d)}\).
Since \(p^2g_{v,r+1}=v^2g_{v,r}\), the contractions give
\[
 c_{r+1}(t)=t^2c_r(t)+O_\delta(\lambda^\delta W^2),
                              \qquad r=0,1.
\]
To check the errors, on this test all four momenta are
\(O(W)\), so \(D_1,D_2\) cost \(O(W^4),O(W^6)\).
Dividing by \(W^2,W^4\), respectively, leaves \(O(W^2)\).

Choose a common subtraction order \(m\) large enough
to annihilate every contact polynomial, including after
the two momentum multiplications. Define
\[
 h_r=(1-\mathcal C_{m,8})(g_{W,r}-g_{2W,r}),\qquad
 d_r=W^{-2r}\operatorname{Re}\mathcal Q_r(h_r).
\]
The zeroth moment of \(g_{W,r}-g_{2W,r}\) is zero.
Its even moment of order \(2l\), \(l\ge1\), is
\[
       (1-2^{2l})W^{2l}\int q^{2l-2r}\sigma(q)\,\dd q=O(W^2).
\]
The fixed-radius correction is therefore a finite sum of
fixed annular profiles with coefficients \(O(W^2)\).
The fixed-annulus estimates give, with
\(e_\lambda=\lambda^\delta W^2\),
\begin{equation}
                    d_r=c_r(1)-c_r(2)+O(e_\lambda).
 \label{eq:angular-correction-price}
\end{equation}

The sign of each \(d_r\) requires the whole corrected
profile. By \eqref{eq:coordinate-positive-residual},
\(L_r(E)=\operatorname{Lap}h_r(E)\) is nonnegative.
Indeed, after scaling \(p=Wq\) or \(2Wq\), its two
integrands have the same nonnegative weight
\(q^{-2r}\sigma(q)\). For fixed correction parameter
\(z\) the remaining factor is
\[
 \frac1{p^2+E^2}
       \prod_{j<m}\frac{(8za_j)^2-p^2}{(8za_j)^2+E^2}.
\]
It decreases with \(p^2\) while \(|p|\le4W\) and
\(W\) is small. The response at \(Wq\) consequently
dominates that at \(2Wq\), at every \(E>0\). At zero
both are zero because these profiles vanish near zero.

The moment subtractions annihilate the polynomials in the
contraction identities. Multiplication by temporal momentum
squared gives \(-E^2\) on the response: division of
\(p^{2r}\) by \(p^2+E^2\) leaves this term and the
vanishing lower moments. Consequently
\[
 d_r=(-1)^rW^{-2r}\int E^{2r}L_r(E)H_\lambda^{(d)}(p)
                                           \,\dd\nu_d^{J_{dd}}.
\]
All these integrals converge by the chosen subtraction
orders and spectral moments. Positivity gives
\(d_0\ge0\), \(d_1\le0\), \(d_2\ge0\).

Put \(A=c_0(1)\), \(B=c_0(2)\). The recurrence and
\eqref{eq:angular-correction-price} give
\[
 d_0=A-B+O(e_\lambda),\qquad
 d_1=A-4B+O(e_\lambda),\qquad
 d_2=A-16B+O(e_\lambda).
\]
Eliminating \(A,B\) yields
\begin{equation}
 4d_0=5d_1-d_2+O(e_\lambda),\qquad
                      0\le d_0\le C_\delta\lambda^\delta W^2.
 \label{eq:angular-positive-diagonal-gain}
\end{equation}

Finally, on \(1\le|q|\le2\),
\[
 \frac1{E^2+W^2q^2}-\frac1{E^2+4W^2q^2}
       =\frac{3W^2q^2}{(E^2+W^2q^2)(E^2+4W^2q^2)}
       \ge\frac{cW^2}{(E+W)^4}.
\]
For \(0<E\le1\), the correction product at \(2Wq\)
has a fixed positive lower bound. Its monotonicity shows
that \(L_0(E)\ge c_mEW^2/(E+W)^4\).
Both sides vanish at zero. Integrate this lower bound and
apply \eqref{eq:angular-positive-diagonal-gain}:
\[
 c_m\int_{E\le1}\frac{EW^2}{(E+W)^4}H_\lambda^{(d)}
                        \,\dd\nu_d^{J_{dd}}
                     \le d_0\le C_\delta\lambda^\delta W^2.
\]
Multiplication by \(W^2/c_m\) proves the third estimate
of the lemma. For the crude estimate use directly
\(\nu_d^{J_{dd}}(E\le1)<\infty\) and
\(\|H_\lambda^{(d)}\|_\infty\le C\lambda^{897}\),
enlarging the exponent to \(900\). This last bound,
rather than the vanishing factor \(E\), will control the
zero-energy atom in the exterior flux.

The constants depend on the fixed words, tensor, profiles,
collision and moment orders, \(\delta\), and the chosen
extension. None depends on \(\lambda\). Their underlying
lattice bounds were uniform in cutoff and volume before
the ordered continuum limits. No new physical limit or
phase has been selected in these distributional calculations.
\end{proof}

Fourier--Laplace tests occurring below satisfy, for every \(L\),
\[
 |\Psi_\lambda(E,p)|\le C_L\lambda^K e^{-\lambda^{-c_d}E}
              \prod_{j\ne d}(1+\lambda^{-c_j}|p_j|)^{-L}
\]
for a fixed finite \(K\). Their integral on
\(\{E>1\}\cup\bigcup_{j\ne d}\{|p_j|>w_j\}\) is
\(O_J(\lambda^{-J})\) for every \(J\). Above energy one
use the exponential and an admissible fixed damping. Beyond
a spatial side of the box use
\(\lambda^{-c_j}|p_j|>\lambda\), arbitrary Schwartz gain,
and \eqref{eq:coordinate-damped-moments}. Mixed measures are
handled by \eqref{eq:coordinate-spectral-schwarz}.
More explicitly, choose one admissible fixed damping \(t_0\).
On \(E>1\), the excess exponential
\(e^{-(\lambda^{-c_d}-t_0)E}\) pays any prescribed
power \(\lambda^{-J-K}\). On a spatial complement
\(|p_j|>w_j\), one factor has
\(1+\lambda^{-c_j}|p_j|>\lambda\); take its Schwartz
order larger than \(J+K\) and keep \(e^{-t_0E}\).
Integration is then bounded by the finite damped spectral
mass. Split the complement into these four regions and use
spectral Schwarz for mixed measures. This proves the stated
inverse power with constants depending on \(J,t_0\) and
the fixed profiles, uniformly in \(\lambda\).

\subsection{An anisotropic cutoff and the vanishing flux}

Define the raw angular Ward sum
\[
 \mathcal W_K(\mathbf A;\mathbf x)
  =\sum_i\left[S(\ldots,Y_KA_i,\ldots;\mathbf x)
                   +\sum_b(Kx_i)_bS(\ldots,p_bA_i,\ldots;\mathbf x)\right].
\]
\begin{theorem}[The collinear exterior boundary vanishes]
\label{thm:collinear-angular-identity}
For every separated configuration on a coordinate line,
\(\mathcal W_K=0\) for all skew-symmetric \(K\), for the
same tensor and raw local charges constructed above.
\end{theorem}
\begin{proof}
It suffices to take \(K=E_{ba}-E_{ab}\). Choose an even smooth
compact cutoff \(\eta=1\) on \([-1,1]\), and put
\(\ell(t)=t\eta(t)\). Define the compact vector field
\[
 f_{ab}(y)=\lambda^{100}\ell(\lambda^{-100}y_a)
                       \prod_{j\ne a}\ell'(\lambda^{c_j}y_j),
 \quad V_b=f_{ab},\quad V_a=-f_{ba},\quad V_j=0\ (j\ne a,b).
\]
It equals \(Ky\) near all fixed insertion points for large
\(\lambda\). Disjoint local windows and conservation give
\(\mathcal W_K=-\int\sum_{\mu,b}(\partial_\mu V_b)
S(J_{\mu b}(y),\mathbf A(\mathbf x))\,\dd y\).
The terms \(\mu=a\) in \(V_b\) and \(\mu=b\) in
\(V_a\) cancel exactly, because
\(\partial_af_{ab}=\partial_bf_{ba}\) and \(J_{ab}=J_{ba}\).

For the other terms of \(V_b\), coordinate reflections replace
the exterior word by its parity projection \(B\) above. The
temporal differentiated profile is supported at
\(|y_\mu|\ge\lambda^{-c_\mu}\), giving the Laplace factor
\(e^{-E/\lambda^{c_\mu}}\). Each transverse Fourier factor
is bounded on its box by \(C\lambda w_j^{-1}h(p_j/w_j)\).
Their product costs \(C\lambda^3H_\lambda^{(\mu)}\).
We estimate before restoring the prefactor \(\lambda^{100}\).
To check that no temporal length is missing, put
\(s=\lambda^{c_\mu}\). The differentiated profile is
\(s\ell''(st)\), with fixed time \(L^1\) norm. Its
support has \(|t|\ge s^{-1}\), and hence
\[
 \int|s\ell''(st)|e^{-E|t|}\,\dd t
                             \le\|\ell''\|_{L^1}e^{-E/s}.
\]
The three transverse bounds give exactly \(C\lambda^3\),
with \(H_\lambda^{(\mu)}\) retained as a normalised profile.

If \(\mu\ne b\), it is also different from \(a\). Put
\(s=\lambda^{c_\mu}\), \(v=\lambda^{-3000}\le s\).
The inequality
\[
                       e^{-E/s}\sqrt{(E+v)(E+s)}\le Cs
\]
and spectral Schwarz bound its low-energy integral by
\[
 C\lambda^{3+c_\mu}
             [I_\mu(B;-3000)I_\mu(J_{\mu b};c_\mu)]^{1/2}.
\]
With \(\delta=1\), its exponent is at most
\(-146\) if \(\mu=l\), and at most \(-221.5\) otherwise.
Indeed \(E+v\le E+s\), and
\(e^{-E/s}(E+s)=s(1+E/s)e^{-E/s}\le Cs\).
The two cases give
\[
 3-150+\tfrac12(1+1)=-146,\qquad
 3-300+\tfrac12(150+1)=-221.5.
\]
When the line axis is one of \(a,b\), the second bound
improves because \(-c_l=100\).

If \(\mu=b\), use the diagonal estimate. Above
\(E=\lambda^{-3000}\),
\(E+\lambda^{-3000}\le2E\) and
\(e^{-E/\lambda^{-100}}\le C[W/(E+W)]^2\).
Spectral Schwarz with the third bound in
\eqref{eq:angular-coordinate-low-energy} gives exponent at most
\(3+(150+1-396)/2=-119.5\).
In detail, on \(E\ge v\),
\[
 e^{-E/\lambda^{-100}}
 \left(\frac{(E+v)(E+W)^4}{EW^4}\right)^{1/2}\le C.
\]
Here \(E+v\le2E\), \(\lambda^{-100}\le W\), and
\(e^{-E/\lambda^{-100}}(1+E/W)^2\le C\).
Thus Schwarz uses the weights \((E+v)^{-1}\) for \(B\)
and \(EW^4/(E+W)^4\) for \(J_{bb}\). The latter integral
costs \(\lambda^\delta W^4=\lambda^{\delta-396}\).
Below that energy, including zero, use the crude diagonal bound
and the factor \((2\lambda^{-3000})^{1/2}\) from the word's
resolvent weight. The exponent is at most
\(3-1500+(150+900)/2=-972\).
For this last use of Schwarz,
\[
 \begin{split}
 \int_{E\le v}H_\lambda^{(b)}\,|\dd\nu_b^{B,J_{bb}}|
 &\le\left(2v I_b(B;-3000)
                  \int_{E\le v}H_\lambda^{(b)}\,\dd\nu_b^{J_{bb}}\right)^{1/2}\\
 &\le C\bigl(v I_b(B;-3000)I_b(J_{bb};0)\bigr)^{1/2}.
 \end{split}
\]
This includes the atom at \(E=0\), where the stronger
diagonal weight itself vanishes.
The complementary energy--momentum region has arbitrary negative
powers by the preceding Schwartz estimate. Exchanging \(a,b\)
gives the identical estimates for \(V_a\). Every remaining term
is therefore \(O(\lambda^{-101})\) before the prefactor and
\(O(\lambda^{-1})\) after it. The angular sum is independent
of \(\lambda\), so it is zero.
\end{proof}

This is an exterior limit of the already constructed infinite-
volume hierarchy. Its order is first the fixed-period cutoff
limit, then the selected volume limit, then \(\lambda\to\infty\).
It is not a claim that an expanding support can be sent to
infinity at fixed finite period. The insertion, shell and moment
inputs had uniform cutoff-tail and volume bounds; the spectral
constants above depend only on the fixed continuum words,
buffers, tensor, profiles and exponents. No mass gap or positive
lower bound on the nonzero spectrum has been used.
% !TeX root = 4DYM_notes.tex
\section{Charge identification and full orthogonal covariance}
\label{sec:rotation-charge-identification}

All statements in this section refer to the same ordered
hierarchy, field quotients and current \(J\). We identify its
charge scale, extend the angular identity beyond coordinate
lines, and construct an actual group representation. The last
step is needed: Lie brackets alone would not specify the
periodicity of a rotation action. The method follows
\cite[Lemmas~16.10--16.11, PDF pp.~184--187]{ArmstrongDraft}.

\subsection{One common translation-charge constant}

The raw charges kill the unit. Indeed remove its slot in a
shell integral; the remaining kernel is smooth through its
centre and conservation makes that shell zero. They preserve
the real structure and all lower degree filtrations, by their
tested laws and injective inclusions.

Translate a collinear configuration on \(\R e_l\) by \(e_l\)
in Theorem~\ref{thm:collinear-angular-identity}, choosing
\(K=E_{bl}-E_{lb}\), \(b\ne l\). Subtract the untranslated
identity. Translation invariance cancels the angular terms,
leaving
\begin{equation}
                   \sum_i S(\ldots,p_bA_i,\ldots;\mathbf x)=0.
 \label{eq:rotation-transverse-vector-ward}
\end{equation}
Its every positional derivative has the same identity with
differentiated labels, by
\eqref{eq:stress-shell-commutators}. Ordinary translation
invariance supplies the identical assertions for \(\partial_b\).

At any sufficiently large degree horizon, transport these two
vector tuples through the coherent frames \(U_k\). They belong
to the finite-dimensional real space \(\mathcal C_d\) of
Section~\ref{sec:free-ward-detection}. Use the finite free
vector test matrix there, and adjoin one real functional
nonzero on \(\partial^{\rm Max}\). The augmented free matrix
is injective. The interacting tested matrices converge to it by
\eqref{eq:stable-smooth-comparison}, so have bounded left
inverses on a tail. In particular the unaugmented interacting
kernel has dimension at most one. It contains both transported
tuples, because their collinear Ward jets vanish. The derivative
tuple is nonzero by
\eqref{eq:stable-derivative-comparison}. Thus
\(p_b=c_d\partial_b\) in this horizon. Compatibility with
inclusions and nonzero differentiation in degree four make all
these constants equal. Embedding each smaller degree in a large
one proves the same identity there.

We determine that this common real constant \(c_\star\) is
nonzero by the local shell, not by choosing a normalization in
advance. Change variables \(y=r_kz\) in both raw shell laws,
and use \eqref{eq:stress-normalised-preimage} and smooth
comparison on its compact shell. Finite separating response
matrices have bounded inverses, so the tested limits give the
operator limits
\begin{align}
 \frac{r_k}{\xi_k}U_{k,d+1}^{-1}p_bU_{k,d}
                          &\longrightarrow\partial_b^{\rm Max},\notag\\
 \xi_k^{-1}U_{k,d}^{-1}Y_KU_{k,d}
                          &\longrightarrow L_d^{\rm Max}(K).
 \label{eq:rotation-raw-charge-limits}
\end{align}
The factors are worth checking. A translation shell contributes
\(r_k^4\) from volume and \(r_k^{-1}\) from its cutoff
derivative, requiring the prefactor \(r_k/\xi_k\). An angular
shell has in addition the factor \(r_k\) from
\(K(y-x_0)\), requiring \(1/\xi_k\). Both then contain
the same normalized tensor \(\xi_k^{-1}r_k^4J\).
Substitute \(p=c_\star\partial\) in the first limit and use
the nonzero derivative comparison. It gives
\(c_\star\ne0\) and \(c_\star/\xi_k\to1\).

Put \(D_K=c_\star^{-1}Y_K\). The commutator identity becomes
\begin{equation}
             \partial_bD_K+\sum_iK_{ib}\partial_i=D_K\partial_b,
 \label{eq:rotation-derivative-commutator}
\end{equation}
and \(U_{k,d}^{-1}D_KU_{k,d}\to L_d^{\rm Max}(K)\).
The normalized collinear angular identity is now
\[
                  \sum_i[S(\ldots,D_KA_i,\ldots)
                                      +(Kx_i)\cdot\nabla_iS]=0.
\]

\subsection{From collinear identities to arbitrary configurations}

Smoothness alone would not continue an identity away from a
line. We prove the analytic property actually needed, using
the coordinate reflection spaces and finite-dimensional charges.
Fix a coordinate plane and \(K=E_{ba}-E_{ab}\). For a finite
label basis write
\[
 \Phi_{u,v}(t,s)=S(u(te_a+se_b),v(0)),\qquad
 \Psi_{u,v}(t,s)=S((e^{-\theta D_K}u)(\rho e_a),
                                      (e^{-\theta D_K}v)(0)),
\]
where \((\rho,\theta)\) are polar coordinates near \(t>0,s=0\).
The axial functions are matrix elements of \(e^{-tH_a}\)
after fixed positive damping of their end vectors. They are
holomorphic for positive real part by
Section~\ref{sec:coordinate-spectral-bounds}. Thus \(\Psi\)
is real analytic near every compact positive axial interval.
It satisfies the angular equation and agrees with \(\Phi\)
on the axis. All angular jets of \(\Phi\) vanish there by the
collinear identity and
\eqref{eq:rotation-derivative-commutator}. For
\(F=\Phi-\Psi\) this gives the recurrence
\begin{equation}
 t\partial_s^{m+1}F_{u,v}(t,0)
 =-\partial_s^m(F_{D_Ku,v}+F_{u,D_Kv})(t,0)
                           +m\partial_t\partial_s^{m-1}F_{u,v}(t,0).
 \label{eq:rotation-transverse-jet-recurrence}
\end{equation}
Omit the last term for \(m=0\). Induction, for the whole
finite label basis at once, makes every transverse jet of \(F\)
zero. At this point no analyticity of \(\Phi\) has been assumed.

In \(\mathcal H_a\), let \(v(s)\) be a single insertion
vector with fixed positive time, translated by \(se_b\).
Its Gram kernel depends on \(s'-s\). The preceding jet
calculation identifies all its diagonal derivatives with those
of the finite analytic matrix \(\Psi\), with the reflected
conjugate label in its first slot. Smoothness of the Gram kernel
implies norm differentiability of the vector: difference quotients
are Cauchy because their squared norm differences are mixed
differences of that kernel. Repeating proves norm derivatives
of all orders. Cauchy estimates on a fixed complex neighbourhood
give
\begin{equation}
 \|v^{(m)}(s)\|^2\le C\rho_0^{-2m}(2m)!
                         \le C(2/\rho_0)^{2m}(m!)^2.
 \label{eq:rotation-vector-analyticity}
\end{equation}
Transverse translation is unitary, so this bound is uniform in
\(s\). The integral Taylor remainder is bounded by
\(C^{1/2}(2|h|/\rho_0)^m\); hence the vector is analytic
for \(|h|<\rho_0/2\). Pairing it with a fixed exterior-word
vector proves analyticity whenever a moving insertion is
strictly earliest or latest in some coordinate. Axial movements
use the heat semigroup. Therefore this assertion needs no
continuous rotation action or spectral cone.

Let \(\mathcal W_K\) now denote the normalized angular sum.
Every positional derivative of it is the same angular sum with
the corresponding differentiated label, by
\eqref{eq:rotation-derivative-commutator}. Consider first a
configuration with distinct first coordinates, a dense set.
Move its points to the first axis by this finite procedure.
Among the points not yet moved, take the one with smallest first
coordinate; it is strictly earliest in axis one. Move it along
axis two past all other second coordinates. Then move it in
axis one past all other first coordinates, while it remains
strictly latest in axis two. Finally set its coordinates two,
three and four to zero, keeping it strictly latest in axis one.
Repeat. Previously moved points have first coordinates beyond
all original ones, so the next point is again strictly earliest.
Every segment stays separated and has an extreme coordinate
giving the analyticity just proved. The final configuration is
collinear with distinct points.

All jets of every angular sum vanish at that final configuration.
Reverse the segments. For every differentiated label list the
one-variable analytic function along a segment has all jets
zero at its endpoint, and hence is zero along the segment.
This preserves all jets at the next endpoint. Induction over
the finite path proves the identity at the original configuration.
Continuity extends it from distinct first coordinates to all
separated configurations. We have obtained the global identity
\begin{equation}
       \sum_i\left[S(\ldots,D_KA_i,\ldots;\mathbf x)
                         +(Kx_i)\cdot\nabla_iS(\mathbf A;\mathbf x)\right]=0.
 \label{eq:rotation-global-angular-ward}
\end{equation}

\subsection{An actual orthogonal action}

First establish the Lie brackets. Commuting two global Ward
identities shows that
\(\Delta(K,L)=[D_K,D_L]-D_{[K,L]}\) has zero summed
responses. It is real, alternating, kills the unit, preserves
every lower degree, and has the required signed covariance.
Transport its finite-horizon maps through \(U_k\). The finite
internal test matrix of
Lemma~\ref{lem:finite-free-internal-tests} is injective, and
smooth comparison makes the interacting matrix injective on a
tail. Thus its transported defect is zero. Faithful inclusions
give \([D_K,D_L]=D_{[K,L]}\) in every degree.

We give a finite-dimensional construction of the group action,
so no integration or period assumption is hidden. Fix a degree
and the six orthonormal plane generators \(K_i\). Write
\(C_i=L^{\rm Max}(K_i)\) and
\(B_{k,i}=U_k^{-1}D_{K_i}U_k\to C_i\).
On endomorphisms of the free label space define
\[
 \rho_{k,i}(X)=B_{k,i}X-XC_i,\qquad
                       \Omega_k=\sum_i\rho_{k,i}^{,2}.
\]
The brackets just proved make \(\rho_k\) a Lie-algebra
representation. The sum of squares commutes with every
\(\rho_{k,i}\): expand the commutator and cancel pairs of
the antisymmetric structure constants. An invariant free inner
product makes \(\Omega_0\), with \(B_i=C_i\), self-adjoint
and nonpositive. Its kernel is the common commutant. A real
polynomial \(z\), equal to one at zero and zero on the other
finitely many eigenvalues, gives its kernel projection \(P\).
Put \(A_k=z(\Omega_k)\to P\).

Turn these near projections into projections without losing
commutation. Iterate \(\varphi(t)=3t^2-2t^3\). The identities
\[
 \varphi(t)-t=(1-2t)(t^2-t),\qquad
 \varphi(t)^2-\varphi(t)=(4(t^2-t)-3)(t^2-t)^2
\]
show that near \(P\) the idempotency error squares and the
successive displacements have a convergent sum. The limits
\(P_k\to P\) are projections commuting with all relative
generators. Their compressed generators tend to zero, since
the limiting commutant is killed by \(\rho_{0,i}\).
If \(m_k\) is their maximum norm, the plane bracket relations
give \(m_k\le2m_k^2\): each plane generator is the bracket
of two others. Since \(m_k\to0\), all compressed generators
are zero eventually. Hence
\[
               f_k=P_k(I)\longrightarrow I,
                          \qquad B_{k,i}f_k=f_kC_i.
\]
The map \(f_k\) is invertible on a tail. All operations are
real and invariant under the signed conjugation action, so
\(f_k\) commutes with the free signed group. For one sufficiently
large \(k\) define
\begin{equation}
 R_d(Q)=U_{k,d}f_kR_d^{\rm Max}(Q)f_k^{-1}U_{k,d}^{-1}.
 \label{eq:rotation-actual-action}
\end{equation}
It is a continuous real representation of \(O(4)\), extends
the original signed action and has generator \(D_K\).
In particular all full-turn periods are inherited from an
actual orthogonal representation.

\begin{theorem}[Full orthogonal covariance and normalized local charges]
\label{thm:full-orthogonal-covariance}
For the same ordered continuum hierarchy there are real
continuous \(O(4)\) actions on \(\mathcal F_d\) such that
\begin{equation}
                    S(R(Q)\mathbf A;Q\mathbf x)=S(\mathbf A;\mathbf x)
 \label{eq:full-orthogonal-covariance}
\end{equation}
on every separated configuration. They preserve units and
inclusions, and
\[
                  R_{d+1}(Q)\partial_a=\sum_bQ_{ba}\partial_bR_d(Q).
\]
There is one nonzero centred real symmetric conserved tensor
\(T_{ab}\in\mathcal F_4\) with
\begin{align}
 R_4(Q)T_{ab}&=\sum_{i,j}Q_{ia}Q_{jb}T_{ij},\notag\\
 \mathcal Q_{e_b}^T(A,\mathbf B)&=S(\partial_bA,\mathbf B),\notag\\
 \mathcal Q_{K(\cdot-x_0)}^T(A,\mathbf B)&=S(L_d(K)A,\mathbf B),
 \label{eq:normalised-interacting-stress-laws}
\end{align}
where \(L_d(K)\) is the generator of \(R_d\).
\end{theorem}
\begin{proof}
Equation~\eqref{eq:rotation-global-angular-ward} and
Proposition~\ref{prop:angular-ward-criterion} give covariance
and the inclusion and derivative laws. Units are fixed because
\(D_K\) kills them and the original signed group fixes them.

Put \(J_0=J/c_\star\). To impose its tensor law without
assuming it, take the finite-dimensional probability Haar integral
\[
 T_{ab}=\int_{O(4)}\sum_{i,j}Q_{ai}Q_{bj}R_4(Q)(J_0)_{ij}\,\dd Q.
\]
Haar invariance gives the tensor law with our column-index
convention. Reality, symmetry and centring pass through the
integral; the derivative vector law gives conservation.
Change variables \(y=Qz\) in each shell. Covariance and
\(R(Q)L(Q^{-1}KQ)=L(K)R(Q)\) show that each transformed
tensor has the same two normalized shell actions as \(J_0\).
Integration preserves them. It cannot give the zero tensor,
because nonzero free differentiation and
\eqref{eq:stable-derivative-comparison} imply that
\(\partial_b\) is nonzero on some degree-four field.
Finally \(T\) and \(J_0\) lie on the same conserved line.
Their identical nonzero translation charge fixes the scalar
ratio to one. Thus the average equals \(J_0\), and the
normalization and full tensor law hold for this one fixed tensor.
\end{proof}

All shell estimates use compact separated kernels. The exterior
estimate was taken only after the ordered physical limits;
neither the ultraviolet family nor its volume extraction has
been changed. Comparison radii were refined once for all finite
requests. Constants in finite-degree matrices may depend on
the degree, fixed tests and calendar; the boundary constants
may also depend on the fixed word, axis, buffer and profile.
They have no residual dependence on a lattice cutoff or period.
The earlier uniform source and derivative estimates supplied
every physical input. No spectral cone, mass gap or prior
continuous rotation covariance was assumed in obtaining the
exterior bound.
% !TeX root = 4DYM_notes.tex
\section{Bounded approximations before phase selection}
\label{sec:bounded-coarse-variance}

Phasewise reflection invariance requires ordinary probability
positivity as well as reflection positivity. Fine observables are
bounded at every cutoff, but that bound may diverge with the
cutoff. The remedy is to stop at one fixed coarse level and
pull a bounded coarse function back to the original probability
space. We give the support and replicated-error arguments
explicitly. This adapts
\cite[Lemmas~14.28--14.29]{ArmstrongDraft}.

\subsection{The probability pullback and its actual dependence}

Distinguish a density push \(\mathcal P_{n,h}\), acting against
coarse Haar measure, from the probability pullback
\[
 (\mathcal K_{n,h}V)(U_n)
    =\int V(U_h)\prod_{i=h}^{n-1}
           K_i(U_{i+1},U_i)\,\dd H_i(U_i).
\]
The parameters in \(K_i\) are those of the same physical Wilson
history. The kernels in \eqref{eq:physical-kernel} have mass one,
so \(\mathcal K_{n,h}1=1\), preserve real values and positivity,
and satisfy \(\|\mathcal K_{n,h}V\|_\infty\le\|V\|_\infty\).
Fubini gives
\begin{equation}
 \int F_n\mathcal K_{n,h}V\,\dd H_n
          =\int(\mathcal P_{n,h}F_n)V\,\dd H_h.
 \label{eq:bounded-push-pull}
\end{equation}
In particular there is no division by an effective density or
by a sector mass.

Here is the required product rule. At one step the actual
kernel can be written
\begin{equation}
 K(U,V)=\sum_{t}\prod_x m_x(U,t_x)
                        \prod_e k_e(U,t_{\partial e};V_e).
 \label{eq:bounded-kernel-parts}
\end{equation}
The site weights satisfy \(\sum_{t_x}m_x(U,t_x)=1\).
The link weight includes the predictor-tag sum and the original
noise normalizer, and \(\int k_e\,\dd h_G=1\).
It uses the shared tags at both endpoints of \(e\).
These identities follow by summing the bond choices in the
likelihood of \eqref{eq:physical-kernel}; the site likelihood
is a product over physical sites.

Suppose \(V_1,\ldots,V_{\mathsf M}\) read only links with both endpoints
in pairwise disjoint site sets \(X_1,\ldots,X_{\mathsf M}\).
For fixed \(t,U\), integrate the unused links in
\eqref{eq:bounded-kernel-parts}. They contribute one.
The remaining integrations split over the \(X_j\).
Each resulting factor reads only the tags in its own \(X_j\).
Summing the unused tags also gives one; summing the remaining,
disjoint sets of tags therefore proves
\begin{equation}
       \mathcal K\Bigl(\prod_jV_j\Bigr)
                            =\prod_j\mathcal K V_j.
 \label{eq:bounded-disjoint-factorisation}
\end{equation}
This proof retains a single tag at a shared endpoint.
The hypothesis is disjoint site sets, not merely distinct links.

For iteration enlarge each site set by all fine sites read by
its frames, predictor, tests and path transport. A step with
ratio \(L_i\) enlarges a native radius \(R\) to at most
\(L_iR+C(J(i)+1)L_i\). In physical units the full ancestor
radius of a fixed level-\(h\) section is consequently bounded by
\begin{equation}
 r_h^{\rm anc}\le {\varepsilon}_hR_d+
                    C\sum_{i\ge h}{\varepsilon}_i(J(i)+1)
                              \longrightarrow0.
 \label{eq:bounded-ancestor-radius}
\end{equation}
The slow profile calendar and \({\varepsilon}_{i+1}\le {\varepsilon}_i/3\) give convergence
of this tail. The estimate includes inactive and fallback tests,
and is uniform in terminal cutoff and period, using periodic
distance. If these ancestor sets are disjoint at every step,
induction proves \eqref{eq:bounded-disjoint-factorisation} for
\(\mathcal K_{n,h}\). Gauge covariance, translation and the finite
signed symmetries pass through the same integrals.

\subsection{Matching at a fixed birth level}

At level \(h\), use the entire compact-group polynomial section
\[
 V_h(x;U_h,y)={\varepsilon}_h^{-d}
                  P_d[\mathsf D_{h_h^{-1}}x](U_h,y/{\varepsilon}_h).
\]
Here \(h_h\) is the previously defined small-field threshold.
The polynomial section is defined on the whole compact link
space. At fixed \(h,d\) its supremum on bounded coefficient
balls is finite; no cutoff parameter occurs in this bound.
On good fields it has the original threshold-scaled derivative
norm. On failures its fixed polynomial price is accompanied by
the original forced weight.

Transport this auxiliary section from \(h\) to the reference
degree coordinate, and denote the resulting linear map by
\(A_{n,M;h,d}\). Denote the selected section's map over the same
interval and unmarked history by \(Z_{n,M;h,d}\).
The auxiliary-section calculation at the end of
Section~\ref{sec:position-local-matching}, now with terminal
birth level \(h\), gives
\begin{equation}
 \|Z_{n,M;h,d}^{-1}A_{n,M;h,d}-I\|_{1\to1}\le\tfrac14,
 \qquad
 \|Z_{n,M;h,d}^{-1}\|_{1\to1}\le E_h,\quad
                         \log E_h=o(s_h),
 \label{eq:coarse-bounded-coordinate-inverse}
\end{equation}
where \(s_h=\log {\varepsilon}_h^{-1}\). To check the application, subtract
the selected section with the same reference readout. Its
coordinate at the reference level is zero. The backward cone
\eqref{eq:position-cone-propagation} applies to this difference
on the same unmarked history. The terminal Taylor read is \(x\);
its regular remainder is bounded by the fixed section column.
The last-step factor \(L_{h-1}^{-1/8}\) pays this column by the
already allowed ratio floor. The forced complement uses its
square-root-normalized norm. This is precisely the auxiliary
comparison argument, without assuming a zero remainder for the
selected section. All constants precede cutoff and volume.

For a reference coordinate \(c\), set
\[
 x_{n,M;h,d}=A_{n,M;h,d}^{-1}c,\qquad
 B_{n,M;h}(c,y)=\mathcal K_{n,h}
                    V_h(x_{n,M;h,d};\,\cdot\,,y).
\]
Neumann inversion gives \(|x|_1\le2E_h|c|_1\). Thus
\begin{equation}
     |B_{n,M;h}(c,y)|\le {\mathsf M}_{h,d,c}<\infty
 \label{eq:coarse-cylinder-supremum}
\end{equation}
uniformly in \(n,M,y\). Real inputs can be replaced by real
parts; conjugation and the root triangle inequality preserve
the estimates below.

\subsection{The replicated discrepancy}

Fix a centred cell side \({\varepsilon}_k\) and a bounded degree-\(d\)
coefficient request. Choose a fixed stopping level \(m\ge k\)
with \(B_m^2\le {\varepsilon}_k/{\varepsilon}_m\) and all stencil margins.
For sufficiently fine \(h\ge m\), the ancestor radius
\eqref{eq:bounded-ancestor-radius} fits strictly inside each
cell. Reflected ancestor sets are then disjoint. Let \(O_j\)
be the original fine insertion and \(V_j\) the matched coarse
section for the copy in cell \(j\).

There is an exact formula for the discrepancy before any norm
is taken. On a mode \(I\) of the square-zero source algebra its
level-\(h\) density is
\begin{equation}
 \sum_{J\subseteq I}(-1)^{|I\setminus J|}
   \mathcal P_{n,h}\left(F_n\prod_{j\in J}O_j\right)
                                      \prod_{j\in I\setminus J}V_j.
 \label{eq:coarse-hybrid-density}
\end{equation}
This is \eqref{eq:common-signed-subset-coefficient} for
\(\mathcal P=\mathcal P_{n,h}\). Its integral is the
expectation of the product of fine discrepancies by
\eqref{eq:common-signed-expectation}:
\eqref{eq:bounded-disjoint-factorisation} verifies the
pullback hypothesis for these disjoint ancestor sets, and
\eqref{eq:bounded-push-pull} supplies adjointness.
The empty mode is the actual unmarked density.

We record how the full-source estimate applies to this signed
birth. A singleton has zero matched reference coordinate;
linearity and the backward cone bound its normalized remainder
at level \(i\le h\) by
\[
        E_h({\varepsilon}_h/{\varepsilon}_i)^{1/8},
\]
after a fixed enlargement of \(E_h\) still satisfying
\(\log E_h=o(s_h)\). This includes regular norms and forced
norms divided by the square root of their original allowance.
There is also a higher-mode birth contribution. Multiplication
by \(\prod_j(1-z_jV_j)\) changes the singleton logarithmic
terms; it does not erase the higher modes of the old source.
Every such mode connects distinct cell anchors and has its full
physical read support.

Here are sufficient explicit prices for those terms. With
\(\ell={\varepsilon}_k\), a connected read carrying a set \(I\) of distinct
cell anchors satisfies
\[
        n_{B_h}(X)\ge
                       c|I|\ell/({\varepsilon}_hB_h).
\]
The doubled spanning-tree argument of
Lemma~\ref{lem:joint-marked-cover} proves this stronger version
with separation \(\ell/{\varepsilon}_h\). One reserved support unit thus
pays \(\exp[-c|I|\ell/({\varepsilon}_hB_h)]\).
Restoring coefficient radii, physical powers, and the finitely
many section columns costs at most
\(\exp[C_d(1+s_h)^2|I|]\).
For this deliberately generous bound use \(h\le C(1+s_h)\),
\(\log\widehat\beta_i\le C+\log(2+s_h)\) for \(i\le h\),
and sum the logarithms of the inverse source radii in
\eqref{eq:joint-radius-transport}. The profile columns are
subpower by the calendar. In particular the bound does not
silently omit the intermediate radius changes.

Here is a choice of weights that closes the mixed induction.
Let \(\epsilon_h\) be the expression inside braces on the
right of \eqref{eq:coarse-mixed-source-error}, and choose
\[
 \rho_m=\frac{\zeta_{d,m}{\varepsilon}_m^d}{4\epsilon_h},\qquad
 \rho_i=\rho_m\prod_{u=m+1}^i\widehat\beta_u^{-1}
                                      \quad(m<i\le h).
\]
The physical singleton envelope is
\(E_h {\varepsilon}_h^{1/8}{\varepsilon}_i^{-d-1/8}\).
The entry inequality already used in
Section~\ref{sec:position-local-matching},
\[
 \widehat\beta_{i+1}^{-1}L_i^{d+1}\zeta_{d,i}
                                             \le\zeta_{d,i+1},
\]
and \(d+1/8\le d+1\) show that its radii-weighted size is
at most \(\zeta_{d,i}/4\) at every intermediate level.
Thus one does not infer a simultaneous source estimate
from singleton cancellation alone.

Compare these radii with the ordinary radii used to transport
the old full source to \(h\). Their ratio is constant across
the interval, since both use
\(\rho_{i+1}=\widehat\beta_{i+1}^{-1}\rho_i\).
Its logarithm is at most \(C_d(1+s_h)^2\), including the
coarse section and normalization columns. For higher modes
the preceding separated-anchor cover estimate pays this
inflation at the signed birth. Keep one fewer stored support
unit there, within the available range \(98\)--\(100\).
Restoration components joining several coarse stencils have
the same named-support property and are paid in the same way.
This supplies the higher-mode income at \(h\), in addition
to the singleton income just checked.

Now apply the proof of
Lemma~\ref{lem:joint-source-attachments} to the complete
signed source on \([m,h]\), using these actual ordinary
radii. At each rooted component keep its distinguished
singleton discrepancy or connected birth, and bound all
other attachments by the small complete income.
The adjacent radius increase is exactly the one in that
lemma, so its separated-name cover estimate applies without
a stronger marker wallet. Cover compression restores the
output reserve at every step. Mixed regular/forced transfers,
all higher modes and their full supports are included.
The descendant sum is at most twice its distinguished income,
as in the original induction. This proves bounded complete
income at \(m\); extraction at radius \(\rho_m\) gives the
claimed small root.

Consequently there are radii for the error names whose
level-\(m\) coefficient estimate gives the reflected-cell root
bound
\begin{equation}
 \|B_{n,M;h}(c,y)-\mathcal O_{n,M}((d,c),y)\|_k
      \le C_{k,m,d,c}\left\{
 E_h({\varepsilon}_h/{\varepsilon}_m)^{1/8}
       +e^{C_d(1+s_h)^2-c {\varepsilon}_k/({\varepsilon}_hB_h)}\right\}
                         =:\eta_h\longrightarrow0.
 \label{eq:coarse-mixed-source-error}
\end{equation}
To see the volume cancellation explicitly, if there are \({\mathsf M}\)
cells, coefficient extraction before the root costs
\[
  \eta_h^{\mathsf M}\exp\{\#\Lambda_mQ_m^{\rm vac}+C {\mathsf M}\}.
 \]
Divide by the unchanged positive partition integral and take
the \({\mathsf M}\)-th root. Since
\(\#\Lambda_m/{\mathsf M}=({\varepsilon}_k/{\varepsilon}_m)^4\), the remaining factor is part of
the fixed constant in \eqref{eq:coarse-mixed-source-error}.
An unreplicated absolute-density estimate would retain an
extensive cost and would not prove volume uniformity.
The second term tends to zero since
\(\log B_h=o(s_h)\). Every mask, branch mass, incoming activity,
Gaussian/Haar normalization and the vacuum scalar remains in
the exact density calculation. Signed and reflected copies use
one finite enlarged collection of sections and coefficient balls.

\subsection{The pair estimate needed by variance}

A real grid word \(w\) has real coefficient representatives at
distinct \({\varepsilon}_k\mathbb Z^4\) points; we may first require distinct
coordinates in each axis. Write \(\vartheta w\) for linear
time reflection of its labels and positions, without scalar
conjugation. Put
\[
 z_j(T)=Tj{\varepsilon}_k(1,1,1,1),\qquad
 f_j=\mathcal O_{n,M}(\tau_{z_j(T)}w)
              -\mathcal O_{n,M}(\tau_{z_j(T)}\vartheta w).
\]
Here \(\mathcal O(w)\) denotes the product of the insertion
functions in the word.

\begin{proposition}[A bounded variance approximation]
\label{prop:bounded-variance-approximation}
For every such fixed word and \(\varepsilon>0\), there are
real fine observables \(g_j\) and a finite \(R_\varepsilon\),
independent of \(n,M,T,j\), with \(|g_j|\le R_\varepsilon\).
For fixed \(i<j\) and every sufficiently large fixed integer
\(T\), eventually in the ordered physical limits,
\begin{equation}
       |\langle g_i g_j\rangle_{n,M}
                      -\langle f_i f_j\rangle_{n,M}|
                                             \le\varepsilon.
 \label{eq:bounded-variance-pairs}
\end{equation}
The bound is chosen before the number of copies in a variance
sum. No common cutoff threshold in \(T,i,j\) is asserted.
\end{proposition}

\begin{proof}
Take a finer fixed tiling with interiors accommodating the
finitely many slots of \(w\) and their signed images.
The original insertions have bounded roots there. By
\eqref{eq:coarse-mixed-source-error}, choose \(h\) so that
each replacement error has root at most \(\delta\le1\).
The replacement has root at most \(C+1\), by the triangle
inequality, where \(C\) bounds the original roots.
Telescope a product of \(r\) separated slots and apply the
chessboard inequality with the boundary margins established
above. Its moment error is at most
\[
                      r\delta(C+1)^{r-1}.
 \]
Indeed the exact product identity is
\[
 \prod_{i=1}^r A_i-\prod_{i=1}^r B_i
  =\sum_{j=1}^r(A_j-B_j)
                         \prod_{i<j}B_i\prod_{i>j}A_i.
\]
Each error slot has root at most \(\delta\), each other
slot has root at most \(C+1\), and all anchors remain
in the disjoint interiors. Taking expectations and using
the chessboard inequality proves the displayed bound.
The number of translated copies does not enter its constant.
The finite translated-tiling argument in
Lemma~\ref{lem:insertion-interior-root} supplies the same
estimate for uncentred anchors.

Define \(g_j\) by replacing the slots in the two terms of
\(f_j\) with their real bounded approximations. Then
\[
          |g_j|\le2\prod_{\text{slots of }w}{\mathsf M}_{h,d,c}
                         =R_\varepsilon.
 \]
For large \(T\), all four pair configurations in the expansion
of \(f_if_j\) have the required fixed separation. Apply the
product estimate with at most \(2|w|\) slots. Choose
\(4(2|w|)\delta(C+1)^{2|w|-1}\le\varepsilon\).
For each fixed finite configuration, sufficiently large
periods remove wrapping identifications, and the cutoff tail
supplies its margins. Intersect the four ordered tails.
Outside tails on which the construction is defined set
\(g_j=0\); the global supremum remains valid. A finite list
of pairs is handled by intersecting finitely many tails.
All choices of \(h,\delta,R_\varepsilon\) precede that list.
Thus a variance with \(q+1\) copies uses one bound
\(R_\varepsilon\) for all \(q\); only the finite
intersection of the configuration-dependent tails changes
with \(q\). The four signed moments are all expectations
under the same Wilson measure, by
Lemma~\ref{lem:complete-subset-identity} and the disjoint
full-cone pullback law.
\end{proof}
% !TeX root = 4DYM_notes.tex
\section{The ground-space algebra and phase decomposition}
\label{sec:phase-compression}

The ordered Wilson limit need not be a single vacuum state.
We first decompose that same limit; we do not take another
lattice subsequence to choose a vacuum. The compression argument
and the variance argument in the next section follow
\cite[Sect.~17.1, Lemma~17.4, PDF pp.~194--195,
and Lemma~17.5, pp.~195--197]{ArmstrongDraft}.

Use the fourth-axis reflection space of
Section~\ref{sec:coordinate-spectral-bounds}, and abbreviate it
as \(\mathcal H\). Write \(E_t=e^{-tH}\), \(H\ge0\),
\(\Omega=[1]\), and \(P_0=\mathbf1_{\{0\}}(H)\).
Thus \(P_0\mathcal H=\ker H\), called the ground space.
It is not yet assumed to be a vacuum line. Spatial translation
is the commuting unitary \(U_y\), \(y\in\mathbb R^3\).
The spectral theorem gives \(E_t\to P_0\) strongly as
\(t\to\infty\), including any mass at zero.

\subsection{Why ground vectors have zero spatial momentum}

\begin{lemma}
\label{lem:phase-spatial-ground}
For every spatial \(y\), \(U_yP_0=P_0\).
\end{lemma}
\begin{proof}
For two fixed words placed in disjoint slabs, a sufficiently
large displacement \(z\), with \(z^4>0\), has correlation
\[
                  \langle v,U_{\vec z}E_{z^4-c}v'\rangle,
\]
for fixed slab vectors \(v,v'\) and width constant \(c\).
Changing \(z^4\) by a fixed admissible constant changes this
by a quantity tending to zero as \(z^4\to\infty\),
uniformly in \(\vec z\). Indeed the norm of the difference
is bounded by
\(\|v\|\|(E_s-I)E_t v'\|\), with an inessential fixed
change of \(t\). Spectral calculus makes this tend to zero.
This argument does not use spatial invariance of \(P_0\).

For fixed \(y\), choose an orthogonal reflection \(K\) taking
\(y+e_4\) to \(|y+e_4|e_4\). Since
\[
       \langle K(y+e_4),Ke_4\rangle=1,
       \qquad (Ke_4)^4=|y+e_4|^{-1}>0,
\]
both rotated long displacements \(tKe_4+K(y+e_4)\) and
\(tKe_4\) have large positive fourth coordinate, and their
difference is purely temporal. Full orthogonal covariance,
already proved in Theorem~\ref{thm:full-orthogonal-covariance},
and the preceding time-shift estimate show that the original
correlations with displacements \(te_4+y+e_4\) and \(te_4\)
have the same limit. Remove the fixed extra \(e_4\) by that
estimate again. Strong heat convergence then gives
\[
           \langle v,U_yP_0v'\rangle
                           =\langle v,P_0v'\rangle.
 \]
Slab word vectors are dense. This proves the operator identity.
\end{proof}

\subsection{Compression and distant clusters}

Let \(\mathsf B_h(w)\) denote the bounded buffered insertion
operator of \eqref{eq:coordinate-buffered-operator}; its width
is \(h\). Define
\[
           \mathsf M^0(w)=P_0\mathsf B_h(w)|_{\ker H}.
\]
Any word can first be translated into a positive buffered slab.
Changing width or translating within time uses the exact rules
\[
 \mathsf B_{h+t}(w)=\mathsf B_h(w)E_t,\qquad
 \mathsf B_{h+t}(\tau_{te_4}w)=E_t\mathsf B_h(w).
\]
Compression removes \(E_t\); Lemma~\ref{lem:phase-spatial-ground}
removes spatial placement as well. Hence \(\mathsf M^0(w)\) is
independent of these choices. Furthermore
\begin{equation}
 \mathsf M^0(1)=I,\quad \mathsf M^0(w)\Omega=P_0[w],\quad
 \langle\Omega,\mathsf M^0(w)\Omega\rangle=S(w),\quad
 \mathsf M^0(w)^*=\mathsf M^0(\vartheta\bar w).
 \label{eq:phase-compression-laws}
\end{equation}
Here \(\bar w\) conjugates labels and coefficients;
\(\vartheta\) is linear fourth-coordinate reflection on both
positions and labels. Reversing a Euclidean word has no effect
on its correlations.

For fixed finite words and displacement lists whose successive
fourth-coordinate gaps tend to infinity, the slab insertion
rule gives
\begin{equation}
 \lim_{j\to\infty}
    S(\tau_{z_1(j)}w_1\cdots\tau_{z_r(j)}w_r)
      =\langle\Omega,\mathsf M^0(w_1)\cdots \mathsf M^0(w_r)\Omega\rangle.
 \label{eq:phase-distant-clusters}
\end{equation}
To justify repeated limits when transverse gaps also vary, if
\(\psi_j\to\psi\) and \(t_j\to\infty\), then
\[
 \|U_{y_j}E_{t_j}\psi_j-P_0\psi\|
       \le\|\psi_j-\psi\|+\|(E_{t_j}-P_0)\psi\|\to0.
\]
This proves the finite induction behind
\eqref{eq:phase-distant-clusters} with arbitrary \(y_j\).
Products of \(\mathsf M^0(w)\)'s describe distant clusters. They do
not assert factorisation of finite-separation correlations.

\begin{proposition}[A commutative algebra on the ground space]
\label{prop:phase-compression-algebra}
The norm-closed unital algebra generated by the \(\mathsf M^0(w)\)
and their adjoints is a separable commutative \(C^*\)-algebra
\(\mathcal A_0\subset\mathcal B(\ker H)\).
The unit vector \(\Omega\) is cyclic and separating for it.
Here cyclic means that \(\mathcal A_0\Omega\) is dense;
separating means \(a\Omega=0\) implies \(a=0\).
\end{proposition}
\begin{proof}
Place four words at
\[
 0,\quad j(e_4+2e_1),\quad j(2e_4+e_1),
                                      \quad3j(e_4+e_1).
\]
Exchanging axes 1 and 4 reverses the order of the middle
two time coordinates. Covariance and
\eqref{eq:phase-distant-clusters} therefore equate the
corresponding compressed products, with all four words also
transformed by that exchange. A second cluster calculation,
now placing all four words on the exchange-invariant line
\(\mathbb R(e_4+e_1)\), removes that simultaneous
transformation. Thus arbitrary outer words give
\[
 \langle\Omega,\mathsf M^0(u)
           [\mathsf M^0(w),\mathsf M^0(v)]\mathsf M^0(u')\Omega\rangle=0.
\]
The span of \(P_0[b]=\mathsf M^0(b)\Omega\) is dense in \(\ker H\):
project the dense positive-word span of \(\mathcal H\).
The adjoint law in \eqref{eq:phase-compression-laws} now
proves commutativity as operators. It also proves cyclicity.
If \(a\Omega=0\), then \(ab\Omega=ba\Omega=0\) for every
\(b\in\mathcal A_0\); density proves \(a=0\).

A \(C^*\)-algebra here means a norm-closed operator algebra
containing the identity and closed under adjoints. It
automatically satisfies \(\|a^*a\|=\|a\|^2\).
Finite degree bases, rational complex coefficients and
configurations with rational coordinates give a countable
collection of words. On compact buffered configuration sets
the operator norm Taylor estimate in
Section~\ref{sec:coordinate-spectral-bounds} approximates all
words by these. They generate a dense subalgebra, proving
separability.
\end{proof}

\subsection{The compact space of characters}

A character of \(\mathcal A_0\) is a complex-linear map
\(\sigma:\mathcal A_0\to\mathbb C\) satisfying
\(\sigma(I)=1\), \(\sigma(ab)=\sigma(a)\sigma(b)\), and
\(\sigma(a^*)=\overline{\sigma(a)}\).
Let \(\Sigma\) be their space with the topology of pointwise
convergence on \(\mathcal A_0\).

We recall the precise elementary functional analysis used
here \cite[Sect.~2.4, Theorem~2.4.1, original preprint p.~25;
proof on pp.~25--27]{Landsman98}.
The character representation of a commutative unital
\(C^*\)-algebra identifies it isometrically with \(C(\Sigma)\):
\[
        a\longmapsto\widehat a,\qquad
                         \widehat a(\sigma)=\sigma(a).
\]
For completeness, maximal ideals give the characters by the
complex Banach-algebra spectrum theorem. Their norms are one,
so their space is closed in the weak-star unit ball and is
compact. The \(C^*\) identity and the spectral-radius formula
for normal elements give
\(\sup_\sigma|\sigma(a)|=\|a\|\). The image is closed, conjugation
invariant, contains constants and separates characters;
the complex Stone--Weierstrass theorem makes it all of
\(C(\Sigma)\). A countable dense set in \(\mathcal A_0\)
metrizes pointwise convergence on its norm-one characters,
so \(\Sigma\) is compact metric.

The positive state
\(\omega(a)=\langle\Omega,a\Omega\rangle\) is a positive
functional on \(C(\Sigma)\) of value one at the constant one.
The Riesz representation theorem
\cite[Sect.~2.8, Theorem~2.8.2, original preprint p.~35]{Landsman98}
gives a probability
measure \(\rho\). It has full support: a nonnegative,
nonzero continuous function supported in any nonempty open
set corresponds to \(a\ge0\), \(a\ne0\).
If \(\omega(a)=0\), then \(a^{1/2}\Omega=0\);
the separating property would force \(a=0\).

Define, for every finite separated word,
\begin{equation}
           S^\sigma(w)=\sigma(\mathsf M^0(w)).
 \label{eq:phase-character-values}
\end{equation}
These are called the phase correlations. They obey
\begin{equation}
            S(w)=\int_\Sigma S^\sigma(w)\,\dd\rho(\sigma).
 \label{eq:phase-mixture}
\end{equation}
Each separated value is continuous in \(\sigma\).
At this point the formula is a character decomposition,
not yet a claim that every \(S^\sigma\) has reflection
positivity, the full field spaces, or a unique vacuum.
Those assertions require the next arguments.
% !TeX root = 4DYM_notes.tex
\section{Reflection-invariant phases and a selected vacuum}
\label{sec:phase-selection}

The variance argument and the construction of phase fields with a
selected vacuum adapt
\cite[Sect.~17.1, Lemma~17.5, PDF pp.~195--197]{ArmstrongDraft}.
The four-moment passage and the common full-measure selection
also follow \cite[Unit V.1, Lemmas V.1.m--n, PDF p.~394,
and Lemma V.1.y, p.~398]{ArmstrongNotes}.

\subsection{Variance makes reflection invariance characterwise}

Reality of a character alone would not imply reflection
invariance. We prove the required operator identity before
evaluating characters.

\begin{lemma}[Reflection and rotation in the compression algebra]
\label{lem:phase-operator-invariance}
For every separated word and \(Q\in O(4)\),
\begin{equation}
       \mathsf M^0(\vartheta w)=\mathsf M^0(w),\qquad
                           \mathsf M^0(Qw)=\mathsf M^0(w).
 \label{eq:phase-operator-invariance}
\end{equation}
\end{lemma}
\begin{proof}
% !TeX root = 4DYM_notes.tex
The empty word is invariant, so fix a nonempty real word
\(w\) on one of the cofinal generic grids of
Proposition~\ref{prop:bounded-variance-approximation}. All
its labels and reflected labels are included in one finite
request. Write
\[
 A=\mathsf M^0(w),\qquad B=\mathsf M^0(\vartheta w),
                            \qquad Z=A-B.
\]
These are elements of the already constructed commutative
algebra \(\mathcal A_0\). The adjoint law, and reality
of the labels, give \(A^*=B\), \(B^*=A\). Hence
\begin{equation}
 Z^*=-Z,\qquad
 a=\langle\Omega,Z^2\Omega\rangle
                         =-\|Z\Omega\|^2\le0.
 \label{eq:phase-variance-negative-square}
\end{equation}
Ordinary positivity of the original Wilson probability
space will supply the reverse inequality.

Fix a tolerance \(\varepsilon>0\). Choose the bounded
approximations of Proposition~\ref{prop:bounded-variance-approximation}
before choosing any number of copies. To recall the relevant
dependence, at a sufficiently fine fixed tiling the original
roots are bounded by \(C\). Choose a slot-replacement
error \(\delta\le1\) so that
\[
        8|w|\delta(C+1)^{2|w|-1}\le\varepsilon.
\]
Then choose its fixed coarse birth level \(h\). The
whole-space coarse sections have suprema
\(\mathsf M_{h,d,c}\), and their probability pullbacks
are contractions in the supremum norm. Thus
\[
         |g_j|\le R_\varepsilon,
       \qquad R_\varepsilon=2\prod_{\text{slots of }w}\mathsf M_{h,d,c}.
\]
The value of \(R_\varepsilon\) may grow as the tolerance
decreases. It is independent of the terminal cutoff,
physical period, displacement parameter \(T\), copy
index \(j\), and the number of copies in the variance.
All constants use the same Wilson history and the original
partition denominator, as proved by the complete signed
source calculation in Section~\ref{sec:bounded-coarse-variance}.

For a fixed integer \(q\ge1\), take copies \(j=0,\ldots,q\)
at \(z_j(T)=Tj\varepsilon_k(1,1,1,1)\). For a fixed
sufficiently large integer \(T\), the finitely many pair
estimates hold together on an ordered tail. On that tail
ordinary variance positivity gives
\begin{equation}
 \begin{split}
 0&\le\left\langle\left(\sum_{j=0}^qg_j\right)^2\right\rangle_{n,M}\\
  &\le(q+1)R_\varepsilon^2+
                  2\sum_{0\le i<j\le q}
                                 \bigl(\langle f_if_j\rangle_{n,M}+\varepsilon\bigr).
 \end{split}
 \label{eq:phase-finite-variance}
\end{equation}
Here \(f_j\) is the difference of the original word and
its reflected word at that copy. The diagonal terms are
bounded using \(g_j\), whose supremum is uniform. No
bound for the coincident product \(f_j^2\) is used.

We take limits only on the right side of
\eqref{eq:phase-finite-variance}. For each fixed \(q,T\),
first take the ultraviolet limit at every fixed sufficiently
large period, then the selected volume limit. Only finitely
many moments occur, so the finite intersection of their tails
suffices. In particular no convergence of the functions
\(g_j\), or of their own correlations, has to be extracted.
The cutoff threshold can depend on the fixed period and
configuration; the bound \(R_\varepsilon\) does not.

For each pair \(i<j\), the right-side correlation is
exactly the signed sum of four separated moments:
\[
 \begin{split}
 \langle f_if_j\rangle_{n,M}
  ={}&\langle w_iw_j\rangle_{n,M}
       -\langle w_i(\vartheta w)_j\rangle_{n,M}\\
    &-\langle(\vartheta w)_iw_j\rangle_{n,M}
       +\langle(\vartheta w)_i(\vartheta w)_j\rangle_{n,M}.
 \end{split}
\]
The subscripts denote the indicated translations of the
physical insertion products. After their ordered limits,
send \(T\to\infty\). Their positive time gap is
\((j-i)T\varepsilon_k\) up to the fixed word widths.
Equation~\eqref{eq:phase-distant-clusters} gives the four
limits, in their displayed order,
\[
 \langle\Omega,A^2\Omega\rangle,quad
 \langle\Omega,AB\Omega\rangle,quad
 \langle\Omega,BA\Omega\rangle,quad
 \langle\Omega,B^2\Omega\rangle.
\]
Their signed sum is
\(\langle\Omega,(A-B)^2\Omega\rangle=a\).
This conclusion holds for every fixed pair; for the fixed
\(q\) there are \(q(q+1)/2\) pairs. Thus
\eqref{eq:phase-finite-variance} implies
\[
        0\le(q+1)R_\varepsilon^2+q(q+1)(a+\varepsilon).
\]
Divide by \(q(q+1)\). Since \(R_\varepsilon\) was
fixed before \(q\), letting \(q\to\infty\) gives
\(a+\varepsilon\ge0\). Now let \(\varepsilon\downarrow0\).
Together with \eqref{eq:phase-variance-negative-square},
this proves \(Z\Omega=0\). The vacuum vector is separating
for \(\mathcal A_0\), so \(Z=0\).

The order of operations was therefore: fix the tolerance
and bounded coarse approximation; fix finitely many copies
and their separation; take the fixed-period ultraviolet
limits and selected volume limit; increase the separation;
increase the number of copies; finally decrease the
tolerance. The supremum bound, not the cutoff thresholds,
is uniform in that copy count.

The cofinal grids approximate every separated real word.
The buffered operator derivative bounds give norm continuity
on each fixed separated compact, so their compressed
operators converge in norm. Multilinearity extends the
result to complex labels by writing each as a real part
plus \(i\) times a real part. This proves the reflection
identity for every word and, in particular,
\(\mathsf M^0(w)^*=\mathsf M^0(\bar w)\).

We explain the passage to all orthogonal maps. For any \(Q\)
admitting a vector \(r\) with \(r^4>0\), \((Qr)^4>0\),
use two-cluster limits along \(tr\). Covariance and the first
identity show
\[
 \langle \mathsf M^0(Qu)\Omega,\mathsf M^0(Qv)\Omega\rangle
              =\langle \mathsf M^0(u)\Omega,\mathsf M^0(v)\Omega\rangle.
\]
Thus the map on this dense word span taking \(w\) to \(Qw\)
is well defined and extends to a ground-space unitary.
Such \(r\) exists except when \(Qe_4=-e_4\); in that case
compose with \(\vartheta\), whose action is already the
identity. For the first case an explicit choice is
\(r=e_4+Q^{-1}e_4\): both required time components equal
\(1+Q_{44}>0\). These word maps form an orthogonal-group
representation. Every conjugate \(Q\vartheta Q^{-1}\)
therefore acts trivially. These are all hyperplane reflections.
An orthogonal map is a product of at most four such reflections:
successively reflect its first basis image to the first basis
vector and restrict to its orthogonal complement.
Consequently every \(Q\) acts trivially on the ground word span.
The second identity in \eqref{eq:phase-operator-invariance}
now follows by separation of \(\Omega\).
\end{proof}

\subsection{Positivity and the common field calculus}

\begin{proposition}[Phase correlations]
\label{prop:phase-correlations}
Every \(S^\sigma\) is a smooth separated Euclidean hierarchy
with permutation symmetry, the unit rule, reality, translations,
full orthogonal covariance and reflection positivity.
It has the same field quotients \(\mathcal F_d\), inclusions,
derivatives and normalized Euclidean stress laws as \(S\).
For each fixed finite degree, arity and coefficient ball it
satisfies the buffered bounds with constants independent of
\(\sigma\). Separated values and their fixed derivatives are
continuous in \(\sigma\).
\end{proposition}
\begin{proof}
The permutation, unit and multilinearity identities hold in
\(\mathcal A_0\): test the corresponding buffered operators
on arbitrary exterior positive words. Smoothness and
derivatives hold in operator norm. Character evaluation is
norm continuous of norm one. Translation and orthogonal
covariance follow from the compression identities; the adjoint
identity and Lemma~\ref{lem:phase-operator-invariance} give
reality.

For positive words \(w_i\) in one common slab, put
\(Y_i=\mathsf B_h(w_i)|_{\ker H}\), as a map
\(\ker H\to\mathcal H\). The exact insertion and adjoint rules
give
\[
 \mathsf M^0(\theta w_iw_j)=Y_i^*Y_j,\qquad
 \sum_{i,j}\bar c_i c_j\mathsf M^0(\theta w_iw_j)
                   =\Big(\sum_i c_iY_i\Big)^*
                                      \Big(\sum_jc_jY_j\Big)\ge0.
\]
Here \(\theta w=\vartheta\bar w\).
This positive operator lies in \(\mathcal A_0\).
Its nonnegative square root also lies in \(\mathcal A_0\):
uniformly approximate the square-root function on its compact
nonnegative spectrum by real polynomials.
Every character therefore evaluates it nonnegatively.
This proves point reflection positivity. Integrating the
continuous positive Gram kernel against compact positive-time
joint tests gives a norm square, proving the smeared version.
Other axes follow by rotation.

A label null for \(S\) has zero matrix elements in every
buffered insertion with that label, since the exterior words
are arbitrary. That insertion operator is zero. Hence
\[
                  \mathcal N_d(S)\subseteq\mathcal N_d(S^\sigma).
\]
Inclusions and derivative identities likewise hold as buffered
operator identities. Compact stress-shell integrals exist in
operator norm: their weights vanish near the marked points
and the fixed spectators. Test on arbitrary exterior words,
use the Euclidean shell laws, and then evaluate the character.
This transfers conservation and both local shell identities
with the same tensor and label maps.

We must still exclude additional phase null labels.
Use the finite singleton/pair free detector rows from
Corollary~\ref{cor:position-finite-free-tests} and their
corrected lifts in
Proposition~\ref{prop:position-relative-gaussian-tests}.
At a fixed comparison trial \(\nu\), denote its one- and
two-slot lifted words by \(u,v\), and its free moments by
\(s,s'\), \(ss'+c\). In \(\mathcal A_0\) put
\[
 \begin{split}
 \Delta_1&=\mathsf M^0(u)-sI,\qquad
 \Delta'_1=\mathsf M^0(v)-s'I,\\
 \Delta_2&=\mathsf M^0(uv)-s\mathsf M^0(v)-s'\mathsf M^0(u)+(ss'-c)I.
 \end{split}
\]
These are exactly the residuals
\eqref{eq:position-relative-pair-errors}, not a pair
subtraction with the means omitted.
Place \(2j\) alternating conjugated copies in distant
comparison cells. Their original Wilson roots are at most
\(\gamma_\nu\). The chessboard estimate, the ordered
physical limits for each fixed \(j\), and then the distant
cluster limit give
\begin{equation}
       \int_\Sigma|\sigma(\Delta)|^{2j}\,\dd\rho(\sigma)
                                         \le\gamma_\nu^{2j}.
 \label{eq:phase-free-row-moments}
\end{equation}
This applies to each of the displayed residuals.
If its continuous absolute value exceeded \(\gamma_\nu\)
at a character, it would exceed some
\(b>\gamma_\nu\) on an open set of positive \(\rho\)-measure.
Letting \(j\to\infty\) contradicts
\eqref{eq:phase-free-row-moments}. Thus these residuals are
bounded by \(\gamma_\nu\) at every character, uniformly.

Expanding \(\Delta_2\) shows that the phase singleton/pair
row matrix of the corrected lifts converges uniformly in
\(\sigma\) to the injective free row matrix. Apply
\eqref{eq:common-fixed-detector-bound} once its uniform
error is less than half the free lower bound. This gives
one detector constant and injectivity for every phase
at the same sufficiently late fixed trial.
These row maps factor through
\(V_d/\mathcal N_d(S^\sigma)\). Their rank is at least
\(\dim V_d^{\rm Max}\). The previously proved dimension
\(\dim\mathcal F_d=\dim V_d^{\rm Max}\), and the null inclusion,
give the reverse inequality. This proves
\begin{equation}
                       \mathcal N_d(S^\sigma)=\mathcal N_d(S).
 \label{eq:phase-equal-null-spaces}
\end{equation}
Thus the field spaces and the transferred calculus are
literally the common ones.

Finally \(|S^\sigma(w)|\le\|\mathsf B_h(w)\|\), and the same
bound holds for each fixed differentiated label list.
This proves phase-uniform separated bounds and continuity.
For use in reconstruction one needs the buffered operator
bound in the phase Hilbert space itself, not just a scalar
bound. Construct its reflection space as before. If \(A_w\)
is its insertion rule and \(v_j=(A_w^*A_w)^jv\), the
mixed-space map \(Y(v_j)=\mathsf B(v_j)|_{\ker H}\) satisfies
\[
 \|[v_j]_\sigma\|^2=\sigma(Y(v_j)^*Y(v_j))
                  \le\|Y(v)\|^2\|\mathsf B_h(w)\|^{4j}.
\]
The repeated Schwarz power argument used for
\eqref{eq:coordinate-buffered-operator} therefore gives
\begin{equation}
        \|\mathsf B_h^\sigma(w)\|\le\|\mathsf B_h(w)\|.
 \label{eq:phase-buffered-operator}
\end{equation}
The insertion, adjoint and operator norm derivative identities
pass to its dense word domain. This proves the claimed common
constants.
\end{proof}

\subsection{Selecting a one-dimensional ground space}

Write \(\mathcal K_\sigma\) for the phase fourth-axis reflection
Hilbert space, \(\Omega_\sigma^{\rm OS}=[1]_\sigma\), and
\(H_\sigma^{\rm OS}\ge0\) for its heat generator.
Let \(P_{0,\sigma}^{\rm OS}=\mathbf1_{\{0\}}(H_\sigma^{\rm OS})\).
These definitions use only the phase Euclidean correlations.

\begin{theorem}[Vacuum selection before reconstruction]
\label{thm:phase-vacuum-selection}
There is a \(\rho\)-full-measure subset of \(\Sigma\) on
which
\[
          \ker H_\sigma^{\rm OS}=\mathbb C\Omega_\sigma^{\rm OS}.
\]
Fix one character \(\sigma_*\) in this set. Its selection changes
neither the Wilson family, the ordered limits defining \(S\),
nor the phase measure in \eqref{eq:phase-mixture}.
\end{theorem}
\begin{proof}
% !TeX root = 4DYM_notes.tex
For a fixed positive word \(b\), write \([b]_\sigma\)
for its class in \(\mathcal K_\sigma\). Its common Gram
bound is finite and uniform in \(\sigma\). At integer
times \(m\ge0\), put
\[
 F_m(\sigma)=\langle[b]_\sigma,
                    e^{-mH_\sigma^{\rm OS}}[b]_\sigma\rangle.
\]
These functions are measurable; in fact each is a continuous
phase correlation of two separated words. By the positive
spectral theorem,
\[
 0\le F_{m+1}(\sigma)\le F_m(\sigma)\le F_0(\sigma),
 \qquad
 F_m(\sigma)\longrightarrow
                       \|P_{0,\sigma}^{\rm OS}[b]_\sigma\|^2.
\]
The function \(F_0\) is continuous on compact \(\Sigma\)
and has a common finite bound. It is integrable against
the probability measure \(\rho\). Dominated convergence,
the mixture identity at each fixed \(m\), and strong heat
convergence in the mixed space therefore give
\begin{align}
 \int_\Sigma\|P_{0,\sigma}^{\rm OS}[b]_\sigma\|^2\,\dd\rho(\sigma)
  &=\lim_{m\to\infty}\int_\Sigma F_m(\sigma)\,\dd\rho(\sigma)\notag\\
  &=\lim_{m\to\infty}\langle[b],e^{-mH}[b]\rangle
   =\|P_0[b]\|^2.
 \label{eq:phase-ground-integral}
\end{align}
Every equality concerns the already fixed hierarchy and its
phase decomposition. The heat limit follows, rather than
interchanges with, the two ordered physical limits.

The compression and Gelfand laws also give
\[
 \int_\Sigma|S^\sigma(b)|^2\,\dd\rho(\sigma)
       =\|\mathsf M^0(b)\Omega\|^2=\|P_0[b]\|^2.
\]
Since \(\Omega_\sigma^{\rm OS}\) is a unit ground vector,
\[
 \langle\Omega_\sigma^{\rm OS},
           P_{0,\sigma}^{\rm OS}[b]_\sigma\rangle=S^\sigma(b).
\]
Expanding the squared distance to its vacuum component yields
\begin{equation}
 \begin{split}
 &\|P_{0,\sigma}^{\rm OS}[b]_\sigma
                       -S^\sigma(b)\Omega_\sigma^{\rm OS}\|^2\\
 &\hspace{12mm}=\|P_{0,\sigma}^{\rm OS}[b]_\sigma\|^2
                                              -|S^\sigma(b)|^2\ge0.
 \end{split}
 \label{eq:phase-ground-defect}
\end{equation}
It is a measurable function with integral zero by
\eqref{eq:phase-ground-integral}. Hence it vanishes outside
a null set which can initially depend on \(b\).

We now remove that dependence with a single countable choice.
For each degree use a fixed finite real basis of labels.
Take all finite positive words with these basis labels at
distinct rational points with positive fourth coordinate,
and include the empty word. This is a countable family
\((b_j)_{j\ge0}\), independent of \(\sigma\).
Multilinearity and Gram continuity show that its complex
linear span is dense in every \(\mathcal K_\sigma\):
first approximate the finitely many positions within their
positive separated neighbourhoods, then expand every label
in its basis. The common buffered bounds justify the Gram
limits for these fixed finite approximations.

Intersect the full-measure sets supplied by
\eqref{eq:phase-ground-defect} for all \(b_j\), obtaining
one measurable set \(E\subseteq\Sigma\) with \(\rho(E)=1\).
For \(\sigma\in E\), the projection of the dense span
of the \([b_j]_\sigma\) lies in
\(\C\Omega_\sigma^{\rm OS}\). The projection is bounded
and this line is closed, so the same holds on the whole
Hilbert space. Conversely the unit vector
\(\Omega_\sigma^{\rm OS}\) belongs to its range.
Thus \(\ker H_\sigma^{\rm OS}=\C\Omega_\sigma^{\rm OS}\)
for every \(\sigma\in E\).

Choose \(\sigma_*\in E\) once. The set is nonempty
because \(\rho\) is a probability measure. The choice is
made using only the Euclidean heat ground space. The later
reconstruction, OPE and trace calculations all use this
same character, the same mixed hierarchy \(S\), and the
same original Wilson family. Auxiliary coefficient limits
do not replace any of these objects. No claim about the
ground-space dimension of an exceptional character is
required here.
\end{proof}

All phase-algebra and vacuum-selection limits take place after
the original ordered physical limits. Constants in the new
variance approximation may depend on its fixed word, tolerance
and coarse stopping levels. They are uniform in cutoff tails
and period, and are fixed before the number of variance copies.
The later phase bounds are uniform over the compact phase
space. No energy gap, clustering rate, or uniqueness for every
exceptional character was needed for this selection.
% !TeX root = 4DYM_notes.tex
\section{Positive-height continuation and tempered boundaries}
\label{sec:wightman-continuation}

We use the positive-buffer method of
\cite[Section~17.2]{ArmstrongDraft}. Its useful feature here
is that the separated Euclidean hierarchy need not first be
extended as a distribution through its collision diagonals.
The boundary distribution will instead be constructed from a
holomorphic function with an explicit common polynomial bound.

\subsection{One label space and its Lorentz action}

Let \(\mathcal L\) be the algebraic direct limit of the
injective inclusions \(\jmath_{d,e}:\mathcal F_d\to\mathcal F_e\).
Explicitly, \((d,A)\) and \((e,B)\) represent the same label if
their images agree in some common higher degree. Addition is
performed after embedding into such a degree. Write
\(\iota_d A\) for the resulting label. Then \(\iota_d\) is
injective and \(\mathcal L=\bigcup_d\iota_d\mathcal F_d\).
Every finite list lies in one finite-dimensional normed space;
no norm on the entire union is needed.

Write \(s_4=\operatorname{diag}(1,1,1,-1)\). The adjoint
involution on labels is
\begin{equation}
                 (\iota_dA)^*=\iota_d(R_d(s_4)\bar A).
 \label{eq:wightman-label-star}
\end{equation}
Reality, \(s_4^2=I\), and compatibility of inclusions show
that this is a well-defined conjugate-linear involution.

For precision we construct the Lorentz action without a
Lie-algebra integration assumption. The free tensor-jet action
is algebraic on the complex orthogonal group: act on every
curvature and derivative index by the complex orthogonal
matrix. Orthogonal contractions and the free equation ideal
are preserved, so this descends to \(V_d^{\rm Max}\).
Conjugate this action by the real intertwiner
\(U_{k,d}f_k\) in \eqref{eq:rotation-actual-action}.
It gives a holomorphic extension \(R_d^{\mathbb C}\) of \(R_d\)
on the connected complex orthogonal group.
Different choices agree there: they agree on real plane
rotations, the one-variable identity theorem gives agreement
on their complex parameters. Products of the six plane
subgroups contain a neighbourhood of the identity: the
differential of their product map at zero is the basis of
antisymmetric matrices, so the inverse function theorem
applies. That neighbourhood generates the connected group.
The same argument extends inclusion compatibility.

Put
\[
 \eta=\operatorname{diag}(-1,-1,-1,1),\qquad
 C=\operatorname{diag}(1,1,1,-i).
\]
If \(\Lambda\in SO^+(1,3)\), then
\((C\Lambda C^{-1})^{\mathsf T}(C\Lambda C^{-1})=I\).
It belongs to the connected complex orthogonal group.
Define
\begin{equation}
 R_{\rm M}(\Lambda)\iota_dA
       =\iota_dR_d^{\mathbb C}(C\Lambda C^{-1})A.
 \label{eq:wightman-lorentz-labels}
\end{equation}
This is a continuous representation on every finite degree.
For the boost generator \(B_i=E_{i4}+E_{4i}\), and
\(K_{4i}=E_{4i}-E_{i4}\),
\[
 CB_iC^{-1}=-iK_{4i},\qquad
          \mathrm dR_{\rm M}(B_i)=-iL_d(K_{4i}).
\]
These equalities fix the boost sign in what follows.
Moreover
\(\overline{C\Lambda C^{-1}}
       =s_4(C\Lambda C^{-1})s_4\).
Using the real extension of \(R_d\), this proves that
\eqref{eq:wightman-label-star} commutes with \(R_{\rm M}\).
Translations act trivially on labels.

\subsection{Strictly ordered imaginary times}

Fix a phase \(\sigma\), and use its reflection space
\(\mathcal K_\sigma\). Abbreviate its vacuum by \(\Omega_\sigma\)
in this and the next two sections. Its heat and spatial
translation operators are \(e^{-tH_\sigma^{\rm OS}}\) and
\(U_y^\sigma\). At this stage their joint spectrum is only known
to have nonnegative energy.

For a label \(a\in\mathcal L\), let
\[
 B_{\sigma,\epsilon}(a)
     =\mathsf B_{2\epsilon}^\sigma(a(\epsilon e_4)).
\]
The notation \(a(x)\) denotes a one-letter Euclidean word.
Equation~\eqref{eq:phase-buffered-operator} gives
\(\|B_{\sigma,\epsilon}(a)\|\le C_a\epsilon^{-q_a}\),
uniformly in \(\sigma\), \(0<\epsilon\le1\).
For a complex time difference \(z\), with
\(\operatorname{Im}z>2\epsilon\), put
\[
 G_{\sigma,\epsilon}(z,y)=
     e^{-(\operatorname{Im}z-2\epsilon)H_\sigma^{\rm OS}}
       e^{i\operatorname{Re}zH_\sigma^{\rm OS}}U_y^\sigma.
\]
For \(r\ge2\), ordered imaginary times, and real spatial
points, define
\begin{equation}
 F_r^\sigma(a;z,\mathbf x)=
 \left\langle\Omega_\sigma,B_{\sigma,\epsilon}(a_1)
  \prod_{j=2}^r
   \left[G_{\sigma,\epsilon}(z_j-z_{j-1},
                \mathbf x_j-\mathbf x_{j-1})
                      B_{\sigma,\epsilon}(a_j)\right]\Omega_\sigma\right\rangle,
 \label{eq:wightman-positive-height}
\end{equation}
where \(4\epsilon<\min_j\operatorname{Im}(z_{j+1}-z_j)\).
Set \(F_0^\sigma=1\); the one-point function is the constant
\(S_1^\sigma(a)\).

At \(z_j=i\tau_j\), with \(\tau_1<\cdots<\tau_r\), the
insertion identities identify \eqref{eq:wightman-positive-height}
with \(S_r^\sigma(a;(\mathbf x_j,\tau_j)_j)\).
Thus the expression is independent of \(\epsilon\) there.
For fixed \(\epsilon\) it is holomorphic in all time
differences: the heat factor is
\(e^{izH_\sigma^{\rm OS}}e^{2\epsilon H_\sigma^{\rm OS}}\)
on its indicated bounded domain. Repeated one-variable
identity theorems extend the agreement between different
\(\epsilon\)'s throughout their common ordered domain.
Consequently \(F_r^\sigma\) is well defined, holomorphic in
the complex times, and continuous in the real spatial
coordinates.

At fixed degree, arity and coefficient ball there are finite
\(C,m\), independent of \(\sigma\), such that
\begin{equation}
 |F_r^\sigma(a;t+iy,\mathbf x)|\le C\delta^{-m},
 \quad 0<\delta\le1,\quad y_{j+1}-y_j\ge\delta.
 \label{eq:wightman-height-bound}
\end{equation}
Take \(\epsilon=\delta/8\); all intermediate heat factors
are contractions, all real translations are unitary, and
the product of the \(r\) insertion bounds proves the claim.
The bound is uniform over all real positions and times.
Fixed spatial derivatives have the same type of bound with
their differentiated labels; Cauchy's estimate gives time
derivative bounds on smaller ordered domains.
Constants depend on the fixed finite request and calendar,
and have no cutoff or volume parameter.

\subsection{An elementary boundary-value calculation}

For a joint Schwartz test
\(f\in\mathcal S((\mathbb R^4)^r)\) define
\[
 \Psi_\sigma(y;f)=
       \int F_r^\sigma(a;t+iy,\mathbf x)f(\mathbf x,t)
                                                   \,\dd\mathbf x\,\dd t.
\]
The integral is absolutely convergent at every strict height
by \eqref{eq:wightman-height-bound}.
Put \(\rho_j=j\), \(1\le j\le r\), and
\(\mathcal A=-i\sum_j\rho_j\partial_{t_j}\).
Holomorphy and integration by parts give
\[
       \frac{\dd}{\dd u}\Psi_\sigma(y+u\rho;f)
                               =\Psi_\sigma(y+u\rho;\mathcal A f).
\]
Choose \(m\) at least the exponent in
\eqref{eq:wightman-height-bound}. Repeated integration on
\([0,1]\) gives the exact identity
\begin{equation}
 \Psi_\sigma(y;f)=
 \sum_{k=0}^m\frac{(-1)^k}{k!}
                \Psi_\sigma(y+\rho;\mathcal A^kf)
   -\int_0^1\frac{(-u)^m}{m!}
                \Psi_\sigma(y+u\rho;\mathcal A^{m+1}f)\,\dd u .
 \label{eq:wightman-boundary-formula}
\end{equation}
For ordered \(y\), consecutive gaps at \(y+u\rho\) are at
least \(u\), so the last integrand is bounded by
\(C\|\mathcal A^{m+1}f\|_{L^1}/m!\), independently of \(y,\sigma\).

The right side at \(y=0\) therefore defines a continuous
linear functional \(W_r^\sigma(a;f)\). For clarity use the
Schwartz seminorm
\[
 p_{\mathsf M}(f)=\max_{|\alpha|\le {\mathsf M}}
             \sup_{x\in(\mathbb R^4)^r}
                    (1+|x|)^{\mathsf M}|\partial^\alpha f(x)|.
\]
For \({\mathsf M}\ge4r+m+2\) the calculation gives
\[
 |W_r^\sigma(a;f)|
      \le C\sum_{k=0}^{m+1}\|\mathcal A^kf\|_{L^1}
                                                  \le C' p_{\mathsf M}(f).
\]
Both constants are uniform in the phase and the fixed
finite coefficient ball. Dominated convergence proves
\begin{equation}
 \lim_{\substack{y\to0\\y_1<\cdots<y_r}}
          \Psi_\sigma(y;f)=W_r^\sigma(a;f).
 \label{eq:wightman-joint-boundary}
\end{equation}
The convergence is for arbitrary strictly ordered approaches,
and includes product tests as a special case. No pointwise
Minkowski boundary value is asserted.
Different choices of \(m,\rho\) give the same functional
by this boundary limit. Set \(W_0^\sigma=1\).

\begin{proposition}[A common phase boundary family]
\label{prop:wightman-tempered-boundary}
The \(W_r^\sigma\) are multilinear in their labels and are joint
tempered distributions. For every fixed finite label request
they obey one Schwartz seminorm bound uniform in \(\sigma\).
Their evaluations are continuous in \(\sigma\), and
\[
       W_r^{\rm mix}(a;f)
                     =\int_\Sigma W_r^\sigma(a;f)\,\dd\rho(\sigma).
\]
The mixed continuation is obtained from the original reflection
space by the same formula.
\end{proposition}
\begin{proof}
Only phase dependence needs further justification.
For fixed spatial points and \(z_j=iy_j+\varsigma t_j\),
the common bound makes the functions of \(\varsigma\)
a normal family in a strip about any compact admissible
segment. If \(\sigma_l\to\sigma\), every convergent subsequence
has the same values on a small imaginary \(\varsigma\)
interval, by continuity of the Euclidean phase values.
The identity theorem makes its limit \(F^\sigma\).
Finite overlapping strips and repeated variables give
continuity throughout the strict-time domain.
The same argument, or Cauchy's integral formula and Fubini,
extends the Euclidean mixture identity to \(F\).
The common integrable bound in
\eqref{eq:wightman-boundary-formula} then proves both assertions
for \(W\). Temperedness and multilinearity were proved in
the boundary calculation.
\end{proof}
% !TeX root = 4DYM_notes.tex
\section{Poincare covariance, the spectral cone and locality}
\label{sec:wightman-spectrum-locality}

The continuation constructed above is initially holomorphic
only in times. We first derive covariance and the momentum
support, then obtain the full spacetime tube from Fourier
inversion. This order avoids assuming the spectral cone in
the Euclidean ground-space argument. The method follows
\cite[Section~17.2]{ArmstrongDraft}.

\subsection{The boost identity with its sign fixed}

Differentiate full Euclidean covariance in the
\((4,i)\)-plane. At ordered imaginary times the infinitesimal
identity is
\[
 \sum_j S(\ldots,L(K_{4i})a_j,\ldots)
 +\sum_j\{x_j^i\partial_{\tau_j}
                         -\tau_j\partial_{x_j^i}\}S=0 .
\]
Under \(\tau=-iz\), \(\partial_\tau=i\partial_z\);
multiply this identity by \(-i\).
The time identity theorem and spatial differentiation under
the buffered operators extend it to strict complex times.
For a joint Schwartz test this says
\begin{equation}
 \sum_j\Psi_\sigma(y;f;\ldots,-iL(K_{4i})a_j,\ldots)
 =\Psi_\sigma\left(y;\sum_j
       [x_j^i\partial_{t_j}+(t_j+iy_j)\partial_{x_j^i}]f\right).
 \label{eq:wightman-boost-ward}
\end{equation}
The boundary formula bounds its test evaluations uniformly
as \(y\to0\). Hence the \(iy_j\) terms tend to zero.
This gives the infinitesimal real boost identity for \(W^\sigma\).
Along a real boost path the derivative of the distribution
with simultaneously transformed labels and tests is therefore
zero. Spatial rotations and all translations retain their
Euclidean identities. Boosts and spatial rotations generate
\(SO^+(1,3)\) (send the final unit vector to its future unit
image by a boost, then rotate its spatial orthogonal
complement). We have proved
\begin{equation}
 W_r^\sigma(R_{\rm M}(\Lambda)a;\,f\circ g^{-1})
                              =W_r^\sigma(a;f),\qquad
                 g x=\Lambda x+b .
 \label{eq:wightman-distribution-covariance}
\end{equation}
This is full proper orthochronous Poincare covariance.

\subsection{Momentum support from the ordered heat damping}

Use the ordinary dual pairing
\(p\cdot x=\mathbf p\cdot\mathbf x+p^4t\).
For a frequency test \(b\) set
\[
 \widehat b_+(x)=\int e^{ip\cdot x}b(p)\,\dd p,\qquad
       \widetilde W_r^\sigma(b)
                    =W_r^\sigma(\widehat b_+(-\,\cdot\,)).
\]
The factor in inverse Fourier formulas is
\((2\pi)^{-4r}\). For \(b\in C_c^\infty\), differentiation
of a height and its Fourier multiplier cancels, by the
time Cauchy--Riemann equation. Integrate this derivative
and include its zero endpoint using the boundary formula:
\begin{equation}
 \widetilde W_r^\sigma(b)=
     \Psi_\sigma\left(y;
       \widehat{\,e^{\sum_j y_jp_j^4}b\,}_+(-\,\cdot\,)\right).
 \label{eq:wightman-laplace}
\end{equation}
Compact frequency support justifies every differentiation.

Translation covariance first gives support in
\(\sum_jp_j=0\). Put \(q_j=\sum_{k=j}^rp_k\).
If \(b\) is supported where \(q_j^4\le-\alpha<0\),
increase all heights with index at least \(j\) by \(T\),
retaining fixed positive gaps elsewhere.
The Fourier multiplier in \eqref{eq:wightman-laplace} then
has all needed seminorms bounded by
\(C(1+T)^{\mathsf M} e^{-\alpha T}\). The strict-height bound,
whose minimum gap is fixed, makes the right side tend to
zero. A partition of unity proves \(q_j^4\ge0\) on the
support, for every \(j\ge2\).

Apply \eqref{eq:wightman-distribution-covariance} to every
Lorentz transformation. The finite-dimensional label action
is invertible. The common support of the finitely many
components in a fixed label tensor space must therefore
satisfy every Lorentz transform of these energy half-space
conditions. A covector has nonnegative time component in all
future Lorentz frames exactly when \(p^4\ge|\mathbf p|\):
if the latter fails, a boost in its spatial direction makes
its energy negative. Thus
\begin{equation}
 \operatorname{supp}\widetilde W_r^\sigma\subseteq
 \mathcal V_r:=
 \left\{p:\sum_jp_j=0,\quad
               q_j\in\overline V_+\ (j\ge2)\right\},
 \qquad \overline V_+=\{p:p^4\ge|\mathbf p|\}.
 \label{eq:wightman-momentum-support}
\end{equation}
The Lorentz action on covectors is the dual action for the
displayed pairing; it preserves this same forward cone.

\subsection{The full spacetime tube and partial boundaries}

Choose a nonnegative smooth unit-integral bump \(\rho_0\)
supported in the unit frequency ball, and put
\[
 \chi=\mathbf1_{\{\operatorname{dist}(p,\mathcal V_r)<2\}}
                                                        *\rho_0 .
\]
Then \(\chi=1\) near \(\mathcal V_r\), has bounded derivatives
of every fixed order, and is zero at distance greater than
three. If complex spacetime points satisfy
\(\operatorname{Im}(z_j-z_{j-1})\in V_+\), where
\(V_+=\{y:y^4>|\mathbf y|\}\), define
\begin{equation}
 \mathcal W_r^\sigma(a;z)
       =(2\pi)^{-4r}\widetilde W_r^\sigma
                 \left(a;\chi(p)e^{i\sum_jp_j\cdot z_j}\right).
 \label{eq:wightman-full-tube}
\end{equation}
This is a legitimate Schwartz test. On \(\mathcal V_r\),
\[
 \sum_jp_j\cdot\operatorname{Im}z_j
       =\sum_{j\ge2}q_j\cdot\operatorname{Im}(z_j-z_{j-1})
                                                   \ge c|p|,
\]
for \(z\) in a fixed compact subset of this tube.
Here the dual pairing of two future covectors/vectors in the
chosen plus-sign Fourier convention is positive: its time
term exceeds the absolute value of the spatial term.
On \(\operatorname{supp}\chi\) the bound holds with a fixed
subtraction. Differentiating either \(z\) or \(p\) only adds
polynomial factors, which the exponential absorbs.
Thus \(\mathcal W\) is holomorphic in all spacetime variables.
Formula~\eqref{eq:wightman-laplace} identifies its complex-time
restriction with \(F\), by compact frequency cutoffs and
Schwartz convergence. The definition is independent of \(\chi\).

There are also continuous partial boundaries as any subset of
the consecutive imaginary differences tends to zero through
the closed cone. To verify this on joint tests, multiply
their Fourier transforms by the corresponding damping
factors and by \(\chi\). On the closed cone the factors have
absolute value at most one; on its fixed neighbourhood they
have a fixed bound. All their derivatives have polynomial
growth uniform on compact height sets. Schwartz convergence
as the heights vary, followed by the common distribution
bound, proves the assertion. This argument permits some
heights to be tied while others remain strict.

\subsection{Locality by an imaginary boost}

\begin{lemma}[Spacelike interchange]
\label{lem:wightman-spacelike-interchange}
The joint distributions \(W_r^\sigma\) are unchanged by
interchanging two adjacent labels and their variables on
the open set where those two variables are spacelike
separated.
\end{lemma}
\begin{proof}
First tie the two adjacent imaginary heights, keeping all
other successive differences strictly positive and temporal.
Their real difference is \((v,t)\), with \(|t|<|v|\).
An imaginary boost of angle \(\beta>0\) along \(v/|v|\)
changes the imaginary part of this difference to
\[
           \sin\beta\,(t\,v/|v|,\ |v|)\in V_+.
\]
Its opposite treats the exchanged ordering. On any fixed
compact set of such tied configurations, choose \(\beta\)
small enough that all the other strict cone inequalities
are retained.

Work in a finite degree containing all labels, so that the
correlation is a holomorphic tensor-valued function.
Let \(\rho_r^\vee(\Lambda)\ell(a)
=\ell(R_{\rm M}(\Lambda)^{-1}a)\) be the dual tensor action.
For the appropriate sign of \(\beta\), the function
\[
       G_\beta(z)=\rho_r^\vee(\Lambda_{i\beta})^{-1}
                                  \mathcal W_r^\sigma(\Lambda_{i\beta}z)
\]
is holomorphic in a neighbourhood of the tied configuration.
It agrees with that ordering's original tube function on
their overlap: real boost covariance, continued in the
one complex boost parameter, proves the agreement.

We spell out why the germs for the two orders agree.
Scale all real time coordinates to zero, retaining the
imaginary heights and spatial coordinates. The adjacent
separation remains spacelike along this compact path;
the other height gaps remain strict. At its terminal point,
the Wick slice has real spatial coordinates and purely
imaginary time coordinates. The imaginary boost is then
a real Euclidean rotation, because
\[
            C\Lambda_{i\beta}C^{-1}
                                      =\exp(\beta K_{4i}).
\]
On an open piece of this slice the two germs both equal
the same smooth Euclidean correlation, by permutation
symmetry and full orthogonal covariance. The slice is
maximally totally real: each complex coordinate ranges
over one real line, and no nonzero complex direction
is contained in its tangent space. Repeated one-variable
identity theorems therefore identify the two holomorphic
germs in a neighbourhood. Cover the compact path by
finitely many overlapping admissible neighbourhoods.
The same identity theorem carries the equality back to
the original tied configuration.

Finally use the continuous partial boundaries proved above,
then let the other heights vanish. Testing on compact
subsets of the spacelike open set gives equality of the
boundary distributions. A partition of unity and Schwartz
cutoffs give the assertion for all tests supported there.
This argument uses no Euclidean distribution at collisions.
\end{proof}

All estimates used here are uniform in the phase at fixed
finite label request. They are continuum estimates obtained
from the earlier cutoff- and volume-uniform buffered bounds.
The spectral cone was derived from ordered heat damping and
Lorentz covariance; it was not an added hypothesis.
% !TeX root = 4DYM_notes.tex
\section{The word domain and the selected Wightman theory}
\label{sec:wightman-reconstruction}

We now turn the boundary distributions into actual smeared
operators on one Hilbert space. Positivity and comparison
with the Euclidean reflection space are proved before
vacuum uniqueness. We continue to follow the buffered
reconstruction of \cite[Section~17.2]{ArmstrongDraft}.
For comparison with reconstruction from nonlocal gauge-invariant
observables, see \cite[Sect.~8(b)--(e), printed pp.~167--179;
Theorems~8.8--8.10, p.~175, and Theorems~8.13--8.14,
pp.~178--179]{Sei82}.
These locators follow the original-page markers in the transcription;
the corresponding original scan pages have not been checked here.
This is a framework comparison: the buffered reconstruction for the
local fields of this manuscript is proved below.

\subsection{Positive boundary words inside the reflection space}

For strictly increasing positive heights \(y\), choose
\(\epsilon<y_1/2\) and \(4\epsilon<y_{j+1}-y_j\).
The following vector is well defined in \(\mathcal K_\sigma\):
\begin{align}
 V_y^\sigma(a;x)={}&
 e^{-(y_1-\epsilon)H_\sigma^{\rm OS}}
       e^{it_1H_\sigma^{\rm OS}}U_{\mathbf x_1}^\sigma
 B_{\sigma,\epsilon}(a_1)\notag\\
 &\quad\cdot\prod_{j=2}^r
    [G_{\sigma,\epsilon}(t_j-t_{j-1}
                          +i(y_j-y_{j-1}),
                        \mathbf x_j-\mathbf x_{j-1})
                        B_{\sigma,\epsilon}(a_j)]\Omega_\sigma .
 \label{eq:wightman-height-vector}
\end{align}
For \(r=0\) it is \(\Omega_\sigma\).
At zero real times it is the Euclidean positive-word vector.
The insertion identities and the vector identity theorem
make it independent of admissible \(\epsilon\).
For fixed heights it is bounded and norm continuous in \(x\),
so \(V_y^\sigma(a;f)=\int V_y^\sigma(a;x)f(x)\,\dd x\) exists
as a Hilbert-space integral for every joint Schwartz test.

Apply the adjoint rule to two such vectors. The negative
half has the first word's order reversed, conjugated tests,
starred labels, and imaginary heights \(-y\), in reversed
order. The positive half has the second word's heights \(y'\).
Their Gram entry is exactly the corresponding
positive-height correlation \(F^\sigma\) integrated against
that joint test. In particular, as \(y,y'\to0\) independently,
\begin{equation}
 \langle V_y^\sigma(b),V_{y'}^\sigma(c)\rangle
                                \longrightarrow W^\sigma(b^*c).
 \label{eq:wightman-vector-gram-limit}
\end{equation}
Here \(b,c\) denote tested words; the star reverses the word
and applies \eqref{eq:wightman-label-star} and test conjugation.
The reflected identity is first checked at imaginary times;
holomorphy in each time parameter proves it for the displayed
real times. The joint boundary theorem justifies the
integrated limit.

Taking \(b=c\), and subtracting the two mixed Gram entries,
shows that these vectors are Cauchy. Denote the limit by
\(v_\sigma(b)\). For any finite tested word combination,
\begin{equation}
                W^\sigma(b^*b)=\|v_\sigma(b)\|^2\ge0.
 \label{eq:wightman-word-positivity}
\end{equation}
The same holds for joint tests by the common seminorm bound.
One way to justify product-test density without a support
assumption is to expand a Schwartz function in tensor
products of one-variable Hermite functions. Integration by
parts with the harmonic oscillator gives coefficients
decaying faster than every polynomial; the finite tensor
sums converge in every Schwartz seminorm.

\subsection{The algebraic word domain}

Let \(\mathcal B\) consist of finite complex sums of words
\((a_1,f_1)\cdots(a_r,f_r)\), modulo multilinearity in
each \(a_j\in\mathcal L\), \(f_j\in\mathcal S(\mathbb R^4)\).
The empty word is the unit; multiplication concatenates.
Define
\[
 ((a_1,f_1)\cdots(a_r,f_r))^*
             =(a_r^*,\bar f_r)\cdots(a_1^*,\bar f_1).
\]
Thus \(\mathcal B\) uses algebraic sums of product tests.
The joint distributions constructed above extend its
expectations, but no completed test space is inserted into
the algebraic field domain.

Define \(\langle b,c\rangle_\sigma=W^\sigma(b^*c)\) and
\(\mathcal I_\sigma=\{b:\langle b,b\rangle_\sigma=0\}\).
Positivity gives Schwarz's inequality, so a null word
pairs to zero with every word.
Left multiplication preserves nullity: for any \(a,c\),
\(\langle ab,c\rangle=\langle b,a^*c\rangle=0\);
take \(c=ab\). The quotient
\(\mathcal D_\sigma=\mathcal B/\mathcal I_\sigma\) is therefore
a pre-Hilbert space and a left word domain.
Complete it to \(\mathcal H_\sigma^{\rm W}\), and let
\(\Omega_\sigma^{\rm W}=[1]\).
Its norm is one and the domain is dense.

Define the smeared field on this common domain by
\begin{equation}
             \Phi_\sigma(a,f)[b]=[(a,f)b].
 \label{eq:wightman-field-definition}
\end{equation}
The domain is invariant, contains the vacuum, and is exactly
the span of finite field words applied to that vacuum.
For \(u,v\in\mathcal D_\sigma\),
\[
       \langle\Phi_\sigma(a,f)u,v\rangle
                   =\langle u,\Phi_\sigma(a^*,\bar f)v\rangle.
\]
Hence \(\Phi_\sigma(a^*,\bar f)\) is contained in the Hilbert
adjoint of \(\Phi_\sigma(a,f)\); in particular the operators
are closable. For fixed \(a,v\), the common joint
distribution bound gives a finite Schwartz seminorm \(p_{\mathsf M}\)
with
\[
        \|\Phi_\sigma(a,f)v\|^2
            =W^\sigma(v^*(a^*,\bar f)(a,f)v)
                                      \le C_{a,v}p_{\mathsf M}(f)^2.
\]
This is temperedness as an operator-valued distribution on
the common word domain. Constants may depend on the fixed
state word and labels. They are uniform in phase when those
words are taken from a fixed finite list.

\subsection{Covariance, locality and the spectral measure}

For \(g=(\Lambda,b)\) let \(g\) act on a letter by
\[
          (a,f)\longmapsto
                 (R_{\rm M}(\Lambda)a,\ f\circ g^{-1}).
\]
The star commutes with this action. Distributional covariance
preserves the word form, so it descends to an isometry on
\(\mathcal D_\sigma\), and \(g^{-1}\) gives its inverse.
It extends to a unitary \(\mathsf U_\sigma(g)\).
The representation law holds on words and hence on the
completion. On a fixed finite label space the action is
continuous, and pullback of Schwartz tests by affine maps
is continuous in their parameters. The norm of a difference
of two transformed word vectors is a finite sum of continuous
Gram entries. This proves strong continuity on the dense
word domain and then on the Hilbert space. The vacuum is
fixed, the domain is invariant, and
\begin{equation}
 \mathsf U_\sigma(g)\Phi_\sigma(a,f)v
      =\Phi_\sigma(R_{\rm M}(\Lambda)a,f\circ g^{-1})
                                      \mathsf U_\sigma(g)v.
 \label{eq:wightman-operator-covariance}
\end{equation}

If every point of \(\operatorname{supp}f\) is spacelike
to every point of \(\operatorname{supp}g\), then
Lemma~\ref{lem:wightman-spacelike-interchange} makes every
word matrix element of
\([\Phi_\sigma(a,f),\Phi_\sigma(b,g)]\) vanish.
Density of word vectors proves that this commutator is zero
on \(\mathcal D_\sigma\). Schwartz tests follow by cutoffs;
their compact truncations retain spacelike separation.
This is smeared locality.

The four commuting strongly continuous translation groups
have a joint projection-valued spectral measure
\(\mathsf E_\sigma\), normalized by \(\mathsf E_\sigma(\mathbb R^4)=I\):
\[
        \mathsf U_\sigma(b)=\int e^{ip\cdot b}\,\dd\mathsf E_\sigma(p).
\]
This is the joint spectral theorem for strongly commuting
unitary groups, applied to their four Stone generators.
Equivalently, integrate their commuting spectral projections
first on rectangles and use countable additivity; the
one-variable spectral integrals determine the group and
its unique joint measure.
For word vectors, their translation matrix coefficients
are Fourier transforms of the word distributions with a
tail momentum at the cut between the two words.
Equation~\eqref{eq:wightman-momentum-support} makes that
momentum lie in \(\overline V_+\). Polarization and density
therefore give
\begin{equation}
                      \mathsf E_\sigma(\overline V_+^c)=0.
 \label{eq:wightman-pvm-cone}
\end{equation}
Indeed the positive scalar measure of this open complement
is zero on every word vector; the corresponding bounded
projection is consequently zero. Its diagonal measures
satisfy the actual characteristic-function identity
\[
  \int e^{ip\cdot b}\,\dd\langle v,\mathsf E_\sigma(p)v\rangle
                           =\langle v,\mathsf U_\sigma(b)v\rangle.
\]
Write \(H_\sigma^{\rm W}\ge0\) for its time generator.

\subsection{The Euclidean comparison isometry}

The Gram identity \eqref{eq:wightman-vector-gram-limit}
defines an isometry
\[
        J_\sigma:\mathcal H_\sigma^{\rm W}\longrightarrow\mathcal K_\sigma,
                         \qquad J_\sigma[b]=v_\sigma(b).
\]
It maps vacuum to vacuum and intertwines spatial translations.
We prove heat intertwining without a surjectivity assertion.
For tested words \(b,c\), the positive-height Gram identity
with one additional gap \(s\) between its two halves, followed
by their boundary limits, gives
\[
 \langle J_\sigma c,e^{-sH_\sigma^{\rm OS}}J_\sigma b\rangle
                  =\langle c,e^{-sH_\sigma^{\rm W}}b\rangle.
\]
On the left use norm convergence of the height vectors.
On the right the Laplace identity and the translation spectral
measure identify the surviving damping as \(e^{-sp^4}\)
at the cut between the words. Compact frequency cutoffs
justify the calculation first; the finite positive spectral
measures and \(0\le e^{-sp^4}\le1\) give its limit.
Thus \(J_\sigma^*e^{-sH_\sigma^{\rm OS}}J_\sigma
=e^{-sH_\sigma^{\rm W}}\) on the whole Hilbert space.

This compression identity is enough for actual intertwining.
For \(v\in\mathcal H_\sigma^{\rm W}\), apply it at \(2s\):
\[
 \|e^{-sH_\sigma^{\rm OS}}J_\sigma v\|^2
                 =\|e^{-sH_\sigma^{\rm W}}v\|^2.
\]
Apply it at \(s\) to the mixed pairing with
\(J_\sigma e^{-sH_\sigma^{\rm W}}v\).
That pairing is the same squared norm. Subtracting the two
vectors therefore gives zero norm, proving
\begin{equation}
       J_\sigma e^{-sH_\sigma^{\rm W}}
                             =e^{-sH_\sigma^{\rm OS}}J_\sigma,\qquad s\ge0 .
 \label{eq:wightman-heat-intertwining}
\end{equation}

\subsection{The chosen vacuum and faithfulness}

\begin{theorem}[Selected Wightman reconstruction]
\label{thm:selected-wightman-reconstruction}
For the fixed phase \(\sigma_*\) of
Theorem~\ref{thm:phase-vacuum-selection}, the Hilbert space
\(\mathcal H_{\sigma_*}^{\rm W}\), word domain
\(\mathcal D_{\sigma_*}\), vacuum, fields
\eqref{eq:wightman-field-definition}, Poincare representation
and spectral measure satisfy W0--W3. More explicitly they
have a normalized cyclic vacuum, common dense invariant
domain, tempered fields with the stated adjoints, full proper
orthochronous covariance, the forward spectral cone, and
spacelike locality. The translation-invariant vacuum is
unique up to scalar. The label action is faithful, each
\(\iota_d\) is injective, and these images exhaust
\(\mathcal L\). The strict-time continuation has Euclidean
root \(S^{\sigma_*}\) and the boundary vacuum moments of this
one theory.
\end{theorem}
\begin{proof}
Only uniqueness and faithfulness remain.
If \(v\in\ker H_{\sigma_*}^{\rm W}\), then
\eqref{eq:wightman-heat-intertwining} puts \(J_{\sigma_*}v\)
in the selected OS ground space, which is
\(\mathbb C\Omega_{\sigma_*}^{\rm OS}\).
Since the isometry maps vacuum to vacuum,
\(J_{\sigma_*}v=cJ_{\sigma_*}\Omega_{\sigma_*}^{\rm W}\)
implies \(v=c\Omega_{\sigma_*}^{\rm W}\).
Thus the full Wightman zero-energy space is a vacuum line.
A translation-invariant vector
has zero energy: its scalar spectral measure satisfies
\(\int|e^{itp^4}-1|^2\,\dd\mu=0\) for every rational \(t\),
so is supported at \(p^4=0\). It is therefore a scalar
multiple of the vacuum.

If a label acts as zero for every test on the word domain,
all vacuum word distributions containing it vanish.
Product-test density, the Laplace identity, and Fourier
inversion imply that all its strictly time-ordered
Euclidean correlations vanish. Permutation and smooth
continuity extend this to every separated configuration.
Equation~\eqref{eq:phase-equal-null-spaces} then makes its
class in the common field quotient zero. Thus the label
action is faithful. Direct-limit injectivity and coverage
were established in Section~\ref{sec:wightman-continuation}.
The unit label acts as \(\Phi(1,f)=(\int f)I\), so the
field system has a nonzero smeared operator.
The root and boundary statements are respectively the
imaginary-time identity and
\eqref{eq:wightman-joint-boundary}, now with the constructed
word expectations. The preceding arguments prove all the
other asserted properties.
\end{proof}

The construction above actually gives faithful covariant
positive-energy local fields for every character. For the
stated theorem, vacuum uniqueness is needed only at the one
selected character. Selecting it in Euclidean space first,
and then using the heat-intertwining isometry, avoids an
additional phase-uniform low-energy argument.
No positive mass gap follows or is used.
% !TeX root = 4DYM_notes.tex
\section{The continued stress tensor and local charges}
\label{sec:wightman-stress}

We retain the normalized Euclidean tensor already constructed.
Continuing its label and shell identities is sufficient;
no new tensor or ultraviolet normalization is chosen.
The conservation-based jump argument is adapted from
\cite[Section~17.2]{ArmstrongDraft}.

Write \(\zeta_a=C_{aa}\), so \(\zeta_i=1\) for \(i\le3\)
and \(\zeta_4=-i\). Define
\[
         T_{ab}^{\rm M}=-\zeta_a\zeta_bT_{ab},\qquad
         T_{ab}^{\rm M}(f)
                   =\Phi_\sigma(\iota_4T_{ab}^{\rm M},f).
\]
In particular
\(T_{44}^{\rm M}=T_{44}\),
\(T_{4i}^{\rm M}=iT_{4i}\), and
\(T_{ij}^{\rm M}=-T_{ij}\).

\begin{proposition}[Minkowski stress and its local action]
\label{prop:wightman-stress-charges}
The \(T_{ab}^{\rm M}\) form a symmetric Hermitian Lorentz
tensor of tempered fields on the common word domain, and
\begin{equation}
           \sum_a\eta_{aa}T_{ab}^{\rm M}(\partial_af)=0.
 \label{eq:wightman-stress-conservation}
\end{equation}
Let \(K(x)=k+Vx\), with \(V^{\mathsf T}\eta+\eta V=0\),
and let \(\Sigma_K=\mathrm dR_{\rm M}(V)\) be its
infinitesimal label action. For real
\(f\in C_c^\infty(\mathbb R)\), \(\int f=1\), and
\(c\in C_c^\infty(\mathbb R^3)\), put
\[
           Q_K(f,c)=\sum_bT_{4b}^{\rm M}(K^b(c\otimes f)).
\]
Suppose \(m\in C_c^\infty(\mathbb R^4)\), and \(c=1\)
on a ball such that all points over its complement at
times in \(\operatorname{supp}f\) are spacelike to
\(\operatorname{supp}m\). Then on the common domain,
\begin{equation}
 i[Q_K(f,c),\Phi_\sigma(a,m)]
       =\Phi_\sigma(\Sigma_Ka,m)
                       -\sum_b\Phi_\sigma(a,\partial_b(K^bm)).
 \label{eq:wightman-local-charge}
\end{equation}
These statements hold in particular in the one selected
Wightman theory.
\end{proposition}
\begin{proof}
The Euclidean derivative identity and the time
Cauchy--Riemann equation continue to the boundary as
\[
 \Phi_\sigma(\iota_{d+1}\partial_a A,f)
                   =-\zeta_a^{-1}\Phi_\sigma(\iota_dA,\partial_af).
\]
One verifies this on arbitrary word matrix elements and
uses their density. Euclidean conservation and
\(\eta_{aa}\zeta_a^2=-1\) prove
\eqref{eq:wightman-stress-conservation}.
The Euclidean tensor law extended through
\eqref{eq:wightman-lorentz-labels} gives the Lorentz law.
Finally \(R(s_4)T_{ab}\) multiplies \(T_{ab}\) by the
product of its time-index signs; this cancels conjugation
of \(\zeta_a\zeta_b\). Thus
\((T_{ab}^{\rm M})^*=T_{ab}^{\rm M}\) as labels.
The adjoint inclusion for fields proves Hermiticity on
the domain. Essential self-adjointness is not needed here.

We supply the jump argument for the charge formula.
In Euclidean coordinates \(x_{\rm E}=(\mathbf x,\tau)\)
put
\[
 K_{\rm E}(x_{\rm E})=Ck+CVC^{-1}x_{\rm E}.
\]
This is a complex Euclidean Killing vector field:
its linear part is antisymmetric. The current
\(J_a=\sum_bK_{\rm E}^bT_{ab}\) is conserved, because
the stress is symmetric and conserved. Near a distinguished
point \(x_0\), its normalized Euclidean shell action is
\[
       \Sigma_K a+\sum_bK_{\rm E}^b(x_0)\partial_ba .
\]
This follows by decomposing \(K_{\rm E}\) into its value
at \(x_0\) and its linear part and using the two phase
shell laws. Complex linearity permits the displayed
complex Killing field.

First take real \(c\). Keep all other Euclidean time
coordinates separated from \(\tau_0\). Integrate the
conservation equation with spatial multiplier \(c\) on
each side of \(\tau_0\). At a spatial boundary its derivative
is supported away from the distinguished point; the
correlations there are uniformly bounded on the resulting
separated compact sets. The spatially integrated current
therefore has limits from both sides. Its difference,
later time minus earlier time, is exactly the shell action:
one can obtain it by inserting a smooth temporal step
function in the shell identity and letting its transition
width tend to zero. Away from the transition conservation
cancels the bulk; the derivative of \(c\) gives the side
boundary; the only enclosed marked point is \(x_0\).

To continue this jump, fix strict heights \(y_j\) for all
noncurrent insertions, including \(y_0\) for the distinguished
one, and real times \(t_j\). Set
\(z_j(\varsigma)=iy_j+\varsigma t_j\).
For \(\epsilon>0\), insert the current immediately before
or after the distinguished field at time
\(z_0(\varsigma)\pm i\epsilon\). Define
\[
 I_\epsilon^\pm(\varsigma)=
 \sum_b\int c(x)K_{\rm E}^b(x,-iz_0(\varsigma)\pm\epsilon)
       F_\pm^\sigma(\ldots,T_{4b}(x,z_0(\varsigma)\pm i\epsilon),
                                               \ldots)\,\dd x .
\]
Take a common complex neighbourhood of \([0,1]\) retaining
the other strict height margins. On \(\operatorname{supp}\nabla c\)
the current and distinguished insertion have positive spatial
separation. The spacelike germs constructed in
Lemma~\ref{lem:wightman-spacelike-interchange} extend the
correlation across their tied height there.
They and their fixed derivatives are bounded on each compact
subset, uniformly for small \(\epsilon\), and uniformly in
the phase. Conservation, integrated in \(x\), gives
\[
 \sup_{\varsigma\in L}
     |I_\epsilon^\pm(\varsigma)-I_{\epsilon'}^\pm(\varsigma)|
                                       \le C_L|\epsilon-\epsilon'|.
\]
Indeed differentiating the current plane leaves only
\(-\sum_i\int(\partial_i c)J_i\); derivatives of
\(K_{\rm E}\) cancel by its Killing property.
The common germ bound controls this integral.
Thus the two limits exist locally uniformly and are
holomorphic. For small imaginary \(\varsigma\) they have
the Euclidean jump just proved. The identity theorem
identifies their jump at \(\varsigma=1\).

There is a useful sign check. On the Wick slice,
\[
 \sum_bK^bT_{4b}^{\rm M}
                      =i\sum_bK_{\rm E}^bT_{4b}.
\]
The earlier-current ordering is \(Q\Phi\), the later one
is \(\Phi Q\). Therefore \(i[Q,\Phi]\) is the later-minus-earlier
Euclidean current jump. The continued right side is
\(\Sigma_Ka+K^b\partial_b a\), with Minkowski derivatives.

For general current time in \(\operatorname{supp}f\), interpret
the commutator of the spatially integrated current as a
distribution in that time. Conservation expresses its
distributional derivative by an integral with \(\nabla c\),
and locality makes that derivative zero. It is therefore
a constant distribution on the allowed time interval;
the continued jump determines that constant.
The support hypothesis ensures spacelike separation along
the intervening time interval as well, since the allowed
times relative to any fixed spatial separation form an
interval. Averaging with \(f\), of integral one, preserves
the jump. Test the distinguished point against \(m\)
and take the continuous partial boundaries of the other
variables. Integration by parts converts \(K^b\partial_ba\)
into \(-\partial_b(K^bm)\). This proves
\eqref{eq:wightman-local-charge} on compact spectator tests.
The common joint Schwartz bounds extend it to word
spectators, hence to the domain.

If \(c\) is complex, its imaginary part vanishes on the
specified ball. Its charge commutes with the distinguished
field by locality and the same support hypothesis. Apply
the real case to its real part. This completes the proof.
\end{proof}

The constants in the jump calculation depend on the fixed
supports, their spacelike and height margins, the finite
degree request and finitely many derivatives of \(c,f,m\).
The common buffered and boundary bounds make them independent
of the phase before selection. No cutoff, period or
additional Gaussian normalization is introduced in this
continuation. The all-orders OPE and the trace analysis
remain separate steps of the construction theorem.
% !TeX root = 4DYM_notes.tex
\section{One coefficient prescription for every finite OPE order}
\label{sec:ope-coefficients}

We now prove the separated expansion stated in
Section~\ref{sec:continuum}. The coefficients must be fixed
before the spectators, the phase and the truncation degree.
A local matching field chosen separately for each estimate
would not meet that requirement. We therefore define the
coefficients directly from the already fixed mixed hierarchy
\(S\), using finitely many correlation measurements in each
degree. The lattice matching fields constructed later are
only tools for estimating this one prescription.

We use the finite-correlation projection of
\cite[Definition~19.3 and Lemmas~19.8--19.10,
PDF pp.~222--230]{ArmstrongDraft}. For the separated statement
we can keep its compatible finite-probe coordinates directly.
The ensuing proof requires no further symmetrization of the
projection or recursion on nested clusters.

\subsection{Finite probes and compatible coordinates}

We identify \(\mathcal F_d\) with its injective image in
\(\mathcal L\); all sums below can equivalently be written
with the inclusions \(\jmath_{d,e}\).
Put \(\mathcal F_{-1}=0\).
A probe word \(B\) is a fixed finite list of fields at
distinct nonzero points. The empty probe is allowed.
It defines a linear functional, on every finite degree, by
\[
                         A\longmapsto S(A(0),B).
\]
These functionals separate \(\mathcal F_d\), by the
all-spectator definition of \(\mathcal N_d\).
Their span is therefore its full dual: otherwise their
common kernel in this finite-dimensional space would
contain a nonzero field.

We make all choices real. The real subspace consists of
fields fixed by the existing conjugation. It spans the
complex space. Real probe words give real values on real
fields; expanding arbitrary words into real fields shows
that these probes separate the real subspace as well.
Thus each finite linear system below can be solved over
\(\mathbb R\).

\begin{lemma}[Compatible finite correlation coordinates]
\label{lem:ope-probe-coordinates}
For each \(j\ge0\) there are a finite set \(I_j\), real fields
\(O_{j,\lambda}\in\mathcal F_j\), \(\lambda\in I_j\), and
linear functionals \(\ell_{j,\lambda}\) given by finite real
combinations of fixed mixed probe correlations, such that
\begin{align}
 \mathcal F_D&=\bigoplus_{j=0}^D
                  \operatorname{span}_{\mathbb C}
                         \{O_{j,\lambda}:\lambda\in I_j\},
                                                        \notag\\
 \ell_{j,\lambda}(O_{i,\mu})
                       &=\delta_{ji}\delta_{\lambda\mu},\qquad
 \ell_{j,\lambda}|_{\mathcal F_{j-1}}=0 .
 \label{eq:ope-probe-coordinates}
\end{align}
Each functional is continued to higher degrees by its same
correlation formula. We may take \(I_0=\{0\}\),
\(O_{0,0}=1\) and \(\ell_{0,0}(A)=S(A(0))\).
\end{lemma}
\begin{proof}
Suppose the fields and functionals below degree \(j\) have
been chosen. Their continued correlation formulas define
the retraction
\[
 \mathsf t_j:\mathcal F_j\longrightarrow\mathcal F_{j-1},
 \qquad
 \mathsf t_j A=\sum_{i<j,\lambda\in I_i}
                                \ell_{i,\lambda}(A)O_{i,\lambda}.
\]
On \(\mathcal F_{j-1}\) it is the identity. Hence
\(\mathcal E_j=\ker\mathsf t_j\) is a real complement to
\(\mathcal F_{j-1}\). Choose a real basis of \(\mathcal E_j\)
and call it \(O_{j,\lambda}\).
The old functionals vanish on this complement, since its
displayed lower basis vectors are independent.
The new dual coordinates on
\(\mathcal F_{j-1}\oplus\mathcal E_j\), zero on the first
summand, are finite real combinations of probe functionals
because those span the dual. Fix such combinations once.
At every later step continue them by those formulas.
The inclusion identities for correlations preserve their
previous restrictions. This proves the induction.
\end{proof}

Concretely, for fixed finite probe lists
\(B_{j,\lambda,\beta}\) and fixed real numbers
\(m_{j,\lambda,\beta}\),
\begin{equation}
 \ell_{j,\lambda}(A)
       =\sum_\beta m_{j,\lambda,\beta}
                             S(A(0),B_{j,\lambda,\beta}).
 \label{eq:ope-probe-functional}
\end{equation}
The same expression makes sense on any single field in
\(\mathcal L\). In particular it is bounded on every fixed
\(\mathcal F_t\). Since it annihilates \(\mathcal F_{j-1}\),
for \(t\ge j\) it satisfies
\begin{equation}
 |\ell_{j,\lambda}(A)|
     \le C_{t,j,\lambda}\,
                 \operatorname{dist}_{\mathcal F_t}
                                  (A,\mathcal F_{j-1}).
 \label{eq:ope-probe-quotient-norm}
\end{equation}
Here and below the distance uses the fixed quotient norm
of \(\mathcal F_t\), with the indicated lower space embedded
by \(\jmath\).

\subsection{Definition of the coefficient functions}

Let \(w=\prod_{i=1}^rA_i(x_i)\) be a separated word and
let \(o\) be its expansion centre. Set
\begin{equation}
 c^{j,\lambda}_{\mathbf A}(\mathbf x;o)
       =\sum_\beta m_{j,\lambda,\beta}
           S\left(\prod_iA_i(x_i-o),B_{j,\lambda,\beta}\right)
 \label{eq:ope-coefficient-definition}
\end{equation}
whenever the combined configuration in every term is
separated. Set this coefficient to zero outside that domain.
This convention defines it at all arguments. For every
fixed \(j,\lambda,K\), all sufficiently small clusters
\(|x_i-o|\le K{\varrho}\) avoid these finitely many nonzero probe
points, so the convention does not affect the required
small-radius estimates.

Define its finite output field by
\begin{equation}
             P_Dw=\sum_{j=0}^D\sum_{\lambda\in I_j}
                    c^{j,\lambda}_{\mathbf A}(\mathbf x;o)
                                                    O_{j,\lambda}.
 \label{eq:ope-raw-truncation}
\end{equation}
It is inserted at \(o\) when used in a correlation.
The maps are multilinear in the input fields, real under
conjugation, invariant under permutation of input slots,
and depend on positions through \(\mathbf x-o\).
On their open separated domains they are smooth.
These facts follow immediately from their finite correlation
formulas.

For a single field \(A\in\mathcal F_t\) at the centre,
\begin{equation}
 P_D[A(o)]=A\quad(D\ge t),\qquad
                  P_D P_H[A(o)]=P_D[A(o)]\quad(D\le H).
 \label{eq:ope-compatible-decoder}
\end{equation}
Indeed \eqref{eq:ope-probe-coordinates} is its basis
reconstruction. Thus increasing the truncation degree never
changes an already chosen coefficient.
For the empty cluster the same formulas give the unit in
degree zero and zero coefficients in positive degrees.

\subsection{The estimate to be proved}

Write \(W=\sum_i d_i\) when \(A_i\in\mathcal F_{d_i}\).
For a spectator word \(B=\prod_bB_b(y_b)\), put
\begin{equation}
 \mathcal R_D(w;B)=\mathsf M^0(wB)-\mathsf M^0((P_Dw)(o)B)\in\mathcal A_0.
 \label{eq:ope-algebra-remainder}
\end{equation}
Here \(\mathsf M^0\) and its operator norm are those of
Section~\ref{sec:phase-compression}; notation for a sum
of words is read linearly. Its character evaluations are
the phase remainders, and its vacuum state is the mixed
remainder.

We shall prove, uniformly over the stated cluster shapes
and over all spectator positions with fixed clearance,
\[
 |c^{j,\lambda}_{\mathbf A}(\mathbf x;o)|
                         \le {\varrho}^{j-W-\varepsilon},
 \qquad
       \|\mathcal R_D(w;B)\|
                         \le {\varrho}^{D+1-W-\varepsilon}
\]
for sufficiently small \({\varrho}\). The second estimate in the
operator norm will give all phases with one threshold.
Every \(D\) is finite; the coefficient prescription is
already fixed for all of them.
The remaining work is analytic: local matching, its
degree-sensitive quotient bound, and control of the full
marked source as the cluster points coalesce.
% !TeX root = 4DYM_notes.tex
\section{Signed matching and the degree-sensitive quotient estimate}
\label{sec:ope-signed-graded}

Two earlier estimates must be used with particular care.
Subtracting a selected column creates a signed local source;
its remainder is controlled before absolute values are
summed. Moreover a coefficient bound in \(V_t\) alone loses
the output degree when converted to the physical quotient.
The following two calculations address these points. They
adapt \cite[Lemmas~19.1--19.2, PDF pp.~219--222]{ArmstrongDraft}.

\subsection{Signed continuation on one unmarked history}

Fix an admitted output degree \(t\). At level \(i\) write
\(\Gamma_i(x)\) for the selected dimensionless one-insertion
column with polynomial coordinate \(x\in V_t\).
The physical insertion has the prefactor \({\varepsilon}_i^{-t}\).
Its remainder norm is
\[
 \|g_i\|_{\rm ins}
      =E_i^{\rm ins}+w_i^{-1/2}K_i^{\rm ins},
\]
exactly \eqref{eq:marked-remainder-norm}.
Regular summands have their original masks and zero radial
jets through weight \(t\) on nonwrapping supports; forced
summands replace the original eligible gas factor.

Temporarily put \(C_i=C_{J(i)}\) and
\(K_i=K_{J(i+1)}\). The constants precede the ratio.
On a sufficiently late interval the local marked estimate
and the existing logarithmic ratio floors give
\(\eta_i\le L_i^{-2/5}\) and
\[
 1+2C_i\le L_i^{1/40},\qquad
 10^5K_i^6(1+C_{i+1})^2\le L_i^{1/8}.
\]
Indeed the logarithm of each displayed adjacent-profile
column product is \(o(\log L_i)\); the ratio is indexed
by its admitted fine profile. Fixed numerical factors are
also absorbed eventually. Thus these reserves hold on
the same calendar, independently of cutoff, volume and
the later cluster.
Set \(\epsilon_i=L_i^{-3/8}\).

A zero-coordinate remainder \(R\) of norm \(q\) has an
exact transported expression \(Z\) whose coefficient \(z\)
and remainder each have size at most \(\eta_iq\), by
Lemma~\ref{lem:insertion-local-prices}. The expression
\[
                         R'=Z-\Gamma_i(z)
\]
has coefficient zero and remainder at most
\((1+C_i)\eta_iq\). This subtraction keeps the full forced
replacement and its eligibility event. Linearity gives
the exact density identity; no estimate on an unsigned
gas has been substituted for it.

For a column plus \(R\), absorb \(z\) into the coarse
column parameter. The change of that column and its
remainder costs at most another \(C_i\eta_iq\).
Thus the extra price is
\((1+2C_i)\eta_iq\le\epsilon_iq\).
This proves that the same signed affine family obeys,
at every step,
\begin{equation}
 \begin{aligned}
 |u-T_i\mathsf S_{L_i}v|_1
                       &\le\epsilon_i(|v^\circ|_1+q),\\
 r&\le K_i|(\mathsf S_{L_i}v)^\circ|_1+\epsilon_iq .
 \end{aligned}
 \label{eq:ope-signed-step}
\end{equation}
The leading map \(T_i\) and the engineering diagonal
\(\mathsf S_{L_i}\) are those of the one-insertion step
in degree \(t\). A zero-coordinate input can have arbitrary
finite norm here; only the unmarked attachments must be
small. Finite truncations and their absolutely convergent
rooted majorants justify the termwise signed identity.

\begin{lemma}[A signed interval estimate]
\label{lem:ope-signed-continuation}
Suppose a local expression at \(m\) has combined coefficient
and remainder size at most \(q_{\rm fin}\), before subtracting
a selected column. On any sufficiently late admitted interval
\([k_*,m]\) there is a column subtraction whose residual has
coefficient zero at \(k_*\). Its coefficients \(x_i\) and
remainder bounds \(R_i\) satisfy
\begin{equation}
 R_i+|x_i|_1\le
     4(1+C_m)q_{\rm fin}({\varepsilon}_m/{\varepsilon}_i)^{1/8},
                     \qquad k_*\le i\le m.
 \label{eq:ope-signed-interval}
\end{equation}
All densities are transported with the same unmarked history:
\(H_i=L_i^t\mathcal P_iH_{i+1}\).
\end{lemma}
\begin{proof}
Continue the residual one step as above, always recentering
its zero coefficient by the selected column. The free
affine parameter is a reference coefficient for that column.
Its map to \(k_*\) is invertible, so it can be chosen to
make \(x_{k_*}=0\).

For completeness write \(U=|u|_1\),
\(Y=|\mathsf S_{L_i}v|_1\), and
\(A_i=32(1+C_i)^2\). Applying \(T_i^{-1}\) to
\eqref{eq:ope-signed-step} gives
\[
               Y\le K_i[U+\epsilon_i(Y+q)],
               \qquad r\le K_iY+\epsilon_iq .
\]
The ratio floors give
\[
 A_i\ge4K_i^2,\quad4K_i\epsilon_i\le1,\quad
 4K_iA_{i+1}\epsilon_i\le1,\quad
              (4K_i^2+1)L_i^{1/8}\epsilon_i\le1.
\]
If \(A_iU\le r\), then
\(4K_iU\le Y+\epsilon_iq\). Substitution gives
\(Y\le4K_i\epsilon_iq\). Consequently
\[
       A_{i+1}|v|_1\le q,\qquad L_i^{1/8}r\le q.
\]
Starting from the zero coordinate at \(k_*\), this is an
induction towards the fine end of the interval.
At the terminal end the original expression and the
subtracted column give
\[
 R_m\le(1+C_m)q_{\rm fin}+C_m|x_m|_1.
\]
The cone implies \(2C_m|x_m|_1\le R_m\), so
\(R_m\le2(1+C_m)q_{\rm fin}\).
Iterating the second cone inequality back to \(i\) gives
\eqref{eq:ope-signed-interval}. Every step was a linear
identity before the majorants were taken.
\end{proof}

\subsection{What an output coordinate means in the quotient}

For a reference label \(v\in V_t\), denote its selected
level-\(m\) polynomial coordinate by \(x_m(v)\), with its
dependence on the finite history understood. Write
\(|v|_m=|x_m(v)|_1\).
We use the fixed auxiliary sequence \((n_l,M_l)\) from
Section~\ref{sec:continuum-field-calculus}, along which
the equation projections converge to \(\Pi_t^\infty\).
This sequence samples the already constructed ordered
hierarchy. Subsequences used for auxiliary coefficient
limits do not define another hierarchy.

\begin{lemma}[Graded quotient estimate]
\label{lem:ope-graded-quotient}
Fix \(0\le i\le t\). There are constants \(C_{t,i}(m)>0\),
independent of auxiliary subsequences and label families,
with
\[
                     \log C_{t,i}(m)=o(s_m),
                         \qquad s_m=\log {\varepsilon}_m^{-1},
\]
such that the following holds. If bounded labels
\(v_l\in V_t\) converge to \(v\) along a subsequence of
the fixed auxiliary sequence and eventually
\(|v_l|_m\le e\), then
\begin{equation}
 \operatorname{dist}_{\mathcal F_t}
           ([v]_t,\jmath_{i-1,t}\mathcal F_{i-1})
              \le C_{t,i}(m){\varepsilon}_m^{-(t-i)}e .
 \label{eq:ope-graded-quotient}
\end{equation}
For \(i=0\), the lower space is zero.
\end{lemma}
\begin{proof}
All cone iterations below take place above one fixed,
sufficiently late admitted level. The finitely many earlier
invertible coordinate transfers are included in the final
fixed constants. First let \(i>0\). Let \(P\) be coordinate projection
in \(V_t\) onto engineering weights at most \(i-1\).
Leading-map triangularity says
\[
 (T_u)_{\alpha\beta}=0
       \quad(\alpha\ne\beta,\ w_\alpha\le w_\beta).
\]
Thus \(T_u\) and \(T_u^{-1}\) preserve \(\ker P\), and
\[
               |Px|_1\le K_u|PT_ux|_1.
\]
The diagonal \(\mathsf S_{L_u}\) expands \(Px\) by at
least \(L_u^{t-i+1}\), and its complementary part by at
most \(L_u^{t-i}\). The actual selected column map differs
from \(T_u\mathsf S_{L_u}\) by at most
\(\epsilon_u(1+C_{u+1})|x|_1\).

These facts imply the two elementary bounds
\begin{equation}
 \begin{aligned}
 |x_{u+1}|_1\le C_{u+1}|Px_{u+1}|_1
 &\ \Longrightarrow\
 \begin{cases}
 |x_u|_1\le C_u|Px_u|_1,\\
 L_u^{t-i+1}|Px_{u+1}|_1\le C_{u+1}|Px_u|_1 ,
 \end{cases}\\
 Px_{u+1}=0
 &\ \Longrightarrow\
       |x_u|_1\le C_{u+1}L_u^{t-i}|x_{u+1}|_1 .
 \end{aligned}
 \label{eq:ope-low-cone}
\end{equation}
Here is the absorption in the first line. Put
\(Y=|P\mathsf S_{L_u}x_{u+1}|_1\).
The hypothesis and \(C_{u+1}\le L_u\) give
\(|\mathsf S_{L_u}x_{u+1}|_1\le2Y\).
After projection, invert the leading map and absorb
\(2K_u\epsilon_u(1+C_{u+1})\le1/2\).
Then \(Y\le2K_u|Px_u|_1\), whereas
\(|x_u|_1\le3K_uY\). The column choice
\(64K_u^2\le C_u,C_{u+1}\) proves the first line.
For the second line use
\(|\mathsf S_{L_u}x_{u+1}|_1\le L_u^{t-i}|x_{u+1}|_1\)
and the same error bound.

At cutoff \(n\), transport the bare low plane \(PV_t\)
along the selected column maps. Its image is in the
first cone at every coarser level. Projection of this
image onto \(PV_t\) is injective by the expansion bound,
and hence onto by dimension. At level \(m\) subtract
the transported low-plane vector with the same low
coordinates as \(x_m(v_l)\). The residual now has zero
low coordinates and size at most \((1+C_m)e\).
Transport it one step, apply the second line of
\eqref{eq:ope-low-cone}, and subtract the low coordinates
again in the same transported plane. Continuing to the
reference level \(\operatorname{ref}(t)\) gives a vector \(y_l\) in the
transported bare low plane with
\begin{equation}
 |v_l-y_l|_1\le
   G_m(1+C_m)({\varepsilon}_{\operatorname{ref}(t)}/{\varepsilon}_m)^{t-i}e,\qquad
             G_m=\prod_{u<m}(1+C_u)C_{u+1}.
 \label{eq:ope-low-plane-distance}
\end{equation}
Each subtracted vector is in that same reference-level
plane; no new plane is chosen during the subtraction.
Since \(v_l\) is bounded, \(y_l\) is bounded at this
fixed \(m\).

We identify every limit \(y\) of this plane family as a
lower-degree physical field. Iterating the expansion
in \eqref{eq:ope-low-cone} shows that its bare coordinate
is bounded by
\[
       {\varepsilon}_n^{t-i+1}\exp(o(s_n))
\]
times a fixed bound for \(y_l\). It uses only weights
at most \(i-1\). Rewriting its physical prefactor
\({\varepsilon}_n^{-t}\) as \({\varepsilon}_n^{-(i-1)}\) leaves a fixed-stencil
local section with subpower coefficient, derivative
and global value prices. Proposition~\ref{prop:position-local-matching}
in degree \(i-1\) gives admissible lower-degree coefficients
\(c_l\) with a root-small comparison error.

Those coefficients need not be bounded. Put
\[
 \lambda_l=(1+|\Pi_{n_l,M_l,i-1}c_l|_1)^{-1},\qquad
                 z_l=\lambda_l\Pi_{n_l,M_l,i-1}c_l .
\]
The normalized equation component
\(\lambda_l(c_l-\Pi_{n_l,M_l,i-1}c_l)\) is admissible;
its original bare subpower bound is only decreased by
\(\lambda_l\le1\), and its projected component is bounded.
The exact equation section and root-small estimates
therefore make all its separated spectator responses
tend to zero. Pass to a further subsequence on which
\(\lambda_l\to\lambda\), \(z_l\to z\) and \(y_l\to y\).
The fixed projection limit gives
\[
 \Pi_{i-1}^\infty z=z,\qquad \lambda+|z|_1=1,\qquad
        S(z(0),B)=\lambda S(y(0),B)
\]
for every separated spectator word \(B\). The last formula
uses the appropriate inclusion into a common degree.
If \(\lambda=0\), then \(z\) is null but lies in the
range of \(\Pi_{i-1}^\infty\), so
Proposition~\ref{prop:continuum-null-kernel} gives \(z=0\).
This contradicts \(|z|_1=1\).
Hence \(\lambda>0\), and \([z/\lambda]_{i-1}\) has
exactly the correlations of \([y]_t\). Faithfulness and
the inclusion law give
\([y]_t\in\jmath_{i-1,t}\mathcal F_{i-1}\).

If necessary, extend the plane labels by zero outside
the set on which their fixed reference bound holds.
They remain in the transported bare low plane and
have a uniform bare subpower envelope. This permits
the local comparison estimates on the original family,
while retaining the chosen auxiliary limit.
The coefficient-ball convergence of the hierarchy and
the root-small estimates justify testing these varying
labels. Passing to a limit in
\eqref{eq:ope-low-plane-distance} now proves
\eqref{eq:ope-graded-quotient}.

For \(i=0\) there is no low plane. The inverse coordinate
bound iterated to the reference level gives directly
\[
       |v|_1\le
           \Big(\prod_{u<m}C_{u+1}\Big)
                            ({\varepsilon}_{\operatorname{ref}(t)}/{\varepsilon}_m)^t e.
\]
Finally the logarithms of both products of column constants
are \(o(s_m)\), by the calendar condition. Fixed reference
spacings and norm equivalences are absorbed in \(C_{t,i}(m)\).
This proves its asserted independence and subpower growth.
\end{proof}

The estimate uses convergence only along the previously
chosen auxiliary sequence, exactly where the equation
projections are known to converge. All finite-history
bounds preceding that limit are uniform in the original
cutoff and volume. The limiting mixed hierarchy has
not been replaced by an auxiliary diagonal limit.
% !TeX root = 4DYM_notes.tex
\section{The full marked source when a cluster coalesces}
\label{sec:ope-cluster-fusion}

Separated insertion estimates cease to apply directly when
several named anchors enter the same coarse neighbourhood.
We retain their original square-zero variables, move their
declared supports to nearby representatives, and pay for
that move by decreasing estimating weights. The actual
inserted functions and the physical measure are unchanged.
This is the grouping argument of
\cite[Lemma~19.8, PDF pp.~226--227]{ArmstrongDraft},
with the support and all-subset bookkeeping made explicit.

\subsection{Weights and physical source norms}

For a finite set of names retain
\(\mathcal A_I=\mathbb C[z_a:a\in I]/(z_a^2)\), and use
\[
       \Big\|\sum_A f_Az_A\Big\|_{\rho}
                   =\sum_A\|f_A\|\prod_{a\in A}\rho_a .
\]
For scalar coefficients this is
\eqref{eq:joint-marker-norm}. For regular and forced
families the coefficient norm means their original
anchored derivative and square-root-normalized forced
majorants, with their covering weights.
The forced normalization is charged once per actual
forced address. A mode with several names keeps that
one replacement and receives the product of its name
weights.
Products use the finite-jet product rule and the scalar
marker norm; the derivative seminorm itself is not being
declared an algebra norm.

During grouping a regular coefficient is allowed its full
radial Taylor polynomial. Its physical size includes the
original \({\varepsilon}_n^{-d}\) insertion prefactors. We postpone
the degree-\(t\) extraction to the common ancestor.
We call this a rough marked coefficient solely to indicate
that no radial jets have yet been removed. Its derivatives,
masks, read support and forced eligibility have exactly
the previous meanings.

All modes are retained, including proper subsets of the
names in a cluster and modes meeting more than one cell.
Grouping is a partition of the names used for estimates.
We never identify two \(z_a\)'s: that would incorrectly
kill their product in the square-zero algebra.

\subsection{One passage before Taylor subtraction}

For clarity we give the weaker passage bound sufficient
here. Let \(\mathcal L\) be an incoming distinguished
physical marked majorant at a fixed anchor. The same
component expansion as in
Lemma~\ref{lem:insertion-local-prices} can be used without
subtracting its radial polynomial.
Separate its unmasked constant unit term, which is
transported exactly. The other value term is the
stationary substitution of the whole remaining function.
Its threshold-scaled derivatives are
bounded by its incoming derivative norm times fixed
profile constants: the interpolation kernels have bounded
row sums, their differentiated rows have inverse spatial
powers, and the ratio of adjacent thresholds is bounded.
For the local Gaussian part the finite chain rule uses
at most the admitted input derivative order. Bounding a
difference by the two values suffices at the highest
order. Measure, cutoff and Haar derivatives are those
already estimated with the common precision reserve;
their remaining fixed-ratio errors tend to zero at entry.
Thus the local regular part has price \(C_j\mathcal L\)
with \(C_j\) preceding the ratio, after entry.

Here is the chain-rule calculation behind this bound.
Let \(\Psi\) be the existing fine chart map of a coarse
coordinate and the held Gaussian variable, and let
\(\operatorname{Part}(q)\) be the partitions of the
\(q\) labelled derivative directions. In the anchored
direction norms, each summand of the finite chain rule
is bounded by
\[
 h_c^q|D^q(f\circ\Psi)|
 \le\sum_{\pi\in\operatorname{Part}(q)}
       h_f^{|\pi|}|D^{|\pi|}f|
           \prod_{B\in\pi}
                    \frac{h_c^{|B|}}{h_f}|D^{|B|}\Psi|.
\]
The first map derivative uses the bounded interpolation
row sum; higher derivatives have the small nonlinear
threshold factors of the stationary graph. Their
polynomial Gaussian moments are bounded with the same
reserve as in
Section~\ref{sec:central-derivative-estimates}.
The finite partition sum is part of the admitted
profile constant. When derivatives hit the normalized
Gaussian measure, its retained scores are the previously
bounded physical-measure derivatives; their fixed-ratio
error is \(\omega_{j,L_i}(\widehat\beta_i)\to0\).
After entry it is at most one. Thus this calculation
uses derivatives of \(f\) through \(q\), and gives
the stated bound without a Taylor-vanishing hypothesis.

The unsubtracted incoming regular term requires its
support conversion as well. If \(s_f,s_c\) are its fine
and coarse connected covering numbers, including all
windows and the anchor, Lemma~\ref{lem:incoming-cover-compression}
gives constants \(K_j^{\rm supp},C_j^{\rm supp}\), before
the ratio, with
\[
        s_c\le K_j^{\rm supp}
                          (1+C_j^{\rm supp}s_f/L_i).
\]
Measuring the source at any fixed exponent
\(98\le\kappa\le100\), choose the two internal support
units as in the previous marked estimates. For the fixed
large ratio this implies
\begin{equation}
 (\kappa+2)s_c
       \le(\kappa+2)K_j^{\rm supp}+\tfrac12\kappa s_f
 \label{eq:ope-rough-cover-conversion}
\end{equation}
whenever the left exponent is among the available internal
estimates. Equivalently take \(\kappa=98\), use \(100\)
internally and store the \(98\)-norm again at the output.
The finite conversion factor is
\(\mathcal A_i=\exp((\kappa+2)K_j^{\rm supp})\).
The unused half of the fine exponent pays the remaining
incidence and anchor factors.

The combined profile constant for this term is at most
\(\mathcal A_iC_j\). These support-conversion and finite
chain-rule constants are included in the common marked
profile bound \(\mathcal C_j\). The logarithmic ratio floors for
the earlier incoming and marked estimates make
\(C_j,\mathcal A_iC_j\le L_i\) eventually on this same
calendar. For forced terms use their original packet
majorants with half the forced suppression retained.
Their normalized cost is at most \(C_j\mathcal L\).
Both the real/imaginary split and the rooted unmarked
descendants have already finite factors in \(C_j\).
For example, allowing a factor two on each regular
part gives the sufficient total bound
\begin{equation}
       2C_j\mathcal L+2\mathcal A_iC_j\mathcal L
                                  +C_j\mathcal L
                                      \le6L_i\mathcal L .
 \label{eq:ope-rough-one-step}
\end{equation}
No extra fine-anchor sum is present for a fixed marked
occurrence.

In a coefficient carrying \({\varrho}\) distinct names, at most
\({\varrho}\) marked occurrences can be multiplied: disjointness
of their name sets follows from \(z_a^2=0\).
The rooted expansion and its unmarked attachment sum
therefore give the conservative factor \((6L_i)^{\varrho}\).
Replacing each estimating weight by \(L_i^{-4}\) times
it pays this passage, since \(6L_i\le L_i^4\) for
\(L_i\ge3\). This proves a complete-source estimate,
including already generated higher modes. It is a
physical marked estimate without the artificial
degree normalization \(L_i^t\); that normalization is
introduced only at the extraction step below.

\subsection{Moving representatives and paying their incidence}

Suppose a marked support \(X\) contains its old anchors.
Joining each new representative to an old one through
at most \(h\) covering balls gives
\begin{equation}
 n_B(X\cup\{\text{new representatives}\})
                        \le n_B(X)+h|A|,
 \label{eq:ope-support-move}
\end{equation}
where \(A\) is the set of names in that coefficient.
The same bound includes the joining paths in the
estimating footprint. Each exact component still records
its original masks, inventories and compatibility tests.
If several gas factors must be expressed as one packet,
use the physical reassembly bijection: partition the
original eligible family into the components meeting
the enlarged packet and multiply their original factors
once, retaining their union inventory and all internal
tests. This is the reassembly of
Section~\ref{sec:physical-reassembly} with the enlarged
distinguished set.

Shrink every affected name weight by
\(\delta e^{-\kappa h}\), including the finite cost of
changing the derivative anchor. The exponential change
in \eqref{eq:ope-support-move} is then paid. In an
anchored sum a new incidence can be charged to the old
marked occurrence on its joining path. If a neighbourhood
contains at most \(b\) representatives and a cell contains
at most \(r\) original names, the multiplicity is at most
\((1+b)r\). Taking \(\delta\le((1+b)r)^{-1}\) pays it.
Polynomial factors from the finite-jet anchor change
are included in the same allowance.

Choose a maximal \(B_i^2\)-separated collection of current
representatives at each passage, and attach discarded
ones to retained neighbours. Previously merged groups
are never split. Within one \(r\)-point cell there are
at most \(r-1\) fusions. When two groups first become
too close after taking parents, the joining radius in
fine units must include \(L_i\), not just \(B_i\).
A sufficient bound for all joining-cover and anchor
costs uses
\[
          H_i=16[1+L_i(B_i+B_{i+1}+1)^2]^4 .
\]
This polynomially excessive bound avoids any assumption
on nonwrapping lifts.

Absorbing unmarked factors in this neighbourhood costs
at most \(\exp(C H_i)\) in the majorant:
the unmarked attachment income per site is bounded,
and its rooted descendants have the same convergent
sum as before. This cost is included at the fusion.
For finite address families the reassembly is a
literal partition of each old eligible family;
absolute convergence then passes it to all addresses.

Its logarithmic weight cost is subpower along the existing
calendar. Indeed \(B_{i+1}^2\le2B_i^2\), and
\eqref{eq:insertion-root-calendar} gives
\[
          (L_iB_i^2)^4Q_{i+1}^{\rm vac}=o(s_{i+1}).
\]
Since the support exponent is fixed and
\(Q_{i+1}^{\rm vac}\ge1\), this gives \(H_i=o(s_i)\).
An anchor polynomial is paid by one further copy of
\(H_i\) and its fixed profile constant. The logarithm
of that constant is \(o(s_i)\), by the same bound on
\(\mathcal C_j\); it is not multiplied by the number
of passages.
Gentleness gives \(s_{i+1}/s_i\to1\).
Consequently one may choose an allowance
\(\alpha_j\le1\), common to passages \(m\le i<j\), with
\begin{equation}
             \log\alpha_j^{-1}=o(s_j).
 \label{eq:ope-fusion-allowance}
\end{equation}
To verify the running minimum, its logarithm is bounded
by the maximum of these finite-passage costs. A fixed
initial segment divided by \(s_j\) vanishes; on a late
segment each cost is at most \(\tau s_i\le\tau s_j\).
Then let \(\tau\downarrow0\).
Only the finite number of fusions spends this allowance;
it is not spent once per level.

Apply this construction simultaneously to the reflected
array of cells. When cell centres are separated by much
more than \({\varepsilon}_iB_i^2\), at most \(r\) representatives
occur in a covering neighbourhood. All incidence bounds
are therefore independent of the number of cells and
of the torus period.
Periodic distances and full declared footprints are used
throughout.

\begin{proposition}[The available weight at the common ancestor]
\label{prop:ope-fusion-weight}
Suppose an \(r\)-point cluster starts at level \(j\) with
individual physical estimating weights \(\rho_a\).
At a common ancestor level \(m+1\), its full-name
coefficient has available product weight at least
\begin{equation}
 \omega\ge
   ({\varepsilon}_j/{\varepsilon}_{m+1})^{4r}
              (\alpha_j/8)^{r(r-1)}\prod_{a=1}^r\rho_a .
 \label{eq:ope-fusion-weight}
\end{equation}
All proper subsets and all intercell modes remain in
the same exact source representation.
\end{proposition}
\begin{proof}
At a passage each original name loses at most \(L_i^4\)
by \eqref{eq:ope-rough-one-step}; using all \(r\) names
is a valid overestimate after some have merged.
Their product over the interval is
\(({\varepsilon}_j/{\varepsilon}_{m+1})^{4r}\).
At a fusion decrease every weight in the merged groups.
There are at most \(r\) such original names, and at
most \(r-1\) fusions, giving the second factor.
The marker norm, the support-move bound and the small
complete attachment income close the induction on the
positive coefficient majorants. Constants for padding,
finite group assignments and conjugation are included
in the factor \(8\) and the allowance \(\alpha_j\).
This does not factorize the physical integral: it bounds
its complete coefficient expansion.
\end{proof}

The calculation uses one unchanged Gaussian law and its
common precision reserve. The vacuum scalar, noise and
Haar normalizations, incoming activities, masks and every
forced inventory stay in the exact density. A change of
representative enlarges a footprint for estimates and
does not change a branch choice or a physical cutoff.
% !TeX root = 4DYM_notes.tex
\section{A uniform lattice cluster replacement}
\label{sec:ope-lattice-replacement}

The product weight from the preceding section bounds a
physical cluster at its common ancestor. We now extract
a degree-\(t\) local field there, continue its signed
residual, and estimate the replicated difference.
The calculation follows
\cite[Lemmas~19.7--19.8, PDF pp.~225--227]{ArmstrongDraft}.
We retain the intercell higher modes explicitly, as in
the signed birth of Section~\ref{sec:bounded-coarse-variance}.

\subsection{Extraction and its continuation weights}

Suppose the full-name physical cluster coefficient at
level \(m+1\) has available weight \(\omega>0\).
Its weighted regular derivative and normalized forced
majorants have already been bounded by the complete
source estimate. Set
\[
 \Xi_m=\exp\sqrt{\widehat\beta_{m+1}},\qquad
 q_m=\frac{4\Xi_m {\varepsilon}_{m+1}^t}{\omega},\qquad
                    E_m=8(1+C_m)q_m .
\]
After entry, \(\Xi_m\) bounds the finite Taylor-reading,
product, section and support-conversion constants, including
the factor \(L_m^t\) and their polynomial coupling losses.
The existing entry bounds for the admitted degrees make
every fixed power of \(L_m\) a coupling-polynomial price.
The logarithms of the remaining prices are
profile constants plus powers of
\(\log\widehat\beta_{m+1}\), whereas
\(\log\Xi_m=\sqrt{\widehat\beta_{m+1}}\).
All these logarithms are \(o(s_m)\).

Read the regular radial Taylor polynomial through weight
\(t\), and use its integral remainder for the regular
complement. Nonwrapping complements have zero jets
through \(t\). Wrapping terms retain the whole function
and their covering decay, exactly as in the earlier
marked reading. Forced terms retain their original
inventories and normalized majorants. Multiplication
by \({\varepsilon}_{m+1}^t\) puts this physical expression into
the dimensionless degree-\(t\) format. Its coefficient
and remainder are each at most \(q_m\).
One passage to \(m\), with its admitted ratio and
profile factors included in \(\Xi_m\), gives this same
bound as a combined bound up to a fixed factor two.

Choose a level-\(k\) reflected cell and let
\[
 k_*=\min\{l\ge k:B_l^2\le {\varepsilon}_k/{\varepsilon}_l\}.
\]
We take \(m\ge k_*\), sufficiently late for degree \(t\).
Subtract the selected column that cancels the residual
coordinate at \(k_*\). Lemma~\ref{lem:ope-signed-continuation}
gives, throughout \([k_*,m]\),
\begin{equation}
 R_l+|x_l|_1\le E_m({\varepsilon}_m/{\varepsilon}_l)^{1/8},\qquad
                  |x_m^{\rm out}|_1\le2E_m .
 \label{eq:ope-extraction-continuation}
\end{equation}
The output coefficient is calculated for the single
cluster before choosing any number of reflected copies.
Every copy uses the literal translated, signed or
conjugated origin coefficient.

Here is a sufficient choice of estimating radii for the
residual. Use the old allowance
\(\zeta_{t,l}=\widehat\beta_l^{-({\varepsilon}_t+4)}\) from
Section~\ref{sec:joint-insertion-bounds}. Its fixed degree
exponent pays the physical section and restoration
prices; the additional finite reading constants are
already included in \(\Xi_m\).
Starting at \(\varpi_{k_*}=\varpi\), set
\(\varpi_{l+1}=\widehat\beta_{l+1}^{-1}\varpi_l\).
It is sufficient that
\begin{equation}
 \begin{aligned}
 \varpi({\varepsilon}_m/{\varepsilon}_{k_*})^{t+1}
           \frac{\zeta_{t,m}}{\zeta_{t,k_*}}
             &\le\frac{\zeta_{t,m}\omega}{16\Xi_m},\\
 32(1+C_m)\varpi {\varepsilon}_{k_*}^{-t}E_m
                   ({\varepsilon}_m/{\varepsilon}_{k_*})^{1/8}
             &\le\zeta_{t,k_*}/2 .
 \end{aligned}
 \label{eq:ope-extraction-weights}
\end{equation}
The extra polynomial allowance in the first line is
harmless for the powers proved below and ensures that
the physical Taylor birth fits the original attachment
budget directly.

Indeed the already established entry inequality
\[
 \widehat\beta_{l+1}^{-1}L_l^{t+1}\zeta_{t,l}
                                               \le\zeta_{t,l+1}
\]
implies
\[
 \varpi_l\le\varpi({\varepsilon}_l/{\varepsilon}_{k_*})^{t+1}
                              \zeta_{t,l}/\zeta_{t,k_*}.
\]
The first line of \eqref{eq:ope-extraction-weights} pays
the extracted physical coefficient at birth. For the
continued signed residual, divide its required weight
by \(\zeta_{t,l}\), use
\eqref{eq:ope-extraction-continuation}, and obtain its
level-\(k_*\) value times
\[
                         ({\varepsilon}_l/{\varepsilon}_{k_*})^{7/8}\le1.
\]
The second line therefore pays the remainder and its
compensating column at every intermediate level.
This is a simultaneous bound along the same history.

\subsection{The higher-mode birth and the replicated root}

% !TeX root = 4DYM_notes.tex
We first write the two-cell calculation without a source
norm. This fixes which cancellations must survive the
subsequent estimates.

Let cells \(C,D\) carry two cluster inputs each, with
physical functions \(O_1,O_2\) and \(O_3,O_4\).
Let \(B_C,B_D\) be the actual selected output insertions,
using the single-cluster coefficient already chosen above.
They are the same functions for every number of copies.
Put \(R_C=O_1O_2-B_C\) and \(R_D=O_3O_4-B_D\).
For any spectator multiplier \(F\), apply
Lemma~\ref{lem:complete-subset-identity} at one configuration
space, with \(\mathcal P\) the identity and the two fine
functions \(O_1O_2,O_3O_4\). The exact numerator is
\begin{equation}
 F R_CR_D=F\{O_1O_2O_3O_4-B_CO_3O_4
                                  -O_1O_2B_D+B_CB_D\}.
 \label{eq:ope-two-cell-signed-numerator}
\end{equation}
Each term is transported as its own complete source.
Its slots are, respectively,
\[
 (1,2;3,4),\quad (C;3,4),\quad (1,2;D),\quad (C;D).
\]
The symbol \(C\) in a slot means its output insertion
alone. Thus an output may be at a cluster point without
requiring an output and that cluster's inputs to be
inserted simultaneously. These four expressions have
the same original density and partition denominator.
Their signs are combined before taking an absolute value.

Here is the bookkeeping that retains the proper subsets.
For example the first source uses commuting markers
\(\theta_1,\ldots,\theta_4\) with \(\theta_i^2=0\).
Before the full coefficient is read it contains every
subset monomial. If \(P,Q\) are two of its factor
polynomials, then
\[
 [\theta_1\theta_2\theta_3\theta_4](PQ)
        =\sum_{A\subset\{1,2,3,4\}}
                 [\theta_A]P\,[\theta_{A^c}]Q.
\]
This is \eqref{eq:common-subset-product}. In particular
the terms with names \(1,2\) in separate
parts are needed for the full coefficient. Replacing
\(\theta_1\theta_2\) by a new cell marker and dropping
the other modes at this point would lose these terms.
The other three sources use only their displayed slots,
with the same full-subset rule. No source with both
versions of one cell has been introduced.

To reorganise the completed coefficient, add one vertex
for each cell, and connect it to every original marked
part carrying a name of that cell. The resulting
connected components are uniquely determined. All
selected names of a cell are then in one component.
The two partitions of two original names are
\(\{\{1,2\}\}\) and \(\{\{1\},\{2\}\}\).
For an output slot there is one partition. Including
the two replacement choices costs a fixed constant
per cell. Unmarked descendants and selection queries
were already in the original components and are not
counted again. For at most \(r\) original names per
cell the partition number is bounded by a constant
depending only on \(r\). This gives \(C_r^{\mathsf M_{\rm cell}}\),
not a partition count depending on all names globally.
It is the grouping allowance of
Proposition~\ref{prop:ope-fusion-weight}.

Only after this regrouping introduce cell markers
\(\xi_C,\xi_D\), with \(\xi_C^2=\xi_D^2=0\).
They describe the exact physical polynomial
\((1+\xi_CR_C)(1+\xi_DR_D)\). It is formed by
combining the preceding complete sources, not by
projecting individual original factors. Its zero
mode is consequently the same unmarked density in
all four replacements, and its singleton modes are
the complete signed residuals. All these operations
are termwise regroupings of absolutely convergent
source series. Long output descendants keep their
actual footprints and covering weights; a descendant
meeting another cell is counted in the intercell
part of the calculation.
For a compatible finite component family \(\Gamma\),
write its factor as
\(a_{\gamma,0}+\xi_Ca_{\gamma,C}
 +\xi_Da_{\gamma,D}+\xi_C\xi_Da_{\gamma,CD}\).
The completed two-cell coefficient has exactly two
types of terms:
\begin{align}
 &\sum_{\gamma\in\Gamma}a_{\gamma,CD}
                             \prod_{\eta\ne\gamma}a_{\eta,0}
       \notag\\
 &\quad+\sum_{\substack{\gamma,\eta\in\Gamma\\\gamma\ne\eta}}
       a_{\gamma,C}a_{\eta,D}
                             \prod_{\nu\notin\{\gamma,\eta\}}a_{\nu,0}.
 \label{eq:ope-two-cell-component-coefficient}
\end{align}
This follows by selecting each marker once in the
product. It uses no division by a zero-mode amplitude,
which may vanish. Sum these identities over the
original allowed families. Enlarged footprints
determine their connected components, while original
supports still determine forced exclusion. All
inventories, coarse prohibitions and normalisations
remain attached to the corresponding terms.

In the second line, the signed replacements in
\eqref{eq:ope-two-cell-signed-numerator} give the
product of the two complete singleton residual
constructions. In the first line they give the
four-term difference for a component containing
both cells. Absolute values are taken only after
these sums. This is a factorisation within the
component expansion, not independence of physical
cell observables in the Wilson measure. The
unmarked gas and its exclusions remain in the
sum and in the partition denominator.

The same argument proves the assertion for any
finite set \(J\) of selected cells. The cell
vertices give a unique component partition of
\(J\); singleton blocks carry the already
cancelled one-cell residual. A block containing
two different cells has a full physical read
joining them. Joining parts inside one cell
does not create an edge between distinct cells.
With birth separation
\(\Delta_m=\varepsilon_k/\varepsilon_m\), the
doubled spanning-tree estimate of
Lemma~\ref{lem:joint-marked-cover} gives
\begin{equation}
 n_{B_m}(X)\ge c\,|J|\,\Delta_m/B_m,
                 \qquad |J|\ge2
                     \quad\text{for a connected cell mode}.
 \label{eq:ope-intercell-cover}
\end{equation}
Indeed an edge of a spanning tree between cell
representatives has length at least their separation.
A tour traverses each physical joining segment at
most twice; one connected \(B_m\)-ball cover can
pay at most \(C B_m\) of that length per ball.
For \(|J|\ge2\), \(|J|-1\ge|J|/2\).
The argument uses periodic distances, so its
constant is independent of the physical volume.

We check explicitly that higher modes fit the
same radii as the signed singleton. Suppose the
scale levels have the fixed positive separated
heights on \(T\to\infty\) used below. For a
fixed output degree \(t\), the output coefficient,
its physical section and every finite coordinate
conversion have price at most
\(\exp(C_tT+o(T))\). This follows from
\eqref{eq:ope-extraction-continuation}, the inverse
coordinate bound and the finite-horizon column
product. Each selected-slot version begins with
its ordinary positive radius, at least
\(\exp(-C_tT-o(T))\) at birth by the rough
passage bound. Increasing it to the prescribed
cell radius therefore costs at most
\(\exp(C_t'T+o(T))\) per cell. A smaller radius
at a finer level causes no enlargement cost.

After summing the fixed local partitions, a
connected mode on \(J\) cells consequently has
the additional price
\[
 \left\{C_r\exp(C_tT+o(T))
                         \exp(-c\Delta_m/B_m)\right\}^{|J|}.
\]
The separation \(\Delta_m\) grows exponentially
in \(T\), while \(B_m\) and the finite conversion
prices have subexponential logarithms. The last
display is smaller than every prescribed inverse
exponential per cell. Thus the intercell terms
in the first line of
\eqref{eq:ope-two-cell-component-coefficient} are
paid at the same radii as the second line.
For the singleton, use exactly the two inequalities
\eqref{eq:ope-extraction-weights}. Subsequent
transport regenerates the ordinary stored reserve
by the source attachment estimate, keeping all
proper subsets until their coefficient is formed.
No logarithm of a signed forced gas is used.

Finally we keep the volume cost visible through
the root. There are \(\mathsf M_{\rm cell}\)
reflected cells of physical side \(\varepsilon_k\).
The normalised source bound before that root is
\[
 \left|\left\langle\prod_C R_C\right\rangle_{n,M}\right|
   \le (C_{t,r}/\varpi)^{\mathsf M_{\rm cell}}
            \exp\{ \#\Lambda_{k_*}Q_{k_*}^{\rm vac}
                                           +C_{t,r}\mathsf M_{\rm cell}\}.
\]
Its denominator is the original positive partition
integral. The common vacuum scalar cancels there;
the scalar-free vacuum ratio of
Lemma~\ref{lem:insertion-vacuum-ratio} gives the
displayed remaining exponential. Since
\[
 \mathsf M_{\rm cell}=(M/\varepsilon_k)^4,
 \quad \#\Lambda_{k_*}=(M/\varepsilon_{k_*})^4,
 \quad \frac{\#\Lambda_{k_*}}{\mathsf M_{\rm cell}}
                          =(\varepsilon_k/\varepsilon_{k_*})^4,
\]
the physical volume cancels after the root. With
\(\mathcal S_k=\exp\{(\varepsilon_k/\varepsilon_{k_*})^4
Q_{k_*}^{\rm vac}+8\}\), we obtain
\begin{equation}
 \left\|\prod_{a=1}^r O_{n,M}(v_a,x_a)
                         -O_{n,M}(v^{\rm out}_{n,M},0)\right\|_k
                                \le C_{t,r}\mathcal S_k/\varpi.
 \label{eq:ope-extracted-root}
\end{equation}
The original minimality argument for
\eqref{eq:uniform-insertion-root} gives
\(\log\mathcal S_k=o(s_k)\) and
\(s_{k_*}/s_k\to1\). This proves the full-source
root bound with the same output coefficient,
incoming activities and physical normalisation.

\subsection{The lattice estimate in scale exponents}

We use only three relative scales for the fixed-clearance
OPE. Choose \(0<\gamma<1/4\), set \(T=-\log {\varrho}\), and
take the first calendar levels satisfying
\[
  s_k\ge\gamma T,\qquad
  s_m\ge(1-\gamma)T,\qquad
  s_j\ge(1+\gamma)T .
\]
Gentleness gives equality up to \(o(T)\), and the same
is true at \(m+1\). Thus \(k_*<m<j\) eventually.
An \(O({\varrho})\) cluster lies in the centred level-\((m+1)\)
block. On a generic spacing-\({\varepsilon}_j\) grid its separation
\(\eta {\varrho}\) exceeds \(B_j^2{\varepsilon}_j\) eventually.
Here generic means that its coordinate values in each
axis are distinct, the grid condition used in the
reflected-cell estimates.

For fixed real \(b\), write \(f\preccurlyeq e^{bT}\)
if every \(\tau>0\) gives \(f\le e^{(b+\tau)T}\)
eventually. For lattice quantities the bound is on
an ordered cutoff tail, uniform in its volumes; the
tail may depend on the fixed geometry and \(T\).
All scalar factors of logarithm \(o(T)\) are thereby
included. Finite sums and products preserve this notation
with the maximum and sum of their exponents.

\begin{theorem}[A same-family lattice cluster replacement]
\label{thm:ope-lattice-replacement}
Fix the input degrees and bounded coefficient balls,
the arity \(r\ge1\), shape bounds \(K,\eta>0\), the
auxiliary degree \(t\), and \(\gamma\) as above. Put
\[
                        Z=(1+\gamma)W+8r\gamma .
\]
For generic spacing-\({\varepsilon}_j\) clusters of extent at most
\(K{\varrho}\), separation at least \(\eta {\varrho}\), and centre zero,
there are output reference labels \(v^{\rm out}_{n,M}\in V_t\)
on the same Wilson family, with
\begin{align}
       |v^{\rm out}_{n,M}|_m
                          &\preccurlyeq e^{[Z-t(1-\gamma)]T},
                                                        \notag\\
 \left\|\prod_aO_{n,M}(v_a,x_a)
                   -O_{n,M}(v^{\rm out}_{n,M},0)\right\|_k
                          &\preccurlyeq
                             e^{[Z-(t+1/16)(1-2\gamma)]T}.
 \label{eq:ope-lattice-replacement}
\end{align}
The error is interior to its cell. These bounds are
uniform in the cutoff tail, period, translated cell
and the fixed input coefficient balls.
\end{theorem}
\begin{proof}
At level \(j\), the complete separated-source estimate
gives individual physical weights with
\(\prod_a\rho_a=\exp[-(1+\gamma)WT+o(T)]\).
It includes the higher modes formed between \(n\) and
\(j\). Proposition~\ref{prop:ope-fusion-weight} gives
\(\omega=\exp[-ZT+o(T)]\) as an available product
weight at the common ancestor. Hence
\[
                  E_m=\exp\{[Z-t(1-\gamma)]T+o(T)\}.
\]
Choose
\[
           \varpi=
            \exp\{[(t+1/16)(1-2\gamma)-Z]T\}.
\]
The exponent in the first left side of
\eqref{eq:ope-extraction-weights} is
\(-Z-\tfrac{15}{16}(1-2\gamma)\), whereas its right
side has exponent \(-Z\). All polynomial allowances
and \(\Xi_m\) are subpower. The exponent in the
second left side is \(-(1-2\gamma)/16\), whereas its
right side has exponent zero. Both inequalities
therefore hold with their strict slack.

The cell separation at birth is
\(\Delta_m=\exp[(1-2\gamma)T+o(T)]\).
This verifies the higher-mode birth condition used
above, including all counterterm-column prices.
Equation~\eqref{eq:ope-extraction-continuation} gives
the coefficient estimate, and
\eqref{eq:ope-extracted-root}, with
\(\log\mathcal S_k=o(T)\), gives the root estimate.
All stencil and ancestor margins are \(o({\varepsilon}_k)\).
The parameters were fixed before the physical limits;
the scalar and Gaussian/Haar normalizations were never
replaced. This proves the theorem.
\end{proof}
% !TeX root = 4DYM_notes.tex
\section{Uniform remainders and the all-orders separated OPE}
\label{sec:ope-separated-bounds}

The output field in
Theorem~\ref{thm:ope-lattice-replacement} depends on its
auxiliary scales and degree. We use it to estimate the
fixed coefficients of Section~\ref{sec:ope-coefficients}.
The degree-sensitive quotient estimate recovers the
power of each retained channel. Finally one finite-order
Taylor formula, using the exact derivative calculus,
removes the grid restriction.
This adapts
\cite[Lemmas~19.9--19.10, PDF pp.~228--230]{ArmstrongDraft},
using only fixed exterior clearance.

\subsection{From replicated roots to the phase-algebra norm}

We first spell out the conversion that makes the result
uniform in the phase. Fix a local lattice difference
\(E_{n,M}\), supported in a level-\(k\) cell, with roots
at most \(b\) for its signed and conjugated images.
Suppose its output reference coefficients are bounded
and converge along a subsequence of the fixed auxiliary
sequence. Let \(E\) be the corresponding difference of
continuum words. Fix a spectator list \(B\) with clearance
\(c>0\) from the cell centre and between its own points.

For \({\varepsilon}_k\) small compared with \(c\), choose one fixed
coarser tiling of side below \(c/(100(|B|+1))\).
Translate its boundaries away from the finitely many
anchors. The proof of
Lemma~\ref{lem:insertion-interior-root} puts the local
difference in its interior with the same root bound:
that proof uses only locality, gauge invariance and the
centred roots, so applies to this difference.
Take the spectators' centred roots at a finer fixed
aligned level to secure the same interior margins.
Their uniform root bounds give a product constant
\(C_{B,c}\), independent
of their positions. Stencil margins hold on the physical
cutoff tail.

Place \(2q\) translated copies of the whole pattern far
apart, alternating it with its conjugate. Their translations
can be integer multiples of the coarser cell side, so
the chosen boundary margins persist for all copies.
The Wilson chessboard estimate gives the bound
\((bC_{B,c})^{2q}\).
At each fixed copy number and translation, pass along
the auxiliary subsequence. Coefficient-ball convergence
of the ordered hierarchy identifies this limit with
the same mixed continuum correlations.
Next send the distances between copies to infinity
in \eqref{eq:phase-distant-clusters}. For
\(U=\mathsf M^0(EB)\), the result is
\[
             \int_\Sigma|\sigma(U)|^{2q}\,\dd\rho(\sigma)
                                       \le(bC_{B,c})^{2q}.
\]
Continuity, full support and \(q\to\infty\), exactly as
in \eqref{eq:phase-free-row-moments}, imply
\begin{equation}
                           \|\mathsf M^0(EB)\|\le bC_{B,c}.
 \label{eq:ope-root-to-algebra}
\end{equation}
All these limits concern the already fixed hierarchy.
In particular the auxiliary coefficient subsequence
does not select a new phase measure or a new physical
volume limit.

The same even-moment argument, now placing each insertion
in its own small cell, gives a useful crude bound.
For a fixed finite word of total degree \(H\), bounded
coefficients and mutual separation at least \(c_0{\varrho}\),
\begin{equation}
                       \|\mathsf M^0(\text{word})\|
                                          \preccurlyeq e^{HT},
                     \qquad T=-\log {\varrho} .
 \label{eq:ope-crude-word}
\end{equation}
Indeed the first sufficiently fine tiling has
\({\varepsilon}_l={\varrho}^{1+o(1)}\), by gentleness. Its individual roots
are \({\varepsilon}_l^{-d-\tau}\) for every fixed \(\tau>0\).
Choose boundaries with the finite margins and apply
chessboard to the repeated arrays before taking limits.
Letting \(\tau\downarrow0\) proves the exponent \(H\).
Constants depend on the finite label request and \(c_0\);
the bound is uniform over positions and phases.
The analogous bound with a fixed positive separation
has a fixed finite constant.

\subsection{Grid estimates for the fixed decoder}

Fix the input arity \(r\ge1\), total degree \(W\),
shape parameters \(K,\eta\), and the fixed spectator
labels and clearance \(c>0\). Translate the centre to
zero. Keep \(0<\gamma<1/4\), and abbreviate
\[
 \theta=1-\gamma,\qquad
              Z=(1+\gamma)W+8r\gamma,\qquad
              G_t=(t+1/16)(1-2\gamma).
\]

For a fixed generic spacing-\({\varepsilon}_{j(T)}\) configuration,
apply Theorem~\ref{thm:ope-lattice-replacement}.
At each fixed \(T\) its output reference coefficients
are bounded uniformly in \(n,M\). To see this directly,
iterate the inverse coordinate bound from \(m\) to
the fixed reference level: the factor is at most
\(\prod_{u<m}C_{u+1}({\varepsilon}_{\operatorname{ref}(t)}/{\varepsilon}_m)^t\).
It is finite at this fixed \(T\).
Choose a subsequence of the fixed auxiliary sequence
on which the coefficients converge, and call their
limit \(v^{\rm out}\), with class \(A^{\rm out}\in\mathcal F_t\).
The same subsequence still converges to \(S\) in every
fixed separated coefficient-ball test.
Equations~\eqref{eq:ope-lattice-replacement} and
\eqref{eq:ope-root-to-algebra} imply
\begin{equation}
       \|\mathsf M^0(wB)-\mathsf M^0(A^{\rm out}(0)B)\|
                                  \preccurlyeq e^{(Z-G_t)T}.
 \label{eq:ope-grid-auxiliary-error}
\end{equation}
Its constants are independent of the auxiliary subsequence.
The same estimate holds for each of the fixed probe words.

For \(i\le t\), apply
Lemma~\ref{lem:ope-graded-quotient} at level \(m\):
\[
 \operatorname{dist}_{\mathcal F_t}
                   (A^{\rm out},\mathcal F_{i-1})
      \preccurlyeq
             e^{(t-i)\theta T}e^{(Z-t\theta)T}
                                      =e^{(Z-i\theta)T}.
\]
The subpower factor \(C_{t,i}(m)\) is included here.
Together with \eqref{eq:ope-probe-quotient-norm}, this
gives
\begin{equation}
       |\ell_{i,\lambda}(A^{\rm out})|
                                  \preccurlyeq e^{(Z-i\theta)T}.
 \label{eq:ope-grid-auxiliary-coordinate}
\end{equation}
Testing \eqref{eq:ope-grid-auxiliary-error} on the
finitely many mixed probes also gives
\begin{equation}
 |c^{i,\lambda}_{\mathbf A}(\mathbf x;0)
                      -\ell_{i,\lambda}(A^{\rm out})|
                                  \preccurlyeq e^{(Z-G_t)T}.
 \label{eq:ope-grid-probe-error}
\end{equation}
Their fixed positive clearances from the origin ensure
that all these tests are permitted at small \({\varrho}\).

For a coefficient in degree \(d\), choose \(t\ge d\)
with \(G_t\ge d\theta\). Then
\begin{equation}
             |c^{d,\lambda}_{\mathbf A}(\mathbf x;0)|
                          \preccurlyeq e^{(Z-d\theta)T}.
 \label{eq:ope-grid-coefficient}
\end{equation}
For a truncation degree \(D\), choose \(t\ge D+1\)
with \(G_t\ge(D+1)\theta\).
Exact basis reconstruction of \(A^{\rm out}\) gives
\[
 A^{\rm out}-P_D[A^{\rm out}(0)]
     =\sum_{i=D+1}^t\sum_{\lambda\in I_i}
                  \ell_{i,\lambda}(A^{\rm out})O_{i,\lambda}.
\]
Insert this expression with the spectators.
Each \(\|\mathsf M^0(O_{i,\lambda}(0)B)\|\) has a fixed bound
depending on \(i,\lambda,B,c\), by the fixed-clearance
version of \eqref{eq:ope-crude-word}.
Use \eqref{eq:ope-grid-auxiliary-coordinate} for the
omitted modes and \eqref{eq:ope-grid-probe-error} for
the retained modes. The three differences, namely
the local replacement, the omitted channels and the
retained-coordinate errors, prove
\begin{equation}
                  \|\mathcal R_D(w;B)\|
                         \preccurlyeq e^{[Z-(D+1)\theta]T}.
 \label{eq:ope-grid-remainder}
\end{equation}
All choices of \(t\) are auxiliary estimates of the
same fixed \(P_D\). They do not change its coefficients.

\subsection{Removing the grid by one finite Taylor formula}

We estimate the required derivatives by their own grid and
crude bounds. We do not infer a derivative estimate by
differentiating an inequality for the original error.
For a coefficient, or for the algebra-valued remainder,
the exact calculus gives
\begin{equation}
 \partial_{x_a^\nu}F_{\mathbf A}
                  =F_{A_1,\ldots,\partial_\nu A_a,\ldots,A_r}.
 \label{eq:ope-exact-derivative}
\end{equation}
The derivative raises the presented total degree by one.
For the remainder this identity follows also for \(P_Dw\),
because its finite mixed probe formulas commute with
input differentiation. The maps are smooth in operator
norm on separated sets, by the buffered operator calculus.
Thus the fundamental theorem of calculus applies in
the Banach space \(\mathcal A_0\).

There is a crude bound with exponent \(W+c_0\), where
\(c_0\) is independent of the input degrees.
For a coefficient take \(c_0\) to be the maximum total
degree of its fixed probe words and use
\eqref{eq:ope-crude-word}. For the remainder use
the same bound on the full word, and its finite
coefficient formula times the fixed-clearance output
word bounds. One may take \(c_0\) to be the maximum
of the spectator total degree and the total degrees
of the probe words through \(D\).
Fixed multipliers and finite sums are included in
\(\preccurlyeq\).

Both desired grid estimates have the form
\[
            \|F_{\mathbf A}\|
                        \preccurlyeq e^{[(1+\gamma)W+d_0]T},
\]
where \(d_0=8r\gamma-d\theta\) for a coefficient and
\(d_0=8r\gamma-(D+1)\theta\) for the remainder.
Choose a positive integer \(q\) so large that
\begin{equation}
                   c_0-q\gamma<\gamma W+d_0.
 \label{eq:ope-taylor-order-choice}
\end{equation}
This is a finite choice for the fixed coefficient or
truncation request. Round each cluster point
to a generic spacing-\({\varepsilon}_{j(T)}\) point at distance
at most \(C_r {\varepsilon}_{j(T)}\). Distinct coordinates in each
axis can be arranged by at most \(r\) extra mesh shifts.
Since
\[
                {\varepsilon}_{j(T)}=e^{-(1+\gamma)T+o(T)}=o({\varrho}),
\]
the straight rounding paths retain separation
\(\eta {\varrho}/2\) and extent at most \((K+1){\varrho}\).
Write the original and rounded configurations as
\(x\) and \(y\), and put \(h=x-y\) in \(\mathbb R^{4r}\).
Taylor's formula in the scalar space or in
\(\mathcal A_0\), with the appropriate norm, is
\begin{align}
 F_{\mathbf A}(x)={}&
  \sum_{a=0}^{q-1}\frac1{a!}
                       D^aF_{\mathbf A}(y)[h^{\otimes a}]\notag\\
 &+\frac1{(q-1)!}\int_0^1(1-t)^{q-1}
                  D^qF_{\mathbf A}(y+th)[h^{\otimes q}]\,\dd t.
 \label{eq:ope-single-taylor-extension}
\end{align}
It follows by integrating the one-variable function
\(t\mapsto F_{\mathbf A}(y+th)\) repeatedly. Smoothness
in the Banach norm justifies that integral.

By \eqref{eq:ope-exact-derivative}, each order-\(a\)
coordinate derivative is the same coefficient or remainder
with input derivatives of total presented degree \(W+a\).
The grid estimate is therefore
\[
 \|D^aF_{\mathbf A}(y)\|
   \preccurlyeq e^{[(1+\gamma)(W+a)+d_0]T},
                         \qquad 0\le a<q.
\]
Since \(|h|\le C_r\varepsilon_{j(T)}\), multiplication
by \(|h|^a\) cancels precisely the additional
\((1+\gamma)aT\). Every term in the finite Taylor
polynomial has exponent \(((1+\gamma)W+d_0)T+o(T)\).
For the remainder use the crude bound on the enlarged
shape class at all points of the segment:
\[
 |h|^q\sup_{0\le t\le1}\|D^qF_{\mathbf A}(y+th)\|
 \preccurlyeq e^{-(1+\gamma)qT}e^{(W+q+c_0)T}
             =e^{(W+c_0-q\gamma)T}.
\]
By \eqref{eq:ope-taylor-order-choice} this has a strictly
smaller exponent than the required one. The finitely many
subpower losses are absorbed in that strict margin.
Formula~\eqref{eq:ope-single-taylor-extension} proves the
desired estimate at \(x\). Only derivatives through
\(q\), degrees through \(W+q\), and their finite
coefficient balls are admitted. Their constants and
cutoff tails are common to the fixed request. No further
induction or new norm is required.

We have proved, for arbitrary admissible configurations,
\begin{align}
 |c^{d,\lambda}_{\mathbf A}(\mathbf x;o)|
 &\preccurlyeq
      e^{[W-d+\gamma(W+8r+d)]T},\notag\\
 \|\mathcal R_D(w;B)\|
 &\preccurlyeq
      e^{[W-D-1+\gamma(W+8r+D+1)]T}.
 \label{eq:ope-arbitrary-configuration-rates}
\end{align}
Translation restores any centre \(o\). The thresholds
are uniform over the input coefficient balls, cluster
shapes and all fixed-clearance spectator configurations.

\subsection{The separated expansion with one coefficient family}

\begin{theorem}[All-orders separated OPE]
\label{thm:all-orders-separated-ope}
Let \(O_{j,\lambda}\) and \(c^{j,\lambda}\) be the single
prescription of Section~\ref{sec:ope-coefficients}.
Fix \(A_i\in\mathcal F_{d_i}\), \(W=\sum_i d_i\).
For every \(K\ge0\), \(\eta,\varepsilon>0\), \(j\ge0\)
and \(\lambda\in I_j\), there is \(r_0>0\) such that
\begin{equation}
 |c^{j,\lambda}_{\mathbf A}(\mathbf x;o)|
                                 \le {\varrho}^{j-W-\varepsilon}
 \label{eq:ope-proved-coefficient}
\end{equation}
for \(0<{\varrho}<r_0\), whenever
\(|x_i-o|\le K{\varrho}\) and
\(|x_i-x_{i'}|\ge\eta {\varrho}\) for \(i\ne i'\).

Fix also a spectator list \(B_b\), \(c>0\) and a finite
\(D\ge0\). There is \(r_1>0\), preceding all positions
and phases, such that on the same cluster geometry and
\[
            |o-y_b|\ge c,\qquad
                 |y_b-y_{b'}|\ge c\quad(b\ne b'),
\]
one has the stronger common algebra estimate
\begin{equation}
                         \|\mathcal R_D(w;B)\|
                                  \le {\varrho}^{D+1-W-\varepsilon}
                            \qquad(0<{\varrho}<r_1).
 \label{eq:ope-proved-algebra-remainder}
\end{equation}
Consequently, for \(\mathcal S=S\) and for every phase
\(\mathcal S=S^\sigma\), the same \(r_1\) gives
\begin{align}
 \left|\mathcal S(\mathbf A,\mathbf x;\mathbf B,\mathbf y)
 -\sum_{j=0}^D\sum_{\lambda\in I_j}
          c^{j,\lambda}_{\mathbf A}(\mathbf x;o)
             \mathcal S(O_{j,\lambda}(o),\mathbf B,\mathbf y)\right|
                               \le {\varrho}^{D+1-W-\varepsilon}.
 \label{eq:ope-proved-separated-remainder}
\end{align}
The thresholds may depend on the fixed finite labels,
coefficient balls, shape and clearance parameters,
the chosen degree or truncation and \(\varepsilon\).
The coefficients themselves precede those choices of
spectators, phase and truncation.
\end{theorem}
\begin{proof}
For \(r\ge1\), apply \eqref{eq:ope-arbitrary-configuration-rates}.
Choose \(\gamma>0\) so small that its displayed loss
is less than \(\varepsilon/4\). The additional
\(\preccurlyeq\) loss can be chosen less than
\(\varepsilon/4\). Thus the estimates eventually have
strictly smaller exponent than \(W-j+\varepsilon\)
and \(W-D-1+\varepsilon\), respectively.
All fixed constants are absorbed in that slack.
Decrease the coefficient threshold until the cluster
avoids its finite probe sets. Decrease the remainder
threshold also so that \(K{\varrho}<c/2\); the cluster and
spectators then form a separated configuration.
For \(K=0\) use the admissible larger extent bound \(1\).
The empty cluster is the exact unit expansion, so its
remainder is zero and its coefficient bound holds for
\({\varrho}<1\).

For precision about uniformity, if either claimed
threshold failed, choose violating configurations at
radii \(r_l\downarrow0\), with the fixed labels and
geometric bounds. Run the preceding estimates along
this sequence with their fixed \(\gamma\) and finite
auxiliary horizon. Their constants and weight budgets
are uniform on exactly this class. The strict exponent
slack contradicts the violation. This also permits
unbounded spectator positions, since every lattice
root and algebra estimate used their clearances and
not their absolute locations.

Finally characters of \(\mathcal A_0\), and its vacuum
state, have norm one. Evaluating
\eqref{eq:ope-proved-algebra-remainder} gives
\eqref{eq:ope-proved-separated-remainder} simultaneously.
This evaluation makes no additional subsequence or
choice of coefficient.
\end{proof}

\begin{corollary}[Tested separated form]
\label{cor:ope-tested-separated}
For fixed finite label and geometric data as above,
integrating the remainder over any measurable admissible
geometry and phase parameters with a complex weight
of total absolute integral at most \(B_0<\infty\) gives
absolute error at most
\(B_0 {\varrho}^{D+1-W-\varepsilon}\), with the same threshold.
\end{corollary}
\begin{proof}
The finite probe formulas and the phase-algebra evaluation
are jointly continuous on the separated domains.
Their fixed-radius bounds make all full and truncated
terms integrable against the stated weight. Integrate
\eqref{eq:ope-proved-separated-remainder} and apply the
triangle inequality.
\end{proof}

The theorem concerns separated Euclidean correlations
of the very same fields used in the selected Wightman
reconstruction. Every finite \(D\) is covered by one
coefficient family. It asserts the specified asymptotic
finite truncations; an infinite sum over \(D\) is not
needed. The finite-history estimates were obtained
before the original ordered limits, with unchanged
cutoff, volume and normalization conventions.
Constants may depend on the fixed group, calendar,
quotient norms and probe prescription, in addition
to the listed finite label and geometric data.
They are uniform in the cutoff and volume on the
physical tails; the final thresholds are uniform
over the phase space.
The trace analysis remains to complete the agreed
construction theorem.
% !TeX root = 4DYM_notes.tex
\section{The centred stress trace and its positive spectral measure}
\label{sec:centered-trace-field}

Keep the Wilson family, the mixed hierarchy, its phase
decomposition and the selected phase \(\sigma_*\) constructed
above. Write \(S_*\) for \(S^{\sigma_*}\). The stress tensor
\(T_{ab}\in\mathcal F_4\) is the translation-normalised tensor
of \eqref{eq:normalised-interacting-stress-laws}. It is centred
in the mixture; its mean in a particular phase need not be zero.
We must subtract the mean in the vacuum actually used for
reconstruction.

\begin{definition}[Centred trace]
Set
\begin{equation}
 e_{\rm tr}=S_{*,1}\left(\sum_{a=1}^4T_{aa}\right),
 \qquad
 \Theta=\sum_{a=1}^4T_{aa}-e_{\rm tr}1_4,
 \qquad
 \tau(r)=S_{*,2}(\Theta,\Theta;0,re_4),\quad r>0.
 \label{eq:centred-trace-definition}
\end{equation}
Here \(1_4\) is the unit field viewed in degree four. Translation
invariance makes the one-point function independent of its
position. The same notation \(\Theta\) will denote its continued
Minkowski label; \(\Phi(\Theta,f)\) is its smeared field on the
common word domain.
\end{definition}

\begin{lemma}[Exact centring and scalar covariance]
\label{lem:centred-trace-identification}
The number \(e_{\rm tr}\) is real. The field \(\Theta\) is real
and \(O(4)\)-invariant, and its continuation is a Hermitian
Lorentz scalar. On the word domain,
\begin{equation}
 \Phi(\Theta,f)=
       \sum_a\Phi(T_{aa},f)
                -e_{\rm tr}\left(\int f(x)\,\dd x\right)I,
 \qquad
 \langle\Omega,\Phi(\Theta,f)\Omega\rangle=0.
 \label{eq:centred-trace-operator-identity}
\end{equation}
The stress labels in the sum are continued as degree-four
labels. Thus this formula uses Euclidean trace labels before
continuation, consistently with the signs of the continued
tensor in Section~\ref{sec:wightman-stress}.
\end{lemma}
\begin{proof}
Reality of the phase hierarchy applied to a real one-point
label gives \(e_{\rm tr}=\overline{e_{\rm tr}}\).
Contract the tensor transformation law with \(\delta_{ab}\);
orthogonality gives
\(\sum_{a,i,j}Q_{ia}Q_{ja}T_{ij}=\sum_iT_{ii}\).
The unit is also fixed. The finite-dimensional analytic
continuation of this fixed vector is fixed by the Lorentz
representation, and its conjugate label equals itself.
The field construction is linear in labels and the unit acts
as \((\int f)I\), giving the first identity. The Euclidean
one-point function of \(\Theta\) is identically zero. Its
positive-height continuation and then its tempered boundary
are zero by the uniqueness argument used in
Section~\ref{sec:wightman-continuation}. This gives the second
identity. All operations use the selected phase and the same
common domain.
\end{proof}

\subsection{Obtaining one measure from the actual field vectors}

A positive spectral representation supplies the monotonicity
needed later to pass from comparison points to an exact
distance. We give its construction, rather than assuming
monotonicity of an interacting correlation.
Let \(P(\dd p)\) be the joint spectral measure of translations
from \eqref{eq:wightman-pvm-cone}, with \(p=(\mathbf p,E)\)
and \(E\ge|\mathbf p|\). Use the Fourier convention
\[
 \widehat f_+(p)=
       \int e^{i(Et-\mathbf p\cdot\mathbf x)}
                                      f(\mathbf x,t)\,\dd\mathbf x\,\dd t.
\]
Put \(v_f=\Phi(\Theta,f)\Omega\) and
\(\mu_f(B)=\langle v_f,P(B)v_f\rangle\).

\begin{lemma}[Trace spectral measure]
\label{lem:trace-spectral-measure}
There is a unique positive measure \(\mu_\Theta\), finite
on compact sets and of polynomial growth, supported in
\(\{E\ge|\mathbf p|\}\), such that
\begin{align}
 \langle v_f,v_g\rangle
     &=\int\overline{\widehat f_+(p)}\widehat g_+(p)
                                             \,\mu_\Theta(\dd p),\label{eq:trace-spectral-inner}\\
 \tau(r)&=\int e^{-rE}\,\mu_\Theta(\dd p),\qquad r>0.
 \label{eq:trace-spectral-laplace}
\end{align}
It is Lorentz invariant. In particular \(\tau\) is real,
nonnegative, smooth and nonincreasing. More precisely,
\[
           (-1)^j\tau^{(j)}(r)
                 =\int E^je^{-rE}\,\mu_\Theta(\dd p)\ge0
                    \quad(j\ge0).
\]
\end{lemma}
\begin{proof}
Translation covariance and Fubini for smeared word vectors
give the convolution identity
\begin{equation}
       \widehat h_+(P)v_f=\widehat f_+(P)v_h.
 \label{eq:trace-spectral-compatibility}
\end{equation}
One may first take Schwartz functions with compact Fourier
support; their spectral multipliers are bounded and all
integrals converge in the Hilbert norm. Continuity extends
the identity to Schwartz tests.
Choose \(f_j\) with \(\widehat f_{j,+}\) smooth, compactly
supported and strictly positive on the ball \(|p|\le j\).
On that ball define
\(\mu_\Theta=|\widehat f_{j,+}|^{-2}\mu_{f_j}\).
For \(i,j\), take the spectral norm squared of
\eqref{eq:trace-spectral-compatibility} on a bounded Borel
set in the overlap. It gives
\[
 |\widehat f_{i,+}|^2\,\dd\mu_{f_j}
                 =|\widehat f_{j,+}|^2\,\dd\mu_{f_i}.
\]
The definitions agree, so they glue to a locally finite
positive measure. The same identity for an arbitrary \(f\)
gives \(\mu_f=|\widehat f_+|^2\mu_\Theta\); polarisation
proves \eqref{eq:trace-spectral-inner} and uniqueness.
Cone support follows from that of \(P\).

Here is also a growth bound. The tempered two-point estimate
already proved for the reconstructed theory bounds
\(\|v_f\|^2\) by a fixed finite sum of Schwartz seminorms of
\(f\), squared. Choose a smooth Fourier cutoff equal to one
on the radius-\(R\) ball, obtained by scaling one fixed
cutoff. The corresponding seminorms grow at most as a fixed
power of \(1+R\). Consequently
\[
             \mu_\Theta(\{|p|\le R\})\le C(1+R)^{\mathsf M}
\]
for some \(C,{\mathsf M}<\infty\). Constants here depend on the fixed
field and selected hierarchy, not on a remaining lattice
cutoff or volume. Scalar covariance, vacuum invariance and
uniqueness of the measure imply Lorentz invariance.

For completeness the measure represents the full relative
two-point distribution: on product tests this is
\eqref{eq:trace-spectral-inner}; finite sums of product tests
are dense in the Schwartz space and both sides are continuous.
Cone support and polynomial growth make its Fourier--Laplace
transform holomorphic at positive imaginary time. On every
compact interval \(r\ge r_0>0\), each differentiated integrand
is bounded by \(C_{j,r_0}e^{-r_0E/2}\), which is integrable
because \(|p|\le\sqrt2E\) on the cone. The positive-height
comparison of Sections~\ref{sec:wightman-continuation}
and~\ref{sec:wightman-reconstruction} identifies this transform
with the continued Euclidean two-point kernel. Equivalently,
test first at compact spatial frequencies, apply the heat
intertwining identity \eqref{eq:wightman-heat-intertwining},
and use the same exponential majorant to remove the frequency
cutoff. At zero spatial separation the result is
\eqref{eq:trace-spectral-laplace}. Differentiation under this
majorant gives the last formula and monotonicity.
\end{proof}

The definitions and spectral reduction follow
\cite[Lemmas~18.1--18.2, PDF pp.~203--204]{ArmstrongDraft}. They use no
short-distance estimate. Positivity here is a consequence of
the reconstructed Hilbert space, and will not be inferred from
the sign of a perturbative coefficient.
% !TeX root = 4DYM_notes.tex
\section{Maxwell comparison at prescribed small radii}
\label{sec:trace-prescribed-radii}

The sequence in Section~\ref{sec:stable-maxwell-comparison}
is adequate for identifying finite-dimensional ranks. A limit
along that sequence alone does not determine a short-distance
asymptotic at every radius. We strengthen the comparison,
using the entry choice \eqref{eq:induction-trace-calendar}.
This is an additional use of the same initial calendar; no
second Wilson family or second continuum hierarchy is chosen.
The argument adapts
\cite[Lemma~18.3, PDF pp.~204--207]{ArmstrongDraft}.
The compatible-window and replicated-source organisation is
also described in \cite[Unit VII.6, Lemmas VII.6.g--h,
PDF p.~569, and Proposition VII.6.2, pp.~563, 573]{ArmstrongNotes}.
We give the error and normalisation calculations below for
the actual family fixed in this manuscript.

All norms below are fixed coefficient norms on
\(V_d^{\rm Max}\). A \emph{stable frame} consists of real,
\(H_4\)-equivariant complex linear isomorphisms
\(U_i(d):V_d^{\rm Max}\to\mathcal F_d\), with exact unit and
inclusion compatibility and convergence of transported
derivatives as in Section~\ref{sec:stable-maxwell-comparison}.
For a point \(x\), \(x^{(i)}\) denotes a nearest point in
\(r_i^{-1}{\varepsilon}_{n_i}\Z^4\). A generic configuration has distinct
coordinates in each coordinate direction; we only round on
compact subsets with a positive lower bound on these coordinate
separations.

\begin{lemma}[Comparison on a prescribed null sequence]
\label{lem:trace-prescribed-radii}
Every positive sequence tending to zero has a subsequence,
denoted \(r_i\), with stable frames \(U_i(d)\) and levels
\(n_i<k_i\), both tending to infinity, such that the following
statements hold simultaneously for every finite degree and
arity request:
\begin{align}
 \beta_i&=\lim_{M}\lim_n\beta_{n,M,k_i}\longrightarrow\infty,
 &\frac{\beta_i}{\widehat\beta_{k_i}}&\longrightarrow1,
 &\frac{\log(1/r_i)}{\beta_i}&\longrightarrow b_0^{-1}.
 \label{eq:trace-prescribed-coupling}
\end{align}
The limits of the bounded coefficient arrays are taken along
further refinements of the already chosen ordered tails;
they preserve all correlations of \(S\).
If \(d_j\le D\), \(A_j^{(i)}\to A_j\) in \(V_{d_j}^{\rm Max}\)
and \(K\) is a compact generic configuration set, then
\begin{equation}
 \sup_{x\in K}\left|
 S_r\bigl(U_i(d_1)A_1^{(i)},\ldots,U_i(d_r)A_r^{(i)};
                                      r_ix^{(i)}\bigr)
                     -S_r^{\rm Max}(\mathbf A;x)\right|\to0.
 \label{eq:trace-prescribed-comparison}
\end{equation}
The corresponding smooth comparison and compact shell-flux
limits hold without rounding. For the two-point comparison
without derivatives the convergence is uniform in phase.
The finite-cutoff estimates underlying these assertions are
uniform in cutoff and period on the specified ordered tails.
\end{lemma}
\begin{proof}
Fix first \(D\), an arity bound \(r\ge2\), a compact position
set and finitely many derivative orders. Put \({\mathsf M}=P_v\) at
a radius level \(v\), and set
\begin{equation}
 \begin{aligned}
 {\varrho}(v,u)&={\varepsilon}_v{\mathsf M}^u,\qquad
 u\in[\alpha_0,\alpha_0+1],\qquad
 H_v=(\log\widehat\beta_v)^{48},\\
 k(v)&=\min\{m>v:\log({\varepsilon}_v/{\varepsilon}_m)\ge H_v\}.
 \end{aligned}
 \label{eq:trace-radius-window}
\end{equation}
The constants \(\alpha_j\) below belong to this finite request.
Choose \(c<0\) and \(\alpha_0,\alpha_1,\alpha_2\) so that
\begin{align}
 \alpha_0-\alpha_1&>6,&
 \alpha_0-\alpha_2&>5-c,\notag\\
 (\alpha_0-\alpha_1)Dr+c&<0,&
 (\alpha_0-\alpha_2)Dr-\alpha_0-c&<0,
 &\alpha_2&\ge Dr+10.
 \label{eq:trace-radius-exponents}
\end{align}
These inequalities are compatible: fix the first difference,
then choose \(-c\), then the second difference, and finally
increase all three \(\alpha_j\) by the same sufficiently large
number.
Include \(\alpha_j>2rD+16\) in this last choice.
The local calibration uses one Taylor order large enough
for the compact exponent intervals
\([\alpha_j,\alpha_j+1]\), with the same full covariance
and coefficient matrix throughout them.

For \(b=\alpha_j+u-\alpha_0\), stop backwards from \(v\)
after logarithmic block length first exceeds
\((b+2)\log {\mathsf M}\), and again after a further \(10\log B_v\).
Every overshoot is \(O(\log\log\widehat\beta_v)\), by
\eqref{eq:induction-trace-calendar}. This is \(o(\log {\mathsf M})\).
The corresponding cells have side at most \({\varepsilon}_v{\mathsf M}^{b+5}\).
Every level in these windows tends to infinity. The change
in nominal coupling is \(O(\log {\mathsf M})\), hence negligible
relative to \(\widehat\beta_v\).
For the column constants of the finite marked request,
\begin{equation}
 \log\prod_{\text{window}}\mathcal C_{J(m)}
 \le\left(\sup_{\text{window}}J(m)^{-2}\right)
                                      \sum_{\text{window}}\log L_m
                  =o(\log {\mathsf M}).
 \label{eq:trace-window-column-cost}
\end{equation}
This bound permits windows crossing profile boundaries.

% !TeX root = 4DYM_notes.tex
We now perform the complete source estimate on these growing
windows. Write \(h_b<l_b<v\) for the two backward stopping
levels, so that
\[
 \log(\varepsilon_{l_b}/\varepsilon_v)\ge(b+2)\log\mathsf M,
 \qquad
 \log(\varepsilon_{h_b}/\varepsilon_{l_b})\ge10\log B_v.
\]
Here \(h_b\) is a comparison cell level; it is distinct from
the later insertion level \(k(v)>v\). All these levels depend
on the finite request, \(v,u\), and the fixed calendar, and
are independent of the terminal cutoff and physical period.

\emph{The norm and the complete correction.}
In one comparison cell retain at most \(r\) names. For a
marker polynomial with activity coefficients use
\[
       \|A(z)\|_{\alpha,m}
          =\sum_{I\subseteq\{1,\ldots,r\}}
                         \alpha^{|I|}\|A_I\|_m,
                 \qquad \alpha=\mathsf M^{1/(8r)}.
\]
In this formula \(\|A_I\|_m\) is the existing rooted
regular family norm, the existing forced family norm divided
by its original square-root allowance, or the scalar price,
as appropriate to that coefficient. Their sum is used when
all three types occur. Supports, inspected reads and covering
weights are the original ones. The finite subset norm is
submultiplicative by \eqref{eq:common-subset-product} and the
already proved family product estimates. At most \(r\) names
in a cell give the single inflation \(\alpha^r=\mathsf M^{1/8}\).
There is no inflation depending on the total number of cells.

Let \(G_I\) be the Maxwell moments of the trial labels on a
subset \(I\), including \(G_\varnothing=1\). As in
\eqref{eq:stable-corrected-source}, define
\[
       G(z)=\sum_I G_Iz_I,\qquad
       R(z)=G(z)^{-1}\prod_{i=1}^r(1+z_iF_i).
\]
The inverse is the finite sum
\(G^{-1}=\sum_{j=0}^r(-(G-1))^j\). Coefficient comparison
in \(GR=\prod_i(1+z_iF_i)\) gives
\[
 R_\varnothing=1,\qquad
 R_I=\prod_{i\in I}F_i-
                \sum_{\varnothing\ne J\subseteq I}G_JR_{I\setminus J}.
\]
For example
\[
 R_{12}=F_1F_2-G_1F_2-G_2F_1+2G_1G_2-G_{12}.
\]
Thus singleton centring alone does not specify the pair.
These are complete physical coefficients. The correction is
multiplied into the full pushed source, with all its proper
subset modes, before the coefficient is estimated.

\emph{The initial error budget.}
The calculation in
\eqref{eq:position-gaussian-source-initial} and
\eqref{eq:position-low-mode-error} applies with \(t=v\).
Its separate contributions, before marker inflation and
column transport, are
\begin{equation}
 \begin{array}{c|c}
 \text{contribution}&\text{price per cell}\\ \hline
 \text{Taylor spatial moments and lattice error}&C\mathsf M^{-b/2}\\
 \text{slow curvature background}&C\mathsf M^{-3}\\
 \text{omitted low covariance}&C\mathsf M^{-8}\\
 \text{finite-interval nonquadratic comparison}&C\mathsf M^A\widehat\beta_v^{-1/16}\\
 \text{localisation and chart failures}&C\mathsf M^Ae^{-c_*B_*} .
 \end{array}
 \label{eq:prescribed-source-error-budget}
\end{equation}
Here \(A<\infty\), and the constants
belong to the fixed finite degree, arity, derivative and compact
position request. The decay constants \(c_*\) and thresholds
\(B_*\) are those of the finitely many profiles met by the
window; their entry conditions pay their actual values. They
are not replaced by a positive infimum over all profiles.

To check the first three lines, the rescaled mesh is
\(\mathsf M^{-b}\). Higher spatial moments of the actual
Taylor polynomials weaken the common first bound to
\(\mathsf M^{-b/2}\), as in
Lemma~\ref{lem:actual-Taylor-window-comparison}.
The slow curvature jet of weight
\(w\ge2\) costs
\(CP_v(\varrho/\varepsilon_{l_b})^w\le C\mathsf M^{-3}\);
and its omitted covariance costs
\((\varrho/\varepsilon_{l_b})^{4+s}\le\mathsf M^{-8}\)
at derivative order \(s\ge0\). A Gaussian moment of fixed
degree is a finite polynomial in its contractions. If
\(|c_j|\le U\) and \(|c_j-\widetilde c_j|\le e\), then
\[
 \left|\prod_{j=1}^q c_j-\prod_{j=1}^q\widetilde c_j\right|
                       \le q e(U+e)^{q-1}.
\]
Apply this to every contraction and every subset, and then
to the finite inverse of \(G\). The resulting constant
depends on \(D,r\) and the separated compact, rather than
on the number of replicated cells. This proves the first
three source prices, including every lower-moment correction.

The finite-interval comparison, localisation and failure
estimates give the last two lines in the same norms.
At fixed finite request, the
calendar entries used for the Gaussian comparison make the
coupling and localisation prices smaller than any prescribed
fixed inverse power of \(\mathsf M\). This uses
\(\log\mathsf M=o(\log\widehat\beta_v)\), with the
profile constants paid before entry. The Taylor polynomial
is constructed from the actual insertion at \(v\);
no terminal Gaussian graph or terminal-memory term occurs.
Fast covariance
between cells has its exponential row sum, while long
activity reads have the reserved covering weight. Summing
either is a rooted sum and introduces no factor counting
cells. Periodic image tails are bounded using the positive
minimum physical period and the fast covariance, as in
Section~\ref{sec:gaussian-position-tests}.

\emph{Closing the transport on a growing window.}
Let \(\Lambda_v\ge1\) be the product of the actual linear
column prices on \([h_b,v]\), including the finitely many
endpoint maps and the three normalisations under consideration.
Use the finite comparison requests included when fixing the
original calendar. Equation
\eqref{eq:trace-window-column-cost}, and the same logarithmic
floor for its endpoint columns, give
\[
                  \Lambda_v\le\mathsf M^{1/64}
\]
eventually. The constant \(1/64\) is only a convenient fixed
reserve: the logarithm of this product is \(o(\log\mathsf M)\).
In particular the product is not treated as a fixed constant.

Let \(a_m\) be the complete corrected income rooted at one
cell, including marked scalar modes. On the outer interval
the same connected expansion as in
\eqref{eq:position-source-iteration} gives
\[
 a_m\le K_{J(m+1)}a_{m+1}
                       +2\mathsf M^{1/2}a_{m+1}^2,
 \qquad \mathcal C_{J(m+1)}\ge K_{J(m+1)}+2.
\]
Every nonlinear term contains at least two corrected sources.
The finite subset partitions use
\eqref{eq:common-subset-product}; Gaussian contractions
and the fixed polynomial packet price are already in the
norm. The remaining powers of \(B_m=\mathsf M^{o(1)}\)
are bounded by \(\mathsf M^{1/2}\) after entry. On failures
one pays the polynomial before dividing by its forced
allowance, using the common precision reserve for that word.

Let \(\mathcal E_{v,n}\) be the sum in
\eqref{eq:prescribed-source-error-budget}. At the entries
fixed for those requests, and then on a full terminal-cutoff
tail, each of its last two lines is at most \(\mathsf M^{-10}\).
The chosen lower bound on \(b\) gives
\begin{equation}
 \alpha^r\Lambda_v\mathcal E_{v,n}
 \le C\mathsf M^{1/8+1/64}
           (\mathsf M^{-b/2}+\mathsf M^{-3}+\mathsf M^{-8}
                                           +2\mathsf M^{-10})
 \le\mathsf M^{-3/4}.
 \label{eq:prescribed-source-entrance}
\end{equation}
The slow-background exponent on the middle expression is
\(-183/64<-3/4\). This leaves room for the fixed constants.
Starting with the jointly integrated packet at \(l_b\),
induct backwards. If \(a_{m+1}\le\mathsf M^{-3/4}\), then
\[
 2\mathsf M^{1/2}a_{m+1}^2\le2a_{m+1},
 \qquad a_m\le\mathcal C_{J(m+1)}a_{m+1}.
\]
Multiplying only the remaining column prices bounds every
intermediate \(a_m\) by the left side of
\eqref{eq:prescribed-source-entrance}. This proves the
induction and the required small source-attachment condition
through the whole window, including profile changes.

\emph{The original denominator and the reflected root.}
Let \(V=\#\Lambda_{h_b}\), the number of comparison cells
in the periodic tiling, and let \(g_{h_b}\) be its fluctuation
coupling. The sharp positive-ratio estimate gives the complete
normalised source norm at most \(e^{Vq_v}\), where
\[
 q_v\le Q_{J(h_b)}+C+
                    \frac{4q_G\log(g_{h_b}^{-1})}{V}
                    \le\frac{\log\mathsf M}{16r}.
\]
The finite comparison entries pay the first two terms. For
the last one,
\(V=(M/\varepsilon_{h_b})^4\ge
 M_{\min}^4\varepsilon_{h_b}^{-4}\), so it tends to zero
uniformly in the physical side \(M\). The vacuum scalar
cancels between the numerator and the same positive partition
integral. No sector partition function is used.

Extract the nonempty mode \(I\) in each reflected cell,
using the coefficient bound at radius \(\alpha\), and
only then take the \(V\)-th root. The result is
\begin{equation}
 \|R_I\|_{\rm root}
       \le e^{q_v}\alpha^{-|I|}
       \le\mathsf M^{1/(16r)-|I|/(8r)}
       \le\mathsf M^{-1/(16r)},\qquad I\ne\varnothing.
 \label{eq:prescribed-complete-source-root}
\end{equation}
Every mask, incoming member, subset mode, forced inventory and
Gaussian/Haar normalisation stays in the complete source.
All constants and tails above are uniform in
\(u\in[\alpha_0,\alpha_0+1]\), by taking the finite
request on this compact exponent interval. Their dependence
on \(D,r\), profiles, derivative orders and position compact
is permitted; no bound is claimed uniformly over these
requests.

We include the position argument because it is what allows
prescribed \(u\). At its own normalisation a singleton has
root at most \(1+|S_1^{\rm Max}(A)|\): extract its centred
coefficient at amplitude \(e^{2q}\) from a norm at most
\(e^q\). Changing normalisation from
\(\alpha_1+u-\alpha_0\) to \(u\) multiplies a homogeneous
weight-\(j\) coefficient by
\({\mathsf M}^{(\alpha_0-\alpha_1)j}\). Thus the root of a degree-\(D\)
label is at most \(C{\mathsf M}^{(\alpha_0-\alpha_1)D}\).
Use also the roots associated with \(\alpha_2\), and
include signed and conjugate basis labels. The chessboard
inequality bounds every reflected product by the product of
these roots.

Let \(\delta={\mathsf M}^c\). The value-bound cells at
normalisation \(\alpha_1+u-\alpha_0\) have side
\(o(\varrho(v,u))\), since their ratio to that radius is
at most \(\mathsf M^{\alpha_1-\alpha_0+5}\).
The displacement-bound cells at normalisation
\(\alpha_2+u-\alpha_0\) have side
\(o(\varrho(v,u)\delta)\), since their ratio to that
smaller margin is at most
\(\mathsf M^{\alpha_2-\alpha_0+5-c}\).
Both conclusions follow from
\eqref{eq:trace-radius-exponents}. The first root bound
controls values throughout the separated compact; only the
second is used with the shrinking displacement margin.
Exclude the finitely many
strips where a coordinate hybrid approaches a coincidence
within \(4\delta\). Their total volume is \(O(\delta)\).
Outside these strips, the displacement estimate already
proved for the physical lattice insertions bounds a
positional difference by \(C/\delta\) times the product of
the roots. The scaled mesh is \(h={\mathsf M}^{-u}\).
The strip and interpolation errors are therefore bounded by
\begin{equation}
 C {\mathsf M}^{(\alpha_0-\alpha_1)Dr+c},
 \qquad
 C {\mathsf M}^{(\alpha_0-\alpha_2)Dr-u-c}.
 \label{eq:trace-radius-position-errors}
\end{equation}
Both tend to zero uniformly over the whole interval of \(u\).
Indeed there are only finitely many coordinate hybrids of
the at most \(r\) slots. In a fixed box, each strip
\(|x_i^a-x_j^a|\le4\delta+O(h)\) has volume at most
\(C(\delta+h)\), by integrating first in \(x_i^a\).
The parameter conditions imply \(h=o(\delta)\).
The value-root product costs
\(C\mathsf M^{(\alpha_0-\alpha_1)Dr}\), giving the
first error. Outside these strips, a mesh interpolation
changes each coordinate by at most \(h\). Apply
Lemma~\ref{lem:position-test-displacement} with physical
margin a fixed multiple of \(\varrho\delta\) and the
second root product. Summing the at most \(4r\) moves
gives the second error. This is also valid for bounded
integrated tests on a separated compact. On a fixed generic
compact the strips are eventually empty, so the same bound
directly gives uniform pointwise comparison there.
All periodic coordinate separations equal ordinary separations
at these shrinking radii. The centred box can be chosen with
mesh below \({\varrho}\delta\); the original insertions and all
reflections are interior.

The corrected-source estimate
\eqref{eq:prescribed-complete-source-root} holds for every
repeated reflected array. To make its phase implication
explicit, let \(\Delta_I\in\mathcal A_0\) be the full
corrected coefficient after taking the ordered limits. For
fixed \(v,u\), take \(2q\) alternating conjugated copies
of that coefficient, take the ordered physical limits for
each fixed array, and then separate its clusters. The same
root estimate gives
\[
 \int_\Sigma|\sigma(\Delta_I)|^{2q}\,\dd\rho(\sigma)
                   \le\mathsf M^{-2q/(16r)},\qquad q\ge1.
\]
The continuous character function has supremum at most
\(\mathsf M^{-1/(16r)}\): a larger value would give an
open set of positive \(\rho\)-measure contradicting this
inequality as \(q\to\infty\). Since the full source
obeys \(\prod_i(1+z_iF_i)=G(z)R(z)\), its moment error
on \(I\) is exactly
\[
              \sum_{J\subsetneq I}G_J\sigma(\Delta_{I\setminus J}).
\]
For \(|I|\le2\) the finite free coefficients are uniformly
bounded on the chosen compact. This proves the asserted
phase-uniform two-point comparison. Finite-dimensional
multilinearity permits moving bounded coefficients and
their limits, using the same estimates on a fixed basis.

The intervals of radii cover every sufficiently small number.
Indeed \(L_v\le\log\widehat\beta_v\le P_{v+1}\) eventually,
and \(P_v\) is increasing, so
\[
        {\varepsilon}_vP_v^{\alpha_0}
                  \le {\varepsilon}_{v+1}P_{v+1}^{\alpha_0+1}.
\]
Their endpoints tend to zero, since \(\log P_v=o(s_v)\).
The first crossing of a prescribed radius and continuity in
\(u\) place it in one interval. Extract successively for
degrees, arities and compact sets. At step \(i\) choose the
level sufficiently late that all previous bounds hold and
\(\alpha_0+1\le\log\log\widehat\beta_{n_i}\).
The real finite-detector construction and recursive inclusion
correction of Section~\ref{sec:stable-maxwell-comparison}
give the frames, with exact units and inclusions.

For derivatives, choose a finite injective matrix of free
responses in degree \(d+1\). Its interacting matrices converge
by the comparison just proved. Equation
\eqref{eq:common-fixed-detector-bound} bounds their left
inverses on the same tails. Testing
\(r_iU_i(d+1)^{-1}\partial_aU_i(d)A\) and integrating by parts
identifies its limit as \(\partial_a^{\rm Max}A\). Iteration
gives bounded coefficients at every fixed derivative order.
The smoothness argument of
Section~\ref{sec:smooth-weak-hierarchy} then gives smooth
comparison on separated compacts. Integrating over a fixed
compact shell gives the stated flux limit.

Minimality of \(k(v)\) and the adjacent ratio bound give
\(\log({\varepsilon}_v/{\varepsilon}_{k(v)})=H_v+O(\log\log\widehat\beta_v)
=o(\widehat\beta_v)\).
The uniform running law gives
\(\beta_i/\widehat\beta_{k_i}\to1\), while
\(\widehat\beta_{k_i}/\widehat\beta_{n_i}\to1\).
Finally
\[
 \log(1/r_i)=s_{n_i}-u_i\log P_{n_i},\qquad
       u_i\log P_{n_i}\le(\log\log\widehat\beta_{n_i})^3
                                    =o(\widehat\beta_{n_i}),
\]
and \(\widehat\beta_r=\bar\beta+b_0s_r\). This proves
\eqref{eq:trace-prescribed-coupling}.
Every finite request was fixed before its ratio floor and
entry. No local bound depends on the number of physical
cells. Refining ordered coefficient tails leaves every
previously convergent continuum correlation unchanged.
\end{proof}
% !TeX root = 4DYM_notes.tex
\section{A regulated Ward calculation for the trace coefficient}
\label{sec:trace-ward-coefficient}

The leading Maxwell stress is traceless. Consequently a
zeroth-order Maxwell comparison cannot determine the trace:
we need one additional order in the inverse coupling.
This section computes the required scalar before transporting
the insertion to continuum coordinates. We use the determinant
calculus already developed for the one-loop drift, retaining
the terms proportional to the action gradient. These terms
are called \emph{off-shell contacts}: they vanish at a critical
background but cannot be discarded before differentiating
with respect to the background.
The organisation follows
\cite[Lemma~18.4, PDF pp.~208--210]{ArmstrongDraft}.

\subsection{The raw tensor coordinates and the scalar read}

Put \(q_G=\dim_{\R}\mathfrak g\), and define the free scalar
and tensors
\begin{equation}
 \mathcal S=\frac14\sum_{a,b,\alpha}:F_{ab}^{\alpha}F_{ab}^{\alpha}:,
 \qquad L_{ab}=\delta_{ab}\mathcal S
                  -\sum_{c,\alpha}:F_{ac}^{\alpha}F_{bc}^{\alpha}:,
 \qquad D_{ab}=\delta_{ab}\mathcal S,
 \label{eq:trace-free-scalar-tensors}
\end{equation}
and \(\Xi_{aa}=0\),
\(\Xi_{ab}=\sum_{c,\alpha}:F_{ac}^{\alpha}F_{bc}^{\alpha}:\)
for \(a\ne b\). The colons mean that there are no Wick
contractions within an insertion. These three tensors span
the real centred symmetric \(H_4\)-tensor space of degree four,
by the counting in Section~\ref{sec:interacting-stress-charges}.

Before imposing the Maxwell equation, use the three vector
coordinates represented in component \(b\) by
\[
 \sum_c\langle F_{bc},\partial_bF_{bc}\rangle,\qquad
 \sum_{a,c}\langle F_{ac},\partial_bF_{ac}\rangle,\qquad
 \sum_{a,c}\langle F_{bc},\partial_aF_{ac}\rangle .
\]
The metric contracts colours. A coordinate triple is written
\((x,y,z)\); this notation refers to coefficients in this
fixed basis. If a tensor has diagonal entries
\(z_0\sum_c\langle F_{ac},F_{ac}\rangle+
 z_1\sum_{b,c}\langle F_{bc},F_{bc}\rangle\)
and off-diagonal entries
\(z_2\sum_c\langle F_{ac},F_{bc}\rangle\), Bianchi and the
product rule give
\[
       \operatorname{div}_{\rm raw}B
                      =\bigl(2(z_0-z_2),\,2z_1+z_2/2,\,z_2\bigr).
\]
Indeed
\(\sum_{a,c}\langle F_{ac},\partial_aF_{bc}\rangle
 =\frac12\sum_{a,c}\langle F_{ac},\partial_bF_{ac}\rangle\);
the remaining term is the third coordinate. Therefore
\begin{equation}
 \rho_{\rm tr}(x,y,z)=-2x-8y,\qquad
 \rho_{\rm tr}(\operatorname{div}_{\rm raw}B)\mathcal S
                         =-\operatorname{tr}B.
 \label{eq:trace-raw-functional}
\end{equation}
In particular its values on \(L,\Xi,D\) are \(0,0,-4\).
It kills the Maxwell equation direction
\(v_{\rm eq}=(0,0,1)\). The raw vector space has dimension
three, whereas its Maxwell quotient has dimension two.

Use the leading analytic path action without its coupling,
denoted \(\mathsf S\). On a finite torus \(\Lambda\), let
\(H\) be its left Hessian and \(d\) the infinitesimal gauge
insertion. All adjoints use counting measure and the fixed
colour metric. For \(\mu>0\) define, near the identity,
\begin{equation}
 \begin{aligned}
 \Gamma_\mu&=\tfrac12\log\det(H+dd^*+\mu^2I)
                       -\log\det(d^*d+\mu^2I),\\
 \mathcal R_\mu&=(H+dd^*+\mu^2I)^{-1},\qquad
 \mathcal G_\mu=(d^*d+\mu^2I)^{-1}.
 \end{aligned}
 \label{eq:trace-scalar-regulator}
\end{equation}
These are ordinary finite matrices. Their flat values are
positive; thus their Taylor coefficients are defined even
when no assertion about distant backgrounds is needed.
The regulator \(\mu\) is an auxiliary infrared mass in the
current lattice units.

Fix a unit colour \(\alpha\). For an anchored chart function
\(F\) write
\[
 \widehat F_{c\alpha,c\alpha}(p,q)=
 \sum_{x,y}e^{i(p\cdot x+q\cdot y)}
          D^2F(0)[e_{x,c,\alpha},e_{y,c,\alpha}].
\]
The vectors \(e_{x,c,\alpha}\) are unit link increments.
Coefficients such as \([z^2t]\) include the factorials.
For an anchored vector family \(F_b\) and an extensive scalar
\(A\), define
\begin{align}
 \mathcal L(F)&=\operatorname{Re}
       \sum_{a\ne c,b}i[z^2t]\,
          \widehat{F_b}_{c\alpha,c\alpha}(ze_a,-ze_a+te_b),
       \notag\\
 \mathcal H(F)&=\mathcal L(F)-\tfrac12\operatorname{Re}
       \sum_{a\ne c}[z^2]\,
          \widehat{F_c}_{c\alpha,c\alpha}(ze_a,-ze_a),
       \notag\\
 \mathcal Q(A)&=(\#\Lambda)^{-1}\operatorname{Re}
       \sum_{a\ne c}[z^2]\,
          \widehat A_{c\alpha,c\alpha}(ze_a,-ze_a),\notag\\
 \mathfrak h(Y)&=\mathcal H(\operatorname{div}_{\rm Haar}Y),
 \qquad
 R_\Lambda(\mu)=\mathcal Q(\Gamma_\mu),\qquad
 \mathcal D_\Lambda(\mu)=-\mu R_\Lambda'(\mu).
 \label{eq:trace-ward-reads}
\end{align}
For periodic kernels the waves read their lifted local
coefficients; the estimates below justify this Taylor read.
The subtraction in \(\mathcal H\) is the correction from the
link-base phase in the lattice difference. The divergence
\(\operatorname{div}_{\rm Haar}\) is divergence with respect
to product normalised Haar measure, not coordinate Lebesgue
measure.

\subsection{Norms and the finite-volume identity}

For \(a>0\) and \(r=1,2,3,4\), the weighted order-\(r\)
seminorm of a function anchored at \(o\) is
\begin{equation}
 |F|_{a,r}=
 \sum_{e_1,\ldots,e_r}
 |D^rF(0)[e_1,\ldots,e_r]|
 \prod_{\ell=1}^r
       e^{a|e_\ell^--o|_\infty} .
 \label{eq:trace-current-seminorm}
\end{equation}
For a displacement sum also over its output link, with the
same anchored weight. Direction and colour indices are included
in the sums. Every norm below is local per anchor.
% !TeX root = 4DYM_notes.tex
It suffices to use exponential weights in this auxiliary
seminorm. To check every polynomial moment explicitly,
for \(p\ge0\) put
\[
 C_{a,p}=\sup_{s\ge0}(1+s)^p e^{-as/2}<\infty.
\]
This follows by differentiating
\(p\log(1+s)-as/2\), or by exponential decay at infinity.
For any distances \(s_1,\ldots,s_r\ge0\),
\[
 \prod_{i=1}^r e^{as_i/2}(1+s_i)^p
                         \le C_{a,p}^r e^{a\sum_i s_i}.
\]
There is one further factor for the output distance
of a displacement. Thus the same seminorm at width
\(a\) controls every fixed polynomial moment with
exponential width \(a/2\). In particular a momentum
derivative multiplies a coefficient by powers of its
input or output displacement; all derivatives needed
in \eqref{eq:trace-ward-reads} are bounded by this
inequality. A periodic-image tail at distance \(R\)
has the additional factor \(e^{-aR/2}\), after its
fixed polynomial moments have been paid.

The widths in the current construction are chosen
before \(L\), below the native kernel widths. Hence
\(C_{a,p}\) is a profile constant, not a ratio loss.
The original regular and forced activity norms are
unchanged. This reduction concerns only the auxiliary
current jets; it will also pay the moment factors in
\eqref{eq:ordinary-common-moment-estimate}.

Fix a coarse level. Consider finite histories and periods
indexed by \(\nu\), with period tending to infinity, and
local current data \(J_{ab},Y_b,\mathcal E_b,q_a,t_{abc}\).
The notation \(A\equiv_4B\) means equality of all field
derivatives through order four at zero. Assume
\begin{align}
 D\mathsf S[Y_b]&\equiv_4\sum_a\nabla_aJ_{ab}-\mathcal E_b,
       \notag\\
 \operatorname{tr}J&=\sum_a\delta_aq_a,\qquad
 J_{ab}-J_{ba}=\sum_c\delta_ct_{abc}.
 \label{eq:trace-current-identities}
\end{align}
Here \(\delta_a=T_a-1\), \(T_aF(x)=F(x+e_a)\), and
\(\nabla_a\) is the improved central difference defined
in the next section. Suppose the seminorms of orders two
through four of \(J,Y,q,t\), and the first-order seminorm
of \(Y\), are bounded by \(C\), and those of
\(\mathcal E\) tend to zero. The scalar currents
\(J,q,t,\mathcal E\) are gauge invariant, and \(Y\) is a
gauge-covariant left link displacement. Their values at
the flat connection vanish. We specify the finite
matching conditions used below. The quadratic momentum
jet of \(D^2J_{ab}(0)\) is that of \(L_{ab}\), with the
free curvature replaced by the linear curl. Its
coefficients may have errors tending to zero.

For the displacement keep its anchor at zero. Define
\[
 \widehat Y_{b,c;d}(p_0;p)
   =\sum_{x,y}e^{-ip_0\cdot x+ip\cdot y}
               D(Y_b^0)_{(x,c)}(0)[e_{y,d}].
\]
This is a colour endomorphism; \(p_0\) and \(p\) are
independent output and input momenta. The required
linear momentum jet is
\begin{equation}
 \widehat Y_{b,c;d}(p_0;p)
    =i(p_b\delta_{cd}-p_c\delta_{bd})I_{\mathfrak g}
                         +O((|p_0|+|p|)^2),
 \label{eq:trace-leading-displacement-read}
\end{equation}
again allowing vanishing errors in the finitely many
displayed coefficients. Higher momentum coefficients
are bounded by the weighted seminorms; they need not
converge to those of one fixed lattice stencil. Gauge
covariance fixes the zero-momentum quadratic completion:
on constant Lie algebra inputs \(A,B\),
\begin{equation}
 \sum_xD^2(Y_b^0)_{(x,c)}(0)[A,B]
       =[A_b,B_c]+[B_b,A_c]+o(1).
 \label{eq:trace-quadratic-displacement-read}
\end{equation}
To see this, differentiate
\(Y(U^\gamma)_e=\operatorname{Ad}(\gamma_{e^-})Y(U)_e\)
once in a constant field and once in a gauge parameter
affine in position. Its linear gauge increment is
an arbitrary constant link field. The variation of
that increment at a constant background, tested with
the leading curl in
\eqref{eq:trace-leading-displacement-read}, gives
the two displayed brackets. A constant increment has
zero linear curl; thus the output adjoint variation
contributes zero at this order. Higher momentum terms
annihilate the affine test. One can first cut off
the affine gauge parameter and then remove the cutoff
using the exponential kernel weights. These are
moment reads of the lifted kernels. Equivalently,
the two brackets are the polarised constant-field
term of \(\partial_bA_c-\partial_cA_b+[A_b,A_c]\).
Thus \eqref{eq:trace-quadratic-displacement-read}
is not an independent renormalisation condition.
The next section constructs these data from the Wilson
action; they are not additional assumptions on the theory.

\begin{lemma}[Regulated Ward coefficient]
\label{lem:trace-regulated-ward}
For these data, \(\mathfrak h(Y)\to-12b_0\).
The estimates used in taking the limit are uniform in
the ultraviolet cutoff and physical period whenever the
displayed local bounds and the stated errors hold.
\end{lemma}
\begin{proof}
For the flow \(\Psi_s(U)=e^{sY(U)}U\) set
\[
 \mathcal Y_U(z)=\left.\partial_s\right|_0
          \log\bigl(\Psi_s(e^zU)\Psi_s(U)^{-1}\bigr),
                   \qquad \mathcal V=D\mathcal Y_U(0).
\]
Let \(\mathcal J_U\) be the mixed \(v,z\) derivative at
zero of
\(\log(\operatorname{gauge}_{e^v}(e^zU)U^{-1})\).
Differentiating gauge invariance once in each variable gives
\(Hd=-\mathcal T\), where
\(\mathcal T=D\mathsf S[\mathcal J_U]\). Hence
\[
          \mathcal R_\mu d
                  =d\mathcal G_\mu+
                              \mathcal R_\mu\mathcal T\mathcal G_\mu.
\]
Expand \(D^2(D\mathsf S[Y])\) by the product rule.
The terms containing \(D H[Y]\), together with the gauge
completion and ghost variations, give \(D\Gamma_\mu[Y]\).
The two terms \(H\mathcal V+\mathcal V^*H\) give
\(\operatorname{tr}\mathcal V-\mu^2\operatorname{tr}
(\mathcal R_\mu\mathcal V)\) and the gauge part just displayed.
The trace of the coordinate-flow differential is
\(\operatorname{div}_{\rm Haar}Y\); the adjoint-frame
commutators have zero trace. The remaining product-rule
terms retain the action gradient. Thus, exactly,
\begin{align}
 \tfrac12\operatorname{tr}
       \{\mathcal R_\mu D^2(D\mathsf S[Y])\}
                       -\operatorname{div}_{\rm Haar}Y
   &=D\Gamma_\mu[Y]-\mu^2\operatorname{tr}
                          (\mathcal R_\mu\mathcal V)
                                  +D\mathsf S[Z],\label{eq:trace-contact-identity}\\
 Z&=\tfrac12\operatorname{contr}
                   (\mathcal R_\mu,D^2\mathcal Y_U(0))
       -\operatorname{contr}\bigl(
          \mathcal G_\mu d^*(\mathcal V+\mathcal V^*)\mathcal R_\mu,
                                         \mathcal J_U\bigr).\notag
\end{align}
Contraction means summation over the indicated finite matrix
indices. In particular the last term has not been set to
zero before its derivatives are read.

We describe the cancellation of action variations under
\(\mathcal L\). At the required momentum order an invariant
quadratic scalar kernel is a colour multiple of
\(|p|^2I-p\otimes p\), by the calculation preceding
\eqref{eq:ordinary-alpha-definition}. Write \(X_{acb}\)
for the momentum-\(e_a\) coefficient of the diagonal
polarisation-\(c\) entry of \(DY_b\), with the scalar
kernel's colour indices contracted. Gauge covariance gives
the mixed coefficient \(-X_{cac}\) and \(X_{aaa}=0\).
The three contributions to the tested variation are
\[
              -2X_{aca}-X_{cac}+\sum_bX_{bcb},\qquad a\ne c.
\]
Their sum is zero: the first two sums equal
\(-3\sum_{b,c}X_{bcb}\), and the last equals its negative.
Higher momentum terms cannot enter this coefficient.
This applies to \(D\Gamma_\mu[Y]\), \(D\mathsf S[Z]\)
and the further action variation below. At fixed
\(\mu\in[\delta,1]\) their Taylor operations have uniform
bounds.

In the first term of \eqref{eq:trace-contact-identity},
insert \eqref{eq:trace-current-identities}. The coefficient
linear in the total momentum \(te_b\) selects \(a=b\).
Summing in \(b\) gives the trace of \(J\), which is a
forward difference and supplies a second total-momentum
factor. Its coefficient linear in \(t\) is zero.
The defect contributes \(o_\nu(1)\), uniformly on each
fixed interval \([\delta,1]\), by
\eqref{eq:trace-current-seminorm}. The same link-base
calculation accounts for the correction in \(\mathcal H\).

To compute the regulator contribution, put
\(\chi(x)=1_{\{x=0\}}\), and on site or link vectors define
\[
 {\mathsf M}_bh(x)=\chi(x)h_b(x),\quad
 (D_bh)_a(x)=(\delta_a\chi)(x)h_b(x),\quad
 T_b^{(r)}h(x)=(\delta_b\chi)(x)h(x).
\]
The superscript \(r=0,1\) specifies site or link vectors.
Set
\[
 W_b^{(1)}=\operatorname{sym}(\mathcal V_b-d{\mathsf M}_b-D_b)
                                      +\tfrac12T_b^{(1)},\qquad
 W_b^{(0)}=\operatorname{sym}({\mathsf M}_bd)-\tfrac12T_b^{(0)},
 \qquad\operatorname{sym}A=(A+A^*)/2.
\]
Use the inverse identity above and cycle the finite traces.
It gives
\begin{align*}
 \operatorname{tr}(\mathcal R_\mu\mathcal V_b)
   &=
       \operatorname{tr}\{\mathcal R_\mu(D_b-\tfrac12T_b^{(1)})\}
       +\tfrac12\operatorname{tr}(\mathcal G_\mu T_b^{(0)})
       +D\mathsf S[Z'_b]+\mathcal W_b,\\
 \mathcal W_b&=\operatorname{tr}(\mathcal R_\mu W_b^{(1)})
                    +\operatorname{tr}(\mathcal G_\mu W_b^{(0)}),\\
 Z'_b&=\operatorname{contr}
                (\mathcal G_\mu {\mathsf M}_b\mathcal R_\mu,\mathcal J_U).
\end{align*}
The symbol of \(\delta_b\chi\) in the total momentum is
\(e^{-it}-1=-it+O(t^2)\). The factor \(i\) in
\(\mathcal L\) therefore turns its linear coefficient
into one. The \(D_b\) term contributes
\(\mathcal Q(\operatorname{tr}\mathcal R_\mu)\);
the four \(T_b^{(1)}\) terms contribute
\(-2\mathcal Q(\operatorname{tr}\mathcal R_\mu)\);
the site terms contribute
\(2\mathcal Q(\operatorname{tr}\mathcal G_\mu)\).
Translation invariance changes an anchored sum to the
volume-normalised trace in \(\mathcal Q\).
Since
\[
 \partial_\mu\Gamma_\mu
       =\mu\operatorname{tr}\mathcal R_\mu
                              -2\mu\operatorname{tr}\mathcal G_\mu,
\]
we obtain
\begin{equation}
       \mathfrak h(Y)=\mathcal D_\Lambda(\mu)
                       +\mu^2\mathcal L(\mathcal W)+o_\nu(1),
          \qquad \delta\le\mu\le1.
 \label{eq:trace-regulated-ward}
\end{equation}

We give the estimates permitting removal of the regulator.
Flat ellipticity and the inverse identity give
\[
        \|\partial_q^jR_\mu(q)\|
                 \le C(|q|+\mu)^{-2-j},\qquad 0\le j\le3,
\]
for both inverse types. The constant depends on the fixed
coarse-profile ellipticity, exponential widths and current
seminorm bounds, but not on period, cutoff or \(\mu\le1\).
We record the current coefficients used in the two
cancellations. At zero field,
\eqref{eq:trace-leading-displacement-read} gives the
leading operator
\((DY_b^0h)_c=\chi(\partial_bh_c-\partial_ch_b)\).
Since \(d=-\partial\) to first momentum order,
subtracting \(d{\mathsf M}_b+D_b\) leaves
\(\chi\partial_b\) on each link component. Its
symmetric part is multiplication by
\(-\partial_b\chi/2\), cancelled by
\(T_b^{(1)}/2\). On sites the symmetric part of
\({\mathsf M}_bd=-\chi\partial_b\) is
\(\partial_b\chi/2\), cancelled by
\(-T_b^{(0)}/2\). Forward differences differ from
these leading symbols only at second momentum order.
For the first field insertion at zero momentum,
\eqref{eq:trace-quadratic-displacement-read} gives
\[
 (\mathcal V_bh)_c
       =\chi\{[A_b,h_c]-[A_c,h_b]\}
                              +\text{terms with a momentum}.
\]
The variation of \(d{\mathsf M}_b\) cancels the second bracket.
The first is \(\chi\operatorname{ad}(A_b)\), whose
symmetric part is zero in the fixed invariant real
metric. The site term is its negative and has the
same vanishing symmetric part. Thus the zero- and
one-field insertions start at momentum orders two
and one, respectively. The two-field insertion only
needs its uniform bound. Errors in the matched
coefficients contribute \(C_\delta\varepsilon_\nu\)
at fixed \(\delta>0\).
Taylor's formula and the weighted seminorms therefore bound
the zero-, one- and two-field insertions by \(Cz^2,Cz,C\),
where \(z=|q|+\mu\). Before extracting the three external
momentum derivatives their paired inverse factors cost
\(Cz^{-4},Cz^{-3},Cz^{-2}\). Each momentum derivative costs
at most one more inverse power of \(z\). Thus every resulting
term is bounded by \(Cz^{-5}\). For example the first
allocation is \(z^2z^{-4}z^{-3}=z^{-5}\); moving derivatives
to the Taylor numerator gives the same or a smaller bound.
The other allocations obey the same count.
In four dimensions
\[
 \int_{|q|\le1}(|q|+\mu)^{-5}\,\dd q\le C/\mu.
\]
The ball \(|q|\le\mu\) costs \(C/\mu\), and successive
dyadic annuli cost \(C\mu^{-1}2^{-j}\).
At fixed positive \(\mu\), Fourier sums converge uniformly
on compact \(\mu\)-intervals to their integrals. We have
\begin{equation}
          |\mu^2\mathcal L(\mathcal W)|
                       \le C\mu+o_\nu(1)
              \quad(\delta\le\mu\le1).
 \label{eq:trace-ward-regulator-error}
\end{equation}

We must also check the logarithmic coefficient for this
regulator, since \eqref{eq:drift-regulator} used a different
one. With scalar \(\mu^2I\), the flat inverse is still
bounded by \(C(|q|+\mu)^{-2}\). On \(|q|\ge\mu\), replacing
it by the massless inverse costs \(C\mu^2|q|^{-4}\).
There are no field derivatives of the added scalar mass.
The two leading vertices are exactly
\eqref{eq:drift-leading-vertices}. The lattice-vertex
replacement and the second-vertex terms have the same
integrable bounds as before. Thus for the coefficient
\(R_\Lambda/12\) the integrand has bounds
\[
 C(|q|+\mu)^{-4},\qquad
 |k_\mu(q)-h_b(q)|
                  \le C\{|q|^{-3}+\mu^2|q|^{-6}\}
                          \quad(|q|\ge\mu).
\]
The factor twelve is exact: there are twelve ordered
transverse pairs \(a\ne c\), and the Hessian is twice the
quadratic-action kernel, which is precisely the factor two
in \eqref{eq:ordinary-alpha-definition}. The small ball
costs \(O(1)\), and the shell errors integrate to \(O(1)\).
The contraction \eqref{eq:drift-pole-integral} consequently gives
\begin{equation}
          |R_\Lambda(\mu)-12b_0\log\mu|
                              \le C+o_\nu(1)
 \label{eq:trace-scalar-regulator-log}
\end{equation}
at each fixed positive \(\mu\). The error from Fourier
sums tends to zero on every compact positive interval;
the constant \(C\) is independent of \(\mu\).
No convergence of the underlying action to a unique
infinite-volume action is required for this uniform bound.

It would be invalid to differentiate the bounded error in
\eqref{eq:trace-scalar-regulator-log}. Instead fix \(u<0\)
and \(A>0\), and apply the ordinary mean value theorem to
\(R_\Lambda(e^t)\) on \([u-A,u]\). Its two endpoint
errors give some \(\mu_\Lambda\in[e^{u-A},e^u]\) with
\[
       |\mathcal D_\Lambda(\mu_\Lambda)+12b_0|
                                       \le(2C+o_\nu(1))/A.
\]
Use \eqref{eq:trace-regulated-ward} uniformly on this
fixed interval and then \eqref{eq:trace-ward-regulator-error}.
It follows that
\[
       \limsup_\nu|\mathfrak h(Y)+12b_0|
                                      \le2C/A+Ce^u.
\]
Let \(u\to-\infty\) and then \(A\to\infty\).
This proves the claim with the required order of limits.
\end{proof}
% !TeX root = 4DYM_notes.tex
\section{Transporting the current and its fourth-order Ward coefficient}
\label{sec:trace-current-transport}

We now supply the current data required in
Lemma~\ref{lem:trace-regulated-ward} and transport its scalar
coefficient along the actual Wilson history. All coefficient
norms in the conclusion are fixed real coefficient norms
in the free spaces. The local estimates use
\eqref{eq:trace-current-seminorm} and the regular and normalised
forced norms of the insertion construction.
The finite-period arguments follow
\cite[Lemma~18.5, PDF pp.~210--213]{ArmstrongDraft}.
We make the last linear-algebra step explicit using two
independent columns of the already known divergence map.

\subsection{Current seeds and averaging kernels}

For a plaquette based at \(x\), put
\[
 f_{ab}(x)=\tfrac12\operatorname{proj}_{\mathfrak g}
                       (U_{\square(x;a,b)}-U_{\square(x;a,b)}^{-1}),
\]
with parallel transport to \(x\) understood in every colour
contraction. At the top level define
\begin{equation}
 J_{ab}=\frac{\delta_{ab}}4\sum_{c,d}\langle f_{cd},f_{cd}\rangle
                         -\sum_c\langle f_{ac},f_{bc}\rangle,
 \qquad
 (Y_b^x)_e=1_{\{e^-=x\}}f_{b,\operatorname{dir}e}(x).
 \label{eq:trace-current-seeds}
\end{equation}
The first seed has trace zero and is symmetric. The second
is a left link displacement anchored at \(x\).
For shifts \(T_a\), define
\begin{equation}
 \nabla_a=\sum_{j=1}^4c_j(T_a^j-T_a^{-j}),\qquad
 (c_1,c_2,c_3,c_4)=(4/5,-1/5,4/105,-1/280).
 \label{eq:trace-central-difference}
\end{equation}
The identities \(2\sum_jjc_j=1\) and
\(\sum_jj^{2v+1}c_j=0\) for \(v=1,2,3\) show that it
agrees with a first derivative through polynomial degree
eight. Define the defect, without omitting any term, by
\(\mathcal E_b=\sum_a\nabla_aJ_{ab}-D\mathsf S[Y_b]\).

Let \(L\) be an odd block ratio. In this paragraph \(B_a\)
is a shift polynomial, distinct from the read radius \(B_m\).
Set
\[
 P_a=L^{-1}\sum_{j=0}^{L-1}T_a^j,\quad
 B_a=\sum_{j=1}^4c_j\sum_{v=-j}^{j-1}T_a^v,\quad
 {\mathsf N}_a=B_a(T_a^L)P_a,\quad Q_0=\prod_aP_a{\mathsf N}_a .
\]
The notation \(B_a(T_a^L)\) replaces every shift by its
\(L\)-th power. With \(F_a=\prod_{c\ne a}P_c{\mathsf N}_c\), put
\(Q_a=B_aP_aF_a\) and \(Q_a^{\rm edge}={\mathsf N}_aF_a\).
The identities \(B_a(T_a-1)=\nabla_a\) and
\(P_a(T_a-1)=L^{-1}(T_a^L-1)\) give
\begin{align}
 Q_0\nabla_a&=L^{-1}\nabla_a^{\rm coarse}Q_a,\qquad
 Q_0\delta_a=L^{-1}\delta_a^{\rm coarse}Q_a^{\rm edge},
       \notag\\
 Q_a-Q_0&=\delta_a^{\rm coarse}Q_a^{\rm tr}.
 \label{eq:trace-averaging-identities}
\end{align}
The last quotient is a finite Laurent polynomial.
To check divisibility, reduce modulo \(T_a^L-1\):
multiplication by \(P_a\) averages all residue classes,
and both \(B_a\) and \(B_a(T_a^L)\) have mass one.
Thus the difference has zero remainder. The kernel \(Q_0\)
has mass one, support radius at most \(56L\), and absolute
row sum at most \(32^4\). These follow by summing the finite
coefficient lists and multiplying their absolute row sums.
Positivity of these averaging coefficients is not needed.

Let \(\mathcal IF=F(U_{\rm stat}(V))\), using the actual
stationary graph of Section~\ref{sec:actual-stationary}.
To push a displacement, differentiate the actual framed
predictor. Write \(R_f\) for its fine differential with
the root frame fixed, and \(R_0\) for the left logarithmic
differential of that frame. These are defined by differentiating
the predictor and the frame functions already constructed
on the same patch; their kernel and derivative bounds are
those of that stationary construction. For the coarse gauge
insertion
\((d_V\phi)_e=\phi(e^-)-\operatorname{Ad}(V_e)\phi(e^+)\),
the product rule gives
\[
             \mathcal I_*Y=(R_f+d_VR_0)Y(U_{\rm stat}(V)).
\]
Keeping the frame term is necessary; gauge invariance only
removes its action variation after this differential is taken.
Transfer the current and its difference potentials by
\begin{align}
 J'_{ab}&=L^4Q_a\mathcal IJ_{ab},&
 Y'_b&=L^5Q_0\mathcal I_*Y_b,&
 \mathcal E'_b&=L^5Q_0\mathcal I\mathcal E_b,\notag\\
 q'_a&=L^3Q_a^{\rm edge}\mathcal Iq_a
                              +L^4Q_a^{\rm tr}\mathcal IJ_{aa}.
 \label{eq:trace-current-transfer}
\end{align}
All kernels act on the lifted anchor translates before their
coarse restriction. Differentiating the stationary equation
eliminates the fine variation. The envelope identity and
classical matching then preserve
\(D\mathsf S[Y_b]\equiv_4\sum_a\nabla_aJ_{ab}-\mathcal E_b\).
The identities \eqref{eq:trace-averaging-identities} preserve
\(\operatorname{tr}J=\sum_a\delta_aq_a\).
Starting from zero, the antisymmetric difference potentials
acquire \(L^4Q_a^{\rm tr}\mathcal IJ_{ab}\) in direction \(a\)
and its negative with \(a,b\) exchanged, besides their
transferred old values. This proves the second identity
in \eqref{eq:trace-current-identities} as well.

\begin{lemma}[Local current bounds at a fixed coarse level]
\label{lem:trace-current-local-bounds}
At a fixed sufficiently late coarse level \(k\), the transferred
currents satisfy the hypotheses of
Lemma~\ref{lem:trace-regulated-ward} as the transfer depth
and then the physical period tend to infinity. Their
weighted seminorm constants may depend on \(k\), but are
independent of the ultraviolet cutoff and period.
\end{lemma}
\begin{proof}
% !TeX root = 4DYM_notes.tex
We work with jets through order four at the flat field.
Separating the matched finite polynomial from its
moment-cancelled remainder is the localisation principle
of \cite[Sect.~1]{BS15b}. The curvature division and
coefficient estimates below give the gauge-theory
version needed for the present seminorms.
The argument uses the free interpolation on curvature
slots, not an \(L^{-2}\) bound for an arbitrary fine
link increment. We first justify that distinction,
then give the finite system of inequalities.

\emph{Locality of the finite jet operations.}
Choose the profile widths \(a_j>0\) smaller than a fixed
fraction of the native path widths, the free-column
width in \eqref{eq:all-scale-linear-difference-bounds},
and the complex strip of the block quotient inverse
in \eqref{eq:ordinary-block-quotient-energy}. All these
widths are chosen before \(L\). Shrink them once for
the finitely many products through order five.
For adjacent widths also require
\(L a_{j_f}\ge32a_{j_c}\), a lower ratio floor.

Here only flat jets, rather than a fixed nonlinear
amplitude ball, require this pre-ratio strip.
Differentiating the stationary equation gives
\(D\rho=-{\mathsf H}^{-1}{\mathsf H_{yx}}\) and the finite partition formula
in the proof of
Proposition~\ref{prop:stationary-response-locality}.
At orders two through four every resulting vertex
is a derivative of the native action through order
five. Its path expansion has the original absolute
budget times a fixed power of path length. Such powers
are paid by shrinking the width.

The inverses in these flat jet formulas can be taken
on the joint gauge quotient, in block energy units,
as in \eqref{eq:ordinary-block-quotient-energy}.
To verify this for each differentiated equation,
differentiate the joint gauge identity
\(R^T\operatorname{grad}S=0\) along the stationary
curve. Every lower derivative of the full joint
fine gradient is zero on that curve. Thus its
order-\(r\) source, after the linear Hessian term
has been separated, annihilates the flat gauge
directions. The quotient inverse applies to this
source. Restoring the tree slice differentiates
the same gauge transformation along the finite
block trees and frame patches. Its coefficients
are finite path sums. This restores the actual
frame terms in \(\mathcal I_*Y\), rather than
discarding a gauge part of that displacement.

The quotient inverse has the fixed complex strip
proved in the classical matching argument. For a
jet with \(r\) inputs reduce that strip by a factor
at least \(r+1\), since an internal momentum is a
sum of those inputs. Fourier inversion gives
exponential coefficient bounds. Passing from
operator bounds to component sums costs only the
number of residue and colour indices in a block,
at most a fixed power of \(L\). Tree lengths are
at most \(C_jL\). For the frame equation the flat
rooted Hessian on \(m-1\) sites has inverse
\(\frac m2(I+\boldsymbol1\boldsymbol1^T)\), since
its Hessian is \(2I/m-2\boldsymbol1\boldsymbol1^T/m^2\).
Here \(m\le C_jL^4\). The predictor Hessian is the
identity. Their finite implicit derivatives therefore
have the same polynomial counting bound.
Consequently all nonlinear finite-jet coefficients
needed below have a bound \(A_{j,L}<\infty\), uniform
in history, cutoff and period. Their exponential
width is the chosen profile width. We reserve the
letter \(C_j\) for the linear-column and geometric
constants that precede the ratio. We shall use
\(A_{j,L}\) only in a later entry condition.

\emph{From gauge invariance to the curvature factors.}
For a gauge-invariant anchored scalar whose lower
field jets vanish, its first nonzero order-\(r\)
kernel annihilates a linear gauge increment in every
input. The input momenta are independent, since the
anchor is retained. The independent-momentum part
of Lemma~\ref{lem:ordinary-curl-division} writes
that kernel with a curl in each input. The Fourier
division replaces a mode by a geometric sum and
costs a fixed power of its displacement. The
exponential reserve just chosen pays this cost.
Complete the resulting coefficients using the
transported curvature slots
\eqref{eq:ordinary-curvature-slot}. Their linear
terms are the required curls. Their derivatives
through order four are bounded by their coefficient
sum times a fixed path-length polynomial.

For example first complete the quadratic jet and
subtract that gauge-invariant function. The new
cubic kernel has no lower jet and is transverse
in its three inputs. Complete it and repeat for
the quartic kernel. This is a linear triangular
reading of the four-jet. A change in the completed
quadratic jet can change orders three and four,
but an order-four jet cannot change order two.
The same construction for a covariant displacement
transports the output colour to its output link
base and uses equivariant, rather than invariant,
colour contractions. Its first nonzero input jet
is again transverse: the output adjoint variation
contains a lower jet, which is zero at that stage.
All constructions are on the integer cover.

This division does not assert that a raw third
derivative by itself is transverse. The lower
completed jets must first be removed. Their
contributions are the off-diagonal terms below.
The finite norm constants follow from the explicit
geometric-sum division and the finite product rule.
In moving a fine coefficient to the coarse lattice,
its exponential weight uses at most \(a_{j_c}/L\)
of its fine displacement. The width floor above
therefore leaves the reserve for all these
divisions and polynomial moments.

\emph{The matched moments and the two-point defect.}
Gauge invariance and simplicity give
\(\mathcal E(1)=D\mathcal E(1)=0\).
The quadratic jet of the Wilson seed defect has
zero momentum coefficients of total degrees at
most three. Indeed each linear curvature is
\(\partial_aA_b-\partial_bA_a\) to lowest order,
and its degree-three vector coefficient is the
polynomial identity
\[
 \sum_a\partial_a\left\{
   \frac{\delta_{ab}}4\sum_{c,d}\langle F_{cd},F_{cd}\rangle
                -\sum_c\langle F_{ac},F_{bc}\rangle\right\}
                  =-\sum_{a,c}\langle F_{bc},\partial_aF_{ac}\rangle.
\]
This follows by the product rule and the linear Bianchi identity.
It is also the leading coefficient of
\(D\mathsf S[Y_b]\), since the action gradient is
\(-\sum_a\partial_aF_{ac}\). Lower momentum degrees
are absent. The difference \(\nabla_a\) has the
same first-derivative coefficient, so this verifies
the seed moment cancellation directly.

It is preserved by transfer. In the principal
alias, the fine momentum is \(p/L\) and the
interpolation symbol is analytic. A nonprincipal
alias has a filter zero of order at least \(j\).
The pole-free formulas
\eqref{eq:linear-E-column}--\eqref{eq:linear-T-column}
show this explicitly on curls; their nonprincipal
denominators are bounded away from zero in rescaled
variables. With \(j\ge7\) no such alias contributes
to a degree-three coefficient. The averaging
polynomials add no pole. Thus a quadratic kernel
vanishing through momentum degree three retains
that property. Fourth-order classical matching
preserves the defect identity through ordinary
field degree four, so this calculation applies to
the actual transferred defect.

After its two curl divisions the quadratic defect kernel
has zero constant and first moments. Apply the common
estimate \eqref{eq:ordinary-common-moment-estimate} with
\((d_1,d_2,s)=(2,2,2)\). The two derivative placements
calculated there both cost \(C_jL^{-6}\); the
dimension-five rescaling gives \(C_j/L\).
The second difference of a curvature input uses
\(\nabla_f^3E\), already among the admitted linear-column
estimates. The kernel moments are paid by the
exponential-width reduction following
\eqref{eq:trace-current-seminorm}. Thus no higher
free-column derivative or additional norm is being used.

For the cubic and quartic diagonal parts there
are respectively three and four curvature
inputs, each with price \(C_jL^{-2}\). Their
dimension-five prices are \(C_jL^{-1}\) and
\(C_jL^{-3}\). A nonlinear chain-rule partition
has fewer derivatives of the incoming function;
the finite-jet locality calculation bounds each
such coefficient by \(A_{j,L}\).

For a tensor family take the maximum of the finitely
many component seminorms. Put
\(\lambda_m=\log\widehat\beta_m\) and
\(u_{m,r}=|\mathcal E_m|_{a_{J(m)},r}/\lambda_m^r\),
for \(r=2,3,4\). Increase entry so that
\(1\le\lambda_m\) and
\(\lambda_{m+1}/\lambda_m\le2\).
The preceding calculations give the explicit
nonnegative system
\begin{equation}
 \begin{pmatrix}u_{m,2}\\u_{m,3}\\u_{m,4}\end{pmatrix}
 \le
 \begin{pmatrix}
 4C_j/L&0&0\\
 4A_{j,L}/\lambda_m&8C_j/L&0\\
 4A_{j,L}/\lambda_m^2&8A_{j,L}/\lambda_m&16C_j/L^3
 \end{pmatrix}
 \begin{pmatrix}u_{m+1,2}\\u_{m+1,3}\\u_{m+1,4}\end{pmatrix}.
 \label{eq:trace-explicit-defect-matrix}
\end{equation}
For instance the middle off-diagonal entry before
the bound is
\(A_{j,L}\lambda_{m+1}^2/\lambda_m^3\).
This displays the inverse logarithm supplied by
the missing field derivative. The order-two
diagonal can equivalently be bounded by
\begin{equation}
 C_j\left\{\lambda_m^{-2}
       +\frac{\lambda_{m+1}^2+1}{L\lambda_m^2}\right\}.
 \label{eq:trace-defect-order-two}
\end{equation}
The extra positive terms in this last bound are
not additive defect sources.

Choose the ratio with \(L\ge128\max(1,C_j)\),
and only then choose entry with
\(\lambda_m\ge64A_{j,L}\). The diagonal row
contribution is at most \(1/4\), and the
off-diagonal contribution is at most \(1/4\).
Thus the maximum of the three components contracts
by \(1/2\). The Wilson seed and its improved
difference operator have fixed finite stencils;
their jets through order four have a fixed bound.
After \(n\) transfers,
\begin{equation}
          |\mathcal E_k|_{a_{J(k)},r}
                        \le C\lambda_k^r2^{-n},\qquad r=2,3,4.
 \label{eq:trace-defect-full-seminorm-decay}
\end{equation}
At fixed \(k\), this tends to zero with the
ultraviolet transfer depth in the actual weighted
seminorm, not only on a set of curvature tests.

\emph{Bounds for the currents that do not vanish.}
Let \(\bar J,\bar Y\) be the same fixed plaquette
seed formulas, installed on the lattice of the
current level. Their lowest curvature coefficients
are preserved by the linear transfer. For \(J\),
the two curvature factors give \(L^{-4}\), cancelled
by \(L^4\), and \(Q_a\) has mass one.
For \(Y\), a useful exact counting identity is
\[
 \sum_{v\in\Z^4}Q_0(Lv-z)=L^{-4}.
\]
Here \(Q_0(z)\) denotes the coefficient of its
shift polynomial. Each factor \(P_a\) is uniform
on residue classes modulo \(L\). The second
\(P_a\) and \(B_a(T_a^L)\), of mass one, preserve
that uniformity. This proves the identity.
The combined framed differential at zero is
\(\widehat Q=Q+d_{0,c}P\) from
\eqref{eq:actual-predictor-linear-frame}.
On a constant fine link increment \(v\),
\(Qv=Lv\) and \(Pv=0\): the mean coordinate
of both the centred patch and \(w\) is zero
in the displayed formula for \(t_z\).
Consequently the leading displacement coefficient
is multiplied by \(L^5L^{-4}LL^{-2}=1\).
More explicitly, summing its coarse output sites
and reindexing the fine output by a coarse
translation replaces its anchor weight by
\(\sum_vQ_0(Lv-f^-)=L^{-4}\).
The remaining sum is one row of \(\widehat Q\)
applied to the constant vector \(L^{-2}F_{b,c}\),
which is \(L^{-1}F_{b,c}\). Multiplication by
\(L^5L^{-4}\) gives \(F_{b,c}\).
Its input transversality precludes a degree-one
term involving output momentum alone. This proves
\eqref{eq:trace-leading-displacement-read}.
Its quadratic completion follows by the gauge
calculation following that equation.

Subtract these fixed stencils to bound
\(J-\bar J\) and \(Y-\bar Y\).
The quadratic curvature kernel of the first has
zero constant moment. The linear curvature kernel
of the second has zero constant moment in the
joint output and input displacements.
Their highest-jet transfer bounds, and those of
\(q,t\), are
\begin{equation}
 \begin{array}{c|c|c}
 \text{remainder}&\text{field order}&
                                \text{diagonal price}\\ \hline
 J-\bar J&2&C_j/L\\
 J-\bar J&r=3,4&C_jL^{4-2r}\\
 Y-\bar Y&1&C_j/L\\
 Y-\bar Y&r=2,3,4&C_jL^{2-2r}\\
 q,t&r=2,3,4&C_jL^{3-2r}.
 \end{array}
 \label{eq:trace-current-diagonal-table}
\end{equation}
The scalar entries follow by inserting the
curvature factors as for the defect, with the
dimension-four or dimension-three prefactor.
The zero constant moment in the first row pays
one additional fine difference.

We detail the displacement count in the other
rows. The combined predictor differential has
row sum at most \(C_jL\): the straight average
has sum \(L\), and every term of its frame
correction is an average of coordinate paths
of length at most \(C_jL\).
The averaging kernel satisfies
\[
 \|Q_0\|_\infty\le C_jL^{-4},\qquad
                       \|\nabla_f Q_0\|_\infty\le C_jL^{-5}.
\]
Indeed \(P_a^2\) is the triangular convolution
of two uniform block weights; its supremum
is at most \(L^{-1}\), and its first difference
at most \(L^{-2}\). The remaining factors have
bounded absolute mass. Hence an order-\(r\)
curvature coefficient has price
\(L^5(C_jL^{-4})(C_jL)(C_jL^{-2})^r
=C_jL^{2-2r}\).
For the residual linear coefficient, its zero
moment moves one difference to \(Q_0\) or to
the interpolated curvature. Both give one
further \(L^{-1}\). To see the full coefficient
sum, fix its finite direction and colour indices
and let \(K_{r,s}\) be its curvature kernel in
the output and input displacements from its
anchor. The coefficient after transfer is
\[
 L^5\sum_{f,r,s}\widehat Q_{ef}
      Q_0(Lx-f^-+r)K_{r,s}(d_{1,f}E)_{f^- -r+s}.
\]
Here \(x\) is the coarse anchor, \(e\) the
coarse output, and the coarse input index of
\(d_{1,f}E\) is suppressed. Since
\(\sum_{r,s}K_{r,s}=0\), subtract the value
at \((r,s)=(0,0)\) from the other two factors.
Its first difference is bounded by
\[
 C_j\{L^{-5}L^{-2}+L^{-4}L^{-3}\}
                          (|r|_1+|s|_1).
\]
The kernel moment sums this last factor.
The prefactor and predictor row sum give
\(L^5(C_jL)L^{-7}=C_j/L\).
The spatially weighted coefficient sums obey
the same estimate, by the width and moment
reserves stated above. This includes the frame
correction; it has not been dropped from the
row sum. Nonlinear derivatives of this
differential act on lower incoming jets and
have price \(A_{j,L}\).

There is no claim of contraction for the
nonzero matched coefficients themselves.
They are the fixed seed coefficients just
computed. Transporting a fixed reference
stencil and subtracting that same stencil
at the new level gives a bounded forcing
with price \(A_{j,L}\).
For the remainders in the first four rows
of \eqref{eq:trace-current-diagonal-table},
divide an order-\(r\) norm by \(\lambda_m^r\).
For \(q,t\), divide by \(\lambda_m^{r+1}\).
The extra logarithm pays their same-order
source \(Q^{\rm tr}\mathcal IJ\); their
homogeneous transfer still has the negative
power shown in the table. The finite Laurent
quotient \(Q^{\rm tr}\) has bounded absolute
mass: on each residue class, division by
\(T_a^L-1\) sums at most a fixed number of
coefficient blocks, since its radius is
\(O(L)\). Its weighted coarse radius is
bounded independently of \(L\).

Every remaining homogeneous off-diagonal
entry now has price at most
\(C A_{j,L}/\lambda_m\): it either lowers
field order or comes from \(J\) into a
difference potential. There are only finitely
many components through order four. Enlarge
\(A_{j,L}\) to include this finite count.
The ratio floor controls the diagonal
entries by \(1/4\), and later entry controls
their off-diagonal row sums by \(1/4\).
It also bounds the normalised forcing by
one. If \(U_m\) is the maximum of these
normalised remainder norms, we obtain
\[
                         U_m\le\tfrac12U_{m+1}+1.
\]
At the Wilson top level the remainder is
zero and \(q=t=0\). Iterating gives
\(U_k\le\sum_{s\ge0}2^{-s}=2\).
Restoring the fixed reference stencils shows
that all required current seminorms are
bounded at fixed \(k\), independently of
transfer depth and period. Only the finite
powers of \(\lambda_k\) enter this bound.

All ratio floors involve the pre-ratio
constants only. The entry condition involving
\(A_{j,L}\) is imposed after that ratio,
as for the other finite-order requests in
the common calendar. Neither depends on
the continuum hierarchy or its selected
phase. These are explicit instances of the
already admitted current estimates, with
the same single Wilson family.

Finally, periodising a rooted kernel with
width \(a_{J(k)}>0\) changes any fixed local
moment read by at most \(C_ke^{-a_{J(k)}\ell_{k,M}/2}\):
the polynomial factors in its tail are paid
by half the width. The local bounds are
uniform before this limit. Thus first taking
the cutoff limit at fixed period, and then
the prescribed period refinement, preserves
\eqref{eq:trace-defect-full-seminorm-decay},
the current bounds and the matched coefficients.
These are precisely the explicit inputs of
Lemma~\ref{lem:trace-regulated-ward}.
\end{proof}

% !TeX root = 4DYM_notes.tex
\subsection{Matching with coefficient control}
\label{subsec:ward-coefficient-control}

The local radial read and the coordinate of a matched physical
source are different quantities. A regular remainder whose radial
read is zero can change the latter coordinate during integration.
We need a bound on that change in addition to the root bound of
Section~\ref{sec:position-local-matching}.

We use the same selected coefficient graphs. At finite terminal
cutoff \(n\) and physical side \(M\), write
\(\mathsf Z_{n,M,k,d}\) for their map from the current
threshold coordinate at level \(k\) to the reference coordinate.
Thus a selected column with reference coordinate \(a\) has
current coordinate \(\mathsf Z_{n,M,k,d}^{-1}a\), including its
regular and forced graph remainders. The norm of a remainder
\(R\) is always \(\|R\|_{\rm ins}\): the sum of its anchored
regular prices and its forced prices divided by the original
\(\sqrt{w_k}\), with all addresses and inventories retained.
The coefficient norm is the \(\ell^1\) norm in the fixed radial
basis. An exact marked representation of a source is the
coefficient of one marker in its full physical density, as in
\eqref{eq:one-insertion-density}.

\begin{lemma}[Matching with coefficient control]
\label{lem:ward-matching-coefficient-control}
Fix one of the finite degree requests \(d=3,4,5\), at sufficiently
late profiles. Suppose that a source born at level \(k<n\)
has current radial read \(y\), a full marked remainder \(R\),
and
\[
                         |y|_1+\|R\|_{\rm ins}\le q.
\]
Use the same backward matching as in
Proposition~\ref{prop:position-local-matching}, and denote its
reference coefficient by \(a\). After imposing a sufficiently
large physical-side threshold for this fixed finite scale
interval, one has
\begin{equation}
 \sqrt{P_k}\,
       |y-\mathsf Z_{n,M,k,d}^{-1}a|_1
                         \le4(1+C_{J(k)})q.
 \label{eq:ward-matching-sharp-bias}
\end{equation}
The coefficient is chosen once, before any lower stopping
level or replicated-cell request. A physical replicated
comparison additionally requires the response data defined
below; its separate budget is not part of the present
coefficient assertion. The physical-side threshold
may depend on the fixed scale interval and requested accuracy,
but not on the terminal cutoff \(n>k\). All local prices are
independent of \(n,M\).
\end{lemma}
\begin{proof}
For \(d=3\), simplicity gives \(V_3=\C1\) and hence
\(\mathcal F_3=\C1_3\):
every nonconstant invariant needs at least two
curvature factors, hence engineering weight at
least four. The nonconstant primary coordinate
is therefore absent. The unit is transported
exactly; a zero-read remainder can nevertheless
generate a unit coefficient under integration.
The following coefficient calculation retains
that coefficient. Its low-order read uses at
most \(d+1=4\) derivatives in this case, within
the already chosen marked order at least six.

\emph{Separate the two one-step prices.}
Let \(c=m\), \(f=m+1\), \(L=L_m\). The proof of
Lemma~\ref{lem:insertion-local-prices} gives a smaller price
for the coefficient than for the whole remainder. More
precisely, for any prescribed \(\xi>0\), the canonical output
of an input \((v,R_f)\), with \(r_f=\|R_f\|_{\rm ins}\),
can be chosen to satisfy
\begin{align}
 |u-T_m\mathsf S_Lv|_1
   &\le \frac{\xi}{\sqrt{P_f}}(|v^\circ|_1+r_f),\notag\\
 r_c&\le K_j|\mathsf S_Lv^\circ|_1+\mu_Lr_f,
                  \qquad \mu_L=L^{-3/8}.
 \label{eq:ward-separated-atom-prices}
\end{align}
Here \(K_j\) is fixed before \(L\). The ratio is made large
enough for the second inequality, and the coupling entry is
then chosen for the first. The unit is transported exactly.

We give the additional coefficient calculation. Only the
finitely many radial derivatives of order at most \(d\) are
read. In the zero-fluctuation, linear-interpolation term the
primary read is exactly \(T_m\mathsf S_Lv\).
For a nonwrapping remainder it is zero: expanding the radial
gauge interpolation through weight \(d\) substitutes a radial
jet into the fine function, and every term of those weights
is killed by \eqref{eq:marked-jet-null}. The Fourier
interpolation calculation in the proof of the leading-map
bound is needed here; an arbitrary interpolation of the
potential would not establish this cancellation.

For clarity, write the fluctuation in fine link units as
\(\widehat\beta_f^{-1/2}\zeta\). In every fixed Gaussian
moment its size in a threshold-normalised direction is
\(P_f^{-1}\zeta\). If \(G\) is one of the functions to be
read and \(D^\alpha\) is a word of at most \(d\) radial
derivatives, the integral formula
\[
 D^\alpha G(A+\widehat\beta_f^{-1/2}\zeta)-D^\alpha G(A)
 =\widehat\beta_f^{-1/2}\int_0^1
       D(D^\alpha G)(A+t\widehat\beta_f^{-1/2}\zeta)[\zeta]
                                                        \,\dd t
\]
costs at most one additional input derivative and one
factor \(P_f^{-1}\). On the central fluctuation region
the entire segment is inside the input good chart by
the central cutoff inclusion; its complement is paid
by the warning estimate below. Integrating \(|\zeta|\) and its fixed
score polynomials uses the retained precision of
Section~\ref{sec:central-derivative-estimates}. Their moments
are bounded by the fixed-order Gaussian moment formula; no
supremum over the fluctuation box is used. Nonlinear
interpolation instead has relative price \(C_{j,L}h_c\),
by its quadratic error and division by one input threshold.

The unmarked density is kept in this calculation. Its
constant background factor cancels against the same
unmarked denominator. A difference of a remaining loop
factor in a fluctuation direction has price at most a
fixed polynomial in its allowance \(\Lambda_f\), times
\(P_f^{-1}\). The smooth interaction and chart differences
give \(\widehat\beta_f^{-c_j}P_f^{b_j}\), for fixed
\(c_j>0,b_j<\infty\). On warning and packet terms the
held physical reads are differentiated through the measure;
the resulting scores are polynomial, while their price
retains \(\exp(-c_{j,L}P_f^2/(1+B_f)^{b_j})\).
For example a forced mark is written
\(k_\alpha=(k_\alpha/\sqrt{w_f})\sqrt{w_f}\);
after dividing an outgoing forced price by \(\sqrt{w_c}\),
the exponent gained is still a negative fixed multiple of
\(P_f^2/(1+B_f)^{b_j}\). Its empty-inventory output is
regular and is estimated by the same integrated score
bound. No forced activity is differentiated.

Summing the rooted words and the finite radial tests
therefore gives, per unit input price, the bound
\begin{equation}
 C_{j,L}(1+B_f)^{b_j}
 \left\{\frac{(1+\Lambda_f)^{a_d}}{P_f}
       +\widehat\beta_f^{-c_j}P_f^{b_j}
       +\widehat\beta_f^{b_j}P_f^{b_j}
           e^{-c_{j,L}P_f^2/(1+B_f)^{b_j}}\right\}.
 \label{eq:ward-sharp-coefficient-income}
\end{equation}
Here \(a_d\) depends only on the fixed degree of the read;
larger derivative orders of the output remainder do not
increase it. Support moments in this formula are summed
with two units of covering reserve. The root is anchored,
so there is no volume factor. With
\(\Lambda_f\le2P_f^{\theta_{J(f)}}\), choose the initial
late profile so that \(a_d\theta_{J(f)}<1/16\).
At each fixed adjacent profile and fixed \(L\), all fixed
constants and powers of \(1+B_f\) are at most
\(P_f^{1/16}\) after entry. The first term is then
\(O(P_f^{-7/8})\); the other two are smaller than every
fixed negative power of \(P_f\). This is at most
\(\xi P_f^{-1/2}\) after a further entry delay.

For regular supports exceeding the local reading radius,
retain the whole function. The radius may be a sufficiently
large fixed-profile multiple of
\((1+B_f)(\log\widehat\beta_f)^2\).
The unused covering exponent makes their coefficient price
smaller than the right side of
\eqref{eq:ward-sharp-coefficient-income}. For the remainder
price one instead keeps the coarser split of
Lemma~\ref{lem:insertion-local-prices}, which gives
\(\mu_L\). This explains why no factor \(L^{-1}\),
independent of the coupling, is needed in the first line
of \eqref{eq:ward-separated-atom-prices}: the free low
read cancels exactly, whereas the whole higher-weight
function does not. The derivative count for the low read
is at most \(d+1\); the remainder count remains
\(\max(q_c,d+1)\).

\emph{Comparison of exact representations on large tori.}
The output in \eqref{eq:ward-separated-atom-prices} need
not be the prescribed coarse representation. Suppose that
two legal marked representations give exactly the same
physical density. Test their difference near the identity
with a radial polynomial field, cut off outside a cover
box of radius \(R\). Take \(\ell_{c,M}>20R\), so that
this test and every short support have a single lift.
All forced terms vanish there and the regular masks are
one. Divide by the common positive unmarked density.
The primary section reads its coefficient, and every
regular support inside the box has zero read. Only
supports reaching its boundary remain. For such a support
\(X\),
\[
 \operatorname{cover}_{B_c}(X)
                 \ge R/(2B_c+1)-1.
\]
The cut-off radial directions grow polynomially in \(R\).
The reading bound and the spare covering exponent thus
bound the difference of coefficients by
\[
 C_{c,d}(1+R)^{b_{c,d}}
          e^{-2R/(2B_c+1)+2}(r_c^{(1)}+r_c^{(2)}).
\]
Given \(\tau>0\), first choose \(R\) to make this at most
\(\tau(r_c^{(1)}+r_c^{(2)})\), then choose the physical
side. Both choices are independent of the terminal cutoff.
In particular, if the prescribed coarse remainder has
norm \(r_c\), comparison with the canonical output gives
\begin{align}
 |u-T_m\mathsf S_Lv|_1
 \le{}&\frac{\xi}{\sqrt{P_f}}(|v^\circ|_1+r_f)\notag\\
 &+\tau\{r_c+K_j|\mathsf S_Lv^\circ|_1+\mu_Lr_f\}.
 \label{eq:ward-exact-representation-price}
\end{align}
This is an arbitrarily small large-volume comparison,
not an exact identification of a periodic field with a
nonperiodic polynomial potential.

\emph{The variable cone.}
Subtract the full selected column with reference coordinate
\(a\). Let \(x_m,r_m\) be the coordinate and remainder
norm of this signed difference. The old fixed-cone
calculation, including its terminal absorption, gives
\begin{equation}
 |x_m|_1+r_m\le \frac{t}{G_m},\qquad
 t=4(1+C_{J(k)})q,\qquad
 G_m=({\varepsilon}_m/{\varepsilon}_k)^{1/8},
       \quad \operatorname{ref}(d)\le m\le k.
 \label{eq:ward-total-residual-budget}
\end{equation}
The coordinate at \(\operatorname{ref}(d)\) is zero.
Put \(z_m=\sqrt{P_m}\) and \(I_m=(G_mz_m)^{-1}\).
We prove \(|x_m|_1\le tI_m\) towards the fine end.
The statement holds at the reference level. On one step,
the inverse leading-map bound, \eqref{eq:ward-exact-representation-price}
and \eqref{eq:ward-total-residual-budget} imply
\begin{align*}
 |x_f|_1
 &\le K_j|x_c|_1
       +K_j\xi\,tI_f
       +K_j\tau t
              \{G_c^{-1}+(K_jL^d+\mu_L)G_f^{-1}\}.
\end{align*}
The factor \(L^d\) is the bound for \(\mathsf S_L\),
including a possible unit. Choose \(\xi\le1/(3K_j)\)
before the ratio, and require \(L^{1/8}\ge4K_j\).
Entry gives \(z_f/z_c\le2\), so
\[
                 K_j I_c\le\tfrac12 I_f.
\]
For this fixed step choose the large-volume comparison
accuracy to satisfy
\[
 \tau\le \frac{I_f}{
   10K_j\{G_c^{-1}+(K_jL^d+\mu_L)G_f^{-1}\}}.
\]
Then
\[
 |x_f|_1\le(\tfrac12+\tfrac13+\tfrac1{10})tI_f
                                    =\tfrac{14}{15}tI_f<tI_f.
\]
There are finitely many steps up to \(k\). Intersecting
their physical-side tails preserves uniformity in \(n>k\).
At \(m=k\), where \(G_k=1\), this is
\eqref{eq:ward-matching-sharp-bias}.

All matching operations are linear before absolute
majorants and preserve the original masks and forced
eligibility. Changing a later stopping clock does not
change this coefficient. The additional construction
needed for its physical root comparison is proved in
Lemma~\ref{lem:physical-response-matching-root}.
\end{proof}

We record the normalisation used below. If
\(q\le C h_k^2P_k^{1/6}\), then
\begin{align}
 h_k^{-2}|y-\mathsf Z_{n,M,k,d}^{-1}a|_1
 &\le4C(1+C_{J(k)})P_k^{-1/3}
 \le P_k^{-1/4}.
 \label{eq:ward-normalised-matching-bias}
\end{align}
The last inequality is precisely the finite entry request
\(4C(1+C_{J(k)})\le P_k^{1/12}\).
It retains the coefficient created by the remainder.

There is also a useful short-window consequence for selected
columns. In the centred degree-four tensor sector and the
degree-five odd vector sector, simplicity leaves only the
critical nonconstant weight, with two curvature factors.
On these sectors
\(\mathsf S_L=I\) and
\(T_m=(h_c/h_f)^2I\). Apply
\eqref{eq:ward-exact-representation-price} to a selected
column, use its graph bound, and then take the period
tail. Its normalised current coordinates satisfy
\begin{equation}
 |h_c^{-2}x_c-h_f^{-2}x_f|_1
 \le \frac{8(1+C_{J(f)})}{\sqrt{P_f}}
                                    |h_f^{-2}x_f|_1.
 \label{eq:ward-critical-column-drift}
\end{equation}
This concerns the nonconstant coordinates; a tensor unit
is removed by centering. Over the window of
\eqref{eq:trace-radius-window}, \(\widehat\beta_m/b\to1\)
and \(P_m/P\to1\), uniformly, with
\(b=\widehat\beta_{n_i}\), \(P=P_{n_i}\).
Entry includes \(C_{J(m)}\le\log\widehat\beta_m\).
There are \(O((\log b)^{48})\) steps, hence the sum of
the relative errors is at most
\[
             C(\log b)^{49}P^{-1/2}\le P^{-2/5}
\]
after entry. Products are bounded by the exponential of
this sum; the reverse products use
\((1-t)^{-1}\le e^{2t}\) for \(0\le t\le1/2\).
Thus the normalised column drift across the whole window
is \(O(P^{-2/5})\). This argument also covers a profile
boundary and will compare current coordinates at \(k_i\)
with the radius-level Maxwell frame, without multiplying
all reference-coordinate losses from earlier scales.

% !TeX root = 4DYM_notes.tex
\subsection{Physical responses and joining a complete source}
\label{subsec:forced-source-response}

An absolute bound for a forced singleton replacement does not
specify its products with other physical sources. We therefore
retain the functions which produce that replacement. This costs
an additional amplitude budget, separate from the coefficient
budget in Lemma~\ref{lem:ward-matching-coefficient-control}.
The finite component regrouping below is done before taking
absolute values. Its positive majorant uses the rooted-family
recursion \eqref{eq:rooted-recursion}. This is the positive-tree
organisation used in \cite[Theorem~1]{FP07}; the displayed
recursion supplies the particular estimate needed here. We do
not apply a convergence criterion for an unconstrained polymer
gas to the original signed, exactly flagged gas.

\begin{definition}[A physical response family]
\label{def:physical-source-response}
Fix a level \(k\), an anchor \(o\), and suppress the level
indices. A response family consists of bounded Borel,
gauge-invariant functions \(L_X\), indexed by site footprints,
such that \(L_X\) reads only the links and final flags in
\(X\), vanishes when \(\chi_X=1\), and is zero as a function
unless \(o\in X\). Reading the final flags includes the halo
\(X^+\) of \eqref{eq:inspected-footprint}. Its reserved price is
\begin{equation}
 e(L)=\sum_X e^{(\kappa+2)n_B(X)}\|L_X\|_\infty.
 \label{eq:physical-response-price}
\end{equation}
The two extra units of support exponent pay for meetings
with the unmarked forced family. No derivative of \(L_X\)
is required.

For a source with forced replacements \(Q_\alpha\), its
forced profile means the functions
\[
 \mathcal Q(Y,I;U)=
 \sum_{\alpha:(Y_\alpha,B_\alpha)=(Y,I)}Q_\alpha(U).
\]
This sum retains each function, including its original
auxiliary tests. Define the response profile
\begin{align}
 (\mathcal H_KL)(Y,I;U)
 ={}&\sum_X\sum_{b}
    L_X(U)\mathbf1_{\{Y\cap\mathcal B_h(U)=I\}}
                                \prod_{\alpha\in b}K_\alpha(U).
 \label{eq:physical-response-profile}
\end{align}
The inner sum runs over nonempty finite families \(b\) of
pairwise disjoint old forced footprints, each meeting \(X\),
with
\[
       Y=X\cup\bigcup_{\alpha\in b}Y_\alpha,
       \qquad I=\bigcup_{\alpha\in b}B_\alpha.
\]
The source has response \(L\) when
\(\mathcal Q=\mathcal H_KL\) as functions of
\((Y,I,U)\), in addition to its exact density identity.
This functional equality is stronger than an equality of
total norms or of integrated densities.
\end{definition}

The labels for \((X,b)\), and for the several-source
components below, are chosen in the initial address book.
Their unmarked activities are zero. A new address has the
displayed union inventory and a tautological additional test;
all old tests remain inside the factors \(K_\alpha\).
The zero-address padding in
Definition~\ref{def:one-insertion-format} permits these
choices before amplitudes or configurations are specified.
Regrouping by a profile is only an intermediate sum; its
terms are recorded again at those fixed addresses.

\begin{lemma}[Exact response birth and its forced price]
\label{lem:physical-response-birth}
Suppose a bounded local physical source decomposes as
\(f=f^{\rm reg}+\sum_XL_X\), where \(f^{\rm reg}\)
already has a masked regular representation and \(L\)
is a response family. Then \(Ff\) has an exact
one-insertion representation whose forced profile is
\(\mathcal H_KL\). If the original weighted forced
incidence is at most \(w\) and \(D_Bw\le1\), its new
forced price satisfies
\begin{equation}
 \sum_\alpha e^{\kappa n_B(Y_\alpha)}
                         \|Q_\alpha\|_\infty
       \le D_Bw\,e(L),\qquad D_B=(2B+1)^4.
 \label{eq:response-forced-birth-price}
\end{equation}
Thus its contribution to \(\|R\|_{\rm ins}\) is at
most \(D_B\sqrt w\,e(L)\). Every constant here is
independent of the cutoff and physical side.
\end{lemma}
\begin{proof}
For an admissible original family \(A\), set
\(b=\{\alpha\in A:Y_\alpha\cap X\ne\varnothing\}\).
If \(L_X(U)\ne0\), \(X\) contains a bad site.
Exact inventory coverage supplies an owner in \(A\),
so \(b\ne\varnothing\). Replace precisely these
members by the recorded address \((X,b)\). Every
remaining footprint avoids \(X\) and the members of
\(b\), hence avoids their union \(Y\).
The new factor is exactly
\[
 L_X\mathbf1_{\{Y\cap\mathcal B_h=I\}}
                                  \prod_{\alpha\in b}K_\alpha.
\]
Conversely, opening this recorded address recovers
the unique old family. The other addresses avoid
\(X\), so the recovered \(b\) is exactly the family
of old members meeting \(X\). Summing this bijection
proves the density identity and
\eqref{eq:physical-response-profile}. It applies also
to threshold equality and wrapping footprints.

Put \(c_\alpha=e^{\kappa n_B(Y_\alpha)}
\|K_\alpha\|_\infty\). Dropping compatibility after
taking absolute values gives
\[
 \sum_{\varnothing\ne b:\,Y_\alpha\cap X\ne\varnothing}
                    \prod_{\alpha\in b}c_\alpha
 \le\prod_{\alpha:Y_\alpha\cap X\ne\varnothing}(1+c_\alpha)-1
 \le e^{|X|w}-1.
\]
The incidence bound was used in the last inequality.
All merged footprints form a connected overlap family,
so \eqref{eq:cover-connected} bounds their covering
weight by the product of the old weights times
\(e^{\kappa n_B(X)}\). Write \(t=n_B(X)\). Since
\(|X|\le D_Bt\), \(D_Bw\le1\), and \(t\le e^t\),
\begin{align*}
 e^{\kappa t}(e^{|X|w}-1)\|L_X\|_\infty
 &\le e^{\kappa t}|X|w e^{|X|w}\|L_X\|_\infty\\
 &\le D_Bw\,t e^{(\kappa+1)t}\|L_X\|_\infty
 \le D_Bw e^{(\kappa+2)t}\|L_X\|_\infty.
\end{align*}
Sum over \(X\). Finite truncations obey the same
bound; positive summation passes to the countable
families. This proves absolute convergence and
\eqref{eq:response-forced-birth-price}.
\end{proof}

For several named sources, expand
\(\prod_i(1+z_i f_i)\) in the coefficient algebra
of Section~\ref{sec:joint-insertion-bounds}, and then
expand each \(f_i=f_i^{\rm reg}+\sum_XL_{i,X}\).
In each term take the overlap components of the old
forced occurrences and the chosen response occurrences.
Every component containing a nonzero response contains
an old forced occurrence, by the same owner argument.
Its weight, name set and inventory are
\begin{equation}
 \left(\prod_{(i,X)\text{ in the component}}L_{i,X}\right)
 \left(\prod_{\alpha\text{ in the component}}K_\alpha\right),
 \qquad \{i\},\qquad \bigcup_\alpha B_\alpha.
 \label{eq:joint-response-component}
\end{equation}
The component also records its union support, exact
union flag test, and every old member. Opening all
recorded components reverses the operation. In
particular a shared old factor occurs once, including
when several names meet it. No division by an old
activity is involved. Components containing no
response remain as their original factors.

The same expansion can be made with a pre-existing
complete source: its forced occurrences are the full
marker polynomials of \eqref{eq:joint-local-source}.
Multiply those polynomials, and retain every disjoint
name partition in their coefficients. Thus inherited
higher modes and the new mixed modes are both kept.
The singleton profile added by a new physical source
is exactly \(\mathcal H_KL\), while the unmarked
coefficient is the original gas. A zero empty
inventory is never relabelled as forced.

\begin{lemma}[Joining with a separate response budget]
\label{lem:complete-source-response-joining}
Let \(\kappa\ge4\). Suppose an existing complete
source has higher regular and forced incidence
majorants at most one at radii \(\rho'_i\), and
every named anchor is different. At the joining
level shrink to \(0<\rho_i\le a\rho'_i\),
where \(0<a\le1\). Suppose its radii-weighted
forced singleton price, before the original
square-root division, is at most \(\sqrt w\)
per name. Let each new physical response have
radii-weighted reserved price at most \(b\ge0\).
Put
\begin{equation}
 r_B=e\{2(16B+9)\}^4,\qquad
 c=w+r_B\sqrt w+a+r_Bb.
 \label{eq:response-joining-income}
\end{equation}
Here the unmarked forced incidence is at most \(w\).
If
\begin{equation}
                  D_Bc\le3,\qquad e^c-1\le1,
 \label{eq:response-joining-smallness}
\end{equation}
the product with the physical sources has an exact
complete representation at exponent \(\kappa-4\).
Its higher regular and forced incidence majorants
are at most one. Its singletons are the sums of
the old singletons and the exact new representations.
The estimates are independent of the number of
names and of the torus side.
\end{lemma}
\begin{proof}
We give the incidence and component calculations.
For a polynomial letter \(A=\sum_Iz_IA_I\), use
\(\sum_I\rho_I\|A_I\|_\infty\), where
\(\rho_I=\prod_{i\in I}\rho_i\). This is the
absolute version of \eqref{eq:joint-marker-norm}.
Shrinking radii multiplies every old higher mode
by at most \(a^{|I|}\le a\). It therefore reduces
its higher incidence bound to \(a\). The old
regular higher modes can be left in the exponential.

To convert an anchored price into a site incidence
bound, spend one support unit. If \(o,y\in X\),
the connected cover gives
\[
 d_\infty(o,y)\le(4B+2)n_B(X)+2B,
 \quad
 e^{-n_B(X)}\le e\,
                  e^{-d_\infty(o,y)/(4B+2)}.
\]
For distinct anchors, summing the last kernel over
the torus is bounded by its sum over \(\Z^4\).
Indeed choose nearest representatives; they are
distinct. Since \(|v|_1\le4|v|_\infty\), this sum
is at most
\[
 e\left(1+\frac{2}{e^{1/(16B+8)}-1}\right)^4
                       \le e\{2(16B+9)\}^4=r_B.
\]
We used \(e^u-1\ge u\). Thus the old forced
singletons and new responses have site incidence
at most \(r_B\sqrt w\) and \(r_Bb\), respectively,
at exponent \(\kappa-1\). Together with the
unmarked and pruned higher forced letters their
complete incidence is at most the number \(c\)
in \eqref{eq:response-joining-income}.

Write these nonnegative letter majorants as \(d_\lambda\),
with their full supports \(X_\lambda\). Apply the distinct-label
part of Lemma~\ref{lem:connected-bound} with output exponent
\(\kappa-4\), covering reserve \(3\), and incidence \(c\).
The input exponent is exactly \((\kappa-4)+3=\kappa-1\),
and its hypothesis is \(D_Bc\le3\). The resulting output
incidence is \(e^c-1\le1\). This controls every name subset
at once. Removing singleton and unmarked modes
from the output only decreases its majorant.
The inherited regular higher part has bound
\(a\le1\), as already observed.

Apply the exact regrouping
\eqref{eq:joint-response-component} to these
letters. Moving the linear response terms from
the exponential into their forced components
leaves the prescribed regular singletons.
The added forced singleton profile is
\eqref{eq:physical-response-profile}; old
profiles and higher polynomials are unchanged
until they occur in the stated products.
All operations retain the original masks,
inventories and payload tests. The positive
majorants justify limits of finite address
truncations. Four support units have been
used once, at this joining: one for the
anchor-to-site sum and three in the component
induction. No derivative or Gaussian precision
is spent on these Borel response functions.
\end{proof}

\subsection{The same matched coefficient in a physical root estimate}
\label{subsec:physical-response-root}

We first check the geometry needed to identify a joined
coefficient with a physical replicated difference. Let a
source at level \(k\) read a ball of native radius \(R\),
including all endpoints and flag halos. At finer levels use
the radii
\[
 R_k=R,\qquad
 R_{j+1}=L_j\{R_j+10(J(j+1)+3)\}.
\]
The second summand bounds the frame, path, predictor and
fallback-test reads of the actual blocking patch. It
includes both endpoints, rather than just a main path.
Put
\(T_j=30\varepsilon_j(J(j)+5)\), a physical length.
The calendar has \(L_j\ge3\) and
\(J(j+1)\le J(j)+1\). Directly,
\begin{align*}
 T_{j+1}&\le\tfrac23T_j,\\
 \varepsilon_{j+1}R_{j+1}+T_{j+1}
 &=\varepsilon_jR_j+10\varepsilon_j(J(j+1)+3)+T_{j+1}
 \le\varepsilon_jR_j+T_j.
\end{align*}
For the second line use
\(10\varepsilon_j(J(j+1)+3)\le T_j/3\).
Consequently the single condition
\begin{equation}
              2(\varepsilon_kR+T_k)<\varepsilon_i
 \label{eq:response-native-room}
\end{equation}
places every finer read cone strictly inside its centred
level-\(i\) cell, for every terminal cutoff \(n\ge k\).
This is a bound on the entire history, uniform in \(n,M\).

To see the associated factorisation, first sum the bond
tags in the physical kernel. Its density can be written
\[
 \mathcal K_j(U,V)=\sum_{\boldsymbol t}
   \left\{\prod_x\lambda_{j,x}(U,t_x)\right\}
   \prod_{e\in E_{\ell_{j,M}}^+}
         p_{j,e}(U,t_{e^-},t_{e^+};V_e).
\]
Here \(t_x\) is the frame tag at \(x\), the nonnegative
masses \(\lambda_{j,x}\) sum to one, and each
\(p_{j,e}\) has Haar integral one. These statements
follow from the tag and noise normalisations in
\eqref{eq:physical-kernel}. Adjacent bonds can share a
site tag. Conditional on all site tags their link
integrations factor. Functions reading distinct cell
interiors use disjoint links and disjoint sets of tag
endpoints. Integrate unused links and sum unused tags,
each with mass one; Fubini then factors the remaining
integrals. Iterating, with the preceding cone bounds,
gives
\begin{equation}
 \mathcal P_{n\to k}^*\left(\prod_{a\in A}f_a\right)
                  =\prod_{a\in A}\mathcal P_{n\to k}^*f_a.
 \label{eq:response-pullback-products}
\end{equation}
The star here denotes the adjoint probability-kernel
pullback, not complex conjugation. The \(f_a\) are
the translated, and when required conjugated, cell
copies. Kernel symmetry commutes with those copies.
Bounded Borel functions are allowed; no differentiability
of their health masks enters Fubini's theorem.

\begin{lemma}[Physical matching with two budgets]
\label{lem:physical-response-matching-root}
Fix \(d\in\{3,4,5\}\) and use the setting of
Lemma~\ref{lem:ward-matching-coefficient-control}.
Suppose in addition that the dimensionless source
\(f\) is a bounded local physical function, its
one-insertion density is exactly \(F_kf\), and its
forced profile has a response family with price
\(e(L)\le x\), where \(x=\widehat\beta_k\ge1\).
Require the finite entry conditions
\begin{equation}
 4(1+C_{J(k)})q\le x^3,\qquad
 4(1+C_{J(k)})\le\sqrt{P_k},\qquad
 D_{B_k}c_k\le3,\quad e^{c_k}-1\le1,
 \label{eq:physical-root-entry}
\end{equation}
where
\[
 c_k=w_k+r_{B_k}\sqrt{w_k}+x^{-1}+r_{B_k}x^{-1}.
\]
Let \(a\) be the same canonical matched coefficient
as in that lemma. For
\(\operatorname{ref}(d)\le m<k\), set
\[
 G_{k,m}=(\varepsilon_m/\varepsilon_k)^{1/8},\qquad
 u_m=\frac{\zeta_{d,m}/2}
                 {\varepsilon_m^{-d}(x^3/G_{k,m})},
             \qquad \rho_m=x^{-1}u_m.
\]
Choose a centred grid of cell side \(\varepsilon_i\)
with \(i\le m\),
\(\varepsilon_i/\varepsilon_m\ge B_m^2\), and
\eqref{eq:response-native-room}. Then
\begin{equation}
 \left\|\varepsilon_k^{-d}
          \mathcal P_{n\to k}^*f-\mathcal O_{n,M}(a)\right\|_i
 \le\frac{\exp\{(\varepsilon_i/\varepsilon_m)^4
                                      Q_m^{\rm vac}+5\}}{\rho_m}
 =\frac{2x^4\varepsilon_m^{-d}}
             {\zeta_{d,m}G_{k,m}}
       e^{(\varepsilon_i/\varepsilon_m)^4Q_m^{\rm vac}+5}.
 \label{eq:physical-response-root-price}
\end{equation}
The sharp coefficient bias remains
\eqref{eq:ward-matching-sharp-bias}, with \(q\),
not \(x^3\). All inequalities are uniform in
\(n>k\) and on the physical-side tail used for
the finite matching interval.
\end{lemma}
\begin{proof}
\emph{One coefficient and two sizes.}
The degree-three case uses the same density and
response identities. Its primary space consists
only of the unmasked unit. No nonconstant primary
function or additional health mask is needed at
that degree. All estimates below retain their
displayed degree \(d\), in particular the physical
factor \(\varepsilon_j^{-d}\), the allowance
\(\zeta_{d,j}\) and the radius-entry power
\(L_j^{d+1}\). Thus \(d=3\) uses \(L_j^4\)
and the existing four-derivative low read; the
joining estimates do not depend on \(d\).

The response identity also fixes the physical
decomposition: applying
Lemma~\ref{lem:physical-response-birth} to the
same regular part gives density
\(F_k(f^{\rm reg}+\sum_XL_X)\).
It equals \(F_kf\) by the prescribed profile
and exact source identity. The actual blocked
Wilson density is strictly positive, so
\(f=f^{\rm reg}+\sum_XL_X\) pointwise.
The canonical coefficient is chosen by exact linear
transport to the reference level and subtraction
of the selected column. The proof of
\eqref{eq:ward-total-residual-budget} gives
\[
 |x_j^{\rm res}|_1+\|R_j^{\rm res}\|_{\rm ins}
       \le \frac{4(1+C_{J(k)})q}{G_{k,j}}
       \le \frac{x^3}{G_{k,j}},\qquad m\le j\le k.
\]
Moreover the sharp bias and the second entry condition
give \(|\mathsf Z^{-1}a-y|_1\le q\), so
\(|\mathsf Z^{-1}a|_1\le2q\). Its selected graph
and the original source at birth together have
size at most
\(q+2(1+C_{J(k)})q\le4(1+C_{J(k)})q\).
Thus \(x^3\) bounds the separate column birth
as well as the residual cone. This enlargement
is used only for source radii.

From \(u_m\), transport estimating weights to
finer levels by
\(u_{j+1}=\widehat\beta_{j+1}^{-1}u_j\), and
put \(\rho_j=x^{-1}u_j\). The common entry
condition
\[
 \widehat\beta_{j+1}^{-1}L_j^{d+1}\zeta_{d,j}
                                       \le\zeta_{d,j+1}
\]
and \(G_{k,j+1}=G_{k,j}L_j^{-1/8}\) imply,
by induction,
\begin{equation}
 u_j\varepsilon_j^{-d}(x^3/G_{k,j})
                              \le\zeta_{d,j}/2.
 \label{eq:response-radius-cone}
\end{equation}
Indeed the left side grows in one fine step by
\(\widehat\beta_{j+1}^{-1}L_j^{d+1/8}\),
which is bounded by the displayed larger factor.
This calculation applies also to the forward
selected-column source, using its leading inverse
bound and the radius entry of Section~43.
Its complete higher modes are transported with it.

At the joining level \(k\) we have two independent
smallness checks. The radii-weighted residual
singleton is at most \(\zeta_{d,k}/2\), and
the response has price at most
\begin{align}
 \rho_k\varepsilon_k^{-d}e(L)
 &\le u_k\varepsilon_k^{-d}e(L)
 \le \frac{e(L)\zeta_{d,k}}{2x^3}
 \le\frac{x^{-1}}2.\notag
 \text{Here }\quad
 e(L)\zeta_{d,k}&\le x\le x^2=x^{-1}x^3.
 \label{eq:response-root-budget-check}
\end{align}
We used \(0<\zeta_{d,k}\le1\). The unnormalised
old forced singleton price is at most \(\sqrt{w_k}\):
its divided price is part of the singleton bound,
and \(\zeta_{d,k}\le1\). Apply
Lemma~\ref{lem:complete-source-response-joining}
with radius reduction and response allowance
both equal to \(x^{-1}\). Its numerical income is
bounded by \(c_k\). Both requirements in
\eqref{eq:physical-root-entry} are eventually
true, since \(r_{B_k},D_{B_k}\) are logarithmic
powers whereas \(x^{-1}\to0\) and
\(w_k=e^{-P_k^{1/5}}\).

\emph{The complete density, before extraction.}
Let \(\mathcal T\) be the set of grid cells and
let \(O_C\) denote the copy of
\(\mathcal O_{n,M}(a)\) in cell \(C\), including
its reflection and conjugation convention.
Begin at the ultraviolet
level with
\(F_n\prod_{C\in\mathcal T}(1-z_CO_C)\).
Transport this entire polynomial to level \(k\).
Multiply there by
\(\prod_{C\in\mathcal T}(1+z_C\varepsilon_k^{-d}f_C)\).
In formulas the resulting density is
\begin{equation}
 \mathcal A_k(z)=
 \mathcal P_{n\to k}
       \left(F_n\prod_{C\in\mathcal T}(1-z_CO_C)\right)
       \prod_{C\in\mathcal T}(1+z_C\varepsilon_k^{-d}f_C).
 \label{eq:physical-response-joined-density}
\end{equation}
Its constant coefficient is \(F_k\). The exact
response regrouping gives its singletons as
the previously computed residuals, including
their complete forced profiles. Rebooking a
singleton profile at its fixed union-inventory
address cannot increase its weighted norm:
the supremum norm of a sum is at most the
sum of the original supremum norms, and all
terms at that address have the same covering
weight. Thus the cone price applies to this
representation as well. Its higher
coefficients include the inherited higher modes,
all old--new products, and all new mixed
responses in \eqref{eq:joint-response-component}.
None is prescribed by multiplying singleton
replacements.

After the join, descend with the original kernels
to \(m\). The attachment proof of
Lemma~\ref{lem:joint-source-attachments} applies:
singleton weights obey
\eqref{eq:response-radius-cone}, and enlarging
radii on a higher mode is paid by its separated
names and the connected-cover estimate. On the
centred grid the separation condition holds at
every finer level, since
\(B_{j+1}^2\le2B_j^2\) and \(L_j\ge3\).
The four exponent units used in the join are
included in the finite support reserve; the
first subsequent cover compression returns to
the prescribed output exponent. This requires
only the same compression inequality with input
exponent \(\kappa-4\), imposed before choosing
the block ratio. It is not a four-unit loss at
every subsequent step. If one stops at the
join itself the nonnegative weaker exponent
already suffices for extraction.

For clarity we also give the final extraction
bound. At \(m\), the radii-weighted singleton
exponential and forced sums are at most two and
one per name. This follows from the cone,
the section allowances \(\zeta_{d,m}\), and
the divided forced norm. Each higher letter
has a support containing at least one named
anchor. A higher incidence bound of one
therefore gives total higher price at most
\(|\mathcal T|\): sum its incidence at all named
anchors, which counts every nonzero letter
at least once. Apply this separately to
the regular and forced higher families.
Consequently, by the submultiplicative marker
norm, their exponential and gas products cost
at most
\[
        \exp\{(2+1+1+1)|\mathcal T|\}=
                                      e^{5|\mathcal T|}.
\]
The unmarked forced products are bounded by
\(e^{\#\Lambda_m w_m}\). Integrating the
resulting absolute numerator and using the
upper and lower bounds in the proof of
Lemma~\ref{lem:insertion-vacuum-ratio} gives
\begin{equation}
 \frac{\left|\int[z_{\mathcal T}]\mathcal A_m\,\dd H_m\right|}
      {Z_{n,M}}\,\rho_m^{|\mathcal T|}
 \le \exp\{\#\Lambda_mQ_m^{\rm vac}+5|\mathcal T|\}.
 \label{eq:response-root-coefficient-bound}
\end{equation}
The harmless fixed enlargement of the vacuum
profile constant in this upper bound is included
in its initial choice. The denominator remains
the original Wilson partition function.

Finally expand \eqref{eq:physical-response-joined-density}
over a subset \(A\subset\mathcal T\) of the new factors.
Adjointness moves its integral to the fine
level; \eqref{eq:response-pullback-products}
then turns that term into
\[
 \left\langle
   \prod_{C\in A}\varepsilon_k^{-d}\mathcal P_{n\to k}^*f_C
   \prod_{C\in\mathcal T\setminus A}(-O_C)
                                                    \right\rangle_{n,M}.
\]
Summing the subsets gives the moment of the
complete product of differences. This uses
no factorisation of the old fine observables.
Mass preservation identifies this moment with
the left side of
\eqref{eq:response-root-coefficient-bound}.
Since
\(\#\Lambda_m/|\mathcal T|=(\varepsilon_i/\varepsilon_m)^4\),
taking its \(|\mathcal T|\)-th root proves
\eqref{eq:physical-response-root-price}.
The same \(a\) has been used throughout.
All local majorants precede cutoff and volume;
the finite matching comparisons require only
their stated physical-side tail.
\end{proof}

\begin{corollary}[A general physical root envelope]
\label{cor:physical-response-general-envelope}
In the preceding argument the positive root envelope
\(x^3\) can be replaced by any \(E_0>0\) satisfying
\[
 4(1+C_{J(k)})q\le E_0,
                \qquad e(L)\zeta_{d,k}\le x^{-1}E_0.
\]
The separate restriction \(e(L)\le x\) is then
unnecessary. All coefficient and joining entry
conditions other than the replaced envelope
inequality are retained.
Keep the radius reduction \(x^{-1}\) and the same
joining entry. The right side of
\eqref{eq:physical-response-root-price} becomes
\[
 \frac{2xE_0\varepsilon_m^{-d}}{\zeta_{d,m}G_{k,m}}
                e^{(\varepsilon_i/\varepsilon_m)^4Q_m^{\rm vac}+5}.
\]
One may also take \(k=n\); the initial transport
in \eqref{eq:physical-response-joined-density}
is then the identity. The canonical matched
coefficient is unchanged.
\end{corollary}
\begin{proof}
Use \(E_0\) in
\eqref{eq:response-radius-cone}. The response
check becomes
\(u_k\varepsilon_k^{-d}e(L)
\le e(L)\zeta_{d,k}/(2E_0)\le x^{-1}/2\).
All component and extraction bounds are the
same. The zero-envelope case, when needed,
is the zero source and has zero residual.
\end{proof}

\subsection{The finite degree-five coefficient}

For an analytic local stencil \(F\), let \(\mathcal P_dF\)
be its radial polynomial read through engineering weight
\(d\): on a radial jet it is
\(\sum_{v=0}^d[\varepsilon^v]F(U(\operatorname{sc}_\varepsilon A))\).
This uses the holonomies and basis of
Section~\ref{sec:local-observable-sections}.
The map \(\mathsf D_h\) scales factor count, as defined
there; it does not scale engineering weight.
For the vector read at level \(k\) use
\[
              \mathcal P_{5,k}=h_k^{-2}\mathsf D_{h_k}\mathcal P_5
\]
and then real \(H_4\) averaging. The raw invariant vectors
in degree five all have two curvature factors, so this
normalisation reads their coefficients in the three-coordinate
basis of the previous section.

% !TeX root = 4DYM_notes.tex
We first derive the equation insertion from the full
physical density. We then perform its local reading.
This order keeps the hard tests out of the differentiation.
The organisation follows
\cite[Lemma~18.5, PDF pp.~210--213]{ArmstrongDraft};
the finite-dimensional identities and the norm conversions
needed here are given explicitly.

\emph{A smooth physical density before any sector expansion.}
Fix a finite terminal level \(n\), a physical side length \(M\)
and a level \(k<n\). Let \(F_{n,k}\) be the original
Wilson density after the exact blocking down to \(k\),
with all extracted scalar factors included, and let
\(Z_{n,k}=\int F_{n,k}\,\dd H_k\). Here \(H_k\) is
product probability Haar measure at this level.
The density is strictly positive and smooth. Indeed,
at the last step it is
\[
 F_{n,k}(V)=\int K_{\beta_{k+1}}(U,V)
                              F_{n,k+1}(U)\,\dd H_{k+1}(U).
\]
For held physical \(U\), the kernel is smooth in \(V\).
All its derivatives are bounded at this fixed finite
cutoff and volume, even though its branch weights are
Borel functions of \(U\). Dominated differentiation
on the compact fine configuration space proves the
claim. Positivity follows from positivity of the
noise and \(\sum_t\lambda_t(U)=1\). Kernel
normalisation gives \(Z_{n,k}=Z_{n,k+1}\).
This argument proves differentiability only; the
uniform estimates below do not use its finite-volume
compactness constants.

We use a smooth finite-range representative of the
current when taking a physical derivative. To construct
it, complete its jets through order four by the
transported curvature atoms used in
Lemma~\ref{lem:trace-current-local-bounds}, and truncate
their kernels and paths at radius
\[
 B_k=\lceil(\log\widehat\beta_k)^2\rceil,\qquad
                 R_{\rm cut}=C_j^{\rm cut}(1+B_k)^{61}.
\]
We use this longer spatial truncation also in the Ward
transport below. Its purpose is to make the discarded
kernel and connected-support tails smaller than
\(\exp\{-(1+B_k)^{40}\}\), even after the change from
level \(k\) to a radius level. A shorter truncation would
suffice for the coefficient read alone.
Each retained atom is a polynomial in group matrices,
their inverses and the fixed Lie-algebra projection.
It is therefore globally smooth and gauge covariant,
not merely a function in a small chart. The finite
jet completion removes the lower jets before the next
one is added, as in the current proof. This gives a
Lie-algebra-valued displacement supported on the
finite set of links read by those atoms. The difference
of its jets through order four from the original
current is at most \(C_k e^{-a_jR_{\rm cut}/2}\); the spare
width pays the path and polynomial factors.

The same construction applies to the scalar current
jets. It is done jointly so that the exact trace and
antisymmetric identities are preserved. To see the
finite algebra, let \(\widetilde J,\widetilde q,
\widetilde t\) be the independently truncated polynomial
representatives. Choose the first coordinate axis and
define, for \(a<b\),
\begin{align*}
 J_{ab}&=\widetilde J_{ab},&
 J_{ba}&=\widetilde J_{ab}-\sum_c\delta_c\widetilde t_{abc},\\
 J_{aa}&=\widetilde J_{aa}\quad(a\ne1),&
 J_{11}&=\sum_a\delta_a\widetilde q_a
                               -\sum_{a\ne1}\widetilde J_{aa}.
\end{align*}
Use \(q_a=\widetilde q_a\),
\(t_{abc}=\widetilde t_{abc}\) for \(a<b\), and
\(t_{bac}=-t_{abc}\), \(t_{aac}=0\).
Summing the diagonal entries, or subtracting the two
off-diagonal entries, proves
\eqref{eq:trace-current-identities} for these finite
functions. Since the original currents obey those
identities, their changes of jets are finite sums of
the individual truncation errors. They obey the same
exponential error bound with an enlarged constant.
The classical defect is defined from these actual
representatives and the chosen \(Y\), retaining its
entire higher-order function. Subsequent real \(H_4\)
averaging is applied to the complete identities.
A reflected forward difference remains a forward
difference of a shifted potential, since
\(T_a^{-1}-1=-(T_a-1)T_a^{-1}\). Thus averaging does
not destroy their exact form.

Every further
fixed derivative of these finite polynomials has
bound \(C_k(1+R_{\rm cut})^q\), from their coefficients and
path lengths. Thus the Taylor estimates below use
these representatives and the native principal
action; flat-jet bounds for the transferred currents
suffice. The constants \(C_k\) have at most fixed powers
of \(\log\widehat\beta_k\), times fixed-profile
constants, by the explicit current inequalities.
The choice of \(C_j^{\rm cut}\) is made after the profile widths
and before the limits. Include it in the finite
current constant \(A_{j,L}\) used in
Section~\ref{sec:nonlinear-induction-step}. The existing
entry condition then gives
\(C_j^{\rm cut}\le\log\widehat\beta_k\), and hence
\(R_{\rm cut}\le C(\log\widehat\beta_k)^{125}\) after entry.
In particular \(\varepsilon_kR_{\rm cut}\to0\) along
the nominal calendar. These are choices before the
Wilson family and its limits, as in the other
current estimates.
In the next identity write \(Y_b\) for this smooth
representative, and keep its approximation error
among \(\varepsilon_{\rm stat}\).

For an orthonormal colour basis, let \(\mathcal L_{e,\alpha}\)
be the left derivative on link \(e\) in colour
\(\alpha\). Haar invariance gives
\(\int\mathcal L_{e,\alpha} f\,\dd H_k=0\) for
every smooth \(f\). Hence
\[
 \operatorname{div}_{\rm Haar}Y_b
                  =\sum_{e,\alpha}\mathcal L_{e,\alpha}(Y_b)_{e,\alpha}.
\]
Define the equation insertion
\begin{equation}
 Q_b^{\rm eq}=F_{n,k}^{-1}
                      \operatorname{div}_{\rm Haar}(F_{n,k}Y_b)
       =\operatorname{div}_{\rm Haar}Y_b+D\log F_{n,k}[Y_b].
 \label{eq:trace-physical-equation-insertion}
\end{equation}
Summing the identities
\(\int\mathcal L_{e,\alpha}(A F_{n,k}(Y_b)_{e,\alpha})=0\)
proves, for every smooth cylinder \(A\),
\begin{equation}
 \frac1{Z_{n,k}}\int A Q_b^{\rm eq}F_{n,k}\,\dd H_k
       =-\frac1{Z_{n,k}}\int DA[Y_b]F_{n,k}\,\dd H_k.
 \label{eq:trace-one-step-weak-ward}
\end{equation}
In particular its exterior moment is zero whenever
the spectator reads avoid the links on which \(Y_b\)
acts. No Maxwell equation has been imposed on a
coefficient in this identity.

The identity can be written before the last fine
integration as
\begin{align*}
 &\int A(V)Q_b^{\rm eq}(V)F_{n,k}(V)\,\dd H_k(V)\\
 &=\int\!\!\int A(V)\{
     (\operatorname{div}_{\rm Haar}Y_b)(V)K(U,V)
                      +D_VK(U,V)[Y_b(V)]\}
               F_{n,k+1}(U)\,\dd H_{k+1}(U)\,\dd H_k(V).
\end{align*}
Here \(K\) is the actual last blocking kernel, including
its branch sum and noise normaliser. Thus physical
fine multipliers are held fixed when the derivative
is taken. Equivalently, for one commuting marker
\(\eta\) with \(\eta^2=0\), the exact one-step
source identity is
\begin{equation}
 F_{n,k}(1+\eta Q_b^{\rm eq})
  =\int\{K+\eta[(\operatorname{div}_{\rm Haar}Y_b)K
                                  +D_VK[Y_b]]\}F_{n,k+1}\,\dd H_{k+1}.
 \label{eq:trace-marked-ward-step}
\end{equation}
The coefficient of \(\eta\) has integral zero,
by \eqref{eq:trace-one-step-weak-ward} with \(A=1\).
The zero mode and its partition normalisation are
unchanged. Sector expansion and physical reassembly
can now be applied to this differentiated smooth
integrand, by linearity. Their masks and selection
queries are retained. This is precisely the order
of operations in
Lemma~\ref{lem:regular-Borel-measure-derivative}.
It never differentiates an outgoing Borel gas one
factor at a time.
There is also a uniform bound for the source away
from its small chart. Boundedness of the first
derivative of the chord square gives
\[
 |D_VK(U,V)[Y_b(V)]|
       \le C_Ga\beta_{k+1}\Bigl(\sum_e|(Y_b)_e(V)|\Bigr)K(U,V).
\]
It holds before the tag sum and survives it because
the weights are nonnegative. Dividing the integrated
inequality by \(F_{n,k}>0\) bounds its logarithmic
derivative by the same local quantity. A fixed
number of further derivatives gives only fixed
powers of the coupling and of the current read
radius. This estimate has no number of sites
outside that radius in it.

\emph{The local germ and its raw coefficient.}
On an all-good neighbourhood of the flat field,
Lemma~\ref{lem:all-good-gas} gives \(Z_K=1\).
All regular masks there are one, so the exact density
representation reads
\[
 \log F_{n,k}=c_{n,k}-\beta_k\mathsf S
                         -\mathcal L_{n,k}-\sum_i E_i.
\]
The constant \(c_{n,k}\) contains the full vacuum
normalisation. It has zero derivative. The sums of
regular germs and their admitted derivatives converge
by their incidence norm. Inserting this expression
in \eqref{eq:trace-physical-equation-insertion} and
using \eqref{eq:trace-current-identities} gives, through
the jets needed for the raw weight-five read,
\begin{align}
 Q_b^{\rm eq}+\beta_k\sum_a\nabla_aJ_{ab}
   ={}&\operatorname{div}_{\rm Haar}Y_b
                  -D\mathcal L_{n,k}[Y_b]
             +\beta_k\mathcal E_b-\sum_iDE_i[Y_b]
                         +\text{retained jet errors}.
 \label{eq:trace-central-ward-germ}
\end{align}
The error from the stated fourth-order classical
identity has zero jets through field order four.
The raw invariant vector space through weight five
has the three quadratic-curvature representatives
of Section~\ref{sec:trace-ward-coefficient}. Its
coefficients are determined by the quadratic field
jet. Thus the higher-field part contributes no new
raw weight-five coefficient. Its physical remainder
is kept in the marked family, not set to zero in an
integral.

Read \eqref{eq:trace-central-ward-germ} with
\(\mathcal P_{5,k}\), keeping all three raw coordinates.
The left side is the compensated Ward source whose local
coefficient is \(\Sigma_k^{\rm read}\); its equation insertion
is retained for the exterior identity used below.
It is not quotiented by \(v_{\rm eq}\) at this stage.
There are four types of discrepancy when this germ
is represented in the original physical source:
stationary-current substitution, kernel truncation,
the classical and regular-action error, and marked
localisation/restoration. Allocating each source
term once gives
\begin{equation}
 \Sigma_k^{\rm read}=
 \mathcal P_{5,k}\{\operatorname{div}_{\rm Haar}Y
                                   -D\mathcal L_{n,k}[Y]\}
       +\varepsilon_{\rm stat}+\varepsilon_{\rm tail}
                    +\varepsilon_{\rm cl}+\varepsilon_{\rm mark}.
 \label{eq:trace-five-coefficient}
\end{equation}
This is the local reading of the exact weak identity
\eqref{eq:trace-one-step-weak-ward}. The physical
remainder, with every source subset mode, accompanies
that reading. In particular no differentiated
partition function is introduced: \(Z_{n,k}\) is
the same number on both sides of the weak identity.

\emph{Quantitative prices after the projection.}
The two-curvature coordinates have factor count two.
The dimensionless read therefore supplies \(h_k^2\),
cancelled by the explicit \(h_k^{-2}\) in
\(\mathcal P_{5,k}\). A regular error with scaled
derivative price \(e_X\) has coefficient price at
most \(C_j h_k^{-2}e_X\), times its fixed support
moment. This is the finite reading inequality in
Section~\ref{sec:insertion-transport}; the exponential
and covering reserves sum those moments.

For example \(DE_i[Y_b]\) has scaled derivative
price \(C_k(1+R_{\rm cut})^q\|E_i\|\), since the current
and its admitted jets have size \(C_kh_k\) in the
physical directions of that norm.
For instance, in ordinary chart operator norms the
order-\(r\) product rule gives, up to the fixed conversion
between chart and left derivatives,
\[
 h_k^r\|D^r(DE_i[Y])\|
 \le C_r\sum_{\ell=0}^r\binom r\ell
     \bigl(h_k^{\ell+1}\|D^{\ell+1}E_i\|\bigr)
     \bigl(h_k^{r-\ell-1}\|D^{r-\ell}Y\|\bigr).
\]
The second factor is bounded by the current constant;
the first uses the unmarked regular estimate through
order \(r+1\). Thus no extra inverse threshold remains.
Only occurrences
whose complete read meets the current can contribute.
Rebase their norms at its anchor once, and spend a
covering reserve. Their summed price is bounded by
\(C_{j} (1+R_{\rm cut})^q(1+B_k)^a\delta_k\), with fixed
additional logarithmic current factors. Consequently
\[
 h_k^{-2}\delta_k
       =\widehat\beta_k^{-1/10}P_k^{-2},\qquad
                       \delta_k=\widehat\beta_k^{-11/10}.
\]
There is no sum over all physical sites in this
inequality.

We also make the Taylor gain precise. A classical
difference with zero ordinary jets through order
four is \(O(C_Bh_k^5)\) on its small local chart.
After its coupling and projection factors its price
is
\[
 h_k^{-2}\beta_k C_Bh_k^5
       \le C_B\widehat\beta_k^{-1/2}P_k^3.
\]
A fourth-jet representative of the displacement
has a fifth-order difference; its divergence costs
one derivative and hence has price \(C_Bh_k^4\).
Its projected price is \(C_Bh_k^2\). Both are
smaller than the preceding regular-action allowance
at entry. Here the gain comes from ordinary field
jets. Engineering weight and factor count remain
distinct: increasing engineering weight alone is
not used to assert an extra factor of \(h_k\).

The classical defect has price \(C_k2^{-s}\) in
the current seminorms after \(s\) ultraviolet
transfers, by
\eqref{eq:trace-defect-full-seminorm-decay}. Its
contribution is bounded by \(C_k\beta_k2^{-s}\)
at fixed \(k\), and vanishes on the original
ultraviolet tail, uniformly in the physical period.
This is a limit at fixed \(k\), not an estimate
obtained by increasing the period first.

For the forced part let \(w_k=e^{-P_k^{1/5}}\)
be the original forced allowance. In a nonzero
marked forced component retain a distinguished
eligible forced occurrence. Split its price as
\(k_\alpha=(k_\alpha/\sqrt{w_k})\sqrt{w_k}\).
The incidence sum of the divided prices is at most
\(\sqrt{w_k}\). The rooted component estimate
sums the other occurrences with the same flags
and leaves one factor \(\sqrt{w_k}\) at the marked
root. The finitely many density derivatives are
estimated by the common score polynomials of
Section~\ref{sec:integrated-regular-derivatives}, using
the one precision reserve. Their covariance moments
and the finite projection cost at most
\(C_j\widehat\beta_k^{C_j}P_k^{C_j}\). No derivative
of a forced activity is taken.

For nonwrapping regular source terms perform the
exact radial reading of
Section~\ref{sec:insertion-transport} and retain
its entire signed residual. Its radial jets through
weight five are zero, so its contribution to the
present raw coefficient is exactly zero. A local
mask restoration with a nonempty bad inventory is
also zero on the all-good germ. Its physical
integral need not be zero: the signed residual
and all forced terms continue in the full source,
with their original norms. The physical Ward
identification below uses backward matching and
its residual bound; it does not infer a small
physical remainder from its zero local read.
The nonzero small regular-action contribution
here is the \(\delta_k\) term already calculated.
Together with the forced restoration bound, this gives
\begin{equation}
 C_j\widehat\beta_k^{-1/10}P_k^{C_j}
        +C_j\widehat\beta_k^{C_j}e^{-P_k^{1/5}/2}
                                             =o(P_k^{-1}).
 \label{eq:trace-projection-errors}
\end{equation}
Indeed the logarithm of the first term divided by
\(P_k^{-1}\) is
\(-\tfrac1{10}\log\widehat\beta_k
 +(C_j+1)(\log\log\widehat\beta_k)^2+O_j(1)\to-\infty\).
For the second use
\(P_k^{1/5}/\log\widehat\beta_k\to\infty\).
All constants depending on the fixed profile are
paid by entry after its ratio has been fixed.
In particular \eqref{eq:trace-projection-errors}
is a bound for the coefficient error. It does not
assert that an arbitrary weight-five marked
remainder has regular norm \(\delta_k\).

Stationary and loop tails have price
\(C_k\widehat\beta_k^{C_j}e^{-a_jR_{\rm cut}/2}\), after
the fixed projection and path moments. Choose
\(C_j^{\rm cut}\) using the pre-existing widths and then
entry. These tails are smaller than \(P_k^{-1}\).
For clarity the four error types have the following
bounds, with their finite sums included in \(C_j\):
\[
\begin{array}{c|l}
 \text{error}&\text{projected bound}\\ \hline
 \varepsilon_{\rm stat}&
 C_B(h_k^2+\widehat\beta_k^{-1/2}P_k^3)
                          +C_k\widehat\beta_k^{C_j}e^{-a_jR_{\rm cut}/2}\\
 \varepsilon_{\rm tail}&C_k\widehat\beta_k^{C_j}e^{-a_jR_{\rm cut}/2}\\
 \varepsilon_{\rm cl}&C_B\widehat\beta_k^{-1/10}P_k^{-2}
                                                  +C_k\beta_k2^{-s}\\
 \varepsilon_{\rm mark}&C_j\widehat\beta_k^{-1/10}P_k^{C_j}
                          +C_j\widehat\beta_k^{C_j}e^{-P_k^{1/5}/2}.
\end{array}
\]
Every bill is local at the marked root. The notation
\(C_B\) in this table includes fixed powers of
\(1+B_k\), \(1+R_{\rm cut}\) and the current logarithmic
bounds; these are bounded by powers of \(P_k\).

A wrapping marked support is treated before making
a radial read. It reaches distance at least
\(\ell_{k,M}/10\), and its covering number is at least
\(\ell_{k,M}/[10(2B_k+1)]-1\). The elementary inequality
\(e^{98c}\le e^{-2q}e^{100c}\) for \(c\ge q\)
therefore gives its full-function price
\begin{equation}
 C_j\widehat\beta_k^{C_j}P_k^{C_j}
       \exp\{-\ell_{k,M}/[5(2B_k+1)]+2\}.
 \label{eq:trace-wrapping-error}
\end{equation}
This term vanishes on the period tails at fixed
\(k\). The ultraviolet limit at each fixed period
and then the prescribed period refinement are the
original ordered limits. All cutoff-dependent
current arrays have the uniform bounds stated
above, so this order is preserved.

\emph{The complete sources to be matched.}
Write \(g_b=\sum_a\nabla_aJ_{ab}\), with the current
representative just fixed. We construct a local physical
representative \(f_b\) of the compensated source
\(Q_b^{\rm eq}+\beta_kg_b\). In the marked kernel identity
\eqref{eq:trace-marked-ward-step}, expand the single mark
with the same connected-component expansion as the
unmarked density. Add \(\beta_kg_b\) before subtracting
any radial jets. Truncate the anchored smooth kernels
and connected reads at \(R_{\rm cut}\); retain their tails
as signed comparison terms. On each retained good chart
keep the entire functions
\[
 \operatorname{div}_{\rm Haar}Y_b-D\mathcal L_{n,k}[Y_b]
 +\beta_k\{g_b-D\mathsf S[Y_b]\}-\sum_iDE_i[Y_b],
\]
rather than their fourth-order Taylor polynomials.
Write \(G_{b,X}\) for these retained anchored connected
germs, with their natural supports \(X\). At this finite
cutoff the spatially truncated component sums converge
absolutely in the admitted regular norms. Put
\[
 y_f=\sum_X\Pi_{5,h_k}G_{b,X},\qquad
 R_{b,X}=G_{b,X}-P_5[\mathsf D_{h_k^{-1}}\Pi_{5,h_k}G_{b,X}],
 \qquad \widetilde X=X\cup\Sigma_5,
\]
where \(\Sigma_5\) is the fixed primary section footprint.
Define the physical function, with each mask in its
original convention, by
\begin{equation}
 f_b=\chi_{\Sigma_5}P_5[\mathsf D_{h_k^{-1}}y_f]
                         +\sum_X\chi_{\widetilde X}R_{b,X}.
 \label{eq:trace-complete-local-source}
\end{equation}
The real vector projection kills the unit. On the common
good chart this function equals the sum of the entire
germs \(G_{b,X}\); the radial subtractions cancel exactly.
Outside that chart it is the displayed bounded Borel
function, and is used as such. No derivative of its
masks is taken. Its birth representation has zero
forced replacement: multiplying the entire unmarked
density by this masked regular function is already exact.
The nonzero forced responses subsequently generated by
RG transport are retained in its full matched column.
For the unmasked physical current \(g\), the usual birth
construction also retains its forced restorations.

Retain each component's natural support and covering
weight. The enclosing box of radius \(R_{\rm cut}\) is
not assigned as the support of a single regular atom:
its covering price would destroy
\eqref{eq:trace-joint-source-price}. The finite current,
the primary-section bound and the summable component
majorants give the global cap
\(\widehat\beta_k^A\), after entry, for this finite
request. The underlying smooth and Borel variables are
treated in the order of
Lemma~\ref{lem:regular-Borel-measure-derivative}.

In particular, \(f_b\) is a bounded local physical source.
We do not assert pointwise equality to the generally
nonlocal score \(Q_b^{\rm eq}+\beta_kg_b\) on the entire
configuration space. Their comparison is made at the
level of the marked physical numerator, retaining the
spatial tails and nonempty inventories just described.
This comparison, including its normalisation, is estimated
in Lemma~\ref{lem:complete-source-ward-transport} below.
Its radial read is \(\Sigma_k^{\rm read}\) from
\eqref{eq:trace-five-coefficient}, with these restoration
and truncation errors included. Both \(f\) and \(g\) are
projected onto real \(H_4\) vector covariance as complete
physical sources, before their prices are taken.

We need the price of their entire birth representations,
not only their reads. Denote the read and remainder of
\(f\) by \(y_f,R_f\). When the source norm uses regular
derivatives through order \(q\), retain the unmarked
regular-action estimate through order \(q+1\) for this
birth calculation. This is one finite request in the
initial derivative catalogue. It pays \(DE_i[Y]\) and
is not an extra derivative lost at each RG step: later
transport uses only \(\max(q_c,6)\) derivatives of the
marked input, as in
Lemma~\ref{lem:insertion-local-prices}. The Haar-divergence term has no
constant vector after this projection. Its linear scalar
gauge-invariant jet is zero by simplicity. The entire
retained analytic germ consequently starts in field
order two. Taylor's formula, with its convergent analytic
majorants, gives the scaled regular price
\(C_j(1+R_{\rm cut}+B_k)^{a_j}h_k^2\).
For \(D\mathcal L[Y]\), a derivative removes one field
factor and \(Y\) restores one. Its price is at most
\(C_j(1+R_{\rm cut}+B_k)^{a_j}h_k^2\Lambda_k\).
The higher classical terms are kept exactly and cost
\(C_j(1+R_{\rm cut}+B_k)^{a_j}\beta_kh_k^5\).
The incoming regular actions cost the already computed
\(C_j(1+R_{\rm cut}+B_k)^{a_j}\delta_k\).
Analytic Taylor's integral remainder also controls every
admitted scaled derivative of these higher terms; it
does not replace the retained function by its read.

The local norm rebasing is at the marked anchor.
The connected tree estimate sums all other unmarked
occurrences, while one eligible forced occurrence supplies
\(\sqrt{w_k}\) after its divided price has been used.
The polynomial score moments and the original Haar/Gaussian
ratios are those already bounded above. Thus
\begin{align}
 |y_f|_1+\|R_f\|_{\rm ins}
 \le{}& C_j(1+R_{\rm cut}+B_k)^{a_j}
       \{h_k^2(1+\Lambda_k)+\beta_kh_k^5+\delta_k\}
          \notag\\
 &+C_j\widehat\beta_k^{a_j}e^{-P_k^{1/5}/2}
   +C_k\beta_k2^{-s}
   +C_j\widehat\beta_k^{a_j}e^{-a_j'R_{\rm cut}}.
 \label{eq:trace-complete-source-birth-price}
\end{align}
Here \(a_j'>0\) is a fixed profile width, unrelated to
the nonnegative exponent \(a_j\). The read and remainder
of \(g\) satisfy the same bound with the loop and
regular-action terms omitted and with \(h_k^2\) as their
principal price. There is no independent price for a
second vacuum or a second Gaussian measure.

Divide \eqref{eq:trace-complete-source-birth-price} by
\(h_k^2\). The last two algebraic small terms become
\(\widehat\beta_k^{-1/2}P_k^3\) and
\(\widehat\beta_k^{-1/10}P_k^{-2}\). The divided forced
and spatial tails tend to zero faster than any fixed
negative power of \(P_k\). Every power of
\(R_{\rm cut},B_k\) is a fixed power of
\(\log\widehat\beta_k\) at a fixed profile. For
\(\theta_j<1/6\) the positive margin in
\(1/6-\theta_j\) absorbs these powers at entry.
First choose the ultraviolet tail at fixed \(k\) to
pay \(C_k\beta_k2^{-s}\). We obtain, uniformly in
the remaining cutoff and physical side,
\begin{equation}
 |y_f|_1+\|R_f\|_{\rm ins}
       +|y_g|_1+\|R_g\|_{\rm ins}
                                 \le C h_k^2P_k^{1/6}.
 \label{eq:trace-joint-source-price}
\end{equation}
All finite constants in this request are included before
the corresponding profile entry. In particular the longer
Ward truncation does not change its power of \(P_k\).

% !TeX root = 4DYM_notes.tex
\subsection{The actual response families of the two Ward sources}
\label{subsec:trace-physical-response-birth}

We verify the additional input of
Lemma~\ref{lem:physical-response-matching-root} for
the sources just constructed. The estimate concerns
their global amplitudes and their natural component
supports. A healthy derivative bound alone would
not give it.

\begin{lemma}[Response prices of the compensated source and current]
\label{lem:trace-physical-response-data}
For the actual finite representatives \(f_b,g_b\)
above, use the exact same density and address book
as in \eqref{eq:trace-joint-source-price}. At birth
\(f_b\) has response zero. The current \(g_b\) has
an explicit response family \(L^g\) with
\begin{equation}
 e(L^g)\le C_j(1+\log\widehat\beta_k)^{A_j}
                                  \le\widehat\beta_k
 \label{eq:trace-actual-response-budget}
\end{equation}
after the common profile entry. Here \(C_j,A_j\)
depend on the admitted finite current jets, the
chosen widths, section and fixed block ratio, but
are independent of \(n>k\) and \(M\). The inequality
holds on the ultraviolet tail on which those
finite currents are used. The same construction
and bounds apply to the finite tensor components
\(J_{ab}\) and their fixed shifts and central
differences, using the primary footprint of
degree four or five as appropriate. Both sources therefore
have the same sharp coefficient budget
\(q\le Ch_k^2P_k^{1/6}\) and the separate physical
root budget of Lemma~\ref{lem:physical-response-matching-root}.
\end{lemma}
\begin{proof}
\emph{The masked compensated source.}
Formula~\eqref{eq:trace-complete-local-source}
is already a sum of a masked primary function
and masked zero-read regular functions. Multiplying
the whole unmarked density by that displayed
function is exactly its one-insertion density,
with every forced replacement zero. Set
\(L^f_X=0\) for all \(X\). This proves both the
response identity and \(e(L^f)=0\). It says
nothing about the forced remainders produced
after integration, which are retained by the
complete transport.

\emph{Global component bounds for the current.}
Use the finite polynomial representative of
\(g_b=\sum_a\nabla_aJ_{ab}\), with its exact trace
and antisymmetric completion, rather than a
nonlinear small-chart continuation of the
untruncated current. Split it into its original
kernel and path terms:
\(g_b=\sum_\lambda G_\lambda\).
Each term is a colour contraction of at most
four transported curvature inputs, or one of
the finitely many shifted terms making the
exact current identities. Its support
\(X_\lambda\) contains the anchor and all the
paths actually read. Complete the real
\(H_4\) orbits as a finite sum of such terms;
do not replace their separate supports by
one enclosing box.

Let \(\ell_\lambda\) denote the total length
of a connected system of those paths, enlarged
by the fixed section and shift lengths. The
curvature division and finite-jet estimates in
the proof of Lemma~\ref{lem:trace-current-local-bounds}
give a spare width \(a_j>0\) and, for every
fixed derivative order \(q_0\) used here,
\begin{equation}
 \sum_\lambda |c_\lambda|
        (1+\ell_\lambda)^{q_0}e^{a_j\ell_\lambda}
                   \le C_j(1+\log\widehat\beta_k)^{A_j}.
 \label{eq:trace-global-current-path-price}
\end{equation}
Here \(c_\lambda\) includes its kernel coefficient
and fixed colour contraction constants. To check
the global use of this estimate, a group matrix
or its inverse has operator norm one. A curvature
atom, a fixed Lie-algebra projection and its
finitely many differentiated products consequently
have norm at most a fixed colour constant times
a power of path length. The product rule gives
at most \(C_{j,q_0}(1+\ell_\lambda)^{q_0}\)
for \(q_0\) derivatives. This proves the required
supremum bounds from the coefficient sum.
No inverse threshold is introduced in this
global estimate. The kernel estimates used to
obtain \eqref{eq:trace-global-current-path-price}
are the exponentially weighted flat-jet
estimates, before finite polynomial completion.
The curvature completion costs just the stated
path moments. Cutting the kernel sums at
\(R_{\rm cut}\) removes terms and cannot
increase this positive majorant.

For clarity, a connected path system of length
\(\ell\), with a fixed finite number of branches,
is covered by at most \(C_j(1+\ell/B_k)\)
connected \(B_k\)-balls. Place their centres
at successive path distances at most \(B_k\),
and add the finitely many branch endpoints.
The same estimate holds for
\(\widetilde X_\lambda=X_\lambda\cup\Sigma_5\)
and its flag halo, after enlarging \(C_j\).
Choose entry so that
\(C_j(\kappa+2)/B_k\le a_j/2\).
The support weights of these terms are then
bounded by a fixed constant times
\(e^{a_j\ell_\lambda/2}\).
The remaining half width pays their finite
path and projection moments. This proves
summability in the \(\kappa+2\) norm uniformly
in the truncation radius. Assigning one atom
the whole radius-\(R_{\rm cut}\) box would
not give this estimate.

\emph{Write the physical complement explicitly.}
Put
\[
 y_\lambda=\Pi_{5,h_k}G_\lambda,\qquad
 A_\lambda=P_5[\mathsf D_{h_k^{-1}}y_\lambda],
                 \qquad R_\lambda=G_\lambda-A_\lambda.
\]
Use the masked regular part
\[
 g^{\rm reg}=
       \chi_{\Sigma_5}\sum_\lambda A_\lambda
               +\sum_\lambda\chi_{\widetilde X_\lambda}R_\lambda.
\]
Its exact complement is the response family
\begin{equation}
 L^g_{\Sigma_5}=(1-\chi_{\Sigma_5})\sum_\lambda A_\lambda,
 \qquad
 L^g_{\widetilde X_\lambda}
                  =(1-\chi_{\widetilde X_\lambda})R_\lambda.
 \label{eq:trace-explicit-current-response}
\end{equation}
When supports coincide, add their functions.
Each response vanishes on its healthy sector,
is Borel and gauge invariant, and has the
stated final-flag locality. Their sum is
exactly \(g_b-g^{\rm reg}\), including on
unhealthy configurations.

There is no hidden factor \(h_k^{-2}\) in
their global amplitude. In fact
\(\mathsf D_{h_k^{-1}}\Pi_{5,h_k}G_\lambda
=\Pi_5G_\lambda\); the factor-count
normalisations cancel exactly before the
primary function is evaluated. The fixed
radial read and primary section therefore
bound \(\|A_\lambda\|_\infty\) by
\(C_j|c_\lambda|(1+\ell_\lambda)^{q_0}\)
for a fixed \(q_0\). The same is true for
\(R_\lambda\), using the preceding bound for
\(G_\lambda\). Combining this fact with
\eqref{eq:trace-global-current-path-price}
and the connected support bound proves the
first inequality in
\eqref{eq:trace-actual-response-budget}.
At each fixed profile its right side is a
logarithmic power, so the second inequality
is a cutoff- and volume-independent entry
condition.

Apply Lemma~\ref{lem:physical-response-birth}
to \eqref{eq:trace-explicit-current-response}.
It produces exactly the forced profile
\(\mathcal H_KL^g\) in the one-source birth
construction, with price at most
\(D_{B_k}\sqrt{w_k}\,e(L^g)\) in the divided
norm. This fits its original coefficient
budget. Indeed, with \(x=\widehat\beta_k\)
and \(h_k=x^{-1/2}P_k\),
\[
 \log\frac{D_{B_k}\sqrt{w_k}
                  C_j(1+\log x)^{A_j}}
                 {h_k^2P_k^{1/6}}
 \le-\tfrac12P_k^{1/5}+\log x
                     +O_j(\log\log x+\log P_k)\longrightarrow-\infty.
\]
Here \(P_k^{1/5}/\log x\to\infty\).
The regular birth estimate remains
\eqref{eq:trace-joint-source-price}.
All these inequalities are local component
estimates; periodising their positive path
majorants does not increase their sums.
They are uniform in the terminal cutoff
and physical side on the prescribed tails.
\end{proof}

For later use the coefficient \(a^f\) or
\(a^g\) is chosen once by the matching map.
Apply Lemma~\ref{lem:physical-response-matching-root}
separately to these two physical sources.
The root comparison of their linear combination
then uses the triangle inequality for Wilson
roots, with the actual coefficient multiplying
each source. This avoids imposing a false
response bound on a coupling-multiplied source.

% !TeX root = 4DYM_notes.tex
\subsection{Degree-three potentials and their physical forward differences}
\label{subsec:degree-three-potentials}

The trace and antisymmetric identities involve the
potentials \(q_a,t_{abc}\), not just the degree-four
current. We need to match a potential in degree three
and its actual forward difference in degree four.
The smallness of the latter field will follow from
translation of the former unit field. It does not
follow from setting a zero radial read equal to its
matched coefficient. We give the response construction
and the two physical comparison estimates before
using that translation argument.

\begin{lemma}[Birth and root bounds for the current potentials]
\label{lem:degree-three-potential-birth}
Let \(H=H_{n,M,k}\) be any real scalar component of
the finite polynomial representatives \(q_a\) or
\(t_{abc}\) constructed above. Fix an axis \(\nu\)
and write \(\delta_\nu=T_\nu-1\), without a lattice
spacing factor. For
\[
              (d,u)=(3,H)\quad\hbox{or}\quad
                         (4,\delta_\nu H),
\]
the dimensionless physical function \(u\) has an
exact one-insertion birth on the same density
\(F_{n,k}\), with radial read zero and
\begin{align}
 \|R_u\|_{\rm ins}
 &\le A_j^{\rm pot}(1+\log x)^{p_j^{\rm pot}}
                   \{h_k^2+D_{B_k}\sqrt{w_k}\}
                    \le h_k^2P_k^{1/6},\notag\\
 e(L^u)&\le A_j^{\rm pot}(1+\log x)^{p_j^{\rm pot}}
                    \le x,\qquad x=\widehat\beta_k.
 \label{eq:degree-three-potential-birth-price}
\end{align}
Here \(L^u\) is an actual physical response family.
The finite constants \(A_j^{\rm pot},p_j^{\rm pot}\)
may depend on the group, profile derivative and spatial
orders, the fixed block ratio, kernel widths and finite
polynomial completion. They are independent of \(n,M\).
These constants are birth constants, distinct from
the selected-column constants \(C_j\).

The inequalities hold after the common profile entry
and on the prescribed ultraviolet and physical-side
tails at each fixed \(k\). The nonwrapping side
threshold includes the full truncated read and its
flag halo, and is uniform in the terminal cutoff.
Let \(a^{d,u}_{n,M,k}\) be the canonical matched
reference coefficient. At fixed grid and stopping
levels \(i,m\), with
\(\operatorname{ref}(d)\le m\) and
\(\varepsilon_i/\varepsilon_m\ge B_m^2\),
one has, for all sufficiently large \(k>m\),
\begin{equation}
 \left\|\varepsilon_k^{-d}\mathcal P_{n\to k}^*u
           -\mathcal O_{n,M}(a^{d,u}_{n,M,k})\right\|_i
            \le C_{i,m,d}\widehat\beta_k^4
                                    \varepsilon_k^{1/8}.
 \label{eq:degree-three-potential-root}
\end{equation}
The constant is independent of \(k,n,M\), once the
fixed grid, stopping level and degree are specified;
the matching side threshold may depend on \(k\).
The degree-three coefficient is a unit coefficient,
which is retained and need not vanish.
\end{lemma}
\begin{proof}
\emph{The actual functions and their low reads.}
Use the polynomial representatives in the exact
trace and antisymmetric completion. Each component
has the path expansion
\(H=\sum_\lambda H_\lambda\), with its natural
anchored supports \(X_\lambda\). The scalar
potentials have no unit or linear field term.
Their completion through field order four uses
products of two, three or four transported
curvature atoms, as in
\eqref{eq:trace-current-diagonal-table}.
A scalar invariant linear in the Lie algebra
would be an invariant linear functional on
\(\mathfrak g\); simplicity makes it zero.
This explains why two is the minimum number
of curvature inputs. Finite shifts and the
exact completion of the identities preserve
this property.

Each curvature atom has engineering weight at
least two. Thus the lifted radial test gives
\(\Pi_3H_\lambda=0\). Its possible weight-four
term is quadratic in the curvature at the
anchor. Shifting the whole stencil by one
native site shifts that anchor by
\(\varepsilon e_\nu\) in the test connection.
The value of that quadratic term changes by
\(O(\varepsilon)\), by Taylor's formula for
the smooth curvature. The two weight-four
coefficients therefore cancel in
\(T_\nu H_\lambda-H_\lambda\), and
\begin{equation}
            \Pi_3H=0,\qquad
                          \Pi_4(\delta_\nu H)=0.
 \label{eq:degree-three-potential-zero-reads}
\end{equation}
This is a statement on the lifted, nonwrapping
read; the fixed-\(k\) side tail supplies that
read. A wrapping comparison term is kept
among the earlier spatial errors.

For the difference use the components
\(H_\lambda^\Delta=T_\nu H_\lambda-H_\lambda\)
on \(X_\lambda\cup(X_\lambda+e_\nu)\).
These footprints contain the original anchor
and both complete reads. There is no factor
\(\varepsilon_k^{-1}\) in this dimensionless
component. Indeed the physical identity is
\begin{equation}
 \varepsilon_k^{-1}(T_\nu-1)
                   (\varepsilon_k^{-3}H)
                            =\varepsilon_k^{-4}\delta_\nu H.
 \label{eq:degree-three-physical-difference}
\end{equation}

\emph{Regular and global prices.}
The spare exponentially weighted path estimate
\eqref{eq:trace-global-current-path-price}
applies also to \(q,t\). Their normalized
four-jet bounds were included in the last
row of \eqref{eq:trace-current-diagonal-table}:
restoring the normalization costs only fixed
powers of \(\log x\). Increasing the path
moment by a fixed amount pays the one-site
shift and any admitted derivative order.

We check why this gives the full healthy
derivative norm of the completed polynomial.
On a healthy support a transported curvature
atom has norm at most \(C_Gh_k\), since its
parallel transports are unitary. An ordered
word of \(b\ge1\) left derivatives has global
bound at most a fixed path polynomial; after
threshold scaling it costs at most
\(C_{j,b}h_k^b(1+\ell_\lambda)^{c_{j,b}}\).
Because \(h_k\le1\), this is bounded by one
factor \(h_k\) times that polynomial.
In the product rule for a term with \(r\)
curvature inputs, every undifferentiated
input supplies \(h_k\), and every
differentiated input supplies at least
\(h_k\). The result is at most
\(C_jh_k^r(1+\ell_\lambda)^{c_j}\), hence
at most \(C_jh_k^2\) times a path polynomial.
The number of assignments of an ordered
derivative word is at most \(4^{q_k}\), a
finite profile constant. The anchored spatial
direction weights add another fixed path
power. The spare exponential width sums
all these powers. Only the four-jet kernel
coefficients, and derivatives of their actual
finite polynomial completion, have been used;
no higher stationary-current jet is required.

For the global supremum drop the factors
\(h_k\) and use unitarity instead. A connected
system of paths of total length
\(\ell_\lambda\), including the shift and
flag halo, has covering number at most
\(C_j(1+\ell_\lambda/B_k)\). Entry leaves
enough width for exponent \(\kappa+2\).
Consequently, for either component family
\(u=\sum_\lambda u_\lambda\), with supports
denoted by \(Z_\lambda\),
\begin{align*}
 \sum_\lambda e^{\kappa n_{B_k}(Z_\lambda)}
     \|u_\lambda\|_{q_k,h_k,s_k;0,Z_\lambda}
    &\le A_j^{\rm pot}h_k^2(1+\log x)^{p_j^{\rm pot}},\\
 \sum_\lambda e^{(\kappa+2)n_{B_k}(Z_\lambda)}
                       \|u_\lambda\|_\infty
    &\le A_j^{\rm pot}(1+\log x)^{p_j^{\rm pot}}.
\end{align*}
The same bounds hold for the truncated sums:
the truncation discards terms of these positive
majorants. Periodization does not increase
their total sums.

\emph{Write the responses.}
No primary subtraction is needed in
\eqref{eq:degree-three-potential-zero-reads}.
Set
\[
 u^{\rm reg}=\sum_\lambda\chi_{Z_\lambda}u_\lambda,
 \qquad
 L^u_Z=\sum_{\lambda:Z_\lambda=Z}
                               (1-\chi_Z)u_\lambda.
\]
Then \(u=u^{\rm reg}+\sum_ZL^u_Z\) exactly,
including at the threshold. The \(L^u_Z\)
are local Borel responses, with the original
flag halos. The masked regular functions
satisfy the requisite zero-jet conditions.
Lemma~\ref{lem:physical-response-birth} gives
their actual forced profile and its divided
price at most \(D_{B_k}\sqrt{w_k}\,e(L^u)\).
This proves the first inequalities in
\eqref{eq:degree-three-potential-birth-price}.

At fixed profile,
\(A_j^{\rm pot}(1+\log x)^{p_j^{\rm pot}}
=o(P_k^{1/6})\) and is at most \(x\) after
entry. The divided forced term fits the
same budget because
\[
 \log\frac{D_{B_k}\sqrt{w_k}\,
       A_j^{\rm pot}(1+\log x)^{p_j^{\rm pot}}}
                    {h_k^2P_k^{1/6}}
 =-\tfrac12P_k^{1/5}+\log x
       -\tfrac{13}{6}\log P_k+O_j(\log\log x)\longrightarrow-\infty.
\]
Choose entry so that the regular and forced
parts each use at most half the stated
budget. These are finitely many requests
before the common family is selected.

\emph{The two root estimates.}
Apply Lemma~\ref{lem:physical-response-matching-root}
with degree three and degree four respectively.
It keeps the actual matched coefficients,
even though the two birth reads are zero.
At fixed \(i,m\) the exact right side is
\[
 \frac{2\widehat\beta_k^4\varepsilon_m^{-d-1/8}}
                        {\zeta_{d,m}}
       e^{(\varepsilon_i/\varepsilon_m)^4Q_m^{\rm vac}+5}
                                     \varepsilon_k^{1/8}.
\]
This gives \eqref{eq:degree-three-potential-root}
with all its constant dependencies displayed.
The full physical read radius is a fixed
profile logarithmic power, and the existing
entry bounds make it at most a fixed power
of \(\log\widehat\beta_k\). The native room
condition therefore holds eventually at
fixed \(i\), uniformly in terminal cutoff.
The local inequalities and their constants
are uniform in \(n,M\) on the stated tails.
Taking limits first in the terminal cutoff
at fixed physical side and then in the
chosen side sequence preserves them. Only
afterwards does \(k\) tend to infinity.
\end{proof}

\begin{corollary}[Matched forward differences of the potentials]
\label{cor:degree-three-forward-difference}
Include the preceding coefficient arrays in
the same countable ordered refinement as
the other current data. Denote their fields
by \(B_k^H\in\mathcal F_3\) and
\(B_k^{\delta_\nu H}\in\mathcal F_4\).
For some fixed \(\eta>0\), which can be
chosen smaller than \(1/8\),
\begin{equation}
 B_k^H=c_k1_3,\qquad
             \|B_k^{\delta_\nu H}\|
                              =O(\varepsilon_k^{\eta/2}).
 \label{eq:degree-three-forward-difference-decay}
\end{equation}
The scalar \(c_k\) is retained; no smallness
of \(c_k\) is assumed.
\end{corollary}
\begin{proof}
Simplicity gives the first assertion, including
when the matching has generated a unit
coefficient. Fix a separated exterior response
row and vary only the potential's anchor along
axis \(\nu\). The root bound and the Wilson
chessboard inequality bound its matching error
by \(E_k=O(\varepsilon_k^\eta)\), and that of
the forward-difference source by
\(E'_k=O(\varepsilon_k^\eta)\). Here
\(\widehat\beta_k^4=e^{o(s_k)}\) is absorbed
by decreasing \(\eta<1/8\). Literal translates
use the same matched coefficient. Kernel
covariance above birth and translation
invariance of the original Wilson measure
give the same bounds throughout a fixed
small interval of anchor translations.

Work first at finite \(n,M\), suppressing those
two indices. Write \(F_k(l)\) for the physical
potential response at anchor
\(l\varepsilon_k e_\nu\), and \(G_k(l)\)
for the response of its independently matched
ordinary degree-four column. The physical
identity \eqref{eq:degree-three-physical-difference}
is exact. The matched potential response is
constant in \(l\); call it \(b_k\). Thus, on
every finite interval in the ordered tails,
\[
 |F_k(l)-b_k|\le E_k,\qquad
 \left|G_k(l)-\frac{F_k(l+1)-F_k(l)}{\varepsilon_k}\right|
                                                       \le E'_k.
\]
The finite-cutoff constant \(b_k\) contains
the unit coefficient and the unchanged
exterior correlation. Put
\(N_k=\lceil\varepsilon_k^{-1+\eta/2}\rceil\).
Telescoping, before any absolute value,
gives
\begin{align*}
 \left|\frac1{N_k}\sum_{l=0}^{N_k-1}G_k(l)\right|
 &\le\frac{|F_k(N_k)-F_k(0)|}{N_k\varepsilon_k}+E'_k\\
 &\le\frac{2E_k}{N_k\varepsilon_k}+E'_k
                                  =O(\varepsilon_k^{\eta/2}).
\end{align*}
Only endpoint errors have been divided by
the averaging length. For each fixed \(k\),
the number \(N_k\) is finite. Pass the matched
responses \(G_k(l)\) through the original
cutoff limit and then the chosen side limit,
including their coefficient arrays in the
same refinement. The bound is uniform on
those tails. No limit of the physical
responses \(F_k(l)\) is required: they have
already disappeared into the endpoint bound.
The averaging length tends to
zero. On each fixed basis vector the
averaged response row therefore converges
to the row at the base point, by smoothness.
Choose finitely many separating rows on
\(\mathcal F_4\). Their averaged matrices
converge to an injective fixed matrix. Equation
\eqref{eq:common-fixed-detector-bound} gives one bounded
left-inverse estimate on their tails. Applying it
to the possibly varying matched field gives
\eqref{eq:degree-three-forward-difference-decay}.
No bound on its norm is used to compare
the averaged and unaveraged rows.
\end{proof}

% !TeX root = 4DYM_notes.tex
\emph{The raw read and the twelve planes.}
We calculate the conversion to \(\rho_{\rm tr}\).
Write \(V^{(1)},V^{(2)},V^{(3)}\) for the three raw
vector basis polynomials in
Section~\ref{sec:trace-ward-coefficient}, in their
displayed order. Fix a unit colour and the external
link polarisation \(c\). Its linear curvature is
\[
                 F_{ij}(p)=i(p_i\delta_{jc}-p_j\delta_{ic}).
\]
The real colour metric is extended complex bilinearly
for these Fourier derivatives.
For \(p=ze_a\), \(q=-ze_a+te_b\), \(a\ne c\), the
three quadratic kernels are respectively
\begin{align*}
 \widehat V_b^{(1)}(p,q)
      &=i(p_b+q_b)\sum_dF_{bd}(p)F_{bd}(q),\\
 \widehat V_b^{(2)}(p,q)
      &=i(p_b+q_b)\sum_{i,d}F_{id}(p)F_{id}(q),\\
 \widehat V_b^{(3)}(p,q)
      &=\sum_{i,d}\{F_{bd}(p)\,iq_iF_{id}(q)
                                  +F_{bd}(q)\,ip_iF_{id}(p)\}.
\end{align*}
These are second derivatives, so both placements of
the two inputs are already included. For one ordered
plane \((a,c)\), direct substitution gives
\[
\begin{array}{c|rrr|r}
 &b=a&b=c&\text{each }b\notin\{a,c\}&\text{sum in }b\\ \hline
 \operatorname{Re}i[z^2t]\widehat V_b^{(1)}&-1&-1&0&-2\\
 \operatorname{Re}i[z^2t]\widehat V_b^{(2)}&-2&-2&-2&-8\\
 \operatorname{Re}i[z^2t]\widehat V_b^{(3)}&-1&-1&1&0.
\end{array}
\]
There are two values in the third column. For example,
in the first row at \(b=a\) the kernel is
\(it(z^2-zt)\), and at \(b=c\) it is \(itz^2\).
For the last row these same two cases have cubic
coefficients \(i\) and \(i\), whereas each other
\(b\) has coefficient \(-i\); their sum is zero.

The link-base correction can also be obtained directly.
For a continuum plane wave in direction \(c\), the
linear variation of its link holonomy contains the
factor
\[
 a_1(p_c)=\int_0^1e^{isp_c}\,\dd s
                                      =1+ip_c/2+O(p_c^2).
\]
The two inputs therefore multiply a left-link kernel
by \(a_1(p_c)a_1(q_c)\). On the tested ray \(p_c=0\)
and \(q_c=t\delta_{bc}\). Its linear total-momentum
factor changes \(\operatorname{Re}i[z^2t]\) by
\(-\tfrac12\delta_{bc}\operatorname{Re}[z^2]\).
There is no second derivative of the holonomy acting
on a first scalar jet, since that jet vanishes by
colour invariance. This is exactly the correction
in \(\mathcal H\), not an independently chosen
subtraction.

The pure weight-five basis kernels have momentum
degree three, so their degree-two correction is
zero. There are \(4\cdot3=12\) ordered planes.
The table gives
\begin{equation}
 \mathcal H(xV^{(1)}+yV^{(2)}+zV^{(3)})
             =-24x-96y=12\rho_{\rm tr}(x,y,z).
 \label{eq:trace-twelve-plane-read}
\end{equation}
Thus \(\mathfrak h(Y)=\mathcal H(\operatorname{div}_{\rm Haar}Y)\)
enters the raw trace functional with the factor
\(1/12\), with its sign fixed by the displayed Fourier
convention.

There is a short direct check for the loop term.
Let \(K_{\rm loop}(p)\) be its quadratic link kernel.
The two moment cancellations of
Proposition~\ref{prop:ordinary-one-loop-matching} give
\(K_{\rm loop}(p)=O(|p|^4)\). Since the loop scalar
has zero first jet and \(Y(1)=0\), the quadratic
kernel of \(D\mathcal L_{n,k}[Y_b]\) is the sum of
the two pairings
\[
 \langle K_{\rm loop}(p)A_p,
                         \widehat{DY_b}(-p;q)A_q\rangle
  +\langle K_{\rm loop}(q)A_q,
                         \widehat{DY_b}(-q;p)A_p\rangle.
\]
The displacement kernel is analytic at zero. Each
pairing has total momentum order at least four,
so both tested coefficients in \(\mathcal H\)
vanish. This proves its zero \(\rho_{\rm tr}\) read
without a separate polynomial-link case expansion.

% !TeX root = 4DYM_notes.tex
Match the two complete dimension-five sources
\(\varepsilon_k^{-5}f\) and \(\varepsilon_k^{-5}g\)
of \eqref{eq:trace-joint-source-price}. At finite cutoff
and physical side let their reference coefficients be
\(a^f_{n,M,k},a^g_{n,M,k}\). Define their normalised
current coordinates by
\begin{equation}
 \Sigma_{n,M,k}=h_k^{-2}\mathsf Z_{n,M,k,5}^{-1}a^f_{n,M,k},
 \qquad
 j_{n,M,k}=h_k^{-2}\mathsf Z_{n,M,k,5}^{-1}a^g_{n,M,k}.
 \label{eq:trace-actual-matched-coordinates}
\end{equation}
Only the real odd vector sector is taken here; it has
two curvature factors in every coordinate. These are
matched coordinates of complete physical sources. Their
local radial reads, divided by \(h_k^2\), are denoted
by \(\Sigma_k^{\rm read},j_k^{\rm read}\).
Lemma~\ref{lem:ward-matching-coefficient-control} and
\eqref{eq:trace-joint-source-price} give on the stated
large-volume tails
\begin{equation}
 \|\Sigma_{n,M,k}-\Sigma_k^{\rm read}\|
       +\|j_{n,M,k}-j_k^{\rm read}\|
                                      \le CP_k^{-1/4}.
 \label{eq:trace-read-matched-difference}
\end{equation}
The coefficient norms are equivalent fixed finite norms;
their conversion is included in the entry request in
\eqref{eq:ward-normalised-matching-bias}.

At fixed \(k\), first use the ultraviolet tail to
remove \(C_k\beta_k2^{-s}\), and then the period
tail to remove \eqref{eq:trace-wrapping-error}.
Entry makes the remaining contribution of each
of the four error types at most \(P_k^{-1}\).
Choose the same tails so that
Lemma~\ref{lem:trace-regulated-ward} has error at
most \(\widehat\beta_k^{-100}\). Its current
seminorm bounds are uniform in precisely these
tails. Kernel truncation and the lifted local
moment read can likewise have error at most
\(\widehat\beta_k^{-30}\), by the exponential
tail bound already proved. Equation
\eqref{eq:trace-twelve-plane-read} therefore gives
\[
 \left|\rho_{\rm tr}(\Sigma_k^{\rm read})-\tfrac1{12}\mathfrak h(Y)\right|
   \le4\|\rho_{\rm tr}\|P_k^{-1}
                                 +\widehat\beta_k^{-30}/12.
\]
These are tail choices in the already convergent
physical rows; they do not redefine the Wilson
family or its ordered hierarchy.

The full coefficient norm also has a uniform
bound. The divergence read uses only a fixed
finite set of current jets and spatial moments.
Their bounds are fixed powers of
\(\log\widehat\beta_k\), by the current proof
and the exponential seminorm reduction. The
loop term has at most its original allowance
\(2P_k^{\theta_{J(k)}}\) times such factors.
Consequently, before increasing entry, the
coefficient norm is bounded by
\[
 C_j(1+\log\widehat\beta_k)^{C_j}
                     (1+2P_k^{\theta_{J(k)}})
                  +\sum_{\bullet}\|\varepsilon_\bullet\|.
\]
The exponent and constant are fixed at a profile.
For \(\theta_j<1/6\), use the positive margin
\((1/6-\theta_j)/2\). Since
\(\log P_k=(\log\log\widehat\beta_k)^2\), entry
makes the first fixed polynomial at most
\(P_k^{(1/6-\theta_j)/2}\). Thus its total exponent is
strictly below \(1/6\); the remaining small errors
fit the same bound. We have proved
\begin{equation}
 \|\Sigma_{n,M,k}\|\le CP_k^{1/6},\qquad
 \left|\rho_{\rm tr}(\Sigma_{n,M,k})
                       -\tfrac1{12}\mathfrak h(Y)\right|
       \le CP_k^{-1/4}+4\|\rho_{\rm tr}\|P_k^{-1}
                                  +\widehat\beta_k^{-30}/12.
 \label{eq:trace-sigma-bounds}
\end{equation}
Constants involving the fixed profile widths,
paths and finite current jets have been used for
entry only. The estimates precede phase evaluation
and are uniform in the cutoff and the physical
period on their stated ordered tails.

The current divergence has local read tending to
\(\operatorname{div}_{\rm raw}L=-v_{\rm eq}\).
Indeed its quadratic curvature jet is that of \(L\),
and the central difference is exact on the relevant
degree-five jet. Moreover its local \(\rho_{\rm tr}\)
read is zero in the exact current identity: the tensor
trace is a forward difference of degree-three scalar
potentials, whose nonconstant radial space is empty by
simplicity. Finite current and kernel approximations
have the prices already charged above. Consequently
\begin{equation}
 j_{n,M,k}\longrightarrow-v_{\rm eq},\qquad
       |\rho_{\rm tr}(j_{n,M,k})|\le CP_k^{-1/4}
 \label{eq:trace-matched-current-bounds}
\end{equation}
on the ordered tails followed by \(k\to\infty\).
The same construction for the full divergences of
the ordinary \(\Xi,D\) stencils gives normalised
matched coordinates \(w_{n,M,k}^\Xi,w_{n,M,k}^D\) with
\begin{equation}
 \|w_{n,M,k}^B-v_B\|\le CP_k^{-1/4},\qquad
     v_B=\operatorname{div}_{\rm raw}B,\quad B=\Xi,D.
 \label{eq:trace-ordinary-matched-bias}
\end{equation}
Their full radial reads are \(v_B\); their higher
functions and forced responses remain in the matching.
All these arrays, their reference coefficients and the
finite coefficient maps are included in the same
countable ordered refinement. We henceforth denote their
limits at fixed \(k\) by \(\Sigma_k,j_k,w_k^B\).

Take joint cluster values of the matched coefficient arrays,
first on ultraviolet tails at each fixed period and then on
the chosen period subsequence. All closed estimates survive.
These refinements preserve the already fixed hierarchy.
The resulting real coefficients satisfy
\begin{equation}
             \rho_{\rm tr}(\Sigma_k)\longrightarrow-b_0,
                      \qquad \|\Sigma_k\|\le CP_k^{1/6}.
 \label{eq:trace-sigma-limit}
\end{equation}

\subsection{Quantitative control of the equation direction}

At the prescribed radii of
Lemma~\ref{lem:trace-prescribed-radii}, match the degree-four
physical current stencil \({\varepsilon}_{k_i}^{-4}J\) to the reference
field space, symmetrise its two indices, average over real
\(H_4\) covariance, and subtract its mixed mean. Denote the
result by \(B_i^L\). Matching the local stencils for
\(\Xi,D\) in exactly the same convention gives
\(B_i^\Xi,B_i^D\). Matching here is the canonical backward
coefficient construction of
Section~\ref{sec:position-local-matching}, with its residual
retained, rather than an assertion that the stencil is
pointwise equal to its finite jet.

For each fixed \(k,d\), let \(\mathsf C_{k,d}\) send a
current threshold coordinate to the continuum field of
the selected column with that coordinate. More explicitly,
it is the reference projection composed with the ordered
limit of \(\mathsf Z_{n,M,k,d}\). The selected column
includes its full regular and forced graph remainder and
its physical factor \(\varepsilon_k^{-d}\).
All these finite maps and their inverses are extracted
simultaneously with the coefficient arrays above; their
uniform bounds at fixed \(k\) justify this extraction.
On the raw three-dimensional vector space define
\begin{equation}
 s_i v=\beta_i r_i^5U_i(5)^{-1}
                      \mathsf C_{k_i,5}(h_{k_i}^2v).
 \label{eq:trace-selected-column-map}
\end{equation}
Its target is the real two-dimensional free tensorial
vector plane. This map acts on the normalised coordinates
of selected columns. It does not mean matching an arbitrary
literal local section with radial read \(v\).
Write \(s\) for the free Maxwell quotient map, with
kernel \(\R v_{\rm eq}\). Below \(j_k,\Sigma_k,w_k^B\)
always mean the matched coordinates of
\eqref{eq:trace-actual-matched-coordinates} and
\eqref{eq:trace-ordinary-matched-bias}.

Put
\begin{equation}
 b_i=\widehat\beta_{k_i},\quad E_i=P_{k_i}^{1/5},\quad
 R_i=P_{k_i}^{1/6},\qquad
 \mathsf A_i=\beta_ir_i^4U_i(4)^{-1},\qquad
 K_i=\beta_ir_i^5U_i(5)^{-1}\operatorname{div}.
 \label{eq:trace-transfer-scales}
\end{equation}
Tensor maps act componentwise. Notice \(R_i/E_i\to0\),
\(E_i/\beta_i\to0\), and \(\beta_i/b_i\to1\).

% !TeX root = 4DYM_notes.tex
\begin{lemma}[Complete-source Ward transport]
\label{lem:complete-source-ward-transport}
Use the complete-source matching and the ordered limits just
defined. Put
\[
             v_i^{\rm eq}=\beta_i^{-1}\Sigma_{k_i}-j_{k_i}.
\]
Then, for every fixed \(p\in\N\),
\begin{equation}
                    b_i^p\|s_i v_i^{\rm eq}\|\longrightarrow0.
 \label{eq:trace-complete-ward-image}
\end{equation}
There are raw coefficient vectors \(y_i^L,y_i^\Xi,y_i^D\)
such that
\begin{align}
 K_iB_i^L&=s_i(j_{k_i}+y_i^L),\notag\\
 K_iB_i^B&=s_i(w_{k_i}^B+y_i^B),\qquad B=\Xi,D,\notag\\
 b_i^p\|y_i^B\|&\longrightarrow0,
                         \qquad B=L,\Xi,D.
 \label{eq:trace-complete-derivative-lift}
\end{align}
The degree-four comparison gives
\(\mathsf A_iB_i^B\to B\), \(B=L,\Xi,D\), and
\(K_iB_i^B\to\operatorname{div}^{\rm Max}B\) for
\(B=\Xi,D\). In addition,
\begin{equation}
 b_i^p\operatorname{tr}(\mathsf A_iB_i^B)\longrightarrow0,
                                      \qquad B=L,\Xi.
 \label{eq:trace-complete-tensor-traces}
\end{equation}
The finite-cutoff bounds are uniform on the specified cutoff
and physical-side tails. No evaluation at a phase, or change
of the selected hierarchy, is involved.
\end{lemma}
\begin{proof}
\emph{The finite physical Ward comparison.}
Work first at fixed \(n,M,k\). Keep the exact smooth density
\(F_{n,k}\), its partition integral \(Z_{n,k}\), and the
smooth displacement \(Y\). The signed kernel numerator in
\eqref{eq:trace-marked-ward-step} is the equation source.
For an exterior physical cylinder \(A\), its integral is
zero, by \eqref{eq:trace-one-step-weak-ward}.
The same assertion holds for a bounded Borel exterior
probe: integrate first over the links on which \(Y\)
acts. The probe is constant in those variables, and
Haar integration by parts applies to the smooth density;
Fubini then integrates the remaining variables.
We compare
\(f-\beta_kg\) with this numerator, before division by
\(Z_{n,k}\). This is why both complete sources, rather
than the local polynomial \(\Sigma_k^{\rm read}
-\beta_kj_k^{\rm read}\), were matched above.

Here are the estimates for that comparison. Set
\[
 D_k=(1+B_k)^{40},\qquad
                 R_{\rm cut}=C_j^{\rm cut}(1+B_k)^{61}.
\]
Expand the differentiated kernel with its one distinguished
mark, together with the physical source used to define
\(f\). On a connected term entirely inside the good local
chart the product rule gives exactly
\eqref{eq:trace-central-ward-germ}, with its entire higher
functions retained. In particular the classical fifth and
higher field terms, the regular action variations, and the
variation of the normalising density all remain in this
identity. They are not small errors in the Ward identity.
Their prices in \eqref{eq:trace-complete-source-birth-price}
are used for matching, not for discarding these functions.

The two numerators therefore differ only by spatial
truncations, the prescribed finite current approximations,
and the chart restorations. For a regular support with
covering number \(c_X\), write its full price as \(v_X\).
The elementary estimate
\begin{equation}
 \sum_{c_X\ge q}e^{98c_X}v_X
                  \le e^{-2q}\sum_Xe^{100c_X}v_X
 \label{eq:trace-spatial-tail}
\end{equation}
holds before or after summing the component addresses.
A term reaching \(R_{\rm cut}\) has
\(q\ge R_{\rm cut}/(2B_k+1)-1\). Its price, including
the finitely many score, path and coefficient factors,
is consequently at most
\[
 C_j\widehat\beta_k^{A_j}(1+R_{\rm cut})^{A_j}
   \left\{e^{-a_j'R_{\rm cut}}
                 +e^{-2R_{\rm cut}/(2B_k+1)+2}\right\}
                                               \le e^{-2D_k}
\]
after entry. Fixed widths and constants are included in
\(C_j^{\rm cut}\) and entry in this order. The powers
\(61,40\) leave a strict margin; their exact values are
immaterial. Finite current approximations at fixed \(k\)
are made on the ultraviolet tail to the same accuracy.
Their bounds are independent of the physical side.

A chart-restoration term has a distinguished physical
failure occurrence. Its complete incoming multiplier and
its original eligible inventory are kept. In a forced
term divide its price by \(\sqrt{w_k}\) once and retain
one \(\sqrt{w_k}\) at that occurrence. For a Gaussian
warning use its stronger integrated warning price instead.
For each fixed source subset the marked component estimate
then gives an income bounded by
\[
 C_j\widehat\beta_k^{A_j}(1+R_{\rm cut})^{A_j}
                     \exp(-P_k^{1/5}/2)
                              \le\exp(-P_k^{1/5}/4).
\]
All other occurrences are summed by the rooted tree
estimate. The divided forced norm, not a derivative of
the forced multiplier, is used here. This calculation
also applies to the chart clipping in the definition of
\(f\): its support of disagreement supplies the same
failure occurrence. It is a marked-numerator comparison;
it does not divide separate component estimates by a
pointwise gas or assume that the logarithm of a Borel
factor can be differentiated.

% !TeX root = 4DYM_notes.tex
\emph{A complete joint test with its physical normalization.}
We spell out one vector/tensor test; finitely many such tests
will detect the two-dimensional vector plane. Keep \(n,M\)
finite throughout. Fix a radius index \(i\), put
\(v=n_i\), \(k=k_i\), and
\(H=(\log\widehat\beta_v)^{48}\). The Ward insertion is
born at the origin of level \(k\). Choose a degree-four
tensor probe at a rounded point a fixed positive multiple
of \(r_i\) away. It is born at the same level, using the
finite curvature section and the critical coefficient
column of the prescribed-radius comparison. Denote its
dimensionless polynomial by \(u\). Its physical insertion
is \(\varepsilon_k^{-4}u\). All finite colour and tensor
recombinations are included in this choice.

We check that choosing this probe does not hide a large
inverse coefficient matrix. On the degree-four critical
block \(w=d=4\), the leading map is exactly its diagonal
threshold ratio: higher-weight entries are absent. After
conjugating by \(h_t^2\), the actual block is identity
plus its coefficient error. The calculation of
\eqref{eq:ward-critical-column-drift} sums these errors
over the present window and gives
\[
 \left\|\prod_{t=v}^{k-1}A_t^{\rm crit}-1\right\|
    \le C(\log\widehat\beta_v)^{49}P_v^{-1/2}=o(1).
\]
Indeed the number of steps is at most \(H/\log3+1\),
and \(P_k/P_v\to1\). The finite basis normalization
at \(v\) has a bounded inverse by its free detector
matrix; the displayed product has inverse at most two
eventually. Thus fixed bounded native curvature
coefficients at \(k\) supply the required probes, up
to these bounded recombinations and the explicit
Gaussian factor \(\beta_k\) used below. An inverse
over the whole reference interval is not used.

Unit terms can first be separated by their exact transport.
For the test below one can center the tensor column by its
original Wilson expectation. A vector Ward field has zero
one-point function by its signed symmetry, so this centering
does not change the test. The subtracted unit remains an
actual constant in its source polynomial. The nonconstant
degree-four labels have two curvature factors, and
\(\|u\|_{\rm healthy}\le C_jh_k^2\) with their admitted
derivatives. The response construction of
Lemma~\ref{lem:trace-physical-response-data}, applied to
this finite section, gives a logarithmic-power global
price. Thus the same \(q\le Ch_k^2P_k^{1/6}\) and
\(e(L)\le\widehat\beta_k\) entries apply. A finite tensor
basis suffices, so all these requests are in the common
initial calendar.

Suppress \(n,M\) in this calculation and write
\(\langle A\rangle_k=Z_{n,k}^{-1}\int A F_{n,k}\,\dd H_k\).
The exact coarse Ward test is
\begin{equation}
 \left\langle
    \varepsilon_k^{-5}Q_b^{\rm eq}\,
                       \varepsilon_k^{-4}u\right\rangle_k=0.
 \label{eq:joint-ward-native-zero}
\end{equation}
Indeed the probe reads no link on which \(Y_b\) acts,
since \(\varepsilon_k R_{\rm cut}=o(r_i)\). The left
side is \(\varepsilon_k^{-9}Z_{n,k}^{-1}\) times
\[
 \sum_{e,\alpha}\int
    u\,\mathcal L_{e,\alpha}
                      (F_{n,k}(Y_b)_{e,\alpha})\,\dd H_k
 =-\sum_{e,\alpha}\int
    (\mathcal L_{e,\alpha}u)F_{n,k}(Y_b)_{e,\alpha}\,\dd H_k=0.
\]
This is integration by parts in the smooth density, before
any sector decomposition. The full equation insertion
need not have a finite read. Its exact vanishing is used
at this level; only the local functions \(f,g,u\) are
subsequently pulled back as separate factors.

\emph{The original kernel and its shared tags.}
For one step write its physical density, with bond tags
summed but site tags still held, as
\[
 K_j(U,V)=\sum_{\boldsymbol t}
    \prod_x\lambda_{j,x}(U,t_x)
    \prod_e p_{j,e}(U,t_{e^-},t_{e^+};V_e).
\]
Here each \(\lambda_{j,x}\) has sum one, and
\begin{align*}
 p_{j,e}(U,t;V_e)&=\sum_s\lambda_{j,e}(U,t;s)
       z_j^{-1}\exp\{-\beta_{j+1}a_j
                            q_G(V_eQ_{j,e}(U,t;s)^{-1})\},\\
 \int_Gp_{j,e}(U,t;V_e)\,\dd h_G(V_e)&=1.
\end{align*}
The number \(z_j\) is the original one-link noise
normalizer. All eligible-tag masses and fallback choices
are the original ones. An endpoint tag in two neighbouring
bonds is the same summation variable in both factors.
Conditional on these site tags, the link integrations are
products; after unused links are integrated, unused site
tags have their original total mass one. Functions on
disjoint full site reads depend on disjoint remaining
tag variables, so their finite tag sums then factor.

Iterating this argument on the disjoint read cones gives
\eqref{eq:response-pullback-products}. Its uniform room
condition is \(2(\varepsilon_kR+T_k)<\varepsilon_m\),
where \(T_k=30\varepsilon_k(J(k)+5)\) and \(m<v\) is
the replication grid. The cone induction proves this
for every terminal cutoff \(n\ge k\), rather than
choosing a new room estimate for each \(n\). Thus,
with a hat denoting \(\mathcal P_{n\to k}^*\),
\begin{equation}
 \langle \varepsilon_k^{-5}(f-\beta_kg),
                           \varepsilon_k^{-4}u\rangle_k
 =\left\langle
   \varepsilon_k^{-5}(\widehat f-\beta_k\widehat g),
                           \varepsilon_k^{-4}\widehat u
                                                    \right\rangle_{n,M}.
 \label{eq:joint-ward-native-pullback}
\end{equation}
Both sides have the same original Wilson partition
denominator: mass preservation gives
\(Z_{n,k}=Z_{n,M}\). No independence of interacting
fields, or of adjacent coarse bonds before their site
tags are fixed, has entered this identity.

\emph{The mixed coefficient is retained explicitly.}
Put
\[
 X_0=\varepsilon_k^{-5}(f-\beta_kg),\quad
 X_1=\varepsilon_k^{-4}u,\qquad
 Y_0=\mathcal O_{n,M}(a^f-\beta_ka^g),\quad
 Y_1=\mathcal O_{n,M}(a^u).
\]
The \(X_a\) are coarse functions and the \(Y_a\) are
fine functions. For two commuting square-zero markers
\(z_0,z_1\), the full joined density is
\[
 \mathcal P_{n\to k}
       \{F_n(1-z_0Y_0)(1-z_1Y_1)\}
                         (1+z_0X_0)(1+z_1X_1).
\]
Its \(z_0z_1\) coefficient is the two-name case of
\eqref{eq:common-signed-subset-coefficient}, with
\(O_i=Y_i\), \(V_i=X_i\), and all marker signs reversed.
Writing out every term gives
\begin{equation}
 F_kX_0X_1-X_0\mathcal P_{n\to k}(F_nY_1)
       -X_1\mathcal P_{n\to k}(F_nY_0)
                         +\mathcal P_{n\to k}(F_nY_0Y_1).
 \label{eq:joint-ward-four-term-density}
\end{equation}
The last term is a joint push, including its inherited
two-name mode. It is not prescribed by two singleton
pushes. On integration, adjointness and the preceding
cone identity identify this coefficient divided by
\(Z_{n,M}\) with
\(\langle(\widehat X_0-Y_0)(\widehat X_1-Y_1)\rangle_{n,M}\).

For the actual estimate first keep distinct markers for
\(f,g,u\), apply their separate response budgets, and
only then substitute the Ward combination. Substituting
\(z_f=z_0\), \(z_g=-\beta_kz_0\) kills the mode
containing both \(f,g\), since \(z_0^2=0\); it keeps
each mixed mode with \(u\) with exactly its signed
coefficient. This substitution explains why no assertion
\(e(L^{\beta_kg})\le\widehat\beta_k\) is needed.

The same finite identity is applied in every reflected
cell before extracting the coefficient of all cell
markers. If a response footprint meets an incoming
forced member \(K_\alpha\), keep that member in the
component record. When several marked footprints meet
it, its factor occurs once, together with all those
markers and its original inventory and payload. Removing
the marks recovers the unmarked gas. Formula
\eqref{eq:joint-response-component} and the reserve-three
case of Lemma~\ref{lem:connected-bound} bound these
components simultaneously for all subsets. Thus
\eqref{eq:joint-ward-four-term-density} is an identity
of the full physical sources, with all mask restorations;
it is not inferred by multiplying replacement norms.

\emph{A rough price for the probes and matched Ward column.}
We need a price before testing the Ward field; assuming
its final coefficient bound here would be circular.
In the vector sector of degree five all nonconstant
primaries have engineering weight five and two curvature
factors. In the centered tensor sector of degree four
they have weight four and two curvature factors.
Thus their engineering diagonal \(\mathsf S_L\) is one.
The leading-map, coefficient-graph and inverse bounds
give, for a selected critical column with current
coordinate \(x_t\),
\[
 h_t^{-2}|x_t^{\rm crit}|_1
 \le\mathcal C_{J(t+1)}h_{t+1}^{-2}|x_{t+1}^{\rm crit}|_1,
                        \qquad v\le t<k.
\]
The regular and divided forced graph prices obey the
same bound after including their column constant in
\(\mathcal C_j\). This follows by inserting
\(w=d\) and factor count two in
\eqref{eq:insertion-leading-inverse} and using
\eqref{eq:insertion-graph-bound}; it does not use an
inverse reference-coordinate matrix.

For the tensor unit subtraction one can see the bounded
graph directly. The exact one-step coefficient matrix
on unit plus critical coordinates has the form
\[
 \binom{u_t}{c_t}
  =\begin{pmatrix}L_t^4&b_t\\0&A_t\end{pmatrix}
                                     \binom{u_{t+1}}{c_{t+1}}.
\]
The unit column is exact. The leading critical block is
\((h_t/h_{t+1})^2I\). Since \(h_t/h_{t+1}\le2\),
\eqref{eq:ward-separated-atom-prices} and
\eqref{eq:insertion-graph-reserves} give the harmless
bounds \(\|A_t\|\le8\), \(\|b_t\|\le1\) on the
stated tails. A centered graph \(u_t=\ell_tc_t\)
consequently satisfies
\[
 \|\ell_{t+1}\|
   \le L_t^{-4}(8\|\ell_t\|+1)
   \le\tfrac18(1+\|\ell_t\|),\qquad L_t\ge3.
\]
This bounded geometric recursion needs no additional
ratio floor. Its initial unit subtraction is
bounded by the fixed-reference one-point insertion
bound. It therefore introduces a fixed column constant,
not a factor growing with the number of scales. For
the odd vector sector the unit coefficient is identically
zero. This accounts for the unit part in the preceding
critical-column estimate.

Let
\(\mathcal C_{v,k}=\prod_{t=v}^{k-1}\mathcal C_{J(t+1)}\).
By the same pre-ratio choice as
\eqref{eq:trace-window-column-cost},
\[
 \log\mathcal C_{v,k}
 \le\sup_{v\le t<k}J(t+1)^{-2}
                    \log(\varepsilon_v/\varepsilon_k)=o(H).
\]
The birth prices are \(Ch_k^2P_k^{1/6}\) for each
unscaled source, and \(C\beta_kh_k^2P_k^{1/6}\) for
their Ward combination. Transporting the selected
columns to \(v\) thus gives dimensionless prices
\[
 q_v^{\rm Ward}\le
 C\mathcal C_{v,k}\beta_k h_v^2P_k^{1/6},\qquad
 q_v^{\rm probe}\le C\mathcal C_{v,k}h_v^2P_k^{1/6}.
\]
The neighboring coupling comparisons give
\(\beta_kh_v^2=P_v^{2+o(1)}\). Every factor on these
right sides, apart from its explicit threshold,
therefore has logarithm \(o(H)\).

Use the scalar-free extraction bound at grid \(m<v\)
with stopping level \(v\). For either column its root
cost is at most
\(C\zeta_{d,v}^{-1}\varepsilon_v^{-d}q_v
 \exp\{(\varepsilon_m/\varepsilon_v)^4Q_v^{\rm vac}+6\}\).
The normalizations of the test are displayed in the table:
\[
\begin{array}{c|c|c|c}
 &\text{physical birth}&\text{stopping factor}&\text{test factor}\\ \hline
 \text{Ward}&\varepsilon_k^{-5}(f-\beta_kg)
                                  &\varepsilon_v^{-5}&r_i^5\\
 \text{tensor probe}&\varepsilon_k^{-4}u
                                  &\varepsilon_v^{-4}&\beta_kr_i^4.
\end{array}
\]
The tensor factor is the canonical Gaussian normalization
of its two curvature inputs. Its remaining bounded change
of critical tensor basis is precisely the finite detector
normalization used to construct \(U_i(4)\). Unit terms
are the ones just treated. Hence the residual powers
are \((r_i/\varepsilon_v)^5\) and
\((r_i/\varepsilon_v)^4\), rather than powers of
\(r_i/\varepsilon_k\). Their logarithms are
\(O((\log\log\widehat\beta_v)^3)\). The grid vacuum
logarithm is \(O((\log\widehat\beta_v)^{22})\).
The inverse fixed-degree allowance \(\zeta_{d,v}^{-1}\)
also has logarithm \(o(H)\), by the same common entry
conditions used in the physical matching bound.
Together with \(\log\mathcal C_{v,k}=o(H)\), these
calculations bound the normalized probe roots and the
rough matched Ward root by \(e^{o(H)}\).

Finally Lemma~\ref{lem:physical-response-matching-root}
gives each normalized matching error the factor
\(e^{-H/8+o(H)}\). For the Ward combination it is
\(\widehat\beta_k^5e^{-H/8}\), times the stopping
factors in the table; for the unscaled probe it is
\(\widehat\beta_k^4e^{-H/8}\), followed by its external
\(\beta_kr_i^4\). All these extra logarithms are
\(o(H)\). In replacing the joint test use the exact
identity
\[
 \widehat X_0\widehat X_1-Y_0Y_1
       =(\widehat X_0-Y_0)\widehat X_1
                                  +Y_0(\widehat X_1-Y_1).
\]
The original Wilson chessboard estimate on the disjoint
interior cells, with the roots just proved, pays the
other factor by \(e^{o(H)}\). It is applicable to
these actual functions because their replicated root
comparisons came from the full joined polynomial above.
The rough bound used here is a bound on these particular
physical columns; it assumes no uniform bound for
\(s_i\) or for its image of a general raw coefficient.

For a spatial or restoration error replace the distinguished
Ward mark in this same calculation by its error numerator.
Its income is respectively \(e^{-2D_k}\) or
\(e^{-P_k^{1/5}/4}\), as proved above. The other marks,
incoming families and one Gaussian precision reserve are
unchanged. Extract the joint coefficient only after their
component sums. The numerator and denominator have the
same vacuum scalar and the same \(Z_{n,M}\); scalar
cancellation is performed before the grid extraction.
The following scale ledger then gives
\eqref{eq:trace-complete-ward-test-price}. All estimates
are uniform on the fixed-\(k\) cutoff tail and its
nonwrapping physical-side tail. Repetition adds no
number-of-cells constant to these local prices.

\emph{Price at the radius level.}
Put \(v=n_i\), \(k=k_i\),
\(H=H_v=(\log\widehat\beta_v)^{48}\).
The first-crossing definition gives
\[
       H\le\log(\varepsilon_v/\varepsilon_k)
              \le H+O(\log\log\widehat\beta_v).
\]
The local read radius at \(k\) is a fixed logarithmic
power, so \(\varepsilon_kR_{\rm cut}=o(\varepsilon_v)\).
In particular an exterior native probe at a positive
multiple of \(r_i\) reads no link displaced by \(Y\).
The Haar identity applies before approximation.

The signed matched reference coefficient of the two
physical sources is exactly
\[
                     a^f_{n,M,k}-\beta_ka^g_{n,M,k}.
\]
Their normalised current difference is
\(\Sigma_{n,M,k}-\beta_kj_{n,M,k}\).
Linearity is used here before taking any norm.
Lemma~\ref{lem:trace-physical-response-data} gives
the actual responses: \(e_f=0\) and
\(e_g\le x=\widehat\beta_k\), with the coefficient
budgets \(q_f+q_g\le Ch_k^2P_k^{1/6}\).
Apply Lemma~\ref{lem:physical-response-matching-root}
separately to \(f\) and \(g\), with their already
chosen coefficients \(a^f\) and \(a^g\).
Each root price uses the broader budget \(x^3\)
and radius reduction \(x^{-1}\), hence has the
factor \(x^4\) in
\eqref{eq:physical-response-root-price}.
The Wilson-root triangle inequality then bounds
the combination by \((1+\beta_k)\) times that
price, which is at most a fixed constant times
\(x^5\), since \(\beta_k\le2x\).
We have not claimed that the response of
\(\beta_kg\) itself has price at most \(x\).
The sharp coefficient comparisons still use
their original \(q_f,q_g\).

Use stopping level \(v\) and the centred grid
level \(m<v\) specified below. Condition
\eqref{eq:response-native-room} holds because
\(\varepsilon_k(R_{\rm cut}+30(J(k)+5))
=o(\varepsilon_v)\). Thus the estimate is for
the product of the actual pullback differences,
with all higher source modes retained.
The inverse-scale gain is \(e^{-H/8}\).
The new factor \(x^5\), the inverse allowance
\(\zeta_{5,v}^{-1}\), and all column losses
have logarithm \(o(H)\), by
\eqref{eq:trace-window-column-cost} and
\(\log x=O(\log\widehat\beta_v)=o(H)\).

For completeness the denominator cost is also local.
For replication take the first backward crossing \(m<v\)
of \(B_v^2\). If \(L\) is the maximum ratio there, then
\[
 B_v^2\le\varepsilon_m/\varepsilon_v\le LB_v^2,
 \qquad
  (\varepsilon_m/\varepsilon_v)^4Q_v^{\rm vac}+6
                  \le L^4B_v^8Q_v^{\rm vac}+6
                              =O((\log\widehat\beta_v)^{22}).
\]
Here \(L\le\log\widehat\beta_v\) and the exponent of the
scalar-free ratio satisfies
\(Q_v^{\rm vac}=O((\log\widehat\beta_v)^2)\), by the
common entry choices. Inverse fixed-degree attachment
allowances cost a fixed coupling power. All these
logarithmic costs are \(o(H)\). Coarsening and scalar
cancellation precede this estimate; charging a separate
fine-site vacuum factor in a radius-level cell would
give a different, unusable bound.

The factor \(\varepsilon_v^{-5}\) in the transported
physical source is cancelled by \(\varepsilon_v^5\) in
the tested Ward coefficient. The remaining change to
\(r_i\) has cost
\[
 (r_i/\varepsilon_v)^5
       \le\exp\{5(\log\log\widehat\beta_v)^3\}.
\]
For a spatial or restoration error still priced at
level \(k\), the larger factor
\((r_i/\varepsilon_k)^5\) has logarithm at most
\(5H+o(H)\). It is paid by the longer truncation above.
Consequently every exterior normalised Ward test is
bounded by
\begin{equation}
 \chi_i=C\left\{
       e^{-H/100+o(H)}
       +e^{10H-D_k/16}
       +e^{10H-P_k^{1/5}/8}\right\}.
 \label{eq:trace-complete-ward-test-price}
\end{equation}
Each term is smaller than \(b_i^{-p}\) for every fixed
\(p\). Indeed \(\widehat\beta_k/\widehat\beta_v\to1\),
\(D_k\asymp(\log\widehat\beta_v)^{80}\), and
\(P_k^{1/5}\) eventually exceeds every fixed power of
\(\log\widehat\beta_v\).

The order of limits in this calculation is important.
At each fixed \(k\), first take a terminal-cutoff tail
for the current approximations. Then take the physical
side large enough for the finitely many lifted reads,
coefficient comparisons and cells. These latter thresholds
are uniform on the terminal tail. The estimates pass
first to the fixed-side cutoff limits and then to the
chosen side subsequence. Thereafter \(i\) tends to
infinity. The current arrays are refined along the
already fixed hierarchy; no correlation limit is changed.

\emph{From tested Ward bounds to coefficient norm.}
In the ordered limit the selected Ward field is
\(\beta_i\mathsf C_{k_i,5}(h_{k_i}^2v_i^{\rm eq})\).
Thus its \(r_i^5U_i(5)^{-1}\) image is exactly
\(s_iv_i^{\rm eq}\); the factor \(\beta_i\) is accounted
for by the definition of \(s_i\).

Degree-four tensor spectators separate the free vector
plane. To verify this, suppose a vector in that plane
has zero correlations with all degree-four tensors at
every separated point. The plane is spanned by
divergences of such tensors. Differentiate in the
spectator position and use their span: the vector has
zero two-point response with itself. Faithfulness of
the free hierarchy makes it zero. Since the plane is
two dimensional, choose two independent scalar rows
from these tensor tests. Their matrix on a fixed basis
of the free vector plane is invertible.

Use the corresponding native physical probes in the
above Ward test, and match them on the same window.
Their reference-coordinate limits are the fields
\(U_i(4)\) applied to the fixed spectator labels. Their
matching errors have the same inverse-scale root price
and were included in \(\chi_i\). The matrices of the
resulting tests on \(U_i(5)\) applied to the fixed vector
basis converge to the invertible free matrix, by
\eqref{eq:trace-prescribed-comparison}. Their inverses
are therefore bounded. Testing the possibly varying
vector \(s_iv_i^{\rm eq}\) gives
\[
                         \|s_iv_i^{\rm eq}\|\le C\chi_i.
\]
Only the matrix entries on a fixed basis were compared;
no bound on \(\|s_i\|\) has been assumed. This proves
\eqref{eq:trace-complete-ward-image} without discarding
the matched coordinate created by a zero-read remainder.

\emph{A quantitative derivative identity.}
We next prove \eqref{eq:trace-complete-derivative-lift}.
At every \(k\), match the complete dimension-four
current and its complete improved central differences
independently. On each fixed exterior response row,
matching stopped at a fixed late tiling level gives
errors bounded by \(E_k=O(\varepsilon_k^\eta)\), for
some \(\eta>0\), uniformly for a small interval of
anchor translations. This follows directly from
\eqref{eq:physical-response-root-price}, applied to
the finite tensor and difference sources covered by
Lemma~\ref{lem:trace-physical-response-data}: the
extra fixed coupling power, and all current,
section and reference-coordinate factors, are subpower
in \(\varepsilon_k^{-1}\). Decrease \(\eta\) once
to include their finite products. A complete source
identity, or any finite sum of them, has the same bound.

Here is the calculation that avoids differentiating
such an error. Let \(H_k(x)\) be the response of the
matched tensor field as its anchor moves in one axis,
and let \(G_k(x)\) be the response of the independently
matched central-difference field. Write \(F_k(l)\) for
the corresponding finite-cutoff physical response at
anchor \(l\varepsilon_k\). At the physical level the
normalised central difference is exactly
\[
 \varepsilon_k^{-1}\sum_{j=1}^4c_j
                         \{F_k(l+j)-F_k(l-j)\}.
\]
The two separately matched responses approximate these
two physical expressions with errors \(E_k\) and
\(E_k'\), both \(O(\varepsilon_k^\eta)\).
For a fixed exterior row, smoothness gives
\[
 \left|H_k'(x)-\varepsilon_k^{-1}\sum_{j=1}^4c_j
               \{H_k(x+j\varepsilon_k)-H_k(x-j\varepsilon_k)\}
       \right|
        \le C\varepsilon_k^8\|B_k^{\rm phys}\|.
\]
In this paragraph \(B_k^{\rm phys}\) denotes the matched
tensor field, not the logarithmic read radius \(B_k\).
The four moment cancellations of
\eqref{eq:trace-central-difference} leave the ninth
derivative in Taylor's integral remainder. Its row
norm is bounded by a fixed constant times the field
norm on this fixed finite-dimensional space. The
reference-coordinate bound gives
\(\|B_k^{\rm phys}\|\le\varepsilon_k^{-4-o(1)}\),
which is more than sufficient for this eighth-order
gain.

Average over \(N_k=\lceil\varepsilon_k^{-1+\eta/2}\rceil\)
adjacent sites, decreasing \(\eta<1\) if necessary.
For each \(j\), the sum of the central difference of
the response error has only \(4j\) endpoint terms.
Thus
\begin{align}
 \left|\frac{1}{N_k}\sum_{l=0}^{N_k-1}
       \{H_k'(l\varepsilon_k)-G_k(l\varepsilon_k)\}\right|
 \le{}&C\varepsilon_k^8\|B_k^{\rm phys}\|
       +\frac{4\sum_{j=1}^4j|c_j|}{N_k\varepsilon_k}E_k
       +E_k'\notag\\
   ={}&O(\varepsilon_k^{\eta/2}).
 \label{eq:trace-averaged-derivative-price}
\end{align}
The averaging length is
\(N_k\varepsilon_k=O(\varepsilon_k^{\eta/2})\to0\).
For each fixed vector in the field space its averaged
response tends to its response at the base point.
Choose a finite matrix of separating rows at that
point. The averaged row matrices converge entrywise
to this fixed invertible matrix, so their left
inverses stay bounded. Applying them to the varying
field difference in
\eqref{eq:trace-averaged-derivative-price} proves a
field-norm error \(O(\varepsilon_k^{\eta/2})\).
This argument does not require an a priori bound
on the difference being detected.

Summing axes gives the divergence identity. The
finite physical sum of the component differences
is exactly \(g\). Comparing its independently
matched column with the sum of the component columns
adds another \(O(\varepsilon_k^\eta)\) error, by
linearity and the same root estimate. Symmetrising
the tensor introduces its antisymmetric part; this
part is a difference of degree-three potentials by
\eqref{eq:trace-current-identities}. Its matched
field is small by the next calculation. Derivation
is a fixed linear map on these finite field spaces,
so its derivative is small too. Real \(H_4\) averaging
commutes with the summed continuum divergence and
does not change any of these orders.

To convert the field error to a raw-coordinate error,
fix a real equivariant right section of the reference
projection onto the vector plane. Compose it with
\(h_k^{-2}\mathsf Z_{k,5}^{-1}\). This is a right
inverse of \(v\mapsto\mathsf C_{k,5}(h_k^2v)\).
Its norm is \(e^{o(s_k)}\), by
\eqref{eq:position-inverse-subpower}; the extra factor
\(h_k^{-2}\) is also subpower. Hence the field error
has a lift \(y_k\) of norm
\(O(\varepsilon_k^{\eta/2-o(1)})\).
The definitions of \(K_i,s_i\) give the exact identities
in \eqref{eq:trace-complete-derivative-lift}, and every
fixed coupling power times \(y_{k_i}\) tends to zero.
This step has not yet used boundedness of \(s_i\).

\emph{Forward differences and tensor traces.}
Apply Lemma~\ref{lem:degree-three-potential-birth}
and Corollary~\ref{cor:degree-three-forward-difference}
to every component of \(q_a,t_{abc}\).
These are separate physical matchings in degrees
three and four, with their actual response families,
the same partition denominator, and all higher
source modes retained. The degree-three matched
coefficient may be a nonzero unit. It cancels
exactly in the translated endpoint difference;
\eqref{eq:degree-three-forward-difference-decay}
therefore gives the degree-four field bound
\(O(\varepsilon_k^{\eta/2})\), before changing
to the moving Maxwell frame. This argument
divides only endpoint errors by the averaging
length, not a numerical matching error by a
single lattice spacing. Apply this bound to
the complete trace and antisymmetric identities
of \eqref{eq:trace-current-identities}. The \(\Xi\)
trace is identically zero at the physical level;
comparison with its full matched column gives the
same rate. Centering and symmetry averaging are
fixed bounded linear maps.

Finally, in degree four there is no nonconstant
equation direction. The selected critical columns
therefore give a uniformly conditioned frame after
prescribed-radius normalisation. The local current
read tends to \(L\), and the ordinary reads are
\(\Xi,D\); \eqref{eq:ward-normalised-matching-bias}
and the short-window drift
\eqref{eq:ward-critical-column-drift} identify their
complete matched tensors with the same frame.
It follows that \(\mathsf A_iB_i^B\to B\).
The transported derivative in that stable frame
converges by
Lemma~\ref{lem:trace-prescribed-radii}; hence also
\(K_iB_i^B\to\operatorname{div}^{\rm Max}B\) for
\(B=\Xi,D\).

The degree-four coordinate conversion of a field
of norm \(O(\varepsilon_k^{\eta/2})\) costs only
\(e^{o(s_k)}\); its subsequent bounded frame
conversion preserves decay faster than every
coupling power. Applying this to the trace errors
proves \eqref{eq:trace-complete-tensor-traces}.
All estimates were made before taking a character
of the ground-space algebra. If a phase bound is
requested, the repeated reflected-source estimates
bound the corresponding multiplication operator,
and therefore every character, with the same
constant. The hierarchy and its selected vacuum
remain the ones already fixed.
\end{proof}

\begin{lemma}[Fourth-order Ward transfer]
\label{lem:trace-ward-transfer}
There are bounded real functionals \(\lambda_i\) on the
free tensorial vector space such that
\begin{align}
 \|\beta_iK_iB_i^L\|&\le CR_i,&
 \lambda_i(\beta_iK_iB_i^L)&\longrightarrow-b_0,\label{eq:trace-transfer-divergence}\\
 K_iB_i^B&\longrightarrow\operatorname{div}^{\rm Max}B,&
 E_i\{\lambda_i(K_iB_i^B)
                    -\rho_{\rm tr}(\operatorname{div}_{\rm raw}B)\}
                      &\longrightarrow0,\qquad B=\Xi,D.\notag
\end{align}
Moreover
\begin{equation}
 \mathsf A_iB_i^B\to B\quad(B=L,\Xi,D),\qquad
 b_i^p\operatorname{tr}(\mathsf A_iB_i^B)\to0
                              \quad(B=L,\Xi;\ p\in\N).
 \label{eq:trace-transfer-tensor}
\end{equation}
The constants depend on the fixed construction and finite
degree request, and not on a remaining cutoff, physical
period or phase.
\end{lemma}
\begin{proof}
All coefficients in this proof are complete-source matched
coordinates. Equations
\eqref{eq:trace-matched-current-bounds} and
\eqref{eq:trace-ordinary-matched-bias} give
\begin{equation}
 j_{k_i}\to-v_{\rm eq},\qquad
 |\rho_{\rm tr}(j_{k_i})|\le CP_{k_i}^{-1/4},\qquad
 \|w_{k_i}^B-v_B\|\le CP_{k_i}^{-1/4},\quad B=\Xi,D.
 \label{eq:trace-current-equation-coefficient}
\end{equation}
With \(v_i^{\rm eq}=\beta_i^{-1}\Sigma_{k_i}-j_{k_i}\)
one obtains
\[
 v_i^{\rm eq}\to v_{\rm eq},\qquad
 E_i|\rho_{\rm tr}(v_i^{\rm eq})|
 \le CE_i/\beta_i+CP_{k_i}^{1/5-1/4}\longrightarrow0.
\]
Here \(\rho_{\rm tr}(\Sigma_k)\) is bounded by
\eqref{eq:trace-sigma-bounds}; the full coefficient norm
may grow as \(R_i\).

We first prove boundedness of \(s_i\), without assuming
that a small field error stays small under this map.
The raw basis
\(w_{k_i}^{\Xi},w_{k_i}^D,v_i^{\rm eq}\) tends to
\(v_{\Xi},v_D,v_{\rm eq}\) and hence has uniformly
bounded inverse matrices. By
Lemma~\ref{lem:complete-source-ward-transport},
\[
 \|s_i\|\le C\{\|K_iB_i^\Xi\|+\|K_iB_i^D\|
       +\|s_iv_i^{\rm eq}\|
       +\|s_i\|(\|y_i^\Xi\|+\|y_i^D\|)\}.
\]
The first two terms are bounded by the stable-frame
comparison, the third tends to zero, and the coefficient
of \(\|s_i\|\) on the right is eventually at most
\(1/2\). Absorption proves \(\sup_i\|s_i\|<\infty\).
The same basis and its three image limits then give
\(s_i\to s\).

The complete-source identities now imply, for every
fixed \(p\),
\begin{align}
 \beta_iK_iB_i^L-s_i\Sigma_{k_i}&=o(b_i^{-p}),\notag\\
 K_iB_i^B-s_iw_{k_i}^B&=o(b_i^{-p}),
                                      \qquad B=\Xi,D.
 \label{eq:trace-transport-identification}
\end{align}
Indeed the first difference is exactly
\(-\beta_i s_iv_i^{\rm eq}+\beta_i s_iy_i^L\), and the
second is \(s_iy_i^B\). The coefficient correction
\(w_{k_i}^B-v_B\) is only \(O(P_{k_i}^{-1/4})\);
\eqref{eq:trace-transport-identification} does not erase it.

The columns \(s_iw_{k_i}^{\Xi},s_iw_{k_i}^D\) converge
to the independent vectors \(sv_{\Xi},sv_D\).
Let \(t_i\) be their inverse coordinate map followed by
\((c_0,c_1)\mapsto c_0w_{k_i}^{\Xi}+c_1w_{k_i}^D\).
It is a bounded right inverse of \(s_i\) onto the vector
plane. Define \(\lambda_i=\rho_{\rm tr}t_i\).
Then \(\lambda_is_i-\rho_{\rm tr}\) is exactly zero on
the first two vectors of our raw basis. On its third
vector it has absolute value at most
\(C\|s_iv_i^{\rm eq}\|+|\rho_{\rm tr}(v_i^{\rm eq})|\).
The uniformly bounded inverse basis therefore gives
\begin{equation}
 \|\lambda_is_i-\rho_{\rm tr}\|
 \le C\{\|s_iv_i^{\rm eq}\|
                       +|\rho_{\rm tr}(v_i^{\rm eq})|\}
                                           =o(E_i^{-1}).
 \label{eq:trace-functional-precision}
\end{equation}
In particular \(\lambda_i\) is bounded. Applying
\eqref{eq:trace-functional-precision} to
\(\|\Sigma_{k_i}\|\le CR_i\), and using
\(R_i/E_i=P_{k_i}^{-1/30}\to0\), yields
\[
 \lambda_i(\beta_iK_iB_i^L)
      =\rho_{\rm tr}(\Sigma_{k_i})+o(1)\longrightarrow-b_0,
 \qquad \|\beta_iK_iB_i^L\|\le CR_i.
\]
For \(B=\Xi,D\), evaluate the second line of
\eqref{eq:trace-transport-identification} instead:
\[
 \lambda_i(K_iB_i^B)-\rho_{\rm tr}(v_B)
       =\rho_{\rm tr}(w_{k_i}^B-v_B)+o(b_i^{-p}).
\]
Multiplication by \(E_i\) gives a bound
\(CP_{k_i}^{-1/20}+o(1)\to0\).
This proves \eqref{eq:trace-transfer-divergence}.
The tensor assertions are exactly the final assertions
of Lemma~\ref{lem:complete-source-ward-transport}.
All coefficient limits, source estimates and finite
matrices refer to the one Wilson family and its already
selected ordered hierarchy. Their bounds precede phase
evaluation and apply to the same selected vacuum.
\end{proof}
% !TeX root = 4DYM_notes.tex
\section{The first nonzero trace coefficient and the full-radius limit}
\label{sec:trace-short-distance}

We first use conservation to identify the scalar coefficient.
This avoids constructing a separate renormalisation scheme
for the trace. The argument uses the unique centred conserved
line and the already normalised translation charge.
It follows the derivation after
\cite[Lemma~18.5, PDF pp.~213--214]{ArmstrongDraft}.

\begin{proposition}[Trace in the moving Maxwell coordinates]
\label{prop:trace-leading-coefficient}
On each prescribed-radius subsequence constructed above,
\begin{equation}
       \beta_ir_i^4U_i(4)^{-1}\Theta
                                     \longrightarrow-b_0\mathcal S.
 \label{eq:trace-leading-coefficient}
\end{equation}
This is convergence in a fixed coefficient norm on
\(V_4^{\rm Max}\). The comparison of the associated
two-point functions is uniform in phase before the
selected phase is evaluated.
\end{proposition}
\begin{proof}
The free divergences of \(\Xi,D\) are independent.
Otherwise a nonzero combination would be in the conserved
line spanned by \(L\), contrary to their tensor coordinates.
The maps
\[
 f_i(c_0,c_1)=c_0K_iB_i^\Xi+c_1K_iB_i^D
\]
are therefore isomorphisms onto the actual divergence plane,
with uniformly bounded inverses. This plane has dimension
two by the interacting rank calculation of
Section~\ref{sec:interacting-stress-charges}.
Set
\[
 (c_{i,0},c_{i,1})=-f_i^{-1}(\beta_iK_iB_i^L),
 \qquad
 X_i=\beta_iB_i^L+c_{i,0}B_i^\Xi+c_{i,1}B_i^D.
\]
Then \(|c_{i,0}|+|c_{i,1}|\le CR_i\) and \(X_i\) is
exactly conserved. Applying \(\lambda_i\) to its divergence
and using \(\rho_{\rm tr}(\operatorname{div}_{\rm raw}D)=-4\)
gives
\begin{align*}
 4c_{i,1}
   ={}&\lambda_i(\beta_iK_iB_i^L)
      +c_{i,0}\lambda_i(K_iB_i^\Xi)\\
      &+c_{i,1}\{\lambda_i(K_iB_i^D)+4\}.
\end{align*}
The last two terms tend to zero because their coefficients
are \(O(R_i)\), their functional errors are \(o(E_i^{-1})\),
and \(R_i/E_i\to0\). Thus
\begin{equation}
                         c_{i,1}\longrightarrow-b_0/4.
 \label{eq:trace-scalar-correction}
\end{equation}
No bound or convergence for \(c_{i,0}\) stronger than
\(O(R_i)\) was required.

Since \(R_i/\beta_i\to0\),
\eqref{eq:trace-transfer-tensor} gives
\[
        r_i^4U_i(4)^{-1}X_i\longrightarrow L.
\]
After taking the trace and multiplying by \(\beta_i\),
the \(L,\Xi\) terms tend to zero by their stronger trace
rates. More explicitly
\[
 \beta_i\operatorname{tr}(r_i^4U_i(4)^{-1}X_i)
    =\beta_i\operatorname{tr}(\mathsf A_iB_i^L)
       +c_{i,0}\operatorname{tr}(\mathsf A_iB_i^\Xi)
       +c_{i,1}\operatorname{tr}(\mathsf A_iB_i^D)
             \longrightarrow-b_0\mathcal S.
\]
The last trace tends to \(4\mathcal S\), and
\eqref{eq:trace-scalar-correction} supplies its coefficient.

The tensors \(X_i\) are centred in the mixture, so uniqueness
of the conserved line gives \(X_i=\nu_iT\) for real \(\nu_i\).
We identify the factor using a local charge, not by a
choice of tensor coordinate. Choose a free label \(A\),
an axis \(b\) and a separating spectator functional \(o\)
with \(y=o(\partial_b^{\rm Max}A)\ne0\). Such a choice
exists by the explicit free Ward calculation.
For a fixed translation shell, smooth prescribed-radius
comparison makes the rescaled flux of \(X_i\) tend to the
flux of \(L\), namely \(y\). The exact translation Ward
identity for \(T\) makes its corresponding rescaled flux
tend to the same \(y\): its right-hand side is the
transported derivative of \(A\). The first flux is
\(\nu_i\) times the second. Since the denominator tends
to a nonzero number, \(\nu_i\to1\).
The trace limit just obtained therefore holds for \(T\).

Finally \(U_i(4)^{-1}1_4=1_4\) exactly and
\(\beta_ir_i^4|e_{\rm tr}|\to0\), by
\eqref{eq:trace-prescribed-coupling}. Subtracting the
selected vacuum mean proves
\eqref{eq:trace-leading-coefficient}.
\end{proof}

\subsection{The free normalisation}

We compute the remaining numerical factor directly.
At \(z=re_4\), the covariance
\eqref{eq:position-free-curvature-covariance} is diagonal
on the six curvature planes \(a<b\). Its three electric
entries \(F_{a4}\) are \(-1/(\pi^2r^4)\), and its three
spatial entries are \(1/(\pi^2r^4)\); different colours
have zero covariance. Since
\(\mathcal S=\frac12\sum_{\alpha,a<b}:F_{ab}^\alpha{}^2:\),
the two cross-pairings give
\begin{equation}
 S_2^{\rm Max}(\mathcal S,\mathcal S;0,z)
   =\frac14\cdot2\cdot6q_G\,\frac1{\pi^4r^8}
   =\frac{3q_G}{\pi^4r^8}.
 \label{eq:trace-free-scalar-normalisation}
\end{equation}
Rotation covariance extends this to every \(z\ne0\).
The signs of the electric entries are squared; no
pointwise positivity of an individual curvature
component was used.

\begin{theorem}[Trace short-distance asymptotic]
\label{thm:trace-short-distance}
For the centred trace of the selected theory,
\begin{equation}
      \boxed{\displaystyle
        \lim_{r\downarrow0}r^8\log^2(1/r)\tau(r)
                                     =\frac{3q_G}{\pi^4}>0 .}
 \label{eq:trace-full-radius-asymptotic}
\end{equation}
Equivalently,
\(\tau(r)=\frac{3q_G}{\pi^4r^8\log^2(1/r)}(1+o(1))\).
The field \(\Theta\) and its reconstructed field are nonzero.
All quantities are those of the single Wilson family and
the selected vacuum already constructed.
\end{theorem}
\begin{proof}
Start with an arbitrary positive null sequence. Extract the
subsequence and frames of
Lemma~\ref{lem:trace-prescribed-radii}. Write
\[
 A_i=\beta_ir_i^4U_i(4)^{-1}\Theta,\qquad
                         A_i\to-b_0\mathcal S.
\]
Multilinearity gives, at a generic pair,
\begin{align*}
 &\beta_i^2r_i^8
    S_{*,2}(\Theta,\Theta;r_ix_0^{(i)},r_ix_1^{(i)})\\
 &\hspace{18mm}
       =S_{*,2}(U_i(4)A_i,U_i(4)A_i;
                                     r_ix_0^{(i)},r_ix_1^{(i)})
       \longrightarrow
                   b_0^2 S_2^{\rm Max}(\mathcal S,\mathcal S;x_0,x_1).
\end{align*}
Here the moving coefficients are bounded, and the phase-uniform
two-point comparison applies before evaluation at \(\sigma_*\).
Since \(\log(1/r_i)/\beta_i\to b_0^{-1}\), the factors
\(b_0^2\) cancel exactly. Translation and \(O(4)\) covariance
identify the left kernel with
\(\tau(r_i|x_1^{(i)}-x_0^{(i)}|)\). Consequently
\begin{equation}
 r_i^8\log^2(1/r_i)
      \tau(r_i|x_1^{(i)}-x_0^{(i)}|)
                      \longrightarrow
                \frac{3q_G}{\pi^4|x_1-x_0|^8}.
 \label{eq:trace-rounded-pair-limit}
\end{equation}

It remains to remove the rounded distance; it is not legitimate
to do this by an unproved derivative bound for a growing
normalisation of \(\Theta\).
Take \(x_0=0\) and
\(x_1=(t/2,t/2,t/2,t/2)\), \(1/2<t<2\).
This is a generic pair of distance \(t\).
For fixed \(t>1\), its rounded distance is eventually
greater than one. Monotonicity from
Lemma~\ref{lem:trace-spectral-measure} then bounds the lower
limit of \(r_i^8\log^2(1/r_i)\tau(r_i)\) from below by
\(3q_G/(\pi^4t^8)\). For fixed \(t<1\), it bounds the
upper limit from above by the same expression.
Let \(t\) approach one from either side. The limit along
this subsequence is \(3q_G/\pi^4\).
Every positive null sequence has such a subsequence. If
the full limit failed, a null sequence staying a fixed
distance from this value would contradict what was just
proved. This proves \eqref{eq:trace-full-radius-asymptotic}.

The nonabelian simple Lie algebra has \(q_G>0\).
A zero label, or an identically zero reconstructed field
on the word domain, would have zero two-point function.
Faithfulness of the reconstruction and the positive
asymptotic exclude both possibilities.
\end{proof}

The finite-cutoff estimates used here have local constants
fixed by the group, derivative request and separation compact;
they are uniform in cutoff and physical period on ordered
tails. First take the ultraviolet limit at each fixed
physical period, then the fixed volume subsequence. The
arbitrary-radius argument is performed on the resulting
kernel. It asserts no exchange of a shrinking radius with
either physical limit, and no volume-dependent renormalisation
of the stress tensor.
% !TeX root = 4DYM_notes.tex
\section{The finite Wick-profile non-freeness test}
\label{sec:trace-nonfreeness}

We finish with the precise comparison class in the construction
theorem. Its exclusion is an elementary consequence of the
inverse logarithm in
\eqref{eq:trace-full-radius-asymptotic}. We prove the needed
expansion rather than assuming that a massive free kernel
has a convergent expansion at zero.
The reduction follows
\cite[Lemma~18.7, PDF p.~215]{ArmstrongDraft}.

\begin{definition}[Elementary kernel and finite Wick profile]
An elementary kernel is specified by \(m\ge0\),
\(s,t\in\N\), and a real continuous function
\(h:\R^3\to\R\) homogeneous of degree \(s\):
\(h(\lambda p)=\lambda^sh(p)\) for \(\lambda>0\).
It is the function, for \(r>0\),
\begin{equation}
 k(r)=\int_{\R^3}
       \frac{E(p)^t h(p)}{2E(p)}e^{-rE(p)}\,\dd p,
                     \qquad E(p)=\sqrt{m^2+|p|^2}.
 \label{eq:trace-elementary-kernel}
\end{equation}
At the single point \(E=0\) the integrand may be defined
arbitrarily. A finite Wick profile is a finite complex linear
combination of finite products of elementary kernels.
The empty product is one.
\end{definition}

This definition allows continuous homogeneous numerators,
not only polynomials. It includes the mass-shell factors
arising from finitely many free masses, finite differential
orders and finite Wick polynomials. The term \emph{finite}
is essential: an arbitrary positive superposition of masses
is not in this comparison class.

\begin{lemma}[Finite Laurent--logarithmic expansions]
\label{lem:trace-wick-expansion}
Every elementary kernel, and every finite Wick profile
\(W\), has an expansion to each finite order in integer
powers of \(r\) and nonnegative integer powers of \(\log r\).
In particular there are finitely many complex polynomials
\(P_a\) and an integer \(a_{\min}\le0\) such that
\begin{equation}
             W(r)=\sum_{a=a_{\min}}^0
                         r^aP_a(\log(1/r))+O(1).
 \label{eq:trace-wick-singular-part}
\end{equation}
The constants may depend on the fixed profile, including its
masses and numerators. No uniformity over the comparison
class is asserted or needed.
\end{lemma}
\begin{proof}
The restriction of \(h\) to the unit sphere is bounded.
In polar coordinates, followed by \(E=\sqrt{m^2+\rho^2}\),
\begin{equation}
 k(r)=c_h\int_m^\infty
               E^t(E^2-m^2)^{(s+1)/2}e^{-rE}\,\dd E,
          \qquad c_h=\tfrac12\int_{S^2}h(\omega)\,\dd\omega .
 \label{eq:trace-radial-kernel}
\end{equation}
Indeed
\(\rho^{s+2}\dd\rho/E=\rho^{s+1}\dd E\).
This also proves absolute convergence for every \(r>0\).
If \(m=0\), it gives exactly
\[
                  k(r)=c_h(s+t+1)!\,r^{-(s+t+2)}.
\]
Thus we may take \(m>0\).

The integral over \([m,2m]\) is analytic in \(r\), with
Taylor remainders of every finite order bounded by the
integral of the usual exponential remainder. On \(E\ge2m\),
Taylor's formula for \((1-u)^{(s+1)/2}\), \(0\le u\le1/4\),
gives, for every \(J\),
\[
 E^t(E^2-m^2)^{(s+1)/2}
     =\sum_{j=0}^Jb_jE^{s+t+1-2j}
                       +O_{m,s,t,J}(E^{s+t-1-2J}).
\]
Fix a requested remainder order \({\mathsf M}\ge0\).
Choose \(J\) so large that the last exponent is less
than \(-{\mathsf M}-1\). The remainder then has an integrable
\({\mathsf M}\)-th moment. Its Laplace transform has a Taylor
polynomial through degree \({\mathsf M}-1\), with error \(O(r^{\mathsf M})\):
use
\[
 \left|e^{-rE}-\sum_{j=0}^{{\mathsf M}-1}\frac{(-rE)^j}{j!}\right|
                         \le C_{\mathsf M}(rE)^{\mathsf M},\qquad rE\ge0,
\]
and integrate. For \({\mathsf M}=0\) boundedness is sufficient.
We are left with finitely many integrals
\[
                         I_a(r)=\int_{2m}^\infty E^ae^{-rE}\,\dd E,
                                      \qquad a\in\Z .
\]
For \(a\ge0\), repeated integration by parts yields
\[
 I_a(r)=e^{-2mr}\sum_{j=0}^a
                       \frac{a!}{j!}(2m)^j r^{-(a-j+1)}.
\]
For \(a=-1\),
\(I_{-1}'(r)=-e^{-2mr}/r\). Since
\((e^{-2mr}-1)/r\) is analytic at zero, integration gives
\(I_{-1}(r)=-\log r+A(r)\), where \(A\) is analytic near
zero. For \(a<-1\), use \(I_a'=-I_{a+1}\) and integrate
successively. Integration of \(r^j\log r\) gives
\[
       \frac{r^{j+1}}{j+1}
                    \left(\log r-\frac1{j+1}\right),\qquad j\ge0.
\]
Thus every \(I_a\) has the asserted expansions. Only
nonnegative logarithmic powers occur.

Finite products preserve expansions to any fixed order.
To justify this with remainders, first note from
\eqref{eq:trace-radial-kernel} that an elementary factor
is \(O(r^{-s-t-2})\). Choose the expansion order of each
factor above the sum of the other factors' pole orders
plus the desired final order, with one further power
to absorb any finite logarithmic powers. In the telescoping
product difference every error is then of the desired
order. Finite sums preserve that property.
Expand far enough that the resulting error is \(O(1)\);
all positive-power logarithmic monomials are bounded.
Collect the remaining monomials by their integer power
to obtain \eqref{eq:trace-wick-singular-part}.
\end{proof}

\begin{lemma}[No positive inverse-logarithmic leading term]
\label{lem:trace-wick-obstruction}
If \(W\) is a finite Wick profile and the limit of
\(r^8\log^2(1/r)W(r)\) exists as a finite complex number,
that limit is zero.
\end{lemma}
\begin{proof}
Combine coefficients of equal powers in
\eqref{eq:trace-wick-singular-part}. If some nonzero
\(P_a\) has \(a\le-8\), take the least such nonzero
power among all terms. Its highest logarithmic coefficient
is nonzero. After multiplication by \(r^8\log^2(1/r)\),
its modulus tends to infinity: an exponential in
\(\log(1/r)\) dominates all polynomials when \(a<-8\),
and a polynomial of degree at least two remains when
\(a=-8\). Larger powers cannot cancel this term.
This contradicts existence of a finite limit.
All remaining powers have \(a>-8\), so their contributions
tend to zero, as does the bounded remainder.
\end{proof}

\begin{theorem}[The stated trace non-freeness test]
\label{thm:trace-nonfreeness}
For every finite Wick profile \(W\) and every \(\delta>0\),
there exists \(r\in(0,\delta)\) such that
\(\tau(r)\ne W(r)\). In particular the centred stress trace
of the selected theory cannot have any such profile as its
two-point function on a punctured neighbourhood of zero.
\end{theorem}
\begin{proof}
If equality held on such a neighbourhood, its scaled limit
would be \(3q_G/\pi^4>0\), by
Theorem~\ref{thm:trace-short-distance}. This contradicts
Lemma~\ref{lem:trace-wick-obstruction}.
\end{proof}

Together with
Lemma~\ref{lem:centred-trace-identification}, this proves
the trace test stated in Section~\ref{sec:target}.
The exclusion concerns finite sums and products of the
specified mass-shell kernels. No assertion excluding
every generalised free field, arbitrary continuous mass
superpositions or every possible Gaussian comparison is
needed for this statement.
% !TeX root = 4DYM_notes.tex
\section{Completion of the joint construction theorem}
\label{sec:construction-assembly}

We assemble the proof of Theorem~\ref{thm:target}. Its
only external model data are the positive representation
dimension, the closed subgroup \(G\subset U(N)\), and the
nonabelian real simplicity condition of
Definition~\ref{def:group}. In particular the RG estimates,
running law, field dimensions, phase properties, vacuum
selection and trace coefficient are not hypotheses of
that theorem.

Three compatibility points deserve an explicit argument:
all finite requests must fit one calendar; reference
coefficients must represent the correctly normalised unit;
and auxiliary coefficient extractions must retain the
already chosen ordered physical limits.

\subsection{One calendar for every finite request}

\begin{lemma}[Admission of later finite requests]
\label{lem:assembly-common-calendar}
The profile constants, odd block ratios, entry thresholds
and reference levels in the preceding construction can be
chosen once, before the ultraviolet cutoffs and any
requested volume sequence. They accommodate every fixed
finite degree, arity, derivative and coefficient request
used later in the proof. A later such request changes only
the tail on which its estimate is used.
\end{lemma}
\begin{proof}
Use the coefficient bases of the \(V_d\), which precede
the continuum quotients. There is a countable cofinal
collection of finite requests: bound degrees, arities,
derivative orders, stencil radii, numbers of terms and
coefficient norms by positive integers; bound compact
position sets by integer boxes and their positive
separations from below by inverse integers. Translation
invariance removes the common centre. Tolerances have the
cofinal values \(1/l\), \(l\in\N_{>0}\).
The estimates above are uniform on the resulting
coefficient balls and separated compacts, so one request
represents all real or complex data in its specified
bounded set. Increasing the bounds gives a cofinal
collection. Every finite list of tests or linear
combinations chosen after the hierarchy has been obtained
belongs to some such request.

At profile \(j\) include the first \(j\) requests and all
their finitely many lower-order prerequisites. Choose
the pre-block geometric, classical and kernel bounds
first. Then choose the odd ratio \(L_j\) to satisfy the
finitely many floors, including
\[
 \frac{d_j+1}{\log L_j}\le\min(b_0/32,j^{-2}),\qquad
             \frac{\log\mathcal C_j}{\log L_j}\le j^{-2}.
\]
Here \(d_j\) is the pre-ratio drift bound and
\(\mathcal C_j\) bounds the finite source and column
constants. Their pre-ratio construction was proved in
Sections~\ref{sec:ordinary-matching}--%
\ref{sec:one-loop-drift} and
Section~\ref{sec:insertion-transport}.
Take the maximum with the preceding ratio.
Choose the noise-sector widths after the ratios, and
only then the entry thresholds. For each fixed adjacent
profile pair and fixed ratio, every required small error
tends to zero as the nominal coupling tends to infinity.
The same is true of the additional conditions
\eqref{eq:induction-trace-calendar} and the local
vacuum-ratio reserve. All these are finite entry
requirements at that profile.

Corollary~\ref{cor:linear-profile-clock} produces one
clock tending to infinity with adjacent increments at
most one. Its proof permits arbitrary further delays
of the first entry of a profile. Include the period
reserves and gentleness conditions in these delays as
in Section~\ref{sec:nonlinear-induction-step}.
At every fixed profile all constants are finite; no
bounded supremum over all profiles is required.
Choose \(\operatorname{ref}(d)\) once the finite degree-\(d\) insertion
request has been admitted, with increasing reference
levels if necessary.

A finite basis or probe matrix chosen from the eventual
hierarchy adds finite coefficients and finitely many
degrees and separation parameters. Its bounds fit a
member of the cofinal collection. Its finite-dimensional
inverse, when used, has a finite norm after its
injectivity has been proved; bounding that norm by a
larger integer likewise changes only the working tail.
This observation is what allows the finite mixed probes
in Section~\ref{sec:ope-coefficients} to be chosen from
the resulting hierarchy without choosing new RG ratios
afterwards. The same reasoning applies to a fixed
spectator list, derivative order or trace comparison
request. Thus every later finite calculation uses this
one initially selected calendar.
\end{proof}

\subsection{The physical unit in the reference coordinates}

\begin{lemma}[Reference normalisation]
\label{lem:assembly-reference-unit}
For \(n\ge \operatorname{ref}(d)\),
\[
 \mathsf Z_{n,M,d}({\varepsilon}_n^d1)=u_d={\varepsilon}_{\operatorname{ref}(d)}^d1,\qquad
             \mathcal O_{n,M}((d,u_d),x;U)=1.
\]
The common unit in \(\mathcal F_d\) is
\(1_d=[u_d]_d\). It is preserved by the inclusions,
annihilated by derivatives, and maps to the single unit
label of the Wightman theory.
\end{lemma}
\begin{proof}
In the exact degree-\(d\) marked step, the density is
multiplied by \(L_m^d\) while its unit remainder remains
zero. Thus \(L_m^{\rm ins}1=L_m^d1\). Since
\({\varepsilon}_m=L_m{\varepsilon}_{m+1}\), iteration gives
\[
 \mathsf Z_{n,M,d}({\varepsilon}_n^d1)
      ={\varepsilon}_n^d\prod_{m=\operatorname{ref}(d)}^{n-1}L_m^d\,1
      ={\varepsilon}_{\operatorname{ref}(d)}^d1 .
\]
Insert its inverse in
\eqref{eq:constructed-local-insertion}. Both
\(\mathsf D_{h_n^{-1}}\) and the local polynomial
section fix the polynomial unit, so the physical
insertion is \({\varepsilon}_n^{-d}{\varepsilon}_n^d=1\).
Consequently its response deletes that insertion
slot in every ordered correlation. The quotient
identifies the physical unit with \([u_d]_d\), and
faithfulness makes the inclusion and derivative
identities for this unit exact.
The direct limit of these inclusions has one common
unit. By the reconstructed unit rule,
\[
                    \Phi(1,f)=\left(\int f\right)I .
\]
The finitely many cutoffs below \(\operatorname{ref}(d)\) have no role
in any of these limiting statements. The original
reference level and \(u_d\) are retained when the
terminal cutoff sequence is refined.
\end{proof}

\subsection{Auxiliary extractions preserve the hierarchy}

\begin{lemma}[Persistence of the ordered limits]
\label{lem:assembly-auxiliary-persistence}
Fix the Wilson family after its common ultraviolet
refinement, the selected volume refinement and the
smooth hierarchy \(S\). A further extraction used to
take limits of bounded finite coefficient arrays does
not change its fixed-period rows, \(S\), the common
field quotients, the phase decomposition or an already
selected phase. Convergent bounded moving labels may
be passed through its separated correlations.
\end{lemma}
\begin{proof}
Every fixed-period row is already a limit along the
full retained mesh sequence. Every outer row value
already converges along the fixed volume refinement.
Subsequences have these same limits.
For clarity, if \(M_j\) is the selected volume
sequence, choose a retained cutoff \(n_j\) so late
that the first \(j\) fixed compact basis tests and
their prescribed derivatives at period \(M_j\)
differ from their row values by at most \(1/j\).
Their row values tend to those of \(S\).
This produces an auxiliary sequence sampling the
same hierarchy; it is not a replacement for either
of its two defining limits.

For a finite multilinear correlation, write a
moving label list as \(\mathbf v_j\to\mathbf v\).
Subtract the two evaluations by replacing one slot
at a time. On a common separated compact the
coefficient-ball bound gives
\[
 \bigl|S_{n_j,M_j}(\mathbf v_j;x)
              -S_{n_j,M_j}(\mathbf v;x)\bigr|
       \le C\sum_i|v_{j,i}-v_i|_1
                      \prod_{\ell\ne i}
                       \max(|v_{j,\ell}|_1,|v_\ell|_1).
\]
The right side tends to zero, uniformly on that
compact. The fixed-label lattice-to-row and
row-to-continuum errors then identify the limit
with \(S(\mathbf v;x)\). The identical argument
works for tested correlations, fixed differences
and the smooth derivatives already constructed.
An auxiliary choice for successively smaller
fixed radii uses this argument at each radius,
after that radius and its finite request are fixed.
It exchanges none of the physical limits with
the radius limit.

Bounded finite arrays have convergent subsequences.
A countable diagonal choice handles every fixed
degree or coarse level. Closed identities and
uniform inequalities pass to their limits. Whenever
an inverse is needed, the preceding detector or
column bound supplies a uniform lower bound, so
the inverse is retained as well. This is the
argument used for \(\Pi_d^\infty\), the field
calculus and the trace Ward functionals.

The null spaces are defined from \(S\), not from
an arbitrary limiting matrix. Their identification
with equation kernels was proved using the same
correlation limits. The phase algebra consists of
operators constructed from \(S\); its character
space and vacuum measure therefore remain the
same. An already chosen character denotes the
same state after every such auxiliary extraction.
\end{proof}

We use the following harmless convention when packaging
the conclusion as total functions. Give the correlations
the value zero at collision configurations; retain the
empty correlation \(S_0=1\). Give a continuation the
value zero outside its ordered tube. The coefficient
prescription already has a total extension outside its
probe domain. All required identities use separated
configurations or the open ordered tube. In the OPE,
its sufficiently small threshold includes \(Kr<c/2\),
so all correlations in the remainder are separated.
These conventions therefore add data without adding
an analytic assertion at collisions or outside a tube.

\subsection{Proof of the main theorem}

\begin{proof}[Proof of Theorem~\ref{thm:target}]
\emph{1. Fix the group and the common construction data.}
Let \(G\subset U(N)\) satisfy
Definition~\ref{def:group}.
Lemma~\ref{lem:group-data-calculus} supplies the local
matrix calculus and probability Haar density used in
the finite integrals. The invariant trace metric is
positive on \(\mathfrak g\). Simplicity gives
\(C_G>0\) and hence
\[
                     b_0=\frac{11C_G}{24\pi^2}>0
\]
by \eqref{eq:drift-casimir}--\eqref{eq:drift-b0}.
Fix the admissible positive kernel parameters, the
cofinal catalogue and one calendar as in
Lemma~\ref{lem:assembly-common-calendar}.
Choose a positive coarse coupling \(\bar\beta\)
large enough for its initial entry. Choose the
strictly increasing physical periods with the
initial divisibility and patch margins.
All these choices precede the requested volume sequence.

\emph{2. Construct the exact Wilson histories.}
The calibrated linear data are supplied by
Theorem~\ref{thm:constructed-calibration}. Their exact
Schur-complement recursion identifies the output of
one step with the input of the next; thus there is
one Gaussian trajectory within each actual history.
The stationary, interacting, component and regular
derivative estimates apply to those data.
In particular \eqref{eq:coercivity-noise-root-identity}
controls the local block-root modes by actual noise rows;
Lemma~\ref{lem:actual-combes-thomas} then gives common
full and principal inverse bounds. The complete one-loop
read \eqref{eq:ordinary-loop-complete-second-read}
keeps the Haar, quartic and two-cubic terms before its
kinetic coefficient is extracted.
Theorem~\ref{thm:central-joint-strength-bound},
Theorem~\ref{thm:sharp-incoming-transport} and
Theorem~\ref{thm:ordinary-auxiliary-action} provide
the central remainder, incoming and matching
estimates required by the nonlinear step.
The drift bound is uniform in its ratio before
that ratio is chosen, by
Theorem~\ref{thm:uniform-one-loop-drift}.

These are precisely the inputs of
Theorem~\ref{thm:nonlinear-one-step}. Wilson
initialisation and the running-window argument then
give Theorem~\ref{thm:finite-cutoff-nonlinear-induction}
for every \(n,M\), with
\[
 F_{n,M,n}=e^{-\beta_{n,M,n}W_{\ell_{n,M}}},\qquad
 \|E_{n,M,m}\|_{\mathcal E,\kappa(m)}\le\delta_m,\qquad
 \|K_{n,M,m}\|_{\mathcal K,\kappa(m)}\le w_m .
\]
The actual couplings lie in their running windows and
\(\beta_{n,M,n}/\log {\varepsilon}_n^{-1}\to b_0\), uniformly in \(M\).
The identities retain all scalar factors, masks,
inventories, incoming activities and Gaussian/Haar
normalisations. In particular the estimates for a
central reference sector have been reassembled into
the full original Wilson density before this theorem
is applied.

\emph{3. Fix one family of local insertions.}
Take the token set to be
\(\mathcal V=\bigsqcup_{d\ge0}V_d\). Birth and exact
one-insertion transport define
\(\mathcal O_{n,M}\) by
\eqref{eq:constructed-local-insertion}; its terminal
coefficient map is chosen once at the origin and
translated literally. Corollary~\ref{cor:local-insertion-regularity}
proves continuity, gauge invariance and the fixed
finite stencil for every token. The stencil is
unchanged by the cutoff-dependent coefficient map.
Lemma~\ref{lem:assembly-reference-unit} gives the unit
normalisation. The running law in step 2 completes
Wilson regularity for this family.

\emph{4. Make the common ultraviolet choice before the volume sequence.}
Theorem~\ref{thm:uniform-separated-insertions} bounds
every fixed finite moment on separated compacts.
At each suitable fixed period its lattice smearings
are therefore bounded, as stated in
Lemma~\ref{lem:uniform-row-smearings}.
Proposition~\ref{prop:common-uv-distribution-rows}
chooses one increasing \(\nu\), simultaneously for
every fixed period and all finite basis tests.
Define the Wilson family \(C\) of the theorem by
reading every terminal datum at the old cutoff
\(\nu(n)\). Its mesh is \({\varepsilon}_{\nu(n)}\); its internal
RG levels and reference units retain their old
indices. Lemma~\ref{lem:fixed-uv-refinement}
preserves Wilson regularity and all cutoff-tail
estimates. This is the final choice of \(C\);
it has been made before any volume request.
For this family the fixed-period mesh limits are
full limits over its cutoff index.

\emph{5. Take the ordered hierarchy for an arbitrary
requested volume sequence.}
Now fix any strictly increasing sequence of admissible side lengths
\(M_j\to\infty\).
Theorem~\ref{thm:ordered-weak-hierarchy} chooses one
strictly increasing \(q\) and the mixed hierarchy \(S\).
For every token list and suitable test, first the
mesh limit is the fixed-period row, and then the
rows along \((M_{q(j)})_j\) converge to its \(S\)-smearing.
One \(q\) works for all finite lists and tests.
Lemma~\ref{lem:common-finite-detection} is applied to the
actual finite-cutoff images before bounding the differentiated
coefficients. Its fixed-domain version then controls the
continuum comparison matrices on these same ordered tails.
Theorem~\ref{thm:smooth-ordered-hierarchy} and
Corollary~\ref{cor:ordered-smooth-local-convergence}
give smooth kernels and identify their local uniform
limits with these same distributions.
This proves the ordered Wilson-limit clause with
the centred placement and cell factor of
\eqref{eq:smear}.

\emph{6. Define the common field and phase data.}
The identification with the finite free detector spaces uses
the actual Taylor construction and finite-interval estimate
of Lemmas~\ref{lem:actual-chart-Taylor}--\ref{lem:actual-Taylor-window-comparison}.
The local coefficient matrix
\eqref{eq:actual-Taylor-calibration-matrix}
has its uniform inverse before the finite source iteration.
The estimates retain the actual regular and divided forced
families on the same physical history; no global comparison
with an evolving Gaussian graph is invoked.
Put \(\mathcal F_d=V_d/\mathcal N_d(S)\), with the
quotient norm of Section~\ref{sec:continuum-field-calculus},
and project a token \((d,v)\) to \((d,[v]_d)\).
This projection is surjective onto the tagged field
spaces. They are finite dimensional, and
Theorem~\ref{thm:continuum-field-calculus} gives their
injective compatible inclusions and derivatives.
The auxiliary coefficient extractions used to
identify these maps retain \(S\), by
Lemma~\ref{lem:assembly-auxiliary-persistence}.

The stress construction and
Theorem~\ref{thm:full-orthogonal-covariance} give a
real centred conserved tensor on these same spaces.
Its translation charges have the prescribed
normalisation, and it has the full orthogonal
tensor law. The coordinate sign transfer
\eqref{eq:coordinate-sign-transfer}, buffered word bound
\eqref{eq:angular-buffered-resolvent} and positive diagonal
gain \eqref{eq:angular-positive-diagonal-gain} remove the
exterior angular boundary before continuous rotations are
identified. The zero-energy part is kept in that estimate.
The compression algebra
\(\mathcal A_0\) of
Proposition~\ref{prop:phase-compression-algebra}
then supplies a compact metric character space
\(\Sigma\), the full-support probability measure
\(\rho\), and \(S^\sigma(w)=\sigma(\mathsf M^0(w))\).
Equation~\eqref{eq:phase-mixture} represents the
original \(S\), not a new physical limit.
Proposition~\ref{prop:phase-correlations} proves
smoothness, permutation symmetry, character
continuity, reflection positivity and the common
field quotients for every phase. The equality
\(\mathcal N_d(S^\sigma)=\mathcal N_d(S)\)
also ensures that no phase uses a smaller
unidentified field space.

\emph{7. Select one vacuum and reconstruct it.}
Choose \(\sigma_*\) once from the full-measure set
in Theorem~\ref{thm:phase-vacuum-selection}.
Reflection invariance uses the finite variance
\eqref{eq:phase-finite-variance}, with bounded coarse
approximations chosen before the copy count. The common
selection set is the countable intersection furnished by
\eqref{eq:phase-ground-integral}--\eqref{eq:phase-ground-defect};
it precedes reconstruction and every later trace request.
Let \(\mathcal L\) be the algebraic direct limit
of the \(\mathcal F_d\), with injective maps
\(\iota_d\). These maps exhaust \(\mathcal L\)
by its definition. Sections~\ref{sec:wightman-continuation}
and~\ref{sec:wightman-reconstruction} construct
the continuation and the common word-domain
fields for \(S^{\sigma_*}\).
Theorem~\ref{thm:selected-wightman-reconstruction}
gives W0--W3, the faithful label action, the
normalised cyclic vacuum, and uniqueness of the
translation-invariant vacuum.

More explicitly, its translation PVM represents
the translation matrix elements and is supported
in the forward cone. Its word domain is dense,
invariant, and exactly the span of field words
applied to the vacuum; the fields are continuous
in Schwartz tests on each domain vector and obey
the smeared adjoint identity. Its proper
orthochronous Poincare action has the ordinary
affine pullback on tests and the compatible
finite-degree Lorentz action on labels. It
preserves that domain and the vacuum. Fields
with spacelike separated supports commute there.
The strict imaginary-time root is
\(S^{\sigma_*}\), and its smeared boundary is
the vacuum moment of these very fields.
The common unit acts as \((\int f)I\), supplying
a nonzero smeared operator.
Thus all field, phase and reconstruction
requirements in the second clause hold for one
selected vacuum.

\emph{8. Fix the coefficients and prove every finite
separated expansion.}
Lemma~\ref{lem:ope-probe-coordinates} constructs the
finite channel sets \(I_j\), actors
\(O_{j,\lambda}\in\mathcal F_j\), and one compatible
family of mixed probe functionals. Their prescribed
cluster evaluations define the coefficients.
They depend on \(S\) and its common field spaces;
they do not depend on the selected character,
spectators or truncation.
Theorem~\ref{thm:all-orders-separated-ope} gives
their coefficient bound and the remainder bound
with precisely the exponents in
\eqref{eq:ope-coefficient}--\eqref{eq:ope-remainder}.
The finite Taylor formula
\eqref{eq:ope-single-taylor-extension} removes the grid
using the existing derivative calculus and finitely many
additional degree requests in the same calendar.
The common algebra norm controls every character
and the mixed state, so one threshold works for
all those states. Empty clusters use the exact
unit expansion. For arbitrary finite index sets,
enumerate the input and spectator slots; symmetry
makes their correlations the grouped evaluations
of this same hierarchy.
No new extraction defines either the coefficients
or the correlations. The auxiliary lattice
matching fields only estimate the fixed finite
decoder, as justified in
Lemma~\ref{lem:assembly-auxiliary-persistence}.

\emph{9. Use the trace of that same selected theory.}
With the tensor from step 6 and the phase selected
in step 7, define
\[
 e_{\rm tr}=S_1^{\sigma_*}\Bigl(\sum_aT_{aa}\Bigr),
       \qquad
 \Theta=\sum_aT_{aa}-e_{\rm tr}1_4 .
\]
Lemma~\ref{lem:centred-trace-identification} gives
its exact centring on the common Wightman domain.
Prescribed-radius comparison and the actual
current estimates are applied along the same
Wilson family. Their auxiliary coefficient
limits preserve \(S\) and \(\sigma_*\).
The source entrance calculation
\eqref{eq:prescribed-source-entrance} pays the growing
column product; the root bound
\eqref{eq:prescribed-complete-source-root} keeps the
original denominator and all subset corrections. Thus the
comparison applies to prescribed radii without changing
the mixed hierarchy or the already selected character.
The joint Ward comparison uses the exact mixed coefficient
\eqref{eq:joint-ward-four-term-density}, with the shared
tags and the one original partition denominator retained.
Its scale factors are paid at the stopping level before
the finite detector bounds the transported vector.
Proposition~\ref{prop:trace-leading-coefficient}
identifies the first nonzero trace coefficient;
Theorem~\ref{thm:trace-short-distance} gives
\[
 \lim_{r\downarrow0}r^8\log^2(1/r)
       S_2^{\sigma_*}(\Theta,\Theta;0,re_4)
                              =\frac{3\dim_\R\mathfrak g}{\pi^4}>0.
\]
Translation invariance identifies this two-point
function with the symmetric placement at times
\(-r/2\) and \(r/2\) in the theorem.
Theorem~\ref{thm:trace-nonfreeness} excludes every
finite profile in the specified class. No further
phase or vacuum selection occurs in this step.

Finally collect the finite-dimensional spaces,
their tagged projection and label maps, the one
phase space and measure, the selected theory,
its continuation, the channel and coefficient
family, and \(\Theta\), with the total-function
conventions above. Step 5 proves its ordered
Wilson origin, steps 6--7 its reconstruction
data, step 8 its separated OPE, and step 9 its
trace test. All four clauses use this one
hierarchy along \((M_{q(j)})_j\). Since \(C\) was
fixed before the arbitrary sequence \((M_j)\), the required
quantifier order is proved.
\end{proof}

\subsection{The shared objects in the four conclusions}

For reference, the following list records which
already constructed object each conclusion uses.
It is part of the identification made in the proof:
\begin{center}
\begin{tabular}{@{}p{.24\linewidth}p{.69\linewidth}@{}}
\hline
Wilson origin &
The mixed \(S\), with full mesh limits for the fixed
family \(C\), followed by the rows along \((M_{q(j)})_j\).\\[3pt]
Field and phase data &
\(\mathcal F_d=V_d/\mathcal N_d(S)\);
the character space \(\Sigma\) of \(\mathcal A_0(S)\)
and its one vacuum probability measure \(\rho\).\\[3pt]
Reconstruction &
The fixed \(S^{\sigma_*}\), its algebraic label union,
common word domain, vacuum and translation PVM.\\[3pt]
Separated OPE &
The mixed-probe coefficient family from \(S\);
the same \(S\) and all \(S^\sigma\) in its remainders.\\[3pt]
Trace test &
The common normalised tensor \(T\), centred using
\(S^{\sigma_*}\), acting in that reconstructed theory.\\
\hline
\end{tabular}
\end{center}

The finite requests have their own constants and
tails; none chooses a new physical limit. The
theorem asserts every finite separated truncation
and the stated finite-profile trace exclusion.
It does not add a positive mass gap or an
infinite-sum convergence assertion for the OPE.

\medskip
\noindent\textbf{Citation conventions.}
For Ba{\l}aban's papers, printed-page locators use the original journal
pagination. Preprint-page and PDF-page locators for
\cite{FP07,BS15a,BS15b,DST23,Landsman98}
use the original arXiv versions linked below, rather than a transcription.
Unverified transcription locators and secondary attributions are identified
at their point of use and in the relevant bibliography entries.

\begin{thebibliography}{99}
\bibitem{OpenAIYM26}
OpenAI,
\emph{Wilson Yang--Mills construction},
OpenAI-YM-20261028, Lean source release (2026).
\href{https://github.com/yang-mills-group/OpenAI-YM-20261028/tree/3456e5e75cf8d4cd558955a5c7638fba2d68f464}{Source version used},
revision \nolinkurl{3456e5e75cf8d4cd558955a5c7638fba2d68f464}.
This is the formal source on which these reading notes are based.
The present exposition has not been mechanically checked against it.
\bibitem{Wilson74}
K.~G. Wilson, \emph{Confinement of quarks},
Physical Review D \textbf{10} (1974), 2445--2459.
\href{https://doi.org/10.1103/PhysRevD.10.2445}{doi:10.1103/PhysRevD.10.2445}.
Historical reference for the lattice formulation; the finite
measure calculations used here are written out in
Section~\ref{sec:wilson}.
\bibitem{FP07}
R. Fern\'andez and A. Procacci,
\emph{Cluster expansion for abstract polymer models.
New bounds from an old approach},
Communications in Mathematical Physics \textbf{274} (2007), 123--140.
\href{https://doi.org/10.1007/s00220-007-0279-2}{doi:10.1007/s00220-007-0279-2}.
Page locators use
\href{https://arxiv.org/abs/math-ph/0605041v2}{arXiv:math-ph/0605041v2}.
The incompatibility-graph convention is described in Sect.~2,
eq.~(2.2), preprint p.~3; the logarithmic and pinned expansions are
eqs.~(2.3) and (2.6)--(2.7), pp.~3--4, and the tree-graph partition
identity is Sect.~4.1, Proposition~5 and eqs.~(4.3)--(4.6),
pp.~10--11. The connected-family bound used
here is proved in Section~\ref{sec:connected-bound}.
\bibitem{Bal87}
T. Ba{\l}aban,
\emph{Renormalization group approach to lattice gauge field theories.
I. Generation of effective actions in a small field approximation
and a coupling constant renormalization in four dimensions},
Communications in Mathematical Physics \textbf{109} (1987), 249--301.
\href{https://doi.org/10.1007/BF01215223}{doi:10.1007/BF01215223}.
Section~2 develops the fluctuation integral, printed pp.~265--268;
eqs.~(2.12)--(2.15), p.~268, retain its normalisation and coupling reset.
The quadratic coefficient is defined in Sect.~1, eqs.~(1.20)--(1.22),
p.~264, and analysed in Sect.~5, pp.~292--297.
\bibitem{BS15a}
D. C. Brydges and G. Slade,
\emph{A renormalisation group method.
I. Gaussian integration and normed algebras},
arXiv:1403.7244v2 (2014).
\href{https://arxiv.org/abs/1403.7244v2}{Preprint version used};
\href{https://doi.org/10.1007/s10955-014-1163-z}{published version}.
The product principle is Proposition~3.7, PDF p.~15;
its proof is in Sect.~5.1, starting on PDF p.~24.
The factorial-weighted adaptation to ordered left
derivatives is proved in
Lemma~\ref{lem:finite-jet-algebra}.
\bibitem{BS15b}
D. C. Brydges and G. Slade,
\emph{A renormalisation group method. II. Approximation by local polynomials},
Journal of Statistical Physics \textbf{159} (2015), 461--491.
\href{https://doi.org/10.1007/s10955-014-1164-y}{doi:10.1007/s10955-014-1164-y}.
Page locators use
\href{https://arxiv.org/abs/1403.7253v2}{arXiv:1403.7253v2}.
The localisation operator is constructed in Sect.~1, Proposition~1.5,
preprint p.~7; its norm and remainder bounds are Propositions~1.11--1.12,
pp.~10--11. This is a comparison with the general extraction framework;
the gauge-invariant affine extraction used here is proved in
Section~\ref{sec:coupling-reset}.
\bibitem{Bal85g}
T. Ba{\l}aban,
\emph{Spaces of regular gauge field configurations on a lattice
and gauge fixing conditions},
Communications in Mathematical Physics \textbf{99} (1985), 75--102.
\href{https://doi.org/10.1007/BF01466594}{doi:10.1007/BF01466594}.
The axial gauge is described in Sect.~B, printed pp.~78--80;
Lemma~1 and eqs.~(1.24)--(1.26), pp.~79--80, give the local path comparison.
The finite-patch path comparisons and derivative bounds used here
are proved in Sections~\ref{sec:actual-stationary}
and~\ref{sec:stationary-localisation}.
\bibitem{Bal85v}
T. Ba{\l}aban,
\emph{The variational problem and background fields in renormalization
group method for lattice gauge theories},
Communications in Mathematical Physics \textbf{102} (1985), 277--309.
\href{https://doi.org/10.1007/BF01229381}{doi:10.1007/BF01229381}.
The bibliographic details and Theorem~1 locator are taken from
\cite[Sect.~1, printed p.~260 and reference~15]{Bal87}.
The original paper has not been checked here. This entry records a
historical attribution, rather than an application of that theorem
to the blocking action used in this manuscript.
\bibitem{DST23}
W. Dybalski, A. Stottmeister and Y. Tanimoto,
\emph{Lattice Green functions for pedestrians: Exponential decay},
arXiv:2303.10754v1 (2023).
\href{https://arxiv.org/abs/2303.10754v1}{Downloaded preprint version}.
The Combes--Thomas bound is Sect.~2.4, Lemma~2.12, preprint p.~14;
its proof is in Appendix~A, pp.~36--37.
The Fourier-analytic kernel argument is Sect.~3.5, Lemma~3.17,
p.~32. These scalar estimates supply methodological comparisons.
The finite-matrix inverse estimate for the present Hessian is proved
in Lemma~\ref{lem:actual-combes-thomas}.
\bibitem{Bal88a}
T. Ba{\l}aban,
\emph{Renormalization group approach to lattice gauge field theories.
II. Cluster expansions},
Communications in Mathematical Physics \textbf{116} (1988), 1--22.
\href{https://doi.org/10.1007/BF01239022}{doi:10.1007/BF01239022}.
Section~1, Lemmas~1--2, printed pp.~9 and~11, localises the fluctuation
interaction. Section~2, eqs.~(2.11)--(2.13), p.~14, passes from its
polymer representation to the connected effective action.
This organisational comparison is distinguished from the quadratic
determinant calculation in Section~\ref{sec:stationary-localisation}.
\bibitem{Bal88b}
T. Ba{\l}aban,
\emph{Convergent renormalization expansions for lattice gauge theories},
Communications in Mathematical Physics \textbf{119} (1988), 243--285.
\href{https://doi.org/10.1007/BF01217741}{doi:10.1007/BF01217741}.
The inductive density representation is Sect.~2, printed pp.~254--258;
eqs.~(2.17)--(2.19), pp.~257--258, separate the small-field tests and
large-field component operations. This is a comparison of the
organisation of the density. Theorem~1, p.~262, has a hypothesis on
the effective coupling sequence and is not used as an input for the
coupling estimates of this manuscript.
\bibitem{Bal89b}
T. Ba{\l}aban,
\emph{Large field renormalization. II. Localization, exponentiation,
and bounds for the R operation},
Communications in Mathematical Physics \textbf{122} (1989), 355--392.
\href{https://doi.org/10.1007/BF01238433}{doi:10.1007/BF01238433}.
Section~1, eqs.~(1.71)--(1.75), printed pp.~379--380, describes the
component integration and its normalisation; pp.~380--383 extract
large-field suppression factors; eqs.~(1.98)--(1.100), p.~390,
give the final exponentiation and activity bounds.
These are comparisons with the large-field organisation.
Theorem~1, p.~355, assumes a uniformly small effective coupling sequence;
the present packet estimates and nonlinear induction are proved here.
\bibitem{Sei82}
E.~Seiler,
\emph{Gauge Theories as a Problem of Constructive Quantum Field
Theory and Statistical Mechanics},
Lecture Notes in Physics \textbf{159}, Springer, 1982.
Reflection positivity and the chessboard estimate are
Theorems~2.1--2.2, printed pp.~14--16.
\href{https://doi.org/10.1007/3-540-11559-5}{doi:10.1007/3-540-11559-5}.
The reconstruction framework for nonlocal gauge-invariant objects is
Sect.~8(b)--(e), pp.~167--179: Theorems~8.8--8.10 are on p.~175,
and Theorems~8.13--8.14 on pp.~178--179.
All printed-page locators for this book follow the original-page markers
in the transcription and have not been checked against the original scans.
The reconstruction reference is a framework comparison; the local-field
construction used here is proved in Sections~\ref{sec:wightman-continuation}--%
\ref{sec:wightman-reconstruction}.
\bibitem{Landsman98}
N.~P. Landsman,
\emph{Lecture notes on C*-algebras, Hilbert C*-modules, and quantum mechanics},
arXiv:math-ph/9807030v1 (1998), draft of 8 April 1998.
\href{https://arxiv.org/abs/math-ph/9807030v1}{Original preprint version used}.
Section~2.4, Theorem~2.4.1, p.~25, with proof on pp.~25--27,
gives the representation of a commutative unital C*-algebra by continuous
functions on its character space. Section~2.8, Theorem~2.8.2, p.~35,
identifies states on that function algebra with probability measures.
These original-preprint page locators have been checked against the
original PDF. The standard results are used in
Section~\ref{sec:phase-compression}; the phase-specific operator estimates
and full-support argument are derived in this manuscript.
\bibitem{MRS93}
J. Magnen, V. Rivasseau and R. S\'en\'eor,
\emph{Construction of YM4 with an infrared cutoff},
Communications in Mathematical Physics \textbf{155} (1993), 325--383.
\href{https://doi.org/10.1007/BF02097397}{doi:10.1007/BF02097397}.
The Abstract and Sect.~I, printed pp.~325--326, describe the scope and
regularisation of the construction. These locators have been checked
against the original journal PDF. The citation is a historical comparison
with a fixed-infrared-cutoff analysis in a regularised axial gauge.
\bibitem{ArmstrongNotes}
Contributors to the Yang--Mills informalization project,
\emph{Four-dimensional quantum Yang--Mills theory: the proof,
outlined (detailed notes)}, 11 November 2026, project
maintained by Scott Armstrong, 595-page working notes.
\href{https://github.com/scottnarmstrong/yang-mills-informalization/blob/1a47ac9007f80d37832291cc12d578874ff6bef9/pdf/yang-mills-detailed-notes.pdf}{Public PDF at the source revision used}
(\texttt{1a47ac9007f8}); page and unit locators below refer to this version.
Unit II.2, Propositions II.2.10--II.2.11, gives the finite
Gaussian contraction and polynomial-step calculus.
Unit III.4, Lemma 4.25, gives the Gaussian density-ratio
calculation; Unit III.11, Lemmas 4.9 and 4.13, organizes
the measured-covariance and local polynomial comparison.
Their use in Section~\ref{sec:gaussian-position-tests}
is expanded in the common regular and divided forced norms.
Unit III.1, PDF pp.~83--98, supplies the branch definitions,
axial layer induction and first-power frame comparison.
The corresponding estimates are derived explicitly here.
Unit III.10, Lemmas 4.4--4.14, PDF pp.~259--262,
gives the row estimates and large-field energy accounting
expanded in Section~\ref{sec:hole-packet-integrals}.
Units V.1, VI.1 and VII.6 give, respectively, the phase
variance and common vacuum selection, exterior angular flux,
and prescribed-radius comparison used later; the relevant
local calculations are proved in the corresponding sections.
\bibitem{ArmstrongDraft}
Contributors to the Yang--Mills informalization project,
\emph{The construction of four-dimensional quantum Yang--Mills theory},
working draft of 11 November 2026, project maintained by
Scott Armstrong.
\href{https://github.com/scottnarmstrong/yang-mills-informalization/blob/1a47ac9007f80d37832291cc12d578874ff6bef9/manuscript/main.pdf}{Public PDF at the source revision used}
(\texttt{1a47ac9007f8}).
Page locators refer to the 241-page main manuscript of that date.
The local-integration organisation is in Lemma~11.13,
PDF p.~73; incoming-term separation and finite-period
localisation are in Lemmas~13.5--13.6, PDF pp.~95--98.
Curl division and classical matching are in
Lemmas~12.5--12.6, PDF pp.~79--81.
The determinant comparison and logarithmic coefficient
are in Lemmas~12.7--12.9, PDF pp.~82--85;
Wilson initialisation is Lemma~13.21, PDF pp.~119--120.
The exact density representation and activity-sum continuity are
Lemmas~13.12--13.13, PDF pp.~107--108; the complete marked-source inputs
are Lemmas~14.12--14.13, pp.~135--137.
Compression is Lemma~17.4, pp.~194--195; phase fields and vacuum
selection are Lemma~17.5, pp.~195--197.
This is a developing exposition. The elementary reductions
adapted from it are proved here; it is not used as a
substitute for unexpanded model-specific estimates.
\end{thebibliography}
\end{document}
